% Auto-generated from data/segments.json + layout.json (+ transforms) — do not edit by hand.
% volume: 38Abhi10
% mode: printing
% segments: 4966
% transform_rules: 8
% toc_marks: 150

% Volume layout defaults (layout.json ``layout``).
\csromanlayoutapply{1}{1.5}{21.6pt}{8pt}{8pt}{65pt}{2.5em}%

\pitaka{อภิธมฺม\-ปิฏก}
\csromanheader{2}{\textbf{ธมฺมานุโลม}}
\csromanmatikaanchor{mtk.0}
\csromantocmarknopage{chapter}{ปฏฺฐานปาฬิ}{mtk.0}
\bookmark[dest={mtk.0},level=0]{ปฏฺฐานปาฬิ}
\gambhira{ทุกปฏฺฐาน\-ปาฬิ}
\namakkaram{นโม ตสฺส ภควโต อรหโต สมฺมา\-สมฺพุทฺธสฺส.}
\csromanmatikaanchor{mtk.1}
\csromanmatikaanchor{mtk.2}
\csromantocmarknopage{section}{ตติยภาค}{mtk.1}
\bookmark[dest={mtk.1},level=1]{ตติยภาค}
\csromantocmarknopage{chapter}{1. เหตุโคจฺฉก}{mtk.2}
\bookmark[dest={mtk.2},level=0]{1. เหตุโคจฺฉก}
\chapterhead{1. \textbf{เหตุโคจฺฉก}}
\csromanmatikaanchor{mtk.3}
\csromantocmarkref{section}{1. เหตุทุก}{mtk.3}
\bookmark[dest={mtk.3},level=1]{1. เหตุทุก}
\csromanmarkh{1. เหตุทุก}\csromanmarkhii{1. ปฏิจฺจวาร}\csromanheadernomark{1}{1. \textbf{เหตุทุก 1. ปฏิจฺจวาร}}
\csromanmarkh{1. ปจฺจยานุโลม}\csromanmarkhii{1. วิภงฺควาร}\csromanheadernomark{2}{1. \textbf{ปจฺจยานุโลม 1. วิภงฺควาร}}
\csromanheader{2}{\textbf{เหตุ}}
\proseitem{1}{\csromanfolio{1}เหตุํ ธมฺมํ ปฏิจฺจ เหตุ ธมฺโม อุปฺปชฺชติ เหตุปจฺจยา. อโลภํ ปฏิจฺจ อโทโส อโมโห. อโทสํ ปฏิจฺจ อโลโภ อโมโห. อโมหํ ปฏิจฺจ อโลโภ อโทโส. โลภํ ปฏิจฺจ โมโห. โมหํ ปฏิจฺจ โลโภ. โทสํ ปฏิจฺจ โมโห. โมหํ ปฏิจฺจ โทโส. ปฏิสนฺธิ\-กฺขเณ ฯเปฯ. (1)}

\prose{เหตุํ ธมฺมํ ปฏิจฺจ นเหตุ ธมฺโม อุปฺปชฺชติ เหตุปจฺจยา. เหตุํ ธมฺมํ ปฏิจฺจ สมฺปยุตฺตกา ขนฺธา จิตฺต\-สมุฏฺฐานญฺจ รูปํ. ปฏิสนฺธิ\-กฺขเณ ฯเปฯ. (2)}

\prose{เหตุํ ธมฺมํ ปฏิจฺจ เหตุ จ นเหตุ จ ธมฺมา อุปฺปชฺชนฺติ เหตุปจฺจยา. อโลภํ ปฏิจฺจ อโทโส อโมโห สมฺปยุตฺตกา จ ขนฺธา จิตฺต\-สมุฏฺฐานญฺจ รูปํ. (จกฺกํ.) โลภํ ปฏิจฺจ โมโห สมฺปยุตฺตกา จ ขนฺธา จิตฺต\-สมุฏฺฐานญฺจ รูปํ ฯเปฯ. ปฏิสนฺธิ\-กฺขเณ ฯเปฯ. (3)}

\proseitem{2}{\csromanfolio{2}นเหตุํ ธมฺมํ ปฏิจฺจ นเหตุ ธมฺโม อุปฺปชฺชติ เหตุปจฺจยา. นเหตุํ เอกํ ขนฺธํ ปฏิจฺจ ตโย ขนฺธา จิตฺต\-สมุฏฺฐานญฺจ รูปํ ฯเปฯ เทฺว ขนฺเธ ปฏิจฺจ เทฺว ขนฺธา จิตฺต\-สมุฏฺฐานญฺจ รูปํ. ปฏิสนฺธิ\-กฺขเณ ฯเปฯ. ขนฺเธ ปฏิจฺจ วตฺถุ, วตฺถุํ ปฏิจฺจ ขนฺธา. เอกํ มหาภูตํ ฯเปฯ. (1)}

\prose{นเหตุํ ธมฺมํ ปฏิจฺจ เหตุ ธมฺโม อุปฺปชฺชติ เหตุปจฺจยา. นเหตู ขนฺเธ ปฏิจฺจ เหตู. ปฏิสนฺธิ\-กฺขเณ ฯเปฯ. วตฺถุํ ปฏิจฺจ เหตู.}

\prose{นเหตุํ ธมฺมํ ปฏิจฺจ เหตุ จ นเหตุ จ ธมฺมา อุปฺปชฺชนฺติ เหตุปจฺจยา. นเหตุํ เอกํ ขนฺธํ ปฏิจฺจ ตโย ขนฺธา เหตุ จ จิตฺต\-สมุฏฺฐานญฺจ รูปํ ฯเปฯ เทฺว ขนฺเธ ปฏิจฺจ เทฺว ขนฺธา เหตุ จ จิตฺต\-สมุฏฺฐานญฺจ รูปํ. ปฏิสนฺธิ\-กฺขเณ ฯเปฯ. ปฏิสนฺธิ\-กฺขเณ วตฺถุํ ปฏิจฺจ เหตู สมฺปยุตฺตกา จ ขนฺธา. (3)}

\proseitem{3}{เหตุญฺจ นเหตุญฺจ ธมฺมํ ปฏิจฺจ เหตุ ธมฺโม อุปฺปชฺชติ เหตุปจฺจยา. อโลภญฺจ สมฺปยุตฺตเก จ ขนฺเธ ปฏิจฺจ อโทโส อโมโห. (จกฺกํ.) โลภญฺจ สมฺปยุตฺตเก จ ขนฺเธ ปฏิจฺจ โมโห. โทสญฺจ สมฺปยุตฺตเก จ ขนฺเธ ปฏิจฺจ โมโห ฯเปฯ. ปฏิสนฺธิ\-กฺขเณ ฯเปฯ. อโลภญฺจ วตฺถุญฺจ ปฏิจฺจ อโทโส อโมโห ฯเปฯ. (1)}

\prose{เหตุญฺจ นเหตุญฺจ ธมฺมํ ปฏิจฺจ นเหตุ ธมฺโม อุปฺปชฺชติ เหตุปจฺจยา. นเหตุํ เอกํ ขนฺธญฺจ เหตุญฺจ ปฏิจฺจ ตโย ขนฺธา จิตฺต\-สมุฏฺฐานญฺจ รูปํ ฯเปฯ เทฺว ขนฺเธ จ เหตุญฺจ ปฏิจฺจ เทฺว ขนฺธา จิตฺต\-สมุฏฺฐานญฺจ รูปํ. ปฏิสนฺธิ\-กฺขเณ ฯเปฯ. ปฏิสนฺธิ\-กฺขเณ วตฺถุญฺจ เหตุญฺจ ปฏิจฺจ สมฺปยุตฺตกา ขนฺธา. (2) ปปฺ}

\prose{เหตุญฺจ นเหตุญฺจ ธมฺมํ ปฏิจฺจ เหตุ จ นเหตุ จ ธมฺมา อุปฺปชฺชนฺติ เหตุปจฺจยา. นเหตุํ เอกํ ขนฺธญฺจ อโลภญฺจ ปฏิจฺจ ตโย ขนฺธา อโทโส อโมโห จ จิตฺต\-สมุฏฺฐานญฺจ รูปํ ฯเปฯ เทฺว ขนฺเธ จ อโลภญฺจ ปฏิจฺจ เทฺว ขนฺธา อโทโส อโมโห จ จิตฺต\-สมุฏฺฐานญฺจ รูปํ. (จกฺกํ.) นเหตุํ เอกํ ขนฺธญฺจ โลภญฺจ ปฏิจฺจ ตโย ขนฺธา โมโห จ จิตฺต\-สมุฏฺฐานญฺจ รูปํ ฯเปฯ เทฺว ขนฺเธ ฯเปฯ. ปฏิสนฺธิ\-กฺขเณ ฯเปฯ. วตฺถุญฺจ อโลภญฺจ ปฏิจฺจ อโทโส อโมโห สมฺปยุตฺตกา จ ขนฺธา. (3)}

\csromanheader{2}{\textbf{อารมฺมณาทิ}}
\proseitem{4}{\csromanfolio{3}เหตุํ ธมฺมํ ปฏิจฺจ เหตุ ธมฺโม อุปฺปชฺชติ อารมฺมณ\-ปจฺจยา. (รูปํ ฉฑฺเฑตฺวา อรูเปเยว นว ปญฺหา.) อธิปติ\-ปจฺจยา. (ปฏิสนฺธิ นตฺถิ. ปริปุณฺณํ.) เอกํ มหาภูตํ ปฏิจฺจ ฯเปฯ มหาภูเต ปฏิจฺจ จิตฺต\-สมุฏฺฐานํ รูปํ อุปาทารูปํ. (อิมํ นานํ.) อนนฺตร\-ปจฺจยา. สมนนฺตร\-ปจฺจยา. สหชาต\-ปจฺจยา. (สพฺเพ มหาภูตา, ยาว อสญฺญสตฺตา.) อญฺญมญฺญ\-ปจฺจยา. นิสฺสยปจฺจยา. อุปนิสฺสย\-ปจฺจยา. ปุเรชาต\-ปจฺจยา. อาเสวนปจฺจยา. (ทฺวีสุปิ ปฏิสนฺธิ นตฺถิ.) กมฺมปจฺจยา. วิปากปจฺจยา. (สํขิตฺตํ.) อวิคตปจฺจยา.}

\csromanmarkh{1. ปจฺจยานุโลม}\csromanmarkhii{2. สงฺขฺยาวาร}\csromanheadernomark{2}{1. \textbf{ปจฺจยานุโลม 2. สงฺขฺยาวาร}}
\proseitem{5}{เหตุยา นว, อารมฺมเณ นว, (สพฺพตฺถ นว.) อวิคเต นว. (เอวํ คเณตพฺพํ.)}

\vspace{\tipitakaparskip}
\csromancenter{\textbf{อนุโลมํ.}}
\csromanmarkh{2. ปจฺจย\-ปจฺจนีย}\csromanmarkhii{1. วิภงฺควาร}\csromanheadernomark{2}{2. ปจฺจย\-ปจฺจนีย 1. วิภงฺควาร}
\proseitem{6}{นเหตุํ ธมฺมํ ปฏิจฺจ นเหตุ ธมฺโม อุปฺปชฺชติ นเหตุปจฺจยา. อเหตุกํ นเหตุํ เอกํ ขนฺธํ ปฏิจฺจ ตโย ขนฺธา จิตฺต\-สมุฏฺฐานญฺจ รูปํ ฯเปฯ เทฺว ขนฺเธ ฯเปฯ. อเหตุก\-ปฏิสนฺธิ\-กฺขเณ ฯเปฯ. ขนฺเธ ปฏิจฺจ วตฺถุ, วตฺถุํ ปฏิจฺจ ขนฺธา. เอกํ มหาภูตํ ฯเปฯ. พาหิรํ. อาหารสมุฏฺฐานํ. อุตุสมุฏฺฐานํ. อสญฺญสตฺตานํ ฯเปฯ. (1)}

\prose{นเหตุํ ธมฺมํ ปฏิจฺจ เหตุ ธมฺโม อุปฺปชฺชติ นเหตุปจฺจยา. วิจิกิจฺฉา\-สหคเต อุทฺธจฺจ\-สหคเต ขนฺเธ ปฏิจฺจ วิจิกิจฺฉา\-สหคโต อุทฺธจฺจ\-สหคโต โมโห. (2)}

\prose{นอารมฺมณาทิ}

\proseitem{7}{เหตุํ ธมฺมํ ปฏิจฺจ นเหตุ ธมฺโม อุปฺปชฺชติ นอารมฺมณ\-ปจฺจยา. เหตุํ ปฏิจฺจ จิตฺต\-สมุฏฺฐานํ รูปํ. ปฏิสนฺธิ\-กฺขเณ ฯเปฯ. (1)}

\prose{\csromanfolio{4}นเหตุํ ธมฺมํ ปฏิจฺจ นเหตุ ธมฺโม อุปฺปชฺชติ นอารมฺมณ\-ปจฺจยา. นเหตู ขนฺเธ ปฏิจฺจ จิตฺต\-สมุฏฺฐานํ รูปํ. ปฏิสนฺธิ\-กฺขเณ ฯเปฯ. สพฺเพ มหาภูตา ฯเปฯ. (1)}

\prose{เหตุญฺจ นเหตุญฺจ ธมฺมํ ปฏิจฺจ นเหตุ ธมฺโม อุปฺปชฺชติ นอารมฺมณ\-ปจฺจยา. เหตุญฺจ นเหตุญฺจ ขนฺเธ ปฏิจฺจ จิตฺต\-สมุฏฺฐานํ รูปํ. ปฏิสนฺธิ\-กฺขเณ ฯเปฯ นอธิปติ\-ปจฺจยา. (ปริปุณฺณํ.) นอนนฺตร\-ปจฺจยา. นสมนนฺตร\-ปจฺจยา. นอญฺญมญฺญ\-ปจฺจยา. นอุปนิสฺสยปจฺจยา.}

\csromanheader{2}{\textbf{นปุเรชาต}}
\proseitem{8}{เหตุํ ธมฺมํ ปฏิจฺจ เหตุ ธมฺโม อุปฺปชฺชติ นปุเรชาต\-ปจฺจยา. อรูเป อโลภํ ปฏิจฺจ อโทโส อโมโห. (จกฺกํ.) โลภํ ปฏิจฺจ โมโห, โมหํ ปฏิจฺจ โลโภ. ปฏิสนฺธิ\-กฺขเณ ฯเปฯ. (1)}

\prose{เหตุํ ธมฺมํ ปฏิจฺจ นเหตุ ธมฺโม อุปฺปชฺชติ นปุเรชาต\-ปจฺจยา. อรูเป เหตุํ ปฏิจฺจ สมฺปยุตฺตกา ขนฺธา. เหตุํ ปฏิจฺจ จิตฺต\-สมุฏฺฐานํ รูปํ. ปฏิสนฺธิ\-กฺขเณ ฯเปฯ. (2)}

\prose{เหตุํ ธมฺมํ ปฏิจฺจ เหตุ จ นเหตุ จ ธมฺมา อุปฺปชฺชนฺติ นปุเรชาต\-ปจฺจยา. อรูเป อโลภํ ปฏิจฺจ อโทโส อโมโห สมฺปยุตฺตกา จ ขนฺธา. (จกฺกํ.) โลภํ ปฏิจฺจ โมโห สมฺปยุตฺตกา จ ขนฺธา. (จกฺกํ.) ปฏิสนฺธิ\-กฺขเณ ฯเปฯ. (3)}

\proseitem{9}{นเหตุํ ธมฺมํ ปฏิจฺจ นเหตุ ธมฺโม อุปฺปชฺชติ นปุเรชาต\-ปจฺจยา. อรูเป นเหตุํ เอกํ ขนฺธํ ปฏิจฺจ ตโย ขนฺธา ฯเปฯ. นเหตู ขนฺเธ ปฏิจฺจ จิตฺต\-สมุฏฺฐานํ รูปํ. ปฏิสนฺธิ\-กฺขเณ ฯเปฯ. เอกํ มหาภูตํ ฯเปฯ. (1)}

\prose{นเหตุํ ธมฺมํ ปฏิจฺจ เหตุ ธมฺโม อุปฺปชฺชติ นปุเรชาต\-ปจฺจยา. อรูเป นเหตู ขนฺเธ ปฏิจฺจ เหตู. ปฏิสนฺธิ\-กฺขเณ ฯเปฯ. (2)}

\prose{นเหตุํ ธมฺมํ ปฏิจฺจ เหตุ จ นเหตุ จ ธมฺมา อุปฺปชฺชนฺติ นปุเรชาต\-ปจฺจยา. อรูเป นเหตุํ เอกํ ขนฺธํ ปฏิจฺจ ตโย ขนฺธา เหตุ จ ฯเปฯ เทฺว ขนฺเธ ฯเปฯ. ปฏิสนฺธิ\-กฺขเณ ฯเปฯ. (3)}

\proseitem{10}{เหตุญฺจ นเหตุญฺจ ธมฺมํ ปฏิจฺจ เหตุ ธมฺโม อุปฺปชฺชติ นปุเรชาต\-ปจฺจยา. อรูเป อโลภญฺจ สมฺปยุตฺตเก จ ขนฺเธ ปฏิจฺจ อโทโส \csromanfolio{5}อโมโห. (จกฺกํ.) อรูเป โลภญฺจ สมฺปยุตฺตเก จ ขนฺเธ ปฏิจฺจ โมโห. (จกฺกํ.) ปฏิสนฺธิ\-กฺขเณ. (1)}

\prose{เหตุญฺจ นเหตุญฺจ ธมฺมํ ปฏิจฺจ นเหตุ ธมฺโม อุปฺปชฺชติ นปุเรชาต\-ปจฺจยา. อรูเป นเหตุํ เอกํ ขนฺธญฺจ เหตุญฺจ ปฏิจฺจ ตโย ขนฺธา ฯเปฯ เทฺว ขนฺเธ ฯเปฯ. นเหตู ขนฺเธ จ เหตุญฺจ ปฏิจฺจ จิตฺต\-สมุฏฺฐานํ รูปํ. เหตุญฺจ มหาภูเต จ ปฏิจฺจ จิตฺต\-สมุฏฺฐานํ รูปํ. ปฏิสนฺธิ\-กฺขเณ ฯเปฯ. (2)}

\prose{เหตุญฺจ นเหตุญฺจ ธมฺมํ ปฏิจฺจ เหตุ จ นเหตุ จ ธมฺมา อุปฺปชฺชนฺติ นปุเรชาต\-ปจฺจยา. อรูเป นเหตุํ เอกํ ขนฺธญฺจ อโลภญฺจ ปฏิจฺจ ตโย ขนฺธา อโทโส อโมโห จ ฯเปฯ เทฺว ขนฺเธ ฯเปฯ. (จกฺกํ.) นเหตุํ เอกํ ขนฺธญฺจ โลภญฺจ ปฏิจฺจ ตโย ขนฺธา โมโห จ. (จกฺกํ.) ปฏิสนฺธิ\-กฺขเณ ฯเปฯ. (3)}

\csromanheader{2}{\textbf{นปจฺฉาชาตาทิ}}
\proseitem{10}{เหตุํ ธมฺมํ ปฏิจฺจ เหตุ ธมฺโม อุปฺปชฺชติ นปจฺฉาชาต\-ปจฺจยา. นอาเสวน\-ปจฺจยา.}

\csromanheader{2}{\textbf{นกมฺมาทิ}}
\proseitem{10}{เหตุํ ธมฺมํ ปฏิจฺจ นเหตุ ธมฺโม อุปฺปชฺชติ นกมฺมปจฺจยา. เหตุํ ปฏิจฺจ สมฺปยุตฺตกา เจตนา. (1)}

\prose{นเหตุํ ธมฺมํ ปฏิจฺจ นเหตุ ธมฺโม อุปฺปชฺชติ นกมฺมปจฺจยา. นเหตู ขนฺเธ ปฏิจฺจ สมฺปยุตฺตกา เจตนา. พาหิรํ. อาหารสมุฏฺฐานํ. อุตุสมุฏฺฐานํ ฯเปฯ. (1)}

\prose{เหตุญฺจ นเหตุญฺจ ธมฺมํ ปฏิจฺจ นเหตุ ธมฺโม อุปฺปชฺชติ นกมฺมปจฺจยา. เหตุญฺจ สมฺปยุตฺตเก จ ขนฺเธ ปฏิจฺจ สมฺปยุตฺตกา เจตนา. (1)}

\prose{เหตุํ ธมฺมํ ปฏิจฺจ เหตุ ธมฺโม อุปฺปชฺชติ นวิปาก\-ปจฺจยา. นว.}

\prose{นอาหาราทิ}

\proseitem{13}{นเหตุํ ธมฺมํ ปฏิจฺจ นเหตุ ธมฺโม อุปฺปชฺชติ นอาหารปจฺจยา. พาหิรํ. อุตุสมุฏฺฐานํ. อสญฺญสตฺตานํ ฯเปฯ นอินฺทฺริยปจฺจยา. พาหิรํ. \csromanfolio{6}อาหารสมุฏฺฐานํ. อุตุสมุฏฺฐานํ ฯเปฯ. อสญฺญสตฺตานํ มหาภูเต ปฏิจฺจ รูปชีวิตินฺทฺริยํ. นฌานปจฺจยา. ปญฺจวิญฺญาณํ ฯเปฯ. พาหิรํ. อาหารสมุฏฺฐานํ. อุตุสมุฏฺฐานํ. อสญฺญสตฺตานํ ฯเปฯ น- มคฺคปจฺจยา. อเหตุกํ นเหตุํ เอกํ ขนฺธํ ปฏิจฺจ ฯเปฯ. พาหิรํ. อาหารสมุฏฺฐานํ. อุตุสมุฏฺฐานํ. อสญฺญสตฺตานํ ฯเปฯ น สมฺปยุตฺต\-ปจฺจยา. นวิปฺปยุตฺต\-ปจฺจยา. (นปุเรชาต\-สทิสํ. อรูปปญฺหาเยว.) โนนตฺถิปจฺจยา. โนวิคต\-ปจฺจยา.}

\csromanmarkh{2. ปจฺจย\-ปจฺจนีย}\csromanmarkhii{2. สงฺขฺยาวาร}\csromanheadernomark{2}{2. ปจฺจย\-ปจฺจนีย 2. สงฺขฺยาวาร}
\proseitem{14}{นเหตุยา เทฺว, นอารมฺมเณ ตีณิ, นอธิปติยา นว, นอนนฺตเร ตีณิ, นสมนนฺตเร ตีณิ, นอญฺญมญฺเญ ตีณิ, นอุปนิสฺสเย ตีณิ, นปุเรชาเต นว, นปจฺฉาชาเต นว, นอาเสวเน นว, นกมฺเม ตีณิ, นวิปาเก นว, นอาหาเร เอกํ, นอินฺทฺริเย เอกํ, นฌาเน เอกํ, นมคฺเค เอกํ, นสมฺปยุตฺเต ตีณิ, นวิปฺปยุตฺเต นว, โน นตฺถิยา ตีณิ, โน วิคเต ตีณิ. (เอวํ คเณตพฺพํ.)}

\vspace{\tipitakaparskip}
\csromancenter{ปจฺจนียํ.}
\proseitem{15}{เหตุปจฺจยา นอารมฺมเณ ตีณิ, นอธิปติยา นว, นอนนฺตเร ตีณิ ฯเปฯ นอุปนิสฺสเย ตีณิ, นปุเรชาเต นว, นปจฺฉาชาเต นว, นอาเสวเน นว, นกมฺเม ตีณิ, นวิปาเก นว, นสมฺปยุตฺเต ตีณิ, นวิปฺปยุตฺเต นว, โนนตฺถิยา ตีณิ, โนวิคเต ตีณิ.}

\vspace{\tipitakaparskip}
\csromancenter{อนุโลม\-ปจฺจนียํ.}
\proseitem{16}{\csromanfolio{7}นเหตุปจฺจยา อารมฺมเณ เทฺว, อนนฺตเร เทฺว ฯเปฯ กมฺเม เทฺว, วิปาเก เอกํ, อาหาเร เทฺว, อินฺทฺริเย เทฺว, ฌาเน เทฺว, มคฺเค เอกํ, สมฺปยุตฺเต เทฺว ฯเปฯ อวิคเต เทฺว. ปจฺจนียา\-นุโลมํ.}

\prose{(สหชาตวาโรปิ ปจฺจยวาโรปิ นิสฺสยวาโรปิ ปฏิจฺจ\-วาร\-สทิสาเยว ปญฺหา. มหาภูเตสุ นิฏฺฐิเตสุ “วตฺถุํ ปจฺจยา”ติ กาตพฺพา. ปญฺจายตนานิ อนุโลเมปิ ปจฺจนีเยปิ ยถา ลพฺภนฺติ, ตถา กาตพฺพา. สํสฏฺฐ\-วาโรปิ สมฺปยุตฺตวาโรปิ ปริปุณฺโณ. รูปํ นตฺถิ, อรูปเมว.)}

\csromanmarkh{1. เหตุทุก}\csromanmarkhii{7. ปญฺหาวาร}\csromanheadernomark{2}{1. เหตุทุก 7. ปญฺหาวาร}
\csromanmarkh{1. ปจฺจยานุโลม}\csromanmarkhii{1. วิภงฺควาร}\csromanheadernomark{2}{1. \textbf{ปจฺจยานุโลม 1. วิภงฺควาร}}
\csromanheader{2}{\textbf{เหตุ}}
\proseitem{17}{เหตุ ธมฺโม เหตุสฺส ธมฺมสฺส เหตุปจฺจเยน ปจฺจโย. อโลโภ อโทสสฺส อโมหสฺส เหตุปจฺจเยน ปจฺจโย. (จกฺกํ.) โลโภ โมหสฺส เหตุปจฺจเยน ปจฺจโย. โทโส โมหสฺส เหตุปจฺจเยน ปจฺจโย. ปฏิสนฺธิ\-กฺขเณ ฯเปฯ. (1)}

\prose{เหตุ ธมฺโม นเหตุสฺส ธมฺมสฺส เหตุปจฺจเยน ปจฺจโย. เหตู สมฺปยุตฺตกานํ ขนฺธานํ จิตฺตสมุฏฺฐานานญฺจ รูปานํ เหตุปจฺจเยน ปจฺจโย. ปฏิสนฺธิ\-กฺขเณ ฯเปฯ. (2)}

\prose{เหตุ ธมฺโม เหตุสฺส จ นเหตุสฺส จ ธมฺมสฺส เหตุปจฺจเยน ปจฺจโย. อโลโภ อโทสสฺส อโมหสฺส สมฺปยุตฺตกานญฺจ ขนฺธานํ จิตฺตสมุฏฺฐานานญฺจ รูปานํ เหตุปจฺจเยน ปจฺจโย. (จกฺกํ.) โลโภ โมหสฺส ฯเปฯ. ปฏิสนฺธิ\-กฺขเณ ฯเปฯ. (3)}

\proseitem{18}{\csromanfolio{8}เหตุ ธมฺโม เหตุสฺส ธมฺมสฺส อารมฺมณ\-ปจฺจเยน ปจฺจโย. เหตุํ อารพฺภ เหตู อุปฺปชฺชนฺติ. (1)}

\prose{เหตุ ธมฺโม นเหตุสฺส ธมฺมสฺส อารมฺมณ\-ปจฺจเยน ปจฺจโย. เหตุํ อารพฺภ นเหตู ขนฺธา อุปฺปชฺชนฺติ. (2)}

\prose{เหตุ ธมฺโม เหตุสฺส จ นเหตุสฺส จ ธมฺมสฺส อารมฺมณ\-ปจฺจเยน ปจฺจโย. เหตุํ อารพฺภ เหตู จ สมฺปยุตฺตกา จ ขนฺธา อุปฺปชฺชนฺติ. (3)}

\proseitem{19}{ธมฺโม นเหตุสฺส ธมฺมสฺส อารมฺมณ\-ปจฺจเยน ปจฺจโย. ทานํ ทตฺวา, สีลํ ฯเปฯ อุโปสถกมฺมํ กตฺวา ตํ ปจฺจเวกฺขติ. ปุพฺเพ สุจิณฺณานิ ปจฺจเวกฺขติ. ฌานา วุฏฺฐหิตฺวา ฯเปฯ. อริยา มคฺคา วุฏฺฐหิตฺวา มคฺคํ ปจฺจเวกฺขนฺติ, ผลํ ฯเปฯ นิพฺพานํ ฯเปฯ. นิพฺพานํ โคตฺรภุสฺส, โวทานสฺส มคฺคสฺส, ผลสฺส, อาวชฺชนาย อารมฺมณ\-ปจฺจเยน ปจฺจโย. อริยา นเหตู ปหีเน กิเลเส ปจฺจเวกฺขนฺติ. วิกฺขมฺภิ เต กิเลเส ปจฺจเวกฺขนฺติ. ปุพฺเพ สมุทาจิณฺเณ กิเลเส ชานนฺติ. จกฺขุํ ฯเปฯ วตฺถุํ, นเหตู ขนฺเธ อนิจฺจโต ฯเปฯ โทมนสฺสํ อุปฺปชฺชติ. ทิพฺเพน จกฺขุนา รูปํ ปสฺสติ. ทิพฺพาย โสตธาตุยา สทฺทํ สุณาติ. เจโตปริย\-ญาเณน นเหตุ\-จิตฺต\-สมงฺคิสฺส จิตฺตํ ชานาติ. อากาสา\-นญฺจา\-ยตนํ\csromansharedfootnote{91c702e073214d7c}{อากาสานญฺจายตนกิริสํ (สฺยา) เอวมุปริปิ ตีสุ ฐาเนสุ.} วิญฺญาณญฺจา\-ยตนสฺส ฯเปฯ. อากิญฺจญฺญา\-ยตนํ เนว\-สญฺญา\-นา\-สญฺญา\-ยตนสฺส ฯเปฯ. รูปายตนํ จกฺขุ\-วิญฺญาณสฺส ฯเปฯ. โผฏฺฐพฺพา\-ยตนํ กายวิญฺญาณสฺส ฯเปฯ. นเหตู ขนฺธา อิทฺธิวิธ\-ญาณสฺส, เจโตปริย\-ญาณสฺส, ปุพฺเพ\-นิวาสา\-นุสฺสติ\-ญาณสฺส, ยถากมฺมูปคญาณสฺส, อนาคตํสญาณสฺส อาวชฺชนาย อารมฺมณ\-ปจฺจเยน ปจฺจโย. (1)}

\prose{ธมฺโม เหตุสฺส ธมฺมสฺส อารมฺมณ\-ปจฺจเยน ปจฺจโย. ทานํ ทตฺวา. (ปฐมคมนํเยว. อาวชฺชนา นตฺถิ. “รูปายตนํ จกฺขุ\-วิญฺญาณสฺส ฯเปฯ. โผฏฺฐพฺพา\-ยตนํ กายวิญฺญาณสฺสา”ติ อิทํ นตฺถิ.)}

\prose{ธมฺโม เหตุสฺส จ นเหตุสฺส จ ธมฺมสฺส อารมฺมณ\-ปจฺจเยน ปจฺจโย. ทานํ ทตฺวา, สีลํ ฯเปฯ อุโปสถกมฺมํ กตฺวา ตํ ปจฺจเวกฺขติ. \csromanfolio{9}ตํ อารพฺภ เหตู จ สมฺปยุตฺตกา จ ขนฺธา อุปฺปชฺชนฺติ. (ตตฺถ ตตฺถ ฐิเตน อิมํ กาตพฺพํ. ทุติยคมน\-สทิสํ.) (3)}

\proseitem{20}{เหตุ จ นเหตุ จ ธมฺมา เหตุสฺส ธมฺมสฺส อารมฺมณ\-ปจฺจเยน ปจฺจโย. เหตุญฺจ สมฺปยุตฺตเก จ ขนฺเธ อารพฺภ เหตู อุปฺปชฺชนฺติ. (1)}

\prose{เหตุ จ นเหตุ จ ธมฺมา นเหตุสฺส ธมฺมสฺส อารมฺมณ\-ปจฺจเยน ปจฺจโย. เหตุญฺจ สมฺปยุตฺตเก จ ขนฺเธ อารพฺภ นเหตู ขนฺธา อุปฺปชฺชนฺติ. (2)}

\prose{เหตุ จ นเหตุ จ ธมฺมา เหตุสฺส จ นเหตุสฺส จ ธมฺมสฺส อารมฺมณ\-ปจฺจเยน ปจฺจโย. เหตุญฺจ สมฺปยุตฺตเก จ ขนฺเธ อารพฺภ เหตู จ สมฺปยุตฺตกา จ ขนฺธา อุปฺปชฺชนฺติ. (3)}

\proseitem{21}{เหตุ ธมฺโม เหตุสฺส ธมฺมสฺส อธิปติ\-ปจฺจเยน ปจฺจโย. อารมฺมณา\-ธิปติ, สหชาตา\-ธิปติ.  อารมฺมณา\-ธิปติ\csromandash{}เหตุํ ครุํ กตฺวา เหตู อุปฺปชฺชนฺติ. สหชาตา\-ธิปติ\csromandash{}เหตุ อธิปติ สมฺปยุตฺตกานํ เหตูนํ อธิปติ\-ปจฺจเยน ปจฺจโย. (1)}

\prose{เหตุ ธมฺโม นเหตุสฺส ธมฺมสฺส อธิปติ\-ปจฺจเยน ปจฺจโย. อารมฺมณา\-ธิปติ, สหชาตา\-ธิปติ.  อารมฺมณา\-ธิปติ\csromandash{}เหตุํ ครุํ กตฺวา นเหตู ขนฺธา อุปฺปชฺชนฺติ. สหชาตา\-ธิปติ\csromandash{}เหตุ อธิปติ สมฺปยุตฺตกานํ ขนฺธานํ จิตฺตสมุฏฺฐานานญฺจ รูปานํ อธิปติ\-ปจฺจเยน ปจฺจโย. (2)}

\prose{เหตุ ธมฺโม เหตุสฺส จ นเหตุสฺส จ ธมฺมสฺส อธิปติ\-ปจฺจเยน ปจฺจโย. อารมฺมณา\-ธิปติ, สหชาตา\-ธิปติ.  อารมฺมณา\-ธิปติ\csromandash{}เหตุํ ครุํ กตฺวา เหตู จ สมฺปยุตฺตกา จ ขนฺธา อุปฺปชฺชนฺติ. สหชาตา\-ธิปติ\csromandash{} เหตุ อธิปติ สมฺปยุตฺตกานํ ขนฺธานํ เหตูนญฺจ จิตฺตสมุฏฺฐานานญฺจ รูปานํ อธิปติ\-ปจฺจเยน ปจฺจโย. (3)}

\proseitem{22}{ธมฺโม นเหตุสฺส ธมฺมสฺส อธิปติ\-ปจฺจเยน ปจฺจโย. อารมฺมณา\-ธิปติ, สหชาตา\-ธิปติ.  อารมฺมณา\-ธิปติ\csromandash{}ทานํ ทตฺวา. (วิตฺถาเรตพฺพํ. ยาว นเหตู ขนฺธา.) สหชาตา\-ธิปติ\csromandash{}นเหตุ อธิปติ สมฺปยุตฺตกานํ ขนฺธานํ จิตฺตสมุฏฺฐานานญฺจ รูปานํ อธิปติ\-ปจฺจเยน ปจฺจโย. (1)}

\prose{\csromanfolio{10}ธมฺโม เหตุสฺส ธมฺมสฺส อธิปติ\-ปจฺจเยน ปจฺจโย. อารมฺมณา\-ธิปติ, สหชาตา\-ธิปติ.  อารมฺมณา\-ธิปติ\csromandash{}ทานํ ทตฺวา. (สํขิตฺตํ. ยาว วตฺถุ นเหตู จ ขนฺธา, ตาว กาตพฺพํ.) สหชาตา\-ธิปติ\csromandash{}นเหตุ อธิปติ สมฺปยุตฺตกานํ เหตูนํ อธิปติ\-ปจฺจเยน ปจฺจโย. (2)}

\prose{ธมฺโม เหตุสฺส จ นเหตุสฺส จ ธมฺมสฺส อธิปติ\-ปจฺจเยน ปจฺจโย. อารมฺมณา\-ธิปติ, สหชาตา\-ธิปติ.  อารมฺมณา\-ธิปติ\csromandash{}ทานํ ทตฺวา, สีลํ ฯเปฯ อุโปสถกมฺมํ กตฺวา ตํ ครุํ กตฺวา ปจฺจเวกฺขติ. ตํ ครุํ กตฺวา นเหตู ขนฺธา จ เหตู จ อุปฺปชฺชนฺติ. ปุพฺเพ สุจิณฺณานิ. (ยาว วตฺถุ นเหตู ขนฺธา จ, ตาว กาตพฺพํ.) สหชาตา\-ธิปติ\csromandash{}นเหตุ อธิปติ สมฺปยุตฺตกานํ ขนฺธานํ เหตูนญฺจ จิตฺตสมุฏฺฐานานญฺจ รูปานํ อธิปติ\-ปจฺจเยน ปจฺจโย. (3)}

\proseitem{23}{เหตุ จ นเหตุ จ ธมฺมา เหตุสฺส ธมฺมสฺส อธิปติ\-ปจฺจเยน ปจฺจโย. อารมฺมณา\-ธิปติ\csromandash{}เหตุญฺจ สมฺปยุตฺตเก จ ขนฺเธ ครุํ กตฺวา เหตู อุปฺปชฺชนฺติ. (1)}

\prose{เหตุ จ นเหตุ จ ธมฺมา นเหตุสฺส ธมฺมสฺส อธิปติ\-ปจฺจเยน ปจฺจโย. อารมฺมณา\-ธิปติ\csromandash{}เหตุญฺจ สมฺปยุตฺตเก จ ขนฺเธ ครุํ กตฺวา นเหตู ขนฺธา อุปฺปชฺชนฺติ. (2)}

\prose{เหตุ จ นเหตุ จ ธมฺมา เหตุสฺส จ นเหตุสฺส จ ธมฺมสฺส อธิปติ\-ปจฺจเยน ปจฺจโย. อารมฺมณา\-ธิปติ\csromandash{}เหตุญฺจ สมฺปยุตฺตเก จ ขนฺเธ ครุํ กตฺวา เหตู จ สมฺปยุตฺตกา จ ขนฺธา อุปฺปชฺชนฺติ. (3)}

\proseitem{22}{เหตุ ธมฺโม เหตุสฺส ธมฺมสฺส อนนฺตร\-ปจฺจเยน ปจฺจโย. ปุริมา ปุริมา เหตู ปจฺฉิมานํ ปจฺฉิมานํ เหตูนํ อนนฺตร\-ปจฺจเยน ปจฺจโย. (1)}

\prose{เหตุ ธมฺโม นเหตุสฺส ธมฺมสฺส อนนฺตร\-ปจฺจเยน ปจฺจโย. ปุริมา ปุริมา เหตู ปจฺฉิมานํ ปจฺฉิมานํ นเหตูนํ ขนฺธานํ อนนฺตร\-ปจฺจเยน ปจฺจโย. (2)}

\prose{เหตุ ธมฺโม เหตุสฺส จ นเหตุสฺส จ ธมฺมสฺส อนนฺตร\-ปจฺจเยน ปจฺจโย. ปุริมา ปุริมา เหตู ปจฺฉิมานํ ปจฺฉิมานํ เหตูนํ สมฺปยุตฺตกานญฺจ ขนฺธานํ อนนฺตร\-ปจฺจเยน ปจฺจโย. (3)}

\proseitem{25}{\csromanfolio{11}ธมฺโม นเหตุสฺส ธมฺมสฺส อนนฺตร\-ปจฺจเยน ปจฺจโย. ปุริมา ปุริมา นเหตู ขนฺธา ปจฺฉิมานํ ปจฺฉิมานํ นเหตูนํ ขนฺธานํ อนนฺตร\-ปจฺจเยน ปจฺจโย. อนุโลมํ โคตฺรภุสฺส. (สํขิตฺตํ.) เนว\-สญฺญา\-นา\-สญฺญา\-ยตนํ ผลสมาปตฺติยา อนนฺตร\-ปจฺจเยน ปจฺจโย. (1)}

\prose{ธมฺโม เหตุสฺส ธมฺมสฺส อนนฺตร\-ปจฺจเยน ปจฺจโย ฯเปฯ.}

\prose{ธมฺโม เหตุสฺส จ นเหตุสฺส จ ธมฺมสฺส อนนฺตร\-ปจฺจเยน ปจฺจโย. (นเหตุมูลกํ ตีณิปิ เอกสทิสํ.) (3)}

\prose{เหตู จ นเหตู จ ธมฺมา เหตุสฺส ธมฺมสฺส อนนฺตร\-ปจฺจเยน ปจฺจโย. ปุริมา ปุริมา เหตู จ สมฺปยุตฺตกา จ ขนฺธา ปจฺฉิมานํ ปจฺฉิมานํ เหตูนํ อนนฺตร\-ปจฺจเยน ปจฺจโย. (1)}

\prose{เหตู จ นเหตู จ ธมฺมา นเหตุสฺส ธมฺมสฺส อนนฺตร\-ปจฺจเยน ปจฺจโย. ปุริมา ปุริมา เหตู จ สมฺปยุตฺตกา จ ขนฺธา ปจฺฉิมานํ ปจฺฉิมานํ นเหตูนํ ขนฺธานํ อนนฺตร\-ปจฺจเยน ปจฺจโย. (2)}

\prose{เหตู จ นเหตู จ ธมฺมา เหตุสฺส จ นเหตุสฺส จ ธมฺมสฺส อนนฺตร\-ปจฺจเยน ปจฺจโย. ปุริมา ปุริมา เหตู จ สมฺปยุตฺตกา จ ขนฺธา ปจฺฉิมานํ ปจฺฉิมานํ เหตูนํ สมฺปยุตฺตกานญฺจ ขนฺธานํ อนนฺตร\-ปจฺจเยน ปจฺจโย. (3)}

\csromanheader{2}{\textbf{สมนนฺตราทิ}}
\proseitem{27}{เหตุ ธมฺโม เหตุสฺส ธมฺมสฺส สมนนฺตร\-ปจฺจเยน ปจฺจโย. (อนนฺตร\-สทิสํ.) สหชาต\-ปจฺจเยน ปจฺจโย. อญฺญมญฺญ\-ปจฺจเยน ปจฺจโย. (อิเม เทฺวปิ ปฏิจฺจสทิสา. นิสฺสยปจฺจโย ปจฺจยวาเร นิสฺสยปจฺจย\-สทิโส.)}

\csromanheader{2}{\textbf{อุปนิสฺสย}}
\proseitem{25}{เหตุ ธมฺโม เหตุสฺส ธมฺมสฺส อุปนิสฺสย\-ปจฺจเยน ปจฺจโย. อารมฺมณู\-ปนิสฺสโย, อนนฺตรู\-ปนิสฺสโย, ปกตู\-ปนิสฺสโย ฯเปฯ. ปกตู\-ปนิสฺสโย\csromandash{}เหตู เหตูนํ อุปนิสฺสย\-ปจฺจเยน ปจฺจโย. (1)}

\prose{เหตุ ธมฺโม นเหตุสฺส ธมฺมสฺส อุปนิสฺสย\-ปจฺจเยน ปจฺจโย. อารมฺมณู\-ปนิสฺสโย, อนนฺตรู\-ปนิสฺสโย, ปกตู\-ปนิสฺสโย ฯเปฯ. ปกตู\-ปนิสฺสโย\csromandash{}เหตู นเหตูนํ ขนฺธานํ อุปนิสฺสย\-ปจฺจเยน ปจฺจโย. (2)}

\proseitem{28}{\csromanfolio{12}เหตุ ธมฺโม เหตุสฺส จ นเหตุสฺส จ ธมฺมสฺส อุปนิสฺสย\-ปจฺจเยน ปจฺจโย. อารมฺมณู\-ปนิสฺสโย, อนนฺตรู\-ปนิสฺสโย, ปกตู\-ปนิสฺสโย ฯเปฯ. ปกตู\-ปนิสฺสโย\csromandash{}เหตู เหตูนํ สมฺปยุตฺตกานญฺจ ขนฺธานํ อุปนิสฺสย\-ปจฺจเยน ปจฺจโย. (3)}

\proseitem{29}{ธมฺโม นเหตุสฺส ธมฺมสฺส อุปนิสฺสย\-ปจฺจเยน ปจฺจโย. อารมฺมณู\-ปนิสฺสโย, อนนฺตรู\-ปนิสฺสโย, ปกตู\-ปนิสฺสโย ฯเปฯ. ปกตู\-ปนิสฺสโย\csromandash{}สทฺธํ อุปนิสฺสาย ทานํ เทติ ฯเปฯ สมาปตฺติํ อุปฺปาเทติ ฯเปฯ ทิฏฺฐิํ คณฺหาติ. สีลํ ฯเปฯ เสนาสนํ อุปนิสฺสาย ทานํ เทติ ฯเปฯ สํฆํ ภินฺทติ. สทฺธา ฯเปฯ เสนาสนํ สทฺธาย ฯเปฯ ผลสมาปตฺถิยา อุปนิสฺสย\-ปจฺจเยน ปจฺจโย. (1)}

\prose{ธมฺโม เหตุสฺส ธมฺมสฺส อุปนิสฺสย\-ปจฺจเยน ปจฺจโย. อารมฺมณู\-ปนิสฺสโย, อนนฺตรู\-ปนิสฺสโย, ปกตู\-ปนิสฺสโย ฯเปฯ. ปกตู\-ปนิสฺสโย\csromandash{}สทฺธํ ฯเปฯ เสนาสนํ อุปนิสฺสาย ทานํ เทติ ฯเปฯ สํฆํ ภินฺทติ. สทฺธา ฯเปฯ เสนาสนํ สทฺธาย ฯเปฯ ปตฺถนาย มคฺคสฺส ผลสมาปตฺติยา อุปนิสฺสย\-ปจฺจเยน ปจฺจโย. (2)}

\prose{ธมฺโม เหตุสฺส จ นเหตุสฺส จ ธมฺมสฺส อุปนิสฺสย\-ปจฺจเยน ปจฺจโย. อารมฺมณู\-ปนิสฺสโย, อนนฺตรู\-ปนิสฺสโย, ปกตู\-ปนิสฺสโย ฯเปฯ. ปกตู\-ปนิสฺสโย. (ทุติยฺ\-เอาปนิสฺสย\-สทิสํ.) (3)}

\proseitem{30}{เหตู จ นเหตู จ ธมฺมา เหตุสฺส ธมฺมสฺส อุปนิสฺสย\-ปจฺจเยน ปจฺจโย. อารมฺมณู\-ปนิสฺสโย, อนนฺตรู\-ปนิสฺสโย, ปกตู\-ปนิสฺสโย ฯเปฯ. ปกตู\-ปนิสฺสโย\csromandash{}เหตู จ สมฺปยุตฺตกา จ ขนฺธา เหตูนํ อุปนิสฺสย\-ปจฺจเยน ปจฺจโย. (1)}

\prose{เหตู จ นเหตู จ ธมฺมา นเหตุสฺส ธมฺมสฺส อุปนิสฺสย\-ปจฺจเยน ปจฺจโย. อารมฺมณู\-ปนิสฺสโย, อนนฺตรู\-ปนิสฺสโย, ปกตู\-ปนิสฺสโย ฯเปฯ. ปกตู\-ปนิสฺสโย\csromandash{}เหตู จ สมฺปยุตฺตกา จ ขนฺธา เนเหตูนํ ขนฺธานํ อุปนิสฺสย\-ปจฺจเยน ปจฺจโย. (2)}

\prose{เหตู จ นเหตู จ ธมฺมา เหตุสฺส จ นเหตุสฺส จ ธมฺมสฺส อุปนิสฺสย\-ปจฺจเยน ปจฺจโย. อารมฺมณู\-ปนิสฺสโย, อนนฺตรู\-ปนิสฺสโย, \csromanfolio{13}ปกตู\-ปนิสฺสโย ฯเปฯ. ปกตู\-ปนิสฺสโย\csromandash{}เหตู จ สมฺปยุตฺตกา จ ขนฺธา เหตูนญฺจ สมฺปยุตฺตกานญฺจ ขนฺธานํ อุปนิสฺสย\-ปจฺจเยน ปจฺจโย. (3)}

\proseitem{31}{ธมฺโม นเหตุสฺส ธมฺมสฺส ปุเรชาต\-ปจฺจเยน ปจฺจโย. อารมฺมณ\-ปุเรชาตํ, วตฺถุ\-ปุเรชาตํ.  อารมฺมณ\-ปุเรชาตํ\csromandash{}จกฺขุํ ฯเปฯ วตฺถุํ อนิจฺจโต ฯเปฯ โทมนสฺสํ อุปฺปชฺชติ. ทิพฺเพน จกฺขุนา รูปํ ปสฺสติ. ทิพฺพาย โสตธาตุยา สทฺทํ สุณาติ. รูปายตนํ จกฺขุ\-วิญฺญาณสฺส ฯเปฯ. โผฏฺฐพฺพา\-ยตนํ กายวิญฺญาณสฺส ฯเปฯ. วตฺถุ\-ปุเรชาตํ\csromandash{}จกฺขายตนํ ฯเปฯ กายายตนํ ฯเปฯ. วตฺถุ นเหตูนํ ขนฺธานํ ปุเรชาต\-ปจฺจเยน ปจฺจโย. (1)}

\prose{ธมฺโม เหตุสฺส ธมฺมสฺส ปุเรชาต\-ปจฺจเยน ปจฺจโย. อารมฺมณ\-ปุเรชาตํ, วตฺถุ\-ปุเรชาตํ.  อารมฺมณ\-ปุเรชาตํ\csromandash{}จกฺขุํ ฯเปฯ วตฺถุํ อนิจฺจโต ฯเปฯ โทมนสฺสํ อุปฺปชฺชติ. ทิพฺเพน จกฺขุนา รูปํ ปสฺสติ. ทิพฺพาย โสตธาตุยา สทฺทํ สุณาติ. วตฺถุ\-ปุเรชาตํ\csromandash{}วตฺถุ เหตูนํ ปุเรชาต\-ปจฺจเยน ปจฺจโย. (2)}

\prose{ธมฺโม เหตุสฺส จ นเหตุสฺส จ ธมฺมสฺส ปุเรชาต\-ปจฺจเยน ปจฺจโย. อารมฺมณ\-ปุเรชาตํ, วตฺถุ\-ปุเรชาตํ.  อารมฺมณ\-ปุเรชาตํ\csromandash{} จกฺขุํ ฯเปฯ วตฺถุํ อนิจฺจโต ฯเปฯ โทมนสฺสํ อุปฺปชฺชติ. วตฺถุ\-ปุเรชาตํ\csromandash{}วตฺถุ เหตูนํ สมฺปยุตฺตกานญฺจ ขนฺธานํ ปุเรชาต\-ปจฺจเยน ปจฺจโย. (3)}

\csromanheader{2}{\textbf{ปจฺฉาชาตาทิ}}
\proseitem{32}{เหตุ ธมฺโม นเหตุสฺส ธมฺมสฺส ปจฺฉาชาต\-ปจฺจเยน ปจฺจโย. ปจฺฉาชาตา เหตู ปุเรชาตสฺส อิมสฺส กายสฺส ปจฺฉาชาต\-ปจฺจเยน ปจฺจโย. (1)}

\prose{ธมฺโม นเหตุสฺส ธมฺมสฺส ปจฺฉาชาต\-ปจฺจเยน ปจฺจโย. ปจฺฉาชาตา นเหตู ขนฺธา ปุเรชาตสฺส อิมสฺส กายสฺส ปจฺฉาชาต\-ปจฺจเยน ปจฺจโย. (1)}

\prose{เหตู จ นเหตู จ ธมฺมา นเหตุสฺส ธมฺมสฺส ปจฺฉาชาต\-ปจฺจเยน. ปจฺฉาชาตา เหตู จ สมฺปยุตฺตกา จ ขนฺธา ปุเรชาตสฺส อิมสฺส กายสฺส ปจฺฉาชาต\-ปจฺจเยน ปจฺจโย. (1)}

\prose{\csromanfolio{14}เหตุ ธมฺโม เหตุสฺส ธมฺมสฺส อาเสวน\-ปจฺจเยน ปจฺจโย. (อนนฺตร\-สทิสํ.)}

\proseitem{33}{ธมฺโม นเหตุสฺส ธมฺมสฺส กมฺมปจฺจเยน ปจฺจโย. สหชาตา, นานากฺขณิกา.  สหชาตา นเหตุ เจตนา สมฺปยุตฺตกานํ ขนฺธานํ จิตฺตสมุฏฺฐานานญฺจ รูปานํ กมฺมปจฺจเยน ปจฺจโย. ปฏิสนฺธิ\-กฺขเณ ฯเปฯ. นานากฺขณิกา นเหตุ เจตนา วิปากานํ ขนฺธานํ กฏตฺตา จ รูปานํ กมฺมปจฺจเยน ปจฺจโย. (1)}

\prose{ธมฺโม เหตุสฺส ธมฺมสฺส กมฺมปจฺจเยน ปจฺจโย. สหชาตา, นานากฺขณิกา.  สหชาตา นเหตุ เจตนา สมฺปยุตฺตกานํ เหตูนํ กมฺมปจฺจเยน ปจฺจโย. ปฏิสนฺธิ\-กฺขเณ ฯเปฯ. นานากฺขณิกา นเหตุ เจตนา วิปากานํ เหตูนํ กมฺมปจฺจเยน ปจฺจโย. (2)}

\prose{ธมฺโม เหตุสฺส จ นเหตุสฺส จ ธมฺมสฺส กมฺมปจฺจเยน ปจฺจโย. สหชาตา, นานากฺขณิกา.  สหชาตา นเหตุ เจตนา สมฺปยุตฺตกานํ ขนฺธานํ เหตูนํ จิตฺตสมุฏฺฐานานญฺจ รูปานํ กมฺมปจฺจเยน ปจฺจโย. ปฏิสนฺธิ\-กฺขเณ ฯเปฯ. นานากฺขณิกา นเหตุ เจตนา วิปากานํ ขนฺธานํ เหตูนํ กฏตฺตา จ รูปานํ กมฺมปจฺจเยน ปจฺจโย. (3)}

\csromanheader{2}{\textbf{วิปาก}}
\proseitem{34}{เหตุ ธมฺโม เหตุสฺส ธมฺมสฺส วิปาก\-ปจฺจเยน ปจฺจโย. วิปาโก อโลโภ อโทสสฺส อโมหสฺส วิปาก\-ปจฺจเยน ปจฺจโย. (ปฏิจฺจ\-วาร\-สทิสํ. วิปากวิภงฺเค นว ปญฺหา.)}

\csromanheader{2}{\textbf{อาหาร}}
\proseitem{35}{ธมฺโม นเหตุสฺส ธมฺมสฺส อาหารปจฺจเยน ปจฺจโย. นเหตู อาหารา สมฺปยุตฺตกานํ ขนฺธานํ จิตฺตสมุฏฺฐานานญฺจ รูปานํ อาหารปจฺจเยน ปจฺจโย. ปฏิสนฺธิ\-กฺขเณ ฯเปฯ. กพฬีกาโร อาหาโร อิมสฺส กายสฺส อาหารปจฺจเยน ปจฺจโย. (1)}

\prose{ธมฺโม เหตุสฺส ธมฺมสฺส อาหารปจฺจเยน ปจฺจโย. นเหตู อาหารา สมฺปยุตฺตกานํ เหตูนํ อาหารปจฺจเยน ปจฺจโย. ปฏิสนฺธิ\-กฺขเณ ฯเปฯ. (2)}

\prose{\csromanfolio{15}ธมฺโม เหตุสฺส จ น เหตุสฺส จ ธมฺมสฺส อาหารปจฺจเยน ปจฺจโย. นเหตู อาหารา สมฺปยุตฺตกานํ ขนฺธานํ, เหตูนํ จิตฺตสมุฏฺฐานานญฺจ รูปานํ อาหารปจฺจเยน ปจฺจโย. ปฏิสนฺธิ\-กฺขเณ ฯเปฯ. (3)}

\csromanheader{2}{\textbf{อินฺทฺริย}}
\proseitem{36}{เหตุ ธมฺโม เหตุสฺส ธมฺมสฺส อินฺทฺริย\-ปจฺจเยน ปจฺจโย ฯเปฯ. (เหตุมูลเก ตีณิ)}

\prose{ธมฺโม นเหตุสฺส ธมฺมสฺส อินฺทฺริย\-ปจฺจเยน ปจฺจโย. นเหตู อินฺทฺริยา สมฺปยุตฺตกานํ ขนฺธานํ จิตฺตสมุฏฺฐานานญฺจ รูปานํ อินฺทฺริย\-ปจฺจเยน ปจฺจโย. ปฏิสนฺธิ\-กฺขเณ ฯเปฯ. จกฺขุนฺทฺริยํ จกฺขุ\-วิญฺญาณสฺส ฯเปฯ. กายินฺทฺริยํ กายวิญฺญาณสฺส ฯเปฯ. รูปชีวิตินฺทฺริยํ กฏตฺตารูปานํ อินฺทฺริย\-ปจฺจเยน ปจฺจโย. (เอวํ อินฺทฺริยปจฺจยา วิตฺถาเรตพฺพา. นว.)}

\csromanheader{2}{\textbf{ฌานาทิ}}
\proseitem{37}{ธมฺโม นเหตุสฺส ธมฺมสฺส ฌานปจฺจเยน ปจฺจโย. ตีณิ.}

\prose{เหตุ ธมฺโม เหตุสฺส ธมฺมสฺส มคฺคปจฺจเยน ปจฺจโย. สมฺปยุตฺต\-ปจฺจเยน ปจฺจโย. (อิเมสุ ทฺวีสุ นว.)}

\proseitem{38}{เหตุ ธมฺโม นเหตุสฺส ธมฺมสฺส วิปฺปยุตฺต\-ปจฺจเยน ปจฺจโย. สหชาตํ, ปจฺฉาชาตํ.  สหชาตา เหตู จิตฺต\-สมุฏฺฐานานํ รูปานํ วิปฺปยุตฺต\-ปจฺจเยน ปจฺจโย. ปฏิสนฺธิ\-กฺขเณ เหตู กฏตฺตารูปานํ วิปฺปยุตฺต\-ปจฺจเยน ปจฺจโย. เหตู วตฺถุสฺส วิปฺปยุตฺต\-ปจฺจเยน ปจฺจโย. ปจฺฉาชาตา เหตู ปุเรชาตสฺส อิมสฺส กายสฺส วิปฺปยุตฺต\-ปจฺจเยน ปจฺจโย. (1)}

\prose{ธมฺโม นเหตุสฺส ธมฺมสฺส วิปฺปยุตฺต\-ปจฺจเยน ปจฺจโย. สหชาตํ, ปุเรชาตํ, ปจฺฉาชาตํ.  สหชาตา นเหตู ขนฺธา จิตฺต\-สมุฏฺฐานานํ รูปานํ วิปฺปยุตฺต\-ปจฺจเยน ปจฺจโย. ปฏิสนฺธิ\-กฺขเณ นเหตู ขนฺธา กฏตฺตารูปานํ วิปฺปยุตฺต\-ปจฺจเยน ปจฺจโย. ขนฺธา วตฺถุสฺส ฯเปฯ วตฺถุ ขนฺธานํ วิปฺปยุตฺต\-ปจฺจเยน ปจฺจโย. ปุเรชาตํ จกฺขายตนํ จกฺขุ\-วิญฺญาณสฺส ฯเปฯ. กายายตนํ กายวิญฺญาณสฺส. วตฺถุ นเหตูนํ ขนฺธานํ วิปฺปยุตฺต\-ปจฺจเยน \csromanfolio{16}ปจฺจโย. ปจฺฉาชาตา นเหตู ขนฺธา ปุเรชาตสฺส อิมสฺส กายสฺส วิปฺปยุตฺต\-ปจฺจเยน ปจฺจโย. (1)}

\prose{ธมฺโม เหตุสฺส ธมฺมสฺส วิปฺปยุตฺต\-ปจฺจเยน ปจฺจโย. สหชาตํ, ปุเรชาตํ.  สหชาตํ ปฏิสนฺธิ\-กฺขเณ วตฺถุ เหตูนํ วิปฺปยุตฺต\-ปจฺจเยน ปจฺจโย. ปุเรชาตํ วตฺถุ เหตูนํ วิปฺปยุตฺต\-ปจฺจเยน ปจฺจโย. (2)}

\prose{ธมฺโม เหตุสฺส จ นเหตุสฺส จ ธมฺมสฺส วิปฺปยุตฺต\-ปจฺจเยน ปจฺจโย. สหชาตํ, ปุเรชาตํ.  สหชาตํ ปฏิสนฺธิ\-กฺขเณ วตฺถุ เหตูนํ สมฺปยุตฺตกานญฺจ ขนฺธานํ วิปฺปยุตฺต\-ปจฺจเยน ปจฺจโย. ปุเรชาตํ วตฺถุ เหตูนํ สมฺปยุตฺตกานญฺจ ขนฺธานํ วิปฺปยุตฺต\-ปจฺจเยน ปจฺจโย. (3)}

\prose{เหตู จ นเหตู จ ธมฺมา นเหตุสฺส ธมฺมสฺส วิปฺปยุตฺต\-ปจฺจเยน ปจฺจโย. สหชาตํ, ปจฺฉาชาตํ.  สหชาตา เหตู จ สมฺปยุตฺตกา จ ขนฺธา จิตฺต\-สมุฏฺฐานานํ รูปานํ วิปฺปยุตฺต\-ปจฺจเยน ปจฺจโย. ปฏิสนฺธิ\-กฺขเณ เหตู จ สมฺปยุตฺตกา จ ขนฺธา กฏตฺตารูปานํ วิปฺปยุตฺต\-ปจฺจเยน ปจฺจโย. ปจฺฉาชาตา เหตู จ สมฺปยุตฺตกา จ ขนฺธา ปุเรชาตสฺส อิมสฺส กายสฺส วิปฺปยุตฺต\-ปจฺจเยน ปจฺจโย. (1)}

\prose{อตฺถิอาทิ}

\proseitem{39}{เหตุ ธมฺโม เหตุสฺส ธมฺมสฺส อตฺถิปจฺจเยน ปจฺจโย. อโลโภ อโทสสฺส อโมหสฺส อตฺถิปจฺจเยน ปจฺจโย. (จกฺกํ.) โลโภ โมหสฺส อตฺถิปจฺจเยน ปจฺจโย. (จกฺกํ.) ปฏิสนฺธิ\-กฺขเณ ฯเปฯ. (1)}

\prose{เหตุ ธมฺโม นเหตุสฺส ธมฺมสฺส อตฺถิปจฺจเยน ปจฺจโย. สหชาตํ, ปจฺฉาชาตํ.  สหชาตา เหตู สมฺปยุตฺตกานํ ขนฺธานํ จิตฺตสมุฏฺฐานานญฺจ รูปานํ อตฺถิปจฺจเยน ปจฺจโย. ปฏิสนฺธิ\-กฺขเณ ฯเปฯ. ปจฺฉาชาตา เหตู ปุเรชาตสฺส อิมสฺส กายสฺส อตฺถิปจฺจเยน ปจฺจโย.}

\prose{เหตุ ธมฺโม เหตุสฺส จ นเหตุสฺส จ ธมฺมสฺส อตฺถิปจฺจเยน ปจฺจโย. อโลโภ อโทสสฺส อโมหสฺส สมฺปยุตฺตกานํ ขนฺธานํ จิตฺตสมุฏฺฐานานญฺจ รูปานํ อตฺถิปจฺจเยน ปจฺจโย. (จกฺกํ) โลโภ โมหสฺส สมฺปยุตฺตกานํ ขนฺธานํ จิตฺตสมุฏฺฐานานญฺจ รูปานํ อตฺถิปจฺจเยน ปจฺจโย. (จกฺกํ.) ปฏิสนฺธิ\-กฺขเณ ฯเปฯ. (3)}

\proseitem{40}{\csromanfolio{17}ธมฺโม นเหตุสฺส ธมฺมสฺส อตฺถิปจฺจเยน ปจฺจโย. สหชาตํ, ปุเรชาตํ, ปจฺฉาชาตํ, อาหารํ, อินฺทฺริยํ.  สหชาโต นเหตุ เอโก ขนฺโธ ติณฺณนฺนํ ขนฺธานํ จิตฺตสมุฏฺฐานานญฺจ รูปานํ อตฺถิปจฺจเยน ปจฺจโย ฯเปฯ เทฺว ขนฺธา ฯเปฯ. ปฏิสนฺธิ\-กฺขเณ นเหตุ เอโก ขนฺโธ ติณฺณนฺนํ ขนฺธานํ กฏตฺตา จ รูปานํ ฯเปฯ. ขนฺธา วตฺถุสฺส อตฺถิปจฺจเยน ปจฺจโย. วตฺถุ ขนฺธานํ อตฺถิปจฺจเยน ปจฺจโย. เอกํ มหาภูตํ ฯเปฯ. พาหิรํ. อาหารสมุฏฺฐานํ. อุตุสมุฏฺฐานํ. อสญฺญสตฺตานํ ฯเปฯ. ปุเรชาตํ จกฺขุํ ฯเปฯ. วตฺถุํ ฯเปฯ. ทิพฺเพน จกฺขุนา รูปํ ปสฺสติ. ทิพฺพาย โสตธาตุยา สทฺทํ สุณาติ. รูปายตนํ จกฺขุ\-วิญฺญาณสฺส ฯเปฯ. โผฏฺฐพฺพา\-ยตนํ กายวิญฺญาณสฺส. จกฺขายตนํ จกฺขุ\-วิญฺญาณสฺส ฯเปฯ. กายายตนํ กายวิญฺญาณสฺส. วตฺถุ นเหตูนํ ขนฺธานํ อตฺถิปจฺจเยน ปจฺจโย. ปจฺฉาชาตา นเหตู ขนฺธา ปุเรชาตสฺส อิมสฺส กายสฺส. อตฺถิปจฺจเยน ปจฺจโย. กพฬีกาโร อาหาโร อิมสฺส กายสฺส อตฺถิปจฺจเยน ปจฺจโย. รูปชีวิตินฺทฺริยํ กฏตฺตารูปานํ อตฺถิปจฺจเยน ปจฺจโย. (1)}

\prose{ธมฺโม เหตุสฺส ธมฺมสฺส อตฺถิปจฺจเยน ปจฺจโย. สหชาตํ, ปุเรชาตํ.  สหชาตา นเหตู ขนฺธา สมฺปยุตฺตกานํ เหตูนํ อตฺถิปจฺจเยน ปจฺจโย. ปฏิสนฺธิ\-กฺขเณ ฯเปฯ. วตฺถุ เหตูนํ อตฺถิปจฺจเยน ปจฺจโย. ปุเรชาตํ จกฺขุํ ฯเปฯ วตฺถุํ อนิจฺจโต ฯเปฯ โทมนสฺสํ อุปฺปชฺชติ. ทิพฺเพน จกฺขุนา รูปํ ปสฺสติ. ทิพฺพาย โสตธาตุยา สทฺทํ สุณาติ. วตฺถุ เหตูนํ อตฺถิปจฺจเยน ปจฺจโย. (2)}

\prose{ธมฺโม เหตุสฺส จ นเหตุสฺส จ ธมฺมสฺส อตฺถิปจฺจเยน ปจฺจโย. สหชาตํ, ปุเรชาตํ.  สหชาโต นเหตุ เอโก ขนฺโธ ติณฺณนฺนํ ขนฺธานํ เหตูนํ จิตฺตสมุฏฺฐานานญฺจ รูปานํ อตฺถิปจฺจเยน ปจฺจโย. ปฏิสนฺธิ\-กฺขเณ ฯเปฯ. วตฺถุ เหตูนํ สมฺปยุตฺตกานญฺจ ขนฺธานํ อตฺถิปจฺจเยน ปจฺจโย. ปุเรชาตํ จกฺขุํ ฯเปฯ วตฺถุํ อนิจฺจโต ฯเปฯ โทมนสฺสํ อุปฺปชฺชติ. ทิพฺเพน จกฺขุนา รูปํ ปสฺสติ. ทิพฺพาย โสตธาตุยา สทฺทํ สุณาติ. วตฺถุ เหตูนํ สมฺปยุตฺตกานญฺจ ขนฺธานํ อตฺถิปจฺจเยน ปจฺจโย. (3)}

\proseitem{41}{เหตุ จ นเหตุ จ ธมฺมา เหตุสฺส ธมฺมสฺส อตฺถิปจฺจเยน ปจฺจโย. \textbf{สหชาตํ, ปุเรชาตํ}.  \textbf{สหชาโต} อโลโภ จ สมฺปยุตฺตกา จ ขนฺธา อโทสสฺส อโมหสฺส อตฺถิปจฺจเยน ปจฺจโย. \csromanfolio{18}(จกฺกํ.) โลโภ จ สมฺปยุตฺตกา จ ขนฺธา โมหสฺส อตฺถิปจฺจเยน ปจฺจโย. (จกฺกํ.) ปฏิสนฺธิ\-กฺขเณ ฯเปฯ. อโลโภ จ วตฺถุ จ อโทสสฺส อโมหสฺส อตฺถิปจฺจเยน ปจฺจโย. (จกฺกํ.) (1)}

\prose{เหตุ จ นเหตุ จ ธมฺมา นเหตุสฺส ธมฺมสฺส อตฺถิปจฺจเยน ปจฺจโย. สหชาตํ, ปจฺฉาชาตํ, อาหารํ, อินฺทฺริยํ.  สหชาโต นเหตุ เอโก ขนฺโธ จ เหตู จ ติณฺณนฺนํ ขนฺธานํ จิตฺตสมุฏฺฐานานญฺจ รูปานํ อตฺถิปจฺจเยน ปจฺจโย ฯเปฯ เทฺว ขนฺธา ฯเปฯ. ปฏิสนฺธิ\-กฺขเณ ฯเปฯ. ปฏิสนฺธิ\-กฺขเณ เหตู จ วตฺถุ จ นเหตูนํ ขนฺธานํ อตฺถิปจฺจเยน ปจฺจโย. สหชาตา เหตู จ มหาภูตา จ จิตฺต\-สมุฏฺฐานานํ รูปานํ อตฺถิปจฺจเยน ปจฺจโย. ปจฺฉาชาตา เหตู จ กพฬีกาโร อาหาโร จ อิมสฺส กายสฺส อตฺถิปจฺจเยน ปจฺจโย. ปจฺฉาชาตา เหตู จ รูปชีวิตินฺทฺริยญฺจ กฏตฺตารูปานํ อตฺถิปจฺจเยน ปจฺจโย. (2)}

\prose{เหตุ จ นเหตุ จ ธมฺมา เหตุสฺส จ นเหตุสฺส จ ธมฺมสฺส อตฺถิปจฺจเยน ปจฺจโย. สหชาตํ, ปุเรชาตํ.  สหชาโต นเหตุ เอโก ขนฺโธ จ อโลโภ จ ติณฺณนฺนํ ขนฺธานํ อโทสสฺส อโมหสฺส จิตฺตสมุฏฺฐานานญฺจ รูปานํ อตฺถิปจฺจเยน ปจฺจโย. (จกฺกํ.) นเหตุ เอโก ขนฺโธ จ โลโภ จ ติณฺณนฺนํ ขนฺธานํ โมหสฺส จิตฺตสมุฏฺฐานานญฺจ รูปานํ อตฺถิปจฺจเยน ปจฺจโย. (จกฺกํ.) ปฏิสนฺธิ\-กฺขเณ นเหตุ เอโก ขนฺโธ จ อโลโภ จ. (จกฺกํ.) ปฏิสนฺธิ\-กฺขเณ ฯเปฯ. อโลโภ จ วตฺถุ จ อโทสสฺส อโมหสฺส สมฺปยุตฺตกานญฺจ ขนฺธานํ อตฺถิปจฺจเยน ปจฺจโย. โลโภ จ วตฺถุ จ โมหสฺส สมฺปยุตฺตกานญฺจ ขนฺธานํ อตฺถิปจฺจเยน ปจฺจโย.}

\prose{นตฺถิ\-ปจฺจเยน ปจฺจโย. วิคต\-ปจฺจเยน ปจฺจโย. อวิคต\-ปจฺจเยน ปจฺจโย. (3)}

\csromanmarkh{1. ปจฺจยานุโลม}\csromanmarkhii{2. สงฺขฺยาวาร}\csromanheadernomark{2}{1. \textbf{ปจฺจยานุโลม 2. สงฺขฺยาวาร}}
\proseitem{42}{เหตุยา ตีณิ, อารมฺมเณ นว, อธิปติยา นว, อนนฺตเร นว, สมนนฺตเร นว, สหชาเต นว, อญฺญมญฺเญ นว, นิสฺสเย นว, อุปนิสฺสเย นว, ปุเรชาเต ตีณิ, ปจฺฉาชาเต ตีณิ, อาเสวเน นว, กมฺเม ตีณิ, วิปาเก นว, อาหาเร ตีณิ, อินฺทฺริเย นว, ฌาเน ตีณิ, \csromanfolio{19}มคฺเค นว, สมฺปยุตฺเต นว, วิปฺปยุตฺเต ปญฺจ, อตฺถิยา นว, นตฺถิยา นว, วิคเต นว, อวิคเต นว. (เอวํ อนุมชฺชนฺเตน คเณตพฺพํ.)}

\vspace{\tipitakaparskip}
\csromancenter{\textbf{อนุโลมํ.}}
\proseitem{43}{เหตุ ธมฺโม เหตุสฺส ธมฺมสฺส อารมฺมณ\-ปจฺจเยน ปจฺจโย. สหชาต\-ปจฺจเยน ปจฺจโย. อุปนิสฺสย\-ปจฺจเยน ปจฺจโย. (1)}

\prose{เหตุ ธมฺโม นเหตุสฺส ธมฺมสฺส อารมฺมณ\-ปจฺจเยน ปจฺจโย. สหชาต\-ปจฺจเยน ปจฺจโย. อุปนิสฺสย\-ปจฺจเยน ปจฺจโย. ปจฺฉาชาต\-ปจฺจเยน ปจฺจโย. (2)}

\prose{เหตุ ธมฺโม เหตุสฺส จ นเหตุสฺส จ ธมฺมสฺส อารมฺมณ\-ปจฺจเยน ปจฺจโย. สหชาต\-ปจฺจเยน ปจฺจโย. อุปนิสฺสย\-ปจฺจเยน ปจฺจโย. (3)}

\proseitem{42}{ธมฺโม นเหตุสฺส ธมฺมสฺส อารมฺมณ\-ปจฺจเยน ปจฺจโย. สหชาต\-ปจฺจเยน ปจฺจโย. อุปนิสฺสย\-ปจฺจเยน ปจฺจโย. ปุเรชาต\-ปจฺจเยน ปจฺจโย. ปจฺฉาชาต\-ปจฺจเยน ปจฺจโย. กมฺมปจฺจเยน ปจฺจโย. อาหารปจฺจเยน ปจฺจโย. อินฺทฺริย\-ปจฺจเยน ปจฺจโย. (1)}

\prose{ธมฺโม เหตุสฺส ธมฺมสฺส อารมฺมณ\-ปจฺจเยน ปจฺจโย. สหชาต\-ปจฺจเยน ปจฺจโย. อุปนิสฺสย\-ปจฺจเยน ปจฺจโย. ปุเรชาต\-ปจฺจเยน ปจฺจโย. กมฺมปจฺจเยน ปจฺจโย. (2)}

\prose{ธมฺโม เหตุสฺส จ นเหตุสฺส จ ธมฺมสฺส อารมฺมณ\-ปจฺจเยน ปจฺจโย. สหชาต\-ปจฺจเยน ปจฺจโย. อุปนิสฺสย\-ปจฺจเยน ปจฺจโย. ปุเรชาต\-ปจฺจเยน ปจฺจโย. กมฺมปจฺจเยน ปจฺจโย. (3)}

\proseitem{45}{เหตุ จ นเหตุ จ ธมฺมา เหตุสฺส ธมฺมสฺส อารมฺมณ\-ปจฺจเยน ปจฺจโย. สหชาต\-ปจฺจเยน ปจฺจโย. อุปนิสฺสย\-ปจฺจเยน ปจฺจโย. (1)}

\prose{เหตุ จ นเหตุ จ ธมฺมา นเหตุสฺส ธมฺมสฺส อารมฺมณ\-ปจฺจเยน ปจฺจโย. สหชาต\-ปจฺจเยน ปจฺจโย. อุปนิสฺสย\-ปจฺจเยน ปจฺจโย. ปจฺฉาชาต\-ปจฺจเยน ปจฺจโย. (2)}

\prose{\csromanfolio{20}เหตุ จ นเหตุ จ ธมฺมา เหตุสฺส จ นเหตุสฺส จ ธมฺมสฺส อารมฺมณ\-ปจฺจเยน ปจฺจโย. สหชาต\-ปจฺจเยน ปจฺจโย. อุปนิสฺสย\-ปจฺจเยน ปจฺจโย. (3)}

\csromanmarkh{2. ปจฺจย\-ปจฺจนีย}\csromanmarkhii{2. สงฺขฺยาวาร}\csromanheadernomark{2}{2. ปจฺจย\-ปจฺจนีย 2. สงฺขฺยาวาร}
\proseitem{46}{นเหตุยา นว, นอารมฺมเณ นว ฯเปฯ โนอวิคเต นว. (เอวํ คเณตพฺพํ.)}

\vspace{\tipitakaparskip}
\csromancenter{\textbf{ปจฺจนียํ.}}
\proseitem{47}{เหตุปจฺจยา นอารมฺมเณ ตีณิ, นอธิปติยา ตีณิ, นอนนฺตเร ตีณิ, นสมนนฺตเร ตีณิ, นอญฺญมญฺเญ เอกํ, นอุปนิสฺสเย ตีณิ ฯเปฯ นมคฺเค ตีณิ, นสมฺปยุตฺเต เอกํ, นวิปฺปยุตฺเต ตีณิ, โนนตฺถิยา ตีณิ, โนวิคเต ตีณิ. (เอวํ คเณตพฺพํ.)}

\vspace{\tipitakaparskip}
\csromancenter{อนุโลม\-ปจฺจนียํ.}
\proseitem{48}{นเหตุปจฺจยา อารมฺมเณ นว, อธิปติยา นว, อนนฺตเร นว, สมนนฺตเร นว, สหชาเต ตีณิ, อญฺญมญฺเญ ตีณิ, นิสฺสเย ตีณิ, อุปนิสฺสเย นว, ปุเรชาเต ตีณิ, ปจฺฉาชาเต ตีณิ, อาเสวเน นว, กมฺเม ตีณิ, วิปาเก ตีณิ, อาหาเร ตีณิ, อินฺทฺริเย ตีณิ, ฌาเน ตีณิ, มคฺเค ตีณิ, สมฺปยุตฺเต ตีณิ, วิปฺปยุตฺเต ตีณิ, อตฺถิยา ตีณิ, นตฺถิยา นว, วิคเต นว, อวิคเต ตีณิ. (เอวํ คเณตพฺพํ.)}

\vspace{\tipitakaparskip}
\csromancenter{ปจฺจนียา\-นุโลมํ.}
\nitthitam{เหตุทุกํ นิฏฺฐิตํ.}
\csromanmatikaanchor{mtk.4}
\csromantocmarkref{section}{2. สเหตุกทุก}{mtk.4}
\bookmark[dest={mtk.4},level=1]{2. สเหตุกทุก}
\csromanmarkh{2. สเหตุทุก}\csromanmarkhii{2. สเหตุกทุก 1. ปฏิจฺจวาร}\csromanheadernomark{1}{2. สเหตุทุก \textbf{2. สเหตุกทุก 1. ปฏิจฺจวาร}}
\csromanmarkh{1. ปจฺจยานุโลม}\csromanmarkhii{1. วิภงฺควาร}\csromanheadernomark{2}{1. \textbf{ปจฺจยานุโลม 1. วิภงฺควาร}}
\csromanheader{2}{\textbf{เหตุ}}
\proseitem{49}{\csromanfolio{21}สเหตุกํ ธมฺมํ ปฏิจฺจ สเหตุโก ธมฺโม อุปฺปชฺชติ เหตุปจฺจยา. สเหตุกํ เอกํ ขนฺธํ ปฏิจฺจ ตโย ขนฺธา ฯเปฯ เทฺว ขนฺเธ ฯเปฯ. ปฏิสนฺธิ\-กฺขเณ ฯเปฯ. (1)}

\prose{สเหตุกํ ธมฺมํ ปฏิจฺจ อเหตุโก ธมฺโม อุปฺปชฺชติ เหตุปจฺจยา. สเหตุเก ขนฺเธ ปฏิจฺจ จิตฺต\-สมุฏฺฐานํ รูปํ. ปฏิสนฺธิ\-กฺขเณ ฯเปฯ.}

\prose{สเหตุกํ ธมฺมํ ปฏิจฺจ สเหตุโก จ อเหตุโก จ ธมฺมา อุปฺปชฺชนฺติ เหตุปจฺจยา. สเหตุกํ เอกํ ขนฺธํ ปฏิจฺจ ตโย ขนฺธา จิตฺต\-สมุฏฺฐานญฺจ รูปํ ฯเปฯ เทฺว ขนฺเธ ฯเปฯ. ปฏิสนฺธิ\-กฺขเณ ฯเปฯ. (3)}

\proseitem{50}{อเหตุกํ ธมฺมํ ปฏิจฺจ อเหตุโก ธมฺโม อุปฺปชฺชติ เหตุปจฺจยา. วิจิกิจฺฉา\-สหคตํ อุทฺธจฺจ\-สหคตํ โมหํ ปฏิจฺจ จิตฺต\-สมุฏฺฐานํ รูปํ. เอกํ มหาภูตํ ปฏิจฺจ ตโย มหาภูตา ฯเปฯ มหาภูเต ปฏิจฺจ จิตฺต\-สมุฏฺฐานํ รูปํ, กฏตฺตารูปํ, อุปาทารูปํ. (1)}

\prose{อเหตุกํ ธมฺมํ ปฏิจฺจ สเหตุโก ธมฺโม อุปฺปชฺชติ เหตุปจฺจยา. วิจิกิจฺฉา\-สหคตํ อุทฺธจฺจ\-สหคตํ โมหํ ปฏิจฺจ สมฺปยุตฺตกา ขนฺธา. ปฏิสนฺธิ\-กฺขเณ วตฺถุํ ปฏิจฺจ สเหตุกา ขนฺธา. (2)}

\prose{อเหตุกํ ธมฺมํ ปฏิจฺจ สเหตุโก จ อเหตุโก จ ธมฺมา อุปฺปชฺชนฺติ อเหตุปจฺจยา. วิจิกิจฺฉา\-สหคตํ อุทฺธจฺจ\-สหคตํ โมหํ ปฏิจฺจ สมฺปยุตฺตกา ขนฺธา จิตฺต\-สมุฏฺฐานญฺจ รูปํ. ปฏิสนฺธิ\-กฺขเณ วตฺถุํ ปฏิจฺจ สเหตุกา ขนฺธา. มหาภูเต ปฏิจฺจ กฏตฺตารูปํ. (3)}

\proseitem{51}{สเหตุกญฺจ อเหตุกญฺจ ธมฺมํ ปฏิจฺจ สเหตุโก ธมฺโม อุปฺปชฺชติ เหตุปจฺจยา. วิจิกิจฺฉา\-สหคตํ อุทฺธจฺจ\-สหคตํ เอกํ ขนฺธญฺจ โมหญฺจ ปฏิจฺจ ตโย ขนฺธา ฯเปฯ เทฺว ขนฺเธ ฯเปฯ. ปฏิสนฺธิ\-กฺขเณ สเหตุกํ เอกํ ขนฺธญฺจ วตฺถุญฺจ ปฏิจฺจ ตโย ขนฺธา ฯเปฯ เทฺว ขนฺเธ ฯเปฯ. (1)}

\prose{\csromanfolio{22}สเหตุกญฺจ อเหตุกญฺจ ธมฺมํ ปฏิจฺจ อเหตุโก ธมฺโม อุปฺปชฺชติ เหตุปจฺจยา. สเหตุเก ขนฺเธ จ มหาภูเต จ ปฏิจฺจ จิตฺต\-สมุฏฺฐานํ รูปํ. วิจิกิจฺฉา\-สหคเต อุทฺธจฺจ\-สหคเต ขนฺเธ จ โมหญฺจ ปฏิจฺจ จิตฺต\-สมุฏฺฐานํ รูปํ. ปฏิสนฺธิ\-กฺขเณ ฯเปฯ. (2)}

\prose{สเหตุกญฺจ อเหตุกญฺจ ธมฺมํ ปฏิจฺจ สเหตุโก จ อเหตุโก จ ธมฺมา อุปฺปชฺชนฺติ เหตุปจฺจยา. วิจิกิจฺฉา\-สหคตํ อุทฺธจฺจ\-สหคตํ เอกํ ขนฺธญฺจ โมหญฺจ ปฏิจฺจ ตโย ขนฺธา จิตฺต\-สมุฏฺฐานญฺจ รูปํ ฯเปฯ เทฺว ขนฺเธ ฯเปฯ. ปฏิสนฺธิ\-กฺขเณ สเหตุกํ เอกํ ขนฺธญฺจ วตฺถุญฺจ ปฏิจฺจ ตโย ขนฺธา ฯเปฯ เทฺว ขนฺเธ ฯเปฯ. สเหตุเก ขนฺเธ จ มหาภูเต จ ปฏิจฺจ กฏตฺตารูปํ. (3)}

\proseitem{52}{สเหตุกํ ธมฺมํ ปฏิจฺจ สเหตุโก ธมฺโม อุปฺปชฺชติ อารมฺมณ\-ปจฺจยา. สเหตุกํ เอกํ ขนฺธํ ปฏิจฺจ ตโย ขนฺธา ฯเปฯ เทฺว ขนฺเธ ฯเปฯ. ปฏิสนฺธิ\-กฺขเณ ฯเปฯ. (1)}

\prose{สเหตุกํ ธมฺมํ ปฏิจฺจ อเหตุโก ธมฺโม อุปฺปชฺชติ อารมฺมณ\-ปจฺจยา. วิจิกิจฺฉา\-สหคเต อุทฺธจฺจ\-สหคเต ขนฺเธ ปฏิจฺจ วิจิกิจฺฉา\-สหคโต อุทฺธจฺจ\-สหคโต โมโห. (2)}

\prose{สเหตุกํ ธมฺมํ ปฏิจฺจ สเหตุโก จ อเหตุโก จ ธมฺมา อุปฺปชฺชนฺติ อารมฺมณ\-ปจฺจยา. วิจิกิจฺฉา\-สหคตํ อุทฺธจฺจ\-สหคตํ เอกํ ขนฺธํ ปฏิจฺจ ตโย ขนฺธา โมโห จ ฯเปฯ เทฺว ขนฺเธ ฯเปฯ. (3)}

\proseitem{53}{อเหตุกํ ธมฺมํ ปฏิจฺจ อเหตุโก ธมฺโม อุปฺปชฺชติ อารมฺมณ\-ปจฺจยา. อเหตุกํ เอกํ ขนฺธํ ปฏิจฺจ ตโย ขนฺธา ฯเปฯ เทฺว ขนฺเธ ฯเปฯ. ปฏิสนฺธิ\-กฺขเณ วตฺถุํ ปฏิจฺจ อเหตุกา ขนฺธา. (1)}

\prose{อเหตุกํ ธมฺมํ ปฏิจฺจ สเหตุโก ธมฺโม อุปฺปชฺชติ อารมฺมณ\-ปจฺจยา. วิจิกิจฺฉา\-สหคตํ อุทฺธจฺจ\-สหคตํ โมหํ ปฏิจฺจ สมฺปยุตฺตกา ขนฺธา. ปฏิสนฺธิ\-กฺขเณ วตฺถุํ ปฏิจฺจ สเหตุกา ขนฺธา. (2)}

\prose{สเหตุกญฺจ อเหตุกญฺจ ธมฺมํ ปฏิจฺจ สเหตุโก ธมฺโม อุปฺปชฺชติ อารมฺมณ\-ปจฺจยา. วิจิกิจฺฉา\-สหคตํ อุทฺธจฺจ\-สหคตํ เอกํ ขนฺธญฺจ โมหญฺจ ปฏิจฺจ \csromanfolio{23}ตโย ขนฺธา ฯเปฯ เทฺว ขนฺเธ ฯเปฯ. ปฏิสนฺธิ\-กฺขเณ สเหตุกํ เอกํ ขนฺธญฺจ วตฺถุญฺจ ปฏิจฺจ ตโย ขนฺธา ฯเปฯ เทฺว ขนฺเธ ฯเปฯ. (1)}

\proseitem{54}{สเหตุกํ ธมฺมํ ปฏิจฺจ สเหตุโก ธมฺโม อุปฺปชฺชติ อธิปติ\-ปจฺจยา. สเหตุกํ เอกํ ขนฺธํ ปฏิจฺจ ตโย ขนฺธา ฯเปฯ เทฺว ขนฺเธ ฯเปฯ. (1)}

\prose{สเหตุกํ ธมฺมํ ปฏิจฺจ อเหตุโก ธมฺโม อุปฺปชฺชติ อธิปติ\-ปจฺจยา. สเหตุเก ขนฺเธ ปฏิจฺจ จิตฺต\-สมุฏฺฐานํ รูปํ. (2)}

\prose{สเหตุกํ ธมฺมํ ปฏิจฺจ สเหตุโก จ อเหตุโก จ ธมฺมา อุปฺปชฺชนฺติ อธิปติ\-ปจฺจยา. สเหตุกํ เอกํ ขนฺธํ ปฏิจฺจ ตโย ขนฺธา จิตฺต\-สมุฏฺฐานญฺจ รูปํ ฯเปฯ เทฺว ขนฺเธ ฯเปฯ. (3)}

\proseitem{53}{อเหตุกํ ธมฺมํ ปฏิจฺจ อเหตุโก ธมฺโม อุปฺปชฺชติ อธิปติ\-ปจฺจยา. เอกํ มหาภูตํ ฯเปฯ มหาภูเต ปฏิจฺจ จิตฺต\-สมุฏฺฐานํ รูปํ, อุปาทารูปํ. (1)}

\prose{สเหตุกญฺจ อเหตุกญฺจ ธมฺมํ ปฏิจฺจ อเหตุโก ธมฺโม อุปฺปชฺชติ อธิปติ\-ปจฺจยา. สเหตุเก ขนฺเธ จ มหาภูเต จ ปฏิจฺจ จิตฺต\-สมุฏฺฐานํ รูปํ. (1)}

\prose{สเหตุกํ ธมฺมํ ปฏิจฺจ สเหตุโก ธมฺโม อุปฺปชฺชติ อนนฺตร\-ปจฺจยา. สมนนฺตร\-ปจฺจยา. สหชาต\-ปจฺจยา. สเหตุกํ เอกํ ขนฺธํ ปฏิจฺจ ตโย ขนฺธา ฯเปฯ. ปฏิสนฺธิ\-กฺขเณ ฯเปฯ. (1)}

\prose{สเหตุกํ ธมฺมํ ปฏิจฺจ อเหตุโก ธมฺโม อุปฺปชฺชติ สหชาต\-ปจฺจยา. สเหตุเก ขนฺเธ ปฏิจฺจ จิตฺต\-สมุฏฺฐานํ รูปํ วิจิกิจฺฉา\-สหคเต อุทฺธจฺจ\-สหคเต ขนฺเธ ปฏิจฺจ วิจิกิจฺฉา\-สหคโต อุทฺธจฺจ\-สหคโต โมโห จิตฺต\-สมุฏฺฐานญฺจ รูปํ. ปฏิสนฺธิ\-กฺขเณ ฯเปฯ.}

\prose{สเหตุกํ ธมฺมํ ปฏิจฺจ สเหตุโก จ อเหตุโก จ ธมฺมา อุปฺปชฺชนฺติ สหชาต\-ปจฺจยา. สเหตุกํ เอกํ ขนฺธํ ปฏิจฺจ ตโย ขนฺธา จิตฺต\-สมุฏฺฐานญฺจ รูปํ เทฺว ขนฺเธ ฯเปฯ. วิจิกิจฺฉา\-สหคตํ อุทฺธจฺจ\-สหคตํ เอกํ ขนฺธํ ปฏิจฺจ ตโย ขนฺธา โมโห จ จิตฺต\-สมุฏฺฐานญฺจ รูปํ ฯเปฯ เทฺว ขนฺเธ ฯเปฯ. ปฏิสนฺธิ\-กฺขเณ ฯเปฯ. (3)}

\proseitem{57}{\csromanfolio{24}อเหตุกํ ธมฺมํ ปฏิจฺจ อเหตุโก ธมฺโม อุปฺปชฺชติ สหชาต\-ปจฺจยา. อเหตุกํ เอกํ ขนฺธํ ปฏิจฺจ ตโย ขนฺธา จิตฺต\-สมุฏฺฐานญฺจ รูปํ ฯเปฯ เทฺว ขนฺเธ ฯเปฯ. วิจิกิจฺฉา\-สหคตํ อุทฺธจฺจ\-สหคตํ โมหํ ปฏิจฺจ จิตฺต\-สมุฏฺฐานํ รูปํ. ปฏิสนฺธิ\-กฺขเณ ฯเปฯ. ขนฺเธ ปฏิจฺจ วตฺถุ, วตฺถุํ ปฏิจฺจ ขนฺธา. เอกํ มหาภูตํ ฯเปฯ. พาหิรํ. อาหารสมุฏฺฐานํ. อุตุสมุฏฺฐานํ. อสญฺญสตฺตานํ เอกํ มหาภูตํ ฯเปฯ. (1)}

\prose{อเหตุกํ ธมฺมํ ปฏิจฺจ สเหตุโก ธมฺโม อุปฺปชฺชติ สหชาต\-ปจฺจยา. (อิเม ปญฺจ ปญฺหา เหตุสทิสา, นินฺนานํ.) (2)}

\csromanheader{2}{\textbf{อญฺญมญฺญ}}
\proseitem{58}{สเหตุกํ ธมฺมํ ปฏิจฺจ สเหตุโก ธมฺโม อุปฺปชฺชติ อญฺญมญฺญ\-ปจฺจยา. สเหตุกํ เอกํ ขนฺธํ ปฏิจฺจ ตโย ขนฺธา ฯเปฯ. ปฏิสนฺธิ\-กฺขเณ ฯเปฯ. (1)}

\prose{สเหตุกํ ธมฺมํ ปฏิจฺจ อเหตุโก ธมฺโม อุปฺปชฺชติ อญฺญมญฺญ\-ปจฺจยา. วิจิกิจฺฉา\-สหคเต อุทฺธจฺจ\-สหคเต ขนฺเธ ปฏิจฺจ วิจิกิจฺฉา\-สหคโต อุทฺธจฺจ\-สหคโต โมโห. ปฏิสนฺธิ\-กฺขเณ สเหตุเก ขนฺเธ ปฏิจฺจ วตฺถุ. (2)}

\prose{สเหตุกํ ธมฺมํ ปฏิจฺจ สเหตุโก จ อเหตุโก จ ธมฺมา อุปฺปชฺชนฺติ อญฺญมญฺญ\-ปจฺจยา. วิจิกิจฺฉา\-สหคตํ อุทฺธจฺจ\-สหคตํ เอกํ ขนฺธํ ปฏิจฺจ ตโย ขนฺธา โมโห จ ฯเปฯ เทฺว ขนฺเธ ฯเปฯ. ปฏิสนฺธิ\-กฺขเณ สเหตุกํ เอกํ ขนฺธํ ปฏิจฺจ ตโย ขนฺธา วตฺถุ จ ฯเปฯ เทฺว ขนฺเธ ฯเปฯ. (3)}

\proseitem{59}{อเหตุกํ ธมฺมํ ปฏิจฺจ อเหตุโก ธมฺโม อุปฺปชฺชติ อญฺญมญฺญ\-ปจฺจยา. อเหตุกํ เอกํ ขนฺธํ ปฏิจฺจ ตโย ขนฺธา ฯเปฯ เทฺว ขนฺเธ ฯเปฯ. ปฏิสนฺธิ\-กฺขเณ อเหตุกํ เอกํ ขนฺธํ ปฏิจฺจ ตโย ขนฺธา วตฺถุ จ ฯเปฯ เทฺว ขนฺเธ ฯเปฯ. (สํขิตฺตํ. ยาว อสญฺญสตฺตา.) (1)}

\prose{อเหตุกํ ธมฺมํ ปฏิจฺจ สเหตุโก ธมฺโม อุปฺปชฺชติ อญฺญมญฺญ\-ปจฺจยา. วิจิกิจฺฉา\-สหคตํ อุทฺธจฺจ\-สหคตํ โมหํ ปฏิจฺจ สมฺปยุตฺตกา ขนฺธา. ปฏิสนฺธิ\-กฺขเณ วตฺถุํ ปฏิจฺจ สเหตุกา ขนฺธา. (2)}

\prose{สเหตุกญฺจ อเหตุกญฺจ ธมฺมํ ปฏิจฺจ สเหตุโก ธมฺโม อุปฺปชฺชติ อญฺญมญฺญ\-ปจฺจยา. วิจิกิจฺฉา\-สหคตํ อุทฺธจฺจ\-สหคตํ เอกํ ขนฺธญฺจ โมหญฺจ \csromanfolio{25}ปฏิจฺจ ตโย ขนฺธา ฯเปฯ เทฺว ขนฺเธ ฯเปฯ. ปฏิสนฺธิ\-กฺขเณ สเหตุกํ เอกํ ขนฺธญฺจ วตฺถุญฺจ ปฏิจฺจ ตโย ขนฺธา ฯเปฯ เทฺว ขนฺเธ ฯเปฯ. (1)}

\csromanheader{2}{\textbf{นิสฺสยาทิ}}
\proseitem{59}{สเหตุกํ ธมฺมํ ปฏิจฺจ สเหตุโก ธมฺโม อุปฺปชฺชติ นิสฺสยปจฺจยา. อุปนิสฺสย\-ปจฺจยา. ปุเรชาต\-ปจฺจยา. อาเสวนปจฺจยา. กมฺมปจฺจยา. วิปากปจฺจยา. อาหารปจฺจยา. อินฺทฺริยปจฺจยา. ฌานปจฺจยา. มคฺคปจฺจยา. (ฌานมฺปิ มคฺคมฺปิ สหชาตปจฺจสสทิสา. พาหิรา มหาภูตา นตฺถิ.) สมฺปยุตฺต\-ปจฺจยา. วิปฺปยุตฺต\-ปจฺจยา. อตฺถิปจฺจยา. นตฺถิปจฺจยา. วิคตปจฺจยา. อวิคตปจฺจยา.}

\csromanmarkh{1. ปจฺจยานุโลม}\csromanmarkhii{2. สงฺขฺยาวาร}\csromanheadernomark{2}{1. \textbf{ปจฺจยานุโลม 2. สงฺขฺยาวาร}}
\proseitem{59}{เหตุยา นว, อารมฺมเณ ฉ, อธิปติยา ปญฺจ, อนนฺตเร ฉ, สมนนฺตเร ฉ, สหชาเต นว, อญฺญมญฺเญ ฉ, นิสฺสเย นว, อุปนิสฺสเย ฉ, ปุเรชาเต ฉ, อาเสวเน ฉ, กมฺเม นว, วิปาเก นว, อาหาเร นว, อินฺทฺริเย นว, ฌาเน นว, มคฺเค นว, สมฺปยุตฺเต ฉ, วิปฺปยุตฺเต นว, อตฺถิยา นว, นตฺถิยา ฉ, วิคเต ฉ, อวิคเต นว. (เอวํ คเณตพฺพํ.)}

\vspace{\tipitakaparskip}
\csromancenter{\textbf{อนุโลมํ.}}
\csromanmarkh{2. ปจฺจย\-ปจฺจนีย}\csromanmarkhii{1. วิภงฺควาร}\csromanheadernomark{2}{2. ปจฺจย\-ปจฺจนีย 1. วิภงฺควาร}
\proseitem{62}{สเหตุกํ ธมฺมํ ปฏิจฺจ อเหตุโก ธมฺโม อุปฺปชฺชติ นเหตุปจฺจยา. วิจิกิจฺฉา\-สหคเต อุทฺธจฺจ\-สหคเต ขนฺเธ ปฏิจฺจ วิจิกิจฺฉา\-สหคโต อุทฺธจฺจ\-สหคโต โมโห. (1)}

\prose{อเหตุกํ ธมฺมํ ปฏิจฺจ อเหตุโก ธมฺโม อุปฺปชฺชติ นเหตุปจฺจยา อเหตุกํ เอกํ ขนฺธํ ปฏิจฺจ ตโย ขนฺธา จิตฺต\-สมุฏฺฐานํ รูปํ ฯเปฯ เทฺว ขนฺเธ ฯเปฯ. ปฏิสนฺธิ\-กฺขเณ ฯเปฯ. (1)}

\vspace{\tipitakaparskip}
\csromancenter{(สพฺพํ ยาว อสญฺญสตฺตา, ตาว กาตพฺพํ.)}
\prose{\csromanfolio{26}นอารมฺมณาทิ}

\proseitem{63}{สเหตุกํ ธมฺมํ ปฏิจฺจ อเหตุโก ธมฺโม อุปฺปชฺชติ นอารมฺมณ\-ปจฺจยา. สเหตุเก ขนฺเธ ปฏิจฺจ จิตฺต\-สมุฏฺฐานํ รูปํ. ปฏิสนฺธิ\-กฺขเณ ฯเปฯ. (1)}

\prose{อเหตุกํ ธมฺมํ ปฏิจฺจ อเหตุโก ธมฺโม อุปฺปชฺชติ นอารมฺมณ\-ปจฺจยา. อเหตุเก ขนฺเธ ปฏิจฺจ จิตฺต\-สมุฏฺฐานํ รูปํ. วิจิกิจฺฉา\-สหคตํ อุทฺธจฺจ\-สหคตํ โมหํ ปฏิจฺจ จิตฺต\-สมุฏฺฐานํ รูปํ. ปฏิสนฺธิ\-กฺขเณ ฯเปฯ. ขนฺเธ ปฏิจฺจ วตฺถุ. เอกํ มหาภูตํ ฯเปฯ. อสญฺญสตฺตานํ เอกํ ฯเปฯ. (1)}

\prose{สเหตุกญฺจ อเหตุกญฺจ ธมฺมํ ปฏิจฺจ อเหตุโก ธมฺโม อุปฺปชฺชติ นอารมฺมณ\-ปจฺจยา. สเหตุเก ขนฺเธ จ มหาภูเต จ ปฏิจฺจจิตฺต\-สมุฏฺฐานํ รูปํ. วิจิกิจฺฉา\-สหคเต อุทฺธจฺจ\-สหคเต ขนฺเธ จ โมหญฺจ ปฏิจฺจ จิตฺต\-สมุฏฺฐานํ รูปํ. ปฏิสนฺธิ\-กฺขเณ ฯเปฯ. (1)}

\prose{สเหตุกํ ธมฺมํ ปฏิจฺจ สเหตุโก ธมฺโม อุปฺปชฺชติ นอธิปติ\-ปจฺจยา. (อนุโลม\-สหชาต\-สทิสา.) นอนนฺตร\-ปจฺจยา. นสมนนฺตร\-ปจฺจยา. นอญฺญมญฺญ\-ปจฺจยา. นอุปนิสฺสยปจฺจยา. นปุเรชาต\-ปจฺจยา. อรูเป สเหตุกํ เอกํ ขนฺธํ ฯเปฯ. ปฏิสนฺธิ\-กฺขเณ ฯเปฯ. (1)}

\proseitem{64}{สเหตุกํ ธมฺมํ ปฏิจฺจ อเหตุโก ธมฺโม อุปฺปชฺชติ นปุเรชาต\-ปจฺจยา. อรูเป วิจิกิจฺฉา\-สหคเต อุทฺธจฺจ\-สหคเต ขนฺเธ ปฏิจฺจ วิจิกิจฺฉา\-สหคโต อุทฺธจฺจ\-สหคโต โมโห. สเหตุเก ขนฺเธ ปฏิจฺจ จิตฺต\-สมุฏฺฐานํ รูปํ. ปฏิสนฺธิ\-กฺขเณ ฯเปฯ. (2)}

\prose{สเหตุกํ ธมฺมํ ปฏิจฺจ สเหตุโก จ อเหตุโก จ ธมฺมา อุปฺปชฺชนฺติ นปุเรชาต\-ปจฺจยา. อรูเป วิจิกิจฺฉา\-สหคตํ อุทฺธจฺจ\-สหคตํ เอกํ ขนฺธํ ปฏิจฺจ ตโย ขนฺธา โมโห จ ฯเปฯ เทฺว ขนฺเธ ฯเปฯ. ปฏิสนฺธิ\-กฺขเณ ฯเปฯ. (3)}

\prose{อเหตุกํ ธมฺมํ ปฏิจฺจ อเหตุโก ธมฺโม อุปฺปชฺชติ นปุเรชาต\-ปจฺจยา. อรูเป อเหตุกํ เอกํ ขนฺธํ ฯเปฯ เทฺว ขนฺเธ ฯเปฯ. อเหตุเก ขนฺเธ ปฏิจฺจ จิตฺต\-สมุฏฺฐานํ รูปํ. วิจิกิจฺฉา\-สหคตํ อุทฺธจฺจ\-สหคตํ โมหํ ปฏิจฺจ จิตฺต\-สมุฏฺฐานํ รูปํ. ปฏิสนฺธิ\-กฺขเณ ฯเปฯ. (ยาว อสญฺญสตฺตา, ตาว วิตฺถาโร.) (1)}

\prose{\csromanfolio{27}อเหตุกํ ธมฺมํ ปฏิจฺจ สเหตุโก ธมฺโม อุปฺปชฺชติ นปุเรชาต\-ปจฺจยา. อรูเป วิจิกิจฺฉา\-สหคตํ อุทฺธจฺจ\-สหคตํ โมหํ ปฏิจฺจ สมฺปยุตฺตกา ขนฺธา. ปฏิสนฺธิ\-กฺขเณ วตฺถุํ ปฏิจฺจ สเหตุกา ขนฺธา. (2)}

\prose{อเหตุกํ ธมฺมํ ปฏิจฺจ สเหตุโก จ อเหตุโก จ ธมฺมา อุปฺปชฺชนฺติ นปุเรชาต\-ปจฺจยา. ปฏิสนฺธิ\-กฺขเณ วตฺถุํ ปฏิจฺจ สเหตุกา ขนฺธา. มหาภูเต ปฏิจฺจ กฏตฺตารูปํ. (3)}

\proseitem{65}{สเหตุกญฺจ อเหตุกญฺจ ธมฺมํ ปฏิจฺจ สเหตุโก ธมฺโม อุปฺปชฺชติ นปุเรชาต\-ปจฺจยา. อรูเป วิจิกิจฺฉา\-สหคตํ อุทฺธจฺจ\-สหคตํ เอกํ ขนฺธญฺจ โมหญฺจ ปฏิจฺจ ตโย ขนฺธา ฯเปฯ เทฺว ขนฺเธ ฯเปฯ. ปฏิสนฺธิ\-กฺขเณ สเหตุกํ เอกํ ขนฺธญฺจ วตฺถุญฺจ ปฏิจฺจ ตโย ขนฺธา ฯเปฯ เทฺว ขนฺเธ ฯเปฯ. (1)}

\prose{สเหตุกญฺจ อเหตุกญฺจ ธมฺมํ ปฏิจฺจ อเหตุโก ธมฺโม อุปฺปชฺชติ นปุเรชาต\-ปจฺจยา. สเหตุเก ขนฺเธ จ มหาภูเต จ ปฏิจฺจ จิตฺต\-สมุฏฺฐานํ รูปํ. วิจิกิจฺฉา\-สหคเต อุทฺธจฺจ\-สหคเต ขนฺเธ จ โมหญฺจ ปฏิจฺจ จิตฺต\-สมุฏฺฐานํ รูปํ. ปฏิสนฺธิ\-กฺขเณ ฯเปฯ. (2)}

\prose{สเหตุกญฺจ อเหตุกญฺจ ธมฺมํ ปฏิจฺจ สเหตุโก จ อเหตุโก จ ธมฺมา อุปฺปชฺชนฺติ นปุเรชาต\-ปจฺจยา. ปฏิสนฺธิ\-กฺขเณ สเหตุกํ เอกํ ขนฺธญฺจ วตฺถุญฺจ ปฏิจฺจ ตโย ขนฺธา ฯเปฯ เทฺว ขนฺเธ ฯเปฯ. สเหตุเก ขนฺเธ จ มหาภูเต จ ปฏิจฺจ กฏตฺตารูปํ. (3)}

\csromanheader{2}{\textbf{นปจฺฉาชาตาทิ}}
\proseitem{66}{สเหตุกํ ธมฺมํ ปฏิจฺจ สเหตุโก ธมฺโม อุปฺปชฺชติ นปจฺฉาชาต\-ปจฺจยา. นอาเสวน\-ปจฺจยา. นกมฺมปจฺจยา. สเหตุเก ขนฺเธ ปฏิจฺจ สเหตุกา เจตนา. (1)}

\prose{อเหตุกํ ธมฺมํ ปฏิจฺจ อเหตุโก ธมฺโม อุปฺปชฺชติ นกมฺมปจฺจยา. อเหตุเก ขนฺเธ ปฏิจฺจ อเหตุกา เจตนา. พาหิรํ. อาหารสมุฏฺฐานํ. อุตุสมุฏฺฐานํ ฯเปฯ. (1)}

\prose{อเหตุกํ ธมฺมํ ปฏิจฺจ สเหตุโก ธมฺโม อุปฺปชฺชติ นกมฺมปจฺจยา. วิจิกิจฺฉา\-สหคตํ อุทฺธจฺจ\-สหคตํ โมหํ ปฏิจฺจ สมฺปยุตฺตกา เจตนา. (2)}

\prose{\csromanfolio{28}สเหตุกญฺจ อเหตุกญฺจ ธมฺมํ ปฏิจฺจ สเหตุโก ธมฺโม อุปฺปชฺชติ นกมฺมปจฺจยา. วิจิกิจฺฉา\-สหคเต อุทฺธจฺจ\-สหคเต ขนฺเธ จ โมหญฺจ ปฏิจฺจ สมฺปยุตฺตกา เจตนา. นวิปาก\-ปจฺจยา. (ปฏิสนฺธิ นตฺถิ.)}

\prose{นอาหาราทิ}

\proseitem{67}{อเหตุกํ ธมฺมํ ปฏิจฺจ อเหตุโก ธมฺโม อุปฺปชฺชติ นอาหารปจฺจยา. นอินฺทฺริยปจฺจยา. นฌานปจฺจยา. นมคฺคปจฺจยา. นสมฺปยุตฺต\-ปจฺจยา.}

\csromanheader{2}{\textbf{นวิปฺปยุตฺตาทิ}}
\proseitem{68}{สเหตุกํ ธมฺมํ ปฏิจฺจ สเหตุโก ธมฺโม อุปฺปชฺชติ นวิปฺปยุตฺต\-ปจฺจยา. อรูเป สเหตุกํ เอกํ ขนฺธํ ฯเปฯ. (1)}

\prose{สเหตุกํ ธมฺมํ ปฏิจฺจ อเหตุโก ธมฺโม อุปฺปชฺชติ นวิปฺปยุตฺต\-ปจฺจยา. อรูเป วิจิกิจฺฉา\-สหคเต อุทฺธจฺจ\-สหคเต ขนฺเธ ปฏิจฺจ วิจิกิจฺฉา\-สหคโต อุทฺธจฺจ\-สหคโต โมโห. (2)}

\prose{สเหตุกํ ธมฺมํ ปฏิจฺจ สเหตุโก จ อเหตุโก จ ธมฺมา อุปฺปชฺชนฺติ นวิปฺปยุตฺต\-ปจฺจยา. อรูเป วิจิกิจฺฉา\-สหคตํ อุทฺธจฺจ\-สหคตํ เอกํ ขนฺธํ ปฏิจฺจ ตโย ขนฺธา โมโห จ ฯเปฯ เทฺว ขนฺเธ ฯเปฯ. (3)}

\prose{อเหตุกํ ธมฺมํ ปฏิจฺจ อเหตุโก ธมฺโม อุปฺปชฺชติ นวิปฺปยุตฺต\-ปจฺจยา. อรูเป อเหตุกํ เอกํ ขนฺธํ ปฏิจฺจ ตโย ขนฺธา ฯเปฯ เทฺว ขนฺเธ ฯเปฯ. (1)}

\prose{อเหตุกํ ธมฺมํ ปฏิจฺจ สเหตุโก ธมฺโม อุปฺปชฺชติ นวิปฺปยุตฺต\-ปจฺจยา. อรูเป วิจิกิจฺฉา\-สหคตํ อุทฺธจฺจ\-สหคตํ โมหํ ปฏิจฺจ สมฺปยุตฺตกา ขนฺธา. (2)}

\prose{สเหตุกญฺจ อเหตุกญฺจ ธมฺมํ ปฏิจฺจ สเหตุโก ธมฺโม อุปฺปชฺชติ นวิปฺปยุตฺต\-ปจฺจยา. อรูเป วิจิกิจฺฉา\-สหคตํ อุทฺธจฺจ\-สหคตํ เอกํ ขนฺธญฺจ โมหญฺจ ปฏิจฺจ ตโย ขนฺธา ฯเปฯ เทฺว ขนฺเธ ฯเปฯ. โนนตฺถิปจฺจยา. โนวิคต\-ปจฺจยา. (1)}

\csromanmarkh{2. ปจฺจย\-ปจฺจนีย}\csromanmarkhii{2. สงฺขฺยาวาร}\csromanheadernomark{2}{2. ปจฺจย\-ปจฺจนีย 2. สงฺขฺยาวาร}
\proseitem{69}{นเหตุยา เทฺว, นอารมฺมเณ ตีณิ, นอธิปติยา นว, นอนนฺตเร ตีณิ, นสมนนฺตเร ตีณิ, นอญฺญมญฺเญ ตีณิ, นอุปนิสฺสเย \csromanfolio{29}ตีณิ, นปุเรชาเต นว, นปจฺฉาชาเต นว, นอาเสวเน นว, นกมฺเม จตฺตาริ, นวิปาเก นว, นอาหาเร เอกํ, ไนนฺทฺริเย เอกํ, นฌาเน เอกํ, นมคฺเค เอกํ, นสมฺปยุตฺเต ตีณิ, นวิปฺปยุตฺเต ฉ, โนนตฺถิยา ตีณิ, โนวิคเต ตีณิ. (เอวํ คเณตพฺพํ.)}

\vspace{\tipitakaparskip}
\csromancenter{ปจฺจนียํ.}
\proseitem{70}{เหตุปจฺจยา นอารมฺมเณ ตีณิ, นอธิปติยา นว, นอนนฺตเร ตีณิ, นสมนนฺตเร ตีณิ, นอญฺญมญฺเญ ตีณิ, นอุปนิสฺสเย ตีณิ, นปุเรชาเต นว, นปจฺฉาชาเต นว, นอาเสวเน นว, นกมฺเม ตีณิ, นวิปาเก นว, นสมฺปยุตฺเต ตีณิ, นวิปฺปยุตฺเต ตีณิ, โนนตฺถิยา ตีณิ, โนวิคเต ตีณิ. (เอวํ คเณตพฺพํ.)}

\vspace{\tipitakaparskip}
\csromancenter{อนุโลม\-ปจฺจนียํ.}
\proseitem{69}{นเหตุปจฺจยา อารมฺมเณ เทฺว, อนนฺตเร เทฺว, สมนนฺตเร เทฺว, (สพฺพตฺถ เทฺว,) วิปาเก เอกํ, อาหาเร เทฺว, อินฺทฺริเย เทฺว, ฌาเน เทฺว, มคฺเค เอกํ, สมฺปยุตฺเต เทฺว ฯเปฯ อวิคเต เทฺว. (เอวํ คเณตพฺพํ.)}

\vspace{\tipitakaparskip}
\csromancenter{ปจฺจนียา\-นุโลมํ.}
\csromancenter{(สหชาตวาโร ปฏิจฺจ\-วาร\-สทิโส.)}
\csromanmarkh{2. สเหตุกทุก}\csromanmarkhii{3. ปจฺจยวาร}\csromanheadernomark{2}{2. \textbf{สเหตุกทุก 3. ปจฺจยวาร}}
\csromanmarkh{1. ปจฺจยานุโลม}\csromanmarkhii{1. วิภงฺควาร}\csromanheadernomark{2}{1. \textbf{ปจฺจยานุโลม 1. วิภงฺควาร}}
\csromanheader{2}{\textbf{เหตุ}}
\proseitem{72}{\csromanfolio{30}สเหตุกํ ธมฺมํ ปจฺจยา สเหตุโก ธมฺโม อุปฺปชฺชติ เหตุปจฺจยา. (สเหตุก\-มูลกํ ปฏิจฺจ\-วาร\-สทิสํ.)}

\prose{อเหตุกํ ธมฺมํ ปจฺจยา อเหตุโก ธมฺโม ฯเปฯ. (ปฏิจฺจวารสทิสํเยว.) (1)}

\prose{อเหตุกํ ธมฺมํ ปจฺจยา สเหตุโก ธมฺโม อุปฺปชฺชติ เหตุปจฺจยา. วตฺถุํ ปจฺจยา สเหตุกา ขนฺธา. วิจิกิจฺฉา\-สหคตํ อุทฺธจฺจ\-สหคตํ โมหํ ปจฺจยา สมฺปยุตฺตกา ขนฺธา. ปฏิสนฺธิ\-กฺขเณ วตฺถุํ ปจฺจยา สเหตุกา ขนฺธา. (2)}

\prose{อเหตุกํ ธมฺมํ ปจฺจยา สเหตุโก จ อเหตุโก จ ธมฺมา อุปฺปชฺชนฺติ เหตุปจฺจยา. วตฺถุํ ปจฺจยา สเหตุกา ขนฺธา. มหาภูเต ปจฺจยา จิตฺต\-สมุฏฺฐานํ รูปํ. วิจิกิจฺฉา\-สหคตํ อุทฺธจฺจ\-สหคตํ โมหํ ปจฺจยา สมฺปยุตฺตกา ขนฺธา จิตฺต\-สมุฏฺฐานญฺจ รูปํ. ปฏิสนฺธิ\-กฺขเณ วตฺถุํ ฯเปฯ. (3)}

\proseitem{73}{สเหตุกญฺจ อเหตุกญฺจ ธมฺมํ ปจฺจยา สเหตุโก ธมฺโม อุปฺปชฺชติ เหตุปจฺจยา. สเหตุกํ เอกํ ขนฺธญฺจ วตฺถุญฺจ ปจฺจยา ตโย ขนฺธา ฯเปฯ เทฺว ขนฺเธ ฯเปฯ. วิจิกิจฺฉา\-สหคตํ อุทฺธจฺจ\-สหคตํ เอกํ ขนฺธญฺจ โมหญฺจ ปจฺจยา ตโย ขนฺธา ฯเปฯ เทฺว ขนฺเธ ฯเปฯ. ปฏิสนฺธิ\-กฺขเณ ฯเปฯ. (1)}

\prose{สเหตุกญฺจ อเหตุกญฺจ ธมฺมํ ปจฺจยา อเหตุโก ธมฺโม อุปฺปชฺชติ เหตุปจฺจยา. สเหตุเก ขนฺเธ จ มหาภูเต จ ปจฺจยา จิตฺต\-สมุฏฺฐานํ รูปํ. วิจิกิจฺฉา\-สหคเต อุทฺธจฺจ\-สหคเต ขนฺเธ จ โมหญฺจ ปจฺจยา จิตฺต\-สมุฏฺฐานํ รูปํ. ปฏิสนฺธิ\-กฺขเณ ฯเปฯ. (2)}

\prose{สเหตุกญฺจ อเหตุกญฺจ ธมฺมํ ปจฺจยา สเหตุโก จ อเหตุโก จ ธมฺมา อุปฺปชฺชนฺติ เหตุปจฺจยา. สเหตุกํ เอกํ ขนฺธญฺจ วตฺถุญฺจ ปจฺจยา ตโย ขนฺธา ฯเปฯ เทฺว ขนฺเธ ฯเปฯ. สเหตุเก ขนฺเธ จ มหาภูเต จ ปจฺจยา จิตฺต\-สมุฏฺฐานํ รูปํ. วิจิกิจฺฉา\-สหคตํ อุทฺธจฺจ\-สหคตํ เอกํ ขนฺธญฺจ \csromanfolio{31}โมหญฺจ ปจฺจยา ตโย ขนฺธา จิตฺต\-สมุฏฺฐานญฺจ รูปํ ฯเปฯ เทฺว ขนฺเธ ฯเปฯ. ปฏิสนฺธิ\-กฺขเณ สเหตุกํ เอกํ ขนฺธญฺจ วตฺถุญฺจ ปจฺจยา ตโย ขนฺธา ฯเปฯ เทฺว ขนฺเธ ฯเปฯ. สเหตุเก ขนฺเธ จ มหาภูเต จ ปจฺจยา กฏตฺตารูปํ. (3)}

\proseitem{74}{สเหตุกํ ธมฺมํ ปจฺจยา สเหตุโก ธมฺโม อุปฺปชฺชติ อารมฺมณ\-ปจฺจยา. สเหตุกํ เอกํ ขนฺธํ ปจฺจยา ตโย ขนฺธา ฯเปฯ. ปฏิสนฺธิ\-กฺขเณ ฯเปฯ. (1)}

\prose{สเหตุกํ ธมฺมํ ปจฺจยา อเหตุโก ธมฺโม อุปฺปชฺชติ อารมฺมณ\-ปจฺจยา. วิจิกิจฺฉา\-สหคเต อุทฺธจฺจ\-สหคเต ขนฺเธ ปจฺจยา วิจิกิจฺฉา\-สหคโต อุทฺธจฺจ\-สหคโต โมโห. (2)}

\prose{สเหตุกํ ธมฺมํ ปจฺจยา สเหตุโก จ อเหตุโก จ ธมฺมา อุปฺปชฺชนฺติ อารมฺมณ\-ปจฺจยา. วิจิกิจฺฉา\-สหคตํ อุทฺธจฺจ\-สหคตํ เอกํ ขนฺธํ ปจฺจยา ตโย ขนฺธา โมโห จ ฯเปฯ เทฺว ขนฺเธ ฯเปฯ. (3)}

\proseitem{75}{อเหตุกํ ธมฺมํ ปจฺจยา อเหตุโก ธมฺโม อุปฺปชฺชติ อารมฺมณ\-ปจฺจยา. อเหตุกํ เอกํ ขนฺธํ ฯเปฯ เทฺว ขนฺเธ ฯเปฯ. ปฏิสนฺธิ\-กฺขเณ วตฺถุํ ปจฺจยา ขนฺธา. จกฺขายตนํ ปจฺจยา จกฺขุวิญฺญาณํ ฯเปฯ. กายายตนํ ปจฺจยา กายวิญฺญาณํ. วตฺถุํ ปจฺจยา อเหตุกา ขนฺธา. (1)}

\prose{อเหตุกํ ธมฺมํ ปจฺจยา สเหตุโก ธมฺโม อุปฺปชฺชติ อารมฺมณ\-ปจฺจยา. วตฺถุํ ปจฺจยา สเหตุกา ขนฺธา. วิจิกิจฺฉา\-สหคตํ อุทฺธจฺจ\-สหคตํ โมหํ ปจฺจยา สมฺปยุตฺตกา ขนฺธา. ปฏิสนฺธิ\-กฺขเณ ฯเปฯ. (2)}

\prose{อเหตุกํ ธมฺมํ ปจฺจยา สเหตุโก จ อเหตุโก จ ธมฺมา อุปฺปชฺชนฺติ อารมฺมณ\-ปจฺจยา. วตฺถุํ ปจฺจยา วิจิกิจฺฉา\-สหคตา อุทฺธจฺจ\-สหคตา ขนฺธา โมโห จ. (3)}

\proseitem{76}{สเหตุกญฺจ อเหตุกญฺจ ธมฺมํ ปจฺจยา สเหตุโก ธมฺโม อุปฺปชฺชติ อารมฺมณ\-ปจฺจยา. สเหตุกํ เอกํ ขนฺธญฺจ วตฺถุญฺจ ปจฺจยา ตโย ขนฺธา ฯเปฯ. วิจิกิจฺฉา\-สหคตํ อุทฺธจฺจ\-สหคตํ เอกํ ขนฺธญฺจ โมหญฺจ ปจฺจยา ตโย ขนฺธา ฯเปฯ เทฺว ขนฺเธ ฯเปฯ. ปฏิสนฺธิ\-กฺขเณ ฯเปฯ. (1)}

\prose{\csromanfolio{32}สเหตุกญฺจ อเหตุกญฺจ ธมฺมํ ปจฺจยา อเหตุโก ธมฺโม อุปฺปชฺชติ อารมฺมณ\-ปจฺจยา. วิจิกิจฺฉา\-สหคเต อุทฺธจฺจ\-สหคเต ขนฺเธ จ วตฺถุญฺจ ปจฺจยา วิจิกิจฺฉา\-สหคโต อุทฺธจฺจ\-สหคโต โมโห. (2)}

\prose{สเหตุกญฺจ อเหตุกญฺจ ธมฺมํ ปจฺจยา สเหตุโก จ อเหตุโก จ ธมฺมา อุปฺปชฺชนฺติ อารมฺมณ\-ปจฺจยา. วิจิกิจฺฉา\-สหคตํ อุทฺธจฺจ\-สหคตํ เอกํ ขนฺธญฺจ วตฺถุญฺจ ปจฺจยา ตโย ขนฺธา โมโห จ ฯเปฯ เทฺว ขนฺเธ ฯเปฯ. (3)}

\proseitem{77}{สเหตุกํ ธมฺมํ ปจฺจยา สเหตุโก ธมฺโม อุปฺปชฺชติ อธิปติ ปจฺจยา. (อธิปติยา นว ปญฺหา ปวตฺเตเยว.)}

\proseitem{78}{สเหตุกํ ธมฺมํ ปจฺจยา สเหตุโก ธมฺโม อุปฺปชฺชติ อนนฺตร\-ปจฺจยา. สมนนฺตร\-ปจฺจยา. สหชาต\-ปจฺจยา. ตีณิ. (ปฏิจฺจ\-วาร\-สทิสา.)}

\prose{อเหตุกํ ธมฺมํ ปจฺจยา อเหตุโก ธมฺโม อุปฺปชฺชติ สหชาต\-ปจฺจยา. อเหตุกํ เอกํ ขนฺธํ ปจฺจยา ตโย ขนฺธา จิตฺต\-สมุฏฺฐานญฺจ รูปํ ฯเปฯ เทฺว ขนฺเธ ฯเปฯ. ปฏิสนฺธิ\-กฺขเณ ฯเปฯ. (ยาว อสญฺญสตฺตา.) จกฺขายตนํ ปจฺจยา จกฺขุวิญฺญาณํ ฯเปฯ. กายายตนํ ปจฺจยา กายวิญฺญาณํ. วตฺถุํ ปจฺจยา อเหตุกา ขนฺธา. (1)}

\prose{อเหตุกํ ธมฺมํ ปจฺจยา สเหตุโก ธมฺโม อุปฺปชฺชติ สหชาต\-ปจฺจยา. วตฺถุํ ปจฺจยา สเหตุกา ขนฺธา. วิจิกิจฺฉา\-สหคตํ อุทฺธจฺจ\-สหคตํ โมหํ ปจฺจยา สมฺปยุตฺตกา ขนฺธา. ปฏิสนฺธิ\-กฺขเณ ฯเปฯ. (2)}

\prose{อเหตุกํ ธมฺมํ ปจฺจยา สเหตุโก จ อเหตุโก จ ธมฺมา อุปฺปชฺชนฺติ สหชาต\-ปจฺจยา. วตฺถุํ ปจฺจยา สเหตุกา ขนฺธา. มหาภูเต ปจฺจยา จิตฺต\-สมุฏฺฐานํ รูปํ. วิจิกิจฺฉา\-สหคตํ อุทฺธจฺจ\-สหคตํ โมหํ ปจฺจยา สมฺปยุตฺตกา ขนฺธา จิตฺต\-สมุฏฺฐานญฺจ รูปํ. ปฏิสนฺธิ\-กฺขเณ วตฺถุํ ฯเปฯ. (3)}

\proseitem{79}{\csromanfolio{33}สเหตุกญฺจ อเหตุกญฺจ ธมฺมํ ปจฺจยา สเหตุโก ธมฺโม อุปฺปชฺชติ สหชาต\-ปจฺจยา. สเหตุกํ เอกํ ขนฺธญฺจ วตฺถุญฺจ ปจฺจยา ตโย ขนฺธา ฯเปฯ เทฺว ขนฺเธ ฯเปฯ. วิจิกิจฺฉา\-สหคตํ อุทฺธจฺจ\-สหคตํ เอกํ ขนฺธญฺจ โมหญฺจ ปจฺจยา ตโย ขนฺธา ฯเปฯ เทฺว ขนฺเธ ฯเปฯ. ปฏิสนฺธิ\-กฺขเณ ฯเปฯ. (1)}

\prose{สเหตุกญฺจ อเหตุกญฺจ ธมฺมํ ปจฺจยา อเหตุโก ธมฺโม อุปฺปชฺชติ สหชาต\-ปจฺจยา. สเหตุเก ขนฺเธ จ มหาภูเต จ ปจฺจยา จิตฺต\-สมุฏฺฐานํ รูปํ. วิจิกิจฺฉา\-สหคเต อุทฺธจฺจ\-สหคเต ขนฺเธ จ โมหญฺจ ปจฺจยา จิตฺต\-สมุฏฺฐานํ รูปํ. ปฏิสนฺธิ\-กฺขเณ ฯเปฯ. วิจิกิจฺฉา\-สหคเต อุทฺธจฺจ\-สหคเต ขนฺเธ จ วตฺถุญฺจ ปจฺจยา วิจิกิจฺฉา\-สหคโต อุทฺธจฺจ\-สหคโต โมโห. (2)}

\prose{สเหตุกญฺจ อเหตุกญฺจ ธมฺมํ ปจฺจยา สเหตุโก จ อเหตุโก จ ธมฺมา อุปฺปชฺชนฺติ สหชาต\-ปจฺจยา. สเหตุกํ เอกํ ขนฺธญฺจ วตฺถุญฺจ ปจฺจยา ตโย ขนฺธา ฯเปฯ เทฺว ขนฺเธ ฯเปฯ. สเหตุเก ขนฺเธ จ มหาภูเต จ ปจฺจยา จิตฺต\-สมุฏฺฐานํ รูปํ. วิจิกิจฺฉา\-สหคตํ อุทฺธจฺจ\-สหคตํ เอกํ ขนฺธญฺจ วตฺถุญฺจ ปจฺจยา ตโย ขนฺธา ฯเปฯ เทฺว ขนฺเธ ฯเปฯ. สเหตุเก ขนฺเธ จ มหาภูเต จ ปจฺจยา จิตฺต\-สมุฏฺฐานํ รูปํ. วิจิกิจฺฉา\-สหคตํ อุทฺธจฺจ\-สหคตํ เอกํ ขนฺธญฺจ โมหญฺจ ปจฺจยา ตโย ขนฺธา จิตฺต\-สมุฏฺฐานญฺจ รูปํ ฯเปฯ เทฺว ขนฺเธ ฯเปฯ. วิจิกิจฺฉา\-สหคตํ อุทฺธจฺจ\-สหคตํ เอกํ ขนฺธญฺจ วตฺถุญฺจ ปจฺจยา ตโย ขนฺธา โมโห จ ฯเปฯ. ปฏิสนฺธิ\-กฺขเณ สเหตุกํ เอกํ ขนฺธญฺจ วตฺถุญฺจ ปจฺจยา ตโย ขนฺธา ฯเปฯ เทฺว ขนฺเธ ฯเปฯ. สเหตุเก ขนฺเธ จ มหาภูเต จ ปจฺจยา กฏตฺตารูปํ. (3)}

\csromanheader{2}{\textbf{อญฺญมญฺญาทิ}}
\proseitem{80}{สเหตุกํ ธมฺมํ ปจฺจยา สเหตุโก ธมฺโม อุปฺปชฺชติ อญฺญมญฺญ\-ปจฺจยา ฯเปฯ อวิคตปจฺจยา.}

\csromanmarkh{1. ปจฺจยานุโลม}\csromanmarkhii{2. สงฺขฺยาวาร}\csromanheadernomark{2}{1. \textbf{ปจฺจยานุโลม 2. สงฺขฺยาวาร}}
\proseitem{81}{เหตุยา นว, อารมฺมเณ นว, อธิปติยา นว, (สพฺพตฺถ นว,) อวิคเต นว. (เอวํ คเณตพฺพํ.)}

\vspace{\tipitakaparskip}
\csromancenter{อนุโลมํ.}
\csromanmarkh{2. ปจฺจย\-ปจฺจนีย}\csromanmarkhii{1. วิภงฺควาร}\csromanheadernomark{2}{2. ปจฺจย\-ปจฺจนีย 1. วิภงฺควาร}
\proseitem{82}{\csromanfolio{34}สเหตุกํ ธมฺมํ ปจฺจยา อเหตุโก ธมฺโม อุปฺปชฺชติ นเหตุปจฺจยา. วิจิกิจฺฉา\-สหคเต อุทฺธจฺจ\-สหคเต ขนฺเธ ปจฺจยา วิจิกิจฺฉา\-สหคโต อุทฺธจฺจ\-สหคโต โมโห. (1)}

\prose{อเหตุกํ ธมฺมํ ปจฺจยา อเหตุโก ธมฺโม อุปฺปชฺชติ นเหตุปจฺจยา. อเหตุกํ เอกํ ขนฺธํ ฯเปฯ. ปฏิสนฺธิ\-กฺขเณ ฯเปฯ. (ยาว อสญฺญสตฺตา.) จกฺขายตนํ ปจฺจยา จกฺขุวิญฺญาณํ ฯเปฯ. กายายตนํ ปจฺจยา กายวิญฺญาณํ. วตฺถุํ ปจฺจยา อเหตุกา ขนฺธา โมโห จ. (1)}

\prose{สเหตุกญฺจ อเหตุกญฺจ ธมฺมํ ปจฺจยา อเหตุโก ธมฺโม อุปฺปชฺชติ นเหตุปจฺจยา. วิจิกิจฺฉา\-สหคเต อุทฺธจฺจ\-สหคเต ขนฺเธ จ วตฺถุญฺจ ปจฺจยา วิจิกิจฺฉา\-สหคโต อุทฺธจฺจ\-สหคโต โมโห. (สํขิตฺตํ.) (1)}

\csromanmarkh{2. ปจฺจย\-ปจฺจนีย}\csromanmarkhii{2. สงฺขฺยาวาร}\csromanheadernomark{2}{2. ปจฺจย\-ปจฺจนีย 2. สงฺขฺยาวาร}
\vspace{\tipitakaparskip}
\csromancenter{\textbf{นเหตุ}ยา ตีณิ, นอารมฺมเณ ตีณิ, นอธิปติยา นว, นอนนฺตเร ตีณิ, นสมนนฺตเร ตีณิ, นอญฺญมญฺเญ ตีณิ, นนิสฺสเย ตีณิ, นฺเอาปนิสฺสเย ตีณิ, นปุเรชาเต นว, นปจฺฉาชาเต นว, นอาเสวเน นว, นกมฺเม จตฺตาริ, นวิปาเก นว, นอาหาเร เอกํ, ไนนฺทฺริเย เอกํ, นฌาเน เอกํ, นมคฺเค เอกํ, นสมฺปยุตฺเต ตีณิ, นวิปฺปยุตฺเต ฉ, โนนตฺถิยา ตีณิ, โนวิคเต ตีณิ. (เอวํ คเณตพฺพํ.)}
\csromancenter{\textbf{ปจฺจนียํ.}}
\proseitem{84}{เหตุปจฺจยา นอารมฺมเณ ตีณิ, นอธิปติยา นว, นอนนฺตเร ตีณิ, นสมนนฺตเร ตีณิ, นอญฺญมญฺเญ ตีณิ, นอุปนิสฺสเย ตีณิ, \csromanfolio{35}นปุเรชาเต นว, นปจฺฉาชาเต นว, นอาเสวเน นว, นกมฺเม ตีณิ, นวิปาเก นว, นสมฺปยุตฺเต ตีณิ, นวิปฺปยุตฺเต ตีณิ, โนนตฺถิยา ตีณิ, โนวิคเต ตีณิ. (เอวํ คเณตพฺพํ.)}

\vspace{\tipitakaparskip}
\csromancenter{อนุโลม\-ปจฺจนียํ.}
\proseitem{85}{นเหตุปจฺจยา อารมฺมเณ ตีณิ, อนนฺตเร ตีณิ ฯเปฯ มคฺเค ตีณิ, สมฺปยุตฺเต ตีณิ ฯเปฯ อวิคเต ตีณิ. (เอวํ คเณตพฺพํ.)}

\vspace{\tipitakaparskip}
\csromancenter{ปจฺจนียา\-นุโลมํ.}
\csromancenter{(นิสฺสยวาโร ปจฺจยวาร\-สทิโส.)}
\csromanmarkh{2. สเหตุกทุก}\csromanmarkhii{5. สํสฏฺฐวาร}\csromanheadernomark{2}{2. \textbf{สเหตุกทุก 5. สํสฏฺฐวาร}}
\csromanmarkh{1. ปจฺจยานุโลม}\csromanmarkhii{1. วิภงฺควาร}\csromanheadernomark{2}{1. \textbf{ปจฺจยานุโลม 1. วิภงฺควาร}}
\csromanheader{2}{\textbf{เหตุ}}
\proseitem{86}{สเหตุกํ ธมฺมํ สํสฏฺโฐ สเหตุโก ธมฺโม อุปฺปชฺชติ เหตุปจฺจยา. สเหตุกํ เอกํ ขนฺธํ สํสฏฺฐา ตโย ขนฺธา ฯเปฯ เทฺว ขนฺเธ ฯเปฯ. ปฏิสนฺธิ\-กฺขเณ ฯเปฯ. (1)}

\prose{อเหตุกํ ธมฺมํ สํสฏฺโฐ สเหตุโก ธมฺโม อุปฺปชฺชติ เหตุปจฺจยา. วิจิกิจฺฉา\-สหคตํ อุทฺธจฺจ\-สหคตํ โมหํ สํสฏฺฐา วิจิกิจฺฉา\-สหคตา อุทฺธจฺจ\-สหคตา ขนฺธา. (1)}

\prose{สเหตุกญฺจ อเหตุกญฺจ ธมฺมํ สํสฏฺโฐ สเหตุโก ธมฺโม อุปฺปชฺชติ เหตุปจฺจยา. วิจิกิจฺฉา\-สหคตํ อุทฺธจฺจ\-สหคตํ เอกํ ขนฺธญฺจ โมหญฺจ สํสฏฺฐา ตโย ขนฺธา ฯเปฯ เทฺว ขนฺเธ ฯเปฯ. (1)}

\proseitem{87}{\csromanfolio{36}สเหตุกํ ธมฺมํ สํสฏฺโฐ อเหตุโก ธมฺโม อุปฺปชฺชติ. อารมฺมณ\-ปจฺจยา. สเหตุกํ เอกํ ขนฺธํ สํสฏฺฐา ตโย ขนฺธา ฯเปฯ เทฺว ขนฺเธ ฯเปฯ. ปฏิสนฺธิ\-กฺขเณ ฯเปฯ. (1)}

\prose{สเหตุกํ ธมฺมํ สํสฏฺโฐ อเหตุโก ธมฺโม อุปฺปชฺชติ อารมฺมณ\-ปจฺจยา. วิจิกิจฺฉา\-สหคเต อุทฺธจฺจ\-สหคเต ขนฺเธ สํสฏฺโฐ วิจิกิจฺฉา\-สหคโต อุทฺธจฺจ\-สหคโต โมโห. (2)}

\prose{สเหตุกํ ธมฺมํ สํสฏฺโฐ สเหตุโก จ อเหตุโก จ ธมฺมา อุปฺปชฺชนฺติ อารมฺมณ\-ปจฺจยา. วิจิกิจฺฉา\-สหคตํ อุทฺธจฺจ\-สหคตํ เอกํ ขนฺธํ สํสฏฺฐา ตโย ขนฺธา โมโห จ ฯเปฯ เทฺว ขนฺเธ ฯเปฯ. (3)}

\proseitem{88}{อเหตุกํ ธมฺมํ สํสฏฺโฐ อเหตุโก ธมฺโม อุปฺปชฺชติ อารมฺมณ\-ปจฺจยา. อเหตุกํ เอกํ ขนฺธํ สํสฏฺฐา ตโย ขนฺธา ฯเปฯ เทฺว ขนฺเธ ฯเปฯ. ปฏิสนฺธิ\-กฺขเณ ฯเปฯ. (1)}

\prose{อเหตุกํ ธมฺมํ สํสฏฺโฐ สเหตุโก ธมฺโม อุปฺปชฺชติ อารมฺมณ\-ปจฺจยา. วิจิกิจฺฉา\-สหคตํ อุทฺธจฺจ\-สหคตํ โมหํ สํสฏฺฐา วิจิกิจฺฉา\-สหคตา อุทฺธจฺจ\-สหคตา ขนฺธา. (2)}

\prose{สเหตุกญฺจ อเหตุกญฺจ ธมฺมํ สํสฏฺโฐ สเหตุโก ธมฺโม อุปฺปชฺชติ อารมฺมณ\-ปจฺจยา. วิจิกิจฺฉา\-สหคตํ อุทฺธจฺจ\-สหคตํ เอกํ ขนฺธญฺจ โมหญฺจ สํสฏฺฐา ตโย ขนฺธา ฯเปฯ เทฺว ขนฺเธ ฯเปฯ. (1)}

\proseitem{89}{สเหตุกํ ธมฺมํ สํสฏฺโฐ สเหตุโก ธมฺโม อุปฺปชฺชติ อธิปติ\-ปจฺจยา. สเหตุกํ เอกํ ขนฺธํ สํสฏฺฐา ตโย ขนฺธา ฯเปฯ เทฺว ขนฺเธ ฯเปฯ. (1)}

\proseitem{90}{สเหตุกํ ธมฺมํ สํสฏฺโฐ สเหตุโก ธมฺโม อุปฺปชฺชติ อนนฺตร\-ปจฺจยา. สมนนฺตร\-ปจฺจยา. สหชาต\-ปจฺจยา ฯเปฯ วิปากปจฺจยา. วิปากํ สเหตุกํ เอกํ ขนฺธํ สํสฏฺฐา ตโย ขนฺธา ฯเปฯ เทฺว ขนฺเธ ฯเปฯ. ปฏิสนฺธิ\-กฺขเณ ฯเปฯ.}

\prose{\csromanfolio{37}อเหตุกํ ธมฺมํ สํสฏฺโฐ อเหตุโก ธมฺโม อุปฺปชฺชติ วิปากปจฺจยา. วิปากํ อเหตุกํ เอกํ ขนฺธํ สํสฏฺฐา ตโย ขนฺธา ฯเปฯ เทฺว ขนฺเธ ฯเปฯ. ปฏิสนฺธิ\-กฺขเณ ฯเปฯ ฌานปจฺจยา ฯเปฯ อวิคตปจฺจยา.}

\csromanmarkh{1. ปจฺจยานุโลม}\csromanmarkhii{2. สงฺขฺยาวาร}\csromanheadernomark{2}{1. \textbf{ปจฺจยานุโลม 2. สงฺขฺยาวาร}}
\proseitem{91}{เหตุยา ตีณิ, อารมฺมเณ ฉ, อธิปติยา เอกํ, อนนฺตเร ฉ, สมนนฺตเร ฉ, สหชาเต ฉ, อญฺญมญฺเญ ฉ, นิสฺสเย ฉ, อุปนิสฺสเย ฉ, ปุเรชาเต ฉ ฯเปฯ วิปาเก เทฺว, อาหาเร ฉ, อินฺทฺริเย ฉ, ฌาเน ฉ, มคฺเค ปญฺจ ฯเปฯ อวิคเต ฉ. (เอวํ คเณตพฺพํ.)}

\vspace{\tipitakaparskip}
\csromancenter{\textbf{อนุโลมํ.}}
\csromanmarkh{2. ปจฺจย\-ปจฺจนีย}\csromanmarkhii{1. วิภงฺควาร}\csromanheadernomark{2}{2. ปจฺจย\-ปจฺจนีย 1. วิภงฺควาร}
\proseitem{92}{สเหตุกํ ธมฺมํ สํสฏฺโฐ อเหตุโก ธมฺโม อุปฺปชฺชติ นเหตุปจฺจยา. วิจิกิจฺฉา\-สหคเต อุทฺธจฺจ\-สหคเต ขนฺเธ สํสฏฺโฐ วิจิกิจฺฉา\-สหคโต อุทฺธจฺจ\-สหคโต โมโห. (1)}

\prose{อเหตุกํ ธมฺมํ สํสฏฺโฐ อเหตุโก ธมฺโม อุปฺปชฺชติ นเหตุปจฺจยา. อเหตุกํ เอกํ ขนฺธํ สํสฏฺฐา ตโย ขนฺธา ฯเปฯ เทฺว ขนฺเธ ฯเปฯ. ปฏิสนฺธิ\-กฺขเณ ฯเปฯ. (สํขิตฺตํ.) (1)}

\csromanmarkh{2. ปจฺจย\-ปจฺจนีย}\csromanmarkhii{2. สงฺขฺยาวาร}\csromanheadernomark{2}{2. ปจฺจย\-ปจฺจนีย 2. สงฺขฺยาวาร}
\vspace{\tipitakaparskip}
\csromancenter{\textbf{นเหตุ}ยา เทฺว, นอธิปติยา ฉ, นปุเรชาเต ฉ, นปจฺฉาชาเต ฉ, นอาเสวเน ฉ, นกมฺเม จตฺตาริ, นวิปาเก ฉ, นฌาเน เอกํ, นมคฺเค เอกํ, นวิปฺปยุตฺเต ฉ. (เอวํ คเณตพฺพํ.)}
\csromancenter{ปจฺจนียํ.}
\csromancenter{\textbf{เหตุ}ปจฺจยา นอธิปติยา ตีณิ, นปุเรชาเต ตีณิ, นปจฺฉาชาเต ตีณิ, นอาเสวเน ตีณิ, นกมฺเม ตีณิ, นวิปาเก ตีณิ, นวิปฺปยุตฺเต ตีณิ.}
\csromancenter{อนุโลม\-ปจฺจนียํ.}
\proseitem{95}{\csromanfolio{38}นเหตุปจฺจยา อารมฺมเณ เทฺว, อนนฺตเร เทฺว ฯเปฯ กมฺเม เทฺว, วิปาเก เอกํ, อาหาเร เทฺว ฯเปฯ มคฺเค เอกํ ฯเปฯ อวิคเต เทฺว. (เอวํ คเณตพฺพํ)}

\vspace{\tipitakaparskip}
\csromancenter{ปจฺจนียา\-นุโลมํ.}
\csromancenter{(สมฺปยุตฺตวาโร สํสฏฺฐวาร\-สทิโส.)}
\csromanmarkh{2. สเหตุกทุก}\csromanmarkhii{7. ปญฺหาวาร}\csromanheadernomark{2}{2. \textbf{สเหตุกทุก 7. ปญฺหาวาร}}
\csromanmarkh{1. ปจฺจยานุโลม}\csromanmarkhii{1. วิภงฺควาร}\csromanheadernomark{2}{1. \textbf{ปจฺจยานุโลม 1. วิภงฺควาร}}
\csromanheader{2}{\textbf{เหตุ}}
\proseitem{96}{สเหตุโก ธมฺโม สเหตุกสฺส ธมฺมสฺส เหตุปจฺจเยน ปจฺจโย. สเหตุกา เหตู สมฺปยุตฺตกานํ ขนฺธานํ เหตุปจฺจเยน ปจฺจโย. ปฏิสนฺธิ\-กฺขเณ ฯเปฯ. (1)}

\prose{สเหตุโก ธมฺโม อเหตุกสฺส ธมฺมสฺส เหตุปจฺจเยน ปจฺจโย. สเหตุกา เหตู จิตฺต\-สมุฏฺฐานานํ รูปานํ เหตุปจฺจเยน ปจฺจโย. ปฏิสนฺธิ\-กฺขเณ ฯเปฯ. (2)}

\prose{\csromanfolio{39}สเหตุโก ธมฺโม สเหตุกสฺส จ อเหตุกสฺส จ ธมฺมสฺส เหตุปจฺจเยน ปจฺจโย. สเหตุกา เหตู สมฺปยุตฺตกานํ ขนฺธานํ จิตฺตสมุฏฺฐานานญฺจ รูปานํ เหตุปจฺจเยน ปจฺจโย. ปฏิสนฺธิ\-กฺขเณ ฯเปฯ. (3)}

\proseitem{97}{อเหตุโก ธมฺโม อเหตุกสฺส ธมฺมสฺส เหตุปจฺจเยน ปจฺจโย. วิจิกิจฺฉา\-สหคโต อุทฺธจฺจ\-สหคโต โมโห จิตฺต\-สมุฏฺฐานานํ รูปานํ เหตุปจฺจเยน ปจฺจโย. (1)}

\prose{อเหตุโก ธมฺโม สเหตุกสฺส ธมฺมสฺส เหตุปจฺจเยน ปจฺจโย. วิจิกิจฺฉา\-สหคโต อุทฺธจฺจ\-สหคโต โมโห สมฺปยุตฺตกานํ ขนฺธานํ เหตุปจฺจเยน ปจฺจโย. (2)}

\prose{อเหตุโก ธมฺโม สเหตุกสฺส จ อเหตุกสฺส จ ธมฺมสฺส เหตุปจฺจเยน ปจฺจโย. วิจิกิจฺฉา\-สหคโต อุทฺธจฺจ\-สหคโต โมโห สมฺปยุตฺตกานํ ขนฺธานํ จิตฺต\-สมุฏฺฐานญฺจ รูปานํ เหตุปจฺจเยน ปจฺจโย. (3)}

\proseitem{98}{สเหตุโก ธมฺโม สเหตุกสฺส ธมฺมสฺส อารมฺมณ\-ปจฺจเยน ปจฺจโย. ทานํ ทตฺวา, สีลํ ฯเปฯ อุโปสถกมฺมํ กตฺวา ตํ ปจฺจเวกฺขติ, ปุพฺเพ สุจิณฺณานิ ปจฺจเวกฺขติ. ฌานา วุฏฺฐหิตฺวา ฌานํ ปจฺจเวกฺขติ. อริยา มคฺคา วุฏฺฐหิตฺวา มคฺคํ ปจฺจเวกฺขนฺติ, ผลํ ปจฺจเวกฺขนฺติ. ปหีเน กิเลเส ฯเปฯ. วิกฺขมฺภิเต กิเลเส ฯเปฯ. ปุพฺเพ ฯเปฯ. สเหตุเก ขนฺเธ อนิจฺจโต ฯเปฯ โทมนสฺสํ อุปฺปชฺชติ. กุสลากุสเล นิรุทฺเธ สเหตุโก วิปาโก ตทารมฺมณตา อุปฺปชฺชติ. เจโตปริย\-ญาเณน สเหตุก\-จิตฺต\-สมงฺคิสฺส จิตฺตํ ชานนฺติ. อากาสา\-นญฺจา\-ยตนํ วิญฺญาณญฺจา\-ยตนสฺส ฯเปฯ. อากิญฺจญฺญา\-ยตนํ เนว\-สญฺญา\-นา\-สญฺญา\-ยตนสฺส ฯเปฯ. สเหตุกา ขนฺธา อิทฺธิวิธ\-ญาณสฺส, เจโตปริย\-ญาณสฺส, ปุพฺเพ\-นิวาสา\-นุสฺสติ\-ญาณสฺส, ยถากมฺมูปคญาณสฺส, อนาคตํสญาณสฺส อารมฺมณ\-ปจฺจเยน ปจฺจโย. สเหตุเก ขนฺเธ อารพฺภ สเหตุกา ขนฺธา อุปฺปชฺชนฺติ. (1)}

\prose{สเหตุโก ธมฺโม อเหตุกสฺส ธมฺมสฺส อารมฺมณ\-ปจฺจเยน ปจฺจโย. สเหตุเก ขนฺเธ อนิจฺจโต ฯเปฯ โทมนสฺสํ อุปฺปชฺชติ. กุสลากุสเล นิรุทฺเธ อเหตุโก วิปาโก ตทารมฺมณตา อุปฺปชฺชติ. สเหตุเก ขนฺเธ อารพฺภ อเหตุกา ขนฺธา จ โมโห จ อุปฺปชฺชนฺติ. (2)}

\prose{\csromanfolio{40}สเหตุโก ธมฺโม สเหตุกสฺส จ อเหตุกสฺส จ ธมฺมสฺส อารมฺมณ\-ปจฺจเยน ปจฺจโย. สเหตุเก ขนฺเธ อารพฺภ วิจิกิจฺฉา\-สหคตา อุทฺธจฺจ\-สหคตา ขนฺธา จ โมโห จ อุปฺปชฺชนฺติ. (3)}

\proseitem{99}{อเหตุโก ธมฺโม อเหตุกสฺส ธมฺมสฺส อารมฺมณ\-ปจฺจเยน ปจฺจโย. นิพฺพานํ อาวชฺชนาย อารมฺมณ\-ปจฺจเยน ปจฺจโย. จกฺขุํ ฯเปฯ. วตฺถุํ. อเหตุเก ขนฺเธ จ โมหญฺจ อนิจฺจโต ฯเปฯ โทมนสฺสํ อุปฺปชฺชติ. กุสลากุสเล นิรุทฺเธ อเหตุโก วิปาโก ตทารมฺมณตา อุปฺปชฺชติ. รูปายตนํ จกฺขุ\-วิญฺญาณสฺส ฯเปฯ. โผฏฺฐพฺพา\-ยตนํ กายวิญฺญาณสฺส ฯเปฯ. อเหตุเก ขนฺเธ จ โมหญฺจ อารพฺภ อเหตุกา ขนฺธา จ โมโห จ อุปฺปชฺชนฺติ. (1)}

\prose{อเหตุโก ธมฺโม สเหตุกสฺส ธมฺมสฺส อารมฺมณ\-ปจฺจเยน ปจฺจโย. อริยา นิพฺพานํ ปจฺจเวกฺขนฺติ. นิพฺพานํ โคตฺรภุสฺส, โวทานสฺส, มคฺคสฺส, ผลสฺส อารมฺมณ\-ปจฺจเยน ปจฺจโย. อริยา อเหตุเก ปหีเน กิเลเส ปจฺจเวกฺขนฺติ. วิกฺขมฺภิเต กิเลเส ฯเปฯ. ปุพฺเพ ฯเปฯ. จกฺขุํ ฯเปฯ วตฺถุํ ฯเปฯ. อเหตุเก ขนฺเธ จ โมหญฺจ อนิจฺจโต ฯเปฯ โทมนสฺสํ อุปฺปชฺชติ. กุสลากุสเล นิรุทฺเธ สเหตุโก วิปาโก ตทารมฺมณตา อุปฺปชฺชติ. ทิพฺเพน จกฺขุนา รูปํ ปสฺสติ. ทิพฺพาย โสตธาตุยา สทฺทํ สุณาติ. เจโตปริย\-ญาเณน อเหตุก\-จิตฺต\-สมงฺคิสฺส จิตฺตํ ชานนฺติ. อเหตุกา ขนฺธา อิทฺธิวิธ\-ญาณสฺส, เจโตปริย\-ญาณสฺส, ปุพฺเพ\-นิวาสา\-นุสฺสติ\-ญาณสฺส, อนาคตํสญาณสฺส อารมฺมณ\-ปจฺจเยน ปจฺจโย. อเหตุเก ขนฺเธ จ โมหญฺจ อารพฺภ สเหตุกา ขนฺธา อุปฺปชฺชนฺติ. (2)}

\prose{อเหตุโก ธมฺโม สเหตุกสฺส จ อเหตุกสฺส จ ธมฺมสฺส อารมฺมณ\-ปจฺจเยน ปจฺจโย. จกฺขุํ อารพฺภ วิจิกิจฺฉา\-สหคตา อุทฺธจฺจ\-สหคตา ขนฺธา จ โมโห จ อุปฺปชฺชนฺติ. โสตํ ฯเปฯ. วตฺถุํ. อเหตุเก ขนฺเธ จ โมหญฺจ อารพฺภ วิจิกิจฺฉา\-สหคตา อุทฺธจฺจ\-สหคตา ขนฺธา จ โมโห จ อุปฺปชฺชนฺติ. (3)}

\proseitem{100}{สเหตุโก จ อเหตุโก จ ธมฺมา สเหตุกสฺส ธมฺมสฺส อารมฺมณ\-ปจฺจเยน ปจฺจโย. วิจิกิจฺฉา\-สหคเต อุทฺธจฺจ\-สหคเต ขนฺเธ จ โมหญฺจ อารพฺภ สเหตุกา ขนฺธา อุปฺปชฺชนฺติ. (1)}

\prose{\csromanfolio{41}สเหตุโก จ อเหตุโก จ ธมฺมา อเหตุกสฺส ธมฺมสฺส อารมฺมณ\-ปจฺจเยน ปจฺจโย. วิจิกิจฺฉา\-สหคเต อุทฺธจฺจ\-สหคเต ขนฺเธ จ โมหญฺจ อารพฺภ อเหตุกา ขนฺธา จ โมโห จ อุปฺปชฺชนฺติ. (2)}

\prose{สเหตุโก จ อเหตุโก จ ธมฺมา สเหตุกสฺส จ อเหตุกสฺส จ ธมฺมสฺส อารมฺมณ\-ปจฺจเยน ปจฺจโย. วิจิกิจฺฉา\-สหคเต อุทฺธจฺจ\-สหคเต ขนฺเธ จ โมหญฺจ อารพฺภ วิจิกิจฺฉา\-สหคตา อุทฺธจฺจ\-สหคตา ขนฺธา จ โมโห จ อุปฺปชฺชนฺติ. (3)}

\proseitem{101}{สเหตุโก ธมฺโม สเหตุกสฺส ธมฺมสฺส อธิปติ\-ปจฺจเยน ปจฺจโย. อารมฺมณา\-ธิปติ, สหชาธิปติ.  อารมฺมณา\-ธิปติ\csromandash{}ทานํ ทตฺวา, สีลํ ฯเปฯ อุโปสถกมฺมํ กตฺวา ตํ ครุํ กตฺวา ปจฺจเวกฺขติ. ฌานา ฯเปฯ. อริยา มคฺคา วุฏฺฐหิตฺวา มคฺคํ ครุํ กตฺวา ฯเปฯ. ผลํ ฯเปฯ. สเหตุเก ขนฺเธ ครุํ กตฺวา อสฺสาเทติ อภินนฺทติ, ตํ ครุํ กตฺวา ราโค อุปฺปชฺชติ, ทิฏฺฐิ อุปฺปชฺชติ. สหชาตา\-ธิปติ\csromandash{}สเหตุกา\-ธิปติ สมฺปยุตฺตกานํ ขนฺธานํ อธิปติ\-ปจฺจเยน ปจฺจโย. (1)}

\prose{สเหตุโก ธมฺโม อเหตุกสฺส ธมฺมสฺส อธิปติ\-ปจฺจเยน ปจฺจโย. สหชาตา\-ธิปติ\csromandash{}สเหตุกา\-ธิปติ จิตฺต\-สมุฏฺฐานานํ รูปานํ อธิปติ\-ปจฺจเยน ปจฺจโย. (2)}

\prose{สเหตุโก ธมฺโม สเหตุกสฺส จ อเหตุกสฺส จ ธมฺมสฺส อธิปติ\-ปจฺจเยน ปจฺจโย. สหชาตา\-ธิปติ\csromandash{}สเหตุกา\-ธิปติ สมฺปยุตฺตกานํ ขนฺธานํ จิตฺตสมุฏฺฐานานญฺจ รูปานํ อธิปติ\-ปจฺจเยน ปจฺจโย. (3)}

\prose{อเหตุโก ธมฺโม สเหตุกสฺส ธมฺมสฺส อธิปติ\-ปจฺจเยน ปจฺจโย. อารมฺมณา\-ธิปติ\csromandash{}อริยา นิพฺพานํ ครุํ กตฺวา ปจฺจเวกฺขนฺติ. นิพฺพานํ โคตฺรภุสฺส, โวทานสฺส, มคฺคสฺส, ผลสฺส อธิปติ\-ปจฺจเยน ปจฺจโย. จกฺขุํ ฯเปฯ วตฺถุํ. อเหตุเก ขนฺเธ ครุํ กตฺวา อสฺสาเทติ อภินนฺทติ, ตํ ครุํ กตฺวา ราโค อุปฺปชฺชติ, ทิฏฺฐิ อุปฺปชฺชติ. (1)}

\proseitem{102}{\csromanfolio{42}สเหตุโก ธมฺโม สเหตุกสฺส ธมฺมสฺส อนนฺตร\-ปจฺจเยน ปจฺจโย. ปุริมา ปุริมา สเหตุกา ขนฺธา ปจฺฉิมานํ ปจฺฉิมานํ สเหตุกานํ ขนฺธานํ อนนฺตร\-ปจฺจเยน ปจฺจโย. อนุโลมํ โคตฺรภุสฺส. อนุโลมํ โวทานสฺส. โคตฺรภุ มคฺคสฺส. โวทานํ มคฺคสฺส. มคฺโค ผลสฺส. ผลํ ผลสฺส. อนุโลมํ ผลสมาปตฺติยา. นิโรธา วุฏฺฐหนฺตสฺส เนว\-สญฺญา\-นา\-สญฺญา\-ยตนํ ผลสมาปตฺติยา อนนฺตร\-ปจฺจเยน ปจฺจโย. (1)}

\prose{สเหตุโก ธมฺโม อเหตุกสฺส ธมฺมสฺส อนนฺตร\-ปจฺจเยน ปจฺจโย. ปุริมา ปุริมา วิจิกิจฺฉา\-สหคตา อุทฺธจฺจ\-สหคตา ขนฺธา ปจฺฉิมสฺส ปจฺฉิมสฺส วิจิกิจฺฉา\-สหคตสฺส อุทฺธจฺจ\-สหคตสฺส โมหสฺส อนนฺตร\-ปจฺจเยน ปจฺจโย. สเหตุกํ จุติจิตฺตํ อเหตุกสฺส อุปปตฺติ\-จิตฺตสฺส อนนฺตร\-ปจฺจเยน ปจฺจโย. สเหตุกํ ภวงฺคํ อาวชฺชนาย อนนฺตร\-ปจฺจเยน ปจฺจโย. สเหตุกํ ภวงฺคํ อเหตุกสฺส ภวงฺคสฺส อนนฺตร\-ปจฺจเยน ปจฺจโย. สเหตุกา ขนฺธา อเหตุกสฺส วุฏฺฐานสฺส อนนฺตร\-ปจฺจเยน ปจฺจโย. (2)}

\prose{สเหตุโก ธมฺโม สเหตุกสฺส จ อเหตุกสฺส จ ธมฺมสฺส อนนฺตร\-ปจฺจเยน ปจฺจโย. ปุริมา ปุริมา วิจิกิจฺฉา\-สหคตา อุทฺธจฺจ\-สหคตา ขนฺธา ปจฺฉิมานํ ปจฺฉิมานํ วิจิกิจฺฉา\-สหคตานํ อุทฺธจฺจ\-สหคตานํ ขนฺธานํ โมหสฺส จ อนนฺตร\-ปจฺจเยน ปจฺจโย. (3)}

\proseitem{103}{อเหตุโก ธมฺโม อเหตุกสฺส ธมฺมสฺส อนนฺตร\-ปจฺจเยน ปจฺจโย. ปุริโม ปุริโม วิจิกิจฺฉา\-สหคโต อุทฺธจฺจ\-สหคโต โมโห ปจฺฉิมสฺส ปจฺฉิมสฺส วิจิกิจฺฉา\-สหคตสฺส อุทฺธจฺจ\-สหคตสฺส โมหสฺส อนนฺตร\-ปจฺจเยน ปจฺจโย. ปุริมา ปุริมา อเหตุกา ขนฺธา ปจฺฉิมานํ ปจฺฉิมานํ อเหตุกานํ ขนฺธานํ อนนฺตร\-ปจฺจเยน ปจฺจโย. อาวชฺชนา ปญฺจนฺนํ วิญฺญาณานํ อนนฺตร\-ปจฺจเยน ปจฺจโย. (1)}

\prose{อเหตุโก ธมฺโม สเหตุกสฺส ธมฺมสฺส อนนฺตร\-ปจฺจเยน ปจฺจโย. ปุริโม ปุริโม วิจิกิจฺฉา\-สหคโต อุทฺธจฺจ\-สหคโต โมโห ปจฺฉิมานํ ปจฺฉิมานํ วิจิกิจฺฉา\-สหคตานํ อุทฺธจฺจ\-สหคตานํ ขนฺธานํ อนนฺตร\-ปจฺจเยน \csromanfolio{43}ปจฺจโย. อเหตุกํ จุติจิตฺตํ สเหตุกสฺส อุปปตฺติ\-จิตฺตสฺส อนนฺตร\-ปจฺจเยน ปจฺจโย. อเหตุกํ ภวงฺคํ สเหตุกสฺส ภวงฺคสฺส อนนฺตร\-ปจฺจเยน ปจฺจโย. อาวชฺชนา สเหตุกานํ ขนฺธานํ อนนฺตร\-ปจฺจเยน ปจฺจโย. อเหตุกา ขนฺธา สเหตุกสฺส วุฏฺฐานสฺส อนนฺตร\-ปจฺจเยน ปจฺจโย. (2)}

\prose{อเหตุโก ธมฺโม สเหตุกสฺส จ อเหตุกสฺส จ ธมฺมสฺส อนนฺตร\-ปจฺจเยน ปจฺจโย. ปุริโม ปุริโม วิจิกิจฺฉา\-สหคโต อุทฺธจฺจ\-สหคโต โมโห ปจฺฉิมานํ ปจฺฉิมานํ วิจิกิจฺฉา\-สหคตานํ อุทฺธจฺจ\-สหคตานํ ขนฺธานํ โมหสฺส จ อนนฺตร\-ปจฺจเยน ปจฺจโย. อาวชฺชนา วิจิกิจฺฉา\-สหคตานํ อุทฺธจฺจ\-สหคตานํ ขนฺธานํ โมหสฺส จ อนนฺตร\-ปจฺจเยน ปจฺจโย. (3)}

\proseitem{104}{สเหตุโก จ อเหตุโก จ ธมฺมา สเหตุกสฺส ธมฺมสฺส อนนฺตร\-ปจฺจเยน ปจฺจโย. ปุริมา ปุริมา วิจิกิจฺฉา\-สหคตา อุทฺธจฺจ\-สหคตา ขนฺธา จ โมโห จ ปจฺฉิมานํ ปจฺฉิมานํ วิจิกิจฺฉา\-สหคตานํ อุทฺธจฺจ\-สหคตานํ ขนฺธานํ อนนฺตร\-ปจฺจเยน ปจฺจโย. (1)}

\prose{สเหตุโก จ อเหตุโก จ ธมฺมา อเหตุกสฺส ธมฺมสฺส อนนฺตร\-ปจฺจเยน ปจฺจโย. ปุริมา ปุริมา วิจิกิจฺฉา\-สหคตา อุทฺธจฺจ\-สหคตา ขนฺธา จ โมโห จ ปจฺฉิมสฺส ปจฺฉิมสฺส วิจิกิจฺฉา\-สหคตสฺส อุทฺธจฺจ\-สหคตสฺส โมหสฺส อนนฺตร\-ปจฺจเยน ปจฺจโย. วิจิกิจฺฉา\-สหคตา อุทฺธจฺจ\-สหคตา ขนฺธา จ โมโห จ อเหตุกสฺส วุฏฺฐานสฺส อนนฺตร\-ปจฺจเยน ปจฺจโย. (2)}

\prose{สเหตุโก จ อเหตุโก จ ธมฺมา สเหตุกสฺส จ อเหตุกสฺส จ ธมฺมสฺส อนนฺตร\-ปจฺจเยน ปจฺจโย. ปุริมา ปุริมา วิจิกิจฺฉา\-สหคตา อุทฺธจฺจสหคถา ขนฺธา จ โมโห จ ปจฺฉิมานํ ปจฺฉิมานํ วิจิกิจฺฉา\-สหคตานํ อุทฺธจฺจ\-สหคตานํ ขนฺธานํ โมหสฺส จ อนนฺตร\-ปจฺจเยน ปจฺจโย. (3)}

\csromanheader{2}{\textbf{สหชาตาทิ}}
\proseitem{105}{สเหตุโก ธมฺโม สเหตุกสฺส ธมฺมสฺส สหชาต\-ปจฺจเยน ปจฺจโย. (ปฏิจฺจวาเร สหชาตสทิสํ, อิห ฆฏนา นตฺถิ.)}

\vspace{\tipitakaparskip}
\csromancenter{(ปฏิจฺจวาเร สหชาต\-สทิสํ, อิห ฆฏนา นตฺถิ.)}
\prosecont{\csromanfolio{44}อญฺญมญฺญ\-ปจฺจเยน ปจฺจโย. (ปฏิจฺจ\-วาร\-สทิสํ.) นิสฺสย\-ปจฺจเยน ปจฺจโย.}

\vspace{\tipitakaparskip}
\csromancenter{(ปฏิจฺจวาเร นิสฺสยปจฺจย\-สทิสํ, อิห ฆฏนา นตฺถิ.)}
\csromanheader{2}{\textbf{อุปนิสฺสย}}
\proseitem{106}{สเหตุโก ธมฺโม สเหตุกสฺส ธมฺมสฺส อุปนิสฺสย\-ปจฺจเยน ปจฺจโย. อารมฺมณู\-ปนิสฺสโย, อนนฺตรู\-ปนิสฺสโย, ปกตู\-ปนิสฺสโย ฯเปฯ. ปกตู\-ปนิสฺสโย\csromandash{}สเหตุกา ขนฺธา สเหตุกานํ ขนฺธานํ อุปนิสฺสย\-ปจฺจเยน ปจฺจโย. (1)}

\prose{สเหตุโก ธมฺโม อเหตุกสฺส ธมฺมสฺส อุปนิสฺสย\-ปจฺจเยน ปจฺจโย. อนนฺตรู\-ปนิสฺสโย, ปกตู\-ปนิสฺสโย ฯเปฯ. ปกตู\-ปนิสฺสโย\csromandash{} สเหตุกา ขนฺธา อเหตุกานํ ขนฺธานํ โมหสฺส จ อุปนิสฺสย\-ปจฺจเยน ปจฺจโย. (2)}

\prose{สเหตุโก ธมฺโม สเหตุกสฺส จ อเหตุกสฺส จ ธมฺมสฺส อุปนิสฺสย\-ปจฺจเยน ปจฺจโย. อนนฺตรู\-ปนิสฺสโย, ปกตู\-ปนิสฺสโย ฯเปฯ. ปกตู\-ปนิสฺสโย\csromandash{}สเหตุกา ขนฺธา วิจิกิจฺฉา\-สหคตานํ อุทฺธจฺจ\-สหคตานํ ขนฺธานํ โมหสฺส จ อุปนิสฺสย\-ปจฺจเยน ปจฺจโย. (3)}

\proseitem{107}{อเหตุโก ธมฺโม อเหตุกสฺส ธมฺมสฺส อุปนิสฺสย\-ปจฺจเยน ปจฺจโย. อนนฺตรู\-ปนิสฺสโย, ปกตู\-ปนิสฺสโย ฯเปฯ. ปกตู\-ปนิสฺสโย\csromandash{} กายิกํ สุขํ กายิกสฺส สุขสฺส, กายิกสฺส ทุกฺขสฺส, โมหสฺส จ อุปนิสฺสย\-ปจฺจเยน ปจฺจโย. กายิกํ ทุกฺขํ. อุตุ. โภชนํ. เสนาสนํ กายิกสฺส สุขสฺส, กายิกสฺส ทุกฺขสฺส, โมหสฺส จ อุปนิสฺสย\-ปจฺจเยน ปจฺจโย. โมโห กายิกสฺส สุขสฺส, กายิกสฺส ทุกฺขสฺส, โมหสฺส จ อุปนิสฺสย\-ปจฺจเยน ปจฺจโย. กายิกํ สุขํ. กายิกํ ทุกฺขํ. อุตุ. โภชนํ เสนาสนํ โมโห จ กายิกสฺส สุขสฺส, กายิกสฺส ทุกฺขสฺส, โมหสฺส จ อุปนิสฺสย\-ปจฺจเยน ปจฺจโย. (1)}

\prose{อเหตุโก ธมฺโม สเหตุกสฺส ธมฺมสฺส อุปนิสฺสย\-ปจฺจเยน ปจฺจโย. อารมฺมณู\-ปนิสฺสโย, อนนฺตรู\-ปนิสฺสโย, ปกตู\-ปนิสฺสโย ฯเปฯ. ปกตู\-ปนิสฺสโย\csromandash{}กายิกํ สุขํ อุปนิสฺสาย ทานํ เทติ ฯเปฯ สํฆํ ภินฺทติ. กายิกํ ทุกฺขํ. อุตุํ. โภชนํ. เสนาสนํ. โมหํ \csromanfolio{45}อุปนิสฺสาย ทานํ เทติ ฯเปฯ สํฆํ ภินฺทติ. กายิกํ สุขํ ฯเปฯ โมโห จ สทฺธาย ฯเปฯ ปญฺญาย. ราคสฺส ฯเปฯ ปตฺถนาย. มคฺคสฺส ผลสมาปตฺติยา อุปนิสฺสย\-ปจฺจเยน ปจฺจโย. (2)}

\prose{อเหตุโก ธมฺโม สเหตุกสฺส จ อเหตุกสฺส จ ธมฺมสฺส อุปนิสฺสย\-ปจฺจเยน ปจฺจโย. อนนฺตรู\-ปนิสฺสโย, ปกตู\-ปนิสฺสโย ฯเปฯ. ปกตู\-ปนิสฺสโย\csromandash{}กายิกํ สุขํ โมโห จ วิจิกิจฺฉา\-สหคตานํ อุทฺธจฺจ\-สหคตานํ ขนฺธานํ โมหสฺส จ อุปนิสฺสย\-ปจฺจเยน ปจฺจโย. (3)}

\proseitem{108}{สเหตุโก จ อเหตุโก จ ธมฺมา สเหตุกสฺส ธมฺมสฺส อุปนิสฺสย\-ปจฺจเยน ปจฺจโย. อนนฺตรู\-ปนิสฺสโย, ปกตู\-ปนิสฺสโย ฯเปฯ. ปกตู\-ปนิสฺสโย\csromandash{}วิจิกิจฺฉา\-สหคตา อุทฺธจฺจ\-สหคตา ขนฺธา จ โมโห จ สเหตุกานํ ขนฺธานํ อุปนิสฺสย\-ปจฺจเยน ปจฺจโย. (1)}

\prose{สเหตุโก จ อเหตุโก จ ธมฺมา อเหตุกสฺส ธมฺมสฺส อุปนิสฺสย\-ปจฺจเยน ปจฺจโย. อนนฺตรู\-ปนิสฺสโย, ปกตู\-ปนิสฺสโย ฯเปฯ. ปกตู\-ปนิสฺสโย\csromandash{}วิจิกิจฺฉา\-สหคตา อุทฺธจฺจ\-สหคตา ขนฺธา จ โมโห จ อเหตุกานํ ขนฺธานํ โมหสฺส จ อุปนิสฺสย\-ปจฺจเยน ปจฺจโย. (2)}

\prose{สเหตุโก จ อเหตุโก จ ธมฺมา สเหตุกสฺส จ อเหตุกสฺส จ ธมฺมสฺส อุปนิสฺสย\-ปจฺจเยน ปจฺจโย. อนนฺตรู\-ปนิสฺสโย, ปกตู\-ปนิสฺสโย ฯเปฯ. ปกตู\-ปนิสฺสโย\csromandash{}วิจิกิจฺฉา\-สหคตา อุทฺธจฺจ\-สหคตา ขนฺธา จ โมโห จ วิจิกิจฺฉา\-สหคตานํ อุทฺธจฺจ\-สหคตานํ ขนฺธานํ โมหสฺส จ อุปนิสฺสย\-ปจฺจเยน ปจฺจโย. (3)}

\proseitem{109}{อเหตุโก ธมฺโม อเหตุกสฺส ธมฺมสฺส ปุเรชาต\-ปจฺจเยน ปจฺจโย. อารมฺมณ\-ปุเรชาตํ, วตฺถุ\-ปุเรชาตํ.  อารมฺมณ\-ปุเรชาตํ\csromandash{} จกฺขุํ ฯเปฯ วตฺถุํ. อนิจฺจโต ฯเปฯ โทมนสฺสํ อุปฺปชฺชติ. กุสลากุสเล นิรุทฺเธ อเหตุโก วิปาโก ตทารมฺมณตา อุปฺปชฺชติ. รูปายตนํ จกฺขุ\-วิญฺญาณสฺส ฯเปฯ. โผฏฺฐพฺพา\-ยตนํ กายวิญฺญาณสฺส ปุเรชาต\-ปจฺจเยน \csromanfolio{46}ปจฺจโย. วตฺถุ\-ปุเรชาตํ\csromandash{}จกฺขายตนํ จกฺขุ\-วิญฺญาณสฺส ฯเปฯ กายายตนํ กายวิญฺญาณสฺส. ปุเรชาตํ วตฺถุ อเหตุกานํ ขนฺธานํ โมหสฺส จ ปุเรชาต\-ปจฺจเยน ปจฺจโย. (1)}

\prose{อเหตุโก ธมฺโม สเหตุกสฺส ธมฺมสฺส ปุเรชาต\-ปจฺจเยน ปจฺจโย. อารมฺมณ\-ปุเรชาตํ, วตฺถุ\-ปุเรชาตํ.  อารมฺมณ\-ปุเรชาตํ\csromandash{}จกฺขุํ ฯเปฯ วตฺถุํ อนิจฺจโต ฯเปฯ โทมนสฺสํ อุปฺปชฺชติ. กุสลากุสเล นิรุทฺเธ สเหตุโก วิปาโก ตทารมฺมณตา อุปฺปชฺชติ. ทิพฺเพน จกฺขุนา รูปํ ปสฺสติ. ทิพฺพาย โสตธาตุยา สทฺทํ สุณาติ. วตฺถุ\-ปุเรชาตํ\csromandash{}วตฺถุ สเหตุกานํ ขนฺธานํ ปุเรชาต\-ปจฺจเยน ปจฺจโย. (2)}

\prose{อเหตุโก ธมฺโม สเหตุกสฺส จ อเหตุกสฺส จ ธมฺมสฺส ปุเรชาต\-ปจฺจเยน ปจฺจโย. อารมฺมณ\-ปุเรชาตํ, วตฺถุ\-ปุเรชาตํ.  อารมฺมณ\-ปุเรชาตํ\csromandash{}จกฺขุํ ฯเปฯ วตฺถุํ อารพฺภ วิจิกิจฺฉา\-สหคตา อุทฺธจฺจ\-สหคตา ขนฺธา จ โมโห จ อุปฺปชฺชนฺติ. วตฺถุ\-ปุเรชาตํ\csromandash{}วตฺถุ วิจิกิจฺฉา\-สหคตานํ อุทฺธจฺจ\-สหคตานํ ขนฺธานํ โมหสฺส จ ปุเรชาต\-ปจฺจเยน ปจฺจโย.}

\csromanheader{2}{\textbf{ปจฺฉาชาตาทิ}}
\proseitem{110}{สเหตุโก ธมฺโม อเหตุกสฺส ธมฺมสฺส ปจฺฉาชาต\-ปจฺจเยน ปจฺจโย. ปจฺฉาชาตา สเหตุกา ขนฺธา ปุเรชาตสฺส อิมสฺส กายสฺส ปจฺฉาชาต\-ปจฺจเยน ปจฺจโย. (1)}

\prose{อเหตุโก ธมฺโม อเหตุกสฺส ธมฺมสฺส ปจฺฉาชาต\-ปจฺจเยน ปจฺจโย. ปจฺฉาชาตา อเหตุกา ขนฺธา จ โมโห จ ปุเรชาตสฺส อิมสฺส กายสฺส ปจฺฉาชาต\-ปจฺจเยน ปจฺจโย. (1)}

\prose{สเหตุโก จ อเหตุโก จ ธมฺมา อเหตุกสฺส ธมฺมสฺส ปจฺฉาชาต\-ปจฺจเยน ปจฺจโย. ปจฺฉาชาตา วิจิกิจฺฉา\-สหคตา อุทฺธจฺจ\-สหคตา ขนฺธา จ โมโห จ ปุเรชาตสฺส อิมสฺส กายสฺส ปจฺฉาชาต\-ปจฺจเยน ปจฺจโย. (1)}

\prose{สเหตุโก ธมฺโม สเหตุกสฺส ธมฺมสฺส อาเสวน\-ปจฺจเยน ปจฺจโย. (อนนฺตร\-สทิสํ. อาวชฺชนมฺปิ ภวงฺคมฺปิ นตฺถิ, อาเสวนปจฺจเย วชฺเชตพฺพา นวปิ.)}

\csromanheader{2}{2. สเหตุทุก \textbf{กมฺม}}
\proseitem{111}{\csromanfolio{47}สเหตุโก ธมฺโม สเหตุกสฺส ธมฺมสฺส กมฺมปจฺจเยน ปจฺจโย. สหชาตา, นานากฺขณิกา.  สหชาตา สเหตุกา เจตนา สมฺปยุตฺตกานํ ขนฺธานํ กมฺมปจฺจเยน ปจฺจโย. ปฏิสนฺธิ\-กฺขเณ ฯเปฯ. นานากฺขณิกา สเหตุกา เจตนา วิปากานํ สเหตุกานํ ขนฺธานํ กมฺมปจฺจเยน ปจฺจโย. (1)}

\prose{สเหตุโก ธมฺโม อเหตุกสฺส ธมฺมสฺส กมฺมปจฺจเยน ปจฺจโย. สหชาตา, นานากฺขณิกา.  สหชาตา สเหตุกา เจตนา จิตฺต\-สมุฏฺฐานานํ รูปานํ กมฺมปจฺจเยน ปจฺจโย. ปฏิสนฺธิ\-กฺขเณ ฯเปฯ. นานากฺขณิกา สเหตุกา เจตนา วิปากานํ อเหตุกานํ ขนฺธานํ กฏตฺตา จ รูปานํ กมฺมปจฺจเยน ปจฺจโย. (2)}

\prose{สเหตุโก ธมฺโม สเหตุกสฺส จ อเหตุกสฺส จ ธมฺมสฺส กมฺมปจฺจเยน ปจฺจโย. สหชาตา, นานากฺขณิกา.  \textbf{สหชาตา} สเหตุกา เจตนา สมฺปยุตฺตกานํ ขนฺธานํ จิตฺตสมุฏฺฐานานญฺจ รูปานํ กมฺมปจฺจเยน ปจฺจโย. นานากฺขณิกา สเหตุกา เจตนา วิปากานํ สเหตุกานํ ขนฺธานํ กฏตฺตา จ รูปานํ กมฺมปจฺจเยน ปจฺจโย. (3)}

\prose{อเหตุโก ธมฺโม อเหตุกสฺส ธมฺมสฺส กมฺมปจฺจเยน ปจฺจโย. อเหตุกา เจตนา สมฺปยุตฺตกานํ ขนฺธานํ จิตฺตสมุฏฺฐานานญฺจ รูปานํ กมฺมปจฺจเยน ปจฺจโย. ปฏิสนฺธิ\-กฺขเณ ฯเปฯ. (1)}

\csromanheader{2}{\textbf{วิปาก}}
\proseitem{112}{สเหตุโก ธมฺโม สเหตุกสฺส ธมฺมสฺส วิปาก\-ปจฺจเยน ปจฺจโย. วิปาโก สเหตุโก เอโก ขนฺโธ ติณฺณนฺนํ ขนฺธานํ ฯเปฯ เทฺว ขนฺธา ฯเปฯ. ปฏิสนฺธิ\-กฺขเณ ฯเปฯ. (1)}

\prose{สเหตุโก ธมฺโม อเหตุกสฺส ธมฺมสฺส วิปาก\-ปจฺจเยน ปจฺจโย. วิปากา สเหตุกา ขนฺธา จิตฺต\-สมุฏฺฐานานํ รูปานํ วิปาก\-ปจฺจเยน ปจฺจโย. ปฏิสนฺธิ\-กฺขเณ ฯเปฯ. (2)}

\prose{สเหตุโก ธมฺโม สเหตุกสฺส จ อเหตุกสฺส จ ธมฺมสฺส วิปาก\-ปจฺจเยน ปจฺจโย. วิปาโก สเหตุโก เอโก ขนฺโธ ติณฺณนฺนํ ขนฺธานํ จิตฺตสมุฏฺฐานานญฺจ รูปานํ ฯเปฯ เทฺว ขนฺธา ฯเปฯ. ปฏิสนฺธิ\-กฺขเณ ฯเปฯ. (3)}

\prose{\csromanfolio{48}อเหตุโก ธมฺโม อเหตุกสฺส ธมฺมสฺส วิปาก\-ปจฺจเยน ปจฺจโย. วิปาโก อเหตุโก เอโก ขนฺโธ ติณฺณนฺนํ ขนฺธานํ จิตฺตสมุฏฺฐานานญฺจ รูปานํ ฯเปฯ เทฺว ขนฺธา ฯเปฯ. ปฏิสนฺธิ\-กฺขเณ ฯเปฯ. ขนฺธา วตฺถุสฺส วิปาก\-ปจฺจเยน ปจฺจโย. (1)}

\csromanheader{2}{\textbf{อาหาร}}
\proseitem{113}{สเหตุโก ธมฺโม สเหตุกสฺส ธมฺมสฺส อาหารปจฺจเยน ปจฺจโย. ตีณิ.}

\prose{อเหตุโก ธมฺโม อเหตุกสฺส ธมฺมสฺส อาหารปจฺจเยน ปจฺจโย. อเหตุกา อาหารา สมฺปยุตฺตกานํ ขนฺธานํ จิตฺตสมุฏฺฐานานญฺจ รูปานํ ฯเปฯ. ปฏิสนฺธิ\-กฺขเณ ฯเปฯ. กพฬีกาโร อาหาโร อิมสฺส กายสฺส อาหารปจฺจเยน ปจฺจโย. (1)}

\csromanheader{2}{\textbf{อินฺทฺริยาทิ}}
\proseitem{114}{สเหตุโก ธมฺโม สเหตุกสฺส ธมฺมสฺส อินฺทฺริย\-ปจฺจเยน ปจฺจโย. ตีณิ.}

\prose{อเหตุโก ธมฺโม อเหตุกสฺส ธมฺมสฺส อินฺทฺริย\-ปจฺจเยน ปจฺจโย. อเหตุกา อินฺทฺริยา สมฺปยุตฺตกานํ ขนฺธานํ จิตฺตสมุฏฺฐานานญฺจ รูปานํ ฯเปฯ. ปฏิสนฺธิ\-กฺขเณ ฯเปฯ. จกฺขุนฺทฺริยํ จกฺขุ\-วิญฺญาณสฺส ฯเปฯ. กายินฺทฺริยํ กายวิญฺญาณสฺส อินฺทฺริย\-ปจฺจเยน ปจฺจโย. รูปชีวิตินฺทฺริยํ กฏตฺตารูปานํ อินฺทฺริย\-ปจฺจเยน ปจฺจโย. (1)}

\prose{สเหตุโก ธมฺโม สเหตุกสฺส ธมฺมสฺส ฌานปจฺจเยน ปจฺจโย. ตีณิ.}

\prose{อเหตุโก ธมฺโม อเหตุกสฺส ธมฺมสฺส ฌานปจฺจเยน ปจฺจโย. อเหตุกานิ ฌานงฺคานิ สมฺปยุตฺตกานํ ขนฺธานํ จิตฺตสมุฏฺฐานานญฺจ รูปานํ ฯเปฯ. ปฏิสนฺธิ\-กฺขเณ ฯเปฯ. (1)}

\prose{สเหตุโก ธมฺโม สเหตุกสฺส ธมฺมสฺส มคฺคปจฺจเยน ปจฺจโย. ตีณิ.}

\prose{สเหตุโก ธมฺโม สเหตุกสฺส ธมฺมสฺส สมฺปยุตฺต\-ปจฺจเยน ปจฺจโย. (ปฏิจฺจวาเร สมฺปยุตฺต\-สทิสา ฉ ปญฺหา.)}

\csromanheader{2}{2. สเหตุทุก \textbf{วิปฺปยุตฺต}}
\proseitem{115}{\csromanfolio{49}สเหตุโก ธมฺโม อเหตุกสฺส ธมฺมสฺส วิปฺปยุตฺต\-ปจฺจเยน ปจฺจโย. สหชาตํ, ปจฺฉาชาตํ.  สหชาตา สเหตุกา ขนฺธา จิตฺต\-สมุฏฺฐานานํ รูปานํ วิปฺปยุตฺต\-ปจฺจเยน ปจฺจโย. ปฏิสนฺธิ\-กฺขเณ สเหตุกา ขนฺธา กฏตฺตารูปานํ วิปฺปยุตฺต\-ปจฺจเยน ปจฺจโย. ปจฺฉาชาตา สเหตุกา ขนฺธา ปุเรชาตสฺส อิมสฺส กายสฺส วิปฺปยุตฺต\-ปจฺจเยน ปจฺจโย. (1)}

\prose{อเหตุโก ธมฺโม อเหตุกสฺส ธมฺมสฺส วิปฺปยุตฺต\-ปจฺจเยน ปจฺจโย. สหชาตํ, ปุเรชาตํ, ปจฺฉาชาตํ.  สหชาตา อเหตุกา ขนฺธา จิตฺต\-สมุฏฺฐานานํ รูปานํ วิปฺปยุตฺต\-ปจฺจเยน ปจฺจโย. ปฏิสนฺธิ\-กฺขเณ อเหตุกา ขนฺธา กฏตฺตารูปานํ วิปฺปยุตฺต\-ปจฺจเยน ปจฺจโย. ขนฺธา วตฺถุสฺส วิปฺปยุตฺต\-ปจฺจเยน ปจฺจโย. วตฺถุ ขนฺธานํ วิปฺปยุตฺต\-ปจฺจเยน ปจฺจโย. ปุเรชาตํ จกฺขายตนํ จกฺขุ\-วิญฺญาณสฺส ฯเปฯ. กายายตนํ กายวิญฺญาณสฺส. วตฺถุ อเหตุกานํ ขนฺธานํ โมหสฺส จ วิปฺปยุตฺต\-ปจฺจเยน ปจฺจโย. ปจฺฉาชาตา อเหตุกา ขนฺธา จ โมโห จ ปุเรชาตสฺส อิมสฺส กายสฺส วิปฺปยุตฺต\-ปจฺจเยน ปจฺจโย. (1)}

\prose{อเหตุโก ธมฺโม สเหตุกสฺส ธมฺมสฺส วิปฺปยุตฺต\-ปจฺจเยน ปจฺจโย. สหชาตํ, ปุเรชาตํ.  สหชาตํ ปฏิสนฺธิ\-กฺขเณ วตฺถุ สเหตุกานํ ขนฺธานํ วิปฺปยุตฺต\-ปจฺจเยน ปจฺจโย. ปุเรชาตํ วตฺถุ สเหตุกานํ ขนฺธานํ วิปฺปยุตฺต\-ปจฺจเยน ปจฺจโย. (2)}

\prose{อเหตุโก ธมฺโม สเหตุกสฺส จ อเหตุกสฺส จ ธมฺมสฺส วิปฺปยุตฺต\-ปจฺจเยน ปจฺจโย. ปุเรชาตํ วตฺถุ วิจิกิจฺฉา\-สหคตานํ อุทฺธจฺจ\-สหคตานํ ขนฺธานํ โมหสฺส จ วิปฺปยุตฺต\-ปจฺจเยน ปจฺจโย. (3)}

\prose{สเหตุโก จ อเหตุโก จ ธมฺมา อเหตุกสฺส ธมฺมสฺส วิปฺปยุตฺต\-ปจฺจเยน ปจฺจโย. สหชาตํ, ปจฺฉาชาตํ.  สหชาตา วิจิกิจฺฉา\-สหคตา อุทฺธจฺจ\-สหคตา ขนฺธา จ โมโห จ จิตฺต\-สมุฏฺฐานานํ รูปานํ วิปฺปยุตฺต\-ปจฺจเยน ปจฺจโย. ปจฺฉาชาตา วิจิกิจฺฉา\-สหคตา อุทฺธจฺจ\-สหคตา ขนฺธา จ โมโห จ ปุเรชาตสฺส อิมสฺส กายสฺส วิปฺปยุตฺต\-ปจฺจเยน ปจฺจโย. (1)}

\proseitem{116}{\csromanfolio{50}สเหตุโก ธมฺโม สเหตุกสฺส ธมฺมสฺส อตฺถิปจฺจเยน ปจฺจโย. สเหตุโก เอโก ขนฺโธ ติณฺณนฺนํ ขนฺธานํ ฯเปฯ เทฺว ขนฺธา ฯเปฯ. ปฏิสนฺธิ\-กฺขเณ ฯเปฯ. (1)}

\prose{สเหตุโก ธมฺโม อเหตุกสฺส ธมฺมสฺส อตฺถิปจฺจเยน ปจฺจโย. สหชาตํ, ปจฺฉาชาตํ.  สหชาตา สเหตุกา ขนฺธา จิตฺต\-สมุฏฺฐานานํ รูปานํ อตฺถิปจฺจเยน ปจฺจโย. วิจิกิจฺฉา\-สหคตา อุทฺธจฺจ\-สหคตา ขนฺธา โมหสฺส จ จิตฺตสมุฏฺฐานานญฺจ รูปานํ อตฺถิปจฺจเยน ปจฺจโย. ปฏิสนฺธิ\-กฺขเณ ฯเปฯ. ปจฺฉาชาตา สเหตุกา ขนฺธา ปุเรชาตสฺส อิมสฺส กายสฺส อตฺถิปจฺจเยน ปจฺจโย. (2)}

\prose{สเหตุโก ธมฺโม สเหตุกสฺส จ อเหตุกสฺส จ ธมฺมสฺส อตฺถิปจฺจเยน ปจฺจโย. สเหตุโก เอโก ขนฺโธ ติณฺณนฺนํ ขนฺธานํ จิตฺตสมุฏฺฐานานญฺจ รูปานํ อตฺถิปจฺจเยน ปจฺจโย. วิจิกิจฺฉา\-สหคโต อุทฺธจฺจ\-สหคโต เอโก ขนฺโธ ติณฺณนฺนํ ขนฺธานํ โมหสฺส จ จิตฺตสมุฏฺฐานานญฺจ รูปานํ อตฺถิปจฺจเยน ปจฺจโย ฯเปฯ เทฺว ขนฺธา ฯเปฯ. ปฏิสนฺธิ\-กฺขเณ สเหตุโก ฯเปฯ. (3)}

\proseitem{117}{อเหตุโก ธมฺโม อเหตุกสฺส ธมฺมสฺส อตฺถิปจฺจเยน ปจฺจโย. สหชาตํ, ปุเรชาตํ, ปจฺฉาชาตํ, อาหารํ, อินฺทฺริยํ.  สหชาโต อเหตุโก เอโก ขนฺโธ ติณฺณนฺนํ ขนฺธานํ จิตฺตสมุฏฺฐานานญฺจ รูปานํ อตฺถิปจฺจเยน ปจฺจโย ฯเปฯ เทฺว ขนฺธา ฯเปฯ. วิจิกิจฺฉา\-สหคโต อุทฺธจฺจ\-สหคโต โมโห จิตฺต\-สมุฏฺฐานานํ รูปานํ อตฺถิปจฺจเยน ปจฺจโย. ปฏิสนฺธิ\-กฺขเณ ฯเปฯ. (ยาว อสญฺญสตฺตา กาตพฺพํ.) ปุเรชาตํ จกฺขุํ ฯเปฯ วตฺถุํ อนิจฺจโต ฯเปฯ โทมนสฺสํ อุปฺปชฺชติ. กุสลากุสเล นิรุทฺเธ อเหตุโก วิปาโก ตทารมฺมณตา อุปฺปชฺชติ. รูปายตนํ จกฺขุ\-วิญฺญาณสฺส ฯเปฯ. โผฏฺฐพฺพา\-ยตนํ กายวิญฺญาณสฺส ฯเปฯ. จกฺขายตนํ จกฺขุ\-วิญฺญาณสฺส ฯเปฯ. กายายตนํ กายวิญฺญาณสฺส อตฺถิปจฺจเยน ปจฺจโย. วตฺถุ อเหตุกานํ ขนฺธานํ โมหสฺส จ อตฺถิปจฺจเยน ปจฺจโย. ปจฺฉาชาตา อเหตุกา ขนฺธา จ โมโห จ ปุเรชาตสฺส อิมสฺส กายสฺส อตฺถิปจฺจเยน ปจฺจโย. กพฬีกาโร \csromanfolio{51}อาหาโร อิมสฺส กายสฺส อตฺถิปจฺจเยน ปจฺจโย. รูปชีวิตินฺทฺริยํ กฏตฺตารูปานํ อตฺถิปจฺจเยน ปจฺจโย. (1)}

\prose{อเหตุโก ธมฺโม สเหตุกสฺส ธมฺมสฺส อตฺถิปจฺจเยน ปจฺจโย. สหชาตํ, ปุเรชาตํ.  สหชาโต วิจิกิจฺฉา\-สหคโต อุทฺธจฺจ\-สหคโต โมโห สมฺปยุตฺตกานํ ขนฺธานํ อตฺถิปจฺจเยน ปจฺจโย. ปฏิสนฺธิ\-กฺขเณ วตฺถุ สเหตุกานํ ขนฺธานํ อตฺถิปจฺจเยน ปจฺจโย. ปุเรชาตํ จกฺขุํ ฯเปฯ วตฺถุํ อนิจฺจโต ฯเปฯ โทมนสฺสํ อุปฺปชฺชติ. กุสลากุสเล นิรุทฺเธ สเหตุโก วิปาโก ตทารมฺมณตา อุปฺปชฺชติ. ทิพฺเพน จกฺขุนา รูปํ ปสฺสติ. ทิพฺพาย โสตธาตุยา สทฺทํ สุณาติ. วตฺถุ สเหตุกานํ ขนฺธานํ อตฺถิปจฺจเยน ปจฺจโย. (2)}

\prose{อเหตุโก ธมฺโม สเหตุกสฺส จ อเหตุกสฺส จ ธมฺมสฺส อตฺถิปจฺจเยน ปจฺจโย. สหชาตํ, ปุเรชาตํ.  สหชาโต วิจิกิจฺฉา\-สหคโต อุทฺธจฺจ\-สหคโต โมโห สมฺปยุตฺตกานํ ขนฺธานํ จิตฺตสมุฏฺฐานานญฺจ รูปานํ อตฺถิปจฺจเยน ปจฺจโย. ปุเรชาตํ จกฺขุํ ฯเปฯ วตฺถุํ อารพฺภ วิจิกิจฺฉา\-สหคตา อุทฺธจฺจ\-สหคตา ขนฺธา จ โมโห จ อุปฺปชฺชนฺติ. วตฺถุ วิจิกิจฺฉา\-สหคตานํ อุทฺธจฺจ\-สหคตานํ ขนฺธานํ โมหสฺส จ อตฺถิปจฺจเยน ปจฺจโย. (3)}

\proseitem{118}{สเหตุโก จ อเหตุโก จ ธมฺมา สเหตุกสฺส ธมฺมสฺส อตฺถิปจฺจเยน ปจฺจโย. สหชาตํ, ปุเรชาตํ.  สหชาโต วิจิกิจฺฉา\-สหคโต อุทฺธจฺจ\-สหคโต เอโก ขนฺโธ จ โมโห จ ติณฺณนฺนํ ขนฺธานํ อตฺถิปจฺจเยน ปจฺจโย ฯเปฯ เทฺว ขนฺธา จ ฯเปฯ. ปฏิสนฺธิ\-กฺขเณ สเหตุโก เอโก ขนฺโธ จ วตฺถุ จ ติณฺณนฺนํ ขนฺธานํ อตฺถิปจฺจเยน ปจฺจโย ฯเปฯ เทฺว ขนฺธา จ ฯเปฯ. สหชาโต สเหตุโก เอโก ขนฺโธ จ วตฺถุ จ ติณฺณนฺนํ ขนฺธานํ อตฺถิปจฺจเยน ปจฺจโย ฯเปฯ เทฺว ขนฺธา จ ฯเปฯ. (1)}

\prose{สเหตุโก จ อเหตุโก จ ธมฺมา อเหตุกสฺส ธมฺมสฺส อตฺถิปจฺจเยน ปจฺจโย. สหชาตํ, ปุเรชาตํ, ปจฺฉาชาตํ, อาหารํ, อินฺทฺริยํ.  สหชาตา สเหตุกา ขนฺธา จ มหาภูตา จ จิตฺต\-สมุฏฺฐานานํ รูปานํ อตฺถิปจฺจเยน ปจฺจโย. วิจิกิจฺฉา\-สหคตา อุทฺธจฺจ\-สหคตา ขนฺธา จ โมโห จ จิตฺต\-สมุฏฺฐานานํ รูปานํ อตฺถิปจฺจเยน ปจฺจโย. ปฏิสนฺธิ\-กฺขเณ สเหตุกา ขนฺธา จ มหาภูตา จ กฏตฺตารูปานํ \csromanfolio{52}อตฺถิปจฺจเยน ปจฺจโย. สหชาตา วิจิกิจฺฉา\-สหคตา อุทฺธจฺจ\-สหคตา ขนฺธา จ วตฺถุ จ โมหสฺส อตฺถิปจฺจเยน ปจฺจโย. ปจฺฉาชาตา วิจิกิจฺฉา\-สหคตา อุทฺธจฺจ\-สหคตา ขนฺธา จ โมโห จ ปุเรชาตสฺส อิมสฺส กายสฺส อตฺถิปจฺจเยน ปจฺจโย. ปจฺฉาชาตา สเหตุกา ขนฺธา จ กพฬีกาโร อาหาโร จ อิมสฺส กายสฺส อตฺถิปจฺจเยน ปจฺจโย. ปจฺฉาชาตา สเหตุกา ขนฺธา จ รูปชีวิตินฺทฺริยญฺจ กฏตฺตา รูปานํ อตฺถิปจฺจเยน ปจฺจโย. (2)}

\prose{สเหตุโก จ อเหตุโก จ ธมฺมา สเหตุกสฺส จ อเหตุกสฺส จ ธมฺมสฺส อตฺถิปจฺจเยน ปจฺจโย. สหชาตํ, ปุเรชาตํ.  สหชาโต วิจิกิจฺฉา\-สหคโต อุทฺธจฺจ\-สหคโต เอโก ขนฺโธ จ โมโห จ ติณฺณนฺนํ ขนฺธานํ จิตฺตสมุฏฺฐานานญฺจ รูปานํ อตฺถิปจฺจเยน ปจฺจโย ฯเปฯ. สหชาโต วิจิกิจฺฉา\-สหคโต อุทฺธจฺจ\-สหคโต เอโก ขนฺโธ จ วตฺถุ จ ติณฺณนฺนํ ขนฺธานํ โมหสฺส จ อตฺถิปจฺจเยน ปจฺจโย ฯเปฯ เทฺว ขนฺธา จ ฯเปฯ. (3)}

\csromanmarkh{1. ปจฺจยานุโลม}\csromanmarkhii{2. สงฺขฺยาวาร}\csromanheadernomark{2}{1. \textbf{ปจฺจยานุโลม 2. สงฺขฺยาวาร}}
\proseitem{119}{เหตุยา ฉ, อารมฺมเณ นว, อธิปติยา จตฺตาริ, อนนฺตเร นว, สมนนฺตเร นว, สหชาเต นว, อญฺญมญฺเญ ฉ, นิสฺสเย นว, อุปนิสฺสเย นว, ปุเรชาเต ตีณิ, ปจฺฉาชาเต ตีณิ อาเสวเน นว, กมฺเม จตฺตาริ, วิปาเก จตฺตาริ, อาหาเร จตฺตาริ, อินฺทฺริเย จตฺตาริ, ฌาเน จตฺตาริ, มคฺเค ตีณิ, สมฺปยุตฺเต ฉ, วิปฺปยุตฺเต ปญฺจ, อตฺถิยา นว, นตฺถิยา นว, วิคเต นว, อวิคเต นว. (เอวํ คเณตพฺพํ.)}

\vspace{\tipitakaparskip}
\csromancenter{อนุโลมํ.}
\proseitem{120}{สเหตุโก ธมฺโม สเหตุกสฺส ธมฺมสฺส อารมฺมณ\-ปจฺจเยน ปจฺจโย. สหชาต\-ปจฺจเยน ปจฺจโย. อุปนิสฺสย\-ปจฺจเยน ปจฺจโย. กมฺมปจฺจเยน ปจฺจโย. (1)}

\prose{\csromanfolio{53}สเหตุโก ธมฺโม อเหตุกสฺส ธมฺมสฺส อารมฺมณ\-ปจฺจเยน ปจฺจโย. สหชาต\-ปจฺจเยน ปจฺจโย. อุปนิสฺสย\-ปจฺจเยน ปจฺจโย. ปจฺฉาชาต\-ปจฺจเยน ปจฺจโย. กมฺมปจฺจเยน ปจฺจโย. (2)}

\prose{สเหตุโก ธมฺโม สเหตุกสฺส จ อเหตุกสฺส จ ธมฺมสฺส อารมฺมณ\-ปจฺจเยน ปจฺจโย. สหชาต\-ปจฺจเยน ปจฺจโย. อุปนิสฺสย\-ปจฺจเยน ปจฺจโย. กมฺมปจฺจเยน ปจฺจโย. (3)}

\proseitem{121}{อเหตุโก ธมฺโม อเหตุกสฺส ธมฺมสฺส อารมฺมณ\-ปจฺจเยน ปจฺจโย. สหชาต\-ปจฺจเยน ปจฺจโย. อุปนิสฺสย\-ปจฺจเยน ปจฺจโย. ปุเรชาต\-ปจฺจเยน ปจฺจโย. ปจฺฉาชาต\-ปจฺจเยน ปจฺจโย. อาหารปจฺจเยน ปจฺจโย. อินฺทฺริย\-ปจฺจเยน ปจฺจโย. (1)}

\prose{อเหตุโก ธมฺโม สเหตุกสฺส ธมฺมสฺส อารมฺมณ\-ปจฺจเยน ปจฺจโย. สหชาต\-ปจฺจเยน ปจฺจโย. อุปนิสฺสย\-ปจฺจเยน ปจฺจโย. ปุเรชาต\-ปจฺจเยน ปจฺจโย. (2)}

\prose{อเหตุโก ธมฺโม สเหตุกสฺส จ อเหตุกสฺส จ ธมฺมสฺส อารมฺมณ\-ปจฺจเยน ปจฺจโย. สหชาต\-ปจฺจเยน ปจฺจโย. อุปนิสฺสย\-ปจฺจเยน ปจฺจโย. ปุเรชาต\-ปจฺจเยน ปจฺจโย. (3)}

\proseitem{122}{สเหตุโก จ อเหตุโก จ ธมฺมา สเหตุกสฺส ธมฺมสฺส อารมฺมณ\-ปจฺจเยน ปจฺจโย. สหชาต\-ปจฺจเยน ปจฺจโย. อุปนิสฺสย\-ปจฺจเยน ปจฺจโย. (1)}

\prose{สเหตุโก จ อเหตุโก จ ธมฺมา อเหตุกสฺส ธมฺมสฺส อารมฺมณ\-ปจฺจเยน ปจฺจโย. สหชาต\-ปจฺจเยน ปจฺจโย. อุปนิสฺสย\-ปจฺจเยน ปจฺจโย. ปจฺฉาชาต\-ปจฺจเยน ปจฺจโย. (2)}

\prose{สเหตุโก จ อเหตุโก จ ธมฺมา สเหตุกสฺส จ อเหตุกสฺส จ ธมฺมสฺส อารมฺมณ\-ปจฺจเยน ปจฺจโย. สหชาต\-ปจฺจเยน ปจฺจโย. อุปนิสฺสย\-ปจฺจเยน ปจฺจโย. (3)}

\csromanmarkh{2. ปจฺจย\-ปจฺจนีย}\csromanmarkhii{2. สงฺขฺยาวาร}\csromanheadernomark{2}{2. ปจฺจย\-ปจฺจนีย 2. สงฺขฺยาวาร}
\proseitem{123}{\csromanfolio{54}นเหตุยา นว ฯเปฯ (สพฺพตฺถ นว) โนอวิคเต นว. (เอวํ คเณตพฺพํ.)}

\vspace{\tipitakaparskip}
\csromancenter{ปจฺจนียํ.}
\proseitem{124}{เหตุปจฺจยา นอารมฺมเณ ฉ, นอธิปติยา ฉ, นอนนฺตเร ฉ, นสมนนฺตเร ฉ, นอญฺญมญฺเญ เทฺว, นอุปนิสฺสเย ฉ ฯเปฯ นมคฺเค ฉ, นสมฺปยุตฺเต เทฺว, นวิปฺปยุตฺเต เทฺว, โนนตฺถิยา ฉ, โนวิคเต ฉ. (เอวํ คเณตพฺพํ.)}

\vspace{\tipitakaparskip}
\csromancenter{อนุโลม\-ปจฺจนียํ.}
\proseitem{125}{นเหตุปจฺจยา อารมฺมเณ นว, อธิปติยา จตฺตาริ, อนนฺตเร นว, สมนนฺตเร นว, สหชาเต นว, อญฺญมญฺเญ ฉ, นิสฺสเย นว, อุปนิสฺสเย นว, ปุเรชาเต ตีณิ, ปจฺฉาชาเต ตีณิ, อาเสวเน นว, กมฺเม จตฺตาริ, วิปาเก จตฺตาริ, อาหาเร จตฺตาริ, อินฺทฺริเย จตฺตาริ, ฌาเน จตฺตาริ, มคฺเค ตีณิ, สมฺปยุตฺเต ฉ, วิปฺปยุตฺเต ปญฺจ อตฺถิยา นว, นตฺถิยา นว, วิคเต นว, อวิคเต นว. (เอวํ คเณตพฺพํ.)}

\vspace{\tipitakaparskip}
\csromancenter{ปจฺจนียา\-นุโลมํ.}
\nitthitam{สเหตุกทุกํ นิฏฺฐิตํ.}
\csromanmatikaanchor{mtk.5}
\csromantocmarkref{section}{3. เหตุสมฺปยุตฺตทุก}{mtk.5}
\bookmark[dest={mtk.5},level=1]{3. เหตุสมฺปยุตฺตทุก}
\csromanmarkh{4. เหตุสเหตุทุก}\csromanmarkhii{3. เหตุ\-สมฺปยุตฺตทุก 1. ปฏิจฺจวาร}\csromanheadernomark{1}{4. เหตุสเหตุทุก 3. เหตุ\-สมฺปยุตฺตทุก 1. ปฏิจฺจวาร}
\csromanheader{2}{\textbf{เหตุ}}
\proseitem{126}{\csromanfolio{55}เหตุ\-สมฺปยุตฺตํ ธมฺมํ ปฏิจฺจ เหตุสมฺปยุตฺโต ธมฺโม อุปฺปชฺชติ เหตุปจฺจยา. เหตุ\-สมฺปยุตฺตํ เอกํ ขนฺธํ ปฏิจฺจ ตโย ขนฺธา ฯเปฯ เทฺว ขนฺเธ ฯเปฯ. ปฏิสนฺธิ\-กฺขเณ ฯเปฯ. (1)}

\prose{เหตุ\-สมฺปยุตฺตํ ธมฺมํ ปฏิจฺจ เหตุวิปฺปยุตฺโต ธมฺโม อุปฺปชฺชติ เหตุปจฺจยา. เหตุสมฺปยุตฺเต ขนฺเธ ปฏิจฺจ จิตฺต\-สมุฏฺฐานํ รูปํ. ปฏิสนฺธิ\-กฺขเณ ฯเปฯ. (2)}

\prosewithcloser{(อิมินา การเณน วิตฺถาเรตพฺพํ. ยถา สเหตุกทุกํ. นินฺนานากรณํ.)}{เหตุ\-สมฺปยุตฺตทุกํ นิฏฺฐิตํ.}
\csromanmatikaanchor{mtk.6}
\csromantocmarkref{section}{4. เหตุสเหตุกทุก}{mtk.6}
\bookmark[dest={mtk.6},level=1]{4. เหตุสเหตุกทุก}
\csromanmarkh{4. เหตุ\-สเหตุกทุก}\csromanmarkhii{1. ปฏิจฺจวาร}\csromanheadernomark{1}{4. เหตุ\-สเหตุกทุก 1. ปฏิจฺจวาร}
\csromanmarkh{1. ปจฺจยานุโลม}\csromanmarkhii{1. วิภงฺควาร}\csromanheadernomark{2}{1. \textbf{ปจฺจยานุโลม 1. วิภงฺควาร}}
\csromanheader{2}{\textbf{เหตุ}}
\proseitem{127}{เหตุญฺเจว สเหตุกญฺจ ธมฺมํ ปฏิจฺจ เหตุ เจว สเหตุโก จ ธมฺโม อุปฺปชฺชติ เหตุปจฺจยา. อโลภํ ปฏิจฺจ อโทโส อโมโห. (จกฺกํ.) โลภํ ปฏิจฺจ โมโห. (จกฺกํ.) ปฏิสนฺธิ\-กฺขเณ อโลภํ ปฏิจฺจ อโทโส อโมโห. (จกฺกํ.) (1)}

\prose{เหตุญฺเจว สเหตุกญฺจ ธมฺมํ ปฏิจฺจ สเหตุโก เจว น จ เหตุ ธมฺโม อุปฺปชฺชติ เหตุปจฺจยา. เหตุํ ปฏิจฺจ สมฺปยุตฺตกา ขนฺธา ปฏิสนฺธิ\-กฺขเณ ฯเปฯ. (2)}

\prose{เหตุญฺเจว สเหตุกญฺจ ธมฺมํ ปฏิจฺจ เหตุ เจว สเหตุโก จ สเหตุโก เจว น จ เหตุ จ ธมฺมา อุปฺปชฺชนฺติ เหตุปจฺจยา. อโลภํ ปฏิจฺจ อโทโส อโมโห สมฺปยุตฺตกา จ ขนฺธา. (จกฺกํ.) โลภํ ปฏิจฺจ โมโห สมฺปยุตฺตกา จ ขนฺธา. (จกฺกํ.) ปฏิสนฺธิ\-กฺขเณ ฯเปฯ. (3)}

\proseitem{128}{\csromanfolio{56}สเหตุกญฺเจว น จ เหตุํ ธมฺมํ ปฏิจฺจ สเหตุโก เจว น จ เหตุ ธมฺโม อุปฺปชฺชติ เหตุปจฺจยา. สเหตุกญฺเจว น จ เหตุํ เอกํ ขนฺธํ ปฏิจฺจ ตโย ขนฺธา ฯเปฯ เทฺว ขนฺเธ ฯเปฯ. ปฏิสนฺธิ\-กฺขเณ ฯเปฯ. (1)}

\prose{สเหตุกญฺเจว น จ เหตุํ ธมฺมํ ปฏิจฺจ เหตุ เจว สเหตุโก จ ธมฺโม อุปฺปชฺชติ เหตุปจฺจยา. สเหตุเก เจว น จ เหตู ขนฺเธ ปฏิจฺจ เหตู. ปฏิสนฺธิ\-กฺขเณ ฯเปฯ. (2)}

\prose{สเหตุกญฺเจว น จ เหตุํ ธมฺมํ ปฏิจฺจ เหตุ เจว สเหตุโก จ สเหตุโก เจว น จ เหตุ จ ธมฺมา อุปฺปชฺชนฺติ เหตุปจฺจยา. สเหตุกญฺเจว น จ เหตุํ เอกํ ขนฺธํ ปฏิจฺจ ตโย ขนฺธา เหตุ จ ฯเปฯ เทฺว ขนฺเธ ฯเปฯ. ปฏิสนฺธิ\-กฺขเณ ฯเปฯ. (3)}

\proseitem{129}{เหตุญฺเจว สเหตุกญฺจ สเหตุกญฺเจว น จ เหตุญฺจ ธมฺมํ ปฏิจฺจ เหตุ เจว สเหตุโก จ ธมฺโม อุปฺปชฺชติ เหตุปจฺจยา. อโลภญฺจ สมฺปยุตฺตเก จ ขนฺเธ ปฏิจฺจ อโทโส อโมโห. (จกฺกํ.) โลภญฺจ สมฺปยุตฺตเก จ ขนฺเธ ปฏิจฺจ โมโห. (จกฺกํ.) ปฏิสนฺธิ\-กฺขเณ ฯเปฯ. (1)}

\prose{เหตุญฺเจว สเหตุกญฺจ สเหตุกญฺเจว น จ เหตุญฺจ ธมฺมํ ปฏิจฺจ สเหตุโก เจว น จ เหตุ ธมฺโม อุปฺปชฺชติ เหตุปจฺจยา. สเหตุกญฺเจว น จ เหตุํ เอกํ ขนฺธญฺจ เหตุญฺจ ปฏิจฺจ ตโย ขนฺธา ฯเปฯ เทฺว ขนฺเธ ฯเปฯ. ปฏิสนฺธิ\-กฺขเณ ฯเปฯ. (2)}

\prose{เหตุญฺเจว สเหตุกญฺจ สเหตุกญฺเจว น จ เหตุญฺจ ธมฺมํ ปฏิจฺจ เหตุ เจว สเหตุโก จ สเหตุโก เจว น จ เหตุ จ ธมฺมา อุปฺปชฺชนฺติ เหตุปจฺจยา. สเหตุกญฺเจว น จ เหตุํ เอกํ ขนฺธญฺจ อโลภญฺจ ปฏิจฺจ ตโย ขนฺธา อโทโส อโมโห จ ฯเปฯ เทฺว ขนฺเธ ฯเปฯ. ปฏิสนฺธิ\-กฺขเณ ฯเปฯ. (3) (สํขิตฺตํ. เอวํ วิตฺถาเรตพฺพํ.)}

\csromanmarkh{1. ปจฺจยานุโลม}\csromanmarkhii{2. สงฺขฺยาวาร}\csromanheadernomark{2}{1. \textbf{ปจฺจยานุโลม 2. สงฺขฺยาวาร}}
\proseitem{130}{เหตุยา นว, อารมฺมเณ นว, อธิปติยา นว, อนนฺตเร นว ฯเปฯ (สพฺพตฺถ นว) อวิคเต นว. (เอวํ คเณตพฺพํ.)}

\vspace{\tipitakaparskip}
\csromancenter{อนุโลมํ.}
\csromanmarkh{4. เหตุสเหตุทุก}\csromanmarkhii{2. ปจฺจย\-ปจฺจนีย 1. วิภงฺควาร}\csromanheadernomark{2}{4. เหตุสเหตุทุก 2. ปจฺจย\-ปจฺจนีย 1. วิภงฺควาร}
\prose{\csromanfolio{57}นอธิปตฺยาทิ}

\proseitem{131}{เหตุญฺเจว สเหตุกญฺจ ธมฺมํ ปฏิจฺจ เหตุ เจว สเหตุโก จ ธมฺโม อุปฺปชฺชติ นอธิปติ\-ปจฺจยา. อโลภํ ปฏิจฺจ อโทโส อโมโห. (จกฺกํ.) ปฏิสนฺธิ\-กฺขเณ ฯเปฯ. (ปริปุณฺณํ นว.) นปุเรชาเต นว, นปจฺฉาชาเต นว, นอาเสวเน นว.}

\csromanheader{2}{\textbf{นกมฺมาทิ}}
\proseitem{132}{เหตุญฺเจว สเหตุกญฺจ ธมฺมํ ปฏิจฺจ สเหตุโก เจว น จ เหตุ ธมฺโม อุปฺปชฺชติ นกมฺมปจฺจยา. เหตุํ ปฏิจฺจ สมฺปยุตฺตกา เจตนา. (1)}

\prose{สเหตุกญฺเจว น จ เหตุํ ธมฺมํ ปฏิจฺจ สเหตุโก เจว น จ เหตุ ธมฺโม อุปฺปชฺชติ นกมฺมปจฺจยา. สเหตุเก เจว น จ เหตู ขนฺเธ ปฏิจฺจ สมฺปยุตฺตกา เจตนา. ปฏิสนฺธิ\-กฺขเณ ฯเปฯ. (1)}

\prose{เหตุญฺเจว สเหตุกญฺจ สเหตุกญฺเจว น จ เหตุญฺจ ธมฺมํ ปฏิจฺจ สเหตุโก เจว น จ เหตุ ธมฺโม อุปฺปชฺชติ นกมฺมปจฺจยา. เหตุญฺจ สมฺปยุตฺตเก จ ขนฺเธ ปฏิจฺจ สมฺปยุตฺตกา เจตนา. นวิปาก\-ปจฺจยา. นวิปฺปยุตฺต\-ปจฺจยา.}

\csromanmarkh{2. ปจฺจย\-ปจฺจนีย}\csromanmarkhii{2. สงฺขฺยาวาร}\csromanheadernomark{2}{2. ปจฺจย\-ปจฺจนีย 2. สงฺขฺยาวาร}
\proseitem{133}{นอธิปติยา นว, นปุเรชาเต นว, นปจฺฉาชาเต นว, นอาเสวเน นว, นกมฺเม ตีณิ, นวิปาเก นว, นวิปฺปยุตฺเต นว. (เอวํ คเณตพฺพํ.)}

\vspace{\tipitakaparskip}
\csromancenter{ปจฺจนียํ.}
\csromancenter{\textbf{เหตุ}ปจฺจยา นอธิปติยา นว, นปุเรชาเต นว, นปจฺฉาชาเต นว, นอาเสวเน นว, นกมฺเม ตีณิ, นวิปาเก นว, นวิปฺปยุตฺเต นว. (เอวํ คเณตพฺพํ.)}
\csromancenter{อนุโลม\-ปจฺจนียํ.}
\prose{\csromanfolio{58}นอธิปติทุก}

\proseitem{135}{นอธิปติ\-ปจฺจยา เหตุยา นว, อารมฺมเณ นว, อนนฺตเร นว ฯเปฯ อวิคเต นว. (เอวํ คเณตพฺพํ.) ปจฺจนียา\-นุโลมํ.}

\prose{(สหชาตวาโรปิ ปจฺจยวาโรปิ นิสฺสยวาโรปิ สํสฏฺฐ\-วาโรปิ สมฺปยุตฺตวาโรปิ ปฏิจฺจ\-วาร\-สทิสา.)}

\csromanmarkh{4. เหตุ\-สเหตุกทุก}\csromanmarkhii{7. ปญฺหาวาร}\csromanheadernomark{2}{4. เหตุ\-สเหตุกทุก 7. ปญฺหาวาร}
\csromanmarkh{1. ปจฺจยานุโลม}\csromanmarkhii{1. วิภงฺควาร}\csromanheadernomark{2}{1. \textbf{ปจฺจยานุโลม 1. วิภงฺควาร}}
\csromanheader{2}{\textbf{เหตุ}}
\proseitem{136}{เหตุ เจว สเหตุโก จ ธมฺโม เหตุสฺส เจว สเหตุกสฺส จ ธมฺมสฺส เหตุปจฺจเยน ปจฺจโย. อโลโภ อโทสสฺส อโมหสฺส เหตุปจฺจเยน ปจฺจโย. (ยถา ปฏิจฺจ\-วาร\-สทิสํ.) (1)}

\prose{เหตุ เจว สเหตุโก จ ธมฺโม สเหตุกสฺส เจว น จ เหตุสฺส ธมฺมสฺส เหตุปจฺจเยน ปจฺจโย. เหตู สมฺปยุตฺตกานํ ขนฺธานํ เหตุปจฺจเยน ปจฺจโย. ปฏิสนฺธิ\-กฺขเณ ฯเปฯ. (2)}

\prose{\csromanfolio{59}เหตุ เจว สเหตุโก จ ธมฺโม เหตุสฺส เจว สเหตุกสฺส จ สเหตุกสฺส เจว น จ เหตุสฺส จ ธมฺมสฺส เหตุปจฺจเยน ปจฺจโย. อโลโภ อโทสสฺส อโมหสฺส สมฺปยุตฺตกานญฺจ ขนฺธานํ เหตุปจฺจเยน ปจฺจโย. (วิตฺถาเรตพฺพํ.) (3)}

\proseitem{137}{เหตุ เจว สเหตุโก จ ธมฺโม เหตุสฺส เจว สเหตุกสฺส จ ธมฺมสฺส อารมฺมณ\-ปจฺจเยน ปจฺจโย. เหตุํ อารพฺภ เหตู อุปฺปชฺชนฺติ. (1)}

\prose{เหตุ เจว สเหตุโก จ ธมฺโม สเหตุกสฺส เจว น จ เหตุสฺส ธมฺมสฺส อารมฺมณ\-ปจฺจเยน ปจฺจโย. เหตุํ อารพฺภ สเหตุกา เจว น จ เหตู ขนฺธา อุปฺปชฺชนฺติ. (2)}

\prose{เหตุ เจว สเหตุโก จ ธมฺโม เหตุสฺส เจว สเหตุกสฺส จ สเหตุกสฺส เจว น จ เหตุสฺส จ ธมฺมสฺส อารมฺมณ\-ปจฺจเยน ปจฺจโย. เหตุํ อารพฺภ เหตู จ สมฺปยุตฺตกา จ ขนฺธา อุปฺปชฺชนฺติ. (3)}

\prose{สเหตุโก เจว น จ เหตุ ธมฺโม สเหตุกสฺส เจว น จ เหตุสฺส ธมฺมสฺส อารมฺมณ\-ปจฺจเยน ปจฺจโย. ทานํ ทตฺวา, สีลํ ฯเปฯ อุโปสถกมฺมํ กตฺวา ตํ ปจฺจเวกฺขติ. ปุพฺเพ สุจิณฺณานิ ปจฺจเวกฺขติ. ฌานา วุฏฺฐหิตฺวา ฯเปฯ. อริยา มคฺคา วุฏฺฐหิตฺวา มคฺคํ ปจฺจเวกฺขนฺติ, ผลํ ปจฺจเวกฺขนฺติ. ปหีเน กิเลเส ฯเปฯ. วิกฺขมฺภิเต กิเลเส ปจฺจเวกฺขนฺติ. ปุพฺเพ สมุทาจิณฺเณ กิเลเส ชานนฺติ. สเหตุเก เจว น จ เหตู ขนฺเธ อนิจฺจโต ฯเปฯ โทมนสฺสํ อุปฺปชฺชติ. เจโตปริย\-ญาเณน สเหตุกา เจว น จ เหตุ\-จิตฺต\-สมงฺคิสฺส จิตฺตํ ชานาติ. อากาสา\-นญฺจา\-ยตนํ วิญฺญาณญฺจา\-ยตนสฺส ฯเปฯ. อากิญฺจญฺญา\-ยตนํ เนว\-สญฺญา\-นา\-สญฺญา\-ยตนสฺส ฯเปฯ. สเหตุกา เจว น จ เหตู ขนฺธา อิทฺธิวิธ\-ญาณสฺส, เจโตปริย\-ญาณสฺส, ปุพฺเพ\-นิวาสา\-นุสฺสติ\-ญาณสฺส, ยถากมฺมูปคญาณสฺส, อนาคตํสญาณสฺส อารมฺมณ\-ปจฺจเยน ปจฺจโย. (1)}

\prose{สเหตุโก เจว น จ เหตุ ธมฺโม เหตุสฺส เจว สเหตุกสฺส จ ธมฺมสฺส อารมฺมณ\-ปจฺจเยน ปจฺจโย. ทานํ ทตฺวา. (ยถา ปฐมคมนํ. เอวํ นินฺนานํ.) (2)}

\prose{\csromanfolio{60}สเหตุโก เจว น จ เหตุ ธมฺโม เหตุสฺส เจว สเหตุกสฺส จ สเหตุกสฺส เจว น จ เหตุสฺส จ ธมฺมสฺส อารมฺมณ\-ปจฺจเยน ปจฺจโย. ทานํ ทตฺวา. (ยถา ปฐมคมนํ. เอวํ นินฺนานํ.) (3)}

\proseitem{138}{เหตุ เจว สเหตุโก จ สเหตุโก เจว น จ เหตุ จ ธมฺมา เหตุสฺส เจว สเหตุกสฺส จ ธมฺมสฺส อารมฺมณ\-ปจฺจเยน ปจฺจโย. เหตุญฺจ สมฺปยุตฺตเก จ ขนฺเธ อารพฺภ เหตู อุปฺปชฺชนฺติ. (1)}

\prose{เหตุ เจว สเหตุโก จ สเหตุโก เจว น จ เหตุ จ ธมฺมา สเหตุกสฺส เจว น จ เหตุสฺส ธมฺมสฺส อารมฺมณ\-ปจฺจเยน ปจฺจโย. เหตุญฺจ สมฺปยุตฺตเก จ ขนฺเธ อารพฺภ สเหตุกา เจว น จ เหตู ขนฺธา อุปฺปชฺชนฺติ. (2)}

\prose{เหตุ เจว สเหตุโก จ สเหตุโก เจว น จ เหตุ จ ธมฺมา เหตุสฺส เจว สเหตุกสฺส จ สเหตุกสฺส เจว น จ เหตุสฺส จ ธมฺมสฺส อารมฺมณ\-ปจฺจเยน ปจฺจโย. เหตุญฺจ สมฺปยุตฺตเก จ ขนฺเธ อารพฺภ เหตู จ สมฺปยุตฺตกา จ ขนฺธา อุปฺปชฺชนฺติ. (3)}

\proseitem{139}{เหตุ เจว สเหตุโก จ ธมฺโม เหตุสฺส เจว สเหตุกสฺส จ ธมฺมสฺส อธิปติ\-ปจฺจเยน ปจฺจโย. อารมฺมณา\-ธิปติ, สหชาตา\-ธิปติ.  อารมฺมณา\-ธิปติ\csromandash{}เหตุํ ครุํ กตฺวา เหตู อุปฺปชฺชนฺติ. สหชาตา\-ธิปติ\csromandash{} เหตุ เจว สเหตุกา\-ธิปติ สมฺปยุตฺตกานํ เหตูนํ อธิปติ\-ปจฺจเยน ปจฺจโย. (1)}

\prose{เหตุ เจว สเหตุโก จ ธมฺโม สเหตุกสฺส เจว น จ เหตุสฺส ธมฺมสฺส อธิปติ\-ปจฺจเยน ปจฺจโย. อารมฺมณา\-ธิปติ, สหชาตา\-ธิปติ.  อารมฺมณา\-ธิปติ\csromandash{}เหตุํ ครุํ กตฺวา สเหตุกา เจว น จ เหตู ขนฺธา อุปฺปชฺชนฺติ. สหชาตา\-ธิปติ\csromandash{}เหตุ เจว สเหตุกา\-ธิปติ สมฺปยุตฺตกานํ ขนฺธานํ อธิปติ\-ปจฺจเยน ปจฺจโย. (2)}

\prose{เหตุ เจว สเหตุโก จ ธมฺโม เหตุสฺส เจว สเหตุกสฺส จ สเหตุกสฺส เจว น จ เหตุสฺส จ ธมฺมสฺส อธิปติ\-ปจฺจเยน ปจฺจโย. อารมฺมณา\-ธิปติ, สหชาตา\-ธิปติ.  อารมฺมณา\-ธิปติ\csromandash{}เหตุํ ครุํ \csromanfolio{61}กตฺวา เหตู จ สมฺปยุตฺตกา จ ขนฺธา อุปฺปชฺชนฺติ. สหชาตา\-ธิปติ\csromandash{}เหตุ เจว สเหตุกา\-ธิปติ สมฺปยุตฺตกานํ ขนฺธานํ เหตูนญฺจ อธิปติ\-ปจฺจเยน ปจฺจโย. (3)}

\proseitem{140}{สเหตุโก เจว น จ เหตุ ธมฺโม สเหตุกสฺส เจว น จ เหตุสฺส ธมฺมสฺส อธิปติ\-ปจฺจเยน ปจฺจโย. อารมฺมณา\-ธิปติ, สหชาตา\-ธิปติ.  อารมฺมณา\-ธิปติ\csromandash{}ทานํ ทตฺวา, สีลํ ฯเปฯ อุโปสถกมฺมํ กตฺวา ตํ ครุํ กตฺวา ปจฺจเวกฺขติ. ปุพฺเพ สุจิณฺณานิ ฯเปฯ. ฌานา วุฏฺฐหิตฺวา ฌานํ ครุํ กตฺวา ปจฺจเวกฺขติ. อริยา มคฺคา วุฏฺฐหิตฺวา มคฺคํ ครุํ กตฺวา ปจฺจเวกฺขนฺติ, ผลํ ครุํ กตฺวา ปจฺจเวกฺขนฺติ. สเหตุเก เจว น จ เหตู ขนฺเธ ครุํ กตฺวา อสฺสาเทติ อภินนฺทติ, ตํ ครุํ กตฺวา ราโค อุปฺปชฺชติ, ทิฏฺฐิ อุปฺปชฺชติ. สหชาตา\-ธิปติ\csromandash{}สเหตุโก เจว น จ เหตุ อธิปติ สมฺปยุตฺตกานํ ขนฺธานํ อธิปติ\-ปจฺจเยน ปจฺจโย. (1)}

\prose{สเหตุโก เจว น จ เหตุ ธมฺโม เหตุสฺส เจว สเหตุกสฺส จ ธมฺมสฺส อธิปติ\-ปจฺจเยน ปจฺจโย. อารมฺมณา\-ธิปติ, สหชาตา\-ธิปติ.  อารมฺมณา\-ธิปติ\csromandash{}ทานํ ทตฺวา. (ปฐมคมนํเยว.) สหชาตา\-ธิปติ\csromandash{}สเหตุโก เจว น จ เหตุ อธิปติ สมฺปยุตฺตกานํ เหตูนํ อธิปติ\-ปจฺจเยน ปจฺจโย. (2)}

\prose{สเหตุโก เจว น จ เหตุ ธมฺโม เหตุสฺส เจว สเหตุกสฺส จ สเหตุกสฺส เจว น จ เหตุสฺส จ ธมฺมสฺส อธิปติ\-ปจฺจเยน ปจฺจโย. อารมฺมณา\-ธิปติ, สหชาตา\-ธิปติ.  อารมฺมณา\-ธิปติ\csromandash{}ทานํ ทตฺวา. (ปฐมคมนํเยว.) สหชาตา\-ธิปติ\csromandash{}สเหตุโก เจว น จ เหตุ อธิปติ สมฺปยุตฺตกานํ ขนฺธานํ เหตูนญฺจ อธิปติ\-ปจฺจเยน ปจฺจโย. (3)}

\proseitem{141}{เหตุ เจว สเหตุโก จ สเหตุโก เจว น จ เหตุ จ ธมฺมา เหตุสฺส เจว สเหตุกสฺส จ ธมฺมสฺส อธิปติ\-ปจฺจเยน ปจฺจโย. อารมฺมณา\-ธิปติ\csromandash{}เหตุญฺจ สมฺปยุตฺตเก จ ขนฺเธ ครุํ กตฺวา เหตู อุปฺปชฺชนฺติ. (1)}

\prose{เหตุ เจว สเหตุโก จ สเหตุโก เจว น จ เหตุ จ ธมฺมา สเหตุกสฺส เจว น จ เหตุสฺส ธมฺมสฺส อธิปติ\-ปจฺจเยน ปจฺจโย. \csromanfolio{62}อารมฺมณา\-ธิปติ\csromandash{}เหตุญฺจ สมฺปยุตฺตเก จ ขนฺเธ ครุํ กตฺวา สเหตุกา เจว น จ เหตู ขนฺธา อุปฺปชฺชนฺติ. (2)}

\prose{เหตุ เจว สเหตุโก จ สเหตุโก เจว น จ เหตุ จ ธมฺมา เหตุสฺส เจว สเหตุกสฺส จ สเหตุกสฺส เจว น จ เหตุสฺส จ ธมฺมสฺส อธิปติ\-ปจฺจเยน ปจฺจโย. อารมฺมณา\-ธิปติ\csromandash{}เหตุญฺจ สมฺปยุตฺตเก จ ขนฺเธ ครุํ กตฺวา เหตู จ สมฺปยุตฺตกา จ ขนฺธา อุปฺปชฺชนฺติ. (3)}

\proseitem{142}{เหตุ เจว สเหตุโก จ ธมฺโม เหตุสฺส เจว สเหตุกสฺส จ ธมฺมสฺส อนนฺตร\-ปจฺจเยน ปจฺจโย. ปุริมา ปุริมา เหตู ปจฺฉิมานํ ปจฺฉิมานํ เหตูนํ อนนฺตร\-ปจฺจเยน ปจฺจโย. (1)}

\prose{เหตุ เจว สเหตุโก จ ธมฺโม สเหตุกสฺส เจว น จ เหตุสฺส ธมฺมสฺส อนนฺตร\-ปจฺจเยน ปจฺจโย. ปุริมา ปุริมา เหตู ปจฺฉิมานํ ปจฺฉิมานํ สเหตุกาน\-ญฺเจว น จ เหตูนํ ขนฺธานํ อนนฺตร\-ปจฺจเยน ปจฺจโย. (2)}

\prose{เหตุ เจว สเหตุโก จ ธมฺโม เหตุสฺส เจว สเหตุกสฺส จ สเหตุกสฺส เจว น จ เหตุสฺส จ ธมฺมสฺส อนนฺตร\-ปจฺจเยน ปจฺจโย. ปุริมา ปุริมา เหตู ปจฺฉิมานํ ปจฺฉิมานํ เหตูนํ สมฺปยุตฺตกานญฺจ ขนฺธานํ อนนฺตร\-ปจฺจเยน ปจฺจโย. (3)}

\proseitem{143}{สเหตุโก เจว น จ เหตุ ธมฺโม สเหตุกสฺส เจว น จ เหตุสฺส ธมฺมสฺส อนนฺตร\-ปจฺจเยน ปจฺจโย. ปุริมา ปุริมา สเหตุกา เจว น จ เหตู ขนฺธา ปจฺฉิมานํ ปจฺฉิมานํ สเหตุกาน\-ญฺเจว น จ เหตูนํ ขนฺธานํ อนนฺตร\-ปจฺจเยน ปจฺจโย. อนุโลมํ โคตฺรภุสฺส. อนุโลมํ โวทานสฺส ฯเปฯ. นิโรธา วุฏฺฐหนฺตสฺส. เนว\-สญฺญา\-นา\-สญฺญา\-ยตนํ ผลสมาปตฺติยา อนนฺตร\-ปจฺจเยน ปจฺจโย. (1)}

\prose{สเหตุโก เจว น จ เหตุ ธมฺโม เหตุสฺส เจว สเหตุกสฺส จ ธมฺมสฺส อนนฺตร\-ปจฺจเยน ปจฺจโย. ปุริมา ปุริมา สเหตุกา เจว น จ เหตู ขนฺธา ปจฺฉิมานํ ปจฺฉิมานํ เหตูนํ อนนฺตร\-ปจฺจเยน ปจฺจโย. อนุโลมํ โคตฺรภุสฺส. (สํขิตฺตํ.) (2)}

\prose{\csromanfolio{63}สเหตุโก เจว น จ เหตุ ธมฺโม เหตุสฺส เจว สเหตุกสฺส จ สเหตุกสฺส เจว น จ เหตุสฺส จ ธมฺมสฺส อนนฺตร\-ปจฺจเยน ปจฺจโย. ปุริมา ปุริมา สเหตุกา เจว น จ เหตู ขนฺธา ปจฺฉิมานํ ปจฺฉิมานํ เหตูนํ สมฺปยุตฺตกานญฺจ ขนฺธานํ อนนฺตร\-ปจฺจเยน ปจฺจโย. อนุโลมํ โคตฺรภุสฺส ฯเปฯ. (3)}

\prose{(สเหตุโก เจว น จ เหตุมูลกํ. ตีณิปิ เอกสทิสา.)}

\proseitem{144}{เหตุ เจว สเหตุโก จ สเหตุโก เจว น จ เหตุ จ ธมฺมา เหตุสฺส เจว สเหตุกสฺส จ ธมฺมสฺส อนนฺตร\-ปจฺจเยน ปจฺจโย. ปุริมา ปุริมา เหตู จ สมฺปยุตฺตกา จ ขนฺธา ปจฺฉิมานํ ปจฺฉิมานํ เหตูนํ อนนฺตร\-ปจฺจเยน ปจฺจโย. (1)}

\prose{เหตุ เจว สเหตุโก จ สเหตุโก เจว น จ เหตุ จ ธมฺมา สเหตุกสฺส เจว น จ เหตุสฺส ธมฺมสฺส อนนฺตร\-ปจฺจเยน ปจฺจโย. ปุริมา ปุริมา เหตู จ สมฺปยุตฺตกา จ ขนฺธา ปจฺฉิมานํ ปจฺฉิมานํ สเหตุกาน\-ญฺเจว น จ เหตูนํ ขนฺธานํ อนนฺตร\-ปจฺจเยน ปจฺจโย. (2)}

\prose{เหตุ เจว สเหตุโก จ สเหตุโก เจว น จ เหตุ จ ธมฺมา เหตุสฺส เจว สเหตุกสฺส จ สเหตุกสฺส เจว น จ เหตุสฺส จ ธมฺมสฺส อนนฺตร\-ปจฺจเยน ปจฺจโย. ปุริมา ปุริมา เหตู จ สมฺปยุตฺตกา จ ขนฺธา ปจฺฉิมานํ ปจฺฉิมานํ เหตูนํ สมฺปยุตฺตกานญฺจ ขนฺธานํ อนนฺตร\-ปจฺจเยน ปจฺจโย. (3)}

\csromanheader{2}{\textbf{สหชาตาทิ}}
\proseitem{145}{เหตุ เจว สเหตุโก จ ธมฺโม เหตุสฺส เจว สเหตุกสฺส จ ธมฺมสฺส สหชาต\-ปจฺจเยน ปจฺจโย. อญฺญมญฺญ\-ปจฺจเยน ปจฺจโย. นิสฺสย\-ปจฺจเยน ปจฺจโย. (ตีณิปิ ปจฺจยา ปฏิจฺจวาเร เหตุสทิสา.)}

\csromanheader{2}{\textbf{อุปนิสฺสย}}
\proseitem{146}{เหตุ เจว สเหตุโก จ ธมฺโม เหตุสฺส เจว สเหตุกสฺส จ ธมฺมสฺส อุปนิสฺสย\-ปจฺจเยน ปจฺจโย. อารมฺมณู\-ปนิสฺสโย, อนนฺตรู\-ปนิสฺสโย, ปกตู\-ปนิสฺสโย ฯเปฯ. ปกตู\-ปนิสฺสโย\csromandash{}เหตู เหตูนํ อุปนิสฺสย\-ปจฺจเยน ปจฺจโย. (มูลํ กาตพฺพํ.) เหตู \csromanfolio{64}สเหตุกาน\-ญฺเจว น จ เหตูนํ ขนฺธานํ อุปนิสฺสย\-ปจฺจเยน ปจฺจโย. (มูลํ กาตพฺพํ.) เหตู เหตูนํ สมฺปยุตฺตกานญฺจ ขนฺธานํ อุปนิสฺสย\-ปจฺจเยน ปจฺจโย.}

\vspace{\tipitakaparskip}
\csromancenter{(อิเมสํ ทฺวินฺนมฺปิ ปญฺหานํ มูลานิ ปุจฺฉิตพฺพานิ.)}
\prose{สเหตุโก เจว น จ เหตุ ธมฺโม สเหตุกสฺส เจว น จ เหตุสฺส จ ธมฺมสฺส อุปนิสฺสย\-ปจฺจเยน ปจฺจโย. อารมฺมณู\-ปนิสฺสโย, อนนฺตรู\-ปนิสฺสโย, ปกตู\-ปนิสฺสโย ฯเปฯ. ปกตู\-ปนิสฺสโย\csromandash{}สทฺธํ อุปนิสฺสาย ทานํ เทติ ฯเปฯ สมาปตฺติํ อุปฺปาเทติ, มานํ ชปฺเปติ, ทิฏฺฐิํ คณฺหาติ, สีลํ ฯเปฯ ปตฺถนํ อุปนิสฺสาย ทานํ เทติ ฯเปฯ สํฆํ ภินฺทติ. สทฺธา ฯเปฯ. ปตฺถนา สทฺธาย ฯเปฯ ปตฺถนาย มคฺคสฺส ผลสมาปตฺติยา อุปนิสฺสย\-ปจฺจเยน ปจฺจโย. (1)}

\prose{(สเหตุโก เจว น จ เหตุมูลเก อิมินากาเรน วิตฺถาเรตพฺพา. อวเสสา เทฺว ปญฺหา.)}

\prose{เหตุ เจว สเหตุโก จ สเหตุโก เจว น จ เหตุ จ ธมฺมา เหตุสฺส เจว สเหตุกสฺส จ ธมฺมสฺส อุปนิสฺสย\-ปจฺจเยน ปจฺจโย. อารมฺมณู\-ปนิสฺสโย, อนนฺตรู\-ปนิสฺสโย, ปกตู\-ปนิสฺสโย ฯเปฯ. ปกตู\-ปนิสฺสโย\csromandash{}เหตู จ สมฺปยุตฺตกา จ ขนฺธา เหตูนํ อุปนิสฺสย\-ปจฺจเยน ปจฺจโย. (เทฺว มูลานิ ปุจฺฉิตพฺพานิ.) เหตู จ สมฺปยุตฺตกา จ ขนฺธา สเหตุกาน\-ญฺเจว น จ เหตูนํ ขนฺธานํ อุปนิสฺสย\-ปจฺจเยน ปจฺจโย. (มูลํ ปุจฺฉิตพฺพํ.) เหตู จ สมฺปยุตฺตกา จ ขนฺธา เหตูนํ สมฺปยุตฺตกานญฺจ ขนฺธานํ อุปนิสฺสย\-ปจฺจเยน ปจฺจโย. (1)}

\csromanheader{2}{\textbf{อาเสวน}}
\proseitem{147}{เหตุ เจว สเหตุโก จ ธมฺโม เหตุสฺส เจว สเหตุกสฺส จ ธมฺมสฺส อาเสวน\-ปจฺจเยน ปจฺจโย. (อนนฺตร\-สทิสํ.)}

\proseitem{148}{สเหตุโก เจว น จ เหตุ ธมฺโม สเหตุกสฺส เจว น จ เหตุสฺส ธมฺมสฺส กมฺมปจฺจเยน ปจฺจโย. สหชาตา, นานากฺขณิกา.  สหชาตา สเหตุกา เจว น จ เหตู เจตนา สมฺปยุตฺตกานํ ขนฺธานํ กมฺมปจฺจเยน ปจฺจโย. ปฏิสนฺธิ\-กฺขเณ ฯเปฯ. \csromanfolio{65}นานากฺขณิกา สเหตุกา เจว น จ เหตู เจตนา วิปากานํ สเหตุกาน\-ญฺเจว น จ เหตูนํ ขนฺธานํ กมฺมปจฺจเยน ปจฺจโย. (1)}

\prose{สเหตุโก เจว น จ เหตุ ธมฺโม เหตุสฺส เจว สเหตุกสฺส จ ธมฺมสฺส กมฺมปจฺจเยน ปจฺจโย. สหชาตา, นานากฺขณิกา.  สหชาตา สเหตุกา เจว น จ เหตู เจตนา สมฺปยุตฺตกานํ เหตูนํ กมฺมปจฺจเยน ปจฺจโย. ปฏิสนฺธิ\-กฺขเณ ฯเปฯ. นานากฺขณิกา สเหตุกา เจว น จ เหตู เจตนา วิปากานํ เหตูนํ กมฺมปจฺจเยน ปจฺจโย. (2)}

\prose{สเหตุโก เจว น จ เหตุ ธมฺโม เหตุสฺส เจว สเหตุกสฺส จ สเหตุกสฺส เจว น จ เหตุสฺส จ ธมฺมสฺส กมฺมปจฺจเยน ปจฺจโย. สหชาตา, นานากฺขณิกา.  สหชาตา สเหตุกา เจว น จ เหตู เจตนา สมฺปยุตฺตกานํ ขนฺธานํ เหตูนญฺจ กมฺมปจฺจเยน ปจฺจโย. ปฏิสนฺธิ\-กฺขเณ ฯเปฯ. นานากฺขณิกา สเหตุกา เจว น จ เหตู เจตนา วิปากานํ ขนฺธานํ เหตูนญฺจ กมฺมปจฺจเยน ปจฺจโย. (3)}

\csromanheader{2}{\textbf{วิปาก}}
\proseitem{149}{เหตุ เจว สเหตุโก จ ธมฺโม เหตุสฺส เจว สเหตุกสฺส จ ธมฺมสฺส วิปาก\-ปจฺจเยน ปจฺจโย. วิปาโก อโลโภ อโทสสฺส อโมหสฺส จ วิปาก\-ปจฺจเยน ปจฺจโย. (จกฺกํ.) ปฏิสนฺธิ\-กฺขเณ อโลโภ. (ยถา เหตุปจฺจยา, เอวํ วิตฺถาเรตพฺพํ. นวปิ ‘วิปากนฺ’ติ นิยาเมตพฺพํ.)}

\csromanheader{2}{\textbf{อาหาราทิ}}
\proseitem{150}{สเหตุโก เจว น จ เหตุ ธมฺโม สเหตุกสฺส เจว น จ เหตุสฺส ธมฺมสฺส อาหารปจฺจเยน ปจฺจโย. ตีณิ.}

\prose{เหตุ เจว สเหตุโก จ ธมฺโม เหตุสฺส เจว สเหตุกสฺส จ ธมฺมสฺส อินฺทฺริย\-ปจฺจเยน ปจฺจโย. (‘อินฺทฺริยนฺ’ติ นิยาเมตพฺพํ. นวปิ ปริปุณฺณํ.)}

\prose{สเหตุโก เจว น จ เหตุ ธมฺโม สเหตุกสฺส เจว น จ เหตุสฺส ธมฺมสฺส ฌานปจฺจเยน ปจฺจโย. ตีณิ.}

\prose{\csromanfolio{66}เหตุ เจว สเหตุโก จ ธมฺโม เหตุสฺส เจว สเหตุกสฺส จ ธมฺมสฺส มคฺคปจฺจเยน ปจฺจโย. สมฺปยุตฺต\-ปจฺจเยน ปจฺจโย. อตฺถิปจฺจเยน ปจฺจโย. นตฺถิ\-ปจฺจเยน ปจฺจโย. วิคต\-ปจฺจเยน ปจฺจโย. อวิคต\-ปจฺจเยน ปจฺจโย.}

\csromanmarkh{1. ปจฺจยานุโลม}\csromanmarkhii{2. สงฺขฺยาวาร}\csromanheadernomark{2}{1. \textbf{ปจฺจยานุโลม 2. สงฺขฺยาวาร}}
\proseitemwithcloserruled{151}{เหตุยา ตีณิ, อารมฺมเณ นว, อธิปติยา นว, อนนฺตเร นว, สมนนฺตเร นว, สหชาเต นว, อญฺญมญฺเญ นว, นิสฺสเย นว, อุปนิสฺสเย นว, อาเสวเน นว, กมฺเม ตีณิ, วิปาเก นว, อาหาเร ตีณิ, อินฺทฺริเย นว, ฌาเน ตีณิ, มคฺเค นว, สมฺปยุตฺเต นว, อตฺถิยา นว, นตฺถิยา นว, วิคเต นว, อวิคเต นว. (เอวํ คเณตพฺพํ.)}{อนุโลมํ.}
\proseitem{152}{เหตุ เจว สเหตุโก จ ธมฺโม เหตุสฺส เจว สเหตุกสฺส จ ธมฺมสฺส อารมฺมณ\-ปจฺจเยน ปจฺจโย. สหชาต\-ปจฺจเยน ปจฺจโย. อุปนิสฺสย\-ปจฺจเยน ปจฺจโย. (1)}

\prose{เหตุ เจว สเหตุโก จ ธมฺโม สเหตุกสฺส เจว น จ เหตุสฺส ธมฺมสฺส อารมฺมณ\-ปจฺจเยน ปจฺจโย. สหชาต\-ปจฺจเยน ปจฺจโย. อุปนิสฺสย\-ปจฺจเยน ปจฺจโย. (2)}

\prose{เหตุ เจว สเหตุโก จ ธมฺโม เหตุสฺส เจว สเหตุกสฺส จ สเหตุกสฺส เจว น จ เหตุสฺส จ ธมฺมสฺส อารมฺมณ\-ปจฺจเยน ปจฺจโย. สหชาต\-ปจฺจเยน ปจฺจโย. อุปนิสฺสย\-ปจฺจเยน ปจฺจโย. (3)}

\proseitem{153}{สเหตุโก เจว น จ เหตุ ธมฺโม สเหตุกสฺส เจว น จ เหตุสฺส ธมฺมสฺส อารมฺมณ\-ปจฺจเยน ปจฺจโย. สหชาต\-ปจฺจเยน ปจฺจโย. อุปนิสฺสย\-ปจฺจเยน ปจฺจโย. กมฺมปจฺจเยน ปจฺจโย. (1)}

\prose{\csromanfolio{67}สเหตุโก เจว น จ เหตุ ธมฺโม เหตุสฺส เจว สเหตุกสฺส จ ธมฺมสฺส อารมฺมณ\-ปจฺจเยน ปจฺจโย. สหชาต\-ปจฺจเยน ปจฺจโย. อุปนิสฺสย\-ปจฺจเยน ปจฺจโย. กมฺมปจฺจเยน ปจฺจโย. (2)}

\prose{สเหตุโก เจว น จ เหตุ ธมฺโม เหตุสฺส เจว สเหตุกสฺส จ สเหตุกสฺส เจว น จ เหตุสฺส จ ธมฺมสฺส อารมฺมณ\-ปจฺจเยน ปจฺจโย. สหชาต\-ปจฺจเยน ปจฺจโย. อุปนิสฺสย\-ปจฺจเยน ปจฺจโย. กมฺมปจฺจเยน ปจฺจโย. (3)}

\proseitem{154}{เหตุ เจว สเหตุโก จ สเหตุโก เจว น จ เหตุ จ ธมฺมา เหตุสฺส เจว สเหตุกสฺส จ ธมฺมสฺส อารมฺมณ\-ปจฺจเยน ปจฺจโย. สหชาต\-ปจฺจเยน ปจฺจโย. อุปนิสฺสย\-ปจฺจเยน ปจฺจโย. (1)}

\prose{เหตุ เจว สเหตุโก จ สเหตุโก เจว น จ เหตุ จ ธมฺมา สเหตุกสฺส เจว น จ เหตุสฺส ธมฺมสฺส อารมฺมณ\-ปจฺจเยน ปจฺจโย. สหชาต\-ปจฺจเยน ปจฺจโย. อุปนิสฺสย\-ปจฺจเยน ปจฺจโย. (2)}

\prose{เหตุ เจว สเหตุโก จ สเหตุโก เจว น จ เหตุ จ ธมฺมา เหตุสฺส เจว สเหตุกสฺส จ สเหตุกสฺส เจว น จ เหตุสฺส จ ธมฺมสฺส อารมฺมณ\-ปจฺจเยน ปจฺจโย. สหชาต\-ปจฺจเยน ปจฺจโย. อุปนิสฺสย\-ปจฺจเยน ปจฺจโย. (3)}

\csromanmarkh{2. ปจฺจย\-ปจฺจนีย}\csromanmarkhii{2. สงฺขฺยาวาร}\csromanheadernomark{2}{2. ปจฺจย\-ปจฺจนีย 2. สงฺขฺยาวาร}
\proseitem{155}{นเหตุยา นว. (สํขิตฺตํ. สพฺพตฺถ นว. เอวํ คเณตพฺพํ.)}

\vspace{\tipitakaparskip}
\csromancenter{ปจฺจนียํ.}
\proseitem{156}{เหตุปจฺจยา นอารมฺมเณ ตีณิ, นอธิปติยา ตีณิ, นอนนฺตเร ตีณิ, นสมนนฺตเร ตีณิ, นอุปนิสฺสเย ตีณิ, (สํขิตฺตํ. สพฺพตฺถ \csromanfolio{68}\textbf{ตีณิ.) นมคฺเค ตีณิ, โนนตฺถิยา ตีณิ, โนวิคเต ตีณิ. (เอวํ คเณตพฺพํ.)}}

\vspace{\tipitakaparskip}
\csromancenter{อนุโลม\-ปจฺจนียํ.}
\proseitem{157}{นเหตุปจฺจยา อารมฺมเณ นว, อธิปติยา นว, อนนฺตเร นว, สมนนฺตเร นว, สหชาเต ตีณิ, อญฺญมญฺเญ ตีณิ, นิสฺสเย ตีณิ, อุปนิสฺสเย นว, อาเสวเน นว, กมฺเม ตีณิ, วิปาเก ตีณิ, อาหาเร ตีณิ, อินฺทฺริเย ตีณิ, ฌาเน ตีณิ, มคฺเค ตีณิ, สมฺปยุตฺเต ตีณิ, อตฺถิยา ตีณิ, นตฺถิยา นว, วิคเต นว, อวิคเต ตีณิ. (เอวํ คเณตพฺพํ.)}

\vspace{\tipitakaparskip}
\csromancenter{ปจฺจนียา\-นุโลมํ.}
\nitthitam{เหตุ\-สเหตุกทุกํ นิฏฺฐิตํ.}
\csromanmatikaanchor{mtk.7}
\csromantocmarkref{section}{5. เหตุเหตุสมฺปยุตฺตทุก}{mtk.7}
\bookmark[dest={mtk.7},level=1]{5. เหตุเหตุสมฺปยุตฺตทุก}
\csromanmarkh{5. เหตุ\-เหตุ\-สมฺปยุตฺตทุก}\csromanmarkhii{1. ปฏิจฺจวาร}\csromanheadernomark{1}{5. เหตุ\-เหตุ\-สมฺปยุตฺตทุก 1. ปฏิจฺจวาร}
\proseitem{158}{เหตุญฺเจว เหตุ\-สมฺปยุตฺตญฺจ ธมฺมํ ปฏิจฺจ เหตุ เจว เหตุสมฺปยุตฺโต จ ธมฺโม อุปฺปชฺชติ เหตุปจฺจยา. อโลภํ ปฏิจฺจ อโทโส อโมโห. (จกฺกํ.) โลภํ ปฏิจฺจ โมโห. (จกฺกํ.) ปฏิสนฺธิ\-กฺขเณ ฯเปฯ.}

\prosewithcloser{(ยถา เหตุ\-สเหตุกทุกํ, เอวํ วิตฺถาเรตพฺพํ, นินฺนานากรณํ.)}{เหตุ\-เหตุ\-สมฺปยุตฺตทุกํ นิฏฺฐิตํ.}
\csromanmatikaanchor{mtk.8}
\csromantocmarkref{section}{6. นเหตุสเหตุกทุก}{mtk.8}
\bookmark[dest={mtk.8},level=1]{6. นเหตุสเหตุกทุก}
\csromanmarkh{6. นเหตุสเหตุทุก}\csromanmarkhii{6. นเหตุ\-สเหตุกทุก 1. ปฏิจฺจวาร}\csromanheadernomark{1}{6. นเหตุสเหตุทุก 6. นเหตุ\-สเหตุกทุก 1. ปฏิจฺจวาร}
\csromanmarkh{1. ปจฺจยานุโลม}\csromanmarkhii{1. วิภงฺควาร}\csromanheadernomark{2}{1. \textbf{ปจฺจยานุโลม 1. วิภงฺควาร}}
\csromanheader{2}{\textbf{เหตุ}}
\proseitem{159}{\csromanfolio{69}นเหตุํ สเหตุกํ ธมฺมํ ปฏิจฺจ นเหตุ สเหตุโก ธมฺโม อุปฺปชฺชติ เหตุปจฺจยา. นเหตุํ สเหตุกํ เอกํ ขนฺธํ ปฏิจฺจ ตโย ขนฺธา ฯเปฯ เทฺว ขนฺเธ ฯเปฯ. ปฏิสนฺธิ\-กฺขเณ ฯเปฯ. (1)}

\prose{นเหตุํ สเหตุกํ ธมฺมํ ปฏิจฺจ นเหตุ อเหตุโก ธมฺโม อุปฺปชฺชติ เหตุปจฺจยา. นเหตู สเหตุเก ขนฺเธ ปฏิจฺจ จิตฺต\-สมุฏฺฐานํ รูปํ. ปฏิสนฺธิ\-กฺขเณ ฯเปฯ. (2)}

\prose{นเหตุํ สเหตุกํ ธมฺมํ ปฏิจฺจ นเหตุ สเหตุโก จ นเหตุ อเหตุโก จ ธมฺมา อุปฺปชฺชนฺติ เหตุปจฺจยา. นเหตุํ สเหตุกํ เอกํ ขนฺธํ ปฏิจฺจ ตโย ขนฺธา จิตฺต\-สมุฏฺฐานญฺจ รูปํ ฯเปฯ เทฺว ขนฺเธ ฯเปฯ. ปฏิสนฺธิ\-กฺขเณ ฯเปฯ. (3)}

\proseitem{160}{นเหตุํ อเหตุกํ ธมฺมํ ปฏิจฺจ นเหตุ อเหตุโก ธมฺโม อุปฺปชฺชติ เหตุปจฺจยา ฯเปฯ. เอกํ มหาภูตํ ปฏิจฺจ ฯเปฯ มหาภูเต ปฏิจฺจ จิตฺต\-สมุฏฺฐานํ รูปํ, กฏตฺตารูปํ, อุปาทารูปํ. (1)}

\prose{นเหตุํ อเหตุกํ ธมฺมํ ปฏิจฺจ นเหตุ สเหตุโก ธมฺโม อุปฺปชฺชติ เหตุปจฺจยา. ปฏิสนฺธิ\-กฺขเณ วตฺถุํ ปฏิจฺจ นเหตู สเหตุกา ขนฺธา. (2)}

\prose{นเหตุํ อเหตุกํ ธมฺมํ ปฏิจฺจ นเหตุ สเหตุโก จ นเหตุ อเหตุโก จ ธมฺมา อุปฺปชฺชนฺติ เหตุปจฺจยา. ปฏิสนฺธิ\-กฺขเณ วตฺถุํ ปฏิจฺจ นเหตู สเหตุกา ขนฺธา. มหาภูเต ปฏิจฺจ กฏตฺตารูปํ. (3)}

\proseitem{161}{นเหตุํ สเหตุกญฺจ นเหตุํ อเหตุกญฺจ ธมฺมํ ปฏิจฺจ นเหตุ สเหตุโก ธมฺโม อุปฺปชฺชติ เหตุปจฺจยา. ปฏิสนฺธิ\-กฺขเณ นเหตุํ สเหตุกํ เอกํ ขนฺธญฺจ วตฺถุญฺจ ปฏิจฺจ ตโย ขนฺธา ฯเปฯ เทฺว ขนฺเธ จ ฯเปฯ. (1)}

\prose{นเหตุํ สเหตุกญฺจ นเหตุํ อเหตุกญฺจ ธมฺมํ ปฏิจฺจ นเหตุ อเหตุโก ธมฺโม อุปฺปชฺชติ เหตุปจฺจยา. นเหตู สเหตุเก ขนฺเธ จ มหาภูเต จ ปฏิจฺจ จิตฺต\-สมุฏฺฐานํ รูปํ. ปฏิสนฺธิ\-กฺขเณ ฯเปฯ. (2)}

\prose{\csromanfolio{70}นเหตุํ สเหตุกญฺจ นเหตุํ อเหตุกญฺจ ธมฺมํ ปฏิจฺจ นเหตุ สเหตุโก จ นเหตุ อเหตุโก จ ธมฺมา อุปฺปชฺชนฺติ เหตุปจฺจยา. ปฏิสนฺธิ\-กฺขเณ นเหตุํ สเหตุกํ เอกํ ขนฺธญฺจ วตฺถุญฺจ ปฏิจฺจ ตโย ขนฺธา ฯเปฯ เทฺว ขนฺเธ ฯเปฯ. นเหตู สเหตุเก ขนฺเธ จ มหาภูเต จ ปฏิจฺจ กฏตฺตารูปํ. (3)}

\proseitem{162}{นเหตุํ สเหตุกํ ธมฺมํ ปฏิจฺจ นเหตุ สเหตุโก ธมฺโม อุปฺปชฺชติ อารมฺมณ\-ปจฺจยา. นเหตุํ สเหตุกํ เอกํ ขนฺธํ ปฏิจฺจ ตโย ขนฺธา ฯเปฯ เทฺว ขนฺเธ ฯเปฯ. ปฏิสนฺธิ\-กฺขเณ ฯเปฯ. (1)}

\prose{นเหตุํ อเหตุกํ ธมฺมํ ปฏิจฺจ นเหตุ อเหตุโก ธมฺโม อุปฺปชฺชติ อารมฺมณ\-ปจฺจยา. นเหตุํ อเหตุกํ เอกํ ขนฺธํ ปฏิจฺจ ตโย ขนฺธา ฯเปฯ. ปฏิสนฺธิ\-กฺขเณ ฯเปฯ. (2)}

\prose{นเหตุํ อเหตุกํ ธมฺมํ ปฏิจฺจ นเหตุ สเหตุโก ธมฺโม อุปฺปชฺชติ อารมฺมณ\-ปจฺจยา. ปฏิสนฺธิ\-กฺขเณ วตฺถุํ ปฏิจฺจ นเหตู สเหตุกา ขนฺธา. (3)}

\prose{นเหตุํ สเหตุกญฺจ นเหตุํ อเหตุกญฺจ ธมฺมํ ปฏิจฺจ นเหตุ สเหตุโก ธมฺโม อุปฺปชฺชติ อารมฺมณ\-ปจฺจยา. ปฏิสนฺธิ\-กฺขเณ นเหตุํ สเหตุกํ เอกํ ขนฺธญฺจ วตฺถุญฺจ ปฏิจฺจ ตโย ขนฺธา ฯเปฯ เทฺว ขนฺเธ ฯเปฯ. (สํขิตฺตํ. เอวํ วิภชิตพฺพํ.)}

\csromanmarkh{1. ปจฺจยานุโลม}\csromanmarkhii{2. สงฺขฺยาวาร}\csromanheadernomark{2}{1. \textbf{ปจฺจยานุโลม 2. สงฺขฺยาวาร}}
\proseitem{163}{เหตุยา นว, อารมฺมเณ จตฺตาริ, อธิปติยา ปญฺจ, อนนฺตเร จตฺตาริ, สมนนฺตเร จตฺตาริ, สหชาเต นว, อญฺญมญฺเญ ฉ, นิสฺสเย นว, อุปนิสฺสเย จตฺตาริ, ปุเรชาเต เทฺว, อาเสวเน เทฺว, กมฺเม นว, วิปาเก นว, อาหาเร นว, (สํขิตฺตํ. สพฺพตฺถ นว.) สมฺปยุตฺเต จตฺตาริ, วิปฺปยุตฺเต นว, อตฺถิยา นว, นตฺถิยา จตฺตาริ, วิคเต จตฺตาริ, อวิคเต นว. (เอวํ คเณตพฺพํ.)}

\vspace{\tipitakaparskip}
\csromancenter{อนุโลมํ.}
\csromanmarkh{6. นเหตุสเหตุทุก}\csromanmarkhii{2. ปจฺจย\-ปจฺจนีย 1. วิภงฺควาร}\csromanheadernomark{2}{6. นเหตุสเหตุทุก 2. ปจฺจย\-ปจฺจนีย 1. วิภงฺควาร}
\proseitem{164}{\csromanfolio{71}นเหตุํ อเหตุกํ ธมฺมํ ปฏิจฺจ นเหตุ อเหตุโก ธมฺโม อุปฺปชฺชติ นเหตุปจฺจยา. นเหตุํ อเหตุกํ เอกํ ขนฺธํ ปฏิจฺจ ตโย ขนฺธา จิตฺต\-สมุฏฺฐานญฺจ รูปํ ฯเปฯ เทฺว ขนฺเธ ฯเปฯ. ปฏิสนฺธิ\-กฺขเณ. (ยาว อสญฺญสตฺตา โมโห นตฺถิ.) (1)}

\prose{นอารมฺมณ}

\proseitem{165}{นเหตุํ สเหตุกํ ธมฺมํ ปฏิจฺจ นเหตุ อเหตุโก ธมฺโม อุปฺปชฺชติ นอารมฺมณ\-ปจฺจยา. นเหตู สเหตุเก ขนฺเธ ปฏิจฺจ จิตฺต\-สมุฏฺฐานํ รูปํ. ปฏิสนฺธิ\-กฺขเณ ฯเปฯ. (1)}

\prose{นเหตุํ อเหตุกํ ธมฺมํ ปฏิจฺจ นเหตุ อเหตุโก ธมฺโม อุปฺปชฺชติ นอารมฺมณ\-ปจฺจยา. นเหตู อเหตุเก ขนฺเธ ปฏิจฺจ จิตฺต\-สมุฏฺฐานํ รูปํ. ปฏิสนฺธิ\-กฺขเณ ฯเปฯ. (ยาว อสญฺญสตฺตา.) (1)}

\prose{นเหตุํ สเหตุกญฺจ นเหตุํ อเหตุกญฺจ ธมฺมํ ปฏิจฺจ นเหตุ อเหตุโก ธมฺโม อุปฺปชฺชติ นอารมฺมณ\-ปจฺจยา. นเหตู สเหตุเก ขนฺเธ จ มหาภูเต จ ปฏิจฺจ จิตฺต\-สมุฏฺฐานํ รูปํ. ปฏิสนฺธิ\-กฺขเณ ฯเปฯ. (สํขิตฺตํ.)}

\csromanmarkh{2. ปจฺจย\-ปจฺจนีย}\csromanmarkhii{2. สงฺขฺยาวาร}\csromanheadernomark{2}{2. ปจฺจย\-ปจฺจนีย 2. สงฺขฺยาวาร}
\vspace{\tipitakaparskip}
\csromancenter{\textbf{นเหตุ}ยา เอกํ, นอารมฺมเณ ตีณิ, นอธิปติยา นว, นอนนฺตเร ตีณิ, นสมนนฺตเร ตีณิ, นอญฺญมญฺเญ ตีณิ, นฺเอาปนิสฺสเย ตีณิ, นปุเรชาเต นว, นปจฺฉาชาเต นว, นอาเสวเน นว, นกมฺเม เทฺว, นวิปาเก ปญฺจ, นอาหาเร เอกํ, ไนนฺทฺริเย เอกํ, นฌาเน เอกํ, นมคฺเค เอกํ, นสมฺปยุตฺเต ตีณิ, นวิปฺปยุตฺเต เทฺว, โนนตฺถิยา ตีณิ, โนวิคเต ตีณิ. (เอวํ คเณตพฺพํ.)}
\csromancenter{ปจฺจนียํ.}
\proseitem{167}{\csromanfolio{72}เหตุปจฺจยา นอารมฺมเณ ตีณิ, นอธิปติยา นว, นอนนฺตเร นว, นสมนนฺตเร นว, นอญฺญมญฺเญ นว, นอุปนิสฺสเย ตีณิ, นปุเรชาเต นว, นปจฺฉาชาเต นว, นอาเสวเน นว, นกมฺเม เอกํ, นวิปาเก ปญฺจ, นสมฺปยุตฺเต ตีณิ, นวิปฺปยุตฺเต เอกํ, โนนตฺถิยา ตีณิ, โนวิคเต ตีณิ. (เอวํ คเณตพฺพํ.)}

\vspace{\tipitakaparskip}
\csromancenter{อนุโลม\-ปจฺจนียํ.}
\proseitem{168}{นเหตุปจฺจยา อารมฺมเณ เอกํ ฯเปฯ อาหาเร เอกํ ฯเปฯ ฌาเน เอกํ, สมฺปยุตฺเต เอกํ, วิปฺปยุตฺเต เอกํ, ฯเปฯ วิคเต เอกํ, อวิคเต เอกํ. (เอวํ คเณตพฺพํ.}

\vspace{\tipitakaparskip}
\csromancenter{ปจฺจนียา\-นุโลมํ.}
\csromancenter{(สหชาตวาเรปิ เอวํ คเณตพฺพํ.)}
\csromanmarkh{6. นเหตุ\-สเหตุกทุก}\csromanmarkhii{3. ปจฺจยวาร}\csromanheadernomark{2}{6. นเหตุ\-สเหตุกทุก 3. ปจฺจยวาร}
\proseitem{169}{นเหตุํ สเหตุกํ ธมฺมํ ปจฺจยา นเหตุ สเหตุโก ธมฺโม อุปฺปชฺชติ เหตุปจฺจยา. ตีณิ.}

\prose{นเหตุํ อเหตุกํ ธมฺมํ ปจฺจยา นเหตุ อเหตุโก ธมฺโม อุปฺปชฺชติ เหตุปจฺจยา. เอกํ มหาภูตํ ปจฺจยา ตโย มหาภูตา. มหาภูเต ปจฺจยา จิตฺต\-สมุฏฺฐานํ รูปํ, กฏตฺตารูปํ, อุปาทารูปํ. (1)}

\prose{\csromanfolio{73}นเหตุํ อเหตุกํ ธมฺมํ ปจฺจยา นเหตุ สเหตุโก ธมฺโม อุปฺปชฺชติ เหตุปจฺจยา. วตฺถุํ ปจฺจยา นเหตู สเหตุกา ขนฺธา. ปฏิสนฺธิ\-กฺขเณ ฯเปฯ. (2)}

\prose{นเหตุํ อเหตุกํ ธมฺมํ ปจฺจยา นเหตุ สเหตุโก จ นเหตุ อเหตุโก จ ธมฺมา อุปฺปชฺชนฺติ เหตุปจฺจยา. วตฺถุํ ปจฺจยา นเหตู สเหตุกา ขนฺธา. มหาภูเต ปจฺจยา จิตฺต\-สมุฏฺฐานํ รูปํ. ปฏิสนฺธิ\-กฺขเณ ฯเปฯ. (3)}

\prose{นเหตุํ สเหตุกญฺจ นเหตุํ อเหตุกญฺจ ธมฺมํ ปจฺจยา นเหตุ สเหตุโก ธมฺโม อุปฺปชฺชติ เหตุปจฺจยา. (ฆฏนา ตีณิ, ปวตฺติ\-ปฏิสนฺธิ ปริปุณฺณํ. สํขิตฺตํ.)}

\proseitem{170}{เหตุยา นว, อารมฺมเณ จตฺตาริ ฯเปฯ อญฺญมญฺเญ ฉ ฯเปฯ ปุเรชาเต, อาเสวเน จตฺตาริ ฯเปฯ อวิคเต นว. (เอวํ คเณตพฺพํ.)}

\vspace{\tipitakaparskip}
\csromancenter{\textbf{อนุโลมํ.}}
\csromancenter{นเหตุยา เอกํ, นอารมฺมเณ ตีณิ ฯเปฯ โนวิคเต ตีณิ.}
\csromancenter{ปจฺจนียํ.}
\csromancenter{(นิสฺสยวาโร ปจฺจยวาร\-สทิโส.)}
\csromanmarkh{6. นเหตุ\-สเหตุกทุก}\csromanmarkhii{5. สํสฏฺฐวาร}\csromanheadernomark{2}{6. นเหตุ\-สเหตุกทุก 5. สํสฏฺฐวาร}
\proseitem{172}{นเหตุํ สเหตุกํ ธมฺมํ สํสฏฺโฐ นเหตุ สเหตุโก ธมฺโม อุปฺปชฺชติ เหตุปจฺจยา. นเหตุํ สเหตุกํ เอกํ ขนฺธํ ฯเปฯ. ปฏิสนฺธิ\-กฺขเณ ฯเปฯ.}

\proseitem{173}{เหตุยา เอกํ, อารมฺมเณ เทฺว, อธิปติยา เอกํ, อนนฺตเร เทฺว, (สพฺพตฺถ เทฺว.) มคฺเค เอกํ ฯเปฯ อวิคเต เทฺว.}

\vspace{\tipitakaparskip}
\csromancenter{\textbf{อนุโลมํ.}}
\proseitem{174}{\csromanfolio{74}นเหตุํ อเหตุกํ ธมฺมํ สํสฏฺโฐ นเหตุ อเหตุโก ธมฺโม อุปฺปชฺชติ นเหตุปจฺจยา. นเหตุํ อเหตุกํ เอกํ ขนฺธํ ฯเปฯ. ปฏิสนฺธิ\-กฺขเณ ฯเปฯ.}

\proseitem{175}{นเหตุยา เอกํ, นอธิปติยา เทฺว, นปุเรชาเต เทฺว, นปจฺฉาชาเต เทฺว, นอาเสวเน เทฺว, นกมฺเม เทฺว, นวิปาเก เทฺว, นฌาเน เอกํ, นมคฺเค เอกํ, นวิปฺปยุตฺเต เทฺว. ปจฺจนียํ. (เอวํ อวเสสาปิ เทฺว คณนา คเณตพฺพา.) (สมฺปยุตฺตวาโร สํสฏฺฐวาร\-สทิโส.)}

\csromanmarkh{6. นเหตุ\-สเหตุกทุก}\csromanmarkhii{7. ปญฺหาวาร}\csromanheadernomark{2}{6. นเหตุ\-สเหตุกทุก 7. ปญฺหาวาร}
\csromanmarkh{1. ปจฺจยานุโลม}\csromanmarkhii{1. วิภงฺควาร}\csromanheadernomark{2}{1. \textbf{ปจฺจยานุโลม 1. วิภงฺควาร}}
\proseitem{176}{สเหตุโก ธมฺโม นเหตุ\-สเหตุกสฺส ธมฺมสฺส อารมฺมณ\-ปจฺจเยน ปจฺจโย. ทานํ ทตฺวา, สีลํ ฯเปฯ อุโปสถกมฺมํ กตฺวา ตํ ปจฺจเวกฺขติ. ปุพฺเพ สุจิณฺณานิ ปจฺจเวกฺขติ. ฌานํ ฯเปฯ. อริยา มคฺคา วุฏฺฐหิตฺวา มคฺคํ ปจฺจเวกฺขนฺติ, ผลํ ปจฺจเวกฺขนฺติ. ปหีเน กิเลเส ฯเปฯ. วิกฺขมฺภิเต กิเลเส ฯเปฯ. ปุพฺเพ ฯเปฯ. นเหตู สเหตุเก ขนฺเธ อนิจฺจโต ฯเปฯ โทมนสฺสํ อุปฺปชฺชติ. กุสลากุสเล นิรุทฺเธ นเหตุ สเหตุโก วิปาโก ตทารมฺมณตา อุปฺปชฺชติ. เจโตปริย\-ญาเณน นเหตุ\-สเหตุก\-จิตฺต\-สมงฺคิสฺส จิตฺตํ ชานาติ. อากาสา\-นญฺจา\-ยตนํ ฯเปฯ. อากิญฺจญฺญา\-ยตนํ ฯเปฯ. นเหตู สเหตุกา ขนฺธา อิทฺธิวิธ\-ญาณสฺส, เจโตปริย\-ญาณสฺส, ปุพฺเพ\-นิวาสา\-นุสฺสติ\-ญาณสฺส, ยถากมฺมูปคญาณสฺส, อนาคตํสญาณสฺส, อารมฺมณ\-ปจฺจเยน ปจฺจโย. นเหตู สเหตุเก ขนฺเธ อารพฺภ นเหตู สเหตุกา ขนฺธา อุปฺปชฺชนฺติ. (1)}

\prose{สเหตุโก ธมฺโม นเหตุ\-อเหตุกสฺส ธมฺมสฺส อารมฺมณ\-ปจฺจเยน ปจฺจโย. นเหตู สเหตุเก ขนฺเธ อนิจฺจโต ฯเปฯ \csromanfolio{75}โทมนสฺสํ อุปฺปชฺชติ. กุสลากุสเล นิรุทฺเธ นเหตุ อเหตุโก วิปาโก ตทารมฺมณตา อุปฺปชฺชติ. นเหตู สเหตุเก ขนฺเธ อารพฺภ นเหตู อเหตุกา ขนฺธา อุปฺปชฺชนฺติ. (2)}

\proseitem{177}{อเหตุโก ธมฺโม นเหตุ\-อเหตุกสฺส ธมฺมสฺส อารมฺมณ\-ปจฺจเยน ปจฺจโย. นิพฺพานํ อาวชฺชนาย อารมฺมณ\-ปจฺจเยน ปจฺจโย. จกฺขุํ ฯเปฯ วตฺถุํ. นเหตู อเหตุเก ขนฺเธ อนิจฺจโต ฯเปฯ โทมนสฺสํ อุปฺปชฺชติ. กุสลากุสเล นิรุทฺเธ นเหตุ อเหตุโก วิปาโก ตทารมฺมณตา อุปฺปชฺชติ. รูปายตนํ จกฺขุ\-วิญฺญาณสฺส ฯเปฯ. โผฏฺฐพฺพา\-ยตนํ กายวิญฺญาณสฺส ฯเปฯ. นเหตู อเหตุเก ขนฺเธ อารพฺภ นเหตู อเหตุกา ขนฺธา อุปฺปชฺชนฺติ. (1)}

\prose{อเหตุโก ธมฺโม นเหตุ\-สเหตุกสฺส ธมฺมสฺส อารมฺมณ\-ปจฺจเยน ปจฺจโย. อริยา นิพฺพานํ ปจฺจเวกฺขนฺติ. นิพฺพานํ โคตฺรภุสฺส, โวทานสฺส, มคฺคสฺส, ผลสฺส อารมฺมณ\-ปจฺจเยน ปจฺจโย. จกฺขุํ ฯเปฯ วตฺถุํ. นเหตู อเหตุเก ขนฺเธ อนิจฺจโต ฯเปฯ โทมนสฺสํ อุปฺปชฺชติ. กุสลากุสเล นิรุทฺเธ นเหตุ สเหตุโก วิปาโก ตทารมฺมณตา อุปฺปชฺชติ. ทิพฺเพน จกฺขุนา รูปํ ปสฺสติ. ทิพฺพาย โสตธาตุยา สทฺทํ สุณาติ. เจโตปริย\-ญาเณน นเหตุ\-อเหตุก\-จิตฺต\-สมงฺคิสฺส จิตฺตํ ชานาติ. นเหตู อเหตุกา ขนฺธา อิทฺธิวิธ\-ญาณสฺส, เจโตปริย\-ญาณสฺส, ปุพฺเพ\-นิวาสา\-นุสฺสติ\-ญาณสฺส, อนาคตํสญาณสฺส, อารมฺมณ\-ปจฺจเยน ปจฺจโย. นเหตู อเหตุเก ขนฺเธ อารพฺภ นเหตู สเหตุกา ขนฺธา อุปฺปชฺชนฺติ. (2)}

\proseitem{178}{สเหตุโก ธมฺโม นเหตุ\-สเหตุกสฺส ธมฺมสฺส อธิปติ\-ปจฺจเยน ปจฺจโย. อารมฺมณา\-ธิปติ, สหชาตา\-ธิปติ.  อารมฺมณา\-ธิปติ\csromandash{}ทานํ ทตฺวา, สีลํ สมาทิยิตฺวา, อุโปสถกมฺมํ กตฺวา ตํ ครุํ กตฺวา ปจฺจเวกฺขติ. ปุพฺเพ สุจิณฺณานิ ครุํ กตฺวา ปจฺจเวกฺขติ. ฌานํ ฯเปฯ. อริยา มคฺคา วุฏฺฐหิตฺวา มคฺคํ ครุํ กตฺวา ปจฺจเวกฺขนฺติ, ผลํ ครุํ กตฺวา ปจฺจเวกฺขนฺติ. นเหตู สเหตุเก ขนฺเธ ครุํ กตฺวา อสฺสาเทติ อภินนฺทติ, ตํ ครุํ กตฺวา ราโค อุปฺปชฺชติ, ทิฏฺฐิ อุปฺปชฺชติ. สหชาตา\-ธิปติ\csromandash{} \csromanfolio{76}นเหตุ\-สเหตุกา\-ธิปติ สมฺปยุตฺตกานํ ขนฺธานํ อธิปติ\-ปจฺจเยน ปจฺจโย. (1)}

\prose{สเหตุโก ธมฺโม นเหตุ\-อเหตุกสฺส ธมฺมสฺส อธิปติ\-ปจฺจเยน ปจฺจโย. สหชาตา\-ธิปติ\csromandash{}นเหตุ สเหตุกา\-ธิปติ จิตฺต\-สมุฏฺฐานานํ รูปานํ อธิปติ\-ปจฺจเยน ปจฺจโย. (2)}

\prose{สเหตุโก ธมฺโม นเหตุ\-สเหตุกสฺส จ นเหตุ\-อเหตุกสฺส จ ธมฺมสฺส อธิปติ\-ปจฺจเยน ปจฺจโย. สหชาตา\-ธิปติ\csromandash{}นเหตุ สเหตุกา\-ธิปติ สมฺปยุตฺตกานํ ขนฺธานํ จิตฺตสมุฏฺฐานานญฺจ รูปานํ อธิปติ\-ปจฺจเยน ปจฺจโย. (3)}

\prose{อเหตุโก ธมฺโม นเหตุ\-สเหตุกสฺส ธมฺมสฺส อธิปติ ปจฺจเยน ปจฺจโย. อารมฺมณา\-ธิปติ\csromandash{}อริยา นิพฺพานํ ครุํ กตฺวา ปจฺจเวกฺขนฺติ. นิพฺพานํ โคตฺรภุสฺส, โวทานสฺส, มคฺคสฺส, ผลสฺส อธิปติ\-ปจฺจเยน ปจฺจโย. จกฺขุํ ฯเปฯ วตฺถุํ. นเหตู อเหตุเก ขนฺเธ ครุํ กตฺวา อสฺสาเทติ อภินนฺทติ, ตํ ครุํ กตฺวา ราโค อุปฺปชฺชติ, ทิฏฺฐิ อุปฺปชฺชติ. (1)}

\proseitem{179}{สเหตุโก ธมฺโม นเหตุ\-สเหตุกสฺส ธมฺมสฺส อนนฺตร\-ปจฺจเยน ปจฺจโย. ปุริมา ปุริมา นเหตู สเหตุกา ขนฺธา ปจฺฉิมานํ ปจฺฉิมานํ นเหตุ\-สเหตุกานํ ขนฺธานํ อนนฺตร\-ปจฺจเยน ปจฺจโย. อนุโลมํ โคตฺรภุสฺส ฯเปฯ. เนว\-สญฺญา\-นา\-สญฺญา\-ยตนํ ผลสมาปตฺติยา อนนฺตร\-ปจฺจเยน ปจฺจโย. (1)}

\prose{สเหตุโก ธมฺโม นเหตุ\-อเหตุกสฺส ธมฺมสฺส อนนฺตร\-ปจฺจเยน ปจฺจโย. นเหตุ สเหตุกํ จุติจิตฺตํ นเหตุ\-อเหตุกสฺส อุปปตฺติ\-จิตฺตสฺส อนนฺตร\-ปจฺจเยน ปจฺจโย. นเหตุ สเหตุกํ ภวงฺคํ อาวชฺชนาย. นเหตุ สเหตุกํ ภวงฺคํ นเหตุ\-อเหตุกสฺส ภวงฺคสฺส. นเหตู สเหตุกา ขนฺธา นเหตุ\-อเหตุกสฺส วุฏฺฐานสฺส อนนฺตร\-ปจฺจเยน ปจฺจโย. (2)}

\prose{อเหตุโก ธมฺโม นเหตุ\-อเหตุกสฺส ธมฺมสฺส อนนฺตร\-ปจฺจเยน ปจฺจโย. ปุริมา ปุริมา นเหตู อเหตุกา ขนฺธา ปจฺฉิมานํ}

\prose{\csromanfolio{77}ปจฺฉิมานํ นเหตุ\-อเหตุกานํ ขนฺธานํ อนนฺตร\-ปจฺจเยน ปจฺจโย. อาวชฺชนา ปญฺจนฺนํ วิญฺญาณานํ อนนฺตร\-ปจฺจเยน ปจฺจโย. (1)}

\prose{อเหตุโก ธมฺโม นเหตุ\-สเหตุกสฺส ธมฺมสฺส อนนฺตร\-ปจฺจเยน ปจฺจโย. นเหตุ อเหตุกํ จุติจิตฺตํ นเหตุ\-สเหตุกสฺส อุปปตฺติ\-จิตฺตสฺส อนนฺตร\-ปจฺจเยน ปจฺจโย. อาวชฺชนา นเหตุ\-สเหตุกานํ ขนฺธานํ อนนฺตร\-ปจฺจเยน ปจฺจโย. นเหตู อเหตุกา ขนฺธา นเหตุ\-สเหตุกสฺส วุฏฺฐานสฺส อนนฺตร\-ปจฺจเยน ปจฺจโย. (2)}

\csromanheader{2}{\textbf{สมนนฺตราทิ}}
\proseitem{180}{สเหตุโก ธมฺโม นเหตุ\-สเหตุกสฺส ธมฺมสฺส สมนนฺตร\-ปจฺจเยน ปจฺจโย. สหชาต\-ปจฺจเยน ปจฺจโย. (อิห ฆฏนา นตฺถิ. สตฺต ปญฺหา.) อญฺญมญฺญ\-ปจฺจเยน ปจฺจโย. ฉ ปญฺหา. นิสฺสย\-ปจฺจเยน ปจฺจโย. (ปวตฺติ\-ปฏิสนฺธิ สตฺต ปญฺหา. อิห ฆฏนา นตฺถิ.)}

\csromanheader{2}{\textbf{อุปนิสฺสย}}
\proseitem{181}{สเหตุโก ธมฺโม นเหตุ\-สเหตุกสฺส ธมฺมสฺส อุปนิสฺสย\-ปจฺจเยน ปจฺจโย. อารมฺมณู\-ปนิสฺสโย, อนนฺตรู\-ปนิสฺสโย, ปกตู\-ปนิสฺสโย ฯเปฯ. ปกตู\-ปนิสฺสโย\csromandash{}สทฺธํ อุปนิสฺสาย ทานํ เทติ ฯเปฯ มานํ ชปฺเปติ. ทิฏฺฐิํ คณฺหาติ. สีลํ ฯเปฯ. ปตฺถนํ อุปนิสฺสาย ทานํ เทติ ฯเปฯ สํฆํ ภินฺทติ. สทฺธา ฯเปฯ. ปตฺถนา สทฺธาย ฯเปฯ ปตฺถนาย มคฺคสฺส ผลสมาปตฺติยา อุปนิสฺสย\-ปจฺจเยน ปจฺจโย. (1)}

\prose{สเหตุโก ธมฺโม นเหตุ\-อเหตุกสฺส ธมฺมสฺส อุปนิสฺสย\-ปจฺจเยน ปจฺจโย. อนนฺตรู\-ปนิสฺสโย, ปกตู\-ปนิสฺสโย ฯเปฯ. ปกตู\-ปนิสฺสโย\csromandash{}สทฺธา กายิกสฺส สุขสฺส, กายิกสฺส ทุกฺขสฺส อุปนิสฺสย\-ปจฺจเยน ปจฺจโย. สีลํ ฯเปฯ. ปตฺถนา กายิกสฺส สุขสฺส, กายิกสฺส ทุกฺขสฺส อุปนิสฺสย\-ปจฺจเยน ปจฺจโย. สทฺธา ฯเปฯ. ปตฺถนา กายิกสฺส สุขสฺส, กายิกสฺส ทุกฺขสฺส อุปนิสฺสย\-ปจฺจเยน ปจฺจโย. (2)}

\proseitem{182}{อเหตุโก ธมฺโม นเหตุ\-อเหตุกสฺส ธมฺมสฺส อุปนิสฺสย\-ปจฺจเยน ปจฺจโย. อนนฺตรู\-ปนิสฺสโย, ปกตู\-ปนิสฺสโย ฯเปฯ. ปกตู\-ปนิสฺสโยกายิกํ สุขํ กายิกสฺส สุขสฺส, กายิกสฺส \csromanfolio{78}ทุกฺขสฺส อุปนิสฺสย\-ปจฺจเยน ปจฺจโย. กายิกํ ทุกฺขํ. อุตุ. โภชนํ. เสนาสนํ กายิกสฺส สุขสฺส, กายิกสฺส ทุกฺขสฺส อุปนิสฺสย\-ปจฺจเยน ปจฺจโย. กายิกํ สุขํ. กายิกํ ทุกฺขํ. อุตุ. โภชนํ. เสนาสนํ กายิกสฺส สุขสฺส, กายิกสฺส ทุกฺขสฺส อุปนิสฺสย\-ปจฺจเยน ปจฺจโย. (1)}

\prose{อเหตุโก ธมฺโม นเหตุ\-สเหตุกสฺส ธมฺมสฺส อุปนิสฺสย\-ปจฺจเยน ปจฺจโย. อารมฺมณู\-ปนิสฺสโย, อนนฺตรู\-ปนิสฺสโย, ปกตู\-ปนิสฺสโย ฯเปฯ. ปกตู\-ปนิสฺสโย\csromandash{}กายิกํ สุขํ อุปนิสฺสาย ทานํ เทติ ฯเปฯ สํฆํ ภินฺทติ. กายิกํ ทุกฺขํ. อุตุํ. โภชนํ. เสนาสนํ อุปนิสฺสาย ทานํ เทติ ฯเปฯ สํฆํ ภินฺทติ. กายิกํ สุขํ ฯเปฯ. เสนาสนํ สทฺธาย ฯเปฯ ปตฺถนาย มคฺคสฺส ผลสมาปตฺติยา อุปนิสฺสย\-ปจฺจเยน ปจฺจโย. (2)}

\proseitem{183}{อเหตุโก ธมฺโม นเหตุ\-อเหตุกสฺส ธมฺมสฺส ปุเรชาต\-ปจฺจเยน ปจฺจโย. อารมฺมณ\-ปุเรชาตํ, วตฺถุ\-ปุเรชาตํ.  อารมฺมณ\-ปุเรชาตํ\csromandash{}จกฺขุํ ฯเปฯ วตฺถุํ อนิจฺจโต ฯเปฯ โทมนสฺสํ อุปฺปชฺชติ. กุสลากุสเล นิรุทฺเธ นเหตุ อเหตุโก วิปาโก ตทารมฺมณตา อุปฺปชฺชติ. รูปายตนํ จกฺขุ\-วิญฺญาณสฺส ปุเรชาต\-ปจฺจเยน ปจฺจโย ฯเปฯ. โผฏฺฐพฺพา\-ยตนํ กายวิญฺญาณสฺส ปุเรชาต\-ปจฺจเยน ปจฺจโย. วตฺถุ\-ปุเรชาตํ\csromandash{}จกฺขายตนํ จกฺขุ\-วิญฺญาณสฺส ฯเปฯ กายายตนํ กายวิญฺญาณสฺส. วตฺถุ นเหตุ\-อเหตุกานํ ขนฺธานํ ปุเรชาต\-ปจฺจเยน ปจฺจโย. (1)}

\prose{อเหตุโก ธมฺโม นเหตุ\-สเหตุกสฺส ธมฺมสฺส ปุเรชาต\-ปจฺจเยน ปจฺจโย. อารมฺมณ\-ปุเรชาตํ, วตฺถุ\-ปุเรชาตํ.  อารมฺมณ\-ปุเรชาตํ\csromandash{}จกฺขุํ ฯเปฯ วตฺถุํ อนิจฺจโต ฯเปฯ โทมนสฺสํ อุปฺปชฺชติ. กุสลากุสเล นิรุทฺเธ นเหตุ สเหตุโก วิปาโก ตทารมฺมณตา อุปฺปชฺชติ. ทิพฺเพน จกฺขุนา รูปํ ปสฺสติ. ทิพฺพาย โสตธาตุยา สทฺทํ สุณาติ. วตฺถุ\-ปุเรชาตํ\csromandash{}วตฺถุ นเหตุ\-สเหตุกานํ ขนฺธานํ ปุเรชาต\-ปจฺจเยน ปจฺจโย. (2)}

\csromanheader{2}{6. นเหตุสเหตุทุก \textbf{ปจฺฉาชาต}}
\proseitem{184}{\csromanfolio{79}สเหตุโก ธมฺโม นเหตุ\-อเหตุกสฺส ธมฺมสฺส ปจฺฉาชาต\-ปจฺจเยน ปจฺจโย. ปจฺฉาชาตา นเหตู สเหตุกา ขนฺธา ปุเรชาตสฺส อิมสฺส กายสฺส ปจฺฉาชาต\-ปจฺจเยน ปจฺจโย. (1)}

\prose{อเหตุโก ธมฺโม นเหตุ\-อเหตุกสฺส ธมฺมสฺส ปจฺฉาชาต\-ปจฺจเยน ปจฺจโย. ปจฺฉาชาตา นเหตู อเหตุกา ขนฺธา ปุเรชาตสฺส อิมสฺส กายสฺส ปจฺฉาชาต\-ปจฺจเยน ปจฺจโย. (1)}

\csromanheader{2}{\textbf{อาเสวน}}
\proseitem{185}{สเหตุโก ธมฺโม นเหตุ\-สเหตุกสฺส ธมฺมสฺส อาเสวน\-ปจฺจเยน ปจฺจโย. ปุริมา ปุริมา นเหตู สเหตุกา ขนฺธา ปจฺฉิมานํ ปจฺฉิมานํ นเหตุ\-สเหตุกานํ ขนฺธานํ อาเสวน\-ปจฺจเยน ปจฺจโย. อนุโลมํ โคตฺรภุสฺส. อนุโลมํ โวทานสฺส. โคตฺรภุ มคฺคสฺส. โวทานํ มคฺคสฺส อาเสวน\-ปจฺจเยน ปจฺจโย. (1)}

\prose{อเหตุโก ธมฺโม นเหตุ\-อเหตุกสฺส ธมฺมสฺส อาเสวน\-ปจฺจเยน ปจฺจโย. ปุริมา ปุริมา นเหตู อเหตุกา ขนฺธา ปจฺฉิมานํ ปจฺฉิมานํ นเหตุ\-อเหตุกานํ ขนฺธานํ อาเสวน\-ปจฺจเยน ปจฺจโย. (1)}

\proseitem{186}{สเหตุโก ธมฺโม นเหตุ\-สเหตุกสฺส ธมฺมสฺส กมฺมปจฺจเยน ปจฺจโย. สหชาตา, นานากฺขณิกา.  สหชาตา นเหตุ สเหตุกา เจตนา สมฺปยุตฺตกานํ ขนฺธานํ กมฺมปจฺจเยน ปจฺจโย. ปฏิสนฺธิ\-กฺขเณ ฯเปฯ. นานากฺขณิกา นเหตุ สเหตุกา เจตนา วิปากานํ นเหตุ\-สเหตุกานํ ขนฺธานํ กมฺมปจฺจเยน ปจฺจโย. (1)}

\prose{สเหตุโก ธมฺโม นเหตุ\-อเหตุกสฺส ธมฺมสฺส กมฺมปจฺจเยน ปจฺจโย. สหชาตา, นานากฺขณิกา.  สหชาตา นเหตุ สเหตุกา เจตนา จิตฺต\-สมุฏฺฐานานํ รูปานํ กมฺมปจฺจเยน ปจฺจโย. ปฏิสนฺธิ\-กฺขเณ ฯเปฯ. นานากฺขณิกา นเหตุ สเหตุกา เจตนา วิปากานํ นเหตุ\-อเหตุกานํ ขนฺธานํ กฏตฺตา จ รูปานํ กมฺมปจฺจเยน ปจฺจโย. (1)}

\prose{\csromanfolio{80}สเหตุโก ธมฺโม นเหตุ\-สเหตุกสฺส จ นเหตุ\-อเหตุกสฺส จ ธมฺมสฺส กมฺมปจฺจเยน ปจฺจโย. สหชาตา, นานากฺขณิกา.  สหชาตา นเหตุ สเหตุกา เจตนา สมฺปยุตฺตกานํ ขนฺธานํ จิตฺตสมุฏฺฐานานญฺจ รูปานํ กมฺมปจฺจเยน ปจฺจโย. ปฏิสนฺธิ\-กฺขเณ ฯเปฯ. นานากฺขณิกา นเหตุ สเหตุกา เจตนา วิปากานํ นเหตุ\-สเหตุกานํ ขนฺธานํ กฏตฺตา จ รูปานํ กมฺมปจฺจเยน ปจฺจโย. (1)}

\prose{อเหตุโก ธมฺโม นเหตุ\-อเหตุกสฺส ธมฺมสฺส กมฺมปจฺจเยน ปจฺจโย. สหชาตา นเหตุ อเหตุกา เจตนา สมฺปยุตฺตกานํ ขนฺธานํ จิตฺตสมุฏฺฐานานญฺจ รูปานํ กมฺมปจฺจเยน ปจฺจโย. ปฏิสนฺธิ\-กฺขเณ นเหตุ อเหตุกา เจตนา สมฺปยุตฺตกานํ ขนฺธานํ กฏตฺตา จ รูปานํ กมฺมปจฺจเยน ปจฺจโย. (1)}

\csromanheader{2}{\textbf{วิปาก}}
\proseitem{187}{สเหตุโก ธมฺโม นเหตุ\-สเหตุกสฺส ธมฺมสฺส วิปาก\-ปจฺจเยน ปจฺจโย. ตีณิ.}

\prose{อเหตุโก ธมฺโม นเหตุ\-อเหตุกสฺส ธมฺมสฺส วิปาก\-ปจฺจเยน ปจฺจโย. เอกํ.}

\csromanheader{2}{\textbf{อาหาร}}
\proseitem{188}{สเหตุโก ธมฺโม นเหตุ\-สเหตุกสฺส ธมฺมสฺส อาหารปจฺจเยน ปจฺจโย. ตีณิ.}

\prose{อเหตุโก ธมฺโม นเหตุ\-อเหตุกสฺส ธมฺมสฺส อาหารปจฺจเยน ปจฺจโย. นเหตู อเหตุกา อาหารา สมฺปยุตฺตกานํ ขนฺธานํ จิตฺตสมุฏฺฐานานญฺจ รูปานํ อาหารปจฺจเยน ปจฺจโย. ปฏิสนฺธิ\-กฺขเณ ฯเปฯ. กพฬีกาโร อาหาโร อิมสฺส กายสฺส อาหารปจฺจเยน ปจฺจโย. (1)}

\csromanheader{2}{\textbf{อินฺทฺริย}}
\vspace{\tipitakaparskip}
\csromancenter{สเหตุโก ธมฺโม นเหตุ\-สเหตุกสฺส ธมฺมสฺส อินฺทฺริย\-ปจฺจเยน ปจฺจโย. ตีณิ.}
\prose{\csromanfolio{81}อเหตุโก ธมฺโม นเหตุ\-อเหตุกสฺส ธมฺมสฺส อินฺทฺริย\-ปจฺจเยน ปจฺจโย. นเหตู อเหตุกา อินฺทฺริยา สมฺปยุตฺตกานํ ขนฺธานํ จิตฺตสมุฏฺฐานานญฺจ รูปานํ อินฺทฺริย\-ปจฺจเยน ปจฺจโย. ปฏิสนฺธิ\-กฺขเณ ฯเปฯ. รูปชีวิตินฺทฺริยํ กฏตฺตารูปานํ อินฺทฺริย\-ปจฺจเยน ปจฺจโย.}

\csromanheader{2}{\textbf{ฌานาทิ}}
\proseitem{190}{สเหตุโก ธมฺโม นเหตุ\-สเหตุกสฺส ธมฺมสฺส ฌานปจฺจเยน ปจฺจโย ฯเปฯ. (จตฺตาริปิ กาตพฺพานิ.) มคฺคปจฺจเยน ปจฺจโย. ตีณิ.}

\csromanheader{2}{\textbf{สมฺปยุตฺต}}
\proseitem{191}{สเหตุโก ธมฺโม นเหตุ\-สเหตุกสฺส ธมฺมสฺส สมฺปยุตฺต\-ปจฺจเยน ปจฺจโย. นเหตุ สเหตุโก เอโก ขนฺโธ ติณฺณนฺนํ ขนฺธานํ ฯเปฯ. ปฏิสนฺธิ\-กฺขเณ ฯเปฯ. (1)}

\prose{อเหตุโก ธมฺโม นเหตุ\-อเหตุกสฺส ธมฺมสฺส สมฺปยุตฺต\-ปจฺจเยน ปจฺจโย. นเหตุ อเหตุโก เอโก ขนฺโธ ติณฺณนฺนํ ขนฺธานํ ฯเปฯ. ปฏิสนฺธิ\-กฺขเณ ฯเปฯ. (2)}

\proseitem{192}{สเหตุโก ธมฺโม นเหตุ\-อเหตุกสฺส ธมฺมสฺส วิปฺปยุตฺต\-ปจฺจเยน ปจฺจโย. สหชาตํ, ปจฺฉาชาตํ.  สหชาตา นเหตู สเหตุกา ขนฺธา จิตฺต\-สมุฏฺฐานานํ รูปานํ วิปฺปยุตฺต\-ปจฺจเยน ปจฺจโย. ปฏิสนฺธิ\-กฺขเณ ฯเปฯ. ปจฺฉาชาตา นเหตู สเหตุกา ขนฺธา ปุเรชาตสฺส อิมสฺส กายสฺส วิปฺปยุตฺต\-ปจฺจเยน ปจฺจโย. (1)}

\prose{อเหตุโก ธมฺโม นเหตุ\-อเหตุกสฺส ธมฺมสฺส วิปฺปยุตฺต\-ปจฺจเยน ปจฺจโย. สหชาตํ, ปุเรชาตํ, ปจฺฉาชาตํ.  สหชาตา นเหตู อเหตุกา ขนฺธา จิตฺต\-สมุฏฺฐานานํ รูปานํ วิปฺปยุตฺต\-ปจฺจเยน ปจฺจโย. ปฏิสนฺธิ\-กฺขเณ ฯเปฯ. ขนฺธา วตฺถุสฺส วิปฺปยุตฺต\-ปจฺจเยน ปจฺจโย. วตฺถุ ขนฺธานํ วิปฺปยุตฺต\-ปจฺจเยน ปจฺจโย. ปุเรชาตํ จกฺขายตนํ จกฺขุ\-วิญฺญาณสฺส ฯเปฯ กายายตนํ กายวิญฺญาณสฺส ฯเปฯ วตฺถุ นเหตุ\-สเหตุกานํ ขนฺธานํ วิปฺปยุตฺต\-ปจฺจเยน ปจฺจโย. ปจฺฉาชาตา นเหตู \csromanfolio{82}อเหตุกา ขนฺธา ปุเรชาตสฺส อิมสฺส กายสฺส วิปฺปยุตฺต\-ปจฺจเยน ปจฺจโย. (1)}

\prose{อเหตุโก ธมฺโม นเหตุ\-สเหตุกสฺส ธมฺมสฺส วิปฺปยุตฺต\-ปจฺจเยน ปจฺจโย. สหชาตํ, ปุเรชาตํ.  สหชาตํ ปฏิสนฺธิ\-กฺขเณ วตฺถุ นเหตุ\-สเหตุกานํ ขนฺธานํ วิปฺปยุตฺต\-ปจฺจเยน ปจฺจโย. ปุเรชาตํ วตฺถุ นเหตุ\-สเหตุกานํ ขนฺธานํ วิปฺปยุตฺต\-ปจฺจเยน ปจฺจโย. (2)}

\proseitem{193}{สเหตุโก ธมฺโม นเหตุ\-สเหตุกสฺส ธมฺมสฺส อตฺถิปจฺจเยน ปจฺจโย. นเหตุ สเหตุโก เอโก ขนฺโธ ติณฺณนฺนํ ขนฺธานํ ฯเปฯ. ปฏิสนฺธิ\-กฺขเณ ฯเปฯ. (1)}

\prose{สเหตุโก ธมฺโม นเหตุ\-อเหตุกสฺส ธมฺมสฺส อตฺถิปจฺจเยน ปจฺจโย. สหชาตํ, ปจฺฉาชาตํ ฯเปฯ. (2)}

\prose{สเหตุโก ธมฺโม นเหตุ\-สเหตุกสฺส จ นเหตุ\-อเหตุกสฺส จ ธมฺมสฺส อตฺถิปจฺจเยน ปจฺจโย. นเหตุ สเหตุโก เอโก ขนฺโธ ติณฺณนฺนํ ขนฺธานํ จิตฺตสมุฏฺฐานานญฺจ รูปานํ อตฺถิปจฺจเยน ปจฺจโย ฯเปฯ. ปฏิสนฺธิ\-กฺขเณ ฯเปฯ. (3)}

\prose{อเหตุโก ธมฺโม นเหตุ\-อเหตุกสฺส ธมฺมสฺส อตฺถิปจฺจเยน ปจฺจโย. สหชาตํ, ปุเรชาตํ, ปจฺฉาชาตํ, อาหารํ, อินฺทฺริยํ.  สหชาโต นเหตุ อเหตุโก เอโก ขนฺโธ ติณฺณนฺนํ ขนฺธานํ จิตฺตสมุฏฺฐานานญฺจ รูปานํ อตฺถิปจฺจเยน ปจฺจโย ฯเปฯ. (ยาว อสญฺญสตฺตา.) ปุเรชาตํ จกฺขุํ ฯเปฯ วตฺถุํ อนิจฺจโต ฯเปฯ โทมนสฺสํ อุปฺปชฺชติ. กุสลากุสเล นิรุทฺเธ นเหตุ อเหตุโก วิปาโก ตทารมฺมณตา อุปฺปชฺชติ. รูปายตนํ จกฺขุ\-วิญฺญาณสฺส ฯเปฯ. โผฏฺฐพฺพา\-ยตนํ กายวิญฺญาณสฺส ฯเปฯ. จกฺขายตนํ ฯเปฯ. กายายตนํ ฯเปฯ วตฺถุ นเหตุ\-อเหตุกานํ ขนฺธานํ อตฺถิปจฺจเยน ปจฺจโย. ปจฺฉาชาตา นเหตู อเหตุกา ขนฺธา ปุเรชาตสฺส ฯเปฯ. กพฬีกาโร อาหาโร อิมสฺส กายสฺส ฯเปฯ. รูปชีวิตินฺทฺริยํ กฏตฺตารูปานํ อตฺถิปจฺจเยน ปจฺจโย. (1)}

\prose{อเหตุโก ธมฺโม นเหตุ\-สเหตุกสฺส ธมฺมสฺส อตฺถิปจฺจเยน ปจฺจโย. สหชาตํ, ปุเรชาตํ.  สหชาตํ ปฏิสนฺธิ\-กฺขเณ \csromanfolio{83}วตฺถุ นเหตุ\-สเหตุกานํ ขนฺธานํ อตฺถิปจฺจเยน ปจฺจโย. ปุเรชาตํ จกฺขุํ ฯเปฯ วตฺถุํ อนิจฺจโต ฯเปฯ โทมนสฺสํ อุปฺปชฺชติ. กุสลากุสเล นิรุทฺเธ นเหตุ สเหตุโก วิปาโก ตทารมฺมณตา อุปฺปชฺชติ. ทิพฺเพน จกฺขุนา รูปํ ปสฺสติ. ทิพฺพาย โสตธาตุยา สทฺทํ สุณาติ. วตฺถุ\-ปุเรชาตํ วตฺถุ นเหตุ\-สเหตุกานํ ขนฺธานํ อตฺถิปจฺจเยน ปจฺจโย. (2)}

\proseitem{194}{สเหตุโก จ นเหตุ อเหตุโก จ ธมฺมา นเหตุ\-สเหตุกสฺส ธมฺมสฺส อตฺถิปจฺจเยน ปจฺจโย. สหชาตํ, ปุเรชาตํ.  สหชาโต นเหตุ สเหตุโก เอโก ขนฺโธ จ วตฺถุ จ ติณฺณนฺนํ ขนฺธานํ ฯเปฯ. ปฏิสนฺธิ\-กฺขเณ ฯเปฯ. นเหตุ สเหตุโก เอโก ขนฺโธ จ วตฺถุ จ ติณฺณนฺนํ ขนฺธานํ ฯเปฯ. (1)}

\prose{สเหตุโก จ นเหตุ อเหตุโก จ ธมฺมา นเหตุ\-อเหตุกสฺส ธมฺมสฺส อตฺถิปจฺจเยน ปจฺจโย. สหชาตํ, ปจฺฉาชาตํ, อาหารํ, อินฺทฺริยํ.  สหชาตา นเหตู สเหตุกา ขนฺธา จ มหาภูตา จ จิตฺต\-สมุฏฺฐานานํ รูปานํ อตฺถิปจฺจเยน ปจฺจโย. ปฏิสนฺธิ\-กฺขเณ ฯเปฯ. ปจฺฉาชาตา นเหตู สเหตุกา ขนฺธา จ กพฬีกาโร อาหาโร จ อิมสฺส กายสฺส อตฺถิปจฺจเยน ปจฺจโย. ปจฺฉาชาตา นเหตู สเหตุกา ขนฺธา จ รูปชีวิตินฺทฺริยญฺจ กฏตฺตารูปานํ อตฺถิปจฺจเยน ปจฺจโย ฯเปฯ. (2)}

\csromanmarkh{1. ปจฺจยานุโลม}\csromanmarkhii{2. สงฺขฺยาวาร}\csromanheadernomark{2}{1. \textbf{ปจฺจยานุโลม 2. สงฺขฺยาวาร}}
\proseitem{195}{อารมฺมเณ จตฺตาริ, อธิปติยา จตฺตาริ, อนนฺตเร จตฺตาริ, สมนนฺตเร จตฺตาริ, สหชาเต สตฺต, อญฺญมญฺเญ ฉ, นิสฺสเย สตฺต, อุปนิสฺสเย จตฺตาริ, ปุเรชาเต เทฺว, ปจฺฉาชาเต เทฺว, อาเสวเน เทฺว, กมฺเม จตฺตาริ, วิปาเก จตฺตาริ, อาหาเร จตฺตาริ, อินฺทฺริเย จตฺตาริ, ฌาเน จตฺตาริ, มคฺเค ตีณิ, สมฺปยุตฺเต เทฺว, วิปฺปยุตฺเต ตีณิ, อตฺถิยา สตฺต, นตฺถิยา จตฺตาริ, วิคเต จตฺตาริ, อวิคเต สตฺต. (เอวํ คเณตพฺพํ)}

\vspace{\tipitakaparskip}
\csromancenter{อนุโลมํ.}
\proseitem{196}{\csromanfolio{84}สเหตุโก ธมฺโม นเหตุ\-สเหตุกสฺส ธมฺมสฺส อารมฺมณ\-ปจฺจเยน ปจฺจโย. สหชาต\-ปจฺจเยน ปจฺจโย. อุปนิสฺสย\-ปจฺจเยน ปจฺจโย. กมฺมปจฺจเยน ปจฺจโย. (1)}

\prose{สเหตุโก ธมฺโม นเหตุ\-อเหตุกสฺส ธมฺมสฺส อารมฺมณ\-ปจฺจเยน ปจฺจโย. สหชาต\-ปจฺจเยน ปจฺจโย. อุปนิสฺสย\-ปจฺจเยน ปจฺจโย. ปจฺฉาชาต\-ปจฺจเยน ปจฺจโย. กมฺมปจฺจเยน ปจฺจโย. (2)}

\prose{สเหตุโก ธมฺโม นเหตุ\-สเหตุกสฺส จ นเหตุ\-อเหตุกสฺส จ ธมฺมสฺส สหชาต\-ปจฺจเยน ปจฺจโย. กมฺมปจฺจเยน ปจฺจโย. (3)}

\proseitem{197}{อเหตุโก ธมฺโม นเหตุ\-อเหตุกสฺส ธมฺมสฺส อารมฺมณ\-ปจฺจเยน ปจฺจโย. สหชาต\-ปจฺจเยน ปจฺจโย. อุปนิสฺสย\-ปจฺจเยน ปจฺจโย. ปุเรชาต\-ปจฺจเยน ปจฺจโย. ปจฺฉาชาต\-ปจฺจเยน ปจฺจโย. อาหารปจฺจเยน ปจฺจโย. อินฺทฺริย\-ปจฺจเยน ปจฺจโย. (1)}

\prose{อเหตุโก ธมฺโม นเหตุ\-สเหตุกสฺส ธมฺมสฺส อารมฺมณ\-ปจฺจเยน ปจฺจโย. สหชาต\-ปจฺจเยน ปจฺจโย. อุปนิสฺสย\-ปจฺจเยน ปจฺจโย. ปุเรชาต\-ปจฺจเยน ปจฺจโย. (2)}

\proseitem{198}{สเหตุโก จ นเหตุ อเหตุโก จ ธมฺมา นเหตุ\-สเหตุกสฺส ธมฺมสฺส สหชาตํ. ปุเรชาตํ. (1)}

\prose{สเหตุโก จ นเหตุ อเหตุโก จ ธมฺมา นเหตุ\-อเหตุกสฺส ธมฺมสฺส สหชาตํ. ปจฺฉาชาตํ. อาหารํ. อินฺทฺริยํ. (2)}

\csromanmarkh{2. ปจฺจย\-ปจฺจนีย}\csromanmarkhii{2. สงฺขฺยาวาร}\csromanheadernomark{2}{2. ปจฺจย\-ปจฺจนีย 2. สงฺขฺยาวาร}
\proseitem{199}{นเหตุยา สตฺต, นอารมฺมเณ สตฺต. (สํขิตฺตํ, สพฺพตฺถ สตฺต, นสหชาเต ฉ, นอญฺญมญฺเญ ฉ, นนิสฺสเย ฉ, (สพฺพตฺถ สตฺต.) \csromanfolio{85}นสมฺปยุตฺเต ฉ, นวิปฺปยุตฺเต ปญฺจ, โนอตฺถิยา ปญฺจ, โนนตฺถิยา สตฺต, โนวิคเต สตฺต, โนอวิคเต ปญฺจ. (เอวํ คเณตพฺพํ.)}

\vspace{\tipitakaparskip}
\csromancenter{ปจฺจนียํ.}
\csromanheader{2}{\textbf{อารมฺมณทุก}}
\proseitem{200}{อารมฺมณ\-ปจฺจยา นเหตุยา จตฺตาริ, นอธิปติยา จตฺตาริ, นอนนฺตเร จตฺตาริ, (สพฺพตฺถ จตฺตาริ.) โนนตฺถิยา จตฺตาริ, โนวิคเต จตฺตาริ, โนอวิคเต จตฺตาริ. (เอวํ คเณตพฺพํ.)}

\vspace{\tipitakaparskip}
\csromancenter{อนุโลม\-ปจฺจนียํ.}
\csromancenter{นเหตุปจฺจยา อารมฺมเณ จตฺตาริ, อธิปติยา จตฺตาริ ฯเปฯ อวิคเต สตฺต.}
\csromancenter{ปจฺจนียา\-นุโลมํ.}
\nitthitam{นเหตุ\-สเหตุกทุกํ นิฏฺฐิตํ.}
\nitthitam{เหตุโคจฺฉกํ นิฏฺฐิตํ.}
\csromanmatikaanchor{mtk.9}
\csromantocmarknopage{chapter}{2. จูฬนฺตรทุก}{mtk.9}
\bookmark[dest={mtk.9},level=0]{2. จูฬนฺตรทุก}
\chapterheadpageread{2. \textbf{จูฬนฺตรทุก}}
\csromanmatikaanchor{mtk.10}
\csromantocmarkref{section}{7. สปฺปจฺจยทุก}{mtk.10}
\bookmark[dest={mtk.10},level=1]{7. สปฺปจฺจยทุก}
\csromanmarkh{7. สปฺปจฺจยทุก}\csromanmarkhii{1. ปฏิจฺจวาร}\csromanheadernomark{1}{7. \textbf{สปฺปจฺจยทุก 1. ปฏิจฺจวาร}}
\proseitem{1}{\csromanfolio{86}สปฺปจฺจยํ ธมฺมํ ปฏิจฺจ สปฺปจฺจโย ธมฺโม อุปฺปชฺชติ เหตุปจฺจยา. สปฺปจฺจยํ เอกํ ขนฺธํ ปฏิจฺจ ตโย ขนฺธา จิตฺต\-สมุฏฺฐานญฺจ รูปํ ฯเปฯ เทฺว ขนฺเธ ฯเปฯ. ปฏิสนฺธิ\-กฺขเณ ฯเปฯ. ขนฺเธ ปฏิจฺจ วตฺถุ, วตฺถุํ ปฏิจฺจ ขนฺธา. เอกํ มหาภูตํ ฯเปฯ มหาภูเต ปฏิจฺจ จิตฺต\-สมุฏฺฐานํ รูปํ, กฏตฺตารูปํ, อุปาทารูปํ.}

\prose{สปฺปจฺจยํ ธมฺมํ ปฏิจฺจ สปฺปจฺจโย ธมฺโม อุปฺปชฺชติ อารมฺมณ\-ปจฺจยา ฯเปฯ อวิคตปจฺจยา.}

\proseitem{2}{เหตุยา เอกํ, อารมฺมเณ เอกํ ฯเปฯ อวิคเต เอกํ.}

\vspace{\tipitakaparskip}
\csromancenter{อนุโลมํ.}
\proseitem{3}{สปฺปจฺจยํ ธมฺมํ ปฏิจฺจ สปฺปจฺจโย ธมฺโม อุปฺปชฺชติ นเหตุปจฺจยา. อเหตุกํ สปฺปจฺจยํ เอกํ ขนฺธํ ปฏิจฺจ ตโย ขนฺธา จิตฺต\-สมุฏฺฐานญฺจ รูปํ. เทฺว ขนฺเธ ฯเปฯ. อเหตุก\-ปฏิสนฺธิ\-กฺขเณ ฯเปฯ. (ยาว อสญฺญสตฺตา.) วิจิกิจฺฉาสเหคเต อุทฺธจฺจ\-สหคเต ขนฺเธ ปฏิจฺจ วิจิกิจฺฉา\-สหคโต อุจฺฉจฺจสหคโต โมโห. (สํขิตฺตํ.)}

\proseitem{4}{นเหตุยา เอกํ, นอารมฺมเณ เอกํ, นอธิปติยา เอกํ ฯเปฯ โนวิคเต เอกํ. ปจฺจนียํ.}

\proseitem{5}{เหตุปจฺจยา นอารมฺมเณ เอกํ, นอธิปติยา เอกํ ฯเปฯ โนวิคเต เอกํ. อนุโลม\-ปจฺจนียํ.}

\proseitem{6}{นเหตุปจฺจยา อารมฺมเณ เอกํ, อนนฺตเร เอกํ ฯเปฯ อวิคเต เอกํ.}

\vspace{\tipitakaparskip}
\csromancenter{ปจฺจนียา\-นุโลมํ.}
\csromancenter{(สหชาตวาโร ปฏิจฺจ\-วาร\-สทิโส.)}
\csromanmarkh{7. สปฺปจฺจยทุก}\csromanmarkhii{7. สปฺปจฺจยทุก 3. ปจฺจยวาร}\csromanheadernomark{2}{7. สปฺปจฺจยทุก \textbf{7. สปฺปจฺจยทุก 3. ปจฺจยวาร}}
\csromanheader{2}{1. \textbf{ปจฺจยานุโลม 1.วิภงฺควาร}}
\proseitem{7}{\csromanfolio{87}สปฺปจฺจยํ ธมฺมํ ปจฺจยา สปฺปจฺจโย ธมฺโม อุปฺปชฺชติ เหตุปจฺจยา. สปฺปจฺจยํ เอกํ ขนฺธํ ปจฺจยา ตโย ขนฺธา จิตฺต\-สมุฏฺฐานญฺจ รูปํ ฯเปฯ. ปฏิสนฺธิ\-กฺขเณ ฯเปฯ. ขนฺเธ ปจฺจยา วตฺถุ, วตฺถุํ ปจฺจยา ขนฺธา. เอกํ มหาภูตํ ฯเปฯ. วตฺถุํ ปจฺจยา สปฺปจฺจยา ขนฺธา.}

\prose{สปฺปจฺจยํ ธมฺมํ ปจฺจยา สปฺปจฺจโย ธมฺโม อุปฺปชฺชติ อารมฺมณ\-ปจฺจยา. (สํขิตฺตํ.) (เอวํ ปจฺจยวาโรปิ นิสฺสยวาโรปิ สํสฏฺฐ\-วาโรปิ สมฺปยุตฺตวาโรปิ วิตฺถาเรตพฺโพ, สพฺพตฺถ เอกาเยว ปญฺหา.)}

\csromanmarkh{7. สปฺปจฺจยทุก}\csromanmarkhii{7. ปญฺหาวาร}\csromanheadernomark{2}{7. สปฺปจฺจยทุก 7. ปญฺหาวาร}
\csromanmarkh{1. ปจฺจยานุโลม}\csromanmarkhii{1. วิภงฺควาร}\csromanheadernomark{2}{1. \textbf{ปจฺจยานุโลม 1. วิภงฺควาร}}
\csromanheader{2}{\textbf{เหตุ}}
\proseitem{8}{สปฺปจฺจโย ธมฺโม สปฺปจฺจยสฺส ธมฺมสฺส เหตุปจฺจเยน ปจฺจโย. สปฺปจฺจยา เหตู สมฺปยุตฺตกานํ ขนฺธานํ จิตฺตสมุฏฺฐานานญฺจ รูปานํ เหตุปจฺจเยน ปจฺจโย. ปฏิสนฺธิ\-กฺขเณ ฯเปฯ.}

\proseitem{9}{สปฺปจฺจโย ธมฺโม สปฺปจฺจยสฺส ธมฺมสฺส อารมฺมณ\-ปจฺจเยน ปจฺจโย. ทานํ ทตฺวา, สีลํ สมาทิยิตฺวา, อุโปสถกมฺมํ กตฺวา ตํ ปจฺจเวกฺขติ. ปุพฺเพ สุจิณฺณานิ ฯเปฯ. ฌานา วุฏฺฐหิตฺวา ฯเปฯ. อริยา มคฺคา วุฏฺฐหิตฺวา มคฺคํ ปจฺจเวกฺขนฺติ, ผลํ ปจฺจเวกฺขนฺติ. ปหีเน กิเลเส ฯเปฯ. วิกฺขมฺภิเต กิเลเส ฯเปฯ. ปุพฺเพ สมุทาจิณฺเณ กิเลเส ชานนฺติ. จกฺขุํ ฯเปฯ วตฺถุํ. สปฺปจฺจเย ขนฺเธ อนิจฺจโต ฯเปฯ โทมนสฺสํ อุปฺปชฺชติ. ทิพฺเพน จกฺขุนา รูปํ ปสฺสติ. ทิพฺพาย โสตธาตุยา สทฺทํ สุณาติ. เจโตปริญาเณน \csromanfolio{88}สปฺปจฺจย\-จิตฺต\-สมงฺคิสฺส จิตฺตํ ชานาติ. อากาสา\-นญฺจา\-ยตนํ วิญฺญาณญฺจา\-ยตนสฺส ฯเปฯ. อากิญฺจญฺญา\-ยตนํ เนว\-สญฺญา\-นา\-สญฺญา\-ยตนสฺส ฯเปฯ. รูปายตนํ จกฺขุ\-วิญฺญาณสฺส ฯเปฯ. โผฏฺฐพฺพา\-ยตนํ กายวิญฺญาณสฺส ฯเปฯ. สปฺปจฺจยา ขนฺธา อิทฺธิวิธ\-ญาณสฺส, เจโตปริย\-ญาณสฺส, ปุพฺเพ\-นิวาสา\-นุสฺสติ\-ญาณสฺส, ยถากมฺมูปคญาณสฺส, อนาคตํสญาณสฺส, อาวชฺชนาย อารมฺมณ\-ปจฺจเยน ปจฺจโย. (1)}

\prose{อปฺปจฺจโย ธมฺโม สปฺปจฺจยสฺส ธมฺมสฺส อารมฺมณ\-ปจฺจเยน ปจฺจโย. อริยา นิพฺพานํ ปจฺจเวกฺขนฺติ. นิพฺพานํ โคตฺรภุสฺส, โวทานสฺส, มคฺคสฺส, ผลสฺส, อาวชฺชนาย อารมฺมณ\-ปจฺจเยน ปจฺจโย. (1)}

\prose{สปฺปจฺจโย ธมฺโม สปฺปจฺจยสฺส ธมฺมสฺส อธิปติ\-ปจฺจเยน ปจฺจโย. อารมฺมณา\-ธิปติ, สหชาตา\-ธิปติ.  อารมฺมณา\-ธิปติ\csromandash{}ทานํ ทตฺวา, สีลํ ฯเปฯ อุโปสถกมฺมํ กตฺวา ตํ ครุํ กตฺวา ปจฺจเวกฺขติ. ปุพฺเพ ฯเปฯ. ฌานา ฯเปฯ. อริยา มคฺคา วุฏฺฐหิตฺวา มคฺคํ ครุํ กตฺวา ปจฺจเวกฺขนฺติ. ผลํ ครุํ กตฺวา ฯเปฯ. จกฺขุํ ฯเปฯ วตฺถุํ. สปฺปจฺจเย ขนฺเธ ครุํ กตฺวา อสฺสาเทติ อภินนฺทติ, ตํ ครุํ กตฺวา ราโค อุปฺปชฺชติ. ทิฏฺฐิ อุปฺปชฺชติ. สหชาตา\-ธิปติ\csromandash{}สปฺปจฺจยา\-ธิปติ สมฺปยุตฺตกานํ ขนฺธานํ จิตฺตสมุฏฺฐานานญฺจ รูปานํ อธิปติ\-ปจฺจเยน ปจฺจโย. (1)}

\prose{อปฺปจฺจโย ธมฺโม สปฺปจฺจยสฺส ธมฺมสฺส อธิปติ\-ปจฺจเยน ปจฺจโย. อารมฺมณา\-ธิปติ\csromandash{}อริยา นิพฺพานํ ครุํ กตฺวา ปจฺจเวกฺขนฺติ. นิพฺพานํ โคตฺรภุสฺส, โวทานสฺส, มคฺคสฺส, ผลสฺส อธิปติ\-ปจฺจเยน ปจฺจโย. (1)}

\prose{สปฺปจฺจโย ธมฺโม สปฺปจฺจยสฺส ธมฺมสฺส อนนฺตร\-ปจฺจเยน ปจฺจโย ฯเปฯ อุปนิสฺสย\-ปจฺจเยน ปจฺจโย ฯเปฯ. (เทฺว ปญฺหา อุปนิสฺสยมูลํ.) ปุเรชาต\-ปจฺจเยน ปจฺจโย ฯเปฯ อวิคต\-ปจฺจเยน ปจฺจโย. (สพฺพตฺถ เอกาเยว ปญฺหา.)}

\csromanmarkh{7. สปฺปจฺจยทุก}\csromanmarkhii{1. ปจฺจยานุโลม 2. สงฺขฺยาวาร}\csromanheadernomark{2}{7. สปฺปจฺจยทุก \textbf{1. ปจฺจยานุโลม 2. สงฺขฺยาวาร}}
\proseitem{12}{\csromanfolio{89}เหตุยา เอกํ, อารมฺมเณ เทฺว, อธิปติยา เทฺว, อนนฺตเร เอกํ, สมนนฺตเร เอกํ, สหชาเต เอกํ, อญฺญมญฺเญ เอกํ, นิสฺสเย เอกํ, อุปนิสฺสเย เทฺว, ปุเรชาเต เอกํ, (สพฺพตฺถ เอกํ.) อวิคเต เอกํ. (เอวํ คเณตพฺพํ.)}

\vspace{\tipitakaparskip}
\csromancenter{\textbf{อนุโลมํ.}}
\proseitem{13}{สปฺปจฺจโย ธมฺโม สปฺปจฺจยสฺส ธมฺมสฺส อารมฺมณ\-ปจฺจเยน ปจฺจโย. สหชาต\-ปจฺจเยน ปจฺจโย. อุปนิสฺสย\-ปจฺจเยน ปจฺจโย. ปุเรชาต\-ปจฺจเยน ปจฺจโย. ปจฺฉาชาต\-ปจฺจเยน ปจฺจโย. กมฺมปจฺจเยน ปจฺจโย. อาหารปจฺจเยน ปจฺจโย. อินฺทฺริย\-ปจฺจเยน ปจฺจโย. (1)}

\prose{อปฺปจฺจโย ธมฺโม สปฺปจฺจยสฺส ธมฺมสฺส อารมฺมณ\-ปจฺจเยน ปจฺจโย. อุปนิสฺสย\-ปจฺจเยน ปจฺจโย. (1)}

\csromanmarkh{2. ปจฺจย\-ปจฺจนีย}\csromanmarkhii{2. สงฺขฺยาวาร}\csromanheadernomark{2}{2. ปจฺจย\-ปจฺจนีย 2. สงฺขฺยาวาร}
\proseitem{14}{นเหตุยา เทฺว, นอารมฺมเณ เอกํ, นอธิปติยา เทฺว, นอนนฺตเร เทฺว, นสมนนฺตเร เทฺว ฯเปฯ นอุปนิสฺสเย เทฺว, นปุเรชาเต เทฺว ฯเปฯ โนวิคเต เทฺว, โนอวิคเต เทฺว. (เอวํ คเณตพฺพํ.)}

\vspace{\tipitakaparskip}
\csromancenter{\textbf{ปจฺจนียํ.}}
\proseitem{15}{เหตุปจฺจยา นอารมฺมเณ เอกํ, นอธิปติยา เอกํ, นอนนฺตเร เอกํ, นสมนนฺตเร เอกํ, นอญฺญมญฺเญ เอกํ, นอุปนิสฺสเย เอกํ ฯเปฯ นสมฺปยุตฺเต เอกํ, นวิปฺปยุตฺเต เอกํ, โนนตฺถิยา เอกํ, โนวิคเต เอกํ. (เอวํ คเณตพฺพํ.)}

\vspace{\tipitakaparskip}
\csromancenter{อนุโลม\-ปจฺจนียํ.}
\proseitem{16}{\csromanfolio{90}นเหตุปจฺจยา อารมฺมเณ เทฺว, อธิปติยา เทฺว, อนนฺตเร เอกํ ฯเปฯ อุปนิสฺสเย เทฺว, ปุเรชาเต เอกํ ฯเปฯ อวิคเต เอกํ. (เอวํ คเณตพฺพํ.)}

\vspace{\tipitakaparskip}
\csromancenter{ปจฺจนียา\-นุโลมํ.}
\nitthitam{สปฺปจฺจยทุกํ นิฏฺฐิตํ.}
\csromanmatikaanchor{mtk.11}
\csromantocmarkref{section}{8. สงฺขตทุก}{mtk.11}
\bookmark[dest={mtk.11},level=1]{8. สงฺขตทุก}
\csromanmarkh{8. สงฺขตทุก}\csromanmarkhii{1. ปฏิจฺจวาร}\csromanheadernomark{1}{8. \textbf{สงฺขตทุก 1. ปฏิจฺจวาร}}
\csromanheader{2}{\textbf{เหตุ}}
\proseitem{17}{สงฺขตํ ธมฺมํ ปฏิจฺจ สงฺขโต ธมฺโม อุปฺปชฺชติ เหตุปจฺจยา. สงฺขตํ เอกํ ขนฺธํ ปฏิจฺจ ตโย ขนฺธา จิตฺต\-สมุฏฺฐานญฺจ รูปํ ฯเปฯ. ปฏิสนฺธิ\-กฺขเณ ฯเปฯ. ขนฺเธ ปฏิจฺจ วตฺถุ, วตฺถุํ ปฏิจฺจ ขนฺธา. เอกํ มหาภูตํ ฯเปฯ มหาภูเต ปฏิจฺจ จิตฺต\-สมุฏฺฐานํ รูปํ, กฏตฺตารูปํ, อุปาทารูปํ.}

\prosewithcloser{(อิมํ ทุกํ ยถา สปฺปจฺจยทุกํ, เอวํ คเณตพฺพํ, นินฺนานากรณํ.)}{สงฺขตทุกํ นิฏฺฐิตํ.}
\csromanmatikaanchor{mtk.12}
\csromantocmarkref{section}{9. สนิทสฺสนทุก}{mtk.12}
\bookmark[dest={mtk.12},level=1]{9. สนิทสฺสนทุก}
\csromanmarkh{9. สนิทสฺสนทุก}\csromanmarkhii{1. ปฏิจฺจวาร}\csromanheadernomark{1}{9. \textbf{สนิทสฺสนทุก 1. ปฏิจฺจวาร}}
\csromanmarkh{1. ปจฺจยานุโลม}\csromanmarkhii{1. วิภงฺควาร}\csromanheadernomark{2}{1. \textbf{ปจฺจยานุโลม 1. วิภงฺควาร}}
\csromanheader{2}{\textbf{เหตุ}}
\proseitem{18}{อนิทสฺสนํ ธมฺมํ ปฏิจฺจ อนิทสฺสโน ธมฺโม อุปฺปชฺชติ เหตุปจฺจยา. อนิทสฺสนํ เอกํ ขนฺธํ ปฏิจฺจ ตโย ขนฺธา อนิทสฺสนํ จิตฺต\-สมุฏฺฐานญฺจ รูปํ ฯเปฯ เทฺว ขนฺเธ ฯเปฯ. ปฏิสนฺธิ\-กฺขเณ อนิทสฺสนํ เอกํ ขนฺธํ ปฏิจฺจ ตโย ขนฺธา อนิทสฺสนํ กฏตฺตา จ รูปํ ฯเปฯ เทฺว ขนฺเธ ฯเปฯ. ขนฺเธ ปฏิจฺจ วตฺถุ, วตฺถุํ ปฏิจฺจ ขนฺธา. เอกํ มหาภูตํ ฯเปฯ มหาภูเต ปฏิจฺจ อนิทสฺสนํ จิตฺต\-สมุฏฺฐานํ รูปํ, กฏตฺตารูปํ, อุปาทารูปํ. (1)}

\prose{\csromanfolio{91}อนิทสฺสนํ ธมฺมํ ปฏิจฺจ สนิทสฺสโน ธมฺโม อุปฺปชฺชติ เหตุปจฺจยา. อนิทสฺสเน ขนฺเธ ปฏิจฺจ สนิทสฺสนํ จิตฺต\-สมุฏฺฐานํ รูปํ. ปฏิสนฺธิ\-กฺขเณ ฯเปฯ มหาภูเต ปฏิจฺจ สนิทสฺสนํ จิตฺต\-สมุฏฺฐานํ รูปํ, กฏตฺตารูปํ, อุปาทารูปํ. (2)}

\prose{อนิทสฺสนํ ธมฺมํ ปฏิจฺจ สนิทสฺสโน จ อนิทสฺสโน จ ธมฺมา อุปฺปชฺชนฺติ เหตุปจฺจยา. อนิทสฺสนํ เอกํ ขนฺธํ ปฏิจฺจ ตโย ขนฺธา สนิทสฺสนญฺจ อนิทสฺสนญฺจ จิตฺต\-สมุฏฺฐานญฺจ รูปํ ฯเปฯ เทฺว ขนฺเธ ฯเปฯ. ปฏิสนฺธิ\-กฺขเณ ฯเปฯ มหาภูเต ปฏิจฺจ สนิทสฺสนญฺจ อนิทสฺสนญฺจ จิตฺต\-สมุฏฺฐานญฺจ รูปํ, กฏตฺตารูปํ, อุปาทารูปํ. (3)}

\proseitem{19}{อนิทสฺสนํ ธมฺมํ ปฏิจฺจ อนิทสฺสโน ธมฺโม อุปฺปชฺชติ อารมฺมณ ปจฺจโย. อนิทสฺสนํ เอกํ ขนฺธํ ปฏิจฺจ ตโย ขนฺธา ฯเปฯ เทฺว ขนฺเธ ฯเปฯ. ปฏิสนฺธิ\-กฺขเณ ฯเปฯ. วตฺถุํ ปฏิจฺจ ขนฺธา. (1)}

\proseitem{20}{อนิทสฺสนํ ธมฺมํ ปฏิจฺจ อนิทสฺสโน ธมฺโม อุปฺปชฺชติ อธิปติ\-ปจฺจยา. อนิทสฺสนํ เอกํ ขนฺธํ ปฏิจฺจ ตโย ขนฺธา อนิทสฺสนํ จิตฺต\-สมุฏฺฐานญฺจ รูปํ ฯเปฯ เทฺว ขนฺเธ ฯเปฯ. เอกํ มหาภูตํ ปฏิจฺจ ตโย มหาภูตา ฯเปฯ เทฺว มหาภูเต ฯเปฯ มหาภูเต ปฏิจฺจ อนิทสฺสนํ จิตฺต\-สมุฏฺฐานํ รูปํ, อุปาทารูปํ. (1)}

\prose{อนิทสฺสนํ ธมฺมํ ปฏิจฺจ สนิทสฺสโน ธมฺโม อุปฺปชฺชติ อธิปติ\-ปจฺจยา. อนิทสฺสเน ขนฺเธ ปฏิจฺจ สนิทสฺสนํ จิตฺต\-สมุฏฺฐานํ รูปํ. มหาภูเต ปฏิจฺจ สนิทสฺสนํ จิตฺต\-สมุฏฺฐานํ รูปํ. (2)}

\prose{อนิทสฺสนํ ธมฺมํ ปฏิจฺจ สนิทสฺสโน จ อนิทสฺสโน จ ธมฺมา อุปฺปชฺชนฺติ อธิปติ\-ปจฺจยา. อนิทสฺสนํ เอกํ ขนฺธํ ปฏิจฺจ ตโย ขนฺธา สนิทสฺสนญฺจ อนิทสฺสนญฺจ จิตฺต\-สมุฏฺฐานํ รูปํ ฯเปฯ เทฺว ขนฺเธ ฯเปฯ มหาภูเต ปฏิจฺจ สนิทสฺสนญฺจ อนิทสฺสนญฺจ จิตฺต\-สมุฏฺฐานํ รูปํ, อุปาทารูปํ. (3) (สํขิตฺตํ. สพฺเพ กาตพฺพา.)}

\csromanmarkh{1. ปจฺจยานุโลม}\csromanmarkhii{2. สงฺขฺยาวาร}\csromanheadernomark{2}{1. \textbf{ปจฺจยานุโลม 2. สงฺขฺยาวาร}}
\proseitem{21}{\csromanfolio{92}เหตุยา ตีณิ, อารมฺมเณ เอกํ, อธิปติยา ตีณิ, อนนฺตเร เอกํ, สมนนฺตเร เอกํ, สหชาเต ตีณิ, อญฺญมญฺเญ เอกํ, นิสฺสเย ตีณิ, อุปนิสฺสเย เอกํ, ปุเรชาเต เอกํ, อาเสวเน เอกํ, กมฺเม ตีณิ, วิปาเก ตีณิ, (สพฺพตฺถ ตีณิ.) มคฺเค ตีณิ, สมฺปยุตฺเต เอกํ, วิปฺปยุตฺเต ตีณิ, อตฺถิยา ตีณิ, นตฺถิยา เอกํ, วิคเต เอกํ, อวิคเต ตีณิ. (เอวํ คเณตพฺพํ.)}

\csromanheader{2}{อนุโลมํ}
\csromanmarkh{2. ปจฺจย\-ปจฺจนีย}\csromanmarkhii{1. วิภงฺควาร}\csromanheadernomark{2}{2. ปจฺจย\-ปจฺจนีย 1. วิภงฺควาร}
\proseitem{22}{อนิทสฺสนํ ธมฺมํ ปฏิจฺจ อนิทสฺสโน ธมฺโม อุปฺปชฺชติ นเหตุปจฺจยา. อเหตุกํ อนิทสฺสนํ เอกํ ขนฺธํ ปฏิจฺจ ตโย ขนฺธา อนิทสฺสนํ จิตฺต\-สมุฏฺฐานญฺจ รูปํ ฯเปฯ เทฺว ขนฺเธ ฯเปฯ. อเหตุก\-ปฏิสนฺธิ\-กฺขเณ ฯเปฯ. ขนฺเธ ปฏิจฺจ วตฺถุ, วตฺถุํ ปฏิจฺจ ขนฺธา. เอกํ มหาภูตํ ฯเปฯ มหาภูเต ปฏิจฺจ อนิทสฺสนํ จิตฺต\-สมุฏฺฐานํ รูปํ, กฏตฺตารูปํ, อุปาทารูปํ. พาหิรํ. อาหารสมุฏฺฐานํ. อุตุสมุฏฺฐานํ. อสญฺญสตฺตานํ ฯเปฯ. วิจิกิจฺฉา\-สหคเต อุทฺธจฺจ\-สหคเต ขนฺเธ ปฏิจฺจ วิจิกิจฺฉา\-สหคโต อุทฺธจฺจ\-สหคโต โมโห. (1)}

\prose{อนิทสฺสนํ ธมฺมํ ปฏิจฺจ สนิทสฺสโน ธมฺโม อุปฺปชฺชติ นเหตุปจฺจยา. อเหตุเก อนิทสฺสเน ขนฺเธ ปฏิจฺจ สนิทสฺสนํ จิตฺต\-สมุฏฺฐานํ รูปํ. ปฏิสนฺธิ\-กฺขเณ ฯเปฯ มหาภูเต ปฏิจฺจ สนิทสฺสนํ จิตฺต\-สมุฏฺฐานํ รูปํ, กฏตฺตารูปํ, อุปาทารูปํ. พาหิเร. อาหารสมุฏฺฐาเน. อุตุสมุฏฺฐาเน. อสญฺญสตฺตานํ มหาภูเต ปฏิจฺจ สนิทสฺสนํ กฏตฺตารูปํ, อุปาทารูปํ. (2)}

\prose{อนิทสฺสนํ ธมฺมํ ปฏิจฺจ สนิทสฺสโน จ อนิทสฺสโน จ ธมฺมา อุปฺปชฺชนฺติ นเหตุปจฺจยา. อเหตุกํ อนิทสฺสนํ เอกํ ขนฺธํ ปฏิจฺจ ตโย ขนฺธา \csromanfolio{93}สนิทสฺสนญฺจ อนิทสฺสนญฺจ จิตฺต\-สมุฏฺฐานํ รูปํ ฯเปฯ เทฺว ขนฺเธ ฯเปฯ. ปฏิสนฺธิ\-กฺขเณ ฯเปฯ มหาภูเต ปฏิจฺจ ฯเปฯ. พาหิเร. อาหารสมุฏฺฐาเน. อุตุสมุฏฺฐาเน. อสญฺญสตฺตานํ มหาภูเต ปฏิจฺจ สนิทสฺสนญฺจ อนิทสฺสนญฺจ กฏตฺตารูปํ, อุปาทารูปํ. (3) (เอวํ สพฺเพ กาตพฺพา.)}

\csromanmarkh{2. ปจฺจย\-ปจฺจนีย}\csromanmarkhii{2. สงฺขฺยาวาร}\csromanheadernomark{2}{2. ปจฺจย\-ปจฺจนีย 2. สงฺขฺยาวาร}
\proseitem{23}{นเหตุยา ตีณิ นอารมฺมเณ ตีณิ นอธิปติยา ตีณิ, นอนนฺตเร ตีณิ, นสมนนฺตเร ตีณิ นอญฺญมญฺเญ ตีณิ, นอุปนิสฺสเย ตีณิ นปุเรชาเต ตีณิ, นปจฺฉาชาเต ตีณิ, นอาเสวเน ตีณิ, นกมฺเม ตีณิ, นวิปาเก ตีณิ, นอาหาเร ตีณิ, นอินฺทฺริเย ตีณิ, นฌาเน ตีณิ, นมคฺเค ตีณิ, นสมฺปยุตฺเต ตีณิ, นวิปฺปยุตฺเต ตีณิ, โนนตฺถิยา ตีณิ, โนวิคเต ตีณิ. (เอวํ คเณตพฺพํ.)}

\vspace{\tipitakaparskip}
\csromancenter{ปจฺจนียํ.}
\proseitem{22}{เหตุปจฺจยา นอารมฺมเณ ตีณิ, นอธิปติยา ตีณิ, (สพฺพตฺถ ตีณิ.) นกมฺเม เอกํ, นวิปาเก ตีณิ, นสมฺปยุตฺเต ตีณิ, นวิปฺปยุตฺเต เอกํ, โนนตฺถิยา ตีณิ, โนวิคเต ตีณิ.}

\vspace{\tipitakaparskip}
\csromancenter{อนุโลม\-ปจฺจนียํ.}
\proseitem{25}{นเหตุปจฺจยา อารมฺมเณ เอกํ, อนนฺตเร เอกํ, สมนนฺตเร เอกํ, สหชาเต ตีณิ, อญฺญมญฺเญ เอกํ, นิสฺสเย ตีณิ, อุปนิสฺสเย \csromanfolio{94}เอกํ, ปุเรชาเต เอกํ, อาเสวเน เอกํ, กมฺเม ตีณิ ฯเปฯ ฌาเน ตีณิ, มคฺเค เอกํ, สมฺปยุตฺเต เอกํ, วิปฺปยุตฺเต ตีณิ, อตฺถิยา ตีณิ, นตฺถิยา เอกํ, วิคเต เอกํ, อวิคเต ตีณิ.}

\vspace{\tipitakaparskip}
\csromancenter{ปจฺจนียา\-นุโลมํ.}
\csromancenter{(สหชาตวาโรปิ ปฏิจฺจ\-วาร\-สทิโส.)}
\csromanmarkh{9. สนิทสฺสนทุก}\csromanmarkhii{3. ปจฺจยวาร}\csromanheadernomark{2}{9. \textbf{สนิทสฺสนทุก 3. ปจฺจยวาร}}
\csromanheader{2}{1. \textbf{ปจฺจยานุโลม}}
\csromanheader{2}{\textbf{เหตุ}}
\proseitem{25}{อนิทสฺสนํ ธมฺมํ ปจฺจยา อนิทสฺสโน ธมฺโม อุปฺปชฺชติ เหตุปจฺจยา. อนิทสฺสนํ เอกํ ขนฺธํ ปจฺจยา ตโย ขนฺธา อนิทสฺสนํ จิตฺต\-สมุฏฺฐานญฺจ รูปํ ฯเปฯ เทฺว ขนฺเธ ฯเปฯ. ปฏิสนฺธิ\-กฺขเณ ฯเปฯ. ขนฺเธ ปจฺจยา วตฺถุ, วตฺถุํ ปจฺจยา ขนฺธา. เอกํ มหาภูตํ ปจฺจยา ฯเปฯ มหาภูเต ปจฺจยา อนิทสฺสนํ จิตฺต\-สมุฏฺฐานํ รูปํ, กฏตฺตารูปํ, อุปาทารูปํ. วตฺถุํ ปจฺจยา อนิทสฺสนา ขนฺธา. (อิตเรปิ เทฺว ปญฺหา กาตพฺพา.)}

\proseitem{27}{อนิทสฺสนํ ธมฺมํ ปจฺจยา อนิทสฺสโน ธมฺโม อุปฺปชฺชติ อารมฺมณ\-ปจฺจยา. อนิทสฺสนํ เอกํ ขนฺธํ ฯเปฯ เทฺว ขนฺเธ ฯเปฯ. ปฏิสนฺธิ\-กฺขเณ ฯเปฯ. วตฺถุํ ปจฺจยา ขนฺธา. จกฺขายตนํ ปจฺจยา จกฺขุวิญฺญาณํ ฯเปฯ. กายายตนํ ปจฺจยา กายวิญฺญาณํ. วตฺถุํ ปจฺจยา อนิทสฺสนา ขนฺธา. (สํขิตฺตํ.)}

\proseitem{25}{\textbf{เหตุ}ยา ตีณิ, อารมฺมเณ เอกํ, อธิปติยา ตีณิ ฯเปฯ อวิคเต ตีณิ.}

\vspace{\tipitakaparskip}
\csromancenter{อนุโลมํ.}
\csromanmarkh{9. สนิทสฺสนทุก}\csromanmarkhii{2. ปจฺจย\-ปจฺจนีย}\csromanheadernomark{2}{9. สนิทสฺสนทุก 2. ปจฺจย\-ปจฺจนีย}
\proseitem{29}{\csromanfolio{95}อนิทสฺสนํ ธมฺมํ ปจฺจยา อนิทสฺสโน ธมฺโม อุปฺปชฺชติ นเหตุปจฺจยา. อเหตุกํ อนิทสฺสนํ เอกํ ขนฺธํ ปจฺจยา ตโย ขนฺธา อนิทสฺสนํ จิตฺต\-สมุฏฺฐานญฺจ รูปํ ฯเปฯ เทฺว ขนฺเธ ฯเปฯ. อเหตุก\-ปฏิสนฺธิ\-กฺขเณ ฯเปฯ. (ยาว อสญฺญสตฺตา.) จกฺขายตนํ ปจฺจยา จกฺขุวิญฺญาณํ ฯเปฯ. กายายตนํ ปจฺจยา กายวิญฺญาณํ. วตฺถุํ ปจฺจยา อเหตุกา อนิทสฺสนา ขนฺธา. วิจิกิจฺฉา\-สหคเต อุทฺธจฺจ\-สหคเต ขนฺเธ จ วตฺถุญฺจ ปจฺจยา วิจิกิจฺฉา\-สหคโต อุทฺธจฺจ\-สหคโต โมโห. (อิตเรปิ เทฺว กาตพฺพา. สํขิตฺตํ.)}

\vspace{\tipitakaparskip}
\csromancenter{\textbf{นเหตุ}ยา ตีณิ, นอารมฺมเณ ตีณิ ฯเปฯ โนวิคเต ตีณิ.}
\csromancenter{\textbf{ปจฺจนียํ.}}
\proseitem{31}{เหตุปจฺจยา นอารมฺมเณ ตีณิ ฯเปฯ นกมฺเม เอกํ ฯเปฯ นวิปฺปยุตฺเต เอกํ, โนนตฺถิยา ตีณิ, โนวิคเต ตีณิ.}

\vspace{\tipitakaparskip}
\csromancenter{อนุโลม\-ปจฺจนียํ.}
\csromancenter{\textbf{นเหตุ}ปจฺจยา อารมฺมเณ เอกํ ฯเปฯ มคฺเค เอกํ ฯเปฯ อวิคเตตีณิ.}
\csromancenter{ปจฺจนียา\-นุโลมํ.}
\csromancenter{(นิสฺสยวาโรปิ เอวํ กาตพฺโพ.)}
\csromanmarkh{9. สนิทสฺสนทุก}\csromanmarkhii{5. สํสฏฺฐวาร 1. ปจฺจยานุโลม}\csromanheadernomark{2}{9. สนิทสฺสนทุก 5. สํสฏฺฐวาร 1. ปจฺจยานุโลม}
\proseitem{33}{อนิทสฺสนํ ธมฺมํ สํสฏฺโฐ อนิทสฺสโน ธมฺโม อุปฺปชฺชติ เหตุปจฺจยา. อนิทสฺสนํ เอกํ ขนฺธํ สํสฏฺฐา ตโย ขนฺธา ฯเปฯ เทฺว ขนฺเธ ฯเปฯ. ปฏิสนฺธิ\-กฺขเณ ฯเปฯ.}

\prose{\csromanfolio{96}อนิทสฺสนํ ธมฺมํ สํสฏฺโฐ อนิทสฺสโน ธมฺโม อุปฺปชฺชติ อารมฺมณ\-ปจฺจยา.}

\vspace{\tipitakaparskip}
\csromancenter{(เอวํ สพฺพํ สปฺปจฺจย\-คณนาหิ สทฺธิํ กาตพฺพํ.)}
\csromancenter{(สมฺปยุตฺตวาโรปิ สํสฏฺฐวาร\-สทิโส.)}
\csromanmarkh{9. สนิทสฺสนทุก}\csromanmarkhii{7. ปญฺหาวาร}\csromanheadernomark{2}{9. \textbf{สนิทสฺสนทุก 7. ปญฺหาวาร}}
\csromanmarkh{1. ปจฺจยานุโลม}\csromanmarkhii{1. วิภงฺควาร}\csromanheadernomark{2}{1. \textbf{ปจฺจยานุโลม 1. วิภงฺควาร}}
\csromanheader{2}{\textbf{เหตุ}}
\proseitem{34}{อนิทสฺสโน ธมฺโม อนิทสฺสนสฺส ธมฺมสฺส เหตุปจฺจเยน ปจฺจโย. อนิทสฺสนา เหตู สมฺปยุตฺตกานํ ขนฺธานํ อนิทสฺสนานํ จิตฺตสมุฏฺฐานานญฺจ รูปานํ เหตุปจฺจเยน ปจฺจโย. ปฏิสนฺธิ\-กฺขเณ ฯเปฯ. (1)}

\prose{อนิทสฺสโน ธมฺโม สนิทสฺสนสฺส ธมฺมสฺส เหตุปจฺจเยน ปจฺจโย. อนิทสฺสนา เหตู สนิทสฺสนานํ จิตฺต\-สมุฏฺฐานานํ รูปานํ เหตุปจฺจเยน ปจฺจโย. ปฏิสนฺธิ\-กฺขเณ ฯเปฯ. (2)}

\prose{อนิทสฺสโน ธมฺโม สนิทสฺสนสฺส จ อนิทสฺสนสฺส จ ธมฺมสฺส เหตุปจฺจเยน ปจฺจโย. อนิทสฺสนา เหตู สมฺปยุตฺตกานํ ขนฺธานํ สนิทสฺส\-นานญฺจ อนิทสฺสนานญฺจ จิตฺตสมุฏฺฐานานญฺจ รูปานํ เหตุปจฺจเยน ปจฺจโย. ปฏิสนฺธิ\-กฺขเณ ฯเปฯ. (3)}

\proseitem{35}{ยนิทสฺสโน ธมฺโม อนิทสฺสนสฺส ธมฺมสฺส อารมฺมณ\-ปจฺจเยน ปจฺจโย. สนิทสฺสนํ รูปํ อนิจฺจโต ฯเปฯ โทมนสฺสํ อุปฺปชฺชติ. ทิพฺเพน จกฺขุนา รูปํ ปสฺสติ. รูปายตนํ จกฺขุ\-วิญฺญาณสฺส อารมฺมณ\-ปจฺจเยน ปจฺจโย. สนิทสฺสนา ขนฺธา อิทฺธิวิธ\-ญาณสฺส, ปุพฺเพ\-นิวาสา\-นุสฺสติ\-ญาณสฺส, อนาคตํสญาณสฺส, อาวชฺชนาย อารมฺมณ\-ปจฺจเยน ปจฺจโย. (1)}

\prose{อนิทสฺสโน ธมฺโม อนิทสฺสนสฺส ธมฺมสฺส อารมฺมณ\-ปจฺจเยน ปจฺจโย. ทานํ ทตฺวา, สีลํ ฯเปฯ อุโปสถกมฺมํ กตฺวา ตํ ปจฺจเวกฺขติ ปุพฺเพ สุจิณฺณานิ ปจฺจเวกฺขติ. ฌานา วุฏฺฐหิตฺวา ฯเปฯ. อริยา มคฺคา วุฏฺฐหิตฺวา มคฺคํ ปจฺจเวกฺขนฺติ, \csromanfolio{97}ผลํ ปจฺจเวกฺขนฺติ, นิพฺพานํ ปจฺจเวกฺขนฺติ. นิพฺพานํ โคตฺรภุสฺส, โวทานสฺส, มคฺคสฺส, ผลสฺส, อาวชฺชนาย อารมฺมณ\-ปจฺจเยน ปจฺจโย. อริยา ปหีเน กิเลเส ปจฺจเวกฺขนฺติ, วิกฺขมฺภิเต กิเลเส ฯเปฯ. ปุพฺเพ ฯเปฯ. จกฺขุํ ฯเปฯ. กายํ. สทฺเท ฯเปฯ. วตฺถุํ อนิทสฺสเน ขนฺเธ อนิจฺจโต ฯเปฯ โทมนสฺสํ อุปฺปชฺชติ. ทิพฺพาย โสตธาตุยา สทฺทํ สุณาติ. เจโตปริย\-ญาเณน อนิทสฺสน\-จิตฺต\-สมงฺคิสฺส จิตฺตํ ชานาติ. อากาสา\-นญฺจา\-ยตนํ วิญฺญาณญฺจา\-ยตนสฺส ฯเปฯ. อากิญฺจญฺญา\-ยตนํ เนว\-สญฺญา\-นา\-สญฺญา\-ยตนสฺส ฯเปฯ. สทฺทายตนํ โสตวิญฺญาณสฺส ฯเปฯ. โผฏฺฐพฺพา\-ยตนํ กายวิญฺญาณสฺส ฯเปฯ. อนิทสฺสนา ขนฺธา อิทฺธิวิธ\-ญาณสฺส, เจโตปริย\-ญาณสฺส, ปุพฺเพ\-นิวาสา\-นุสฺสติ\-ญาณสฺส, ยถากมฺมูปคญาณสฺส, อนาคตํสญาณสฺส, อาวชฺชนาย อารมฺมณ\-ปจฺจเยน ปจฺจโย. (1)}

\proseitem{36}{สนิทสฺสโน ธมฺโม อนิทสฺสนสฺส ธมฺมสฺส อธิปติ\-ปจฺจเยน ปจฺจโย. อารมฺมณา\-ธิปติ\csromandash{}สนิทสฺสนํ รูปํ ครุํ กตฺวา อสฺสาเทติ อภินนฺทติ, ตํ ครุํ กตฺวา ราโค อุปฺปชฺชติ, ทิฏฺฐิ อุปฺปชฺชติ. (1)}

\prose{อนิทสฺสโน ธมฺโม อนิทสฺสนสฺส ธมฺมสฺส อธิปติ\-ปจฺจเยน ปจฺจโย. อารมฺมณา\-ธิปติ, สหชาตา\-ธิปติ.  อารมฺมณา\-ธิปติ\csromandash{}ทานํ ทตฺวา, สีลํ ฯเปฯ อุโปสถกมฺมํ กตฺวา ตํ ครุํ กตฺวา ปจฺจเวกฺขติ. ปุพฺเพ สุจิณฺณานิ ฯเปฯ. ฌานา วุฏฺฐหิตฺวา ฯเปฯ. อริยา มคฺคา วุฏฺฐหิตฺวา มคฺคํ ครุํ กตฺวา ปจฺจเวกฺขนฺติ, ผลํ ครุํ กตฺวา ปจฺจเวกฺขนฺติ, นิพฺพานํ ครุํ กตฺวา ปจฺจเวกฺขนฺติ. นิพฺพานํ โคตฺรภุสฺส, โวทานสฺส, มคฺคสฺส, ผลสฺส อธิปติ\-ปจฺจเยน ปจฺจโย. จกฺขุํ ฯเปฯ วตฺถุํ อนิทสฺสเน ขนฺเธ ครุํ กตฺวา อสฺสาเทติ อภินนฺทติ, ตํ ครุํ กตฺวา ราโค อุปฺปชฺชติ, ทิฏฺฐิ อุปฺปชฺชติ. สหชาตา\-ธิปติ\csromandash{}อนิทสฺสนา\-ธิปติ สมฺปยุตฺตกานํ ขนฺธานํ อนิทสฺสนานํ จิตฺตสมุฏฺฐานานญฺจ รูปานํ อธิปติ\-ปจฺจเยน ปจฺจโย. (1)}

\prose{อนิทสฺสโน ธมฺโม สนิทสฺสนสฺส ธมฺมสฺส อธิปติ\-ปจฺจเยน ปจฺจโย. สหชาตา\-ธิปติ\csromandash{}อนิทสฺสนา\-ธิปติ สนิทสฺสนานํ จิตฺต\-สมุฏฺฐานานํ รูปานํ อธิปติ\-ปจฺจเยน ปจฺจโย. (2)}

\prose{อนิทสฺสโน ธมฺโม สนิทสฺสนสฺส จ อนิทสฺสนสฺส จ ธมฺมสฺส อธิปติ\-ปจฺจเยน ปจฺจโย. สหชาตา\-ธิปติ\csromandash{}อนิทสฺสนา\-ธิปติ สมฺปยุตฺตกานํ \csromanfolio{98}ขนฺธานํ สนิทสฺส\-นานญฺจ อนิทสฺสนานญฺจ จิตฺตสมุฏฺฐานานญฺจ รูปานํ อธิปติ\-ปจฺจเยน ปจฺจโย.(3)}

\proseitem{37}{อนิทสฺสโน ธมฺโม อนิทสฺสนสฺส ธมฺมสฺส อนนฺตร\-ปจฺจเยน ปจฺจโย. ปุริมา ปุริมา อนิทสฺสนา ขนฺธา ปจฺฉิมานํ ปจฺฉิมานํ อนิทสฺสนานํ ขนฺธานํ อนนฺตร\-ปจฺจเยน ปจฺจโย. อนุโลมํ โคตฺรภุสฺส ฯเปฯ. โคตฺรภุ มคฺคสฺส ฯเปฯ. เนว\-สญฺญา\-นา\-สญฺญา\-ยตนํ ผลสมาปตฺติยา อนนฺตร\-ปจฺจเยน ปจฺจโย. (1)}

\csromanheader{2}{\textbf{สมนนฺตราทิ}}
\proseitem{38}{อนิทสฺสโน ธมฺโม อนิทสฺสนสฺส ธมฺมสฺส สมนนฺตร\-ปจฺจเยน ปจฺจโย. สหชาต\-ปจฺจเยน ปจฺจโย. ตีณิ. อญฺญมญฺญ\-ปจฺจเยน ปจฺจโย. เอกํ. นิสฺสย\-ปจฺจเยน ปจฺจโย. ตีณิ.}

\csromanheader{2}{\textbf{อุปนิสฺสย}}
\proseitem{39}{สนิทสฺสโน ธมฺโม อนิทสฺสนสฺส ธมฺมสฺส อุปนิสฺสย\-ปจฺจเยน ปจฺจโย. อารมฺมณู\-ปนิสฺสโย, ปกตู\-ปนิสฺสโย ฯเปฯ. ปกตู\-ปนิสฺสโย\csromandash{}วณฺณสมฺปทํ ปตฺถยมาโน ทานํ ฯเปฯ สีลํ ฯเปฯ อุโปสถกมฺมํ กโรติ. วณฺณสมฺปทา สทฺธาย ฯเปฯ ปตฺถนาย. กายิกสฺส สุขสฺส, กายิกสฺส ทุกฺขสฺส, มคฺคสฺส, ผลสมาปตฺติยา อุปนิสฺสย\-ปจฺจเยน ปจฺจโย. (1)}

\prose{อนิทสฺสโน ธมฺโม อนิทสฺสนสฺส ธมฺมสฺส อุปนิสฺสย\-ปจฺจเยน ปจฺจโย. อารมฺมณู\-ปนิสฺสโย, อนนฺตรู\-ปนิสฺสโย, ปกตู\-ปนิสฺสโย ฯเปฯ. ปกตู\-ปนิสฺสโย\csromandash{}สทฺธํ อุปนิสฺสาย ทานํ เทติ ฯเปฯ สมาปตฺติํ อุปฺปาเทติ. มานํ ชปฺเปติ. ทิฏฺฐํ คณฺหาติ. สีลํ ฯเปฯ. เสนาสนํ อุปนิสฺสาย ทานํ เทติ ฯเปฯ สํฆํ ภินฺทติ. สทฺธา ฯเปฯ. เสนาสนํ สทฺธาย ฯเปฯ ผลสมาปตฺติยา อุปนิสฺสย\-ปจฺจเยน ปจฺจโย. (1)}

\proseitem{40}{สนิทสฺสโน ธมฺโม อนิทสฺสนสฺส ธมฺมสฺส ปุเรชาต\-ปจฺจเยน ปจฺจโย. สนิทสฺสนํ รูปํ อนิจฺจโต ฯเปฯ โทมนสฺสํ อุปฺปชฺชติ. ทิพฺเพน \csromanfolio{99}จกฺขุนา รูปํ ปสฺสติ. รูปายตนํ จกฺขุ\-วิญฺญาณสฺส ปุเรชาต\-ปจฺจเยน ปจฺจโย. (1)}

\prose{อนิทสฺสโน ธมฺโม อนิทสฺสนสฺส ธมฺมสฺส ปุเรชาต\-ปจฺจเยน ปจฺจโย. อารมฺมณ\-ปุเรชาตํ, วตฺถุ\-ปุเรชาตํ.  อารมฺมณ\-ปุเรชาตํ\csromandash{} จกฺขุํ ฯเปฯ วตฺถุํ อนิจฺจโต ฯเปฯ โทมนสฺสํ อุปฺปชฺชติ. ทิพฺพาย โสตธาตุยา สทฺทํ สุณาติ. สทฺทายตนํ โสตวิญฺญาณสฺส ฯเปฯ. โผฏฺฐพฺพา\-ยตนํ กายวิญฺญาณสฺส ปุเรชาต\-ปจฺจเยน ปจฺจโย. วตฺถุ\-ปุเรชาตํ\csromandash{}จกฺขายตนํ จกฺขุ\-วิญฺญาณสฺส ฯเปฯ. กายายตนํ กายวิญฺญาณสฺส ฯเปฯ. วตฺถุ อนิทสฺสนานํ ขนฺธานํ ปุเรชาต\-ปจฺจเยน ปจฺจโย. (1)}

\prose{สนิทสฺสโน จ อนิทสฺสโน จ ธมฺโม อนิทสฺสนสฺส ธมฺมสฺส ปุเรชาต\-ปจฺจเยน ปจฺจโย. อารมฺมณ\-ปุเรชาตํ, วตฺถุ\-ปุเรชาตํ\csromandash{} รูปายตนญฺจ วตฺถุ จ อนิทสฺสนานํ ขนฺธานํ ปุเรชาต\-ปจฺจเยน ปจฺจโย. รูปายตนญฺจ จกฺขายตนญฺจ จกฺขุ\-วิญฺญาณสฺส ปุเรชาต\-ปจฺจเยน ปจฺจโย. (1)}

\csromanheader{2}{\textbf{ปจฺฉาชาต}}
\proseitem{41}{อนิทสฺสโน ธมฺโม อนิทสฺสนสฺส ธมฺมสฺส ปจฺฉาชาต\-ปจฺจเยน ปจฺจโย. ปจฺฉาชาตา อนิทสฺสนา ขนฺธา ปุเรชาตสฺส อิมสฺส อนิทสฺสนสฺส กายสฺส ปจฺฉาชาต\-ปจฺจเยน ปจฺจโย. (1)}

\prose{อนิทสฺสโน ธมฺโม สนิทสฺสนสฺส ธมฺมสฺส ปจฺฉาชาต\-ปจฺจเยน ปจฺจโย. ปจฺฉาชาตา อนิทสฺสนา ขนฺธา ปุเรชาตสฺส อิมสฺส สนิทสฺสนสฺส กายสฺส ปจฺฉาชาต\-ปจฺจเยน ปจฺจโย. (2)}

\prose{อนิทสฺสโน ธมฺโม สนิทสฺสนสฺส จ อนิทสฺสนสฺส จ ธมฺมสฺส ปจฺฉาชาต\-ปจฺจเยน ปจฺจโย. ปจฺฉาชาตา อนิทสฺสนา ขนฺธา ปุเรชาตสฺส อิมสฺส สนิทสฺสนสฺส จ อนิทสฺสนสฺส จ กายสฺส ปจฺฉาชาต ปจฺจเยน ปจฺจโย. (3)}

\csromanheader{2}{\textbf{อาเสวน}}
\proseitem{42}{อนิทสฺสโน ธมฺโม อนิทสฺสนสฺส ธมฺมสฺส อาเสวน\-ปจฺจเยน ปจฺจโย. ปุริมา ปุริมา อนิทสฺสนา ขนฺธา ปจฺฉิมานํ ปจฺฉิมานํ อนิทสฺสนานํ ขนฺธานํ ฯเปฯ. อนุโลมํ โคตฺรภุสฺส. อนุโลมํ โวทานสฺส. โคตฺรภุ มคฺคสฺส. โวทานํ มคฺคสฺส อาเสวน\-ปจฺจเยน ปจฺจโย. (1)}

\proseitem{43}{\csromanfolio{100}อนิทสฺสโน ธมฺโม อนิทสฺสนสฺส ธมฺมสฺส กมฺมปจฺจเยน ปจฺจโย. สหชาตา, นานากฺขณิกา.  สหชาตา อนิทสฺสนา เจตนา สมฺปยุตฺตกานํ ขนฺธานํ อนิทสฺสนานญฺจ จิตฺต\-สมุฏฺฐานานํ รูปานํ กมฺมปจฺจเยน ปจฺจโย. ปฏิสนฺธิ\-กฺขเณ ฯเปฯ. นานากฺขณิกา อนิทสฺสนา เจตนา วิปากานํ ขนฺธานํ อนิทสฺสนานญฺจ กฏตฺตารูปานํ กมฺมปจฺจเยน ปจฺจโย. (1)}

\prose{อนิทสฺสโน ธมฺโม สนิทสฺสนสฺส ธมฺมสฺส กมฺมปจฺจเยน ปจฺจโย. \textbf{สหชาตา, นานากฺขณิกา}. (วิตฺถาเรตพฺพํ.) (2)}

\prose{อนิทสฺสโน ธมฺโม สนิทสฺสนสฺส จ อนิทสฺสนสฺส จ ธมฺมสฺส กมฺมปจฺจเยน ปจฺจโย. \textbf{สหชาตา, นานา กฺขณิกา}. (วิตฺถาเรตพฺพํ.) (3)}

\csromanheader{2}{\textbf{วิปากาทิ}}
\proseitem{44}{อนิทสฺสโน ธมฺโม อนิทสฺสนสฺส ธมฺมสฺส วิปาก\-ปจฺจเยน ปจฺจโย. ตีณิ. อาหารปจฺจเยน ปจฺจโย. ตีณิ. (ตีสุปิ กพฬีกาโร อาหาโร กาตพฺโพ.) อินฺทฺริย\-ปจฺจเยน ปจฺจโย. ตีณิ. (ตีสุปิ รูปชีวิตินฺทฺริยํ.) ฌานปจฺจเยน ปจฺจโย. ตีณิ. มคฺคปจฺจเยน ปจฺจโย. ตีณิ. สมฺปยุตฺต\-ปจฺจเยน ปจฺจโย. เอกํ.}

\proseitem{45}{อนิทสฺสโน ธมฺโม อนิทสฺสนสฺส ธมฺมสฺส วิปฺปยุตฺต\-ปจฺจเยน ปจฺจโย. สหชาตํ, ปุเรชาตํ, ปจฺฉาชาตํ.  สหชาตา อนิทสฺสนา ขนฺธา อนิทสฺสนานํ จิตฺต\-สมุฏฺฐานานํ รูปานํ วิปฺปยุตฺต\-ปจฺจเยน ปจฺจโย. ปฏิสนฺธิ\-กฺขเณ อนิทสฺสนา ขนฺธา อนิทสฺสนานํ กฏตฺตารูปานํ วิปฺปยุตฺต\-ปจฺจเยน ปจฺจโย. ขนฺธา วตฺถุสฺส วิปฺปยุตฺต\-ปจฺจเยน ปจฺจโย. วตฺถุ ขนฺธานํ วิปฺปยุตฺต\-ปจฺจเยน ปจฺจโย. ปุเรชาตํ จกฺขายตนํ จกฺขุ\-วิญฺญาณสฺส ฯเปฯ. กายายตนํ กายวิญฺญาณสฺส ฯเปฯ. วตฺถุ อนิทสฺสนานํ ขนฺธานํ วิปฺปยุตฺต\-ปจฺจเยน ปจฺจโย. ปจฺฉาชาตา อนิทสฺสนา ขนฺธา ปุเรชาตสฺส อิมสฺส อนิทสฺสนสฺส กายสฺส วิปฺปยุตฺต\-ปจฺจเยน ปจฺจโย. (1)}

\prose{อนิทสฺสโน ธมฺโม สนิทสฺสนสฺส ธมฺมสฺส วิปฺปยุตฺต\-ปจฺจเยน ปจฺจโย. สหชาตํ, ปจฺฉาชาตํ.  สหชาตา อนิทสฺสนา ขนฺธา สนิทสฺสนานํ \csromanfolio{101}จิตฺต\-สมุฏฺฐานานํ รูปานํ วิปฺปยุตฺต\-ปจฺจเยน ปจฺจโย. ปฏิสนฺธิ\-กฺขเณ ฯเปฯ. ปจฺฉาชาตา อนิทสฺสนา ขนฺธา ปุเรชาตสฺส อิมสฺส สนิทสฺสนสฺส กายสฺส วิปฺปยุตฺต\-ปจฺจเยน ปจฺจโย. (2)}

\prose{อนิทสฺสโน ธมฺโม สนิทสฺสนสฺส จ อนิทสฺสนสฺส จ ธมฺมสฺส วิปฺปยุตฺต\-ปจฺจเยน ปจฺจโย. สหชาตํ, ปจฺฉาชาตํ.  สหชาตา อนิทสฺสนา ขนฺธา สนิทสฺส\-นานญฺจ อนิทสฺสนานญฺจ จิตฺต\-สมุฏฺฐานานํ รูปานํ วิปฺปยุตฺต\-ปจฺจเยน ปจฺจโย. ปฏิสนฺธิ\-กฺขเณ ฯเปฯ. ปจฺฉาชาตา อนิทสฺสนา ขนฺธา ปุเรชาตสฺส อิมสฺส สนิทสฺสนสฺส จ อนิทสฺสนสฺส จ กายสฺส วิปฺปยุตฺต\-ปจฺจเยน ปจฺจโย. (3)}

\prose{อตฺถิอาทิ}

\proseitem{45}{สนิทสฺสโน ธมฺโม อนิทสฺสนสฺส ธมฺมสฺส อตฺถิปจฺจเยน ปจฺจโย. สนิทสฺสนํ รูปํ อนิจฺจโต ฯเปฯ โทมนสฺสํ อุปฺปชฺชติ. ทิพฺเพน จกฺขุนา รูปํ ปสฺสติ. รูปายตนํ จกฺขุ\-วิญฺญาณสฺส อตฺถิปจฺจเยน ปจฺจโย. (1)}

\prose{อนิทสฺสโน ธมฺโม อนิทสฺสนสฺส ธมฺมสฺส อตฺถิปจฺจเยน ปจฺจโย. สหชาตํ, ปุเรชาตํ, ปจฺฉาชาตํ, อาหารํ, อินฺทฺริยํ.  สหชาโต อนิทสฺสโน เอโก ขนฺโธ ติณฺณนฺนํ ขนฺธานํ อนิทสฺสนานญฺจ จิตฺต\-สมุฏฺฐานานํ รูปานํ อตฺถิปจฺจเยน ปจฺจโย ฯเปฯ เทฺว ขนฺเธ ฯเปฯ. (สํขิตฺตํ. ยาว อสญฺญสตฺตา.) ปุเรชาตํ จกฺขุํ ฯเปฯ วตฺถุํ อนิจฺจโต ฯเปฯ โทมนสฺสํ อุปฺปชฺชติ. ทิพฺพาย โสตธาตุยา สทฺทํ สุณาติ. สทฺทายตนํ โสตวิญฺญาณสฺส ฯเปฯ. โผฏฺฐพฺพา\-ยตนํ กายวิญฺญาณสฺส ฯเปฯ. จกฺขายตนํ จกฺขุ\-วิญฺญาณสฺส ฯเปฯ. กายายตนํ กายวิญฺญาณสฺส ฯเปฯ. วตฺถุ อนิทสฺสนานํ ขนฺธานํ อตฺถิปจฺจเยน ปจฺจโย. ปจฺฉาชาตา อนิทสฺสนา ขนฺธา ปุเรชาตสฺส อิมสฺส อนิทสฺสนสฺส กายสฺส อตฺถิปจฺจเยน ปจฺจโย.  กพฬีกาโร อาหาโร อิมสฺส อนิทสฺสนสฺส กายสฺส อตฺถิปจฺจเยน ปจฺจโย.  รูปชีวิตินฺทฺริยํ อนิทสฺสนานํ กฏตฺตารูปานํ อตฺถิปจฺจเยน ปจฺจโย. (1)}

\prose{อนิทสฺสโน ธมฺโม สนิทสฺสนสฺส ธมฺมสฺส อตฺถิปจฺจเยน ปจฺจโย. สหชาตํ, ปจฺฉาชาตํ, อาหารํ, อินฺทฺริยํ.  สหชาตา อนิทสฺสนา ขนฺธา สนิทสฺสนานํ จิตฺต\-สมุฏฺฐานานํ รูปานํ อตฺถิปจฺจเยน ปจฺจโย. ปฏิสนฺธิ\-กฺขเณ ฯเปฯ. มหาภูตา สนิทสฺสนานํ จิตฺต\-สมุฏฺฐานานํ รูปานํ, \csromanfolio{102}กฏตฺตารูปานํ, อุปาทารูปานํ อตฺถิปจฺจเยน ปจฺจโย. พาหิรา. อาหารสมุฏฺฐานา. อุตุสมุฏฺฐานา. อสญฺญสตฺตานํ มหาภูตา สนิทสฺสนานํ กฏตฺตารูปานํ, อุปาทารูปานํ อตฺถิปจฺจเยน ปจฺจโย. \textbf{ปจฺฉาชาตา} อนิทสฺสนา ขนฺธา ปุเรชาตสฺส อิมสฺส สนิทสฺสนสฺส กายสฺส อตฺถิปจฺจเยน ปจฺจโย.  กพฬีกาโร อาหาโร อิมสฺส สนิทสฺสนสฺส กายสฺส อตฺถิปจฺจเยน ปจฺจโย. รูปชีวิตินฺทฺริยํ สนิทสฺสนานํ กฏตฺตารูปานํ อตฺถิปจฺจเยน ปจฺจโย. (2)}

\proseitem{46}{อนิทสฺสโน ธมฺโม สนิทสฺสนสฺส จ อนิทสฺสนสฺส จ ธมฺมสฺส อตฺถิปจฺจเยน ปจฺจโย. สหชาตํ, ปจฺฉาชาตํ, อาหารํ, อินฺทฺริยํ.  สหชาโต อนิทสฺสโน เอโก ขนฺโธ ติณฺณนฺนํ ขนฺธานํ สนิทสฺส\-นานญฺจ อนิทสฺสนานญฺจ จิตฺต\-สมุฏฺฐานานํ รูปานํ อตฺถิปจฺจเยน ปจฺจโย ฯเปฯ เทฺว ขนฺธา ฯเปฯ. ปฏิสนฺธิ\-กฺขเณ ฯเปฯ มหาภูตา สนิทสฺส\-นานญฺจ อนิทสฺสนานญฺจ จิตฺต\-สมุฏฺฐานานํ รูปานํ, กฏตฺตารูปานํ, อุปาทารูปานํ อตฺถิปจฺจเยน ปจฺจโย. พาหิรา. อาหารสมุฏฺฐานา. อุตุสมุฏฺฐานา. อสญฺญสตฺตานํ มหาภูตา สนิทสฺส\-นานญฺจ อนิทสฺสนานญฺจ กฏตฺตารูปานํ, อุปาทารูปานํ อตฺถิปจฺจเยน ปจฺจโย ฯเปฯ. (3)}

\proseitem{47}{สนิทสฺสโน จ อนิทสฺสโน จ ธมฺมา อนิทสฺสนสฺส ธมฺมสฺส อตฺถิปจฺจเยน ปจฺจโย. ปุเรชาตํ รูปายตนญฺจ วตฺถุ จ อนิทสฺสนานํ ขนฺธานํ อตฺถิปจฺจเยน ปจฺจโย. รูปายตนญฺจ จกฺขายตนญฺจ จกฺขุ\-วิญฺญาณสฺส อตฺถิปจฺจเยน ปจฺจโย.}

\prose{นตฺถิ\-ปจฺจเยน ปจฺจโย. วิคต\-ปจฺจเยน ปจฺจโย. อวิคต\-ปจฺจเยน ปจฺจโย.}

\csromanmarkh{1. ปจฺจยานุโลม}\csromanmarkhii{2. สงฺขฺยาวาร}\csromanheadernomark{2}{1. \textbf{ปจฺจยานุโลม 2. สงฺขฺยาวาร}}
\proseitem{46}{เหตุยา ตีณิ, อารมฺมเณ เทฺว, อธิปติยา จตฺตาริ, อนนฺตเร เอกํ, สมนนฺตเร เอกํ, สหชาเต ตีณิ, อญฺญมญฺเญ เอกํ, นิสฺสเย ตีณิ, อุปนิสฺสเย เทฺว, ปุเรชาเต ตีณิ, ปจฺฉาชาเต ตีณิ, อาเสวเน เอกํ, กมฺเม ตีณิ, วิปาเก ตีณิ, อาหาเร ตีณิ, \csromanfolio{103}อินฺทฺริเย ตีณิ, ฌาเน ตีณิ, มคฺเค ตีณิ, สมฺปยุตฺเต เอกํ, วิปฺปยุตฺเต ตีณิ, อตฺถิยา ปญฺจ, นตฺถิยา เอกํ, วิคเต เอกํ, อวิคเต ปญฺจ. (เอวํ คเณตพฺพํ.)}

\vspace{\tipitakaparskip}
\csromancenter{\textbf{อนุโลมํ.}}
\proseitem{49}{สนิทสฺสโน ธมฺโม อนิทสฺสนสฺส ธมฺมสฺส อารมฺมณ\-ปจฺจเยน ปจฺจโย. อุปนิสฺสายปจฺจเยน ปจฺจโย. (1)}

\prose{อนิทสฺสโน ธมฺโม อนิทสฺสนสฺส ธมฺมสฺส อารมฺมณ\-ปจฺจเยน ปจฺจโย. สหชาต\-ปจฺจเยน ปจฺจโย. อุปนิสฺสย\-ปจฺจเยน ปจฺจโย. ปุเรชาต\-ปจฺจเยน ปจฺจโย. ปจฺฉาชาต\-ปจฺจเยน ปจฺจโย. กมฺมปจฺจเยน ปจฺจโย. อาหารปจฺจเยน ปจฺจโย. อินฺทฺริย\-ปจฺจเยน ปจฺจโย. (1)}

\prose{อนิทสฺสโน ธมฺโม สนิทสฺสนสฺส ธมฺมสฺส สหชาต\-ปจฺจเยน ปจฺจโย. ปจฺฉาชาต\-ปจฺจเยน ปจฺจโย. กมฺมปจฺจเยน ปจฺจโย. อาหารปจฺจเยน ปจฺจโย. อินฺทฺริย\-ปจฺจเยน ปจฺจโย. (2)}

\prose{อนิทสฺสโน ธมฺโม สนิทสฺสนสฺส จ อนิทสฺสนสฺส จ ธมฺมสฺส สหชาต\-ปจฺจเยน ปจฺจโย. ปจฺฉาชาต\-ปจฺจเยน ปจฺจโย. กมฺมปจฺจเยน ปจฺจโย. อาหารปจฺจเยน ปจฺจโย. อินฺทฺริย\-ปจฺจเยน ปจฺจโย. (3)}

\prose{สนิทสฺสโน จ อนิทสฺสโน จ ธมฺมา อนิทสฺสนสฺส ธมฺมสฺส ปุเรชาต\-ปจฺจเยน ปจฺจโย. (1)}

\csromanmarkh{2. ปจฺจย\-ปจฺจนีย}\csromanmarkhii{2. สงฺขฺยาวาร}\csromanheadernomark{2}{2. ปจฺจย\-ปจฺจนีย 2. สงฺขฺยาวาร}
\proseitem{48}{นเหตุยา ปญฺจ, นอารมฺมเณ จตฺตาริ, นอธิปติยา ปญฺจ, นอนนฺตเร ปญฺจ, นสมนนฺตเร ปญฺจ, นสหชาเต ปญฺจ, นอญฺญมญฺเญ ปญฺจ, นนิสฺสเย จตฺตาริ, นอุปนิสฺสเย ปญฺจ, นปุเรชาเต จตฺตาริ, นปจฺฉาชาเต ปญฺจ, (สพฺพตฺถ ปญฺจ.) นสมฺปยุตฺเต ปญฺจ, นวิปฺปยุตฺเต จตฺตาริ, โนอตฺถิยา จตฺตาริ, โนนตฺถิยา ปญฺจ, โนวิคเต ปญฺจ, โนอวิคเต จตฺตาริ.}

\vspace{\tipitakaparskip}
\csromancenter{ปจฺจนียํ.}
\csromancenter{\textbf{เหตุ}ปจฺจยา นอารมฺมเณ ตีณิ, นอธิปติยา ตีณิ, นอนนฺตเร ตีณิ, นสมนนฺตเร ตีณิ, นอญฺญมญฺเญ ตีณิ, นฺเอาปนิสฺสเย ตีณิ ฯเปฯ นสมฺปยุตฺเต ตีณิ, นวิปฺปยุตฺเต เอกํ, โนนตฺถิยา ตีณิ, โนวิคเต ตีณิ.}
\csromancenter{อนุโลม\-ปจฺจนียํ.}
\csromanheader{2}{4. \textbf{ปจฺจปจฺจนียานุโลม}}
\proseitem{52}{\csromanfolio{104}นเหตุปจฺจยา อารมฺมเณ เทฺว, อธิปติยา จตฺตาริ, อนนฺตเร เอกํ, สมนนฺตเร เอกํ, สหชาเต ตีณิ, อญฺญมญฺเญ เอกํ, นิสฺสเย ตีณิ, อุปนิสฺสเย เทฺว, ปุเรชาเต ตีณิ, ปจฺฉาชาเต ตีณิ, อาเสวเน เอกํ, กมฺเม ตีณิ ฯเปฯ มคฺเค ตีณิ, สมฺปยุตฺเต เอกํ, วิปฺปยุตฺเต ตีณิ, อตฺถิยา ปญฺจ, นตฺถิยา เอกํ, วิคเต เอกํ, อวิคเต ปญฺจ.}

\vspace{\tipitakaparskip}
\csromancenter{ปจฺจนียา\-นุโลมํ.}
\nitthitam{สนิทสฺสนทุกํ นิฏฺฐิตํ.}
\csromanmatikaanchor{mtk.13}
\csromantocmarkref{section}{10. สปฺปฏิฆทุก}{mtk.13}
\bookmark[dest={mtk.13},level=1]{10. สปฺปฏิฆทุก}
\csromanmarkh{10. สปฺปฏิฆทุก}\csromanmarkhii{1. ปฏิจฺจวาร}\csromanheadernomark{1}{10. \textbf{สปฺปฏิฆทุก 1. ปฏิจฺจวาร}}
\csromanmarkh{1. ปจฺจยานุโลม}\csromanmarkhii{1. วิภงฺควาร}\csromanheadernomark{2}{1. \textbf{ปจฺจยานุโลม 1. วิภงฺควาร}}
\csromanheader{2}{\textbf{เหตุ}}
\proseitem{53}{สปฺปฏิฆํ ธมฺมํ ปฏิจฺจ สปฺปฏิโฆ ธมฺโม อุปฺปชฺชติ เหตุปจฺจยา. สปฺปฏิฆํ เอกํ มหาภูตํ ปฏิจฺจ เทฺว มหาภูตา, เทฺว มหาภูเต ปฏิจฺจ เอกํ \csromanfolio{105}มหาภูตํ. สปฺปฏิเฆ มหาภูเต ปฏิจฺจ สปฺปฏิฆํ จิตฺต\-สมุฏฺฐานํ รูปํ, กฏตฺตารูปํ, อุปาทารูปํ. โผฏฺฐพฺพา\-ยตนํ ปฏิจฺจ จกฺขายตนํ ฯเปฯ รสายตนํ. (1)}

\prose{สปฺปฏิฆํ ธมฺมํ ปฏิจฺจ อปฺปฏิโฆ ธมฺโม อุปฺปชฺชติ เหตุปจฺจยา. สปฺปฏิเฆ มหาภูเต ปฏิจฺจ อาโปธาตุ. สปฺปฏิเฆ มหาภูเต ปฏิจฺจ อปฺปฏิฆํ จิตฺต\-สมุฏฺฐานํ รูปํ, กฏตฺตารูปํ, อุปาทารูปํ. โผฏฺฐพฺพา\-ยตนํ ปฏิจฺจ อาโปธาตุ. อิตฺถินฺทฺริยํ ฯเปฯ กพฬีกาโร อาหาโร. (2)}

\prose{สปฺปฏิฆํ ธมฺมํ ปฏิจฺจ สปฺปฏิโฆ จ อปฺปฏิโฆ จ ธมฺมา อุปฺปชฺชนฺติ เหตุปจฺจยา. สปฺปฏิฆํ เอกํ มหาภูตํ ปฏิจฺจ เทฺว มหาภูตา อาโปธาตุ จ, เทฺว มหาภูเต ฯเปฯ สปฺปฏิเฆ มหาภูเต ปฏิจฺจ สปฺปฏิฆญฺจ อปฺปฏิฆญฺจ จิตฺต\-สมุฏฺฐานํ รูปํ, กฏตฺตารูปํ, อุปาทารูปํ. โผฏฺฐพฺพา\-ยตนํ ปฏิจฺจ จกฺขายตนํ ฯเปฯ รสายตนํ อาโปธาตุ อิตฺถินฺทฺริยํ ฯเปฯ กพฬีกาโร อาหาโร. (3)}

\proseitem{54}{อปฺปฏิฆํ ธมฺมํ ปฏิจฺจ อปฺปฏิโฆ ธมฺโม อุปฺปชฺชติ เหตุปจฺจยา. อปฺปฏิฆํ เอกํ ขนฺธํ ปฏิจฺจ ตโย ขนฺธา อปฺปฏิฆํ จิตฺต\-สมุฏฺฐานญฺจ รูปํ ฯเปฯ เทฺว ขนฺเธ ฯเปฯ. ปฏิสนฺธิ\-กฺขเณ ฯเปฯ. ขนฺเธ ปฏิจฺจ วตฺถุ, วตฺถุํ ปฏิจฺจ ขนฺธา. อาโปธาตุํ ปฏิจฺจ อปฺปฏิฆํ จิตฺต\-สมุฏฺฐานํ รูปํ, กฏตฺตารูปํ, อุปาทารูปํ. อาโปธาตุํ ปฏิจฺจ อิตฺถินฺทฺริยํ ฯเปฯ กพฬีกาโร อาหาโร. (1)}

\prose{อปฺปฏิฆํ ธมฺมํ ปฏิจฺจ สปฺปฏิโฆ ธมฺโม อุปฺปชฺชติ เหตุปจฺจยา. อปฺปฏิเฆ ขนฺเธ ปฏิจฺจ สปฺปฏิฆํ จิตฺต\-สมุฏฺฐานํ รูปํ. ปฏิสนฺธิ\-กฺขเณ ฯเปฯ. อาโปธาตุํ ปฏิจฺจ สปฺปฏิฆา มหาภูตา. อาโปธาตุํ ปฏิจฺจ สปฺปฏิฆํ จิตฺต\-สมุฏฺฐานํ รูปํ, กฏตฺตารูปํ, อุปาทารูปํ. อาโปธาตุํ ปฏิจฺจ จกฺขายตนํ ฯเปฯ โผฏฺฐพฺพา\-ยตนํ.}

\prose{อปฺปฏิฆํ ธมฺมํ ปฏิจฺจ สปฺปฏิโฆ จ อปฺปฏิโฆ จ ธมฺมา อุปฺปชฺชนฺติ เหตุปจฺจยา. อปฺปฏิฆํ เอกํ ขนฺธํ ปฏิจฺจ ตโย ขนฺธา สปฺปฏิฆญฺจ อปฺปฏิฆญฺจ จิตฺต\-สมุฏฺฐานํ รูปํ ฯเปฯ เทฺว ขนฺเธ ฯเปฯ. ปฏิสนฺธิ\-กฺขเณ ฯเปฯ. อาโปธาตุํ ปฏิจฺจ สปฺปฏิฆญฺจ อปฺปฏิฆญฺจ จิตฺต\-สมุฏฺฐานํ รูปํ, กฏตฺตารูปํ, อุปาทารูปํ. อาโปธาตุํ ปฏิจฺจ จกฺขายตนํ ฯเปฯ โผฏฺฐพฺพา\-ยตนํ, อิตฺถินฺทฺริยํ ฯเปฯ กพฬีกาโร อาหาโร. (3)}

\proseitem{55}{\csromanfolio{106}สปฺปฏิฆญฺจ อปฺปฏิฆญฺจ ธมฺมํ ปฏิจฺจ สปฺปฏิโฆ ธมฺโม อุปฺปชฺชติ เหตุปจฺจยา. อปฺปฏิเฆ ขนฺเธ จ มหาภูเต จ ปฏิจฺจ สปฺปฏิฆํ จิตฺต\-สมุฏฺฐานํ รูปํ. ปฏิสนฺธิ\-กฺขเณ ฯเปฯ. สปฺปฏิฆํ เอกํ มหาภูตญฺจ อาโปธาตุญฺจ ปฏิจฺจ เทฺว มหาภูตา ฯเปฯ สปฺปฏิเฆ มหาภูเต จ อาโปธาตุญฺจ ปฏิจฺจ สปฺปฏิฆํ จิตฺต\-สมุฏฺฐานํ รูปํ, กฏตฺตารูปํ, อุปาทารูปํ. โผฏฺฐพฺพายตนญฺจ อาโปธาตุญฺจ ปฏิจฺจ จกฺขายตนํ ฯเปฯ รสายตนํ. (1)}

\prose{สปฺปฏิฆญฺจ อปฺปฏิฆญฺจ ธมฺมํ ปฏิจฺจ อปฺปฏิโฆ ธมฺโม อุปฺปชฺชติ เหตุปจฺจยา. อปฺปฏิเฆ ขนฺเธ จ มหาภูเต จ ปฏิจฺจ อปฺปฏิฆํ จิตฺต\-สมุฏฺฐานํ รูปํ. ปฏิสนฺธิ\-กฺขเณ อปฺปฏิเฆ ขนฺเธ จ มหาภูเต จ ปฏิจฺจ อปฺปฏิฆํ กฏตฺตารูปํ. โผฏฺฐพฺพายตนญฺจ อาโปธาตุญฺจ ปฏิจฺจ อปฺปฏิฆํ จิตฺต\-สมุฏฺฐานํ รูปํ, กฏตฺตารูปํ, อุปาทารูปํ. โผฏฺฐพฺพายตนญฺจ อาโปธาตุญฺจ ปฏิจฺจ อิตฺถินฺทฺริยํ ฯเปฯ กพฬีกาโร อาหาโร. (2)}

\prose{สปฺปฏิฆญฺจ อปฺปฏิฆญฺจ ธมฺมํ ปฏิจฺจ สปฺปฏิโฆ จ อปฺปฏิโฆ จ ธมฺมา อุปฺปชฺชนฺติ เหตุปจฺจยา. อปฺปฏิเฆ ขนฺเธ จ มหาภูเต จ ปฏิจฺจ สปฺปฏิฆญฺจ อปฺปฏิฆญฺจ จิตฺต\-สมุฏฺฐานํ รูปํ. ปฏิสนฺธิ\-กฺขเณ ฯเปฯ. อปฺปฏิเฆ ขนฺเธ จ มหาภูเต จ ปฏิจฺจ สปฺปฏิฆญฺจ อปฺปฏิฆญฺจ กฏตฺตารูปํ. โผฏฺฐพฺพายตนญฺจ อาโปธาตุญฺจ ปฏิจฺจ สปฺปฏิฆญฺจ อปฺปฏิฆญฺจ จิตฺต\-สมุฏฺฐานํ รูปํ, กฏตฺตารูปํ, อุปาทารูปํ. โผฏฺฐพฺพายตนญฺจ อาโปธาตุญฺจ ปฏิจฺจ จกฺขายตนํ ฯเปฯ รสายตนํ, อิตฺถินฺทฺริยํ ฯเปฯ กพฬีกาโร อาหาโร. (3)}

\proseitem{56}{อปฺปฏิฆํ ธมฺมํ ปฏิจฺจ อปฺปฏิโฆ ธมฺโม อุปฺปชฺชติ อารมฺมณ\-ปจฺจยา. อปฺปฏิฆํ เอกํ ขนฺธํ ปฏิจฺจ ตโย ขนฺธา ฯเปฯ เทฺว ขนฺเธ ฯเปฯ. ปฏิสนฺธิ\-กฺขเณ ฯเปฯ. วตฺถุํ ปฏิจฺจ ขนฺธา.}

\csromanheader{2}{\textbf{อธิปตฺยาทิ}}
\proseitem{57}{สปฺปฏิฆํ ธมฺมํ ปฏิจฺจ อปฺปฏิโฆ ธมฺโม อุปฺปชฺชติ อธิปติ\-ปจฺจยา. (ปฏิสนฺธิ วชฺเชตพฺพา, กฏตฺตารูปา จ.) อนนฺตร\-ปจฺจยา. สมนนฺตร\-ปจฺจยา. สหชาต\-ปจฺจยา. (สพฺเพ มหาภูตา กาตพฺพา.) อญฺญมญฺญ\-ปจฺจยา. สปฺปฏิฆํ เอกํ มหาภูตํ ปฏิจฺจ เทฺว มหาภูตา. เทฺว มหาภูเต ฯเปฯ. (1)}

\prose{\csromanfolio{107}สปฺปฏิฆํ ธมฺมํ ปฏิจฺจ อปฺปฏิโฆ ธมฺโม อุปฺปชฺชติ อญฺญมญฺญ\-ปจฺจยา. สปฺปฏิเฆ มหาภูเต ปฏิจฺจ อาโปธาตุ ฯเปฯ. (2)}

\prose{สปฺปฏิฆํ ธมฺมํ ปฏิจฺจ สปฺปฏิโฆ จ อปฺปฏิโฆ จ ธมฺมา อุปฺปชฺชนฺติ อญฺญมญฺญ\-ปจฺจยา. สปฺปฏิฆํ เอกํ มหาภูตํ ปฏิจฺจ เทฺว มหาภูตา อาโปธาตุ จ. เทฺว มหาภูเต ฯเปฯ. (3)}

\proseitem{58}{อปฺปฏิฆํ ธมฺมํ ปฏิจฺจ อปฺปฏิโฆ ธมฺโม อุปฺปชฺชติ อญฺญมญฺญ\-ปจฺจยา. อปฺปฏิฆํ เอกํ ขนฺธํ ปฏิจฺจ ตโย ขนฺธา ฯเปฯ เทฺว ขนฺเธ ฯเปฯ. ปฏิสนฺธิ\-กฺขเณ ฯเปฯ. ขนฺเธ ปฏิจฺจ วตฺถุ, วตฺถุํ ปฏิจฺจ ขนฺธา. (1)}

\prose{อปฺปฏิฆํ ธมฺมํ ปฏิจฺจ สปฺปฏิโฆ ธมฺโม อุปฺปชฺชติ อญฺญมญฺญ\-ปจฺจยา. อาโปธาตุํ ปฏิจฺจ สปฺปฏิฆา มหาภูตา. (อิเม อชฺฌตฺติก\-พาหิรา มหาภูตา กาตพฺพา.) (2)}

\prose{สปฺปฏิฆญฺจ อปฺปฏิฆญฺจ ธมฺมํ ปฏิจฺจ สปฺปฏิโฆ ธมฺโม อุปฺปชฺชติ อญฺญมญฺญ\-ปจฺจยา. สปฺปฏิฆํ เอกํ มหาภูตญฺจ อาโปธาตุญฺจ ปฏิจฺจ เทฺว มหาภูตา ฯเปฯ. นิสฺสยปจฺจยา ฯเปฯ อวิคตปจฺจยา.}

\csromanmarkh{1. ปจฺจยานุโลม}\csromanmarkhii{2. สงฺขฺยาวาร}\csromanheadernomark{2}{1. \textbf{ปจฺจยานุโลม 2. สงฺขฺยาวาร}}
\proseitem{59}{เหตุยา นว, อารมฺมเณ เอกํ, อธิปติยา นว, อนนฺตเร เอกํ, สมนนฺตเร เอกํ, สหชาเต นว, อญฺญมญฺเญ ฉ, นิสฺสเย นว, อุปนิสฺสเย เอกํ, ปุเรชาเต เอกํ, อาเสวเน เอกํ, กมฺเม นว, วิปาเก นว, อาหาเร นว, อินฺทฺริเย นว, ฌาเน นว, มคฺเค นว, สมฺปยุตฺเต เอกํ, วิปฺปยุตฺเต นว, อตฺถิยา นว, นตฺถิยา เอกํ, วิคเต เอกํ, อวิคเต นว.}

\vspace{\tipitakaparskip}
\csromancenter{อนุโลมํ.}
\csromanmarkh{2. ปจฺจย\-ปจฺจนีย}\csromanmarkhii{1. วิภงฺควาร}\csromanheadernomark{2}{2. ปจฺจย\-ปจฺจนีย 1. วิภงฺควาร}
\proseitem{60}{\csromanfolio{108}สปฺปฏิฆํ ธมฺมํ ปฏิจฺจ สปฺปฏิโฆ ธมฺโม อุปฺปชฺชติ นเหตุปจฺจยา ตีณิ.}

\prose{อปฺปฏิฆํ ธมฺมํ ปฏิจฺจ อปฺปฏิโฆ ธมฺโม อุปฺปชฺชติ นเหตุปจฺจยา. อเหตุกํ อปฺปฏิฆํ เอกํ ขนฺธํ ปฏิจฺจ ตโย ขนฺธา อปฺปฏิฆํ จิตฺต\-สมุฏฺฐานญฺจ รูปํ ฯเปฯ เทฺว ขนฺเธ ฯเปฯ. อเหตุก\-ปฏิสนฺธิ\-กฺขเณ ฯเปฯ. ขนฺเธ ปฏิจฺจ วตฺถุ, วตฺถุํ ปฏิจฺจ ขนฺธา. อาโปธาตุํ ปฏิจฺจ อปฺปฏิฆํ จิตฺต\-สมุฏฺฐานํ รูปํ, กฏตฺตารูปํ, อุปาทารูปํ. อาโปธาตุํ ปฏิจฺจ อิตฺถินฺทฺริยํ ฯเปฯ กพฬีกาโร อาหาโร. พาหิรํ. อาหารสมุฏฺฐานํ. อุตุสมุฏฺฐานํ. อสญฺญสตฺตานํ. อาโปธาตุํ ปฏิจฺจ อปฺปฏิฆํ กฏตฺตารูปํ, อุปาทารูปํ. วิจิกิจฺฉา\-สหคเต อุทฺธจฺจ\-สหคเต ขนฺเธ ปฏิจฺจ วิจิกิจฺฉา\-สหคโต อุทฺธจฺจ\-สหคโต โมโห. (1)}

\prose{(อปฺปฏิฆ\-มูลกํ อิตเรปิ เทฺว ปญฺหา กาตพฺพา. ฆฏเนปิ ตีณิ ปญฺหา กาตพฺพา. อชฺฌตฺติกา พาหิรา มหาภูตา สพฺเพ ชานิตฺวา กาตพฺพา.)}

\prose{นอารมฺมณาทิ}

\proseitem{61}{สปฺปฏิฆํ ธมฺมํ ปฏิจฺจ สปฺปฏิโฆ ธมฺโม อุปฺปชฺชติ นอารมฺมณ\-ปจฺจยา. (สพฺพํ ขิตฺตํ.) โนวิคต\-ปจฺจยา.}

\csromanmarkh{2. ปจฺจย\-ปจฺจนีย}\csromanmarkhii{2. สงฺขฺยาวาร}\csromanheadernomark{2}{2. ปจฺจย\-ปจฺจนีย 2. สงฺขฺยาวาร}
\vspace{\tipitakaparskip}
\csromancenter{\textbf{นเหตุ}ยา นว, นอารมฺมเณ นว, นอธิปติยา นว, นอนนฺตเร นว, นสมนนฺตเร นว, นอญฺญมญฺเญ นว, นฺเอาปนิสฺสเย นว, นปุเรชาเต นว, นปจฺฉาชาเต นว, นอาเสวเน นว, นกมฺเม นว, นวิปาเก นว, นอาหาเร นว, ไนนฺทฺริเย นว, นฌาเน นว, นมคฺเค นว, นสมฺปยุตฺเต นว, นวิปฺปยุตฺเต นว, โนนตฺถิยา นว, โนวิคเต นว.}
\csromancenter{ปจฺจนียํ.}
\csromanmarkh{10. สปฺปฏิฆทุก}\csromanmarkhii{3. ปจฺจยา\-นุโลม\-ปจฺจนีย}\csromanheadernomark{2}{10. สปฺปฏิฆทุก 3. ปจฺจยา\-นุโลม\-ปจฺจนีย}
\vspace{\tipitakaparskip}
\csromancenter{\textbf{เหตุ}ปจฺจยา นอารมฺมเณ นว, นอธิปติยา นว, นอนนฺตเร นว, นสมนนฺตเร นว, นอญฺญมญฺเญ นว, นฺเอาปนิสฺสเย นว, นปุเรชาเต นว, นปจฺฉาชาเต นว, นอาเสวเน นว, นกมฺเม เอกํ, นวิปาเก นว, นสมฺปยุตฺเต นว, นวิปฺปยุตฺเต เอกํ, โนนตฺถิยา นว, โนวิคเต นว.}
\csromancenter{อนุโลม\-ปจฺจนียํ.}
\proseitem{64}{\csromanfolio{109}นเหตุปจฺจยา อารมฺมเณ เอกํ, อนนฺตเร เอกํ, สมนนฺตเร เอกํ, สหชาเต นว, อญฺญมญฺเญ ฉ, นิสฺสเย นว, อุปนิสฺสเย เอกํ, ปุเรชาเต เอกํ, อาเสวเน เอกํ, กมฺเม นว ฯเปฯ มคฺเค เอกํ, สมฺปยุตฺเต เอกํ, วิปฺปยุตฺเต นว, อตฺถิยา นว, นตฺถิยา เอกํ, วิคเต เอกํ, อวิคเต นว.}

\vspace{\tipitakaparskip}
\csromancenter{ปจฺจนียา\-นุโลมํ.}
\csromancenter{(สหชาตวาโรปิ ปฏิจฺจ\-วาร\-สทิโส.)}
\csromanmarkh{10. สปฺปฏิฆทุก}\csromanmarkhii{3. ปจฺจยวาร}\csromanheadernomark{2}{10. สปฺปฏิฆทุก 3. ปจฺจยวาร}
\csromanmarkh{1. ปจฺจยานุโลม}\csromanmarkhii{1. วิภงฺควาร}\csromanheadernomark{2}{1. \textbf{ปจฺจยานุโลม 1. วิภงฺควาร}}
\csromanheader{2}{\textbf{เหตุ}}
\proseitem{65}{สปฺปฏิฆํ ธมฺมํ ปจฺจยา สปฺปฏิโฆ ธมฺโม อุปฺปชฺชติ เหตุปจฺจยา. ตีณิ.}

\prose{\csromanfolio{110}อปฺปฏิฆํ ธมฺมํ ปจฺจยา อปฺปฏิโฆ ธมฺโม อุปฺปชฺชติ เหตุปจฺจยา. อปฺปฏิฆํ เอกํ ขนฺธํ ปจฺจยา ตโย ขนฺธา อปฺปฏิฆํ จิตฺต\-สมุฏฺฐานญฺจ รูปํ ฯเปฯ เทฺว ขนฺเธ ฯเปฯ. ปฏิสนฺธิ\-กฺขเณ ฯเปฯ. อาโปธาตุํ ปจฺจยา อปฺปฏิฆํ จิตฺต\-สมุฏฺฐานํ รูปํ กฏตฺตารูปํ. อุปาทารูปํ. อาโปธาตุํ ปจฺจยา อิตฺถินฺทฺริยํ ฯเปฯ กพฬีกาโร อาหาโร. วตฺถุํ ปจฺจยา อปฺปฏิฆา ขนฺธา. (1) (อวเสสา ปญฺจ ปญฺหา ปฏิจฺจ\-วาร\-สทิสา.)}

\csromanheader{2}{\textbf{อารมฺมณาทิ}}
\proseitem{66}{สปฺปฏิฆํ ธมฺมํ ปจฺจยา อปฺปฏิโฆ ธมฺโม อุปฺปชฺชติ อารมฺมณ\-ปจฺจยา. จกฺขายตนํ ปจฺจยา จกฺขุวิญฺญาณํ ฯเปฯ กายายตนํ ปจฺจยา กายวิญฺญาณํ. (1)}

\prose{อปฺปฏิฆํ ธมฺมํ ปจฺจยา อปฺปฏิโฆ ธมฺโม อุปฺปชฺชติ อารมฺมณ\-ปจฺจยา. อปฺปฏิฆํ เอกํ ขนฺธํ ปจฺจยา ตโย ขนฺธา ฯเปฯ เทฺว ขนฺเธ ฯเปฯ. ปฏิสนฺธิ\-กฺขเณ ฯเปฯ. วตฺถุํ ปจฺจยา อปฺปฏิฆา ขนฺธา. (1)}

\prose{สปฺปฏิฆญฺจ อปฺปฏิฆญฺจ ธมฺมํ ปจฺจยา อปฺปฏิโฆ ธมฺโม อุปฺปชฺชติ อารมฺมณ\-ปจฺจยา. จกฺขุ\-วิญฺญาณสหคตํ เอกํ ขนฺธญฺจ จกฺขายตนญฺจ ปจฺจยา ตโย ขนฺธา ฯเปฯ. กาย\-วิญฺญาณสหคตํ เอกํ ขนฺธญฺจ กายายตนญฺจ ปจฺจยา ตโย ขนฺธา ฯเปฯ. อธิปติ\-ปจฺจยา. (สํขิตฺตํ.) อวิคตปจฺจยา.}

\csromanmarkh{1. ปจฺจยานุโลม}\csromanmarkhii{2. สงฺขฺยาวาร}\csromanheadernomark{2}{1. \textbf{ปจฺจยานุโลม 2. สงฺขฺยาวาร}}
\proseitem{67}{เหตุยา นว, อารมฺมเณ ตีณิ, อธิปติยา นว, อนนฺตเร ตีณิ, สมนนฺตเร ตีณิ, สหชาเต นว, อญฺญมญฺเญ ฉ, นิสฺสเย นว, อุปนิสฺสเย ตีณิ, ปุเรชาเต ตีณิ, อาเสวเน เอกํ, กมฺเม นว ฯเปฯ อวิคเต นว. (เอวํ ปจฺจนีย\-คณนาปิ กาตพฺพา.) (นิสฺสยวาโรปิ ปจฺจยวาร\-สทิโส.)}

\prose{(สํสฏฺฐ\-วาเรปิ สพฺพตฺถ เอกํ. สํขิตฺตํ. อวิคตปจฺจยา, เอกาเยว ปญฺหา. เทฺวปิ วารา กาตพฺพา.)}

\csromanmarkh{10. สปฺปฏิฆทุก}\csromanmarkhii{10. สปฺปฏิฆทุก 7. ปญฺหาวาร}\csromanheadernomark{2}{10. สปฺปฏิฆทุก \textbf{10. สปฺปฏิฆทุก 7. ปญฺหาวาร}}
\csromanmarkh{1. ปจฺจยานุโลม}\csromanmarkhii{1. วิภงฺควาร}\csromanheadernomark{2}{1. \textbf{ปจฺจยานุโลม 1. วิภงฺควาร}}
\csromanheader{2}{\textbf{เหตุ}}
\proseitem{68}{\csromanfolio{111}อปฺปฏิโฆ ธมฺโม อปฺปฏิฆสฺส ธมฺมสฺส เหตุปจฺจเยน ปจฺจโย. อปฺปฏิฆา เหตู สมฺปยุตฺตกานํ ขนฺธานํ อปฺปฏิฆานญฺจ จิตฺต\-สมุฏฺฐานานํ รูปานํ เหตุปจฺจเยน ปจฺจโย. ปฏิสนฺธิ\-กฺขเณ ฯเปฯ. (1)}

\prose{อปฺปฏิโฆ ธมฺโม สปฺปฏิฆสฺส ธมฺมสฺส เหตุปจฺจเยน ปจฺจโย. อปฺปฏิฆา เหตู สปฺปฏิฆานํ จิตฺต\-สมุฏฺฐานานํ รูปานํ เหตุปจฺจเยน ปจฺจโย. ปฏิสนฺธิ\-กฺขเณ ฯเปฯ. (2)}

\prose{อปฺปฏิโฆ ธมฺโม สปฺปฏิฆสฺส จ อปฺปฏิฆสฺส จ ธมฺมสฺส เหตุปจฺจเยน ปจฺจโย. อปฺปฏิฆา เหตู สมฺปยุตฺตกานํ ขนฺธานํ สปฺปฏิฆานญฺจ อปฺปฏิฆานญฺจ จิตฺต\-สมุฏฺฐานานํ รูปานํ เหตุปจฺจเยน ปจฺจโย. ปฏิสนฺธิ\-กฺขเณ ฯเปฯ. (3)}

\proseitem{69}{สปฺปฏิโฆ ธมฺโม อปฺปฏิฆสฺส ธมฺมสฺส อารมฺมณ\-ปจฺจเยน ปจฺจโย. จกฺขุํ ฯเปฯ. โผฏฺฐพฺเพ อนิจฺจโต ฯเปฯ โทมนสฺสํ อุปฺปชฺชติ. ทิพฺเพน จกฺขุนา รูปํ ปสฺสติ. ทิพฺพาย โสตธาตุยา สทฺทํ สุณาติ. รูปายตนํ จกฺขุ\-วิญฺญาณสฺส ฯเปฯ. โผฏฺฐพฺพา\-ยตนํ กายวิญฺญาณสฺส ฯเปฯ. สปฺปฏิฆา ขนฺธา อิทฺธิวิธ\-ญาณสฺส, ปุพฺเพ\-นิวาสา\-นุสฺสติ\-ญาณสฺส, อนาคตํสญาณสฺส, อาวชฺชนาย อารมฺมณ\-ปจฺจเยน ปจฺจโย. (1)}

\prose{อปฺปฏิโฆ ธมฺโม อปฺปฏิฆสฺส ธมฺมสฺส อารมฺมณ\-ปจฺจเยน ปจฺจโย. ทานํ ฯเปฯ สีลํ ฯเปฯ อุโปสถกมฺมํ กตฺวา ตํ ปจฺจเวกฺขติ. ปุพฺเพ สุจิณฺณานิ ปจฺจเวกฺขติ. ฌานา วุฏฺฐหิตฺวา ฌานํ ปจฺจเวกฺขติ. อริยา มคฺคา วุฏฺฐหิตฺวา มคฺคํ ปจฺจเวกฺขนฺติ, ผลํ ปจฺจเวกฺขนฺติ, นิพฺพานํ ปจฺจเวกฺขนฺติ. นิพฺพานํ โคตฺรภุสฺส, โวทานสฺส, มคฺคสฺส, ผลสฺส, อาวชฺชนาย อารมฺมณ\-ปจฺจเยน ปจฺจโย. อริยา ปหีเน กิเลเส ปจฺจเวกฺขนฺติ, วิกฺขมฺภิเต กิเลเส ปจฺจเวกฺขนฺติ. ปุพฺเพ สมุทาจิณฺเณ กิเลเส ชานนฺติ. วตฺถุํ ฯเปฯ. อิตฺถินฺทฺริยํ. ปุริสินฺทฺริยํ. ชีวิตินฺทฺริยํ. อาโปธาตุํ กพฬีการํ อาหารํ อนิจฺจโต ฯเปฯ โทมนสฺสํ อุปฺปชฺชติ. เจโตปริย\-ญาเณน อปฺปฏิฆ\-จิตฺต\-สมงฺคิสฺส จิตฺตํ ชานาติ. \csromanfolio{112}อากาสา\-นญฺจา\-ยตนํ วิญฺญาณญฺจา\-ยตนสฺส ฯเปฯ. อากิญฺจญฺญา\-ยตนํ เนว\-สญฺญา\-นา\-สญฺญา\-ยตนสฺส ฯเปฯ. อปฺปฏิฆา ขนฺธา อิทฺธิวิธ\-ญาณสฺส, เจโตปริย\-ญาณสฺส, ปุพฺเพ\-นิวาสา\-นุสฺสติ\-ญาณสฺส ยถากมฺมูปคญาณสฺส, อนาคตํสญาณสฺส, อาวชฺชนาย อารมฺมณ\-ปจฺจเยน ปจฺจโย. (1)}

\proseitem{70}{สปฺปฏิโฆ ธมฺโม อปฺปฏิฆสฺส ธมฺมสฺส อธิปติ\-ปจฺจเยน ปจฺจโย. อารมฺมณา\-ธิปติ\csromandash{}จกฺขุํ ฯเปฯ โผฏฺฐพฺเพ ครุํ กตฺวา อสฺสาเทติ อภินนฺทติ, ตํ ครุํ กตฺวา ราโค อุปฺปชฺชติ, ทิฏฺฐิ อุปฺปชฺชติ. (1)}

\prose{อปฺปฏิโฆ ธมฺโม อปฺปฏิฆสฺส ธมฺมสฺส อธิปติ\-ปจฺจเยน ปจฺจโย. อารมฺมณา\-ธิปติ, สหชาตา\-ธิปติ.  อารมฺมณา\-ธิปติ\csromandash{}ทานํ ฯเปฯ สีลํ ฯเปฯ อุโปสถกมฺมํ กตฺวา ตํ ครุํ กตฺวา ปจฺจเวกฺขติ. ปุพฺเพ สุจิณฺณานิ ครุํ กตฺวา ปจฺจเวกฺขติ. ฌานา วุฏฺฐหิตฺวา ฌานํ ครุํ กตฺวา ปจฺจเวกฺขติ. อริยา มคฺคา วุฏฺฐหิตฺวา มคฺคํ ครุํ กตฺวา ปจฺจเวกฺขนฺติ, ผลํ ครุํ กตฺวา ปจฺจเวกฺขนฺติ, นิพฺพานํ ครุํ กตฺวา ปจฺจเวกฺขนฺติ. นิพฺพานํ โคตฺรภุสฺส, โวทานสฺส, มคฺคสฺส, ผลสฺส อธิปติ\-ปจฺจเยน ปจฺจโย. วตฺถุํ ฯเปฯ. อิตฺถินฺทฺริยํ. ปุริสินฺทฺริยํ. ชีวิตินฺทฺริยํ. อาโปธาตุํ, กพฬีการํ อาหารํ ครุํ กตฺวา อสฺสาเทติ อภินนฺทติ, ตํ ครุํ กตฺวา ราโค อุปฺปชฺชติ, ทิฏฺฐิ อุปฺปชฺชติ. สหชาตา\-ธิปติ\csromandash{}อปฺปฏิฆา\-ธิปติ สมฺปยุตฺตกานํ ขนฺธานํ อปฺปฏิฆานญฺจ จิตฺต\-สมุฏฺฐานานํ รูปานํ อธิปติ\-ปจฺจเยน ปจฺจโย. (1)}

\prose{อปฺปฏิโฆ ธมฺโม สปฺปฏิฆสฺส ธมฺมสฺส อธิปติ\-ปจฺจเยน ปจฺจโย. สหชาตา\-ธิปติ\csromandash{}อปฺปฏิฆา\-ธิปติ สปฺปฏิฆานํ จิตฺต\-สมุฏฺฐานานํ รูปานํ อธิปติ\-ปจฺจเยน ปจฺจโย. (2)}

\prose{อปฺปฏิโฆ ธมฺโม สปฺปฏิฆสฺส จ อปฺปฏิฆสฺส จ ธมฺมสฺส อธิปติ\-ปจฺจเยน ปจฺจโย. สหชาตา\-ธิปติ\csromandash{}อปฺปฏิฆา\-ธิปติ สมฺปยุตฺตกานํ ขนฺธานํ สปฺปฏิฆานญฺจ อปฺปฏิฆานญฺจ จิตฺต\-สมุฏฺฐานานํ รูปานํ อธิปติ\-ปจฺจเยน ปจฺจโย. (3)}

\proseitem{71}{อปฺปฏิโฆ ธมฺโม อปฺปฏิฆสฺส ธมฺมสฺส อนนฺตร\-ปจฺจเยน ปจฺจโย. ปุริมา ปุริมา อปฺปฏิฆา ขนฺธา ฯเปฯ ผลสมาปตฺติยา อนนฺตร\-ปจฺจเยน ปจฺจโย.}

\csromanheader{2}{10. สปฺปฏิฆทุก \textbf{สมนนฺตร}}
\vspace{\tipitakaparskip}
\csromancenter{อปฺปฏิโฆ ธมฺโม อปฺปฏิฆสฺส ธมฺมสฺส สมนนฺตร\-ปจฺจเยน ปจฺจโย ฯเปฯ.}
\csromanheader{2}{\textbf{สหชาตาทิ}}
\proseitem{73}{\csromanfolio{113}สปฺปฏิโฆ ธมฺโม สปฺปฏิฆสฺส ธมฺมสฺส สหชาต\-ปจฺจเยน ปจฺจโย. นว. อญฺญมญฺญ\-ปจฺจเยน ปจฺจโย. ฉ. นิสฺสย\-ปจฺจเยน ปจฺจโย. นว.}

\csromanheader{2}{\textbf{อุปนิสฺสย}}
\proseitem{74}{สปฺปฏิโฆ ธมฺโม อปฺปฏิฆสฺส ธมฺมสฺส อุปนิสฺสย\-ปจฺจเยน ปจฺจโย. อารมฺมณู\-ปนิสฺสโย, ปกตู\-ปนิสฺสโย ฯเปฯ. ปกตู\-ปนิสฺสโย\csromandash{} อุตุํ. เสนาสนํ อุปนิสฺสาย ทานํ เทติ ฯเปฯ สํฆํ ภินฺทติ. อุตุ. เสนาสนํ สทฺธาย ฯเปฯ ผลสมาปตฺติยา อุปนิสฺสย\-ปจฺจเยน ปจฺจโย. (1)}

\prose{อปฺปฏิโฆ ธมฺโม อปฺปฏิฆสฺส ธมฺมสฺส อุปนิสฺสย\-ปจฺจเยน ปจฺจโย. อารมฺมณู\-ปนิสฺสโย, อนนฺตรู\-ปนิสฺสโย, ปกตู\-ปนิสฺสโย ฯเปฯ. ปกตู\-ปนิสฺสโย\csromandash{}สทฺธํ อุปนิสฺสาย ทานํ เทติ ฯเปฯ ทิฏฺฐิํ คณฺหาติ. สีลํ ฯเปฯ. กายิกํ สุขํ. กายิกํ ทุกฺขํ. โภชนํ อุปนิสฺสาย ทานํ เทติ ฯเปฯ สํฆํ ภินฺทติ. สทฺธา ฯเปฯ. โภชนํ สทฺธาย ฯเปฯ ผลสมาปตฺติยา อุปนิสฺสย\-ปจฺจเยน ปจฺจโย. (1)}

\proseitem{75}{สปฺปฏิโฆ ธมฺโม อปฺปฏิฆสฺส ธมฺมสฺส ปุเรชาต\-ปจฺจเยน ปจฺจโย. อารมฺมณ\-ปุเรชาตํ, วตฺถุ\-ปุเรชาตํ.  อารมฺมณ\-ปุเรชาตํ\csromandash{} จกฺขุํ ฯเปฯ. โผฏฺฐพฺเพ อนิจฺจโต ฯเปฯ โทมนสฺสํ อุปฺปชฺชติ. ทิพฺเพน จกฺขุนา รูปํ ปสฺสติ. ทิพฺพาย โสตธาตุยา สทฺทํ สุณาติ. รูปายตนํ จกฺขุ\-วิญฺญาณสฺส ฯเปฯ. โผฏฺฐพฺพา\-ยตนํ กายวิญฺญาณสฺส ปุเรชาต\-ปจฺจเยน ปจฺจโย. วตฺถุ\-ปุเรชาตํ\csromandash{}จกฺขายตนํ จกฺขุ\-วิญฺญาณสฺส ฯเปฯ กายายตนํ กายวิญฺญาณสฺส ปุเรชาต\-ปจฺจเยน ปจฺจโย. (1)}

\prose{อปฺปฏิโฆ ธมฺโม อปฺปฏิฆสฺส ธมฺมสฺส ปุเรชาต\-ปจฺจเยน ปจฺจโย. อารมฺมณ\-ปุเรชาตํ, วตฺถุ\-ปุเรชาตํ.  อารมฺมณ\-ปุเรชาตํ\csromandash{} วตฺถุํ ฯเปฯ อิตฺถินฺทฺริยํ. ปุริสินฺทฺริยํ. ชีวิตินฺทฺริยํ. อาโปธาตุํ. กพฬีการํ อาหารํ \csromanfolio{114}อนิจฺจโต ฯเปฯ โทมนสฺสํ อุปฺปชฺชติ.  วตฺถุ\-ปุเรชาตํ\csromandash{}วตฺถุ อปฺปฏิฆานํ ขนฺธานํ ปุเรชาต\-ปจฺจเยน ปจฺจโย. (1)}

\prose{สปฺปฏิโฆ จ อปฺปฏิโฆ จ ธมฺมา อปฺปฏิฆสฺส ธมฺมสฺส ปุเรชาต\-ปจฺจเยน ปจฺจโย. อารมฺมณ\-ปุเรชาตํ, วตฺถุ\-ปุเรชาตํ.  จกฺขายตนญฺจ วตฺถุ จ ฯเปฯ. โผฏฺฐพฺพายตนญฺจ วตฺถุ จ อปฺปฏิฆานํ ขนฺธานํ ปุเรชาต\-ปจฺจเยน ปจฺจโย. (1)}

\csromanheader{2}{\textbf{ปจฺฉาชาต}}
\proseitem{76}{อปฺปฏิโฆ ธมฺโม อปฺปฏิฆสฺส ธมฺมสฺส ปจฺฉาชาต\-ปจฺจเยน ปจฺจโย. ปจฺฉาชาตา อปฺปฏิฆา ขนฺธา ปุเรชาตสฺส อิมสฺส อปฺปฏิฆสฺส กายสฺส ปจฺฉาชาต\-ปจฺจเยน ปจฺจโย. (มูลํ กาตพฺพํ.) ปจฺฉาชาตา อปฺปฏิฆา ขนฺธา ปุเรชาตสฺส อิมสฺส สปฺปฏิฆสฺส กายสฺส ปจฺฉาชาต\-ปจฺจเยน ปจฺจโย. (มูลํ กาตพฺพํ.) ปจฺฉาชาตา อปฺปฏิฆา ขนฺธา ปุเรชาตสฺส อิมสฺส สปฺปฏิฆสฺส จ อปฺปฏิฆสฺส จ กายสฺส ปจฺฉาชาต\-ปจฺจเยน ปจฺจโย. (ทฺวินฺนมฺปิ มูลา กาตพฺพา.) (3)}

\csromanheader{2}{\textbf{อาเสวน}}
\proseitem{77}{อปฺปฏิโฆ ธมฺโม อปฺปฏิฆสฺส ธมฺมสฺส อาเสวน\-ปจฺจเยน ปจฺจโย. ปุริมา ปุริมา อปฺปฏิฆา ขนฺธา ฯเปฯ. โวทานํ มคฺคสฺส อาเสวน\-ปจฺจเยน ปจฺจโย. (1)}

\proseitem{78}{อปฺปฏิโฆ ธมฺโม อปฺปฏิฆสฺส ธมฺมสฺส กมฺมปจฺจเยน ปจฺจโย. สหชาตา, นานากฺขณิกา.  สหชาตา อปฺปฏิฆา เจตนา สมฺปยุตฺตกานํ ขนฺธานํ อปฺปฏิฆานญฺจ จิตฺต\-สมุฏฺฐานานํ รูปานํ กมฺมปจฺจเยน ปจฺจโย. ปฏิสนฺธิ\-กฺขเณ ฯเปฯ. นานากฺขณิกา อปฺปฏิฆา เจตนา วิปากานํ ขนฺธานํ อปฺปฏิฆานญฺจ กฏตฺตารูปานํ กมฺมปจฺจเยน ปจฺจโย. (1)}

\prose{อปฺปฏิโฆ ธมฺโม สปฺปฏิฆสฺส ธมฺมสฺส กมฺมปจฺจเยน ปจฺจโย. สหชาตา, นานากฺขณิกา.  สหชาตา อปฺปฏิฆา เจตนา สปฺปฏิฆานํ จิตฺต\-สมุฏฺฐานานํ รูปานํ กมฺมปจฺจเยน ปจฺจโย. ปฏิสนฺธิ\-กฺขเณ ฯเปฯ. นานากฺขณิกา อปฺปฏิฆา เจตนา สปฺปฏิฆานํ กฏตฺตารูปานํ กมฺมปจฺจเยน ปจฺจโย. (2)}

\prose{\csromanfolio{115}อปฺปฏิโฆ ธมฺโม สปฺปฏิฆสฺส จ อปฺปฏิฆสฺส จ ธมฺมสฺส กมฺมปจฺจเยน ปจฺจโย. สหชาตา, นานากฺขณิกา.  สหชาตา อปฺปฏิฆา เจตนา สมฺปยุตฺตกานํ ขนฺธานํ สปฺปฏิฆานญฺจ อปฺปฏิฆานญฺจ จิตฺต\-สมุฏฺฐานานํ รูปานํ กมฺมปจฺจเยน ปจฺจโย. ปฏิสนฺธิ\-กฺขเณ ฯเปฯ. นานากฺขณิกา อปฺปฏิฆา เจตนา วิปากานํ ขนฺธานํ สปฺปฏิฆานญฺจ อปฺปฏิฆานญฺจ กฏตฺตารูปานํ กมฺมปจฺจเยน ปจฺจโย. (3)}

\csromanheader{2}{\textbf{วิปาก}}
\proseitem{79}{อปฺปฏิโฆ ธมฺโม อปฺปฏิฆสฺส ธมฺมสฺส วิปาก\-ปจฺจเยน ปจฺจโย. วิปาโก อปฺปฏิโฆ ฯเปฯ. ตีณิ.}

\csromanheader{2}{\textbf{อาหาร}}
\proseitem{80}{อปฺปฏิโฆ ธมฺโม อปฺปฏิฆสฺส ธมฺมสฺส อาหารปจฺจเยน ปจฺจโย. อปฺปฏิฆา อาหารา สมฺปยุตฺตกานํ ขนฺธานํ อปฺปฏิฆานญฺจ จิตฺต\-สมุฏฺฐานานํ รูปานํ อาหารปจฺจเยน ปจฺจโย. ปฏิสนฺธิ\-กฺขเณ ฯเปฯ. กพฬีกาโร อาหาโร อิมสฺส อปฺปฏิฆสฺส กายสฺส อาหารปจฺจเยน ปจฺจโย. (อวเสสา เทฺวปิ ปญฺหา กาตพฺพา. ปฏิสนฺธิ กพฬีกาโร อาหาโร ทฺวีสุปิ กาตพฺโพ อคฺเค.) (3)}

\csromanheader{2}{\textbf{อินฺทฺริยาทิ}}
\proseitem{81}{สปฺปฏิโฆ ธมฺโม อปฺปฏิฆสฺส ธมฺมสฺส อินฺทฺริย\-ปจฺจเยน ปจฺจโย. จกฺขุนฺทฺริยํ จกฺขุ\-วิญฺญาณสฺส ฯเปฯ. กายินฺทฺริยํ กายวิญฺญาณสฺส อินฺทฺริย\-ปจฺจเยน ปจฺจโย. (1)}

\prose{อปฺปฏิโฆ ธมฺโม อปฺปฏิฆสฺส ธมฺมสฺส อินฺทฺริย\-ปจฺจเยน ปจฺจโย. ตีณิ. (ตีสุปิ ชีวิตินฺทฺริยํ อคฺเค กาตพฺพํ.) (3)}

\prose{สปฺปฏิโฆ จ อปฺปฏิโฆ จ ธมฺมา อปฺปฏิฆสฺส ธมฺมสฺส อินฺทฺริย\-ปจฺจเยน ปจฺจโย. จกฺขุนฺทฺริยญฺจ จกฺขุวิญฺญาณญฺจ จกฺขุ\-วิญฺญาณสหคตานํ ขนฺธานํ อินฺทฺริย\-ปจฺจเยน ปจฺจโย ฯเปฯ. กายินฺทฺริยญฺจ กายวิญฺญาณญฺจ กาย\-วิญฺญาณสหคตานํ ขนฺธานํ อินฺทฺริย\-ปจฺจเยน ปจฺจโย. ฌานปจฺจเยน ปจฺจโย. ตีณิ. มคฺคปจฺจเยน ปจฺจโย. ตีณิ. สมฺปยุตฺต\-ปจฺจเยน ปจฺจโย. เอกํ.}

\proseitem{82}{\csromanfolio{116}สปฺปฏิโฆ ธมฺโม อปฺปฏิฆสฺส ธมฺมสฺส วิปฺปยุตฺต\-ปจฺจเยน ปจฺจโย. ปุเรชาตํ จกฺขายตนํ จกฺขุ\-วิญฺญาณสฺส ฯเปฯ. กายายตนํ กายวิญฺญาณสฺส วิปฺปยุตฺต\-ปจฺจเยน ปจฺจโย. (1)}

\prose{อปฺปฏิโฆ ธมฺโม อปฺปฏิฆสฺส ธมฺมสฺส วิปฺปยุตฺต\-ปจฺจเยน ปจฺจโย. สหชาตํ, ปุเรชาตํ, ปจฺฉาชาตํ.  สหชาตา อปฺปฏิฆา ขนฺธา อปฺปฏิฆานํ จิตฺต\-สมุฏฺฐานานํ รูปานํ วิปฺปยุตฺต\-ปจฺจเยน ปจฺจโย. ปฏิสนฺธิ\-กฺขเณ ฯเปฯ. ขนฺธา วตฺถุสฺส วิปฺปยุตฺต\-ปจฺจเยน ปจฺจโย. วตฺถุ ขนฺธานํ วิปฺปยุตฺต\-ปจฺจเยน ปจฺจโย. ปุเรชาตํ วตฺถุ อปฺปฏิฆานํ ขนฺธานํ วิปฺปยุตฺต\-ปจฺจเยน ปจฺจโย. ปจฺฉาชาตา อปฺปฏิฆา ขนฺธา ปุเรชาตสฺส อิมสฺส อปฺปฏิฆสฺส กายสฺส วิปฺปยุตฺต\-ปจฺจเยน ปจฺจโย. (1)}

\prose{อปฺปฏิโฆ ธมฺโม สปฺปฏิฆสฺส ธมฺมสฺส วิปฺปยุตฺต\-ปจฺจเยน ปจฺจโย. สหชาตํ, ปจฺฉาชาตํ.  สหชาตา อปฺปฏิฆา ขนฺธา สปฺปฏิฆานํ จิตฺต\-สมุฏฺฐานานํ รูปานํ วิปฺปยุตฺต\-ปจฺจเยน ปจฺจโย. ปฏิสนฺธิ\-กฺขเณ ฯเปฯ. ปจฺฉาชาตา อปฺปฏิฆา ขนฺธา ปุเรชาตสฺส อิมสฺส สปฺปฏิฆสฺส กายสฺส วิปฺปยุตฺต\-ปจฺจเยน ปจฺจโย. (2)}

\prose{อปฺปฏิโฆ ธมฺโม สปฺปฏิฆสฺส จ อปฺปฏิฆสฺส จ ธมฺมสฺส วิปฺปยุตฺต\-ปจฺจเยน ปจฺจโย. สหชาตํ, ปจฺฉาชาตํ.  สหชาตา อปฺปฏิฆา ขนฺธา สปฺปฏิฆานญฺจ อปฺปฏิฆานญฺจ จิตฺต\-สมุฏฺฐานานํ รูปานํ วิปฺปยุตฺต\-ปจฺจเยน ปจฺจโย. ปฏิสนฺธิ\-กฺขเณ ฯเปฯ. ปจฺฉาชาตา อปฺปฏิฆา ขนฺธา ปุเรชาตสฺส อิมสฺส สปฺปฏิฆสฺส จ อปฺปฏิฆสฺส จ กายสฺส วิปฺปยุตฺต\-ปจฺจเยน ปจฺจโย. (3)}

\proseitem{83}{สปฺปฏิโฆ ธมฺโม สปฺปฏิฆสฺส ธมฺมสฺส อตฺถิปจฺจเยน ปจฺจโย. เอกํ. (ปฏิจฺจสทิสา ปฐมปญฺหา.) (1)}

\prose{สปฺปฏิโฆ ธมฺโม อปฺปฏิฆสฺส ธมฺมสฺส อตฺถิปจฺจเยน ปจฺจโย. สหชาตํ, ปุเรชาตํ.  สหชาตา สปฺปฏิฆา มหาภูตา อาโปธาตุยา อตฺถิปจฺจเยน ปจฺจโย. สปฺปฏิฆา มหาภูตา อปฺปฏิฆานํ จิตฺต\-สมุฏฺฐานานํ รูปานํ, กฏตฺตารูปานํ, อุปาทารูปานํ อตฺถิปจฺจเยน ปจฺจโย. โผฏฺฐพฺพา\-ยตนํ อิตฺถินฺทฺริยสฺส ฯเปฯ. กพฬีการสฺส อาหารสฺส อตฺถิปจฺจเยน \csromanfolio{117}ปจฺจโย. พาหิรํ. อาหารสมุฏฺฐานํ. อุตุสมุฏฺฐานํ. อสญฺญสตฺตานํ ฯเปฯ. ปุเรชาตํ จกฺขุํ ฯเปฯ. โผฏฺฐพฺเพ อนิจฺจโต ฯเปฯ โทมนสฺสํ อุปฺปชฺชติ. ทิพฺเพน จกฺขุนา รูปํ ปสฺสติ. ทิพฺพาย โสตธาตุยา สทฺทํ สุณาติ. รูปายตนญฺจ จกฺขายตนญฺจ จกฺขุ\-วิญฺญาณสฺส ฯเปฯ. โผฏฺฐพฺพายตนญฺจ กายายตนญฺจ กายวิญฺญาณสฺส อตฺถิปจฺจเยน ปจฺจโย. (2)}

\prose{สปฺปฏิโฆ ธมฺโม สปฺปฏิฆสฺส จ อปฺปฏิฆสฺส จ ธมฺมสฺส อตฺถิปจฺจเยน ปจฺจโย. สปฺปฏิฆํ เอกํ มหาภูตํ ทฺวินฺนํ มหาภูตานํ อาโปธาตุยา จ อตฺถิปจฺจเยน ปจฺจโย ฯเปฯ. (ปฏิจฺจสทิสํ ยาว อสญฺญสตฺตา.) (3)}

\proseitem{84}{อปฺปฏิโฆ ธมฺโม อปฺปฏิฆสฺส ธมฺมสฺส อตฺถิปจฺจเยน ปจฺจโย. สหชาตํ, ปุเรชาตํ, ปจฺฉาชาตํ, อาหารํ, อินฺทฺริยํ.  \textbf{สหชาโต} อปฺปฏิโฆ เอโก ขนฺโธ ติณฺณนฺนํ ขนฺธานํ ฯเปฯ. (ยาว อสญฺญสตฺตา.) \textbf{ปุเรชาตํ} วตฺถุํ ฯเปฯ อิตฺถินฺทฺริยํ. ปุริสินฺทฺริยํ. ชีวิตินฺทฺริยํ. อาโปธาตุํ. กพฬีการํ อาหารํ อนิจฺจโต ฯเปฯ โทมนสฺสํ อุปฺปชฺชติ. วตฺถุ อปฺปฏิฆานํ ขนฺธานํ อตฺถิปจฺจเยน ปจฺจโย. \textbf{ปจฺฉาชาตา} อปฺปฏิฆา ขนฺธา ปุเรชาตสฺส อิมสฺส อปฺปฏิฆสฺส กายสฺส อตฺถิปจฺจเยน ปจฺจโย. \textbf{กพฬีกาโร} \textbf{อาหาโร} อิมสฺส อปฺปฏิฆสฺส กายสฺส อตฺถิปจฺจเยน ปจฺจโย. \textbf{รูปชีวิตินฺทฺริยํ} อปฺปฏิฆานํ กฏตฺตารูปานํ อตฺถิปจฺจเยน ปจฺจโย. (1)}

\prose{อปฺปฏิโฆ ธมฺโม สปฺปฏิฆสฺส ธมฺมสฺส อตฺถิปจฺจเยน ปจฺจโย. สหชาตํ, ปจฺฉาชาตํ, อาหารํ, อินฺทฺริยํ.  สหชาตา อปฺปฏิฆา ขนฺธา สปฺปฏิฆานํ จิตฺต\-สมุฏฺฐานานํ รูปานํ อตฺถิปจฺจเยน ปจฺจโย. ปฏิสนฺธิ\-กฺขเณ ฯเปฯ. อาโปธาตุ สปฺปฏิฆานํ มหาภูตานํ อตฺถิปจฺจเยน ปจฺจโย. อาโปธาตุ สปฺปฏิฆานํ จิตฺต\-สมุฏฺฐานานํ รูปานํ, กฏตฺตารูปานํ, อุปาทารูปานํ อตฺถิปจฺจเยน ปจฺจโย. อาโปธาตุ จกฺขา\-ยตนสฺส ฯเปฯ โผฏฺฐพฺพา\-ยตนสฺส อตฺถิปจฺจเยน ปจฺจโย. พาหิรํ. อาหารสมุฏฺฐานํ. อุตุสมุฏฺฐานํ. อสญฺญสตฺตานํ ฯเปฯ. ปจฺฉาชาตา อปฺปฏิฆา ขนฺธา ปุเรชาตสฺส อิมสฺส สปฺปฏิฆสฺส กายสฺส อตฺถิปจฺจเยน ปจฺจโย. กพฬีกาโร อาหาโร อิมสฺส สปฺปฏิฆสฺส กายสฺส อตฺถิปจฺจเยน ปจฺจโย. รูปชีวิตินฺทฺริยํ สปฺปฏิฆานํ กฏตฺตารูปานํ อตฺถิปจฺจเยน ปจฺจโย. (2)}

\prose{อปฺปฏิโฆ ธมฺโม สปฺปฏิฆสฺส จ อปฺปฏิฆสฺส จ ธมฺมสฺส อตฺถิปจฺจเยน ปจฺจโย. สหชาตํ, ปจฺฉาชาตํ, อาหารํ, อินฺทฺริยํ.  สหชาโต \csromanfolio{118}อปฺปฏิโฆ เอโก ขนฺโธ ติณฺณนฺนํ ขนฺธานํ สปฺปฏิฆานญฺจ อปฺปฏิฆานญฺจ ฯเปฯ. (ปฏิจฺจสทิสํ ยาว อสญฺญสตฺตา.) \textbf{ปจฺฉาชาตา} อปฺปฏิฆา ขนฺธา ปุเรชาตสฺส อิมสฺส สปฺปฏิฆสฺส จ อปฺปฏิฆสฺส จ กายสฺส อตฺถิปจฺจเยน ปจฺจโย. \textbf{กพฬีกาโร อาหาโร} อิมสฺส สปฺปฏิฆสฺส จ อปฺปฏิฆสฺส จ กายสฺส อตฺถิปจฺจเยน ปจฺจโย. \textbf{รูปชีวิตินฺทฺริยํ} สปฺปฏิฆานญฺจ อปฺปฏิฆานญฺจ กฏตฺตารูปานํ อตฺถิปจฺจเยน ปจฺจโย. (3)}

\proseitem{85}{สปฺปฏิโฆ จ อปฺปฏิโฆ จ ธมฺมา สปฺปฏิฆสฺส ธมฺมสฺส อตฺถิปจฺจเยน ปจฺจโย. (ปฏิจฺจสทิสํ ยาว อสญฺญสตฺตา.) (1)}

\prose{สปฺปฏิโฆ จ อปฺปฏิโฆ จ ธมฺมา อปฺปฏิฆสฺส ธมฺมสฺส อตฺถิปจฺจเยน ปจฺจโย. สหชาตํ, ปุเรชาตํ.  สหชาตา อปฺปฏิฆา ขนฺธา จ มหาภูตา จ อปฺปฏิฆานํ จิตฺต\-สมุฏฺฐานานํ รูปานํ ฯเปฯ. (ปฏิจฺจสทิสํ ยาว อสญฺญสตฺตา.) สหชาโต จกฺขุ\-วิญฺญาณ\-สหคโต เอโก ขนฺโธ จ จกฺขายตนญฺจ ติณฺณนฺนํ ขนฺธานํ ฯเปฯ เทฺว ขนฺธา จ ฯเปฯ. กาย\-วิญฺญาณ\-สหคโต เอโก ขนฺโธ จ กายายตนญฺจ ติณฺณนฺนํ ขนฺธานํ อตฺถิปจฺจเยน ปจฺจโย ฯเปฯ เทฺว ขนฺธา จ ฯเปฯ. (2)}

\prose{สปฺปฏิโฆ จ อปฺปฏิโฆ จ ธมฺมา สปฺปฏิฆสฺส จ อปฺปฏิฆสฺส จ ธมฺมสฺส อตฺถิปจฺจเยน ปจฺจโย. (ปฏิจฺจสทิสํ.) (3)}

\csromanmarkh{1. ปจฺจยานุโลม}\csromanmarkhii{2. สงฺขฺยาวาร}\csromanheadernomark{2}{1. \textbf{ปจฺจยานุโลม 2. สงฺขฺยาวาร}}
\proseitem{86}{เหตุยา ตีณิ, อารมฺมเณ เทฺว, อธิปติยา จตฺตาริ, อนนฺตเร เอกํ, สมนนฺตเร เอกํ, สหชาเต นว, อญฺญมญฺเญ ฉ, นิสฺสเย นว, อุปนิสฺสเย เทฺว, ปุเรชาเต ตีณิ, ปจฺฉาชาเต ตีณิ, อาเสวเน เอกํ, กมฺเม ตีณิ, วิปาเก ตีณิ, อาหาเร ตีณิ, อินฺทฺริเย ปญฺจ, ฌาเน ตีณิ, มคฺเค ตีณิ, สมฺปยุตฺเต เอกํ, วิปฺปยุตฺเต จตฺตาริ, อตฺถิยา นว, นตฺถิยา เอกํ, วิคเต เอกํ, อวิคเต นว. (เอวํ คเณตพฺพํ.)}

\vspace{\tipitakaparskip}
\csromancenter{อนุโลมํ.}
\csromanmarkh{10. สปฺปฏิฆทุก}\csromanmarkhii{2. ปจฺจนียุ\-ทฺธาร}\csromanheadernomark{2}{10. สปฺปฏิฆทุก 2. ปจฺจนียุ\-ทฺธาร}
\proseitem{87}{\csromanfolio{119}สปฺปฏิโฆ ธมฺโม สปฺปฏิฆสฺส ธมฺมสฺส สหชาต\-ปจฺจเยน ปจฺจโย. (1)}

\prose{สปฺปฏิโฆ ธมฺโม อปฺปฏิฆสฺส ธมฺมสฺส อารมฺมณ\-ปจฺจเยน ปจฺจโย. สหชาต\-ปจฺจเยน ปจฺจโย. อุปนิสฺสย\-ปจฺจเยน ปจฺจโย. ปุเรชาต\-ปจฺจเยน ปจฺจโย. (2)}

\prose{สปฺปฏิโฆ ธมฺโม สปฺปฏิฆสฺส จ อปฺปฏิฆสฺส จ ธมฺมสฺส สหชาต\-ปจฺจเยน ปจฺจโย. (3)}

\proseitem{88}{อปฺปฏิโฆ ธมฺโม อปฺปฏิฆสฺส ธมฺมสฺส อารมฺมณ\-ปจฺจเยน ปจฺจโย. สหชาต\-ปจฺจเยน ปจฺจโย. อุปนิสฺสย\-ปจฺจเยน ปจฺจโย. ปุเรชาต\-ปจฺจเยน ปจฺจโย. ปจฺฉาชาต\-ปจฺจเยน ปจฺจโย. กมฺมปจฺจเยน ปจฺจโย. อาหารปจฺจเยน ปจฺจโย. อินฺทฺริย\-ปจฺจเยน ปจฺจโย. (1)}

\prose{อปฺปฏิโฆ ธมฺโม สปฺปฏิฆสฺส ธมฺมสฺส สหชาต\-ปจฺจเยน ปจฺจโย. ปจฺฉาชาต\-ปจฺจเยน ปจฺจโย. กมฺมปจฺจเยน ปจฺจโย. อาหารปจฺจเยน ปจฺจโย. อินฺทฺริย\-ปจฺจเยน ปจฺจโย. (2)}

\prose{อปฺปฏิโฆ ธมฺโม สปฺปฏิฆสฺส จ อปฺปฏิฆสฺส จ ธมฺมสฺส สหชาต\-ปจฺจเยน ปจฺจโย. ปจฺฉาชาต\-ปจฺจเยน ปจฺจโย. กมฺมปจฺจเยน ปจฺจโย. อาหารปจฺจเยน ปจฺจโย. อินฺทฺริย\-ปจฺจเยน ปจฺจโย. (3)}

\proseitem{89}{สปฺปฏิโฆ จ อปฺปฏิโฆ จ ธมฺมา สปฺปฏิฆสฺส ธมฺมสฺส สหชาต\-ปจฺจเยน ปจฺจโย. (1)}

\prose{สปฺปฏิโฆ จ อปฺปฏิโฆ จ ธมฺมา อปฺปฏิฆสฺส ธมฺมสฺส สหชาตํ. ปุเรชาตํ. (2)}

\prose{สปฺปฏิโฆ จ อปฺปฏิโฆ จ ธมฺมา สปฺปฏิฆสฺส จ อปฺปฏิฆสฺส จ ธมฺมสฺส สหชาต\-ปจฺจเยน ปจฺจโย. (3)}

\csromanmarkh{2. ปจฺจย\-ปจฺจนีย}\csromanmarkhii{2. สงฺขฺยาวาร}\csromanheadernomark{2}{2. ปจฺจย\-ปจฺจนีย 2. สงฺขฺยาวาร}
\proseitem{90}{\csromanfolio{120}นเหตุยา นว ฯเปฯ นอนนฺตเร นว, นสมนนฺตเร นว, นสหชาเต จตฺตาริ, นอญฺญมญฺเญ นว, นนิสฺสเย จตฺตาริ, นอุปนิสฺสเย นว, นปุเรชาเต นว ฯเปฯ นสมฺปยุตฺเต นว, นวิปฺปยุตฺเต นว, โนอตฺถิยา จตฺตาริ, โนนตฺถิยา นว, โนวิคเต นว, โนอวิคเต จตฺตาริ.}

\vspace{\tipitakaparskip}
\csromancenter{ปจฺจนียํ.}
\proseitem{91}{เหตุปจฺจยา นอารมฺมเณ ตีณิ ฯเปฯ นอนนฺตเร ตีณิ, นสมนนฺตเร ตีณิ, นอญฺญมญฺเญ ตีณิ, นอุปนิสฺสเย ตีณิ ฯเปฯ นสมฺปยุตฺเต ตีณิ, นวิปฺปยุตฺเต เอกํ, โนนตฺถิยา ตีณิ, โนวิคเต ตีณิ.}

\vspace{\tipitakaparskip}
\csromancenter{อนุโลม\-ปจฺจนียํ.}
\proseitem{92}{นเหตุปจฺจยา อารมฺมเณ เทฺว, อธิปติยา จตฺตาริ. (อนุโลมมาติกา คเณตพฺพา.) อวิคเต นว.}

\vspace{\tipitakaparskip}
\csromancenter{ปจฺจนียา\-นุโลมํ.}
\nitthitam{สปฺปฏิฆทุกํ นิฏฺฐิตํ.}
\csromanmatikaanchor{mtk.14}
\csromantocmarkref{section}{11. รูปีทุก}{mtk.14}
\bookmark[dest={mtk.14},level=1]{11. รูปีทุก}
\csromanmarkh{11. รูปีทุก}\csromanmarkhii{11. รูปีทุก 1. ปฏิจฺจวาร}\csromanheadernomark{1}{11. รูปีทุก \textbf{11. รูปีทุก 1. ปฏิจฺจวาร}}
\csromanmarkh{1. ปจฺจยานุโลม}\csromanmarkhii{1. วิภงฺควาร}\csromanheadernomark{2}{1. \textbf{ปจฺจยานุโลม 1. วิภงฺควาร}}
\csromanheader{2}{\textbf{เหตุ}}
\proseitem{93}{\csromanfolio{121}รูปิํ ธมฺมํ ปฏิจฺจ รูปี ธมฺโม อุปฺปชฺชติ เหตุปจฺจยา. เอกํ มหาภูตํ ปฏิจฺจ ตโย มหาภูตา ฯเปฯ เทฺว มหาภูเต ฯเปฯ มหาภูเต ปฏิจฺจ จิตฺต\-สมุฏฺฐานํ รูปํ, กฏตฺตารูปํ, อุปาทารูปํ. (1)}

\prose{รูปิํ ธมฺมํ ปฏิจฺจ อรูปี ธมฺโม อุปฺปชฺชติ เหตุปจฺจยา. ปฏิสนฺธิ\-กฺขเณ วตฺถุํ ปฏิจฺจ อรูปิโน ขนฺธา. (2)}

\prose{รูปิํ ธมฺมํ ปฏิจฺจ รูปี จ อรูปี จ ธมฺมา อุปฺปชฺชนฺติ เหตุปจฺจยา. ปฏิสนฺธิ\-กฺขเณ วตฺถุํ ปฏิจฺจ อรูปิโน ขนฺธา. มหาภูเต ปฏิจฺจ กฏตฺตารูปํ. (3)}

\proseitem{94}{อรูปิํ ธมฺมํ ปฏิจฺจ อรูปี ธมฺโม อุปฺปชฺชติ เหตุปจฺจยา. อรูปิํ เอกํ ขนฺธํ ปฏิจฺจ ตโย ขนฺธา ฯเปฯ เทฺว ขนฺเธ ฯเปฯ. ปฏิสนฺธิ\-กฺขเณ ฯเปฯ. (1)}

\prose{อรูปิํ ธมฺมํ ปฏิจฺจ รูปี ธมฺโม อุปฺปชฺชติ เหตุปจฺจยา. อรูปิโน ขนฺเธ ปฏิจฺจ จิตฺต\-สมุฏฺฐานํ รูปํ. ปฏิสนฺธิ\-กฺขเณ ฯเปฯ.}

\prose{อรูปิํ ธมฺมํ ปฏิจฺจ รูปี จ อรูปี จ ธมฺมา อุปฺปชฺชนฺติ เหตุปจฺจยา. อรูปิํ เอกํ ขนฺธํ ปฏิจฺจ ตโย ขนฺธา จิตฺต\-สมุฏฺฐานญฺจ รูปํ ฯเปฯ เทฺว ขนฺเธ ฯเปฯ. ปฏิสนฺธิ\-กฺขเณ ฯเปฯ. (3)}

\proseitem{95}{รูปิญฺจ อรูปิญฺจ ธมฺมํ ปฏิจฺจ รูปี ธมฺโม อุปฺปชฺชติ เหตุปจฺจยา. อรูปิโน ขนฺเธ จ มหาภูเต จ ปฏิจฺจ จิตฺต\-สมุฏฺฐานํ รูปํ. ปฏิสนฺธิ\-กฺขเณ ฯเปฯ. (1)}

\prose{รูปิญฺจ อรูปิญฺจ ธมฺมํ ปฏิจฺจ อรูปี ธมฺโม อุปฺปชฺชติ เหตุปจฺจยา. ปฏิสนฺธิ\-กฺขเณ อรูปิํ เอกํ ขนฺธญฺจ วตฺถุญฺจ ปฏิจฺจ ตโย ขนฺธา ฯเปฯ เทฺว ขนฺเธ ฯเปฯ. (2)}

\prose{รูปิญฺจ อรูปิญฺจ ธมฺมํ ปฏิจฺจ รูปี จ อรูปี จ ธมฺมา อุปฺปชฺชนฺติ เหตุปจฺจยา. ปฏิสนฺธิ\-กฺขเณ อรูปิํ เอกํ ขนฺธญฺจ วตฺถุญฺจ ปฏิจฺจ ตโย ขนฺธา ฯเปฯ เทฺว ขนฺเธ ฯเปฯ. อรูปิโน ขนฺเธ จ มหาภูเต จ ปฏิจฺจ กฏตฺตารูปํ. (สํขิตฺตํ.) (3)}

\csromanmarkh{1. ปจฺจยานุโลม}\csromanmarkhii{2. สงฺขฺยาวาร}\csromanheadernomark{2}{1. \textbf{ปจฺจยานุโลม 2. สงฺขฺยาวาร}}
\proseitem{96}{\csromanfolio{122}เหตุยา นว, อารมฺมเณ ตีณิ, อธิปติยา ปญฺจ, อนนฺตเร ตีณิ, สมนนฺตเร ตีณิ, สหชาเต นว, อญฺญมญฺเญ ฉ, นิสฺสเย นว, อุปนิสฺสเย ตีณิ, ปุเรชาเต เอกํ, อาเสวเน เอกํ, กมฺเม นว, วิปาเก นว, อาหาเร นว, อินฺทฺริเย นว, ฌาเน นว, มคฺเค นว, สมฺปยุตฺเต ตีณิ, วิปฺปยุตฺเต นว, อตฺถิยา นว, นตฺถิยา ตีณิ, วิคเต ตีณิ, อวิคเต นว.}

\vspace{\tipitakaparskip}
\csromancenter{อนุโลมํ.}
\csromanmarkh{2. ปจฺจย\-ปจฺจนีย}\csromanmarkhii{1. วิภงฺควาร}\csromanheadernomark{2}{2. ปจฺจย\-ปจฺจนีย 1. วิภงฺควาร}
\proseitem{97}{รูปิํ ธมฺมํ ปฏิจฺจ รูปี ธมฺโม อุปฺปชฺชติ นเหตุปจฺจยา. ตีณิ.}

\prose{อรูปิํ ธมฺมํ ปฏิจฺจ อรูปี ธมฺโม อุปฺปชฺชติ นเหตุปจฺจยา. อเหตุกํ อรูปิํ เอกํ ขนฺธํ ปฏิจฺจ ตโย ขนฺธา ฯเปฯ เทฺว ขนฺเธ ฯเปฯ. อเหตุก\-ปฏิสนฺธิ\-กฺขเณ ฯเปฯ. วิจิกิจฺฉา\-สหคเต อุทฺธจฺจ\-สหคเต ขนฺเธ ปฏิจฺจ วิจิกิจฺฉา\-สหคโต อุทฺธจฺจ\-สหคโต โมโห.}

\vspace{\tipitakaparskip}
\csromancenter{(นเหตุปจฺจยา นว ปญฺหา, “อเหตุกนฺ”ติ นิยาเมตพฺพํ.)}
\csromanmarkh{2. ปจฺจย\-ปจฺจนีย}\csromanmarkhii{2. สงฺขฺยาวาร}\csromanheadernomark{2}{2. ปจฺจย\-ปจฺจนีย 2. สงฺขฺยาวาร}
\vspace{\tipitakaparskip}
\csromancenter{\textbf{นเหตุ}ยา นว, นอารมฺมเณ ตีณิ, นอธิปติยา นว, นอนนฺตเร ตีณิ, นสมนนฺตเร ตีณิ, นอญฺญมญฺเญ ตีณิ, นฺเอาปนิสฺสเย ตีณิ, นปุเรชาเต นว, นปจฺฉาชาเต นว, นอาเสวเน นว, นกมฺเม เทฺว, นวิปาเก ปญฺจ, นอาหาเร เอกํ, ไนนฺทฺริเย เอกํ, นฌาเน เทฺว, นมคฺเค นว, นสมฺปยุตฺเต ตีณิ, นวิปฺปยุตฺเต เทฺว, โนนตฺถิยา ตีณิ, โนวิคเต ตีณิ.}
\csromancenter{ปจฺจนียํ.}
\csromanmarkh{11. รูปีทุก}\csromanmarkhii{3. ปจฺจยา\-นุโลม\-ปจฺจนีย}\csromanheadernomark{2}{11. รูปีทุก 3. ปจฺจยา\-นุโลม\-ปจฺจนีย}
\vspace{\tipitakaparskip}
\csromancenter{\textbf{เหตุ}ปจฺจยา นอารมฺมเณ ตีณิ, นอธิปติยา นว, นอนนฺตเร ตีณิ, นสมนนฺตเร ตีณิ, นอญฺญมญฺเญ ตีณิ, นฺเอาปนิสฺสเย ตีณิ, นปุเรชาเต นว, นปจฺฉาชาเต นว, นอาเสวเน นว, นกมฺเม เอกํ, นวิปาเก ปญฺจ, นสมฺปยุตฺเต ตีณิ, นวิปฺปยุตฺเต เอกํ, โนนตฺถิยา ตีณิ, โนวิคเต ตีณิ.}
\csromancenter{อนุโลม\-ปจฺจนียํ.}
\proseitem{100}{\csromanfolio{123}นเหตุปจฺจยา อารมฺมเณ ตีณิ, อนนฺตเร ตีณิ, สมนนฺตเร ตีณิ, สหชาเต นว, อญฺญมญฺเญ ฉ, นิสฺสเย นว, อุปนิสฺสเย ตีณิ, ปุเรชาเต เอกํ, อาเสวเน เอกํ, กมฺเม นว, วิปาเก นว, อาหาเร นว, อินฺทฺริเย นว, ฌาเน นว, มคฺเค เอกํ, สมฺปยุตฺเต ตีณิ, วิปฺปยุตฺเต นว, อตฺถิยา นว, นตฺถิยา ตีณิ, วิคเต ตีณิ, อวิคเต นว.}

\vspace{\tipitakaparskip}
\csromancenter{ปจฺจนียา\-นุโลมํ.}
\csromancenter{(สหชาตวาโรปิ ปฏิจฺจ\-วาร\-สทิโส.)}
\csromanmarkh{11. รูปีทุก}\csromanmarkhii{3. ปจฺจยวาร}\csromanheadernomark{2}{11. รูปีทุก 3. ปจฺจยวาร}
\csromanmarkh{1. ปจฺจยานุโลม}\csromanmarkhii{1. วิภงฺควาร}\csromanheadernomark{2}{1. \textbf{ปจฺจยานุโลม 1. วิภงฺควาร}}
\csromanheader{2}{\textbf{เหตุ}}
\proseitem{101}{รูปิํ ธมฺมํ ปจฺจยา รูปี ธมฺโม อุปฺปชฺชติ เหตุปจฺจยา. เอกํ มหาภูตํ ฯเปฯ. (ปฏิจฺจสทิสํ.) (1)}

\prose{\csromanfolio{124}รูปิํ ธมฺมํ ปจฺจยา อรูปี ธมฺโม อุปฺปชฺชติ เหตุปจฺจยา. วตฺถุํ ปจฺจยา อรูปิโน ขนฺธา. ปฏิสนฺธิ\-กฺขเณ ฯเปฯ. (2)}

\prose{รูปิํ ธมฺมํ ปจฺจยา รูปี จ อรูปี จ ธมฺมา อุปฺปชฺชนฺติ เหตุปจฺจยา. วตฺตุํ ปจฺจยา อรูปิโน ขนฺธา. มหาภูเต ปจฺจยา จิตฺต\-สมุฏฺฐานํ รูปํ. ปฏิสนฺธิ\-กฺขเณ ฯเปฯ. (เอวํ อวเสสา ปญฺหา, ปวตฺติ\-ปฏิสนฺธิ วิภชิตพฺพา.) (3)}

\proseitem{102}{รูปิํ ธมฺมํ ปจฺจยา อรูปี ธมฺโม อุปฺปชฺชติ อารมฺมณ\-ปจฺจยา. จกฺขายตนํ ปจฺจยา จกฺขุวิญฺญาณํ ฯเปฯ. กายายตนํ ปจฺจยา กายวิญฺญาณํ. วตฺถุํ ปจฺจยา อรูปิโน ขนฺธา. ปฏิสนฺธิ\-กฺขเณ ฯเปฯ. (1)}

\prose{อรูปิํ ธมฺมํ ปจฺจยา อรูปี ธมฺโม อุปฺปชฺชติ อารมฺมณ\-ปจฺจยา. อรูปิํ เอกํ ขนฺธํ ฯเปฯ เทฺว ขนฺเธ ฯเปฯ. (1)}

\prose{รูปิญฺจ อรูปิญฺจ ธมฺมํ ปจฺจยา อรูปี ธมฺโม อุปฺปชฺชติ อารมฺมณ\-ปจฺจยา. จกฺขุ\-วิญฺญาณสหคตํ เอกํ ขนฺธญฺจ จกฺขายตนญฺจ ปจฺจยา ตโย ขนฺธา ฯเปฯ เทฺว ขนฺเธ ฯเปฯ. กาย\-วิญฺญาณสหคตํ ฯเปฯ. อรูปิํ เอกํ ขนฺธญฺจ วตฺถุญฺจ ปจฺจยา ตโย ขนฺธา ฯเปฯ เทฺว ขนฺเธ ฯเปฯ. (สํขิตฺตํ.) (3)}

\csromanmarkh{1. ปจฺจยานุโลม}\csromanmarkhii{2. สงฺขฺยาวาร}\csromanheadernomark{2}{1. \textbf{ปจฺจยานุโลม 2. สงฺขฺยาวาร}}
\proseitem{103}{เหตุยา นว, อารมฺมเณ ตีณิ, อธิปติยา นว, อนนฺตเร ตีณิ, สมนนฺตเร ตีณิ, สหชาเต นว, อญฺญมญฺเญ ฉ, นิสฺสเย นว, อุปนิสฺสเย ตีณิ, ปุเรชาเต ตีณิ, อาเสวเน ตีณิ, กมฺเม นว ฯเปฯ มคฺเค นว, สมฺปยุตฺเต ตีณิ, วิปฺปยุตฺเต นว, อตฺถิยา นว, นตฺถิยา ตีณิ, วิคเต ตีณิ, อวิคเต นว.}

\vspace{\tipitakaparskip}
\csromancenter{อนุโลมํ.}
\csromanmarkh{11. รูปีทุก}\csromanmarkhii{2. ปจฺจย\-ปจฺจนีย 1. วิภงฺควาร}\csromanheadernomark{2}{11. รูปีทุก 2. ปจฺจย\-ปจฺจนีย 1. วิภงฺควาร}
\proseitem{104}{\csromanfolio{125}รูปิํ ธมฺมํ ปจฺจยา รูปี ธมฺโม อุปฺปชฺชติ นเหตุปจฺจยา. เอกํ มหาภูตํ ฯเปฯ. อสญฺญสตฺตานํ เอกํ มหาภูตํ ฯเปฯ. (1)}

\prose{รูปิํ ธมฺมํ ปจฺจยา อรูปี ธมฺโม อุปฺปชฺชติ นเหตุปจฺจยา. จกฺขายตนํ ปจฺจยา จกฺขุวิญฺญาณํ ฯเปฯ. กายายตนํ ปจฺจยา กายวิญฺญาณํ. วตฺถุํ ปจฺจยา อเหตุกา อรูปิโน ขนฺธา. อเหตุก\-ปฏิสนฺธิ\-กฺขเณ ฯเปฯ. วตฺถุํ ปจฺจยา วิจิกิจฺฉา\-สหคโต อุทฺธจฺจ\-สหคโต โมโห. (2)}

\prose{รูปิํ ธมฺมํ ปจฺจยา รูปี จ อรูปี จ ธมฺมา อุปฺปชฺชนฺติ นเหตุปจฺจยา. (ปวตฺติ\-ปฏิสนฺธิ กาตพฺพา.) (3)}

\proseitem{105}{อรูปิํ ธมฺมํ ปจฺจยา อรูปี ธมฺโม อุปฺปชฺชติ นเหตุปจฺจยา. อเหตุกํ อรูปิํ เอกํ ขนฺธํ ฯเปฯ. ปฏิสนฺธิ\-กฺขเณ ฯเปฯ. วิจิกิจฺฉา\-สหคเต อุทฺธจฺจ\-สหคเต ขนฺเธ ปจฺจยา วิจิกิจฺฉา\-สหคโต อุทฺธจฺจ\-สหคโต โมโห. (1)}

\prose{อรูปิํ ธมฺมํ ปจฺจยา รูปี ธมฺโม อุปฺปชฺชติ นเหตุปจฺจยา. อเหตุเก อรูปิโน ขนฺเธ ปจฺจยา จิตฺต\-สมุฏฺฐานํ รูปํ. ปฏิสนฺธิ\-กฺขเณ ฯเปฯ. (2)}

\prose{อรูปิํ ธมฺมํ ปจฺจยา รูปี จ อรูปี จ ธมฺมา อุปฺปชฺชนฺติ นเหตุปจฺจยา. อเหตุกํ อรูปิํ เอกํ ขนฺธํ ปจฺจยา ตโย ขนฺธา จิตฺต\-สมุฏฺฐานญฺจ รูปํ ฯเปฯ เทฺว ขนฺเธ ฯเปฯ. ปฏิสนฺธิ\-กฺขเณ ฯเปฯ. (3)}

\proseitem{106}{รูปิญฺจ อรูปิญฺจ ธมฺมํ ปจฺจยา รูปี ธมฺโม อุปฺปชฺชติ นเหตุปจฺจยา. อเหตุเก อรูปิโน ขนฺเธ จ มหาภูเต จ ปจฺจยา จิตฺต\-สมุฏฺฐานํ รูปํ. ปฏิสนฺธิ\-กฺขเณ ฯเปฯ. (1)}

\prose{รูปิญฺจ อรูปิญฺจ ธมฺมํ ปจฺจยา อรูปี ธมฺโม อุปฺปชฺชติ นเหตุปจฺจยา. จกฺขุ\-วิญฺญาณสหคตํ เอกํ ขนฺธญฺจ จกฺขายตนญฺจ ปจฺจยา ตโย ขนฺธา ฯเปฯ เทฺว ขนฺเธ ฯเปฯ. กาย\-วิญฺญาณสหคตํ ฯเปฯ. อรูปิํ เอกํ ขนฺธญฺจ วตฺถุญฺจ ปจฺจยา ตโย ขนฺธา ฯเปฯ เทฺว ขนฺเธ ฯเปฯ. ปฏิสนฺธิ\-กฺขเณ ฯเปฯ. วิจิกิจฺฉา\-สหคเต อุทฺธจฺจ\-สหคเต ขนฺเธ จ วตฺถุญฺจ ปจฺจยา วิจิกิจฺฉา\-สหคโต อุทฺธจฺจ\-สหคโต โมโห. (2)}

\prose{\csromanfolio{126}รูปิญฺจ อรูปิญฺจ ธมฺมํ ปจฺจยา รูปี จ อรูปี จ ธมฺมา อุปฺปชฺชนฺติ นเหตุปจฺจยา. อเหตุกํ อรูปิํ เอกํ ขนฺธญฺจ วตฺถุญฺจ ปจฺจยา ตโย ขนฺธา ฯเปฯ เทฺว ขนฺเธ ฯเปฯ. อรูปิโน ขนฺเธ จ มหาภูเต จ ปจฺจยา จิตฺต\-สมุฏฺฐานํ รูปํ. ปฏิสนฺธิ\-กฺขเณ. (3)}

\csromanmarkh{2. ปจฺจย\-ปจฺจนีย}\csromanmarkhii{2. สงฺขฺยาวาร}\csromanheadernomark{2}{2. ปจฺจย\-ปจฺจนีย 2. สงฺขฺยาวาร}
\proseitem{107}{นเหตุยา นว, นอารมฺมเณ ตีณิ, นอธิปติยา นว, นอนนฺตเร ตีณิ, นสมนนฺตเร ตีณิ, นอญฺญมญฺเญ ตีณิ, นอุปนิสฺสเย ตีณิ, นปุเรชาเต นว, นปจฺฉาชาเต นว, นอาเสวเน นว, นกมฺเม จตฺตาริ, นวิปาเก นว, นอาหาเร เอกํ, นอินฺทฺริเย เอกํ, นฌาเน จตฺตาริ, นมคฺเค นว, นสมฺปยุตฺเต ตีณิ, นวิปฺปยุตฺเต เทฺว, โนนตฺถิยา ตีณิ, โนวิคเต ตีณิ.}

\vspace{\tipitakaparskip}
\csromancenter{\textbf{ปจฺจนียํ.}}
\proseitem{108}{เหตุปจฺจยา นอารมฺมเณ ตีณิ. (สํขิตฺตํ. สพฺเพ กาตพฺพา.) นกมฺเม ตีณิ, นวิปาเก นว, นสมฺปยุตฺเต ตีณิ, นวิปฺปยุตฺเต เอกํ, โนนตฺถิยา ตีณิ, โนวิคเต ตีณิ.}

\vspace{\tipitakaparskip}
\csromancenter{อนุโลม\-ปจฺจนียํ.}
\proseitem{109}{นเหตุปจฺจยา อารมฺมเณ ตีณิ, (สพฺเพ กาตพฺพา.) ฯเปฯ ฌาเน นว, มคฺเค ตีณิ ฯเปฯ อวิคเต นว.}

\vspace{\tipitakaparskip}
\csromancenter{ปจฺจนียา\-นุโลมํ.}
\csromancenter{(นิสฺสยวาโรปิ ปจฺจยวาร\-สทิโส.)}
\csromanmarkh{11. รูปีทุก}\csromanmarkhii{11. รูปีทุก 5. สํสฏฺฐวาร}\csromanheadernomark{2}{11. รูปีทุก \textbf{11. รูปีทุก 5. สํสฏฺฐวาร}}
\proseitem{110}{\csromanfolio{127}อรูปิํ ธมฺมํ สํสฏฺโฐ อรูปี ธมฺโม อุปฺปชฺชติ เหตุปจฺจยา. อรูปิํ เอกํ ขนฺธํ สํสฏฺฐา ตโย ขนฺธา ฯเปฯ เทฺว ขนฺเธ ฯเปฯ. ปฏิสนฺธิ\-กฺขเณ ฯเปฯ.}

\prose{เหตุยา เอกํ ฯเปฯ อวิคเต เอกํ. (เอวํ ปจฺจนียาทีนิ คณนาปิ สมฺปยุตฺตวาเรปิ สพฺเพ กาตพฺพา. เอโกเยว ปโญฺห.)}

\csromanmarkh{11. รูปีทุก}\csromanmarkhii{7. ปญฺหาวาร}\csromanheadernomark{2}{11. รูปีทุก 7. ปญฺหาวาร}
\csromanmarkh{1. ปจฺจยานุโลม}\csromanmarkhii{1. วิภงฺควาร}\csromanheadernomark{2}{1. \textbf{ปจฺจยานุโลม 1. วิภงฺควาร}}
\vspace{\tipitakaparskip}
\csromancenter{\textbf{เหตุ}}
\proseitem{111}{อรูปี ธมฺโม อรูปิสฺส ธมฺมสฺส เหตุปจฺจเยน ปจฺจโย. อรูปี เหตู สมฺปยุตฺตกานํ ขนฺธานํ เหตุปจฺจเยน ปจฺจโย. ปฏิสนฺธิ\-กฺขเณ ฯเปฯ. (1)}

\prose{อรูปี ธมฺโม รูปิสฺส ธมฺมสฺส เหตุปจฺจเยน ปจฺจโย. อรูปี เหตู จิตฺต\-สมุฏฺฐานานํ รูปานํ เหตุปจฺจเยน ปจฺจโย. ปฏิสนฺธิ\-กฺขเณ ฯเปฯ. (2)}

\prose{อรูปี ธมฺโม รูปิสฺส จ อรูปิสฺส จ ธมฺมสฺส เหตุปจฺจเยน ปจฺจโย. อรูปี เหตู สมฺปยุตฺตกานํ ขนฺธานํ จิตฺตสมุฏฺฐานานญฺจ รูปานํ เหตุปจฺจเยน ปจฺจโย. ปฏิสนฺธิ\-กฺขเณ ฯเปฯ. (3)}

\proseitem{112}{รูปี ธมฺโม อรูปิสฺส ธมฺมสฺส อารมฺมณ\-ปจฺจเยน ปจฺจโย. จกฺขุํ ฯเปฯ วตฺถุํ. อิตฺถินฺทฺริยํ. ปุริสินฺทฺริยํ. ชีวิตินฺทฺริยํ. อาโปธาตุํ. กพฬีการํ อาหารํ อนิจฺจโต ฯเปฯ โทมนสฺสํ อุปฺปชฺชติ. ทิพฺเพน จกฺขุนา รูปํ ปสฺสติ. ทิพฺพาย โสตธาตุยา สทฺทํ สุณาติ. รูปายตนํ จกฺขุ\-วิญฺญาณสฺส อารมฺมณ\-ปจฺจเยน ปจฺจโย ฯเปฯ. โผฏฺฐพฺพา\-ยตนํ กายวิญฺญาณสฺส อารมฺมณ\-ปจฺจเยน ปจฺจโย. รูปิโน ขนฺธา อิทฺธิวิธ\-ญาณสฺส, ปุพฺเพ\-นิวาสา\-นุสฺสติ\-ญาณสฺส, อนาคตํสญาณสฺส, อาวชฺชนาย อารมฺมณ\-ปจฺจเยน ปจฺจโย. (1)}

\proseitem{113}{\csromanfolio{128}อรูปี ธมฺโม อรูปิสฺส ธมฺมสฺส อารมฺมณ\-ปจฺจเยน ปจฺจโย. ทานํ ฯเปฯ สีลํ ฯเปฯ อุโปสถกมฺมํ กตฺวา ตํ ปจฺจเวกฺขติ. ปุพฺเพ สุจิณฺณานิ ปจฺจเวกฺขติ. ฌานา ฯเปฯ. อริยา มคฺคา วุฏฺฐหิตฺวา มคฺคํ ปจฺจเวกฺขนฺติ, ผลํ ปจฺจเวกฺขนฺติ, นิพฺพานํ ปจฺจเวกฺขนฺติ. นิพฺพานํ โคตฺรภุสฺส, โวทานสฺส, มคฺคสฺส, ผลสฺส, อาวชฺชนาย อารมฺมณ\-ปจฺจเยน ปจฺจโย. อริยา ปหีเน กิเลเส ฯเปฯ. วิกฺขมฺภิเต กิเลเส ฯเปฯ. ปุพฺเพ ฯเปฯ. อรูปิโน ขนฺเธ อนิจฺจโต ฯเปฯ โทมนสฺสํ อุปฺปชฺชติ. เจโตปริย\-ญาเณน อรูปิ\-จิตฺต\-สมงฺคิสฺส จิตฺตํ ชานาติ. อากาสา\-นญฺจา\-ยตนํ ฯเปฯ. เนว\-สญฺญา\-นา\-สญฺญา\-ยตนสฺส ฯเปฯ. อรูปิโน ขนฺธา อิทฺธิวิธ\-ญาณสฺส, เจโตปริย\-ญาณสฺส, ปุพฺเพ\-นิวาสา\-นุสฺสติ\-ญาณสฺส, ยถากมฺมูปคญาณสฺส, อนาคตํสญาณสฺส, อาวชฺชนาย อารมฺมณ\-ปจฺจเยน ปจฺจโย. (1)}

\proseitem{114}{รูปี ธมฺโม อรูปิสฺส ธมฺมสฺส อธิปติ\-ปจฺจเยน ปจฺจโย. อารมฺมณา\-ธิปติ\csromandash{}จกฺขุํ ฯเปฯ. กพฬีการํ อาหารํ ครุํ กตฺวา อสฺสาเทติ อภินนฺทติ, ตํ ครุํ กตฺวา ราโค อุปฺปชฺชติ. ทิฏฺฐิ อุปฺปชฺชติ. (1)}

\prose{อรูปี ธมฺโม อรูปิสฺส ธมฺมสฺส อธิปติ\-ปจฺจเยน ปจฺจโย. อารมฺมณา\-ธิปติ, สหชาตา\-ธิปติ.  อารมฺมณา\-ธิปติ\csromandash{}ทานํ ฯเปฯ. (สํขิตฺตํ.) นิพฺพานํ มคฺคสฺส, ผลสฺส อธิปติ\-ปจฺจเยน ปจฺจโย. อรูปิโน ขนฺเธ ครุํ กตฺวา อสฺสาเทติ ฯเปฯ ราโค อุปฺปชฺชติ. ทิฏฺฐิ อุปฺปชฺชติ. สหชาตา\-ธิปติ\csromandash{}อรูปี อธิปติ สมฺปยุตฺตกานํ ขนฺธานํ อธิปติ\-ปจฺจเยน ปจฺจโย.}

\prose{อรูปี ธมฺโม รูปิสฺส ธมฺมสฺส อธิปติ\-ปจฺจเยน ปจฺจโย. สหชาตา\-ธิปติ\csromandash{}อรูปี อธิปติ จิตฺต\-สมุฏฺฐานานํ รูปานํ อธิปติ\-ปจฺจเยน ปจฺจโย. (2)}

\prose{อรูปี ธมฺโม รูปิสฺส จ อรูปิสฺส จ ธมฺมสฺส อธิปติ\-ปจฺจเยน ปจฺจโย. สหชาตา\-ธิปติ\csromandash{}อรูปี อธิปติ สมฺปยุตฺตกานํ ขนฺธานํ จิตฺตสมุฏฺฐานานญฺจ รูปานํ อธิปติ\-ปจฺจเยน ปจฺจโย. (3)}

\csromanheader{2}{11. รูปีทุก \textbf{อนนฺตราทิ}}
\proseitem{115}{\csromanfolio{129}อรูปี ธมฺโม อรูปิสฺส ธมฺมสฺส อนนฺตร\-ปจฺจเยน ปจฺจโย. ปุริมา ปุริมา อรูปิโน ขนฺโธ ปจฺฉิมานํ ปจฺฉิมานํ อรูปีนํ ขนฺธานํ ฯเปฯ ผลสมาปตฺติยา อนนฺตร\-ปจฺจเยน ปจฺจโย. สมนนฺตร\-ปจฺจเยน ปจฺจโย. (สหชาต\-ปจฺจเย สตฺต. อิห ฆฏนา นตฺถิ. อญฺญมญฺญ\-ปจฺจเย ฉ. นิสฺสยปจฺจเย สตฺต ปญฺหา. อิห ฆฏนา นตฺถิ.)}

\csromanheader{2}{\textbf{อุปนิสฺสย}}
\proseitem{116}{รูปี ธมฺโม อรูปิสฺส ธมฺมสฺส อุปนิสฺสย\-ปจฺจเยน ปจฺจโย. อารมฺมณู\-ปนิสฺสโย, ปกตู\-ปนิสฺสโย ฯเปฯ. ปกตู\-ปนิสฺสโย\csromandash{}อุตุํ. โภชนํ. เสนาสนํ อุปนิสฺสาย ทานํ เทติ ฯเปฯ สํฆํ ภินฺทติ. อุตุ. โภชนํ. เสนาสนํ สทฺธาย ฯเปฯ ผลสมาปตฺติยา อุปนิสฺสย\-ปจฺจเยน ปจฺจโย. (1)}

\prose{อรูปี ธมฺโม อรูปิสฺส ธมฺมสฺส อุปนิสฺสย\-ปจฺจเยน ปจฺจโย. อารมฺมณู\-ปนิสฺสโย, อนนฺตรู\-ปนิสฺสโย, ปกตู\-ปนิสฺสโย ฯเปฯ. ปกตู\-ปนิสฺสโย\csromandash{}สทฺธํ อุปนิสฺสาย ทานํ เทติ. สีลํ ฯเปฯ. กายิกํ ทุกฺขํ อุปนิสฺสาย ทานํ เทติ ฯเปฯ ปํฆํ ภินฺทติ. สทฺธา ฯเปฯ. กายิกํ ทุกฺขํ. สทฺธาย ฯเปฯ ผลสมาปตฺติยา อุปนิสฺสย\-ปจฺจเยน ปจฺจโย. (1)}

\proseitem{117}{รูปี ธมฺโม อรูปิสฺส ธมฺมสฺส ปุเรชาต\-ปจฺจเยน ปจฺจโย. อารมฺมณ\-ปุเรชาตํ, วตฺถุ\-ปุเรชาตํ.  อารมฺมณ\-ปุเรชาตํ\csromandash{}จกฺขุํ ฯเปฯ วตฺถุํ ฯเปฯ. กพฬีการํ อาหารํ อนิจฺจโต ฯเปฯ โทมนสฺสํ อุปฺปชฺชติ. ทิพฺเพน จกฺขุนา รูปํ ปสฺสติ. ทิพฺพาย โสตธาตุยา สทฺทํ สุณาติ. รูปายตนํ จกฺขุ\-วิญฺญาณสฺส ฯเปฯ. โผฏฺฐพฺพา\-ยตนํ กายวิญฺญาณสฺส ฯเปฯ. วตฺถุ\-ปุเรชาตํ\csromandash{}จกฺขายตนํ จกฺขุ\-วิญฺญาณสฺส ฯเปฯ. กายายตนํ กายวิญฺญาณสฺส ฯเปฯ. วตฺถุ อรูปีนํ ขนฺธานํ ปุเรชาต\-ปจฺจเยน ปจฺจโย. (1)}

\csromanheader{2}{\textbf{ปจฺฉาชาต}}
\proseitem{118}{\csromanfolio{130}อรูปี ธมฺโม รูปิสฺส ธมฺมสฺส ปจฺฉาชาต\-ปจฺจเยน ปจฺจโย. ปจฺฉาชาตา อรูปิโน ขนฺธา ปุเรชาตสฺส อิมสฺส กายสฺส ปจฺฉาชาต\-ปจฺจเยน ปจฺจโย. (1)}

\csromanheader{2}{\textbf{อาเสวน}}
\proseitem{119}{อรูปี ธมฺโม อรูปิสฺส ธมฺมสฺส อาเสวน\-ปจฺจเยน ปจฺจโย. ปุริมา ปุริมา อรูปิโน ขนฺธา ปจฺฉิมานํ ปจฺฉิมานํ อรูปีนํ ขนฺธานํ อาเสวน ปจฺจเยน ปจฺจโย. (1)}

\proseitem{120}{อรูปี ธมฺโม อรูปิสฺส ธมฺมสฺส กมฺมปจฺจเยน ปจฺจโย. สหชาตา, นานากฺขณิกา.  สหชาตา อรูปี เจตนา สมฺปยุตฺตกานํ ขนฺธานํ กมฺมปจฺจเยน ปจฺจโย. ปฏิสนฺธิ\-กฺขเณ ฯเปฯ. นานากฺขณิกา อรูปี เจตนา วิปากานํ ขนฺธานํ กมฺมปจฺจเยน ปจฺจโย. (1)}

\prose{อรูปี ธมฺโม รูปิสฺส ธมฺมสฺส กมฺมปจฺจเยน ปจฺจโย. สหชาตา, นานากฺขณิกา.  สหชาตา อรูปี เจตนา จิตฺต\-สมุฏฺฐานานํ รูปานํ กมฺมปจฺจเยน ปจฺจโย. ปฏิสนฺธิ\-กฺขเณ ฯเปฯ. นานากฺขณิกา อรูปี เจตนา กฏตฺตารูปานํ กมฺมปจฺจเยน ปจฺจโย. (2)}

\prose{อรูปี ธมฺโม รูปิสฺส จ อรูปิสฺส จ ธมฺมสฺส กมฺมปจฺจเยน ปจฺจโย. สหชาตา, นานากฺขณิกา.  สหชาตา อรูปี เจตนา สมฺปยุตฺตกานํ ขนฺธานํ จิตฺตสมุฏฺฐานานญฺจ รูปานํ กมฺมปจฺจเยน ปจฺจโย. ปฏิสนฺธิ\-กฺขเณ ฯเปฯ. นานากฺขณิกา อรูปี เจตนา วิปากานํ ขนฺธานํ กฏตฺตา จ รูปานํ กมฺมปจฺจเยน ปจฺจโย. (3)}

\csromanheader{2}{\textbf{วิปากาหาร}}
\proseitem{121}{อรูปี ธมฺโม อรูปิสฺส ธมฺมสฺส วิปาก\-ปจฺจเยน ปจฺจโย. ตีณิ.}

\prose{รูปี ธมฺโม รูปิสฺส ธมฺมสฺส อาหารปจฺจเยน ปจฺจโย. กพฬีกาโร อาหาโร อิมสฺส กายสฺส อาหารปจฺจเยน ปจฺจโย. (1)}

\prose{อรูปี ธมฺโม อรูปิสฺส ธมฺมสฺส อาหารปจฺจเยน ปจฺจโย. ตีณิ.}

\csromanheader{2}{11. รูปีทุก \textbf{อินฺทฺริย}}
\proseitem{122}{\csromanfolio{131}รูปี ธมฺโม รูปิสฺส ธมฺมสฺส อินฺทฺริย\-ปจฺจเยน ปจฺจโย. รูปชีวิตินฺทฺริยํ กฏตฺตารูปานํ อินฺทฺริย\-ปจฺจเยน ปจฺจโย. (1)}

\prose{รูปี ธมฺโม อรูปิสฺส ธมฺมสฺส อินฺทฺริย\-ปจฺจเยน ปจฺจโย. จกฺขุนฺทฺริยํ ฯเปฯ. กายินฺทฺริยํ กายวิญฺญาณสฺส อินฺทฺริย\-ปจฺจเยน ปจฺจโย. (2)}

\prose{อรูปี ธมฺโม อรูปิสฺส ธมฺมสฺส อินฺทฺริย\-ปจฺจเยน ปจฺจโย. ตีณิ.}

\prose{รูปี จ อรูปี จ ธมฺมา อรูปิสฺส ธมฺมสฺส อินฺทฺริย\-ปจฺจเยน ปจฺจโย. จกฺขุนฺทฺริยญฺจ จกฺขุวิญฺญาณญฺจ จกฺขุ\-วิญฺญาณสหคตานํ ขนฺธานํ อินฺทฺริย\-ปจฺจเยน ปจฺจโย ฯเปฯ. กายินฺทฺริยญฺจ ฯเปฯ.}

\csromanheader{2}{\textbf{ฌานาทิ}}
\proseitem{123}{อรูปี ธมฺโม อรูปิสฺส ธมฺมสฺส ฌานปจฺจเยน ปจฺจโย. ตีณิ. มคฺคปจฺจเยน ปจฺจโย. ตีณิ. สมฺปยุตฺต\-ปจฺจเยน ปจฺจโย. เอกํ.}

\proseitem{124}{รูปี ธมฺโม อรูปิสฺส ธมฺมสฺส วิปฺปยุตฺต\-ปจฺจเยน ปจฺจโย. สหชาตํ, ปุเรชาตํ.  สหชาตํ ปฏิสนฺธิ\-กฺขเณ วตฺถุ อรูปีนํ ขนฺธานํ วิปฺปยุตฺต\-ปจฺจเยน ปจฺจโย. ปุเรชาตํ จกฺขายตนํ จกฺขุ\-วิญฺญาณสฺส ฯเปฯ กายายตนํ กายวิญฺญาณสฺส ฯเปฯ. วตฺถุ อรูปีนํ ขนฺธานํ วิปฺปยุตฺต\-ปจฺจเยน ปจฺจโย. (1)}

\prose{อรูปี ธมฺโม รูปิสฺส ธมฺมสฺส วิปฺปยุตฺต\-ปจฺจเยน ปจฺจโย. สหชาตํ ปจฺฉาชาตํ.  สหชาตา อรูปิโน ขนฺธา จิตฺต\-สมุฏฺฐานานํ รูปานํ วิปฺปยุตฺต\-ปจฺจเยน ปจฺจโย. ปฏิสนฺธิ\-กฺขเณ อรูปิโน ขนฺธา กฏตฺตารูปานํ วิปฺปยุตฺต\-ปจฺจเยน ปจฺจโย. อรูปิโน ขนฺธา วตฺถุสฺส วิปฺปยุตฺต\-ปจฺจเยน ปจฺจโย. ปจฺฉาชาตา อรูปิโน ขนฺธา ปุเรชาตสฺส อิมสฺส กายสฺส วิปฺปยุตฺต\-ปจฺจเยน ปจฺจโย. (1)}

\prose{\csromanfolio{132}อตฺถิอาทิ}

\proseitem{125}{รูปี ธมฺโม รูปิสฺส ธมฺมสฺส อตฺถิปจฺจเยน ปจฺจโย. สหชาตํ, อาหารํ, อินฺทฺริยํ.  \textbf{สหชาตํ} เอกํ มหาภูตํ ฯเปฯ. (ยาว อสญฺญสตฺตา.) \textbf{กพฬีกาโร อาหาโร} อิมสฺส กายสฺส อตฺถิปจฺจเยน ปจฺจโย. \textbf{รูปชีวินฺทฺริยํ} กฏตฺตารูปานํ อตฺถิปจฺจเยน ปจฺจโย. (1)}

\prose{รูปี ธมฺโม อรูปิสฺส ธมฺมสฺส อตฺถิปจฺจเยน ปจฺจโย. สหชาตํ, ปุเรชาตํ.  สหชาตํ ปฏิสนฺธิ\-กฺขเณ วตฺถุ อรูปีนํ ขนฺธานํ อตฺถิปจฺจเยน ปจฺจโย. ปุเรชาตํ จกฺขุํ ฯเปฯ. กพฬีการํ อาหารํ อนิจฺจโต ฯเปฯ โทมนสฺสํ อุปฺปชฺชติ. ทิพฺเพน จกฺขุนา รูปํ ปสฺสติ. ทิพฺพาย โสตธาตุยา สทฺทํ สุณาติ. รูปายตนํ จกฺขุ\-วิญฺญาณสฺส ฯเปฯ. โผฏฺฐพฺพา\-ยตนํ กายวิญฺญาณสฺส ฯเปฯ. จกฺขายตนํ จกฺขุ\-วิญฺญาณสฺส ฯเปฯ. กายายตนํ กายวิญฺญาณสฺส ฯเปฯ. วตฺถุ อรูปีนํ ขนฺธานํ อตฺถิปจฺจเยน ปจฺจโย. (2)}

\proseitem{126}{อรูปี ธมฺโม อรูปิสฺส ธมฺมสฺส อตฺถิปจฺจเยน ปจฺจโย. อรูปี เอโก ขนฺโธ ติณฺณนฺนํ ขนฺธานํ ฯเปฯ เทฺว ขนฺธา ฯเปฯ. ปฏิสนฺธิ\-กฺขเณ ฯเปฯ. (1)}

\prose{อรูปี ธมฺโม รูปิสฺส ธมฺมสฺส อตฺถิปจฺจเยน ปจฺจโย. สหชาตํ, ปจฺฉาชาตํ.  สหชาตา อรูปิโน ขนฺธา จิตฺต\-สมุฏฺฐานานํ รูปานํ อตฺถิปจฺจเยน ปจฺจโย. ปฏิสนฺธิ\-กฺขเณ ฯเปฯ. ปจฺฉาชาตา อรูปิโน ขนฺธา ปุเรชาตสฺส อิมสฺส กายสฺส อตฺถิปจฺจเยน ปจฺจโย. (2)}

\prose{อรูปี ธมฺโม รูปิสฺส จ อรูปิสฺส จ ธมฺมสฺส อตฺถิปจฺจเยน ปจฺจโย. อรูปี เอโก ขนฺโธ ติณฺณนฺนํ ขนฺธานํ จิตฺตสมุฏฺฐานานญฺจ รูปานํ อตฺถิปจฺจเยน ปจฺจโย ฯเปฯ เทฺว ขนฺธา ฯเปฯ. ปฏิสนฺธิ\-กฺขเณ ฯเปฯ. (3)}

\proseitem{127}{รูปี จ อรูปี จ ธมฺมา รูปิสฺส ธมฺมสฺส อตฺถิปจฺจเยน ปจฺจโย. สหชาตํ, ปจฺฉาชาตํ, อาหารํ, อินฺทฺริยํ.  สหชาตา อรูปี ขนฺธา จ มหาภูตา จ จิตฺต\-สมุฏฺฐานานํ รูปานํ อตฺถิปจฺจเยน ปจฺจโย. ปฏิสนฺธิ\-กฺขเณ ฯเปฯ. ปจฺฉาชาตา อรูปิโน ขนฺธา จ กพฬีกาโร อาหาโร จ อิมสฺส กายสฺส อตฺถิปจฺจเยน ปจฺจโย. ปจฺฉาชาตา อรูปิโน ขนฺธา จ รูปชีวิตินฺทฺริยญฺจ กฏตฺตารูปานํ อตฺถิปจฺจเยน ปจฺจโย. (1)}

\prose{\csromanfolio{133}รูปี จ อรูปี จ ธมฺมา อรูปิสฺส ธมฺมสฺส อตฺถิปจฺจเยน ปจฺจโย. สหชาตํ, ปุเรชาตํ.  สหชาโต จกฺขุ\-วิญฺญาณ\-สหคโต เอโก ขนฺโธ จ จกฺขายตนญฺจ ติณฺณนฺนํ ขนฺธานํ อตฺถิปจฺจเยน ปจฺจโย ฯเปฯ เทฺว ขนฺธา จ ฯเปฯ. กาย\-วิญฺญาณ\-สหคโต ฯเปฯ. อรูปี เอโก ขนฺโธ จ วตฺถุ จ ติณฺณนฺนํ ขนฺธานํ อตฺถิปจฺจเยน ปจฺจโย ฯเปฯ เทฺว ขนฺธา จ ฯเปฯ. ปฏิสนฺธิ\-กฺขเณ อรูปี เอโก ขนฺโธ จ วตฺถุ จ ติณฺณนฺนํ ขนฺธานํ ฯเปฯ เทฺว ขนฺธา จ ฯเปฯ. (2)}

\prose{นตฺถิ\-ปจฺจเยน ปจฺจโย. วิคต\-ปจฺจเยน ปจฺจโย. อวิคต\-ปจฺจเยน ปจฺจโย.}

\csromanmarkh{1. ปจฺจยานุโลม}\csromanmarkhii{2. สงฺขฺยาวาร}\csromanheadernomark{2}{1. \textbf{ปจฺจยานุโลม 2. สงฺขฺยาวาร}}
\proseitem{128}{เหตุยา ตีณิ, อารมฺมเณ เทฺว, อธิปติยา จตฺตาริ, อนนฺตเร เอกํ, สมนนฺตเร เอกํ, สหชาเต สตฺต, อญฺญมญฺเญ ฉ, นิสฺสเย สตฺต, อุปนิสฺสเย เทฺว, ปุเรชาเต เอกํ, ปจฺฉาชาเต เอกํ, อาเสวเน เอกํ, กมฺเม ตีณิ, วิปาเก ตีณิ, อาหาเร จตฺตาริ, อินฺทฺริเย ฉ, ฌาเน ตีณิ, มคฺเค ตีณิ, สมฺปยุตฺเต เอกํ, วิปฺปยุตฺเต เทฺว, อตฺถิยา สตฺต, นตฺถิยา เอกํ, วิคเต เอกํ, อวิคเต สตฺต.}

\vspace{\tipitakaparskip}
\csromancenter{อนุโลมํ.}
\proseitem{129}{รูปี ธมฺโม รูปิสฺส ธมฺมสฺส สหชาต\-ปจฺจเยน ปจฺจโย. อาหารปจฺจเยน ปจฺจโย. อินฺทฺริย\-ปจฺจเยน ปจฺจโย. (1)}

\prose{รูปี ธมฺโม อรูปิสฺส ธมฺมสฺส อารมฺมณ\-ปจฺจเยน ปจฺจโย. สหชาต\-ปจฺจเยน ปจฺจโย. อุปนิสฺสย\-ปจฺจเยน ปจฺจโย. ปุเรชาต\-ปจฺจเยน ปจฺจโย. (2)}

\prose{อรูปี ธมฺโม อรูปิสฺส ธมฺมสฺส อารมฺมณ\-ปจฺจเยน ปจฺจโย. สหชาต\-ปจฺจเยน ปจฺจโย. อุปนิสฺสย\-ปจฺจเยน ปจฺจโย. กมฺมปจฺจเยน ปจฺจโย. (1)}

\prose{\csromanfolio{134}อรูปี ธมฺโม รูปิสฺส ธมฺมสฺส สหชาต\-ปจฺจเยน ปจฺจโย. ปจฺฉาชาต\-ปจฺจเยน ปจฺจโย. กมฺมปจฺจเยน ปจฺจโย. (2)}

\prose{อรูปี ธมฺโม รูปิสฺส จ อรูปิสฺส จ ธมฺมสฺส สหชาต\-ปจฺจเยน ปจฺจโย. กมฺมปจฺจเยน ปจฺจโย. (3)}

\prose{รูปี จ อรูปี จ ธมฺมา รูปิสฺส ธมฺมสฺส สหชาตํ, ปจฺฉาชาตํ, อาหารํ, อินฺทฺริยํ. (1)}

\prose{รูปี จ อรูปี จ ธมฺมา อรูปิสฺส ธมฺมสฺส สหชาตํ, ปุเรชาตํ.}

\csromanmarkh{2. ปจฺจย\-ปจฺจนีย}\csromanmarkhii{2. สงฺขฺยาวาร}\csromanheadernomark{2}{2. ปจฺจย\-ปจฺจนีย 2. สงฺขฺยาวาร}
\proseitem{130}{นเหตุยา สตฺต, นอารมฺมเณ สตฺต, นอธิปติยา สตฺต นอนนฺตเร สตฺต, นสมนนฺตเร สตฺต, นสหชาเต ฉ, นอญฺญมญฺเญ ฉ, นนิสฺสเย ฉ, นอุปนิสฺสเย สตฺต, นปุเรชาเต สตฺต ฯเปฯ นมคฺเค สตฺต, นสมฺปยุตฺเต ฉ, นวิปฺปยุตฺเต ปญฺจ, โนอตฺถิยา จตฺตาริ, โนนตฺถิยา สตฺต, โนวิคเต สตฺต, โนอวิคเต จตฺตาริ.}

\vspace{\tipitakaparskip}
\csromancenter{ปจฺจนียํ.}
\proseitem{131}{เหตุปจฺจยา นอารมฺมเณ ตีณิ, นอธิปติยา ตีณิ, นอนนฺตเร ตีณิ, นสมนนฺตเร ตีณิ, นอญฺญมญฺเญ เอกํ, นอุปนิสฺสเย ตีณิ, (สพฺพตฺถ ตีณิ.) นสมฺปยุตฺเต เอกํ, นวิปฺปยุตฺเต เอกํ, โนนตฺถิยา ตีณิ, โนวิคเต ตีณิ.}

\vspace{\tipitakaparskip}
\csromancenter{อนุโลม\-ปจฺจนียํ.}
\csromanmatikaanchor{mtk.15}
\csromantocmarkref{section}{12. โลกิยทุก}{mtk.15}
\bookmark[dest={mtk.15},level=1]{12. โลกิยทุก}
\csromanmarkh{12. โลกิยทุก}\csromanmarkhii{4. ปจฺจย\-ปจฺจนียา\-นุโลม}\csromanheadernomark{1}{12. โลกิยทุก 4. ปจฺจย\-ปจฺจนียา\-นุโลม}
\proseitem{132}{\csromanfolio{135}นเหตุปจฺจยา อารมฺมเณ เทฺว, อธิปติยา จตฺตาริ, (อนุโลมมาติกา กาตพฺพา.) อวิคเต สตฺต.}

\vspace{\tipitakaparskip}
\csromancenter{ปจฺจนียา\-นุโลมํ.}
\nitthitam{รูปีทุกํ นิฏฺฐิตํ.}
\csromanmarkh{12. โลกิยทุก}\csromanmarkhii{1. ปฏิจฺจวาร}\csromanheadernomark{2}{12. โลกิยทุก 1. ปฏิจฺจวาร}
\csromanmarkh{1. ปจฺจยานุโลม}\csromanmarkhii{1. วิภงฺควาร}\csromanheadernomark{2}{1. \textbf{ปจฺจยานุโลม 1. วิภงฺควาร}}
\csromanheader{2}{\textbf{เหตุ}}
\proseitem{133}{โลกิยํ ธมฺมํ ปฏิจฺจ โลกิโย ธมฺโม อุปฺปชฺชติ เหตุปจฺจยา. โลกิยํ เอกํ ขนฺธํ ปฏิจฺจ ตโย ขนฺธา จิตฺต\-สมุฏฺฐานญฺจ รูปํ ฯเปฯ เทฺว ขนฺเธ ฯเปฯ. ปฏิสนฺธิ\-กฺขเณ ฯเปฯ. ขนฺเธ ปฏิจฺจ วตฺถุ, วตฺถุํ ปฏิจฺจ ขนฺธา. เอกํ มหาภูตํ ฯเปฯ มหาภูเต ปฏิจฺจ จิตฺต\-สมุฏฺฐานํ รูปํ, กฏตฺตารูปํ, อุปาทารูปํ. (1)}

\proseitem{134}{โลกุตฺตรํ ธมฺมํ ปฏิจฺจ โลกุตฺตโร ธมฺโม อุปฺปชฺชติ เหตุปจฺจยา. โลกุตฺตรํ เอกํ ขนฺธํ ปฏิจฺจ ตโย ขนฺธา ฯเปฯ เทฺว ขนฺเธ ฯเปฯ. (1)}

\prose{โลกุตฺตรํ ธมฺมํ ปฏิจฺจ โลกิโย ธมฺโม อุปฺปชฺชติ เหตุปจฺจยา. โลกุตฺตเร ขนฺเธ ปฏิจฺจ จิตฺต\-สมุฏฺฐานํ รูปํ. (2)}

\prose{โลกุตฺตรํ ธมฺมํ ปฏิจฺจ โลกิโย จ โลกุตฺตโร จ ธมฺมา อุปฺปชฺชนฺติ เหตุปจฺจยา. โลกุตฺตรํ เอกํ ขนฺธํ ปฏิจฺจ ตโย ขนฺธา จิตฺต\-สมุฏฺฐานญฺจ รูปํ ฯเปฯ เทฺว ขนฺเธ ฯเปฯ. (3)}

\prose{โลกิยญฺจ โลกุตฺตรญฺจ ธมฺมํ ปฏิจฺจ โลกิโย ธมฺโม อุปฺปชฺชติ เหตุปจฺจยา. โลกุตฺตเร ขนฺเธ จ มหาภูเต จ ปฏิจฺจ จิตฺต\-สมุฏฺฐานํ รูปํ. (สํขิตฺตํ.) (1)}

\csromanmarkh{1. ปจฺจยานุโลม}\csromanmarkhii{2. สงฺขฺยาวาร}\csromanheadernomark{2}{1. \textbf{ปจฺจยานุโลม 2. สงฺขฺยาวาร}}
\proseitem{135}{\csromanfolio{136}เหตุยา ปญฺจ, อารมฺมเณ, เทฺว, อธิปติยา ปญฺจ, อนนฺตเร เทฺว, สมนนฺตเร เทฺว, สหชาเต ปญฺจ, อญฺญมญฺเญ เทฺว, นิสฺสเย ปญฺจ, อุปนิสฺสเย เทฺว, ปุเรชาเต เทฺว, อาเสวเน เทฺว, กมฺเม ปญฺจ, วิปาเก ปญฺจ, อาหาเร ปญฺจ, อินฺทฺริเย ปญฺจ, ฌาเน ปญฺจ, มคฺเค ปญฺจ, สมฺปยุตฺเต เทฺว, วิปฺปยุตฺเต ปญฺจ, อตฺถิยา ปญฺจ, นตฺถิยา เทฺว, วิคเต เทฺว, อวิคเต ปญฺจ.}

\vspace{\tipitakaparskip}
\csromancenter{อนุโลมํ.}
\csromanmarkh{2. ปจฺจย\-ปจฺจนีย}\csromanmarkhii{1. วิภงฺควาร}\csromanheadernomark{2}{2. ปจฺจย\-ปจฺจนีย 1. วิภงฺควาร}
\proseitem{136}{โลกิยํ ธมฺมํ ปฏิจฺจ โลกิโย ธมฺโม อุปฺปชฺชติ นเหตุปจฺจยา. อเหตุกํ โลกิยํ เอกํ ขนฺธํ ปฏิจฺจ ตโย ขนฺธา จิตฺต\-สมุฏฺฐานญฺจ รูปํ ฯเปฯ เทฺว ขนฺเธ ฯเปฯ. อเหตุก\-ปฏิสนฺธิ\-กฺขเณ ฯเปฯ. (ยาว อสญฺญสตฺตา.) วิจิกิจฺฉา\-สหคเต อุทฺธจฺจ\-สหคเต ขนฺเธ ปฏิจฺจ วิจิกิจฺฉา\-สหคโต อุทฺธจฺจ\-สหคโต โมโห. (สํขิตฺตํ.)}

\csromanmarkh{2. ปจฺจย\-ปจฺจนีย}\csromanmarkhii{2. สงฺขฺยาวาร}\csromanheadernomark{2}{2. ปจฺจย\-ปจฺจนีย 2. สงฺขฺยาวาร}
\vspace{\tipitakaparskip}
\csromancenter{นเหตุยา เอกํ, นอารมฺมเณ ตีณิ, นอธิปติยา เทฺว, นอนนฺตเร ตีณิ, นสมนนฺตเร ตีณิ, นอญฺญมญฺเญ ตีณิ, นฺเอาปนิสฺสเย ตีณิ, นปุเรชาเต จตฺตาริ, นปจฺฉาชาเต ปญฺจ, นอาเสวเน ปญฺจ, (นอาเสวน\-มูลเก โลกุตฺตเร สุทฺธเก “วิปาโก”ติ นิยาเมตพฺพํ. อวเสสา ปกติกาเยว.) นกมฺเม เทฺว, นวิปาเก ปญฺจ, นอาหาเร เอกํ, ไนนฺทฺริเย เอกํ, นฌาเน เอกํ, นมคฺเค เอกํ, นสมฺปยุตฺเต ตีณิ, นวิปฺปยุตฺเต เทฺว, โนนตฺถิยา ตีณิ, โนวิคเต ตีณิ.}
\csromancenter{ปจฺจนียํ.}
\csromanmarkh{12. โลกิยทุก}\csromanmarkhii{3. ปจฺจยา\-นุโลม\-ปจฺจนีย}\csromanheadernomark{2}{12. โลกิยทุก 3. ปจฺจยา\-นุโลม\-ปจฺจนีย}
\vspace{\tipitakaparskip}
\csromancenter{เหตุปจฺจยา นอารมฺมเณ ตีณิ, นอธิปติยา เทฺว, (นอนนฺตร\-ปทาที ปจฺจนีย\-สทิสา.) ฯเปฯ นวิปาเก ปญฺจ, นสมฺปยุตฺเต ตีณิ, นวิปฺปยุตฺเต เทฺว, โนนตฺถิยา ตีณิ, โนวิคเต ตีณิ.}
\csromancenter{อนุโลม\-ปจฺจนียํ.}
\csromancenter{นเหตุปจฺจยา อารมฺมเณ เอกํ, อนนฺตเร เอกํ ฯเปฯ อวิคเต เอกํ.}
\csromancenter{ปจฺจนียา\-นุโลมํ.}
\csromancenter{(สหชาตวาโร ปฏิจฺจ\-วาร\-สทิโส.)}
\csromanmarkh{12. โลกิยทุก}\csromanmarkhii{3. ปจฺจยวาร}\csromanheadernomark{2}{12. โลกิยทุก 3. ปจฺจยวาร}
\csromanmarkh{1. ปจฺจยานุโลม}\csromanmarkhii{1. วิภงฺควาร}\csromanheadernomark{2}{1. \textbf{ปจฺจยานุโลม 1. วิภงฺควาร}}
\csromanheader{2}{\textbf{เหตุ}}
\proseitem{140}{\csromanfolio{137}โลกิยํ ธมฺมํ ปจฺจยา โลกิโย ธมฺโม อุปฺปชฺชติ เหตุปจฺจยา. โลกิยํ เอกํ ขนฺธํ ปจฺจยา ฯเปฯ. เอกํ มหาภูตํ ฯเปฯ มหาภูเต ปจฺจยา จิตฺต\-สมุฏฺฐานํ รูปํ, กฏตฺตารูปํ, อุปาทารูปํ. วตฺถุํ ปจฺจยา โลกิยา ขนฺธา. (1)}

\prose{โลกิยํ ธมฺมํ ปจฺจยา โลกุตฺตโร ธมฺโม อุปฺปชฺชติ เหตุปจฺจยา. วตฺถุํ ปจฺจยา โลกุตฺตรา ขนฺธา. (2)}

\prose{\csromanfolio{138}โลกิยํ ธมฺมํ ปจฺจยา โลกิโย จ โลกุตฺตโร จ ธมฺมา อุปฺปชฺชนฺติ เหตุปจฺจยา. วตฺถุํ ปจฺจยา โลกุตฺตรา ขนฺธา. มหาภูเต ปจฺจยา จิตฺต\-สมุฏฺฐานํ รูปํ. (3)}

\proseitem{141}{โลกุตฺตรํ ธมฺมํ ปจฺจยา โลกุตฺตโร ธมฺโม อุปฺปชฺชติ เหตุปจฺจยา. ตีณิ.}

\prose{โลกิยญฺจ โลกุตฺตรญฺจ ธมฺมํ ปจฺจยา โลกิโย ธมฺโม อุปฺปชฺชติ เหตุปจฺจยา. โลกุตฺตเร ขนฺเธ จ มหาภูเต จ ปจฺจยา จิตฺต\-สมุฏฺฐานํ รูปํ. (1)}

\prose{โลกิยญฺจ โลกุตฺตรญฺจ ธมฺมํ ปจฺจยา โลกุตฺตโร ธมฺโม อุปฺปชฺชติ เหตุปจฺจยา. โลกุตฺตรํ เอกํ ขนฺธญฺจ วตฺถุญฺจ ปจฺจยา ตโย ขนฺธา ฯเปฯ เทฺว ขนฺเธ ฯเปฯ. (2)}

\prose{โลกิยญฺจ โลกุตฺตรญฺจ ธมฺมํ ปจฺจยา โลกิโย จ โลกุตฺตโร จ ธมฺมา อุปฺปชฺชนฺติ เหตุปจฺจยา. โลกุตฺตรํ เอกํ ขนฺธญฺจ วตฺถุญฺจ ปจฺจยา ตโย ขนฺธา ฯเปฯ เทฺว ขนฺเธ ฯเปฯ. โลกุตฺตเร ขนฺเธ จ มหาภูเต จ ปจฺจยา จิตฺต\-สมุฏฺฐานํ รูปํ. (สํขิตฺตํ.) (3)}

\csromanmarkh{1. ปจฺจยานุโลม}\csromanmarkhii{2. สงฺขฺยาวาร}\csromanheadernomark{2}{1. \textbf{ปจฺจยานุโลม 2. สงฺขฺยาวาร}}
\proseitem{142}{เหตุยา นว, อารมฺมเณ จตฺตาริ, อธิปติยา นว, อนนฺตเร จตฺตาริ, สมนนฺตเร จตฺตาริ, สหชาเต นว, อญฺญมญฺเญ จตฺตาริ, นิสฺสเย นว, อุปนิสฺสเย จตฺตาริ, ปุเรชาเต จตฺตาริ, อาเสวเน จตฺตาริ, กมฺเม นว, วิปาเก นว ฯเปฯ มคฺเค นว, สมฺปยุตฺเต จตฺตาริ วิปฺปยุตฺเต นว, อตฺถิยา นว, นตฺถิยา จตฺตาริ, วิคเต จตฺตาริ, อวิคเต นว.}

\vspace{\tipitakaparskip}
\csromancenter{อนุโลมํ.}
\csromanmarkh{12. โลกิยทุก}\csromanmarkhii{2. ปจฺจย\-ปจฺจนีย 1. วิภงฺควาร}\csromanheadernomark{2}{12. โลกิยทุก 2. ปจฺจย\-ปจฺจนีย 1. วิภงฺควาร}
\proseitem{143}{\csromanfolio{139}โลกิยํ ธมฺมํ ปจฺจยา โลกิโย ธมฺโม อุปฺปชฺชติ นเหตุปจฺจยา. อเหตุกํ โลกิยํ เอกํ ขนฺธํ ฯเปฯ. (ยาว อสญฺญสตฺตา.) จกฺขายตนํ ปจฺจยา จกฺขุวิญฺญาณํ ฯเปฯ. กายายตนํ ปจฺจยา กายวิญฺญาณํ. วตฺถุํ ปจฺจยา อเหตุกา โลกิยา ขนฺธา. วิจิกิจฺฉา\-สหคเต อุทฺธจฺจ\-สหคเต ขนฺเธ จ วตฺถุญฺจ ปจฺจยา วิจิกิจฺฉา\-สหคโต อุทฺธจฺจ\-สหคโต โมโห. (สํขิตฺตํ.)}

\csromanmarkh{2. ปจฺจย\-ปจฺจนีย}\csromanmarkhii{2. สงฺขฺยาวาร}\csromanheadernomark{2}{2. ปจฺจย\-ปจฺจนีย 2. สงฺขฺยาวาร}
\vspace{\tipitakaparskip}
\csromancenter{\textbf{นเหตุ}ยา เอกํ, นอารมฺมเณ ตีณิ, นอธิปติยา จตฺตาริ, นอนนฺตเร ตีณิ ฯเปฯ นฺเอาปนิสฺสเย ตีณิ, นปุเรชาเต จตฺตาริ, นปจฺฉาชาเต นว, นอาเสวเน นว, (โลกุตฺตเร อรูเป วิปากนฺติ นิยาเมตพฺพํ.) นกมฺเม จตฺตาริ, นวิปาเก นว, นอาหาเร เอกํ, ไนนฺทฺริเย เอกํ, นฌาเน เอกํ, นมคฺเค เอกํ, นสมฺปยุตฺเต ตีณิ, นวิปฺปยุตฺเต เทฺว, โนนตฺถิยา ตีณิ, โนวิคเต ตีณิ.}
\csromancenter{ปจฺจนียํ.}
\proseitem{145}{เหตุปจฺจโย นอารมฺมเณ ตีณิ, นอธิปติยา จตฺตาริ, นอนนฺตเร ตีณิ. (นสมนนฺตร\-ปทาที ปจฺจนีย\-สทิสา.) นวิปาเก นว, นสมฺปยุตฺเต ตีณิ, นวิปฺปยุตฺเต เทฺว, โนนตฺถิยา ตีณิ, โนวิคเต ตีณิ.}

\vspace{\tipitakaparskip}
\csromancenter{อนุโลม\-ปจฺจนียํ.}
\csromancenter{นเหตุปจฺจยา อารมฺมเณ เอกํ, อนนฺตเร เอกํ ฯเปฯ อวิคเต เอกํ.}
\csromancenter{ปจฺจนียา\-นุโลมํ.}
\csromancenter{(นิสฺสยวาโร ปจฺจยวาร\-สทิโส)}
\csromanmarkh{12. โลกิยทุก}\csromanmarkhii{5. สํสฏฺฐวาร}\csromanheadernomark{2}{12. \textbf{โลกิยทุก 5. สํสฏฺฐวาร}}
\csromanmarkh{1. ปจฺจยานุโลม}\csromanmarkhii{1. วิภงฺควาร}\csromanheadernomark{2}{1. \textbf{ปจฺจยานุโลม 1. วิภงฺควาร}}
\csromanheader{2}{\textbf{เหตุ}}
\proseitem{147}{\csromanfolio{140}โลกิยํ ธมฺมํ สํสฏฺโฐ โลกิโย ธมฺโม อุปฺปชฺชติ เหตุปจฺจยา. โลกิยํ เอกํ ขนฺธํ สํสฏฺฐา ตโย ขนฺธา ฯเปฯ เทฺว ขนฺเธ ฯเปฯ. ปฏิสนฺธิ\-กฺขเณ ฯเปฯ. (1)}

\prose{โลกุตฺตรํ ธมฺมํ สํสฏฺโฐ โลกุตฺตโร ธมฺโม อุปฺปชฺชติ เหตุปจฺจยา. โลกุตฺตรํ เอกํ ขนฺธํ สํสฏฺฐา ตโย ขนฺธา ฯเปฯ เทฺว ขนฺเธ ฯเปฯ. (1)}

\prose{(สํสฏฺฐวาโร เอวํ วิตฺถาเรตพฺโพ, สห คณนาหิ เทฺว ปญฺหา.)}

\vspace{\tipitakaparskip}
\csromancenter{(สมฺปยุตฺตวาโร สํสฏฺฐวาร\-สทิโส.)}
\csromanmarkh{12. โลกิยทุก}\csromanmarkhii{7. ปญฺหาวาร}\csromanheadernomark{2}{12. \textbf{โลกิยทุก 7. ปญฺหาวาร}}
\csromanmarkh{1. ปจฺจยานุโลม}\csromanmarkhii{1. วิภงฺควาร}\csromanheadernomark{2}{1. \textbf{ปจฺจยานุโลม 1. วิภงฺควาร}}
\csromanheader{2}{\textbf{เหตุ}}
\proseitem{148}{โลกิโย ธมฺโม โลกิยสฺส ธมฺมสฺส เหตุปจฺจเยน ปจฺจโย. โลกิยา เหตู สมฺปยุตฺตกานํ ขนฺธานํ จิตฺตสมุฏฺฐานานญฺจ รูปานํ เหตุปจฺจเยน ปจฺจโย. ปฏิสนฺธิ\-กฺขเณ ฯเปฯ. (1)}

\vspace{\tipitakaparskip}
\csromancenter{โลกุตฺตโร ธมฺโม โลกุตฺตรสฺส ธมฺมสฺส เหตุปจฺจเยน ปจฺจโย. ตีณิ.}
\proseitem{149}{\csromanfolio{141}โลกิโย ธมฺโม โลกิยสฺส ธมฺมสฺส อารมฺมณ\-ปจฺจเยน ปจฺจโย. ทานํ ฯเปฯ สีลํ ฯเปฯ อุโปสถกมฺมํ กตฺวา ตํ ปจฺจเวกฺขติ. ปุพฺเพ สุจิณฺณานิ ปจฺจเวกฺขติ. ฌานา ฯเปฯ. อริยา โคตฺรภุํ ปจฺจเวกฺขนฺติ. โวทานํ ปจฺจเวกฺขนฺติ. ปหีเน กิเลเส ปจฺจเวกฺขนฺติ. วิกฺขมฺภิเต กิเลเส ฯเปฯ. ปุพฺเพ สมุทาจิณฺเณ กิเลเส ชานนฺติ. จกฺขุํ ฯเปฯ วตฺถุํ โลกิเย ขนฺเธ อนิจฺจโต ฯเปฯ โทมนสฺสํ อุปฺปชฺชติ. ทิพฺเพน จกฺขุนา รูปํ ปสฺสติ. ทิพฺพาย โสตธาตุยา สทฺทํ สุณาติ. เจโตปริย\-ญาเณน โลกิยจิตฺต\-สมงฺคิสฺส จิตฺตํ ชานาติ. อากาสา\-นญฺจา\-ยตนํ วิญฺญาณญฺจา\-ยตนสฺส ฯเปฯ. อากิญฺจญฺญา\-ยตนํ เนว\-สญฺญา\-นา\-สญฺญา\-ยตนสฺส ฯเปฯ. รูปายตนํ จกฺขุ\-วิญฺญาณสฺส ฯเปฯ. โผฏฺฐพฺพา\-ยตนํ กายวิญฺญาณสฺส. โลกิยา ขนฺธา อิทฺธิวิธ\-ญาณสฺส, เจโตปริย\-ญาณสฺส, ปุพฺเพ\-นิวาสา\-นุสฺสติ\-ญาณสฺส, ยถากมฺมูปคญาณสฺส, อนาคตํสญาณสฺส, อาวชฺชนาย อารมฺมณ\-ปจฺจเยน ปจฺจโย. (1)}

\proseitem{150}{โลกุตฺตโร ธมฺโม โลกุตฺตรสฺส ธมฺมสฺส อารมฺมณ\-ปจฺจเยน ปจฺจโย. นิพฺพานํ มคฺคสฺส, ผลสฺส อารมฺมณ\-ปจฺจเยน ปจฺจโย. (1)}

\prose{โลกุตฺตโร ธมฺโม โลกิยสฺส ธมฺมสฺส อารมฺมณ\-ปจฺจเยน ปจฺจโย. อริยา มคฺคา วุฏฺฐหิตฺวา มคฺคํ ปจฺจเวกฺขนฺติ, ผลํ ปจฺจเวกฺขนฺติ, นิพฺพานํ ปจฺจเวกฺขนฺติ. นิพฺพานํ โคตฺรภุสฺส, โวทานสฺส, อาวชฺชนาย, อารมฺมณ\-ปจฺจเยน ปจฺจโย. อริยา เจโตปริย\-ญาเณน โลกุตฺตรจิตฺต\-สมงฺคิสฺส จิตฺตํ ชานนฺติ. โลกุตฺตรา ขนฺธา เจโตปริย\-ญาณสฺส, ปุพฺเพ\-นิวาสา\-นุสฺสติ\-ญาณสฺส, อนาคตํสญาณสฺส, อาวชฺชนาย อารมฺมณ\-ปจฺจเยน ปจฺจโย. (2)}

\proseitem{151}{โลกิโย ธมฺโม โลกิยสฺส ธมฺมสฺส อธิปติ\-ปจฺจเยน ปจฺจโย. อารมฺมณา\-ธิปติ, สหชาตา\-ธิปติ.  อารมฺมณา\-ธิปติ\csromandash{}ทานํ ทตฺวา, สีลํ ฯเปฯ อุโปสถกมฺมํ ฯเปฯ. ปุพฺเพ ฯเปฯ. ฌานา ฯเปฯ. เสกฺขา โคตฺรภุํ ครุํ กตฺวา ปจฺจเวกฺขนฺติ, โวทานํ ครุํ กตฺวา ปจฺจเวกฺขนฺติ. จกฺขุํ ฯเปฯ. \csromanfolio{142}วตฺถุํ โลกิเย ขนฺเธ ครุํ กตฺวา อสฺสาเทติ อภินนฺทติ, ตํ ครุํ กตฺวา ราโค อุปฺปชฺชติ, ทิฏฺฐิ อุปฺปชฺชติ.  สหชาตา\-ธิปติ\csromandash{}โลกิยาธิปติ สมฺปยุตฺตกานํ ขนฺธานํ จิตฺตสมุฏฺฐานานญฺจ รูปานํ อธิปติ\-ปจฺจเยน ปจฺจโย. (1)}

\proseitem{152}{โลกุตฺตโร ธมฺโม โลกุตฺตรสฺส ธมฺมสฺส อธิปติ\-ปจฺจเยน ปจฺจโย. อารมฺมณา\-ธิปติ, สหชาตา\-ธิปติ.  อารมฺมณา\-ธิปติ\csromandash{} นิพฺพานํ มคฺคสฺส, ผลสฺส อธิปติ\-ปจฺจเยน ปจฺจโย. สหชาตา\-ธิปติ\csromandash{} โลกุตฺตรา\-ธิปติ สมฺปยุตฺตกานํ ขนฺธานํ อธิปติ\-ปจฺจเยน ปจฺจโย. (1)}

\prose{โลกุตฺตโร ธมฺโม โลกิยสฺส ธมฺมสฺส อธิปติ\-ปจฺจเยน ปจฺจโย อารมฺมณา\-ธิปติ, สหชาตา\-ธิปติ.  อารมฺมณา\-ธิปติ\csromandash{}อริยา มคฺคา วุฏฺฐหิตฺวา มคฺคํ ครุํ กตฺวา ปจฺจเวกฺขนฺติ, ผลํ ครุํ กตฺวา ปจฺจเวกฺขนฺติ, นิพฺพานํ ครุํ กตฺวา ปจฺจเวกฺขนฺติ. นิพฺพานํ โคตฺรภุสฺส, โวทานสฺส อธิปติ\-ปจฺจเยน ปจฺจโย. สหชาตา\-ธิปติ\csromandash{} โลกุตฺตรา\-ธิปติ จิตฺต\-สมุฏฺฐานานํ รูปานํ อธิปติ\-ปจฺจเยน ปจฺจโย.}

\prose{โลกุตฺตโร ธมฺโม โลกิยสฺส จ โลกุตฺตรสฺส จ ธมฺมสฺส อธิปติ\-ปจฺจเยน ปจฺจโย.  สหชาตา\-ธิปติ\csromandash{}โลกุตฺตรา\-ธิปติ สมฺปยุตฺตกานํ ขนฺธานํ จิตฺตสมุฏฺฐานานญฺจ รูปานํ อธิปติ\-ปจฺจเยน ปจฺจโย. (3)}

\proseitem{153}{โลกิโย ธมฺโม โลกิยสฺส ธมฺมสฺส อนนฺตร\-ปจฺจเยน ปจฺจโย. ปุริมา ปุริมา โลกิยา ขนฺธา ปจฺฉิมานํ ปจฺฉิมานํ โลกิยานํ ขนฺธานํ ฯเปฯ. อนุโลมํ โคตฺรภุสฺส. อนุโลมํ โวทานสฺส อนนฺตร\-ปจฺจเยน ปจฺจโย. (1)}

\prose{โลกิโย ธมฺโม โลกุตฺตรสฺส ธมฺมสฺส อนนฺตร\-ปจฺจเยน ปจฺจโย. โคตฺรภุ มคฺคสฺส. โวทานํ มคฺคสฺส. อนุโลมํ ผลสมาปตฺติยา. นิโรธา วุฏฺฐหนฺตสฺส เนว\-สญฺญา\-นา\-สญฺญา\-ยตนํ ผลสมาปตฺติยา อนนฺตร\-ปจฺจเยน ปจฺจโย. (2)}

\proseitem{154}{โลกุตฺตโร ธมฺโม โลกุตฺตรสฺส ธมฺมสฺส อนนฺตร\-ปจฺจเยน ปจฺจโย. ปุริมา ปุริมา โลกุตฺตรา ขนฺธา ปจฺฉิมานํ ปจฺฉิมานํ โลกุตฺตรานํ \csromanfolio{143}ขนฺธานํ อนนฺตร\-ปจฺจเยน ปจฺจโย. มคฺโค ผลสฺส. ผลํ ผลสฺส อนนฺตร\-ปจฺจเยน ปจฺจโย. (1)}

\prose{โลกุตฺตโร ธมฺโม โลกิยสฺส ธมฺมสฺส อนนฺตร\-ปจฺจเยน ปจฺจโย. ผลํ วุฏฺฐานสฺส อนนฺตร\-ปจฺจเยน ปจฺจโย. (2)}

\csromanheader{2}{\textbf{สมนนฺตราทิ}}
\proseitem{155}{โลกิโย ธมฺโม โลกิยสฺส ธมฺมสฺส สมนนฺตร\-ปจฺจเยน ปจฺจโย. สหชาต\-ปจฺจเยน ปจฺจโย. (ปญฺจ ปญฺหา ฆฏนา นตฺถิ.) อญฺญมญฺญ\-ปจฺจเยน ปจฺจโย. เทฺว. นิสฺสย\-ปจฺจเยน ปจฺจโย.}

\csromanheader{2}{\textbf{อุปนิสฺสย}}
\proseitem{156}{โลกิโย ธมฺโม โลกิยสฺส ธมฺมสฺส อุปนิสฺสย\-ปจฺจเยน ปจฺจโย. อารมฺมณู\-ปนิสฺสโย, อนนฺตรู\-ปนิสฺสโย, ปกตู\-ปนิสฺสโย ฯเปฯ. ปกตู\-ปนิสฺสโย\csromandash{}โลกิยํ สทฺธํ อุปนิสฺสาย ทานํ เทติ ฯเปฯ วิปสฺสนํ อุปฺปาเทติ, อภิญฺญํ อุปฺปาเทติ, สมาปตฺติํ อุปฺปาเทติ. มานํ ชปฺเปติ. ทิฏฺฐิํ คณฺหาติ. โลกิยํ สีลํ ฯเปฯ. เสนาสนํ อุปนิสฺสาย ทานํ เทติ ฯเปฯ สํฆํ ภินฺทติ. โลกิยา สทฺธา ฯเปฯ. เสนาสนํ โลกิยาย สทฺธาย ฯเปฯ กายิกสฺส ทุกฺขสฺส อุปนิสฺสย\-ปจฺจเยน ปจฺจโย. กุสลากุสลํ กมฺมํ วิปากสฺส อุปนิสฺสย\-ปจฺจเยน ปจฺจโย. (1)}

\prose{โลกิโย ธมฺโม โลกุตฺตรสฺส ธมฺมสฺส อุปนิสฺสย\-ปจฺจเยน ปจฺจโย. อนนฺตรู\-ปนิสฺสโย, ปกตู\-ปนิสฺสโย ฯเปฯ. ปกตู\-ปนิสฺสโย\csromandash{} ปฐมสฺส มคฺคสฺส ปริกมฺมํ ปฐมสฺส มคฺคสฺส อุปนิสฺสย\-ปจฺจเยน ปจฺจโย ฯเปฯ. จตุตฺถสฺส มคฺคสฺส ปริกมฺมํ จตุตฺถสฺส มคฺคสฺส อุปนิสฺสย\-ปจฺจเยน ปจฺจโย. (2)}

\proseitem{157}{โลกุตฺตโร ธมฺโม โลกุตฺตรสฺส ธมฺมสฺส อุปนิสฺสย\-ปจฺจเยน ปจฺจโย. อารมฺมณู\-ปนิสฺสโย, อนนฺตรู\-ปนิสฺสโย, ปกตู\-ปนิสฺสโย ฯเปฯ. ปกตู\-ปนิสฺสโย\csromandash{}ปฐโม มคฺโค ทุติยสฺส มคฺคสฺส อุปนิสฺสย\-ปจฺจเยน ปจฺจโย ฯเปฯ. ตติโย มคฺโค จตุตฺถสฺส มคฺคสฺส อุปนิสฺสย\-ปจฺจเยน ปจฺจโย. (1)}

\prose{\csromanfolio{144}โลกุตฺตโร ธมฺโม โลกิยสฺส ธมฺมสฺส อุปนิสฺสย\-ปจฺจเยน ปจฺจโย. อารมฺมณู\-ปนิสฺสโย, อนนฺตรู\-ปนิสฺสโย, ปกตู\-ปนิสฺสโย ฯเปฯ. ปกตู\-ปนิสฺสโย\csromandash{}อริยา มคฺคํ อุปนิสฺสาย อนุปฺปนฺนํ สมาปตฺติํ อุปฺปาเทนฺติ, อุปฺปนฺนํ สมาปชฺชนฺติ. สงฺขาเร อนิจฺจโต ทุกฺขโต อนตฺตโต วิปสฺสนฺติ. อริยานํ มคฺโค ฯเปฯ. ฐานาฐาน\-โกสลฺลสฺส อุปนิสฺสย\-ปจฺจเยน ปจฺจโย. ผลสมาปตฺติ กายิกสฺส สุขสฺส อุปนิสฺสย\-ปจฺจเยน ปจฺจโย. (2)}

\proseitem{158}{โลกิโย ธมฺโม โลกิยสฺส ธมฺมสฺส ปุเรชาต\-ปจฺจเยน ปจฺจโย. อารมฺมณ\-ปุเรชาตํ, วตฺถุ\-ปุเรชาตํ.  อารมฺมณ\-ปุเรชาตํ\csromandash{}จกฺขุํ ฯเปฯ วตฺถุํ อนิจฺจโต ทุกฺขโต อนตฺตโต ฯเปฯ โทมนสฺสํ อุปฺปชฺชติ. ทิพฺเพน จกฺขุนา รูปํ ปสฺสติ. ทิพฺพาย โสตธาตุยา สทฺทํ สุณาติ. รูปายตนํ จกฺขุ\-วิญฺญาณสฺส ฯเปฯ. โผฏฺฐพฺพา\-ยตนํ กายวิญฺญาณสฺส ปุเรชาต\-ปจฺจเยน ปจฺจโย. วตฺถุ\-ปุเรชาตํ\csromandash{}จกฺขายตนํ จกฺขุ\-วิญฺญาณสฺส ฯเปฯ. กายาตนํ กายวิญฺญาณสฺส ฯเปฯ. วตฺถุ โลกิยานํ ขนฺธานํ ปุเรชาต\-ปจฺจเยน ปจฺจโย. (1)}

\prose{โลกิโย ธมฺโม โลกุตฺตรสฺส ธมฺมสฺส ปุเรชาต\-ปจฺจเยน ปจฺจโย. วตฺถุ\-ปุเรชาตํ\csromandash{}วตฺถุ โลกุตฺตรานํ ขนฺธานํ ปุเรชาต\-ปจฺจเยน ปจฺจโย. (2)}

\csromanheader{2}{\textbf{ปจฺฉาชาต}}
\proseitem{159}{โลกิโย ธมฺโม โลกิยสฺส ธมฺมสฺส ปจฺฉาชาต\-ปจฺจเยน ปจฺจโย. ปจฺฉาชาตา โลกิยา ขนฺธา ปุเรชาตสฺส อิมสฺส กายสฺส ปจฺฉาชาต\-ปจฺจเยน ปจฺจโย. (1)}

\prose{โลกุตฺตโร ธมฺโม โลกิยสฺส ธมฺมสฺส ปจฺฉาชาต\-ปจฺจเยน ปจฺจโย. ปจฺฉาชาตา โลกุตฺตรา ขนฺธา ปุเรชาตสฺส อิมสฺส กายสฺส ปจฺฉาชาต\-ปจฺจเยน ปจฺจโย. (1)}

\csromanheader{2}{12. โลกิยทุก \textbf{อาเสวน}}
\proseitem{160}{\csromanfolio{145}โลกิโย ธมฺโม โลกิยสฺส ธมฺมสฺส อาเสวน\-ปจฺจเยน ปจฺจโย. ปุริมา ปุริมา โลกิยา ขนฺธา ปจฺฉิมานํ ปจฺฉิมานํ โลกิยานํ ขนฺธานํ อาเสวน\-ปจฺจเยน ปจฺจโย. อนุโลมํ โคตฺรภุสฺส. อนุโลมํ โวทานสฺส อาเสวน\-ปจฺจเยน ปจฺจโย. (1)}

\prose{โลกิโย ธมฺโม โลกุตฺตรสฺส ธมฺมสฺส อาเสวน\-ปจฺจเยน ปจฺจโย. โคตฺรภุ มคฺคสฺส. โวทานํ มคฺคสฺส อาเสวน\-ปจฺจเยน ปจฺจโย. (2)}

\proseitem{161}{โลกิโย ธมฺโม โลกิยสฺส ธมฺมสฺส กมฺมปจฺจเยน ปจฺจโย. สหชาตา, นานากฺขณิกา.  สหชาตา โลกิยา เจตนา สมฺปยุตฺตกานํ ขนฺธานํ จิตฺตสมุฏฺฐานานญฺจ รูปานํ กมฺมปจฺจเยน ปจฺจโย. ปฏิสนฺธิ\-กฺขเณ ฯเปฯ. นานากฺขณิกา โลกิยา เจตนา วิปากานํ ขนฺธานํ กฏตฺตา จ รูปานํ กมฺมปจฺจเยน ปจฺจโย. (1)}

\prose{โลกุตฺตโร ธมฺโม โลกุตฺตรสฺส ธมฺมสฺส กมฺมปจฺจเยน ปจฺจโย. สหชาตา, นานากฺขณิกา.  สหชาตา โลกุตฺตรา เจตนา สมฺปยุตฺตกานํ ขนฺธานํ กมฺมปจฺจเยน ปจฺจโย. นานากฺขณิกา โลกุตฺตรา เจตนา วิปากานํ ขนฺธานํ กมฺมปจฺจเยน ปจฺจโย. (1)}

\prose{โลกุตฺตโร ธมฺโม โลกิยสฺส ธมฺมสฺส กมฺมปจฺจเยน ปจฺจโย. โลกุตฺตรา เจตนา จิตฺต\-สมุฏฺฐานานํ รูปานํ กมฺมปจฺจเยน ปจฺจโย.}

\prose{โลกุตฺตโร ธมฺโม โลกิยสฺส จ โลกุตฺตรสฺส จ ธมฺมสฺส กมฺมปจฺจเยน ปจฺจโย. โลกุตฺตรา เจตนา สมฺปยุตฺตกานํ ขนฺธานํ จิตฺตสมุฏฺฐานานญฺจ รูปานํ กมฺมปจฺจเยน ปจฺจโย. (3)}

\csromanheader{2}{\textbf{วิปาก}}
\proseitem{162}{โลกิโย ธมฺโม โลกิยสฺส ธมฺมสฺส วิปาก\-ปจฺจเยน ปจฺจโย. วิปาโก โลกิโย เอโก ขนฺโธ ติณฺณนฺนํ ขนฺธานํ จิตฺตสมุฏฺฐานานญฺจ รูปานํ วิปาก\-ปจฺจเยน ปจฺจโย ฯเปฯ เทฺว ขนฺธา ฯเปฯ. ปฏิสนฺธิ\-กฺขเณ ฯเปฯ. (1)}

\prose{\csromanfolio{146}โลกุตฺตโร ธมฺโม โลกุตฺตรสฺส ธมฺมสฺส วิปาก\-ปจฺจเยน ปจฺจโย. ตีณิ.}

\csromanheader{2}{\textbf{อาหาร}}
\proseitem{163}{โลกิโย ธมฺโม โลกิยสฺส ธมฺมสฺส อาหารปจฺจเยน ปจฺจโย. โลกิยา อาหารา สมฺปยุตฺตกานํ ขนฺธานํ จิตฺต\-สมุฏฺฐานญฺจ รูปานํ อาหารปจฺจเยน ปจฺจโย. ปฏิสนฺธิ\-กฺขเณ ฯเปฯ. กพฬีกาโร, อาหาโร อิมสฺส กายสฺส อาหารปจฺจเยน ปจฺจโย. (1)}

\prose{โลกุตฺตโร ธมฺโม โลกุตฺตรสฺส ธมฺมสฺส อาหารปจฺจเยน ปจฺจโย. ตีณิ.}

\csromanheader{2}{\textbf{อินฺทฺริย}}
\proseitem{164}{โลกิโย ธมฺโม โลกิยสฺส ธมฺมสฺส อินฺทฺริย\-ปจฺจเยน ปจฺจโย. (ปฏิสนฺธิ กาตพฺพา.) จกฺขุนฺทฺริยํ จกฺขุ\-วิญฺญาณสฺส ฯเปฯ. กายินฺทฺริยํ กายวิญฺญาณสฺส อินฺทฺริย\-ปจฺจเยน ปจฺจโย. รูปชีวิตินฺทฺริยํ กฏตฺตารูปานํ อินฺทฺริย\-ปจฺจเยน ปจฺจโย. (1)}

\prose{โลกุตฺตโร ธมฺโม โลกุตฺตรสฺส ธมฺมสฺส อินฺทฺริย\-ปจฺจเยน ปจฺจโย. ตีณิ.}

\csromanheader{2}{\textbf{ฌานาทิ}}
\proseitem{165}{โลกิโย ธมฺโม โลกิยสฺส ธมฺมสฺส ฌานปจฺจเยน ปจฺจโย. เอกํ. โลกุตฺตโร ธมฺโม ฯเปฯ. ตีณิ. มคฺคปจฺจเยน ปจฺจโย. โลกิเย เอกํ. โลกุตฺตเร ตีณิ.}

\prose{โลกิโย ธมฺโม โลกิยสฺส ธมฺมสฺส สมฺปยุตฺต\-ปจฺจเยน ปจฺจโย. เอกํ. โลกุตฺตโร ธมฺโม ฯเปฯ. เอกํ.}

\proseitem{166}{โลกิโย ธมฺโม โลกิยสฺส ธมฺมสฺส วิปฺปยุตฺต\-ปจฺจเยน ปจฺจโย. สหชาตํ, ปุเรชาตํ, ปจฺฉาชาตํ.  สหชาตา โลกิยา ขนฺธา จิตฺต\-สมุฏฺฐานานํ รูปานํ วิปฺปยุตฺต\-ปจฺจเยน ปจฺจโย. ปฏิสนฺธิ\-กฺขเณ ฯเปฯ. ขนฺธา วตฺถุสฺส วิปฺปยุตฺต\-ปจฺจเยน ปจฺจโย. วตฺถุ ขนฺธานํ วิปฺปยุตฺต\-ปจฺจเยน ปจฺจโย. ปุเรชาตํ จกฺขายตนํ จกฺขุ\-วิญฺญาณสฺส ฯเปฯ. \csromanfolio{147}กายายตนํ กายวิญฺญาณสฺส ฯเปฯ. วตฺถุ โลกิยานํ ขนฺธานํ วิปฺปยุตฺต\-ปจฺจเยน ปจฺจโย. ปจฺฉาชาตา โลกิยา ขนฺธา ปุเรชาตสฺส อิมสฺส กายสฺส วิปฺปยุตฺต\-ปจฺจเยน ปจฺจโย. (1)}

\prose{โลกิโย ธมฺโม โลกุตฺตรสฺส ธมฺมสฺส วิปฺปยุตฺต\-ปจฺจเยน ปจฺจโย. ปุเรชาตํ วตฺถุ โลกุตฺตรานํ ขนฺธานํ วิปฺปยุตฺต\-ปจฺจเยน ปจฺจโย.}

\prose{โลกุตฺตโร ธมฺโม โลกิยสฺส ธมฺมสฺส วิปฺปยุตฺต\-ปจฺจเยน ปจฺจโย. สหชาตํ, ปจฺฉาชาตํ.  สหชาตา โลกุตฺตรา ขนฺธา จิตฺต\-สมุฏฺฐานานํ รูปานํ วิปฺปยุตฺต\-ปจฺจเยน ปจฺจโย. ปจฺฉาชาตา โลกุตฺตรา ขนฺธา ปุเรชาตสฺส อิมสฺส กายสฺส วิปฺปยุตฺต\-ปจฺจเยน ปจฺจโย. (1)}

\prose{อตฺถิอาทิ}

\proseitem{167}{โลกิโย ธมฺโม โลกิยสฺส ธมฺมสฺส อตฺถิปจฺจเยน ปจฺจโย. สหชาตํ, ปุเรชาตํ, ปจฺฉาชาตํ, อาหารํ, อินฺทฺริยํ.  สหชาโต โลกิโย เอโก ขนฺโธ ติณฺณนฺนํ ขนฺธานํ จิตฺตสมุฏฺฐานานญฺจ รูปานํ ฯเปฯ. (ยาว อสญฺญสตฺตา.) ปุเรชาตํ จกฺขุํ ฯเปฯ วตฺถุํ ฯเปฯ. (ปุเรชาต\-สทิสํ.) วตฺถุ โลกิยานํ ขนฺธานํ อตฺถิปจฺจเยน ปจฺจโย. ปจฺฉาชาตา โลกิยา ขนฺธา ปุเรชาตสฺส อิมสฺส กายสฺส อตฺถิปจฺจเยน ปจฺจโย. กพฬีกาโร อาหาโร อิมสฺส กายสฺส อตฺถิปจฺจเยน ปจฺจโย. รูปชีวิตินฺทฺริยํ กฏตฺตารูปานํ อตฺถิปจฺจเยน ปจฺจโย. (1)}

\prose{โลกิโย ธมฺโม โลกุตฺตรสฺส ธมฺมสฺส อตฺถิปจฺจเยน ปจฺจโย. \textbf{ปุเรชาตํ} วตฺถุ โลกุตฺตรานํ ขนฺธานํ อตฺถิปจฺจเยน ปจฺจโย. (2)}

\proseitem{168}{โลกุตฺตโร ธมฺโม โลกุตฺตรสฺส ธมฺมสฺส อตฺถิปจฺจเยน ปจฺจโย. โลกุตฺตโร เอโก ขนฺโธ ติณฺณนฺนํ ขนฺธานํ อตฺถิปจฺจเยน ปจฺจโย ฯเปฯ เทฺว ขนฺธา ฯเปฯ. (1)}

\prose{โลกุตฺตโร ธมฺโม โลกิยสฺส ธมฺมสฺส อตฺถิปจฺจเยน ปจฺจโย. สหชาตํ, ปจฺฉาชาตํ.  สหชาตา โลกุตฺตรา ขนฺธา จิตฺต\-สมุฏฺฐานานํ รูปานํ อตฺถิปจฺจเยน ปจฺจโย. ปจฺฉาชาตา โลกุตฺตรา ขนฺธา ปุเรชาตสฺส อิมสฺส กายสฺส อตฺถิปจฺจเยน ปจฺจโย. (2)}

\prose{\csromanfolio{148}โลกุตฺตโร ธมฺโม โลกิยสฺส จ โลกุตฺตรสฺส จ ธมฺมสฺส อตฺถิปจฺจเยน ปจฺจโย. โลกุตฺตโร เอโก ขนฺโธ ติณฺณนฺนํ ขนฺธานํ จิตฺตสมุฏฺฐานานญฺจ รูปานํ อตฺถิปจฺจเยน ปจฺจโย ฯเปฯ เทฺว ขนฺธา ฯเปฯ. (3)}

\proseitem{169}{โลกิโย จ โลกุตฺตโร จ ธมฺมา โลกิยสฺส ธมฺมสฺส อตฺถิปจฺจเยน ปจฺจโย. สหชาตํ, ปจฺฉาชาตํ, อาหารํ, อินฺทฺริยํ. สหชาตา โลกุตฺตรา ขนฺธา จ มหาภูตา จ จิตฺต\-สมุฏฺฐานานํ รูปานํ อตฺถิปจฺจเยน ปจฺจโย. ปจฺฉาชาตา โลกุตฺตรา ขนฺธา จ กพฬีกาโร อาหาโร จ อิมสฺส กายสฺส อตฺถิปจฺจเยน ปจฺจโย. ปจฺฉาชาตา โลกุตฺตรา ขนฺธา จ รูปชีวิตินฺทฺริยญฺจ กฏตฺตารูปานํ อตฺถิปจฺจเยน ปจฺจโย. (1)}

\prose{โลกิโย จ โลกุตฺตโร จ ธมฺมา โลกุตฺตรสฺส ธมฺมสฺส อตฺถิปจฺจเยน ปจฺจโย. สหชาตํ, ปุเรชาตํ.  สหชาโต โลกุตฺตโร เอโก ขนฺโธ จ วตฺถุ จ ติณฺณนฺนํ ขนฺธานํ อตฺถิปจฺจเยน ปจฺจโย ฯเปฯ เทฺว ขนฺธา จ ฯเปฯ. (2)}

\prose{นตฺถิ\-ปจฺจเยน ปจฺจโย. วิคต\-ปจฺจเยน ปจฺจโย. อวิคต\-ปจฺจเยน ปจฺจโย. (2)}

\csromanmarkh{1. ปจฺจยานุโลม}\csromanmarkhii{2. สงฺขฺยาวาร}\csromanheadernomark{2}{1. \textbf{ปจฺจยานุโลม 2. สงฺขฺยาวาร}}
\proseitem{170}{เหตุยา จตฺตาริ, อารมฺมเณ ตีณิ, อธิปติยา จตฺตาริ, อนนฺตเร จตฺตาริ, สมนนฺตเร จตฺตาริ, สหชาเต ปญฺจ, อญฺญมญฺเญ เทฺว, นิสฺสเย สตฺต, อุปนิสฺสเย จตฺตาริ, ปุเรชาเต เทฺว, ปจฺฉาชาเต เทฺว, อาเสวเน เทฺว, กมฺเม จตฺตาริ, วิปาเก จตฺตาริ, อาหาเร จตฺตาริ, อินฺทฺริเย จตฺตาริ, ฌาเน จตฺตาริ, มคฺเค จตฺตาริ, สมฺปยุตฺเต เทฺว, วิปฺปยุตฺเต ตีณิ, อตฺถิยา สตฺต, นตฺถิยา จตฺตาริ, วิคเต จตฺตาริ, อวิคเต สตฺต.}

\vspace{\tipitakaparskip}
\csromancenter{\textbf{อนุโลมํ.}}
\proseitem{171}{โลกิโย ธมฺโม โลกิยสฺส ธมฺมสฺส อารมฺมณ\-ปจฺจเยน ปจฺจโย. สหชาต\-ปจฺจเยน ปจฺจโย. อุปนิสฺสย\-ปจฺจเยน ปจฺจโย. \csromanfolio{149}ปุเรชาต\-ปจฺจเยน ปจฺจโย. ปจฺฉาชาต\-ปจฺจเยน ปจฺจโย. กมฺมปจฺจเยน ปจฺจโย. อาหารปจฺจเยน ปจฺจโย. อินฺทฺริย\-ปจฺจเยน ปจฺจโย. (1)}

\prose{โลกิโย ธมฺโม โลกุตฺตรสฺส ธมฺมสฺส อุปนิสฺสย\-ปจฺจเยน ปจฺจโย. ปุเรชาต\-ปจฺจเยน ปจฺจโย. (2)}

\proseitem{172}{โลกุตฺตโร ธมฺโม โลกุตฺตรสฺส ธมฺมสฺส อารมฺมณ\-ปจฺจเยน ปจฺจโย. สหชาต\-ปจฺจเยน ปจฺจโย. อุปนิสฺสย\-ปจฺจเยน ปจฺจโย. (1)}

\prose{โลกุตฺตโร ธมฺโม โลกิยสฺส ธมฺมสฺส อารมฺมณ\-ปจฺจเยน ปจฺจโย. สหชาต\-ปจฺจเยน ปจฺจโย. อุปนิสฺสย\-ปจฺจเยน ปจฺจโย. ปจฺฉาชาต\-ปจฺจเยน ปจฺจโย. (2)}

\prose{โลกุตฺตโร ธมฺโม โลกิยสฺส จ โลกุตฺตรสฺส จ ธมฺมสฺส สหชาต\-ปจฺจเยน ปจฺจโย. (3)}

\prose{โลกิโย จ โลกุตฺตโร จ ธมฺมา โลกิยสฺส ธมฺมสฺส สหชาตํ. ปจฺฉาชาตํ. อาหารํ. อินฺทฺริยํ. (1)}

\prose{โลกิโย จ โลกุตฺตโร จ ธมฺมา โลกุตฺตรสฺส ธมฺมสฺส สหชาตํ. ปุเรชาตํ. (2)}

\csromanmarkh{2. ปจฺจย\-ปจฺจนีย}\csromanmarkhii{2. สงฺขฺยาวาร}\csromanheadernomark{2}{2. ปจฺจย\-ปจฺจนีย 2. สงฺขฺยาวาร}
\proseitem{173}{นเหตุยา สตฺต ฯเปฯ นสมนนฺตเร สตฺต, นสหชาเต ปญฺจ, นอญฺญมญฺเญ ปญฺจ, นนิสฺสเย ปญฺจ, นอุปนิสฺสเย สตฺต, นปุเรชาเต ฉ, นปจฺฉาชาเต สตฺต ฯเปฯ นมคฺเค สตฺต, นสมฺปยุตฺเต ปญฺจ, นวิปฺปยุตฺเต จตฺตาริ, โนอตฺถิยา จตฺตาริ, โนนตฺถิยา สตฺต, โนวิคเต สตฺต, โนอวิคเต จตฺตาริ.}

\vspace{\tipitakaparskip}
\csromancenter{ปจฺจนียํ.}
\csromancenter{\textbf{เหตุ}ปจฺจยา นอารมฺมเณ จตฺตาริ ฯเปฯ นสมนนฺตเร จตฺตาริ, นอญฺญมญฺเญ เทฺว, นฺเอาปนิสฺสเย จตฺตาริ ฯเปฯ นมคฺเค จตฺตาริ, นสมฺปยุตฺเต เทฺว, นวิปฺปยุตฺเต เทฺว, โนนตฺถิยา จตฺตาริ, โนวิคเต จตฺตาริ.}
\csromancenter{อนุโลม\-ปจฺจนียํ.}
\proseitem{175}{\csromanfolio{150}นเหตุปจฺจยา อารมฺมเณ ตีณิ, อธิปติยา จตฺตาริ, (อนุโลมมาติกา กาตพฺพา.) อวิคเต สตฺต.}

\vspace{\tipitakaparskip}
\csromancenter{ปจฺจนียา\-นุโลมํ.}
\nitthitam{โลกิยทุกํ นิฏฺฐิตํ.}
\csromanmatikaanchor{mtk.16}
\csromantocmarkref{section}{13. เกนจิวิญฺเญยฺยทุก}{mtk.16}
\bookmark[dest={mtk.16},level=1]{13. เกนจิวิญฺเญยฺยทุก}
\csromanmarkh{13. เกนจิ\-วิญฺเญยฺยทุก}\csromanmarkhii{1. ปฏิจฺจวาร}\csromanheadernomark{1}{13. เกนจิ\-วิญฺเญยฺยทุก 1. ปฏิจฺจวาร}
\csromanmarkh{1. ปจฺจยานุโลม}\csromanmarkhii{1. วิภงฺควาร}\csromanheadernomark{2}{1. \textbf{ปจฺจยานุโลม 1. วิภงฺควาร}}
\csromanheader{2}{\textbf{เหตุ}}
\proseitem{176}{เกนจิ วิญฺเญยฺยํ ธมฺมํ ปฏิจฺจ เกนจิ วิญฺเญยฺโย ธมฺโม อุปฺปชฺชติ เหตุปจฺจยา. เกนจิ วิญฺเญยฺยํ เอกํ ขนฺธํ ปฏิจฺจ ตโย ขนฺธา จิตฺต\-สมุฏฺฐานญฺจ รูปํ ฯเปฯ เทฺว ขนฺเธ ฯเปฯ. ปฏิสนฺธิ\-กฺขเณ ฯเปฯ. ขนฺเธ ปฏิจฺจ วตฺถุ, วตฺถุํ ปฏิจฺจ ขนฺธา. เอกํ มหาภูตํ ปฏิจฺจ ตโย มหาภูตา ฯเปฯ มหาภูเต ปฏิจฺจ จิตฺต\-สมุฏฺฐานํ รูปํ, กฏตฺตารูปํ, อุปาทารูปํ. (1)}

\prose{เกนจิ วิญฺเญยฺยํ ธมฺมํ ปฏิจฺจ เกนจิ นวิญฺเญยฺโย ธมฺโม อุปฺปชฺชติ เหตุปจฺจยา. เกนจิ วิญฺเญยฺยํ เอกํ ขนฺธํ ปฏิจฺจ เกนจิ นวิญฺเญยฺยา}

\vspace{\tipitakaparskip}
\csromancenter{เกนจิ\-วิญฺเญยฺยทุก ตโย ขนฺธา จิตฺต\-สมุฏฺฐานญฺจ รูปํ ฯเปฯ เทฺว ขนฺเธ ฯเปฯ. ปฏิสนฺธิ\-กฺขเณ ฯเปฯ. ขนฺเธ ปฏิจฺจ วตฺถุ, วตฺถุํ ปฏิจฺจ ขนฺธา. เอกํ มหาภูตํ ฯเปฯ มหาภูเต ปฏิจฺจ จิตฺต\-สมุฏฺฐานํ รูปํ, กฏตฺตารูปํ, อุปาทารูปํ. (2)}
\prose{\csromanfolio{151}เกนจิ วิญฺเญยฺยํ ธมฺมํ ปฏิจฺจ เกนจิ วิญฺเญยฺโย จ เกนจิ นวิญฺเญยฺโย จ ธมฺมา อุปฺปชฺชนฺติ เหตุปจฺจยา. เกนจิ วิญฺเญยฺยํ เอกํ ขนฺธํ ปฏิจฺจ เกนจิ วิญฺเญยฺยา จ เกนจิ นวิญฺเญยฺยา จ ตโย ขนฺธา จิตฺต\-สมุฏฺฐานญฺจ รูปํ ฯเปฯ เทฺว ขนฺเธ ฯเปฯ. ปฏิสนฺธิ\-กฺขเณ ฯเปฯ. ขนฺเธ ปฏิจฺจ วตฺถุ, วตฺถุํ ปฏิจฺจ ขนฺธา. เอกํ มหาภูตํ ฯเปฯ มหาภูเต ปฏิจฺจ จิตฺต\-สมุฏฺฐานํ รูปํ, กฏตฺตารูปํ, อุปาทารูปํ. (3)}

\proseitem{177}{เกนจิ นวิญฺเญยฺยํ ธมฺมํ ปฏิจฺจ เกนจิ นวิญฺเญยฺโย ธมฺโม อุปฺปชฺชติ เหตุปจฺจยา. เกนจิ นวิญฺเญยฺยํ เอกํ ขนฺธํ ปฏิจฺจ เกนจิ นวิญฺเญยฺยา ตโย ขนฺธา จิตฺต\-สมุฏฺฐานญฺจ รูปํ ฯเปฯ เทฺว ขนฺเธ ฯเปฯ. ปฏิสนฺธิ\-กฺขเณ ฯเปฯ. ขนฺเธ ปฏิจฺจ วตฺถุ, วตฺถุํ ปฏิจฺจ ขนฺธา. เอกํ มหาภูตํ ฯเปฯ มหาภูเต ปฏิจฺจ จิตฺต\-สมุฏฺฐานํ รูปํ, กฏตฺตารูปํ, อุปาทารูปํ. (1)}

\prose{เกนจิ นวิญฺเญยฺยํ ธมฺมํ ปฏิจฺจ เกนจิ วิญฺเญยฺโย ธมฺโม อุปฺปชฺชติ เหตุปจฺจยา. เกนจิ นวิญฺเญยฺยํ เอกํ ขนฺธํ ปฏิจฺจ เกนจิ วิญฺเญยฺยา ตโย ขนฺธา จิตฺต\-สมุฏฺฐานญฺจ รูปํ ฯเปฯ เทฺว ขนฺเธ ฯเปฯ. ปฏิสนฺธิ\-กฺขเณ ฯเปฯ. ขนฺเธ ปฏิจฺจ วตฺถุ, วตฺถุํ ปฏิจฺจ ขนฺธา. เอกํ มหาภูตํ ฯเปฯ มหาภูเต ปฏิจฺจ จิตฺต\-สมุฏฺฐานํ รูปํ, กฏตฺตารูปํ, อุปาทารูปํ. (2)}

\prose{เกนจิ นวิญฺเญยฺยํ ธมฺมํ ปฏิจฺจ เกนจิ วิญฺเญยฺโย จ เกนจิ นวิญฺเญยฺโย จ ธมฺมา อุปฺปชฺชนฺติ เหตุปจฺจยา. เกนจิ นวิญฺเญยฺยํ เอกํ ขนฺธํ ปฏิจฺจ เกนจิ วิญฺเญยฺยา จ เกนจิ นวิญฺเญยฺยา จ ตโย ขนฺธา จิตฺต\-สมุฏฺฐานญฺจ รูปํ ฯเปฯ เทฺว ขนฺเธ ฯเปฯ. ปฏิสนฺธิ\-กฺขเณ ฯเปฯ. ขนฺเธ ปฏิจฺจ วตฺถุ, วตฺถุํ ปฏิจฺจ ขนฺธา. เอกํ มหาภูตํ ฯเปฯ มหาภูเต ปฏิจฺจ จิตฺต\-สมุฏฺฐานํ รูปํ, กฏตฺตารูปํ, อุปาทารูปํ. (3)}

\proseitem{178}{เกนจิ วิญฺเญยฺยญฺจ เกนจิ นวิญฺเญยฺยญฺจ ธมฺมํ ปฏิจฺจ เกนจิ วิญฺเญยฺโย ธมฺโม อุปฺปชฺชติ เหตุปจฺจยา. เกนจิ วิญฺเญยฺยญฺจ เกนจิ นวิญฺเญยฺยญฺจ เอกํ ขนฺธํ ปฏิจฺจ เกนจิ วิญฺเญยฺยา ตโย ขนฺธา จิตฺต\-สมุฏฺฐานญฺจ รูปํ ฯเปฯ เทฺว ขนฺเธ ฯเปฯ. ปฏิสนฺธิ\-กฺขเณ ฯเปฯ. ขนฺเธ ปฏิจฺจ วตฺถุ, วตฺถุํ ปฏิจฺจ \csromanfolio{152}ขนฺธา. เอกํ มหาภูตํ ฯเปฯ มหาภูเต ปฏิจฺจ จิตฺต\-สมุฏฺฐานํ รูปํ, กฏตฺตารูปํ, อุปาทารูปํ. (1)}

\prose{เกนจิ วิญฺเญยฺยญฺจ เกนจิ นวิญฺเญยฺยญฺจ ธมฺมํ ปฏิจฺจ เกนจิ นวิญฺเญยฺโย ธมฺโม อุปฺปชฺชติ เหตุปจฺจยา. เกนจิ วิญฺเญยฺยญฺจ เกนจิ นวิญฺเญยฺยญฺจ เอกํ ขนฺธํ ปฏิจฺจ เกนจิ นวิญฺเญยฺยา ตโย ขนฺธา จิตฺต\-สมุฏฺฐานญฺจ รูปํ ฯเปฯ เทฺว ขนฺเธ ฯเปฯ. ปฏิสนฺธิ\-กฺขเณ ฯเปฯ. ขนฺเธ ปฏิจฺจ วตฺถุ, วตฺถุํ ปฏิจฺจ ขนฺธา. เอกํ มหาภูตํ ฯเปฯ มหาภูเต ปฏิจฺจ จิตฺต\-สมุฏฺฐานํ รูปํ, กฏตฺตารูปํ, อุปาทารูปํ. (2)}

\prose{เกนจิ วิญฺเญยฺยญฺจ เกนจิ นวิญฺเญยฺยญฺจ ธมฺมํ ปฏิจฺจ เกนจิ วิญฺเญยฺโย จ เกนจิ นวิญฺเญยฺโย จ ธมฺมา อุปฺปชฺชนฺติ เหตุปจฺจยา. เกนจิ วิญฺเญยฺยญฺจ เกนจิ นวิญฺเญยฺยญฺจ เอกํ ขนฺธํ ปฏิจฺจ เกนจิ วิญฺเญยฺยา จ เกนจิ นวิญฺเญยฺยา จ ตโย ขนฺธา จิตฺต\-สมุฏฺฐานญฺจ รูปํ ฯเปฯ เทฺว ขนฺเธ ฯเปฯ. ปฏิสนฺธิ\-กฺขเณ ฯเปฯ. ขนฺเธ ปฏิจฺจ วตฺถุ, วตฺถุํ ปฏิจฺจ ขนฺธา. เอกํ มหาภูตํ ฯเปฯ มหาภูเต ปฏิจฺจ จิตฺต\-สมุฏฺฐานํ รูปํ, กฏตฺตารูปํ, อุปาทารูปํ. (สํขิตฺตํ.) (3)}

\csromanmarkh{1. ปจฺจยานุโลม}\csromanmarkhii{2. สงฺขฺยาวาร}\csromanheadernomark{2}{1. \textbf{ปจฺจยานุโลม 2. สงฺขฺยาวาร}}
\vspace{\tipitakaparskip}
\csromancenter{เหตุยา นว, อารมฺมเณ นว ฯเปฯ อวิคเต นว.}
\csromancenter{อนุโลมํ.}
\csromanmarkh{2. ปจฺจย\-ปจฺจนีย}\csromanmarkhii{2. สงฺขฺยาวาร}\csromanheadernomark{2}{2. ปจฺจย\-ปจฺจนีย 2. สงฺขฺยาวาร}
\proseitem{180}{นเหตุยา นว, นอารมฺมเณ นว ฯเปฯ โนวิคเต นว. (เอวํ จตฺตาริปิ คณนา ปริปุณฺณา.) (สหชาตวาโรปิ ปจฺจยวาโรปิ นิสฺสยวาโรปิ สํสฏฺฐ\-วาโรปิ สมฺปยุตฺตวาโรปิ เอวํ วิตฺถาเรตพฺพา. ปจฺจยวาเร วตฺถุ จ ปญฺจายตนานิ จ ทสฺเสตพฺพานิ. ยถา ยถา ลพฺภติ. ตํ ตํ กาตพฺพํ.)}

\csromanmarkh{13. เกนจิ\-วิญฺเญยฺยทุก}\csromanmarkhii{13. เกนจิ\-วิญฺเญยฺยทุก 7. ปญฺหาวาร}\csromanheadernomark{2}{13. เกนจิ\-วิญฺเญยฺยทุก 13. เกนจิ\-วิญฺเญยฺยทุก 7. ปญฺหาวาร}
\proseitem{181}{\csromanfolio{153}เกนจิ วิญฺเญยฺโย ธมฺโม เกนจิ วิญฺเญยฺยสฺส ธมฺมสฺส เหตุปจฺจเยน ปจฺจโย. เกนจิ วิญฺเญยฺยา เหตู สมฺปยุตฺตกานํ ขนฺธานํ จิตฺตสมุฏฺฐานานญฺจ รูปานํ เหตุปจฺจเยน ปจฺจโย. ปฏิสนฺธิ\-กฺขเณ ฯเปฯ. (สํขิตฺตํ.)}

\prose{เหตุยา นว, อารมฺมเณ นว ฯเปฯ อวิคเต นว. อนุโลมํ.}

\prose{นเหตุยา นว ฯเปฯ โนวิคเต นว. (เอวํ จตฺตาริปิ คณนา ปริปุณฺณา.)}

\vspace{\tipitakaparskip}
\csromancenter{ปจฺจนียํ.}
\nitthitam{เกนจิ\-วิญฺเญยฺยทุกํ นิฏฺฐิตํ.}
\nitthitam{จูฬนฺตรทุกํ นิฏฺฐิตํ.}
\csromanmatikaanchor{mtk.17}
\csromantocmarknopage{chapter}{3. อาสวโคจฺฉก}{mtk.17}
\bookmark[dest={mtk.17},level=0]{3. อาสวโคจฺฉก}
\chapterheadpageread{3. \textbf{อาสวโคจฺฉก}}
\csromanmatikaanchor{mtk.18}
\csromantocmarkref{section}{14. อาสวทุก}{mtk.18}
\bookmark[dest={mtk.18},level=1]{14. อาสวทุก}
\csromanmarkh{14. อาสวทุก}\csromanmarkhii{1. ปฏิจฺจวาร}\csromanheadernomark{1}{14. \textbf{อาสวทุก 1. ปฏิจฺจวาร}}
\csromanmarkh{1. ปจฺจยานุโลม}\csromanmarkhii{1. วิภงฺควาร}\csromanheadernomark{2}{1. \textbf{ปจฺจยานุโลม 1. วิภงฺควาร}}
\csromanheader{2}{\textbf{เหตุ}}
\proseitem{1}{\csromanfolio{154}อาสวํ ธมฺมํ ปฏิจฺจ อาสโว ธมฺโม อุปฺปชฺชติ เหตุปจฺจยา. กามาสวํ ปฏิจฺจ ทิฏฺฐาสโว อวิชฺชาสโว. ทิฏฺฐาสวํ ปฏิจฺจ กามาสโว อวิชฺชาสโว. อวิชฺชาสวํ ปฏิจฺจ กามาสโว ทิฏฺฐาสโว. ภวาสวํ ปฏิจฺจ อวิชฺชาสโว. ทิฏฺฐาสวํ ปฏิจฺจ อวิชฺชาสโว. (เอเกกมฺปิ จกฺกํ กาตพฺพํ.) (1)}

\prose{อาสวํ ธมฺมํ ปฏิจฺจ โนอาสโว ธมฺโม อุปฺปชฺชติ เหตุปจฺจยา. อาสวํ ปฏิจฺจ อาสว\-สมฺปยุตฺตกา ขนฺธา จิตฺต\-สมุฏฺฐานญฺจ รูปํ. (2)}

\prose{อาสวํ ธมฺมํ ปฏิจฺจ อาสโว จ โนอาสโว จ ธมฺมา อุปฺปชฺชนฺติ เหตุปจฺจยา. กามาสวํ ปฏิจฺจ ทิฏฺฐาสโว อวิชฺชาสโว สมฺปยุตฺตกา จ ขนฺธา จิตฺต\-สมุฏฺฐานญฺจ รูปํ. (จกฺกํ.) (3)}

\proseitem{2}{โนอาสวํ ธมฺมํ ปฏิจฺจ โนอาสโว ธมฺโม อุปฺปชฺชติ เหตุปจฺจยา. โนอาสวํ เอกํ ขนฺธํ ปฏิจฺจ ตโย ขนฺธา จิตฺต\-สมุฏฺฐานํ รูปํ ฯเปฯ เทฺว ขนฺเธ ฯเปฯ. ปฏิสนฺธิ\-กฺขเณ ฯเปฯ. ขนฺเธ ปฏิจฺจ วตฺถุ, วตฺถุํ ปฏิจฺจ ขนฺธา. เอกํ มหาภูตํ ฯเปฯ มหาภูเต ปฏิจฺจ จิตฺต\-สมุฏฺฐานํ รูปํ, กฏตฺตารูปํ, อุปาทารูปํ. (1)}

\prose{โนอาสวํ ธมฺมํ ปฏิจฺจ อาสโว ธมฺโม อุปฺปชฺชติ เหตุปจฺจยา. โนอาสเว ขนฺเธ ปฏิจฺจ อาสวา. (2)}

\prose{โนอาสวํ ธมฺมํ ปฏิจฺจ อาสโว จ โนอาสโว จ ธมฺมา อุปฺปชฺชนฺติ เหตุปจฺจยา. โนอาสวํ เอกํ ขนฺธํ ปฏิจฺจ ตโย ขนฺธา อาสวา จ จิตฺต\-สมุฏฺฐานญฺจ รูปํ ฯเปฯ เทฺว ขนฺเธ ฯเปฯ. (3)}

\proseitem{3}{อาสวญฺจ โนอาสวญฺจ ธมฺมํ ปฏิจฺจ อาสโว ธมฺโม อุปฺปชฺชติ เหตุปจฺจยา. กามาสวญฺจ สมฺปยุตฺตเก จ ขนฺเธ ปฏิจฺจ ทิฏฺฐาสโว อวิชฺชาสโว. (จกฺกํ พนฺธิตพฺพํ.) (1)}

\prose{\csromanfolio{155}อาสวญฺจ โนอาสวญฺจ ธมฺมํ ปฏิจฺจ โนอาสโว ธมฺโม อุปฺปชฺชติ เหตุปจฺจยา. โนอาสวํ เอกํ ขนฺธญฺจ อาสเว จ ปฏิจฺจ ตโย ขนฺธา จิตฺต\-สมุฏฺฐานญฺจ รูปํ ฯเปฯ เทฺว ขนฺเธ ฯเปฯ. (2)}

\prose{อาสวญฺจ โนอาสวญฺจ ธมฺมํ ปฏิจฺจ อาสโว จ โนอาสโว จ ธมฺมา อุปฺปชฺชนฺติ เหตุปจฺจยา. โนอาสวํ เอกํ ขนฺธญฺจ กามาสวญฺจ ปฏิจฺจ ตโย ขนฺธา ทิฏฺฐาสโว อวิชฺชาสโว จิตฺต\-สมุฏฺฐานญฺจ รูปํ ฯเปฯ เทฺว ขนฺเธ จ ฯเปฯ. (จกฺกํ. สํขิตฺตํ.) (3)}

\csromanmarkh{1. ปจฺจยานุโลม}\csromanmarkhii{2. สงฺขฺยาวาร}\csromanheadernomark{2}{1. \textbf{ปจฺจยานุโลม 2. สงฺขฺยาวาร}}
\proseitem{4}{เหตุยา นว, อารมฺมเณ นว, (สพฺพตฺถ นว.) วิปาเก เอกํ อาหาเร นว ฯเปฯ อวิคเต นว.}

\vspace{\tipitakaparskip}
\csromancenter{อนุโลมํ.}
\csromanmarkh{2. ปจฺจย\-ปจฺจนีย}\csromanmarkhii{1. วิภงฺควาร}\csromanheadernomark{2}{2. ปจฺจย\-ปจฺจนีย 1. วิภงฺควาร}
\proseitem{5}{โนอาสวํ ธมฺมํ ปฏิจฺจ โนอาสโว ธมฺโม อุปฺปชฺชติ นเหตุปจฺจยา. อเหตุกํ โนอาสวํ เอกํ ขนฺธํ ปฏิจฺจ ตโย ขนฺธา จิตฺต\-สมุฏฺฐานญฺจ รูปํ ฯเปฯ เทฺว ขนฺเธ ฯเปฯ. อเหตุก\-ปฏิสนฺธิ\-กฺขเณ ฯเปฯ. ขนฺเธ ปฏิจฺจ วตฺถุ, วตฺถุํ ปฏิจฺจ ขนฺธา. เอกํ มหาภูตํ ฯเปฯ. (ยาว อสญฺญสตฺตา.) (1)}

\prose{โนอาสวํ ธมฺมํ ปฏิจฺจ อาสโว ธมฺโม อุปฺปชฺชติ นเหตุปจฺจยา. วิจิกิจฺฉา\-สหคเต อุทฺธจฺจ\-สหคเต ขนฺเธ ปฏิจฺจ วิจิกิจฺฉา\-สหคโต อุทฺธจฺจ\-สหคโต โมโห. (2)}

\prose{นอารมฺมณ}

\proseitem{6}{อาสวํ ธมฺมํ ปฏิจฺจ โนอาสโว ธมฺโม อุปฺปชฺชติ นอารมฺมณ\-ปจฺจยา. อาสเว ปฏิจฺจ จิตฺต\-สมุฏฺฐานํ รูปํ. (1)}

\prose{\csromanfolio{156}โนอาสวํ ธมฺมํ ปฏิจฺจ โนอาสโว ธมฺโม อุปฺปชฺชติ นอารมฺมณ\-ปจฺจยา. โนอาสเว ขนฺเธ ปฏิจฺจ จิตฺต\-สมุฏฺฐานํ รูปํ. ปฏิสนฺธิ\-กฺขเณ ฯเปฯ. ขนฺเธ ปฏิจฺจ วตฺถุ. เอกํ มหาภูตํ ฯเปฯ. (ยาว อสญฺญสตฺตา.) (1)}

\prose{อาสวญฺจ โนอาสวญฺจ ธมฺมํ ปฏิจฺจ โนอาสโว ธมฺโม อุปฺปชฺชติ นอารมฺมณ\-ปจฺจยา. อาสเว จ สมฺปยุตฺตเก จ ขนฺเธ ปฏิจฺจ จิตฺต\-สมุฏฺฐานํ รูปํ. (สํขิตฺตํ.) (1)}

\csromanmarkh{2. ปจฺจย\-ปจฺจนีย}\csromanmarkhii{2. สงฺขฺยาวาร}\csromanheadernomark{2}{2. ปจฺจย\-ปจฺจนีย 2. สงฺขฺยาวาร}
\proseitem{7}{นเหตุยา เทฺว, นอารมฺมเณ ตีณิ, นอธิปติยา นว, นอนนฺตเร ตีณิ, นสมนนฺตเร ตีณิ, นอญฺญมญฺเญ ตีณิ, นอุปนิสฺสเย ตีณิ, นปุเรชาเต นว, นปจฺฉาชาเต นว, นอาเสวเน นว, นกมฺเม ตีณิ, นวิปาเก นว, นอาหาเร เอกํ, นอินฺทฺริเย เอกํ, นฌาเน เอกํ, นมคฺเค เอกํ, นสมฺปยุตฺเต ตีณิ, นวิปฺปยุตฺเต นว, โนนตฺถิยา ตีณิ, โนวิคเต ตีณิ.}

\vspace{\tipitakaparskip}
\csromancenter{ปจฺจนียํ.}
\proseitem{8}{เหตุปจฺจยา นอารมฺมเณ ตีณิ, นอธิปติยา นว, นอนนฺตเร ตีณิ, นสมนนฺตเร ตีณิ, นอญฺญมญฺเญ ตีณิ, นอุปนิสฺสเย ตีณิ, นปุเรชาเต นว, นปจฺฉาชาเต นว, นอาเสวเน นว, นกมฺเม ตีณิ, นวิปาเก นว, นสมฺปยุตฺเต ตีณิ, นวิปฺปยุตฺเต นว, โนนตฺถิยา ตีณิ, โนวิคเต ตีณิ.}

\vspace{\tipitakaparskip}
\csromancenter{อนุโลม\-ปจฺจนียํ.}
\csromanmarkh{14. อาสวทุก}\csromanmarkhii{4. ปจฺจย\-ปจฺจนียา\-นุโลม}\csromanheadernomark{2}{14. อาสวทุก 4. ปจฺจย\-ปจฺจนียา\-นุโลม}
\proseitem{9}{\csromanfolio{157}นเหตุปจฺจยา อารมฺมเณ เทฺว, อนนฺตเร เทฺว ฯเปฯ วิปาเก เอกํ ฯเปฯ มคฺเค เอกํ ฯเปฯ อวิคเต เทฺว.}

\vspace{\tipitakaparskip}
\csromancenter{ปจฺจนียา\-นุโลมํ.}
\csromancenter{(สหชาตวาโร ปฏิจฺจ\-วาร\-สทิโส.)}
\csromanmarkh{14. อาสวทุก}\csromanmarkhii{3. ปจฺจยวาร}\csromanheadernomark{2}{14. อาสวทุก 3. ปจฺจยวาร}
\csromanmarkh{1. ปจฺจยานุโลม}\csromanmarkhii{1. วิภงฺควาร}\csromanheadernomark{2}{1. \textbf{ปจฺจยานุโลม 1. วิภงฺควาร}}
\csromanheader{2}{\textbf{เหตุ}}
\proseitem{10}{อาสวํ ธมฺมํ ปจฺจยา อาสโว ธมฺโม อุปฺปชฺชติ เหตุปจฺจยา. (อาสวมูลกํ ตีณิ. ปฏิจฺจสทิสา.)}

\prose{โนอาสวํ ธมฺมํ ปจฺจยา โนอาสโว ธมฺโม อุปฺปชฺชติ เหตุปจฺจยา. โนอาสวํ เอกํ ขนฺธํ ปจฺจยา ตโย ขนฺธา จิตฺต\-สมุฏฺฐานญฺจ รูปํ ฯเปฯ เทฺว ขนฺเธ ฯเปฯ. ปฏิสนฺธิ\-กฺขเณ ฯเปฯ. ขนฺเธ ปจฺจยา วตฺถุ, วตฺถุํ ปจฺจยา ขนฺธา. มหาภูเต ปจฺจยา จิตฺต\-สมุฏฺฐานํ รูปํ, กฏตฺตารูปํ, อุปาทารูปํ. วตฺถุํ ปจฺจยา โนอาสวา ขนฺธา. (1)}

\prose{โนอาสวํ ธมฺมํ ปจฺจยา อาสโว ธมฺโม อุปฺปชฺชติ เหตุปจฺจยา. โนอาสเว ขนฺเธ ปจฺจยา อาสวา. วตฺถุํ ปจฺจยา อาสวา. (2)}

\prose{โนอาสวํ ธมฺมํ ปจฺจยา อาสโว จ โนอาสโว จ ธมฺมา อุปฺปชฺชนฺติ เหตุปจฺจยา. โนอาสวํ เอกํ ขนฺธํ ปจฺจยา ตโย ขนฺธา อาสวา จ จิตฺต\-สมุฏฺฐานญฺจ รูปํ ฯเปฯ เทฺว ขนฺเธ ฯเปฯ. วตฺถุํ ปจฺจยา อาสวา สมฺปยุตฺตกา จ ขนฺธา. (3)}

\proseitem{11}{อาสวญฺจ โนอาสวญฺจ ธมฺมํ ปจฺจยา อาสโว ธมฺโม อุปฺปชฺชติ เหตุปจฺจยา. กามาสวญฺจ สมฺปยุตฺตเก จ ขนฺเธ ปจฺจยา ทิฏฺฐาสโว \csromanfolio{158}\textbf{อวิชฺชาสโว. (จกฺกํ.) กามาสวญฺจ วตฺถุญฺจ ปจฺจยา ทิฏฺฐาสโว} อวิชฺชาสโว. (จกฺกํ.) (1)}

\prose{อาสวญฺจ โนอาสวญฺจ ธมฺมํ ปจฺจยา โนอาสโว ธมฺโม อุปฺปชฺชติ เหตุปจฺจยา. โนอาสวํ เอกํ ขนฺธญฺจ อาสเว จ ปจฺจยา ตโย ขนฺธา จิตฺต\-สมุฏฺฐานญฺจ รูปํ ฯเปฯ เทฺว ขนฺเธ ฯเปฯ. อาสวญฺจ วตฺถุญฺจ ปจฺจยา โนอาสวา ขนฺธา. (2)}

\prose{อาสวญฺจ โนอาสวญฺจ ธมฺมํ ปจฺจยา อาสโว จ โนอาสโว จ ธมฺมา อุปฺปชฺชนฺติ เหตุปจฺจยา. โนอาสวํ เอกํ ขนฺธญฺจ กามาสวญฺจ ปจฺจยา ตโย ขนฺธา ทิฏฺฐาสโว อวิชฺชาสโว จิตฺต\-สมุฏฺฐานญฺจ รูปํ ฯเปฯ เทฺว ขนฺเธ ฯเปฯ. (จกฺกํ.) กามาสวญฺจ วตฺถุญฺจ ปจฺจยา ทิฏฺฐาสโว อวิชฺชาสโว สมฺปยุตฺตกา จ ขนฺธา. (จกฺกํ. สํขิตฺตํ.) (3)}

\csromanmarkh{1. ปจฺจยานุโลม}\csromanmarkhii{2. สงฺขฺยาวาร}\csromanheadernomark{2}{1. \textbf{ปจฺจยานุโลม 2. สงฺขฺยาวาร}}
\vspace{\tipitakaparskip}
\csromancenter{เหตุยา นว, อารมฺมเณ นว, อธิปติยา นว ฯเปฯ วิปาเก เอกํ ฯเปฯ อวิคเต นว.}
\csromancenter{อนุโลมํ.}
\csromanmarkh{2. ปจฺจย\-ปจฺจนีย}\csromanmarkhii{1. วิภงฺควาร}\csromanheadernomark{2}{2. ปจฺจย\-ปจฺจนีย 1. วิภงฺควาร}
\proseitem{13}{โนอาสวํ ธมฺมํ ปจฺจยา โนอาสโว ธมฺโม อุปฺปชฺชติ นเหตุปจฺจยา. อเหตุกํ โนอาสวํ เอกํ ขนฺธํ ปจฺจยา ตโย ขนฺธา จิตฺต\-สมุฏฺฐานญฺจ รูปํ ฯเปฯ เทฺว ขนฺเธ ฯเปฯ. อเหตุก\-ปฏิสนฺธิ\-กฺขเณ ฯเปฯ. (ยาว อสญฺญสตฺตา.) จกฺขายตนํ ปจฺจยา จกฺขุวิญฺญาณํ ฯเปฯ. กายายตนํ ปจฺจยา กายวิญฺญาณํ. วตฺถุํ ปจฺจยา อเหตุกา โนอาสวา ขนฺธา. (1)}

\prose{โนอาสวํ ธมฺมํ ปจฺจยา อาสโว ธมฺโม อุปฺปชฺชติ นเหตุปจฺจยา. วิจิกิจฺฉา\-สหคเต อุทฺธจฺจ\-สหคเต ขนฺเธ จ วตฺถุญฺจ ปจฺจยา วิจิกิจฺฉา\-สหคโต อุทฺธสหคโต โมโห. (2)}

\csromanmarkh{14. อาสวทุก}\csromanmarkhii{2. ปจฺจย\-ปจฺจนีย 2. สงฺขฺยาวร}\csromanheadernomark{2}{14. อาสวทุก 2. ปจฺจย\-ปจฺจนีย 2. สงฺขฺยาวร}
\proseitem{14}{\csromanfolio{159}นเหตุยา เทฺว, นอารมฺมเณ ตีณิ, นอธิปติยา นว ฯเปฯ นกมฺเม ตีณิ ฯเปฯ นวิปฺปยุตฺเต นว, โนนตฺถิยา ตีณิ, โนวิคเต ตีณิ.(เอวํ สพฺเพ คณนา คเณตพฺพา.)}

\vspace{\tipitakaparskip}
\csromancenter{(นิสฺสยวาโร ปจฺจยวาร\-สทิโส.)}
\csromanmarkh{14. อาสวทุก}\csromanmarkhii{5. สํสฏฺฐวาร}\csromanheadernomark{2}{14. อาสวทุก 5. สํสฏฺฐวาร}
\proseitem{15}{อาสวํ ธมฺมํ สํสฏฺโฐ อาสโว ธมฺโม อุปฺปชฺชติ เหตุปจฺจยา. กามาสวํ สํสฏฺโฐ ทิฏฺฐาสโว อวิชฺชาสโว. (จกฺกํ. สํขิตฺตํ.)}

\prose{เหตุยา นว, อารมฺมเณ นว, (สพฺพตฺถ นว.) วิปาเก เอกํ ฯเปฯ อวิคเต นว.}

\prose{นเหตุยา เทฺว, นอธิปติยา นว, นปุเรชาเต นว, นปจฺฉาชาเต นว, นอาเสวเน นว, นกมฺเม ตีณิ, นวิปาเก นว, นฌาเน เอกํ, นมคฺเค เอกํ, นวิปฺปยุตฺเต นว.}

\vspace{\tipitakaparskip}
\csromancenter{(คณนาปิ สมฺปยุตฺตวาโรปิ สํสฏฺฐวาร\-สทิโส.)}
\csromanmarkh{14. อาสวทุก}\csromanmarkhii{7. ปญฺหาวาร}\csromanheadernomark{2}{14. อาสวทุก 7. ปญฺหาวาร}
\csromanmarkh{1. ปจฺจยานุโลม}\csromanmarkhii{1. วิภงฺควาร}\csromanheadernomark{2}{1. \textbf{ปจฺจยานุโลม 1. วิภงฺควาร}}
\vspace{\tipitakaparskip}
\csromancenter{\textbf{เหตุ}}
\proseitem{16}{อาสโว ธมฺโม อาสวสฺส ธมฺมสฺส เหตุปจฺจเยน ปจฺจโย. กามาสโว ทิฏฺฐาสวสฺส อวิชฺชาสวสฺส เหตุปจฺจเยน ปจฺจโย. ภวาสโว อวิชฺชาสวสฺส เหตุปจฺจเยน ปจฺจโย. (จกฺกํ.) (1)}

\prose{อาสโว ธมฺโม โนอาสวสฺส ธมฺมสฺส เหตุปจฺจเยน ปจฺจโย. อาสวา เหตู สมฺปยุตฺตกานํ ขนฺธานํ จิตฺตสมุฏฺฐานานญฺจ รูปานํ เหตุปจฺจเยน ปจฺจโย. (2)}

\prose{\csromanfolio{160}อาสโว ธมฺโม อาสวสฺส จ โนอาสวสฺส จ ธมฺมสฺส เหตุปจฺจเยน ปจฺจโย. กามาสโว ทิฏฺฐาสวสฺส อวิชฺชาสวสฺส สมฺปยุตฺตกานญฺจ ขนฺธานํ จิตฺตสมุฏฺฐานานญฺจ รูปานํ เหตุปจฺจเยน ปจฺจโย. (3)}

\proseitem{17}{โนอาสโว ธมฺโม โนอาสวสฺส ธมฺมสฺส เหตุปจฺจเยน ปจฺจโย. โนอาสวา เหตู สมฺปยุตฺตกานํ ขนฺธานํ จิตฺตสมุฏฺฐานานญฺจ รูปานํ เหตุปจฺจเยน ปจฺจโย. ปฏิสนฺธิ\-กฺขเณ ฯเปฯ. (1)}

\prose{โนอาสโว ธมฺโม อาสวสฺส ธมฺมสฺส เหตุปจฺจเยน ปจฺจโย. โนอาสวา เหตู สมฺปยุตฺตกานํ อาสวานํ เหตุปจฺจเยน ปจฺจโย. (2)}

\prose{โนอาสโว ธมฺโม อาสวสฺส จ โนอาสวสฺส จ ธมฺมสฺส เหตุปจฺจเยน ปจฺจโย. โนอาสวา เหตู สมฺปยุตฺตกานํ ขนฺธานํ อาสวานํ จิตฺตสมุฏฺฐานานญฺจ รูปานํ เหตุปจฺจเยน ปจฺจโย. (3)}

\prose{อาสโว จ โนอาสโว จ ธมฺมา โนอาสวสฺส ธมฺมสฺส เหตุปจฺจเยน ปจฺจโย. อาสวา จ โนอาสวา จ เหตู สมฺปยุตฺตกานํ ขนฺธานํ จิตฺตสมุฏฺฐานานญฺจ รูปานํ เหตุปจฺจเยน ปจฺจโย. (1)}

\proseitem{18}{อาสโว ธมฺโม อาสวสฺส ธมฺมสฺส อารมฺมณ\-ปจฺจเยน ปจฺจโย. อาสเว อารพฺภ อาสวา อุปฺปชฺชนฺติ. (1)}

\prose{อาสโว ธมฺโม โนอาสวสฺส ธมฺมสฺส อารมฺมณ\-ปจฺจเยน ปจฺจโย. อาสเว อารพฺภ โนอาสวา ขนฺธา อุปฺปชฺชนฺติ. (2)}

\prose{อาสโว ธมฺโม อาสวสฺส จ โนอาสวสฺส จ ธมฺมสฺส อารมฺมณ\-ปจฺจเยน ปจฺจโย. อาสเว อารพฺภ อาสวา จ สมฺปยุตฺตกา จ ขนฺธา อุปฺปชฺชนฺติ. (3)}

\proseitem{19}{โนอาสโว ธมฺโม โนอาสวสฺส ธมฺมสฺส อารมฺมณ\-ปจฺจเยน ปจฺจโย. ทานํ ฯเปฯ สีลํ ฯเปฯ อุโปสถกมฺมํ ฯเปฯ. ปุพฺเพ ฯเปฯ. ฌานา ฯเปฯ. อริยา มคฺคา วุฏฺฐหิตฺวา มคฺคํ ปจฺจเวกฺขนฺติ, ผลํ ปจฺจเวกฺขนฺติ, นิพฺพานํ ปจฺจเวกฺขนฺติ. นิพฺพานํ โคตฺรภุสฺส, โวทานสฺส, มคฺคสฺส, ผลสฺส, อาวชฺชนาย อารมฺมณ\-ปจฺจเยน ปจฺจโย. อริยา โนอาสเว ปหีเน กิเลเส ฯเปฯ. \csromanfolio{161}วิกฺขมฺภิเต กิเลเส ฯเปฯ. ปุพฺเพ สมุทาจิณฺเณ กิเลเส ชานนฺติ. จกฺขุํ ฯเปฯ วตฺถุํ โนอาสเว ขนฺเธ อนิจฺจโต ฯเปฯ โทมนสฺสํ อุปฺปชฺชติ. ทิพฺเพน จกฺขุนา รูปํ ปสฺสติ. ทิพฺพาย โสตธาตุยา สทฺทํ สุณาติ. เจโตปริย\-ญาเณน โนอาสวจิตฺตสมงฺคิสฺส จิตฺตํ ชานาติ.  อากาสา\-นญฺจา\-ยตนํ วิญฺญาณญฺจา\-ยตนสฺส ฯเปฯ. อากิญฺจญฺญา\-ยตนํ เนว\-สญฺญา\-นา\-สญฺญา\-ยตนสฺส ฯเปฯ. รูปายตนํ จกฺขุ\-วิญฺญาณสฺส ฯเปฯ. โผฏฺฐพฺพา\-ยตนํ กายวิญฺญาณสฺส ฯเปฯ. โนอาสวา ขนฺธา อิทฺธิวิธ\-ญาณสฺส, เจโตปริย\-ญาณสฺส, ปุพฺเพ\-นิวาสา\-นุสฺสติ\-ญาณสฺส, ยถากมฺมูปคญาณสฺส, อนาคตํสญาณสฺส, อาวชฺชนาย อารมฺมณ\-ปจฺจเยน ปจฺจโย. (1)}

\prose{โนอาสโว ธมฺโม อาสวสฺส ธมฺมสฺส อารมฺมณ\-ปจฺจเยน ปจฺจโย. ทานํ ทตฺวา ฯเปฯ ตํ อสฺสาเทติ อภินนฺทติ, ตํ อารพฺภ อาสวา อุปฺปชฺชนฺติ. สีลํ ฯเปฯ อุโปสถกมฺมํ ฯเปฯ. ปุพฺเพ ฯเปฯ. ฌานา ฯเปฯ. จกฺขุํ ฯเปฯ วตฺถุํ โนอาสเว ขนฺเธ อสฺสาเทติ อภินนฺทติ, ตํ อารพฺภ อาสวา อุปฺปชฺชนฺติ. (2)}

\prose{โนอาสโว ธมฺโม อาสวสฺส จ โนอาสวสฺส จ ธมฺมสฺส อารมฺมณ\-ปจฺจเยน ปจฺจโย. ทานํ ทตฺวา ฯเปฯ. (ทุติยคมนํ.) โนอาสเว ขนฺเธ อสฺสาเทติ อภินนฺทติ, ตํ อารพฺภ อาสวา จ สมฺปยุตฺตกา จ ขนฺธา อุปฺปชฺชนฺติ. (3)}

\proseitem{20}{อาสโว จ โนอาสโว จ ธมฺมา อาสวสฺส ธมฺมสฺส อารมฺมณ\-ปจฺจเยน ปจฺจโย. อาสเว จ สมฺปยุตฺตเก จ ขนฺเธ อารพฺภ อาสวา อุปฺปชฺชนฺติ. (1)}

\prose{อาสโว จ โนอาสโว จ ธมฺมา โนอาสวสฺส ธมฺมสฺส อารมฺมณ\-ปจฺจเยน ปจฺจโย. อาสเว จ สมฺปยุตฺตเก จ ขนฺเธ อารพฺภ โนอาสวา ขนฺธา อุปฺปชฺชนฺติ. (2)}

\prose{อาสโว จ โนอาสโว จ ธมฺมา อาสวสฺส จ โนอาสวสฺส จ ธมฺมสฺส อารมฺมณ\-ปจฺจเยน ปจฺจโย. อาสเว จ สมฺปยุตฺตเก จ ขนฺเธ อารพฺภ อาสวา จ สมฺปยุตฺตกา จ ขนฺธา อุปฺปชฺชนฺติ. (3)}

\proseitem{21}{\csromanfolio{162}อาสโว ธมฺโม อาสวสฺส ธมฺมสฺส อธิปติ\-ปจฺจเยน ปจฺจโย. อารมฺมณาธิปตฺอิ\csromandash{}อาสเว ครุํ กตฺวา อาสวา อุปฺปชฺชนฺติ. (ตีณิ อารมฺมณสทิสา, ครุการมฺมณา กาตพฺพา.)}

\vspace{\tipitakaparskip}
\csromancenter{(ตีณิ อารมฺมณสทิสา, ครุการมฺมณา กาตพฺพา.)}
\prose{โนอาสโว ธมฺโม โนอาสวสฺส ธมฺมสฺส อธิปติ\-ปจฺจเยน ปจฺจโย. อารมฺมณา\-ธิปติ, สหชาตา\-ธิปติ.  อารมฺมณา\-ธิปติ\csromandash{}ทานํ ฯเปฯ สีลํ ฯเปฯ อุโปสถกมฺมํ ฯเปฯ. ปุพฺเพ ฯเปฯ. ฌานา ฯเปฯ. อริยา มคฺคา ฯเปฯ ผลํ ฯเปฯ. นิพฺพานํ ครุํ ฯเปฯ. นิพฺพานํ โคตฺรภุสฺส, โวทานสฺส, มคฺคสฺส, ผลสฺส อธิปติ\-ปจฺจเยน ปจฺจโย. จกฺขุํ ฯเปฯ วตฺถุํ โนอาสเว ขนฺเธ ครุํ กตฺวา อสฺสาเทติ อภินนฺทติ, ตํ ครุํ กตฺวา ราโค อุปฺปชฺชติ, ทิฏฺฐิ อุปฺปชฺชติ. สหชาตา\-ธิปติ\csromandash{}โนอาสวา อธิปติ สมฺปยุตฺตกานํ ขนฺธานํ จิตฺตสมุฏฺฐานานญฺจ รูปานํ อธิปติ\-ปจฺจเยน ปจฺจโย. (1)}

\prose{โนอาสโว ธมฺโม อาสวสฺส ธมฺมสฺส อธิปติ\-ปจฺจเยน ปจฺจโย. อารมฺมณา\-ธิปติ, สหชาตา\-ธิปติ.  อารมฺมณา\-ธิปติ\csromandash{}ทานํ ฯเปฯ โนอาสเว ขนฺเธ ครุํ กตฺวา อสฺสาเทติ อภินนฺทติ, ตํ ครุํ กตฺวา อาสวา อุปฺปชฺชนฺติ. สหชาตา\-ธิปติ\csromandash{}โนอาสวา อธิปติ สมฺปยุตฺตกานํ อาสวานํ อธิปติ\-ปจฺจเยน ปจฺจโย. (2)}

\prose{โนอาสโว ธมฺโม อาสวสฺส จ โนอาสวสฺส จ ธมฺมสฺส อธิปติ\-ปจฺจเยน ปจฺจโย. อารมฺมณา\-ธิปติ, สหชาตา\-ธิปติ.  อารมฺมณา\-ธิปติ\csromandash{}ทานํ ฯเปฯ โนอาสเว ขนฺเธ ครุํ กตฺวา อสฺสาเทติ อภินนฺทติ, ตํ ครุํ กตฺวา ราโค อุปฺปชฺชติ, ทิฏฺฐิ อุปฺปชฺชติ. สหชาตา\-ธิปติ\csromandash{}โนอาสวา อธิปติ สมฺปยุตฺตกานํ ขนฺธานํ อาสวานํ จิตฺตสมุฏฺฐานานญฺจ รูปานํ อธิปติ\-ปจฺจเยน ปจฺจโย. (3)}

\prose{อาสโว จ โนอาสโว จ ธมฺมา อาสวสฺส ธมฺมสฺส อธิปติ\-ปจฺจเยน ปจฺจโย. อารมฺมณา\-ธิปติ\csromandash{}อาสเว จ สมฺปยุตฺตเก จ ขนฺเธ ครุํ กตฺวา อสฺสาเทติ ฯเปฯ. อาสวา อุปฺปชฺชนฺติ. (ตีณิ. ครุการมฺมณา.)}

\csromanheader{2}{14. อาสวทุก \textbf{อนนฺตร}}
\proseitem{22}{\csromanfolio{163}อาสโว ธมฺโม อาสวสฺส ธมฺมสฺส อนนฺตร\-ปจฺจเยน ปจฺจโย. ปุริมา ปุริมา อาสวา ปจฺฉิมานํ ปจฺฉิมานํ อาสวานํ อนนฺตร\-ปจฺจเยน ปจฺจโย. (1)}

\prose{อาสโว ธมฺโม โนอาสวสฺส ธมฺมสฺส อนนฺตร\-ปจฺจเยน ปจฺจโย. ปุริมา ปุริมา อาสวา ปจฺฉิมานํ ปจฺฉิมานํ โนอาสวานํ ขนฺธานํ อนนฺตร\-ปจฺจเยน ปจฺจโย. อาสวา วุฏฺฐานสฺส อนนฺตร\-ปจฺจเยน ปจฺจโย. (2)}

\prose{อาสโว ธมฺโม อาสวสฺส จ โนอาสวสฺส จ ธมฺมสฺส อนนฺตร\-ปจฺจเยน ปจฺจโย. ปุริมา ปุริมา อาสวา ปจฺฉิมานํ ปจฺฉิมานํ อาสวานํ สมฺปยุตฺตกานญฺจ ขนฺธานํ อนนฺตร\-ปจฺจเยน ปจฺจโย. (3)}

\proseitem{23}{โนอาสโว ธมฺโม โนอาสวสฺส ธมฺมสฺส อนนฺตร\-ปจฺจเยน ปจฺจโย. ปุริมา ปุริมา โนอาสวา ขนฺธา ปจฺฉิมานํ ปจฺฉิมานํ โนอาสวานํ ขนฺธานํ อนนฺตร\-ปจฺจเยน ปจฺจโย. อนุโลมํ โคตฺรภุสฺส ฯเปฯ ผลสมาปตฺติยา อนนฺตร\-ปจฺจเยน ปจฺจโย. (1)}

\prose{โนอาสโว ธมฺโม อาสวสฺส ธมฺมสฺส อนนฺตร\-ปจฺจเยน ปจฺจโย. ปุริมา ปุริมา โนอาสวา ขนฺธา ปจฺฉิมานํ ปจฺฉิมานํ อาสวานํ อนนฺตร\-ปจฺจเยน ปจฺจโย. อาวชฺชนา อาสวานํ อนนฺตร\-ปจฺจเยน ปจฺจโย. (2)}

\prose{โนอาสโว ธมฺโม อาสวสฺส จ โนอาสวสฺส จ ธมฺมสฺส อนนฺตร\-ปจฺจเยน ปจฺจโย. ปุริมา ปุริมา โนอาสวา ขนฺธา ปจฺฉิมานํ ปจฺฉิมานํ อาสวานํ สมฺปยุตฺตกานญฺจ ขนฺธานํ อนนฺตร\-ปจฺจเยน ปจฺจโย. (3)}

\proseitem{24}{อาสโว จ โนอาสโว จ ธมฺมา อาสวสฺส ธมฺมสฺส อนนฺตร\-ปจฺจเยน ปจฺจโย. ตีณิ.}

\csromanheader{2}{\textbf{สมนนฺตราทิ}}
\proseitem{25}{อาสโว ธมฺโม อาสวสฺส ธมฺมสฺส สมนนฺตร\-ปจฺจเยน ปจฺจโย. สหชาต\-ปจฺจเยน ปจฺจโย. นว. อญฺญมญฺญ\-ปจฺจเยน ปจฺจโย. นว. นิสฺสย\-ปจฺจเยน ปจฺจโย. นว. (วตฺถุ จ ทสฺเสตพฺพํ.)}

\csromanheader{2}{\textbf{อุปนิสฺสย}}
\proseitem{26}{\csromanfolio{164}อาสโว ธมฺโม อาสวสฺส ธมฺมสฺส อุปนิสฺสย\-ปจฺจเยน ปจฺจโย. อารมฺมณู\-ปนิสฺสโย, อนนฺตรู\-ปนิสฺสโย, ปกตู\-ปนิสฺสโย ฯเปฯ. ปกตู\-ปนิสฺสโย\csromandash{}อาสวา อาสวานํ อุปนิสฺสย\-ปจฺจเยน ปจฺจโย. (ตีณิ.)}

\proseitem{27}{โนอาสโว ธมฺโม โนอาสวสฺส ธมฺมสฺส อุปนิสฺสย\-ปจฺจเยน ปจฺจโย. อารมฺมณู\-ปนิสฺสโย, อนนฺตรู\-ปนิสฺสโย, ปกตู\-ปนิสฺสโย ฯเปฯ. ปกตู\-ปนิสฺสโย\csromandash{}สทฺธํ อุปนิสฺสาย ทานํ เทติ ฯเปฯ สมาปตฺติํ อุปฺปาเทติ. มานํ ชปฺเปติ. ทิฏฺฐิํ คณฺหาติ. สีลํ ฯเปฯ. เสนาสนํ อุปนิสฺสาย ทานํ เทติ ฯเปฯ สํฆํ ภินฺทติ. สทฺธา ฯเปฯ. เสนาสนํ สทฺธาย ฯเปฯ ผลสมาปตฺติยา อุปนิสฺสย\-ปจฺจเยน ปจฺจโย. (1)}

\prose{โนอาสโว ธมฺโม อาสวสฺส ธมฺมสฺส อุปนิสฺสย\-ปจฺจเยน ปจฺจโย. อารมฺมณู\-ปนิสฺสโย, อนนฺตรู\-ปนิสฺสโย, ปกตู\-ปนิสฺสโย ฯเปฯ. ปกตู\-ปนิสฺสโย\csromandash{}สทฺธํ อุปนิสฺสาย มานํ ชปฺเปติ. ทิฏฺฐิํ คณฺหาติ. สีลํ ฯเปฯ สํฆํ ภินฺทติ. สทฺธา ฯเปฯ. เสนาสนํ ราคสฺส ฯเปฯ ปตฺถนาย อุปนิสฺสย\-ปจฺจเยน ปจฺจโย. (2)}

\prose{โนอาสโว ธมฺโม อาสวสฺส จ โนอาสวสฺส จ ธมฺมสฺส อุปนิสฺสย\-ปจฺจเยน ปจฺจโย. อารมฺมณู\-ปนิสฺสโย, อนนฺตรู\-ปนิสฺสโย, ปกตู\-ปนิสฺสโย ฯเปฯ. ปกตู\-ปนิสฺสโย\csromandash{}สทฺธํ อุปนิสฺสาย มานํ ชปฺเปติ. ทิฏฺฐิํ คณฺหาติ. สีลํ ฯเปฯ. เสนาสนํ อุปนิสฺสาย ฯเปฯ ปาณํ หนติ ฯเปฯ. สํฆํ ภินฺทติ. สทฺธา ฯเปฯ. เสนาสนํ ราคสฺส ฯเปฯ ปตฺถนาย อุปนิสฺสย\-ปจฺจเยน ปจฺจโย. (3)}

\prose{อาสโว จ โนอาสโว จ ธมฺมา อาสวสฺส ธมฺมสฺส อุปนิสฺสย\-ปจฺจเยน ปจฺจโย. อารมฺมณู\-ปนิสฺสโย, อนนฺตรู\-ปนิสฺสโย, ปกตู\-ปนิสฺสโย ฯเปฯ. ปกตู\-ปนิสฺสโย\csromandash{}ตีณิ.}

\proseitem{28}{โนอาสโว ธมฺโม โนอาสวสฺส ธมฺมสฺส ปุเรชาต\-ปจฺจเยน ปจฺจโย. อารมฺมณ\-ปุเรชาตํ, วตฺถุ\-ปุเรชาตํ.  อารมฺมณ\-ปุเรชาตํ\csromandash{} \csromanfolio{165}จกฺขุํ ฯเปฯ วตฺถุํ. (เอวํ วิตฺถาเรตพฺพํ.) โผฏฺฐพฺพา\-ยตนํ กายวิญฺญาณสฺส ปุเรชาต\-ปจฺจเยน ปจฺจโย. วตฺถุ\-ปุเรชาตํ\csromandash{} จกฺขายตนํ จกฺขุ\-วิญฺญาณสฺส ฯเปฯ กายายตนํ กายวิญฺญาณสฺส. วตฺถุ โนอาสวานํ ขนฺธานํ ปุเรชาต\-ปจฺจเยน ปจฺจโย. (1)}

\prose{โนอาสโว ธมฺโม อาสวสฺส ธมฺมสฺส ปุเรชาต\-ปจฺจเยน ปจฺจโย. อารมฺมณ\-ปุเรชาตํ, วตฺถุ\-ปุเรชาตํ.  อารมฺมณ\-ปุเรชาตํ\csromandash{}จกฺขุํ ฯเปฯ วตฺถุํ อสฺสาเทติ อภินนฺทติ, ตํ อารพฺภ อาสวา อุปฺปชฺชนฺติ. วตฺถุ\-ปุเรชาตํ\csromandash{}วตฺถุ อาสวานํ ปุเรชาต\-ปจฺจเยน ปจฺจโย. (2)}

\prose{โนอาสโว ธมฺโม อาสวสฺส จ โนอาสวสฺส จ ธมฺมสฺส ปุเรชาต\-ปจฺจเยน ปจฺจโย. อารมฺมณ\-ปุเรชาตํ, วตฺถุ\-ปุเรชาตํ.  อารมฺมณ\-ปุเรชาตํ\csromandash{}จกฺขุํ ฯเปฯ วตฺถุํ อสฺสาเทติ อภินนฺทติ, ตํ อารพฺภ อาสวา จ สมฺปยุตฺตกา จ ขนฺธา อุปฺปชฺชนฺติ. วตฺถุ\-ปุเรชาตํ\csromandash{} วตฺถุ อาสวานญฺจ อาสว\-สมฺปยุตฺตกานญฺจ ขนฺธานํ ปุเรชาต\-ปจฺจเยน ปจฺจโย. (3)}

\csromanheader{2}{\textbf{ปจฺฉาชาต}}
\proseitem{29}{อาสโว ธมฺโม โนอาสวสฺส ธมฺมสฺส ปจฺฉาชาต\-ปจฺจเยน ปจฺจโย. ปจฺฉาชาตา อาสวา ปุเรชาตสฺส อิมสฺส กายสฺส ปจฺฉาชาต\-ปจฺจเยน ปจฺจโย. (1)}

\prose{โนอาสโว ธมฺโม โนอาสวสฺส ธมฺมสฺส ปจฺฉาชาต\-ปจฺจเยน ปจฺจโย. ปจฺฉาชาตา โนอาสวา ขนฺธา ปุเรชาตสฺส อิมสฺส กายสฺส ปจฺฉาชาต\-ปจฺจเยน ปจฺจโย. (1)}

\prose{อาสโว จ โนอาสโว จ ธมฺมา โนอาสวสฺส ธมฺมสฺส ปจฺฉาชาต\-ปจฺจเยน ปจฺจโย. ปจฺฉาชาตา อาสวา จ สมฺปยุตฺตกา จ ขนฺธา ปุเรชาตสฺส อิมสฺส กายสฺส ปจฺฉาชาต\-ปจฺจเยน ปจฺจโย. (1)}

\csromanheader{2}{\textbf{อาเสวน}}
\vspace{\tipitakaparskip}
\csromancenter{อาสโว ธมฺโม อาสวสฺส ธมฺมสฺส อาเสวน\-ปจฺจเยน ปจฺจโย. นว.}
\proseitem{31}{\csromanfolio{166}โนอาสโว ธมฺโม โนอาสวสฺส ธมฺมสฺส กมฺมปจฺจเยน ปจฺจโย. สหชาตา, นานากฺขณิกา.  สหชาตา โนอาสวา เจตนา สมฺปยุตฺตกานํ ขนฺธานํ จิตฺตสมุฏฺฐานานญฺจ รูปานํ กมฺมปจฺจเยน ปจฺจโย. ปฏิสนฺธิ\-กฺขเณ ฯเปฯ. นานากฺขณิกา โนอาสวา เจตนา วิปากานํ ขนฺธานํ กฏตฺตา จ รูปานํ กมฺมปจฺจเยน ปจฺจโย. (1)}

\prose{โนอาสโว ธมฺโม อาสวสฺส ธมฺมสฺส กมฺมปจฺจเยน ปจฺจโย. โนอาสวา เจตนา สมฺปยุตฺตกานํ อาสวานํ กมฺมปจฺจเยน ปจฺจโย.}

\prose{โนอาสโว ธมฺโม อาสวสฺส จ โนอาสวสฺส จ ธมฺมสฺส กมฺมปจฺจเยน ปจฺจโย. โนอาสวา เจตนา สมฺปยุตฺตกานํ ขนฺธานํ อาสวานํ จิตฺตสมุฏฺฐานานญฺจ รูปานํ กมฺมปจฺจเยน ปจฺจโย. (3)}

\csromanheader{2}{\textbf{วิปากาหาร}}
\proseitem{32}{โนอาสโว ธมฺโม โนอาสวสฺส ธมฺมสฺส วิปาก\-ปจฺจเยน ปจฺจโย. เอกํ. อาหารปจฺจเยน ปจฺจโย. โนอาสวา อาหารา สมฺปยุตฺตกานํ ขนฺธานํ จิตฺตสมุฏฺฐานานญฺจ รูปานํ อาหารปจฺจเยน ปจฺจโย. ปฏิสนฺธิ\-กฺขเณ ฯเปฯ. กพฬีกาโร อาหาโร อิมสฺส กายสฺส อาหารปจฺจเยน ปจฺจโย. (1)}

\prose{โนอาสโว ธมฺโม อาสวสฺส ธมฺมสฺส อาหารปจฺจเยน ปจฺจโย. โนอาสวา อาหารา สมฺปยุตฺตกานํ อาสวานํ อาหารปจฺจเยน ปจฺจโย. (2)}

\prose{โนอาสโว ธมฺโม อาสวสฺส จ โนอาสวสฺส จ ธมฺมสฺส อาหารปจฺจเยน ปจฺจโย. โนอาสวา อาหารา สมฺปยุตฺตกานํ ขนฺธานํ อาสวานํ จิตฺตสมุฏฺฐานานญฺจ รูปานํ อาหารปจฺจเยน ปจฺจโย. (3)}

\csromanheader{2}{\textbf{อินฺทฺริยาทิ}}
\proseitem{33}{โนอาสโว ธมฺโม โนอาสวสฺส ธมฺมสฺส อินฺทฺริย\-ปจฺจเยน ปจฺจโย. โนอาสวา อินฺทฺริยา สมฺปยุตฺตกานํ ขนฺธานํ จิตฺตสมุฏฺฐานานญฺจ \csromanfolio{167}รูปานํ อินฺทฺริย\-ปจฺจเยน ปจฺจโย. ปฏิสนฺธิ\-กฺขเณ ฯเปฯ. จกฺขุนฺทฺริยํ จกฺขุ\-วิญฺญาณสฺส ฯเปฯ. กายินฺทฺริยํ กายวิญฺญาณสฺส ฯเปฯ. รูปชีวิตินฺทฺริยํ กฏตฺตารูปานํ อินฺทฺริย\-ปจฺจเยน ปจฺจโย. ตีณิ. ฌานปจฺจเยน ปจฺจโย. ตีณิ. มคฺคปจฺจเยน ปจฺจโย. นว. สมฺปยุตฺต\-ปจฺจเยน ปจฺจโย. นว.}

\proseitem{34}{อาสโว ธมฺโม โนอาสวสฺส ธมฺมสฺส วิปฺปยุตฺต\-ปจฺจเยน ปจฺจโย. สหชาตํ, ปจฺฉาชาตํ.  สหชาตา อาสวา จิตฺต\-สมุฏฺฐานานํ รูปานํ วิปฺปยุตฺต\-ปจฺจเยน ปจฺจโย. ปจฺฉาชาตา อาสวา ปุเรชาตสฺส อิมสฺส กายสฺส วิปฺปยุตฺต\-ปจฺจเยน ปจฺจโย. (1)}

\proseitem{35}{โนอาสโว ธมฺโม โนอาสวสฺส ธมฺมสฺส วิปฺปยุตฺต\-ปจฺจเยน ปจฺจโย. สหชาตํ, ปุเรชาตํ, ปจฺฉาชาตํ.  สหชาตา โนอาสวา ขนฺธา จิตฺต\-สมุฏฺฐานานํ รูปานํ วิปฺปยุตฺต\-ปจฺจเยน ปจฺจโย. ปฏิสนฺธิ\-กฺขเณ ฯเปฯ. ขนฺธา วตฺถุสฺส วิปฺปยุตฺต\-ปจฺจเยน ปจฺจโย. วตฺถุ ขนฺธานํ วิปฺปยุตฺต\-ปจฺจเยน ปจฺจโย. ปุเรชาตํ จกฺขายตนํ จกฺขุ\-วิญฺญาณสฺส ฯเปฯ กายายตนํ กายวิญฺญาณสฺส ฯเปฯ วตฺถุ โนอาสวานํ ขนฺธานํ วิปฺปยุตฺต\-ปจฺจเยน ปจฺจโย. ปจฺฉาชาตา โนอาสวา ขนฺธา ปุเรชาตสฺส อิมสฺส กายสฺส วิปฺปยุตฺต\-ปจฺจเยน ปจฺจโย. (1)}

\prose{โนอาสโว ธมฺโม อาสวสฺส ธมฺมสฺส วิปฺปยุตฺต\-ปจฺจเยน ปจฺจโย. ปุเรชาตํ วตฺถุ อาสวานํ วิปฺปยุตฺต\-ปจฺจเยน ปจฺจโย. (2)}

\prose{โนอาสโว ธมฺโม อาสวสฺส จ โนอาสวสฺส จ ธมฺมสฺส วิปฺปยุตฺต\-ปจฺจเยน ปจฺจโย. ปุเรชาตํ วตฺถุ อาสวานํ สมฺปยุตฺตกานญฺจ ขนฺธานํ วิปฺปยุตฺต\-ปจฺจเยน ปจฺจโย. (3)}

\proseitem{36}{อาสโว จ โนอาสโว จ ธมฺมา โนอาสวสฺส ธมฺมสฺส วิปฺปยุตฺต\-ปจฺจเยน ปจฺจโย. สหชาตํ, ปจฺฉาชาตํ.  สหชาตา อาสวา จ สมฺปยุตฺตกา จ ขนฺธา จิตฺต\-สมุฏฺฐานานํ รูปานํ วิปฺปยุตฺต\-ปจฺจเยน ปจฺจโย. ปจฺฉาชาตา อาสวา จ สมฺปยุตฺตกา จ ขนฺธา ปุเรชาตสฺส อิมสฺส กายสฺส วิปฺปยุตฺต\-ปจฺจเยน ปจฺจโย. (1)}

\proseitem{37}{\csromanfolio{168}อาสโว ธมฺโม อาสวสฺส ธมฺมสฺส อตฺถิปจฺจเยน ปจฺจโย. กามาสโว ทิฏฺฐาสวสฺส อวิชฺชาสวสฺส อตฺถิปจฺจเยน ปจฺจโย. (จกฺกํ.) (1)}

\prose{อาสโว ธมฺโม โนอาสวสฺส ธมฺมสฺส อตฺถิปจฺจเยน ปจฺจโย. สหชาตํ, ปจฺฉาชาตํ.  \textbf{สหชาตา} อาสวา สมฺปยุตฺตกานํ ขนฺธานํ จิตฺตสมุฏฺฐานานญฺจ รูปานํ อตฺถิปจฺจเยน ปจฺจโย. \textbf{ปจฺฉาชาตา} อาสวา ปุเรชาตสฺส อิมสฺส กายสฺส อตฺถิปจฺจเยน ปจฺจโย. (2)}

\prose{อาสโว ธมฺโม อาสวสฺส จ โนอาสวสฺส จ ธมฺมสฺส อตฺถิปจฺจเยน ปจฺจโย. กามาสโว ทิฏฺฐาสวสฺส อวิชฺชาสวสฺส สมฺปยุตฺตกานํ ขนฺธานํ จิตฺตสมุฏฺฐานานญฺจ รูปานํ อตฺถิปจฺจเยน ปจฺจโย. (จกฺกํ.) (3)}

\proseitem{38}{โนอาสโว ธมฺโม โนอาสวสฺส ธมฺมสฺส อตฺถิปจฺจเยน ปจฺจโย. สหชาตํ, ปุเรชาตํ, ปจฺฉาชาตํ, อาหารํ, อินฺทฺริยํ.  สหชาโต โนอาสโว เอโก ขนฺโธ ติณฺณนฺนํ ขนฺธานํ ฯเปฯ. (ยาว อสญฺญสตฺตา.) ปุเรชาตํ จกฺขุํ ฯเปฯ วตฺถุํ อนิจฺจโต ฯเปฯ โทมนสฺสํ อุปฺปชฺชติ. ทิพฺเพน จกฺขุนา รูปํ ปสฺสติ. ทิพฺพาย โสตธาตุยา สทฺทํ สุณาติ. รูปายตนํ จกฺขุ\-วิญฺญาณสฺส ฯเปฯ. โผฏฺฐพฺพา\-ยตนํ กายวิญฺญาณสฺส ฯเปฯ. จกฺขายตนํ จกฺขุ\-วิญฺญาณสฺส ฯเปฯ. กายายตนํ กายวิญฺญาณสฺส ฯเปฯ. วตฺถุ โนอาสวานํ ขนฺธานํ อตฺถิปจฺจเยน ปจฺจโย. ปจฺฉาชาตา โนอาสวา ขนฺธา ปุเรชาตสฺส อิมสฺส กายสฺส อตฺถิปจฺจเยน ปจฺจโย. กพฬีกาโร อาหาโร อิมสฺส กายสฺส อตฺถิปจฺจเยน ปจฺจโย. รูปชีวิตินฺทฺริยํ กฏตฺตารูปานํ อตฺถิปจฺจเยน ปจฺจโย. (1)}

\prose{โนอาสโว ธมฺโม อาสวสฺส ธมฺมสฺส อตฺถิปจฺจเยน ปจฺจโย. สหชาตํ, ปุเรชาตํ.  \textbf{สหชาตา} โนอาสวา ขนฺธา อาสวานํ อตฺถิปจฺจเยน ปจฺจโย. \textbf{ปุเรชาตํ} จกฺขุํ ฯเปฯ วตฺถุํ อสฺสาเทติ อภินนฺทติ, ตํ อารพฺภ อาสวา อุปฺปชฺชนฺติ. วตฺถุ อาสวานํ อตฺถิปจฺจเยน ปจฺจโย. (2)}

\prose{\csromanfolio{169}โนอาสโว ธมฺโม อาสวสฺส จ โนอาสวสฺส จ ธมฺมสฺส อตฺถิปจฺจเยน ปจฺจโย. สหชาตํ, ปุเรชาตํ.  \textbf{สหชาโต} โนอาสโว เอโก ขนฺโธ ติณฺณนฺนํ ขนฺธานํ อาสวานํ จิตฺตสมุฏฺฐานานญฺจ รูปานํ อตฺถิปจฺจเยน ปจฺจโย ฯเปฯ เทฺว ขนฺธา ฯเปฯ. (จกฺกํ.) (3)}

\proseitem{39}{อาสโว จ โนอาสโว จ ธมฺมา อาสวสฺส ธมฺมสฺส อตฺถิปจฺจเยน ปจฺจโย. สหชาตํ, ปุเรชาตํ.  \textbf{สหชาโต} กามาสโว จ สมฺปยุตฺตกา จ ขนฺธา ทิฏฺฐาสวสฺส อวิชฺชาสวสฺส อตฺถิปจฺจเยน ปจฺจโย. (จกฺกํ.) กามาสโว จ วตฺถุ จ ทิฏฺฐาสวสฺส อวิชฺชาสวสฺส อตฺถิปจฺจเยน ปจฺจโย. (จกฺกํ.) (1)}

\prose{อาสโว จ โนอาสโว จ ธมฺมา โนอาสวสฺส ธมฺมสฺส อตฺถิปจฺจเยน ปจฺจโย. สหชาตํ, ปุเรชาตํ, ปจฺฉาชาตํ, อาหารํ, อินฺทฺริยํ.  สหชาโต โนอาสโว เอโก ขนฺโธ จ อาสวา จ ติณฺณนฺนํ ขนฺธานํ จิตฺตสมุฏฺฐานานญฺจ รูปานํ อตฺถิปจฺจเยน ปจฺจโย. สหชาตา อาสวา จ มหาภูตา จ จิตฺต\-สมุฏฺฐานานํ รูปานํ อตฺถิปจฺจเยน ปจฺจโย. อาสวา จ วตฺถุ จ โนอาสวานํ ขนฺธานํ อตฺถิปจฺจเยน ปจฺจโย. ปจฺฉาชาตา อาสวา จ กพฬีกาโร อาหาโร จ อิมสฺส กายสฺส อตฺถิปจฺจเยน ปจฺจโย. ปจฺฉาชาตา อาสวา จ รูปชีวิตินฺทฺริยญฺจ กฏตฺตารูปานํ อตฺถิปจฺจเยน ปจฺจโย. (2)}

\prose{อาสโว จ โนอาสโว จ ธมฺมา อาสวสฺส จ โนอาสวสฺส จ ธมฺมสฺส อตฺถิปจฺจเยน ปจฺจโย. สหชาตํ, ปุเรชาตํ.  สหชาโต โนอาสโว เอโก ขนฺโธ จ กามาสโว จ ติณฺณนฺนํ ขนฺธานํ ทิฏฺฐาสวสฺส อวิชฺชาสวสฺส จิตฺตสมุฏฺฐานานญฺจ รูปานํ อตฺถิปจฺจเยน ปจฺจโย ฯเปฯ เทฺว ขนฺธา จ ฯเปฯ. (จกฺกํ.) สหชาโต กามาสโว จ วตฺถุ จ ทิฏฺฐาสวสฺส อวิชฺชาสวสฺส สมฺปยุตฺตกานญฺจ ขนฺธานํ อตฺถิปจฺจเยน ปจฺจโย. (จกฺกํ.) (3)}

\csromanmarkh{1. ปจฺจยานุโลม}\csromanmarkhii{2. สงฺขฺยาวาร}\csromanheadernomark{2}{1. \textbf{ปจฺจยานุโลม 2. สงฺขฺยาวาร}}
\proseitem{40}{เหตุยา สตฺต, อารมฺมเณ นว, อธิปติยา นว, อนนฺตเร นว, สมนนฺตเร นว, สหชาเต นว, อญฺญมญฺเญ นว, นิสฺสเย นว, \csromanfolio{170}อุปนิสฺสเย นว, ปุเรชาเต ตีณิ, ปจฺฉาชาเต ตีณิ, อาเสวเน นว, กมฺเม ตีณิ, วิปาเก เอกํ, อาหาเร ตีณิ ฯเปฯ มคฺเค นว, สมฺปยุตฺเต นว, วิปฺปยุตฺเต ปญฺจ, อตฺถิยา นว, นตฺถิยา นว, วิคเต นว, อวิคเต นว.}

\vspace{\tipitakaparskip}
\csromancenter{อนุโลมํ.}
\proseitem{41}{อาสโว ธมฺโม อาสวสฺส ธมฺมสฺส อารมฺมณ\-ปจฺจเยน ปจฺจโย. สหชาต\-ปจฺจเยน ปจฺจโย. อุปนิสฺสย\-ปจฺจเยน ปจฺจโย. (1)}

\prose{อาสโว ธมฺโม โนอาสวสฺส ธมฺมสฺส อารมฺมณ\-ปจฺจเยน ปจฺจโย. สหชาต\-ปจฺจเยน ปจฺจโย. อุปนิสฺสย\-ปจฺจเยน ปจฺจโย. ปจฺฉาชาต\-ปจฺจเยน ปจฺจโย. (2)}

\prose{อาสโว ธมฺโม อาสวสฺส จ โนอาสวสฺส จ ธมฺมสฺส อารมฺมณ\-ปจฺจเยน ปจฺจโย. สหชาต\-ปจฺจเยน ปจฺจโย. อุปนิสฺสย\-ปจฺจเยน ปจฺจโย. (3)}

\proseitem{42}{โนอาสโว ธมฺโม โนอาสวสฺส ธมฺมสฺส อารมฺมณ\-ปจฺจเยน ปจฺจโย. สหชาต\-ปจฺจเยน ปจฺจโย. อุปนิสฺสย\-ปจฺจเยน ปจฺจโย. ปุเรชาต\-ปจฺจเยน ปจฺจโย. ปจฺฉาชาต\-ปจฺจเยน ปจฺจโย. กมฺมปจฺจเยน ปจฺจโย. อาหารปจฺจเยน ปจฺจโย. อินฺทฺริย\-ปจฺจเยน ปจฺจโย. (1)}

\prose{โนอาสโว ธมฺโม อาสวสฺส ธมฺมสฺส อารมฺมณ\-ปจฺจเยน ปจฺจโย. สหชาต\-ปจฺจเยน ปจฺจโย. อุปนิสฺสย\-ปจฺจเยน ปจฺจโย. ปุเรชาต\-ปจฺจเยน ปจฺจโย. (2)}

\prose{โนอาสโว ธมฺโม อาสวสฺส จ โนอาสวสฺส จ ธมฺมสฺส อารมฺมณ\-ปจฺจเยน ปจฺจโย. สหชาต\-ปจฺจเยน ปจฺจโย. อุปนิสฺสย\-ปจฺจเยน ปจฺจโย. ปุเรชาต\-ปจฺจเยน ปจฺจโย. (3)}

\proseitem{43}{อาสโว จ โนอาสโว จ ธมฺมา อาสวสฺส ธมฺมสฺส อารมฺมณ\-ปจฺจเยน ปจฺจโย. สหชาต\-ปจฺจเยน ปจฺจโย. อุปนิสฺสย\-ปจฺจเยน ปจฺจโย. (1)}

\prose{\csromanfolio{171}อาสโว จ โนอาสโว จ ธมฺมา โนอาสวสฺส ธมฺมสฺส อารมฺมณ\-ปจฺจเยน ปจฺจโย. สหชาต\-ปจฺจเยน ปจฺจโย. อุปนิสฺสย\-ปจฺจเยน ปจฺจโย. ปจฺฉาชาต\-ปจฺจเยน ปจฺจโย. (2)}

\prose{อาสโว จ โนอาสโว จ ธมฺมา อาสวสฺส จ โนอาสวสฺส จ ธมฺมสฺส อารมฺมณ\-ปจฺจเยน ปจฺจโย. สหชาต\-ปจฺจเยน ปจฺจโย. อุปนิสฺสย\-ปจฺจเยน ปจฺจโย. (3)}

\csromanmarkh{2. ปจฺจย\-ปจฺจนีย}\csromanmarkhii{2. สงฺขฺยาวาร}\csromanheadernomark{2}{2. ปจฺจย\-ปจฺจนีย 2. สงฺขฺยาวาร}
\proseitem{44}{นเหตุยา นว, นอารมฺมเณ นว, นอธิปติยา นว, (สพฺพตฺถ นว,) โนอวิคเต นว.}

\vspace{\tipitakaparskip}
\csromancenter{\textbf{ปจฺจนียํ.}}
\proseitem{45}{เหตุปจฺจยา นอารมฺมเณ สตฺต, นอธิปติยา สตฺต, นอนนฺตเร สตฺต, นสมนนฺตเร สตฺต, นอญฺญมญฺเญ ตีณิ, นอุปนิสฺสเย สตฺต, (สพฺพตฺถ สตฺต,) นมคฺเค สตฺต, นสมฺปยุตฺเต ตีณิ, นวิปฺปยุตฺเต สตฺต, โนนตฺถิยา สตฺต, โนวิคเต สตฺต.}

\vspace{\tipitakaparskip}
\csromancenter{อนุโลม\-ปจฺจนียํ.}
\proseitem{46}{นเหตุปจฺจยา อารมฺมเณ นว, อธิปติยา นว. (อนุโลมปทา ปริปุณฺณา) อวิคเต นว.}

\vspace{\tipitakaparskip}
\csromancenter{ปจฺจนียา\-นุโลมํ.}
\nitthitam{อาสวทุกํ นิฏฺฐิตํ.}
\csromanmatikaanchor{mtk.19}
\csromantocmarkref{section}{15. สาสวทุก}{mtk.19}
\bookmark[dest={mtk.19},level=1]{15. สาสวทุก}
\csromanmarkh{15. สาสวทุก}\csromanmarkhii{1. ปฏิจฺจวาร}\csromanheadernomark{1}{15. \textbf{สาสวทุก 1. ปฏิจฺจวาร}}
\csromanheader{2}{\textbf{เหตุ}}
\proseitem{47}{\csromanfolio{172}สาสวํ ธมฺมํ ปฏิจฺจ สาสโว ธมฺโม อุปฺปชฺชติ เหตุปจฺจยา. สาสวํ เอกํ ขนฺธํ ปฏิจฺจ ตโย ขนฺธา จิตฺต\-สมุฏฺฐานญฺจ รูปํ ฯเปฯ เทฺว ขนฺเธ ฯเปฯ. ปฏิสนฺธิ\-กฺขเณ ฯเปฯ. ขนฺเธ ปฏิจฺจ วตฺถุ, วตฺถุํ ปฏิจฺจ ขนฺธา. เอกํ มหาภูตํ ฯเปฯ มหาภูเต ปฏิจฺจ จิตฺต\-สมุฏฺฐานํ รูปํ, กฏตฺตารูปํ, อุปาทารูปํ. (1)}

\prose{อนาสวํ ธมฺมํ ปฏิจฺจ อนาสโว ธมฺโม อุปฺปชฺชติ เหตุปจฺจยา. อนาสวํ เอกํ ขนฺธํ ปฏิจฺจ ตโย ขนฺธา ฯเปฯ เทฺว ขนฺเธ ฯเปฯ. (1)}

\prose{อนาสวํ ธมฺมํ ปฏิจฺจ สาสโว ธมฺโม อุปฺปชฺชติ เหตุปจฺจยา. อนาสเว ขนฺเธ ปฏิจฺจ จิตฺต\-สมุฏฺฐานํ รูปํ. (2)}

\prose{อนาสวํ ธมฺมํ ปฏิจฺจ สาสโว จ อนาสโว จ ธมฺมา อุปฺปชฺชนฺติ เหตุปจฺจยา. อนาสวํ เอกํ ขนฺธํ ปฏิจฺจ ตโย ขนฺธา จิตฺต\-สมุฏฺฐานญฺจ รูปํ ฯเปฯ เทฺว ขนฺเธ ฯเปฯ. (3)}

\prose{สาสวญฺจ อนาสวญฺจ ธมฺมํ ปฏิจฺจ สาสโว ธมฺโม อุปฺปชฺชติ เหตุปจฺจยา. อนาสเว ขนฺเธ จ มหาภูเต จ ปฏิจฺจ จิตฺต\-สมุฏฺฐานํ รูปํ. (1)}

\prosewithcloser{(ยถา จูฬนฺตรทุเก โลกิยทุกํ, เอวํ กาตพฺพํ, นินฺนานากรณํ.)}{สาสวทุกํ นิฏฺฐิตํ.}
\csromanmatikaanchor{mtk.20}
\csromantocmarkref{section}{16. อาสวสมฺปยุตฺตทุก}{mtk.20}
\bookmark[dest={mtk.20},level=1]{16. อาสวสมฺปยุตฺตทุก}
\csromanmarkh{16. อาสว\-สมฺปยุตฺตทุก}\csromanmarkhii{1. ปฏิจฺจวาร}\csromanheadernomark{1}{16. อาสว\-สมฺปยุตฺตทุก 1. ปฏิจฺจวาร}
\csromanmarkh{1. ปจฺจยานุโลม}\csromanmarkhii{1. วิภงฺควาร}\csromanheadernomark{2}{1. \textbf{ปจฺจยานุโลม 1. วิภงฺควาร}}
\csromanheader{2}{\textbf{เหตุ}}
\proseitem{48}{อาสว\-สมฺปยุตฺตํ ธมฺมํ ปฏิจฺจ อาสว\-สมฺปยุตฺโต ธมฺโม อุปฺปชฺชติ เหตุปจฺจยา. อาสว\-สมฺปยุตฺตํ เอกํ ขนฺธํ ปฏิจฺจ ตโย ขนฺธา ฯเปฯ เทฺว ขนฺเธ ฯเปฯ. (1)}

\prose{\csromanfolio{173}อาสว\-สมฺปยุตฺตํ ธมฺมํ ปฏิจฺจ อาสว\-วิปฺปยุตฺโต ธมฺโม อุปฺปชฺชติ เหตุปจฺจยา. อาสว\-สมฺปยุตฺเต ขนฺเธ ปฏิจฺจ จิตฺต\-สมุฏฺฐานํ รูปํ. โทมนสฺส\-สหคเต ขนฺเธ ปฏิจฺจ โมโห จิตฺต\-สมุฏฺฐานญฺจ รูปํ. (2)}

\prose{อาสว\-สมฺปยุตฺตํ ธมฺมํ ปฏิจฺจ อาสว\-สมฺปยุตฺโต จ อาสว\-วิปฺปยุตฺโต จ ธมฺมา อุปฺปชฺชนฺติ เหตุปจฺจยา. อาสว\-สมฺปยุตฺตํ เอกํ ขนฺธํ ปฏิจฺจ ตโย ขนฺธา จิตฺต\-สมุฏฺฐานญฺจ รูปํ ฯเปฯ เทฺว ขนฺเธ ฯเปฯ. โทมนสฺส\-สหคตํ เอกํ ขนฺธํ ปฏิจฺจ ตโย ขนฺธา โมโห จิตฺต\-สมุฏฺฐานญฺจ รูปํ ฯเปฯ เทฺว ขนฺเธ ฯเปฯ. (3)}

\proseitem{49}{อาสว\-วิปฺปยุตฺตํ ธมฺมํ ปฏิจฺจ อาสว\-วิปฺปยุตฺโต ธมฺโม อุปฺปชฺชติ เหตุปจฺจยา. อาสว\-วิปฺปยุตฺตํ เอกํ ขนฺธํ ปฏิจฺจ ตโย ขนฺธา จิตฺต\-สมุฏฺฐานญฺจ รูปํ ฯเปฯ เทฺว ขนฺเธ ฯเปฯ. โทมนสฺส\-สหคตํ วิจิกิจฺฉา\-สหคตํ อุทฺธจฺจ\-สหคตํ โมหํ ปฏิจฺจ จิตฺต\-สมุฏฺฐานํ รูปํ. ปฏิสนฺธิ\-กฺขเณ ฯเปฯ. ขนฺเธ ปฏิจฺจ วตฺถุ, วตฺถุํ ปฏิจฺจ ขนฺธา. เอกํ มหาภูตํ ปฏิจฺจ ตโย มหาภูตา, ตโย มหาภูเต ปฏิจฺจ เอกํ มหาภูตํ, เทฺว มหาภูเต ปฏิจฺจ เทฺว มหาภูตา. มหาภูเต ปฏิจฺจ จิตฺต\-สมุฏฺฐานํ รูปํ, กฏตฺตารูปํ, อุปาทารูปํ. (1)}

\prose{อาสว\-วิปฺปยุตฺตํ ธมฺมํ ปฏิจฺจ อาสว\-สมฺปยุตฺโต ธมฺโม อุปฺปชฺชติ เหตุปจฺจยา. โทมนสฺส\-สหคตํ วิจิกิจฺฉา\-สหคตํ อุทฺธจฺจ\-สหคตํ โมหํ ปฏิจฺจ สมฺปยุตฺตกา ขนฺธา. (2)}

\prose{อาสว\-วิปฺปยุตฺตํ ธมฺมํ ปฏิจฺจ อาสว\-สมฺปยุตฺโต จ อาสว\-วิปฺปยุตฺโต จ ธมฺมา อุปฺปชฺชนฺติ เหตุปจฺจยา. โทมนสฺส\-สหคตํ วิจิกิจฺฉา\-สหคตํ อุทฺธจฺจ\-สหคตํ โมหํ ปฏิจฺจ สมฺปยุตฺตกา ขนฺธา จิตฺต\-สมุฏฺฐานญฺจ รูปํ. (3)}

\proseitem{50}{อาสว\-สมฺปยุตฺตญฺจ อาสว\-วิปฺปยุตฺตญฺจ ธมฺมํ ปฏิจฺจ อาสว\-สมฺปยุตฺโต ธมฺโม อุปฺปชฺชติ เหตุปจฺจยา. โทมนสฺส\-สหคตํ วิจิกิจฺฉา\-สหคตํ อุทฺธจฺจ\-สหคตํ เอกํ ขนฺธญฺจ โมหญฺจ ปฏิจฺจ ตโย ขนฺธา ฯเปฯ เทฺว ขนฺเธ ฯเปฯ. (1)}

\prose{อาสว\-สมฺปยุตฺตญฺจ อาสว\-วิปฺปยุตฺตญฺจ ธมฺมํ ปฏิจฺจ อาสว\-วิปฺปยุตฺโต ธมฺโม อุปฺปชฺชติ เหตุปจฺจยา. อาสว\-สมฺปยุตฺเต ขนฺเธ จ มหาภูเต จ ปฏิจฺจ จิตฺต\-สมุฏฺฐานํ รูปํ. โทมนสฺส\-สหคเต วิจิกิจฺฉา\-สหคเต อุทฺธจฺจ\-สหคเต ขนฺเธ จ โมหญฺจ ปฏิจฺจ จิตฺต\-สมุฏฺฐานํ รูปํ. (2)}

\prose{\csromanfolio{174}อาสว\-สมฺปยุตฺตญฺจ อาสว\-วิปฺปยุตฺตญฺจ ธมฺมํ ปฏิจฺจ อาสว\-สมฺปยุตฺโต จ อาสว\-วิปฺปยุตฺโต จ ธมฺมา อุปฺปชฺชนฺติ เหตุปจฺจยา. โทมนสฺส\-สหคตํ วิจิกิจฺฉา\-สหคตํ อุทฺธจฺจ\-สหคตํ เอกํ ขนฺธญฺจ โมหญฺจ ปฏิจฺจ ตโย ขนฺธา จิตฺต\-สมุฏฺฐานญฺจ รูปํ. เทฺว ขนฺเธ ฯเปฯ. (3)}

\proseitem{51}{อาสว\-สมฺปยุตฺตํ ธมฺมํ ปฏิจฺจ อาวสมฺปยุตฺโต ธมฺโม อุปฺปชฺชติ อารมฺมณ\-ปจฺจยา. อาสว\-สมฺปยุตฺตํ เอกํ ขนฺธํ ปฏิจฺจ ตโย ขนฺธา ฯเปฯ เทฺว ขนฺเธ ฯเปฯ. (1)}

\prose{อาสว\-สมฺปยุตฺตํ ธมฺมํ ปฏิจฺจ อาสว\-วิปฺปยุตฺโต ธมฺโม อุปฺปชฺชติ อารมฺมณ\-ปจฺจยา. โทมนสฺส\-สหคเต วิจิกิจฺฉา\-สหคเต อุทฺธจฺจ\-สหคเต ขนฺเธ ปฏิจฺจ โมโห. (2)}

\prose{อาสว\-สมฺปยุตฺตํ ธมฺมํ ปฏิจฺจ อาสว\-สมฺปยุตฺโต จ อาสว\-วิปฺปยุตฺโต จ ธมฺมา อุปฺปชฺชนฺติ อารมฺมณ\-ปจฺจยา. โทมนสฺส\-สหคตํ วิจิกิจฺฉา\-สหคตํ อุทฺธจฺจ\-สหคตํ เอกํ ขนฺธํ ปฏิจฺจ ตโย ขนฺธา โมโห จ ฯเปฯ เทฺว ขนฺเธ ฯเปฯ. (3)}

\proseitem{52}{อาสว\-วิปฺปยุตฺตํ ธมฺมํ ปฏิจฺจ อาสว\-วิปฺปยุตฺโต ธมฺโม อุปฺปชฺชติ อารมฺมณ\-ปจฺจยา. อาสว\-วิปฺปยุตฺตํ เอกํ ขนฺธํ ปฏิจฺจ ตโย ขนฺธา ฯเปฯ เทฺว ขนฺเธ ฯเปฯ. ปฏิสนฺธิ\-กฺขเณ ฯเปฯ. วตฺถุํ ปฏิจฺจ ขนฺธา. (1)}

\prose{อาสว\-วิปฺปยุตฺตํ ธมฺมํ ปฏิจฺจ อาสว\-สมฺปยุตฺโต ธมฺโม อุปฺปชฺชติ อารมฺมณ\-ปจฺจยา. โทมนสฺส\-สหคตํ วิจิกิจฺฉา\-สหคตํ อุทฺธจฺจ\-สหคตํ โมหํ ปฏิจฺจ สมฺปยุตฺตกา ขนฺธา. (2)}

\prose{อาสว\-สมฺปยุตฺตญฺจ อาสว\-วิปฺปยุตฺตญฺจ ธมฺมํ ปฏิจฺจ อาสว\-สมฺปยุตฺโต ธมฺโม อุปฺปชฺชติ อารมฺมณ\-ปจฺจยา. โทมนสฺส\-สหคตํ วิจิกิจฺฉา\-สหคตํ อุทฺธจฺจ\-สหคตํ เอกํ ขนฺธญฺจ โมหญฺจ ปฏิจฺจ ตโย ขนฺธา ฯเปฯ เทฺว ขนฺเธ ฯเปฯ. (1)}

\proseitem{53}{อาสว\-สมฺปยุตฺตํ ธมฺมํ ปฏิจฺจ อาสว\-สมฺปยุตฺโต ธมฺโม อุปฺปชฺชติ อธิปติ\-ปจฺจยา. ตีณิ.}

\prose{\csromanfolio{175}อาสว\-วิปฺปยุตฺตํ ธมฺมํ ปฏิจฺจ อาสว\-วิปฺปยุตฺโต ธมฺโม อุปฺปชฺชติ อธิปติ\-ปจฺจยา. อาสว\-วิปฺปยุตฺตํ เอกํ ขนฺธํ ปฏิจฺจ ตโย ขนฺธา จิตฺต\-สมุฏฺฐานญฺจ รูปํ ฯเปฯ เทฺว ขนฺเธ ฯเปฯ. โทมนสฺส\-สหคตํ โมหํ ปฏิจฺจ จิตฺต\-สมุฏฺฐานํ รูปํ. เอกํ มหาภูตํ ฯเปฯ มหาภูเต ปฏิจฺจ จิตฺต\-สมุฏฺฐานํ รูปํ, อุปาทารูปํ. (1)}

\prose{อาสว\-วิปฺปยุตฺตํ ธมฺมํ ปฏิจฺจ อาสว\-สมฺปยุตฺโต ธมฺโม อุปฺปชฺชติ อธิปติ\-ปจฺจยา. โทมนสฺส\-สหคตํ โมหํ ปฏิจฺจ สมฺปยุตฺตกา ขนฺธา. (2)}

\prose{อาสว\-วิปฺปยุตฺตํ ธมฺมํ ปฏิจฺจ อาสว\-สมฺปยุตฺโต จ อาสว\-วิปฺปยุตฺโต จ ธมฺมา อุปฺปชฺชนฺติ อธิปติ\-ปจฺจยา. โทมนสฺส\-สหคตํ โมหํ ปฏิจฺจ สมฺปยุตฺตกา ขนฺธา จิตฺต\-สมุฏฺฐานญฺจ รูปํ. (3)}

\proseitem{54}{อาสว\-สมฺปยุตฺตญฺจ อาสว\-วิปฺปยุตฺตญฺจ ธมฺมํ ปฏิจฺจ อาสว\-สมฺปยุตฺโต ธมฺโม อุปฺปชฺชติ อธิปติ\-ปจฺจยา. โทมนสฺส\-สหคตํ เอกํ ขนฺธญฺจ โมหญฺจ ปฏิจฺจ ตโย ขนฺธา ฯเปฯ เทฺว ขนฺเธ ฯเปฯ. (1)}

\prose{อาสว\-สมฺปยุตฺตญฺจ อาสว\-วิปฺปยุตฺตญฺจ ธมฺมํ ปฏิจฺจ อาสว\-วิปฺปยุตฺโต ธมฺโม อุปฺปชฺชติ อธิปติ\-ปจฺจยา. อาสว\-สมฺปยุตฺเต ขนฺเธ จ มหาภูเต จ ปฏิจฺจ จิตฺต\-สมุฏฺฐานํ รูปํ. โทมนสฺส\-สหคเต ขนฺเธ จ โมหญฺจ ปฏิจฺจ จิตฺต\-สมุฏฺฐานํ รูปํ.}

\prose{อาสว\-สมฺปยุตฺตญฺจ อาสว\-วิปฺปยุตฺตญฺจ ธมฺมํ ปฏิจฺจ อาสว\-สมฺปยุตฺโต จ อาสว\-วิปฺปยุตฺโต จ ธมฺมา อุปฺปชฺชนฺติ อธิปติ\-ปจฺจยา. โทมนสฺส\-สหคตํ เอกํ ขนฺธญฺจ โมหญฺจ ปฏิจฺจ ตโย ขนฺธา จิตฺต\-สมุฏฺฐานญฺจ รูปํ ฯเปฯ เทฺว ขนฺเธ ฯเปฯ. (3) (เอวํ สพฺเพ ปจฺจยา วิตฺถาเรตพฺพา. สํขิตฺตํ.)}

\csromanmarkh{1. ปจฺจยานุโลม}\csromanmarkhii{2. สงฺขฺยาวาร}\csromanheadernomark{2}{1. \textbf{ปจฺจยานุโลม 2. สงฺขฺยาวาร}}
\proseitem{55}{เหตุยา นว, อารมฺมเณ ฉ, อธิปติยา นว, อนนฺตเร ฉ, สมนนฺตเร ฉ, สหชาเต นว, อญฺญมญฺเญ ฉ, นิสฺสเย นว, อุปนิสฺสเย ฉ, ปุเรชาเต ฉ, อาเสวเน ฉ, กมฺเม นว, วิปาเก \csromanfolio{176}เอกํ, อาหาเร นว, อินฺทฺริเย นว, ฌาเน นว, มคฺเค นว, สมฺปยุตฺเต ฉ, วิปฺปยุตฺเต นว, อตฺถิยา นว, นตฺถิยา ฉ, วิคเต ฉ, อวิคเต นว.}

\vspace{\tipitakaparskip}
\csromancenter{อนุโลมํ.}
\csromanmarkh{2. ปจฺจย\-ปจฺจนีย}\csromanmarkhii{1. วิภงฺควาร}\csromanheadernomark{2}{2. ปจฺจย\-ปจฺจนีย 1. วิภงฺควาร}
\proseitem{55}{อาสว\-สมฺปยุตฺตํ ธมฺมํ ปฏิจฺจ อาสว\-วิปฺปยุตฺโต ธมฺโม อุปฺปชฺชติ นเหตุปจฺจยา. วิจิกิจฺฉา\-สหคเต อุทฺธจฺจ\-สหคเต ขนฺเธ ปฏิจฺจ วิจิกิจฺฉา\-สหคโต อุทฺธจฺจ\-สหคโต โมโห. (1)}

\prose{อาสว\-วิปฺปยุตฺตํ ธมฺมํ ปฏิจฺจ อาสว\-วิปฺปยุตฺโต ธมฺโม อุปฺปชฺชติ นเหตุปจฺจยา. อเหตุกํ อาสว\-วิปฺปยุตฺตํ เอกํ ขนฺธํ ปฏิจฺจ ตโย ขนฺธา จิตฺต\-สมุฏฺฐานญฺจ รูปํ ฯเปฯ เทฺว ขนฺเธ ฯเปฯ. อเหตุก\-ปฏิสนฺธิ\-กฺขเณ ฯเปฯ. (ยาว อสญฺญสตฺตา.) (1)}

\prose{นอารมฺมณ}

\proseitem{55}{อาสว\-สมฺปยุตฺตํ ธมฺมํ ปฏิจฺจ อาสว\-วิปฺปยุตฺโต ธมฺโม อุปฺปชฺชติ นอารมฺมณ\-ปจฺจโย. อาสว\-สมฺปยุตฺเต ขนฺเธ ปฏิจฺจ จิตฺต\-สมุฏฺฐานํ รูปํ. (1)}

\prose{อาสว\-วิปฺปยุตฺตํ ธมฺมํ ปฏิจฺจ อาสว\-วิปฺปยุตฺโต ธมฺโม อุปฺปชฺชติ นอารมฺมณ\-ปจฺจยา. อาสววิปยุตฺเต ขนฺเธ ปฏิจฺจ จิตฺต\-สมุฏฺฐานํ รูปํ. โทมนสฺส\-สหคตํ วิจิกิจฺฉา\-สหคตํ อุทฺธจฺจ\-สหคตํ โมหํ ปฏิจฺจ จิตฺต\-สมุฏฺฐานํ รูปํ. ปฏิสนฺธิ\-กฺขเณ ฯเปฯ. ขนฺเธ ปฏิจฺจ วตฺถุ. เอกํ มหาภูตํ ฯเปฯ. (ยาว อสญฺญสตฺตา.) (1)}

\prose{อาสว\-สมฺปยุตฺตญฺจ อาสว\-วิปฺปยุตฺตญฺจ ธมฺมํ ปฏิจฺจ อาสว\-วิปฺปยุตฺโต ธมฺโม อุปฺปชฺชติ นอารมฺมณ\-ปจฺจยา. อาสว\-สมฺปยุตฺเต ขนฺเธ จ มหาภูเต จ ปฏิจฺจ จิตฺต\-สมุฏฺฐานํ รูปํ. โทมนสฺส\-สหคเต วิจิกิจฺฉา\-สหคเต อุทฺธจฺจ\-สหคเต ขนฺเธ จ โมหญฺจ ปฏิจฺจ จิตฺต\-สมุฏฺฐานํ รูปํ. (1)}

\csromanheader{2}{16. อาสว\-สมฺปยุตฺตทุก นอธิปติ}
\proseitem{58}{\csromanfolio{177}อาสว\-สมฺปยุตฺตํ ธมฺมํ ปฏิจฺจ อาสว\-สมฺปยุตฺโต ธมฺโม อุปฺปชฺชติ นอธิปติ\-ปจฺจยา. (สํขิตฺตํ.)}

\csromanheader{2}{\textbf{นปุเรชาตาทิ}}
\proseitem{59}{อาสว\-สมฺปยุตฺตํ ธมฺมํ ปฏิจฺจ อาสว\-สมฺปยุตฺโต ธมฺโม อุปฺปชฺชติ นปุเรชาต\-ปจฺจยา. อรูเป อาสว\-สมฺปยุตฺตํ เอกํ ขนฺธํ ปฏิจฺจ ตโย ขนฺธา ฯเปฯ เทฺว ขนฺเธ ฯเปฯ. (1)}

\prose{อาสว\-สมฺปยุตฺตํ ธมฺมํ ปฏิจฺจ อาสว\-วิปฺปยุตฺโต ธมฺโม อุปฺปชฺชติ นปุเรชาต\-ปจฺจยา. อรูเป วิจิกิจฺฉา\-สหคเต อุทฺธจฺจ\-สหคเต ขนฺเธ ปฏิจฺจ วิจิกิจฺฉา\-สหคโต อุทฺธจฺจ\-สหคโต โมโห. อาสว\-สมฺปยุตฺเต ขนฺเธ ปฏิจฺจ จิตฺต\-สมุฏฺฐานํ รูปํ. (2)}

\prose{อาสว\-สมฺปยุตฺตํ ธมฺมํ ปฏิจฺจ อาสว\-สมฺปยุตฺโต จ อาสว\-วิปฺปยุตฺโต จ ธมฺมา อุปฺปชฺชนฺติ นปุเรชาต\-ปจฺจยา. อรูเป วิจิกิจฺฉา\-สหคตํ อุทฺธจฺจ\-สหคตํ เอกํ ขนฺธํ ปฏิจฺจ ตโย ขนฺธา โมโห จ ฯเปฯ เทฺว ขนฺเธ ฯเปฯ. (3)}

\proseitem{60}{อาสว\-วิปฺปยุตฺตํ ธมฺมํ ปฏิจฺจ อาสว\-วิปฺปยุตฺโต ธมฺโม อุปฺปชฺชติ นปุเรชาต\-ปจฺจยา. อรูเป อาสว\-วิปฺปยุตฺตํ เอกํ ขนฺธํ ปฏิจฺจ ตโย ขนฺธา ฯเปฯ เทฺว ขนฺเธ ฯเปฯ. อาสว\-วิปฺปยุตฺเต ขนฺเธ ปฏิจฺจ จิตฺต\-สมุฏฺฐานํ รูปํ. โทมนสฺส\-สหคตํ วิจิกิจฺฉา\-สหคตํ อุทฺธจฺจ\-สหคตํ โมหํ ปฏิจฺจ จิตฺต\-สมุฏฺฐานํ รูปํ. ปฏิสนฺธิ\-กฺขเณ ฯเปฯ. (ยาว อสญฺญสตฺตา.) (1)}

\prose{อาสว\-วิปฺปยุตฺตํ ธมฺมํ ปฏิจฺจ อาสว\-สมฺปยุตฺโต ธมฺโม อุปฺปชฺชติ นปุเรชาต\-ปจฺจยา. อรูเป วิจิกิจฺฉา\-สหคตํ อุทฺธจฺจ\-สหคตํ โมหํ ปฏิจฺจ สมฺปยุตฺตกา ขนฺธา. (2)}

\proseitem{61}{อาสว\-สมฺปยุตฺตญฺจ อาสว\-วิปฺปยุตฺตญฺจ ธมฺมํ ปฏิจฺจ อาสว\-สมฺปยุตฺโต ธมฺโม อุปฺปชฺชติ นปุเรชาต\-ปจฺจยา. อรูเป วิจิกิจฺฉา\-สหคตํ อุทฺธจฺจ\-สหคตํ เอกํ ขนฺธญฺจ โมหญฺจ ปฏิจฺจ ตโย ขนฺธา ฯเปฯ เทฺว ขนฺเธ จ ฯเปฯ. (1)}

\prose{อาสว\-สมฺปยุตฺตญฺจ อาสว\-วิปฺปยุตฺตญฺจ ธมฺมํ ปฏิจฺจ อาสว\-วิปฺปยุตฺโต ธมฺโม อุปฺปชฺชติ นปุเรชาต\-ปจฺจยา. อาสว\-สมฺปยุตฺเต ขนฺเธ จ มหาภูเต \csromanfolio{178}จ ปฏิจฺจ จิตฺต\-สมุฏฺฐานํ รูปํ. โทมนสฺส\-สหคเต วิจิกิจฺฉา\-สหคเต อุทฺธจฺจ\-สหคเต ขนฺเธ จ โมหญฺจ ปฏิจฺจ จิตฺต\-สมุฏฺฐานํ รูปํ. (2) (นปจฺฉาชาต\-ปจฺจยา นว, นอาเสวน\-ปจฺจยา นว.)}

\csromanheader{2}{\textbf{นกมฺมาทิ}}
\proseitem{62}{อาสว\-สมฺปยุตฺตํ ธมฺมํ ปฏิจฺจ อาสว\-สมฺปยุตฺโต ธมฺโม อุปฺปชฺชติ นกมฺมปจฺจยา. อาสว\-สมฺปยุตฺเต ขนฺเธ ปฏิจฺจ สมฺปยุตฺตกา เจตนา. (1)}

\prose{อาสว\-วิปฺปยุตฺตํ ธมฺมํ ปฏิจฺจ อาสว\-วิปฺปยุตฺโต ธมฺโม อุปฺปชฺชติ นกมฺมปจฺจยา. อาสว\-วิปฺปยุตฺเต ขนฺเธ ปฏิจฺจ วิปฺปยุตฺตกา เจตนา. (1)}

\prose{อาสว\-วิปฺปยุตฺตํ ธมฺมํ ปฏิจฺจ อาสว\-สมฺปยุตฺโต ธมฺโม อุปฺปชฺชติ นกมฺมปจฺจยา. โทมนสฺส\-สหคตํ วิจิกิจฺฉา\-สหคตํ อุทฺธจฺจ\-สหคตํ โมหํ ปฏิจฺจ สมฺปยุตฺตกา เจตนา. (2)}

\prose{อาสว\-สมฺปยุตฺตญฺจ อาสว\-วิปฺปยุตฺตญฺจ ธมฺมํ ปฏิจฺจ อาสว\-สมฺปยุตฺโต ธมฺโม อุปฺปชฺชติ นกมฺมปจฺจยา. โทมนสฺส\-สหคเต วิจิกิจฺฉา\-สหคเต อุทฺธจฺจ\-สหคเต ขนฺเธ จ โมหญฺจ ปฏิจฺจ สมฺปยุตฺตกา เจตนา. นวิปาก\-ปจฺจยา. นอาหารปจฺจยา. นอินฺทฺริยปจฺจยา. นฌานปจฺจยา. นมคฺคปจฺจยา. นสมฺปยุตฺต\-ปจฺจยา. นวิปฺปยุตฺต\-ปจฺจยา. โนนตฺถิปจฺจยา. โนวิคต\-ปจฺจยา.}

\csromanmarkh{2. ปจฺจย\-ปจฺจนีย}\csromanmarkhii{2. สงฺขฺยาวาร}\csromanheadernomark{2}{2. ปจฺจย\-ปจฺจนีย 2. สงฺขฺยาวาร}
\proseitem{61}{นเหตุยา เทฺว, นอารมฺมเณ ตีณิ, นอธิปติยา นว, นอนนฺตเร ตีณิ, นสมนนฺตเร ตีณิ, นอญฺญมญฺเญ ตีณิ, นอุปนิสฺสเย ตีณิ, นปุเรชาเต สตฺต, นปจฺฉาชาเต นว, นอาเสวเน นว, นกมฺเม จตฺตาริ, นวิปาเก นว, นอาหาเร เอกํ, นอินฺทฺริเย เอกํ, นฌาเน เอกํ, นมคฺเค เอกํ, นสมฺปยุตฺเต ตีณิ, นวิปฺปยุตฺเต ฉ, โนนตฺถิยา ตีณิ, โนวิคเต ตีณิ.}

\vspace{\tipitakaparskip}
\csromancenter{ปจฺจนียํ.}
\csromanmarkh{16. อาสว\-สมฺปยุตฺตทุก}\csromanmarkhii{3. ปจฺจยา\-นุโลม\-ปจฺจนีย}\csromanheadernomark{2}{16. อาสว\-สมฺปยุตฺตทุก 3. ปจฺจยา\-นุโลม\-ปจฺจนีย}
\vspace{\tipitakaparskip}
\csromancenter{\textbf{เหตุ}ปจฺจยา นอารมฺมเณ ตีณิ ฯเปฯ นปุเรชาเต ฉ ฯเปฯ นวิปาเก นว ฯเปฯ นสมฺปยุตฺเต ตีณิ, นวิปฺปยุตฺเต จตฺตาริ, โนนตฺถิยา ตีณิ, โนวิคเต ตีณิ.}
\csromancenter{อนุโลม\-ปจฺจนียํ.}
\proseitem{65}{\csromanfolio{179}นเหตุปจฺจยา อารมฺมเณ เทฺว, อนนฺตเร เทฺว, สมนนฺตเร เทฺว ฯเปฯ กมฺเม เทฺว, วิปาเก เอกํ, อาหาเร เทฺว ฯเปฯ มคฺเค เอกํ, สมฺปยุตฺเต เทฺว, วิปฺปยุตฺเต เทฺว ฯเปฯ อวิคเต เทฺว.}

\vspace{\tipitakaparskip}
\csromancenter{(สหชาตวาโร ปฏิจฺจ\-วาร\-สทิโส.)}
\csromanmarkh{16. อาสว\-สมฺปยุตฺตทุก}\csromanmarkhii{3. ปจฺจยวาร}\csromanheadernomark{2}{16. อาสว\-สมฺปยุตฺตทุก 3. ปจฺจยวาร}
\csromanmarkh{1. ปจฺจยานุโลม}\csromanmarkhii{1. วิภงฺควาร}\csromanheadernomark{2}{1. \textbf{ปจฺจยานุโลม 1. วิภงฺควาร}}
\csromanheader{2}{\textbf{เหตุ}}
\proseitem{66}{อาสว\-สมฺปยุตฺตํ ธมฺมํ ปจฺจยา อาสว\-สมฺปยุตฺโต ธมฺโม อุปฺปชฺชติ เหตุปจฺจยา. ตีณิ. (ปฏิจฺจ\-วาร\-สทิโส.)}

\prose{อาสว\-วิปฺปยุตฺตํ ธมฺมํ ปจฺจยา อาสว\-วิปฺปยุตฺโต ธมฺโม อุปฺปชฺชติ เหตุปจฺจยา. อาสว\-วิปฺปยุตฺตํ เอกํ ขนฺธํ ปจฺจยา ตโย ขนฺธา จิตฺต\-สมุฏฺฐานญฺจ รูปํ ฯเปฯ เทฺว ขนฺเธ ฯเปฯ. โทมนสฺส\-สหคตํ วิจิกิจฺฉา\-สหคตํ อุทฺธจฺจ\-สหคตํ โมหํ ปจฺจยา จิตฺต\-สมุฏฺฐานํ รูปํ. ปฏิสนฺธิ\-กฺขเณ ฯเปฯ. ขนฺเธ ปจฺจยา วตฺถุ, วตฺถุํ ปจฺจยา ขนฺธา. เอกํ มหาภูตํ ฯเปฯ มหาภูเต ปจฺจยา จิตฺต\-สมุฏฺฐานํ \csromanfolio{180}รูปํ, กฏตฺตารูปํ, อุปาทารูปํ. วตฺถุํ ปจฺจยา อาสว\-วิปฺปยุตฺตา ขนฺธา. วตฺถุํ ปจฺจยา โทมนสฺส\-สหคโต โมโห. (1)}

\prose{อาสว\-วิปฺปยุตฺตํ ธมฺมํ ปจฺจยา อาสว\-สมฺปยุตฺโต ธมฺโม อุปฺปชฺชติ เหตุปจฺจยา. วตฺถุํ ปจฺจยา อาสว\-สมฺปยุตฺตกา ขนฺธา. โทมนสฺส\-สหคตํ วิจิกิจฺฉา\-สหคตํ อุทฺธจฺจ\-สหคตํ โมหํ ปจฺจยา สมฺปยุตฺตกา ขนฺธา. (2)}

\prose{อาสว\-วิปฺปยุตฺตํ ธมฺมํ ปจฺจยา อาสว\-สมฺปยุตฺโต จ อาสว\-วิปฺปยุตฺโต จ ธมฺมา อุปฺปชฺชนฺติ เหตุปจฺจยา. วตฺถุํ ปจฺจยา อาสว\-สมฺปยุตฺตกา ขนฺธา. มหาภูเต ปจฺจยา จิตฺต\-สมุฏฺฐานํ รูปํ. โทมนสฺส\-สหคตํ วิจิกิจฺฉา\-สหคตํ อุทฺธจฺจ\-สหคตํ โมหํ ปจฺจยา สมฺปยุตฺตกา ขนฺธา จิตฺต\-สมุฏฺฐานญฺจ รูปํ. วตฺถุํ ปจฺจยา โทมนสฺส\-สหคตา ขนฺธา จ โมโห จ. (3)}

\prose{อาสว\-สมฺปยุตฺตญฺจ อาสว\-วิปฺปยุตฺตญฺจ ธมฺมํ ปจฺจยา อาสว\-สมฺปยุตฺโต ธมฺโม อุปฺปชฺชติ เหตุปจฺจยา. อาสว\-สมฺปยุตฺตํ เอกํ ขนฺธญฺจ วตฺถุญฺจ ปจฺจยา ตโย ขนฺธา ฯเปฯ เทฺว ขนฺเธ จ ฯเปฯ. โทมนสฺส\-สหคตํ วิจิกิจฺฉา\-สหคตํ อุทฺธจฺจ\-สหคตํ เอกํ ขนฺธญฺจ โมหญฺจ ปจฺจยา ตโย ขนฺธา ฯเปฯ เทฺว ขนฺเธ จ ฯเปฯ. (1)}

\prose{อาสว\-สมฺปยุตฺตญฺจ อาสว\-วิปฺปยุตฺตญฺจ ธมฺมํ ปจฺจยา อาสว\-วิปฺปยุตฺโต ธมฺโม อุปฺปชฺชติ เหตุปจฺจยา. อาสว\-สมฺปยุตฺเต ขนฺเธ จ มหาภูเต จ ปจฺจยา จิตฺต\-สมุฏฺฐานํ รูปํ. โทมนสฺส\-สหคเต วิจิกิจฺฉา\-สหคเต อุทฺธจฺจ\-สหคเต ขนฺเธ จ โมหญฺจ ปจฺจยา จิตฺต\-สมุฏฺฐานํ รูปํ. โทมนสฺส\-สหคเต ขนฺเธ จ วตฺถุญฺจ ปจฺจยา โทมนสฺส\-สหคโต โมโห. (2)}

\prose{อาสว\-สมฺปยุตฺตญฺจ อาสว\-วิปฺปยุตฺตญฺจ ธมฺมํ ปจฺจยา อาสว\-สมฺปยุตฺโต จ อาสว\-วิปฺปยุตฺโต จ ธมฺมา อุปฺปชฺชนฺติ เหตุปจฺจยา. อาสว\-สมฺปยุตฺตํ เอกํ ขนฺธญฺจ วตฺถุญฺจ ปจฺจยา ตโย ขนฺธา ฯเปฯ เทฺว ขนฺเธ ฯเปฯ. อาสว\-สมฺปยุตฺเต ขนฺเธ จ มหาภูเต จ ปจฺจยา จิตฺต\-สมุฏฺฐานํ รูปํ. โทมนสฺส\-สหคตํ วิจิกิจฺฉา\-สหคตํ อุทฺธจฺจ\-สหคตํ เอกํ ขนฺธญฺจ โมหญฺจ ปจฺจยา ตโย ขนฺธา จิตฺต\-สมุฏฺฐานญฺจ รูปํ ฯเปฯ เทฺว ขนฺเธ ฯเปฯ. โทมนสฺส\-สหคตํ เอกํ ขนฺธญฺจ วตฺถุญฺจ ปจฺจยา ตโย ขนฺธา โมโห จ ฯเปฯ เทฺว ขนฺเธ ฯเปฯ. (3)}

\csromanheader{2}{16. อาสว\-สมฺปยุตฺตทุก อารมฺมณาทิ}
\proseitem{68}{\csromanfolio{181}อาสว\-สมฺปยุตฺตํ ธมฺมํ ปจฺจยา อาสว\-สมฺปยุตฺโต ธมฺโม อุปฺปชฺชติ อารมฺมณ\-ปจฺจยา. ตีณิ. (ปฏิจฺจ\-วาร\-สทิสา.)}

\prose{อาสว\-วิปฺปยุตฺตํ ธมฺมํ ปจฺจยา อาสว\-วิปฺปยุตฺโต ธมฺโม อุปฺปชฺชติ อารมฺมณ\-ปจฺจยา. อาสว\-วิปฺปยุตฺตํ เอกํ ขนฺธํ ปจฺจยา ตโย ขนฺธา ฯเปฯ เทฺว ขนฺเธ ฯเปฯ. ปฏิสนฺธิ\-กฺขเณ ฯเปฯ. วตฺถุํ ปจฺจยา ขนฺธา. จกฺขายตนํ ปจฺจยา จกฺขุวิญฺญาณํ ฯเปฯ. กายายตนํ ปจฺจยา กายวิญฺญาณํ. วตฺถุํ ปจฺจยา อาสว\-วิปฺปยุตฺตา ขนฺธา. วตฺถุํ ปจฺจยา โทมนสฺส\-สหคโต วิจิกิจฺฉา\-สหคโต อุทฺธจฺจ\-สหคโต โมโห. (1)}

\prose{อาสว\-วิปฺปยุตฺตํ ธมฺมํ ปจฺจยา อาสว\-สมฺปยุตฺโต ธมฺโม อุปฺปชฺชติ อารมฺมณ\-ปจฺจยา. วตฺถุํ ปจฺจยา อาสว\-สมฺปยุตฺตกา ขนฺธา. โทมนสฺส\-สหคตํ วิจิกิจฺฉา\-สหคตํ อุทฺธจฺจ\-สหคตํ โมหํ ปจฺจยา สมฺปยุตฺตกา ขนฺธา. (2)}

\prose{อาสว\-วิปฺปยุตฺตํ ธมฺมํ ปจฺจยา อาสว\-สมฺปยุตฺโต จ อาสว\-วิปฺปยุตฺโต จ ธมฺมา อุปฺปชฺชนฺติ อารมฺมณ\-ปจฺจยา. วตฺถุํ ปจฺจยา โทมนสฺส\-สหคตา วิจิกิจฺฉา\-สหคตา อุทฺธจฺจ\-สหคตา ขนฺธา โมโห จ. (3)}

\proseitem{69}{อาสว\-สมฺปยุตฺตญฺจ อาสว\-วิปฺปยุตฺตญฺจ ธมฺมํ ปจฺจยา อาสว\-สมฺปยุตฺโต ธมฺโม อุปฺปชฺชติ อารมฺมณ\-ปจฺจยา. อาสว\-สมฺปยุตฺตํ เอกํ ขนฺธญฺจ วตฺถุญฺจ ปจฺจยา ตโย ขนฺธา ฯเปฯ เทฺว ขนฺเธ ฯเปฯ. (1)}

\prose{อาสว\-สมฺปยุตฺตญฺจ อาสว\-วิปฺปยุตฺตญฺจ ธมฺมํ ปจฺจยา อาสว\-วิปฺปยุตฺโต ธมฺโม อุปฺปชฺชติ อารมฺมณ\-ปจฺจยา. โทมนสฺส\-สหคเต วิจิกิจฺฉา\-สหคเต อุทฺธจฺจ\-สหคเต ขนฺเธ จ วตฺถุญฺจ ปจฺจยา โทมนสฺส\-สหคโต วิจิกิจฺฉา\-สหคโต อุทฺธจฺจ\-สหคโต โมโห. (2)}

\prose{อาสว\-สมฺปยุตฺตญฺจ อาสว\-วิปฺปยุตฺตญฺจ ธมฺมํ ปจฺจยา อาสว\-สมฺปยุตฺโต จ อาสว\-วิปฺปยุตฺโต จ ธมฺมา อุปฺปชฺชนฺติ อารมฺมณ\-ปจฺจยา. โทมนสฺส\-สหคตํ วิจิกิจฺฉา\-สหคตํ อุทฺธจฺจ\-สหคตํ เอกํ ขนฺธญฺจ วตฺถุญฺจ ปจฺจยา ตโย ขนฺธา โมโห จ ฯเปฯ เทฺว ขนฺเธ ฯเปฯ. อธิปติ\-ปจฺจยา. อนนฺตร\-ปจฺจยา ฯเปฯ อวิคตปจฺจยา. (3)}

\csromanmarkh{1. ปจฺจยานุโลม}\csromanmarkhii{2. สงฺขฺยาวาร}\csromanheadernomark{2}{1. \textbf{ปจฺจยานุโลม 2. สงฺขฺยาวาร}}
\proseitem{70}{\csromanfolio{182}เหตุยา นว, อารมฺมเณ นว, (สพฺพตฺถ นว.) กมฺเม นว, วิปาเก เอกํ ฯเปฯ อวิคเต นว.}

\vspace{\tipitakaparskip}
\csromancenter{อนุโลมํ.}
\csromanmarkh{2. ปจฺจย\-ปจฺจนีย}\csromanmarkhii{1. วิภงฺควาร}\csromanheadernomark{2}{2. ปจฺจย\-ปจฺจนีย 1. วิภงฺควาร}
\proseitem{71}{อาสว\-สมฺปยุตฺตํ ธมฺมํ ปจฺจยา อาสว\-วิปฺปยุตฺโต ธมฺโม อุปฺปชฺชติ นเหตุปจฺจยา. วิจิกิจฺฉา\-สหคเต อุทฺธจฺจ\-สหคเต ขนฺเธ ปจฺจยา วิจิกิจฺฉา\-สหคโต อุทฺธจฺจ\-สหคโต โมโห. (1)}

\prose{อาสว\-วิปฺปยุตฺตํ ธมฺมํ ปจฺจยา อาสว\-วิปฺปยุตฺโต ธมฺโม อุปฺปชฺชติ นเหตุปจฺจยา. อเหตุกํ อาสว\-วิปฺปยุตฺตํ เอกํ ขนฺธํ ปจฺจยา ตโย ขนฺธา จิตฺต\-สมุฏฺฐานญฺจ รูปํ ฯเปฯ เทฺว ขนฺเธ ฯเปฯ. อเหตุก\-ปฏิสนฺธิ\-กฺขเณ ฯเปฯ. (ยาว อสญฺญสตฺตา.) จกฺขายตนํ ปจฺจยา จกฺขุวิญฺญาณํ ฯเปฯ. กายายตนํ ปจฺจยา กายวิญฺญาณํ. วตฺถุํ ปจฺจยา อเหตุกา อาสว\-วิปฺปยุตฺตา ขนฺธา. วตฺถุํ ปจฺจยา วิจิกิจฺฉา\-สหคโต อุทฺธจฺจ\-สหคโต โมโห. (1)}

\prose{อาสว\-สมฺปยุตฺตญฺจ อาสว\-วิปฺปยุตฺตญฺจ ธมฺมํ ปจฺจยา อาสว\-วิปฺปยุตฺโต ธมฺโม อุปฺปชฺชติ นเหตุปจฺจยา. วิจิกิจฺฉา\-สหคโต อุทฺธจฺจ\-สหคเต ขนฺเธ จ วตฺถุญฺจ ปจฺจยา วิจิกิจฺฉา\-สหคโต อุทฺธจฺจ\-สหคโต โมโห. (1) (สํขิตฺตํ.)}

\csromanmarkh{2. ปจฺจย\-ปจฺจนีย}\csromanmarkhii{2. สงฺขฺยาวาร}\csromanheadernomark{2}{2. ปจฺจย\-ปจฺจนีย 2. สงฺขฺยาวาร}
\vspace{\tipitakaparskip}
\csromancenter{\textbf{นเหตุ}ยา ตีณิ, นอารมฺมเณ ตีณิ, นอธิปติยา นว, นอนนฺตเร ตีณิ, นสมนนฺตเร ตีณิ, นอญฺญมญฺเญ ตีณิ, นฺเอาปนิสฺสเย ตีณิ, นปุเรชาเต สตฺต, นปจฺฉาชาเต นว, นอาเสวเน นว, นกมฺเม}
\prose{\csromanfolio{183}จตฺตาริ, นวิปาเก นว, นอาหาเร เอกํ, ไนนฺทฺริเย เอกํ, นฌาเน เอกํ, นมคฺเค เอกํ, นสมฺปยุตฺเต ตีณิ, นวิปฺปยุตฺเต ฉ, โนนตฺถิยา ตีณิ, โนวิคเต ตีณิ. (เอวํ อิตเรปิ เทฺว คณนา กาตพฺพา.)}

\vspace{\tipitakaparskip}
\csromancenter{(นิสฺสยวาโรปิ ปจฺจยวาร\-สทิโส.)}
\csromanmarkh{16. อาสว\-สมฺปยุตฺตทุก}\csromanmarkhii{5. สํสฏฺฐวาร}\csromanheadernomark{2}{16. อาสว\-สมฺปยุตฺตทุก 5. สํสฏฺฐวาร}
\proseitem{73}{อาสว\-สมฺปยุตฺตํ ธมฺมํ สํสฏฺโฐ อาสว\-สมฺปยุตฺโต ธมฺโม อุปฺปชฺชติ เหตุปจฺจยา. (สํขิตฺตํ.)}

\prose{เหตุยา ฉ, อารมฺมเณ ฉ, อธิปติยา ฉ, (สพฺพตฺถ ฉ.) วิปาเก เอกํ ฯเปฯ อวิคเต ฉ.}

\prose{อาสว\-สมฺปยุตฺตํ ธมฺมํ สํสฏฺโฐ อาสว\-วิปฺปยุตฺโต ธมฺโม อุปฺปชฺชติ นเหตุปจฺจยา. วิจิกิจฺฉา\-สหคเต อุทฺธจฺจ\-สหคเต ขนฺเธ สํสฏฺโฐ วิจิกิจฺฉา\-สหคโต อุทฺธจฺจ\-สหคโต โมโห. (1)}

\prose{อาสว\-วิปฺปยุตฺตํ ธมฺมํ สํสฏฺโฐ อาสว\-วิปฺปยุตฺโต ธมฺโม อุปฺปชฺชติ นเหตุปจฺจยา ฯเปฯ.}

\prose{นเหตุยา เทฺว, นอธิปติยา ฉ, นปุเรชาเต ฉ, นปจฺฉาชาเต ฉ, นอาเสวเน ฉ, นกมฺเม จตฺตาริ, นวิปาเก ฉ, นฌาเน เอกํ, นมคฺเค เอกํ, นวิปฺปยุตฺเต ฉ. (เอวํ อิตเรปิ เทฺว คณนา กาตพฺพา.)}

\vspace{\tipitakaparskip}
\csromancenter{(สมฺปยุตฺตวาโรปิ สํสฏฺฐวาร\-สทิโส.)}
\csromanmarkh{16. อาสว\-สมฺปยุตฺตทุก}\csromanmarkhii{7. ปญฺหาวาร}\csromanheadernomark{2}{16. อาสว\-สมฺปยุตฺตทุก 7. ปญฺหาวาร}
\csromanmarkh{1. ปจฺจยานุโลม}\csromanmarkhii{1. วิภงฺควาร}\csromanheadernomark{2}{1. \textbf{ปจฺจยานุโลม 1. วิภงฺควาร}}
\vspace{\tipitakaparskip}
\csromancenter{\textbf{เหตุ}}
\proseitem{74}{อาสว\-สมฺปยุตฺโต ธมฺโม อาสว\-สมฺปยุตฺตสฺส ธมฺมสฺส เหตุปจฺจเยน ปจฺจโย. อาสว\-สมฺปยุตฺตา เหตู สมฺปยุตฺตกานํ ขนฺธานํ เหตุปจฺจเยน ปจฺจโย. (1)}

\prose{\csromanfolio{184}อาสว\-สมฺปยุตฺโต ธมฺโม อาสว\-วิปฺปยุตฺตสฺส ธมฺมสฺส เหตุปจฺจเยน ปจฺจโย. อาสว\-สมฺปยุตฺตา เหตู จิตฺต\-สมุฏฺฐานานํ รูปานํ เหตุปจฺจเยน ปจฺจโย. โทโส โมหสฺส จิตฺตสมุฏฺฐานานญฺจ รูปานํ เหตุปจฺจเยน ปจฺจโย. (2)}

\prose{อาสว\-สมฺปยุตฺโต ธมฺโม อาสว\-สมฺปยุตฺตสฺส จ อาสว\-วิปฺปยุตฺตสฺส จ ธมฺมสฺส เหตุปจฺจเยน ปจฺจโย. อาสว\-สมฺปยุตฺตา เหตู สมฺปยุตฺตกานํ ขนฺธานํ จิตฺตสมุฏฺฐานานญฺจ รูปานํ เหตุปจฺจเยน ปจฺจโย. โทโส สมฺปยุตฺตกานํ ขนฺธานํ โมหสฺส จิตฺตสมุฏฺฐานานญฺจ รูปานํ เหตุปจฺจเยน ปจฺจโย. (3)}

\proseitem{75}{อาสว\-วิปฺปยุตฺโต ธมฺโม อาสว\-วิปฺปยุตฺตสฺส ธมฺมสฺส เหตุปจฺจเยน ปจฺจโย. อาสว\-วิปฺปยุตฺตา เหตู สมฺปยุตฺตกานํ ขนฺธานํ จิตฺตสมุฏฺฐานานญฺจ รูปานํ เหตุปจฺจเยน ปจฺจโย. โทมนสฺส\-สหคโต วิจิกิจฺฉา\-สหคโต อุทฺธจฺจ\-สหคโต โมโห จิตฺต\-สมุฏฺฐานานํ รูปานํ เหตุปจฺจเยน ปจฺจโย. ปฏิสนฺธิ\-กฺขเณ ฯเปฯ. (1)}

\prose{อาสว\-วิปฺปยุตฺโต ธมฺโม อาสว\-สมฺปยุตฺตสฺส ธมฺมสฺส เหตุปจฺจเยน ปจฺจโย. โทมนสฺส\-สหคโต วิจิกิจฺฉา\-สหคโต อุทฺธจฺจ\-สหคโต โมโห สมฺปยุตฺตกานํ ขนฺธานํ เหตุปจฺจเยน ปจฺจโย. (2)}

\prose{อาสว\-วิปฺปยุตฺโต ธมฺโม อาสว\-สมฺปยุตฺตสฺส จ อาสว\-วิปฺปยุตฺตสฺส จ ธมฺมสฺส เหตุปจฺจเยน ปจฺจโย. โทมนสฺส\-สหคโต วิจิกิจฺฉา\-สหคโต อุทฺธจฺจ\-สหคโต โมโห สมฺปยุตฺตกานํ ขนฺธานํ จิตฺตสมุฏฺฐานานญฺจ รูปานํ เหตุปจฺจเยน ปจฺจโย. (3)}

\proseitem{76}{อาสว\-สมฺปยุตฺโต จ อาสว\-วิปฺปยุตฺโต จ ธมฺมา อาสว\-สมฺปยุตฺตสฺส ธมฺมสฺส เหตุปจฺจเยน ปจฺจโย. โทโส จ โมโห จ อาสว\-สมฺปยุตฺตกานํ ขนฺธานํ เหตุปจฺจเยน ปจฺจโย. (1)}

\prose{อาสว\-สมฺปยุตฺโต จ อาสว\-วิปฺปยุตฺโต จ ธมฺมา อาสว\-วิปฺปยุตฺตสฺส ธมฺมสฺส เหตุปจฺจเยน ปจฺจโย. โทโส จ โมโห จ จิตฺต\-สมุฏฺฐานานํ รูปานํ เหตุปจฺจเยน ปจฺจโย. (2)}

\prose{อาสว\-สมฺปยุตฺโต จ อาสว\-วิปฺปยุตฺโต จ ธมฺมา อาสว\-สมฺปยุตฺตสฺส จ อาสว\-วิปฺปยุตฺตสฺส จ ธมฺมสฺส เหตุปจฺจเยน ปจฺจโย. โทโส จ \csromanfolio{185}โมโห จ สมฺปยุตฺตกานํ ขนฺธานํ จิตฺตสมุฏฺฐานานญฺจ รูปานํ เหตุปจฺจเยน ปจฺจโย. (3)}

\proseitem{77}{อาสว\-สมฺปยุตฺโต ธมฺโม อาสว\-สมฺปยุตฺตสฺส ธมฺมสฺส อารมฺมณ\-ปจฺจเยน ปจฺจโย. อาสว\-สมฺปยุตฺเต ขนฺเธ อารพฺภ อาสว\-สมฺปยุตฺตกา ขนฺธา อุปฺปชฺชนฺติ. (1)}

\prose{อาสว\-สมฺปยุตฺโต ธมฺโม อาสว\-วิปฺปยุตฺตสฺส ธมฺมสฺส อารมฺมณ\-ปจฺจเยน ปจฺจโย. อาสว\-สมฺปยุตฺเต ขนฺเธ อารพฺภ อาสว\-วิปฺปยุตฺตา ขนฺธา จ โมโห จ อุปฺปชฺชนฺติ. (2)}

\prose{อาสว\-สมฺปยุตฺโต ธมฺโม อาสว\-สมฺปยุตฺตสฺส จ อาสว\-วิปฺปยุตฺตสฺส จ ธมฺมสฺส อารมฺมณ\-ปจฺจเยน ปจฺจโย. อาสว\-สมฺปยุตฺเต ขนฺเธ อารพฺภ โทมนสฺส\-สหคตา วิจิกิจฺฉา\-สหคตา อุทฺธจฺจ\-สหคตา ขนฺธา จ โมโห จ อุปฺปชฺชนฺติ. (3)}

\proseitem{78}{อาสว\-วิปฺปยุตฺโต ธมฺโม อาสว\-วิปฺปยุตฺตสฺส ธมฺมสฺส อารมฺมณ\-ปจฺจเยน ปจฺจโย. ทานํ ทตฺวา, สีลํ ฯเปฯ อุโปสถกมฺมํ กตฺวา ตํ ปจฺจเวกฺขติ. ปุพฺเพ สุจิณฺณานิ ฯเปฯ. ฌานา ฯเปฯ. อริยา มคฺคํ ฯเปฯ. ผลํ ฯเปฯ. นิพฺพานํ ฯเปฯ. นิพฺพานํ โคตฺรภุสฺส, โวทานสฺส, มคฺคสฺส, ผลสฺส, อาวชฺชนาย อารมฺมณ\-ปจฺจเยน ปจฺจโย. อริยา อาสว\-วิปฺปยุตฺเต ปหีเน กิเลเส ปจฺจเวกฺขนฺติ. วิกฺขมฺภิเต ฯเปฯ. ปุพฺเพ ฯเปฯ. จกฺขุํ ฯเปฯ วตฺถุํ อาสว\-วิปฺปยุตฺเต ขนฺเธ อนิจฺจโต ทุกฺขโต อนตฺตโต วิปสฺสติ. (อิธ อสฺสาทนา นตฺถิ.) ทิพฺเพน จกฺขุนา รูปํ ปสฺสติ. ทิพฺพาย โสตธาตุยา สทฺทํ สุณาติ. เจโตปริย\-ญาเณน อาสว\-วิปฺปยุตฺต\-จิตฺต\-สมงฺคิสฺส จิตฺตํ ชานาติ. อากาสา\-นญฺจา\-ยตนํ วิญฺญาณญฺจา\-ยตนสฺส ฯเปฯ. อากิญฺจญฺญา\-ยตนํ เนว\-สญฺญา\-นา\-สญฺญา\-ยตนสฺส ฯเปฯ. รูปายตนํ จกฺขุ\-วิญฺญาณสฺส ฯเปฯ. โผฏฺฐพฺพา\-ยตนํ กายวิญฺญาณสฺส ฯเปฯ. อาสว\-วิปฺปยุตฺตา ขนฺธา อิทฺธิวิธ\-ญาณสฺส, เจโตปริย\-ญาณสฺส, ปุพฺเพ\-นิวาสา\-นุสฺสติ\-ญาณสฺส, ยถากมฺมูปคญาณสฺส, อนาคตํสญาณสฺส, อาวชฺชนาย โมหสฺส อารมฺมณ\-ปจฺจเยน ปจฺจโย. (1) \csromanfolio{186}อาสว\-วิปฺปยุตฺโต ธมฺโม อาสว\-สมฺปยุตฺตสฺส ธมฺมสฺส อารมฺมณ\-ปจฺจเยน ปจฺจโย. ทานํ ทตฺวา, สีลํ ฯเปฯ อุโปสถกมฺมํ กตฺวา ตํ อสฺสาเทติ อภินนฺทติ, ตํ อารพฺภ ราโค อุปฺปชฺชติ. ทิฏฺฐิ ฯเปฯ วิจิกิจฺฉา ฯเปฯ อุทฺธจฺจํ ฯเปฯ โทมนสฺสํ อุปฺปชฺชติ. ปุพฺเพ สุจิณฺณานิ ฯเปฯ. ฌานา วุฏฺฐหิตฺวา ฌานํ ฯเปฯ. จกฺขุํ ฯเปฯ วตฺถุํ อาสว\-วิปฺปยุตฺเต ขนฺเธ อสฺสาเทติ อภินนฺทติ, ตํ อารพฺภ ราโค ฯเปฯ. ทิฏฺฐิ. โทมนสฺสํ. วิจิกิจฺฉา. อุทฺธจฺจํ อุปฺปชฺชติ. (2)}

\prose{อาสว\-วิปฺปยุตฺโต ธมฺโม อาสว\-สมฺปยุตฺตสฺส จ อาสว\-วิปฺปยุตฺตสฺส จ ธมฺมสฺส อารมฺมณ\-ปจฺจเยน ปจฺจโย. จกฺขุํ ฯเปฯ วตฺถุํ อาสว\-วิปฺปยุตฺเต ขนฺเธ อารพฺภ โทมนสฺส\-สหคตา วิจิกิจฺฉา\-สหคตา อุทฺธจฺจ\-สหคตา ขนฺธา จ โมโห จ อุปฺปชฺชนฺติ. (3)}

\proseitem{79}{อาสว\-สมฺปยุตฺโต จ อาสว\-วิปฺปยุตฺโต จ ธมฺมา อาสว\-สมฺปยุตฺตสฺส ธมฺมสฺส อารมฺมณ\-ปจฺจเยน ปจฺจโย. โทมนสฺส\-สหคเต วิจิกิจฺฉา\-สหคเต อุทฺธจฺจ\-สหคเต ขนฺเธ จ โมหญฺจ อารพฺภ อาสว\-สมฺปยุตฺตา ขนฺธา อุปฺปชฺชนฺติ. (1)}

\prose{อาสว\-สมฺปยุตฺโต จ อาสว\-วิปฺปยุตฺโต จ ธมฺมา อาสว\-วิปฺปยุตฺตสฺส ธมฺมสฺส อารมฺมณ\-ปจฺจเยน ปจฺจโย. โทมนสฺส\-สหคเต วิจิกิจฺฉา\-สหคเต อุทฺธจฺจ\-สหคเต ขนฺเธ จ โมหญฺจ อารพฺภ อาสว\-วิปฺปยุตฺตา ขนฺธา จ โมโห จ อุปฺปชฺชนฺติ. (2)}

\prose{อาสว\-สมฺปยุตฺโต จ อาสว\-วิปฺปยุตฺโต จ ธมฺมา อาสว\-สมฺปยุตฺตสฺส จ อาสว\-วิปฺปยุตฺตสฺส จ ธมฺมสฺส อารมฺมณ\-ปจฺจเยน ปจฺจโย. โทมนสฺส\-สหคเต วิจิกิจฺฉา\-สหคเต อุทฺธจฺจ\-สหคเต ขนฺเธ จ โมหญฺจ อารพฺภ โทมนสฺส\-สหคตา วิจิกิจฺฉา\-สหคตา อุทฺธจฺจ\-สหคตา ขนฺธา จ โมโห จ อุปฺปชฺชนฺติ. (3)}

\proseitem{80}{อาสว\-สมฺปยุตฺโต ธมฺโม อาสว\-สมฺปยุตฺตสฺส ธมฺมสฺส อธิปติ\-ปจฺจเยน ปจฺจโย. อารมฺมณา\-ธิปติ, สหชาตา\-ธิปติ.  อารมฺมณา\-ธิปติ\csromandash{}อาสว\-สมฺปยุตฺเต ขนฺเธ ครุํ กตฺวา อาสว\-สมฺปยุตฺตกา ขนฺธา อุปฺปชฺชนฺติ. สหชาตา\-ธิปติ\csromandash{} อาสว\-สมฺปยุตฺตา\-ธิปติ สมฺปยุตฺตกานํ ขนฺธานํ อธิปติ\-ปจฺจเยน ปจฺจโย. (1)}

\prose{\csromanfolio{187}อาสว\-สมฺปยุตฺโต ธมฺโม อาสว\-วิปฺปยุตฺตสฺส ธมฺมสฺส อธิปติ\-ปจฺจเยน ปจฺจโย. สหชาตา\-ธิปติ\csromandash{}อาสว\-สมฺปยุตฺตา\-ธิปติ จิตฺต\-สมุฏฺฐานานํ รูปานํ อธิปติ\-ปจฺจเยน ปจฺจโย. โทมนสฺส\-สหคตาธิปติ โมหสฺส จิตฺตสมุฏฺฐานานญฺจ รูปานํ อธิปติ\-ปจฺจเยน ปจฺจโย. (2)}

\prose{อาสว\-สมฺปยุตฺโต ธมฺโม อาสว\-สมฺปยุตฺตสฺส จ อาสว\-วิปฺปยุตฺตสฺส จ ธมฺมสฺส อธิปติ\-ปจฺจเยน ปจฺจโย. สหชาตา\-ธิปติ\csromandash{} อาสว\-สมฺปยุตฺตา\-ธิปติ สมฺปยุตฺตกานํ ขนฺธานํ จิตฺตสมุฏฺฐานานญฺจ รูปานํ อธิปติ\-ปจฺจเยน ปจฺจโย. โทมนสฺส\-สหคตาธิปติ สมฺปยุตฺตกานํ ขนฺธานํ โมหสฺส จ จิตฺตสมุฏฺฐานานญฺจ รูปานํ อธิปติ\-ปจฺจเยน ปจฺจโย. (3)}

\proseitem{81}{อาสว\-วิปฺปยุตฺโต ธมฺโม อาสว\-วิปฺปยุตฺตสฺส ธมฺมสฺส อธิปติ\-ปจฺจเยน ปจฺจโย. อารมฺมณา\-ธิปติ, สหชาตา\-ธิปติ.  อารมฺมณา\-ธิปติ\csromandash{}ทานํ ทตฺวา, สีลํ ฯเปฯ อุโปสถกมฺมํ กตฺวา ตํ ครุํ กตฺวา ปจฺจเวกฺขติ. ปุพฺเพ สุจิณฺณานิ ฯเปฯ. ฌานา ฯเปฯ. อริยา มคฺคา วุฏฺฐหิตฺวา มคฺคํ ฯเปฯ. ผลํ ฯเปฯ. นิพฺพานํ ครุํ กตฺวา ปจฺจเวกฺขนฺติ. นิพฺพานํ โคตฺรภุสฺส, โวทานสฺส, มคฺคสฺส, ผลสฺส อธิปติ\-ปจฺจเยน ปจฺจโย. สหชาตา\-ธิปติ\csromandash{}อาสว\-วิปฺปยุตฺตา\-ธิปติ สมฺปยุตฺตกานํ ขนฺธานํ จิตฺตสมุฏฺฐานานญฺจ รูปานํ อธิปติ\-ปจฺจเยน ปจฺจโย. (1)}

\prose{อาสว\-วิปฺปยุตฺโต ธมฺโม อาสว\-สมฺปยุตฺตสฺส ธมฺมสฺส อธิปติ\-ปจฺจเยน ปจฺจโย. อารมฺมณา\-ธิปติ\csromandash{}ทานํ ทตฺวา, สีลํ ฯเปฯ อุโปสถกมฺมํ กตฺวา ตํ ครุํ กตฺวา อสฺสาเทติ อภินนฺทติ, ตํ ครุํ กตฺวา ราโค อุปฺปชฺชติ. ทิฏฺฐิ อุปฺปชฺชติ. ปุพฺเพ สุจิณฺณานิ ฯเปฯ. ฌานา ฯเปฯ. จกฺขุํ ฯเปฯ วตฺถุํ อาสว\-วิปฺปยุตฺเต ขนฺเธ ครุํ กตฺวา อสฺสาเทติ อภินนฺทติ, ตํ ครุํ กตฺวา ราโค อุปฺปชฺชติ. ทิฏฺฐิ อุปฺปชฺชติ. (2)}

\proseitem{82}{อาสว\-สมฺปยุตฺโต ธมฺโม อาสว\-สมฺปยุตฺตสฺส ธมฺมสฺส อนนฺตร\-ปจฺจเยน ปจฺจโย. ปุริมา ปุริมา อาสว\-สมฺปยุตฺตกา ขนฺธา ปจฺฉิมานํ ปจฺฉิมานํ อาสว\-สมฺปยุตฺตกานํ ขนฺธานํ อนนฺตร\-ปจฺจเยน ปจฺจโย. (1)}

\prose{อาสว\-สมฺปยุตฺโต ธมฺโม อาสว\-วิปฺปยุตฺตสฺส ธมฺมสฺส อนนฺตร\-ปจฺจเยน ปจฺจโย. ปุริมา ปุริมา โทมนสฺส\-สหคตา วิจิกิจฺฉา\-สหคตา \csromanfolio{188}อุทฺธจฺจ\-สหคตา ขนฺธา ปจฺฉิมสฺส ปจฺฉิมสฺส โทมนสฺส\-สหคตสฺส วิจิกิจฺฉา\-สหคตสฺส อุทฺธจฺจ\-สหคตสฺส โมหสฺส อนนฺตร\-ปจฺจเยน ปจฺจโย. อาสว\-สมฺปยุตฺตา ขนฺธา วุฏฺฐานสฺส อนนฺตร\-ปจฺจเยน ปจฺจโย. (2)}

\prose{อาสว\-สมฺปยุตฺโต ธมฺโม อาสว\-สมฺปยุตฺตสฺส จ อาสว\-วิปฺปยุตฺตสฺส จ ธมฺมสฺส อนนฺตร\-ปจฺจเยน ปจฺจโย. ปุริมา ปุริมา โทมนสฺส\-สหคตา วิจิกิจฺฉา\-สหคตา อุทฺธจฺจ\-สหคตา ขนฺธา ปจฺฉิมานํ ปจฺฉิมานํ โทมนสฺส\-สหคตานํ วิจิกิจฺฉา\-สหคตานํ อุทฺธจฺจ\-สหคตานํ ขนฺธานํ โมหสฺส จ อนนฺตร\-ปจฺจเยน ปจฺจโย. (3)}

\proseitem{83}{อาสว\-วิปฺปยุตฺโต ธมฺโม อาสว\-วิปฺปยุตฺตสฺส ธมฺมสฺส อนนฺตร\-ปจฺจเยน ปจฺจโย. ปุริโม ปุริโม โทมนสฺส\-สหคโต วิจิกิจฺฉา\-สหคโต อุทฺธจฺจ\-สหคโต โมโห ปจฺฉิมสฺส ปจฺฉิมสฺส โทมนสฺส\-สหคตสฺส วิจิกิจฺฉา\-สหคตสฺส อุทฺธจฺจ\-สหคตสฺส โมหสฺส อนนฺตร\-ปจฺจเยน ปจฺจโย. ปุริมา ปุริมา อาสว\-วิปฺปยุตฺตา ขนฺธา ปจฺฉิมานํ ปจฺฉิมานํ อาสว\-วิปฺปยุตฺตานํ ขนฺธานํ อนนฺตร\-ปจฺจเยน ปจฺจโย. อนุโลมํ โคตฺรภุสฺส ฯเปฯ ผลสมาปตฺติยา อนนฺตร\-ปจฺจเยน ปจฺจโย. (1)}

\prose{อาสว\-วิปฺปยุตฺโต ธมฺโม อาสว\-สมฺปยุตฺตสฺส ธมฺมสฺส อนนฺตร\-ปจฺจเยน ปจฺจโย. ปุริโม ปุริโม โทมนสฺส\-สหคโต วิจิกิจฺฉา\-สหคโต อุทฺธจฺจ\-สหคโต โมโห ปจฺฉิมานํ ปจฺฉิมานํ โทมนสฺส\-สหคตานํ วิจิกิจฺฉา\-สหคตานํ อุทฺธจฺจ\-สหคตานํ ขนฺธานํ อนนฺตร\-ปจฺจเยน ปจฺจโย. อาวชฺชนา อาสว\-สมฺปยุตฺตกานํ ขนฺธานํ อนนฺตร\-ปจฺจเยน ปจฺจโย. (2)}

\prose{อาสว\-วิปฺปยุตฺโต ธมฺโม อาสว\-สมฺปยุตฺตสฺส จ อาสว\-วิปฺปยุตฺตสฺส จ ธมฺมสฺส อนนฺตร\-ปจฺจเยน ปจฺจโย. ปุริโม ปุริโม โทมนสฺส\-สหคโต วิจิกิจฺฉา\-สหคโต อุทฺธจฺจ\-สหคโต โมโห ปจฺฉิมานํ ปจฺฉิมานํ โทมนสฺส\-สหคตานํ วิจิกิจฺฉา\-สหคตานํ อุทฺธจฺจ\-สหคตานํ ขนฺธานํ โมหสฺส จ อนนฺตร\-ปจฺจเยน ปจฺจโย. อาวชฺชนา โทมนสฺส\-สหคตานํ วิจิกิจฺฉา\-สหคตานํ อุทฺธจฺจ\-สหคตานํ ขนฺธานํ โมหสฺส จ อนนฺตร\-ปจฺจเยน ปจฺจโย. (3)}

\proseitem{84}{\csromanfolio{189}อาสว\-สมฺปยุตฺโต จ อาสว\-วิปฺปยุตฺโต จ ธมฺมา อาสว\-สมฺปยุตฺตสฺส ธมฺมสฺส อนนฺตร\-ปจฺจเยน ปจฺจโย. ปุริมา ปุริมา โทมนสฺส\-สหคตา วิจิกิจฺฉา\-สหคตา อุทฺธจฺจ\-สหคตา ขนฺธา จ โมโห จ ปจฺฉิมานํ ปจฺฉิมานํ โทมนสฺส\-สหคตานํ วิจิกิจฺฉา\-สหคตานํ อุทฺธจฺจ\-สหคตานํ ขนฺธานํ อนนฺตร\-ปจฺจเยน ปจฺจโย. (1)}

\prose{อาสว\-สมฺปยุตฺโต จ อาสว\-วิปฺปยุตฺโต จ ธมฺมา อาสว\-วิปฺปยุตฺตสฺส ธมฺมสฺส อนนฺตร\-ปจฺจเยน ปจฺจโย. ปุริมา ปุริมา โทมนสฺส\-สหคตา วิจิกิจฺฉา\-สหคตา อุทฺธจฺจ\-สหคตา ขนฺธา จ โมโห จ ปจฺฉิมสฺส ปจฺฉิมสฺส โทมนสฺส\-สหคตสฺส วิจิกิจฺฉา\-สหคตสฺส อุทฺธจฺจ\-สหคตสฺส โมหสฺส อนนฺตร\-ปจฺจเยน ปจฺจโย. โทมนสฺส\-สหคตา วิจิกิจฺฉา\-สหคตา อุทฺธจฺจ\-สหคตา ขนฺธา จ โมโห จ วุฏฺฐานสฺส อนนฺตร\-ปจฺจเยน ปจฺจโย. (2)}

\prose{อาสว\-สมฺปยุตฺโต จ อาสว\-วิปฺปยุตฺโต จ ธมฺมา อาสว\-สมฺปยุตฺตสฺส จ อาสว\-วิปฺปยุตฺตสฺส จ ธมฺมสฺส อนนฺตร\-ปจฺจเยน ปจฺจโย. ปุริมา ปุริมา โทมนสฺส\-สหคตา วิจิกิจฺฉา\-สหคตา อุทฺธจฺจ\-สหคตา ขนฺธา จ โมโห จ ปจฺฉิมานํ ปจฺฉิมานํ โทมนสฺส\-สหคตานํ วิจิกิจฺฉา\-สหคตานํ อุทฺธจฺจ\-สหคตานํ ขนฺธานํ โมหสฺส จ อนนฺตร\-ปจฺจเยน ปจฺจโย. (3)}

\csromanheader{2}{\textbf{สมนนฺตราทิ}}
\proseitem{85}{อาสว\-สมฺปยุตฺโต ธมฺโม อาสว\-สมฺปยุตฺตสฺส ธมฺมสฺส สมนนฺตร\-ปจฺจเยน ปจฺจโย. สหชาต\-ปจฺจเยน ปจฺจโย. นว. อญฺญมญฺญ\-ปจฺจเยน ปจฺจโย. ฉ. นิสฺสย\-ปจฺจเยน ปจฺจโย. นว.}

\csromanheader{2}{\textbf{อุปนิสฺสย}}
\proseitem{86}{อาสว\-สมฺปยุตฺโต ธมฺโม อาสว\-สมฺปยุตฺตสฺส ธมฺมสฺส อุปนิสฺสย\-ปจฺจเยน ปจฺจโย. อารมฺมณู\-ปนิสฺสโย, อนนฺตรู\-ปนิสฺสโย, ปกตู\-ปนิสฺสโย ฯเปฯ. ปกตู\-ปนิสฺสโย\csromandash{}อาสว\-สมฺปยุตฺตา ขนฺธา อาสว\-สมฺปยุตฺตานํ ขนฺธานํ อุปนิสฺสย\-ปจฺจเยน ปจฺจโย. (1)}

\prose{อาสว\-สมฺปยุตฺโต ธมฺโม อาสว\-วิปฺปยุตฺตสฺส ธมฺมสฺส อุปนิสฺสย\-ปจฺจเยน ปจฺจโย. อนนฺตรู\-ปนิสฺสโย, ปกตู\-ปนิสฺสโย ฯเปฯ. \csromanfolio{190}ปกตู\-ปนิสฺสโย\csromandash{}อาสว\-สมฺปยุตฺตา ขนฺธา อาสว\-วิปฺปยุตฺตานํ ขนฺธานํ โมหสฺส จ อุปนิสฺสย\-ปจฺจเยน ปจฺจโย. (2)}

\prose{อาสว\-สมฺปยุตฺโต ธมฺโม อาสว\-สมฺปยุตฺตสฺส จ อาสว\-วิปฺปยุตฺตสฺส จ ธมฺมสฺส อุปนิสฺสย\-ปจฺจเยน ปจฺจโย. อนนฺตรู\-ปนิสฺสโย, ปกตู\-ปนิสฺสโย ฯเปฯ. ปกตู\-ปนิสฺสโย\csromandash{}อาสว\-สมฺปยุตฺตกา ขนฺธา โทมนสฺส\-สหคตานํ วิจิกิจฺฉา\-สหคตานํ อุทฺธจฺจ\-สหคตานํ ขนฺธานํ โมหสฺส จ อุปนิสฺสย\-ปจฺจเยน ปจฺจโย. (3)}

\proseitem{87}{อาสว\-วิปฺปยุตฺโต ธมฺโม อาสว\-วิปฺปยุตฺตสฺส ธมฺมสฺส อุปนิสฺสย\-ปจฺจเยน ปจฺจโย. อารมฺมณู\-ปนิสฺสโย, อนนฺตรู\-ปนิสฺสโย, ปกตู\-ปนิสฺสโย ฯเปฯ. ปกตู\-ปนิสฺสโย\csromandash{}สทฺธํ อุปนิสฺสาย ทานํ เทติ ฯเปฯ สีลํ ฯเปฯ ปญฺญํ. กายิกํ สุขํ ฯเปฯ. เสนาสนํ. โมหํ อุปนิสฺสาย ทานํ เทติ ฯเปฯ สมาปตฺติํ อุปฺปาเทติ, สทฺธา ฯเปฯ ปญฺญา. กายิกํ สุขํ. กายิกํ ทุกฺขํ. โมโห จ สทฺธาย ฯเปฯ โมหสฺส จ อุปนิสฺสย\-ปจฺจเยน ปจฺจโย. (1)}

\prose{อาสว\-วิปฺปยุตฺโต ธมฺโม อาสว\-สมฺปยุตฺตสฺส ธมฺมสฺส อุปนิสฺสย\-ปจฺจเยน ปจฺจโย. อารมฺมณู\-ปนิสฺสโย, อนนฺตรู\-ปนิสฺสโย, ปกตู\-ปนิสฺสโย ฯเปฯ. ปกตู\-ปนิสฺสโย\csromandash{}สทฺธํ อุปนิสฺสาย มานํ ชปฺเปติ, ทิฏฺฐํ คณฺหาติ, สีลํ ฯเปฯ ปญฺญํ. กายิกํ สุขํ. กายิกํ ทุกฺขํ. อุตุํ. โภชนํ. เสนาสนํ. โมหํ อุปนิสฺสาย ปาณํ หนติ ฯเปฯ สํฆํ ภินฺทติ. สทฺธา ฯเปฯ โมโห จ ราคสฺส ฯเปฯ ปตฺถนาย อุปนิสฺสย\-ปจฺจเยน ปจฺจโย. (2)}

\prose{อาสว\-วิปฺปยุตฺโต ธมฺโม อาสว\-สมฺปยุตฺตสฺส จ อาสว\-วิปฺปยุตฺตสฺส จ ธมฺมสฺส อุปนิสฺสย\-ปจฺจเยน ปจฺจยฺโอ. อนนฺตรู\-ปนิสฺสโย, ปกตู\-ปนิสฺสโย ฯเปฯ. ปกตู\-ปนิสฺสโย\csromandash{}สทฺธา. สีลํ ฯเปฯ โมโห โทมนสฺส\-สหคตานํ วิจิกิจฺฉา\-สหคตานํ อุทฺธจฺจ\-สหคตานํ ขนฺธานํ โมหสฺส จ อุปนิสฺสย\-ปจฺจเยน ปจฺจโย. (3)}

\proseitem{88}{อาสว\-สมฺปยุตฺโต จ อาสว\-วิปฺปยุตฺโต จ ธมฺมา อาสว\-สมฺปยุตฺตสฺส ธมฺมสฺส อุปนิสฺสย\-ปจฺจเยน ปจฺจโย. อนนฺตรู\-ปนิสฺสโย, \csromanfolio{191}ปกตู\-ปนิสฺสโย ฯเปฯ. ปกตู\-ปนิสฺสโย\csromandash{}โทมนสฺส\-สหคตา วิจิกิจฺฉา\-สหคตา อุทฺธจฺจ\-สหคตา ขนฺธา จ โมโห จ อาสว\-สมฺปยุตฺตกานํ ขนฺธานํ อุปนิสฺสย\-ปจฺจเยน ปจฺจโย. (ปุจฺฉิตพฺพํ มูลํ.) โทมนสฺส\-สหคตา วิจิกิจฺฉา\-สหคตา อุทฺธจฺจ\-สหคตา ขนฺธา จ โมโห จ อาสว\-สมฺปยุตฺตกานํ ขนฺธานํ อุปนิสฺสย\-ปจฺจเยน ปจฺจโย. (1)}

\prose{อาสว\-สมฺปยุตฺโต จ อาสว\-วิปฺปยุตฺโต จ ธมฺมา อาสว\-สมฺปยุตฺตสฺส จ อาสววิปฺปยุสฺส จ ธมฺมสฺส อุปนิสฺสย\-ปจฺจเยน ปจฺจโย. อนนฺตรู\-ปนิสฺสโย, ปกตู\-ปนิสฺสโย ฯเปฯ. ปกตู\-ปนิสฺสโย\csromandash{} โทมนสฺส\-สหคตา วิจิกิจฺฉา\-สหคตา อุทฺธจฺจ\-สหคตา ขนฺธา จ โมโห จ โทมนสฺส\-สหคตานํ วิจิกิจฺฉา\-สหคตานํ อุทฺธจฺจ\-สหคตานํ ขนฺธานํ โมหสฺส จ อุปนิสฺสย\-ปจฺจเยน ปจฺจโย. (2)}

\proseitem{89}{อาสว\-วิปฺปยุตฺโต ธมฺโม อาสว\-วิปฺปยุตฺตสฺส ธมฺมสฺส ปุเรชาต\-ปจฺจเยน ปจฺจโย. อารมฺมณ\-ปุเรชาตํ, วตฺถุ\-ปุเรชาตํ.  อารมฺมณ\-ปุเรชาตํ\csromandash{}จกฺขุํ ฯเปฯ วตฺถุํ อนิจฺจโต ทุกฺขโต อนตฺตโต วิปสฺสติ. ทิพฺเพน จกฺขุนา รูปํ ปสฺสติ. ทิพฺพาย โสตธาตุยา สทฺทํ สุณาติ. รูปายตนํ จกฺขุ\-วิญฺญาณสฺส ฯเปฯ. โผฏฺฐพฺพา\-ยตนํ กายวิญฺญาณสฺส ปุเรชาต\-ปจฺจเยน ปจฺจโย. วตฺถุ\-ปุเรชาตํ\csromandash{} จกฺขายตนํ จกฺขุ\-วิญฺญาณสฺส ฯเปฯ. กายายตนํ กายวิญฺญาณสฺส ฯเปฯ. วตฺถุ อาสว\-วิปฺปยุตฺตานํ ขนฺธานํ โมหสฺส จ ปุเรชาต\-ปจฺจเยน ปจฺจโย. (1)}

\prose{อาสว\-วิปฺปยุตฺโต ธมฺโม อาสว\-สมฺปยุตฺตสฺส ธมฺมสฺส ปุเรชาต\-ปจฺจเยน ปจฺจโย. อารมฺมณ\-ปุเรชาตํ, วตฺถุ\-ปุเรชาตํ.  อารมฺมณ\-ปุเรชาตํ\csromandash{}จกฺขุํ ฯเปฯ วตฺถุํ อสฺสาเทติ อภินนฺทติ, ตํ อารพฺภ ราโค อุปฺปชฺชติ ฯเปฯ โทมนสฺสํ อุปฺปชฺชติ. วตฺถุ\-ปุเรชาตํ\csromandash{}วตฺถุ อาสว\-สมฺปยุตฺตกานํ ขนฺธานํ ปุเรชาต\-ปจฺจเยน ปจฺจโย. (2)}

\prose{อาสว\-วิปฺปยุตฺโต ธมฺโม อาสว\-สมฺปยุตฺตสฺส จ อาสว\-วิปฺปยุตฺตสฺส จ ธมฺมสฺส ปุเรชาต\-ปจฺจเยน ปจฺจโย. อารมฺมณ\-ปุเรชาตํ, วตฺถุ\-ปุเรชาตํ.  อารมฺมณ\-ปุเรชาตํ\csromandash{}จกฺขุํ ฯเปฯ วตฺถุํ อารพฺภ โทมนสฺส\-สหคตา วิจิกิจฺฉา\-สหคตา อุทฺธจฺจ\-สหคตา ขนฺธา จ โมโห จ \csromanfolio{192}อุปฺปชฺชนฺติ. วตฺถุ\-ปุเรชาตํ\csromandash{}วตฺถุ โทมนสฺส\-สหคตานํ วิจิกิจฺฉา\-สหคตานํ อุทฺธจฺจ\-สหคตานํ ขนฺธานํ โมหสฺส จ ปุเรชาต\-ปจฺจเยน ปจฺจโย. (3)}

\csromanheader{2}{\textbf{ปจฺฉาชาต}}
\proseitem{90}{อาสว\-สมฺปยุตฺโต ธมฺโม อาสว\-วิปฺปยุตฺตสฺส ธมฺมสฺส ปจฺฉาชาต\-ปจฺจเยน ปจฺจโย. ปจฺฉาชาตา อาสว\-สมฺปยุตฺตกา ขนฺธา ปุเรชาตสฺส อิมสฺส กายสฺส ปจฺฉาชาต\-ปจฺจเยน ปจฺจโย. (1)}

\prose{อาสว\-วิปฺปยุตฺโต ธมฺโม อาสว\-วิปฺปยุตฺตสฺส ธมฺมสฺส ปจฺฉาชาต\-ปจฺจเยน ปจฺจโย. ปจฺฉาชาตา อาสว\-วิปฺปยุตฺโต ขนฺธา จ โมโห จ ปุเรชาตสฺส อิมสฺส กายสฺส ปจฺฉาชาต\-ปจฺจเยน ปจฺจโย. (1)}

\prose{อาสว\-สมฺปยุตฺโต จ อาสว\-วิปฺปยุตฺโต จ ธมฺมา อาสว\-วิปฺปยุตฺตสฺส ธมฺมสฺส ปจฺฉาชาต\-ปจฺจเยน ปจฺจโย. ปจฺฉาชาตา โทมนสฺส\-สหคตา วิจิกิจฺฉา\-สหคตา อุทฺธจฺจ\-สหคตา ขนฺธา จ โมโห จ ปุเรชาตสฺส อิมสฺส กายสฺส ปจฺฉาชาต\-ปจฺจเยน ปจฺจโย. (1)}

\csromanheader{2}{\textbf{อาเสวน}}
\proseitem{91}{อาสว\-สมฺปยุตฺโต ธมฺโม อาสว\-สมฺปยุตฺตสฺส ธมฺมสฺส อาเสวน\-ปจฺจเยน ปจฺจโย. นว. (อาวชฺชนาปิ วุฏฺฐานมฺปิ นตฺถิ.)}

\proseitem{92}{อาสว\-สมฺปยุตฺโต ธมฺโม อาสว\-สมฺปยุตฺตสฺส ธมฺมสฺส กมฺมปจฺจเยน ปจฺจโย. อาสว\-สมฺปยุตฺตกา เจตนา สมฺปยุตฺตกานํ ขนฺธานํ กมฺมปจฺจเยน ปจฺจโย. (1)}

\prose{อาสว\-สมฺปยุตฺโต ธมฺโม อาสว\-วิปฺปยุตฺตสฺส ธมฺมสฺส กมฺมปจฺจเยน ปจฺจโย. สหชาตา, นานากฺขณิกา.  สหชาตา อาสว\-สมฺปยุตฺตา เจตนา จิตฺต\-สมุฏฺฐานานํ รูปานํ กมฺมปจฺจเยน ปจฺจโย. โทมนสฺส\-สหคตา วิจิกิจฺฉา\-สหคตา อุทฺธจฺจ\-สหคตา เจตนา โมหสฺส จิตฺตสมุฏฺฐานานญฺจ รูปานํ กมฺมปจฺจเยน ปจฺจโย. นานากฺขณิกา อาสว\-สมฺปยุตฺตกา เจตนา วิปากานํ ขนฺธานํ กฏตฺตา จ รูปานํ กมฺมปจฺจเยน ปจฺจโย. (2)}

\prose{\csromanfolio{193}อาสว\-สมฺปยุตฺโต ธมฺโม อาสว\-สมฺปยุตฺตสฺส จ อาสว\-วิปฺปยุตฺตสฺส จ ธมฺมสฺส กมฺมปจฺจเยน ปจฺจโย. อาสว\-สมฺปยุตฺตา เจตนา สมฺปยุตฺตกานํ ขนฺธานํ จิตฺตสมุฏฺฐานานญฺจ รูปานํ กมฺมปจฺจเยน ปจฺจโย. โทมนสฺส\-สหคตา วิจิกิจฺฉา\-สหคตา อุทฺธจฺจ\-สหคตา เจตนา สมฺปยุตฺตกานํ ขนฺธานํ โมหสฺส จ จิตฺตสมุฏฺฐานานญฺจ รูปานํ กมฺมปจฺจเยน ปจฺจโย. (3)}

\prose{อาสว\-วิปฺปยุตฺโต ธมฺโม อาสว\-วิปฺปยุตฺตสฺส ธมฺมสฺส กมฺมปจฺจเยน ปจฺจโย. สหชาตา, นานากฺขณิกา.  สหชาตา อาสว\-วิปฺปยุตฺตา เจตนา สมฺปยุตฺตกานํ ขนฺธานํ จิตฺตสมุฏฺฐานานญฺจ รูปานํ กมฺมปจฺจเยน ปจฺจโย.  ปฏิสนฺธิ\-กฺขเณ ฯเปฯ. นานากฺขณิกา อาสว\-วิปฺปยุตฺตา เจตนา วิปากานํ ขนฺธานํ กฏตฺตา จ รูปานํ กมฺมปจฺจเยน ปจฺจโย. (1)}

\csromanheader{2}{\textbf{วิปาก}}
\proseitem{93}{อาสว\-วิปฺปยุตฺโต ธมฺโม อาสว\-วิปฺปยุตฺตสฺส ธมฺมสฺส วิปาก\-ปจฺจเยน ปจฺจโย. เอกํ.}

\csromanheader{2}{\textbf{อาหาร}}
\proseitem{94}{อาสว\-สมฺปยุตฺโต ธมฺโม อาสว\-สมฺปยุตฺตสฺส ธมฺมสฺส อาหารปจฺจเยน ปจฺจโย. อาสว\-สมฺปยุตฺตา อาหารา สมฺปยุตฺตกานํ ขนฺธานํ อาหารปจฺจเยน ปจฺจโย. (1)}

\prose{อาสว\-สมฺปยุตฺโต ธมฺโม อาสว\-วิปฺปยุตฺตสฺส ธมฺมสฺส อาหารปจฺจเยน ปจฺจโย. อาสว\-สมฺปยุตฺตา อาหารา จิตฺต\-สมุฏฺฐานานํ รูปานํ อาหารปจฺจเยน ปจฺจโย. โทมนสฺส\-สหคตา วิจิกิจฺฉา\-สหคตา อุทฺธจฺจ\-สหคตา อาหารา โมหสฺส จิตฺตสมุฏฺฐานานญฺจ รูปานํ อาหารปจฺจเยน ปจฺจโย. (2)}

\prose{อาสว\-สมฺปยุตฺโต ธมฺโม อาสว\-สมฺปยุตฺตสฺส จ อาสว\-วิปฺปยุตฺตสฺส จ ธมฺมสฺส อาหารปจฺจเยน ปจฺจโย. อาสว\-สมฺปยุตฺตา อาหารา สมฺปยุตฺตกานํ ขนฺธานํ จิตฺตสมุฏฺฐานานญฺจ รูปานํ อาหารปจฺจเยน ปจฺจโย. โทมนสฺส\-สหคตา วิจิกิจฺฉา\-สหคตา อุทฺธจฺจ\-สหคตา อาหารา สมฺปยุตฺตกานํ ขนฺธานํ โมหสฺส จ จิตฺตสมุฏฺฐานานญฺจ รูปานํ อาหารปจฺจเยน ปจฺจโย. (3)}

\proseitem{95}{\csromanfolio{194}อาสว\-วิปฺปยุตฺโต ธมฺโม อาสว\-วิปฺปยุตฺตสฺส ธมฺมสฺส อาหารปจฺจเยน ปจฺจโย. อาสว\-วิปฺปยุตฺตา อาหารา สมฺปยุตฺตกานํ ขนฺธานํ จิตฺตสมุฏฺฐานานญฺจ รูปานํ อาหารปจฺจเยน ปจฺจโย. ปฏิสนฺธิ\-กฺขเณ ฯเปฯ. กพฬีกาโร อาหาโร อิมสฺส กายสฺส อาหารปจฺจเยน ปจฺจโย. (1)}

\csromanheader{2}{\textbf{อินฺทฺริยาทิ}}
\proseitem{96}{อาสว\-สมฺปยุตฺโต ธมฺโม อาสว\-สมฺปยุตฺตสฺส ธมฺมสฺส อินฺทฺริย\-ปจฺจเยน ปจฺจโย. จตฺตาริ. ฌานปจฺจเยน ปจฺจโย. จตฺตาริ. มคฺคปจฺจเยน ปจฺจโย. จตฺตาริ. สมฺปยุตฺต\-ปจฺจเยน ปจฺจโย. ฉ.}

\proseitem{97}{อาสว\-สมฺปยุตฺโต ธมฺโม อาสว\-วิปฺปยุตฺตสฺส ธมฺมสฺส วิปฺปยุตฺต\-ปจฺจเยน ปจฺจโย. สหชาตํ, ปจฺฉาชาตํ.  สหชาตา อาสว\-สมฺปยุตฺโต ขนฺธา จิตฺต\-สมุฏฺฐานานํ รูปานํ วิปฺปยุตฺต\-ปจฺจเยน ปจฺจโย. ปจฺฉาชาตา อาสว\-สมฺปยุตฺตา ขนฺธา ปุเรชาตสฺส อิมสฺส กายสฺส วิปฺปยุตฺต\-ปจฺจเยน ปจฺจโย. (1)}

\prose{อาสว\-วิปฺปยุตฺโต ธมฺโม อาสว\-วิปฺปยุตฺตสฺส ธมฺมสฺส วิปฺปยุตฺต\-ปจฺจเยน ปจฺจโย. สหชาตํ. ปุเรชาตํ. ปจฺฉาชาตํ. (สํขิตฺตํ. วิตฺถาเรตพฺพํ.) (1)}

\proseitem{98}{อาสว\-วิปฺปยุตฺโต ธมฺโม อาสว\-สมฺปยุตฺตสฺส ธมฺมสฺส วิปฺปยุตฺต\-ปจฺจเยน ปจฺจโย. ปุเรชาตํ วตฺถุ อาสว\-สมฺปยุตฺตกานํ ขนฺธานํ วิปฺปยุตฺต\-ปจฺจเยน ปจฺจโย. (1)}

\prose{อาสว\-วิปฺปยุตฺโต ธมฺโม อาสว\-สมฺปยุตฺตสฺส จ อาสว\-วิปฺปยุตฺตสฺส จ ธมฺมสฺส วิปฺปยุตฺต\-ปจฺจเยน ปจฺจโย. ปุเรชาตํ วตฺถุ โทมนสฺส\-สหคตานํ วิจิกิจฺฉา\-สหคตานํ อุทฺธจฺจ\-สหคตานํ ขนฺธานํ โมหสฺส จ วิปฺปยุตฺต\-ปจฺจเยน ปจฺจโย. (2)}

\prose{อาสว\-สมฺปยุตฺโต จ อาสววิยุตฺโต จ ธมฺมา อาสว\-วิปฺปยุตฺตสฺส ธมฺมสฺส วิปฺปยุตฺต\-ปจฺจเยน ปจฺจโย. สหชาตํ, ปจฺฉาชาตํ.  สหชาตา โทมนสฺส\-สหคตา วิจิกิจฺฉา\-สหคตา อุทฺธจฺจ\-สหคตา \csromanfolio{195}ขนฺธา จ โมโห จ จิตฺต\-สมุฏฺฐานํ รูปานํ วิปฺปยุตฺต\-ปจฺจเยน ปจฺจโย. ปจฺฉาชาตา โทมนสฺส\-สหคตา วิจิกิจฺฉา\-สหคตา อุทฺธจฺจ\-สหคตา ขนฺธา จ โมโห จ ปุเรชาตสฺส อิมสฺส กายสฺส วิปฺปยุตฺต\-ปจฺจเยน ปจฺจโย. (1)}

\proseitem{99}{อาสว\-สมฺปยุตฺโต ธมฺโม อาสว\-สมฺปยุตฺตสฺส ธมฺมสฺส อตฺถิปจฺจเยน ปจฺจโย. เอกํ. (1)}

\prose{อาสว\-สมฺปยุตฺโต ธมฺโม อาสว\-วิปฺปยุตฺตสฺส ธมฺมสฺส อตฺถิปจฺจเยน ปจฺจโย. สหชาตํ, ปจฺฉาชาตํ.  สหชาตา อาสว\-สมฺปยุตฺตา ขนฺธา จิตฺต\-สมุฏฺฐานานํ รูปานํ อตฺถิปจฺจเยน ปจฺจโย. โทมนสฺส\-สหคตา วิจิกิจฺฉา\-สหคตา อุทฺธจฺจ\-สหคตา ขนฺธา โมหสฺส จ จิตฺตสมุฏฺฐานานญฺจ รูปานํ อตฺถิปจฺจเยน ปจฺจโย. ปจฺฉาชาตา อาสว\-สมฺปยุตฺตา ขนฺธา ปุเรชาตสฺส อิมสฺส กายสฺส อตฺถิปจฺจเยน ปจฺจโย. (2)}

\prose{อาสว\-สมฺปยุตฺโต ธมฺโม อาสว\-สมฺปยุตฺตสฺส จ อาสว\-วิปฺปยุตฺตสฺส จ ธมฺมสฺส อตฺถิปจฺจเยน ปจฺจโย. (สหชาต\-สทิสํ.) (3)}

\proseitem{100}{อาสว\-วิปฺปยุตฺโต ธมฺโม อาสว\-วิปฺปยุตฺตสฺส ธมฺมสฺส อตฺถิปจฺจเยน ปจฺจโย. สหชาตํ, ปุเรชาตํ, ปจฺฉาชาตํ, อาหารํ, อินฺทฺริยํ. (สํขิตฺตํ. วิตฺถาเรตพฺพํ.) (1)}

\prose{อาสว\-วิปฺปยุตฺโต ธมฺโม อาสว\-สมฺปยุตฺตสฺส ธมฺมสฺส อตฺถิปจฺจเยน ปจฺจโย. สหชาตํ, ปุเรชาตํ.  สหชาโต โทนมสฺสสหคโต วิจิกิจฺฉา\-สหคโต อุทฺธจฺจ\-สหคโต โมโห สมฺปยุตฺตกานํ ขนฺธานํ อตฺถิปจฺจเยน ปจฺจโย. ปุเรชาตํ จกฺขุํ ฯเปฯ วตฺถุํ อสฺสาเทติ อภินนฺทติ, ตํ อารพฺภ ราโค อุปฺปชฺชติ, ทิฏฺฐิ อุปฺปชฺชติ, วิจิกิจฺฉา อุปฺปชฺชติ, อุทฺธจฺจํ อุปฺปชฺชติ, โทมนสฺสํ อุปฺปชฺชติ. วตฺถุ อาสว\-สมฺปยุตฺตกานํ ขนฺธานํ อตฺถิปจฺจเยน ปจฺจโย. (2)}

\prose{อาสว\-วิปฺปยุตฺโต ธมฺโม อาสว\-สมฺปยุตฺตสฺส จ อาสว\-วิปฺปยุตฺตสฺส จ ธมฺมสฺส อตฺถิปจฺจเยน ปจฺจโย. สหชาตํ, ปุเรชาตํ.  สหชาโต \csromanfolio{196}โทมนสฺส\-สหคโต วิจิกิจฺฉา\-สหคโต อุทฺธจฺจ\-สหคโต โมโห สมฺปยุตฺตกานํ ขนฺธานํ จิตฺตสมุฏฺฐานานญฺจ รูปานํ อตฺถิปจฺจเยน ปจฺจโย. ปุเรชาตํ จกฺขุํ ฯเปฯ วตฺถุํ อารพฺภ โทมนสฺส\-สหคตา วิจิกิจฺฉา\-สหคตา อุทฺธจฺจ\-สหคตา ขนฺธา จ โมโห จ อุปฺปชฺชนฺติ. (3)}

\proseitem{101}{อาสว\-สมฺปยุตฺโต จ อาสว\-วิปฺปยุตฺโต จ ธมฺมา อาสว\-สมฺปยุตฺตสฺส ธมฺมสฺส อตฺถิปจฺจเยน ปจฺจโย. สหชาตํ, ปุเรชาตํ.  สหชาโต อาสว\-สมฺปยุตฺโต เอโก ขนฺโธ จ วตฺถุ จ ติณฺณนฺนํ ขนฺธานํ อตฺถิปจฺจเยน ปจฺจโย ฯเปฯ เทฺว ขนฺธา จ ฯเปฯ. โทมนสฺส\-สหคโต วิจิกิจฺฉา\-สหคโต อุทฺธจฺจ\-สหคโต เอโก ขนฺโธ จ โมโห จ ติณฺณนฺนํ ขนฺธานํ อตฺถิปจฺจเยน ปจฺจโย. เทฺว ขนฺธา จ ฯเปฯ. (1)}

\prose{อาสว\-สมฺปยุตฺโต จ อาสว\-วิปฺปยุตฺโต จ ธมฺมา อาสว\-วิปฺปยุตฺตสฺส ธมฺมสฺส อตฺถิปจฺจเยน ปจฺจโย. สหชาตํ, ปุเรชาตํ, ปจฺฉาชาตํ, อาหารํ, อินฺทฺริยํ.  สหชาตา อาสว\-สมฺปยุตฺตา ขนฺธา จ มหาภูตา จ จิตฺต\-สมุฏฺฐานานํ รูปานํ อตฺถิปจฺจเยน ปจฺจโย. โทมนสฺส\-สหคตา วิจิกิจฺฉา\-สหคตา อุทฺธจฺจ\-สหคตา ขนฺธา จ โมโห จ จิตฺต\-สมุฏฺฐานานํ รูปานํ อตฺถิปจฺจเยน ปจฺจโย. โทมนสฺส\-สหคตา วิจิกิจฺฉา\-สหคตา อุทฺธจฺจ\-สหคตา ขนฺธา จ วตฺถุ จ โมหสฺส อตฺถิปจฺจเยน ปจฺจโย. ปจฺฉาชาตา อาสว\-สมฺปยุตฺตา ขนฺธา จ กพฬีกาโร อาหาโร จ อิมสฺส กายสฺส อตฺถิปจฺจเยน ปจฺจโย. ปจฺฉาชาตา อาสว\-สมฺปยุตฺตา ขนฺธา จ รูปชีวิตินฺทฺริยญฺจ กฏตฺตารูปานํ อตฺถิปจฺจเยน ปจฺจโย. (2)}

\prose{อาสว\-สมฺปยุตฺโต จ อาสว\-วิปฺปยุตฺโต จ ธมฺมา อาสว\-สมฺปยุตฺตสฺส จ อาสว\-วิปฺปยุตฺตสฺส จ ธมฺมสฺส อตฺถิปจฺจเยน ปจฺจโย. สหชาตํ, ปุเรชาตํ.  สหชาโต โทมนสฺส\-สหคโต วิจิกิจฺฉา\-สหคโต อุทฺธจฺจ\-สหคโต เอโก ขนฺโธ จ โมโห จ ติณฺณนฺนํ ขนฺธานํ จิตฺตสมุฏฺฐานานญฺจ รูปานํ อตฺถิปจฺจเยน ปจฺจโย ฯเปฯ เทฺว ขนฺธา จ ฯเปฯ. สหชาโต โทมนสฺส\-สหคโต วิจิกิจฺฉา\-สหคโต อุทฺธจฺจ\-สหคโต เอโก ขนฺโธ จ วตฺถุ จ ติณฺณนฺนํ ขนฺธานํ โมหสฺส จ ฯเปฯ เทฺว ขนฺธา จ ฯเปฯ. (3)}

\csromanmarkh{16. อาสว\-สมฺปยุตฺตทุก}\csromanmarkhii{1. ปจฺจยานุโลม 2. สงฺขฺยาวาร}\csromanheadernomark{2}{16. อาสว\-สมฺปยุตฺตทุก 1. ปจฺจยานุโลม 2. สงฺขฺยาวาร}
\proseitem{102}{\csromanfolio{197}เหตุยา นว, อารมฺมเณ นว, อธิปติยา ปญฺจ, อนนฺตเร นว, สมนนฺตเร นว, สหชาเต นว, อญฺญมญฺเญ ฉ, นิสฺสเย นว, อุปนิสฺสเย นว, ปุเรชาเต ตีณิ, ปจฺฉาชาเต ตีณิ, อาเสวเน นว, กมฺเม จตฺตาริ, วิปาเก เอกํ, อาหาเร จตฺตาริ, อินฺทฺริเย จตฺตาริ, ฌาเน จตฺตาริ, มคฺเค จตฺตาริ, สมฺปยุตฺเต ฉ, วิปฺปยุตฺเต ปญฺจ, อตฺถิยา นว, นตฺถิยา นว, วิคเต นว, อวิคเต นว.}

\proseitem{103}{อาสว\-สมฺปยุตฺโต ธมฺโม อาสว\-สมฺปยุตฺตสฺส ธมฺมสฺส อารมฺมณ\-ปจฺจเยน ปจฺจโย. สหชาต\-ปจฺจเยน ปจฺจโย. อุปนิสฺสย\-ปจฺจเยน ปจฺจโย. (1)}

\prose{อาสว\-สมฺปยุตฺโต ธมฺโม อาสว\-วิปฺปยุตฺตสฺส ธมฺมสฺส อารมฺมณ\-ปจฺจเยน ปจฺจโย. สหชาต\-ปจฺจเยน ปจฺจโย. อุปนิสฺสย\-ปจฺจเยน ปจฺจโย. ปจฺฉาชาต\-ปจฺจเยน ปจฺจโย. กมฺมปจฺจเยน ปจฺจโย. (2)}

\prose{อาสว\-สมฺปยุตฺโต ธมฺโม อาสว\-สมฺปยุตฺตสฺส จ อาสว\-วิปฺปยุตฺตสฺส จ ธมฺมสฺส อารมฺมณ\-ปจฺจเยน ปจฺจโย. สหชาต\-ปจฺจเยน ปจฺจโย. อุปนิสฺสย\-ปจฺจเยน ปจฺจโย. (3)}

\proseitem{104}{อาสว\-วิปฺปยุตฺโต ธมฺโม อาสว\-วิปฺปยุตฺตสฺส ธมฺมสฺส อารมฺมณ\-ปจฺจเยน ปจฺจโย. สหชาต\-ปจฺจเยน ปจฺจโย. อุปนิสฺสย\-ปจฺจเยน ปจฺจโย. ปุเรชาต\-ปจฺจเยน ปจฺจโย. ปจฺฉาชาต\-ปจฺจเยน ปจฺจโย. กมฺมปจฺจเยน ปจฺจโย. อาหารปจฺจเยน ปจฺจโย. อินฺทฺริย\-ปจฺจเยน ปจฺจโย. (1)}

\prose{อาสว\-วิปฺปยุตฺโต ธมฺโม อาสว\-สมฺปยุตฺตสฺส ธมฺมสฺส อารมฺมณ\-ปจฺจเยน ปจฺจโย. สหชาต\-ปจฺจเยน ปจฺจโย. อุปนิสฺสย\-ปจฺจเยน ปจฺจโย. ปุเรชาต\-ปจฺจเยน ปจฺจโย. (2)}

\prose{\csromanfolio{198}อาสว\-วิปฺปยุตฺโต ธมฺโม อาสว\-สมฺปยุตฺตสฺส จ อาสว\-วิปฺปยุตฺตสฺส จ ธมฺมสฺส อารมฺมณ\-ปจฺจเยน ปจฺจโย. สหชาต\-ปจฺจเยน ปจฺจโย. อุปนิสฺสย\-ปจฺจเยน ปจฺจโย. ปุเรชาต\-ปจฺจเยน ปจฺจโย. (3)}

\proseitem{105}{อาสว\-สมฺปยุตฺโต จ อาสว\-วิปฺปยุตฺโต จ ธมฺมา อาสว\-สมฺปยุตฺตสฺส ธมฺมสฺส อารมฺมณ\-ปจฺจเยน ปจฺจโย. สหชาต\-ปจฺจเยน ปจฺจโย. อุปนิสฺสย\-ปจฺจเยน ปจฺจโย. (1)}

\prose{อาสว\-สมฺปยุตฺโต จ อาสว\-วิปฺปยุตฺโต จ ธมฺมา อาสว\-วิปฺปยุตฺตสฺส ธมฺมสฺส อารมฺมณ\-ปจฺจเยน ปจฺจโย. สหชาต\-ปจฺจเยน ปจฺจโย. อุปนิสฺสย\-ปจฺจเยน ปจฺจโย. ปจฺฉาชาต\-ปจฺจเยน ปจฺจโย. (2)}

\prose{อาสว\-สมฺปยุตฺโต จ อาสว\-วิปฺปยุตฺโต จ ธมฺมา อาสว\-สมฺปยุตฺตสฺส จ อาสว\-วิปฺปยุตฺตสฺส จ ธมฺมสฺส อารมฺมณ\-ปจฺจเยน ปจฺจโย. สหชาต\-ปจฺจเยน ปจฺจโย. อุปนิสฺสย\-ปจฺจเยน ปจฺจโย. (3)}

\csromanmarkh{2. ปจฺจย\-ปจฺจนีย}\csromanmarkhii{2. สงฺขฺยาวาร}\csromanheadernomark{2}{2. ปจฺจย\-ปจฺจนีย 2. สงฺขฺยาวาร}
\proseitem{106}{นเหตุยา นว, นอารมฺมเณ นว, นอธิปติยา นว, (สพฺพตฺถ นว,) โนอวิคเต นว.}

\proseitem{107}{เหตุปจฺจยา นอารมฺมเณ นว, นอธิปติยา นว ฯเปฯ นสมนนฺตเร นว, นอญฺญมญฺเญ ตีณิ, นอุปนิสฺสเย นว ฯเปฯ นมคฺเค นว, นสมฺปยุตฺเต ตีณิ, นวิปฺปยุตฺเต ฉ, โนนตฺถิยา นว, โนวิคเต นว.}

\csromanmatikaanchor{mtk.21}
\csromantocmarkref{section}{17. อาสวสาสวทุก}{mtk.21}
\bookmark[dest={mtk.21},level=1]{17. อาสวสาสวทุก}
\csromanmarkh{17. อาสว\-สาสวทุก}\csromanmarkhii{4. ปจฺจย\-ปจฺจนียา\-นุโลม}\csromanheadernomark{1}{17. อาสว\-สาสวทุก 4. ปจฺจย\-ปจฺจนียา\-นุโลม}
\vspace{\tipitakaparskip}
\csromancenter{นเหตุปจฺจยา อารมฺมเณ นว, อธิปติยา ปญฺจ ฯเปฯ อวิคเต นว.}
\csromancenter{ปจฺจนียา\-นุโลมํ.}
\nitthitam{อาสว\-สมฺปยุตฺตทุกํ นิฏฺฐิตํ.}
\csromanmarkh{17. อาสว\-สาสวทุก}\csromanmarkhii{1. ปฏิจฺจวาร}\csromanheadernomark{2}{17. อาสว\-สาสวทุก 1. ปฏิจฺจวาร}
\csromanmarkh{1. ปจฺจยานุโลม}\csromanmarkhii{1. วิภงฺควาร}\csromanheadernomark{2}{1. \textbf{ปจฺจยานุโลม 1. วิภงฺควาร}}
\csromanheader{2}{\textbf{เหตุ}}
\proseitem{109}{\csromanfolio{199}อาสวญฺเจว สาสวญฺจ ธมฺมํ ปฏิจฺจ อาสโว เจว สาสโว จ ธมฺโม อุปฺปชฺชติ เหตุปจฺจยา. กามาสวํ ปฏิจฺจ ทิฏฺฐาสโว อวิชฺชาสโว. (จกฺกํ พนฺธิตพฺพํ.) ภวาสวํ ปฏิจฺจ อวิชฺชาสโว. (จกฺกํ พนฺธิตพฺพํ.) ทิฏฺฐาสวํ ปฏิจฺจ อวิชฺชาสโว. (1)}

\prose{อาสวญฺเจว สาสวญฺจ ธมฺมํ ปฏิจฺจ สาสโว เจว โน จ อาสโว ธมฺโม อุปฺปชฺชติ เหตุปจฺจยา. อาสเว ปฏิจฺจ สมฺปยุตฺตกา ขนฺธา จิตฺต\-สมุฏฺฐานญฺจ รูปํ. (2)}

\prose{อาสวญฺเจว สาสวญฺจ ธมฺมํ ปฏิจฺจ อาสโว เจว สาสโว จ สาสโว เจว โน จ อาสโว จ ธมฺมา อุปฺปชฺชนฺติ เหตุปจฺจยา. กามาสวํ ปฏิจฺจ ทิฏฺฐาสโว อวิชฺชาสโว สมฺปยุตฺตกา ขนฺธา จิตฺต\-สมุฏฺฐานญฺจ รูปํ. ภวาสวํ (จกฺกํ พนฺธิตพฺพํ.) (3)}

\proseitem{110}{สาสวญฺเจว โน จ อาสวํ ธมฺมํ ปฏิจฺจ สาสโว เจว โน จ อาสโว ธมฺโม อุปฺปชฺชติ เหตุปจฺจยา. สาสวญฺเจว โน จ อาสวํ เอกํ ขนฺธํ ปฏิจฺจ ตโย ขนฺธา จิตฺต\-สมุฏฺฐานญฺจ รูปํ ฯเปฯ เทฺว ขนฺเธ ฯเปฯ. ปฏิสนฺธิ\-กฺขเณ ฯเปฯ. ขนฺเธ ปฏิจฺจ วตฺถุ, วตฺถุํ ปฏิจฺจ ขนฺธา. เอกํ มหาภูตํ ฯเปฯ. (1)}

\prose{\csromanfolio{200}สาสวญฺเจว โน จ อาสวํ ธมฺมํ ปฏิจฺจ อาสโว เจว สาสโว จ ธมฺโม อุปฺปชฺชติ เหตุปจฺจยา. สาสเว เจว โน จ อาสเว ขนฺเธ ปฏิจฺจ อาสวา. (2)}

\prose{สาสวญฺเจว โน จ อาสวํ ธมฺมํ ปฏิจฺจ อาสโว เจว สาสโว จ สาสโว เจว โน จ อาสโว จ ธมฺมา อุปฺปชฺชนฺติ เหตุปจฺจยา. สาสวญฺเจว โน จ อาสวํ เอกํ ขนฺธํ ปฏิจฺจ ตโย ขนฺธา อาสวา จ จิตฺต\-สมุฏฺฐานญฺจ รูปํ. เทฺว ขนฺเธ ฯเปฯ. (3)}

\proseitem{111}{อาสวญฺเจว สาสวญฺจ สาสวญฺเจว โน จ อาสวญฺจ ธมฺมํ ปฏิจฺจ อาสโว เจว สาสโว จ ธมฺโม อุปฺปชฺชติ เหตุปจฺจยา. กามาสวญฺจ สมฺปยุตฺตเก จ ขนฺเธ ปฏิจฺจ ทิฏฺฐาสโว อวิชฺชาสโว. (เอวํ จกฺกํ พนฺธิตพฺพํ.) (1)}

\prose{อาสวญฺเจว สาสวญฺจ สาสวญฺเจว โน จ อาสวญฺจ ธมฺมํ ปฏิจฺจ สาสโว เจว โน จ อาสโว ธมฺโม อุปฺปชฺชติ เหตุปจฺจยา. สาสวญฺเจว โน จ อาสวํ เอกํ ขนฺธญฺจ อาสเว จ ปฏิจฺจ ตโย ขนฺธา จิตฺต\-สมุฏฺฐานญฺจ รูปํ. เทฺว ขนฺเธ จ ฯเปฯ. (2)}

\prose{อาสวญฺเจว สาสวญฺจ สาสวญฺเจว โน จ อาสวญฺจ ธมฺมํ ปฏิจฺจ อาสโว เจว สาสโว จ สาสโว เจว โน จ อาสโว จ ธมฺมา อุปฺปชฺชนฺติ เหตุปจฺจยา. สาสวญฺเจว โน จ อาสวํ เอกํ ขนฺธญฺจ กามาสวญฺจ ปฏิจฺจ ตโย ขนฺธา ทิฏฺฐาสโว อวิชฺชาสโว จิตฺต\-สมุฏฺฐานญฺจ รูปํ. เทฺว ขนฺเธ จ ฯเปฯ. (จกฺกํ สํขิตฺตํ.) (3)}

\prose{(เอวํ ปฏิจฺจวาโรปิ สหชาตวาโรปิ ปจฺจยวาโรปิ นิสฺสายวาโรปิ สํสฏฺฐ\-วาโรปิ สมฺปยุตฺตวาโรปิ ยถา อาสวทุกํ, เอวํ กาตพฺพํ. นินฺนานํ.)}

\csromanmarkh{17. อาสว\-สาสวทุก}\csromanmarkhii{7. ปญฺหาวาร}\csromanheadernomark{2}{17. อาสว\-สาสวทุก 7. ปญฺหาวาร}
\vspace{\tipitakaparskip}
\csromancenter{ปจฺจยานุโลม 1. วิภงฺควาร (ปญฺหาวาเร เหตุปจฺจเยปิ อารมฺมณ\-ปจฺจเยปิ โลกุตฺตรํ นกาตพฺพํ. เสกฺขา โคตฺรภุํ ปจฺจเวกฺขนฺติ. โวทานํ ปจฺจเวกฺขนฺตีติ กาตพฺพา.  อธิปติปจฺจยมฺปิ สพฺพํ ชานิตฺวา กาตพฺพํ.)}
\csromanheader{2}{17. อาสว\-สาสวทุก อนนฺตร}
\proseitem{112}{\csromanfolio{201}อาสโว เจว สาสโว จ ธมฺโม อาสวสฺส เจว สาสวสฺส จ ธมฺมสฺส อนนฺตร\-ปจฺจเยน ปจฺจโย. ปุริมา ปุริมา อาสวา ปจฺฉิมานํ ปจฺฉิมานํ อาสวานํ อนนฺตร\-ปจฺจเยน ปจฺจโย. (1)}

\prose{อาสโว เจว สาสโว จ ธมฺโม สาสวสฺส เจว โน จ อาสวสฺส ธมฺมสฺส อนนฺตร\-ปจฺจเยน ปจฺจโย. ปุริมา ปุริมา อาสวา ปจฺฉิมานํ ปจฺฉิมานํ สาสวานญฺเจว โน จ อาสวานํ ขนฺธานํ อนนฺตร\-ปจฺจเยน ปจฺจโย. อาสวา วุฏฺฐานสฺส อนนฺตร\-ปจฺจเยน ปจฺจโย.}

\prose{อาสโว เจว สาสโว จ ธมฺโม อาสวสฺส เจว สาสวสฺส จ สาสวสฺส เจว โน จ อาสวสฺส จ ธมฺมสฺส อนนฺตร\-ปจฺจเยน ปจฺจโย. ปุริมา ปุริมา อาสวา ปจฺฉิมานํ ปจฺฉิมานํ อาสวานํ สมฺปยุตฺตกานญฺจ ขนฺธานํ อนนฺตร\-ปจฺจเยน ปจฺจโย. (3)}

\proseitem{113}{สาสโว เจว โน จ อาสโว ธมฺโม สาสวสฺส เจว โน จ อาสวสฺส ธมฺมสฺส อนนฺตร\-ปจฺจเยน ปจฺจโย. ปุริมา ปุริมา สาสวา เจว โน จ อาสวา ขนฺธา ปจฺฉิมานํ ปจฺฉิมานํ สาสวานญฺเจว โน จ อาสวานํ ขนฺธานํ อนนฺตร\-ปจฺจเยน ปจฺจโย. อนุโลมํ โคตฺรภุสฺส. อนุโลมํ โวทานสฺส. อาวชฺชนา สาสวานญฺเจว โน จ อาสวานํ ขนฺธานํ อนนฺตร\-ปจฺจเยน ปจฺจโย. (1)}

\prose{สาสโว เจว โน จ อาสโว ธมฺโม อาสวสฺส เจว สาสวสฺส จ ธมฺมสฺส อนนฺตร\-ปจฺจเยน ปจฺจโย. ปุริมา ปุริมา สาสวา เจว โน จ อาสวา ขนฺธา ปจฺฉิมานํ ปจฺฉิมานํ อาสวานํ อนนฺตร\-ปจฺจเยน ปจฺจโย. อาวชฺชนา อาสวานํ อนนฺตร\-ปจฺจเยน ปจฺจโย. (2)}

\prose{สาสโว เจว โน จ อาสโว ธมฺโม อาสวสฺส เจว สาสวสฺส จ สาสวสฺส เจว โน จ อาสวสฺส จ ธมฺมสฺส อนนฺตร\-ปจฺจเยน ปจฺจโย. ปุริมา ปุริมา สาสวา เจว โน จ อาสวา ขนฺธา ปจฺฉิมานํ ปจฺฉิมานํ อาสวานํ สมฺปยุตฺตกานญฺจ ขนฺธานํ อนนฺตร\-ปจฺจเยน ปจฺจโย. อาวชฺชนา อาสวานํ สมฺปยุตฺตกานญฺจ ขนฺธานํ อนนฺตร\-ปจฺจเยน ปจฺจโย. (3)}

\csromanmatikaanchor{mtk.22}
\csromantocmarkref{section}{18. อาสวอาสวสมฺปยุตฺตทุก}{mtk.22}
\bookmark[dest={mtk.22},level=1]{18. อาสวอาสวสมฺปยุตฺตทุก}
\proseitem{114}{อาสโว เจว สาสโว จ สาสโว เจว โน จ อาสโว จ ธมฺมา อาสวสฺส เจว สาสวสฺส จ ธมฺมสฺส อนนฺตร\-ปจฺจเยน \csromanfolio{202}ปจฺจโย. ปุริมา ปุริมา อาสวา จ สมฺปยุตฺตกา จ ขนฺธา ปจฺฉิมานํ ปจฺฉิมานํ อาสวานํ อนนฺตร\-ปจฺจเยน ปจฺจโย. (1)}

\prose{อาสโว เจว สาสโว จ สาสโว เจว โน จ อาสโว จ ธมฺมา สาสวสฺส เจว โน จ อาสวสฺส ธมฺมสฺส อนนฺตร\-ปจฺจเยน ปจฺจโย. ปุริมา ปุริมา อาสวา จ สมฺปยุตฺตกา จ ขนฺธา ปจฺฉิมานํ ปจฺฉิมานํ สาสวานญฺเจว โน จ อาสวานํ ขนฺธานํ อนนฺตร\-ปจฺจเยน ปจฺจโย. อาสวา จ สมฺปยุตฺตกา จ ขนฺธา วุฏฺฐานสฺส อนนฺตร\-ปจฺจเยน ปจฺจโย.}

\prose{อาสโว เจว สาสโว จ สาสโว เจว โน จ อาสโว จ ธมฺมา อาสวสฺส เจว สาสวสฺส จ สาสวสฺส เจว โน จ อาสวสฺส จ ธมฺมสฺส อนนฺตร\-ปจฺจเยน ปจฺจโย. ปุริมา ปุริมา อาสวา จ สมฺปยุตฺตกา จ ขนฺธา ปจฺฉิมานํ ปจฺฉิมานํ อาสวานํ สมฺปยุตฺตกานญฺจ ขนฺธานํ อนนฺตร\-ปจฺจเยน ปจฺจโย. (เอวํ สพฺพํ วิตฺถาเรตพฺพํ.) (3)}

\prosewithcloser{(อาสวทุเกปิ อนนฺตรํ อิมินา สทิสํ กาตพฺพํ. อาวชฺชนาปิ วุฏฺฐานมฺปิ เอวํ สมุทฺทิฏฺฐํ สํขิตฺตํ. สพฺพํ ปริปุณฺณํ. อาสว\-ทุก\-สทิสํ กาตพฺพํ. นินฺนานํ.)}{อาสว\-สาสวทุกํ นิฏฺฐิตํ.}
\csromanmarkh{18. อาสว\-อาสว\-สมฺปยุตฺตทุก}\csromanmarkhii{1. ปฏิจฺจวาร}\csromanheadernomark{2}{18. อาสว\-อาสว\-สมฺปยุตฺตทุก 1. ปฏิจฺจวาร}
\csromanmarkh{1. ปจฺจยานุโลม}\csromanmarkhii{1. วิภงฺควาร}\csromanheadernomark{2}{1. \textbf{ปจฺจยานุโลม 1. วิภงฺควาร}}
\csromanheader{2}{\textbf{เหตุ}}
\proseitem{115}{อาสวญฺเจว อาสว\-สมฺปยุตฺตญฺจ ธมฺมํ ปฏิจฺจ อาสโว เจว อาสว\-สมฺปยุตฺโต จ ธมฺโม อุปฺปชฺชติ เหตุปจฺจยา. กามาสวํ ปฏิจฺจ ทิฏฺฐาสโว อวิชฺชาสโว. (จกฺกํ พนฺธิตพฺพํ.) ภวาสวํ ปฏิจฺจ อวิชฺชาสโว. (จกฺกํ พนฺธิตพฺพํ.) ทิฏฺฐาสวํ ปฏิจฺจ อวิชฺชาสโว. (1)}

\prose{อาสวญฺเจว อาสว\-สมฺปยุตฺตญฺจ ธมฺมํ ปฏิจฺจ อาสว\-สมฺปยุตฺโต เจว โน จ อาสโว ธมฺโม อุปฺปชฺชติ เหตุปจฺจยา. อาสเว ปฏิจฺจ สมฺปยุตฺตกา ขนฺธา. (2)}

\prose{\csromanfolio{203}18. อาสว\-อาสว\-สมฺปยุตฺตทุก}

\prose{อาสวญฺเจว อาสว\-สมฺปยุตฺตญฺจ ธมฺมํ ปฏิจฺจ อาสโว เจว อาสว\-สมฺปยุตฺโต จ อาสว\-สมฺปยุตฺโต เจว โน จ อาสโว จ ธมฺมา อุปฺปชฺชนฺติ เหตุปจฺจยา. อาสว\-สมฺปยุตฺตํ กามาสวํ ปฏิจฺจ ทิฏฺฐาสโว อวิชฺชาสโว สมฺปยุตฺตกา จ ขนฺธา. (สพฺพํ จกฺกํ.) (3)}

\proseitem{116}{อาสว\-สมฺปยุตฺต\-ญฺเจว โน จ อาสวํ ธมฺมํ ปฏิจฺจ อาสว\-สมฺปยุตฺโต เจว โน จ อาสโว ธมฺโม อุปฺปชฺชติ เหตุปจฺจยา. อาสว\-สมฺปยุตฺต\-ญฺเจว โน จ อาสวํ เอกํ ขนฺธํ ปฏิจฺจ ตโย ขนฺธา ฯเปฯ เทฺว ขนฺเธ ฯเปฯ. (1)}

\prose{อาสว\-สมฺปยุตฺต\-ญฺเจว โน จ อาสวํ ธมฺมํ ปฏิจฺจ อาสโว เจว อาสว\-สมฺปยุตฺโต จ ธมฺโม อุปฺปชฺชติ เหตุปจฺจยา. อาสว\-สมฺปยุตฺต\-ญฺเจว โน จ อาสเว ขนฺเธ ปฏิจฺจ อาสวา. (2)}

\prose{อาสว\-สมฺปยุตฺต\-ญฺเจว โน จ อาสวํ ธมฺมํ ปฏิจฺจ อาสโว เจว อาสว\-สมฺปยุตฺโต จ อาสว\-สมฺปยุตฺโต เจว โน จ อาสโว จ ธมฺมา อุปฺปชฺชนฺติ เหตุปจฺจยา. อาสว\-สมฺปยุตฺต\-ญฺเจว โน จ อาสวํ เอกํ ขนฺธํ ปฏิจฺจ ตโย ขนฺธา อาสวา จ ฯเปฯ เทฺว ขนฺเธ ฯเปฯ. (3)}

\proseitem{117}{อาสวญฺเจว อาสว\-สมฺปยุตฺตญฺจ อาสว\-สมฺปยุตฺต\-ญฺเจว โน จ อาสวญฺจ ธมฺมํ ปฏิจฺจ อาสโว เจว อาสว\-สมฺปยุตฺโต จ ธมฺโม อุปฺปชฺชติ เหตุปจฺจยา. กามาสวญฺจ สมฺปยุตฺตเก จ ขนฺเธ ปฏิจฺจ ทิฏฺฐาสโว อวิชฺชาสโว. (สพฺพํ จกฺกํ.) (1)}

\prose{อาสวญฺเจว อาสว\-สมฺปยุตฺตญฺจ อาสว\-สมฺปยุตฺต\-ญฺเจว โน จ อาสวญฺจ ธมฺมํ ปฏิจฺจ อาสว\-สมฺปยุตฺโต เจว โน จ อาสโว ธมฺโม อุปฺปชฺชติ เหตุปจฺจยา. อาสว\-สมฺปยุตฺต\-ญฺเจว โน จ อาสวํ เอกํ ขนฺธญฺจ อาสเว จ ปฏิจฺจ ตโย ขนฺธา ฯเปฯ เทฺว ขนฺเธ ฯเปฯ. (2)}

\prose{อาสวญฺเจว อาสว\-สมฺปยุตฺตญฺจ อาสว\-สมฺปยุตฺต\-ญฺเจว โน จ อาสวญฺจ ธมฺมํ ปฏิจฺจ อาสโว เจว อาสว\-สมฺปยุตฺโต จ อาสว\-สมฺปยุตฺโต เจว โน จ อาสโว จ ธมฺมา อุปฺปชฺชนฺติ เหตุปจฺจยา. อาสว\-สมฺปยุตฺต\-ญฺเจว โน จ อาสวํ เอกํ ขนฺธญฺจ กามาสวญฺจ ปฏิจฺจ ตโย ขนฺธา ทิฏฺฐาสโว อวิชฺชาสโว ฯเปฯ เทฺว ขนฺเธ ฯเปฯ. (3)}

\vspace{\tipitakaparskip}
\csromancenter{(จกฺกํ. เอวํ สพฺเพ ปจฺจยา กาตพฺพา.)}
\csromanmarkh{1. ปจฺจยานุโลม}\csromanmarkhii{2. สงฺขฺยาวาร}\csromanheadernomark{2}{1. \textbf{ปจฺจยานุโลม 2. สงฺขฺยาวาร}}
\proseitem{118}{\csromanfolio{204}เหตุยา นว, อารมฺมเณ นว, อธิปติยา นว, (สพฺพตฺถ นว, สํขิตฺตํ.) กมฺเม นว, (วิปากํ นตฺถิ.) อาหาเร นว ฯเปฯ อวิคเต นว.}

\vspace{\tipitakaparskip}
\csromancenter{\textbf{อนุโลมํ.}}
\csromanmarkh{2. ปจฺจย\-ปจฺจนีย}\csromanmarkhii{1. วิภงฺควาร}\csromanheadernomark{2}{2. ปจฺจย\-ปจฺจนีย 1. วิภงฺควาร}
\prose{นอธิปตฺยาทิ}

\proseitem{119}{อาสวญฺเจว อาสว\-สมฺปยุตฺตํ ธมฺมํ ปฏิจฺจ อาสโว เจว อาสว\-สมฺปยุตฺโต จ ธมฺโม อุปฺปชฺชติ นอธิปติ\-ปจฺจยา. (นเหตุมูลกํ นตฺถิ.) นปุเรชาต\-ปจฺจยา. นปจฺฉาชาต\-ปจฺจยา. (สํขิตฺตํ.)}

\csromanmarkh{2. ปจฺจย\-ปจฺจนีย}\csromanmarkhii{2. สงฺขฺยาวาร}\csromanheadernomark{2}{2. ปจฺจย\-ปจฺจนีย 2. สงฺขฺยาวาร}
\proseitem{120}{นอธิปติยา นว, นปุเรชาเต นว, นปจฺฉาชาเต นว, นอาเสวเน นว, นกมฺเม ตีณิ, นวิปาเก นว, นวิปฺปยุตฺเต นว.}

\prose{(เอวํ อิตเร เทฺว คณนาปิ สหชาตวารมฺปิ ปจฺจยวารมฺปิ นิสฺสยวารมฺปิ สํสฏฺฐ\-วารมฺปิ สมฺปยุตฺตวารมฺปิ ปริปุณฺณํ ปฏิจฺจสทิสา.)}

\prose{18. อาสว\-อาสว\-สมฺปยุตฺตทุก 7. ปญฺหาวาร}

\csromanmarkh{1. ปจฺจยานุโลม}\csromanmarkhii{1. วิภงฺควาร}\csromanheadernomark{2}{1. \textbf{ปจฺจยานุโลม 1. วิภงฺควาร}}
\vspace{\tipitakaparskip}
\csromancenter{\textbf{เหตุ}}
\proseitem{121}{อาสโว เจว อาสว\-สมฺปยุตฺโต จ ธมฺโม อาสวสฺส เจว อาสว\-สมฺปยุตฺตสฺส จ ธมฺมสฺส เหตุปจฺจเยน ปจฺจโย. ตีณิ.}

\prose{\csromanfolio{205}18. อาสว\-อาสว\-สมฺปยุตฺตทุก}

\prose{อาสว\-สมฺปยุตฺโต เจว โน จ อาสโว ธมฺโม อาสว\-สมฺปยุตฺตสฺส เจว โน จ อาสวสฺส ธมฺมสฺส เหตุปจฺจเยน ปจฺจโย. อาสว\-สมฺปยุตฺตา เจว โน จ อาสวา เหตู สมฺปยุตฺตกานํ ขนฺธานํ เหตุปจฺจเยน ปจฺจโย. (1)}

\csromanheader{2}{\textbf{อารมฺมณาทิ}}
\proseitem{122}{อาสโว เจว อาสว\-สมฺปยุตฺโต จ ธมฺโม อาสวสฺส เจว อาสว\-สมฺปยุตฺตสฺส จ ธมฺมสฺส อารมฺมณ\-ปจฺจเยน ปจฺจโย. ตีณิ.}

\prose{อาสว\-สมฺปยุตฺโต เจว โน จ อาสโว ธมฺโม อาสว\-สมฺปยุตฺตสฺส เจว โน จ อาสวสฺส ธมฺมสฺส อารมฺมณ\-ปจฺจเยน ปจฺจโย. อาสว\-สมฺปยุตฺเต เจว โน จ อาสเว ขนฺเธ อารพฺภ อาสว\-สมฺปยุตฺตา เจว โน จ อาสวา ขนฺธา อุปฺปชฺชนฺติ. (1)}

\prose{อาสว\-สมฺปยุตฺโต เจว โน จ อาสโว ธมฺโม อาสวสฺส เจว อาสว\-สมฺปยุตฺตสฺส จ ธมฺมสฺส อารมฺมณ\-ปจฺจเยน ปจฺจโย. อาสว\-สมฺปยุตฺเต เจว โน จ อาสเว ขนฺเธ อารพฺภ อาสวา อุปฺปชฺชนฺติ. (2)}

\prose{อาสว\-สมฺปยุตฺโต เจว โน จ อาสโว ธมฺโม อาสวสฺส เจว อาสว\-สมฺปยุตฺตสฺส จ อาสว\-สมฺปยุตฺตสฺส เจว โน จ อาสวสฺส จ ธมฺมสฺส อารมฺมณ\-ปจฺจเยน ปจฺจโย. อาสว\-สมฺปยุตฺเต เจว โน จ อาสเว ขนฺเธ อารพฺภ อาสวา จ อาสว\-สมฺปยุตฺตกา จ ขนฺธา อุปฺปชฺชนฺติ. (3)}

\prose{อาสโว เจว อาสว\-สมฺปยุตฺโต จ อาสว\-สมฺปยุตฺโต เจว โน จ อาสโว จ ธมฺมา อาสวสฺส เจว อาสว\-สมฺปยุตฺตสฺส จ ธมฺมสฺส อารมฺมณ\-ปจฺจเยน ปจฺจโย. ตีณิ.}

\prose{อธิปติ\-ปจฺจยา. (อารมฺมณสทิสา, ครุการมฺมณา.) อนนฺตร\-ปจฺจยา. (อารมฺมณ\-สทิสาเยว, “ปุริมา ปุริมา”ติ กาตพฺพา.) สมนนฺตร\-ปจฺจยา.สหชาต\-ปจฺจยา. อญฺญมญฺญ\-ปจฺจยา. นิสฺสยปจฺจยา. อุปนิสฺสย\-ปจฺจยา. (อารมฺมณสทิสํเยว. วิภชนา นตฺถิ. ตีณิ. อุปนิสฺสยํ สพฺพํ กาตพฺพํ.)}

\csromanheader{2}{\textbf{กมฺมาทิ}}
\proseitem{123}{\csromanfolio{206}อาสว\-สมฺปยุตฺโต เจว โน จ อาสโว ธมฺโม อาสว\-สมฺปยุตฺตสฺส เจว โน จ อาสวสฺส ธมฺมสฺส กมฺมปจฺจเยน ปจฺจโย. ตีณิ. อาหารปจฺจเยน ปจฺจโย. ตีณิ. อินฺทฺริย\-ปจฺจเยน ปจฺจโย. ตีณิ. ฌานปจฺจเยน ปจฺจโย. ตีณิ. มคฺคปจฺจเยน ปจฺจโย. นว. สมฺปยุตฺต\-ปจฺจเยน ปจฺจโย. นว. อตฺถิปจฺจเยน ปจฺจโย. นตฺถิ\-ปจฺจเยน ปจฺจโย. วิคต\-ปจฺจเยน ปจฺจโย. อวิคต\-ปจฺจเยน ปจฺจโย. นว.}

\csromanmarkh{1. ปจฺจยานุโลม}\csromanmarkhii{2. สงฺขฺยาวาร}\csromanheadernomark{2}{1. \textbf{ปจฺจยานุโลม 2. สงฺขฺยาวาร}}
\proseitem{124}{เหตุยา จตฺตาริ, อารมฺมเณ นว, อธิปติยา นว, อนนฺตเร นว, สมนนฺตเร นว, สหชาเต นว, อญฺญมญฺเญ นว, นิสฺสเย นว, อุปนิสฺสเย นว, อาเสวเน นว, กมฺเม ตีณิ, อาหาเร ตีณิ, อินฺทฺริเย ตีณิ, ฌาเน ตีณิ, มคฺเค นว, สมฺปยุตฺเต นว, อตฺถิยา นว, นตฺถิยา นว, วิคเต นว, อวิคเต นว.}

\proseitem{125}{อาสโว เจว อาสว\-สมฺปยุตฺโต จ ธมฺโม อาสวสฺส เจว อาสว\-สมฺปยุตฺตสฺส จ ธมฺมสฺส อารมฺมณ\-ปจฺจเยน ปจฺจโย. สหชาต\-ปจฺจเยน ปจฺจโย. อุปนิสฺสย\-ปจฺจเยน ปจฺจโย. (1)}

\prose{อาสโว เจว อาสว\-สมฺปยุตฺโต จ ธมฺโม อาสว\-สมฺปยุตฺตสฺส เจว โน จ อาสวสฺส ธมฺมสฺส อารมฺมณ\-ปจฺจเยน ปจฺจโย. สหชาต\-ปจฺจเยน ปจฺจโย. อุปนิสฺสย\-ปจฺจเยน ปจฺจโย. (2)}

\prose{อาสโว เจว อาสว\-สมฺปยุตฺโต จ ธมฺโม อาสวสฺส เจว อาสว\-สมฺปยุตฺตสฺส จ อาสว\-สมฺปยุตฺตสฺส เจว โน จ อาสวสฺส จ ธมฺมสฺส อารมฺมณ\-ปจฺจเยน ปจฺจโย. สหชาต\-ปจฺจเยน ปจฺจโย. อุปนิสฺสย\-ปจฺจเยน ปจฺจโย. (3)}

\prose{\csromanfolio{207}18. อาสว\-อาสว\-สมฺปยุตฺตทุก}

\proseitem{126}{อาสว\-สมฺปยุตฺโต เจว โน จ อาสโว ธมฺโม อาสว\-สมฺปยุตฺตสฺส เจว โน จ อาสวสฺส ธมฺมสฺส อารมฺมณ\-ปจฺจเยน ปจฺจโย. สหชาต\-ปจฺจเยน ปจฺจโย. อุปนิสฺสย\-ปจฺจเยน ปจฺจโย. (1)}

\prose{อาสว\-สมฺปยุตฺโต เจว โน จ อาสโว ธมฺโม อาสวสฺส เจว อาสว\-สมฺปยุตฺตสฺส จ ธมฺมสฺส อารมฺมณ\-ปจฺจเยน ปจฺจโย. สหชาต\-ปจฺจเยน ปจฺจโย. อุปนิสฺสย\-ปจฺจเยน ปจฺจโย. (2)}

\prose{อาสว\-สมฺปยุตฺโต เจว โน จ อาสโว ธมฺโม อาสวสฺส เจว อาสว\-สมฺปยุตฺตสฺส จ อาสว\-สมฺปยุตฺตสฺส เจว โน จ อาสวสฺส จ ธมฺมสฺส อารมฺมณ\-ปจฺจเยน ปจฺจโย. สหชาต\-ปจฺจเยน ปจฺจโย. อุปนิสฺสย\-ปจฺจเยน ปจฺจโย. (3)}

\proseitem{127}{อาสโว เจว อาสว\-สมฺปยุตฺโต จ อาสว\-สมฺปยุตฺโต เจว โน จ อาสโว จ ธมฺมา อาสวสฺส เจว อาสว\-สมฺปยุตฺตสฺส จ ธมฺมสฺส อารมฺมณ\-ปจฺจเยน ปจฺจโย. สหชาต\-ปจฺจเยน ปจฺจโย. อุปนิสฺสย\-ปจฺจเยน ปจฺจโย. (1)}

\prose{อาสโว เจว อาสว\-สมฺปยุตฺโต จ อาสว\-สมฺปยุตฺโต เจว โน จ อาสโว จ ธมฺมา อาสว\-สมฺปยุตฺตสฺส เจว โน จ อาสวสฺส ธมฺมสฺส อารมฺมณ\-ปจฺจเยน ปจฺจโย. สหชาต\-ปจฺจเยน ปจฺจโย. อุปนิสฺสย\-ปจฺจเยน ปจฺจโย. (2)}

\prose{อาสโว เจว อาสว\-สมฺปยุตฺโต จ อาสว\-สมฺปยุตฺโต เจว โน จ อาสโว จ ธมฺมา อาสวสฺส เจว อาสว\-สมฺปยุตฺตสฺส จ อาสว\-สมฺปยุตฺตสฺส เจว โน จ อาสวสฺส จ ธมฺมสฺส อารมฺมณ\-ปจฺจเยน ปจฺจโย. สหชาต\-ปจฺจเยน ปจฺจโย. อุปนิสฺสย\-ปจฺจเยน ปจฺจโย. (3)}

\csromanmarkh{2. ปจฺจย\-ปจฺจนีย}\csromanmarkhii{2. สงฺขฺยาวาร}\csromanheadernomark{2}{2. ปจฺจย\-ปจฺจนีย 2. สงฺขฺยาวาร}
\vspace{\tipitakaparskip}
\csromancenter{นเหตุยา นว, นอารมฺมเณ นว, (สพฺพตฺถ นว.) โนอวิคเต นว.}
\csromancenter{\textbf{เหตุ}ปจฺจยา นอารมฺมเณ จตฺตาริ ฯเปฯ นสมนนฺตเร จตฺตาริ, นฺเอาปนิสฺสเย จตฺตาริ ฯเปฯ นมคฺเค จตฺตาริ ฯเปฯ นวิปฺปยุตฺเต จตฺตาริ, โนนตฺถิยา จตฺตาริ, โนวิคเต จตฺตาริ.}
\proseitemwithcloser{130}{\csromanfolio{208}นเหตุปจฺจยา อารมฺมเณ นว, อธิปติยา นว, (อนุโลมปทานิ คณิตพฺพานิ.) ฯเปฯ อวิคเต นว.}{อาสว\-อาสว\-สมฺปยุตฺตทุกํ นิฏฺฐิตํ.}
\csromanmatikaanchor{mtk.23}
\csromantocmarkref{section}{19. อาสววิปฺปยุตฺตสาสวทุก}{mtk.23}
\bookmark[dest={mtk.23},level=1]{19. อาสววิปฺปยุตฺตสาสวทุก}
\csromanmarkh{19. อาสว\-วิปฺปยุตฺต\-สาสวทุก}\csromanmarkhii{1. ปฏิจฺจวาร}\csromanheadernomark{1}{19. อาสว\-วิปฺปยุตฺต\-สาสวทุก 1. ปฏิจฺจวาร}
\csromanheader{2}{\textbf{เหตุ}}
\proseitem{131}{อาสว\-วิปฺปยุตฺตํ สาสวํ ธมฺมํ ปฏิจฺจ อาสว\-วิปฺปยุตฺโต สาสโว ธมฺโม อุปฺปชฺชติ เหตุปจฺจยา. อาสว\-วิปฺปยุตฺตํ สาสวํ เอกํ ขนฺธํ ปฏิจฺจ ตโย ขนฺธา จิตฺต\-สมุฏฺฐานญฺจ รูปํ ฯเปฯ เทฺว ขนฺเธ ฯเปฯ. ปฏิสนฺธิ\-กฺขเณ ฯเปฯ. ขนฺเธ ปฏิจฺจ วตฺถุ, วตฺถุํ ปฏิจฺจ ขนฺธา. เอกํ มหาภูตํ ฯเปฯ. (1)}

\prose{อาสว\-วิปฺปยุตฺตํ อนาสวํ ธมฺมํ ปฏิจฺจ อาสว\-วิปฺปยุตฺโต อนาสโว ธมฺโม อุปฺปชฺชติ เหตุปจฺจยา. ตีณิ.}

\prosewithcloser{(ยถา จูฬนฺตรทุเก โลกิยทุกํ, เอวํ วิตฺถาเรตพฺพํ. นินฺนานากรณํ. สํขิตฺตํ.)}{อาสว\-วิปฺปยุตฺต\-สาสวทุกํ นิฏฺฐิตํ.}
\nitthitam{อาสวโคจฺฉกํ นิฏฺฐิตํ.}
\csromanmatikaanchor{mtk.24}
\csromanmatikaanchor{mtk.25}
\csromantocmarknopage{chapter}{4. สญฺโญชนโคจฺฉก}{mtk.24}
\bookmark[dest={mtk.24},level=0]{4. สญฺโญชนโคจฺฉก}
\csromantocmarkref{section}{20. สญฺโญชนทุก}{mtk.25}
\bookmark[dest={mtk.25},level=1]{20. สญฺโญชนทุก}
\csromanmarkh{20. สญฺโญชนทุก}\csromanmarkhii{4. สญฺโญชน\-โคจฺฉก}\csromanheadernomark{1}{20. สญฺโญชนทุก 4. สญฺโญชน\-โคจฺฉก}
\csromanmarkh{20. สญฺโญชนทุก}\csromanmarkhii{1. ปฏิจฺจวาร}\csromanheadernomark{2}{20. สญฺโญชนทุก 1. ปฏิจฺจวาร}
\csromanmarkh{1. ปจฺจยานุโลม}\csromanmarkhii{1. วิภงฺควาร}\csromanheadernomark{2}{1. \textbf{ปจฺจยานุโลม 1. วิภงฺควาร}}
\csromanheader{2}{\textbf{เหตุ}}
\proseitem{1}{\csromanfolio{209}สญฺโญชนํ ธมฺมํ ปฏิจฺจ สญฺโญชโน ธมฺโม อุปฺปชฺชติ เหตุปจฺจยา. กามราค\-สญฺโญชนํ ปฏิจฺจ ทิฏฺฐิ\-สญฺโญชนํ อวิชฺชา\-สญฺโญชนํ. กามราค\-สญฺโญชนํ ปฏิจฺจ สีลพฺพตปรามาส\-สญฺโญชนํ อวิชฺชา\-สญฺโญชนํ. กามราค\-สญฺโญชนํ ปฏิจฺจ มานสญฺโญชนํ อวิชฺชา\-สญฺโญชนํ. กามราค\-สญฺโญชนํ ปฏิจฺจ อวิชฺชา\-สญฺโญชนํ. ปฏิฆ\-สญฺโญชนํ ปฏิจฺจ อิสฺสาสญฺโญชนํ อวิชฺชา\-สญฺโญชนํ. ปฏิฆ\-สญฺโญชนํ ปฏิจฺจ มจฺฉริย\-สญฺโญชนํ อวิชฺชา\-สญฺโญชนํ. ปฏิฆ\-สญฺโญชนํ ปฏิจฺจ อวิชฺชา\-สญฺโญชนํ. มานสญฺโญชนํ ปฏิจฺจ ภวราค\-สญฺโญชนํ อวิชฺชา\-สญฺโญชนํ. ภวราค\-สญฺโญชนํ ปฏิจฺจ อวิชฺชา\-สญฺโญชนํ. วิจิกิจฺฉา\-สญฺโญชนํ ปฏิจฺจ อวิชฺชา\-สญฺโญชนํ. (1)}

\prose{สญฺโญชนํ ธมฺมํ ปฏิจฺจ โนสญฺโญชโน ธมฺโม อุปฺปชฺชติ เหตุปจฺจยา. สญฺโญชเน ปฏิจฺจ สมฺปยุตฺตกา ขนฺธา จิตฺต\-สมุฏฺฐานญฺจ รูปํ. (2)}

\prose{สญฺโญชนํ ธมฺมํ ปฏิจฺจ สญฺโญชโน จ โนสญฺโญชโน จ ธมฺมา อุปฺปชฺชนฺติ เหตุปจฺจยา. กามราค\-สญฺโญชนํ ปฏิจฺจ ทิฏฺฐิ\-สญฺโญชนํ อวิชฺชา\-สญฺโญชนํ สมฺปยุตฺตกา จ ขนฺธา จ จิตฺต\-สมุฏฺฐานญฺจ รูปํ. (จกฺกํ พนฺธิตพฺพํ.) (3)}

\proseitem{2}{โนสญฺโญชนํ ธมฺมํ ปฏิจฺจ โนสญฺโญชโน ธมฺโม อุปฺปชฺชติ เหตุปจฺจยา. โนสญฺโญชนํ เอกํ ขนฺธํ ปฏิจฺจ ตโย ขนฺธา จิตฺต\-สมุฏฺฐานญฺจ รูปํ ฯเปฯ เทฺว ขนฺเธ ฯเปฯ. ปฏิสนฺธิ\-กฺขเณ ฯเปฯ. ขนฺเธ ปฏิจฺจ วตฺถุ, วตฺถุํ ปฏิจฺจ ขนฺธา. เอกํ มหาภูตํ ฯเปฯ มหาภูเต ปฏิจฺจ จิตฺต\-สมุฏฺฐานํ รูปํ, กฏตฺตารูปํ, อุปาทารูปํ. (1)}

\prose{โนสญฺโญชนํ ธมฺมํ ปฏิจฺจ สญฺโญชโน ธมฺโม อุปฺปชฺชติ เหตุปจฺจยา. โนสญฺโญชเน ขนฺเธ ปฏิจฺจ สญฺโญชนา. (2)}

\prose{\csromanfolio{210}โนสญฺโญชนํ ธมฺมํ ปฏิจฺจ สญฺโญชโน จ โนสญฺโญชโน จ ธมฺมา อุปฺปชฺชนฺติ เหตุปจฺจยา. โนสญฺโญชนํ เอกํ ขนฺธํ ปฏิจฺจ ตโย ขนฺธา สญฺโญชนา จ จิตฺต\-สมุฏฺฐานญฺจ รูปํ. เทฺว ขนฺเธ ฯเปฯ. (3)}

\proseitem{3}{สญฺโญชนญฺจ โนสญฺโญชนญฺจ ธมฺมํ ปฏิจฺจ สญฺโญชโน ธมฺโม อุปฺปชฺชติ เหตุปจฺจยา. กามราคสญฺโญชนญฺจ สมฺปยุตฺตเก จ ขนฺเธ ปฏิจฺจ ทิฏฺฐิ\-สญฺโญชนํ อวิชฺชา\-สญฺโญชนํ. (จกฺกํ พนฺธิตพฺพํ.) (1)}

\prose{สญฺโญชนญฺจ โนสญฺโญชนญฺจ ธมฺมํ ปฏิจฺจ โนสญฺโญชโน ธมฺโม อุปฺปชฺชติ เหตุปจฺจยา. โนสญฺโญชนํ เอกํ ขนฺธญฺจ สญฺโญชเน จ ปฏิจฺจ ตโย ขนฺธา จิตฺต\-สมุฏฺฐานญฺจ รูปํ ฯเปฯ เทฺว ขนฺเธ จ ฯเปฯ.}

\prose{สญฺโญชนญฺจ โนสญฺโญชนญฺจ ธมฺมํ ปฏิจฺจ สญฺโญชโน จ โนสญฺโญชโน จ ธมฺมา อุปฺปชฺชนฺติ เหตุปจฺจยา. โนสญฺโญชนํ เอกํ ขนฺธญฺจ กามราคสญฺโญชนญฺจ ปฏิจฺจ ตโย ขนฺธา ทิฏฺฐิ\-สญฺโญชนํ อวิชฺชา\-สญฺโญชนํ จิตฺต\-สมุฏฺฐานญฺจ รูปํ ฯเปฯ. (จกฺกํ พนฺธิตพฺพํ.) (3)}

\prose{(อารมฺมณ\-ปจฺจเย รูปํ นตฺถิ. อธิปติ\-ปจฺจโย เหตุสทิโส. วิจิกิจฺฉา\-สญฺโญชนํ นตฺถิ.) อนนฺตร\-ปจฺจยา ฯเปฯ อวิคตปจฺจยา.}

\csromanmarkh{1. ปจฺจยานุโลม}\csromanmarkhii{2. สงฺขฺยาวาร}\csromanheadernomark{2}{1. \textbf{ปจฺจยานุโลม 2. สงฺขฺยาวาร}}
\proseitem{4}{เหตุยา นว, อารมฺมเณ นว, อธิปติยา, นว, อนนฺตเร นว, (สพฺพตฺถ นว.) วิปาเก เอกํ, อาหาเร นว ฯเปฯ อวิคเต นว.}

\vspace{\tipitakaparskip}
\csromancenter{\textbf{อนุโลมํ.}}
\csromanmarkh{2. ปจฺจย\-ปจฺจนีย}\csromanmarkhii{1. วิภงฺควาร}\csromanheadernomark{2}{2. ปจฺจย\-ปจฺจนีย 1. วิภงฺควาร}
\proseitem{5}{สญฺโญชนํ ธมฺมํ ปฏิจฺจ สญฺโญชโน ธมฺโม อุปฺปชฺชติ นเหตุปจฺจยา. วิจิกิจฺฉา\-สญฺโญชนํ ปฏิจฺจ อวิชฺชา\-สญฺโญชนํ. (1)}

\prose{\csromanfolio{211}โนสญฺโญชนํ ธมฺมํ ปฏิจฺจ โนสญฺโญชโน ธมฺโม อุปฺปชฺชติ นเหตุปจฺจยา. อเหตุกํ โนสญฺโญชนํ เอกํ ขนฺธํ ปฏิจฺจ ตโย ขนฺธา จิตฺต\-สมุฏฺฐานญฺจ รูปํ ฯเปฯ เทฺว ขนฺเธ ฯเปฯ. (ยาว อสญฺญสตฺตา.) (1)}

\prose{โนสญฺโญชนํ ธมฺมํ ปฏิจฺจ สญฺโญชโน ธมฺโม อุปฺปชฺชติ นเหตุปจฺจยา. วิจิกิจฺฉา\-สหคเต อุทฺธจฺจ\-สหคเต ขนฺเธ ปฏิจฺจ อวิชฺชา\-สญฺโญชนํ. (2)}

\prose{สญฺโญชนญฺจ โนสญฺโญชนญฺจ ธมฺมํ ปฏิจฺจ สญฺโญชโน ธมฺโม อุปฺปชฺชติ นเหตุปจฺจยา. วิจิกิจฺฉาสญฺโญชนญฺจ สมฺปยุตฺตเก จ ขนฺเธ ปฏิจฺจ อวิชฺชา\-สญฺโญชนํ. (3) (สํขิตฺตํ. อาสว\-โคจฺฉก\-สทิสํ. นอารมฺมณาปิ สพฺเพ อุทฺธริตพฺพา.)}

\csromanmarkh{2. ปจฺจย\-ปจฺจนีย}\csromanmarkhii{2. สงฺขฺยาวาร}\csromanheadernomark{2}{2. ปจฺจย\-ปจฺจนีย 2. สงฺขฺยาวาร}
\proseitem{6}{นเหตุยา จตฺตาริ, นอารมฺมเณ ตีณิ, นอธิปติยา นว, นอนนฺตเร ตีณิ, นสมนนฺตเร ตีณิ, นอญฺญมญฺเญ ตีณิ, นอุปนิสฺสเย ตีณิ, นปุเรชาเต นว, นปจฺฉาชาเต นว, นอาเสวเน นว, นกมฺเม ตีณิ, นวิปาเก นว, นอาหาเร เอกํ, นอินฺทฺริเย เอกํ, นฌาเน เอกํ, นมคฺเค เอกํ, นสมฺปยุตฺเต ตีณิ, นวิปฺปยุตฺเต นว, โนนตฺถิยา ตีณิ, โนวิคเต ตีณิ.}

\vspace{\tipitakaparskip}
\csromancenter{ปจฺจนียํ.}
\proseitem{7}{เหตุปจฺจยา นอารมฺมเณ ตีณิ, นอธิปติยา นว. (เอวํ สพฺพํ คเณตพฺพํ.)}

\proseitem{8}{\csromanfolio{212}นเหตุปจฺจยา อารมฺมเณ จตฺตาริ, (สพฺพตฺถ จตฺตาริ.) วิปาเก เอกํ, อาหาเร จตฺตาริ ฯเปฯ มคฺเค ตีณิ, สมฺปยุตฺเต จตฺตาริ ฯเปฯ อวิคเต จตฺตาริ.}

\vspace{\tipitakaparskip}
\csromancenter{ปจฺจนียา\-นุโลมํ.}
\csromanmarkh{20. สญฺโญชนทุก}\csromanmarkhii{2. สหชาตวาร}\csromanheadernomark{2}{20. \textbf{สญฺโญชนทุก 2. สหชาตวาร}}
\proseitem{9}{สญฺโญชนํ ธมฺมํ สหชาโต สญฺโญชโน ธมฺโม อุปฺปชฺชติ เหตุปจฺจยา. (ปฏิจฺจ\-วาร\-สทิสํ.)}

\csromanmarkh{20. สญฺโญชนทุก}\csromanmarkhii{3. ปจฺจยวาร}\csromanheadernomark{2}{20. \textbf{สญฺโญชนทุก 3. ปจฺจยวาร}}
\csromanheader{2}{1. ปจฺจยา\-นุโลมาทิ}
\csromanheader{2}{\textbf{เหตุ}}
\proseitem{10}{สญฺโญชนํ ธมฺมํ ปจฺจยา สญฺโญชโน ธมฺโม อุปฺปชฺชติ เหตุปจฺจยา. ตีณิ. (ปฏิจฺจสทิสํ.)}

\prose{โนสญฺโญชนํ ธมฺมํ ปจฺจยา โนสญฺโญชโน ธมฺโม อุปฺปชฺชติ เหตุปจฺจยา. โนสญฺโญชนํ เอกํ ขนฺธํ ปจฺจยา ตโย ขนฺธา จิตฺต\-สมุฏฺฐานญฺจ รูปํ ฯเปฯ เทฺว ขนฺเธ ฯเปฯ. ปฏิสนฺธิ\-กฺขเณ ฯเปฯ. ขนฺเธ ปจฺจยา วตฺถุ, วตฺถุํ ปจฺจยา ขนฺธา. เอกํ มหาภูตํ ฯเปฯ มหาภูเต ปจฺจยา จิตฺต\-สมุฏฺฐานํ รูปํ, กฏตฺตารูปํ, อุปาทารูปํ. วตฺถุํ ปจฺจยา โนสญฺโญชนา ขนฺธา. (1)}

\prose{โนสญฺโญชนํ ธมฺมํ ปจฺจยา สญฺโญชโน ธมฺโม อุปฺปชฺชติ เหตุปจฺจยา. โนสญฺโญชเน ขนฺเธ ปจฺจยา สญฺโญชนา. วตฺถุํ ปจฺจยา สญฺโญชนา. (2)}

\prose{โนสญฺโญชนํ ธมฺมํ ปจฺจยา สญฺโญชโน จ โนสญฺโญชโน จ ธมฺมา อุปฺปชฺชนฺติ เหตุปจฺจยา. โนสญฺโญชนํ เอกํ ขนฺธํ ปจฺจยา ตโย ขนฺธา สญฺโญชนญฺจ จิตฺต\-สมุฏฺฐานญฺจ รูปํ ฯเปฯ เทฺว ขนฺเธ ฯเปฯ. วตฺถุํ ปจฺจยา สญฺโญชนา. มหาภูเต ปจฺจยา จิตฺต\-สมุฏฺฐานํ รูปํ. วตฺถุํ ปจฺจยา สญฺโญชนา สมฺปยุตฺตกา จ ขนฺธา. (3)}

\proseitem{11}{\csromanfolio{213}สญฺโญชนญฺจ โนสญฺโญชนญฺจ ธมฺมํ ปจฺจยา สญฺโญชโน ธมฺโม อุปฺปชฺชติ เหตุปจฺจยา. กามราคสญฺโญชนญฺจ สมฺปยุตฺตเก จ ขนฺเธ ปจฺจยา ทิฏฺฐิ\-สญฺโญชนํ อวิชฺชา\-สญฺโญชนํ. กามราคสญฺโญชนญฺจ วตฺถุญฺจ ปจฺจยา ทิฏฺฐิ\-สญฺโญชนํ อวิชฺชา\-สญฺโญชนํ. (จกฺกํ.) (1)}

\prose{สญฺโญชนญฺจ โนสญฺโญชนญฺจ ธมฺมํ ปจฺจยา โนสญฺโญชโน ธมฺโม อุปฺปชฺชติ เหตุปจฺจยา. โนสญฺโญชนํ เอกํ ขนฺธญฺจ สญฺโญชเน จ ปจฺจยา ตโย ขนฺธา จิตฺต\-สมุฏฺฐานญฺจ รูปํ ฯเปฯ เทฺว ขนฺเธ ฯเปฯ. (จกฺกํ.) สญฺโญชเน จ วตฺถุญฺจ ปจฺจยา โนสญฺโญชนา ขนฺธา. (2)}

\prose{สญฺโญชนญฺจ โนสญฺโญชนญฺจ ธมฺมํ ปจฺจยา สญฺโญชโน จ โนสญฺโญชโน จ ธมฺมา อุปฺปชฺชนฺติ เหตุปจฺจยา. โนสญฺโญชนํ เอกํ ขนฺธญฺจ กามราคสญฺโญชนญฺจ ปจฺจยา ตโย ขนฺธา ทิฏฺฐิ\-สญฺโญชนํ อวิชฺชา\-สญฺโญชนํ จิตฺต\-สมุฏฺฐานญฺจ รูปํ. เทฺว ขนฺเธ ฯเปฯ. (จกฺกํ.) กามราคสญฺโญชนญฺจ วตฺถุญฺจ ปจฺจยา ทิฏฺฐิ\-สญฺโญชนํ อวิชฺชา\-สญฺโญชนํ สมฺปยุตฺตกา จ ขนฺธา. (จกฺกํ. สํขิตฺตํ.) (3)}

\proseitem{12}{เหตุยา นว, อารมฺมเณ นว, (สพฺพตฺถ นว.) วิปาเก เอกํ ฯเปฯ อวิคเต นว.}

\proseitem{13}{นเหตุยา จตฺตาริ, (ยตฺถ ยตฺถ วตฺถุ ลพฺภติ, ตตฺถ ตตฺถ นินฺเนตพฺพํ.) นอารมฺมเณ ตีณิ ฯเปฯ โนวิคเต ตีณิ.}

\vspace{\tipitakaparskip}
\csromancenter{(เอวํ อิตเรปิ เทฺว คณนา จ นิสฺสยวาโร จ กาตพฺโพ.)}
\csromanmarkh{20. สญฺโญชนทุก}\csromanmarkhii{5. สํสฏฺฐวาร}\csromanheadernomark{2}{20. สญฺโญชนทุก 5. สํสฏฺฐวาร}
\csromanheader{2}{1. \textbf{ปจฺจยานุโลม}}
\vspace{\tipitakaparskip}
\csromancenter{\textbf{เหตุ}}
\proseitem{14}{สญฺโญชนํ ธมฺมํ สํสฏฺโฐ สญฺโญชโน ธมฺโม อุปฺปชฺชติ เหตุปจฺจยา. กามราค\-สญฺโญชนํ สํสฏฺฐํ ทิฏฺฐิ\-สญฺโญชนํ อวิชฺชา\-สญฺโญชนํ. (เอวํ นว ปญฺหา. อรูปาเยว กาตพฺพา.)}

\vspace{\tipitakaparskip}
\csromancenter{(สํสฏฺฐ\-วาโรปิ สมฺปยุตฺตวาโรปิ เอวํ กาตพฺพา.)}
\csromanmarkh{20. สญฺโญชนทุก}\csromanmarkhii{7. ปญฺหาวาร}\csromanheadernomark{2}{20. \textbf{สญฺโญชนทุก 7. ปญฺหาวาร}}
\csromanmarkh{1. ปจฺจยานุโลม}\csromanmarkhii{1. วิภงฺควาร}\csromanheadernomark{2}{1. \textbf{ปจฺจยานุโลม 1. วิภงฺควาร}}
\csromanheader{2}{\textbf{เหตุ}}
\proseitem{15}{\csromanfolio{214}สญฺโญชโน ธมฺโม สญฺโญชนสฺส ธมฺมสฺส เหตุปจฺจเยน ปจฺจโย. สญฺโญชนา เหตู สมฺปยุตฺตกานํ สญฺโญชนานํ เหตุปจฺจเยน ปจฺจโย. (1)}

\prose{สญฺโญชโน ธมฺโม โนสญฺโญชนสฺส ธมฺมสฺส เหตุปจฺจเยน ปจฺจโย. สญฺโญชนา เหตู สมฺปยุตฺตกานํ ขนฺธานํ จิตฺตสมุฏฺฐานานญฺจ รูปานํ เหตุปจฺจเยน ปจฺจโย. (2)}

\prose{สญฺโญชนา ธมฺโม สญฺโญชนสฺส จ โนสญฺโญชนสฺส จ ธมฺมสฺส เหตุปจฺจเยน ปจฺจโย. สญฺโญชนา เหตู สมฺปยุตฺตกานํ ขนฺธานํ สญฺโญชนานํ จิตฺตสมุฏฺฐานานญฺจ รูปานํ เหตุปจฺจเยน ปจฺจโย. (3)}

\proseitem{16}{โนสญฺโญชโน ธมฺโม โนสญฺโญชนสฺส ธมฺมสฺส เหตุปจฺจเยน ปจฺจโย. โนสญฺโญชนา เหตู สมฺปยุตฺตกานํ ขนฺธานํ จิตฺตสมุฏฺฐานานญฺจ รูปานํ เหตุปจฺจเยน ปจฺจโย. ปฏิสนฺธิ\-กฺขเณ ฯเปฯ. (1)}

\proseitem{17}{สญฺโญชโน ธมฺโม สญฺโญชนสฺส ธมฺมสฺส อารมฺมณ\-ปจฺจเยน ปจฺจโย. สญฺโญชเน อารพฺภ สญฺโญชนา อุปฺปชฺชนฺติ. (มูลํ กาตพฺพํ.) สญฺโญชเน อารพฺภ โนสญฺโญชนา ขนฺธา อุปฺปชฺชนฺติ. (มูลํ กาตพฺพํ.) สญฺโญชเน อารพฺภ สญฺโญชนา จ สมฺปยุตฺตกา จ ขนฺธา อุปฺปชฺชนฺติ. (3)}

\proseitem{18}{โนสญฺโญชโน ธมฺโม โนสญฺโญชนสฺส ธมฺมสฺส อารมฺมณ\-ปจฺจเยน ปจฺจโย. ทานํ, ทตฺวา, สีลํ ฯเปฯ อุโปสถกมฺมํ กตฺวา ตํ ปจฺจเวกฺขติ. ปุพฺเพ สุจิณฺณานิ ฯเปฯ. ฌานา ฯเปฯ. อริยา มคฺคา วุฏฺฐหิตฺวา มคฺคํ ปจฺจเวกฺขนฺติ, ผลํ ฯเปฯ. นิพฺพานํ ปจฺจเวกฺขนฺติ. นิพฺพานํ โคตฺรภุสฺส, โวทานสฺส, มคฺคสฺส, ผลสฺส, อาวชฺชนาย อารมฺมณ\-ปจฺจเยน ปจฺจโย. อริยา โนสญฺโญชเน ปหีเน กิเลเส ฯเปฯ. วิกฺขมฺภิเต กิเลเส ฯเปฯ. ปุพฺเพ ฯเปฯ. \csromanfolio{215}จกฺขุํ ฯเปฯ วตฺถุํ โนสญฺโญชเน ขนฺเธ อนิจฺจโต ฯเปฯ โทมนสฺสํ อุปฺปชฺชติ. ทิพฺเพน จกฺขุนา รูปํ ปสฺสติ. ทิพฺพาย โสตธาตุยา สทฺทํ สุณาติ. เจโตปริย\-ญาเณน โนสญฺโญชนจิตฺตสมงฺคิสฺส จิตฺตํ ชานาติ. อากาสา\-นญฺจา\-ยตนํ วิญฺญาณญฺจา\-ยตนสฺส ฯเปฯ. อากิญฺจญฺญา\-ยตนํ เนว\-สญฺญา\-นา\-สญฺญา\-ยตนสฺส ฯเปฯ. รูปายตนํ จกฺขุ\-วิญฺญาณสฺส ฯเปฯ. โผฏฺฐพฺพา\-ยตนํ กายวิญฺญาณสฺส ฯเปฯ. โนสญฺโญชนา ขนฺธา อิทฺธิวิธ\-ญาณสฺส, เจโตปริย\-ญาณสฺส, ปุพฺเพ\-นิวาสา\-นุสฺสติ\-ญาณสฺส, ยถากมฺมูปคญาณสฺส, อนาคตํสญาณสฺส, อาวชฺชนาย อารมฺมณ\-ปจฺจเยน ปจฺจโย. (1)}

\prose{โนสญฺโญชโน ธมฺโม สญฺโญชนสฺส ธมฺมสฺส อารมฺมณ\-ปจฺจเยน ปจฺจโย. ทานํ ฯเปฯ สีลํ ฯเปฯ อุโปสถกมฺมํ ฯเปฯ. ปุพฺเพ ฯเปฯ. ฌานา ฯเปฯ. จกฺขุํ ฯเปฯ วตฺถุํ โนสญฺโญชเน ขนฺเธ อสฺสาเทติ อภินนฺทติ, ตํ อารพฺภ ราโค อุปฺปชฺชติ ฯเปฯ โทมนสฺสํ อุปฺปชฺชติ. (2)}

\prose{โนสญฺโญชโน ธมฺโม สญฺโญชนสฺส จ โนสญฺโญชนสฺส จ ธมฺมสฺส อารมฺมณ\-ปจฺจเยน ปจฺจโย. ทานํ ทตฺวา, สีลํ ฯเปฯ อุโปสถกมฺมํ ฯเปฯ. ปุพฺเพ ฯเปฯ. ฌานา ฯเปฯ. จกฺขุํ ฯเปฯ วตฺถุํ โนสญฺโญชเน ขนฺเธ อสฺสาเทติ อภินนฺทติ, ตํ อารพฺภ สญฺโญชนา จ สญฺโญชน\-สมฺปยุตฺตกา จ ขนฺธา อุปฺปชฺชนฺติ. (3)}

\prose{สญฺโญชโน จ โนสญฺโญชโน จ ธมฺมา สญฺโญชนสฺส ธมฺมสฺส อารมฺมณ\-ปจฺจเยน ปจฺจโย. ตีณิ. (อารพฺภเยว กาตพฺพา.)}

\proseitem{19}{สญฺโญชโน ธมฺโม สญฺโญชนสฺส ธมฺมสฺส อธิปติ\-ปจฺจเยน ปจฺจโย. อารมฺมณา\-ธิปติ\csromandash{}สญฺโญชนํ ครุํ กตฺวา ฯเปฯ. ตีณิ. (ครุการมฺมณา)}

\prose{โนสญฺโญชโน ธมฺโม โนสญฺโญชนสฺส ธมฺมสฺส อธิปติ\-ปจฺจเยน ปจฺจโย. อารมฺมณา\-ธิปติ, สหชาตา\-ธิปติ.  อารมฺมณา\-ธิปติ\csromandash{}ทานํ ทตฺวา, สีลํ ฯเปฯ. ตีณิ. (ติณฺณมฺปิ อารมฺมณา\-ธิปติ. สหชาตา\-ธิปติปิ กาตพฺพา. วิภชิตพฺพา. ตีณิปิ.)}

\prose{\csromanfolio{216}สญฺโญชโน จ โนสญฺโญชโน จ ธมฺมา สญฺโญชนสฺส ธมฺมสฺส อธิปติ\-ปจฺจเยน ปจฺจโย. อารมฺมณา\-ธิปติ\csromandash{}สญฺโญชเน จ สมฺปยุตฺตเก จ ขนฺเธ ครุํ กตฺวา ฯเปฯ. ตีณิ.}

\proseitem{20}{สญฺโญชโน ธมฺโม สญฺโญชนสฺส ธมฺมสฺส อนนฺตร\-ปจฺจเยน ปจฺจโย. ปุริมา ปุริมา สญฺโญชนา ปจฺฉิมานํ ปจฺฉิมานํ สญฺโญชนานํ อนนฺตร\-ปจฺจเยน ปจฺจโย. ตีณิ.}

\prose{โนสญฺโญชโน ธมฺโม โนสญฺโญชนสฺส ธมฺมสฺส อนนฺตร\-ปจฺจเยน ปจฺจโย. ปุริมา ปุริมา โนสญฺโญชนา ขนฺธา ปจฺฉิมานํ ฯเปฯ ผลสมาปตฺติยา อนนฺตร\-ปจฺจเยน ปจฺจโย. (1)}

\prose{โนสญฺโญชโน ธมฺโม สญฺโญชนสฺส ธมฺมสฺส อนนฺตร\-ปจฺจเยน ปจฺจโย. ปุริมา ปุริมา โนสญฺโญชนา ขนฺธา ปจฺฉิมานํ ปจฺฉิมานํ สญฺโญชนานํ อนนฺตร\-ปจฺจเยน ปจฺจโย. อาวชฺชนา สญฺโญชนานํ อนนฺตร\-ปจฺจเยน ปจฺจโย. (เอวํ เทฺวปิ กาตพฺพา.) (3)}

\prose{สญฺโญชโน จ โนสญฺโญชโน จ ธมฺมา สญฺโญชนสฺส ธมฺมสฺส อนนฺตร\-ปจฺจเยน ปจฺจโย. ตีณิ.}

\prose{\textbf{สมนนฺตราทิ}}

\proseitem{21}{สญฺโญชโน ธมฺโม สญฺโญชนสฺส ธมฺมสฺส สมนนฺตร\-ปจฺจเยน ปจฺจโย. นว. สหชาต\-ปจฺจเยน ปจฺจโย. นว. อญฺญมญฺญ\-ปจฺจเยน ปจฺจโย. นว. นิสฺสย\-ปจฺจเยน ปจฺจโย. นว.}

\csromanheader{2}{\textbf{อุปนิสฺสย}}
\proseitem{22}{สญฺโญชโน ธมฺโม สญฺโญชนสฺส ธมฺมสฺส อุปนิสฺสย\-ปจฺจเยน ปจฺจโย. อารมฺมณู\-ปนิสฺสโย, อนนฺตรู\-ปนิสฺสโย, ปกตู\-ปนิสฺสโย ฯเปฯ. ปกตู\-ปนิสฺสโย\csromandash{}สญฺโญชนา สญฺโญชนานํ อุปนิสฺสย\-ปจฺจเยน ปจฺจโย. (เอวํ ตีณิปิ.)}

\proseitem{23}{\csromanfolio{217}โนสญฺโญชโน ธมฺโม โนสญฺโญชนสฺส ธมฺมสฺส อุปนิสฺสย\-ปจฺจเยน ปจฺจโย. อารมฺมณู\-ปนิสฺสโย, อนนฺตรู\-ปนิสฺสโย, ปกตู\-ปนิสฺสโย ฯเปฯ. ปกตู\-ปนิสฺสโย\csromandash{}สทฺธํ อุปนิสฺสาย ทานํ เทติ ฯเปฯ สมาปตฺติํ อุปฺปาเทติ. มานํ ชปฺเปติ. ทิฏฺฐิํ คณฺหาติ. สีลํ ฯเปฯ. ปญฺญํ. ราคํ ฯเปฯ ปตฺถนํ ฯเปฯ. เสนาสนํ อุปนิสฺสาย ทานํ เทติ ฯเปฯ สํฆํ ภินฺทติ. สทฺธา ฯเปฯ. เสนาสนํ สทฺธาย ฯเปฯ ผลสมาปตฺติยา อุปนิสฺสย\-ปจฺจเยน ปจฺจโย. (1)}

\prose{โนสญฺโญชโน ธมฺโม สญฺโญชนสฺส ธมฺมสฺส อุปนิสฺสย\-ปจฺจเยน ปจฺจโย. อารมฺมณู\-ปนิสฺสโย, อนนฺตรู\-ปนิสฺสโย, ปกตู\-ปนิสฺสโย ฯเปฯ. ปกตู\-ปนิสฺสโย\csromandash{}สทฺธํ อุปนิสฺสาย มานํ ชปฺเปติ. ทิฏฺฐิํ คณฺหาติ. สีลํ ฯเปฯ. เสนาสนํ อุปนิสฺสาย ปาณํ หนติ ฯเปฯ สํฆํ ภินฺทติ. สทฺธา ฯเปฯ. เสนาสนํ ราคสฺส ฯเปฯ ปตฺถนาย อุปนิสฺสย\-ปจฺจเยน ปจฺจโย. (2)}

\prose{โนสญฺโญชโน ธมฺโม สญฺโญชนสฺส จ โนสญฺโญชนสฺส จ ธมฺมสฺส อุปนิสฺสย\-ปจฺจเยน ปจฺจโย. อารมฺมณู\-ปนิสฺสโย, อนนฺตรู\-ปนิสฺสโย, ปกตู\-ปนิสฺสโย ฯเปฯ. ปกตู\-ปนิสฺสโย\csromandash{}สทฺธํ อุปนิสฺสาย มานํ ชปฺเปติ. ทิฏฺฐิํ คณฺหาติ. สีลํ ฯเปฯ. เสนาสนํ อุปนิสฺสาย ปาณํ หนติ ฯเปฯ สํฆํ ภินฺทติ. สทฺธา ฯเปฯ. เสนาสนํ สญฺโญชนานํ สมฺปยุตฺตกานญฺจ ขนฺธานํ อุปนิสฺสย\-ปจฺจเยน ปจฺจโย. (3)}

\prose{สญฺโญชโน จ โนสญฺโญชโน จ ธมฺมา สญฺโญชนสฺส ธมฺมสฺส อุปนิสฺสย\-ปจฺจเยน ปจฺจโย. ตีณิ.}

\proseitem{24}{โนสญฺโญชโน ธมฺโม โนสญฺโญชนสฺส ธมฺมสฺส ปุเรชาต\-ปจฺจเยน ปจฺจโย. อารมฺมณ\-ปุเรชาตํ, วตฺถุ\-ปุเรชาตํ.  อารมฺมณ\-ปุเรชาตํ\csromandash{}จกฺขุํ ฯเปฯ วตฺถุํ อนิจฺจโต ฯเปฯ โทมนสฺสํ อุปฺปชฺชติ. ทิพฺเพน จกฺขุนา รูปํ ปสฺสติ. ทิพฺพาย โสตธาตุยา สทฺทํ สุณาติ. รูปายตนํ จกฺขุ\-วิญฺญาณสฺส ฯเปฯ. โผฏฺฐพฺพา\-ยตนํ กายวิญฺญาณสฺส ฯเปฯ. วตฺถุ\-ปุเรชาตํ\csromandash{}จกฺขายตนํ จกฺขุ\-วิญฺญาณสฺส ฯเปฯ กายายตนํ กายวิญฺญาณสฺส ฯเปฯ. วตฺถุ โนสญฺโญชนานํ ขนฺธานํ ปุเรชาต\-ปจฺจเยน ปจฺจโย. (1)}

\prose{\csromanfolio{218}โนสญฺโญชโน ธมฺโม สญฺโญชนสฺส ธมฺมสฺส ปุเรชาต\-ปจฺจเยน ปจฺจโย. อารมฺมณ\-ปุเรชาตํ, วตฺถุ\-ปุเรชาตํ.  อารมฺมณ\-ปุเรชาตํ\csromandash{} จกฺขุํ ฯเปฯ วตฺถุํ อสฺสาเทติ อภินนฺทติ, ตํ อารพฺภ ราโค อุปฺปชฺชติ ฯเปฯ โทมนสฺสํ อุปฺปชฺชติ. วตฺถุ\-ปุเรชาตํ\csromandash{}วตฺถุ สญฺโญชนานํ ปุเรชาต\-ปจฺจเยน ปจฺจโย. (2)}

\prose{โนสญฺโญชโน ธมฺโม สญฺโญชนสฺส จ โนสญฺโญชนสฺส จ ธมฺมสฺส ปุเรชาต\-ปจฺจเยน ปจฺจโย. อารมฺมณ\-ปุเรชาตํ, วตฺถุ\-ปุเรชาตํ.  อารมฺมณ\-ปุเรชาตํ\csromandash{}จกฺขุํ ฯเปฯ วตฺถุํ อสฺสาเทติ อภินนฺทติ, ตํ อารพฺภ สญฺโญชนา จ สมฺปยุตฺตกา จ ขนฺธา อุปฺปชฺชนฺติ. วตฺถุ\-ปุเรชาตํ\csromandash{}วตฺถุ สญฺโญชนานํ สมฺปยุตฺตกานญฺจ ขนฺธานํ ปุเรชาต\-ปจฺจเยน ปจฺจโย. (3)}

\csromanheader{2}{\textbf{ปจฺฉาชาต}}
\proseitem{25}{สญฺโญชโน ธมฺโม โนสญฺโญชนสฺส ธมฺมสฺส ปจฺฉาชาต\-ปจฺจเยน ปจฺจโย. เอกํ.}

\prose{โนสญฺโญชโน ธมฺโม โนสญฺโญชนสฺส ธมฺมสฺส ปจฺฉาชาต\-ปจฺจเยน ปจฺจโย ฯเปฯ. (1)}

\prose{สญฺโญชโน จ โนสญฺโญชโน จ ธมฺมา โนสญฺโญชนสฺส ธมฺมสฺส ปจฺฉาชาต\-ปจฺจเยน ปจฺจโย ฯเปฯ. (1)}

\csromanheader{2}{\textbf{อาเสวน}}
\vspace{\tipitakaparskip}
\csromancenter{สญฺโญชโน ธมฺโม สญฺโญชนสฺส ธมฺมสฺส อาเสวน\-ปจฺจเยน ปจฺจโย. นว.}
\csromanheader{2}{\textbf{กมฺมาทิ}}
\proseitem{27}{โนสญฺโญชโน ธมฺโม โนสญฺโญชนสฺส ธมฺมสฺส กมฺมปจฺจเยน ปจฺจโย. ตีณิ. วิปาก\-ปจฺจเยน ปจฺจโย. เอกํ. อาหารปจฺจเยน ปจฺจโย. ตีณิ. อินฺทฺริย\-ปจฺจเยน ปจฺจโย. ตีณิ. ฌานปจฺจเยน ปจฺจโย. ตีณิ. มคฺคปจฺจเยน ปจฺจโย. นว. สมฺปยุตฺต\-ปจฺจเยน ปจฺจโย. นว.}

\csromanheader{2}{20. สญฺโญชนทุก \textbf{วิปฺปยุตฺต}}
\proseitem{28}{\csromanfolio{219}สญฺโญชโน ธมฺโม โนสญฺโญชนสฺส ธมฺมสฺส วิปฺปยุตฺต\-ปจฺจเยน ปจฺจโย. สหชาตํ, ปจฺฉาชาตํ. (วิภชิตพฺพํ.) (1)}

\prose{โนสญฺโญชโน ธมฺโม โนสญฺโญชนสฺส ธมฺมสฺส วิปฺปยุตฺต\-ปจฺจเยน ปจฺจโย. สหชาตํ, ปุเรชาตํ, ปจฺฉาชาตํ. (วิภชิตพฺพํ.) (1)}

\prose{โนสญฺโญชโน ธมฺโม สญฺโญชนสฺส ธมฺมสฺส วิปฺปยุตฺต\-ปจฺจเยน ปจฺจโย. ปุเรชาตํ วตฺถุ สญฺโญชนานํ วิปฺปยุตฺต\-ปจฺจเยน ปจฺจโย. (2)}

\prose{โนสญฺโญชโน ธมฺโม สญฺโญชนสฺส จ โนสญฺโญชนสฺส จ ธมฺมสฺส วิปฺปยุตฺต\-ปจฺจเยน ปจฺจโย. ปุเรชาตํ วตฺถุ สญฺโญชนานํ สมฺปยุตฺตกานญฺจ ขนฺธานํ วิปฺปยุตฺต\-ปจฺจเยน ปจฺจโย. (3)}

\prose{สญฺโญชโน จ โนสญฺโญชโน จ ธมฺมา โนสญฺโญชนสฺส ธมฺมสฺส วิปฺปยุตฺต\-ปจฺจเยน ปจฺจโย. สหชาตํ, ปจฺฉาชาตํ. (วิภชิตพฺพํ.) (1)}

\proseitem{29}{สญฺโญชโน ธมฺโม สญฺโญชนสฺส ธมฺมสฺส อตฺถิปจฺจเยน ปจฺจโย. เอกํ. (ปฏิจฺจสทิสํ.) (1)}

\prose{สญฺโญชโน ธมฺโม โนสญฺโญชนสฺส ธมฺมสฺส อตฺถิปจฺจเยน ปจฺจโย. สหชาตํ, ปจฺฉาชาตํ. (สํขิตฺตํ.) (2)}

\prose{สญฺโญชโน ธมฺโม สญฺโญชนสฺส จ โนสญฺโญชนสฺส จ ธมฺมสฺส อตฺถิปจฺจเยน ปจฺจโย. (ปฏิจฺจสทิสํ.) (3)}

\proseitem{30}{โนสญฺโญชโน ธมฺโม โนสญฺโญชนสฺส ธมฺมสฺส อตฺถิปจฺจเยน ปจฺจโย. สหชาตํ, ปุเรชาตํ, ปจฺฉาชาตํ, อาหารํ, อินฺทฺริยํ (สํขิตฺตํ.) (1)}

\prose{โนสญฺโญชโน ธมฺโม สญฺโญชนสฺส ธมฺมสฺส อตฺถิปจฺจเยน ปจฺจโย. สหชาตํ, ปุเรชาตํ.  สหชาตา โนสญฺโญชนา ขนฺธา สมฺปยุตฺตกานํ สญฺโญชนานํ อตฺถิปจฺจเยน ปจฺจโย. ปุเรชาตํ จกฺขุํ ฯเปฯ วตฺถุํ อสฺสาเทติ อภินนฺทติ, ตํ อารพฺภ ราโค อุปฺปชฺชติ ฯเปฯ โทมนสฺสํ อุปฺปชฺชติ. วตฺถุ อญฺโญชนานํ อตฺถิปจฺจเยน ปจฺจโย. (2)}

\prose{\csromanfolio{220}โนสญฺโญชโน ธมฺโม สญฺโญชนสฺส จ โนสญฺโญชนสฺส จ ธมฺมสฺส อตฺถิปจฺจเยน ปจฺจโย. สหชาตํ, ปุเรชาตํ.  สหชาโต โนสญฺโญชโน ฯเปฯ. (สํขิตฺตํ. อาสวสทิสํ.) (3)}

\proseitem{31}{สญฺโญชโน จ โนสญฺโญชโน จ ธมฺมา สญฺโญชนสฺส ธมฺมสฺส อตฺถิปจฺจเยน ปจฺจโย. สหชาตํ, ปุเรชาตํ. (อาสวสทิสํ.) (1)}

\prose{สญฺโญชโน จ โนสญฺโญชโน จ ธมฺมา โนสญฺโญชนสฺส ธมฺมสฺส อตฺถิปจฺจเยน ปจฺจโย. สหชาตํ, ปุเรชาตํ, ปจฺฉาชาตํ, อาหารํ, อินฺทฺริยํ. (วิภชิตพฺพํ. อาสวสทิสํ.) (2)}

\prose{สญฺโญชโน จ โนสญฺโญชโน จ ธมฺมา สญฺโญชนสฺส จ โนสญฺโญชนสฺส จ ธมฺมสฺส อตฺถิปจฺจเยน ปจฺจโย. สหชาตํ, ปุเรชาตํ. (วิภชิตพฺพํ. อาสวสทิสํ.) (3)}

\csromanmarkh{1. ปจฺจยานุโลม}\csromanmarkhii{2. สงฺขฺยาวาร}\csromanheadernomark{2}{1. \textbf{ปจฺจยานุโลม 2. สงฺขฺยาวาร}}
\proseitem{32}{เหตุยา จตฺตาริ, อารมฺมเณ นว, อธิปติยา นว, อนนฺตเร นว, สมนนฺตเร นว, สหชาเต นว, อญฺญมญฺเญ นว, นิสฺสเย นว, อุปนิสฺสเย นว, ปุเรชาเต ตีณิ, ปจฺฉาชาเต ตีณิ, อาเสวเน นว, กมฺเม ตีณิ, วิปาเก เอกํ, อาหาเร ตีณิ, อินฺทฺริเย ตีณิ, ฌาเน ตีณิ, มคฺเค นว, สมฺปยุตฺเต นว, วิปฺปยุตฺเต ปญฺจ, อตฺถิยา นว, นตฺถิยา นว, วิคเต นว, อวิคเต นว.}

\vspace{\tipitakaparskip}
\csromancenter{อนุโลมํ.}
\proseitem{33}{สญฺโญชโน ธมฺโม สญฺโญชนสฺส ธมฺมสฺส อารมฺมณ\-ปจฺจเยน ปจฺจโย. สหชาต\-ปจฺจเยน ปจฺจโย. อุปนิสฺสย\-ปจฺจเยน ปจฺจโย. (1)}

\prose{สญฺโญชโน ธมฺโม โนสญฺโญชนสฺส ธมฺมสฺส อารมฺมณ\-ปจฺจเยน ปจฺจโย. สหชาต\-ปจฺจเยน ปจฺจโย. อุปนิสฺสย\-ปจฺจเยน ปจฺจโย. ปจฺฉาชาต\-ปจฺจเยน ปจฺจโย. (2)}

\prose{\csromanfolio{221}สญฺโญชโน ธมฺโม สญฺโญชนสฺส จ โนสญฺโญชนสฺส จ ธมฺมสฺส อารมฺมณ\-ปจฺจเยน ปจฺจโย. สหชาต\-ปจฺจเยน ปจฺจโย. อุปนิสฺสย\-ปจฺจเยน ปจฺจโย. (3)}

\proseitem{34}{โนสญฺโญชโน ธมฺโม โนสญฺโญชนสฺส ธมฺมสฺส อารมฺมณ\-ปจฺจเยน ปจฺจโย. สหชาต\-ปจฺจเยน ปจฺจโย. อุปนิสฺสย\-ปจฺจเยน ปจฺจโย. ปุเรชาต\-ปจฺจเยน ปจฺจโย. ปจฺฉาชาต\-ปจฺจเยน ปจฺจโย. กมฺมปจฺจเยน ปจฺจโย. อาหารปจฺจเยน ปจฺจโย. อินฺทฺริย\-ปจฺจเยน ปจฺจโย. (1)}

\prose{โนสญฺโญชโน ธมฺโม สญฺโญชนสฺส ธมฺมสฺส อารมฺมณ\-ปจฺจเยน ปจฺจโย. สหชาต\-ปจฺจเยน ปจฺจโย. อุปนิสฺสย\-ปจฺจเยน ปจฺจโย. ปุเรชาต\-ปจฺจเยน ปจฺจโย. (2)}

\prose{โนสญฺโญชโน ธมฺโม สญฺโญชนสฺส จ โนสญฺโญชนสฺส จ ธมฺมสฺส อารมฺมณ\-ปจฺจเยน ปจฺจโย. สหชาต\-ปจฺจเยน ปจฺจโย. อุปนิสฺสย\-ปจฺจเยน ปจฺจโย. ปุเรชาต\-ปจฺจเยน ปจฺจโย. (3)}

\proseitem{35}{สญฺโญชโน จ โนสญฺโญชโน จ ธมฺมา สญฺโญชนสฺส ธมฺมสฺส อารมฺมณ\-ปจฺจเยน ปจฺจโย. สหชาต\-ปจฺจเยน ปจฺจโย. อุปนิสฺสย\-ปจฺจเยน ปจฺจโย. (1)}

\prose{สญฺโญชโน จ โนสญฺโญชโน จ ธมฺมา โนสญฺโญชนสฺส ธมฺมสฺส อารมฺมณ\-ปจฺจเยน ปจฺจโย. สหชาต\-ปจฺจเยน ปจฺจโย. อุปนิสฺสย\-ปจฺจเยน ปจฺจโย. ปจฺฉาชาต\-ปจฺจเยน ปจฺจโย. (2)}

\prose{สญฺโญชโน จ โนสญฺโญชโน จ ธมฺมา สญฺโญชนสฺส จ โนสญฺโญชนสฺส จ ธมฺมสฺส อารมฺมณ\-ปจฺจเยน ปจฺจโย. สหชาต\-ปจฺจเยน ปจฺจโย. อุปนิสฺสย\-ปจฺจเยน ปจฺจโย. (3)}

\csromanmarkh{2. ปจฺจย\-ปจฺจนีย}\csromanmarkhii{2. สงฺขฺยาวาร}\csromanheadernomark{2}{2. ปจฺจย\-ปจฺจนีย 2. สงฺขฺยาวาร}
\vspace{\tipitakaparskip}
\csromancenter{นเหตุยา นว, นอารมฺมเณ นว, (สพฺพตฺถ นว.) โนอวิคเต นว.}
\csromancenter{\textbf{เหตุ}ปจฺจยา นอารมฺมเณ จตฺตาริ ฯเปฯ นสมนนฺตเร จตฺตาริ, นอญฺญมญฺเญ เทฺว, นฺเอาปนิสฺสเย จตฺตาริ ฯเปฯ นมคฺเค จตฺตาริ, นสมฺปยุตฺเต เทฺว, นวิปฺปยุตฺเต จตฺตาริ, โนนตฺถิยา จตฺตาริ, โนวิคเต จตฺตาริ.}
\csromanheader{2}{\textbf{นเหตุทุก}}
\csromanheader{2}{38. นเหตุปจฺจยา อารมฺมเณ นว, อธิปติยา นว, (อนุโลมมาติกา) อวิคเต นว.}
\nitthitam{สญฺโญชนทุกํ นิฏฺฐิตํ.}
\csromanmatikaanchor{mtk.26}
\csromantocmarkref{section}{21. สญฺโญชนิยทุก}{mtk.26}
\bookmark[dest={mtk.26},level=1]{21. สญฺโญชนิยทุก}
\csromanheader{1}{21. \textbf{สญฺโญชนิยทุก 1-7. ปฏิจฺจาทิวาร}}
\proseitem{39}{\csromanfolio{222}สญฺโญชนิยํ ธมฺมํ ปฏิจฺจ สญฺโญชนิโย ธมฺโม อุปฺปชฺชติ เหตุปจฺจยา. สญฺโญชนิยํ เอกํ ขนฺธํ ปฏิจฺจ ตโย ขนฺธา จิตฺต\-สมุฏฺฐานญฺจ รูปํ ฯเปฯ เทฺว ขนฺเธ ฯเปฯ. ปฏิสนฺธิ\-กฺขเณ ฯเปฯ. ขนฺเธ ปฏิจฺจ วตฺถุ, วตฺถุํ ปฏิจฺจ ขนฺธา. เอกํ มหาภูตํ ฯเปฯ.}

\vspace{\tipitakaparskip}
\csromancenter{(จูฬนฺตรทุเก โลกิย\-ทุก\-สทิสํ, นินฺนานากรณํ.)}
\nitthitam{สญฺโญชนิยทุกํ นิฏฺฐิตํ.}
\csromanmatikaanchor{mtk.27}
\csromantocmarkref{section}{22. สญฺโญชนสมฺปยุตฺตทุก}{mtk.27}
\bookmark[dest={mtk.27},level=1]{22. สญฺโญชนสมฺปยุตฺตทุก}
\csromanmarkh{22. สญฺโญชน\-สมฺปยุตฺตทุก}\csromanmarkhii{1. ปฏิจฺจวาร}\csromanheadernomark{1}{22. สญฺโญชน\-สมฺปยุตฺตทุก 1. ปฏิจฺจวาร}
\csromanmarkh{1. ปจฺจยานุโลม}\csromanmarkhii{1. วิภงฺควาร}\csromanheadernomark{2}{1. \textbf{ปจฺจยานุโลม 1. วิภงฺควาร}}
\csromanheader{2}{\textbf{เหตุ}}
\proseitem{40}{สญฺโญชน\-สมฺปยุตฺตํ ธมฺมํ ปฏิจฺจ สญฺโญชน\-สมฺปยุตฺโต ธมฺโม อุปฺปชฺชติ เหตุปจฺจยา. สญฺโญชน\-สมฺปยุตฺตํ เอกํ ขนฺธํ ปฏิจฺจ ตโย ขนฺธา ฯเปฯ เทฺว ขนฺเธ ฯเปฯ. (1)}

\prose{\csromanfolio{223}สญฺโญชน\-สมฺปยุตฺตํ ธมฺมํ ปฏิจฺจ สญฺโญชน\-วิปฺปยุตฺโต ธมฺโม อุปฺปชฺชติ เหตุปจฺจยา. สญฺโญชน\-สมฺปยุตฺเต ขนฺเธ ปฏิจฺจ จิตฺต\-สมุฏฺฐานํ รูปํ. (2)}

\prose{สญฺโญชน\-สมฺปยุตฺตํ ธมฺมํ ปฏิจฺจ สญฺโญชน\-สมฺปยุตฺโต จ สญฺโญชน\-วิปฺปยุตฺโต จ ธมฺมา อุปฺปชฺชนฺติ เหตุปจฺจยา. สญฺโญชน\-สมฺปยุตฺตํ เอกํ ขนฺธํ ปฏิจฺจ ตโย ขนฺธา จิตฺต\-สมุฏฺฐานญฺจ รูปํ ฯเปฯ เทฺว ขนฺเธ ฯเปฯ. (3)}

\proseitem{41}{สญฺโญชน\-วิปฺปยุตฺตํ ธมฺมํ ปฏิจฺจ สญฺโญชน\-วิปฺปยุตฺโต ธมฺโม อุปฺปชฺชติ เหตุปจฺจยา. สญฺโญชน\-วิปฺปยุตฺตํ เอกํ ขนฺธํ ปฏิจฺจ ตโย ขนฺธา จิตฺต\-สมุฏฺฐานญฺจ รูปํ ฯเปฯ เทฺว ขนฺเธ ฯเปฯ. อุทฺธจฺจ\-สหคตํ โมหํ ปฏิจฺจ จิตฺต\-สมุฏฺฐานํ รูปํ. ปฏิสนฺธิ\-กฺขเณ ฯเปฯ. เอกํ มหาภูตํ ฯเปฯ. (1)}

\prose{สญฺโญชน\-วิปฺปยุตฺตํ ธมฺมํ ปฏิจฺจ สญฺโญชน\-สมฺปยุตฺโต ธมฺโม อุปฺปชฺชติ เหตุปจฺจยา. อุทฺธจฺจ\-สหคตํ โมหํ ปฏิจฺจ สมฺปยุตฺตกา ขนฺธา. (2)}

\prose{สญฺโญชน\-วิปฺปยุตฺตํ ธมฺมํ ปฏิจฺจ สญฺโญชน\-สมฺปยุตฺโต จ สญฺโญชน\-วิปฺปยุตฺโต จ ธมฺมา อุปฺปชฺชนฺติ เหตุปจฺจยา. อุทฺธจฺจ\-สหคตํ โมหํ ปฏิจฺจ สมฺปยุตฺตกา ขนฺธา จิตฺต\-สมุฏฺฐานญฺจ รูปํ. (3)}

\proseitem{42}{สญฺโญชน\-สมฺปยุตฺตญฺจ สญฺโญชน\-วิปฺปยุตฺตญฺจ ธมฺมํ ปฏิจฺจ สญฺโญชน\-สมฺปยุตฺโต ธมฺโม อุปฺปชฺชติ เหตุปจฺจยา. อุทฺธจฺจ\-สหคตํ เอกํ ขนฺธญฺจ โมหญฺจ ปฏิจฺจ ตโย ขนฺธา ฯเปฯ เทฺว ขนฺเธ ฯเปฯ. (1)}

\prose{สญฺโญชน\-สมฺปยุตฺตญฺจ สญฺโญชน\-วิปฺปยุตฺตญฺจ ธมฺมํ ปฏิจฺจ สญฺโญชน\-วิปฺปยุตฺโต ธมฺโม อุปฺปชฺชติ เหตุปจฺจยา. สญฺโญชน\-สมฺปยุตฺเต ขนฺเธ จ มหาภูเต จ ปฏิจฺจ จิตฺต\-สมุฏฺฐานํ รูปํ. อุทฺธจฺจ\-สหคเต ขนฺเธ จ โมหญฺจ ปฏิจฺจ จิตฺต\-สมุฏฺฐานํ รูปํ.}

\prose{สญฺโญชน\-สมฺปยุตฺตญฺจ สญฺโญชน\-วิปฺปยุตฺตญฺจ ธมฺมํ ปฏิจฺจ สญฺโญชน\-สมฺปยุตฺโต จ สญฺโญชน\-วิปฺปยุตฺโต จ ธมฺมา อุปฺปชฺชนฺติ เหตุปจฺจยา. อุทฺธจฺจ\-สหคตํ เอกํ ขนฺธญฺจ โมหญฺจ ปฏิจฺจ ตโย ขนฺธา จิตฺต\-สมุฏฺฐานญฺจ รูปํ ฯเปฯ เทฺว ขนฺเธ ฯเปฯ. (3)}

\proseitem{43}{\csromanfolio{224}สญฺโญชน\-สมฺปยุตฺตํ ธมฺมํ ปฏิจฺจ สญฺโญชน\-สมฺปยุตฺโต ธมฺโม อุปฺปชฺชติ อารมฺมณ\-ปจฺจยา. สญฺโญชน\-สมฺปยุตฺตํ เอกํ ขนฺธํ ฯเปฯ เทฺว ขนฺเธ ฯเปฯ. (1)}

\prose{สญฺโญชน\-สมฺปยุตฺตํ ธมฺมํ ปฏิจฺจ สญฺโญชน\-วิปฺปยุตฺโต ธมฺโม อุปฺปชฺชติ อารมฺมณ\-ปจฺจยา. อุทฺธจฺจ\-สหคเต ขนฺเธ ปฏิจฺจ อุทฺธจฺจ\-สหคโต โมโห. (2)}

\prose{สญฺโญชน\-สมฺปยุตฺตํ ธมฺมํ ปฏิจฺจ สญฺโญชน\-สมฺปยุตฺโต จ สญฺโญชน\-วิปฺปยุตฺโต จ ธมฺมา อุปฺปชฺชนฺติ อารมฺมณ\-ปจฺจยา. อุทฺธจฺจ\-สหคตํ เอกํ ขนฺธํ ปฏิจฺจ ตโย ขนฺธา โมโห จ ฯเปฯ เทฺว ขนฺเธ ฯเปฯ. (3)}

\proseitem{44}{สญฺโญชน\-วิปฺปยุตฺตํ ธมฺมํ ปฏิจฺจ สญฺโญชน\-วิปฺปยุตฺโต ธมฺโม อุปฺปชฺชติ อารมฺมณ\-ปจฺจยา. สญฺโญชน\-วิปฺปยุตฺตํ เอกํ ขนฺธํ ปฏิจฺจ ตโย ขนฺธา ฯเปฯ เทฺว ขนฺเธ ฯเปฯ. ปฏิสนฺธิ\-กฺขเณ ฯเปฯ. วตฺถุํ ปฏิจฺจ ขนฺธา. (1)}

\prose{สญฺโญชน\-วิปฺปยุตฺตํ ธมฺมํ ปฏิจฺจ สญฺโญชน\-สมฺปยุตฺโต ธมฺโม อุปฺปชฺชติ อารมฺมณ\-ปจฺจยา. อุทฺธจฺจ\-สหคตํ โมหํ ปฏิจฺจ สมฺปยุตฺตกา ขนฺธา. (2)}

\prose{สญฺโญชน\-สมฺปยุตฺตญฺจ สญฺโญชน\-วิปฺปยุตฺตญฺจ ธมฺมํ ปฏิจฺจ สญฺโญชน\-สมฺปยุตฺโต ธมฺโม อุปฺปชฺชติ อารมฺมณ\-ปจฺจยา. อุทฺธจฺจ\-สหคตํ เอกํ ขนฺธญฺจ โมหญฺจ ปฏิจฺจ ตโย ขนฺธา ฯเปฯ เทฺว ขนฺเธ ฯเปฯ. (1)}

\proseitem{45}{สญฺโญชน\-สมฺปยุตฺตํ ธมฺมํ ปฏิจฺจ สญฺโญชน\-สมฺปยุตฺโต ธมฺโม อุปฺปชฺชติ อธิปติ\-ปจฺจยา. ตีณิ.}

\prose{สญฺโญชน\-วิปฺปยุตฺตํ ธมฺมํ ปฏิจฺจ สญฺโญชน\-วิปฺปยุตฺโต ธมฺโม อุปฺปชฺชติ อธิปติ\-ปจฺจยา. เอกํ.}

\prose{สญฺโญชน\-สมฺปยุตฺตญฺจ สญฺโญชน\-วิปฺปยุตฺตญฺจ ธมฺมํ ปฏิจฺจ สญฺโญชน\-วิปฺปยุตฺโต ธมฺโม อุปฺปชฺชติ อธิปติ\-ปจฺจยา. สญฺโญชน\-สมฺปยุตฺตเก ขนฺเธ จ มหาภูเต จ ปฏิจฺจ จิตฺต\-สมุฏฺฐานํ รูปํ. (1) (สํขิตฺตํ.)}

\csromanmarkh{22. สญฺโญชน\-สมฺปยุตฺตทุก}\csromanmarkhii{1. ปจฺจยานุโลม 2. สงฺขฺยาวาร}\csromanheadernomark{2}{22. สญฺโญชน\-สมฺปยุตฺตทุก 1. ปจฺจยานุโลม 2. สงฺขฺยาวาร}
\proseitem{46}{\csromanfolio{225}เหตุยา นว, อารมฺมเณ ฉ, อธิปติยา ปญฺจ, อนนฺตเร ฉ, สมนนฺตเร ฉ, สหชาเต นว, อญฺญมญฺเญ ฉ, นิสฺสเย นว, อุปนิสฺสเย ฉ, ปุเรชาเต ฉ, อาเสวเน ฉ, กมฺเม นว, วิปาเก เอกํ, อาหาเร นว, อินฺทฺริเย นว, ฌาเน นว, มคฺเค นว, สมฺปยุตฺเต ฉ, วิปฺปยุตฺเต นว, อตฺถิยา นว, นตฺถิยา ฉ, วิคเต ฉ, อวิคเต นว.}

\vspace{\tipitakaparskip}
\csromancenter{อนุโลมํ.}
\csromanmarkh{2. ปจฺจย\-ปจฺจนีย}\csromanmarkhii{1. วิภงฺควาร}\csromanheadernomark{2}{2. ปจฺจย\-ปจฺจนีย 1. วิภงฺควาร}
\proseitem{47}{สญฺโญชน\-สมฺปยุตฺตํ ธมฺมํ ปฏิจฺจ สญฺโญชน\-สมฺปยุตฺโต ธมฺโม อุปฺปชฺชติ นเหตุปจฺจยา. วิจิกิจฺฉา\-สหคเต ขนฺเธ ปฏิจฺจ วิจิกิจฺฉา\-สหคโต โมโห. (1)}

\prose{สญฺโญชน\-สมฺปยุตฺตํ ธมฺมํ ปฏิจฺจ สญฺโญชน\-วิปฺปยุตฺโต ธมฺโม อุปฺปชฺชติ นเหตุปจฺจยา. อุทฺธจฺจ\-สหคเต ขนฺเธ ปฏิจฺจ อุทฺธจฺจ\-สหคโต โมโห. (2)}

\prose{สญฺโญชน\-วิปฺปยุตฺตํ ธมฺมํ ปฏิจฺจ สญฺโญชน\-วิปฺปยุตฺโต ธมฺโม อุปฺปชฺชติ นเหตุปจฺจยา. อเหตุกํ สญฺโญชน\-วิปฺปยุตฺตํ เอกํ ขนฺธํ ฯเปฯ. (ยาว อสญฺญสตฺตา. สํขิตฺตํ.) (1)}

\csromanmarkh{2. ปจฺจย\-ปจฺจนีย}\csromanmarkhii{2. สงฺขฺยาวาร}\csromanheadernomark{2}{2. ปจฺจย\-ปจฺจนีย 2. สงฺขฺยาวาร}
\vspace{\tipitakaparskip}
\csromancenter{\textbf{นเหตุ}ยา ตีณิ, นอารมฺมเณ ตีณิ, นอธิปติยา นว, นอนนฺตเร ตีณิ, นสมนนฺตเร ตีณิ, นอญฺญมญฺเญ ตีณิ, นฺเอาปนิสฺสเย ตีณิ, นปุเรชาเต สตฺต, นปจฺฉาชาเต นว, นอาเสวเน นว, นกมฺเม}
\prosecont{\csromanfolio{226}จตฺตาริ, นวิปาเก นว, นอาหาเร เอกํ, นอินฺทฺริเย เอกํ, นฌาเน เอกํ, นมคฺเค เอกํ, นสมฺปยุตฺเต ตีณิ, นวิปฺปยุตฺเต ฉ, โนนตฺถิยา ตีณิ, โนวิคเต ตีณิ.}

\vspace{\tipitakaparskip}
\csromancenter{ปจฺจนียํ.}
\csromancenter{\textbf{เหตุ}ปจฺจยา นอารมฺมเณ ตีณิ, นอธิปติยา นว ฯเปฯ นฺเอาปนิสฺสเย ตีณิ, นปุเรชาเต ฉ, นปจฺฉาชาเต นว, นอาเสวเน นว, นกมฺเม จตฺตาริ, นวิปาเก นว, นสมฺปยุตฺเต ตีณิ, นวิปฺปยุตฺเต จตฺตาริ, โนนตฺถิยา ตีณิ, โนวิคเต ตีณิ.}
\proseitem{48}{นเหตุปจฺจยา อารมฺมเณ ตีณิ ฯเปฯ วิปาเก เอกํ, อาหาเร ตีณิ ฯเปฯ มคฺเค เทฺว ฯเปฯ อวิคเต ตีณิ.}

\vspace{\tipitakaparskip}
\csromancenter{(สหชาตวาโร ปฏิจฺจ\-วาร\-สทิโส.)}
\csromanmarkh{22. สญฺโญชน\-สมฺปยุตฺตทุก}\csromanmarkhii{3. ปจฺจยวาร}\csromanheadernomark{2}{22. สญฺโญชน\-สมฺปยุตฺตทุก 3. ปจฺจยวาร}
\csromanmarkh{1. ปจฺจยานุโลม}\csromanmarkhii{1. วิภงฺควาร}\csromanheadernomark{2}{1. \textbf{ปจฺจยานุโลม 1. วิภงฺควาร}}
\csromanheader{2}{\textbf{เหตุ}}
\proseitem{51}{สญฺโญชน\-สมฺปยุตฺตํ ธมฺมํ ปจฺจยา สญฺโญชน\-สมฺปยุตฺโต ธมฺโม อุปฺปชฺชติ เหตุปจฺจยา. ตีณิ. (ปฏิจฺจสทิสํ.)}

\prose{\csromanfolio{227}สญฺโญชน\-วิปฺปยุตฺตํ ธมฺมํ ปจฺจยา สญฺโญชน\-วิปฺปยุตฺโต ธมฺโม อุปฺปชฺชติ เหตุปจฺจยา. (ยาว ปฏิสนฺธิ.) เอกํ มหาภูตํ ฯเปฯ วตฺถุํ ปจฺจยา สญฺโญชน\-วิปฺปยุตฺตา ขนฺธา. วตฺถุํ ปจฺจยา อุทฺธจฺจ\-สหคโต โมโห. (1)}

\prose{สญฺโญชน\-วิปฺปยุตฺตํ ธมฺมํ ปจฺจยา สญฺโญชน\-สมฺปยุตฺโต ธมฺโม อุปฺปชฺชติ เหตุปจฺจยา. วตฺถุํ ปจฺจยา สญฺโญชน\-สมฺปยุตฺตกา ขนฺธา. อุทฺธจฺจ\-สหคตํ โมหํ ปจฺจยา สมฺปยุตฺตกา ขนฺธา. (2)}

\prose{สญฺโญชน\-วิปฺปยุตฺตํ ธมฺมํ ปจฺจยา สญฺโญชน\-สมฺปยุตฺโต จ สญฺโญชน\-วิปฺปยุตฺโต จ ธมฺมา อุปฺปชฺชนฺติ เหตุปจฺจยา. วตฺถุํ ปจฺจยา สญฺโญชน\-สมฺปยุตฺตกา ขนฺธา. มหาภูเต ปจฺจยา จิตฺต\-สมุฏฺฐานํ รูปํ. อุทฺธจฺจ\-สหคตํ โมหํ ปจฺจยา สมฺปยุตฺตกา ขนฺธา จิตฺต\-สมุฏฺฐานญฺจ รูปํ. (3)}

\proseitem{52}{สญฺโญชน\-สมฺปยุตฺตญฺจ สญฺโญชน\-วิปฺปยุตฺตญฺจ ธมฺมํ ปจฺจยา สญฺโญชน\-สมฺปยุตฺโต ธมฺโม อุปฺปชฺชติ เหตุปจฺจยา. สญฺโญชน\-สมฺปยุตฺตํ เอกํ ขนฺธญฺจ วตฺถุญฺจ ปจฺจยา ตโย ขนฺธา ฯเปฯ เทฺว ขนฺเธ ฯเปฯ. อุทฺธจฺจ\-สหคตํ เอกํ ขนฺธญฺจ โมหญฺจ ปจฺจยา ตโย ขนฺธา ฯเปฯ เทฺว ขนฺเธ ฯเปฯ. (1)}

\prose{สญฺโญชน\-สมฺปยุตฺตญฺจ สญฺโญชน\-วิปฺปยุตฺตญฺจ ธมฺมํ ปจฺจยา สญฺโญชน\-วิปฺปยุตฺโต ธมฺโม อุปฺปชฺชติ เหตุปจฺจยา. สญฺโญชน\-สมฺปยุตฺเต ขนฺเธ จ มหาภูเต จ ปจฺจยา จิตฺต\-สมุฏฺฐานํ รูปํ. อุทฺธจฺจ\-สหคเต ขนฺเธ จ โมหญฺจ ปจฺจยา จิตฺต\-สมุฏฺฐานํ รูปํ.}

\prose{สญฺโญชน\-สมฺปยุตฺตญฺจ สญฺโญชน\-วิปฺปยุตฺตญฺจ ธมฺมํ ปจฺจยา สญฺโญชน\-สมฺปยุตฺโต จ สญฺโญชน\-วิปฺปยุตฺโต จ ธมฺมา อุปฺปชฺชนฺติ เหตุปจฺจยา. สญฺโญชน\-สมฺปยุตฺตํ เอกํ ขนฺธญฺจ วตฺถุญฺจ ปจฺจยา ตโย ขนฺธา ฯเปฯ เทฺว ขนฺเธ ฯเปฯ. สญฺโญชน\-สมฺปยุตฺเต ขนฺเธ จ มหาภูเต จ ปจฺจยา จิตฺต\-สมุฏฺฐานํ รูปํ. อุทฺธจฺจ\-สหคตํ เอกํ ขนฺธญฺจ โมหญฺจ ปจฺจยา ตโย ขนฺธา จิตฺต\-สมุฏฺฐานญฺจ รูปํ ฯเปฯ เทฺว ขนฺเธ ฯเปฯ. (3)}

\proseitem{53}{สญฺโญชน\-สมฺปยุตฺตํ ธมฺมํ ปจฺจยา สญฺโญชน\-สมฺปยุตฺโต ธมฺโม อุปฺปชฺชติ อารมฺมณ\-ปจฺจยา. ตีณิ. (ปฏิจฺจสทิสา.)}

\prose{สญฺโญชน\-วิปฺปยุตฺตํ ธมฺมํ ปจฺจยา สญฺโญชน\-วิปฺปยุตฺโต ธมฺโม อุปฺปชฺชติ อารมฺมณ\-ปจฺจยา. (ยาว ปฏิสนฺธิ.) วตฺถุํ ปจฺจยา สญฺโญชน\-วิปฺปยุตฺตา \csromanfolio{228}ขนฺธา. จกฺขายตนํ ปจฺจยา จกฺขุวิญฺญาณํ ฯเปฯ. กายายตนํ ปจฺจยา กายวิญฺญาณํ. วตฺถุํ ปจฺจยา สญฺโญชน\-วิปฺปยุตฺตา ขนฺธา. วตฺถุํ ปจฺจยา อุทฺธจฺจ\-สหคโต โมโห. (1)}

\prose{สญฺโญชน\-วิปฺปยุตฺตํ ธมฺมํ ปจฺจยา สญฺโญชน\-สมฺปยุตฺโต ธมฺโม อุปฺปชฺชติ อารมฺมณ\-ปจฺจยา. วตฺถุํ ปจฺจยา สญฺโญชน\-สมฺปยุตฺตกา ขนฺธา. อุทฺธจฺจ\-สหคตํ โมหํ ปจฺจยา สมฺปยุตฺตกา ขนฺธา. (2)}

\prose{สญฺโญชน\-วิปฺปยุตฺตํ ธมฺมํ ปจฺจยา สญฺโญชน\-สมฺปยุตฺโต จ สญฺโญชน\-วิปฺปยุตฺโต จ ธมฺมา อุปฺปชฺชนฺติ อารมฺมณ\-ปจฺจยา. วตฺถุํ ปจฺจยา อุทฺธจฺจ\-สหคตา ขนฺธา จ โมโห จ. (3)}

\proseitem{54}{สญฺโญชน\-สมฺปยุตฺตญฺจ สญฺโญชน\-วิปฺปยุตฺตญฺจ ธมฺมํ ปจฺจยา สญฺโญชน\-สมฺปยุตฺโต ธมฺโม อุปฺปชฺชติ อารมฺมณ\-ปจฺจยา. สญฺโญชน\-สมฺปยุตฺตํ เอกํ ขนฺธญฺจ วตฺถุญฺจ ปจฺจยา ตโย ขนฺธา ฯเปฯ เทฺว ขนฺเธ ฯเปฯ. อุทฺธจฺจ\-สหคตํ เอกํ ขนฺธญฺจ โมหญฺจ ปจฺจยา ตโย ขนฺธา ฯเปฯ เทฺว ขนฺเธ ฯเปฯ. (1)}

\prose{สญฺโญชน\-สมฺปยุตฺตญฺจ สญฺโญชน\-วิปฺปยุตฺตญฺจ ธมฺมํ ปจฺจยา สญฺโญชน\-วิปฺปยุตฺโต ธมฺโม อุปฺปชฺชติ อารมฺมณ\-ปจฺจยา. อุทฺธจฺจ\-สหคเต ขนฺเธ จ วตฺถุญฺจ ปจฺจยา อุทฺธจฺจ\-สหคโต โมโห.}

\prose{สญฺโญชน\-สมฺปยุตฺตญฺจ สญฺโญชน\-วิปฺปยุตฺตญฺจ ธมฺมํ ปจฺจยา สญฺโญชน\-สมฺปยุตฺโต จ สญฺโญชน\-วิปฺปยุตฺโต จ ธมฺมา อุปฺปชฺชนฺติ อารมฺมณ\-ปจฺจยา. อุทฺทจฺจสหคตํ เอกํ ขนฺธญฺจ วตฺถุญฺจ ปจฺจยา ตโย ขนฺธา โมโห จ ฯเปฯ เทฺว ขนฺเธ ฯเปฯ. (3)}

\csromanheader{2}{\textbf{อธิปตฺยาทิ}}
\proseitem{53}{สญฺโญชน\-สมฺปยุตฺตํ ธมฺมํ ปจฺจยา สญฺโญชน\-สมฺปยุตฺโต ธมฺโม อุปฺปชฺชติ อธิปติ\-ปจฺจยา ฯเปฯ อวิคตปจฺจยา ฯเปฯ.}

\csromanmarkh{1. ปจฺจยานุโลม}\csromanmarkhii{2. สงฺขฺยาวาร}\csromanheadernomark{2}{1. \textbf{ปจฺจยานุโลม 2. สงฺขฺยาวาร}}
\proseitem{53}{เหตุยา นว, อารมฺมเณ นว, อธิปติยา นว, (สพฺพตฺถ นว.) วิปาเก เอกํ, อาหาเร นว ฯเปฯ อวิคเต นว.}

\vspace{\tipitakaparskip}
\csromancenter{อนุโลมํ.}
\csromanmarkh{22. สญฺโญชน\-สมฺปยุตฺตทุก}\csromanmarkhii{2. ปจฺจย\-ปจฺจนีย 1. วิภงฺควาร}\csromanheadernomark{2}{22. สญฺโญชน\-สมฺปยุตฺตทุก 2. ปจฺจย\-ปจฺจนีย 1. วิภงฺควาร}
\proseitem{57}{\csromanfolio{229}สญฺโญชน\-สมฺปยุตฺตํ ธมฺมํ ปจฺจยา สญฺโญชน\-สมฺปยุตฺโต ธมฺโม อุปฺปชฺชติ นเหตุปจฺจยา. วิจิกิจฺฉา\-สหคเต ขนฺเธ ปจฺจยา วิจิกิจฺฉา\-สหคโต โมโห. (1)}

\prose{สญฺโญชน\-สมฺปยุตฺตํ ธมฺมํ ปจฺจยา สญฺโญชน\-วิปฺปยุตฺโต ธมฺโม อุปฺปชฺชติ นเหตุปจฺจยา. อุทฺธจฺจ\-สหคเต ขนฺเธ ปจฺจยา อุทฺธจฺจ\-สหคโต โมโห. (2)}

\proseitem{58}{สญฺโญชน\-วิปฺปยุตฺตํ ธมฺมํ ปจฺจยา สญฺโญชน\-วิปฺปยุตฺโต ธมฺโม อุปฺปชฺชติ นเหตุปจฺจยา. อเหตุกํ สญฺโญชน\-วิปฺปยุตฺตํ เอกํ ขนฺธํ ฯเปฯ. (ยาว อสญฺญสตฺตา.) จกฺขายตนํ ปจฺจยา จกฺขุวิญฺญาณํ ฯเปฯ. กายายตนํ ปจฺจยา กายวิญฺญาณํ. วตฺถุํ ปจฺจยา อเหตุกา สญฺโญชน\-วิปฺปยุตฺตา ขนฺธา จ อุทฺธจฺจ\-สหคโต โมโห จ. (1)}

\prose{สญฺโญชน\-วิปฺปยุตฺตํ ธมฺมํ ปจฺจยา สญฺโญชน\-สมฺปยุตฺโต ธมฺโม อุปฺปชฺชติ นเหตุปจฺจยา. วตฺถุํ ปจฺจยา วิจิกิจฺฉา\-สหคโต โมโห. (2)}

\proseitem{59}{สญฺโญชน\-สมฺปยุตฺตญฺจ สญฺโญชน\-วิปฺปยุตฺตญฺจ ธมฺมํ ปจฺจยา สญฺโญชน\-สมฺปยุตฺโต ธมฺโม อุปฺปชฺชติ นเหตุปจฺจยา. วิจิกิจฺฉา\-สหคเต ขนฺเธ จ วตฺถุญฺจ ปจฺจยา วิจิกิจฺฉา\-สหคโต โมโห. (1)}

\prose{สญฺโญชน\-สมฺปยุตฺตญฺจ สญฺโญชน\-วิปฺปยุตฺตญฺจ ธมฺมํ ปจฺจยา สญฺโญชน\-วิปฺปยุตฺโต ธมฺโม อุปฺปชฺชติ นเหตุปจฺจยา. อุทฺธจฺจ\-สหคเต ขนฺเธ จ วตฺถุญฺจ ปจฺจยา อุทฺธจฺจ\-สหคโต โมโห. (สํขิตฺตํ.) (2)}

\csromanmarkh{2. ปจฺจย\-ปจฺจนีย}\csromanmarkhii{2. สงฺขฺยาวาร}\csromanheadernomark{2}{2. ปจฺจย\-ปจฺจนีย 2. สงฺขฺยาวาร}
\vspace{\tipitakaparskip}
\csromancenter{\textbf{นเหตุ}ยา ฉ, นอารมฺมเณ ตีณิ, นอธิปติยา นว, นอนนฺตเร ตีณิ, นสมนนฺตเร ตีณิ, นอญฺญมญฺเญ ตีณิ, นฺเอาปนิสฺสเย ตีณิ, นปุเรชาเต สตฺต, นปจฺฉาชาเต นว, นอาเสวเน นว, นกมฺเม จตฺตาริ,}
\prosecont{\csromanfolio{230}นวิปาเก นว, นอาหาเร เอกํ, นอินฺทฺริเย เอกํ, นฌาเน เอกํ, นมคฺเค เอกํ, นสมฺปยุตฺเต ตีณิ, นวิปฺปยุตฺเต ฉ, โนนตฺถิยา ตีณิ, โนวิคเต ตีณิ.}

\prose{เหตุปจฺจยา นอารมฺมเณ ตีณิ, นอธิปติยา นว. (สํขิตฺตํ. เอวํ กาตพฺพํ.)}

\proseitem{62}{นเหตุปจฺจยา อารมฺมเณ ฉ, (สพฺพตฺถ ฉ,) วิปาเก เอกํ, อาหาเร ฉ ฯเปฯ มคฺเค ฉ ฯเปฯ อวิคเต ฉ.}

\vspace{\tipitakaparskip}
\csromancenter{(นิสฺสยวาโรปิ ปจฺจยวาร\-สทิโส.)}
\csromanmarkh{22. สญฺโญชน\-สมฺปยุตฺตทุก}\csromanmarkhii{5. สํสฏฺฐวาร}\csromanheadernomark{2}{22. สญฺโญชน\-สมฺปยุตฺตทุก 5. สํสฏฺฐวาร}
\proseitem{60}{สญฺโญชน\-สมฺปยุตฺตํ ธมฺมํ สํสฏฺโฐ สญฺโญชน\-สมฺปยุตฺโต ธมฺโม อุปฺปชฺชติ เหตุปจฺจยา. สญฺโญชน\-สมฺปยุตฺตํ เอกํ ขนฺธํ สํสฏฺฐา ตโย ขนฺธา ฯเปฯ เทฺว ขนฺเธ ฯเปฯ. (1)}

\prose{สญฺโญชน\-วิปฺปยุตฺตํ ธมฺมํ สํสฏฺโฐ สญฺโญชน\-วิปฺปยุตฺโต ธมฺโม อุปฺปชฺชติ เหตุปจฺจยา. สญฺโญชน\-วิปฺปยุตฺตํ เอกํ ขนฺธํ สํสฏฺฐา ตโย ขนฺธา ฯเปฯ เทฺว ขนฺเธ ฯเปฯ. ปฏิสนฺธิ\-กฺขเณ ฯเปฯ. (1)}

\prose{สญฺโญชน\-วิปฺปยุตฺตํ ธมฺมํ สํสฏฺโฐ สญฺโญชน\-สมฺปยุตฺโต ธมฺโม อุปฺปชฺชติ เหตุปจฺจยา. อุทฺธจฺจ\-สหคตํ โมหํ สํสฏฺฐา สมฺปยุตฺตกา ขนฺธา. (2)}

\prose{\csromanfolio{231}สญฺโญชน\-สมฺปยุตฺตญฺจ สญฺโญชน\-วิปฺปยุตฺตญฺจ ธมฺมํ สํสฏฺโฐ สญฺโญชน\-สมฺปยุตฺโต ธมฺโม อุปฺปชฺชติ เหตุปจฺจยา. อุทฺธจฺจ\-สหคตํ เอกํ ขนฺธญฺจ โมหญฺจ สํสฏฺฐา ตโย ขนฺธา ฯเปฯ เทฺว ขนฺเธ ฯเปฯ. (1) (สํขิตฺตํ.)}

\proseitem{64}{เหตุยา จตฺตาริ, อารมฺมเณ ฉ, อธิปติยา เทฺว, อนนฺตเร ฉ, สมนนฺตเร ฉ ฯเปฯ อวิคเต ฉ.}

\proseitem{65}{นเหตุยา ตีณิ, นอธิปติยา ฉ, นปุเรชาเต ฉ, นปจฺฉาชาเต ฉ, นอาเสวเน ฉ, นกมฺเม จตฺตาริ, นวิปาเก ฉ, นฌาเน เอกํ, นมคฺเค เอกํ, นวิปฺปยุตฺเต ฉ.}

\vspace{\tipitakaparskip}
\csromancenter{(อิตเร เทฺว คณนาปิ สมฺปยุตฺตวาโรปิ กาตพฺโพ.)}
\csromanmarkh{22. สญฺโญชน\-สมฺปยุตฺตทุก}\csromanmarkhii{7. ปญฺหาวาร}\csromanheadernomark{2}{22. สญฺโญชน\-สมฺปยุตฺตทุก 7. ปญฺหาวาร}
\csromanmarkh{1. ปจฺจยานุโลม}\csromanmarkhii{1. วิภงฺควาร}\csromanheadernomark{2}{1. \textbf{ปจฺจยานุโลม 1. วิภงฺควาร}}
\vspace{\tipitakaparskip}
\csromancenter{\textbf{เหตุ}}
\proseitem{66}{สญฺโญชน\-สมฺปยุตฺโต ธมฺโม สญฺโญชน\-สมฺปยุตฺตสฺส ธมฺมสฺส เหตุปจฺจเยน ปจฺจโย. สญฺโญชน\-สมฺปยุตฺตา เหตู สมฺปยุตฺตกานํ ขนฺธานํ เหตุปจฺจเยน ปจฺจโย. (1)}

\prose{สญฺโญชน\-สมฺปยุตฺโต ธมฺโม สญฺโญชน\-วิปฺปยุตฺตสฺส ธมฺมสฺส เหตุปจฺจเยน ปจฺจโย. สญฺโญชน\-สมฺปยุตฺตา เหตู จิตฺต\-สมุฏฺฐานานํ รูปานํ เหตุปจฺจเยน ปจฺจโย. (2)}

\prose{สญฺโญชน\-สมฺปยุตฺโต ธมฺโม สญฺโญชน\-สมฺปยุตฺตสฺส จ สญฺโญชน\-วิปฺปยุตฺตสฺส จ ธมฺมสฺส เหตุปจฺจเยน ปจฺจโย. สญฺโญชน\-สมฺปยุตฺตา เหตู สมฺปยุตฺตกานํ ขนฺธานํ จิตฺตสมุฏฺฐานานญฺจ รูปานํ เหตุปจฺจเยน ปจฺจโย. (3)}

\proseitem{67}{สญฺโญชน\-วิปฺปยุตฺโต ธมฺโม สญฺโญชน\-วิปฺปยุตฺตสฺส ธมฺมสฺส เหตุปจฺจเยน ปจฺจโย. สญฺโญชน\-วิปฺปยุตฺโต เหตู สมฺปยุตฺตกานํ ขนฺธานํ จิตฺตสมุฏฺฐานานญฺจ รูปานํ เหตุปจฺจเยน ปจฺจโย. อุทฺธจฺจ\-สหคโต \csromanfolio{232}โมโห จิตฺต\-สมุฏฺฐานานํ รูปานํ เหตุปจฺจเยน ปจฺจโย. ปฏิสนฺธิ\-กฺขเณ ฯเปฯ. (1)}

\prose{สญฺโญชน\-วิปฺปยุตฺโต ธมฺโม สญฺโญชน\-สมฺปยุตฺตสฺส ธมฺมสฺส เหตุปจฺจเยน ปจฺจโย. อุทฺธจฺจ\-สหคโต โมโห สมฺปยุตฺตกานํ ขนฺธานํ เหตุปจฺจเยน ปจฺจโย. (2)}

\prose{สญฺโญชน\-วิปฺปยุตฺโต ธมฺโม สญฺโญชน\-สมฺปยุตฺตสฺส จ สญฺโญชน\-วิปฺปยุตฺตสฺส จ ธมฺมสฺส เหตุปจฺจเยน ปจฺจโย. อุทฺธจฺจ\-สหคโต โมโห สมฺปยุตฺตกานํ ขนฺธาน จิตฺตสมุฏฺฐานาญฺจ รูปานํ เหตุปจฺจเยน ปจฺจโย. (3)}

\prose{สญฺโญชน\-สมฺปยุตฺโต ธมฺโม สญฺโญชน\-สมฺปยุตฺตสฺส ธมฺมสฺส อารมฺมณ\-ปจฺจเยน ปจฺจโย. สญฺโญชน\-สมฺปยุตฺเต ขนฺเธ อารพฺภ สญฺโญชน\-สมฺปยุตฺตกา ขนฺธา อุปฺปชฺชนฺติ. (มูลํ กาตพฺพํ.) สญฺโญชน\-สมฺปยุตฺเต ขนฺเธ อารพฺภ สญฺโญชน\-วิปฺปยุตฺตา ขนฺธา จ โมโห จ อุปฺปชฺชนฺติ. (มูลํ กาตพฺพํ.) สญฺโญชน\-สมฺปยุตฺเต ขนฺเธ อารพฺภ อุทฺธจฺจ\-สหคตา ขนฺธา จ โมโห จ อุปฺปชฺชนฺติ. (3)}

\prose{สญฺโญชน\-วิปฺปยุตฺโต ธมฺโม สญฺโญชน\-วิปฺปยุตฺตสฺส ธมฺมสฺส อารมฺมณ\-ปจฺจเยน ปจฺจโย. ทานํ ทตฺวา, สีลํ ฯเปฯ อุโปสถกมฺมํ กตฺวา ตํ ปจฺจเวกฺขติ. ปุพฺเพ สุจิณฺณานิ ฯเปฯ. ฌานา ฯเปฯ. อริยา มคฺคา วุฏฺฐหิตฺวา มคฺคํ ปจฺจเวกฺขนฺติ, ผลํ ปจฺจเวกฺขนฺติ, นิพฺพานํ ปจฺจเวกฺขนฺติ. นิพฺพานํ โคตฺรภุสฺส, โวทานสฺส, มคฺคสฺส, ผลสฺส, อาวชฺชนาย อารมฺมณ\-ปจฺจเยน ปจฺจโย. อริยา สญฺโญชน\-วิปฺปยุตฺเต ปหีเน กิเลเส ปจฺจเวกฺขนฺติ, วิกฺขมฺภิเต กิเลเส ปจฺจเวกฺขนฺติ. ปุพฺเพ สมุทาจิณฺเณ กิเลเส ชานนฺติ. จกฺขุํ ฯเปฯ วตฺถุํ สญฺโญชน\-วิปฺปยุตฺเต ขนฺเธ จ โมหญฺจ อนิจฺจโต ฯเปฯ วิปสฺสติ ฯเปฯ. ทิพฺเพน จกฺขุนา รูปํ ปสฺสติ. ทิพฺพาย โสตธาตุยา สทฺทํ สุณาติ. เจโตปริย\-ญาเณน สญฺโญชน\-วิปฺปยุตฺต\-จิตฺต\-สมงฺคิสฺส จิตฺตํ ชานาติ. อากาสา\-นญฺจา\-ยตนํ วิญฺญาณญฺจา\-ยตนสฺส ฯเปฯ. รูปายตนํ ฯเปฯ. โผฏฺฐพฺพา\-ยตนํ ฯเปฯ. สญฺโญชน\-วิปฺปยุตฺตา ขนฺธา อิทฺธิวิธ\-ญาณสฺส, เจโตปริย\-ญาณสฺส, ปุพฺเพ\-นิวาสา\-นุสฺสติ\-ญาณสฺส, ยถากมฺมูปคญาณสฺส, \csromanfolio{233}อนาคตํสญาณสฺส, อาวชฺชนาย โมหสฺส จ อารมฺมณ\-ปจฺจเยน ปจฺจโย. (1)}

\proseitem{69}{สญฺโญชน\-วิปฺปยุตฺโต ธมฺโม สญฺโญชน\-สมฺปยุตฺตสฺส ธมฺมสฺส อารมฺมณ\-ปจฺจเยน ปจฺจโย. ทานํ ทตฺวา, สีลํ ฯเปฯ อุโปสถกมฺมํ ฯเปฯ. ปุพฺเพ สุจิณฺณานิ ฯเปฯ. ฌานา วุฏฺฐหิตฺวา ฌานํ ฯเปฯ. จกฺขุํ ฯเปฯ วตฺถุํ สญฺโญชน\-วิปฺปยุตฺเต ขนฺเธ จ โมหญฺจ อสฺสาเทติ อภินนฺทติ, ตํ อารพฺภ ราโค อุปฺปชฺชติ ฯเปฯ โทมนสฺสํ อุปฺปชฺชติ. (2)}

\prose{สญฺโญชน\-วิปฺปยุตฺโต ธมฺโม สญฺโญชน\-สมฺปยุตฺตสฺส จ สญฺโญชน\-วิปฺปยุตฺตสฺส จ ธมฺมสฺส อารมฺมณ\-ปจฺจเยน ปจฺจโย. จกฺขุํ ฯเปฯ วตฺถุํ สญฺโญชน\-วิปฺปยุตฺเต ขนฺเธ จ โมหญฺจ อารพฺภ อุทฺธจฺจ\-สหคตา ขนฺธา จ โมโห จ อุปฺปชฺชนฺติ. (3)}

\proseitem{70}{สญฺโญชน\-สมฺปยุตฺโต จ สญฺโญชน\-วิปฺปยุตฺโต จ ธมฺมา สญฺโญชน\-สมฺปยุตฺตสฺส ธมฺมสฺส อารมฺมณ\-ปจฺจเยน ปจฺจโย. อุทฺธจฺจ\-สหคเต ขนฺเธ จ โมหญฺจ อารพฺภ สญฺโญชน\-สมฺปยุตฺตกา ขนฺธา อุปฺปชฺชนฺติ. (มูลํ กาตพฺพํ.) อุทฺธจฺจ\-สหคเต ขนฺเธ จ โมหญฺจ อารพฺภ สญฺโญชน\-วิปฺปยุตฺตา ขนฺธา จ โมโห จ อุปฺปชฺชนฺติ. (มูลํ กาตพฺพํ.) อุทฺธจฺจ\-สหคเต ขนฺเธ จ โมหญฺจ อารพฺภ อุทฺธจฺจ\-สหคตา ขนฺธา จ โมโห จ อุปฺปชฺชนฺติ. (3)}

\proseitem{69}{สญฺโญชน\-สมฺปยุตฺโต ธมฺโม สญฺโญชน\-สมฺปยุตฺตสฺส ธมฺมสฺส อธิปติ\-ปจฺจเยน ปจฺจโย. อารมฺมณา\-ธิปติ, สหชาตา\-ธิปติ.  อารมฺมณา\-ธิปติ\csromandash{}ราคํ ฯเปฯ. ทิฏฺฐิํ ครุํ กตฺวา อสฺสาเทติ อภินนฺทติ, ตํ ครุํ กตฺวา ราโค อุปฺปชฺชติ, ทิฏฺฐิ อุปฺปชฺชติ. สหชาตา\-ธิปติ\csromandash{} สญฺโญชน\-สมฺปยุตฺตา\-ธิปติ สมฺปยุตฺตกานํ ขนฺธานํ อธิปติ\-ปจฺจเยน ปจฺจโย. (1)}

\prose{สญฺโญชน\-สมฺปยุตฺโต ธมฺโม สญฺโญชน\-วิปฺปยุตฺตสฺส ธมฺมสฺส อธิปติ\-ปจฺจเยน ปจฺจโย. สหชาตา\-ธิปติ\csromandash{}สญฺโญชน\-สมฺปยุตฺตา\-ธิปติ จิตฺต\-สมุฏฺฐานานํ รูปานํ อธิปติ\-ปจฺจเยน ปจฺจโย. (2)}

\proseitem{71}{\csromanfolio{234}สญฺโญชน\-สมฺปยุตฺโต ธมฺโม สญฺโญชน\-สมฺปยุตฺตสฺส จ สญฺโญชน\-วิปฺปยุตฺตสฺส จ ธมฺมสฺส อธิปติ\-ปจฺจเยน ปจฺจโย. สหชาตา\-ธิปติ\csromandash{}สญฺโญชน\-สมฺปยุตฺตา\-ธิปติ สมฺปยุตฺตกานํ ขนฺธานํ จิตฺตสมุฏฺฐานานญฺจ รูปานํ อธิปติ\-ปจฺจเยน ปจฺจโย. (3)}

\proseitem{72}{สญฺโญชน\-วิปฺปยุตฺโต ธมฺโม สญฺโญชน\-วิปฺปยุตฺตสฺส ธมฺมสฺส อธิปติ\-ปจฺจเยน ปจฺจโย. อารมฺมณา\-ธิปติ, สหชาตา\-ธิปติ.  อารมฺมณา\-ธิปติ\csromandash{}ทานํ ทตฺวา, สีลํ ฯเปฯ อุโปสถกมฺมํ กตฺวา ตํ ครุํ กตฺวา ปจฺจเวกฺขติ. ปุพฺเพ สุจิณฺณานิ ฯเปฯ. ฌานา ฯเปฯ. อริยา มคฺคา ฯเปฯ. ผลํ ฯเปฯ. นิพฺพานํ ฯเปฯ. นิพฺพานํ โคตฺรภุสฺส, โวทานสฺส, มคฺคสฺส, ผลสฺส อธิปติ\-ปจฺจเยน ปจฺจโย. สหชาตา\-ธิปติ\csromandash{} สญฺโญชน\-วิปฺปยุตฺตา\-ธิปติ สมฺปยุตฺตกานํ ขนฺธานํ จิตฺตสมุฏฺฐานานญฺจ รูปานํ อธิปติ\-ปจฺจเยน ปจฺจโย. (1)}

\prose{สญฺโญชน\-วิปฺปยุตฺโต ธมฺโม สญฺโญชน\-สมฺปยุตฺตสฺส ธมฺมสฺส อธิปติ\-ปจฺจเยน ปจฺจโย. อารมฺมณา\-ธิปติ\csromandash{}ทานํ ทตฺวา, สีลํ ฯเปฯ อุโปสถกมฺมํ กตฺวา ตํ ครุํ กตฺวา อสฺสาเทติ อภินนฺทติ, ตํ ครุํ กตฺวา ราโค อุปฺปชฺชติ, ทิฏฺฐิ อุปฺปชฺชติ. ปุพฺเพ สุจิณฺณานิ ฯเปฯ. ฌานา ฯเปฯ. จกฺขุํ ฯเปฯ วตฺถุํ สญฺโญชน\-วิปฺปยุตฺเต ขนฺเธ ครุํ กตฺวา อสฺสาเทติ อภินนฺทติ, ตํ ครุํ กตฺวา ราโค อุปฺปชฺชติ, ทิฏฺฐิ อุปฺปชฺชติ. (2)}

\proseitem{73}{สญฺโญชน\-สมฺปยุตฺโต ธมฺโม สญฺโญชน\-สมฺปยุตฺตสฺส ธมฺมสฺส อนนฺตร\-ปจฺจเยน ปจฺจโย. ปุริมา ปุริมา สญฺโญชน\-สมฺปยุตฺตา ขนฺธา ปจฺฉิมานํ ปจฺฉิมานํ สญฺโญชน\-สมฺปยุตฺตกานํ ขนฺธานํ อนนฺตร\-ปจฺจเยน ปจฺจโย. (1)}

\prose{สญฺโญชน\-สมฺปยุตฺโต ธมฺโม สญฺโญชน\-วิปฺปยุตฺตสฺส ธมฺมสฺส อนนฺตร\-ปจฺจเยน ปจฺจโย. ปุริมา ปุริมา อุทฺธจฺจ\-สหคตา ขนฺธา ปจฺฉิมสฺส ปจฺฉิมสฺส อุทฺธจฺจ\-สหคตสฺส โมหสฺส อนนฺตร\-ปจฺจเยน ปจฺจโย. สญฺโญชน\-สมฺปยุตฺตา ขนฺธา วุฏฺฐานสฺส อนนฺตร\-ปจฺจเยน ปจฺจโย. (2)}

\prose{สญฺโญชน\-สมฺปยุตฺโต ธมฺโม สญฺโญชน\-สมฺปยุตฺตสฺส จ สญฺโญชน\-วิปฺปยุตฺตสฺส จ ธมฺมสฺส อนนฺตร\-ปจฺจเยน ปจฺจโย. ปุริมา ปุริมา \csromanfolio{235}อุทฺธจฺจ\-สหคตา ขนฺธา ปจฺฉิมานํ ปจฺฉิมานํ อุทฺธจฺจ\-สหคตานํ ขนฺธานํ โมหสฺส จ อนนฺตร\-ปจฺจเยน ปจฺจโย. (3)}

\proseitem{74}{สญฺโญชน\-วิปฺปยุตฺโต ธมฺโม สญฺโญชน\-วิปฺปยุตฺตสฺส ธมฺมสฺส อนนฺตร\-ปจฺจเยน ปจฺจโย. ปุริโม ปุริโม อุทฺธจฺจ\-สหคโต โมโห ปจฺฉิมสฺส ปจฺฉิมสฺส อุทฺธจฺจ\-สหคตสฺส โมหสฺส อนนฺตร\-ปจฺจเยน ปจฺจโย. ปุริมา ปุริมา สญฺโญชน\-วิปฺปยุตฺตา ขนฺธา ปจฺฉิมานํ ฯเปฯ ผลสมาปตฺติยา อนนฺตร\-ปจฺจเยน ปจฺจโย. (1)}

\prose{สญฺโญชน\-วิปฺปยุตฺโต ธมฺโม สญฺโญชน\-สมฺปยุตฺตสฺส ธมฺมสฺส อนนฺตร\-ปจฺจเยน ปจฺจโย. ปุริโม ปุริโม อุทฺธจฺจ\-สหคโต โมโห ปจฺฉิมานํ ปจฺฉิมานํ อุทฺธจฺจ\-สหคตานํ ขนฺธานํ อนนฺตร\-ปจฺจเยน ปจฺจโย. อาวชฺชนา สญฺโญชน\-สมฺปยุตฺตกานํ ขนฺธานํ อนนฺตร\-ปจฺจเยน ปจฺจโย. (2)}

\prose{สญฺโญชน\-วิปฺปยุตฺโต ธมฺโม สญฺโญชน\-สมฺปยุตฺตสฺส จ สญฺโญชน\-วิปฺปยุตฺตสฺส จ ธมฺมสฺส อนนฺตร\-ปจฺจเยน ปจฺจโย. ปุริโม ปุริโม อุทฺธจฺจ\-สหคโต โมโห ปจฺฉิมานํ ปจฺฉิมานํ อุทฺธจฺจ\-สหคตานํ ขนฺธานํ โมหสฺส จ อนนฺตร\-ปจฺจเยน ปจฺจโย. อาวชฺชนา อุทฺธจฺจ\-สหคตานํ ขนฺธานํ โมหสฺส จ อนนฺตร\-ปจฺจเยน ปจฺจโย. (3)}

\proseitem{75}{สญฺโญชน\-สมฺปยุตฺโต จ สญฺโญชน\-วิปฺปยุตฺโต จ ธมฺมา สญฺโญชน\-สมฺปยุตฺตสฺส ธมฺมสฺส อนนฺตร\-ปจฺจเยน ปจฺจโย. ปุริมา ปุริมา อุทฺธจฺจ\-สหคตา ขนฺธา จ โมโห จ ปจฺฉิมานํ ปจฺฉิมานํ อุทฺธจฺจ\-สหคตานํ ขนฺธานํ อนนฺตร\-ปจฺจเยน ปจฺจโย. (1)}

\prose{สญฺโญชน\-สมฺปยุตฺโต จ สญฺโญชน\-วิปฺปยุตฺโต จ ธมฺมา สญฺโญชน\-วิปฺปยุตฺตสฺส ธมฺมสฺส อนนฺตร\-ปจฺจเยน ปจฺจโย. ปุริมา ปุริมา อุทฺธจฺจ\-สหคตา ขนฺธา จ โมโห จ ปจฺฉิมสฺส ปจฺฉิมสฺส อุทฺธจฺจ\-สหคตสฺส โมหสฺส อนนฺตร\-ปจฺจเยน ปจฺจโย. อุทฺธจฺจ\-สหคตา ขนฺธา จ โมโห จ วุฏฺฐานสฺส อนนฺตร\-ปจฺจเยน ปจฺจโย. (2)}

\prose{สญฺโญชน\-สมฺปยุตฺโต จ สญฺโญชน\-วิปฺปยุตฺโต จ ธมฺมา สญฺโญชน\-สมฺปยุตฺตสฺส จ สญฺโญชน\-วิปฺปยุตฺตสฺส จ ธมฺมสฺส อนนฺตร\-ปจฺจเยน ปจฺจโย. ปุริมา ปุริมา อุทฺธจฺจ\-สหคตา ขนฺธา จ โมโห จ ปจฺฉิมานํ \csromanfolio{236}ปจฺฉิมานํ อุทฺธจฺจ\-สหคตานํ ขนฺธานํ โมหสฺส จ อนนฺตร\-ปจฺจเยน ปจฺจโย. (3)}

\csromanheader{2}{\textbf{สมนนฺตราทิ}}
\proseitem{76}{สญฺโญชน\-สมฺปยุตฺโต ธมฺโม สญฺโญชน\-สมฺปยุตฺตสฺส ธมฺมสฺส สมนนฺตร\-ปจฺจเยน ปจฺจโย. นว. สหชาต\-ปจฺจเยน ปจฺจโย. นว. อญฺญมญฺญ\-ปจฺจเยน ปจฺจโย. ฉ. นิสฺสย\-ปจฺจเยน ปจฺจโย. นว.}

\csromanheader{2}{\textbf{อุปนิสฺสย}}
\proseitem{77}{สญฺโญชน\-สมฺปยุตฺโต ธมฺโม สญฺโญชน\-สมฺปยุตฺตสฺส ธมฺมสฺส อุปนิสฺสย\-ปจฺจเยน ปจฺจโย. อารมฺมณู\-ปนิสฺสโย, อนนฺตรู\-ปนิสฺสโย, ปกตู\-ปนิสฺสโย ฯเปฯ. ปกตู\-ปนิสฺสโย\csromandash{}สญฺโญชน\-สมฺปยุตฺตา ขนฺธา สญฺโญชน\-สมฺปยุตฺตกานํ ขนฺธานํ อุปนิสฺสย\-ปจฺจเยน ปจฺจโย. (1)}

\prose{สญฺโญชน\-สมฺปยุตฺโต ธมฺโม สญฺโญชน\-วิปฺปยุตฺตสฺส ธมฺมสฺส อุปนิสฺสย\-ปจฺจเยน ปจฺจโย. อนนฺตรู\-ปนิสฺสโย, ปกตู\-ปนิสฺสโย ฯเปฯ. ปกตู\-ปนิสฺสโย\csromandash{}สญฺโญชน\-สมฺปยุตฺตา ขนฺธา สญฺโญชน\-วิปฺปยุตฺตานํ ขนฺธานํ โมหสฺส จ อุปนิสฺสย\-ปจฺจเยน ปจฺจโย. (2)}

\prose{สญฺโญชน\-สมฺปยุตฺโต ธมฺโม สญฺโญชน\-สมฺปยุตฺตสฺส จ สญฺโญชน\-วิปฺปยุตฺตสฺส จ ธมฺมสฺส อุปนิสฺสย\-ปจฺจเยน ปจฺจโย. อนนฺตรู\-ปนิสฺสโย. ปกตู\-ปนิสฺสโย ฯเปฯ. ปกตู\-ปนิสฺสโย\csromandash{} สญฺโญชน\-สมฺปยุตฺตา ขนฺธา อุทฺธจฺจ\-สหคตานํ ขนฺธานํ โมหสฺส จ อุปนิสฺสย\-ปจฺจเยน ปจฺจโย. (3)}

\proseitem{78}{สญฺโญชน\-วิปฺปยุตฺโต ธมฺโม สญฺโญชน\-วิปฺปยุตฺตสฺส ธมฺมสฺส อุปนิสฺสย\-ปจฺจเยน ปจฺจโย. อารมฺมณู\-ปนิสฺสโย, อนนฺตรู\-ปนิสฺสโย, ปกตู\-ปนิสฺสโย ฯเปฯ. ปกตู\-ปนิสฺสโย\csromandash{}สทฺธํ อุปนิสฺสาย ทานํ เทติ ฯเปฯ สมาปตฺติํ อุปฺปาเทติ. สีลํ ฯเปฯ. ปญฺญํ. กายิกํ สุขํ. กายิกํ ทุกฺขํ. อุตุํ. โภชนํ. เสนาสนํ โมหํ อุปนิสฺสาย ทานํ เทติ ฯเปฯ สมาปตฺติํ อุปฺปาเทติ. สทฺธา ฯเปฯ. เสนาสนํ โมโห จ สทฺธาย ฯเปฯ ผลสมาปตฺติยา โมหสฺส จ อุปนิสฺสย\-ปจฺจเยน ปจฺจโย. (1)}

\prose{\csromanfolio{237}สญฺโญชน\-วิปฺปยุตฺโต ธมฺโม สญฺโญชน\-สมฺปยุตฺตสฺส ธมฺมสฺส อุปนิสฺสย\-ปจฺจเยน ปจฺจโย. อารมฺมณู\-ปนิสฺสโย, อนนฺตรู\-ปนิสฺสโย, ปกตู\-ปนิสฺสโย ฯเปฯ. ปกตู\-ปนิสฺสโย\csromandash{}สทฺธํ อุปนิสฺสาย มานํ ชปฺเปติ. ทิฏฺฐิํ คณฺหาติ. สีลํ ฯเปฯ. ปญฺญํ. กายิกํ สุขํ. กายิกํ ทุกฺขํ. อุตุํ. โภชนํ. เสนาสนํ โมหํ อุปนิสฺสาย ปาณํ หนติ ฯเปฯ สํฆํ ภินฺทติ. สทฺธา ฯเปฯ. เสนาสนํ โมโห ราคสฺส ฯเปฯ ปตฺถนาย อุปนิสฺสย\-ปจฺจเยน ปจฺจโย. (2)}

\prose{สญฺโญชน\-วิปฺปยุตฺโต ธมฺโม สญฺโญชน\-สมฺปยุตฺตสฺส จ สญฺโญชน\-วิปฺปยุตฺตสฺส จ ธมฺมสฺส อุปนิสฺสย\-ปจฺจเยน ปจฺจยฺโอ. อนนฺตรู\-ปนิสฺสโย, ปกตู\-ปนิสฺสโย ฯเปฯ ปกตู\-ปนิสฺสโย\csromandash{}สทฺธา ฯเปฯ ปญฺญา. กายิกํ สุขํ ฯเปฯ. เสนาสนํ โมโห จ อุทฺธจฺจ\-สหคตานํ ขนฺธานํ โมหสฺส จ อุปนิสฺสย\-ปจฺจเยน ปจฺจโย. (3)}

\proseitem{79}{สญฺโญชน\-สมฺปยุตฺโต จ สญฺโญชน\-วิปฺปยุตฺโต จ ธมฺมา สญฺโญชน\-สมฺปยุตฺตสฺส ธมฺมสฺส อุปนิสฺสย\-ปจฺจเยน ปจฺจโย. อนนฺตรู\-ปนิสฺสโย, ปกตู\-ปนิสฺสโย ฯเปฯ. ปกตู\-ปนิสฺสโย\csromandash{} อุทฺธจฺจ\-สหคตา ขนฺธา จ โมโห จ สญฺโญชน\-สมฺปยุตฺตกานํ ขนฺธานํ อุปนิสฺสย\-ปจฺจเยน ปจฺจโย. (1)}

\prose{สญฺโญชน\-สมฺปยุตฺโต จ สญฺโญชน\-วิปฺปยุตฺโต จ ธมฺมา สญฺโญชน\-วิปฺปยุตฺตสฺส ธมฺมสฺส อุปนิสฺสย\-ปจฺจเยน ปจฺจโย. อนนฺตรู\-ปนิสฺสโย, ปกตู\-ปนิสฺสโย ฯเปฯ. ปกตู\-ปนิสฺสโย\csromandash{} อุทฺธจฺจ\-สหคตา ขนฺธา จ โมโห จ สญฺโญชน\-วิปฺปยุตฺตานํ ขนฺธานํ โมหสฺส จ อุปนิสฺสย\-ปจฺจเยน ปจฺจโย. (2)}

\prose{สญฺโญชน\-สมฺปยุตฺโต จ สญฺโญชน\-วิปฺปยุตฺโต จ ธมฺมา สญฺโญชน\-สมฺปยุตฺตสฺส จ สญฺโญชน\-วิปฺปยุตฺตสฺส จ ธมฺมสฺส อุปนิสฺสย\-ปจฺจเยน ปจฺจโย. อนนฺตรู\-ปนิสฺสโย, ปกตู\-ปนิสฺสโย ฯเปฯ. ปกตู\-ปนิสฺสโย\csromandash{}อุทฺธจฺจ\-สหคตา ขนฺธา จ โมโห จ อุทฺธจฺจ\-สหคตานํ ขนฺธานํ โมหสฺส จ อุปนิสฺสย\-ปจฺจเยน ปจฺจโย. (3)}

\proseitem{80}{\csromanfolio{238}สญฺโญชน\-วิปฺปยุตฺโต ธมฺโม สญฺโญชน\-วิปฺปยุตฺตสฺส ธมฺมสฺส ปุเรชาต\-ปจฺจเยน ปจฺจโย. อารมฺมณ\-ปุเรชาตํ, วตฺถุ\-ปุเรชาตํ.  อารมฺมณ\-ปุเรชาตํ\csromandash{}จกฺขุํ ฯเปฯ วตฺถุํ อนิจฺจโต ฯเปฯ วิปสฺสติ. ทิพฺเพน จกฺขุนา รูปํ ปสฺสติ. ทิพฺพาย โสตธาตุยา สทฺทํ สุณาติ. รูปายตนํ จกฺขุ\-วิญฺญาณสฺส ฯเปฯ. โผฏฺฐพฺพา\-ยตนํ กายวิญฺญาณสฺส ฯเปฯ. วตฺถุ\-ปุเรชาตํ\csromandash{}จกฺขายตนํ จกฺขุ\-วิญฺญาณสฺส ฯเปฯ กายายตนํ กายวิญฺญาณสฺส ฯเปฯ. วตฺถุ สญฺโญชน\-วิปฺปยุตฺตานํ ขนฺธานํ โมหสฺส จ ปุเรชาต\-ปจฺจเยน ปจฺจโย. (1)}

\prose{สญฺโญชน\-วิปฺปยุตฺโต ธมฺโม สญฺโญชน\-สมฺปยุตฺตสฺส ธมฺมสฺส ปุเรชาต\-ปจฺจเยน ปจฺจโย. อารมฺมณ\-ปุเรชาตํ, วตฺถุ\-ปุเรชาตํ.  อารมฺมณ\-ปุเรชาตํ\csromandash{}จกฺขุํ ฯเปฯ วตฺถุํ อสฺสาเทติ อภินนฺทติ, ตํ อารพฺภ ราโค อุปฺปชฺชติ ฯเปฯ โทมนสฺสํ อุปฺปชฺชติ. วตฺถุ\-ปุเรชาตํ\csromandash{}วตฺถุ สญฺโญชน\-สมฺปยุตฺตกานํ ขนฺธานํ ปุเรชาต\-ปจฺจเยน ปจฺจโย. (2)}

\prose{สญฺโญชน\-วิปฺปยุตฺโต ธมฺโม สญฺโญชน\-สมฺปยุตฺตสฺส จ สญฺโญชน\-วิปฺปยุตฺตสฺส จ ธมฺมสฺส ปุเรชาต\-ปจฺจเยน ปจฺจโย. อารมฺมณ\-ปุเรชาตํ, วตฺถุ\-ปุเรชาตํ.  อารมฺมณ\-ปุเรชาตํ\csromandash{}จกฺขุํ ฯเปฯ วตฺถุํ อารพฺภ อุทฺธจฺจ\-สหคตา ขนฺธา จ โมโห จ อุปฺปชฺชนฺติ. วตฺถุ\-ปุเรชาตํ\csromandash{}วตฺถุ อุทฺธจฺจ\-สหคตานํ ขนฺธานํ โมหสฺส จ ปุเรชาต\-ปจฺจเยน ปจฺจโย. (3)}

\csromanheader{2}{\textbf{ปจฺฉาชาต}}
\proseitem{81}{สญฺโญชน\-สมฺปยุตฺโต ธมฺโม สญฺโญชน\-วิปฺปยุตฺตสฺส ธมฺมสฺส ปจฺฉาชาต\-ปจฺจเยน ปจฺจโย. ปจฺฉาชาตา สญฺโญชน\-สมฺปยุตฺตา ขนฺธา ปุเรชาตสฺส อิมสฺส กายสฺส ปจฺฉาชาต\-ปจฺจเยน ปจฺจโย. (1)}

\prose{สญฺโญชน\-วิปฺปยุตฺโต ธมฺโม สญฺโญชน\-วิปฺปยุตฺตสฺส ธมฺมสฺส ปจฺฉาชาต\-ปจฺจเยน ปจฺจโย. ปจฺฉาชาตา สญฺโญชน\-วิปฺปยุตฺตา ขนฺธา จ โมโห จ ปุเรชาตสฺส อิมสฺส กายสฺส ปจฺฉาชาต\-ปจฺจเยน ปจฺจโย. (1)}

\prose{\csromanfolio{239}สญฺโญชน\-สมฺปยุตฺโต จ สญฺโญชน\-วิปฺปยุตฺโต จ ธมฺมา สญฺโญชน\-วิปฺปยุตฺตสฺส ธมฺมสฺส ปจฺฉาชาต\-ปจฺจเยน ปจฺจโย. ปจฺฉาชาตา อุทฺธจฺจ\-สหคตา ขนฺธา จ โมโห จ ปุเรชาตสฺส อิมสฺส กายสฺส ปจฺฉาชาต\-ปจฺจเยน ปจฺจโย. (1)}

\csromanheader{2}{\textbf{อาเสวน}}
\proseitem{82}{สญฺโญชน\-สมฺปยุตฺโต ธมฺโม สญฺโญชน\-สมฺปยุตฺตสฺส ธมฺมสฺส อาเสวน\-ปจฺจเยน ปจฺจโย. นว. (อาวชฺชนาปิ วุฏฺฐานมฺปิ นตฺถิ.)}

\proseitem{83}{สญฺโญชน\-สมฺปยุตฺโต ธมฺโม สญฺโญชน\-สมฺปยุตฺตสฺส ธมฺมสฺส กมฺมปจฺจเยน ปจฺจโย. สญฺโญชน\-สมฺปยุตฺตา เจตนา สญฺโญชน\-สมฺปยุตฺตกานํ ขนฺธานํ กมฺมปจฺจเยน ปจฺจโย. (1)}

\prose{สญฺโญชน\-สมฺปยุตฺโต ธมฺโม สญฺโญชน\-วิปฺปยุตฺตสฺส ธมฺมสฺส กมฺมปจฺจเยน ปจฺจโย. สหชาตา, นานากฺขณิกา.  สหชาตา สญฺโญชน\-สมฺปยุตฺตา เจตนา จิตฺต\-สมุฏฺฐานานํ รูปานํ กมฺมปจฺจเยน ปจฺจโย. อุทฺธจฺจ\-สหคตา เจตนา โมหสฺส จิตฺตสมุฏฺฐานานญฺจ รูปานํ กมฺมปจฺจเยน ปจฺจโย. นานากฺขณิกา สญฺโญชน\-สมฺปยุตฺตกา เจตนา วิปากานํ ขนฺธานํ กฏตฺตา จ รูปานํ กมฺมปจฺจเยน ปจฺจโย. (2)}

\prose{สญฺโญชน\-สมฺปยุตฺโต ธมฺโม สญฺโญชน\-สมฺปยุตฺตสฺส จ สญฺโญชน\-วิปฺปยุตฺตสฺส จ ธมฺมสฺส กมฺมปจฺจเยน ปจฺจโย. สญฺโญชน\-สมฺปยุตฺตา เจตนา สมฺปยุตฺตกานํ ขนฺธานํ จิตฺตสมุฏฺฐานานญฺจ รูปานํ กมฺมปจฺจเยน ปจฺจโย. อุทฺธจฺจ\-สหคตา เจตนา สมฺปยุตฺตกานํ ขนฺธานํ โมหสฺส จ จิตฺตสมุฏฺฐานานญฺจ รูปานํ กมฺมปจฺจเยน ปจฺจโย. (3)}

\prose{สญฺโญชน\-วิปฺปยุตฺโต ธมฺโม สญฺโญชน\-วิปฺปยุตฺตสฺส ธมฺมสฺส กมฺมปจฺจเยน ปจฺจโย. สหชาตา, นานากฺขณิกา.  สหชาตา สญฺโญชน\-วิปฺปยุตฺตา เจตนา สมฺปยุตฺตกานํ ขนฺธานํ จิตฺตสมุฏฺฐานานญฺจ รูปานํ กมฺมปจฺจเยน ปจฺจโย. ปฏิสนฺธิ\-กฺขเณ ฯเปฯ. นานากฺขณิกา สญฺโญชน\-วิปฺปยุตฺตา เจตนา วิปากานํ ขนฺธานํ กฏตฺตา จ รูปานํ กมฺมปจฺจเยน ปจฺจโย.}

\csromanheader{2}{\textbf{วิปาก}}
\proseitem{84}{\csromanfolio{240}สญฺโญชน\-วิปฺปยุตฺโต ธมฺโม สญฺโญชน\-วิปฺปยุตฺตสฺส ธมฺมสฺส วิปาก\-ปจฺจเยน ปจฺจโย. เอกํ.}

\csromanheader{2}{\textbf{อาหาราทิ}}
\proseitem{85}{สญฺโญชน\-สมฺปยุตฺโต ธมฺโม สญฺโญชน\-สมฺปยุตฺตสฺส ธมฺมสฺส อาหารปจฺจเยน ปจฺจโย. จตฺตาริ. อินฺทฺริย\-ปจฺจเยน ปจฺจโย. จตฺตาริ. ฌานปจฺจเยน ปจฺจโย. จตฺตาริ. มคฺคปจฺจเยน ปจฺจโย. จตฺตาริ. สมฺปยุตฺต\-ปจฺจเยน ปจฺจโย. ฉ.}

\proseitem{86}{สญฺโญชน\-สมฺปยุตฺโต ธมฺโม สญฺโญชน\-วิปฺปยุตฺตสฺส ธมฺมสฺส วิปฺปยุตฺต\-ปจฺจเยน ปจฺจโย. สหชาตํ, ปจฺฉาชาตํ. (สํขิตฺตํ.) (1)}

\prose{สญฺโญชน\-วิปฺปยุตฺโต ธมฺโม สญฺโญชน\-วิปฺปยุตฺตสฺส ธมฺมสฺส วิปฺปยุตฺต\-ปจฺจเยน ปจฺจโย. สหชาตํ, ปุเรชาตํ, ปจฺฉาชาตํ. (สํขิตฺตํ.) (1)}

\prose{สญฺโญชน\-วิปฺปยุตฺโต ธมฺโม สญฺโญชน\-สมฺปยุตฺตสฺส ธมฺมสฺส วิปฺปยุตฺต\-ปจฺจเยน ปจฺจโย. ปุเรชาตํ วตฺถุ สญฺโญชน\-สมฺปยุตฺตกานํ ขนฺธานํ วิปฺปยุตฺต\-ปจฺจเยน ปจฺจโย. (2)}

\prose{สญฺโญชน\-วิปฺปยุตฺโต ธมฺโม สญฺโญชน\-สมฺปยุตฺตสฺส จ สญฺโญชน\-วิปฺปยุตฺตสฺส จ ธมฺมสฺส วิปฺปยุตฺต\-ปจฺจเยน ปจฺจโย. ปุเรชาตํ วตฺถุ อุทฺธจฺจ\-สหคตานํ ขนฺธานํ โมหสฺส จ วิปฺปยุตฺต\-ปจฺจเยน ปจฺจโย. (3)}

\prose{สญฺโญชน\-สมฺปยุตฺโต จ สญฺโญชน\-วิปฺปยุตฺโต จ ธมฺมา สญฺโญชน\-วิปฺปยุตฺตสฺส ธมฺมสฺส วิปฺปยุตฺต\-ปจฺจเยน ปจฺจโย. สหชาตํ, ปจฺฉาชาตํ. (สํขิตฺตํ.) (1)}

\prose{อตฺถิอาทิ}

\proseitem{87}{สญฺโญชน\-สมฺปยุตฺโต ธมฺโม สญฺโญชน\-สมฺปยุตฺตสฺส ธมฺมสฺส อตฺถิปจฺจเยน ปจฺจโย. เอกํ. (ปฏิจฺจ\-วาร\-สทิสํ.) (1)}

\prose{\csromanfolio{241}สญฺโญชน\-สมฺปยุตฺโต ธมฺโม สญฺโญชน\-วิปฺปยุตฺตสฺส ธมฺมสฺส อตฺถิปจฺจเยน ปจฺจโย. สหชาตํ, ปจฺฉาชาตํ.  สหชาตา สญฺโญชน\-สมฺปยุตฺตา ขนฺธา จิตฺต\-สมุฏฺฐานานํ รูปานํ อตฺถิปจฺจเยน ปจฺจโย. อุทฺธจฺจ\-สหคตา ขนฺธา โมหสฺส จิตฺตสมุฏฺฐานานญฺจ รูปานํ อตฺถิปจฺจเยน ปจฺจโย. ปจฺฉาชาตา สญฺโญชน\-สมฺปยุตฺตา ขนฺธา ปุเรชาตสฺส อิมสฺส กายสฺส อตฺถิปจฺจเยน ปจฺจโย. (2)}

\prose{สญฺโญชน\-สมฺปยุตฺโต ธมฺโม สญฺโญชน\-สมฺปยุตฺตสฺส จ สญฺโญชน\-วิปฺปยุตฺตสฺส จ ธมฺมสฺส อตฺถิปจฺจเยน ปจฺจโย. สญฺโญชน\-สมฺปยุตฺโต เอโก ขนฺโธ ติณฺณินฺนํ ขนฺธานํ จิตฺตสมุฏฺฐานานญฺจ รูปานํ อตฺถิปจฺจเยน ปจฺจโย ฯเปฯ เทฺว ขนฺธา ฯเปฯ. อุทฺธจฺจ\-สหคโต เอโก ขนฺโธ ติณฺณนฺนํ ขนฺธานํ โมหสฺส จ จิตฺตสมุฏฺฐานานญฺจ รูปานํ อตฺถิปจฺจเยน ปจฺจโย ฯเปฯ. (3)}

\proseitem{88}{สญฺโญชน\-วิปฺปยุตฺโต ธมฺโม สญฺโญชน\-วิปฺปยุตฺตสฺส ธมฺมสฺส อตฺถิปจฺจเยน ปจฺจโย. สหชาตํ, ปุเรชาตํ, ปจฺฉาชาตํ, อาหารํ อินฺทฺริยํ. (สํขิตฺตํ. วิตฺถาเรตพฺพํ.) (1)}

\prose{สญฺโญชน\-วิปฺปยุตฺโต ธมฺโม สญฺโญชน\-สมฺปยุตฺตสฺส ธมฺมสฺส อตฺถิปจฺจเยน ปจฺจโย. สหชาตํ, ปุเรชาตํ.  สหชาโต อุทฺธจฺจ\-สหคโต โมโห สมฺปยุตฺตกานํ ขนฺธานํ อตฺถิปจฺจเยน ปจฺจโย. ปุเรชาตํ จกฺขุํ ฯเปฯ วตฺถุํ อสฺสาเทติ อภินนฺทติ, ตํ อารพฺภ ราโค อุปฺปชฺชติ ฯเปฯ โทมนสฺสํ อุปฺปชฺชติ. วตฺถุ สญฺโญชน\-สมฺปยุตฺตกานํ ขนฺธานํ อตฺถิปจฺจเยน ปจฺจโย. (2)}

\prose{สญฺโญชน\-วิปฺปยุตฺโต ธมฺโม สญฺโญชน\-สมฺปยุตฺตสฺส จ สญฺโญชน\-วิปฺปยุตฺตสฺส จ ธมฺมสฺส อตฺถิปจฺจเยน ปจฺจโย. สหชาตํ, ปุเรชาตํ.  สหชาโต อุทฺธจฺจ\-สหคโต โมโห สมฺปยุตฺตกานํ ขนฺธานํ จิตฺตสมุฏฺฐานานญฺจ รูปานํ อตฺถิปจฺจเยน ปจฺจโย. ปุเรชาตํ จกฺขุํ ฯเปฯ วตฺถุํ อารพฺภ อุทฺธจฺจ\-สหคตา ขนฺธา จ โมโห จ อุปฺปชฺชนฺติ. วตฺถุ อุทฺธจฺจ\-สหคตานํ ขนฺธานํ โมหสฺส จ อตฺถิปจฺจเยน ปจฺจโย. (3)}

\proseitem{89}{สญฺโญชน\-สมฺปยุตฺโต จ สญฺโญชน\-วิปฺปยุตฺโต จ ธมฺมา สญฺโญชน\-สมฺปยุตฺตสฺส ธมฺมสฺส อตฺถิปจฺจเยน ปจฺจโย. สหชาตํ, ปุเรชาตํ.  สหชาโต สญฺโญชน\-สมฺปยุตฺโต เอโก ขนฺโธ จ \csromanfolio{242}วตฺถุ จ ติณฺณนฺนํ ขนฺธานํ อตฺถิปจฺจเยน ปจฺจโย ฯเปฯ เทฺว ขนฺธา จ ฯเปฯ. อุทฺธจฺจ\-สหคโต เอโก ขนฺโธ จ โมโห จ ติณฺณนฺนํ ขนฺธานํ อตฺถิปจฺจเยน ปจฺจโย ฯเปฯ เทฺว ขนฺธา จ ฯเปฯ. (1)}

\prose{สญฺโญชน\-สมฺปยุตฺโต จ สญฺโญชน\-วิปฺปยุตฺโต จ ธมฺมา สญฺโญชน\-วิปฺปยุตฺตสฺส ธมฺมสฺส อตฺถิปจฺจเยน ปจฺจโย. สหชาตํ, ปุเรชาตํ, ปจฺฉาชาตํ, อาหารํ อินฺทฺริยํ.  สหชาตา สญฺโญชน\-สมฺปยุตฺตา ขนฺธา จ โมโห จ มหาภูตา จ จิตฺต\-สมุฏฺฐานานํ รูปานํ อตฺถิปจฺจเยน ปจฺจโย. อุทฺธจฺจ\-สหคตา ขนฺธา จ โมโห จ จิตฺต\-สมุฏฺฐานานํ รูปานํ อตฺถิปจฺจเยน ปจฺจโย. อุทฺธจฺจ\-สหคตา ขนฺธา จ วตฺถุ จ โมหสฺส อตฺถิปจฺจเยน ปจฺจโย. ปจฺฉาชาตา อุทฺธจฺจ\-สหคตา ขนฺธา จ โมโห จ ปุเรชาตสฺส อิมสฺส กายสฺส อตฺถิปจฺจเยน ปจฺจโย. ปจฺฉาชาตา สญฺโญชน\-สมฺปยุตฺตา ขนฺธา จ กพฬีกาโร อาหาโร จ อิมสฺส กายสฺส อตฺถิปจฺจเยน ปจฺจโย. ปจฺฉาชาตา สญฺโญชน\-สมฺปยุตฺตา ขนฺธา จ รูปชีวิตินฺทฺริยญฺจ กฏตฺตารูปานํ อตฺถิปจฺจเยน ปจฺจโย. (2)}

\prose{สญฺโญชน\-สมฺปยุตฺโต จ สญฺโญชน\-วิปฺปยุตฺโต จ ธมฺมา สญฺโญชน\-สมฺปยุตฺตสฺส จ สญฺโญชน\-วิปฺปยุตฺตสฺส จ ธมฺมสฺส อตฺถิปจฺจเยน ปจฺจโย. สหชาตํ, ปุเรชาตํ.  สหชาโต อุทฺธจฺจ\-สหคโต เอโก ขนฺโธ จ โมโห จ ติณฺณนฺนํ ขนฺธานํ จิตฺตสมุฏฺฐานานญฺจ รูปานํ อตฺถิปจฺจเยน ปจฺจโย ฯเปฯ เทฺว ขนฺธา จ ฯเปฯ. สหชาโต อุทฺธจฺจ\-สหคโต เอโก ขนฺโธ จ วตฺถุ จ ติณฺณนฺนํ ขนฺธานํ โมหสฺส จ อตฺถิปจฺจเยน ปจฺจโย ฯเปฯ เทฺว ขนฺธา ฯเปฯ. (3)}

\prose{นตฺถิ\-ปจฺจเยน ปจฺจโย. วิคต\-ปจฺจเยน ปจฺจโย. อวิคต\-ปจฺจเยน ปจฺจโย.}

\csromanmarkh{1. ปจฺจยานุโลม}\csromanmarkhii{2. สงฺขฺยาวาร}\csromanheadernomark{2}{1. \textbf{ปจฺจยานุโลม 2. สงฺขฺยาวาร}}
\proseitem{90}{เหตุยา ฉ, อารมฺมเณ นว, อธิปติยา ปญฺจ, อนนฺตเร นว, สมนนฺตเร นว, สหชาเต นว, อญฺญมญฺเญ ฉ, นิสฺสเย นว, อุปนิสฺสเย นว, ปุเรชาเต ตีณิ, ปจฺฉาชาเต ตีณิ, อาเสวเน นว, กมฺเม จตฺตาริ, วิปาเก เอกํ, อาหาเร จตฺตาริ, อินฺทฺริเย จตฺตาริ, ฌาเน จตฺตาริ, \csromanfolio{243}มคฺเค จตฺตาริ, สมฺปยุตฺเต ฉ, วิปฺปยุตฺเต ปญฺจ, อตฺถิยา นว, นตฺถิยา นว, วิคเต นว, อวิคเต นว.}

\vspace{\tipitakaparskip}
\csromancenter{\textbf{อนุโลมํ.}}
\proseitem{91}{สญฺโญชน\-สมฺปยุตฺโต ธมฺโม สญฺโญชน\-สมฺปยุตฺตสฺส ธมฺมสฺส อารมฺมณ\-ปจฺจเยน ปจฺจโย. สหชาต\-ปจฺจเยน ปจฺจโย. อุปนิสฺสย\-ปจฺจเยน ปจฺจโย. (1)}

\prose{สญฺโญชน\-สมฺปยุตฺโต ธมฺโม สญฺโญชน\-วิปฺปยุตฺตสฺส ธมฺมสฺส อารมฺมณ\-ปจฺจเยน ปจฺจโย. สหชาต\-ปจฺจเยน ปจฺจโย. อุปนิสฺสย\-ปจฺจเยน ปจฺจโย. ปจฺฉาชาต\-ปจฺจเยน ปจฺจโย. กมฺมปจฺจเยน ปจฺจโย. (2)}

\prose{สญฺโญชน\-สมฺปยุตฺโต ธมฺโม สญฺโญชน\-สมฺปยุตฺตสฺส จ สญฺโญชน\-วิปฺปยุตฺตสฺส จ ธมฺมสฺส อารมฺมณ\-ปจฺจเยน ปจฺจโย. สหชาต\-ปจฺจเยน ปจฺจโย. อุปนิสฺสย\-ปจฺจเยน ปจฺจโย. (3)}

\proseitem{92}{สญฺโญชน\-วิปฺปยุตฺโต ธมฺโม สญฺโญชน\-วิปฺปยุตฺตสฺส ธมฺมสฺส อารมฺมณ\-ปจฺจเยน ปจฺจโย. สหชาต\-ปจฺจเยน ปจฺจโย. อุปนิสฺสย\-ปจฺจเยน ปจฺจโย. ปุเรชาต\-ปจฺจเยน ปจฺจโย. ปจฺฉาชาต\-ปจฺจเยน ปจฺจโย. กมฺมปจฺจเยน ปจฺจโย. อาหารปจฺจเยน ปจฺจโย. อินฺทฺริย\-ปจฺจเยน ปจฺจโย. (1)}

\prose{สญฺโญชน\-วิปฺปยุตฺโต ธมฺโม สญฺโญชน\-สมฺปยุตฺตสฺส ธมฺมสฺส อารมฺมณ\-ปจฺจเยน ปจฺจโย. สหชาต\-ปจฺจเยน ปจฺจโย. อุปนิสฺสย\-ปจฺจเยน ปจฺจโย. ปุเรชาต\-ปจฺจเยน ปจฺจโย. (2)}

\prose{สญฺโญชน\-วิปฺปยุตฺโต ธมฺโม สญฺโญชน\-สมฺปยุตฺตสฺส จ สญฺโญชน\-วิปฺปยุตฺตสฺส จ ธมฺมสฺส อารมฺมณ\-ปจฺจเยน ปจฺจโย. สหชาต\-ปจฺจเยน ปจฺจโย. อุปนิสฺสย\-ปจฺจเยน ปจฺจโย. ปุเรชาต\-ปจฺจเยน ปจฺจโย. (3)}

\proseitem{93}{สญฺโญชน\-สมฺปยุตฺโต จ สญฺโญชน\-วิปฺปยุตฺโต จ ธมฺมา สญฺโญชน\-สมฺปยุตฺตสฺส ธมฺมสฺส อารมฺมณ\-ปจฺจเยน ปจฺจโย. สหชาต\-ปจฺจเยน ปจฺจโย. อุปนิสฺสย\-ปจฺจเยน ปจฺจโย. (1)}

\prose{\csromanfolio{244}สญฺโญชน\-สมฺปยุตฺโต จ สญฺโญชน\-วิปฺปยุตฺโต จ ธมฺมา สญฺโญชน\-วิปฺปยุตฺตสฺส ธมฺมสฺส อารมฺมณ\-ปจฺจเยน ปจฺจโย. สหชาต\-ปจฺจเยน ปจฺจโย. อุปนิสฺสย\-ปจฺจเยน ปจฺจโย. ปจฺฉาชาต\-ปจฺจเยน ปจฺจโย. (2)}

\prose{สญฺโญชน\-สมฺปยุตฺโต จ สญฺโญชน\-วิปฺปยุตฺโต จ ธมฺมา สญฺโญชน\-สมฺปยุตฺตสฺส จ สญฺโญชน\-วิปฺปยุตฺตสฺส จ ธมฺมสฺส อารมฺมณ\-ปจฺจเยน ปจฺจโย. สหชาต\-ปจฺจเยน ปจฺจโย. อุปนิสฺสย\-ปจฺจเยน ปจฺจโย. (3)}

\csromanmarkh{2. ปจฺจย\-ปจฺจนีย}\csromanmarkhii{2. สงฺขฺยาวาร}\csromanheadernomark{2}{2. ปจฺจย\-ปจฺจนีย 2. สงฺขฺยาวาร}
\vspace{\tipitakaparskip}
\csromancenter{นเหตุยา นว, นอารมฺมเณ นว, (สพฺพตฺถ นว.) โนวิคเต นว, โนอวิคเต นว.}
\csromancenter{\textbf{เหตุ}ปจฺจยา นอารมฺมเณ ฉ ฯเปฯ นสมนนฺตเร ฉ. นอญฺญมญฺเญ เทฺว, นฺเอาปนิสฺสเย ฉ ฯเปฯ นมคฺเค ฉ, นสมฺปยุตฺเต เทฺว, นวิปฺปยุตฺเต ตีณิ, โนนตฺถิยา ฉ, โนวิคเต ฉ.}
\proseitemwithcloser{96}{นเหตุปจฺจยา อารมฺมเณ นว, อธิปติยา ปญฺจ, (อนุโลมปทานิ คณิตพฺพานิ.) อวิคเต นว.}{สญฺโญชน\-สมฺปยุตฺตทุกํ นิฏฺฐิตํ.}
\csromanmatikaanchor{mtk.28}
\csromantocmarkref{section}{23. สญฺโญชนสญฺโญชนิยทุก}{mtk.28}
\bookmark[dest={mtk.28},level=1]{23. สญฺโญชนสญฺโญชนิยทุก}
\csromanmarkh{23. สญฺโญชน\-สญฺโญชนิยทุก}\csromanmarkhii{1. ปฏิจฺจวาร}\csromanheadernomark{1}{23. สญฺโญชน\-สญฺโญชนิยทุก 1. ปฏิจฺจวาร}
\csromanmarkh{1. ปจฺจยานุโลม}\csromanmarkhii{1. วิภงฺควาร}\csromanheadernomark{2}{1. \textbf{ปจฺจยานุโลม 1. วิภงฺควาร}}
\csromanheader{2}{\textbf{เหตุ}}
\proseitem{97}{สญฺโญชนญฺเจว สญฺโญชนิยญฺจ ธมฺมํ ปฏิจฺจ สญฺโญชโน เจว สญฺโญชนิโย จ ธมฺโม อุปฺปชฺชติ เหตุปจฺจยา. กามราค\-สญฺโญชนํ ปฏิจฺจ ทิฏฺฐิ\-สญฺโญชนํ อวิชฺชา\-สญฺโญชนํ. (จกฺกํ.) (1)}

\csromanheader{2}{23. สญฺโญชน\-สญฺโญชนิยทุก}
\prose{\csromanfolio{245}สญฺโญชนญฺเจว สญฺโญชนิยญฺจ ธมฺมํ ปฏิจฺจ สญฺโญชนิโย เจว โน จ สญฺโญชโน ธมฺโม อุปฺปชฺชติ เหตุปจฺจยา. สญฺโญชเน ปฏิจฺจ สมฺปยุตฺตกา ขนฺธา จิตฺต\-สมุฏฺฐานญฺจ รูปํ. (2)}

\prose{สญฺโญชนญฺเจว สญฺโญชนิยญฺจ ธมฺมํ ปฏิจฺจ สญฺโญชโน เจว สญฺโญชนิโย จ สญฺโญชนิโย เจว โน จ สญฺโญชโน จ ธมฺมา อุปฺปชฺชนฺติ เหตุปจฺจยา. กามราค\-สญฺโญชนํ ปฏิจฺจ ทิฏฺฐิ\-สญฺโญชนํ อวิชฺชา\-สญฺโญชนํ สมฺปยุตฺตกา ขนฺธา จิตฺต\-สมุฏฺฐานญฺจ รูปํ. (จกฺกํ.) (3)}

\proseitem{98}{สญฺโญชนิย\-ญฺเจว โน จ สญฺโญชนํ ธมฺมํ ปฏิจฺจ สญฺโญชนิโย เจว โน จ สญฺโญชโน ธมฺโม อุปฺปชฺชติ เหตุปจฺจยา. สญฺโญชนิย\-ญฺเจว โน จ สญฺโญชนํ เอกํ ขนฺธํ ปฏิจฺจ ตโย ขนฺธา จิตฺต\-สมุฏฺฐานญฺจ รูปํ ฯเปฯ เทฺว ขนฺเธ ฯเปฯ. ปฏิสนฺธิ\-กฺขเณ ฯเปฯ. (ยาว มหาภูตา.) (1)}

\prose{สญฺโญชนิย\-ญฺเจว โน จ สญฺโญชนํ ธมฺมํ ปฏิจฺจ สญฺโญชโน เจว สญฺโญชนิโย จ ธมฺโม อุปฺปชฺชติ เหตุปจฺจยา. สญฺโญชนิเย เจว โน จ สญฺโญชเน ขนฺเธ ปฏิจฺจ สญฺโญชนา. (2)}

\prose{สญฺโญชนิย\-ญฺเจว โน จ สญฺโญชนํ ธมฺมํ ปฏิจฺจ สญฺโญชโน เจว สญฺโญชนิโย จ สญฺโญชนิโย เจว โน จ สญฺโญชโน จ ธมฺมา อุปฺปชฺชนฺติ เหตุปจฺจยา. สญฺโญชนิย\-ญฺเจว โน จ สญฺโญชนํ เอกํ ขนฺธํ ปฏิจฺจ ตโย ขนฺธา สญฺโญชนา จ จิตฺต\-สมุฏฺฐานญฺจ รูปํ ฯเปฯ เทฺว ขนฺเธ ฯเปฯ. (3)}

\proseitem{99}{สญฺโญชนญฺเจว สญฺโญชนิยญฺจ สญฺโญชนิย\-ญฺเจว โน จ สญฺโญชนํ ธมฺมํ ปฏิจฺจ สญฺโญชโน เจว สญฺโญชนิโย จ ธมฺโม อุปฺปชฺชติ เหตุปจฺจยา. กามราคสญฺโญชนญฺจ สมฺปยุตฺตเก จ ขนฺเธ ปฏิจฺจ ทิฏฺฐิ\-สญฺโญชนํ อวิชฺชา\-สญฺโญชนํ. (จกฺกํ.) (1)}

\prose{สญฺโญชนญฺเจว สญฺโญชนิยญฺจ สญฺโญชนิย\-ญฺเจว โน จ สญฺโญชนญฺจ ธมฺมํ ปฏิจฺจ สญฺโญชนิโย เจว โน จ สญฺโญชโน ธมฺโม อุปฺปชฺชติ เหตุปจฺจยา. สญฺโญชนิย\-ญฺเจว โน จ สญฺโญชนํ เอกํ ขนฺธญฺจ สญฺโญชเน จ ปฏิจฺจ ตโย ขนฺธา จิตฺต\-สมุฏฺฐานญฺจ รูปํ ฯเปฯ เทฺว ขนฺเธ ฯเปฯ. (2)}

\prose{สญฺโญชนญฺเจว สญฺโญชนิยญฺจ สญฺโญชนิย\-ญฺเจว โน จ สญฺโญชนญฺจ ธมฺมํ ปฏิจฺจ สญฺโญชโน เจว สญฺโญชนิโย จ สญฺโญชนิโย \csromanfolio{246}เจว โน จ สญฺโญชโน จ ธมฺมา อุปฺปชฺชนฺติ เหตุปจฺจยา. สญฺโญชนิย\-ญฺเจว โน จ สญฺโญชนํ เอกํ ขนฺธญฺจ กามราคสญฺโญชนญฺจ ปฏิจฺจ ตโย ขนฺธา ทิฏฺฐิ\-สญฺโญชนํ อวิชฺชา\-สญฺโญชนํ จิตฺต\-สมุฏฺฐานญฺจ รูปํ ฯเปฯ เทฺว ขนฺเธ จ ฯเปฯ. (จกฺกํ พนฺธิตพฺพํ.) (3) (สญฺโญชน\-โคจฺฉเก ปฐม\-ทุก\-สทิสํ.)}

\prosewithcloser{(เอวํ อิมมฺปิ ทุกํ วิตฺถาเรตพฺพํ.}{นินฺนานากรณํ ฐเปตฺวา โลกุตฺตรํ.) สญฺโญชน\-สญฺโญชนิยทุกํ นิฏฺฐิตํ.}
\csromanmatikaanchor{mtk.29}
\csromantocmarkref{section}{24. สญฺโญชนสญฺโญชนสมฺปยุตฺตทุก}{mtk.29}
\bookmark[dest={mtk.29},level=1]{24. สญฺโญชนสญฺโญชนสมฺปยุตฺตทุก}
\csromanmarkh{24. สญฺโญชน\-สญฺโญชน\-สมฺปยุตฺตทุก}\csromanmarkhii{1. ปฏิจฺจวาร}\csromanheadernomark{1}{24. สญฺโญชน\-สญฺโญชน\-สมฺปยุตฺตทุก 1. ปฏิจฺจวาร}
\csromanmarkh{1. ปจฺจยานุโลม}\csromanmarkhii{1. วิภงฺควาร}\csromanheadernomark{2}{1. \textbf{ปจฺจยานุโลม 1. วิภงฺควาร}}
\csromanheader{2}{\textbf{เหตุ}}
\proseitem{100}{สญฺโญชนญฺเจว สญฺโญชน\-สมฺปยุตฺตญฺจ ธมฺมํ ปฏิจฺจ สญฺโญชโน เจว สญฺโญชน\-สมฺปยุตฺโต จ ธมฺโม อุปฺปชฺชติ เหตุปจฺจยา. กามราค\-สญฺโญชนํ ปฏิจฺจ ทิฏฺฐิ\-สญฺโญชนํ อวิชฺชา\-สญฺโญชนํ. (จกฺกํ.) (1)}

\prose{สญฺโญชนญฺเจว สญฺโญชน\-สมฺปยุตฺตญฺจ ธมฺมํ ปฏิจฺจ สญฺโญชน\-สมฺปยุตฺโต เจว โน จ สญฺโญชโน ธมฺโม อุปฺปชฺชติ เหตุปจฺจยา. สญฺโญชเน ปฏิจฺจ สมฺปยุตฺตกา ขนฺธา. (2)}

\prose{สญฺโญชนญฺเจว สญฺโญชน\-สมฺปยุตฺตญฺจ ธมฺมํ ปฏิจฺจ สญฺโญชโน เจว สญฺโญชน\-สมฺปยุตฺโต จ สญฺโญชน\-สมฺปยุตฺโต เจว โน จ สญฺโญชโน จ ธมฺมา อุปฺปชฺชนฺติ เหตุปจฺจยา. กามราค\-สญฺโญชนํ ปฏิจฺจ ทิฏฺฐิ\-สญฺโญชนํ อวิชฺชา\-สญฺโญชนํ สมฺปยุตฺตกา จ ขนฺธา. (จกฺกํ.) (3)}

\proseitem{101}{สญฺโญชน\-สมฺปยุตฺต\-ญฺเจว โน จ สญฺโญชนํ ธมฺมํ ปฏิจฺจ สญฺโญชน\-สมฺปยุตฺโต เจว โน จ สญฺโญชโน ธมฺโม อุปฺปชฺชติ เหตุปจฺจยา. สญฺโญชน\-สมฺปยุตฺต\-ญฺเจว โน จ สญฺโญชนํ เอกํ ขนฺธํ ปฏิจฺจ ตโย ขนฺธา ฯเปฯ เทฺว ขนฺเธ ฯเปฯ. (1)}

\prose{\csromanfolio{247}สญฺโญชน\-สมฺปยุตฺต\-ญฺเจว โน จ สญฺโญชนํ ธมฺมํ ปฏิจฺจ สญฺโญชโน เจว สญฺโญชน\-สมฺปยุตฺโต จ ธมฺโม อุปฺปชฺชติ เหตุปจฺจยา. สญฺโญชน\-สมฺปยุตฺเต เจว โน จ สญฺโญชเน ขนฺเธ ปฏิจฺจ สญฺโญชนา. (2)}

\prose{สญฺโญชน\-สมฺปยุตฺต\-ญฺเจว โน จ สญฺโญชนํ ธมฺมํ ปฏิจฺจ สญฺโญชโน เจว สญฺโญชน\-สมฺปยุตฺโต จ สญฺโญชน\-สมฺปยุตฺโต เจว โน จ สญฺโญชโน จ ธมฺมา อุปฺปชฺชนฺติ เหตุปจฺจยา. สญฺโญชน\-สมฺปยุตฺต\-ญฺเจว โน จ สญฺโญชนํ เอกํ ขนฺธํ ปฏิจฺจ ตโย ขนฺธา สญฺโญชนา จ ฯเปฯ เทฺว ขนฺเธ. (3)}

\proseitem{102}{สญฺโญชนญฺเจว สญฺโญชน\-สมฺปยุตฺตญฺจ สญฺโญชน\-สมฺปยุตฺต\-ญฺเจว โน จ สญฺโญชนญฺจ ธมฺมํ ปฏิจฺจ สญฺโญชโน เจว สญฺโญชน\-สมฺปยุตฺโต จ ธมฺโม อุปฺปชฺชติ เหตุปจฺจยา. กามราคสญฺโญชนญฺจ สมฺปยุตฺตเก จ ขนฺเธ ปฏิจฺจ ทิฏฺฐิ\-สญฺโญชนํ อวิชฺชา\-สญฺโญชนํ. (จกฺกํ.) (1)}

\prose{สญฺโญชนญฺเจว สญฺโญชน\-สมฺปยุตฺตญฺจ สญฺโญชน\-สมฺปยุตฺต\-ญฺเจว โน จ สญฺโญชนญฺจ ธมฺมํ ปฏิจฺจ สญฺโญชน\-สมฺปยุตฺโต เจว โน จ สญฺโญชโน ธมฺโม อุปฺปชฺชติ เหตุปจฺจยา. สญฺโญชน\-สมฺปยุตฺต\-ญฺเจว โน จ สญฺโญชนํ เอกํ ขนฺธญฺจ สญฺโญชเน จ ปฏิจฺจ ตโย ขนฺธา ฯเปฯ เทฺว ขนฺเธ จ ฯเปฯ. (2)}

\prose{สญฺโญชนญฺเจว สญฺโญชน\-สมฺปยุตฺตญฺจ สญฺโญชน\-สมฺปยุตฺต\-ญฺเจว โน จ สญฺโญชนญฺจ ธมฺมํ ปฏิจฺจ สญฺโญชโน เจว สญฺโญชน\-สมฺปยุตฺโต จ สญฺโญชน\-สมฺปยุตฺโต เจว โน จ สญฺโญชโน จ ธมฺมา อุปฺปชฺชนฺติ เหตุปจฺจยา. สญฺโญชน\-สมฺปยุตฺต\-ญฺเจว โน จ สญฺโญชนํ เอกํ ขนฺธญฺจ กามราคสญฺโญชนญฺจ ปฏิจฺจ ตโย ขนฺธา ทิฏฺฐิ\-สญฺโญชนํ อวิชฺชา\-สญฺโญชนํ ฯเปฯ เทฺว ขนฺเธ จ ฯเปฯ. (จกฺกํ.) (3)}

\csromanmarkh{1. ปจฺจยานุโลม}\csromanmarkhii{2. สงฺขฺยาวาร}\csromanheadernomark{2}{1. \textbf{ปจฺจยานุโลม 2. สงฺขฺยาวาร}}
\proseitem{103}{เหตุยา นว, อารมฺมเณ นว, อธิปติยา นว, (สพฺพตฺถ นว.) กมฺเม นว, อาหาเร นว ฯเปฯ อวิคเต นว.}

\vspace{\tipitakaparskip}
\csromancenter{อนุโลมํ.}
\csromanmarkh{2. ปจฺจย\-ปจฺจนีย}\csromanmarkhii{1. วิภงฺควาร}\csromanheadernomark{2}{2. ปจฺจย\-ปจฺจนีย 1. วิภงฺควาร}
\proseitem{104}{\csromanfolio{248}สญฺโญชนญฺเจว สญฺโญชน\-สมฺปยุตฺตญฺจ ธมฺมํ ปฏิจฺจ สญฺโญชโน เจว สญฺโญชน\-สมฺปยุตฺโต จ ธมฺโม อุปฺปชฺชติ นเหตุปจฺจยา. วิจิกิจฺฉา\-สญฺโญชนํ ปฏิจฺจ อวิชฺชา\-สญฺโญชนํ. (1)}

\prose{สญฺโญชน\-สมฺปยุตฺต\-ญฺเจว โน จ สญฺโญชนํ ธมฺมํ ปฏิจฺจ สญฺโญชโน เจว สญฺโญชน\-สมฺปยุตฺโต จ ธมฺโม อุปฺปชฺชติ นเหตุปจฺจยา. วิจิกิจฺฉา\-สหคเต ขนฺเธ ปฏิจฺจ วิจิกิจฺฉา\-สหคโต โมโห. (1)}

\prose{สญฺโญชนญฺเจว สญฺโญชน\-สมฺปยุตฺตญฺจ สญฺโญชน\-สมฺปยุตฺต\-ญฺเจว โน จ สญฺโญชนญฺจ ธมฺมํ ปฏิจฺจ สญฺโญชโน เจว สญฺโญชน\-สมฺปยุตฺโต จ ธมฺโม อุปฺปชฺชติ นเหตุปจฺจยา. วิจิกิจฺฉาสญฺโญชนญฺจ สมฺปยุตฺตเก จ ขนฺเธ ปฏิจฺจ อวิชฺชา\-สญฺโญชนํ. (1)}

\csromanmarkh{2. ปจฺจย\-ปจฺจนีย}\csromanmarkhii{2. สงฺขฺยาวาร}\csromanheadernomark{2}{2. ปจฺจย\-ปจฺจนีย 2. สงฺขฺยาวาร}
\vspace{\tipitakaparskip}
\csromancenter{\textbf{นเหตุ}ยา ตีณิ, นอธิปติยา นว, นปุเรชาเต นว, นปจฺฉาชาเต นว, นอาเสวเน นว, นกมฺเม ตีณิ, นวิปาเก นว, นวิปฺปยุตฺเต นว.}
\prose{(เอวํ อิตเร เทฺว คณนาปิ สหชาตวาโรปิ กาตพฺโพ. ปจฺจยวาโรปิ นิสฺสยวาโรปิ สํสฏฺฐ\-วาโรปิ สมฺปยุตฺตวาโรปิ ปฏิจฺจ\-วาร\-สทิสา.)}

\prose{24. สญฺโญชน\-สญฺโญชน\-สมฺปยุตฺตทุก 7. ปญฺหาวาร}

\csromanmarkh{1. ปจฺจยานุโลม}\csromanmarkhii{1. วิภงฺควาร}\csromanheadernomark{2}{1. \textbf{ปจฺจยานุโลม 1. วิภงฺควาร}}
\csromanheader{2}{\textbf{เหตุ}}
\proseitem{106}{สญฺโญชโน เจว สญฺโญชน\-สมฺปยุตฺโต จ ธมฺโม สญฺโญชนสฺส เจว สญฺโญชน\-สมฺปยุตฺตสฺส จ ธมฺมสฺส เหตุปจฺจเยน ปจฺจโย \csromanfolio{249}กามราค\-สญฺโญ\-ชโน ทิฏฺฐิ\-สญฺโญชนสฺส อวิชฺชา\-สญฺโญชนสฺส เหตุปจฺจเยน ปจฺจโย. (จกฺกํ.) (1)}

\prose{สญฺโญชโน เจว สญฺโญชน\-สมฺปยุตฺโต จ ธมฺโม สญฺโญชน\-สมฺปยุตฺตสฺส เจว โน จ สญฺโญชนสฺส ธมฺมสฺส เหตุปจฺจเยน ปจฺจโย. สญฺโญชนา เจว สญฺโญชน\-สมฺปยุตฺตา จ เหตู สมฺปยุตฺตกานํ ขนฺธานํ เหตุปจฺจเยน ปจฺจโย. (2)}

\prose{สญฺโญชโน เจว สญฺโญชน\-สมฺปยุตฺโต จ ธมฺโม สญฺโญชนสฺส เจว สญฺโญชน\-สมฺปยุตฺตสฺส จ สญฺโญชน\-สมฺปยุตฺตสฺส เจว โน จ สญฺโญชนสฺส จ ธมฺมสฺส เหตุปจฺจเยน ปจฺจโย. กามราค\-สญฺโญ\-ชโน ทิฏฺฐิ\-สญฺโญชนสฺส อวิชฺชา\-สญฺโญชนสฺส สมฺปยุตฺตกานญฺจ ขนฺธานํ เหตุปจฺจเยน ปจฺจโย. (จกฺกํ.) (3)}

\proseitem{107}{สญฺโญชโน เจว สญฺโญชน\-สมฺปยุตฺโต จ ธมฺโม สญฺโญชนสฺส เจว สญฺโญชน\-สมฺปยุตฺตสฺส จ ธมฺมสฺส อารมฺมณ\-ปจฺจเยน ปจฺจโย. สญฺโญชเน อารพฺภ สญฺโญชนา อุปฺปชฺชนฺติ. (มูลํ กาตพฺพํ.) สญฺโญชเน อารพฺภ สญฺโญชน\-สมฺปยุตฺตา เจว โน จ สญฺโญชนา ขนฺธา อุปฺปชฺชนฺติ. (มูลํ กาตพฺพํ.) สญฺโญชเน อารพฺภ สญฺโญชนา จ สญฺโญชน\-สมฺปยุตฺตกา จ ขนฺธา อุปฺปชฺชนฺติ. (1)}

\prose{สญฺโญชน\-สมฺปยุตฺโต เจว โน จ สญฺโญชโน ธมฺโม สญฺโญชน\-สมฺปยุตฺตสฺส เจว โน จ สญฺโญชนสฺส ธมฺมสฺส อารมฺมณ\-ปจฺจเยน ปจฺจโย. สญฺโญชน\-สมฺปยุตฺเต เจว โน จ สญฺโญชเน ขนฺเธ อารพฺภ สญฺโญชน\-สมฺปยุตฺตา เจว โน จ สญฺโญชนา ขนฺธา อุปฺปชฺชนฺติ. (1)}

\prose{สญฺโญชน\-สมฺปยุตฺโต เจว โน จ สญฺโญชโน ธมฺโม สญฺโญชนสฺส เจว สญฺโญชน\-สมฺปยุตฺตสฺส จ ธมฺมสฺส อารมฺมณ\-ปจฺจเยน ปจฺจโย. สญฺโญชน\-สมฺปยุตฺเต เจว โน จ สญฺโญชเน ขนฺเธ อารพฺภ สญฺโญชนา อุปฺปชฺชนฺติ. (2)}

\prose{สญฺโญชน\-สมฺปยุตฺโต เจว โน จ สญฺโญชโน ธมฺโม สญฺโญชนสฺส เจว สญฺโญชน\-สมฺปยุตฺตสฺส จ สญฺโญชน\-สมฺปยุตฺตสฺส เจว โน จ สญฺโญชนสฺส จ ธมฺมสฺส อารมฺมณ\-ปจฺจเยน ปจฺจโย. สญฺโญชน\-สมฺปยุตฺเต \csromanfolio{250}เจว โน จ สญฺโญชเน ขนฺเธ อารพฺภ สญฺโญชนา จ สญฺโญชน\-สมฺปยุตฺตกา จ ขนฺธา อุปฺปชฺชนฺติ. (3)}

\prose{สญฺโญชโน เจว สญฺโญชน\-สมฺปยุตฺโต จ สญฺโญชน\-สมฺปยุตฺโต เจว โน จ สญฺโญชโน จ ธมฺมา สญฺโญชนสฺส เจว สญฺโญชน\-สมฺปยุตฺตสฺส จ ธมฺมสฺส อารมฺมณ\-ปจฺจเยน ปจฺจโย. ตีณิ.}

\proseitem{108}{สญฺโญชโน เจว สญฺโญชน\-สมฺปยุตฺโต จ ธมฺโม สญฺโญชนสฺส เจว สญฺโญชน\-สมฺปยุตฺตสฺส จ ธมฺมสฺส อธิปติ\-ปจฺจเยน ปจฺจโย. อารมฺมณา\-ธิปติ. ตีณิ.}

\prose{สญฺโญชน\-สมฺปยุตฺโต เจว โน จ สญฺโญชโน ธมฺโม สญฺโญชน\-สมฺปยุตฺตสฺส เจว โน จ สญฺโญชนสฺส ธมฺมสฺส อธิปติ\-ปจฺจเยน ปจฺจโย. อารมฺมณา\-ธิปติ, สหชาตา\-ธิปติ. ตีณิ.}

\vspace{\tipitakaparskip}
\csromancenter{(อิมาสุ ตีสุปิ ปญฺหาสุ อารมฺมณาธิปติปิ สหชาตาธิปติปิ กาตพฺพา.)}
\prose{สญฺโญชโน เจว สญฺโญชน\-สมฺปยุตฺโต จ สญฺโญชน\-สมฺปยุตฺโต เจว โน จ สญฺโญชโน จ ธมฺมา สญฺโญชนสฺส เจว สญฺโญชน\-สมฺปยุตฺตสฺส จ ธมฺมสฺส อธิปติ\-ปจฺจเยน ปจฺจโย. อารมฺมณา\-ธิปติ. ตีณิ.}

\proseitem{109}{สญฺโญชโน เจว สญฺโญชน\-สมฺปยุตฺโต จ ธมฺโม สญฺโญชนสฺส เจว สญฺโญชน\-สมฺปยุตฺตสฺส จ ธมฺมสฺส อนนฺตร\-ปจฺจเยน ปจฺจโย. นว. (นินฺนานากรณํ, วิภชนา นตฺถิ. อารมฺมณสทิสา.) สมนนฺตร\-ปจฺจเยน ปจฺจโย. นว. สหชาต\-ปจฺจเยน ปจฺจโย. นว. อญฺญมญฺญ\-ปจฺจเยน ปจฺจโย. นว. นิสฺสย\-ปจฺจเยน ปจฺจโย. นว. อุปนิสฺสย\-ปจฺจเยน ปจฺจโย. นว. (อารมฺมณนเยน กาตพฺพา.) อาเสวน\-ปจฺจเยน ปจฺจโย. นว.}

\csromanheader{2}{\textbf{กมฺมาทิ}}
\proseitem{110}{สญฺโญชน\-สมฺปยุตฺโต เจว โน จ สญฺโญชโน ธมฺโม สญฺโญชน\-สมฺปยุตฺตสฺส เจว โน จ สญฺโญชนสฺส ธมฺมสฺส กมฺมปจฺจเยน ปจฺจโย. ตีณิ. อาหารปจฺจเยน ปจฺจโย. ตีณิ. อินฺทฺริย\-ปจฺจเยน ปจฺจโย. ตีณิ. ฌานปจฺจเยน ปจฺจโย. ตีณิ. มคฺคปจฺจเยน ปจฺจโย. \csromanfolio{251}นว. สมฺปยุตฺต\-ปจฺจเยน ปจฺจโย. นว. อตฺถิปจฺจเยน ปจฺจโย. นว. นตฺถิ\-ปจฺจเยน ปจฺจโย. นว. วิคต\-ปจฺจเยน ปจฺจโย. นว. อวิคต\-ปจฺจเยน ปจฺจโย. นว.}

\csromanmarkh{1. ปจฺจยานุโลม}\csromanmarkhii{2. สงฺขฺยาวาร}\csromanheadernomark{2}{1. \textbf{ปจฺจยานุโลม 2. สงฺขฺยาวาร}}
\proseitem{111}{เหตุยา ตีณิ, อารมฺมเณ นว, อธิปติยา นว, อนนฺตเร นว, สมนนฺตเร นว, สหชาเต นว, อญฺญมญฺเญ นว, นิสฺสเย นว, อุปนิสฺสเย นว, อาเสวเน นว, กมฺเม ตีณิ, อาหาเร ตีณิ, อินฺทฺริเย ตีณิ, ฌาเน ตีณิ, มคฺเค นว, สมฺปยุตฺเต นว, อตฺถิยา นว, นตฺถิยา นว, วิคเต นว, อวิคเต นว.}

\vspace{\tipitakaparskip}
\csromancenter{\textbf{อนุโลมํ.}}
\proseitem{112}{สญฺโญชโน เจว สญฺโญชน\-สมฺปยุตฺโต จ ธมฺโม สญฺโญชนสฺส เจว สญฺโญชน\-สมฺปยุตฺตสฺส จ ธมฺมสฺส อารมฺมณ\-ปจฺจเยน ปจฺจโย. สหชาต\-ปจฺจเยน ปจฺจโย. อุปนิสฺสย\-ปจฺจเยน ปจฺจโย. (สํขิตฺตํ. เอวํ นว ปญฺหา กาตพฺพา. ตีสุเยว ปเทสุ ปริวตฺเตตพฺพา. นานากฺขณิกา นตฺถิ.)}

\csromanmarkh{2. ปจฺจย\-ปจฺจนีย}\csromanmarkhii{2. สงฺขฺยาวาร}\csromanheadernomark{2}{2. ปจฺจย\-ปจฺจนีย 2. สงฺขฺยาวาร}
\vspace{\tipitakaparskip}
\csromancenter{นเหตุยา นว, นอารมฺมเณ นว, (สพฺพตฺถ นว.) โนอวิคเต นว.}
\proseitem{114}{เหตุปจฺจยา นอารมฺมเณ ตีณิ ฯเปฯ นสมนนฺตเร ตีณิ, นอุปนิสฺสเย ตีณิ, นปุเรชาเต ตีณิ ฯเปฯ นมคฺเค ตีณิ, นสมฺปยุตฺเต ตีณิ, นวิปฺปยุตฺเต ตีณิ, โนนตฺถิยา ตีณิ, โนวิคเต ตีณิ.}

\proseitemwithcloser{115}{\csromanfolio{252}นเหตุปจฺจยา อารมฺมเณ นว, อธิปติยา นว, (อนุโลมปทานิ กาตพฺพานิ.) ฯเปฯ อวิคเต นว.}{สญฺโญชน\-สญฺโญชน\-สมฺปยุตฺตทุกํ นิฏฺฐิตํ.}
\csromanmatikaanchor{mtk.30}
\csromantocmarkref{section}{25. สญฺโญชนวิปฺปยุตฺตสญฺโญชนิยทุก}{mtk.30}
\bookmark[dest={mtk.30},level=1]{25. สญฺโญชนวิปฺปยุตฺตสญฺโญชนิยทุก}
\csromanmarkh{25. สญฺโญชน\-วิปฺปยุตฺต\-สญฺโญชนิยทุก}\csromanmarkhii{1. ปฏิจฺจวาร}\csromanheadernomark{1}{25. สญฺโญชน\-วิปฺปยุตฺต\-สญฺโญชนิยทุก 1. ปฏิจฺจวาร}
\csromanmarkh{1. ปจฺจยานุโลม}\csromanmarkhii{1. วิภงฺควาร}\csromanheadernomark{2}{1. \textbf{ปจฺจยานุโลม 1. วิภงฺควาร}}
\proseitem{116}{สญฺโญชน\-วิปฺปยุตฺตํ สญฺโญชนิยํ ธมฺมํ ปฏิจฺจ สญฺโญชน\-วิปฺปยุตฺโต สญฺโญชนิโย ธมฺโม อุปฺปชฺชติ เหตุปจฺจยา. สญฺโญชน\-วิปฺปยุตฺตํ สญฺโญชนิยํ เอกํ ขนฺธํ ปฏิจฺจ ตโย ขนฺธา จิตฺต\-สมุฏฺฐานญฺจ รูปํ ฯเปฯ เทฺว ขนฺเธ ฯเปฯ. ปฏิสนฺธิ\-กฺขเณ ฯเปฯ. (ยาว อสญฺญสตฺตา. มหาภูตา.)}

\prose{สญฺโญชน\-วิปฺปยุตฺตํ อสญฺโญชนิยํ ธมฺมํ ปฏิจฺจ สญฺโญชน\-วิปฺปยุตฺโต อสญฺโญชนิโย ธมฺโม อุปฺปชฺชติ เหตุปจฺจยา. สญฺโญชน\-วิปฺปยุตฺตํ อสญฺโญชนิยํ เอกํ ขนฺธํ ปฏิจฺจ ตโย ขนฺธา ฯเปฯ เทฺว ขนฺเธ ฯเปฯ.}

\prose{สญฺโญชน\-วิปฺปยุตฺตํ อสญฺโญชนิยํ ธมฺมํ ฯเปฯ. (เทฺว ปญฺหา กาตพฺพา.)}

\prosewithcloser{(อิมํ ทุกํ จูฬนฺตรทุเก โลกิย\-ทุก\-สทิสํ, นินฺนานากรณํ.)}{สญฺโญชน\-วิปฺปยุตฺต\-สญฺโญชนิยทุกํ นิฏฺฐิตํ.}
\nitthitam{สญฺโญชน\-โคจฺฉกํ นิฏฺฐิตํ.}
\csromanmatikaanchor{mtk.31}
\csromanmatikaanchor{mtk.32}
\csromantocmarknopage{chapter}{5. คนฺถโคจฺฉก}{mtk.31}
\bookmark[dest={mtk.31},level=0]{5. คนฺถโคจฺฉก}
\csromantocmarkref{section}{26. คนฺถทุก}{mtk.32}
\bookmark[dest={mtk.32},level=1]{26. คนฺถทุก}
\csromanmarkh{26. คนฺถทุก}\csromanmarkhii{5. คนฺถโคจฺฉก}\csromanheadernomark{1}{26. คนฺถทุก \textbf{5. คนฺถโคจฺฉก}}
\csromanmarkh{26. คนฺถทุก}\csromanmarkhii{1. ปฏิจฺจวาร}\csromanheadernomark{2}{26. คนฺถทุก 1. ปฏิจฺจวาร}
\csromanmarkh{1. ปจฺจยานุโลม}\csromanmarkhii{1. วิภงฺควาร}\csromanheadernomark{2}{1. \textbf{ปจฺจยานุโลม 1. วิภงฺควาร}}
\csromanheader{2}{\textbf{เหตุ}}
\proseitem{1}{\csromanfolio{253}คนฺถํ ธมฺมํ ปฏิจฺจ คนฺโถ ธมฺโม อุปฺปชฺชติ เหตุปจฺจยา. สีลพฺพต\-ปรามาสํ กายคนฺถํ ปฏิจฺจ อภิชฺฌา\-กายคนฺโถ. อภิชฺฌา\-กายคนฺถํ ปฏิจฺจ สีลพฺพต\-ปรามาโส กายคนฺโถ. อิทํสจฺจา\-ภินิเวสํ กายคนฺถํ ปฏิจฺจ อภิชฺฌา\-กายคนฺโถ. อภิชฺฌา\-กายคนฺถํ ปฏิจฺจ อิทํสจฺจา\-ภินิเวโส กายคนฺโถ. (1)}

\prose{คนฺถํ ธมฺมํ ปฏิจฺจ โนคนฺโถ ธมฺโม อุปฺปชฺชติ เหตุปจฺจยา. คนฺเถ ปฏิจฺจ สมฺปยุตฺตกา ขนฺธา จิตฺต\-สมุฏฺฐานญฺจ รูปํ. (2)}

\prose{คนฺถํ ธมฺมํ ปฏิจฺจ คนฺโถ จ โนคนฺโถ จ ธมฺมา อุปฺปชฺชนฺติ เหตุปจฺจยา. สีลพฺพต\-ปรามาสํ กายคนฺถํ ปฏิจฺจ อภิชฺฌา\-กายคนฺโถ สมฺปยุตฺตกา จ ขนฺธา จิตฺต\-สมุฏฺฐานญฺจ รูปํ. (จกฺกํ.) (3)}

\proseitem{2}{โนคนฺถํ ธมฺมํ ปฏิจฺจ โนคนฺโถ ธมฺโม อุปฺปชฺชติ เหตุปจฺจยา. โนคนฺถํ เอกํ ขนฺธํ ปฏิจฺจ ตโย ขนฺธา จิตฺต\-สมุฏฺฐานญฺจ รูปํ ฯเปฯ เทฺว ขนฺเธ ฯเปฯ. ปฏิสนฺธิ\-กฺขเณ ฯเปฯ. ขนฺเธ ปฏิจฺจ วตฺถุ, วตฺถุํ ปฏิจฺจ ขนฺธา. เอกํ มหาภูตํ ฯเปฯ. (1)}

\prose{โนคนฺถํ ธมฺมํ ปฏิจฺจ คนฺโถ ธมฺโม อุปฺปชฺชติ เหตุปจฺจยา. โนคนฺเถ ขนฺเธ ปฏิจฺจ คนฺถา. (2)}

\prose{โนคนฺถํ ธมฺมํ ปฏิจฺจ คนฺโถ จ โนคนฺโถ จ ธมฺมา อุปฺปชฺชนฺติ เหตุปจฺจยา. โนคนฺถํ เอกํ ขนฺธํ ปฏิจฺจ ตโย ขนฺธา คนฺถา จ จิตฺต\-สมุฏฺฐานญฺจ รูปํ ฯเปฯ เทฺว ขนฺเธ ฯเปฯ. (3)}

\proseitem{3}{คนฺถญฺจ โนคนฺถญฺจ ธมฺมํ ปฏิจฺจ คนฺโถ ธมฺโม อุปฺปชฺชติ เหตุปจฺจยา. สีลพฺพต\-ปรามาสํ กายคนฺถญฺจ สมฺปยุตฺตเก จ ขนฺเธ ปฏิจฺจ อภิชฺฌา\-กายคนฺโถ. (จกฺกํ.) (1)}

\prose{\csromanfolio{254}คนฺถญฺจ โนคนฺถญฺจ ธมฺมํ ปฏิจฺจ โนคนฺโถ ธมฺโม อุปฺปชฺชติ เหตุปจฺจยา. โนคนฺถํ เอกํ ขนฺธญฺจ คนฺเถ จ ปฏิจฺจ ตโย ขนฺธา จิตฺต\-สมุฏฺฐานญฺจ รูปํ ฯเปฯ เทฺว ขนฺเธ จ ฯเปฯ. คนฺเถ จ สมฺปยุตฺตเก ขนฺเธ จ ปฏิจฺจ จิตฺต\-สมุฏฺฐานํ รูปํ. (2)}

\prose{คนฺถญฺจ โนคนฺถญฺจ ธมฺมํ ปฏิจฺจ คนฺโถ จ โนคนฺโถ จ ธมฺมา อุปฺปชฺชนฺติ เหตุปจฺจยา. โนคนฺถํ เอกํ ขนฺธญฺจ สีลพฺพตปรามาส\-กายคนฺถญฺจ ปฏิจฺจ ตโย ขนฺธา อภิชฺฌา\-กายคนฺโถ จิตฺต\-สมุฏฺฐานญฺจ รูปํ ฯเปฯ เทฺว ขนฺเธ จ ฯเปฯ. (3) (จกฺกํ. สํขิตฺตํ.) อารมฺมณ\-ปจฺจยา ฯเปฯ อวิคตปจฺจยา.}

\csromanmarkh{1. ปจฺจยานุโลม}\csromanmarkhii{2. สงฺขฺยาวาร}\csromanheadernomark{2}{1. \textbf{ปจฺจยานุโลม 2. สงฺขฺยาวาร}}
\proseitem{4}{เหตุยา นว, อารมฺมเณ นว, อธิปติยา นว, (สพฺพตฺถ นว.) วิปาเก เอกํ, อาหาเร นว ฯเปฯ อวิคเต นว.}

\vspace{\tipitakaparskip}
\csromancenter{\textbf{อนุโลมํ.}}
\csromanmarkh{2. ปจฺจย\-ปจฺจนีย}\csromanmarkhii{1. วิภงฺควาร}\csromanheadernomark{2}{2. ปจฺจย\-ปจฺจนีย 1. วิภงฺควาร}
\proseitem{5}{โนคนฺถํ ธมฺมํ ปฏิจฺจ โนคนฺโถ ธมฺโม อุปฺปชฺชติ นเหตุปจฺจยา. อเหตุกํ โนคนฺถํ เอกํ ขนฺธํ ปฏิจฺจ ตโย ขนฺธา จิตฺต\-สมุฏฺฐานญฺจ รูปํ ฯเปฯ เทฺว ขนฺเธ ฯเปฯ. อเหตุก\-ปฏิสนฺธิ\-กฺขเณ. (ยาว อสญฺญสตฺตา.) วิจิกิจฺฉา\-สหคเต อุทฺธจฺจ\-สหคเต ขนฺเธ ปฏิจฺจ วิจิกิจฺฉา\-สหคโต อุทฺธจฺจ\-สหคโต โมโห. (1)}

\prose{นอารมฺมณาทิ}

\proseitem{6}{คนฺถํ ธมฺมํ ปฏิจฺจ โนคนฺโถ ธมฺโม อุปฺปชฺชติ นอารมฺมณ\-ปจฺจยา. คนฺเถ ปฏิจฺจ จิตฺต\-สมุฏฺฐานํ รูปํ. (1)}

\prose{โนคนฺถํ ธมฺมํ ปฏิจฺจ โนคนฺโถ ธมฺโม อุปฺปชฺชติ นอารมฺมณ\-ปจฺจยา. โนคนฺเถ ขนฺเธ ปฏิจฺจ จิตฺต\-สมุฏฺฐานํ รูปํ. ปฏิสนฺธิ\-กฺขเณ ฯเปฯ. (ยาว อสญฺญสตฺตา.) (1)}

\prose{\csromanfolio{255}คนฺถญฺจ โนคนฺถญฺจ ธมฺมํ ปฏิจฺจ โนคนฺโถ ธมฺโม อุปฺปชฺชติ นอารมฺมณ\-ปจฺจยา. คนฺเถ จ สมฺปยุตฺตเก จ ขนฺเธ ปฏิจฺจ จิตฺต\-สมุฏฺฐานํ รูปํ. (สํขิตฺตํ.)}

\prose{นอธิปติ\-ปจฺจยา. นว. นอนนฺตร\-ปจฺจยา. ตีณิ. นสมนนฺตร\-ปจฺจยา. ตีณิ. นอญฺญมญฺญ\-ปจฺจยา. ตีณิ. นอุปนิสฺสยปจฺจยา. ตีณิ.}

\csromanheader{2}{\textbf{นปุเรชาตาทิ}}
\proseitem{7}{คนฺถํ ธมฺมํ ปฏิจฺจ คนฺโถ ธมฺโม อุปฺปชฺชติ นปุเรชาต\-ปจฺจยา. อรูเป อิทํสจฺจา\-ภินิเวสํ กายคนฺถํ ปฏิจฺจ อภิชฺฌา\-กายคนฺโถ. อภิชฺฌา\-กายคนฺถํ ปฏิจฺจ อิทํสจฺจา\-ภินิเวโส กายคนฺโถ. (อรูเป สีลพฺพต\-ปรามาโส นตฺถิ, เอวํ นว ปญฺหา กาตพฺพา.) นปจฺฉาชาเต นว, นอาเสวเน นว, นกมฺเม ตีณิ, นวิปาเก นว, นอาหาเร เอกํ, นอินฺทฺริเย เอกํ, นฌาเน เอกํ, นมคฺเค เอกํ, นสมฺปยุตฺเต ตีณิ, นวิปฺปยุตฺเต นว, โนนตฺถิยา ตีณิ, โนวิคเต ตีณิ.}

\csromanmarkh{2. ปจฺจย\-ปจฺจนีย}\csromanmarkhii{2. สงฺขฺยาวาร}\csromanheadernomark{2}{2. ปจฺจย\-ปจฺจนีย 2. สงฺขฺยาวาร}
\proseitem{8}{นเหตุยา เอกํ, นอารมฺมเณ ตีณิ, นอธิปติยา นว, นอนนฺตเร ตีณิ, นสมนนฺตเร ตีณิ, นอญฺญมญฺเญ ตีณิ, นอุปนิสฺสเย ตีณิ, นปุเรชาเต นว, นปจฺฉาชาเต นว, นอาเสวเน นว, นกมฺเม ตีณิ, นวิปาเก นว, นอาหาเร เอกํ, นอินฺทฺริเย เอกํ, นฌาเน เอกํ, นมคฺเค เอกํ, นสมฺปยุตฺเต ตีณิ, นวิปฺปยุตฺเต นว, โนนตฺถิยา ตีณิ, โนวิคเต ตีณิ.}

\proseitem{9}{เหตุปจฺจยา นอารมฺมเณ ตีณิ, นอธิปติยา นว. (สํขิตฺตํ. เอวํ คเณตพฺพํ.)}

\vspace{\tipitakaparskip}
\csromancenter{นเหตุปจฺจยา อารมฺมเณ เอกํ ฯเปฯ อนนฺตเร เอกํ ฯเปฯ อวิคเต เอกํ.}
\csromancenter{(สหชาตวาโร ปฏิจฺจ\-วาร\-สทิโส.)}
\csromanmarkh{26. คนฺถทุก}\csromanmarkhii{3. ปจฺจยวาร}\csromanheadernomark{2}{26. \textbf{คนฺถทุก 3. ปจฺจยวาร}}
\csromanmarkh{1. ปจฺจยานุโลม}\csromanmarkhii{1. วิภงฺควาร}\csromanheadernomark{2}{1. \textbf{ปจฺจยานุโลม 1. วิภงฺควาร}}
\csromanheader{2}{\textbf{เหตุ}}
\proseitem{11}{\csromanfolio{256}คนฺถํ ธมฺมํ ปจฺจยา คนฺโถ ธมฺโม อุปฺปชฺชติ เหตุปจฺจยา. ตีณิ. (ปฏิจฺจ\-วาร\-สทิโส.)}

\prose{โนคนฺถํ ธมฺมํ ปจฺจยา โนคนฺโถ ธมฺโม อุปฺปชฺชติ เหตุปจฺจยา. โนคนฺถํ เอกํ ขนฺธํ ปจฺจยา ตโย ขนฺธา จิตฺต\-สมุฏฺฐานญฺจ รูปํ ฯเปฯ เทฺว ขนฺเธ ฯเปฯ. ปฏิสนฺธิ\-กฺขเณ ฯเปฯ. ขนฺเธ ปจฺจยา วตฺถุ, วตฺถุํ ปจฺจยา ขนฺธา. เอกํ มหาภูตํ ฯเปฯ. วตฺถุํ ปจฺจยา โนคนฺถา ขนฺธา. (1)}

\prose{โนคนฺถํ ธมฺมํ ปจฺจยา คนฺโถ ธมฺโม อุปฺปชฺชติ เหตุปจฺจยา. โนคนฺเถ ขนฺเธ ปจฺจยา คนฺถา. วตฺถุํ ปจฺจยา คนฺถา. (2)}

\prose{โนคนฺถํ ธมฺมํ ปจฺจยา คนฺโถ จ โนคนฺโถ จ ธมฺมา อุปฺปชฺชนฺติ เหตุปจฺจยา. โนคนฺถํ เอกํ ขนฺธํ ปจฺจยา ตโย ขนฺธา คนฺถา จ จิตฺต\-สมุฏฺฐานญฺจ รูปํ ฯเปฯ เทฺว ขนฺเธ ฯเปฯ. วตฺถุํ ปจฺจยา คนฺถา สมฺปยุตฺตกา จ ขนฺธา. (3)}

\proseitem{12}{คนฺถญฺจ โนคนฺถญฺจ ธมฺมํ ปจฺจยา คนฺโถ ธมฺโม อุปฺปชฺชติ เหตุปจฺจยา. สีลพฺพต\-ปรามาสํ กายคนฺถญฺจ สมฺปยุตฺตเก จ ขนฺเธ ปจฺจยา อภิชฺฌา\-กายคนฺโถ. (จกฺกํ.) สีลพฺพต\-ปรามาสํ กายคนฺถญฺจ วตฺถุญฺจ ปจฺจยา อภิชฺฌา\-กายคนฺโถ. (จกฺกํ.) (1)}

\prose{คนฺถญฺจ โนคนฺถญฺจ ธมฺมํ ปจฺจยา โนคนฺโถ ธมฺโม อุปฺปชฺชติ เหตุปจฺจยา. โนคนฺถํ เอกํ ขนฺธญฺจ คนฺเถ จ ปจฺจยา ตโย ขนฺธา จิตฺต\-สมุฏฺฐานญฺจ รูปํ ฯเปฯ เทฺว ขนฺเธ ฯเปฯ. คนฺเถ จ วตฺถุญฺจ ปจฺจยา โนคนฺถา ขนฺธา. (2)}

\prose{คนฺถญฺจ โนคนฺถญฺจ ธมฺมํ ปจฺจยา คนฺโถ จ โนคนฺโถ จ ธมฺมา อุปฺปชฺชนฺติ เหตุปจฺจยา. โนคนฺถํ เอกํ ขนฺธญฺจ สีลพฺพตรามาสํ กายคนฺถญฺจ ปจฺจยา ตโย ขนฺธา อภิชฺฌา\-กายคนฺโถ จ จิตฺต\-สมุฏฺฐานญฺจ รูปํ ฯเปฯ เทฺว ขนฺเธ ฯเปฯ. (จกฺกํ.) สีลพฺพต\-ปรามาสํ กายคนฺถญฺจ วตฺถุญฺจ ปจฺจยา อภิชฺฌา\-กายคนฺโถ จ สมฺปยุตฺตกา จ ขนฺธา. (จกฺกํ.) (3)}

\csromanmarkh{26. คนฺถทุก}\csromanmarkhii{1. ปจฺจยานุโลม 2. สงฺขฺยาวาร}\csromanheadernomark{2}{26. คนฺถทุก \textbf{1. ปจฺจยานุโลม 2. สงฺขฺยาวาร}}
\vspace{\tipitakaparskip}
\csromancenter{เหตุยา นว, อารมฺมเณ นว, อธิปติยา นว ฯเปฯ อวิคเต นว.}
\csromancenter{อนุโลมํ.}
\csromanmarkh{2. ปจฺจย\-ปจฺจนีย}\csromanmarkhii{1. วิภงฺควาร}\csromanheadernomark{2}{2. ปจฺจย\-ปจฺจนีย 1. วิภงฺควาร}
\proseitem{14}{\csromanfolio{257}โนคนฺถํ ธมฺมํ ปจฺจยา โนคนฺโถ ธมฺโม อุปฺปชฺชติ นเหตุปจฺจยา. อเหตุกํ โนคนฺถํ เอกํ ขนฺธํ ปจฺจยา ตโย ขนฺธา จิตฺต\-สมุฏฺฐานญฺจ รูปํ ฯเปฯ เทฺว ขนฺเธ ฯเปฯ. ปฏิสนฺธิ\-กฺขเณ ฯเปฯ. (ยาว อสญฺญสตฺตา.) จกฺขายตนํ ปจฺจยา จกฺขุวิญฺญาณํ ฯเปฯ. กายายตนํ ปจฺจยา กายวิญฺญาณํ. วตฺถุํ ปจฺจยา อเหตุกา โนคนฺถา ขนฺธา. วิจิกิจฺฉา\-สหคเต อุทฺธจฺจ\-สหคเต ขนฺเธ จ วตฺถุญฺจ ปจฺจยา วิจิกิจฺฉา\-สหคโต อุทฺธจฺจ\-สหคโต โมโห (สํขิตฺตํ.)}

\csromanmarkh{2. ปจฺจย\-ปจฺจนีย}\csromanmarkhii{2. สงฺขฺยาวาร}\csromanheadernomark{2}{2. ปจฺจย\-ปจฺจนีย 2. สงฺขฺยาวาร}
\vspace{\tipitakaparskip}
\csromancenter{\textbf{นเหตุ}ยา เอกํ, นอารมฺมเณ ตีณิ, นอธิปติยา นว. (เอวํ คเณตพฺพํ.)}
\csromanheader{2}{3. \textbf{ปจฺจยานุโลมปจฺจนิย}}
\vspace{\tipitakaparskip}
\csromancenter{เหตุปจฺจยา นอารมฺมเณ ตีณิ, นอธิปติยา นว ฯเปฯ โนวิคเต ตีณิ.}
\proseitem{17}{\textbf{นเหตุ}ปจฺจยา อารมฺมเณ เอกํ ฯเปฯ อวิคเต เอกํ.}

\prose{(นิสฺสยวาโร ปจฺจยวารสทิโสว. สํสฏฺฐ\-วาโรปิ สมฺปยุตฺตวาโรปิ นว ปญฺหา กาตพฺพา. รูปํ นตฺถิ.)}

\csromanmarkh{26. คนฺถทุก}\csromanmarkhii{7. ปญฺหาวาร}\csromanheadernomark{2}{26. \textbf{คนฺถทุก 7. ปญฺหาวาร}}
\csromanmarkh{1. ปจฺจยานุโลม}\csromanmarkhii{1. วิภงฺควาร}\csromanheadernomark{2}{1. \textbf{ปจฺจยานุโลม 1. วิภงฺควาร}}
\csromanheader{2}{\textbf{เหตุ}}
\proseitem{18}{\csromanfolio{258}คนฺโถ ธมฺโม คนฺถสฺส ธมฺมสฺส เหตุปจฺจเยน ปจฺจโย. คนฺถา เหตู สมฺปยุตฺตกานํ คนฺถานํ เหตุปจฺจเยน ปจฺจโย. (1)}

\prose{คนฺโถ ธมฺโม โนคนฺถสฺส ธมฺมสฺส เหตุปจฺจเยน ปจฺจโย. คนฺถา เหตู สมฺปยุตฺตกานํ ขนฺธานํ จิตฺตสมุฏฺฐานานญฺจ รูปานํ เหตุปจฺจเยน ปจฺจโย. (2)}

\prose{คนฺโถ ธมฺโม คนฺถสฺส จ โนคนฺถสฺส จ ธมฺมสฺส เหตุปจฺจเยน ปจฺจโย. คนฺถา เหตู สมฺปยุตฺตกานํ ขนฺธานํ คนฺถานญฺจ จิตฺต\-สมุฏฺฐานญฺจ รูปานํ เหตุปจฺจเยน ปจฺจโย. (3)}

\proseitem{19}{โนคนฺโถ ธมฺโม โนคนฺถสฺส ธมฺมสฺส เหตุปจฺจเยน ปจฺจโย. โนคนฺถา เหตู สมฺปยุตฺตกานํ ขนฺธานํ จิตฺตสมุฏฺฐานานญฺจ รูปานํ เหตุปจฺจเยน ปจฺจโย. ปฏิสนฺธิ\-กฺขเณ ฯเปฯ. (1)}

\prose{โนคนฺโถ ธมฺโม คนฺถสฺส ธมฺมสฺส เหตุปจฺจเยน ปจฺจโย. โนคนฺถา เหตู สมฺปยุตฺตกานํ คนฺถานํ เหตุปจฺจเยน ปจฺจโย. (2)}

\prose{โนคนฺโถ ธมฺโม คนฺถสฺส จ โนคนฺถสฺส จ ธมฺมสฺส เหตุปจฺจเยน ปจฺจโย. โนคนฺถา เหตู สมฺปยุตฺตกานํ ขนฺธานํ คนฺถานญฺจ จิตฺตสมุฏฺฐานานญฺจ รูปานํ เหตุปจฺจเยน ปจฺจโย. (3)}

\proseitem{20}{คนฺโถ จ โนคนฺโถ จ ธมฺมา คนฺถสฺส ธมฺมสฺส เหตุปจฺจเยน ปจฺจโย. คนฺถา จ โนคนฺถา จ เหตู สมฺปยุตฺตกานํ คนฺถานํ เหตุปจฺจเยน ปจฺจโย. (1)}

\prose{คนฺโถ จ โนคนฺโถ จ ธมฺมา โนคนฺถสฺส ธมฺมสฺส เหตุปจฺจเยน ปจฺจโย. คนฺถา จ โนคนฺถา จ เหตู สมฺปยุตฺตกานํ ขนฺธานํ จิตฺตสมุฏฺฐานานญฺจ รูปานํ เหตุปจฺจเยน ปจฺจโย. (2)}

\prose{\csromanfolio{259}คนฺโถ จ โนคนฺโถ จ ธมฺมา คนฺถสฺส จ โนคนฺถสฺส จ ธมฺมสฺส เหตุปจฺจเยน ปจฺจโย. คนฺถา จ โนคนฺถา จ เหตู สมฺปยุตฺตกานํ ขนฺธานํ คนฺถานญฺจ จิตฺต\-สมุฏฺฐานญฺจ รูปานํ เหตุปจฺจเยน ปจฺจโย. (3)}

\proseitem{21}{คนฺโถ ธมฺโม คนฺถสฺส ธมฺมสฺส อารมฺมณ\-ปจฺจเยน ปจฺจโย. คนฺเถ อารพฺภ คนฺถา อุปฺปชฺชนฺติ. (มูลํ กาตพฺพํ.) คนฺเถ อารพฺภ โนคนฺถา ขนฺธา อุปฺปชฺชนฺติ. (มูลํ กาตพฺพํ.) คนฺเถ อารพฺภ คนฺถา จ สมฺปยุตฺตกา จ ขนฺธา อุปฺปชฺชนฺติ. (3)}

\proseitem{22}{โนคนฺโถ ธมฺโม โนคนฺถสฺส ธมฺมสฺส อารมฺมณ\-ปจฺจเยน ปจฺจโย. ทานํ ทตฺวา, สีลํ ฯเปฯ อุโปสถกมฺมํ กตฺวา ตํ ปจฺจเวกฺขติ. ปุพฺเพ สุจิณฺณานิ ฯเปฯ. ฌานา ฯเปฯ. อริยา มคฺคา วุฏฺฐหิตฺวา มคฺคํ ปจฺจเวกฺขนฺติ, ผลํ, ปจฺจเวกฺขนฺติ, นิพฺพานํ ปจฺจเวกฺขนฺติ. นิพฺพานํ โคตฺรภุสฺส, โวทานสฺส, มคฺคสฺส, ผลสฺส, อาวชฺชนาย อารมฺมณ\-ปจฺจเยน ปจฺจโย. อริยา โนคนฺเถ ปหีเน กิเลเส ฯเปฯ. วิกฺขมฺภิเต กิเลเส ปจฺจเวกฺขนฺติ. ปุพฺเพ ฯเปฯ. จกฺขุํ ฯเปฯ วตฺถุํ โนคนฺเถ ขนฺเธ อนิจฺจโต ฯเปฯ โทมนสฺสํ อุปฺปชฺชติ. ทิพฺเพน จกฺขุนา รูปํ ปสฺสติ. ทิพฺพาย โสตธาตุยา สทฺทํ สุณาติ. เจโตปริย\-ญาเณน โนคนฺถ\-จิตฺต\-สมงฺคิสฺส จิตฺตํ ชานาติ. อากาสา\-นญฺจา\-ยตนํ วิญฺญาณญฺจา\-ยตนสฺส ฯเปฯ. อากิญฺจญฺญา\-ยตนํ เนว\-สญฺญา\-นา\-สญฺญา\-ยตนสฺส ฯเปฯ. รูปายตนํ จกฺขุ\-วิญฺญาณสฺส ฯเปฯ. โผฏฺฐพฺพา\-ยตนํ กายวิญฺญาณสฺส ฯเปฯ. โนคนฺถา ขนฺธา อิทฺธิวิธ\-ญาณสฺส, เจโตปริย\-ญาณสฺส, ปุพฺเพ\-นิวาสา\-นุสฺสติ\-ญาณสฺส, ยถากมฺมูปคญาณสฺส, อนาคตํสญาณสฺส, อาวชฺชนาย อารมฺมณ\-ปจฺจเยน ปจฺจโย. (1)}

\prose{โนคนฺโถ ธมฺโม คนฺถสฺส ธมฺมสฺส อารมฺมณ\-ปจฺจเยน ปจฺจโย. ทานํ ทตฺวา, สีลํ ฯเปฯ อุโปสถกมฺมํ กตฺวา ตํ อสฺสาเทติ อภินนฺทติ, ตํ อารพฺภ ราโค อุปฺปชฺชติ, ทิฏฺฐิ อุปฺปชฺชติ, โทมนสฺสํ อุปฺปชฺชติ. ปุพฺเพ สุจิณฺณานิ ฯเปฯ. ฌานา ฯเปฯ. จกฺขุํ ฯเปฯ วตฺถุํ โนคนฺเถ ขนฺเธ อสฺสาเทติ อภินนฺทติ, ตํ อารพฺภ ราโค อุปฺปชฺชติ ฯเปฯ โทมนสฺสํ อุปฺปชฺชติ. (2)}

\prose{โนคนฺโถ ธมฺโม คนฺถสฺส จ โนคนฺถสฺส จ ธมฺมสฺส อารมฺมณ\-ปจฺจเยน ปจฺจโย. ทานํ ทตฺวา, สีลํ ฯเปฯ อุโปสถกมฺมํ กตฺวา ตํ \csromanfolio{260}อสฺสาเทติ อภินนฺทติ, ตํ อารพฺภ คนฺถา จ สมฺปยุตฺตกา จ ขนฺธา อุปฺปชฺชนฺติ. ปุพฺเพ สุจิณฺณานิ ฯเปฯ. ฌานา ฯเปฯ. จกฺขุํ ฯเปฯ วตฺถุํ โนคนฺเถ ขนฺเธ อสฺสาเทติ อภินนฺทติ, ตํ อารพฺภ คนฺถา จ สมฺปยุตฺตกา จ ขนฺธา อุปฺปชฺชนฺติ. (3)}

\prose{คนฺโถ จ โนคนฺโถ จ ธมฺมา คนฺถสฺส ธมฺมสฺส อารมฺมณ\-ปจฺจเยน ปจฺจโย. ตีณิ. (อารพฺภ กาตพฺพา.)}

\proseitem{23}{คนฺโถ ธมฺโม คนฺถสฺส ธมฺมสฺส อธิปติ\-ปจฺจเยน ปจฺจโย. ตีณิ. (อารมฺมณสทิสา. ครุการมฺมณา กาตพฺพา.)}

\prose{โนคนฺโถ ธมฺโม โนคนฺถสฺส ธมฺมสฺส อธิปติ\-ปจฺจเยน ปจฺจโย. อารมฺมณา\-ธิปติ, สหชาตา\-ธิปติ.  อารมฺมณา\-ธิปติ\csromandash{}ทานํ ฯเปฯ สีลํ ฯเปฯ อุโปสถกมฺมํ ฯเปฯ. ปุพฺเพ ฯเปฯ. ฌานา ฯเปฯ. อริยา มคฺคา ฯเปฯ. ผลํ ฯเปฯ. นิพฺพานํ ฯเปฯ. นิพฺพานํ โคตฺรภุสฺส, โวทานสฺส, มคฺคสฺส, ผลสฺส, อธิปติ\-ปจฺจเยน ปจฺจโย. จกฺขุํ ฯเปฯ วตฺถุํ โนคนฺเถ ขนฺเธ ครุํ กตฺวา อสฺสาเทติ อภินนฺทติ, ตํ ครุํ กตฺวา ราโค อุปฺปชฺชติ, ทิฏฺฐิ อุปฺปชฺชติ. สหชาตาธิปตฺอิ\csromandash{}โนคนฺถา\-ธิปติ สมฺปยุตฺตกานํ ขนฺธานํ จิตฺตสมุฏฺฐานานญฺจ รูปานํ อธิปติ\-ปจฺจเยน ปจฺจโย. (1)}

\prose{โนคนฺโถ ธมฺโม คนฺถสฺส ธมฺมสฺส อธิปติ\-ปจฺจเยน ปจฺจโย. อารมฺมณา\-ธิปติ, สหชาตา\-ธิปติ.  อารมฺมณา\-ธิปติ\csromandash{}ทานํ ฯเปฯ สีลํ ฯเปฯ อุโปสถกมฺมํ ฯเปฯ. ปุพฺเพ สุจิณฺณานิ ฯเปฯ. ฌานา ฯเปฯ. จกฺขุํ ฯเปฯ วตฺถุํ โนคนฺเถ ขนฺเธ ครุํ กตฺวา อสฺสาเทติ อภินนฺทติ, ตํ ครุํ กตฺวา ราโค อุปฺปชฺชติ, ทิฏฺฐิ อุปฺปชฺชติ. สหชาตา\-ธิปติ\csromandash{} โนคนฺถา\-ธิปติ คนฺถานํ อธิปติ\-ปจฺจเยน ปจฺจโย. (2)}

\prose{โนคนฺโถ ธมฺโม คนฺถสฺส จ โนคนฺถสฺส จ ธมฺมสฺส อธิปติ\-ปจฺจเยน ปจฺจโย. อารมฺมณา\-ธิปติ, สหชาตา\-ธิปติ.  อารมฺมณา\-ธิปติ\csromandash{}ทานํ ฯเปฯ โนคนฺเถ ขนฺเธ ครุํ กตฺวา อสฺสาเทติ อภินนฺทติ, ตํ ครุํ กตฺวา คนฺถา จ สมฺปยุตฺตกา จ ขนฺธา อุปฺปชฺชนฺติ. สหชาตา\-ธิปติ\csromandash{}โนคนฺถา\-ธิปติ สมฺปยุตฺตกานํ ขนฺธานํ คนฺถานญฺจ จิตฺตสมุฏฺฐานานญฺจ รูปานํ อธิปติ\-ปจฺจเยน ปจฺจโย. (3)}

\proseitem{24}{\csromanfolio{261}คนฺโถ จ โนคนฺโถ จ ธมฺมา คนฺถสฺส ธมฺมสฺส อธิปติ\-ปจฺจเยน ปจฺจโย. อารมฺมณา\-ธิปติ. ตีณิ.}

\proseitem{25}{คนฺโถ ธมฺโม คนฺถสฺส ธมฺมสฺส อนนฺตร\-ปจฺจเยน ปจฺจโย. ปุริมา ปุริมา คนฺถา ปจฺฉิมานํ ปจฺฉิมานํ คนฺถานํ อนนฺตร\-ปจฺจเยน ปจฺจโย. (1)}

\prose{คนฺโถ ธมฺโม โนคนฺถสฺส ธมฺมสฺส อนนฺตร\-ปจฺจเยน ปจฺจโย. ปุริมา ปุริมา คนฺถา ปจฺฉิมานํ ปจฺฉิมานํ โนคนฺถานํ ขนฺธานํ อนนฺตร\-ปจฺจเยน ปจฺจโย. คนฺถา วุฏฺฐานสฺส อนนฺตร\-ปจฺจเยน ปจฺจโย. (2)}

\prose{คนฺโถ ธมฺโม คนฺถสฺส จ โนคนฺถสฺส จ ธมฺมสฺส อนนฺตร\-ปจฺจเยน ปจฺจโย. ปุริมา ปุริมา คนฺถา ปจฺฉิมานํ ปจฺฉิมานํ คนฺถา นํ สมฺปยุตฺตกานญฺจ ขนฺธานํ อนนฺตร\-ปจฺจเยน ปจฺจโย. (3)}

\proseitem{26}{โนคนฺโถ ธมฺโม โนคนฺถสฺส ธมฺมสฺส อนนฺตร\-ปจฺจเยน ปจฺจโย. ตีณิ. (เทฺว อาวชฺชนา กาตพฺพา. ปฐโม นตฺถิ.)}

\prose{คนฺโถ จ โนคนฺโถ จ ธมฺมา คนฺถสฺส ธมฺมสฺส อนนฺตร\-ปจฺจเยน ปจฺจโย. ตีณิ. (เอกมฺปิ วุฏฺฐานํ กาตพฺพํ. มชฺเฌ.)}

\csromanheader{2}{\textbf{สมนนฺตราทิ}}
\proseitem{27}{คนฺโถ ธมฺโม คนฺถสฺส ธมฺมสฺส สมนนฺตร\-ปจฺจเยน ปจฺจโย. นว. สหชาต\-ปจฺจเยน ปจฺจโย. นว. อญฺญมญฺญ\-ปจฺจเยน ปจฺจโย. นว. นิสฺสย\-ปจฺจเยน ปจฺจโย. นว.}

\csromanheader{2}{\textbf{อุปนิสฺสย}}
\proseitem{28}{คนฺโถ ธมฺโม คนฺถสฺส ธมฺมสฺส อุปนิสฺสย\-ปจฺจเยน ปจฺจโย. อารมฺมณู\-ปนิสฺสโย, อนนฺตรู\-ปนิสฺสโย, ปกตู\-ปนิสฺสโย ฯเปฯ. ปกตู\-ปนิสฺสโย\csromandash{}คนฺถา คนฺถานํ อุปนิสฺสย\-ปจฺจเยน ปจฺจโย. ตีณิ.}

\proseitem{29}{โนคนฺโถ ธมฺโม โนคนฺถสฺส ธมฺมสฺส อุปนิสฺสย\-ปจฺจเยน ปจฺจโย. อารมฺมณู\-ปนิสฺสโย, อนนฺตรู\-ปนิสฺสโย, ปกตู\-ปนิสฺสโย \csromanfolio{262}ฯเปฯ. ปกตู\-ปนิสฺสโย\csromandash{}สทฺธํ อุปนิสฺสาย ทานํ เทติ ฯเปฯ มานํ ชปฺเปติ. ทิฏฺฐิํ คณฺหาติ. สีลํ ฯเปฯ. เสนาสนํ อุปนิสฺสาย ทานํ เทติ ฯเปฯ สํฆํ ภินฺทติ. สทฺธา ฯเปฯ. เสนาสนํ สทฺธาย ฯเปฯ ผลสมาปตฺติยา อุปนิสฺสย\-ปจฺจเยน ปจฺจโย. (1)}

\prose{โนคนฺโถ ธมฺโม คนฺถสฺส ธมฺมสฺส อุปนิสฺสย\-ปจฺจเยน ปจฺจโย. อารมฺมณู\-ปนิสฺสโย, อนนฺตรู\-ปนิสฺสโย, ปกตู\-ปนิสฺสโย ฯเปฯ. ปกตู\-ปนิสฺสโย\csromandash{}สทฺธํ อุปนิสฺสาย ทิฏฺฐิํ คณฺหาติ. สีลํ ฯเปฯ. เสนาสนํ อุปนิสฺสาย ปาณํ หนติ ฯเปฯ สํฆํ ภินฺทติ. สทฺธา ฯเปฯ. เสนาสนํ ราคสฺส ฯเปฯ ปตฺถนาย อุปนิสฺสย\-ปจฺจเยน ปจฺจโย. (2)}

\prose{โนคนฺโถ ธมฺโม คนฺถสฺส จ โนคนฺถสฺส จ ธมฺมสฺส อุปนิสฺสย\-ปจฺจเยน ปจฺจโย. อารมฺมณู\-ปนิสฺสโย, อนนฺตรู\-ปนิสฺสโย, ปกตู\-ปนิสฺสโย ฯเปฯ. ปกตู\-ปนิสฺสโย\csromandash{}สทฺธํ อุปนิสฺสาย มานํ ชปฺเปติ. ทิฏฺฐิํ คณฺหาติ. สีลํ ฯเปฯ. เสนาสนํ อุปนิสฺสาย ปาณํ หนติ ฯเปฯ สํฆํ ภินฺทติ. สทฺธา ฯเปฯ. เสนาสนํ คนฺถานํ สมฺปยุตฺตกานญฺจ ขนฺธานํ อุปนิสฺสย\-ปจฺจเยน ปจฺจโย. (3)}

\prose{คนฺโถ จ โนคนฺโถ จ ธมฺมา คนฺถสฺส ธมฺมสฺส อุปนิสฺสย\-ปจฺจเยน ปจฺจโย. อารมฺมณู\-ปนิสฺสโย, อนนฺตรู\-ปนิสฺสโย, ปกตู\-ปนิสฺสโย ฯเปฯ. ปกตู\-ปนิสฺสโย. ตีณิ. (อารมฺมณนเยน กาตพฺพา.)}

\csromanheader{2}{\textbf{ปุเรชาตาทิ}}
\proseitem{30}{โนคนฺโถ ธมฺโม โนคนฺถสฺส ธมฺมสฺส ปุเรชาต\-ปจฺจเยน ปจฺจโย. อารมฺมณ\-ปุเรชาตํ, วตฺถุ\-ปุเรชาตํ.  อารมฺมณ\-ปุเรชาตํ\csromandash{} จกฺขุํ ฯเปฯ วตฺถุํ อนิจฺจโต ฯเปฯ โทมนสฺสํ อุปฺปชฺชติ. ทิพฺเพน จกฺขุนา รูปํ ปสฺสติ. ทิพฺพาย โสตธาตุยา สทฺทํ สุณาติ. รูปายตนํ จกฺขุ\-วิญฺญาณสฺส ฯเปฯ. โผฏฺฐพฺพา\-ยตนํ กายวิญฺญาณสฺส ฯเปฯ. วตฺถุ\-ปุเรชาตํ\csromandash{}จกฺขายตนํ จกฺขุ\-วิญฺญาณสฺส ฯเปฯ กายายตนํ กายวิญฺญาณสฺส ฯเปฯ. วตฺถุ โนคนฺถานํ ขนฺธานํ ปุเรชาต\-ปจฺจเยน ปจฺจโย. (1)}

\prose{โนคนฺโถ ธมฺโม คนฺถสฺส ธมฺมสฺส ปุเรชาต\-ปจฺจเยน ปจฺจโย. อารมฺมณ\-ปุเรชาตํ, วตฺถุ\-ปุเรชาตํ.  อารมฺมณ\-ปุเรชาตํ\csromandash{}จกฺขุํ ฯเปฯ วตฺถุํ อสฺสาเทติ อภินนฺทติ, ตํ อารพฺภ ราโค อุปฺปชฺชติ, ทิฏฺฐิ อุปฺปชฺชติ, \csromanfolio{263}โทมนสฺสํ อุปฺปชฺชติ. วตฺถุ\-ปุเรชาตํ\csromandash{}วตฺถุ คนฺถานํ ปุเรชาต\-ปจฺจเยน ปจฺจโย. (2)}

\prose{โนคนฺโถ ธมฺโม คนฺถสฺส จ โนคนฺถสฺส จ ธมฺมสฺส ปุเรชาต\-ปจฺจเยน ปจฺจโย. อารมฺมณ\-ปุเรชาตํ, วตฺถุ\-ปุเรชาตํ.  อารมฺมณ\-ปุเรชาตํ\csromandash{}จกฺขุํ ฯเปฯ วตฺถุํ อสฺสาเทติ อภินนฺทติ, ตํ อารพฺภ คนฺถา จ สมฺปยุตฺตกา จ ขนฺธา อุปฺปชฺชนฺติ. วตฺถุ\-ปุเรชาตํ\csromandash{} วตฺถุ คนฺถานํ สมฺปยุตฺตกานญฺจ ขนฺธานํ ปุเรชาต\-ปจฺจเยน ปจฺจโย. (3) ปจฺฉาชาต\-ปจฺจเยน ปจฺจโย. ตีณิ. อาเสวน\-ปจฺจเยน ปจฺจโย. นว.}

\proseitem{31}{โนคนฺโถ ธมฺโม โนคนฺถสฺส ธมฺมสฺส กมฺมปจฺจเยน ปจฺจโย. สหชาตา, นานากฺขณิกา.  สหชาตา โนคนฺถา เจตนา สมฺปยุตฺตกานํ ขนฺธานํ จิตฺตสมุฏฺฐานานญฺจ รูปานํ กมฺมปจฺจเยน ปจฺจโย. ปฏิสนฺธิ\-กฺขเณ ฯเปฯ. นานากฺขณิกา โนคนฺถา เจตนา วิปากานํ ขนฺธานํ กฏตฺตา จ รูปานํ กมฺมปจฺจเยน ปจฺจโย. (1)}

\prose{โนคนฺโถ ธมฺโม คนฺถสฺส ธมฺมสฺส กมฺมปจฺจเยน ปจฺจโย. โนคนฺถา เจตนา สมฺปยุตฺตกานํ คนฺถานํ กมฺมปจฺจเยน ปจฺจโย.}

\prose{โนคนฺโถ ธมฺโม คนฺถสฺส จ โนคนฺถสฺส จ ธมฺมสฺส กมฺมปจฺจเยน ปจฺจโย. โนคนฺถา เจตนา สมฺปยุตฺตกานํ ขนฺธานํ คนฺถานญฺจ จิตฺตสมุฏฺฐานานญฺจ รูปานํ กมฺมปจฺจเยน ปจฺจโย. (3)}

\csromanheader{2}{\textbf{วิปากาทิ}}
\proseitem{32}{โนคนฺโถ ธมฺโม โนคนฺถสฺส ธมฺมสฺส วิปาก\-ปจฺจเยน ปจฺจโย. เอกํ. อาหารปจฺจเยน ปจฺจโย. ตีณิ. อินฺทฺริย\-ปจฺจเยน ปจฺจโย. ตีณิ. ฌานปจฺจเยน ปจฺจโย. ตีณิ. มคฺคปจฺจเยน ปจฺจโย. นว. สมฺปยุตฺต\-ปจฺจเยน ปจฺจโย. นว.}

\proseitem{33}{คนฺโถ ธมฺโม โนคนฺถสฺส ธมฺมสฺส วิปฺปยุตฺต\-ปจฺจเยน ปจฺจโย. สหชาตํ, ปจฺฉาชาตํ. (สํขิตฺตํ.) (1)}

\prose{\csromanfolio{264}โนคนฺโถ ธมฺโม โนคนฺถสฺส ธมฺมสฺส วิปฺปยุตฺต\-ปจฺจเยน ปจฺจโย. สหชาตํ, ปุเรชาตํ, ปจฺฉาชาตํ. (สํขิตฺตํ.) (1)}

\prose{โนคนฺโถ ธมฺโม คนฺถสฺส ธมฺมสฺส วิปฺปยุตฺต\-ปจฺจเยน ปจฺจโย. ปุเรชาตํ วตฺถุ คนฺถานํ วิปฺปยุตฺต\-ปจฺจเยน ปจฺจโย. (2)}

\prose{โนคนฺโถ ธมฺโม คนฺถสฺส จ โนคนฺถสฺส จ ธมฺมสฺส วิปฺปยุตฺต\-ปจฺจเยน ปจฺจโย. ปุเรชาตํ วตฺถุ คนฺถานํ สมฺปยุตฺตกานญฺจ ขนฺธานํ วิปฺปยุตฺต\-ปจฺจเยน ปจฺจโย. (3)}

\prose{คนฺโธ จ โนคนฺโถ จ ธมฺมา โนคนฺถสฺส ธมฺมสฺส วิปฺปยุตฺต\-ปจฺจเยน ปจฺจโย. สหชาตํ, ปจฺฉาชาตํ. (สํขิตฺตํ.) (1)}

\proseitem{34}{คนฺโถ ธมฺโม คนฺถสฺส ธมฺมสฺส อตฺถิปจฺจเยน ปจฺจโย. เอกํ. (ปฏิจฺจสทิสํ.)}

\prose{คนฺโถ ธมฺโม โนคนฺถสฺส ธมฺมสฺส อตฺถิปจฺจเยน ปจฺจโย. สหชาตํ, ปจฺฉาชาตํ.  สหชาตา คนฺถา สมฺปยุตฺตกานํ ขนฺธานํ จิตฺตสมุฏฺฐานานญฺจ รูปานํ อตฺถิปจฺจเยน ปจฺจโย. ปจฺฉาชาตา คนฺถา ปุเรชาตสฺส อิมสฺส กายสฺส อตฺถิปจฺจเยน ปจฺจโย. (2)}

\prose{คนฺโถ ธมฺโม คนฺถสฺส จ โนคนฺถสฺส จ ธมฺมสฺส อตฺถิปจฺจเยน ปจฺจโย. เอกํ. (ปฏิจฺจสทิสํ.) (3)}

\proseitem{35}{โนคนฺโถ ธมฺโม โนคนฺถสฺส ธมฺมสฺส อตฺถิปจฺจเยน ปจฺจโย. สหชาตํ, ปุเรชาตํ, ปจฺฉาชาตํ, อาหารํ, อินฺทฺริยํ. (สํขิตฺตํ.) (1)}

\prose{โนคนฺโถ ธมฺโม คนฺถสฺส ธมฺมสฺส อตฺถิปจฺจเยน ปจฺจโย. สหชาตํ, ปุเรชาตํ.  สหชาตา โนคนฺถา ขนฺธา คนฺถานํ อตฺถิปจฺจเยน ปจฺจโย. ปุเรชาตํ จกฺขุํ ฯเปฯ วตฺถุํ อสฺสาเทติ อภินนฺทติ, ตํ อารพฺภ ราโค อุปฺปชฺชติ, ทิฏฺฐิ อุปฺปชฺชติ, โทมนสฺสํ อุปฺปชฺชติ. วตฺถุ คนฺถานํ อตฺถิปจฺจเยน ปจฺจโย. (2)}

\prose{\csromanfolio{265}โนคนฺโถ ธมฺโม คนฺถสฺส จ โนคนฺถสฺส จ ธมฺมสฺส อตฺถิปจฺจเยน ปจฺจโย. สหชาตํ, ปุเรชาตํ.  \textbf{สหชาโต} โนคนฺโถ เอโก ขนฺโธ ติณฺณนฺนํ ขนฺธานํ คนฺถานญฺจ จิตฺตสมุฏฺฐานานญฺจ รูปานํ อตฺถิปจฺจเยน ปจฺจโย ฯเปฯ. \textbf{ปุเรชาตํ} จกฺขุํ ฯเปฯ วตฺถุํ อสฺสาเทติ อภินนฺทติ, ตํ อารพฺภ คนฺถา จ สมฺปยุตฺตกา จ ขนฺธา อุปฺปชฺชนฺติ. วตฺถุ คนฺถานํ สมฺปยุตฺตกานญฺจ ขนฺธานํ อตฺถิปจฺจเยน ปจฺจโย. (3)}

\proseitem{36}{คนฺโถ จ โนคนฺโถ จ ธมฺมา คนฺถสฺส ธมฺมสฺส อตฺถิปจฺจเยน ปจฺจโย. สหชาตํ, ปุเรชาตํ.  สหชาโต สีลพฺพต\-ปรามาโส กายคนฺโถ จ สมฺปยุตฺตกา จ ขนฺธา อภิชฺโฌกายคนฺถสฺส อตฺถิปจฺจเยน ปจฺจโย. (จกฺกํ.) สหชาโต สีลพฺพต\-ปรามาโส กายคนฺโถ จ วตฺถุ จ อภิชฺฌา\-กายคนฺถสฺส อตฺถิปจฺจเยน ปจฺจโย. (จกฺกํ.) (1)}

\prose{คนฺโถ จ โนคนฺโถ จ ธมฺมา โนคนฺถสฺส ธมฺมสฺส อตฺถิปจฺจเยน ปจฺจโย. สหชาตํ, ปุเรชาตํ, ปจฺฉาชาตํ, อาหารํ, อินฺทฺริยํ.  สหชาโต โนคนฺโถ เอโก ขนฺโธ จ คนฺโถ จ ติณฺณนฺนํ ขนฺธานํ จิตฺตสมุฏฺฐานานญฺจ รูปานํ อตฺถิปจฺจเยน ปจฺจโย ฯเปฯ เทฺว ขนฺธา จ ฯเปฯ. สหชาตา คนฺถา จ วตฺถุ จ โนคนฺถานํ ขนฺธานํ อตฺถิปจฺจเยน ปจฺจโย. สหชาตา คนฺถา จ มหาภูตา จ จิตฺตสมุฏฺฐานานญฺจ รูปานํ อตฺถิปจฺจเยน ปจฺจโย. ปจฺฉาชาตา คนฺถา จ สมฺปยุตฺตกา จ ขนฺธา ปุเรชาตสฺส อิมสฺส กายสฺส อตฺถิปจฺจเยน ปจฺจโย. ปจฺฉาชาตา คนฺถา จ กพฬีกาโร อาหาโร จ อิมสฺส กายสฺส อตฺถิปจฺจเยน ปจฺจโย. ปจฺฉาชาตา คนฺถา จ รูปชีวิตินฺทฺริยญฺจ กฏตฺตารูปานํ อตฺถิปจฺจเยน ปจฺจโย. (2)}

\prose{คนฺโถ จ โนคนฺโถ จ ธมฺมา คนฺถสฺส จ โนคนฺถสฺส จ ธมฺมสฺส อตฺถิปจฺจเยน ปจฺจโย. สหชาตํ, ปุเรชาตํ.  สหชาโต โนคนฺโถ เอโก ขนฺโธ จ สีลพฺพต\-ปรามาโส กายคนฺโถ จ ติณฺณนฺนํ ขนฺธานํ อภิชฺฌา\-กายคนฺถสฺส จ จิตฺตสมุฏฺฐานานญฺจ รูปานํ อตฺถิปจฺจเยน ปจฺจโย ฯเปฯ. สหชาโต สีลพฺพต\-ปรามาโส กายคนฺโถ จ วตฺถุ จ อภิชฺฌา\-กายคนฺถสฺส จ สมฺปยุตฺตกานญฺจ ขนฺธานํ อตฺถิปจฺจเยน ปจฺจโย. (จกฺกํ.) (3)}

\csromanmarkh{1. ปจฺจยานุโลม}\csromanmarkhii{2. สงฺขฺยาวาร}\csromanheadernomark{2}{1. \textbf{ปจฺจยานุโลม 2. สงฺขฺยาวาร}}
\proseitem{37}{\csromanfolio{266}เหตุยา นว, อารมฺมเณ นว, (สพฺพตฺถ นว,) อุปนิสฺสเย นว, ปุเรชาเต ตีณิ, ปจฺฉาชาเต ตีณิ, อาเสวเน นว, กมฺเม ตีณิ, วิปาเก เอกํ, อาหาเร ตีณิ, อินฺทฺริเย ตีณิ, ฌาเน ตีณิ, มคฺเค นว, สมฺปยุตฺเต นว, วิปฺปยุตฺเต ปญฺจ, อตฺถิยา นว, นตฺถิยา นว, วิคเต นว, อวิคเต นว.}

\vspace{\tipitakaparskip}
\csromancenter{อนุโลมํ.}
\proseitem{38}{คนฺโถ ธมฺโม คนฺถสฺส ธมฺมสฺส อารมฺมณ\-ปจฺจเยน ปจฺจโย. สหชาต\-ปจฺจเยน ปจฺจโย. อุปนิสฺสย\-ปจฺจเยน ปจฺจโย. (1)}

\prose{คนฺโถ ธมฺโม โนคนฺถสฺส ธมฺมสฺส อารมฺมณ\-ปจฺจเยน ปจฺจโย. สหชาต\-ปจฺจเยน ปจฺจโย. อุปนิสฺสย\-ปจฺจเยน ปจฺจโย. ปจฺฉาชาต\-ปจฺจเยน ปจฺจโย. (2)}

\prose{คนฺโถ ธมฺโม คนฺถสฺส จ โนคนฺถสฺส จ ธมฺมสฺส อารมฺมณ\-ปจฺจเยน ปจฺจโย. สหชาต\-ปจฺจเยน ปจฺจโย. อุปนิสฺสย\-ปจฺจเยน ปจฺจโย. (3)}

\proseitem{39}{โนคนฺโถ ธมฺโม โนคนฺถสฺส ธมฺมสฺส อารมฺมณ\-ปจฺจเยน ปจฺจโย. สหชาต\-ปจฺจเยน ปจฺจโย. อุปนิสฺสย\-ปจฺจเยน ปจฺจโย. ปุเรชาต\-ปจฺจเยน ปจฺจโย. ปจฺฉาชาต\-ปจฺจเยน ปจฺจโย. กมฺมปจฺจเยน ปจฺจโย. อาหารปจฺจเยน ปจฺจโย. อินฺทฺริย\-ปจฺจเยน ปจฺจโย. (1)}

\prose{โนคนฺโถ ธมฺโม คนฺถสฺส ธมฺมสฺส อารมฺมณ\-ปจฺจเยน ปจฺจโย. สหชาต\-ปจฺจเยน ปจฺจโย. อุปนิสฺสย\-ปจฺจเยน ปจฺจโย. ปุเรชาต\-ปจฺจเยน ปจฺจโย. (2)}

\prose{โนคนฺโถ ธมฺโม คนฺถสฺส จ โนคนฺถสฺส จ ธมฺมสฺส อารมฺมณ\-ปจฺจเยน ปจฺจโย. สหชาต\-ปจฺจเยน ปจฺจโย. อุปนิสฺสย\-ปจฺจเยน ปจฺจโย. ปุเรชาต\-ปจฺจเยน ปจฺจโย. (3)}

\csromanmatikaanchor{mtk.33}
\csromantocmarkref{section}{27. คนฺถนิยทุก}{mtk.33}
\bookmark[dest={mtk.33},level=1]{27. คนฺถนิยทุก}
\csromanheader{1}{27. คนฺถนิยทุก}
\proseitem{40}{\csromanfolio{267}คนฺโถ จ โนคนฺโถ จ ธมฺมา คนฺถสฺส ธมฺมสฺส อารมฺมณ\-ปจฺจเยน ปจฺจโย. สหชาต\-ปจฺจเยน ปจฺจโย. อุปนิสฺสย\-ปจฺจเยน ปจฺจโย. (1)}

\prose{คนฺโถ จ โนคนฺโถ จ ธมฺมา โนคนฺถสฺส ธมฺมสฺส อารมฺมณ\-ปจฺจเยน ปจฺจโย. สหชาต\-ปจฺจเยน ปจฺจโย. อุปนิสฺสย\-ปจฺจเยน ปจฺจโย. ปจฺฉาชาต\-ปจฺจเยน ปจฺจโย. (2)}

\prose{คนฺโถ จ โนคนฺโถ จ ธมฺมา คนฺถสฺส จ โนคนฺถสฺส จ ธมฺมสฺส อารมฺมณ\-ปจฺจเยน ปจฺจโย. สหชาต\-ปจฺจเยน ปจฺจโย. อุปนิสฺสย\-ปจฺจเยน ปจฺจโย. (3)}

\csromanmarkh{2. ปจฺจย\-ปจฺจนีย}\csromanmarkhii{2. สงฺขฺยาวาร}\csromanheadernomark{2}{2. ปจฺจย\-ปจฺจนีย 2. สงฺขฺยาวาร}
\vspace{\tipitakaparskip}
\csromancenter{นเหตุยา นว, นอารมฺมเณ นว, (สพฺพตฺถ นว,) โนอวิคเต นว.}
\proseitem{42}{เหตุปจฺจยา นอารมฺมเณ นว ฯเปฯ นสมนนฺตเร นว, นอญฺญมญฺเญ ตีณิ, นอุปนิสฺสเย นว, (สพฺพตฺถ นว,) นมคฺเค นว, นสมฺปยุตฺเต ตีณิ, นวิปฺปยุตฺเต นว, โนนตฺถิยา นว, โนวิคเต นว.}

\proseitemwithcloser{43}{นเหตุปจฺจยา อารมฺมเณ นว, อธิปติยา นว, (อนุโลมปทานิ ปริปุณฺณานิ กาตพฺพานิ.) อวิคเต นว.}{คนฺถทุกํ นิฏฺฐิตํ.}
\csromanheader{2}{27. คนฺถนิยทุก 1-7. วารสตฺตก}
\proseitem{44}{คนฺถนิยํ ธมฺมํ ปฏิจฺจ คนฺถนิโย ธมฺโม อุปฺปชฺชติ เหตุปจฺจยา. คนฺถนิยํ เอกํ ขนฺธํ. (สํขิตฺตํ.)}

\prosewithcloser{(ยถา จูฬนฺตรทุเก โลกิยทุกํ, เอวํ วิภชิตพฺพํ. นินฺนานากรณํ.)}{คนฺถนิยทุกํ นิฏฺฐิตํ.}
\csromanmatikaanchor{mtk.34}
\csromantocmarkref{section}{28. คนฺถสมฺปยุตฺตทุก}{mtk.34}
\bookmark[dest={mtk.34},level=1]{28. คนฺถสมฺปยุตฺตทุก}
\csromanmarkh{28. คนฺถ\-สมฺปยุตฺตทุก}\csromanmarkhii{1. ปฏิจฺจวาร}\csromanheadernomark{1}{28. คนฺถ\-สมฺปยุตฺตทุก 1. ปฏิจฺจวาร}
\csromanmarkh{1. ปจฺจยานุโลม}\csromanmarkhii{1. วิภงฺควาร}\csromanheadernomark{2}{1. \textbf{ปจฺจยานุโลม 1. วิภงฺควาร}}
\csromanheader{2}{\textbf{เหตุ}}
\proseitem{45}{\csromanfolio{268}คนฺถ\-สมฺปยุตฺตํ ธมฺมํ ปฏิจฺจ คนฺถ\-สมฺปยุตฺโต ธมฺโม อุปฺปชฺชติ เหตุปจฺจยา คนฺถ\-สมฺปยุตฺตํ เอกํ ขนฺธํ ปฏิจฺจ ตโย ขนฺธา ฯเปฯ เทฺว ขนฺเธ ฯเปฯ. (1)}

\prose{คนฺถ\-สมฺปยุตฺตํ ธมฺมํ ปฏิจฺจ คนฺถ\-วิปฺปยุตฺโต ธมฺโม อุปฺปชฺชติ เหตุปจฺจยา. คนฺถ\-สมฺปยุตฺเต ขนฺเธ ปฏิจฺจ จิตฺต\-สมุฏฺฐานํ รูปํ. ทิฏฺฐิคต\-วิปฺปยุตฺต\-โลภ\-สหคเต ขนฺเธ ปฏิจฺจ โลโภ จิตฺต\-สมุฏฺฐานญฺจ รูปํ. โทมนสฺส\-สหคเต ขนฺเธ ปฏิจฺจ ปฏิฆํ จิตฺต\-สมุฏฺฐานญฺจ รูปํ. (2)}

\prose{คนฺถ\-สมฺปยุตฺตํ ธมฺมํ ปฏิจฺจ คนฺถ\-สมฺปยุตฺโต จ คนฺถ\-วิปฺปยุตฺโต จ ธมฺมา อุปฺปชฺชนฺติ เหตุปจฺจยา. คนฺถ\-สมฺปยุตฺตํ เอกํ ขนฺธํ ปฏิจฺจ ตโย ขนฺธา จิตฺต\-สมุฏฺฐานญฺจ รูปํ ฯเปฯ เทฺว ขนฺเธ ฯเปฯ. ทิฏฺฐิคต\-วิปฺปยุตฺต\-โลภ\-สหคตํ เอกํ ขนฺธํ ปฏิจฺจ ตโย ขนฺธา โลโภ จ จิตฺต\-สมุฏฺฐานญฺจ รูปํ ฯเปฯ เทฺว ขนฺเธ ฯเปฯ. โทมนสฺส\-สหคตํ เอกํ ขนฺธํ ปฏิจฺจ ตโย ขนฺธา ปฏิฆญฺจ จิตฺต\-สมุฏฺฐานญฺจ รูปํ ฯเปฯ เทฺว ขนฺเธ ฯเปฯ. (3)}

\proseitem{46}{คนฺถ\-วิปฺปยุตฺตํ ธมฺมํ ปฏิจฺจ คนฺถ\-วิปฺปยุตฺโต ธมฺโม อุปฺปชฺชติ เหตุปจฺจยา. คนฺถ\-วิปฺปยุตฺตํ เอกํ ขนฺธํ ปฏิจฺจ ตโย ขนฺธา จิตฺต\-สมุฏฺฐานญฺจ รูปํ ฯเปฯ เทฺว ขนฺเธ ฯเปฯ. ทิฏฺฐิคต\-วิปฺปยุตฺตํ โลภํ ปฏิจฺจ จิตฺต\-สมุฏฺฐานํ รูปํ. ปฏิฆํ ปฏิจฺจ จิตฺต\-สมุฏฺฐานํ รูปํ. ปฏิสนฺธิ\-กฺขเณ ฯเปฯ ขนฺเธ ปฏิจฺจ. (สํขิตฺตํ.) (1)}

\prose{คนฺถ\-วิปฺปยุตฺตํ ธมฺมํ ปฏิจฺจ คนฺถ\-สมฺปยุตฺโต ธมฺโม อุปฺปชฺชติ เหตุปจฺจยา. ทิฏฺฐิคต\-วิปฺปยุตฺตํ โลภํ ปฏิจฺจ สมฺปยุตฺตกา ขนฺธา. ปฏิฆํ ปฏิจฺจ สมฺปยุตฺตกา ขนฺธา. (2)}

\prose{คนฺถ\-วิปฺปยุตฺตํ ธมฺมํ ปฏิจฺจ คนฺถ\-สมฺปยุตฺโต จ คนฺถ\-วิปฺปยุตฺโต จ ธมฺมา อุปฺปชฺชนฺติ เหตุปจฺจยา. ทิฏฺฐคตวิปฺปยุตฺตํ โลภํ ปฏิจฺจ สมฺปยุตฺตกา ขนฺธา จิตฺต\-สมุฏฺฐานญฺจ รูปํ. ปฏิฆํ ปฏิจฺจ สมฺปยุตฺตกา ขนฺธา จิตฺต\-สมุฏฺฐานญฺจ รูปํ. (3)}

\proseitem{47}{คนฺถ\-สมฺปยุตฺตญฺจ คนฺถ\-วิปฺปยุตฺตญฺจ ธมฺมํ ปฏิจฺจ คนฺถ\-สมฺปยุตฺโต ธมฺโม อุปฺปชฺชติ เหตุปจฺจยา. ทิฏฺฐิคต\-วิปฺปยุตฺต\-โลภ\-สหคตํ เอกํ ขนฺธญฺจ \csromanfolio{269}โลภญฺจ ปฏิจฺจ ตโย ขนฺธา ฯเปฯ เทฺว ขนฺเธ ฯเปฯ. โทมนสฺส\-สหคตํ เอกํ ขนฺธญฺจ ปฏิฆญฺจ ปฏิจฺจ ตโย ขนฺธา ฯเปฯ เทฺว ขนฺเธ ฯเปฯ. (1)}

\prose{คนฺถ\-สมฺปยุตฺตญฺจ คนฺถ\-วิปฺปยุตฺตญฺจ ธมฺมํ ปฏิจฺจ คนฺถ\-วิปฺปยุตฺโต ธมฺโม อุปฺปชฺชติ เหตุปจฺจยา. คนฺถ\-สมฺปยุตฺเต ขนฺเธ จ มหาภูเต จ ปฏิจฺจ จิตฺต\-สมุฏฺฐานํ รูปํ. ทิฏฺฐิคต\-วิปฺปยุตฺต\-โลภ\-สหคเต ขนฺเธ จ โลภญฺจ ปฏิจฺจ จิตฺต\-สมุฏฺฐานํ รูปํ. โทมนสฺส\-สหคเต ขนฺเธ จ ปฏิฆญฺจ ปฏิจฺจ จิตฺต\-สมุฏฺฐานํ รูปํ. (2)}

\prose{คนฺถ\-สมฺปยุตฺตญฺจ คนฺถ\-วิปฺปยุตฺตญฺจ ธมฺมํ ปฏิจฺจ คนฺถ\-สมฺปยุตฺโต จ คนฺถ\-วิปฺปยุตฺโต จ ธมฺมา อุปฺปชฺชนฺติ เหตุปจฺจยา. ทิฏฺฐิคต\-วิปฺปยุตฺต\-โลภ\-สหคตํ เอกํ ขนฺธญฺจ โลภญฺจ ปฏิจฺจ ตโย ขนฺธา จิตฺต\-สมุฏฺฐานญฺจ รูปํ ฯเปฯ เทฺว ขนฺเธ ฯเปฯ. โทมนสฺส\-สหคตํ เอกํ ขนฺธญฺจ ปฏิฆญฺจ ปฏิจฺจ ตโย ขนฺธา จิตฺต\-สมุฏฺฐานญฺจ รูปํ ฯเปฯ เทฺว ขนฺเธ ฯเปฯ. (3)}

\prose{คนฺถ\-สมฺปยุตฺตํ ธมฺมํ ปฏิจฺจ คนฺถ\-สมฺปยุตฺโต ธมฺโม อุปฺปชฺชติ อารมฺมณ\-ปจฺจยา. คนฺถ\-สมฺปยุตฺตํ เอกํ ขนฺธํ ปฏิจฺจ ตโย ขนฺธา ฯเปฯ เทฺว ขนฺเธ ฯเปฯ. (1)}

\prose{คนฺถ\-สมฺปยุตฺตํ ธมฺมํ ปฏิจฺจ คนฺถ\-วิปฺปยุตฺโต ธมฺโม อุปฺปชฺชติ อารมฺมณ\-ปจฺจยา. ทิฏฺฐิคต\-วิปฺปยุตฺต\-โลภ\-สหคเต ขนฺเธ ปฏิจฺจ โลโภ. โทมนสฺส\-สหคเต ขนฺเธ ปฏิจฺจ ปฏิฆํ. (2)}

\prose{คนฺถ\-สมฺปยุตฺตํ ธมฺมํ ปฏิจฺจ คนฺถ\-สมฺปยุตฺโต จ คนฺถ\-วิปฺปยุตฺโต จ ธมฺมา อุปฺปชฺชนฺติ อารมฺมณ\-ปจฺจยา. ทิฏฺฐิคต\-วิปฺปยุตฺต\-โลภ\-สหคตํ เอกํ ขนฺธํ ปฏิจฺจ ตโย ขนฺธา โลโภ จ ฯเปฯ เทฺว ขนฺเธ ฯเปฯ. โทมนสฺส\-สหคตํ เอกํ ขนฺธํ ปฏิจฺจ ตโย ขนฺธา ปฏิฆญฺจ ฯเปฯ เทฺว ขนฺเธ ฯเปฯ. (3)}

\proseitem{49}{คนฺถ\-วิปฺปยุตฺตํ ธมฺมํ ปฏิจฺจ คนฺถ\-วิปฺปยุตฺโต ธมฺโม อุปฺปชฺชติ อารมฺมณ\-ปจฺจยา. คนฺถ\-วิปฺปยุตฺตํ เอกํ ขนฺธํ ปฏิจฺจ ตโย ขนฺธา ฯเปฯ เทฺว ขนฺเธ ฯเปฯ. ปฏิสนฺธิ\-กฺขเณ ฯเปฯ วตฺถุํ ปฏิจฺจ ขนฺธา. (1)}

\prose{คนฺถ\-วิปฺปยุตฺตํ ธมฺมํ ปฏิจฺจ คนฺถ\-สมฺปยุตฺโต ธมฺโม อุปฺปชฺชติ อารมฺมณ\-ปจฺจยา. ทิฏฺฐิคต\-วิปฺปยุตฺตํ โลภํ ปฏิจฺจ สมฺปยุตฺตกา ขนฺธา. ปฏิฆํ ปฏิจฺจ สมฺปยุตฺตกา ขนฺธา. (2)}

\prose{\csromanfolio{270}คนฺถ\-สมฺปยุตฺตญฺจ คนฺถ\-วิปฺปยุตฺตญฺจ ธมฺมํ ปฏิจฺจ คนฺถ\-สมฺปยุตฺโต ธมฺโม อุปฺปชฺชติ อารมฺมณ\-ปจฺจยา. ทิฏฺฐิคต\-วิปฺปยุตฺต\-โลภ\-สหคตํ เอกํ ขนฺธญฺจ โลภญฺจ ปฏิจฺจ ตโย ขนฺธา ฯเปฯ เทฺว ขนฺเธ ฯเปฯ. โทมนสฺส\-สหคตํ เอกํ ขนฺธญฺจ ปฏิฆญฺจ ปฏิจฺจ ตโย ขนฺธา ฯเปฯ เทฺว ขนฺเธ ฯเปฯ. (สํขิตฺตํ.) (1)}

\csromanmarkh{1. ปจฺจยานุโลม}\csromanmarkhii{2. สงฺขฺยาวาร}\csromanheadernomark{2}{1. \textbf{ปจฺจยานุโลม 2. สงฺขฺยาวาร}}
\proseitem{50}{เหตุยา นว, อารมฺมเณ ฉ, อธิปติยา นว, อนนฺตเร ฉ, สมนนฺตเร ฉ, สหชาเต นว, อญฺญมญฺเญ ฉ, นิสฺสเย นว, อุปนิสฺสเย ฉ, ปุเรชาเต ฉ, อาเสวเน ฉ, กมฺเม นว, วิปาเก เอกํ, อาหาเร นว, อินฺทฺริเย นว, ฌาเน นว, มคฺเค นว, สมฺปยุตฺเต ฉ, วิปฺปยุตฺเต นว, อตฺถิยา นว, นตฺถิยา ฉ, วิคเต ฉ, อวิคเต นว.}

\vspace{\tipitakaparskip}
\csromancenter{\textbf{อนุโลมํ.}}
\csromanmarkh{2. ปจฺจย\-ปจฺจนีย}\csromanmarkhii{1. วิภงฺควาร}\csromanheadernomark{2}{2. ปจฺจย\-ปจฺจนีย 1. วิภงฺควาร}
\proseitem{51}{คนฺถ\-วิปฺปยุตฺตํ ธมฺมํ ปฏิจฺจ คนฺถ\-วิปฺปยุตฺโต ธมฺโม อุปฺปชฺชติ นเหตุปจฺจยา อเหตุกํ คนฺถ\-วิปฺปยุตฺตํ เอกํ ขนฺธํ ปฏิจฺจ ตโย ขนฺธา จิตฺต\-สมุฏฺฐานญฺจ รูปํ ฯเปฯ เทฺว ขนฺเธ ฯเปฯ. อเหตุก\-ปฏิสนฺธิ\-กฺขเณ ฯเปฯ. (ยาว อสญฺญสตฺตา.) วิจิกิจฺฉา\-สหคเต อุทฺธจฺจ\-สหคเต ขนฺเธ ปฏิจฺจ วิจิกิจฺฉา\-สหคโต อุทฺธจฺจ\-สหคโต โมโห. (สํขิตฺตํ.) (1)}

\csromanmarkh{2. ปจฺจย\-ปจฺจนีย}\csromanmarkhii{2. สงฺขฺยาวาร}\csromanheadernomark{2}{2. ปจฺจย\-ปจฺจนีย 2. สงฺขฺยาวาร}
\vspace{\tipitakaparskip}
\csromancenter{\textbf{นเหตุ}ยา เอกํ, นอารมฺมเณ ตีณิ, นอธิปติยา นว, นอนนฺตเร ตีณิ, นสมนนฺตเร ตีณิ, นอญฺญมญฺเญ ตีณิ, นฺเอาปนิสฺสเย ตีณิ, นปุเรชาเต สตฺต. (นปุเรชาเต วิภชนฺเตน อรูปํ ปฐมํ}
\prosecont{\csromanfolio{271}กาตพฺพํ. รูปํ ยตฺถ ลพฺภติ, ปจฺฉา กาตพฺพํ. ปฏิฆญฺจ อรูเป นตฺถิ.) นปจฺฉาชาเต นว, นอาเสวเน นว, นกมฺเม จตฺตาริ, นวิปาเก นว, นอาหาเร เอกํ, ไนนฺทฺริเย เอกํ, นฌาเน เอกํ, นมคฺเค เอกํ, นสมฺปยุตฺเต ตีณิ, นวิปฺปยุตฺเต ฉ, โนนตฺถิยา ตีณิ, โนวิคเต ตีณิ.}

\vspace{\tipitakaparskip}
\csromancenter{\textbf{เหตุ}ปจฺจยา นอารมฺมเณ ตีณิ, นอธิปติยา นว, (เอวํ คเณตพฺพํ. สํขิตฺตํ.) โนวิคเต ตีณิ.}
\csromancenter{นเหตุปจฺจยา อารมฺมเณ เอกํ ฯเปฯ อวิคเต เอกํ.}
\csromancenter{(สหชาตวาโรปิ เอวํ กาตพฺโพ.)}
\csromanmarkh{28. คนฺถ\-สมฺปยุตฺตทุก}\csromanmarkhii{3. ปจฺจยวาร}\csromanheadernomark{2}{28. คนฺถ\-สมฺปยุตฺตทุก 3. ปจฺจยวาร}
\csromanmarkh{1. ปจฺจยานุโลม}\csromanmarkhii{1. วิภงฺควาร}\csromanheadernomark{2}{1. \textbf{ปจฺจยานุโลม 1. วิภงฺควาร}}
\vspace{\tipitakaparskip}
\csromancenter{\textbf{เหตุ}}
\proseitem{52}{คนฺถ\-สมฺปยุตฺตํ ธมฺมํ ปจฺจยา คนฺถ\-สมฺปยุตฺโต ธมฺโม อุปฺปชฺชติ เหตุปจฺจยา. ตีณิ. (ปฏิจฺจสทิสา.)}

\prose{คนฺถ\-วิปฺปยุตฺตํ ธมฺมํ ปจฺจยา คนฺถ\-วิปฺปยุตฺโต ธมฺโม อุปฺปชฺชติ เหตุปจฺจยา. คนฺถ\-วิปฺปยุตฺตํ เอกํ ขนฺธํ ปจฺจยา ฯเปฯ. ปฏิสนฺธิ\-กฺขเณ ฯเปฯ. ขนฺเธ ปจฺจยา วตฺถุ, วตฺถุํ ปจฺจยา ขนฺธา. เอกํ มหาภูตํ ฯเปฯ วตฺถุํ ปจฺจยา คนฺถ\-วิปฺปยุตฺตา ขนฺธา. (1)}

\proseitem{55}{\csromanfolio{272}คนฺถ\-วิปฺปยุตฺตํ ธมฺมํ ปจฺจยา คนฺถ\-สมฺปยุตฺโต ธมฺโม อุปฺปชฺชติ เหตุปจฺจยา. วตฺถุํ ปจฺจยา คนฺถ\-สมฺปยุตฺตกา ขนฺธา. ทิฏฺฐิคต\-วิปฺปยุตฺตํ โลภํ ปจฺจยา สมฺปยุตฺตกา ขนฺธา. ปฏิฆํ ปจฺจยา สมฺปยุตฺตกา ขนฺธา. (2)}

\prose{คนฺถ\-วิปฺปยุตฺตํ ธมฺมํ ปจฺจยา คนฺถ\-สมฺปยุตฺโต จ คนฺถ\-วิปฺปยุตฺโต จ ธมฺมา อุปฺปชฺชนฺติ เหตุปจฺจยา. วตฺถุํ ปจฺจยา คนฺถ\-สมฺปยุตฺตกา ขนฺธา. มหาภูเต ปจฺจยา จิตฺต\-สมุฏฺฐานํ รูปํ. ทิฏฺฐิคต\-วิปฺปยุตฺตํ โลภํ ปจฺจยา สมฺปยุตฺตกา ขนฺธา จิตฺต\-สมุฏฺฐานญฺจ รูปํ. ปฏิฆํ ปจฺจยา สมฺปยุตฺตกา ขนฺธา จิตฺต\-สมุฏฺฐานญฺจ รูปํ. วตฺถุํ ปจฺจยา ทิฏฺฐิคต\-วิปฺปยุตฺต\-โลภ\-สหคตา ขนฺธา จ โลโภ จ. วตฺถุํ ปจฺจยา โทมนสฺส\-สหคตา ขนฺธา จ ปฏิฆญฺจ. (3)}

\proseitem{56}{คนฺถ\-สมฺปยุตฺตญฺจ คนฺถ\-วิปฺปยุตฺตญฺจ ธมฺมํ ปจฺจยา คนฺถ\-สมฺปยุตฺโต ธมฺโม อุปฺปชฺชติ เหตุปจฺจยา. คนฺถ\-สมฺปยุตฺตํ เอกํ ขนฺธญฺจ วตฺถุญฺจ ปจฺจยา ตโย ขนฺธา ฯเปฯ เทฺว ขนฺเธ ฯเปฯ. ทิฏฺฐิคต\-วิปฺปยุตฺต\-โลภ\-สหคตํ เอกํ ขนฺธญฺจ วตฺถุญฺจ โลภญฺจ ปจฺจยา ตโย ขนฺธา ฯเปฯ เทฺว ขนฺเธ ฯเปฯ. โทมนสฺส\-สหคตํ เอกํ ขนฺธญฺจ วตฺถุญฺจ ปฏิฆญฺจ ปจฺจยา ตโย ขนฺธา ฯเปฯ เทฺว ขนฺเธ ฯเปฯ. (1)}

\prose{คนฺถ\-สมฺปยุตฺตญฺจ คนฺถ\-วิปฺปยุตฺตญฺจ ธมฺมํ ปจฺจยา คนฺถ\-วิปฺปยุตฺโต ธมฺโม อุปฺปชฺชติ เหตุปจฺจยา. คนฺถ\-สมฺปยุตฺเต ขนฺเธ จ มหาภูเต จ ปจฺจยา จิตฺต\-สมุฏฺฐานํ รูปํ. ทิฏฺฐิคต\-วิปฺปยุตฺต\-โลภ\-สหคเต ขนฺเธ จ โลภญฺจ ปจฺจยา จิตฺต\-สมุฏฺฐานํ รูปํ. โทมนสฺส\-สหคเต ขนฺเธ จ ปฏิฆญฺจ ปจฺจยา จิตฺต\-สมุฏฺฐานํ รูปํ. ทิฏฺฐิคต\-วิปฺปยุตฺต\-โลภ\-สหคเต ขนฺเธ จ วตฺถุญฺจ ปจฺจยา โลโภ. โทมนสฺส\-สหคเต ขนฺเธ จ วตฺถุญฺจ ปจฺจยา ปฏิฆํ. (2)}

\prose{คนฺถ\-สมฺปยุตฺตญฺจ คนฺถ\-วิปฺปยุตฺตญฺจ ธมฺมํ ปจฺจยา คนฺถ\-สมฺปยุตฺโต จ คนฺถ\-วิปฺปยุตฺโต จ ธมฺมา อุปฺปชฺชนฺติ เหตุปจฺจยา. คนฺถ\-สมฺปยุตฺตํ เอกํ ขนฺธญฺจ วตฺถุญฺจ ปจฺจยา ตโย ขนฺธา ฯเปฯ เทฺว ขนฺเธ ฯเปฯ. คนฺถ\-สมฺปยุตฺเต ขนฺเธ จ มหาภูเต จ ปจฺจยา จิตฺต\-สมุฏฺฐานํ รูปํ. ทิฏฺฐิคต\-วิปฺปยุตฺต\-โลภ\-สหคตํ เอกํ ขนฺธญฺจ โลภญฺจ ปจฺจยา ตโย ขนฺธา จิตฺต\-สมุฏฺฐานญฺจ รูปํ ฯเปฯ เทฺว ขนฺเธ ฯเปฯ. โทมนสฺส\-สหคตํ เอกํ ขนฺธญฺจ ปฏิฆญฺจ ปจฺจยา ตโย ขนฺธา จิตฺต\-สมุฏฺฐานญฺจ รูปํ ฯเปฯ เทฺว ขนฺเธ ฯเปฯ. ทิฏฺฐิคต\-วิปฺปยุตฺต\-โลภ\-สหคตํ เอกํ ขนฺธญฺจ วตฺถุญฺจ ปจฺจยา ตโย ขนฺธา โลโภ จ ฯเปฯ เทฺว ขนฺเธ ฯเปฯ. โทมนสฺส\-สหคตํ เอกํ ขนฺธญฺจ วตฺถุญฺจ ปจฺจยา ตโย ขนฺธา ปฏิฆญฺจ ฯเปฯ เทฺว ขนฺเธ ฯเปฯ. (สํขิตฺตํ.) (3)}

\csromanmarkh{28. คนฺถ\-สมฺปยุตฺตทุก}\csromanmarkhii{1. ปจฺจยานุโลม 2. สงฺขฺยาวาร}\csromanheadernomark{2}{28. คนฺถ\-สมฺปยุตฺตทุก 1. ปจฺจยานุโลม 2. สงฺขฺยาวาร}
\proseitem{57}{\csromanfolio{273}เหตุยา นว, อารมฺมเณ นว, อธิปติยา นว, อนนฺตเร นว, สมนนฺตเร นว, (สพฺพตฺถ นว.) วิปาเก เอกํ, อาหาเร นว ฯเปฯ อวิคเต นว.}

\vspace{\tipitakaparskip}
\csromancenter{\textbf{อนุโลมํ.}}
\csromanmarkh{2. ปจฺจย\-ปจฺจนีย}\csromanmarkhii{1. วิภงฺควาร}\csromanheadernomark{2}{2. ปจฺจย\-ปจฺจนีย 1. วิภงฺควาร}
\proseitem{58}{คนฺถ\-วิปฺปยุตฺตํ ธมฺมํ ปจฺจยา คนฺถ\-วิปฺปยุตฺโต ธมฺโม อุปฺปชฺชติ นเหตุปจฺจยา. อเหตุกํ คนฺถ\-วิปฺปยุตฺตํ ฯเปฯ. อเหตุก\-ปฏิสนฺธิ\-กฺขเณ. (ยาว อสญฺญสตฺตา.) จกฺขายตนํ ปจฺจยา จกฺขุวิญฺญาณํ ฯเปฯ. กายายตนํ ปจฺจยา กายวิญฺญาณํ. วตฺถุํ ปจฺจยา อเหตุกา คนฺถ\-วิปฺปยุตฺตา ขนฺธา. วิจิกิจฺฉา\-สหคเต อุทฺธจฺจ\-สหคเต ขนฺเธ จ วตฺถุญฺจ ปจฺจยา วิจิกิจฺฉา\-สหคโต อุทฺธจฺจ\-สหคโต โมโห. (สํขิตฺตํ.)}

\csromanmarkh{2. ปจฺจย\-ปจฺจนีย}\csromanmarkhii{2. สงฺขฺยาวาร}\csromanheadernomark{2}{2. ปจฺจย\-ปจฺจนีย 2. สงฺขฺยาวาร}
\vspace{\tipitakaparskip}
\csromancenter{\textbf{นเหตุ}ยา เอกํ, นอารมฺมเณ ตีณิ, นอธิปติยา นว, นอนนฺตเร ตีณิ ฯเปฯ นฺเอาปนิสฺสเย ตีณิ, นปุเรชาเต สตฺต, นปจฺฉาชาเต นว, นอาเสวเน นว, นกมฺเม จตฺตาริ, นวิปาเก นว, นอาหาเร เอกํ, ไนนฺทฺริเย เอกํ, นฌาเน เอกํ, นมคฺเค เอกํ, นสมฺปยุตฺเต ตีณิ, นวิปฺปยุตฺเต ฉ, โนนตฺถิยา ตีณิ, โนวิคเต ตีณิ.}
\proseitem{60}{เหตุปจฺจยา นอารมฺมเณ ตีณิ, นอธิปติยา นว. (เอวํ คเณตพฺพํ.)}

\vspace{\tipitakaparskip}
\csromancenter{นเหตุปจฺจยา อารมฺมเณ เอกํ ฯเปฯ อวิคเต เอกํ.}
\csromancenter{(นิสฺสยวาโร ปจฺจยวาร\-สทิโส)}
\csromanmarkh{28. คนฺถ\-สมฺปยุตฺตทุก}\csromanmarkhii{5. สํสฏฺฐวาร}\csromanheadernomark{2}{28. คนฺถ\-สมฺปยุตฺตทุก 5. สํสฏฺฐวาร}
\csromanmarkh{1. ปจฺจยานุโลม}\csromanmarkhii{1. วิภงฺควาร}\csromanheadernomark{2}{1. \textbf{ปจฺจยานุโลม 1. วิภงฺควาร}}
\csromanheader{2}{\textbf{เหตุ}}
\proseitem{62}{\csromanfolio{274}คนฺถ\-สมฺปยุตฺตํ ธมฺมํ สํสฏฺโฐ คนฺถ\-สมฺปยุตฺโต ธมฺโม อุปฺปชฺชติ เหตุปจฺจยา. คนฺถ\-สมฺปยุตฺตํ เอกํ ขนฺธํ สํสฏฺฐา ตโย ขนฺธา ฯเปฯ เทฺว ขนฺเธ ฯเปฯ. (1)}

\prose{คนฺถ\-สมฺปยุตฺตํ ธมฺมํ สํสฏฺโฐ คนฺถ\-วิปฺปยุตฺโต ธมฺโม อุปฺปชฺชติ เหตุปจฺจยา. ทิฏฺฐิคต\-วิปฺปยุตฺต\-โลภ\-สหคเต ขนฺเธ สํสฏฺโฐ โลโภ. โทมนสฺส\-สหคเต ขนฺเธ สํสฏฺฐํ ปฏิฆํ. (2)}

\prose{คนฺถ\-สมฺปยุตฺตํ ธมฺมํ สํสฏฺโฐ คนฺถ\-สมฺปยุตฺโต จ คนฺถ\-วิปฺปยุตฺโต จ ธมฺมา อุปฺปชฺชนฺติ เหตุปจฺจยา. ทิฏฺฐิคต\-วิปฺปยุตฺต\-โลภ\-สหคตํ เอกํ ขนฺธํ สํสฏฺฐา ตโย ขนฺธา โลโภ จ ฯเปฯ เทฺว ขนฺเธ ฯเปฯ. โทมนสฺส\-สหคตํ เอกํ ขนฺธํ สํสฏฺฐา ตโย ขนฺธา ปฏิฆญฺจ. เทฺว ขนฺเธ ฯเปฯ. (3)}

\proseitem{63}{คนฺถ\-วิปฺปยุตฺตํ ธมฺมํ สํสฏฺโฐ คนฺถ\-วิปฺปยุตฺโต ธมฺโม อุปฺปชฺชติ เหตุปจฺจยา. คนฺถ\-วิปฺปยุตฺตํ เอกํ ขนฺธํ สํสฏฺฐา ตโย ขนฺธา ฯเปฯ เทฺว ขนฺเธ ฯเปฯ. ปฏิสนฺธิ\-กฺขเณ ฯเปฯ. (1)}

\prose{คนฺถ\-วิปฺปยุตฺตํ ธมฺมํ สํสฏฺโฐ คนฺถ\-สมฺปยุตฺโต ธมฺโม อุปฺปชฺชติ เหตุปจฺจยา. ทิฏฺฐิคต\-วิปฺปยุตฺตํ โลภํ สํสฏฺฐา สมฺปยุตฺตกา ขนฺธา. ปฏิฆํ สํสฏฺฐา สมฺปยุตฺตกา ขนฺธา. (2)}

\prose{คนฺถสมฺปยุตญฺจ คนฺถ\-วิปฺปยุตฺตญฺจ ธมฺมํ สํสฏฺโฐ คนฺถ\-สมฺปยุตฺโต ธมฺโม อุปฺปชฺชติ เหตุปจฺจยา. ทิฏฺฐิคต\-วิปฺปยุตฺตํ โลภสหคตํ เอกํ ขนฺธญฺจ \csromanfolio{275}โลภญฺจ สํสฏฺฐา ตโย ขนฺธา ฯเปฯ เทฺว ขนฺเธ ฯเปฯ. โทมนสฺส\-สหคตํ เอกํ ขนฺธญฺจ ปฏิฆญฺจ สํสฏฺฐา ตโย ขนฺธา ฯเปฯ เทฺว ขนฺเธ ฯเปฯ. (สํขิตฺตํ.)}

\csromanmarkh{1. ปจฺจยานุโลม}\csromanmarkhii{2. สงฺขฺยาวาร}\csromanheadernomark{2}{1. \textbf{ปจฺจยานุโลม 2. สงฺขฺยาวาร}}
\proseitem{64}{เหตุยา ฉ, อารมฺมเณ ฉ, อธิปติยา ฉ, (สพฺพตฺถ ฉ.) วิปาเก เอกํ, อาหาเร ฉ ฯเปฯ อวิคเต ฉ.}

\vspace{\tipitakaparskip}
\csromancenter{อนุโลมํ.}
\csromanmarkh{2. ปจฺจย\-ปจฺจนีย}\csromanmarkhii{1. วิภงฺควาร}\csromanheadernomark{2}{2. ปจฺจย\-ปจฺจนีย 1. วิภงฺควาร}
\proseitem{65}{คนฺถ\-วิปฺปยุตฺตํ ธมฺมํ สํสฏฺโฐ คนฺถ\-วิปฺปยุตฺโต ธมฺโม อุปฺปชฺชติ นเหตุปจฺจยา. อเหตุกํ คนฺถ\-วิปฺปยุตฺตํ ฯเปฯ. อเหตุก\-ปฏิสนฺธิ\-กฺขเณ ฯเปฯ. วิจิกิจฺฉา\-สหคเต อุทฺธจฺจ\-สหคเต ขนฺเธ สํสฏฺโฐ วิจิกิจฺฉา\-สหคโต อุทฺธจฺจ\-สหคโต โมโห. (สํขิตฺตํ.)}

\csromanmarkh{2. ปจฺจย\-ปจฺจนีย}\csromanmarkhii{2. สงฺขฺยาวาร}\csromanheadernomark{2}{2. ปจฺจย\-ปจฺจนีย 2. สงฺขฺยาวาร}
\vspace{\tipitakaparskip}
\csromancenter{\textbf{นเหตุ}ยา เอกํ, นอธิปติยา ฉ, นปุเรชาเต ฉ, นปจฺฉาชาเต ฉ, นอาเสวเน ฉ, นกมฺเม จตฺตาริ, นวิปาเก ฉ, นฌาเน เอกํ, นมคฺเค เอกํ, นวิปฺปยุตฺเต ฉ.}
\prose{เหตุปจฺจยา นอธิปติยา ฉ, นปุเรชาเต ฉ, นปจฺฉาชาเต ฉ, นอาเสวเน ฉ, นกมฺเม จตฺตาริ, นวิปาเก ฉ, นวิปฺปยุตฺเต ฉ.}

\vspace{\tipitakaparskip}
\csromancenter{นเหตุปจฺจยา อารมฺมเณ เอกํ, อนนฺตเร เอกํ ฯเปฯ อวิคเต เอกํ.}
\csromancenter{(สมฺปยุตฺตวาโร สํสฏฺฐวาร\-สทิโส)}
\csromanmarkh{28. คนฺถ\-สมฺปยุตฺตทุก}\csromanmarkhii{7. ปญฺหาวาร}\csromanheadernomark{2}{28. คนฺถ\-สมฺปยุตฺตทุก 7. ปญฺหาวาร}
\csromanmarkh{1. ปจฺจยานุโลม}\csromanmarkhii{1. วิภงฺควาร}\csromanheadernomark{2}{1. \textbf{ปจฺจยานุโลม 1. วิภงฺควาร}}
\csromanheader{2}{\textbf{เหตุ}}
\proseitem{69}{\csromanfolio{276}คนฺถ\-สมฺปยุตฺโต ธมฺโม คนฺถ\-สมฺปยุตฺตสฺส ธมฺมสฺส เหตุปจฺจเยน ปจฺจโย. คนฺถ\-สมฺปยุตฺตา เหตู สมฺปยุตฺตกานํ ขนฺธานํ เหตุปจฺจเยน ปจฺจโย. (1)}

\prose{คนฺถ\-สมฺปยุตฺโต ธมฺโม คนฺถ\-วิปฺปยุตฺตสฺส ธมฺมสฺส เหตุปจฺจเยน ปจฺจโย. คนฺถ\-สมฺปยุตฺตา เหตู จิตฺต\-สมุฏฺฐานานํ รูปานํ เหตุปจฺจเยน ปจฺจโย. ทิฏฺฐิคต\-วิปฺปยุตฺต\-โลภ\-สหคโต เหตุ โลภสฺส จิตฺตสมุฏฺฐานานญฺจ รูปานํ เหตุปจฺจเยน ปจฺจโย. โทมนสฺส\-สหคโต เหตุ ปฏิฆสฺส จิตฺตสมุฏฺฐานานญฺจ รูปานํ เหตุปจฺจเยน ปจฺจโย. (2)}

\prose{คนฺถ\-สมฺปยุตฺโต ธมฺโม คนฺถ\-สมฺปยุตฺตสฺส จ คนฺถ\-วิปฺปยุตฺตสฺส จ ธมฺมสฺส เหตุปจฺจเยน ปจฺจโย. คนฺถ\-สมฺปยุตฺตา เหตู สมฺปยุตฺตกานํ ขนฺธานํ จิตฺตสมุฏฺฐานานญฺจ รูปานํ เหตุปจฺจเยน ปจฺจโย. ทิฏฺฐิคต\-วิปฺปยุตฺต\-โลภ\-สหคโต เหตุ สมฺปยุตฺตกานํ ขนฺธานํ โลภสฺส จ จิตฺตสมุฏฺฐานานญฺจ รูปานํ เหตุปจฺจเยน ปจฺจโย. โทมนสฺส\-สหคโต เหตุ สมฺปยุตฺตกานํ ขนฺธานํ ปฏิฆสฺส จ จิตฺตสมุฏฺฐานานญฺจ รูปานํ เหตุปจฺจเยน ปจฺจโย. (3)}

\proseitem{70}{คนฺถ\-วิปฺปยุตฺโต ธมฺโม คนฺถ\-วิปฺปยุตฺตสฺส ธมฺมสฺส เหตุปจฺจเยน ปจฺจโย. คนฺถ\-วิปฺปยุตฺตา เหตู สมฺปยุตฺตกานํ ขนฺธานํ จิตฺตสมุฏฺฐานานญฺจ \csromanfolio{277}รูปานํ เหตุปจฺจเยน ปจฺจโย. ทิฏฺฐิคต\-วิปฺปยุตฺโต โลโภ จิตฺต\-สมุฏฺฐานานํ รูปานํ เหตุปจฺจเยน ปจฺจโย. ปฏิฆํ จิตฺต\-สมุฏฺฐานานํ รูปานํ เหตุปจฺจเยน ปจฺจโย. ปฏิสนฺธิ\-กฺขเณ ฯเปฯ. (1)}

\prose{คนฺถ\-วิปฺปยุตฺโต ธมฺโม คนฺถ\-สมฺปยุตฺตสฺส ธมฺมสฺส เหตุปจฺจเยน ปจฺจโย. ทิฏฺฐิคต\-วิปฺปยุตฺโต โลโภ สมฺปยุตฺตกานํ ขนฺธานํ เหตุปจฺจเยน ปจฺจโย. ปฏิฆํ สมฺปยุตฺตกานํ ขนฺธานํ เหตุปจฺจเยน ปจฺจโย. (2)}

\prose{คนฺถ\-วิปฺปยุตฺโต ธมฺโม คนฺถ\-สมฺปยุตฺตสฺส จ คนฺถ\-วิปฺปยุตฺตสฺส จ ธมฺมสฺส เหตุปจฺจเยน ปจฺจโย. ทิฏฺฐิคต\-วิปฺปยุตฺโต โลโภ สมฺปยุตฺตกานํ ขนฺธานํ จิตฺตสมุฏฺฐานานญฺจ รูปานํ เหตุปจฺจเยน ปจฺจโย. ปฏิฆํ สมฺปยุตฺตกานํ ขนฺธานํ จิตฺตสมุฏฺฐานานญฺจ รูปานํ เหตุปจฺจเยน ปจฺจโย. (3)}

\prose{คนฺถ\-สมฺปยุตฺโต จ คนฺถ\-วิปฺปยุตฺโต จ ธมฺมา คนฺถ\-สมฺปยุตฺตสฺส ธมฺมสฺส เหตุปจฺจเยน ปจฺจโย. ทิฏฺฐิคต\-วิปฺปยุตฺต\-โลภ\-สหคโต เหตุ จ โลโภ จ สมฺปยุตฺตกานํ ขนฺธานํ เหตุปจฺจเยน ปจฺจโย. โทมนสฺส\-สหคโต เหตุ จ ปฏิฆญฺจ สมฺปยุตฺตกานํ ขนฺธานํ เหตุปจฺจเยน ปจฺจโย. (1)}

\prose{คนฺถ\-สมฺปยุตฺโต จ คนฺถ\-วิปฺปยุตฺโต จ ธมฺมา คนฺถ\-วิปฺปยุตฺตสฺส ธมฺมสฺส เหตุปจฺจเยน ปจฺจโย. ทิฏฺฐิคต\-วิปฺปยุตฺต\-โลภ\-สหคโต เหตุ จ โลโภ จ จิตฺต\-สมุฏฺฐานานํ รูปานํ เหตุปจฺจเยน ปจฺจโย. โทมนสฺส\-สหคโต เหตุ จ ปฏิฆญฺจ จิตฺต\-สมุฏฺฐานานํ รูปานํ เหตุปจฺจเยน ปจฺจโย. (2)}

\prose{คนฺถ\-สมฺปยุตฺโต จ คนฺถ\-วิปฺปยุตฺโต จ ธมฺมา คนฺถ\-สมฺปยุตฺตสฺส จ คนฺถ\-วิปฺปยุตฺตสฺส จ ธมฺมสฺส เหตุปจฺจเยน ปจฺจโย. ทิฏฺฐิคต\-วิปฺปยุตฺต\-โลภ\-สหคโต เหตุ จ โลโภ จ สมฺปยุตฺตกานํ ขนฺธานํ จิตฺตสมุฏฺฐานานญฺจ รูปานํ เหตุปจฺจเยน ปจฺจโย. โทมนสฺส\-สหคโต เหตุ จ ปฏิฆญฺจ สมฺปยุตฺตกานํ ขนฺธานํ จิตฺตสมุฏฺฐานานญฺจ รูปานํ เหตุปจฺจเยน ปจฺจโย. (3)}

\proseitem{72}{คนฺถ\-สมฺปยุตฺโต ธมฺโม คนฺถ\-สมฺปยุตฺตสฺส ธมฺมสฺส อารมฺมณ\-ปจฺจเยน ปจฺจโย. คนฺถ\-สมฺปยุตฺเต ขนฺเธ อารพฺภ คนฺถ\-สมฺปยุตฺตกา ขนฺธา \csromanfolio{278}อุปฺปชฺชนฺติ. (ตีสุปิ มูลา ปุจฺฉิตพฺพา.) คนฺถ\-สมฺปยุตฺเต ขนฺเธ อารพฺภ คนฺถ\-วิปฺปยุตฺตา ขนฺธา อุปฺปชฺชนฺติ. คนฺถ\-สมฺปยุตฺเต ขนฺเธ อารพฺภ ทิฏฺฐิคต\-วิปฺปยุตฺต\-โลภ\-สหคตา ขนฺธา จ โลโภ จ อุปฺปชฺชนฺติ. โทมนสฺส\-สหคตา ขนฺธา จ ปฏิฆญฺจ อุปฺปชฺชนฺติ. (3)}

\proseitem{73}{คนฺถ\-วิปฺปยุตฺโต ธมฺโม คนฺถ\-วิปฺปยุตฺตสฺส ธมฺมสฺส อารมฺมณ\-ปจฺจเยน ปจฺจโย. ทานํ ทตฺวา, สีลํ ฯเปฯ อุโปสถกมฺมํ กตฺวา ตํ ปจฺจเวกฺขติ, ปุพฺเพ สุจิณฺณานิ ปจฺจเวกฺขติ, ฌานา วุฏฺฐหิตฺวา ฌานํ ปจฺจเวกฺขติ. อริยา มคฺคา วุฏฺฐหิตฺวา มคฺคํ ปจฺจเวกฺขนฺติ, ผลํ ปจฺจเวกฺขนฺติ, นิพฺพานํ ปจฺจเวกฺขนฺติ. นิพฺพานํ โคตฺรภุสฺส, โวทานสฺส, มคฺคสฺส, ผลสฺส, อาวชฺชนาย อารมฺมณ\-ปจฺจเยน ปจฺจโย. อริยา คนฺถ\-วิปฺปยุตฺเต ปหีเน กิเลเส ปจฺจเวกฺขนฺติ, วิกฺขมฺภิเต กิเลเส ปจฺจเวกฺขนฺติ. ปุพฺเพ สมุทาจิณฺเณ กิเลเส ชานนฺติ. จกฺขุํ ฯเปฯ วตฺถุํ คนฺถ\-วิปฺปยุตฺเต ขนฺเธ จ โลภญฺจ ปฏิฆญฺจ อนิจฺจโต ฯเปฯ วิปสฺสติ, อสฺสาเทติ อภินนฺทติ, ตํ อารพฺภ คนฺถ\-วิปฺปยุตฺโต ราโค อุปฺปชฺชติ. วิจิกิจฺฉา ฯเปฯ อุทฺธจฺจํ ฯเปฯ โทมนสฺสํ อุปฺปชฺชติ. ทิพฺเพน จกฺขุนา รูปํ ปสฺสติ. ทิพฺพาย โสตธาตุยา สทฺทํ สุณาติ. เจโตปริย\-ญาเณน คนฺถ\-วิปฺปยุตฺต\-จิตฺต\-สมงฺคิสฺส จิตฺตํ ชานาติ. อากาสา\-นญฺจา\-ยตนํ วิญฺญาณญฺจา\-ยตนสฺส ฯเปฯ. อากิญฺจญฺญา\-ยตนํ เนว\-สญฺญา\-นา\-สญฺญา\-ยตนสฺส ฯเปฯ. รูปายตนํ จกฺขุ\-วิญฺญาณสฺส ฯเปฯ. โผฏฺฐพฺพา\-ยตนํ กายวิญฺญาณสฺส ฯเปฯ. คนฺถ\-วิปฺปยุตฺตา ขนฺธา อิทฺธิวิธ\-ญาณสฺส, เจโตปริย\-ญาณสฺส, ปุพฺเพ\-นิวาสา\-นุสฺสติ\-ญาณสฺส, ยถากมฺมูปคญาณสฺส, อนาคตํสญาณสฺส, อาวชฺชนาย อารมฺมณ\-ปจฺจเยน ปจฺจโย. (1)}

\prose{คนฺถ\-วิปฺปยุตฺโต ธมฺโม คนฺถ\-สมฺปยุตฺตสฺส ธมฺมสฺส อารมฺมณ\-ปจฺจเยน ปจฺจโย. ทานํ ทตฺวา, สีลํ ฯเปฯ อุโปสถกมฺมํ ฯเปฯ. ปุพฺเพ ฯเปฯ. ฌานา ฯเปฯ ตํ อสฺสาเทติ อภินนฺทติ, ตํ อารพฺภ คนฺถ\-สมฺปยุตฺโต ราโค อุปฺปชฺชติ. ทิฏฺฐิ อุปฺปชฺชติ. โทมนสฺสํ อุปฺปชฺชติ. (สํขิตฺตํ.) (2)}

\prose{คนฺถ\-วิปฺปยุตฺโต ธมฺโม คนฺถ\-สมฺปยุตฺตสฺส จ คนฺถ\-วิปฺปยุตฺตสฺส จ ธมฺมสฺส อารมฺมณ\-ปจฺจเยน ปจฺจโย. จกฺขุํ ฯเปฯ วตฺถุํ คนฺถ\-วิปฺปยุตฺเต ขนฺเธ จ โลภญฺจ ปฏิฆญฺจ อารพฺภ ทิฏฺฐิคต\-วิปฺปยุตฺต\-โลภ\-สหคตา ขนฺธา จ โลโภ จ โทมนสฺส\-สหคตา ขนฺธา จ ปฏิฆญฺจ อุปฺปชฺชนฺติ. (3)}

\proseitem{74}{\csromanfolio{279}คนฺถ\-สมฺปยุตฺโต จ คนฺถ\-วิปฺปยุตฺโต จ ธมฺมา คนฺถ\-สมฺปยุตฺตสฺส ธมฺมสฺส อารมฺมณ\-ปจฺจเยน ปจฺจโย. ทิฏฺฐิคต\-วิปฺปยุตฺต\-โลภ\-สหคเต ขนฺเธ จ โลภญฺจ โทมนสฺส\-สหคเต ขนฺเธ จ ปฏิฆญฺจ อารพฺภ คนฺถ\-สมฺปยุตฺตกา ขนฺธา อุปฺปชฺชนฺติ. (มูลํ ปุจฺฉิตพฺพํ.) ทิฏฺฐิคต\-วิปฺปยุตฺต\-โลภ\-สหคเต ขนฺเธ จ โลภญฺจ โทมนสฺส\-สหคเต ขนฺเธ จ ปฏิฆญฺจ อารพฺภ คนฺถ\-วิปฺปยุตฺตา ขนฺธา อุปฺปชฺชนฺติ. ทิฏฺฐิคต\-วิปฺปยุตฺต\-โลภ\-สหคเต ขนฺเธ จ โลภญฺจ โทมนสฺส\-สหคเต ขนฺเธ จ ปฏิฆญฺจ อารพฺภ ทิฏฺฐิคต\-วิปฺปยุตฺต\-โลภ\-สหคตา ขนฺธา จ โลโภ จ โทมนสฺส\-สหคตา ขนฺธา จ ปฏิฆญฺจ อุปฺปชฺชนฺติ. (3)}

\proseitem{75}{คนฺถ\-สมฺปยุตฺโต ธมฺโม คนฺถ\-สมฺปยุตฺตสฺส ธมฺมสฺส อธิปติ\-ปจฺจเยน ปจฺจโย. อารมฺมณา\-ธิปติ, สหชาตา\-ธิปติ.  อารมฺมณา\-ธิปติ\csromandash{}คนฺถ\-สมฺปยุตฺเต ขนฺเธ ครุํ กตฺวา คนฺถ\-สมฺปยุตฺตกา ขนฺธา อุปฺปชฺชนฺติ.  สหชาตา\-ธิปติ\csromandash{} คนฺถ\-สมฺปยุตฺตา\-ธิปติ สมฺปยุตฺตกานํ ขนฺธานํ อธิปติ\-ปจฺจเยน ปจฺจโย. (1)}

\prose{คนฺถ\-สมฺปยุตฺโต ธมฺโม คนฺถ\-วิปฺปยุตฺตสฺส ธมฺมสฺส อธิปติ\-ปจฺจเยน ปจฺจโย. อารมฺมณา\-ธิปติ, สหชาตา\-ธิปติ.  อารมฺมณา\-ธิปติ\csromandash{}คนฺถ\-สมฺปยุตฺเต ขนฺเธ ครุํ กตฺวา ทิฏฺฐิคต\-วิปฺปยุตฺโต โลโภ อุปฺปชฺชติ.  สหชาตา\-ธิปติ\csromandash{} คนฺถ\-สมฺปยุตฺตา\-ธิปติ จิตฺต\-สมุฏฺฐานานํ รูปานํ อธิปติ\-ปจฺจเยน ปจฺจโย. ทิฏฺฐิคต\-วิปฺปยุตฺต\-โลภ\-สหคตา\-ธิปติ โลภสฺส จ จิตฺตสมุฏฺฐานานญฺจ รูปานํ อธิปติ\-ปจฺจเยน ปจฺจโย. โทมนสฺส\-สหคตาธิปติ ปฏิฆสฺส จ จิตฺตสมุฏฺฐานานญฺจ รูปานํ อธิปติ\-ปจฺจเยน ปจฺจโย. (2)}

\prose{คนฺถ\-สมฺปยุตฺโต ธมฺโม คนฺถ\-สมฺปยุตฺตสฺส จ คนฺถ\-วิปฺปยุตฺตสฺส จ ธมฺมสฺส อธิปติ\-ปจฺจเยน ปจฺจโย. อารมฺมณา\-ธิปติ, สหชาตา\-ธิปติ.  อารมฺมณา\-ธิปติ\csromandash{}คนฺถ\-สมฺปยุตฺเต ขนฺเธ ครุํ กตฺวา ทิฏฺฐิคต\-วิปฺปยุตฺต\-โลภ\-สหคตา ขนฺธา จ โลโภ จ อุปฺปชฺชนฺติ.  สหชาตา\-ธิปติ\csromandash{}ทิฏฺฐิคต\-วิปฺปยุตฺต\-โลภ\-สหคตา\-ธิปติ สมฺปยุตฺตกานํ ขนฺธานํ โลภสฺส จ จิตฺตสมุฏฺฐานานญฺจ รูปานํ อธิปติ\-ปจฺจเยน ปจฺจโย. โทมนสฺส\-สหคตาธิปติ สมฺปยุตฺตกานํ ขนฺธานํ ปฏิฆสฺส จ จิตฺตสมุฏฺฐานานญฺจ รูปานํ อธิปติ\-ปจฺจเยน ปจฺจโย. (3)}

\proseitem{76}{\csromanfolio{280}คนฺถ\-วิปฺปยุตฺโต ธมฺโม คนฺถ\-วิปฺปยุตฺตสฺส ธมฺมสฺส อธิปติ\-ปจฺจเยน ปจฺจโย. อารมฺมณา\-ธิปติ, สหชาตา\-ธิปติ.  อารมฺมณา\-ธิปติ\csromandash{}ทานํ ฯเปฯ สีลํ ฯเปฯ อุโปสถกมฺมํ กตฺวา ตํ ครุํ กตฺวา ปจฺจเวกฺขติ, ปุพฺเพ สุจิณฺณานิ ครุํ กตฺวา ปจฺจเวกฺขติ. ฌานา วุฏฺฐหิตฺวา ฌานํ ครุํ กตฺวา ปจฺจเวกฺขติ. อริยา มคฺคา วุฏฺฐหิตฺวา มคฺคํ ครุํ กตฺวา ปจฺจเวกฺขนฺติ, ผลา วุฏฺฐหิตฺวา ผลํ ครุํ กตฺวา ปจฺจเวกฺขนฺติ, นิพฺพานํ ครุํ กตฺวา ปจฺจเวกฺขนฺติ. นิพฺพานํ โคตฺรภุสฺส, โวทานสฺส, มคฺคสฺส, ผลสฺส อธิปติ\-ปจฺจเยน ปจฺจโย. จกฺขุํ ฯเปฯ วตฺถุํ คนฺถ\-วิปฺปยุตฺเต ขนฺเธ จ โลภญฺจ ครุํ กตฺวา อสฺสาเทติ อภินนฺทติ, ตํ ครุํ กตฺวา คนฺถ\-วิปฺปยุตฺโต ราโค อุปฺปชฺชติ.  สหชาตา\-ธิปติ\csromandash{}คนฺถ\-วิปฺปยุตฺตา\-ธิปติ สมฺปยุตฺตกานํ ขนฺธานํ จิตฺตสมุฏฺฐานานญฺจ รูปานํ อธิปติ\-ปจฺจเยน ปจฺจโย. (1)}

\prose{คนฺถ\-วิปฺปยุตฺโต ธมฺโม คนฺถ\-สมฺปยุตฺตสฺส ธมฺมสฺส อธิปติ\-ปจฺจเยน ปจฺจโย. อารมฺมณา\-ธิปติ\csromandash{}ทานํ ฯเปฯ สีลํ ฯเปฯ อุโปสถกมฺมํ ฯเปฯ. ปุพฺเพ สุจิณฺณานิ ฯเปฯ. ฌานา วุฏฺฐหิตฺวา ฌานํ ฯเปฯ. จกฺขุํ ฯเปฯ วตฺถุํ คนฺถ\-วิปฺปยุตฺเต ขนฺเธ จ โลภญฺจ ครุํ กตฺวา อสฺสาเทติ อภินนฺทติ, ตํ ครุํ กตฺวา คนฺถ\-สมฺปยุตฺโต ราโค อุปฺปชฺชติ. ทิฏฺฐิ อุปฺปชฺชติ. (2)}

\prose{คนฺถ\-วิปฺปยุตฺโต ธมฺโม คนฺถ\-สมฺปยุตฺตสฺส จ คนฺถ\-วิปฺปยุตฺตสฺส จ ธมฺมสฺส อธิปติ\-ปจฺจเยน ปจฺจโย. อารมฺมณา\-ธิปติ\csromandash{}จกฺขุํ ฯเปฯ วตฺถุํ คนฺถ\-วิปฺปยุตฺเต ขนฺเธ จ โลภญฺจ ครุํ กตฺวา ทิฏฺฐิคต\-วิปฺปยุตฺต\-โลภ\-สหคตา ขนฺธา จ โลโภ จ อุปฺปชฺชนฺติ. (3)}

\proseitem{77}{คนฺถ\-สมฺปยุตฺโต จ คนฺถ\-วิปฺปยุตฺโต จ ธมฺมา คนฺถ\-สมฺปยุตฺตสฺส ธมฺมสฺส อธิปติ\-ปจฺจเยน ปจฺจโย. อารมฺมณา\-ธิปติ\csromandash{}ทิฏฺฐิคต\-วิปฺปยุตฺต\-โลภ\-สหคเต ขนฺเธ จ โลภญฺจ ครุํ กตฺวา คนฺถ\-สมฺปยุตฺตกา ขนฺธา อุปฺปชฺชนฺติ. (มูลา ปุจฺฉิตพฺพา.) ทิฏฺฐิคต\-วิปฺปยุตฺต\-โลภ\-สหคเต ขนฺเธ จ โลภญฺจ ครุํ กตฺวา ทิฏฺฐิคต\-วิปฺปยุตฺโต โลโภ อุปฺปชฺชติ. (มูลา ปุจฺฉิตพฺพา.) ทิฏฺฐิคต\-วิปฺปยุตฺต\-โลภ\-สหคเต ขนฺเธ จ โลภญฺจ ครุํ กตฺวา ทิฏฺฐิคต\-วิปฺปยุตฺต\-โลภ\-สหคตา ขนฺธา จ โลโภ จ อุปฺปชฺชนฺติ. (3)}

\csromanheader{2}{28. คนฺถ\-สมฺปยุตฺตทุก อนนฺตร}
\proseitem{78}{\csromanfolio{281}คนฺถ\-สมฺปยุตฺโต ธมฺโม คนฺถ\-สมฺปยุตฺตสฺส ธมฺมสฺส อนนฺตร\-ปจฺจเยน ปจฺจโย. ปุริมา ปุริมา คนฺถ\-สมฺปยุตฺตา ขนฺธา ปจฺฉิมานํ ปจฺฉิมานํ คนฺถ\-สมฺปยุตฺตกานํ ขนฺธานํ อนนฺตร\-ปจฺจเยน ปจฺจโย. (มูลา ปุจฺฉิตพฺพา.) ปุริมา ปุริมา ทิฏฺฐิคต\-วิปฺปยุตฺต\-โลภ\-สหคตา ขนฺธา ปจฺฉิมสฺส ปจฺฉิมสฺส ทิฏฺฐิคกวิปฺปยุตฺตสฺส โลภสฺส อนนฺตร\-ปจฺจเยน ปจฺจโย. ปุริมา ปุริมา โทมนสฺส\-สหคตา ขนฺธา ปจฺฉิมสฺส ปจฺฉิมสฺส ปฏิฆสฺส อนนฺตร\-ปจฺจเยน ปจฺจโย. คนฺถ\-สมฺปยุตฺตา ขนฺธา วุฏฺฐานสฺส อนนฺตร\-ปจฺจเยน ปจฺจโย. (มูลา ปุจฺฉิตพฺพา.) ปุริมา ปุริมา ทิฏฺฐิคต\-วิปฺปยุตฺต\-โลภ\-สหคตา ขนฺธา ปจฺฉิมานํ ปจฺฉิมานํ ทิฏฺฐิคต\-วิปฺปยุตฺต\-โลภ\-สหคตานํ ขนฺธานํ โลภสฺส จ อนนฺตร\-ปจฺจเยน ปจฺจโย. ปุริมา ปุริมา โทมนสฺส\-สหคตา ขนฺธา ปจฺฉิมานํ ปจฺฉิมานํ โทมนสฺส\-สหคตานํ ขนฺธานํ ปฏิฆสฺส จ อนนฺตร\-ปจฺจเยน ปจฺจโย. (3)}

\proseitem{79}{คนฺถ\-วิปฺปยุตฺโต ธมฺโม คนฺถ\-วิปฺปยุตฺตสฺส ธมฺมสฺส อนนฺตร\-ปจฺจเยน ปจฺจโย. ปุริโม ปุริโม ทิฏฺฐิคต\-วิปฺปยุตฺโต โลโภ ปจฺฉิมสฺส ปจฺฉิมสฺส ทิฏฺฐิคต\-วิปฺปยุตฺตสฺส โลภสฺส อนนฺตร\-ปจฺจเยน ปจฺจโย. ปุริมํ ปุริมํ ปฏิฆํ ปจฺฉิมสฺส ปจฺฉิมสฺส ปฏิฆสฺส อนนฺตร\-ปจฺจเยน ปจฺจโย. ปุริมา ปุริมา คนฺถ\-วิปฺปยุตฺตา ขนฺธา ปจฺฉิมานํ ปจฺฉิมานํ คนฺถ\-วิปฺปยุตฺตานํ ขนฺธานํ อนนฺตร\-ปจฺจเยน ปจฺจโย. อนุโลมํ โคตฺรภุสฺส ฯเปฯ ผลสมาปตฺติยา อนนฺตร\-ปจฺจเยน ปจฺจโย. (1)}

\prose{คนฺถ\-วิปฺปยุตฺโต ธมฺโม คนฺถ\-สมฺปยุตฺตสฺส ธมฺมสฺส อนนฺตร\-ปจฺจเยน ปจฺจโย. ปุริโม ปุริโม ทิฏฺฐิคต\-วิปฺปยุตฺโต โลโภ ปจฺฉิมานํ ปจฺฉิมานํ ทิฏฺฐิคต\-วิปฺปยุตฺต\-โลภ\-สหคตานํ ขนฺธานํ อนนฺตร\-ปจฺจเยน ปจฺจโย. ปุริมํ ปุริมํ ปฏิฆํ ปจฺฉิมานํ ปจฺฉิมานํ โทมนสฺส\-สหคตานํ ขนฺธานํ อนนฺตร\-ปจฺจเยน ปจฺจโย. อาวชฺชนา คนฺถ\-สมฺปยุตฺตกานํ ขนฺธานํ อนนฺตร\-ปจฺจเยน ปจฺจโย. (2)}

\prose{คนฺถ\-วิปฺปยุตฺโต ธมฺโม คนฺถ\-สมฺปยุตฺตสฺส จ คนฺถ\-วิปฺปยุตฺตสฺส จ ธมฺมสฺส อนนฺตร\-ปจฺจเยน ปจฺจโย. ปุริโม ปุริโม ทิฏฺฐิคต\-วิปฺปยุตฺโต โลโภ ปจฺฉิมานํ ปจฺฉิมานํ ทิฏฺฐิคต\-วิปฺปยุตฺต\-โลภ\-สหคตานํ ขนฺธานํ โลภสฺส จ \csromanfolio{282}อนนฺตร\-ปจฺจเยน ปจฺจโย. ปุริมํ ปุริมํ ปฏิฆํ ปจฺฉิมานํ ปจฺฉิมานํ โทมนสฺส\-สหคตานํ ขนฺธานํ ปฏิฆสฺส จ อนนฺตร\-ปจฺจเยน ปจฺจโย. อาวชฺชนา ทิฏฺฐิคต\-วิปฺปยุตฺต\-โลภ\-สหคตานํ ขนฺธานํ โลภสฺส จ, โทมนสฺส\-สหคตานํ ขนฺธานํ ปฏิฆสฺส จ อนนฺตร\-ปจฺจเยน ปจฺจโย. (3)}

\proseitem{80}{คนฺถ\-สมฺปยุตฺโต จ คนฺถ\-วิปฺปยุตฺโต จ ธมฺมา คนฺถ\-สมฺปยุตฺตสฺส ธมฺมสฺส อนนฺตร\-ปจฺจเยน ปจฺจโย. ปุริมา ปุริมา ทิฏฺฐิคต\-วิปฺปยุตฺต\-โลภ\-สหคตา ขนฺธา จ โลโภ จ ปจฺฉิมานํ ปจฺฉิมานํ ทิฏฺฐิคต\-วิปฺปยุตฺต\-โลภ\-สหคตานํ ขนฺธานํ อนนฺตร\-ปจฺจเยน ปจฺจโย. ปุริมา ปุริมา โทมนสฺส\-สหคตา ขนฺธา จ ปฏิฆญฺจ ปจฺฉิมานํ ปจฺฉิมานํ โทมนสฺส\-สหคตานํ ขนฺธานํ อนนฺตร\-ปจฺจเยน ปจฺจโย. (มูลา ปุจฺฉิตพฺพา.) ปุริมา ปุริมา ทิฏฺฐิคต\-วิปฺปยุตฺต\-โลภ\-สหคตา ขนฺธา จ โลโภ จ ปจฺฉิมสฺส ปจฺฉิมสฺส ทิฏฺฐิคต\-วิปฺปยุตฺตสฺส โลภสฺส อนนฺตร\-ปจฺจเยน ปจฺจโย. ปุริมา ปุริมา โทมนสฺส\-สหคตา ขนฺธา จ ปฏิฆญฺจ ปจฺฉิมสฺส ปจฺฉิมสฺส ปฏิฆสฺส อนนฺตร\-ปจฺจเยน ปจฺจโย. ทิฏฺฐิคต\-วิปฺปยุตฺต\-โลภ\-สหคตา ขนฺธา จ โลโภ จ, โทมนสฺส\-สหคตา ขนฺธา จ ปฏิฆญฺจ วุฏฺฐานสฺส อนนฺตร\-ปจฺจเยน ปจฺจโย. (มูลา ปุจฺฉิตพฺพา.) ปุริมา ปุริมา ทิฏฺฐิคต\-วิปฺปยุตฺต\-โลภ\-สหคตา ขนฺธา จ โลโภ จ ปจฺฉิมานํ ปจฺฉิมานํ ทิฏฺฐิคต\-วิปฺปยุตฺต\-โลภ\-สหคตานํ ขนฺธานํ โลภสฺส จ อนนฺตร\-ปจฺจเยน ปจฺจโย. ปุริมา ปุริมา โทมนสฺส\-สหคตา ขนฺธา จ ปฏิฆญฺจ ปจฺฉิมานํ ปจฺฉิมานํ โทมนสฺส\-สหคตานํ ขนฺธานํ ปฏิฆสฺส จ อนนฺตร\-ปจฺจเยน ปจฺจโย. (3)}

\csromanheader{2}{\textbf{สมนนฺตราทิ}}
\proseitem{81}{คนฺถ\-สมฺปยุตฺโต ธมฺโม คนฺถ\-สมฺปยุตฺตสฺส ธมฺมสฺส สมนนฺตร\-ปจฺจเยน ปจฺจโย. สหชาต\-ปจฺจเยน ปจฺจโย. อญฺญมญฺญ\-ปจฺจเยน ปจฺจโย. นิสฺสย\-ปจฺจเยน ปจฺจโย.}

\csromanheader{2}{\textbf{อุปนิสฺสย}}
\proseitem{82}{คนฺถ\-สมฺปยุตฺโต ธมฺโม คนฺถ\-สมฺปยุตฺตสฺส ธมฺมสฺส อุปนิสฺสย\-ปจฺจเยน ปจฺจโย. อารมฺมณู\-ปนิสฺสโย, อนนฺตรู\-ปนิสฺสโย, ปกตู\-ปนิสฺสโย ฯเปฯ. ปกตู\-ปนิสฺสโย\csromandash{}คนฺถ\-สมฺปยุตฺตา ขนฺธา คนฺถ\-สมฺปยุตฺตกานํ ขนฺธานํ อุปนิสฺสย\-ปจฺจเยน ปจฺจโย. (มูลํ ปุจฺฉิตพฺพํ. \csromanfolio{283}ตีณิปิ อุปนิสฺสยา.) คนฺถ\-สมฺปยุตฺตา ขนฺธา คนฺถ\-วิปฺปยุตฺตานํ ขนฺธานํ อุปนิสฺสย\-ปจฺจเยน ปจฺจโย. (มูลํ. ตีณิปิ อุปนิสฺสยา.) คนฺถ\-สมฺปยุตฺตา ขนฺธา ทิฏฺฐิคต\-วิปฺปยุตฺต\-โลภ\-สหคตานํ ขนฺธานํ โลภสฺส จ อุปนิสฺสย\-ปจฺจเยน ปจฺจโย. โทมนสฺส\-สหคตานํ ขนฺธานํ ปฏิฆสฺส จ อุปนิสฺสย\-ปจฺจเยน ปจฺจโย. (3)}

\proseitem{83}{คนฺถ\-วิปฺปยุตฺโต ธมฺโม คนฺถ\-วิปฺปยุตฺตสฺส ธมฺมสฺส อุปนิสฺสย\-ปจฺจเยน ปจฺจโย. อารมฺมณู\-ปนิสฺสโย, อนนฺตรู\-ปนิสฺสโย, ปกตู\-ปนิสฺสโย ฯเปฯ. ปกตู\-ปนิสฺสโย\csromandash{}สทฺธํ อุปนิสฺสาย ทานํ เทติ ฯเปฯ สมาปตฺติํ อุปฺปาเทติ. มานํ ชปฺเปติ. สีลํ ฯเปฯ. ปญฺญํ. ราคํ. โทสํ. โมหํ. มานํ. ปตฺถนํ. กายิกํ สุขํ. กายิกํ ทุกฺขํ. อุตุํ. โภชนํ. เสนาสนํ อุปนิสฺสาย ทานํ เทติ ฯเปฯ สมาปตฺติํ อุปฺปาเทติ. ปาณํ หนติ ฯเปฯ สํฆํ ภินฺทติ. สทฺธา ฯเปฯ. ปญฺญา. ราโค ฯเปฯ. ปตฺถนา ฯเปฯ. เสนาสนํ สทฺธาย ฯเปฯ ปญฺญาย. ราคสฺส. โทสสฺส. โมหสฺส. มานสฺส. ปตฺถนาย ฯเปฯ ผลสมาปตฺติยา อุปนิสฺสย\-ปจฺจเยน ปจฺจโย. (1)}

\prose{คนฺถ\-วิปฺปยุตฺโต ธมฺโม คนฺถ\-สมฺปยุตฺตสฺส ธมฺมสฺส อุปนิสฺสย\-ปจฺจเยน ปจฺจโย. อารมฺมณู\-ปนิสฺสโย, อนนฺตรู\-ปนิสฺสโย, ปกตู\-ปนิสฺสโย ฯเปฯ. ปกตู\-ปนิสฺสโย\csromandash{}สทฺธํ อุปนิสฺสาย มานํ ชปฺเปติ. ทิฏฺฐิํ คณฺหาติ. สีลํ ฯเปฯ. ปญฺญํ. ราคํ ฯเปฯ. มานํ. ปตฺถนํ ฯเปฯ. เสนาสนํ อุปนิสฺสาย ปาณํ หนติ ฯเปฯ สํฆํ ภินฺทติ. สทฺธา ฯเปฯ. เสนาสนํ ราคสฺส. โทสสฺส. โมหสฺส. มานสฺส. ทิฏฺฐิยา. ปตฺถนาย อุปนิสฺสย\-ปจฺจเยน ปจฺจโย. (2)}

\prose{คนฺถ\-วิปฺปยุตฺโต ธมฺโม คนฺถ\-สมฺปยุตฺตสฺส จ คนฺถ\-วิปฺปยุตฺตสฺส จ ธมฺมสฺส อุปนิสฺสย\-ปจฺจเยน ปจฺจโย. (ตีณิ อุปนิสฺสยา.) สทฺธํ อุปนิสฺสาย มานํ ชปฺเปติ. สีลํ ฯเปฯ. ปญฺญํ. ราคํ. โทสํ. โมหํ. มานํ. ปตฺถนํ. กายิกํ สุขํ ฯเปฯ. เสนาสนํ อุปนิสฺสาย ปาณํ หนติ. สํฆํ ภินฺทติ. สทฺธา ฯเปฯ. ปญฺญา. ราโค. โทโส. โมโห. มาโน. ปตฺถนา ฯเปฯ. เสนาสนํ ทิฏฺฐิคต\-วิปฺปยุตฺต\-โลภ\-สหคตานํ ขนฺธานํ โลภสฺส จ, โทมนสฺส\-สหคตานํ ขนฺธานํ ปฏิฆสฺส จ อุปนิสฺสย\-ปจฺจเยน ปจฺจโย. (3)}

\proseitem{84}{คนฺถ\-สมฺปยุตฺโต จ คนฺถ\-วิปฺปยุตฺโต จ ธมฺมา คนฺถ สมฺปยุตฺตสฺส ธมฺมสฺส อุปนิสฺสย\-ปจฺจเยน ปจฺจโย. อารมฺมณู\-ปนิสฺสโย, อนนฺตรู\-ปนิสฺสโย, \csromanfolio{284}ปกตู\-ปนิสฺสโย ฯเปฯ. ปกตู\-ปนิสฺสโย\csromandash{}ทิฏฺฐิคต\-วิปฺปยุตฺต\-โลภ\-สหคตา ขนฺธา จ โลโภ จ, โทมนสฺส\-สหคตา ขนฺธา จ ปฏิฆญฺจ คนฺถ\-สมฺปยุตฺตกานํ ขนฺธานํ อุปนิสฺสย\-ปจฺจเยน ปจฺจโย. (มูลํ กาตพฺพํ.) ทิฏฺฐิคต\-วิปฺปยุตฺต\-โลภ\-สหคตา ขนฺธา จ โลโภ จ, โทมนสฺส\-สหคตา ขนฺธา จ ปฏิฆญฺจ คนฺถ\-วิปฺปยุตฺตานํ ขนฺธานํ ทิฏฺฐิคต\-วิปฺปยุตฺต\-โลภสฺส จ ปฏิฆสฺส จ อุปนิสฺสย\-ปจฺจเยน ปจฺจโย. (มูลํ ปุจฺฉิตพฺพํ.) ทิฏฺฐิคต\-วิปฺปยุตฺต\-โลภ\-สหคตา ขนฺธา จ โลโภ จ, โทมนสฺส\-สหคตา ขนฺธา จ ปฏิฆญฺจ ทิฏฺฐิคต\-วิปฺปยุตฺต\-โลภ\-สหคตานํ ขนฺธานํ โลภสฺส จ, โทมนสฺส\-สหคตานํ ขนฺธานํ ปฏิฆสฺส จ อุปนิสฺสย\-ปจฺจเยน ปจฺจโย. (3)}

\proseitem{85}{คนฺถ\-วิปฺปยุตฺโต ธมฺโม คนฺถ\-วิปฺปยุตฺตสฺส ธมฺมสฺส ปุเรชาต\-ปจฺจเยน ปจฺจโย. อารมฺมณ\-ปุเรชาตํ, วตฺถุ\-ปุเรชาตํ.  อารมฺมณ\-ปุเรชาตํ\csromandash{}จกฺขุํ ฯเปฯ วตฺถุํ อนิจฺจโต วิปสฺสติ, อสฺสาเทติ อภินนฺทติ, ตํ อารพฺภ คนฺถ\-วิปฺปยุตฺโต ราโค อุปฺปชฺชติ. วิจิกิจฺฉา ฯเปฯ อุทฺธจฺจํ ฯเปฯ โทมนสฺสํ อุปฺปชฺชติ. ทิพฺเพน จกฺขุนา รูปํ ปสฺสติ. ทิพฺพาย โสตธาตุยา สทฺทํ สุณาติ. รูปายตนํ จกฺขุ\-วิญฺญาณสฺส ฯเปฯ. โผฏฺฐพฺพา\-ยตนํ กายวิญฺญาณสฺส ฯเปฯ. วตฺถุ\-ปุเรชาตํ\csromandash{}จกฺขายตนํ จกฺขุ\-วิญฺญาณสฺส ฯเปฯ กายายตนํ กายวิญฺญาณสฺส ฯเปฯ. วตฺถุ คนฺถ\-วิปฺปยุตฺตานํ ขนฺธานํ ทิฏฺฐิคต\-วิปฺปยุตฺตสฺส โลภสฺส จ ปฏิฆสฺส จ ปุเรชาต\-ปจฺจเยน ปจฺจโย. (1)}

\prose{คนฺถ\-วิปฺปยุตฺโต ธมฺโม คนฺถ\-สมฺปยุตฺตสฺส ธมฺมสฺส ปุเรชาต\-ปจฺจเยน ปจฺจโย. อารมฺมณ\-ปุเรชาตํ, วตฺถุ\-ปุเรชาตํ.  อารมฺมณ\-ปุเรชาตํ\csromandash{}จกฺขุํ ฯเปฯ วตฺถุํ อสฺสาเทติ อภินนฺทติ, ตํ อารพฺภ คนฺถ\-สมฺปยุตฺโต ราโค อุปฺปชฺชติ. ทิฏฺฐิ ฯเปฯ โทมนสฺสํ อุปฺปชฺชติ.  วตฺถุ\-ปุเรชาตํ\csromandash{}วตฺถุ คนฺถ\-สมฺปยุตฺตกานํ ขนฺธานํ ปุเรชาต\-ปจฺจเยน ปจฺจโย. (2)}

\prose{คนฺถ\-วิปฺปยุตฺโต ธมฺโม คนฺถ\-สมฺปยุตฺตสฺส จ คนฺถ\-วิปฺปยุตฺตสฺส จ ธมฺมสฺส ปุเรชาต\-ปจฺจเยน ปจฺจโย. อารมฺมณ\-ปุเรชาตํ, วตฺถุ\-ปุเรชาตํ.  อารมฺมณ\-ปุเรชาตํ\csromandash{}จกฺขุํ ฯเปฯ วตฺถุํ อารพฺภ ทิฏฺฐิคต\-วิปฺปยุตฺต\-โลภ\-สหคตา ขนฺธา จ โลโภ จ, โทมนสฺส\-สหคตา ขนฺธา จ ปฏิฆญฺจ อุปฺปชฺชนฺติ.  วตฺถุ\-ปุเรชาตํ\csromandash{}วตฺถุ ทิฏฺฐิคต\-วิปฺปยุตฺต\-โลภ\-สหคตานํ \csromanfolio{285}ขนฺธานํ โลภสฺส จ, โทมนสฺส\-สหคตานํ ขนฺธานํ ปฏิฆสฺส จ ปุเรชาต\-ปจฺจเยน ปจฺจโย. (3)}

\csromanheader{2}{\textbf{ปจฺฉาชาต}}
\proseitem{86}{คนฺถ\-สมฺปยุตฺโต ธมฺโม คนฺถ\-วิปฺปยุตฺตสฺส ธมฺมสฺส ปจฺฉาชาต\-ปจฺจเยน ปจฺจโย. เอกํ.}

\prose{คนฺถ\-วิปฺปยุตฺโต ธมฺโม คนฺถ\-วิปฺปยุตฺตสฺส ธมฺมสฺส ปจฺฉาชาต\-ปจฺจเยน ปจฺจโย. เอกํ.}

\prose{คนฺถ\-สมฺปยุตฺโต จ คนฺถ\-วิปฺปยุตฺโต จ ธมฺมา คนฺถ\-วิปฺปยุตฺตสฺส ธมฺมสฺส ปจฺฉาชาต\-ปจฺจเยน ปจฺจโย. ปจฺฉาชาตา ทิฏฺฐิคต\-วิปฺปยุตฺต\-โลภ\-สหคตา ขนฺธา จ โลโภ จ, โทมนสฺส\-สหคตา ขนฺธา จ ปฏิฆญฺจ ปุเรชาตสฺส อิมสฺส กายสฺส ปจฺฉาชาต\-ปจฺจเยน ปจฺจโย. (1)}

\csromanheader{2}{\textbf{อาเสวน}}
\proseitem{87}{คนฺถ\-สมฺปยุตฺโต ธมฺโม คนฺถ\-สมฺปยุตฺตสฺส ธมฺมสฺส อาเสวน\-ปจฺจเยน ปจฺจโย. (อนนฺตร\-สทิสํ. อาวชฺชนาปิ วุฏฺฐานมฺปิ นตฺถิ.)}

\proseitem{88}{คนฺถ\-สมฺปยุตฺโต ธมฺโม คนฺถ\-สมฺปยุตฺตสฺส ธมฺมสฺส กมฺมปจฺจเยน ปจฺจโย. คนฺถ\-สมฺปยุตฺตา เจตนา สมฺปยุตฺตกานํ ขนฺธานํ กมฺมปจฺจเยน ปจฺจโย. (1)}

\prose{คนฺถ\-สมฺปยุตฺโต ธมฺโม คนฺถ\-วิปฺปยุตฺตสฺส ธมฺมสฺส กมฺมปจฺจเยน ปจฺจโย. สหชาตา, นานากฺขณิกา.  สหชาตา คนฺถ\-สมฺปยุตฺตา เจตนา จิตฺต\-สมุฏฺฐานานํ รูปานํ กมฺมปจฺจเยน ปจฺจโย. ทิฏฺฐิคต\-วิปฺปยุตฺต\-โลภ\-สหคตา เจตนา โลภสฺส จ จิตฺตสมุฏฺฐานานญฺจ รูปานํ กมฺมปจฺจเยน ปจฺจโย. โทมนสฺส\-สหคตา เจตนา ปฏิฆสฺส จ จิตฺตสมุฏฺฐานานญฺจ รูปานํ กมฺมปจฺจเยน ปจฺจโย.  นานากฺขณิกา คนฺถ\-สมฺปยุตฺตา เจตนา วิปากานํ ขนฺธานํ กฏตฺตา จ รูปานํ กมฺมปจฺจเยน ปจฺจโย. (2)}

\prose{คนฺถ\-สมฺปยุตฺโต ธมฺโม คนฺถ\-สมฺปยุตฺตสฺส จ คนฺถ\-วิปฺปยุตฺตสฺส จ ธมฺมสฺส กมฺมปจฺจเยน ปจฺจโย. คนฺถ\-สมฺปยุตฺตา เจตนา สมฺปยุตฺตกานํ ขนฺธานํ \csromanfolio{286}จิตฺตสมุฏฺฐานานญฺจ รูปานํ กมฺมปจฺจเยน ปจฺจโย. ทิฏฺฐิคต\-วิปฺปยุตฺต\-โลภ\-สหคตา เจตนา สมฺปยุตฺตกานํ ขนฺธานํ โลภสฺส จ จิตฺตสมุฏฺฐานานญฺจ รูปานํ กมฺมปจฺจเยน ปจฺจโย. โทมนสฺส\-สหคตา เจตนา สมฺปยุตฺตกานํ ขนฺธานํ ปฏิฆสฺส จ จิตฺตสมุฏฺฐานานญฺจ รูปานํ กมฺมปจฺจเยน ปจฺจโย. (3)}

\proseitem{89}{คนฺถ\-วิปฺปยุตฺโต ธมฺโม คนฺถ\-วิปฺปยุตฺตสฺส ธมฺมสฺส กมฺมปจฺจเยน ปจฺจโย. สหชาตา, นานากฺขณิกา.  สหชาตา คนฺถ\-วิปฺปยุตฺตา เจตนา สมฺปยุตฺตกานํ ขนฺธานํ จิตฺตสมุฏฺฐานานญฺจ รูปานํ กมฺมปจฺจเยน ปจฺจโย. ปฏิสนฺธิ\-กฺขเณ ฯเปฯ. นานากฺขณิกา คนฺถ\-วิปฺปยุตฺตา เจตนา วิปากานํ ขนฺธานํ กฏตฺตา จ รูปานํ กมฺมปจฺจเยน ปจฺจโย. (1)}

\csromanheader{2}{\textbf{วิปากาทิ}}
\proseitem{90}{คนฺถ\-วิปฺปยุตฺโต ธมฺโม คนฺถ\-วิปฺปยุตฺตสฺส ธมฺมสฺส วิปาก\-ปจฺจเยน ปจฺจโย. เอกํ.}

\prose{คนฺถ\-สมฺปยุตฺโต ธมฺโม คนฺถ\-สมฺปยุตฺตสฺส ธมฺมสฺส อาหารปจฺจเยน ปจฺจโย. จตฺตาริ. อินฺทฺริย\-ปจฺจเยน ปจฺจโย. จตฺตาริ. ฌานปจฺจเยน ปจฺจโย. จตฺตาริ. มคฺคปจฺจเยน ปจฺจโย. จตฺตาริ. สมฺปยุตฺต\-ปจฺจเยน ปจฺจโย. ฉ.}

\proseitem{91}{คนฺถ\-สมฺปยุตฺโต ธมฺโม คนฺถ\-วิปฺปยุตฺตสฺส ธมฺมสฺส วิปฺปยุตฺต\-ปจฺจเยน ปจฺจโย. สหชาตํ, ปจฺฉาชาตํ. (สํขิตฺตํ. วิภชิตพฺพํ.) (1)}

\prose{คนฺถ\-วิปฺปยุตฺโต ธมฺโม คนฺถ\-วิปฺปยุตฺตสฺส ธมฺมสฺส วิปฺปยุตฺต\-ปจฺจเยน ปจฺจโย. สหชาตํ, ปุเรชาตํ, ปจฺฉาชาตํ. (สํขิตฺตํ.) (1)}

\prose{คนฺถ\-วิปฺปยุตฺโต ธมฺโม คนฺถ\-สมฺปยุตฺตสฺส ธมฺมสฺส วิปฺปยุตฺต\-ปจฺจเยน ปจฺจโย. ปุเรชาตํ วตฺถุ คนฺถ\-สมฺปยุตฺตกานํ ขนฺธานํ วิปฺปยุตฺต\-ปจฺจเยน ปจฺจโย.  (2)}

\prose{คนฺถ\-วิปฺปยุตฺโต ธมฺโม คนฺถ\-สมฺปยุตฺตสฺส จ คนฺถ\-วิปฺปยุตฺตสฺส จ ธมฺมสฺส วิปฺปยุตฺต\-ปจฺจเยน ปจฺจโย. ปุเรชาตํ วตฺถุ ทิฏฺฐิคต\-วิปฺปยุตฺต\-โลภ\-สหคตานํ ขนฺธานํ โลภสฺส จ, โทมนสฺส\-สหคตานํ ขนฺธานํ ปฏิฆสฺส จ วิปฺปยุตฺต\-ปจฺจเยน ปจฺจโย. (3)}

\proseitem{92}{\csromanfolio{287}คนฺถ\-สมฺปยุตฺโต จ คนฺถ\-วิปฺปยุตฺโต จ ธมฺมา คนฺถ\-วิปฺปยุตฺตสฺส ธมฺมสฺส วิปฺปยุตฺต\-ปจฺจเยน ปจฺจโย. สหชาตํ, ปจฺฉาชาตํ.  สหชาตา ทิฏฺฐิคต\-วิปฺปยุตฺต\-โลภ\-สหคตา ขนฺธา จ โลโภ จ จิตฺต\-สมุฏฺฐานานํ รูปานํ วิปฺปยุตฺต\-ปจฺจเยน ปจฺจโย. โทมนสฺส\-สหคตา ขนฺธา จ ปฏิฆญฺจ จิตฺต\-สมุฏฺฐานานํ รูปานํ วิปฺปยุตฺต\-ปจฺจเยน ปจฺจโย.  ปจฺฉาชาตา ทิฏฺฐิคต\-วิปฺปยุตฺต\-โลภ\-สหคตา ขนฺธา จ โลโภ จ, โทมนสฺส\-สหคตา ขนฺธา จ ปฏิฆญฺจ ปุเรชาตสฺส อิมสฺส กายสฺส วิปฺปยุตฺต\-ปจฺจเยน ปจฺจโย.}

\proseitem{93}{คนฺถ\-สมฺปยุตฺโต ธมฺโม คนฺถ\-สมฺปยุตฺตสฺส ธมฺมสฺส อตฺถิปจฺจเยน ปจฺจโย. เอกํ. (ปฏิจฺจสทิสํ.) (1)}

\prose{คนฺถ\-สมฺปยุตฺโต ธมฺโม คนฺถ\-วิปฺปยุตฺตสฺส ธมฺมสฺส อตฺถิปจฺจเยน ปจฺจโย. สหชาตํ, ปจฺฉาชาตํ. (สํขิตฺตํ. วิภชิตพฺพํ.) (2)}

\prose{คนฺถ\-สมฺปยุตฺโต ธมฺโม คนฺถ\-สมฺปยุตฺตสฺส จ คนฺถ\-วิปฺปยุตฺตสฺส จ ธมฺมสฺส อตฺถิปจฺจเยน ปจฺจโย. เอกํ. (ปฏิจฺจสทิสํ.) (3)}

\proseitem{94}{คนฺถ\-วิปฺปยุตฺโต ธมฺโม คนฺถ\-วิปฺปยุตฺตสฺส ธมฺมสฺส อตฺถิปจฺจเยน ปจฺจโย. สหชาตํ, ปุเรชาตํ, ปจฺฉาชาตํ, อาหารํ, อินฺทฺริยํ. (สํขิตฺตํ. วิภชิตพฺพํ.) (1)}

\prose{คนฺถ\-วิปฺปยุตฺโต ธมฺโม คนฺถ\-สมฺปยุตฺตสฺส ธมฺมสฺส อตฺถิปจฺจเยน ปจฺจโย. สหชาตํ, ปุเรชาตํ.  สหชาโต ทิฏฺฐิคต\-วิปฺปยุตฺโต โลโภ สมฺปยุตฺตกานํ ขนฺธานํ อตฺถิปจฺจเยน ปจฺจโย. โทมนสฺส\-สหคตํ ปฏิฆํ สมฺปยุตฺตกานํ ขนฺธานํ อตฺถิปจฺจเยน ปจฺจโย.  ปุเรชาตํ จกฺขุํ ฯเปฯ วตฺถุํ อสฺสาเทติ อภินนฺทติ, ตํ อารพฺภ คนฺถ\-สมฺปยุตฺโต ราโค ฯเปฯ ทิฏฺฐิ ฯเปฯ โทมนสฺสํ อุปฺปชฺชติ. วตฺถุ คนฺถ\-สมฺปยุตฺตกานํ ขนฺธานํ อตฺถิปจฺจเยน ปจฺจโย. (2)}

\prose{คนฺถ\-วิปฺปยุตฺโต ธมฺโม คนฺถ\-สมฺปยุตฺตสฺส จ คนฺถ\-วิปฺปยุตฺตสฺส จ ธมฺมสฺส อตฺถิปจฺจเยน ปจฺจโย. สหชาตํ, ปุเรชาตํ.  สหชาโต ทิฏฺฐิคต\-วิปฺปยุตฺโต โลโภ สมฺปยุตฺตกานํ ขนฺธานํ จิตฺตสมุฏฺฐานานญฺจ รูปานํ อตฺถิปจฺจเยน ปจฺจโย. ปฏิฆํ สมฺปยุตฺตกานํ ขนฺธานํ จิตฺตสมุฏฺฐานานญฺจ \csromanfolio{288}รูปานํ อตฺถิปจฺจเยน ปจฺจโย. ปุเรชาตํ จกฺขุํ ฯเปฯ วตฺถุํ อารพฺภ ทิฏฺฐิคต\-วิปฺปยุตฺต\-โลภ\-สหคตา ขนฺธา จ โลโภ จ, โทมนสฺส\-สหคตา ขนฺธา จ ปฏิฆญฺจ อุปฺปชฺชนฺติ. วตฺถุ ทิฏฺฐิคต\-วิปฺปยุตฺต\-โลภ\-สหคตานํ ขนฺธานํ โลภสฺส จ, โทมนสฺส\-สหคตานํ ขนฺธานํ ปฏิฆสฺส จ อตฺถิปจฺจเยน ปจฺจโย. (3)}

\proseitem{95}{คนฺถ\-สมฺปยุตฺโต จ คนฺถ\-วิปฺปยุตฺโต จ ธมฺมา คนฺถ\-สมฺปยุตฺตสฺส ธมฺมสฺส อตฺถิปจฺจเยน ปจฺจโย. สหชาตํ, ปุเรชาตํ.  สหชาโต คนฺถ\-สมฺปยุตฺโต เอโก ขนฺโธ จ วตฺถุ จ ติณฺณนฺนํ ขนฺธานํ อตฺถิปจฺจเยน ปจฺจโย ฯเปฯ. เทฺว ขนฺธา จ ฯเปฯ. สหชาโต ทิฏฺฐิคต\-วิปฺปยุตฺต\-โลภ\-สหคโต เอโก ขนฺโธ จ โลโภ จ ติณฺณนฺนํ ขนฺธานํ อตฺถิปจฺจเยน ปจฺจโย ฯเปฯ เทฺว ขนฺธา จ ฯเปฯ. โทมนสฺส\-สหคโต เอโก ขนฺโธ จ ปฏิฆญฺจ ติณฺณนฺนํ ขนฺธานํ อตฺถิปจฺจเยน ปจฺจโย ฯเปฯ เทฺว ขนฺธา จ ฯเปฯ. (1)}

\prose{คนฺถ\-สมฺปยุตฺโต จ คนฺถ\-วิปฺปยุตฺโต จ ธมฺมา คนฺถ\-วิปฺปยุตฺตสฺส ธมฺมสฺส อตฺถิปจฺจเยน ปจฺจโย. สหชาตํ, ปุเรชาตํ, ปจฺฉาชาตํ, อาหารํ, อินฺทฺริยํ.  สหชาตา คนฺถ\-สมฺปยุตฺตา ขนฺธา จ มหาภูตา จ จิตฺต\-สมุฏฺฐานานํ รูปานํ อตฺถิปจฺจเยน ปจฺจโย. สหชาตา ทิฏฺฐิคต\-วิปฺปยุตฺต\-โลภ\-สหคตา ขนฺธา จ โลโภ จ จิตฺต\-สมุฏฺฐานานํ รูปานํ อตฺถิปจฺจเยน ปจฺจโย. สหชาตา โทมนสฺส\-สหคตา ขนฺธา จ ปฏิฆญฺจ จิตฺต\-สมุฏฺฐานานํ รูปานํ อตฺถิปจฺจเยน ปจฺจโย. ทิฏฺฐิคต\-วิปฺปยุตฺต\-โลภ\-สหคตา ขนฺธา จ วตฺถุ จ โลภสฺส อตฺถิปจฺจเยน ปจฺจโย. โทมนสฺส\-สหคตา ขนฺธา จ วตฺถุ จ ปฏิฆสฺส อตฺถิปจฺจเยน ปจฺจโย. ปจฺฉาชาตา ทิฏฺฐิคต\-วิปฺปยุตฺต\-โลภ\-สหคตา ขนฺธา จ โลโภ จ โทมนสฺส\-สหคตา ขนฺธา จ ปฏิฆญฺจ ปุเรชาตสฺส อิมสฺส กายสฺส อตฺถิปจฺจเยน ปจฺจโย. ปจฺฉาชาตา คนฺถ\-สมฺปยุตฺตา ขนฺธา จ กพฬีกาโร อาหาโร จ อิมสฺส กายสฺส อตฺถิปจฺจเยน ปจฺจโย. ปจฺฉาชาตา คนฺถ\-สมฺปยุตฺตา ขนฺธา จ รูปชีวิตินฺทฺริยญฺจ กฏตฺตารูปานํ อตฺถิปจฺจเยน ปจฺจโย. (2)}

\prose{คนฺถ\-สมฺปยุตฺโต จ คนฺถ\-วิปฺปยุตฺโต จ ธมฺมา คนฺถ\-สมฺปยุตฺตสฺส จ คนฺถ\-วิปฺปยุตฺตสฺส จ ธมฺมสฺส อตฺถิปจฺจเยน ปจฺจโย. สหชาตํ, ปุเรชาตํ.  สหชาโต ทิฏฺฐิคต\-วิปฺปยุตฺต\-โลภ\-สหคโต เอโก ขนฺโธ จ \csromanfolio{289}โลโภ จ ติณฺณนฺนํ ขนฺธานํ จิตฺตสมุฏฺฐานานญฺจ รูปานํ อตฺถิปจฺจเยน ปจฺจโย ฯเปฯ เทฺว ขนฺธา จ ฯเปฯ. สหชาโต โทมนสฺส\-สหคโต เอโก ขนฺโธ จ ปฏิฆญฺจ ติณฺณนฺนํ ขนฺธานํ จิตฺตสมุฏฺฐานานญฺจ รูปานํ อตฺถิปจฺจเยน ปจฺจโย ฯเปฯ เทฺว ขนฺธา จ ฯเปฯ. สหชาโต ทิฏฺฐิคต\-วิปฺปยุตฺต\-โลภ\-สหคโต เอโก ขนฺโธ จ วตฺถุ จ ติณฺณนฺนํ ขนฺธานํ โลภสฺส จ อตฺถิปจฺจเยน ปจฺจโย ฯเปฯ เทฺว ขนฺธา จ ฯเปฯ. สหชาโต โทมนสฺส\-สหคโต เอโก ขนฺโธ จ วตฺถุ จ ติณฺณนฺนํ ขนฺธานํ ปฏิฆสฺส จ อตฺถิปจฺจเยน ปจฺจโย ฯเปฯ เทฺว ขนฺธา จ ฯเปฯ. (3)}

\csromanmarkh{1. ปจฺจยานุโลม}\csromanmarkhii{2. สงฺขฺยาวาร}\csromanheadernomark{2}{1. \textbf{ปจฺจยานุโลม 2. สงฺขฺยาวาร}}
\proseitem{96}{เหตุยา นว, อารมฺมเณ นว, อธิปติยา นว, อนนฺตเร นว, สมนนฺตเร นว, สหชาเต นว, อญฺญมญฺเญ ฉ, นิสฺสเย นว, อุปนิสฺสเย นว, ปุเรชาเต ตีณิ, ปจฺฉาชาเต ตีณิ, อาเสวเน นว, กมฺเม จตฺตาริ, วิปาเก เอกํ, อาหาเร จตฺตาริ, อินฺทฺริเย จตฺตาริ, ฌาเน จตฺตาริ, มคฺเค จตฺตาริ, สมฺปยุตฺเต ฉ, วิปฺปยุตฺเต ปญฺจ, อตฺถิยา นว, นตฺถิยา นว, วิคเต นว, อวิคเต นว.}

\vspace{\tipitakaparskip}
\csromancenter{อนุโลมํ.}
\proseitem{97}{คนฺถ\-สมฺปยุตฺโต ธมฺโม คนฺถ\-สมฺปยุตฺตสฺส ธมฺมสฺส อารมฺมณ\-ปจฺจเยน ปจฺจโย. สหชาต\-ปจฺจเยน ปจฺจโย. อุปนิสฺสย\-ปจฺจเยน ปจฺจโย. (1)}

\prose{คนฺถ\-สมฺปยุตฺโต ธมฺโม คนฺถ\-วิปฺปยุตฺตสฺส ธมฺมสฺส อารมฺมณ\-ปจฺจเยน ปจฺจโย. สหชาต\-ปจฺจเยน ปจฺจโย. อุปนิสฺสย\-ปจฺจเยน ปจฺจโย. ปจฺฉาชาต\-ปจฺจเยน ปจฺจโย. กมฺมปจฺจเยน ปจฺจโย. (2)}

\prose{คนฺถ\-สมฺปยุตฺโต ธมฺโม คนฺถ\-สมฺปยุตฺตสฺส จ คนฺถ\-วิปฺปยุตฺตสฺส จ ธมฺมสฺส อารมฺมณ\-ปจฺจเยน ปจฺจโย. สหชาต\-ปจฺจเยน ปจฺจโย. อุปนิสฺสย\-ปจฺจเยน ปจฺจโย. (93)}

\proseitem{98}{\csromanfolio{290}คนฺถ\-วิปฺปยุตฺโต ธมฺโม คนฺถ\-วิปฺปยุตฺตสฺส ธมฺมสฺส อารมฺมณ\-ปจฺจเยน ปจฺจโย. สหชาต\-ปจฺจเยน ปจฺจโย. อุปนิสฺสย\-ปจฺจเยน ปจฺจโย. ปุเรชาต\-ปจฺจเยน ปจฺจโย. ปจฺฉาชาต\-ปจฺจเยน ปจฺจโย. กมฺมปจฺจเยน ปจฺจโย. อาหารปจฺจเยน ปจฺจโย. อินฺทฺริย\-ปจฺจเยน ปจฺจโย. (1)}

\prose{คนฺถ\-วิปฺปยุตฺโต ธมฺโม คนฺถ\-สมฺปยุตฺตสฺส ธมฺมสฺส อารมฺมณ\-ปจฺจเยน ปจฺจโย. สหชาต\-ปจฺจเยน ปจฺจโย. อุปนิสฺสย\-ปจฺจเยน ปจฺจโย. ปุเรชาต\-ปจฺจเยน ปจฺจโย. (2)}

\prose{คนฺถ\-วิปฺปยุตฺโต ธมฺโม คนฺถ\-สมฺปยุตฺตสฺส จ คนฺถ\-วิปฺปยุตฺตสฺส จ ธมฺมสฺส อารมฺมณ\-ปจฺจเยน ปจฺจโย. สหชาต\-ปจฺจเยน ปจฺจโย. อุปนิสฺสย\-ปจฺจเยน ปจฺจโย. ปุเรชาต\-ปจฺจเยน ปจฺจโย. (3)}

\proseitem{99}{คนฺถ\-สมฺปยุตฺโต จ คนฺถ\-วิปฺปยุตฺโต จ ธมฺมา คนฺถ\-สมฺปยุตฺตสฺส ธมฺมสฺส อารมฺมณ\-ปจฺจเยน ปจฺจโย. สหชาต\-ปจฺจเยน ปจฺจโย. อุปนิสฺสย\-ปจฺจเยน ปจฺจโย. (1)}

\prose{คนฺถ\-สมฺปยุตฺโต จ คนฺถ\-วิปฺปยุตฺโต จ ธมฺมา คนฺถ\-วิปฺปยุตฺตสฺส ธมฺมสฺส อารมฺมณ\-ปจฺจเยน ปจฺจโย. สหชาต\-ปจฺจเยน ปจฺจโย. อุปนิสฺสย\-ปจฺจเยน ปจฺจโย. ปจฺฉาชาต\-ปจฺจเยน ปจฺจโย. (2)}

\prose{คนฺถ\-สมฺปยุตฺโต จ คนฺถ\-วิปฺปยุตฺโต จ ธมฺมา คนฺถ\-สมฺปยุตฺตสฺส จ คนฺถ\-วิปฺปยุตฺตสฺส จ ธมฺมสฺส อารมฺมณ\-ปจฺจเยน ปจฺจโย. สหชาต\-ปจฺจเยน ปจฺจโย. อุปนิสฺสย\-ปจฺจเยน ปจฺจโย. (3)}

\csromanmarkh{2. ปจฺจย\-ปจฺจนีย}\csromanmarkhii{2. สงฺขฺยาวาร}\csromanheadernomark{2}{2. ปจฺจย\-ปจฺจนีย 2. สงฺขฺยาวาร}
\vspace{\tipitakaparskip}
\csromancenter{นเหตุยา นว, นอารมฺมเณ นว, (สพฺพตฺถ นว,) โนอวิคเต นว.}
\proseitem{101}{เหตุปจฺจยา นอารมฺมเณ นว, นอธิปติยา นว ฯเปฯ นสมนนฺตเร นว, นอญฺญมญฺเญ ตีณิ, นอุปนิสฺสเย นว ฯเปฯ นมคฺเค}

\vspace{\tipitakaparskip}
\csromancenter{คนฺถ\-คนฺถนิยทุก นว, นสมฺปยุตฺเต ตีณิ, นวิปฺปยุตฺเต ฉ, โนนตฺถิยา นว, โนวิคเต นว.}
\proseitemwithcloser{102}{\csromanfolio{291}นเหตุปจฺจยา อารมฺมเณ นว, อธิปติยา นว. (อนุโลมมาติกา วิตฺถาเรตพฺพา.) ฯเปฯ อวิคเต นว.}{คนฺถ\-สมฺปยุตฺตทุกํ นิฏฺฐิตํ.}
\csromanmatikaanchor{mtk.35}
\csromantocmarkref{section}{29. คนฺถคนฺถนิยทุก}{mtk.35}
\bookmark[dest={mtk.35},level=1]{29. คนฺถคนฺถนิยทุก}
\csromanmarkh{29. คนฺถ\-คนฺถนิยทุก}\csromanmarkhii{1. ปฏิจฺจวาร}\csromanheadernomark{1}{29. คนฺถ\-คนฺถนิยทุก 1. ปฏิจฺจวาร}
\csromanheader{2}{\textbf{เหตุ}}
\proseitem{103}{คนฺถญฺเจว คนฺถนิยญฺจ ธมฺมํ ปฏิจฺจ คนฺโถ เจว คนฺถนิโย จ ธมฺโม อุปฺปชฺชติ เหตุปจฺจยา. สีลพฺพต\-ปรามาสํ กายคนฺถํ ปฏิจฺจ อภิชฺฌา\-กายคนฺโถ, อภิชฺฌา\-กายคนฺถํ ปฏิจฺจ สีลพฺพต\-ปรามาโส กายคนฺโถ, อิทํสจฺจาภินิเวส\-กายคนฺถํ ปฏิจฺจ อภิชฺฌา\-กายคนฺโถ, อภิชฺฌา\-กายคนฺถํ ปฏิจฺจ อิทํสจฺจาภินิเวส\-กายคนฺโถ. (1)}

\prose{คนฺถญฺเจว คนฺถนิยญฺจ ธมฺมํ ปฏิจฺจ คนฺถนิโย เจว โน จ คนฺโถ ธมฺโม อุปฺปชฺชติ เหตุปจฺจยา. คนฺเถ ปฏิจฺจ สมฺปยุตฺตกา ขนฺธา จิตฺต\-สมุฏฺฐานญฺจ รูปํ. (2)}

\prose{คนฺถญฺเจว คนฺถนิยญฺจ ธมฺมํ ปฏิจฺจ คนฺโถ เจว คนฺถนิโย จ คนฺถนิโย เจว โน จ คนฺโถ จ ธมฺมา อุปฺปชฺชนฺติ เหตุปจฺจยา. (1)}

\prose{(ปฏิจฺจวารมฺปิ สหชาตวารมฺปิ ปจฺจยวารมฺปิ นิสฺสยวารมฺปิ สํสฏฺฐ\-วารมฺปิ สมฺปยุตฺตวารมฺปิ คนฺถ\-ทุก\-สทิสํ, นินฺนานากรณํ.)}

\csromanmarkh{29. คนฺถ\-คนฺถนิยทุก}\csromanmarkhii{7. ปญฺหาวาร}\csromanheadernomark{2}{29. คนฺถ\-คนฺถนิยทุก 7. ปญฺหาวาร}
\csromanmarkh{1. ปจฺจยานุโลม}\csromanmarkhii{1. วิภงฺควาร}\csromanheadernomark{2}{1. \textbf{ปจฺจยานุโลม 1. วิภงฺควาร}}
\csromanheader{2}{\textbf{เหตุ}}
\proseitem{104}{\csromanfolio{292}คนฺโถ เจว คนฺถนิโย จ ธมฺโม คนฺถสฺส เจว คนฺถนิยสฺส จ ธมฺมสฺส เหตุปจฺจเยน ปจฺจโย. คนฺถา เหตู สมฺปยุตฺตกานํ คนฺถานํ เหตุปจฺจเยน ปจฺจโย. (เอวํ นว ปญฺหา วิตฺถาเรตพฺพา.)}

\proseitem{105}{คนฺโถ เจว คนฺถนิโย จ ธมฺโม คนฺถสฺส เจว คนฺถนิยสฺส จ ธมฺมสฺส อารมฺมณ\-ปจฺจเยน ปจฺจโย. คนฺเถ อารพฺภ คนฺถา อุปฺปชฺชนฺติ. (มูลํ ปุจฺฉิตพฺพํ.) คนฺเถ อารพฺภ คนฺถนิยา เจว โน จ คนฺถา ขนฺธา อุปฺปชฺชนฺติ. (มูลํ ปุจฺฉิตพฺพํ.) คนฺเถ อารพฺภ คนฺถา จ สมฺปยุตฺตกา จ ขนฺธา อุปฺปชฺชนฺติ. (3)}

\proseitem{106}{คนฺถนิโย เจว โน จ คนฺโถ ธมฺโม คนฺถนิยสฺส เจว โน จ คนฺถสฺส จ ธมฺมสฺส อารมฺมณ\-ปจฺจเยน ปจฺจโย. ทานํ ฯเปฯ สีลํ ฯเปฯ อุโปสถกมฺมํ กตฺวา ตํ ปจฺจเวกฺขติ. ปุพฺเพ สุจิณฺณานิ ฯเปฯ. ฌานา ฯเปฯ. อริยา โคตฺรภุํ ปจฺจเวกฺขนฺติ, โวทานํ ปจฺจเวกฺขนฺติ. ปหีเน กิเลเส ฯเปฯ. วิกฺขมฺภิเต กิเลเส ปจฺจเวกฺขนฺติ. ปุพฺเพ สมุทาจิณฺเณ กิเลเส ชานนฺติ. จกฺขุํ ฯเปฯ วตฺถุํ คนฺถนิเย เจว โน จ คนฺเถ ขนฺเธ อนิจฺจโต ฯเปฯ โทมนสฺสํ อุปฺปชฺชติ. ทิพฺเพน จกฺขุนา รูปํ ปสฺสติ. ทิพฺพาย โสตธาตุยา สทฺทํ สุณาติ. (สพฺพํ วิตฺถาเรตพฺพํ.) อาวชฺชนาย อารมฺมณ\-ปจฺจเยน ปจฺจโย. (1)}

\prose{คนฺถนิโย เจว โน จ คนฺโถ ธมฺโม คนฺถสฺส เจว คนฺถนิยสฺส จ ธมฺมสฺส อารมฺมณ\-ปจฺจเยน ปจฺจโย. ทานํ ฯเปฯ สีลํ ฯเปฯ อุโปสถกมฺมํ กตฺวา ตํ อสฺสาเทติ อภินนฺทติ, ตํ อารพฺภ ราโค อุปฺปชฺชติ, ทิฏฺฐิ ฯเปฯ วิจิกิจฺฉา ฯเปฯ โทมนสฺสํ อุปฺปชฺชติ. ปุพฺเพ สุจิณฺณานิ ฯเปฯ. ฌานา วุฏฺฐหิตฺวา ฌานํ ฯเปฯ. จกฺขุํ ฯเปฯ วตฺถุํ คนฺถนิเย เจว โน จ คนฺเถ ขนฺเธ อสฺสาเทติ อภินนฺทติ, ตํ อารพฺภ ราโค อุปฺปชฺชติ, ทิฏฺฐิ ฯเปฯ โทมนสฺสํ อุปฺปชฺชติ. (2)}

\csromanmatikaanchor{mtk.36}
\csromantocmarkref{section}{30. คนฺถคนฺถสมฺปยุตฺตทุก}{mtk.36}
\bookmark[dest={mtk.36},level=1]{30. คนฺถคนฺถสมฺปยุตฺตทุก}
\csromanheader{1}{30. คนฺถ\-คนฺถ\-สมฺปยุตฺตทุก}
\prose{\csromanfolio{293}คนฺถนิโย เจว โน จ คนฺโถ ธมฺโม คนฺถสฺส เจว คนฺถนิยสฺส จ คนฺถนิยสฺส เจว โน จ คนฺถสฺส จ ธมฺมสฺส อารมฺมณ\-ปจฺจเยน ปจฺจโย. ทานํ ฯเปฯ สีลํ ฯเปฯ อุโปสถกมฺมํ ฯเปฯ. ปุพฺเพ สุจิณฺณานิ ฯเปฯ. ฌานา วุฏฺฐหิตฺวา ฌานํ ฯเปฯ. จกฺขุํ ฯเปฯ วตฺถุํ คนฺถนิเย เจว โน จ คนฺเถ ขนฺเธ อสฺสาเทติ อภินนฺทติ, ตํ อารพฺภ คนฺถา จ สมฺปยุตฺตกา จ ขนฺธา อุปฺปชฺชนฺติ. (เอวํ อิตเรปิ ตีณิ วิตฺถาเรตพฺพา.) (3)}

\prosewithcloser{(อารพฺภ กาตพฺพา. อิมสฺมิํ ทุเก โลกุตฺตรํ นตฺถิ, คนฺถ\-ทุก\-สทิสํ, นินฺนานากรณํ. “คนฺถนิยนฺ”ติ นิยาเมตพฺพํ.  มคฺเค นว ปญฺหา กาตพฺพา.)}{คนฺถ\-คนฺถนิยทุกํ นิฏฺฐิตํ.}
\csromanmarkh{30. คนฺถ\-คนฺถ\-สมฺปยุตฺตทุก}\csromanmarkhii{1. ปฏิจฺจวาร}\csromanheadernomark{2}{30. คนฺถ\-คนฺถ\-สมฺปยุตฺตทุก 1. ปฏิจฺจวาร}
\csromanheader{2}{1-4. ปจฺจยจตุกฺก}
\csromanheader{2}{\textbf{เหตุ}}
\proseitem{107}{คนฺถญฺเจว คนฺถ\-สมฺปยุตฺตญฺจ ธมฺมํ ปฏิจฺจ คนฺโถ เจว คนฺถ\-สมฺปยุตฺโต จ ธมฺโม อุปฺปชฺชติ เหตุปจฺจยา. สีลพฺพต\-ปรามาสํ กายคนฺถํ ปฏิจฺจ อภิชฺฌา\-กายคนฺโถ, อภิชฺฌา\-กายคนฺถํ ปฏิจฺจ สีลพฺพต\-ปรามาโส กายคนฺโถ, อิทํสจฺจาภินิเวส\-กายคนฺถํ ปฏิจฺจ อภิชฺฌา\-กายคนฺโถ, อภิชฺฌา\-กายคนฺถํ ปฏิจฺจ อิทํสจฺจา\-ภินิเวโส กายคนฺโถ. (1)}

\prose{คนฺถญฺเจว คนฺถ\-สมฺปยุตฺตญฺจ ธมฺมํ ปฏิจฺจ คนฺถ\-สมฺปยุตฺโต เจว โน จ คนฺโถ ธมฺโม อุปฺปชฺชติ เหตุปจฺจยา. คนฺเถ ปฏิจฺจ สมฺปยุตฺตกา ขนฺธา. (2)}

\prose{คนฺถญฺเจว คนฺถ\-สมฺปยุตฺตญฺจ ธมฺมํ ปฏิจฺจ คนฺโถ เจว คนฺถ\-สมฺปยุตฺโต จ คนฺถ\-สมฺปยุตฺโต เจว โน จ คนฺโถ จ ธมฺมา อุปฺปชฺชนฺติ เหตุปจฺจยา. สีลพฺพต\-ปรามาสํ กายคนฺถํ ปฏิจฺจ อภิชฺฌา\-กายคนฺโถ สมฺปยุตฺตกา จ ขนฺธา. (จกฺกํ.) (3)}

\proseitem{108}{คนฺถ\-สมฺปยุตฺต\-ญฺเจว โน จ คนฺถํ ธมฺมํ ปฏิจฺจ คนฺถ\-สมฺปยุตฺโต เจว โน จ คนฺโถ ธมฺโม อุปฺปชฺชติ เหตุปจฺจยา. คนฺถ\-สมฺปยุตฺต\-ญฺเจว โน จ คนฺถํ เอกํ ขนฺธํ ปฏิจฺจ ตโย ขนฺธา ฯเปฯ เทฺว ขนฺเธ ฯเปฯ. (1)}

\prose{\csromanfolio{294}คนฺถ\-สมฺปยุตฺต\-ญฺเจว โน จ คนฺถํ ธมฺมํ ปฏิจฺจ คนฺโถ เจว คนฺถ\-สมฺปยุตฺโต จ ธมฺโม อุปฺปชฺชติ เหตุปจฺจยา. คนฺถ\-สมฺปยุตฺเต เจว โน จ คนฺเถ ขนฺเธ ปฏิจฺจ คนฺถา. (2)}

\prose{คนฺถ\-สมฺปยุตฺต\-ญฺเจว โน จ คนฺถํ ธมฺมํ ปฏิจฺจ คนฺโถ เจว คนฺถ\-สมฺปยุตฺโต จ คนฺถ\-สมฺปยุตฺโต เจว โน จ คนฺโถ จ ธมฺมา อุปฺปชฺชนฺติ เหตุปจฺจยา. คนฺถ\-สมฺปยุตฺต\-ญฺเจว โน จ คนฺถํ เอกํ ขนฺธํ ปฏิจฺจ ตโย ขนฺธา คนฺถา จ ฯเปฯ เทฺว ขนฺเธ ฯเปฯ. (3)}

\proseitem{109}{คนฺถญฺเจว คนฺถ\-สมฺปยุตฺตญฺจ คนฺถ\-สมฺปยุตฺต\-ญฺเจว โน จ คนฺถญฺจ ธมฺมํ ปฏิจฺจ คนฺโถ เจว คนฺถ\-สมฺปยุตฺโต จ ธมฺโม อุปฺปชฺชติ เหตุปจฺจยา. คนฺเถ จ สมฺปยุตฺตเก จ ขนฺเธ ปฏิจฺจ คนฺถา. (1)}

\prose{คนฺถญฺเจว คนฺถ\-สมฺปยุตฺตญฺจ คนฺถ\-สมฺปยุตฺต\-ญฺเจว โน จ คนฺถญฺจ ธมฺมํ ปฏิจฺจ คนฺถ\-สมฺปยุตฺโต เจว โน จ คนฺโถ ธมฺโม อุปฺปชฺชติ เหตุปจฺจยา. คนฺถ\-สมฺปยุตฺต\-ญฺเจว โน จ คนฺถํ เอกํ ขนฺธญฺจ คนฺเถ จ ปฏิจฺจ ตโย ขนฺธา ฯเปฯ เทฺว ขนฺเธ จ ฯเปฯ. (2)}

\prose{คนฺถญฺเจว คนฺถ\-สมฺปยุตฺตญฺจ คนฺถ\-สมฺปยุตฺต\-ญฺเจว โน จ คนฺถญฺจ ธมฺมํ ปฏิจฺจ คนฺโถ เจว คนฺถ\-สมฺปยุตฺโต จ คนฺถ\-สมฺปยุตฺโต เจว โน จ คนฺโถ จ ธมฺมา อุปฺปชฺชนฺติ เหตุปจฺจยา. คนฺถ\-สมฺปยุตฺต\-ญฺเจว โน จ คนฺถํ เอกํ ขนฺธญฺจ สีลพฺพต\-ปรามาสํ กายคนฺถญฺจ ปฏิจฺจ ตโย ขนฺธา อภิชฺฌา\-กายคนฺโถ จ ฯเปฯ เทฺว ขนฺเธ จ ฯเปฯ. (จกฺกํ พนฺธิตพฺพํ. สํขิตฺตํ.) (3)}

\proseitem{110}{เหตุยา นว, อารมฺมเณ นว, (สพฺพตฺถ นว,) กมฺเม นว, อาหาเร นว ฯเปฯ อวิคเต นว.}

\vspace{\tipitakaparskip}
\csromancenter{อนุโลมํ.}
\csromanheader{2}{2. \textbf{ปจฺจนีย}}
\proseitem{111}{คนฺถญฺเจว คนฺถ\-สมฺปยุตฺตญฺจ ธมฺมํ ปฏิจฺจ คนฺโถ เจว คนฺถ\-สมฺปยุตฺโต จ ธมฺโม อุปฺปชฺชติ นอธิปติ\-ปจฺจยา. (สํขิตฺตํ.)}

\csromanheader{2}{30. คนฺถ\-คนฺถ\-สมฺปยุตฺตทุก}
\prose{\csromanfolio{295}(อิธ นเหตุปจฺจโย นตฺถิ.) นอธิปติยา นว, นปุเรชาเต นว, นปจฺฉาชาเต นว, นอาเสวเน นว, นกมฺเม ตีณิ, นวิปาเก นว, นวิปฺปยุตฺเต นว.}

\prose{(เอวํ อิตเร เทฺว คณนาปิ สหชาตวาโรปิ ปจฺจยวาโรปิ นิสฺสยวาโรปิ สํสฏฺฐ\-วาโรปิ สมฺปยุตฺตวาโรปิ ปฏิจฺจ\-วาร\-สทิสา.)}

\csromanmarkh{30. คนฺถ\-คนฺถ\-สมฺปยุตฺตทุก}\csromanmarkhii{7. ปญฺหาวาร}\csromanheadernomark{2}{30. คนฺถ\-คนฺถ\-สมฺปยุตฺตทุก 7. ปญฺหาวาร}
\csromanmarkh{1. ปจฺจยานุโลม}\csromanmarkhii{1. วิภงฺควาร}\csromanheadernomark{2}{1. \textbf{ปจฺจยานุโลม 1. วิภงฺควาร}}
\csromanheader{2}{\textbf{เหตุ}}
\proseitem{112}{คนฺโถ เจว คนฺถ\-สมฺปยุตฺโต จ ธมฺโม คนฺถสฺส เจว คนฺถ\-สมฺปยุตฺตสฺส จ ธมฺมสฺส เหตุปจฺจเยน ปจฺจโย. คนฺถา เจว คนฺถ\-สมฺปยุตฺตา จ เหตู สมฺปยุตฺตกานํ คนฺถานํ เหตุปจฺจเยน ปจฺจโย. (1)}

\prose{คนฺโถ เจว คนฺถ\-สมฺปยุตฺโต จ ธมฺโม คนฺถ\-สมฺปยุตฺตสฺส เจว โน จ คนฺถสฺส ธมฺมสฺส เหตุปจฺจเยน ปจฺจโย. คนฺถา เจว คนฺถ\-สมฺปยุตฺตา จ เหตู สมฺปยุตฺตกานํ ขนฺธานํ เหตุปจฺจเยน ปจฺจโย. (2)}

\prose{คนฺโถ เจว คนฺถ\-สมฺปยุตฺโต จ ธมฺโม คนฺถสฺส เจว คนฺถ\-สมฺปยุตฺตสฺส จ คนฺถ\-สมฺปยุตฺตสฺส เจว โน จ คนฺถสฺส จ ธมฺมสฺส เหตุปจฺจเยน ปจฺจโย. คนฺถา เจว คนฺถ\-สมฺปยุตฺตา จ เหตู สมฺปยุตฺตกานํ ขนฺธานํ คนฺถานญฺจ เหตุปจฺจเยน ปจฺจโย. (3)}

\proseitem{113}{คนฺถ\-สมฺปยุตฺโต เจว โน จ คนฺโถ ธมฺโม คนฺถ\-สมฺปยุตฺตสฺส เจว โน จ คนฺถสฺส ธมฺมสฺส เหตุปจฺจเยน ปจฺจโย. คนฺถ\-สมฺปยุตฺตา เจว โน จ คนฺถา เหตู สมฺปยุตฺตกานํ คนฺถานํ เหตุปจฺจเยน ปจฺจโย. (1)}

\prose{คนฺถ\-สมฺปยุตฺโต เจว โน จ คนฺโถ ธมฺโม คนฺถสฺส เจว คนฺถ\-สมฺปยุตฺตสฺส จ ธมฺมสฺส เหตุปจฺจเยน ปจฺจโย. คนฺถ\-สมฺปยุตฺตา เจว โนจ คนฺถา เหตู สมฺปยุตฺตกานํ คนฺถานํ เหตุปจฺจเยน ปจฺจโย. (2)}

\prose{คนฺถ\-สมฺปยุตฺโต เจว โน จ คนฺโถ ธมฺโม คนฺถสฺส เจว คนฺถ\-สมฺปยุตฺตสฺส จ คนฺถ\-สมฺปยุตฺตสฺส เจว โน จ คนฺถสฺส จ ธมฺมสฺส \csromanfolio{296}เหตุปจฺจเยน ปจฺจโย. คนฺถ\-สมฺปยุตฺตา เจว โน จ คนฺถา เหตู สมฺปยุตฺตกานํ ขนฺธานํ คนฺถานญฺจ เหตุปจฺจเยน ปจฺจโย. (3)}

\proseitem{114}{คนฺโถ เจว คนฺถ\-สมฺปยุตฺโต จ คนฺถ\-สมฺปยุตฺโต เจว โน จ คนฺโถ จ ธมฺมา คนฺถสฺส เจว คนฺถ\-สมฺปยุตฺตสฺส จ ธมฺมสฺส เหตุปจฺจเยน ปจฺจโย. คนฺถา เจว คนฺถ\-สมฺปยุตฺตา จ คนฺถ\-สมฺปยุตฺตา เจว โน จ คนฺถา จ เหตู สมฺปยุตฺตกานํ ขนฺธานํ เหตุปจฺจเยน ปจฺจโย. (1)}

\prose{คนฺโถ เจว คนฺถ\-สมฺปยุตฺโต จ คนฺถ\-สมฺปยุตฺโต เจว โน จ คนฺโถ จ ธมฺมา คนฺถ\-สมฺปยุตฺตสฺส เจว โน จ คนฺถสฺส ธมฺมสฺส เหตุปจฺจเยน ปจฺจโย. คนฺถา เจว คนฺถ\-สมฺปยุตฺตา จ คนฺถ\-สมฺปยุตฺตา เจว โน จ คนฺถา จ เหตู สมฺปยุตฺตกานํ ขนฺธานํ เหตุปจฺจเยน ปจฺจโย. (2)}

\prose{คนฺโถ เจว คนฺถ\-สมฺปยุตฺโต จ คนฺถ\-สมฺปยุตฺโต เจว โน จ คนฺโถ จ ธมฺมา คนฺถสฺส เจว คนฺถ\-สมฺปยุตฺตสฺส จ คนฺถ\-สมฺปยุตฺตสฺส เจว โน จ คนฺถสฺส จ ธมฺมสฺส เหตุปจฺจเยน ปจฺจโย. คนฺถา เจว คนฺถ\-สมฺปยุตฺตา จ คนฺถ\-สมฺปยุตฺตา เจว โน จ คนฺถา จ เหตู สมฺปยุตฺตกานํ ขนฺธานํ คนฺถานญฺจ เหตุปจฺจเยน ปจฺจโย. (3)}

\csromanheader{2}{\textbf{อารมฺมณาทิ}}
\proseitem{115}{คนฺโถ เจว คนฺถ\-สมฺปยุตฺโต จ ธมฺโม คนฺถสฺส เจว คนฺถ\-สมฺปยุตฺตสฺส จ ธมฺมสฺส อารมฺมณ\-ปจฺจเยน ปจฺจโย. คนฺเถ อารพฺภ คนฺถา อุปฺปชฺชนฺติ. (มูลํ กาตพฺพํ.) คนฺเถ อารพฺภ คนฺถ\-สมฺปยุตฺตา เจว โน จ คนฺถา ขนฺธา อุปฺปชฺชนฺติ. (มูลํ กาตพฺพํ.). คนฺเถ อารพฺภ คนฺถา จ คนฺถ\-สมฺปยุตฺตกา จ ขนฺธา อุปฺปชฺชนฺติ. (3)}

\prose{คนฺถ\-สมฺปยุตฺโต เจว โน จ คนฺโถ ธมฺโม คนฺถ\-สมฺปยุตฺตสฺส เจว โน จ คนฺถสฺส ธมฺมสฺส อารมฺมณ\-ปจฺจเยน ปจฺจโย. คนฺถ\-สมฺปยุตฺเต เจว โน จ คนฺเถ ขนฺเธ อารพฺภ คนฺถ\-สมฺปยุตฺตา เจว โน จ คนฺถา ขนฺธา อุปฺปชฺชนฺติ. (มูลํ กาตพฺพํ.) คนฺถ\-สมฺปยุตฺเต เจว โน จ คนฺเถ ขนฺเธ อารพฺภ คนฺถา อุปฺปชฺชนฺติ. (มูลํ กาตพฺพํ.) คนฺถ\-สมฺปยุตฺเต เจว โน จ คนฺเถ ขนฺเธ อารพฺภ คนฺถา จ คนฺธสมฺปยุตฺตกา จ ขนฺธา อุปฺปชฺชนฺติ. (3)}

\prose{(เอวํ อิตเรปิ ตีณิ ปญฺหา กาตพฺพา. อารมฺมณ\-สทิสํ เยว. อธิปติยาปิ อนนฺตเรปิ อุปนิสฺสเยปิ วิภาโค นตฺถิ.)}

\csromanmatikaanchor{mtk.37}
\csromanmatikaanchor{mtk.39}
\csromantocmarkref{section}{31. คนฺถวิปฺปยุตฺตคนฺถนิยทุก}{mtk.37}
\bookmark[dest={mtk.37},level=1]{31. คนฺถวิปฺปยุตฺตคนฺถนิยทุก}
\csromanmarkh{32-43. โอฆาทิทุก}\csromanmarkhii{1. ปจฺจยานุโลม 2. สงฺขฺยาวาร}\csromanheadernomark{2}{32-43. โอฆาทิทุก \textbf{1. ปจฺจยานุโลม 2. สงฺขฺยาวาร}}
\proseitemwithcloser{116}{\csromanfolio{297}เหตุยา นว, อารมฺมเณ นว, อธิปติยา นว, อนนฺตเร นว, สมนนฺตเร นว, สหชาเต นว, อญฺญมญฺเญ นว, นิสฺสเย นว, อุปนิสฺสเย นว, อาเสวเน นว, กมฺเม ตีณิ, อาหาเร ตีณิ, อินฺทฺริเย ตีณิ, ฌาเน ตีณิ, มคฺเค นว, สมฺปยุตฺเต นว, อตฺถิยา นว, นตฺถิยา นว, วิคเต นว, อวิคเต นว. (อรูปํ เยว ปจฺจยํ, เอเกกสฺส ตีณิ ตีณิ กาตพฺพา.  อารมฺมณญฺจ สหชาตญฺจ อุปนิสฺสยญฺจ นวสุปิ ปริวตฺเต\-ตพฺพํ. เอวํ ปญฺหาวาเรปิ สพฺพํ กาตพฺพํ.)}{คนฺถ\-คนฺถ\-สมฺปยุตฺตทุกํ นิฏฺฐิตํ.}
\csromanmarkh{31. คนฺถ\-วิปฺปยุตฺต\-คนฺถนิยทุก}\csromanmarkhii{1. ปฏิจฺจวาร}\csromanheadernomark{1}{31. คนฺถ\-วิปฺปยุตฺต\-คนฺถนิยทุก 1. ปฏิจฺจวาร}
\vspace{\tipitakaparskip}
\csromancenter{\textbf{เหตุ}}
\proseitemwithcloser{117}{คนฺถ\-วิปฺปยุตฺต\-คนฺถนิยํ ธมฺมํ ปฏิจฺจ คนฺถ\-วิปฺปยุตฺต\-คนฺถนิโย ธมฺโม อุปฺปชฺชติ เหตุปจฺจยา. คนฺถ\-วิปฺปยุตฺตํ คนฺถนิยํ เอกํ ขนฺธํ ปฏิจฺจ ตโย ขนฺธา จิตฺต\-สมุฏฺฐานญฺจ รูปํ ฯเปฯ เทฺว ขนฺเธ ฯเปฯ. ปฏิสนฺธิ\-กฺขเณ ฯเปฯ. เอกํ มหาภูตํ ฯเปฯ. (ยถา จูฬนฺตรทุเก โลกิยทุกํ, เอวํ วิตฺถาเรตพฺพํ, นินฺนานากรณํ.) คนฺถ\-วิปฺปยุตฺต\-คนฺถนิยทุกํ นิฏฺฐิตํ.}{คนฺถ\-โคจฺฉกํ นิฏฺฐิตํ.}
\csromanmatikaanchor{mtk.38}
\csromantocmarknopage{section}{6-7. โอฆโยคโคจฺฉก}{mtk.38}
\bookmark[dest={mtk.38},level=1]{6-7. โอฆโยคโคจฺฉก}
\csromantocmarkref{subsection}{32-43. โอฆาทิทุก}{mtk.39}
\bookmark[dest={mtk.39},level=2]{32-43. โอฆาทิทุก}
\csromanheader{1}{6-7. โอฆ\-โยค\-โคจฺฉก}
\csromanheader{2}{32-43. โอฆาทิทุก}
\proseitem{1}{โอฆํ ธมฺมํ ปฏิจฺจ โอโฆ ธมฺโม อุปฺปชฺชติ เหตุปจฺจยา ฯเปฯ.}

\proseitem{2}{โยคํ ธมฺมํ ปฏิจฺจ โยโค ธมฺโม อุปฺปชฺชติ เหตุปจฺจยา ฯเปฯ.}

\vspace{\tipitakaparskip}
\csromancenter{(เทฺวปิ โคจฺฉกา อาสว\-โคจฺฉก\-สทิสา, นินฺนานากรณา.)}
\nitthitam{โอฆ\-โยค\-โคจฺฉกํ นิฏฺฐิตํ.}
\csromanmatikaanchor{mtk.40}
\csromantocmarknopage{section}{8. นีวรณโคจฺฉก}{mtk.40}
\bookmark[dest={mtk.40},level=1]{8. นีวรณโคจฺฉก}
\csromanheader{1}{8. นีวรณ\-โคจฺฉก}
\csromanmatikaanchor{mtk.41}
\csromantocmarkref{subsection}{44. นีวรณทุก}{mtk.41}
\bookmark[dest={mtk.41},level=2]{44. นีวรณทุก}
\csromanmarkh{44. นีวรณทุก}\csromanmarkhii{1. ปฏิจฺจวาร}\csromanheadernomark{2}{44. \textbf{นีวรณทุก 1. ปฏิจฺจวาร}}
\csromanmarkh{1. ปจฺจยานุโลม}\csromanmarkhii{1. วิภงฺควาร}\csromanheadernomark{2}{1. \textbf{ปจฺจยานุโลม 1. วิภงฺควาร}}
\csromanheader{2}{\textbf{เหตุ}}
\proseitem{1}{\csromanfolio{298}นีวรณํ ธมฺมํ ปฏิจฺจ นีวรโณ ธมฺโม อุปฺปชฺชติ เหตุปจฺจยา. กามจฺฉนฺท\-นีวรณํ ปฏิจฺจ ถินมิทฺธ\-นีวรณํ\csromansharedfootnote{fe5b891e8c0e7adc}{ถีนมิทฺธนีวรณํ (สฺยา)}, อุทฺธจฺจ\-นีวรณํ, อวิชฺชา\-นีวรณํ.  กามจฺฉนฺท\-นีวรณํ ปฏิจฺจ อุทฺธจฺจ\-นีวรณํ, อวิชฺชา\-นีวรณํ.  พฺยาปาท\-นีวรณํ ปฏิจฺจ ถินมิทฺธ\-นีวรณํ, อุทฺธจฺจ\-นีวรณํ, อวิชฺชา\-นีวรณํ.  พฺยาปาท\-นีวรณํ ปฏิจฺจ อุทฺธจฺจ\-นีวรณํ, อวิชฺชา\-นีวรณํ.  พฺยาปาท\-นีวรณํ ปฏิจฺจ ถินมิทฺธ\-นีวรณํ, อุทฺธจฺจ\-นีวรณํ, กุกฺกุจฺจ\-นีวรณํ, อวิชฺชา\-นีวรณํ.  พฺยาปาท\-นีวรณํ ปฏิจฺจ อุทฺธจฺจ\-นีวรณํ, กุกฺกุจฺจ\-นีวรณํ, อวิชฺชา\-นีวรณํ.  วิจิกิจฺฉา\-นีวรณํ ปฏิจฺจ อุทฺธจฺจ\-นีวรณํ.  อุทฺธจฺจ\-นีวรณํ ปฏิจฺจ อวิชฺชา\-นีวรณํ. (1)}

\prose{นีวรณํ ธมฺมํ ปฏิจฺจ โนนีวรโณ ธมฺโม อุปฺปชฺชติ เหตุปจฺจยา. นีวรเณ ปฏิจฺจ สมฺปยุตฺตกา ขนฺธา จิตฺต\-สมุฏฺฐานญฺจ รูปํ. (2)}

\prose{นีวรณํ ธมฺมํ ปฏิจฺจ นีวรโณ จ โนนีวรโณ จ ธมฺมา อุปฺปชฺชนฺติ เหตุปจฺจยา. กามจฺฉนฺท\-นีวรณํ ปฏิจฺจ ถินมิทฺธ\-นีวรณํ, อุทฺธจฺจ\-นีวรณํ, อวิชฺชา\-นีวรณํ สมฺปยุตฺตกา จ ขนฺธา จิตฺต\-สมุฏฺฐานญฺจ รูปํ. (จกฺกํ.) (3)}

\proseitem{2}{โนนีวรณํ ธมฺมํ ปฏิจฺจ โนนีวรโณ ธมฺโม อุปฺปชฺชติ เหตุปจฺจยา. โนนีวรณํ เอกํ ขนฺธํ ปฏิจฺจ ตโย ขนฺธา จิตฺต\-สมุฏฺฐานญฺจ รูปํ ฯเปฯ เทฺว ขนฺเธ ฯเปฯ. ปฏิสนฺธิ\-กฺขเณ ฯเปฯ. (1)}

\prose{โนนีวรณํ ธมฺมํ ปฏิจฺจ นีวรโณ ธมฺโม อุปฺปชฺชติ เหตุปจฺจยา. โนนีวรเณ ขนฺเธ ปฏิจฺจ นีวรณา. (2)}

\prose{โนนีวรณํ ธมฺมํ ปฏิจฺจ นีวรโณ จ โนนีวรโณ จ ธมฺมา อุปฺปชฺชนฺติ เหตุปจฺจยา. โนนีวรณํ เอกํ ขนฺธํ ปฏิจฺจ ตโย ขนฺธา นีวรณา จ จิตฺต\-สมุฏฺฐานญฺจ รูปํ ฯเปฯ เทฺว ขนฺเธ ฯเปฯ. (3)}

\proseitem{3}{\csromanfolio{299}นีวรณญฺจ โนนีวรณญฺจ ธมฺมํ ปฏิจฺจ นีวรโณ ธมฺโม อุปฺปชฺชติ เหตุปจฺจยา. กามจฺฉนฺทนีวรณญฺจ สมฺปยุตฺตเก จ ขนฺเธ ปฏิจฺจ ถินมิทฺธ\-นีวรณํ, อุทฺธจฺจ\-นีวรณํ, อวิชฺชา\-นีวรณํ. (จกฺกํ.) (1)}

\prose{นีวรณญฺจ โนนีวรณญฺจ ธมฺมํ ปฏิจฺจ โนนีวรโณ ธมฺโม อุปฺปชฺชติ เหตุปจฺจยา. โนนีวรณํ เอกํ ขนฺธญฺจ นีวรณญฺจ ปฏิจฺจ ตโย ขนฺธา จิตฺต\-สมุฏฺฐานญฺจ รูปํ ฯเปฯ เทฺว ขนฺเธ จ ฯเปฯ. นีวรเณ จ มหาภูเต จ ปฏิจฺจ จิตฺต\-สมุฏฺฐานํ รูปํ. (2)}

\prose{นีวรณญฺจ โนนีวรณญฺจ ธมฺมํ ปฏิจฺจ นีวรโณ จ โนนีวรโณ จ ธมฺมา อุปฺปชฺชนฺติ เหตุปจฺจยา. โนนีวรณํ เอกํ ขนฺธญฺจ กามจฺฉนฺทนีวรณญฺจ ปฏิจฺจ ตโย ขนฺธา ถินมิทฺธ\-นีวรณํ, อุทฺธจฺจ\-นีวรณํ, อวิชฺชา\-นีวรณํ ฯเปฯ เทฺว ขนฺเธ จ ฯเปฯ. (จกฺกํ. สํขิตฺตํ.) (3)}

\csromanmarkh{1. ปจฺจยานุโลม}\csromanmarkhii{2. สงฺขฺยาวาร}\csromanheadernomark{2}{1. \textbf{ปจฺจยานุโลม 2. สงฺขฺยาวาร}}
\proseitem{4}{เหตุยา นว, อารมฺมเณ นว, อธิปติยา นว, อนนฺตเร นว, สมนนฺตเร นว, (สพฺพตฺถ นว,) วิปาเก เอกํ, อาหาเร นว ฯเปฯ อวิคเต นว.}

\vspace{\tipitakaparskip}
\csromancenter{\textbf{อนุโลมํ.}}
\csromanmarkh{2. ปจฺจย\-ปจฺจนีย}\csromanmarkhii{1. วิภงฺควาร}\csromanheadernomark{2}{2. ปจฺจย\-ปจฺจนีย 1. วิภงฺควาร}
\proseitem{5}{นีวรณํ ธมฺมํ ปฏิจฺจ นีวรโณ ธมฺโม อุปฺปชฺชติ นเหตุปจฺจยา. วิจิกิจฺฉา\-นีวรณํ ปฏิจฺจ อวิชฺชา\-นีวรณํ. อุทฺธจฺจ\-นีวรณํ ปฏิจฺจ อวิชฺชา\-นีวรณํ. (1)}

\prose{โนนีวรณํ ธมฺมํ ปฏิจฺจ โนนีวรโณ ธมฺโม อุปฺปชฺชติ นเหตุปจฺจยา. อเหตุกํ โนนีวรณํ เอกํ ขนฺธํ ปฏิจฺจ ตโย ขนฺธา จิตฺต\-สมุฏฺฐานญฺจ รูปํ ฯเปฯ เทฺว ฯเปฯ. อเหตุก\-ปฏิสนฺธิ\-กฺขเณ ฯเปฯ. (ยาว อสญฺญสตฺตา.) (1)}

\prose{\csromanfolio{300}โนนีวรณํ ธมฺมํ ปฏิจฺจ นีวรโณ ธมฺโม อุปฺปชฺชติ นเหตุปจฺจยา. วิจิกิจฺฉา\-สหคเต อุทฺธจฺจ\-สหคเต ขนฺเธ ปฏิจฺจ อวิชฺชา\-นีวรณํ. (2)}

\proseitem{6}{นีวรณญฺจ โนนีวรณญฺจ ธมฺมํ ปฏิจฺจ นีวรโณ ธมฺโม อุปฺปชฺชติ นเหตุปจฺจยา. วิจิกิจฺฉานีวรณญฺจ สมฺปยุตฺตเก จ ขนฺเธ ปฏิจฺจ อวิชฺชา\-นีวรณํ. อุทฺธจฺจ\-นีวรณญฺจ สมฺปยุตฺตเก จ ขนฺเธ ปฏิจฺจ อวิชฺชา\-นีวรณํ. (1)}

\prose{นอารมฺมณาทิ}

\proseitem{7}{นีวรณํ ธมฺมํ ปฏิจฺจ โนนีวรโณ ธมฺโม อุปฺปชฺชติ นอารมฺมณ\-ปจฺจยา. นีวรเณ ปฏิจฺจ จิตฺต\-สมุฏฺฐานํ รูปํ. (1)}

\prose{โนนีวรณํ ธมฺมํ ปฏิจฺจ โนนีวรโณ ธมฺโม อุปฺปชฺชติ นอารมฺมณ\-ปจฺจยา. โนนีวรเณ ขนฺเธ ปฏิจฺจ จิตฺต\-สมุฏฺฐานํ รูปํ. ปฏิสนฺธิ\-กฺขเณ ฯเปฯ. (ยาว อสญฺญสตฺตา.) (1)}

\prose{นีวรณญฺจ โนนีวรณญฺจ ธมฺมํ ปฏิจฺจ โนนีวรโณ ธมฺโม อุปฺปชฺชติ นอารมฺมณ\-ปจฺจยา. นีวรเณ จ สมฺปยุตฺตเก จ ขนฺเธ ปฏิจฺจ จิตฺต\-สมุฏฺฐานํ รูปํ. (สํขิตฺตํ.) นอธิปติ\-ปจฺจยา. นอนนฺตร\-ปจฺจยา. นสมนนฺตร\-ปจฺจยา. นอญฺญมญฺญ\-ปจฺจยา. นอุปนิสฺสยปจฺจยา.}

\csromanheader{2}{\textbf{นปุเรชาต}}
\proseitem{8}{นีวรณํ ธมฺมํ ปฏิจฺจ นีวรโณ ธมฺโม อุปฺปชฺชติ นปุเรชาต\-ปจฺจยา. อรูเป กามจฺฉนฺท\-นีวรณํ ปฏิจฺจ ถินมิทฺธ\-นีวรณํ, อุทฺธจฺจ\-นีวรณํ, อวิชฺชา\-นีวรณํ. อรูเป กามจฺฉนฺท\-นีวรณํ ปฏิจฺจ อุทฺธจฺจ\-นีวรณํ, อวิชฺชา\-นีวรณํ. อรูเป วิจิกิจฺฉา\-นีวรณํ ปฏิจฺจ อุทฺธจฺจ\-นีวรณํ, อวิชฺชา\-นีวรณํ. อรูเป อุทฺธจฺจ\-นีวรณํ ปฏิจฺจ อวิชฺชา\-นีวรณํ. (1)}

\prose{นีวรณํ ธมฺมํ ปฏิจฺจ โนนีวรโณ ธมฺโม อุปฺปชฺชติ นปุเรชาต\-ปจฺจยา. อรูเป นีวรเณ ปฏิจฺจ สมฺปยุตฺตกา ขนฺธา. นีวรเณ ปฏิจฺจ จิตฺต\-สมุฏฺฐานํ รูปํ. (อวเสสา ปญฺหา สพฺเพ วิตฺถาเรตพฺพา. อรูปํ ปฐมํ กาตพฺพํ. รูปํ ปจฺฉา ยถา ลภติ.)}

\proseitem{9}{\csromanfolio{301}นีวรณญฺจ โนนีวรณญฺจ ธมฺมํ ปฏิจฺจ นีวรโณ ธมฺโม อุปฺปชฺชติ นปุเรชาต\-ปจฺจยา. อรูเป โนนีวรเณ ขนฺเธ จ กามจฺฉนฺทนีวรณญฺจ ปฏิจฺจ ถินมิทฺธ\-นีวรณํ, อุทฺธจฺจ\-นีวรณํ. (จกฺกํ.) (1)}

\prose{นีวรณญฺจ โนนีวรณญฺจ ธมฺมํ ปฏิจฺจ โนนีวรโณ ธมฺโม อุปฺปชฺชติ นปุเรชาต\-ปจฺจยา. อรูเป โนนีวรณํ เอกํ ขนฺธญฺจ นีวรเณ จ ปฏิจฺจ ตโย ขนฺธา ฯเปฯ เทฺว ขนฺเธ จ ฯเปฯ. นีวรเณ จ สมฺปยุตฺตเก จ ขนฺเธ ปฏิจฺจ จิตฺต\-สมุฏฺฐานํ รูปํ. นีวรเณ จ มหาภูเต จ ปฏิจฺจ จิตฺต\-สมุฏฺฐานํ รูปํ. (2)}

\prose{นีวรณญฺจ โนนีวรณญฺจ ธมฺมํ ปฏิจฺจ นีวรโณ จ โนนีวรโณ จ ธมฺมา อุปฺปชฺชนฺติ นปุเรชาต\-ปจฺจยา. อรูเป โนนีวรณํ เอกํ ขนฺธญฺจ กมจฺฉนฺทนีวรณญฺจ ปฏิจฺจ ตโย ขนฺธา ถินมิทฺธ\-นีวรณํ, อุทฺธจฺจ\-นีวรณํ, อวิชฺชา\-นีวรณํ. (จกฺกํ. สํขิตฺตํ.) (3)}

\prose{2. ปจฺจย\-ปจฺจนีย 2. สงฺขฺยาวาร}

\proseitem{10}{นเหตุยา จตฺตาริ, นอารมฺมเณ ตีณิ, นอธิปติยา นว, นอนนฺตเร ตีณิ, นสมนนฺตเร ตีณิ, นอญฺญมญฺเญ ตีณิ, นอุปนิสฺสเย ตีณิ, นปุเรชาเต นว, นปจฺฉาชาเต นว, นอาเสวเน นว, นกมฺเม ตีณิ, นวิปาเก นว, นอาหาเร เอกํ, นอินฺทฺริเย เอกํ, นฌาเน เอกํ, นมคฺเค เอกํ, นสมฺปยุตฺเต ตีณิ, นวิปฺปยุตฺเต นว, โนนตฺถิยา ตีณิ, โนวิคเต ตีณิ.}

\proseitem{11}{เหตุปจฺจยา นอารมฺมเณ ตีณิ, นอธิปติยา นว. (สํขิตฺตํ.)}

\vspace{\tipitakaparskip}
\csromancenter{นเหตุปจฺจยา อารมฺมเณ จตฺตาริ ฯเปฯ มคฺเค ตีณิ ฯเปฯ อวิคเต จตฺตาริ.}
\csromancenter{(สหชาตวาโรปิ เอวํ วิตฺถาเรตพฺโพ.)}
\csromanmarkh{44. นีวรณทุก}\csromanmarkhii{3. ปจฺจยวาร}\csromanheadernomark{2}{44. \textbf{นีวรณทุก 3. ปจฺจยวาร}}
\csromanmarkh{1. ปจฺจยานุโลม}\csromanmarkhii{1. วิภงฺควาร}\csromanheadernomark{2}{1. \textbf{ปจฺจยานุโลม 1. วิภงฺควาร}}
\csromanheader{2}{\textbf{เหตุ}}
\proseitem{13}{\csromanfolio{302}นีวรณํ ธมฺมํ ปจฺจยา นีวรโณ ธมฺโม อุปฺปชฺชติ เหตุปจฺจยา. ตีณิ.}

\prose{โนนีวรณํ ธมฺมํ ปจฺจยา โนนีวรโณ ธมฺโม อุปฺปชฺชติ เหตุปจฺจยา. โนนีวรณํ เอกํ ขนฺธํ ปจฺจยา ตโย ขนฺธา จิตฺต\-สมุฏฺฐานญฺจ รูปํ ฯเปฯ. (ยาว อชฺฌตฺติกา มหาภูตา.) วตฺถุํ ปจฺจยา โนนีวรณา ขนฺธา. (1)}

\prose{โนนีวรณํ ธมฺมํ ปจฺจยา นีวรโณ ธมฺโม อุปฺปชฺชติ เหตุปจฺจยา. โนนีวรเณ ขนฺเธ ปจฺจยา นีวรณา. วตฺถุํ ปจฺจยา นีวรณา. (2)}

\prose{โนนีวรณํ ธมฺมํ ปจฺจยา นีวรโณ จ โนนีวรโณ จ ธมฺมา อุปฺปชฺชนฺติ เหตุปจฺจยา. โนนีวรณํ เอกํ ขนฺธํ ปจฺจยา ตโย ขนฺธา นีวรณา จ จิตฺต\-สมุฏฺฐานญฺจ รูปํ ฯเปฯ เทฺว ขนฺเธ ฯเปฯ. วตฺถุํ ปจฺจยา นีวรณา. มหาภูเต ปจฺจยา จิตฺต\-สมุฏฺฐานํ รูปํ. วตฺถุํ ปจฺจยา นีวรณา จ สมฺปยุตฺตกา จ ขนฺธา. (3)}

\proseitem{14}{นีวรณญฺจ โนนีวรณญฺจ ธมฺมํ ปจฺจยา นีวรโณ ธมฺโม อุปฺปชฺชติ เหตุปจฺจยา. กามจฺฉนฺทนีวรณญฺจ สมฺปยุตฺตเก จ ขนฺเธ ปจฺจยา ถินมิทฺธ\-นีวรณํ อุทฺธจฺจ\-นีวรณํ อวิชฺชา\-นีวรณํ. (จกฺกํ.) กามจฺฉนฺทนีวรณญฺจ วตฺถุญฺจ ปจฺจยา ถินมิทฺธ\-นีวรณํ อุทฺธจฺจ\-นีวรณํ อวิชฺชา\-นีวรณํ. (จกฺกํ.) (1)}

\prose{นีวรณญฺจ โนนีวรณญฺจ ธมฺมํ ปจฺจยา โนนีวรโณ ธมฺโม อุปฺปชฺชติ เหตุปจฺจยา. โนนีวรณํ เอกํ ขนฺธญฺจ นีวรณญฺจ ปจฺจยา ตโย ขนฺธา จิตฺต\-สมุฏฺฐานญฺจ รูปํ ฯเปฯ เทฺว ขนฺเธ ฯเปฯ. นีวรณญฺจ วตฺถุญฺจ ปจฺจยา สมฺปยุตฺตกา ขนฺธา. นีวรณญฺจ สมฺปยุตฺตเก จ ขนฺเธ ปจฺจยา จิตฺต\-สมุฏฺฐานํ รูปํ. นีวรเณ จ มหาภูเต จ ปจฺจยา จิตฺต\-สมุฏฺฐานํ รูปํ. (2)}

\prose{นีวรณญฺจ โนนีวรณญฺจ ธมฺมํ ปจฺจยา นีวรโณ จ โนนีวรโณ จ ธมฺมา อุปฺปชฺชนฺติ เหตุปจฺจยา. โนนีวรณํ เอกํ ขนฺธญฺจ กามจฺฉนฺทนีวรณญฺจ ปจฺจยา ตโย ขนฺธา ถินมิทฺธ\-นีวรณํ อุทฺธจฺจ\-นีวรณํ อวิชฺชา\-นีวรณํ ฯเปฯ เทฺว \csromanfolio{303}ขนฺเธ ฯเปฯ. (จกฺกํ.) กามจฺฉนฺทนีวรณญฺจ วตฺถุญฺจ ปจฺจยา ถินมิทฺธ\-นีวรณํ อุทฺธจฺจ\-นีวรณํ อวิชฺชา\-นีวรณํ สมฺปยุตฺตกา จ ขนฺธา. (จกฺกํ. สํขิตฺตํ.) (3)}

\csromanmarkh{1. ปจฺจยานุโลม}\csromanmarkhii{2. สงฺขฺยาวาร}\csromanheadernomark{2}{1. \textbf{ปจฺจยานุโลม 2. สงฺขฺยาวาร}}
\vspace{\tipitakaparskip}
\csromancenter{เหตุยา นว, อารมฺมเณ นว, (สพฺพตฺถ นว,) วิปาเก เอกํ ฯเปฯ อวิคเต นว.}
\csromancenter{อนุโลมํ.}
\csromanmarkh{2. ปจฺจย\-ปจฺจนีย}\csromanmarkhii{1. วิภงฺควาร}\csromanheadernomark{2}{2. ปจฺจย\-ปจฺจนีย 1. วิภงฺควาร}
\proseitem{14}{นีวรณํ ธมฺมํ ปจฺจยา นีวรโณ ธมฺโม อุปฺปชฺชติ นเหตุปจฺจยา. วิจิกิจฺฉา\-นีวรณํ ปจฺจยา อวิชฺชา\-นีวรณํ. อุทฺธจฺจ\-นีวรณํ ปจฺจยา อวิชฺชา\-นีวรณํ. (1)}

\prose{โนนีวรณํ ธมฺมํ ปจฺจยา โนนีวรโณ ธมฺโม อุปฺปชฺชติ นเหตุปจฺจยา. อเหตุกํ โนนีวรณํ เอกํ ขนฺธํ ปจฺจยา ตโย ขนฺธา จิตฺต\-สมุฏฺฐานญฺจ รูปํ ฯเปฯ เทฺว ขนฺเธ ฯเปฯ. (ยาว อสญฺญสตฺตา.) จกฺขายตนํ ปจฺจยา จกฺขุวิญฺญาณํ ฯเปฯ. กายายตนํ ปจฺจยา กายวิญฺญาณํ. วตฺถุํ ปจฺจยา อเหตุกา โนนีวรณา ขนฺธา. (1)}

\prose{โนนีวรณํ ธมฺมํ ปจฺจยา นีวรโณ ธมฺโม อุปฺปชฺชติ นเหตุปจฺจยา. วิจิกิจฺฉา\-สหคเต อุทฺธจฺจ\-สหคเต ขนฺเธ ปจฺจยา อวิชฺชา\-นีวรณํ. วตฺถุํ ปจฺจยา อวิชฺชา\-นีวรณํ. (2)}

\proseitem{17}{นีวรณญฺจ โนนีวรณญฺจ ธมฺมํ ปจฺจยา นีวรโณ ธมฺโม อุปฺปชฺชติ นเหตุปจฺจยา. วิจิกิจฺฉานีวรณญฺจ สมฺปยุตฺตเก จ ขนฺเธ ปจฺจยา อวิชฺชา\-นีวรณํ. อุทฺธจฺจ\-นีวรณญฺจ สมฺปยุตฺตเก จ ขนฺเธ ปจฺจยา อวิชฺชา\-นีวรณํ. วิจิกิจฺฉานีวรณญฺจ วตฺถุญฺจ ปจฺจยา อวิชฺชา\-นีวรณํ. อุทฺธจฺจ\-นีวรณญฺจ วตฺถุญฺจ ปจฺจยา อวิชฺชา\-นีวรณํ. (สํขิตฺตํ.) (1)}

\csromanmarkh{2. ปจฺจย\-ปจฺจนีย}\csromanmarkhii{2. สงฺขฺยาวาร}\csromanheadernomark{2}{2. ปจฺจย\-ปจฺจนีย 2. สงฺขฺยาวาร}
\proseitem{18}{\csromanfolio{304}นเหตุยา จตฺตาริ, นอารมฺมเณ ตีณิ, นอธิปติยา นว, นอนนฺตเร ตีณิ, นสมนนฺตเร ตีณิ, นอญฺญมญฺเญ ตีณิ, นอุปนิสฺสเย ตีณิ, นปุเรชาเต นว, นปจฺฉาชาเต นว, นอาเสวเน นว, นกมฺเม ตีณิ, นวิปาเก นว, นอาหาเร เอกํ, นอินฺทฺริเย เอกํ, นฌาเน เอกํ, นมคฺเค เอกํ, นสมฺปยุตฺเต ตีณิ, นวิปฺปยุตฺเต นว, โนนตฺถิยา ตีณิ, โนวิคเต ตีณิ.}

\proseitem{19}{เหตุปจฺจยา นอารมฺมเณ ตีณิ, นอธิปติยา นว. (สํขิตฺตํ.)}

\proseitem{20}{นเหตุปจฺจยา อารมฺมเณ จตฺตาริ, อนนฺตเร จตฺตาริ, สมนนฺตเร จตฺตาริ ฯเปฯ มคฺเค ตีณิ, อวิคเต จตฺตาริ.}

\csromanmarkh{44. นีวรณทุก}\csromanmarkhii{5. สํสฏฺฐวาร}\csromanheadernomark{2}{44. \textbf{นีวรณทุก 5. สํสฏฺฐวาร}}
\csromanheader{2}{1-4. ปจฺจยจตุกฺก}
\proseitem{21}{นีวรณํ ธมฺมํ สํสฏฺโฐ นีวรโณ ธมฺโม อุปฺปชฺชติ เหตุปจฺจยา. กามจฺฉนฺท\-นีวรณํ สํสฏฺฐํ ถินมิทฺธ\-นีวรณํ อุทฺธจฺจ\-นีวรณํ อวิชฺชา\-นีวรณํ. (จกฺกํ. สพฺพํ นีวรณํ วิตฺถาเรตพฺพํ.)}

\proseitem{22}{เหตุยา นว, อารมฺมเณ นว, (สพฺพตฺถ นว,) วิปาเก เอกํ ฯเปฯ อวิคเต นว.}

\vspace{\tipitakaparskip}
\csromancenter{อนุโลมํ.}
\proseitem{23}{\csromanfolio{305}นีวรณํ ธมฺมํ สํสฏฺโฐ นีวรโณ ธมฺโม อุปฺปชฺชติ นเหตุปจฺจยา. วิจิกิจฺฉา\-นีวรณํ สํสฏฺฐํ อวิชฺชา\-นีวรณํ. อุทฺธจฺจ\-นีวรณํ สํสฏฺฐํ อวิชฺชา\-นีวรณํ. (สํขิตฺตํ.)}

\proseitem{24}{นเหตุยา จตฺตาริ, นอธิปติยา นว, นปุเรชาเต นว, นปจฺฉาชาเต นว, นอาเสวเน นว, นกมฺเม ตีณิ, นวิปาเก นว, นฌาเน เอกํ, นมคฺเค เอกํ, นวิปฺปยุตฺเต นว.}

\vspace{\tipitakaparskip}
\csromancenter{ปจฺจนียํ.}
\csromancenter{(เอวํ อิตเร เทฺว คณนาปิ สมฺปยุตฺตวาโรปิ กาตพฺโพ.)}
\csromanmarkh{44. นีวรณทุก}\csromanmarkhii{7. ปญฺหาวาร}\csromanheadernomark{2}{44. นีวรณทุก 7. ปญฺหาวาร}
\csromanmarkh{1. ปจฺจยานุโลม}\csromanmarkhii{1. วิภงฺควาร}\csromanheadernomark{2}{1. \textbf{ปจฺจยานุโลม 1. วิภงฺควาร}}
\csromanheader{2}{\textbf{เหตุ}}
\proseitem{25}{นีวรโณ ธมฺโม นีวรณสฺส ธมฺมสฺส เหตุปจฺจเยน ปจฺจโย. นีวรณา เหตู สมฺปยุตฺตกานํ นีวรณานํ เหตุปจฺจเยน ปจฺจโย. (1)}

\prose{นีวรโณ ธมฺโม โนนีวรณสฺส ธมฺมสฺส เหตุปจฺจเยน ปจฺจโย. นีวรณา เหตู สมฺปยุตฺตกานํ ขนฺธานํ จิตฺตสมุฏฺฐานานญฺจ รูปานํ เหตุปจฺจเยน ปจฺจโย. (2)}

\prose{นีวรโณ ธมฺโม นีวรณสฺส จ โนนีวรณสฺส จ ธมฺมสฺส เหตุปจฺจเยน ปจฺจโย. นีวรณา เหตู สมฺปยุตฺตกานํ ขนฺธานํ นีวรณานญฺจ จิตฺตสมุฏฺฐานานญฺจ รูปานํ เหตุปจฺจเยน ปจฺจโย. (3)}

\prose{โนนีวรโณ ธมฺโม โนนีวรณสฺส ธมฺมสฺส เหตุปจฺจเยน ปจฺจโย. โนนีวรณา เหตู สมฺปยุตฺตกานํ ขนฺธานํ จิตฺตสมุฏฺฐานานญฺจ รูปานํ เหตุปจฺจเยน ปจฺจโย. ปฏิสนฺธิ\-กฺขเณ ฯเปฯ. (1)}

\proseitem{26}{\csromanfolio{306}นีวรโณ ธมฺโม นีวรณสฺส ธมฺมสฺส อารมฺมณ\-ปจฺจเยน ปจฺจโย. นีวรเณ อารพฺภ นีวรณา อุปฺปชฺชนฺติ. (มูลํ ปุจฺฉิตพฺพํ.) นีวรเณ อารพฺภ โนนีวรณา ขนฺธา อุปฺปชฺชนฺติ. (มูลํ กาตพฺพํ.) นีวรเณ อารพฺภ นีวรณา จ สมฺปยุตฺตกา จ ขนฺธา อุปฺปชฺชนฺติ. (3)}

\proseitem{27}{โนนีวรโณ ธมฺโม โนนีวรณสฺส ธมฺมสฺส อารมฺมณ\-ปจฺจเยน ปจฺจโย. ทานํ ฯเปฯ สีลํ ฯเปฯ อุโปสถกมฺมํ กตฺวา ตํ ปจฺจเวกฺขติ, ปุพฺเพ สุจิณฺณานิ ฯเปฯ ฌานา วุฏฺฐหิตฺวา ฌานํ ฯเปฯ. อริยา มคฺคา วุฏฺฐหิตฺวา มคฺคํ ฯเปฯ. ผลํ ฯเปฯ. นิพฺพานํ ฯเปฯ. นิพฺพานํ โคตฺรภุสฺส, โวทานสฺส, มคฺคสฺส, ผลสฺส, อาวชฺชนาย อารมฺมณ\-ปจฺจเยน ปจฺจโย. อริยา โนนีวรเณ ปหีเน กิเลเส ปจฺจเวกฺขนฺติ. วิกฺขมฺภิเต กิเลเส ฯเปฯ. ปุพฺเพ สมุทาจิณฺเณ ฯเปฯ. จกฺขุํ ฯเปฯ วตฺถุํ โนนีวรเณ ขนฺเธ อนิจฺจโต ฯเปฯ โทมนสฺสํ อุปฺปชฺชติ. ทิพฺเพน จกฺขุนา รูปํ ปสฺสติ. ทิพฺพาย โสตธาตุยา สทฺทํ สุณาติ. เจโตปริย\-ญาเณน โนนีวรณจิตฺตสมงฺคิสฺส จิตฺตํ ชานาติ. อากาสา\-นญฺจา\-ยตนํ วิญฺญาณญฺจา\-ยตนสฺส ฯเปฯ. อากิญฺจญฺญา\-ยตนํ เนว\-สญฺญา\-นา\-สญฺญา\-ยตนสฺส ฯเปฯ. รูปายตนํ จกฺขุ\-วิญฺญาณสฺส ฯเปฯ. โผฏฺฐพฺพา\-ยตนํ กายวิญฺญาณสฺส ฯเปฯ. โนนีวรณา ขนฺธา อิทฺธิวิธ\-ญาณสฺส, เจโตปริย\-ญาณสฺส, ปุพฺเพ\-นิวาสา\-นุสฺสติ\-ญาณสฺส, ยถากมฺมูปคญาณสฺส, อนาคตํสญาณสฺส, อาวชฺชนาย อารมฺมณ\-ปจฺจเยน ปจฺจโย. (1)}

\prose{โนนีวรโณ ธมฺโม นีวรณสฺส ธมฺมสฺส อารมฺมณ\-ปจฺจเยน ปจฺจโย. ทานํ ฯเปฯ สีลํ ฯเปฯ อุโปสถกมฺมํ ฯเปฯ. ปุพฺเพ สุจิณฺณานิ ฯเปฯ. ฌานา ฯเปฯ. จกฺขุํ ฯเปฯ วตฺถุํ โนนีวรเณ ขนฺเธ อสฺสาเทติ อภินนฺทติ, ตํ อารพฺภ ราโค อุปฺปชฺชติ. ทิฏฺฐิ ฯเปฯ. วิจิกิจฺฉา ฯเปฯ. อุทฺธจฺจํ ฯเปฯ. โทมนสฺสํ อุปฺปชฺชติ. (2)}

\prose{โนนีวรโณ ธมฺโม นีวรณสฺส จ โนนีวรณสฺส จ ธมฺมสฺส อารมฺมณ\-ปจฺจเยน ปจฺจโย. ทานํ ฯเปฯ สีลํ ฯเปฯ อุโปสถกมฺมํ ฯเปฯ. ปุพฺเพ สุจิณฺณานิ ฯเปฯ. ฌานา ฯเปฯ. จกฺขุํ ฯเปฯ วตฺถุํ โนนีวรเณ ขนฺเธ อสฺสาเทติ อภินนฺทติ ตํ อารพฺภ นีวรณา จ สมฺปยุตฺตกา จ ขนฺธา อุปฺปชฺชนฺติ. (3)}

\prose{\csromanfolio{307}นีวรโณ จ โนนีวรโณ จ ธมฺมา นีวรณสฺส ธมฺมสฺส อารมฺมณ\-ปจฺจเยน ปจฺจโย. (อารพฺภ กาตพฺพา.)}

\proseitem{28}{นีวรโณ ธมฺโม นีวรณสฺส ธมฺมสฺส อธิปติ\-ปจฺจเยน ปจฺจโย. อารมฺมณา\-ธิปติ\csromandash{}นีวรเณ ครุํ กตฺวา นีวรณา อุปฺปชฺชนฺติ. ตีณิ. (อารมฺมณ\-สทิสํ.) (3)}

\proseitem{29}{โนนีวรโณ ธมฺโม โนนีวรณสฺส ธมฺมสฺส อธิปติ\-ปจฺจเยน ปจฺจโย. อารมฺมณา\-ธิปติ, สหชาตา\-ธิปติ.  อารมฺมณา\-ธิปติ\csromandash{}ทานํ ฯเปฯ สีลํ ฯเปฯ อุโปสถกมฺมํ กตฺวา ตํ ครุํ กตฺวา ปจฺจเวกฺขติ, ปุพฺเพ สุจิณฺณานิ ฯเปฯ. ฌานา ฯเปฯ. อริยา มคฺคา วุฏฺฐหิตฺวา มคฺคํ ฯเปฯ. ผลํ ฯเปฯ. นิพฺพานํ ฯเปฯ. นิพฺพานํ โคตฺรภุสฺส, โวทานสฺส, มคฺคสฺส, ผลสฺส อธิปติ\-ปจฺจเยน ปจฺจโย. จกฺขุํ ฯเปฯ วตฺถุํ โนนีวรเณ ขนฺเธ ครุํ กตฺวา อสฺสาเทติ อภินนฺทติ, ตํ ครุํ กตฺวา ราโค อุปฺปชฺชติ, ทิฏฺฐิ อุปฺปชฺชติ. สหชาตา\-ธิปติ\csromandash{}โนนีวรณา\-ธิปติ สมฺปยุตฺตกานํ ขนฺธานํ จิตฺตสมุฏฺฐานานญฺจ รูปานํ อธิปติ\-ปจฺจเยน ปจฺจโย. (1)}

\prose{โนนีวรโณ ธมฺโม นีวรณสฺส ธมฺมสฺส อธิปติ\-ปจฺจเยน ปจฺจโย. อารมฺมณา\-ธิปติ, สหชาตา\-ธิปติ.  อารมฺมณา\-ธิปติ\csromandash{}ทานํ ฯเปฯ สีลํ ฯเปฯ อุโปสถกมฺมํ ฯเปฯ. ปุพฺเพ ฯเปฯ. ฌานา ฯเปฯ. จกฺขุํ ฯเปฯ วตฺถุํ โนนีวรเณ ขนฺเธ ครุํ กตฺวา อสฺสาเทติ อภินนฺทติ, ตํ ครุํ กตฺวา ราโค อุปฺปชฺชติ, ทิฏฺฐิ อุปฺปชฺชติ. สหชาตา\-ธิปติ\csromandash{} โนนีวรณา\-ธิปติ สมฺปยุตฺตกานํ นีวรณานํ อธิปติ\-ปจฺจเยน ปจฺจโย. (2)}

\prose{โนนีวรโณ ธมฺโม นีวรณสฺส จ โนนีวรณสฺส จ ธมฺมสฺส อธิปติ\-ปจฺจเยน ปจฺจโย. อารมฺมณา\-ธิปติ, สหชาตา\-ธิปติ.  อารมฺมณา\-ธิปติ\csromandash{}ทานํ ฯเปฯ สีลํ ฯเปฯ อุโปสถกมฺมํ ฯเปฯ. ปุพฺเพ ฯเปฯ. ฌานา ฯเปฯ. จกฺขุํ ฯเปฯ วตฺถุํ โนนีวรโณ ขนฺเธ ครุํ กตฺวา อสฺสาเทติ อภินนฺทติ, ตํ ครุํ กตฺวา นีวรณา จ สมฺปยุตฺตกา จ ขนฺธา อุปฺปชฺชนฺติ. สหชาตา\-ธิปติ\csromandash{}โนนีวรณา\-ธิปติ สมฺปยุตฺตกานํ ขนฺธานํ นีวรณานญฺจ จิตฺตสมุฏฺฐานานญฺจ รูปานํ อธิปติ\-ปจฺจเยน ปจฺจโย. (3)}

\prose{\csromanfolio{308}นีวรโณ จ โนนีวรโณ จ ธมฺมา นีวรณสฺส ธมฺมสฺส อธิปติ\-ปจฺจเยน ปจฺจโย. อารมฺมณา\-ธิปติ. ตีณิ. (อารมฺมณา\-ธิปติเยว.)}

\proseitem{30}{นีวรโณ ธมฺโม นีวรณสฺส ธมฺมสฺส อนนฺตร\-ปจฺจเยน ปจฺจโย. ปุริมา ปุริมา นีวรณา ปจฺฉิมานํ ปจฺฉิมานํ นีวรณานํ อนนฺตร\-ปจฺจเยน ปจฺจโย. (มูลํ ปุจฺฉิตพฺพํ.) ปุริมา ปุริมา นีวรณา ปจฺฉิมานํ ปจฺฉิมานํ โนนีวรณานํ ขนฺธานํ อนนฺตร\-ปจฺจเยน ปจฺจโย. นีวรณา วุฏฺฐานสฺส อนนฺตร\-ปจฺจเยน ปจฺจโย. (มูลํ ปุจฺฉิตพฺพํ.) ปุริมา ปุริมา นีวรณา ปจฺฉิมานํ ปจฺฉิมานํ นีวรณานํ สมฺปยุตฺตกานญฺจ ขนฺธานํ อนนฺตร\-ปจฺจเยน ปจฺจโย. (3)}

\proseitem{31}{โนนีวรโณ ธมฺโม โนนีวรณสฺส ธมฺมสฺส อนนฺตร\-ปจฺจเยน ปจฺจโย. ปุริมา ปุริมา โนนีวรณา ขนฺธา ปจฺฉิมานํ ปจฺฉิมานํ โนนีวรณานํ ขนฺธานํ อนนฺตร\-ปจฺจเยน ปจฺจโย ฯเปฯ. นิโรธา วุฏฺฐหนฺตสฺส เนว\-สญฺญา\-นา\-สญฺญา\-ยตนํ ผลสมาปตฺติยา อนนฺตร\-ปจฺจเยน ปจฺจโย. (มูลํ ปุจฺฉิตพฺพํ.) ปุริมา ปุริมา โนนีวรณา ขนฺธา ปจฺฉิมานํ ปจฺฉิมานํ นีวรณานํ อนนฺตร\-ปจฺจเยน ปจฺจโย. อาวชฺชนา นีวรณานํ อนนฺตร\-ปจฺจเยน ปจฺจโย. (มูลํ ปุจฺฉิตพฺพํ.) ปุริม ปุริมา โนนีวรณา ขนฺธา ปจฺฉิมานํ ปจฺฉิมานํ นีวรณานํ สมฺปยุตฺตกานญฺจ ขนฺธานํ อนนฺตร\-ปจฺจเยน ปจฺจโย. อาวชฺชนา นีวรณานํ สมฺปยุตฺตกานญฺจ ขนฺธานํ อนนฺตร\-ปจฺจเยน ปจฺจโย. (3)}

\proseitem{32}{นีวรโณ จ โนนีวรโณ จ ธมฺมา นีวรณสฺส ธมฺมสฺส อนนฺตร\-ปจฺจเยน ปจฺจโย. ปุริมา ปุริมา นีวรณา จ สมฺปยุตฺตกา จ ขนฺธา ปจฺฉิมานํ ปจฺฉิมานํ นีวรณานํ อนนฺตร\-ปจฺจเยน ปจฺจโย. (มูลํ ปุจฺฉิตพฺพํ.) ปุริมา ปุริมา นีวรณา จ สมฺปยุตฺตกา จ ขนฺธา ปจฺฉิมานํ ปจฺฉิมานํ โนนีวรณานํ ขนฺธานํ อนนฺตร\-ปจฺจเยน ปจฺจโย. นีวรณา จ สมฺปยุตฺตกา จ ขนฺธา วุฏฺฐานสฺส อนนฺตร\-ปจฺจเยน ปจฺจโย. (มูลํ ปุจฺฉิตพฺพํ.) ปุริมา ปุริมา นีวรณา จ สมฺปยุตฺตกา จ ขนฺธา ปจฺฉิมานํ ปจฺฉิมานํ นีวรณานญฺจ สมฺปยุตฺตกานญฺจ ขนฺธานํ อนนฺตร\-ปจฺจเยน ปจฺจโย. (3)}

\csromanheader{2}{44. นีวรณทุก \textbf{สมนนฺตราทิ}}
\proseitem{33}{\csromanfolio{309}นีวรโณ ธมฺโม นีวรณสฺส ธมฺมสฺส สมนนฺตร\-ปจฺจเยน ปจฺจโย. สหชาต\-ปจฺจเยน ปจฺจโย. อญฺญมญฺญ\-ปจฺจเยน ปจฺจโย. นิสฺสย\-ปจฺจเยน ปจฺจโย.}

\csromanheader{2}{\textbf{อุปนิสฺสย}}
\proseitem{34}{นีวรโณ ธมฺโม นีวรณสฺส ธมฺมสฺส อุปนิสฺสย\-ปจฺจเยน ปจฺจโย. อารมฺมณู\-ปนิสฺสโย, อนนฺตรู\-ปนิสฺสโย, ปกตู\-ปนิสฺสโย ฯเปฯ. ปกตู\-ปนิสฺสโย\csromandash{}นีวรณานิ นีวรณานํ อุปนิสฺสย\-ปจฺจเยน ปจฺจโย. ตีณิ.}

\proseitem{35}{โนนีวรโณ ธมฺโม โนนีวรณสฺส ธมฺมสฺส อุปนิสฺสย\-ปจฺจเยน ปจฺจโย.  อารมฺมณู\-ปนิสฺสโย, อนนฺตรู\-ปนิสฺสโย, ปกตู\-ปนิสฺสโย ฯเปฯ. ปกตู\-ปนิสฺสโย\csromandash{}สทฺธํ อุปนิสฺสาย ทานํ เทติ ฯเปฯ สีลํ ฯเปฯ อุโปสถกมฺมํ ฯเปฯ. ฌานํ อุปฺปาเทติ, วิปสฺสนํ ฯเปฯ มคฺคํ ฯเปฯ อภิญฺญํ ฯเปฯ สมาปตฺติํ อุปฺปาเทติ, มานํ ชปฺเปติ, ทิฏฺฐิํ คณฺหาติ. สีลํ ฯเปฯ. ปญฺญํ ฯเปฯ มานํ. ทิฏฺฐิํ. ปตฺถนํ. กายิกํ สุขํ. กายิกํ ทุกฺขํ. อุตุํ. โภชนํ. เสนาสนํ อุปนิสฺสาย ทานํ เทติ ฯเปฯ สํฆํ ภินฺทติ. สทฺธา ฯเปฯ. เสนาสนํ สทฺธาย ฯเปฯ ปญฺญาย. ราคสฺส ฯเปฯ ปตฺถนาย. กายิกสฺส สุขสฺส, กายิกสฺส ทุกฺขสฺส, มคฺคสฺส, ผลสมาปตฺติยา อุปนิสฺสย\-ปจฺจเยน ปจฺจโย. (1)}

\prose{โนนีวรโณ ธมฺโม นีวรณสฺส ธมฺมสฺส อุปนิสฺสย\-ปจฺจเยน ปจฺจโย. (ตีณิ อุปนิสฺสยา.) ปกตู\-ปนิสฺสโย\csromandash{}สทฺธํ อุปนิสฺสาย มานํ ชปฺเปติ, ทิฏฺฐิํ คณฺหาติ. สีลํ ฯเปฯ. เสนาสนํ อุปนิสฺสาย ปาณํ หนติ ฯเปฯ สํฆํ ภินฺทติ. สทฺธา ฯเปฯ. เสนาสนํ ราคสฺส ฯเปฯ ปตฺถนาย อุปนิสฺสย\-ปจฺจเยน ปจฺจโย. (2)}

\prose{โนนีวรโณ ธมฺโม นีวรณสฺส จ โนนีวรณสฺส จ ธมฺมสฺส อุปนิสฺสย\-ปจฺจเยน ปจฺจโย. อารมฺมณู\-ปนิสฺสโย, อนนฺตรู\-ปนิสฺสโย, ปกตู\-ปนิสฺสโย ฯเปฯ. ปกตู\-ปนิสฺสโย\csromandash{}สทฺธํ อุปนิสฺสาย มานํ ชปฺเปติ, ทิฏฺฐิํ คณฺหาติ. สีลํ ฯเปฯ. เสนาสนํ อุปนิสฺสาย ปาณํ หนติ ฯเปฯ สํฆํ ภินฺทติ. สทฺธา ฯเปฯ. เสนาสนํ นีวรณานํ สมฺปยุตฺตกานญฺจ ขนฺธานํ อุปนิสฺสย\-ปจฺจเยน ปจฺจโย. (3)}

\prose{\csromanfolio{310}นีวรโณ จ โนนีวรโณ จ ธมฺมา นีวรณสฺส ธมฺมสฺส อุปนิสฺสย\-ปจฺจเยน ปจฺจโย. อารมฺมณู\-ปนิสฺสโย, อนนฺตรู\-ปนิสฺสโย, ปกตู\-ปนิสฺสโย ฯเปฯ. ปกตู\-ปนิสฺสโย\csromandash{}นีวรณา จ สมฺปยุตฺตกา จ ขนฺธา นีวรณานํ อุปนิสฺสย\-ปจฺจเยน ปจฺจโย. ตีณิ.}

\proseitem{36}{โนนีวรโณ ธมฺโม โนนีวรณสฺส ธมฺมสฺส ปุเรชาต\-ปจฺจเยน ปจฺจโย. อารมฺมณ\-ปุเรชาตํ, วตฺถุ\-ปุเรชาตํ.  อารมฺมณ\-ปุเรชาตํ\csromandash{} จกฺขุํ ฯเปฯ วตฺถุํ อนิจฺจโต ฯเปฯ โทมนสฺสํ อุปฺปชฺชติ ฯเปฯ. ตีณิ.}

\vspace{\tipitakaparskip}
\csromancenter{(ปุเรชาตํ อารมฺมณ\-สทิสํ, กุสลา\-กุสลสฺส วิภชิตพฺพํ.)}
\csromanheader{2}{\textbf{ปจฺฉาชาต อาเสวน}}
\proseitem{37}{นีวรโณ ธมฺโม โนนีวรณสฺส ธมฺมสฺส ปจฺฉาชาต\-ปจฺจเยน ปจฺจโย. ตีณิ. อาเสวน\-ปจฺจเยน ปจฺจโย. นว.}

\proseitem{38}{โนนีวรโณ ธมฺโม โนนีวรณสฺส ธมฺมสฺส กมฺมปจฺจเยน ปจฺจโย. สหชาตา, นานากฺขณิกา.  สหชาตา โนนีวรณา เจตนา สมฺปยุตฺตกานํ ขนฺธานํ จิตฺตสมุฏฺฐานานญฺจ รูปานํ กมฺมปจฺจเยน ปจฺจโย. ปฏิสนฺธิ\-กฺขเณ ฯเปฯ. นานากฺขณิกา โนนีวรณา เจตนา วิปากานํ ขนฺธานํ กฏตฺตา จ รูปานํ กมฺมปจฺจเยน ปจฺจโย. (มูลํ ปุจฺฉิตพฺพํ.) โนนีวรณา เจตนา สมฺปยุตฺตกานํ นีวรณานํ กมฺมปจฺจเยน ปจฺจโย. (มูลํ ปุจฺฉิตพฺพํ.) โนนีวรณา เจตนา สมฺปยุตฺตกานํ ขนฺธานํ นีวรณานญฺจ จิตฺตสมุฏฺฐานานญฺจ รูปานํ กมฺมปจฺจเยน ปจฺจโย. (3)}

\csromanheader{2}{\textbf{วิปากาทิ}}
\proseitem{39}{โนนีวรโณ ธมฺโม โนนีวรณสฺส ธมฺมสฺส วิปาก\-ปจฺจเยน ปจฺจโย. เอกํ. อาหารปจฺจเยน ปจฺจโย. อินฺทฺริย\-ปจฺจเยน ปจฺจโย. ฌานปจฺจเยน ปจฺจโย. มคฺคปจฺจเยน ปจฺจโย. สมฺปยุตฺต\-ปจฺจเยน ปจฺจโย.}

\csromanheader{2}{44. นีวรณทุก \textbf{วิปฺปยุตฺต}}
\proseitem{40}{\csromanfolio{311}นีวรโณ ธมฺโม โนนีวรณสฺส ธมฺมสฺส วิปฺปยุตฺต\-ปจฺจเยน ปจฺจโย. สหชาตํ, ปจฺฉาชาตํ. (เอวํ อวเสสา จตฺตาริ ปญฺหา กาตพฺพา.)}

\proseitem{41}{นีวรโณ ธมฺโม นีวรณสฺส ธมฺมสฺส อตฺถิปจฺจเยน ปจฺจโย. กามจฺฉนฺท\-นีวรณํ ถินมิทฺธ\-นีวรณสฺส, อุทฺธจฺจ\-นีวรณสฺส, อวิชฺชา\-นีวรณสฺส อตฺถิปจฺจเยน ปจฺจโย. (จกฺกํ.) (1)}

\prose{นีวรโณ ธมฺโม โนนีวรณสฺส ธมฺมสฺส อตฺถิปจฺจเยน ปจฺจโย. สหชาตํ, ปจฺฉาชาตํ. (เอวํ นีวรณมูเล ตีณิ.) (3)}

\proseitem{42}{โนนีวรโณ ธมฺโม โนนีวรณสฺส ธมฺมสฺส อตฺถิปจฺจเยน ปจฺจโย. สหชาตํ, ปุเรชาตํ, ปจฺฉาชาตํ, อาหารํ, อินฺทฺริยํ. (สํขิตฺตํ.) (1)}

\prose{โนนีวรโณ ธมฺโม นีวรณสฺส ธมฺมสฺส อตฺถิปจฺจเยน ปจฺจโย. สหชาตํ ปุเรชาตํ. (สํขิตฺตํ.) (2)}

\prose{โนนีวรโณ ธมฺโม นีวรณสฺส จ โนนีวรณสฺส จ ธมฺมสฺส อตฺถิปจฺจเยน ปจฺจโย. สหชาตํ, ปุเรชาตํ. (สํขิตฺตํ.) (3)}

\proseitem{43}{นีวรโณ จ โนนีวรโณ จ ธมฺมา นีวรณสฺส ธมฺมสฺส อตฺถิปจฺจเยน ปจฺจโย. สหชาตํ, ปุเรชาตํ.  สหชาตํ กามจฺฉนฺทนีวรณญฺจ สมฺปยุตฺตกา จ ขนฺธา ถินมิทฺธ\-นีวรณสฺส, อุทฺธจฺจ\-นีวรณสฺส, อวิชฺชา\-นีวรณสฺส อตฺถิปจฺจเยน ปจฺจโย. กามจฺฉนฺทนีวรณญฺจ วตฺถุ จ ถินมิทฺธ\-นีวรณสฺส, อุทฺธจฺจ\-นีวรณสฺส, อวิชฺชา\-นีวรณสฺส อตฺถิปจฺจเยน ปจฺจโย. (1)}

\prose{นีวรโณ จ โนนีวรณฺโอ จ ธมฺมา โนนีวรณสฺส ธมฺมสฺส อตฺถิปจฺจเยน ปจฺจโย. สหชาตํ, ปุเรชาตํ, ปจฺฉาชาตํ, อาหารํ, อินฺทฺริยํ.  สหชาโต โนนีวรโณ เอโก ขนฺโธ จ นีวรณา จ ติณฺณนฺนํ ขนฺธานํ จิตฺตสมุฏฺฐานานญฺจ รูปานํ อตฺถิปจฺจเยน ปจฺจโย ฯเปฯ เทฺว ขนฺธา ฯเปฯ. นีวรณา จ วตฺถุ จ โนนีวรณานํ ขนฺธานํ อตฺถิปจฺจเยน ปจฺจโย. นีวรณา จ สมฺปยุตฺตกา จ ขนฺธา จิตฺต\-สมุฏฺฐานานํ รูปานํ อตฺถิปจฺจเยน \csromanfolio{312}ปจฺจโย. นีวรณา จ มหาภูตา จ จิตฺต\-สมุฏฺฐานานํ รูปานํ อตฺถิปจฺจเยน ปจฺจโย. ปจฺฉาชาตา นีวรณา จ สมฺปยุตฺตกา จ ขนฺธา ปุเรชาตสฺส อิมสฺส กายสฺส อตฺถิปจฺจเยน ปจฺจโย. ปจฺฉาชาตา นีวรณา จ สมฺปยุตฺตกา ขนฺธา จ กพฬีกาโร อาหาโร จ อิมสฺส กายสฺส อตฺถิปจฺจเยน ปจฺจโย. ปจฺฉาชาตา นีวรณา จ สมฺปยุตฺตกา ขนฺธา จ รูปชีวิตินฺทฺริยญฺจ กฏตฺตารูปานํ อตฺถิปจฺจเยน ปจฺจโย. (2)}

\prose{นีวรโณ จ โนนีวรโณ จ ธมฺมา นีวรณสฺส จ โนนีวรณสฺส จ ธมฺมสฺส อตฺถิปจฺจเยน ปจฺจโย. สหชาตํ, ปุเรชาตํ.  สหชาโต โนนีวรโณ เอโก ขนฺโธ จ กามจฺฉนฺทนีวรณญฺจ ติณฺณนฺนํ ขนฺธานํ ถินมิทฺธ\-นีวรณสฺส, อุทฺธจฺจ\-นีวรณสฺส, อวิชฺชา\-นีวรณสฺส จ จิตฺตสมุฏฺฐานานญฺจ รูปานํ อตฺถิปจฺจเยน ปจฺจโย. เทฺว ขนฺธา จ ฯเปฯ. กามจฺฉนฺทนีวรณญฺจ วตฺถุ จ ถินมิทฺธ\-นีวรณสฺส, อุทฺธจฺจ\-นีวรณสฺส, อวิชฺชา\-นีวรณสฺส จ สมฺปยุตฺตกานญฺจ ขนฺธานํ อตฺถิปจฺจเยน ปจฺจโย. (จกฺกํ.) (3)}

\csromanmarkh{1. ปจฺจยานุโลม}\csromanmarkhii{2. สงฺขฺยาวาร}\csromanheadernomark{2}{1. \textbf{ปจฺจยานุโลม 2. สงฺขฺยาวาร}}
\proseitem{43}{เหตุยา จตฺตาริ, อารมฺมเณ นว, อธิปติยา นว, อนนฺตเร นว, สมนนฺตเร นว, สหชาเต นว, อญฺญมญฺเญ นว, นิสฺสเย นว, อุปนิสฺสเย นว, ปุเรชาเต ตีณิ, ปจฺฉาชาเต ตีณิ, อาเสวเน นว, กมฺเม ตีณิ, วิปาเก เอกํ, อาหาเร ตีณิ, อินฺทฺริเย ตีณิ, ฌาเน ตีณิ, มคฺเค ตีณิ, สมฺปยุตฺเต นว, วิปฺปยุตฺเต ปญฺจ, อตฺถิยา นว, นตฺถิยา นว, วิคเต นว, อวิคเต นว.}

\vspace{\tipitakaparskip}
\csromancenter{อนุโลมํ.}
\proseitem{45}{นีวรโณ ธมฺโม นีวรณสฺส ธมฺมสฺส อารมฺมณ\-ปจฺจเยน ปจฺจโย. สหชาต\-ปจฺจเยน ปจฺจโย. อุปนิสฺสย\-ปจฺจเยน ปจฺจโย. (1)}

\prose{\csromanfolio{313}นีวรโณ ธมฺโม โนนีวรณสฺส ธมฺมสฺส อารมฺมณ\-ปจฺจเยน ปจฺจโย. สหชาต\-ปจฺจเยน ปจฺจโย. อุปนิสฺสย\-ปจฺจเยน ปจฺจโย. ปจฺฉาชาต\-ปจฺจเยน ปจฺจโย. (2)}

\prose{นีวรโณ ธมฺโม นีวรณสฺส จ โนนีวรณสฺส จ ธมฺมสฺส อารมฺมณ\-ปจฺจเยน ปจฺจโย. สหชาต\-ปจฺจเยน ปจฺจโย. อุปนิสฺสย\-ปจฺจเยน ปจฺจโย. (3)}

\proseitem{46}{โนนีวรโณ ธมฺโม โนนีวรณสฺส ธมฺมสฺส อารมฺมณ\-ปจฺจเยน ปจฺจโย. สหชาต\-ปจฺจเยน ปจฺจโย. อุปนิสฺสย\-ปจฺจเยน ปจฺจโย. ปุเรชาต\-ปจฺจเยน ปจฺจโย. ปจฺฉาชาต\-ปจฺจเยน ปจฺจโย. กมฺมปจฺจเยน ปจฺจโย. อาหารปจฺจเยน ปจฺจโย. อินฺทฺริย\-ปจฺจเยน ปจฺจโย. (1)}

\prose{โนนีวรโณ ธมฺโม นีวรณสฺส ธมฺมสฺส อารมฺมณ\-ปจฺจเยน ปจฺจโย. สหชาต\-ปจฺจเยน ปจฺจโย. อุปนิสฺสย\-ปจฺจเยน ปจฺจโย. ปุเรชาต\-ปจฺจเยน ปจฺจโย. (2)}

\prose{โนนีวรโณ ธมฺโม นีวรณสฺส จ โนนีวรณสฺส จ ธมฺมสฺส อารมฺมณ\-ปจฺจเยน ปจฺจโย. สหชาต\-ปจฺจเยน ปจฺจโย. อุปนิสฺสย\-ปจฺจเยน ปจฺจโย. ปุเรชาต\-ปจฺจเยน ปจฺจโย. (3)}

\proseitem{47}{นีวรโณ จ โนนีวรโณ จ ธมฺมา นีวรณสฺส ธมฺมสฺส อารมฺมณ\-ปจฺจเยน ปจฺจโย. สหชาต\-ปจฺจเยน ปจฺจโย. อุปนิสฺสย\-ปจฺจเยน ปจฺจโย. (1)}

\prose{นีวรโณ จ โนนีวรโณ จ ธมฺมา โนนีวรณสฺส ธมฺมสฺส อารมฺมณ\-ปจฺจเยน ปจฺจโย. สหชาต\-ปจฺจเยน ปจฺจโย. อุปนิสฺสย\-ปจฺจเยน ปจฺจโย. ปจฺฉาชาต\-ปจฺจเยน ปจฺจโย. (2)}

\prose{นีวรโณ จ โนนีวรโณ จ ธมฺมา นีวรณสฺส จ โนนีวรณสฺส จ ธมฺมสฺส อารมฺมณ\-ปจฺจเยน ปจฺจโย. สหชาต\-ปจฺจเยน ปจฺจโย. อุปนิสฺสย\-ปจฺจเยน ปจฺจโย. (3)}

\csromanmarkh{2. ปจฺจย\-ปจฺจนีย}\csromanmarkhii{2. สงฺขฺยาวาร}\csromanheadernomark{2}{2. ปจฺจย\-ปจฺจนีย 2. สงฺขฺยาวาร}
\proseitem{48}{นเหตุยา นว, นอารมฺมเณ นว, นอธิปติยา นว, (สพฺพตฺถ นว.) โนวิคเต นว, โนอวิคเต นว.}

\vspace{\tipitakaparskip}
\csromancenter{\textbf{เหตุ}ปจฺจยา นอารมฺมเณ จตฺตาริ ฯเปฯ นสมนนฺตเร จตฺตาริ, นอญฺญมญฺเญ เทฺว, นฺเอาปนิสฺสเย จตฺตาริ ฯเปฯ นมคฺเค จตฺตาริ, นสมฺปยุตฺเต เทฺว, นวิปฺปยุตฺเต จตฺตาริ, โนนตฺถิยา จตฺตาริ, โนวิคเต จตฺตาริ.}
\proseitemwithcloser{50}{\csromanfolio{314}นเหตุปจฺจยา อารมฺมเณ นว, อธิปติยา นว, (อนุโลมมาติกา.) ฯเปฯ อวิคเต นว.}{นีวรณทุกํ นิฏฺฐิตํ.}
\csromanmatikaanchor{mtk.42}
\csromantocmarkref{subsection}{45. นีวรณิยทุก}{mtk.42}
\bookmark[dest={mtk.42},level=2]{45. นีวรณิยทุก}
\csromanmarkh{45. นีวรณิยทุก}\csromanmarkhii{1. ปฏิจฺจวาร}\csromanheadernomark{2}{45. \textbf{นีวรณิยทุก 1. ปฏิจฺจวาร}}
\proseitem{51}{นีวรณิยํ ธมฺมํ ปฏิจฺจ นีวรณิโย ธมฺโม อุปฺปชฺชติ เหตุปจฺจยา ฯเปฯ. (นีวรณิยทุกํ ยถา โลกิยทุกํ, เอวํ กาตพฺพํ. นินฺนานากรณํ.)}

\setlength{\csromangathapagewidth}{0pt}
\setlength{\csromangathawakwidth}{0pt}
\csromangathameasuregroup{ทฺวิกฺขตฺตุํ กามจฺฉนฺเทน, \\ อุทฺธจฺจํ วิจิกิจฺฉา จ,}{\csromangathabat{ทฺวิกฺขตฺตุํ กามจฺฉนฺเทน,}{จตุกฺขตฺตุํ ปฏิเฆน จ.} \\ \csromangathabat{อุทฺธจฺจํ วิจิกิจฺฉา จ,}{อุโภเปเต สกิํ สกิํ.}}
\csromangathasetleft{ทฺวิกฺขตฺตุํ กามจฺฉนฺเทน, \\ อุทฺธจฺจํ วิจิกิจฺฉา จ,}
\csromangathagroup{\csromangathastanza{\csromangathabat{ทฺวิกฺขตฺตุํ กามจฺฉนฺเทน,}{จตุกฺขตฺตุํ ปฏิเฆน จ.} \\ \csromangathabat{อุทฺธจฺจํ วิจิกิจฺฉา จ,}{อุโภเปเต สกิํ สกิํ.}}}

\prosewithcloser{นีวรณานํ นีวรเณหิ, อฏฺฐวิธํ ปโยชนํ. (นีวรณ\-ทุกสฺส มาติกา อิธ กตา.)}{นีวรณิยทุกํ นิฏฺฐิตํ.}
\csromanmatikaanchor{mtk.43}
\csromantocmarkref{subsection}{46. นีวรณสมฺปยุตฺตทุก}{mtk.43}
\bookmark[dest={mtk.43},level=2]{46. นีวรณสมฺปยุตฺตทุก}
\csromanmarkh{46. นีวรณ\-สมฺปยุตฺตทุก}\csromanmarkhii{1. ปฏิจฺจวาร}\csromanheadernomark{2}{46. นีวรณ\-สมฺปยุตฺตทุก 1. ปฏิจฺจวาร}
\csromanmarkh{1. ปจฺจยานุโลม}\csromanmarkhii{1. วิภงฺควาร}\csromanheadernomark{2}{1. \textbf{ปจฺจยานุโลม 1. วิภงฺควาร}}
\csromanheader{2}{\textbf{เหตุ}}
\proseitem{52}{นีวรณ\-สมฺปยุตฺตํ ธมฺมํ ปฏิจฺจ นีวรณ\-สมฺปยุตฺโต ธมฺโม อุปฺปชฺชติ เหตุปจฺจยา. นีวรณ\-สมฺปยุตฺตํ เอกํ ขนฺธํ ปฏิจฺจ ตโย ขนฺธา ฯเปฯ เทฺว ขนฺเธ ฯเปฯ. (1)}

\prose{\csromanfolio{315}นีวรณ\-สมฺปยุตฺตํ ธมฺมํ ปฏิจฺจ นีวรณ\-วิปฺปยุตฺโต ธมฺโม อุปฺปชฺชติ เหตุปจฺจยา. นีวรณ\-สมฺปยุตฺเต ขนฺเธ ปฏิจฺจ จิตฺต\-สมุฏฺฐานํ รูปํ. (2)}

\prose{นีวรณ\-สมฺปยุตฺตํ ธมฺมํ ปฏิจฺจ นีวรณ\-สมฺปยุตฺโต จ นีวรณ\-วิปฺปยุตฺโต จ ธมฺมา อุปฺปชฺชนฺติ เหตุปจฺจยา. นีวรณ\-สมฺปยุตฺตํ เอกํ ขนฺธํ ปฏิจฺจ ตโย ขนฺธา จิตฺต\-สมุฏฺฐานญฺจ รูปํ ฯเปฯ เทฺว ขนฺเธ ฯเปฯ. (3)}

\proseitem{53}{นีวรณ\-วิปฺปยุตฺตํ ธมฺมํ ปฏิจฺจ นีวรณ\-วิปฺปยุตฺโต ธมฺโม อุปฺปชฺชติ เหตุปจฺจยา. นีวรณ\-วิปฺปยุตฺตํ เอกํ ขนฺธํ ปฏิจฺจ ตโย ขนฺธา จิตฺต\-สมุฏฺฐานญฺจ รูปํ ฯเปฯ เทฺว ขนฺเธ ฯเปฯ. ปฏิสนฺธิ\-กฺขเณ ฯเปฯ. (ยาว อชฺฌตฺติกา มหาภูตา.) (1)}

\prose{นีวรณ\-สมฺปยุตฺตญฺจ นีวรณ\-วิปฺปยุตฺตญฺจ ธมฺมํ ปฏิจฺจ นีวรณ\-วิปฺปยุตฺโต ธมฺโม อุปฺปชฺชติ เหตุปจฺจยา. นีวรณ\-สมฺปยุตฺเต ขนฺเธ จ มหาภูเต จ ปฏิจฺจ จิตฺต\-สมุฏฺฐานํ รูปํ. (สํขิตฺตํ.)}

\csromanmarkh{1. ปจฺจยานุโลม}\csromanmarkhii{2. สงฺขฺยาวาร}\csromanheadernomark{2}{1. \textbf{ปจฺจยานุโลม 2. สงฺขฺยาวาร}}
\proseitem{54}{เหตุยา ปญฺจ, อารมฺมเณ เทฺว, อธิปติยา ปญฺจ ฯเปฯ นตฺถิยา เทฺว ฯเปฯ อวิคเต ปญฺจ.}

\vspace{\tipitakaparskip}
\csromancenter{อนุโลมํ.}
\csromanmarkh{2. ปจฺจย\-ปจฺจนีย}\csromanmarkhii{1. วิภงฺควาร}\csromanheadernomark{2}{2. ปจฺจย\-ปจฺจนีย 1. วิภงฺควาร}
\proseitem{55}{นีวรณ\-สมฺปยุตฺตํ ธมฺมํ ปฏิจฺจ นีวรณ\-สมฺปยุตฺโต ธมฺโม อุปฺปชฺชติ นเหตุปจฺจยา. วิจิกิจฺฉา\-สหคเต อุทฺธจฺจ\-สหคเต ขนฺเธ ปฏิจฺจ อวิชฺชา\-นีวรณํ. (1)}

\prose{นีวรณ\-วิปฺปยุตฺตํ ธมฺมํ ปฏิจฺจ นีวรณ\-วิปฺปยุตฺโต ธมฺโม อุปฺปชฺชติ นเหตุปจฺจยา. อเหตุกํ นีวรณ\-วิปฺปยุตฺตํ ฯเปฯ. (ยาว อสญฺญสตฺตา. สํขิตฺตํ.)}

\csromanmarkh{2. ปจฺจย\-ปจฺจนีย}\csromanmarkhii{2. สงฺขฺยาวาร}\csromanheadernomark{2}{2. ปจฺจย\-ปจฺจนีย 2. สงฺขฺยาวาร}
\proseitem{56}{\csromanfolio{316}นเหตุยา เทฺว, นอารมฺมเณ ตีณิ, นอธิปติยา ปญฺจ, นอนนฺตเร, นสมนนฺตเร, นอญฺญมญฺเญ, นอุปนิสฺสเย ตีณิ, นปุเรชาเต จตฺตาริ, นปจฺฉาชาเต ปญฺจ, นอาเสวเน ปญฺจ, นกมฺเม เทฺว, นวิปาเก ปญฺจ, นอาหาเร เอกํ, นอินฺทฺริเย เอกํ, นฌาเน เอกํ, นมคฺเค เอกํ, นสมฺปยุตฺเต ตีณิ, นวิปฺปยุตฺเต เทฺว, โนนตฺถิยา ตีณิ, โนวิคเต ตีณิ.}

\vspace{\tipitakaparskip}
\csromancenter{\textbf{เหตุ}ปจฺจยา นอารมฺมเณ ตีณิ, นอธิปติยา ปญฺจ, นปุเรชาเต จตฺตาริ ฯเปฯ นกมฺเม เทฺว, นวิปาเก ปญฺจ, นสมฺปยุตฺเต ตีณิ, นวิปฺปยุตฺเต เทฺว, โนนตฺถิยา ตีณิ, โนวิคเต ตีณิ.}
\csromancenter{นเหตุปจฺจยา อารมฺมเณ เทฺว ฯเปฯ มคฺเค เอกํ ฯเปฯ อวิคเต เทฺว.}
\csromancenter{(สหชาตวาโรปิ เอวํ กาตพฺโพ.)}
\csromanmarkh{46. นีวรณ\-สมฺปยุตฺตทุก}\csromanmarkhii{3. ปจฺจยวาร}\csromanheadernomark{2}{46. นีวรณ\-สมฺปยุตฺตทุก 3. ปจฺจยวาร}
\csromanmarkh{1. ปจฺจยานุโลม}\csromanmarkhii{1. วิภงฺควาร}\csromanheadernomark{2}{1. \textbf{ปจฺจยานุโลม 1. วิภงฺควาร}}
\csromanheader{2}{\textbf{เหตุ}}
\proseitem{59}{นีวรณ\-สมฺปยุตฺตํ ธมฺมํ ปจฺจยา นีวรณ\-สมฺปยุตฺโต ธมฺโม อุปฺปชฺชติ เหตุปจฺจยา. ตีณิ.}

\prose{นีวรณ\-วิปฺปยุตฺตํ ธมฺมํ ปจฺจยา นีวรณ\-วิปฺปยุตฺโต ธมฺโม อุปฺปชฺชติ เหตุปจฺจยา. นีวรณ\-วิปฺปยุตฺตํ เอกํ ขนฺธํ ปจฺจยา ตโย ขนฺธา จิตฺต\-สมุฏฺฐานญฺจ \csromanfolio{317}รูปํ ฯเปฯ เทฺว ขนฺเธ ฯเปฯ. ปฏิสนฺธิ\-กฺขเณ ฯเปฯ. เอกํ มหาภูตํ ฯเปฯ. วตฺถุํ ปจฺจยา นีวรณ\-วิปฺปยุตฺตา ขนฺธา. (1)}

\prose{นีวรณ\-วิปฺปยุตฺตํ ธมฺมํ ปจฺจยา นีวรณ\-สมฺปยุตฺโต ธมฺโม อุปฺปชฺชติ เหตุปจฺจยา. วตฺถุํ ปจฺจยา นีวรณ\-สมฺปยุตฺตา ขนฺธา. (2)}

\prose{นีวรณ\-วิปฺปยุตฺตํ ธมฺมํ ปจฺจยา นีวรณ\-สมฺปยุตฺโต จ นีวรณ\-วิปฺปยุตฺโต จ ธมฺมา อุปฺปชฺชนฺติ เหตุปจฺจยา. วตฺถุํ ปจฺจยา นีวรณ\-สมฺปยุตฺตา ขนฺธา. มหาภูเต ปจฺจยา จิตฺต\-สมุฏฺฐานํ รูปํ. (3)}

\prose{นีวรณ\-สมฺปยุตฺตญฺจ นีวรณ\-วิปฺปยุตฺตญฺจ ธมฺมํ ปจฺจยา นีวรณ\-สมฺปยุตฺโต ธมฺโม อุปฺปชฺชติ เหตุปจฺจยา. นีวรณ\-สมฺปยุตฺตํ เอกํ ขนฺธญฺจ วตฺถุญฺจ ปจฺจยา ตโย ขนฺธา ฯเปฯ เทฺว ขนฺเธ ฯเปฯ. (1)}

\prose{นีวรณ\-สมฺปยุตฺตญฺจ นีวรณ\-วิปฺปยุตฺตญฺจ ธมฺมํ ปจฺจยา นีวรณ\-วิปฺปยุตฺโต ธมฺโม อุปฺปชฺชติ เหตุปจฺจยา. นีวรณ\-สมฺปยุตฺเต ขนฺเธ จ มหาภูเต จ ปจฺจยา จิตฺต\-สมุฏฺฐานํ รูปํ. (2)}

\prose{นีวรณ\-สมฺปยุตฺตญฺจ นีวรณ\-วิปฺปยุตฺตญฺจ ธมฺมํ ปจฺจยา นีวรณ\-สมฺปยุตฺโต จ นีวรณ\-วิปฺปยุตฺโต จ ธมฺมา อุปฺปชฺชนฺติ เหตุปจฺจยา. นีวรณ\-สมฺปยุตฺตํ เอกํ ขนฺธญฺจ วตฺถุญฺจ ปจฺจยา ตโย ขนฺธา ฯเปฯ เทฺว ขนฺเธ ฯเปฯ. นีวรณ\-สมฺปยุตฺเต ขนฺเธ จ มหาภูเต จ ปจฺจยา จิตฺต\-สมุฏฺฐานํ รูปํ. (สํขิตฺตํ.) (3)}

\csromanmarkh{1. ปจฺจยานุโลม}\csromanmarkhii{2. สงฺขฺยาวาร}\csromanheadernomark{2}{1. \textbf{ปจฺจยานุโลม 2. สงฺขฺยาวาร}}
\proseitem{59}{เหตุยา นว, อารมฺมเณ จตฺตาริ, อธิปติยา นว, อนนฺตเร จตฺตาริ ฯเปฯ วิปาเก เอกํ ฯเปฯ อวิคเต นว.}

\vspace{\tipitakaparskip}
\csromancenter{อนุโลมํ.}
\csromanmarkh{2. ปจฺจย\-ปจฺจนีย}\csromanmarkhii{1. วิภงฺควาร}\csromanheadernomark{2}{2. ปจฺจย\-ปจฺจนีย 1. วิภงฺควาร}
\proseitem{62}{\csromanfolio{318}นีวรณ\-สมฺปยุตฺตํ ธมฺมํ ปจฺจยา นีวรณ\-สมฺปยุตฺโต ธมฺโม อุปฺปชฺชติ นเหตุปจฺจยา. วิจิกิจฺฉา\-สหคเต อุทฺธจฺจ\-สหคเต ขนฺเธ ปจฺจยา อวิชฺชา\-นีวรณํ. (1)}

\prose{นีวรณ\-วิปฺปยุตฺตํ ธมฺมํ ปจฺจยา นีวรณ\-วิปฺปยุตฺโต ธมฺโม อุปฺปชฺชติ นเหตุปจฺจยา. อเหตุกํ นีวรณ\-วิปฺปยุตฺตํ เอกํ ขนฺธํ ปจฺจยา ตโย ขนฺธา จิตฺต\-สมุฏฺฐานญฺจ รูปํ ฯเปฯ. (ยาว อสญฺญสตฺตา.) จกฺขายตนํ ปจฺจยา จกฺขุวิญฺญาณํ ฯเปฯ. กายายตนํ ปจฺจยา กายวิญฺญาณํ. วตฺถุํ ปจฺจยา อเหตุกา นีวรณ\-วิปฺปยุตฺตา ขนฺธา. (1)}

\prose{นีวรณ\-วิปฺปยุตฺตํ ธมฺมํ ปจฺจยา นีวรณ\-สมฺปยุตฺโต ธมฺโม อุปฺปชฺชติ นเหตุปจฺจยา. วตฺถุํ ปจฺจยา วิจิกิจฺฉา\-สหคโต อุทฺธจฺจ\-สหคโต โมโห. (2)}

\prose{นีวรณ\-สมฺปยุตฺตญฺจ นีวรณ\-วิปฺปยุตฺตญฺจ ธมฺมํ ปจฺจยา นีวรณ\-สมฺปยุตฺโต ธมฺโม อุปฺปชฺชติ นเหตุปจฺจยา. วิจิกิจฺฉา\-สหคเต อุทฺธจฺจ\-สหคเต ขนฺเธ จ วตฺถุญฺจ ปจฺจยา วิจิกิจฺฉา\-สหคโต อุทฺธจฺจ\-สหคโต โมโห. (สํขิตฺตํ.) (1)}

\csromanmarkh{2. ปจฺจย\-ปจฺจนีย}\csromanmarkhii{2. สงฺขฺยาวาร}\csromanheadernomark{2}{2. ปจฺจย\-ปจฺจนีย 2. สงฺขฺยาวาร}
\vspace{\tipitakaparskip}
\csromancenter{\textbf{นเหตุ}ยา จตฺตาริ, นอารมฺมเณ ตีณิ ฯเปฯ นปุเรชาเต จตฺตาริ, นปจฺฉาชาเต นว, นอาเสวเน นว, นกมฺเม จตฺตาริ, นวิปาเก นว ฯเปฯ นสมฺปยุตฺเต ตีณิ, นวิปฺปยุตฺเต เทฺว, โนนตฺถิยา ตีณิ, โนวิคเต ตีณิ.}
\csromancenter{(เอวํ อิตเร เทฺว คณนาปิ นิสฺสยวาโรปิ กาตพฺโพ.)}
\csromancenter{46. นีวรณ\-สมฺปยุตฺตทุก 5. สํสฏฺฐวาร}
\csromanheader{2}{1-4. ปจฺจยจตุกฺก}
\proseitem{64}{\csromanfolio{319}นีวรณ\-สมฺปยุตฺตํ ธมฺมํ สํสฏฺโฐ นีวรณ\-สมฺปยุตฺโต ธมฺโม อุปฺปชฺชติ เหตุปจฺจยา ฯเปฯ.}

\prose{เหตุยา เทฺว, อารมฺมเณ เทฺว, (สพฺพตฺถ เทฺว.) วิปาเก เอกํ ฯเปฯ อวิคเต เทฺว.}

\vspace{\tipitakaparskip}
\csromancenter{\textbf{อนุโลมํ.}}
\proseitem{65}{นีวรณ\-สมฺปยุตฺตํ ธมฺมํ สํสฏฺโฐ นีวรณ\-สมฺปยุตฺโต ธมฺโม อุปฺปชฺชติ นเหตุปจฺจยา. วิจิกิจฺฉา\-สหคเต อุทฺธจฺจ\-สหคเต ขนฺเธ สํสฏฺโฐ วิจิกิจฺฉา\-สหคโต อุทฺธจฺจ\-สหคโต โมโห.}

\prose{นีวรณ\-วิปฺปยุตฺตํ ธมฺมํ สํสฏฺโฐ. (สํขิตฺตํ.)}

\prose{นเหตุยา เทฺว, นอธิปติยา เทฺว, นปุเรชาเต เทฺว, นกมฺเม เทฺว, นวิปาเก เทฺว, นฌาเน เอกํ, นมคฺเค เอกํ, นวิปฺปยุตฺเต เทฺว.}

\vspace{\tipitakaparskip}
\csromancenter{ปจฺจนียํ.}
\csromancenter{(เอวํ อิตเร เทฺว คณนาปิ สมฺปยุตฺตวาโรปิ กาตพฺโพ.)}
\csromanmarkh{46. นีวรณ\-สมฺปยุตฺตทุก}\csromanmarkhii{7. ปญฺหาวาร}\csromanheadernomark{2}{46. นีวรณ\-สมฺปยุตฺตทุก 7. ปญฺหาวาร}
\csromanmarkh{1. ปจฺจยานุโลม}\csromanmarkhii{1. วิภงฺควาร}\csromanheadernomark{2}{1. \textbf{ปจฺจยานุโลม 1. วิภงฺควาร}}
\vspace{\tipitakaparskip}
\csromancenter{\textbf{เหตุ}}
\proseitem{66}{นีวรณ\-สมฺปยุตฺโต ธมฺโม นีวรณ\-สมฺปยุตฺตสฺส ธมฺมสฺส เหตุปจฺจเยน ปจฺจโย. นีวรณ\-สมฺปยุตฺตา เหตู สมฺปยุตฺตกานํ ขนฺธานํ เหตุปจฺจเยน ปจฺจโย. (มูลํ ปุจฺฉิตพฺพํ.) นีวรณ\-สมฺปยุตฺตา เหตู จิตฺต\-สมุฏฺฐานานํ รูปานํ เหตุปจฺจเยน ปจฺจโย. (มูลํ ปุจฺฉิตพฺพํ.) \csromanfolio{320}นีวรณ\-สมฺปยุตฺตา เหตู สมฺปยุตฺตกานํ ขนฺธานํ จิตฺตสมุฏฺฐานานญฺจ รูปานํ เหตุปจฺจเยน ปจฺจโย. (3)}

\prose{นีวรณ\-วิปฺปยุตฺโต ธมฺโม นีวรณ\-วิปฺปยุตฺตสฺส ธมฺมสฺส เหตุปจฺจเยน ปจฺจโย. นีวรณ\-วิปฺปยุตฺตา เหตู สมฺปยุตฺตกานํ ขนฺธานํ จิตฺตสมุฏฺฐานานญฺจ รูปานํ เหตุปจฺจเยน ปจฺจโย. ปฏิสนฺธิ\-กฺขเณ ฯเปฯ. (1)}

\proseitem{67}{นีวรณ\-สมฺปยุตฺโต ธมฺโม นีวรณ\-สมฺปยุตฺตสฺส ธมฺมสฺส อารมฺมณ\-ปจฺจเยน ปจฺจโย. ราคํ อสฺสาเทติ อภินนฺทติ, ตํ อารพฺภ ราโค อุปฺปชฺชติ, ทิฏฺฐิ ฯเปฯ วิจิกิจฺฉา ฯเปฯ อุทฺธจฺจํ ฯเปฯ โทมนสฺสํ อุปฺปชฺชติ. ทิฏฺฐิํ อสฺสาเทติ อภินนฺทติ, ตํ อารพฺภ ราโค อุปฺปชฺชติ. ทิฏฺฐิ ฯเปฯ วิจิกิจฺฉา ฯเปฯ อุทฺธจฺจํ ฯเปฯ โทมนสฺสํ อุปฺปชฺชติ. วิจิกิจฺฉํ อารพฺภ วิจิกิจฺฉา ฯเปฯ ทิฏฺฐิ ฯเปฯ อุทฺธจฺจํ ฯเปฯ โทมนสฺสํ อุปฺปชฺชติ. อุทฺธจฺจํ อารพฺภ อุทฺธจฺจํ อุปฺปชฺชติ. ทิฏฺฐิ ฯเปฯ วิจิกิจฺฉา. โทมนสฺสํ อุปฺปชฺชติ. โทมนสฺสํ อารพฺภ โทมนสฺสํ อุปฺปชฺชติ. ทิฏฺฐิ ฯเปฯ วิจิกิจฺฉา ฯเปฯ อุทฺธจฺจํ อุปฺปชฺชติ. (1)}

\prose{นีวรณ\-สมฺปยุตฺโต ธมฺโม นีวรณ\-วิปฺปยุตฺตสฺส ธมฺมสฺส อารมฺมณ\-ปจฺจเยน ปจฺจโย. อริยา นีวรณ\-สมฺปยุตฺเต ปหีเน กิเลเส ปจฺจเวกฺขนฺติ, วิกฺขมฺภิเต กิเลเส ปจฺจเวกฺขนฺติ. ปุพฺเพ สมุทาจิณฺเณ กิเลเส ชานนฺติ. นีวรณ\-สมฺปยุตฺเต ขนฺเธ อนิจฺจโต ทุกฺขโต อนตฺตโต วิปสฺสติ. เจโตปริย\-ญาเณน นีวรณ\-สมฺปยุตฺตจิตฺต\-สมงฺคิสฺส จิตฺตํ ชานาติ. นีวรณ\-สมฺปยุตฺตา ขนฺธา เจโตปริย\-ญาณสฺส, ปุพฺเพ\-นิวาสา\-นุสฺสติ\-ญาณสฺส, ยถากมฺมูปคญาณสฺส, อนาคตํสญาณสฺส, อาวชฺชนาย อารมฺมณ\-ปจฺจเยน ปจฺจโย. (2)}

\proseitem{68}{นีวรณ\-วิปฺปยุตฺโต ธมฺโม นีวรณ\-วิปฺปยุตฺตสฺส ธมฺมสฺส อารมฺมณ\-ปจฺจเยน ปจฺจโย. ทานํ ฯเปฯ สีลํ ฯเปฯ อุโปสถกมฺมํ กตฺวา ตํ ปจฺจเวกฺขติ. ปุพฺเพ สุจิณฺณานิ ฯเปฯ. ฌานา วุฏฺฐหิตฺวา ฌานํ ฯเปฯ. อริยา มคฺคา วุฏฺฐหิตฺวา มคฺคํ ฯเปฯ. ผลํ ฯเปฯ. นิพฺพานํ โคตฺรภุสฺส, โวทานสฺส, มคฺคสฺส, ผลสฺส, อาวชฺชนาย อารมฺมณ\-ปจฺจเยน ปจฺจโย. จกฺขุํ ฯเปฯ วตฺถุํ นีวรณ\-วิปฺปยุตฺเต ขนฺเธ อนิจฺจโต ทุกฺขโต อนตฺตโต วิปสฺสติ. ทิพฺเพน จกฺขุนา รูปํ ปสฺสติ. (ยาว อาวชฺชนาย.) (1)}

\prose{\csromanfolio{321}นีวรณาวิปฺปยุตฺโต ธมฺโม นีวรณ\-สมฺปยุตฺตสฺส ธมฺมสฺส อารมฺมณ\-ปจฺจเยน ปจฺจโย. ทานํ ฯเปฯ สีลํ ฯเปฯ อุโปสถกมฺมํ ฯเปฯ. ปุพฺเพ ฯเปฯ. ฌานา ฯเปฯ. จกฺขุํ ฯเปฯ วตฺถุํ นีวรณ\-วิปฺปยุตฺเต ขนฺเธ อสฺสาเทติ อภินนฺทติ, ตํ อารพฺภ ราโค อุปฺปชฺชติ, ทิฏฺฐิ ฯเปฯ วิจิกิจฺฉา ฯเปฯ อุทฺธจฺจํ ฯเปฯ โทมนสฺสํ อุปฺปชฺชติ. (2)}

\proseitem{69}{นีวรณ\-สมฺปยุตฺโต ธมฺโม นีวรณ\-สมฺปยุตฺตสฺส ธมฺมสฺส อธิปติ\-ปจฺจเยน ปจฺจโย. อารมฺมณา\-ธิปติ, สหชาตา\-ธิปติ.  อารมฺมณา\-ธิปติ\csromandash{}ราคํ ครุํ กตฺวา อสฺสาเทติ อภินนฺทติ, ตํ ครุํ กตฺวา ราโค อุปฺปชฺชติ, ทิฏฺฐิ อุปฺปชฺชติ. ทิฏฺฐิํ ครุํ กตฺวา อสฺสาเทติ ฯเปฯ. สหชาตา\-ธิปติ\csromandash{}นีวรณ\-สมฺปยุตฺตา\-ธิปติ นีวรณ\-สมฺปยุตฺตกานํ ขนฺธานํ อธิปติ\-ปจฺจเยน ปจฺจโย. (1)}

\prose{นีวรณ\-สมฺปยุตฺโต ธมฺโม นีวรณ\-วิปฺปยุตฺตสฺส ธมฺมสฺส อธิปติ\-ปจฺจเยน ปจฺจโย. สหชาตา\-ธิปติ\csromandash{}นีวรณ\-สมฺปยุตฺตา\-ธิปติ จิตฺต\-สมุฏฺฐานานํ รูปานํ อธิปติ\-ปจฺจเยน ปจฺจโย. (มูลํ ปุจฺฉิตพฺพํ.) นีวรณ\-สมฺปยุตฺตา\-ธิปติ สมฺปยุตฺตกานํ ขนฺธานํ จิตฺตสมุฏฺฐานานญฺจ รูปานํ อธิปติ\-ปจฺจเยน ปจฺจโย. (3)}

\proseitem{70}{นีวรณ\-วิปฺปยุตฺโต ธมฺโม นีวรณ\-วิปฺปยุตฺตสฺส ธมฺมสฺส อธิปติ\-ปจฺจเยน ปจฺจโย. อารมฺมณา\-ธิปติ, สหชาตา\-ธิปติ.  อารมฺมณา\-ธิปติ\csromandash{}ทานํ ฯเปฯ สีลํ ฯเปฯ อุโปสถกมฺมํ กตฺวา ตํ ครุํ กตฺวา ปจฺจเวกฺขติ. ปุพฺเพ สุจิณฺณานิ ฯเปฯ. ฌานา วุฏฺฐหิตฺวา ฌานํ ฯเปฯ. อริยา มคฺคา วุฏฺฐหิตฺวา มคฺคํ ครุํ กตฺวา ฯเปฯ. ผลํ ฯเปฯ. นิพฺพานํ ฯเปฯ. นิพฺพานํ โคตฺรภุสฺส, โวทานสฺส, มคฺคสฺส, ผลสฺส อธิปติ\-ปจฺจเยน ปจฺจโย. สหชาตา\-ธิปติ\csromandash{}นีวรณ\-วิปฺปยุตฺตา\-ธิปติ สมฺปยุตฺตกานํ ขนฺธานํ จิตฺตสมุฏฺฐานานญฺจ รูปานํ อธิปติ\-ปจฺจเยน ปจฺจโย. (1)}

\prose{นีวรณ\-วิปฺปยุตฺโต ธมฺโม นีวรณ\-สมฺปยุตฺตสฺส ธมฺมสฺส อธิปติ\-ปจฺจเยน ปจฺจโย. อารมฺมณา\-ธิปติ\csromandash{}ทานํ ฯเปฯ สีลํ ฯเปฯ. ปุพฺเพ ฯเปฯ. ฌานํ ฯเปฯ. จกฺขุํ ฯเปฯ วตฺถุํ นีวรณ\-วิปฺปยุตฺเต ขนฺเธ ครุํ กตฺวา อสฺสาเทติ อภินนฺทติ, ตํ ครุํ กตฺวา ราโค อุปฺปชฺชติ, ทิฏฺฐิ อุปฺปชฺชติ. (2)}

\proseitem{71}{\csromanfolio{322}นีวรณ\-สมฺปยุตฺโต ธมฺโม นีวรณ\-สมฺปยุตฺตสฺส ธมฺมสฺส อนนฺตร\-ปจฺจเยน ปจฺจโย. ปุริมา ปุริมา นีวรณ\-สมฺปยุตฺตา ขนฺธา ปจฺฉิมานํ ปจฺฉิมานํ นีวรณ\-สมฺปยุตฺตกานํ ขนฺธานํ อนนฺตร\-ปจฺจเยน ปจฺจโย. (มูลํ ปุจฺฉิตพฺพํ.) นีวรณ\-สมฺปยุตฺตา ขนฺธา วุฏฺฐานสฺส อนนฺตร\-ปจฺจเยน ปจฺจโย. (อิธ “ปุริมา ปุริมา”ติ นตฺถิ.) (2)}

\prose{(มูลํ ปุจฺฉิตพฺพํ.) ปุริมา ปุริมา นีวรณ\-วิปฺปยุตฺตา ขนฺธา ปจฺฉิมานํ ปจฺฉิมานํ นีวรณ\-วิปฺปยุตฺตานํ ขนฺธานํ อนนฺตร\-ปจฺจเยน ปจฺจโย. อนุโลมํ โคตฺรภุสฺส ฯเปฯ ผลสมาปตฺติยา อนนฺตร\-ปจฺจเยน ปจฺจโย. (1)}

\prose{นีวรณ\-วิปฺปยุตฺโต ธมฺโม นีวรณ\-สมฺปยุตฺตสฺส ธมฺมสฺส อนนฺตร\-ปจฺจเยน ปจฺจโย. อาวชฺชนา นีวรณ\-สมฺปยุตฺตกานํ ขนฺธานํ อนนฺตร\-ปจฺจเยน ปจฺจโย. (2)}

\csromanheader{2}{\textbf{สมนนฺตราทิ}}
\proseitem{72}{นีวรณ\-สมฺปยุตฺโต ธมฺโม นีวรณ\-สมฺปยุตฺตสฺส ธมฺมสฺส สมนนฺตร\-ปจฺจเยน ปจฺจโย. สหชาต\-ปจฺจเยน ปจฺจโย. อญฺญมญฺญ\-ปจฺจเยน ปจฺจโย. นิสฺสย\-ปจฺจเยน ปจฺจโย.}

\csromanheader{2}{\textbf{อุปนิสฺสย}}
\proseitem{73}{นีวรณ\-สมฺปยุตฺโต ธมฺโม นีวรณ\-สมฺปยุตฺตสฺส ธมฺมสฺส อุปนิสฺสย\-ปจฺจเยน ปจฺจโย. อารมฺมณู\-ปนิสฺสโย, อนนฺตรู\-ปนิสฺสโย, ปกตู\-ปนิสฺสโย ฯเปฯ. ปกตู\-ปนิสฺสโย\csromandash{}ราคํ อุปนิสฺสาย ปาณํ หนติ ฯเปฯ สํฆํ ภินฺทติ. โทสํ. โมหํ. มานํ. ทิฏฺฐิํ. ปตฺถนํ อุปนิสฺสาย ปาณํ หนติ ฯเปฯ สํฆํ ภินฺทติ. ราโค ฯเปฯ. ปตฺถนา ราคสฺส ฯเปฯ ปตฺถนาย อุปนิสฺสย\-ปจฺจเยน ปจฺจโย. (1)}

\prose{นีวรณ\-สมฺปยุตฺโต ธมฺโม นีวรณ\-วิปฺปยุตฺตสฺส ธมฺมสฺส อุปนิสฺสย\-ปจฺจเยน ปจฺจโย. อนนฺตรู\-ปนิสฺสโย, ปกตู\-ปนิสฺสโย ฯเปฯ. ปกตู\-ปนิสฺสโย\csromandash{}ราคํ อุปนิสฺสาย ทานํ เทติ ฯเปฯ สีลํ ฯเปฯ อุโปสถกมฺมํ ฯเปฯ. ฌานํ ฯเปฯ. วิปสฺสนํ. มคฺคํ. อภิญฺญํ. สมาปตฺติํ อุปฺปาเทติ. โทสํ ฯเปฯ. ปตฺถนํ อุปนิสฺสาย ทานํ เทติ ฯเปฯ สมาปตฺติํ อุปฺปาเทติ. ราโค ฯเปฯ. ปตฺถนา \csromanfolio{323}สทฺธาย ฯเปฯ ปญฺญาย. กายิกสฺส สุขสฺส, กายิกสฺส ทุกฺขสฺส, มคฺคสฺส, ผลสมาปตฺติยา อุปนิสฺสย\-ปจฺจเยน ปจฺจโย. (2)}

\proseitem{74}{นีวรณ\-วิปฺปยุตฺโต ธมฺโม นีวรณ\-วิปฺปยุตฺตสฺส ธมฺมสฺส อุปนิสฺสย\-ปจฺจเยน ปจฺจโย. อารมฺมณู\-ปนิสฺสโย, อนนฺตรู\-ปนิสฺสโย, ปกตู\-ปนิสฺสโย ฯเปฯ. ปกตู\-ปนิสฺสโย\csromandash{}สทฺธํ อุปนิสฺสาย ทานํ เทติ. สีลํ ฯเปฯ มคฺคํ ฯเปฯ อภิญฺญํ ฯเปฯ สมาปตฺติํ อุปฺปาเทติ. สีลํ ฯเปฯ. ปญฺญํ ฯเปฯ. เสนาสนํ อุปนิสฺสาย ทานํ เทติ ฯเปฯ สมาปตฺติํ อุปฺปาเทติ. สทฺธา ฯเปฯ. เสนาสนํ สทฺธาย ฯเปฯ ปญฺญาย ฯเปฯ มคฺคสฺส ผลสมาปตฺติยา อุปนิสฺสย\-ปจฺจเยน ปจฺจโย. (1)}

\prose{นีวรณ\-วิปฺปยุตฺโต ธมฺโม นีวรณ\-สมฺปยุตฺตสฺส ธมฺมสฺส อุปนิสฺสย\-ปจฺจเยน ปจฺจโย. อารมฺมณู\-ปนิสฺสโย, อนนฺตรู\-ปนิสฺสโย, ปกตู\-ปนิสฺสโย ฯเปฯ. ปกตู\-ปนิสฺสโย\csromandash{}สทฺธํ อุปนิสฺสาย มานํ ชปฺเปติ. ทิฏฺฐิํ คณฺหาติ. สีลํ ฯเปฯ. ปญฺญํ. กายิกํ สุขํ ฯเปฯ. เสนาสนํ อุปนิสฺสาย ปาณํ หนติ ฯเปฯ. สํฆํ ภินฺทติ. สทฺธา ฯเปฯ. เสนาสนํ ราคสฺส. โทสสฺส. โมหสฺส. มานสฺส. ทิฏฺฐิยา. ปตฺถนาย อุปนิสฺสย\-ปจฺจเยน ปจฺจโย. (2)}

\proseitem{75}{นีวรณ\-วิปฺปยุตฺโต ธมฺโม นีวรณ\-วิปฺปยุตฺตสฺส ธมฺมสฺส ปุเรชาต\-ปจฺจเยน ปจฺจโย. อารมฺมณ\-ปุเรชาตํ, วตฺถุ\-ปุเรชาตํ.  อารมฺมณ\-ปุเรชาตํ\csromandash{}จกฺขุํ ฯเปฯ วตฺถุํ อนิจฺจโต ทุกฺขโต อนตฺตโต วิปฺปสฺสติ. ทิพฺเพน จกฺขุนา รูปํ ปสฺสติ. ทิพฺพาย โสตธาตุยา สทฺทํ สุณาติ. รูปายตนํ จกฺขุ\-วิญฺญาณสฺส ฯเปฯ. โผฏฺฐพฺพา\-ยตนํ กายวิญฺญาณสฺส ฯเปฯ. วตฺถุ\-ปุเรชาตํ\csromandash{}จกฺขายตนํ จกฺขุ\-วิญฺญาณสฺส ฯเปฯ. กายายตนํ กายวิญฺญาณสฺส ฯเปฯ. วตฺถุ นีวรณ\-วิปฺปยุตฺตานํ ขนฺธานํ ปุเรชาต\-ปจฺจเยน ปจฺจโย. (1)}

\prose{นีวรณ\-วิปฺปยุตฺโต ธมฺโม นีวรณ\-สมฺปยุตฺตสฺส ธมฺมสฺส ปุเรชาต\-ปจฺจเยน ปจฺจโย. อารมฺมณ\-ปุเรชาตํ, วตฺถุ\-ปุเรชาตํ.  อารมฺมณ\-ปุเรชาตํ\csromandash{}จกฺขุํ ฯเปฯ วตฺถุํ อสฺสาเทติ อภินนฺทติ, ตํ อารพฺภ ราโค. โทโส. โมโห. วตฺถุ\-ปุเรชาตํ\csromandash{}วตฺถุ นีวรณ\-สมฺปยุตฺตานํ ขนฺธานํ ปุเรชาต\-ปจฺจเยน ปจฺจโย. (2)}

\proseitem{76}{\csromanfolio{324}นีวรณ\-สมฺปยุตฺโต ธมฺโม นีวรณ\-วิปฺปยุตฺตสฺส ธมฺมสฺส ปจฺฉาชาต\-ปจฺจเยน ปจฺจโย. เทฺว. อาเสวน\-ปจฺจเยน ปจฺจโย. เทฺว.}

\csromanheader{2}{\textbf{กมฺมาทิ}}
\proseitem{77}{นีวรณ\-สมฺปยุตฺโต ธมฺโม นีวรณ\-สมฺปยุตฺตสฺส ธมฺมสฺส กมฺมปจฺจเยน ปจฺจโย. นีวรณ\-สมฺปยุตฺตา เจตนา สมฺปยุตฺตกานํ ขนฺธานํ กมฺมปจฺจเยน ปจฺจโย. (มูลํ กาตพฺพํ.) สหชาตา, นานากฺขณิกา.  สหชาตา นีวรณ\-สมฺปยุตฺตา เจตนา จิตฺต\-สมุฏฺฐานานํ รูปานํ กมฺมปจฺจเยน ปจฺจโย. นานากฺขณิกา นีวรณ\-สมฺปยุตฺตา เจตนา วิปากานํ ขนฺธานํ กฏตฺตา จ รูปานํ กมฺมปจฺจเยน ปจฺจโย. (มูลํ กาตพฺพํ.) นีวรณ\-สมฺปยุตฺตา เจตนา สมฺปยุตฺตกานํ ขนฺธานํ จิตฺตสมุฏฺฐานานญฺจ รูปานํ กมฺมปจฺจเยน ปจฺจโย. (3)}

\prose{นีวรณ\-วิปฺปยุตฺโต ธมฺโม นีวรณ\-วิปฺปยุตฺตสฺส ธมฺมสฺส กมฺมปจฺจเยน ปจฺจโย. สหชาตา, นานากฺขณิกา. (สํขิตฺตํ.) วิปากปจฺจยา เอกํ.}

\csromanheader{2}{\textbf{อาหาราทิ}}
\proseitem{78}{นีวรณ\-สมฺปยุตฺโต ธมฺโม นีวรณ\-สมฺปยุตฺตสฺส ธมฺมสฺส อาหารปจฺจเยน ปจฺจโย. อินฺทฺริย\-ปจฺจเยน ปจฺจโย. ฌานปจฺจเยน ปจฺจโย. มคฺคปจฺจเยน ปจฺจโย. สมฺปยุตฺต\-ปจฺจเยน ปจฺจโย. เทฺว.}

\proseitem{79}{นีวรณ\-สมฺปยุตฺโต ธมฺโม นีวรณ\-วิปฺปยุตฺตสฺส ธมฺมสฺส วิปฺปยุตฺต\-ปจฺจเยน ปจฺจโย. สหชาตํ, ปจฺฉาชาตํ. (สํขิตฺตํ.) (1)}

\prose{นีวรณ\-วิปฺปยุตฺโต ธมฺโม นีวรณ\-วิปฺปยุตฺตสฺส ธมฺมสฺส วิปฺปยุตฺต\-ปจฺจเยน ปจฺจโย. สหชาตํ, ปุเรชาตํ, ปจฺฉาชาตํ. (สํขิตฺตํ.) (1)}

\prose{นีวรณ\-วิปฺปยุตฺโต ธมฺโม นีวรณ\-สมฺปยุตฺตสฺส ธมฺมสฺส วิปฺปยุตฺต\-ปจฺจเยน ปจฺจโย. ปุเรชาตํ วตฺถุ นีวรณ\-สมฺปยุตฺตกานํ ขนฺธานํ วิปฺปยุตฺต\-ปจฺจเยน ปจฺจโย. (2)}

\csromanheader{2}{46. นีวรณ\-สมฺปยุตฺตทุก อตฺถิ}
\proseitem{80}{\csromanfolio{325}นีวรณ\-สมฺปยุตฺโต ธมฺโม นีวรณ\-สมฺปยุตฺตสฺส ธมฺมสฺส อตฺถิปจฺจเยน ปจฺจโย. นีวรณ\-สมฺปยุตฺโต เอโก ขนฺโธ ติณฺณนฺนํ ขนฺธานํ อตฺถิปจฺจเยน ปจฺจโย ฯเปฯ. (1)}

\prose{นีวรณ\-สมฺปยุตฺโต ธมฺโม นีวรณ\-วิปฺปยุตฺตสฺส ธมฺมสฺส อตฺถิปจฺจเยน ปจฺจโย. สหชาตํ, ปจฺฉาชาตํ.  สหชาตา นีวรณ\-สมฺปยุตฺตา ขนฺธา จิตฺตสมุฏฺฐานานญฺจ รูปานํ อตฺถิปจฺจเยน ปจฺจโย. ปจฺฉาชาตา นีวรณ\-สมฺปยุตฺตา ขนฺธา ปุเรชาตสฺส อิมสฺส กายสฺส อตฺถิปจฺจเยน ปจฺจโย. (2)}

\prose{นีวรณ\-สมฺปยุตฺโต ธมฺโม นีวรณ\-สมฺปยุตฺตสฺส จ นีวรณ\-วิปฺปยุตฺตสฺส จ ธมฺมสฺส อตฺถิปจฺจเยน ปจฺจโย. นีวรณ\-สมฺปยุตฺโต เอโก ขนฺโธ ติณฺณนฺนํ ขนฺธานํ จิตฺตสมุฏฺฐานานญฺจ รูปานํ อตฺถิปจฺจเยน ปจฺจโย ฯเปฯ เทฺว ขนฺธา ทฺวินฺนํ ขนฺธานํ จิตฺตสมุฏฺฐานานญฺจ รูปานํ อตฺถิปจฺจเยน ปจฺจโย. (3)}

\proseitem{81}{นีวรณ\-วิปฺปยุตฺโต ธมฺโม นีวรณ\-วิปฺปยุตฺตสฺส ธมฺมสฺส อตฺถิปจฺจเยน ปจฺจโย. สหชาตํ, ปุเรชาตํ, ปจฺฉาชาตํ, อาหารํ, อินฺทฺริยํ. (สํขิตฺตํ.) (1)}

\prose{นีวรณ\-วิปฺปยุตฺโต ธมฺโม นีวรณ\-สมฺปยุตฺตสฺส ธมฺมสฺส อตฺถิปจฺจเยน ปจฺจโย. ปุเรชาตํ จกฺขุํ ฯเปฯ วตฺถุํ อสฺสาเทติ อภินนฺทติ, ตํ อารพฺภ ราโค ฯเปฯ ทิฏฺฐิ ฯเปฯ วิจิกิจฺฉา ฯเปฯ อุทฺธจฺจํ ฯเปฯ โทมนสฺสํ อุปฺปชฺชติ. วตฺถุ นีวรณ\-สมฺปยุตฺตกานํ ขนฺธานํ อตฺถิปจฺจเยน ปจฺจโย. (2)}

\proseitem{82}{นีวรณ\-สมฺปยุตฺโต จ นีวรณ\-วิปฺปยุตฺโต จ ธมฺมา นีวรณ\-สมฺปยุตฺตสฺส ธมฺมสฺส อตฺถิปจฺจเยน ปจฺจโย. สหชาตํ, ปุเรชาตํ.  สหชาโต นีวรณ\-สมฺปยุตฺโต เอโก ขนฺโธ จ วตฺถุ จ ติณฺณนฺนํ ขนฺธานํ อตฺถิปจฺจเยน ปจฺจโย ฯเปฯ เทฺว ขนฺธา จ วตฺถุ จ ทฺวินฺนํ ขนฺธานํ อตฺถิปจฺจเยน ปจฺจโย. (1)}

\prose{นีวรณ\-สมฺปยุตฺโต จ นีวรณ\-วิปฺปยุตฺโต จ ธมฺมา นีวรณ\-วิปฺปยุตฺตสฺส ธมฺมสฺส อตฺถิปจฺจเยน ปจฺจโย. สหชาตํ, ปจฺฉาชาตํ, อาหารํ, อินฺทฺริยํ.  สหชาตา นีวรณ\-สมฺปยุตฺตา ขนฺธา จ มหาภูตา จ \csromanfolio{326}จิตฺต\-สมุฏฺฐานานํ รูปานํ อตฺถิปจฺจเยน ปจฺจโย. ปจฺฉาชาตา นีวรณ\-สมฺปยุตฺตา ขนฺธา จ กพฬีกาโร อาหาโร จ อิมสฺส กายสฺส อตฺถิปจฺจเยน ปจฺจโย. ปจฺฉาชาตา นีวรณ\-สมฺปยุตฺตา ขนฺธา จ รูปาชีวิตินฺทฺริยญฺจ กฏตฺตารูปานํ อตฺถิปจฺจเยน ปจฺจโย. (2)}

\csromanmarkh{1. ปจฺจยานุโลม}\csromanmarkhii{2. สงฺขฺยาวาร}\csromanheadernomark{2}{1. \textbf{ปจฺจยานุโลม 2. สงฺขฺยาวาร}}
\proseitem{83}{เหตุยา จตฺตาริ, อารมฺมเณ จตฺตาริ, อธิปติยา ปญฺจ, อนนฺตเร จตฺตาริ, สมนนฺตเร จตฺตาริ, สหชาเต ปญฺจ, อญฺญมญฺเญ เทฺว, นิสฺสเย สตฺต, อุปนิสฺสเย จตฺตาริ, ปุเรชาเต เทฺว, ปจฺฉาชาเต เทฺว, อาเสวเน เทฺว, กมฺเม จตฺตาริ, วิปาเก เอกํ, อาหาเร จตฺตาริ, อินฺทฺริเย จตฺตาริ, ฌาเน จตฺตาริ, มคฺเค จตฺตาริ, สมฺปยุตฺเต เทฺว, วิปฺปยุตฺเต ตีณิ, อตฺถิยา สตฺต, นตฺถิยา จตฺตาริ, วิคเต จตฺตาริ, อวิคเต สตฺต.}

\vspace{\tipitakaparskip}
\csromancenter{อนุโลมํ.}
\proseitem{84}{นีวรณ\-สมฺปยุตฺโต ธมฺโม นีวรณ\-สมฺปยุตฺตสฺส ธมฺมสฺส อารมฺมณ\-ปจฺจเยน ปจฺจโย. สหชาต\-ปจฺจเยน ปจฺจโย. อุปนิสฺสย\-ปจฺจเยน ปจฺจโย. (1)}

\prose{นีวรณ\-สมฺปยุตฺโต ธมฺโม นีวรณ\-วิปฺปยุตฺตสฺส ธมฺมสฺส อารมฺมณ\-ปจฺจเยน ปจฺจโย. สหชาต\-ปจฺจเยน ปจฺจโย. อุปนิสฺสย\-ปจฺจเยน ปจฺจโย. ปจฺฉาชาต\-ปจฺจเยน ปจฺจโย. กมฺมปจฺจเยน ปจฺจโย. (2)}

\prose{นีวรณ\-สมฺปยุตฺโต ธมฺโม นีวรณ\-สมฺปยุตฺตสฺส จ นีวรณ\-วิปฺปยุตฺตสฺส จ ธมฺมสฺส สหชาต\-ปจฺจเยน ปจฺจโย. (3)}

\proseitem{85}{นีวรณ\-วิปฺปยุตฺโต ธมฺโม นีวรณ\-วิปฺปยุตฺตสฺส ธมฺมสฺส อารมฺมณ\-ปจฺจเยน ปจฺจโย. สหชาต\-ปจฺจเยน ปจฺจโย. อุปนิสฺสย\-ปจฺจเยน ปจฺจโย. ปุเรชาต\-ปจฺจเยน ปจฺจโย. ปจฺฉาชาต\-ปจฺจเยน ปจฺจโย. \csromanfolio{327}กมฺมปจฺจเยน ปจฺจโย. อาหารปจฺจเยน ปจฺจโย. อินฺทฺริย\-ปจฺจเยน ปจฺจโย. (1)}

\prose{นีวรณ\-วิปฺปยุตฺโต ธมฺโม นีวรณ\-สมฺปยุตฺตสฺส ธมฺมสฺส อารมฺมณ\-ปจฺจเยน ปจฺจโย. อุปนิสฺสย\-ปจฺจเยน ปจฺจโย. ปุเรชาต\-ปจฺจเยน ปจฺจโย. (2)}

\prose{นีวรณ\-สมฺปยุตฺโต จ นีวรณ\-วิปฺปยุตฺโต จ ธมฺมา นีวรณ\-สมฺปยุตฺตสฺส ธมฺมสฺส สหชาตํ, ปุเรชาตํ. (1)}

\prose{นีวรณ\-สมฺปยุตฺโต จ นีวรณ\-วิปฺปยุตฺโต จ ธมฺมา นีวรณ\-วิปฺปยุตฺตสฺส ธมฺมสฺส สหชาตํ, ปจฺฉาชาตํ, อาหารํ, อินฺทฺริยํ. (2)}

\csromanmarkh{2. ปจฺจย\-ปจฺจนีย}\csromanmarkhii{2. สงฺขฺยาวาร}\csromanheadernomark{2}{2. ปจฺจย\-ปจฺจนีย 2. สงฺขฺยาวาร}
\proseitem{86}{นเหตุยา สตฺต, นอารมฺมเณ สตฺต, นอธิปติยา สตฺต, นอนนฺตเร สตฺต, นสมนนฺตเร สตฺต, นสหชาเต ปญฺจ, นอญฺญมญฺเญ ปญฺจ, นนิสฺสเย ปญฺจ, นอุปนิสฺสเย สตฺต, นปุเรชาเต ฉ ฯเปฯ นมคฺเค สตฺต, นสมฺปยุตฺเต ปญฺจ, นวิปฺปยุตฺเต จตฺตาริ, โนอตฺถิยา จตฺตาริ, โนนตฺถิยา สตฺต, โนวิคเต สตฺต, โนอวิคเต จตฺตาริ.}

\vspace{\tipitakaparskip}
\csromancenter{ปจฺจนียํ.}
\proseitem{87}{เหตุปจฺจยา นอารมฺมเณ จตฺตาริ, นอธิปติยา จตฺตาริ, นอนนฺตเร จตฺตาริ, นสมนนฺตเร จตฺตาริ, นอญฺญมญฺเญ เทฺว, นอุปนิสฺสเย จตฺตาริ ฯเปฯ นมคฺเค จตฺตาริ, นสมฺปยุตฺเต เทฺว, นวิปฺปยุตฺเต เทฺว, โนนตฺถิยา จตฺตาริ, โนวิคเต จตฺตาริ.}

\proseitemwithcloser{88}{\csromanfolio{328}นเหตุปจฺจยา อารมฺมเณ จตฺตาริ, อธิปติยา ปญฺจ, (อนุโลมคณนา) ฯเปฯ อวิคเต สตฺต.}{นีวรณ\-สมฺปยุตฺตทุกํ นิฏฺฐิตํ.}
\csromanmatikaanchor{mtk.44}
\csromantocmarkref{subsection}{47. นีวรณนีวรณิยทุก}{mtk.44}
\bookmark[dest={mtk.44},level=2]{47. นีวรณนีวรณิยทุก}
\csromanmarkh{47. นีวรณ\-นีวรณิยทุก}\csromanmarkhii{1. ปฏิจฺจวาร}\csromanheadernomark{2}{47. นีวรณ\-นีวรณิยทุก 1. ปฏิจฺจวาร}
\proseitem{89}{นีวรณญฺเจว นีวรณิยญฺจ ธมฺมํ ปฏิจฺจ นีวรโณ เจว นีวรณิโย จ ธมฺโม อุปฺปชฺชติ เหตุปจฺจยา. กามจฺฉนฺท\-นีวรณํ ปฏิจฺจ ถินมิทฺธ\-นีวรณํ, อุทฺธจฺจ\-นีวรณํ, อวิชฺชา\-นีวรณํ. (เอวํ สพฺเพ คณนา วิภชิตพฺพา. นีวรณ\-ทุก\-สทิสํ. นินฺนานากรณํ.)}

\csromanmarkh{47. นีวรณ\-นีวรณิยทุก}\csromanmarkhii{7. ปญฺหาวาร}\csromanheadernomark{2}{47. นีวรณ\-นีวรณิยทุก 7. ปญฺหาวาร}
\csromanmarkh{1. ปจฺจยานุโลม}\csromanmarkhii{1. วิภงฺควาร}\csromanheadernomark{2}{1. \textbf{ปจฺจยานุโลม 1. วิภงฺควาร}}
\csromanheader{2}{\textbf{เหตุ}}
\proseitem{90}{นีวรโณ เจว นีวรณิโย จ ธมฺโม นีวรณสฺส เจว นีวรณิยสฺส จ ธมฺมสฺส เหตุปจฺจเยน ปจฺจโย. นีวรณา เจว นีวรณิยา จ เหตู สมฺปยุตฺตกานํ นีวรณานํ เหตุปจฺจเยน ปจฺจโย. (1)}

\prose{นีวรโณ เจว นีวรณิโย จ ธมฺโม นีวรณิยสฺส เจว โน จ นีวรณสฺส ธมฺมสฺส เหตุปจฺจเยน ปจฺจโย. นีวรณา เจว นีวรณิยา จ เหตู สมฺปยุตฺตกานํ ขนฺธานํ จิตฺตสมุฏฺฐานานญฺจ รูปานํ เหตุปจฺจเยน ปจฺจโย. (2)}

\prose{นีวรโณ เจว นีวรณิโย จ ธมฺโม นีวรณสฺส เจว นีวรณิยสฺส จ นีวรณิยสฺส เจว โน จ นีวรณสฺส จ ธมฺมสฺส เหตุปจฺจเยน ปจฺจโย.}

\vspace{\tipitakaparskip}
\csromancenter{นีวรณ\-นีวรณิยทุก นีวรณา เจว นีวรณิยา จ เหตู สมฺปยุตฺตกานํ ขนฺธานํ นีวรณานญฺจ จิตฺตสมุฏฺฐานานญฺจ รูปานํ เหตุปจฺจเยน ปจฺจโย. (3)}
\prose{\csromanfolio{329}นีวรณิโย เจว โน จ นีวรโณ ธมฺโม นีวรณิยสฺส เจว โน จ นีวรณสฺส ธมฺมสฺส เหตุปจฺจเยน ปจฺจโย. นีวรณิยา เจว โน จ นีวรณา เหตู สมฺปยุตฺตกานํ ขนฺธานํ จิตฺตสมุฏฺฐานานญฺจ รูปานํ เหตุปจฺจเยน ปจฺจโย. ปฏิสนฺธิ\-กฺขเณ ฯเปฯ. (1)}

\proseitem{91}{นีวรโณ เจว นีวรณิโย จ ธมฺโม นีวรณสฺส เจว นีวรณิยสฺส จ ธมฺมสฺส อารมฺมณ\-ปจฺจเยน ปจฺจโย. นีวรเณ อารพฺภ นีวรณา อุปฺปชฺชนฺติ. (มูลํ ปุจฺฉิตพฺพํ.) นีวรเณ อารพฺภ นีวรณิยา เจว โน จ นีวรณา ขนฺธา อุปฺปชฺชนฺติ. (มูลํ ปุจฺฉิตพฺพํ.) นีวรเณ อารพฺภ นีวรณา จ สมฺปยุตฺตกา จ ขนฺธา อุปฺปชฺชนฺติ. (3)}

\prose{นีวรณิโย เจว โน จ นีวรโณ ธมฺโม นีวรณิยสฺส เจว โน จ นีวรณสฺส ธมฺมสฺส อารมฺมณ\-ปจฺจเยน ปจฺจโย. ทานํ ฯเปฯ สีลํ ฯเปฯ อุโปสถกมฺมํ กตฺวา ตํ ปจฺจเวกฺขติ, อสฺสาเทติ อภินนฺทติ, ตํ อารพฺภ ราโค อุปฺปชฺชติ. ทิฏฺฐิ. วิจิกิจฺฉา. อุทฺธจฺจํ. โทมนสฺสํ อุปฺปชฺชติ. ปุพฺเพ สุจิณฺณานิ ฯเปฯ. ฌานา ฯเปฯ. อริยา โคตฺรภุํ ปจฺจเวกฺขนฺติ. โวทานํ ฯเปฯ. ปหีเน กิเลเส ฯเปฯ. วิกฺขมฺภิเต กิเลเส ฯเปฯ. ปุพฺเพ สมุทาจิณฺเณ ฯเปฯ. จกฺขุํ ฯเปฯ วตฺถุํ นีวรณิเย เจว โน จ นีวรเณ ขนฺเธ อนิจฺจโต ฯเปฯ วิปสฺสติ, อสฺสาเทติ อภินนฺทติ ฯเปฯ โทมนสฺสํ อุปฺปชฺชติ. ทิพฺเพน จกฺขุนา รูปํ ปสฺสติ ฯเปฯ. (ยาว อาวชฺชนา, ตาว กาตพฺพา.) (1)}

\prose{นีวรณิโย เจว โน จ นีวรโณ ธมฺโม นีวรณสฺส เจว นีวรณิยสฺส จ ธมฺมสฺส อารมฺมณ\-ปจฺจเยน ปจฺจโย. ทานํ ฯเปฯ สีลํ ฯเปฯ อุโปสถกมฺมํ ฯเปฯ. ปุพฺเพ สุจิณฺณานิ ฯเปฯ. ฌานา ฯเปฯ. จกฺขุํ ฯเปฯ วตฺถุํ นีวรณิเย เจว โน จ นีวรเณ ขนฺเธ อสฺสาเทติ อภินนฺทติ, ตํ อารพฺภ ราโค ฯเปฯ ทิฏฺฐิ. วิจิกิจฺฉา. อุทฺธจฺจํ. โทมนสฺสํ อุปฺปชฺชติ. (2)}

\prose{นีวรณิโย เจว โน จ นีวรโณ ธมฺโม นีวรณสฺส เจว นีวรณิยสฺส จ นีวรณิยสฺส เจว โน จ นีวรณสฺส จ ธมฺมสฺส อารมฺมณ\-ปจฺจเยน ปจฺจโย. ทานํ ฯเปฯ สีลํ ฯเปฯ อุโปสถกมฺมํ ฯเปฯ. ปุพฺเพ สุจิณฺณานิ ฯเปฯ. \csromanfolio{330}ฌานา ฯเปฯ. จกฺขุํ ฯเปฯ วตฺถุํ นีวรณิเย เจว โน จ นีวรเณ ขนฺเธ อสฺสาเทติ อภินนฺทติ, ตํ อารพฺภ นีวรณา จ สมฺปยุตฺตกา จ ขนฺธา อุปฺปชฺชนฺติ. (เอวํ อิตเรปิ ตีณิ กาตพฺพา.) (3)}

\prosewithcloser{(อารมฺมณสทิสา. อธิปติ\-ปจฺจยา. ปุเรชาตมฺปิ อารมฺมณ\-สทิสํ. อุปนิสฺสเยปิ โลกุตฺตรํ น กาตพฺพํ. สํขิตฺตํ. เอวํ วิตฺถาเรตพฺพํ, ยถา นีวรณทุกํ, เอวํ ปจฺจเวกฺขิตฺวา กาตพฺพํ.)}{นีวรณ\-นีวรณิยทุกํ นิฏฺฐิตํ.}
\csromanmatikaanchor{mtk.45}
\csromantocmarkref{subsection}{48. นีวรณนีวรณสมฺปยุตฺตทุก}{mtk.45}
\bookmark[dest={mtk.45},level=2]{48. นีวรณนีวรณสมฺปยุตฺตทุก}
\csromanmarkh{48. นีวรณ\-นีวรณ\-สมฺปยุตฺตทุก}\csromanmarkhii{1. ปฏิจฺจวาร}\csromanheadernomark{2}{48. นีวรณ\-นีวรณ\-สมฺปยุตฺตทุก 1. ปฏิจฺจวาร}
\csromanmarkh{1. ปจฺจยานุโลม}\csromanmarkhii{1. วิภงฺควาร}\csromanheadernomark{2}{1. \textbf{ปจฺจยานุโลม 1. วิภงฺควาร}}
\csromanheader{2}{\textbf{เหตุ}}
\proseitem{92}{นีวรณญฺเจว นีวรณ\-สมฺปยุตฺตญฺจ ธมฺมํ ปฏิจฺจ นีวรโณ เจว นีวรณ\-สมฺปยุตฺโต จ ธมฺโม อุปฺปชฺชติ เหตุปจฺจยา. กามจฺฉนฺท\-นีวรณํ ปฏิจฺจ ถินมิทฺธ\-นีวรณํ, อุทฺธจฺจ\-นีวรณํ, อวิชฺชา\-นีวรณํ. (จกฺกํ. สพฺเพปิ นีวรณา กาตพฺพา.) (1)}

\prose{นีวรณญฺเจว นีวรณ\-สมฺปยุตฺตญฺจ ธมฺมํ ปฏิจฺจ นีวรณ\-สมฺปยุตฺโต เจว โน จ นีวรโณ ธมฺโม อุปฺปชฺชติ เหตุปจฺจยา. นีวรเณ ปฏิจฺจ สมฺปยุตฺตกา ขนฺธา. (2)}

\prose{นีวรณญฺเจว นีวรณ\-สมฺปยุตฺตญฺจ ธมฺมํ ปฏิจฺจ นีวรโณ เจว นีวรณ\-สมฺปยุตฺโต จ นีวรณ\-สมฺปยุตฺโต เจว โน จ นีวรโณ จ ธมฺมา อุปฺปชฺชนฺติ เหตุปจฺจยา. กามจฺฉนฺท\-นีวรณํ ปฏิจฺจ ถินมิทฺธ\-นีวรณํ, อุทฺธจฺจ\-นีวรณํ, อวิชฺชา\-นีวรณํ สมฺปยุตฺตกา จ ขนฺธา. (จกฺกํ.) (3)}

\proseitem{93}{นีวรณ\-สมฺปยุตฺต\-ญฺเจว โน จ นีวรณํ ธมฺมํ ปฏิจฺจ นีวรณ\-สมฺปยุตฺโต เจว โน จ นีวรโณ ธมฺโม อุปฺปชฺชติ เหตุปจฺจยา. นีวรณ\-สมฺปยุตฺต\-ญฺเจว โน จ นีวรณํ เอกํ ขนฺธํ ปฏิจฺจ ตโย ขนฺธา ฯเปฯ เทฺว ขนฺเธ ฯเปฯ. (1)}

\prose{\csromanfolio{331}นีวรณ\-สมฺปยุตฺต\-ญฺเจว โน จ นีวรณํ ธมฺมํ ปฏิจฺจ นีวรโณ เจว นีวรณ\-สมฺปยุตฺโต จ ธมฺโม อุปฺปชฺชติ เหตุปจฺจยา. นีวรณ\-สมฺปยุตฺเต เจว โน จ นีวรเณ ขนฺเธ ปฏิจฺจ นีวรณา. (2)}

\prose{นีวรณ\-สมฺปยุตฺต\-ญฺเจว โน จ นีวรณํ ธมฺมํ ปฏิจฺจ นีวรโณ เจว นีวรณ\-สมฺปยุตฺโต จ นีวรณ\-สมฺปยุตฺโต เจว โน จ นีวรโณ จ ธมฺมา อุปฺปชฺชนฺติ เหตุปจฺจยา. นีวรณ\-สมฺปยุตฺต\-ญฺเจว โน จ นีวรณํ เอกํ ขนฺธํ ปฏิจฺจ ตโย ขนฺธา นีวรณา จ ฯเปฯ เทฺว ขนฺเธ ฯเปฯ. (3)}

\proseitem{94}{นีวรณญฺเจว นีวรณ\-สมฺปยุตฺตญฺจ นีวรณ\-สมฺปยุตฺต\-ญฺเจว โน จ นีวรณญฺจ ธมฺมํ ปฏิจฺจ นีวรโณ เจว นีวรณ\-สมฺปยุตฺโต จ ธมฺโม อุปฺปชฺชติ เหตุปจฺจยา. กามจฺฉนฺทนีวรณญฺจ สมฺปยุตฺตเก จ ขนฺเธ ปฏิจฺจ ถินมิทฺธ\-นีวรณํ, อุทฺธจฺจ\-นีวรณํ, อวิชฺชา\-นีวรณํ. (จกฺกํ.) (1)}

\prose{นีวรณญฺเจว นีวรณ\-สมฺปยุตฺตญฺจ นีวรณ\-สมฺปยุตฺต\-ญฺเจว โน จ นีวรณญฺจ ธมฺมํ ปฏิจฺจ นีวรณ\-สมฺปยุตฺโต เจว โน จ นีวรโณ ธมฺโม อุปฺปชฺชติ เหตุปจฺจยา. นีวรณ\-สมฺปยุตฺต\-ญฺเจว โน จ นีวรณํ เอกํ ขนฺธญฺจ นีวรณญฺจ ปฏิจฺจ ตโย ขนฺธา ฯเปฯ เทฺว ขนฺเธ จ ฯเปฯ. (2)}

\prose{นีวรณญฺเจว นีวรณ\-สมฺปยุตฺตญฺจ นีวรณ\-สมฺปยุตฺต\-ญฺเจว โน จ นีวรณญฺจ ธมฺมํ ปฏิจฺจ นีวรโณ เจว นีวรณ\-สมฺปยุตฺโต จ นีวรณ\-สมฺปยุตฺโต เจว โน จ นีวรโณ จ ธมฺมา อุปฺปชฺชนฺติ เหตุปจฺจยา. นีวรณ\-สมฺปยุตฺต\-ญฺเจว โน จ นีวรณํ เอกํ ขนฺธญฺจ กามจฺฉนฺทนีวรณญฺจ ปฏิจฺจ ตโย ขนฺธา ถินมิทฺธ\-นีวรณํ, อุทฺธจฺจ\-นีวรณํ, อวิชฺชา\-นีวรณํ ฯเปฯ เทฺว ขนฺเธ จ ฯเปฯ. (จกฺกํ. สํขิตฺตํ.) (3)}

\csromanmarkh{1. ปจฺจยานุโลม}\csromanmarkhii{2. สงฺขฺยาวาร}\csromanheadernomark{2}{1. \textbf{ปจฺจยานุโลม 2. สงฺขฺยาวาร}}
\proseitem{95}{เหตุยา นว, อารมฺมเณ นว, อธิปติยา นว, (สพฺพตฺถ นว.) กมฺเม นว, อาหาเร นว ฯเปฯ อวิคเต นว.}

\vspace{\tipitakaparskip}
\csromancenter{อนุโลมํ.}
\csromanmarkh{2. ปจฺจย\-ปจฺจนีย}\csromanmarkhii{1. วิภงฺควาร}\csromanheadernomark{2}{2. ปจฺจย\-ปจฺจนีย 1. วิภงฺควาร}
\proseitem{96}{\csromanfolio{332}นีวรณญฺเจว นีวรณ\-สมฺปยุตฺตญฺจ ธมฺมํ ปฏิจฺจ นีวรโณ เจว นีวรณ\-สมฺปยุตฺโต จ ธมฺโม อุปฺปชฺชติ นเหตุปจฺจยา. วิจิกิจฺฉา\-นีวรณํ ปฏิจฺจ อวิชฺชา\-นีวรณํ. อุทฺธจฺจ\-นีวรณํ ปฏิจฺจ อวิชฺชา\-นีวรณํ. (1)}

\prose{นีวรณ\-สมฺปยุตฺต\-ญฺเจว โน จ นีวรณํ ธมฺมํ ปฏิจฺจ นีวรโณ เจว นีวรณ\-สมฺปยุตฺโต จ ธมฺโม อุปฺปชฺชติ นเหตุปจฺจยา. วิจิกิจฺฉา\-สหคเต อุทฺธจฺจ\-สหคเต ขนฺเธ ปฏิจฺจ อวิชฺชา\-นีวรณํ. (1)}

\prose{นีวรณญฺเจว นีวรณ\-สมฺปยุตฺตญฺจ นีวรณ\-สมฺปยุตฺต\-ญฺเจว โน จ นีวรณญฺจ ธมฺมํ ปฏิจฺจ นีวรโณ เจว นีวรณ\-สมฺปยุตฺโต จ ธมฺโม อุปฺปชฺชติ นเหตุปจฺจยา. วิจิกิจฺฉานีวรณญฺจ สมฺปยุตฺตเก จ ขนฺเธ ปฏิจฺจ อวิชฺชา\-นีวรณํ. อุทฺธจฺจ\-นีวรณญฺจ สมฺปยุตฺตเก จ ขนฺเธ ปฏิจฺจ อวิชฺชา\-นีวรณํ (สํขิตฺตํ.) (1)}

\csromanmarkh{2. ปจฺจย\-ปจฺจนีย}\csromanmarkhii{2. สงฺขฺยาวาร}\csromanheadernomark{2}{2. ปจฺจย\-ปจฺจนีย 2. สงฺขฺยาวาร}
\vspace{\tipitakaparskip}
\csromancenter{\textbf{นเหตุ}ยา ตีณิ, นอธิปติยา นว, นปุเรชาเต นว, นปจฺฉาชาเต นว, นอาเสวเน นว, นกมฺเม ตีณิ, นวิปาเก นว, นวิปฺปยุตฺเต นว.}
\csromancenter{เหตุปจฺจยา นอธิปติยา นว ฯเปฯ นวิปฺปยุตฺเต นว.}
\proseitem{99}{\textbf{นเหตุ}ปจฺจยา อารมฺมเณ ตีณิ, อนนฺตเร ตีณิ, สมนนฺตเร ตีณิ, (สพฺพตฺถ ตีณิ.) มคฺเค ตีณิ ฯเปฯ อวิคเต ตีณิ.}

\prose{(เอวํ สหชาตวาโรปิ ปจฺจยวาโรปิ นิสฺสยวาโรปิ สํสฏฺฐ\-วาโรปิ สมฺปยุตฺตวาโรปิ ปฏิจฺจ\-วาร\-สทิสา, นินฺนานากรณา.)}

\vspace{\tipitakaparskip}
\csromancenter{48. นีวรณ\-นีวรณ\-สมฺปยุตฺตทุก 7. ปญฺหาวาร}
\csromanmarkh{1. ปจฺจยานุโลม}\csromanmarkhii{1. วิภงฺควาร}\csromanheadernomark{2}{1. \textbf{ปจฺจยานุโลม 1. วิภงฺควาร}}
\csromanheader{2}{\textbf{เหตุ}}
\proseitem{100}{\csromanfolio{333}นีวรโณ เจว นีวรณ\-สมฺปยุตฺโต จ ธมฺโม นีวรณสฺส เจว นีวรณ\-สมฺปยุตฺตสฺส จ ธมฺมสฺส เหตุปจฺจเยน ปจฺจโย. นีวรณา เจว นีวรณ\-สมฺปยุตฺตา จ เหตู สมฺปยุตฺตกานํ นีวรณานํ เหตุปจฺจเยน ปจฺจโย. (1)}

\prose{นีวรโณ เจว นีวรณ\-สมฺปยุตฺโต จ ธมฺโม นีวรณ\-สมฺปยุตฺตสฺส เจว โน จ นีวรณสฺส ธมฺมสฺส เหตุปจฺจเยน ปจฺจโย. นีวรณา เจว นีวรณ\-สมฺปยุตฺตา จ เหตู สมฺปยุตฺตกานํ ขนฺธานํ เหตุปจฺจเยน ปจฺจโย. (2)}

\prose{นีวรโณ เจว นีวรณ\-สมฺปยุตฺโต จ ธมฺโม นีวรณสฺส เจว นีวรณ\-สมฺปยุตฺตสฺส จ นีวรณ\-สมฺปยุตฺตสฺส เจว โน จ นีวรณสฺส จ ธมฺมสฺส เหตุปจฺจเยน ปจฺจโย. นีวรณา เจว นีวรณ\-สมฺปยุตฺตา จ เหตู สมฺปยุตฺตกานํ ขนฺธานํ นีวรณานญฺจ เหตุปจฺจเยน ปจฺจโย. (3)}

\proseitem{111}{นีวรโณ เจว นีวรณ\-สมฺปยุตฺโต จ ธมฺโม นีวรณสฺส เจว นีวรณ\-สมฺปยุตฺตสฺส จ ธมฺมสฺส อารมฺมณ\-ปจฺจเยน ปจฺจโย. นีวรเณ อารพฺภ นีวรณา อุปฺปชฺชนฺติ. (มูลํ กาตพฺพํ.) นีวรเณ อารพฺภ นีวรณ\-สมฺปยุตฺตา เจว โน จ นีวรณา ขนฺธา อุปฺปชฺชนฺติ. (มูลํ กาตพฺพํ.) นีวรเณ อารพฺภ นีวรณา จ สมฺปยุตฺตกา จ ขนฺธา อุปฺปชฺชนฺติ. (3)}

\prose{นีวรณ\-สมฺปยุตฺโต เจว โน จ นีวรโณ ธมฺโม นีวรณ\-สมฺปยุตฺตสฺส เจว โน จ นีวรณสฺส ธมฺมสฺส อารมฺมณ\-ปจฺจเยน ปจฺจโย. นีวรณ\-สมฺปยุตฺเต เจว โน จ นีวรเณ ขนฺเธ อารพฺภ นีวรณ\-สมฺปยุตฺตา เจว โน จ นีวรณา ขนฺธา อุปฺปชฺชนฺติ. (มูลํ ปุจฺฉิตพฺพํ.) นีวรณ\-สมฺปยุตฺเต เจว โน จ นีวรเณ ขนฺเธ อารพฺภ นีวรณา อุปฺปชฺชนฺติ. (มูลํ ปุจฺฉิตพฺพํ.) นีวรณ\-สมฺปยุตฺเต เจว โน จ นีวรเณ ขนฺเธ อารพฺภ นีวรณา จ สมฺปยุตฺตกา จ ขนฺธา อุปฺปชฺชนฺติ. (3)}

\prose{\csromanfolio{334}นีวรโณ เจว นีวรณ\-สมฺปยุตฺโต จ นีวรณ\-สมฺปยุตฺโต เจว โน จ นีวรโณ จ ธมฺมา นีวรณสฺส เจว นีวรณ\-สมฺปยุตฺตสฺส จ ธมฺมสฺส อารมฺมณ\-ปจฺจเยน ปจฺจโย. ตีณิ. (3)}

\csromanheader{2}{\textbf{อธิปตฺยาทิ}}
\proseitem{102}{นีวรโณ เจว นีวรณ\-สมฺปยุตฺโต จ ธมฺโม นีวรณสฺส เจว นีวรณ\-สมฺปยุตฺตสฺส จ ธมฺมสฺส อธิปติ\-ปจฺจเยน ปจฺจโย. อารมฺมณา\-ธิปติ. ตีณิ. (ครุกา\-รมฺมณาเยว.)}

\prose{นีวรณ\-สมฺปยุตฺโต เจว โน จ นีวรโณ ธมฺโม นีวรณ\-สมฺปยุตฺตสฺส เจว โน จ นีวรณสฺส ธมฺมสฺส อธิปติ\-ปจฺจเยน ปจฺจโย. อารมฺมณา\-ธิปติ, สหชาตา\-ธิปติ.  อารมฺมณา\-ธิปติ\csromandash{}นีวรณ\-สมฺปยุตฺเต เจว โน จ นีวรเณ ขนฺเธ ครุํ กตฺวา นีวรณ\-สมฺปยุตฺตา เจว โน จ นีวรณา ขนฺธา อุปฺปชฺชนฺติ. สหชาตา\-ธิปติ\csromandash{}นีวรณ\-สมฺปยุตฺตา เจว โน จ นีวรณา\-ธิปติ สมฺปยุตฺตกานํ ขนฺธานํ อธิปติ\-ปจฺจเยน ปจฺจโย. (มูลํ กาตพฺพํ.) อารมฺมณา\-ธิปติ\csromandash{}นีวรณ\-สมฺปยุตฺเต เจว โน จ นีวรเณ ขนฺเธ ครุํ กตฺวา นีวรณา อุปฺปชฺชนฺติ. สหชาตา\-ธิปติ\csromandash{} นีวรณ\-สมฺปยุตฺตา เจว โน จ นีวรณา\-ธิปติ สมฺปยุตฺตกานํ นีวรณานํ อธิปติ\-ปจฺจเยน ปจฺจโย. (มูลํ กาตพฺพํ.) อารมฺมณา\-ธิปติ\csromandash{}นีวารณสมฺปยุตฺเต เจว โน จ นีวรเณ ขนฺเธ ครุํ กตฺวา นีวรณา จ สมฺปยุตฺตกา จ ขนฺธา อุปฺปชฺชนฺติ. สหชาตา\-ธิปติ\csromandash{}นีวรณ\-สมฺปยุตฺตา เจว โน จ นีวรณา\-ธิปติ สมฺปยุตฺตกานํ ขนฺธานํ นีวรณานญฺจ อธิปติ\-ปจฺจเยน ปจฺจโย. (3)}

\proseitem{103}{นีวรโณ เจว นีวรณ\-สมฺปยุตฺโต จ นีวรณ\-สมฺปยุตฺโต เจว โน จ นีวรโณ จ ธมฺมา นีวรณสฺส เจว นีวรณ\-สมฺปยุตฺตสฺส จ ธมฺมสฺส อธิปติ\-ปจฺจเยน ปจฺจโย. อารมฺมณา\-ธิปติ. ตีณิ. (3)}

\prose{อนนฺตร\-ปจฺจเยน ปจฺจโย. (อาวชฺชนาปิ วุฏฺฐานมฺปิ นตฺถิ, สพฺพตฺถ ปุริมา ปุริมา กาตพฺพา.) สมนนฺตร\-ปจฺจเยน ปจฺจโย. นว. สหชาต\-ปจฺจเยน ปจฺจโย. นว. อญฺญมญฺญ\-ปจฺจเยน ปจฺจโย. นว. นิสฺสย\-ปจฺจเยน ปจฺจโย. นว. อุปนิสฺสย\-ปจฺจเยน ปจฺจโย. นว. (อารมฺมณ\-สทิสํ, วิปาโก นตฺถิ.) อาเสวน\-ปจฺจเยน ปจฺจโย. ปญฺจ.}

\csromanheader{2}{48. นีวรณ\-นีวรณ\-สมฺปยุตฺตทุก กมฺม}
\proseitem{104}{\csromanfolio{335}นีวรณ\-สมฺปยุตฺโต เจว โน จ นีวรโณ ธมฺโม นีวรณ\-สมฺปยุตฺตสฺส เจว โน จ นีวรณสฺส ธมฺมสฺส กมฺมปจฺจเยน ปจฺจโย. นีวรณ\-สมฺปยุตฺตา เจว โน จ นีวรณา เจตนา สมฺปยุตฺตกานํ ขนฺธานํ กมฺมปจฺจเยน ปจฺจโย. (มูลํ กาตพฺพํ.) นีวรณ\-สมฺปยุตฺตา เจว โน จ นีวรณา เจตนา สมฺปยุตฺตกานํ นีวรณานํ กมฺมปจฺจเยน ปจฺจโย. (มูลํ กาตพฺพํ.) นีวรณ\-สมฺปยุตฺตา เจว โน จ นีวรณา เจตนา สมฺปยุตฺตกานํ ขนฺธานํ นีวรณานญฺจ กมฺมปจฺจเยน ปจฺจโย. (3)}

\csromanheader{2}{\textbf{อาหาราทิ}}
\proseitem{105}{นีวรณ\-สมฺปยุตฺโต เจว โน จ นีวรโณ ธมฺโม นีวรณ\-สมฺปยุตฺตสฺส เจว โน จ นีวรณสฺส ธมฺมสฺส อาหารปจฺจเยน ปจฺจโย. ตีณิ. อินฺทฺริย\-ปจฺจเยน ปจฺจโย. ตีณิ. ฌานปจฺจเยน ปจฺจโย. ตีณิ. มคฺคปจฺจเยน ปจฺจโย. ตีณิ. สมฺปยุตฺต\-ปจฺจเยน ปจฺจโย. นว. อตฺถิปจฺจเยน ปจฺจโย. นว. นตฺถิ\-ปจฺจเยน ปจฺจโย. นว. วิคต\-ปจฺจเยน ปจฺจโย. นว. อวิคต\-ปจฺจเยน ปจฺจโย. นว.}

\csromanmarkh{1. ปจฺจยานุโลม}\csromanmarkhii{2. สงฺขฺยาวาร}\csromanheadernomark{2}{1. \textbf{ปจฺจยานุโลม 2. สงฺขฺยาวาร}}
\proseitem{106}{เหตุยา ตีณิ, อารมฺมเณ นว, อธิปติยา นว, อนนฺตเร นว, สมนนฺตเร นว, สหชาเต นว, อญฺญมญฺเญ นว, นิสฺสเย นว, อุปนิสฺสเย นว, อาเสวเน นว, กมฺเม ตีณิ, อาหาเร ตีณิ, อินฺทฺริเย ตีณิ, ฌาเน ตีณิ, มคฺเค ตีณิ, สมฺปยุตฺเต นว, อตฺถิยา นว, นตฺถิยา นว, วิคเต นว, อวิคเต นว.}

\proseitem{107}{นีวรโณ เจว นีวรณ\-สมฺปยุตฺโต จ ธมฺโม นีวรณสฺส เจว นีวรณ\-สมฺปยุตฺตสฺส จ ธมฺมสฺส อารมฺมณ\-ปจฺจเยน ปจฺจโย. สหชาต\-ปจฺจเยน ปจฺจโย. อุปนิสฺสย\-ปจฺจเยน ปจฺจโย. (เอวํ นวปิ ตีสุ ปเทสุ ปริวตฺเตตพฺพา.)}

\csromanmarkh{2. ปจฺจย\-ปจฺจนีย}\csromanmarkhii{2. สงฺขฺยาวาร}\csromanheadernomark{2}{2. ปจฺจย\-ปจฺจนีย 2. สงฺขฺยาวาร}
\vspace{\tipitakaparskip}
\csromancenter{นเหตุยา นว, นอารมฺมเณ นว ฯเปฯ โนอวิคเต นว.}
\proseitem{109}{\csromanfolio{336}เหตุปจฺจยา นอารมฺมเณ ตีณิ, นอธิปติยา ตีณิ, นอนนฺตเร ตีณิ, นสมนนฺตเร ตีณิ, นอุปนิสฺสเย ตีณิ ฯเปฯ นมคฺเค ตีณิ, นวิปฺปยุตฺเต ตีณิ, โนนตฺถิยา ตีณิ, โนวิคเต ตีณิ.}

\csromanheader{2}{110. นเหตุปจฺจยา อารมฺมเณ นว, อธิปติยา นว, (อนุโลมมาติกา) ฯเปฯ อวิคเต นว.}
\nitthitam{นีวรณ\-นีวรณ\-สมฺปยุตฺตทุกํ นิฏฺฐิตํ.}
\csromanmatikaanchor{mtk.46}
\csromantocmarkref{subsection}{49. นีวรณวิปฺปยุตฺตนีวรณิยทุก}{mtk.46}
\bookmark[dest={mtk.46},level=2]{49. นีวรณวิปฺปยุตฺตนีวรณิยทุก}
\csromanmarkh{49. นีวรณ\-วิปฺปยุตฺต\-นีวรณิยทุก}\csromanmarkhii{1. ปฏิจฺจวาร}\csromanheadernomark{2}{49. นีวรณ\-วิปฺปยุตฺต\-นีวรณิยทุก 1. ปฏิจฺจวาร}
\proseitem{111}{นีวรณ\-วิปฺปยุตฺตํ นีวรณิยํ ธมฺมํ ปฏิจฺจ นีวรณ\-วิปฺปยุตฺโต นีวรณิโย ธมฺโม อุปฺปชฺชติ เหตุปจฺจยา. นีวรณ\-วิปฺปยุตฺตํ นีวรณิยํ เอกํ ขนฺธํ ปฏิจฺจ ตโย ขนฺธา จิตฺต\-สมุฏฺฐานญฺจ รูปํ ฯเปฯ เทฺว ขนฺเธ ฯเปฯ. ปฏิสนฺธิ\-กฺขเณ ฯเปฯ. (สํขิตฺตํ.)}

\prosewithcloser{(ยถา จูฬนฺตรทุเก โลกิยทุกํ, เอวํ กาตพฺพํ, นินฺนานากรณํ.)}{นีวรณ\-วิปฺปยุตฺต\-นีวรณิยทุกํ นิฏฺฐิตํ.}
\nitthitam{นีวรณ\-โคจฺฉกํ นิฏฺฐิตํ.}
\csromanmatikaanchor{mtk.47}
\csromanmatikaanchor{mtk.48}
\csromantocmarknopage{section}{9. ปรามาสโคจฺฉก}{mtk.47}
\bookmark[dest={mtk.47},level=1]{9. ปรามาสโคจฺฉก}
\csromantocmarkref{subsection}{50. ปรามาสทุก}{mtk.48}
\bookmark[dest={mtk.48},level=2]{50. ปรามาสทุก}
\csromanmarkh{50. ปรามาสทุก}\csromanmarkhii{9. ปรามาส\-โคจฺฉก}\csromanheadernomark{2}{50. ปรามาสทุก 9. ปรามาส\-โคจฺฉก}
\csromanmarkh{50. ปรามาสทุก}\csromanmarkhii{1. ปฏิจฺจวาร}\csromanheadernomark{2}{50. ปรามาสทุก 1. ปฏิจฺจวาร}
\csromanmarkh{1. ปจฺจยานุโลม}\csromanmarkhii{1. วิภงฺควาร}\csromanheadernomark{2}{1. \textbf{ปจฺจยานุโลม 1. วิภงฺควาร}}
\csromanheader{2}{\textbf{เหตุ}}
\proseitem{1}{\csromanfolio{337}ปรามาสํ ธมฺมํ ปฏิจฺจ โนปรามาโส ธมฺโม อุปฺปชฺชติ เหตุปจฺจยา. ปรามาสํ ปฏิจฺจ สมฺปยุตฺตกา ขนฺธา จิตฺต\-สมุฏฺฐานญฺจ รูปํ. (1)}

\proseitem{2}{โนปรามาสํ ธมฺมํ ปฏิจฺจ โนปรามาโส ธมฺโม อุปฺปชฺชติ เหตุปจฺจยา. โนปรามาสํ เอกํ ขนฺธํ ปฏิจฺจ ตโย ขนฺธา จิตฺต\-สมุฏฺฐานญฺจ รูปํ ฯเปฯ เทฺว ขนฺเธ ฯเปฯ. ปฏิสนฺธิ\-กฺขเณ ฯเปฯ. (ยาว อสญฺญสตฺตา มหาภูตา.) (1)}

\prose{โนปรามาสํ ธมฺมํ ปฏิจฺจ ปรามาโส ธมฺโม อุปฺปชฺชติ เหตุปจฺจยา. โนปรามาเส ขนฺเธ ปฏิจฺจ ปรามาโส. (2)}

\prose{โนปรามาสํ ธมฺมํ ปฏิจฺจ ปรามาโส จ โนปรามาโส จ ธมฺมา อุปฺปชฺชนฺติ เหตุปจฺจยา. โนปรามาสํ เอกํ ขนฺธํ ปฏิจฺจ ตโย ขนฺธา ปรามาโส จ จิตฺต\-สมุฏฺฐานญฺจ รูปํ ฯเปฯ เทฺว ขนฺเธ ฯเปฯ. (3)}

\proseitem{3}{ปรามาสญฺจ โนปรามาสญฺจ ธมฺมํ ปฏิจฺจ โนปรามาโส ธมฺโม อุปฺปชฺชติ เหตุปจฺจยา. โนปรามาสํ เอกํ ขนฺธญฺจ ปรามาสญฺจ ปฏิจฺจ ตโย ขนฺธา จิตฺต\-สมุฏฺฐานญฺจ รูปํ ฯเปฯ เทฺว ขนฺเธ จ ฯเปฯ. ปรามาเส จ มหาภูเต จ ปฏิจฺจ จิตฺต\-สมุฏฺฐานํ รูปํ. (1)}

\csromanmarkh{1. ปจฺจยานุโลม}\csromanmarkhii{2. สงฺขฺยาวาร}\csromanheadernomark{2}{1. \textbf{ปจฺจยานุโลม 2. สงฺขฺยาวาร}}
\vspace{\tipitakaparskip}
\csromancenter{\textbf{เหตุ}ยา ปญฺจ, อารมฺมเณ ปญฺจ, (สพฺพตฺถ ปญฺจ.) วิปาเก เอกํ ฯเปฯ อวิคเต ปญฺจ.}
\csromancenter{อนุโลมํ.}
\csromanmarkh{2. ปจฺจย\-ปจฺจนีย}\csromanmarkhii{1. วิภงฺควาร}\csromanheadernomark{2}{2. ปจฺจย\-ปจฺจนีย 1. วิภงฺควาร}
\proseitem{5}{\csromanfolio{338}โนปรามาสํ ธมฺมํ ปฏิจฺจ โนปรามาโส ธมฺโม อุปฺปชฺชติ นเหตุปจฺจยา. อเหตุกํ โนปรามาสํ เอกํ ขนฺธํ ปฏิจฺจ ตโย ขนฺธา จิตฺต\-สมุฏฺฐานญฺจ รูปํ ฯเปฯ เทฺว ขนฺเธ ฯเปฯ. อเหตุก\-ปฏิสนฺธิ\-กฺขเณ ฯเปฯ. (ยาว อสญฺญสตฺตา.) วิจิกิจฺฉา\-สหคเต อุทฺธจฺจ\-สหคเต ขนฺเธ ปฏิจฺจ วิจิกิจฺฉา\-สหคโต อุทฺธจฺจ\-สหคโต โมโห. (1)}

\prose{นอารมฺมณ}

\proseitem{6}{ปรามาสํ ธมฺมํ ปฏิจฺจ โนปรามาโส ธมฺโม อุปฺปชฺชติ นอารมฺมณ\-ปจฺจยา. ปรามาสํ ปฏิจฺจ จิตฺต\-สมุฏฺฐานํ รูปํ. (1)}

\prose{โนปรามาสํ ธมฺมํ ปฏิจฺจ โนปรามาโส ธมฺโม อุปฺปชฺชติ นอารมฺมณ\-ปจฺจยา. โนปรามาเส ขนฺเธ ปฏิจฺจ จิตฺต\-สมุฏฺฐานํ รูปํ. ปฏิสนฺธิ\-กฺขเณ ฯเปฯ. (ยาว อสญฺญสตฺตา.) (1)}

\prose{ปรามาสญฺจ โนปรามาสญฺจ ธมฺมํ ปฏิจฺจ โนปรามาโส ธมฺโม อุปฺปชฺชติ. นอารมฺมณ\-ปจฺจยา. ปรามาสญฺจ สมฺปยุตฺตเก จ ขนฺเธ ปฏิจฺจ จิตฺต\-สมุฏฺฐานํ รูปํ. (1)}

\csromanmarkh{2. ปจฺจย\-ปจฺจนีย}\csromanmarkhii{2. สงฺขฺยาวาร}\csromanheadernomark{2}{2. ปจฺจย\-ปจฺจนีย 2. สงฺขฺยาวาร}
\vspace{\tipitakaparskip}
\csromancenter{\textbf{นเหตุ}ยา เอกํ, นอารมฺมเณ ตีณิ, นอธิปติยา ปญฺจ, นอนนฺตเร ตีณิ, นสมนนฺตเร ตีณิ, นอญฺญมญฺเญ ตีณิ, นฺเอาปนิสฺสเย ตีณิ, นปุเรชาเต ปญฺจ, นปจฺฉาชาเต ปญฺจ, นอาเสวเน ปญฺจ, นกมฺเม ตีณิ, นวิปาเก ปญฺจ, นอาหาเร เอกํ, ไนนฺทฺริเย เอกํ, นฌาเน เอกํ, นมคฺเค เอกํ, นสมฺปยุตฺเต ตีณิ, นวิปฺปยุตฺเต ปญฺจ, โนนตฺถิยา ตีณิ, โนวิคเต ตีณิ.}
\csromanmarkh{50. ปรามาสทุก}\csromanmarkhii{3. ปจฺจยา\-นุโลม\-ปจฺจนีย}\csromanheadernomark{2}{50. ปรามาสทุก 3. ปจฺจยา\-นุโลม\-ปจฺจนีย}
\vspace{\tipitakaparskip}
\csromancenter{\textbf{เหตุ}ปจฺจยา นอารมฺมเณ ตีณิ, นอธิปติยา ปญฺจ ฯเปฯ นวิปาเก ปญฺจ ฯเปฯ นสมฺปยุตฺเต ตีณิ, นวิปฺปยุตฺเต ปญฺจ, โนนตฺถิยา ตีณิ, โนวิคเต ตีณิ.}
\csromancenter{นเหตุปจฺจยา อารมฺมเณ เอกํ, อนนฺตเร เอกํ, (สพฺพตฺถ เอกํ.) อวิคเต เอกํ.}
\csromancenter{(สหชาตวาโรปิ ปฏิจฺจ\-วาร\-สทิโส.)}
\csromanmarkh{50. ปรามาสทุก}\csromanmarkhii{3. ปจฺจยวาร}\csromanheadernomark{2}{50. ปรามาสทุก 3. ปจฺจยวาร}
\csromanmarkh{1. ปจฺจยานุโลม}\csromanmarkhii{1. วิภงฺควาร}\csromanheadernomark{2}{1. \textbf{ปจฺจยานุโลม 1. วิภงฺควาร}}
\csromanheader{2}{\textbf{เหตุ}}
\proseitem{10}{\csromanfolio{339}ปรามาสํ ธมฺมํ ปจฺจยา โนปรามาโส ธมฺโม อุปฺปชฺชติ เหตุปจฺจยา. ปรามาสํ ปจฺจยา สมฺปยุตฺตกา ขนฺธา จิตฺต\-สมุฏฺฐานญฺจ รูปํ. (1)}

\prose{โนปรามาสํ ธมฺมํ ปจฺจยา โนปรามาโส ธมฺโม อุปฺปชฺชติ เหตุปจฺจยา. โนปรามาสํ เอกํ ขนฺธํ ปจฺจยา ตโย ขนฺธา จิตฺต\-สมุฏฺฐานญฺจ รูปํ ฯเปฯ เทฺว ขนฺเธ ฯเปฯ. ปฏิสนฺธิ\-กฺขเณ ฯเปฯ. (ยาว อชฺฌตฺติกา มหาภูตา.) วตฺถุํ ปจฺจยา โนปรามาสา ขนฺธา. (1)}

\prose{โนปรามาสํ ธมฺมํ ปจฺจยา ปรามาโส ธมฺโม อุปฺปชฺชติ เหตุปจฺจยา. โนปรามาเส ขนฺเธ ปจฺจยา ปรามาโส. วตฺถุํ ปจฺจยา ปรามาโส. (2)}

\prose{โนปรามาสํ ธมฺมํ ปจฺจยา ปรามาโส จ โนปรามาโส จ ธมฺมา อุปฺปชฺชนฺติ เหตุปจฺจยา. โนปรามาสํ เอกํ ขนฺธํ ปจฺจยา ตโย ขนฺธา \csromanfolio{340}ปรามาโส จ จิตฺต\-สมุฏฺฐานญฺจ รูปํ ฯเปฯ เทฺว ขนฺเธ ฯเปฯ. วตฺถุํ ปจฺจยา ปรามาโส. มหาภูเต ปจฺจยา จิตฺต\-สมุฏฺฐานํ รูปํ. วตฺถุํ ปจฺจยา ปรามาโส จ สมฺปยุตฺตกา จ ขนฺธา. (3)}

\prose{ปรามาสญฺจ โนปรามาสญฺจ ธมฺมํ ปจฺจยา โนปรามาโส ธมฺโม อุปฺปชฺชติ เหตุปจฺจยา. โนปรามาสํ เอกํ ขนฺธญฺจ ปรามาสญฺจ ปจฺจยา ตโย ขนฺธา จิตฺต\-สมุฏฺฐานญฺจ รูปํ ฯเปฯ เทฺว ขนฺเธ จ ฯเปฯ. ปรามาสญฺจ สมฺปยุตฺตเก จ ขนฺเธ ปจฺจยา จิตฺต\-สมุฏฺฐานํ รูปํ. ปรามาสญฺจ มหาภูเต จ ปจฺจยา จิตฺต\-สมุฏฺฐานํ รูปํ. ปรามาสญฺจ วตฺถุญฺจ ปจฺจยา โนปรามาสา ขนฺธา. (สํขิตฺตํ.) (1)}

\csromanmarkh{1. ปจฺจยานุโลม}\csromanmarkhii{2. สงฺขฺยาวาร}\csromanheadernomark{2}{1. \textbf{ปจฺจยานุโลม 2. สงฺขฺยาวาร}}
\proseitem{10}{เหตุยา ปญฺจ, อารมฺมเณ ปญฺจ, อธิปติยา ปญฺจ, (สพฺพตฺถ ปญฺจ.) วิปาเก เอกํ ฯเปฯ อวิคเต ปญฺจ.}

\csromanmarkh{2. ปจฺจย\-ปจฺจนีย}\csromanmarkhii{1. วิภงฺควาร}\csromanheadernomark{2}{2. ปจฺจย\-ปจฺจนีย 1. วิภงฺควาร}
\proseitem{13}{โนปรามาสํ ธมฺมํ ปจฺจยา โนปรามาโส ธมฺโม อุปฺปชฺชติ นเหตุปจฺจยา. อเหตุกํ โนปรามาสํ เอกํ ขนฺธํ ปจฺจยา ตโย ขนฺธา จิตฺต\-สมุฏฺฐานญฺจ รูปํ ฯเปฯ เทฺว ขนฺเธ ฯเปฯ. อเหตุก\-ปฏิสนฺธิ\-กฺขเณ ฯเปฯ. (ยาว อสญฺญสตฺตา.) จกฺขายตนํ ปจฺจยา จกฺขุวิญฺญาณํ ฯเปฯ. กายายตนํ ปจฺจยา กายวิญฺญาณํ. วตฺถุํ ปจฺจยา อเหตุกา โนปรามาสา ขนฺธา. วิจิกิจฺฉา\-สหคเต อุทฺธจฺจ\-สหคเต ขนฺเธ จ วตฺถุญฺจ ปจฺจยา วิจิกิจฺฉา\-สหคโต อุทฺธจฺจ\-สหคโต โมโห. (สํขิตฺตํ.)}

\csromanmarkh{2. ปจฺจย\-ปจฺจนีย}\csromanmarkhii{2. สงฺขฺยาวาร}\csromanheadernomark{2}{2. ปจฺจย\-ปจฺจนีย 2. สงฺขฺยาวาร}
\vspace{\tipitakaparskip}
\csromancenter{\textbf{นเหตุ}ยา เอกํ, นอารมฺมเณ ตีณิ, นอธิปติยา ปญฺจ, นอนนฺตเร ตีณิ ฯเปฯ นฺเอาปนิสฺสเย ตีณิ, นปุเรชาเต ปญฺจ, นปจฺฉาชาเต}
\prosecont{\csromanfolio{341}ปญฺจ, อาเสวเน ปญฺจ, นกมฺเม ตีณิ, นวิปาเก ปญฺจ, นอาหาเร เอกํ, ไนนฺทฺริเย เอกํ, นฌาเน เอกํ, นมคฺเค เอกํ, นสมฺปยุตฺเต ตีณิ, นวิปฺปยุตฺเต ปญฺจ, โนนตฺถิยา ตีณิ, โนวิคเต ตีณิ.}

\vspace{\tipitakaparskip}
\csromancenter{\textbf{เหตุ}ปจฺจยา นอารมฺมเณ ตีณิ, นอธิปติยา ปญฺจ. (เอวํ สพฺพตฺถ กาตพฺพํ.)}
\csromancenter{นเหตุปจฺจยา อารมฺมเณ เอกํ ฯเปฯ อวิคเต เอกํ.}
\csromancenter{(นิสฺสยวาโร ปจฺจยวาร\-สทิโส.)}
\csromanmarkh{50. ปรามาสทุก}\csromanmarkhii{5. สํสฏฺฐวาร}\csromanheadernomark{2}{50. ปรามาสทุก 5. สํสฏฺฐวาร}
\prose{ปรามาสํ ธมฺมํ สํสฏฺโฐ โนปรามาโส ธมฺโม อุปฺปชฺชติ เหตุปจฺจยา. ปรามาสํ สํสฏฺฐา สมฺปยุตฺตกา ขนฺธา.}

\vspace{\tipitakaparskip}
\csromancenter{(เอวํ ปญฺจ ปญฺหา กาตพฺพา, อรูเปเยว, สํสฏฺฐ\-วาโรปิ สมฺปยุตฺตวาโรปิ เอวํ กาตพฺพา.)}
\csromanmarkh{50. ปรามาสทุก}\csromanmarkhii{7. ปญฺหาวาร}\csromanheadernomark{2}{50. ปรามาสทุก 7. ปญฺหาวาร}
\csromanmarkh{1. ปจฺจยานุโลม}\csromanmarkhii{1. วิภงฺควาร}\csromanheadernomark{2}{1. \textbf{ปจฺจยานุโลม 1. วิภงฺควาร}}
\vspace{\tipitakaparskip}
\csromancenter{\textbf{เหตุ}}
\proseitem{17}{โนปรามาโส ธมฺโม โนปรามาสสฺส ธมฺมสฺส เหตุปจฺจเยน ปจฺจโย. โนปรามาสา เหตู สมฺปยุตฺตกานํ ขนฺธานํ จิตฺตสมุฏฺฐานานญฺจ รูปานํ เหตุปจฺจเยน ปจฺจโย. ปฏิสนฺธิ\-กฺขเณ ฯเปฯ. (1)}

\prose{โนปรามาโส ธมฺโม ปรามาสสฺส ธมฺมสฺส เหตุปจฺจเยน ปจฺจโย. โนปรามาโส เหตู ปรามาสสฺส เหตุปจฺจเยน ปจฺจโย. (2)}

\prose{โนปรามาโส ธมฺโม ปรามาสสฺส จ โนปรามาสสฺส จ ธมฺมสฺส เหตุปจฺจเยน ปจฺจโย. โนปรามาสา เหตู สมฺปยุตฺตกานํ ขนฺธานํ ปรามาสสฺส จ จิตฺตสมุฏฺฐานานญฺจ รูปานํ เหตุปจฺจเยน ปจฺจโย. (3)}

\proseitem{18}{\csromanfolio{342}ปรามาโส ธมฺโม ปรามาสสฺส ธมฺมสฺส อารมฺมณ\-ปจฺจเยน ปจฺจโย. ปรามาสํ อารพฺภ ปรามาโส อุปฺปชฺชติ. (มูลํ กาตพฺพํ.) ปรามาสํ อารพฺภ โนปรามาสา ขนฺธา อุปฺปชฺชนฺติ. (มูลํ กาตพฺพํ.) ปรามาสํ อารพฺภ ปรามาโส จ สมฺปยุตฺตกา จ ขนฺธา อุปฺปชฺชนฺติ. (3)}

\proseitem{19}{โนปรามาโส ธมฺโม โนปรามาสสฺส ธมฺมสฺส อารมฺมณ\-ปจฺจเยน ปจฺจโย. ทานํ ฯเปฯ สีลํ ฯเปฯ อุโปสถกมฺมํ กตฺวา ตํ ปจฺจเวกฺขติ, อสฺสาเทติ อภินนฺทติ, ตํ อารพฺภ ราโค ฯเปฯ วิจิกิจฺฉา. อุทฺธจฺจํ. โทมนสฺสํ อุปฺปชฺชติ. ปุพฺเพ สุจิณฺณานิ ฯเปฯ. ฌานา วุฏฺฐหิตฺวา ฌานํ ฯเปฯ. อริยา มคฺคา วุฏฺฐหิตฺวา มคฺคํ ปจฺจเวกฺขนฺติ, ผลํ ปจฺจเวกฺขนฺติ, นิพฺพานํ ปจฺจเวกฺขนฺติ. นิพฺพานํ โคตฺรภุสฺส, โวทานสฺส, มคฺคสฺส, ผลสฺส, อาวชฺชนาย อารมฺมณ\-ปจฺจเยน ปจฺจโย. อริยา โนปรามาเส ปหีเน กิเลเส ฯเปฯ. วิกฺขมฺภิเต กิเลเส ฯเปฯ. ปุพฺเพ ฯเปฯ. จกฺขุํ ฯเปฯ วตฺถุํ โนปรามาเส ขนฺเธ อนิจฺจโต ฯเปฯ วิปสฺสติ, อสฺสาเทติ อภินนฺทติ, ตํ อารพฺภ ราโค ฯเปฯ วิจิกิจฺฉา. อุทฺธจฺจํ. โทมนสฺสํ อุปฺปชฺชติ. ทิพฺเพน จกฺขุนา รูปํ ปสฺสติ. ทิพฺพาย โสตธาตุยา สทฺทํ สุณาติ. เจโตปริย\-ญาเณน โนปรามาสจิตฺตสมงฺคิสฺส จิตฺตํ ชานาติ. อากาสา\-นญฺจา\-ยตนํ วิญฺญาณญฺจา\-ยตนสฺส ฯเปฯ. อากิญฺจญฺญา\-ยตนํ เนว\-สญฺญา\-นา\-สญฺญา\-ยตนสฺส ฯเปฯ. รูปายตนํ จกฺขุ\-วิญฺญาณสฺส ฯเปฯ. โผฏฺฐพฺพา\-ยตนํ กายวิญฺญาณสฺส ฯเปฯ โนปรามาสา ขนฺธา อิทฺธิวิธ\-ญาณสฺส, เจโตปริย\-ญาณสฺส, ปุพฺเพ\-นิวาสา\-นุสฺสติ\-ญาณสฺส, ยถากมฺมูปคญาณสฺส อนาคตํสญาณสฺส, อาวชฺชนาย อารมฺมณ\-ปจฺจเยน ปจฺจโย. (1)}

\prose{โนปรามาโส ธมฺโม ปรามาสสฺส ธมฺมสฺส อารมฺมณ\-ปจฺจเยน ปจฺจโย. ทานํ ฯเปฯ สีลํ ฯเปฯ อุโปสถกมฺมํ ฯเปฯ ตํ อสฺสาเทติ, อภินนฺทติ, ตํ อารพฺภ ทิฏฺฐิ อุปฺปชฺชติ. ปุพฺเพ ฯเปฯ. ฌานา ฯเปฯ. จกฺขุํ ฯเปฯ วตฺถุํ โนปรามาเส ขนฺเธ อสฺสาเทติ อภินนฺทติ. ตํ อารพฺภ ทิฏฺฐิ อุปฺปชฺชติ. (2)}

\prose{โนปรามาโส ธมฺโม ปรามาสสฺส จ โนปรามาสสฺส จ ธมฺมสฺส อารมฺมณ\-ปจฺจเยน ปจฺจโย. ทานํ ฯเปฯ สีลํ ฯเปฯ อุโปสถกมฺมํ ฯเปฯ ตํ อสฺสาเทติ อภินนฺทติ, ตํ อารพฺภ ปรามาโส จ สมฺปยุตฺตกา จ ขนฺธา อุปฺปชฺชนฺติ. ปุพฺเพ ฯเปฯ. ฌานา ฯเปฯ. จกฺขุํ ฯเปฯ วตฺถุํ โนปรามาเส ขนฺเธ \csromanfolio{343}อสฺสาเทติ อภินนฺทติ, ตํ อารพฺภ ปรามาโส จ สมฺปยุตฺตกา จ ขนฺธา อุปฺปชฺชนฺติ. (3)}

\prose{ปรามาโส จ โนปรามาโส จ ธมฺมา ปรามาสสฺส ธมฺมสฺส อารมฺมณ\-ปจฺจเยน ปจฺจโย. ปรามาสญฺจ สมฺปยุตฺตเก จ ขนฺเธ อารพฺภ ปรามาโส อุปฺปชฺชติ. ตีณิ.}

\proseitem{21}{ปรามาโส ธมฺโม ปรามาสสฺส ธมฺมสฺส อธิปติ\-ปจฺจเยน ปจฺจโย. ปรามาสํ ครุํ กตฺวา ปรามาโส อุปฺปชฺชติ. ตีณิ. (อารมฺมณา\-ธิปติเยว กาตพฺพา.)}

\proseitem{19}{โนปรามาโส ธมฺโม โนปรามาสสฺส ธมฺมสฺส อธิปติ\-ปจฺจเยน ปจฺจโย. อารมฺมณา\-ธิปติ, สหชาตา\-ธิปติ.  อารมฺมณา\-ธิปติ\csromandash{}ทานํ ฯเปฯ สีลํ ฯเปฯ อุโปสถกมฺมํ ฯเปฯ ตํ ครุํ กตฺวา ปจฺจเวกฺขติ, อสฺสาเทติ อภินนฺทติ, ตํ ครุํ กตฺวา ราโค อุปฺปชฺชติ. ปุพฺเพ ฯเปฯ. ฌานา ฯเปฯ. อริยา มคฺคา วุฏฺฐหิตฺวา มคฺคํ ครุํ กตฺวา ฯเปฯ. ผลํ ครุํ กตฺวา ฯเปฯ. นิพฺพานํ ครุํ กตฺวา ฯเปฯ. นิพฺพานํ โคตฺรภุสฺส, โวทานสฺส, มคฺคสฺส, ผลสฺส อธิปติ\-ปจฺจเยน ปจฺจโย. จกฺขุํ ฯเปฯ วตฺถุํ โนปรามาเส ขนฺเธ ครุํ กตฺวา อสฺสาเทติ อภินนฺทติ, ตํ ครุํ กตฺวา ราโค อุปฺปชฺชติ. สหชาตา\-ธิปติ\csromandash{} โนปรามาสา\-ธิปติ สมฺปยุตฺตกานํ ขนฺธานํ จิตฺตสมุฏฺฐานานญฺจ รูปานํ อธิปติ\-ปจฺจเยน ปจฺจโย. (1)}

\prose{โนปรามาโส ธมฺโม ปรามาสสฺส ธมฺมสฺส อธิปติ\-ปจฺจเยน ปจฺจโย. อารมฺมณา\-ธิปติ, สหชาตา\-ธิปติ.  อารมฺมณา\-ธิปติ\csromandash{}ทานํ ฯเปฯ สีลํ ฯเปฯ อุโปสถกมฺมํ ฯเปฯ. ปุพฺเพ ฯเปฯ. ฌานา ฯเปฯ. จกฺขุํ ฯเปฯ วตฺถุํ โนปรามาเส ขนฺเธ ครุํ กตฺวา อสฺสาเทติ อภินนฺทติ, ตํ ครุํ กตฺวา ทิฏฺฐิ อุปฺปชฺชติ. สหชาตา\-ธิปติ\csromandash{}โนปรามาสา\-ธิปติ ปรามาสสฺส อธิปติ\-ปจฺจเยน ปจฺจโย. (2)}

\prose{โนปรามาโส ธมฺโม ปรามาสสฺส จ โนปรามาสสฺส จ ธมฺมสฺส อธิปติ\-ปจฺจเยน ปจฺจโย. อารมฺมณา\-ธิปติ, สหชาตา\-ธิปติ.  อารมฺมณา\-ธิปติ\csromandash{}ทานํ ฯเปฯ สีลํ ฯเปฯ อุโปสถกมฺมํ ฯเปฯ. ปุพฺเพ ฯเปฯ. \csromanfolio{344}ฌานา ฯเปฯ. จกฺขุํ ฯเปฯ วตฺถุํ โนปรามาเส ขนฺเธ ครุํ กตฺวา อสฺสาเทติ อภินนฺทติ, ตํ ครุํ กตฺวา ปรามาโส จ สมฺปยุตฺตกา จ ขนฺธา อุปฺปชฺชนฺติ. สหชาตา\-ธิปติ\csromandash{}โนปรามาสา\-ธิปติ สมฺปยุตฺตกานํ ขนฺธานํ ปรามาสสฺส จ จิตฺตสมุฏฺฐานานญฺจ รูปานํ อธิปติ\-ปจฺจเยน ปจฺจโย. (3)}

\proseitem{22}{ปรามาโส จ โนปรามาโส จ ธมฺมา ปรามาสสฺส ธมฺมสฺส อธิปติ\-ปจฺจเยน ปจฺจโย. อารมฺมณา\-ธิปติ\csromandash{}ปรามาสญฺจ สมฺปยุตฺตเก จ ขนฺเธ ครุํ กตฺวา ปรามาโส. ตีณิ.}

\proseitem{23}{ปรามาโส ธมฺโม ปรามาสสฺส ธมฺมสฺส อนนฺตร\-ปจฺจเยน ปจฺจโย. ปุริโม ปุริโม ปรามาโส ปจฺฉิมสฺส ปจฺฉิมสฺส ปรามาสสฺส อนนฺตร\-ปจฺจเยน ปจฺจโย. (มูลํ กาตพฺพํ.) ปุริโม ปุริโม ปรามาโส ปจฺฉิมานํ ปจฺฉิมานํ โนปรามาสานํ ขนฺธานํ อนนฺตร\-ปจฺจเยน ปจฺจโย. ปรามาโส วุฏฺฐานสฺส อนนฺตร\-ปจฺจเยน ปจฺจโย. (มูลํ กาตพฺพํ.) ปุริโม ปุริโม ปรามาโส ปจฺฉิมสฺส ปจฺฉิมสฺส ปรามาสสฺส สมฺปยุตฺตกานญฺจ ขนฺธานํ อนนฺตร\-ปจฺจเยน ปจฺจโย. (3)}

\proseitem{24}{โนปรามาโส ธมฺโม โนปรามาสสฺส ธมฺมสฺส อนนฺตร\-ปจฺจเยน ปจฺจโย. ปุริมา ปุริมา โนปรามาสา ขนฺธา ปจฺฉิมานํ ปจฺฉิมานํ โนปรามาสานํ ขนฺธานํ อนนฺตร\-ปจฺจเยน ปจฺจโย ฯเปฯ. อนุโลมํ ผลสมาปตฺติยา อนนฺตร\-ปจฺจเยน ปจฺจโย. (มูลํ กาตพฺพํ.) ปุริมา ปุริมา โนปรามาสา ขนฺธา ปจฺฉิมสฺส ปจฺฉิมสฺส ปรามาสสฺส อนนฺตร\-ปจฺจเยน ปจฺจโย. อาวชฺชนา ปรามาสสฺส อนนฺตร\-ปจฺจเยน ปจฺจโย. (มูลํ กาตพฺพํ.) ปุริมา ปุริมา โนปรามาสา ขนฺธา ปจฺฉิมสฺส ปจฺฉิมสฺส ปรามาสสฺส จ สมฺปยุตฺตกานญฺจ ขนฺธานํ อนนฺตร\-ปจฺจเยน ปจฺจโย. อาวชฺชนา ปรามาสสฺส จ สมฺปยุตฺตกานญฺจ ขนฺธานํ อนนฺตร\-ปจฺจเยน ปจฺจโย. (3)}

\proseitem{25}{ปรามาโส จ โนปรามาโส จ ธมฺมา ปรามาสสฺส ธมฺมสฺส อนนฺตร\-ปจฺจเยน ปจฺจโย. ปุริโม ปุริโม ปรามาโส จ สมฺปยุตฺตกา จ ขนฺธา ปจฺฉิมสฺส ปจฺฉิมสฺส ปรามาสสฺส อนนฺตร\-ปจฺจเยน ปจฺจโย. (มูลํ \csromanfolio{345}กาตพฺพํ.) ปุริโม ปุริโม ปรามาโส จ สมฺปยุตฺตกา จ ขนฺธา ปจฺฉิมานํ ปจฺฉิมานํ โนปรามาสานํ ขนฺธานํ อนนฺตร\-ปจฺจเยน ปจฺจโย. ปรามาโส จ สมฺปยุตฺตกา จ ขนฺธา วุฏฺฐานสฺส อนนฺตร\-ปจฺจเยน ปจฺจโย. (มูลํ กาตพฺพํ.) ปุริโม ปุริโม ปรามาโส จ สมฺปยุตฺตกา จ ขนฺธา ปจฺฉิมสฺส ปจฺฉิมสฺส ปรามาสสฺส สมฺปยุตฺตกานญฺจ ขนฺธานํ อนนฺตร\-ปจฺจเยน ปจฺจโย. (3)}

\csromanheader{2}{\textbf{สมนนฺตราทิ}}
\proseitem{25}{ปรามาโส ธมฺโม ปรามาสสฺส ธมฺมสฺส สมนนฺตร\-ปจฺจเยน ปจฺจโย. สหชาต\-ปจฺจเยน ปจฺจโย. อญฺญมญฺญ\-ปจฺจเยน ปจฺจโย. นิสฺสย\-ปจฺจเยน ปจฺจโย. ปญฺจ.}

\csromanheader{2}{\textbf{อุปนิสฺสย}}
\proseitem{27}{ปรามาโส ธมฺโม ปรามาสสฺส ธมฺมสฺส อุปนิสฺสย\-ปจฺจเยน ปจฺจโย. อารมฺมณู\-ปนิสฺสโย, อนนฺตรู\-ปนิสฺสโย, ปกตู\-ปนิสฺสโย ฯเปฯ. ปกตู\-ปนิสฺสโย\csromandash{}ปรามาโส ปรามาสสฺส อุปนิสฺสย\-ปจฺจเยน ปจฺจโย. ตีณิ.}

\proseitem{25}{โนปรามาโส ธมฺโม โนปรามาสสฺส ธมฺมสฺส อุปนิสฺสย\-ปจฺจเยน ปจฺจโย. อารมฺมณู\-ปนิสฺสโย, อนนฺตรู\-ปนิสฺสโย, ปกตู\-ปนิสฺสโย ฯเปฯ. ปกตู\-ปนิสฺสโย\csromandash{}สทฺธํ อุปนิสฺสาย ทานํ เทติ ฯเปฯ สมาปตฺติํ อุปฺปาเทติ. มานํ ชปฺเปติ. สีลํ ฯเปฯ. ปญฺญํ. ราคํ. โทสํ. โมหํ. มานํ. ปตฺถนํ. กายิกํ สุขํ ฯเปฯ. เสนาสนํ อุปนิสฺสาย ทานํ เทติ ฯเปฯ สมาปตฺติํ อุปฺปาเทติ. ปาณํ หนติ ฯเปฯ สํฆํ ภินฺทติ. สทฺธา ฯเปฯ. เสนาสนํ. สทฺธาย ฯเปฯ. ปญฺญาย. ราคสฺส ฯเปฯ ปตฺถนาย. กายิกสฺส สุขสฺส ฯเปฯ มคฺคสฺส, ผลสมาปตฺติยา อุปนิสฺสย\-ปจฺจเยน ปจฺจโย. (1)}

\prose{โนปรามาโส ธมฺโม ปรามาสสฺส ธมฺมสฺส อุปนิสฺสย\-ปจฺจเยน ปจฺจโย. อารมฺมณู\-ปนิสฺสโย, อนนฺตรู\-ปนิสฺสโย, ปกตู\-ปนิสฺสโย ฯเปฯ. ปกตู\-ปนิสฺสโย\csromandash{}สทฺธํ อุปนิสฺสาย ทิฏฺฐิํ คณฺหาติ. สีลํ ฯเปฯ. ปญฺญํ. ราคํ ฯเปฯ. ปตฺถนํ. กายิกํ สุขํ ฯเปฯ. เสนาสนํ อุปนิสฺสาย ทิฏฺฐิํ คณฺหาติ. สทฺธา ฯเปฯ. เสนาสนํ ปรามาสสฺส อุปนิสฺสย\-ปจฺจเยน ปจฺจโย. (2)}

\proseitem{28}{\csromanfolio{346}โนปรามาโส ธมฺโม ปรามาสสฺส จ โนปรามาสสฺส จ ธมฺมสฺส อุปนิสฺสย\-ปจฺจเยน ปจฺจโย. อารมฺมณู\-ปนิสฺสโย, อนนฺตรู\-ปนิสฺสโย, ปกตู\-ปนิสฺสโย ฯเปฯ. ปกตู\-ปนิสฺสโย\csromandash{}สทฺธํ อุปนิสฺสาย ทิฏฺฐิํ คณฺหาติ. สีลํ ฯเปฯ. เสนาสนํ อุปนิสฺสาย ทิฏฺฐิํ คณฺหาติ. สทฺธา ฯเปฯ. เสนาสนํ ปรามาสสฺส สมฺปยุตฺตกานญฺจ ขนฺธานํ อุปนิสฺสย\-ปจฺจเยน ปจฺจโย. (3)}

\proseitem{29}{ปรามาโส จ โนปรามาโส จ ธมฺมา ปรามาสสฺส ธมฺมสฺส อุปนิสฺสย\-ปจฺจเยน ปจฺจโย. ตีณิ อุปนิสฺสยา ปรามาโส จ สมฺปยุตฺตกา จ ขนฺธา ปรามาสสฺส อุปนิสฺสย\-ปจฺจเยน ปจฺจโย. ตีณิ.}

\proseitem{30}{โนปรามาโส ธมฺโม โนปรามาสสฺส ธมฺมสฺส ปุเรชาต\-ปจฺจเยน ปจฺจโย. อารมฺมณ\-ปุเรชาตํ, วตฺถุ\-ปุเรชาตํ.  อารมฺมณ\-ปุเรชาตํ\csromandash{}จกฺขุํ ฯเปฯ วตฺถุํ อนิจฺจโต ฯเปฯ วิปสฺสติ, อสฺสาเทติ อภินนฺทติ, ตํ อารพฺภ ราโค ฯเปฯ วิจิกิจฺฉา ฯเปฯ อุทฺธจฺจํ ฯเปฯ โทมนสฺสํ อุปฺปชฺชติ. ทิพฺเพน จกฺขุนา รูปํ ปสฺสติ. ทิพฺพาย โสตธาตุยา สทฺทํ สุณาติ. รูปายตนํ จกฺขุ\-วิญฺญาณสฺส ฯเปฯ. โผฏฺฐพฺพา\-ยตนํ กายวิญฺญาณสฺส ฯเปฯ. วตฺถุ\-ปุเรชาตํ\csromandash{}จกฺขายตนํ จกฺขุ\-วิญฺญาณสฺส ฯเปฯ กายายตนํ กายวิญฺญาณสฺส ฯเปฯ. วตฺถุ โนปรามาสานํ ขนฺธานํ ปุเรชาต\-ปจฺจเยน ปจฺจโย. (1)}

\prose{โนปรามาโส ธมฺโม ปรามาสสฺส ธมฺมสฺส ปุเรชาต\-ปจฺจเยน ปจฺจโย. อารมฺมณ\-ปุเรชาตํ, วตฺถุ\-ปุเรชาตํ.  อารมฺมณ\-ปุเรชาตํ\csromandash{} จกฺขุํ ฯเปฯ วตฺถุํ อสฺสาเทติ อภินนฺทติ, ตํ อารพฺภ ทิฏฺฐิ อุปฺปชฺชติ. วตฺถุ\-ปุเรชาตํ\csromandash{}วตฺถุ ปรามาสสฺส ปุเรชาต\-ปจฺจเยน ปจฺจโย. (2)}

\prose{โนปรามาโส ธมฺโม ปรามาสสฺส จ โนปรามาสสฺส จ ธมฺมสฺส ปุเรชาต\-ปจฺจเยน ปจฺจโย. อารมฺมณ\-ปุเรชาตํ, วตฺถุ\-ปุเรชาตํ.  อารมฺมณ\-ปุเรชาตํ\csromandash{}จกฺขุํ ฯเปฯ วตฺถุํ อสฺสาเทติ อภินนฺทติ, ตํ อารพฺภ ปรามาโส จ สมฺปยุตฺตกา จ ขนฺธา อุปฺปชฺชนฺติ. วตฺถุ\-ปุเรชาตํ\csromandash{}วตฺถุ ปรามาสสฺส จ สมฺปยุตฺตกานญฺจ ขนฺธานํ ปุเรชาต\-ปจฺจเยน ปจฺจโย. (3)}

\csromanheader{2}{50. ปรามาสทุก ปจฺฉาชาตา\-เสวน}
\proseitem{31}{\csromanfolio{347}ปรามาโส ธมฺโม โนปรามาสสฺส ธมฺมสฺส ปจฺฉาชาต\-ปจฺจเยน ปจฺจโย. ตีณิ. (ปจฺฉาชาตา กาตพฺพา.) อาเสวน\-ปจฺจเยน ปจฺจโย. นว.}

\proseitem{32}{โนปรามาโส ธมฺโม โนปรามาสสฺส ธมฺมสฺส กมฺมปจฺจเยน ปจฺจโย. สหชาตา, นานากฺขณิกา.  สหชาตา โนปรามาสา เจตนา สมฺปยุตฺตกานํ ขนฺธานํ จิตฺตสมุฏฺฐานานญฺจ รูปานํ กมฺมปจฺจเยน ปจฺจโย. ปฏิสนฺธิ\-กฺขเณ ฯเปฯ. นานากฺขณิกา โนปรามาสา เจตนา วิปากานํ ขนฺธานํ กฏตฺตา จ รูปานํ กมฺมปจฺจเยน ปจฺจโย. (มูลํ กาตพฺพํ.) โนปรามาสา เจตนา ปรามาสสฺส กมฺมปจฺจเยน ปจฺจโย. (มูลํ กาตพฺพํ.) โนปรามาสา เจตนา ปรามาสสฺส จ สมฺปยุตฺตกานญฺจ ขนฺธานํ จิตฺตสมุฏฺฐานานญฺจ รูปานํ กมฺมปจฺจเยน ปจฺจโย. (3)}

\csromanheader{2}{\textbf{วิปากาทิ}}
\proseitem{33}{โนปรามาโส ธมฺโม โนปรามาสสฺส ธมฺมสฺส วิปาก\-ปจฺจเยน ปจฺจโย. เอกํ. อาหารปจฺจเยน ปจฺจโย. ตีณิ. อินฺทฺริย\-ปจฺจเยน ปจฺจโย. ตีณิ. ฌานปจฺจเยน ปจฺจโย. ตีณิ.}

\prose{ปรามาโส ธมฺโม โนปรามาสสฺส ธมฺมสฺส มคฺคปจฺจเยน ปจฺจโย. ปรามาสานิ มคฺคงฺคานิ ฯเปฯ. (เอวํ ปญฺจ ปญฺหา กาตพฺพา.) สมฺปยุตฺต\-ปจฺจเยน ปจฺจโย. ปญฺจ}

\proseitem{34}{ปรามาโส ธมฺโม โนปรามาสสฺส ธมฺมสฺส วิปฺปยุตฺต\-ปจฺจเยน ปจฺจโย. สหชาตํ, ปจฺฉาชาตํ. (สํขิตฺตํ.) (1)}

\prose{โนปรามาโส ธมฺโม โนปรามาสสฺส ธมฺมสฺส วิปฺปยุตฺต\-ปจฺจเยน ปจฺจโย. สหชาตํ, ปุเรชาตํ, ปจฺฉาชาตํ. (สํขิตฺตํ.) (1)}

\prose{โนปรามาโส ธมฺโม ปรามาสสฺส ธมฺมสฺส วิปฺปยุตฺต\-ปจฺจเยน ปจฺจโย. ปุเรชาตํ วตฺถุ ปรามาสสฺส วิปฺปยุตฺต\-ปจฺจเยน ปจฺจโย. (2)}

\prose{\csromanfolio{348}โนปรามาโส ธมฺโม ปรามาสสฺส จ โนปรามาสสฺส จ ธมฺมสฺส วิปฺปยุตฺต\-ปจฺจเยน ปจฺจโย. ปุเรชาตํ วตฺถุ ปรามาสสฺส จ สมฺปยุตฺตกานญฺจ ขนฺธานํ วิปฺปยุตฺต\-ปจฺจเยน ปจฺจโย. (3)}

\prose{ปรามาโส จ โนปรามาโส จ ธมฺมา โนปรามาสสฺส ธมฺมสฺส วิปฺปยุตฺต\-ปจฺจเยน ปจฺจโย. สหชาตํ, ปจฺฉาชาตํ. (สํขิตฺตํ.) (1)}

\proseitem{35}{ปรามาโส ธมฺโม โนปรามาสสฺส ธมฺมสฺส อตฺถิปจฺจเยน ปจฺจโย. สหชาตํ, ปจฺฉาชาตํ.  \textbf{สหชาโต} ปรามาโส สมฺปยุตฺตกานญฺจ ขนฺธานํ จิตฺตสมุฏฺฐานานญฺจ รูปานํ อตฺถิปจฺจเยน ปจฺจโย. \textbf{ปจฺฉาชาโต} ปรามาโส ปุเรชาตสฺส อิมสฺส กายสฺส อตฺถิปจฺจเยน ปจฺจโย. (1)}

\proseitem{36}{โนปรามาโส ธมฺโม โนปรามาสสฺส ธมฺมสฺส อตฺถิปจฺจเยน ปจฺจโย. สหชาตํ, ปุเรชาตํ, ปจฺฉาชาตํ, อาหารํ, อินฺทฺริยํ. (สํขิตฺตํ.) (1)}

\prose{โนปรามาโส ธมฺโม ปรามาสสฺส ธมฺมสฺส อตฺถิปจฺจเยน ปจฺจโย. สหชาตํ, ปุเรชาตํ.  สหชาตา โนปรามาสา ขนฺธา ปรามาสสฺส อตฺถิปจฺจเยน ปจฺจโย. ปุเรชาตํ จกฺขุํ ฯเปฯ วตฺถุํ อสฺสาเทติ อภินนฺทติ, ตํ อารพฺภ ทิฏฺฐิ อุปฺปชฺชติ. วตฺถุ ปรามาสสฺส อตฺถิปจฺจเยน ปจฺจโย. (2)}

\prose{โนปรามาโส ธมฺโม ปรามาสสฺส จ โนปรามาสสฺส จ ธมฺมสฺส อตฺถิปจฺจเยน ปจฺจโย. สหชาตํ, ปุเรชาตํ.  สหชาโต โนปรามาโส เอโก ขนฺโธ ติณฺณนฺนํ ขนฺธานํ ปรามาสสฺส จ จิตฺตสมุฏฺฐานานญฺจ รูปานํ อตฺถิปจฺจเยน ปจฺจโย ฯเปฯ. ปุเรชาตํ จกฺขุํ ฯเปฯ วตฺถุํ อสฺสาเทติ อภินนฺทติ, ตํ อารพฺภ ปรามาโส จ สมฺปยุตฺตกา จ ขนฺธา อุปฺปชฺชนฺติ. วตฺถุ ปรามาสสฺส จ สมฺปยุตฺตกานญฺจ ขนฺธานํ อตฺถิปจฺจเยน ปจฺจโย. (3)}

\proseitem{37}{ปรามาโส จ โนปรามาโส จ ธมฺมา โนปรามาสสฺส ธมฺมสฺส อตฺถิปจฺจเยน ปจฺจโย. สหชาตํ, ปุเรชาตํ, ปจฺฉาชาตํ}

\prose{\csromanfolio{349}อาหารํ, อินฺทฺริยํ.  สหชาโต โนปรามาโส เอโก ขนฺโธ จ ปรามาโส จ ติณฺณนฺนํ ขนฺธานํ จิตฺตสมุฏฺฐานานญฺจ รูปานํ อตฺถิปจฺจเยน ปจฺจโย ฯเปฯ เทฺว ขนฺธา จ ฯเปฯ. ปรามาโส จ สมฺปยุตฺตกา จ ขนฺธา จิตฺต\-สมุฏฺฐานานํ รูปานํ อตฺถิปจฺจเยน ปจฺจโย. ปรามาโส จ มหาภูตา จ จิตฺต\-สมุฏฺฐานานํ รูปานํ อตฺถิปจฺจเยน ปจฺจโย. ปรามาโส จ วตฺถุ จ โนปรามาสานํ ขนฺธานํ อตฺถิปจฺจเยน ปจฺจโย. ปจฺฉาชาโต ปรามาโส จ สมฺปยุตฺตกา จ ขนฺธา ปุเรชาตสฺส อิมสฺส กายสฺส อตฺถิปจฺจเยน ปจฺจโย. ปจฺฉาชาโต ปรามาโส จ สมฺปยุตฺตกา จ ขนฺธา กพฬีกาโร อาหาโร จ อิมสฺส กายสฺส อตฺถิปจฺจเยน ปจฺจโย. ปจฺฉาชาโต ปรามาโส จ สมฺปยุตฺตกา จ ขนฺธา รูปชีวิตินฺทฺริยญฺจ กฏตฺตารูปานํ อตฺถิปจฺจเยน ปจฺจโย. (1)}

\csromanmarkh{1. ปจฺจยานุโลม}\csromanmarkhii{2. สงฺขฺยาวาร}\csromanheadernomark{2}{1. \textbf{ปจฺจยานุโลม 2. สงฺขฺยาวาร}}
\proseitem{38}{เหตุยา ตีณิ, อารมฺมเณ นว, อธิปติยา นว, อนนฺตเร นว, สมนนฺตเร นว, สหชาเต ปญฺจ, อญฺญมญฺเญ ปญฺจ, นิสฺสเย ปญฺจ, อุปนิสฺสเย นว, ปุเรชาเต ตีณิ, ปจฺฉาชาเต ตีณิ, อาเสวเน นว, กมฺเม ตีณิ, วิปาเก เอกํ, อาหาเร ตีณิ, อินฺทฺริเย ตีณิ, ฌาเน ตีณิ, มคฺเค ปญฺจ, สมฺปยุตฺเต ปญฺจ, วิปฺปยุตฺเต ปญฺจ, อตฺถิยา ปญฺจ, นตฺถิยา นว, วิคเต นว, อวิคเต ปญฺจ.}

\vspace{\tipitakaparskip}
\csromancenter{อนุโลมํ.}
\proseitem{39}{ปรามาโส ธมฺโม ปรามาสสฺส ธมฺมสฺส อารมฺมณ\-ปจฺจเยน ปจฺจโย. อุปนิสฺสย\-ปจฺจเยน ปจฺจโย. (1)}

\prose{ปรามาโส ธมฺโม โนปรามาสสฺส ธมฺมสฺส อารมฺมณ\-ปจฺจเยน ปจฺจโย. สหชาต\-ปจฺจเยน ปจฺจโย. อุปนิสฺสย\-ปจฺจเยน ปจฺจโย. ปจฺฉาชาต\-ปจฺจเยน ปจฺจโย. (2)}

\prose{\csromanfolio{350}ปรามาโส ธมฺโม ปรามาสสฺส จ โนปรามาสสฺส จ ธมฺมสฺส อารมฺมณ\-ปจฺจเยน ปจฺจโย. อุปนิสฺสย\-ปจฺจเยน ปจฺจโย. (3)}

\proseitem{40}{โนปรามาโส ธมฺโม โนปรามาสสฺส ธมฺมสฺส อารมฺมณ\-ปจฺจเยน ปจฺจโย. สหชาต\-ปจฺจเยน ปจฺจโย. อุปนิสฺสย\-ปจฺจเยน ปจฺจโย. ปุเรชาต\-ปจฺจเยน ปจฺจโย. ปจฺฉาชาต\-ปจฺจเยน ปจฺจโย. กมฺมปจฺจเยน ปจฺจโย. อาหารปจฺจเยน ปจฺจโย. อินฺทฺริย\-ปจฺจเยน ปจฺจโย. (1)}

\prose{โนปรามาโส ธมฺโม ปรามาสสฺส ธมฺมสฺส อารมฺมณ\-ปจฺจเยน ปจฺจโย. สหชาต\-ปจฺจเยน ปจฺจโย. อุปนิสฺสย\-ปจฺจเยน ปจฺจโย. ปุเรชาต\-ปจฺจเยน ปจฺจโย. (2)}

\prose{โนปรามาโส ธมฺโม ปรามาสสฺส จ โนปรามาสสฺส จ ธมฺมสฺส อารมฺมณ\-ปจฺจเยน ปจฺจโย. สหชาต\-ปจฺจเยน ปจฺจโย. อุปนิสฺสย\-ปจฺจเยน ปจฺจโย. ปุเรชาต\-ปจฺจเยน ปจฺจโย. (3)}

\proseitem{41}{ปรามาโส จ โนปรามาโส จ ธมฺมา ปรามาสสฺส ธมฺมสฺส อารมฺมณ\-ปจฺจเยน ปจฺจโย. อุปนิสฺสย\-ปจฺจเยน ปจฺจโย. (1)}

\prose{ปรามาโส จ โนปรามาโส จ ธมฺมา โนปรามาสสฺส ธมฺมสฺส อารมฺมณ\-ปจฺจเยน ปจฺจโย. สหชาต\-ปจฺจเยน ปจฺจโย. อุปนิสฺสย\-ปจฺจเยน ปจฺจโย. ปจฺฉาชาต\-ปจฺจเยน ปจฺจโย. (2)}

\prose{ปรามาโส จ โนปรามาโส จ ธมฺมา ปรามาสสฺส จ โนปรามาสสฺส จ ธมฺมสฺส อารมฺมณ\-ปจฺจเยน ปจฺจโย. อุปนิสฺสย\-ปจฺจเยน ปจฺจโย. (3)}

\csromanmarkh{2. ปจฺจย\-ปจฺจนีย}\csromanmarkhii{2. สงฺขฺยาวาร}\csromanheadernomark{2}{2. ปจฺจย\-ปจฺจนีย 2. สงฺขฺยาวาร}
\vspace{\tipitakaparskip}
\csromancenter{นเหตุยา นว, นอารมฺมเณ นว, (สพฺพตฺถ นว.) โนอวิคเต นว.}
\csromanmatikaanchor{mtk.49}
\csromanmatikaanchor{mtk.50}
\csromantocmarkref{subsection}{51. ปรามฏฺฐทุก}{mtk.49}
\bookmark[dest={mtk.49},level=2]{51. ปรามฏฺฐทุก}
\csromantocmarkref{subsection}{52. ปรามาสสมฺปยุตฺตทุก}{mtk.50}
\bookmark[dest={mtk.50},level=2]{52. ปรามาสสมฺปยุตฺตทุก}
\csromanmarkh{52. ปรามาส\-สมฺปยุตฺตทุก}\csromanmarkhii{3. ปจฺจยา\-นุโลม\-ปจฺจนีย}\csromanheadernomark{2}{52. ปรามาส\-สมฺปยุตฺตทุก 3. ปจฺจยา\-นุโลม\-ปจฺจนีย}
\vspace{\tipitakaparskip}
\csromancenter{\textbf{เหตุ}ปจฺจยา นอารมฺมเณ ตีณิ, นอธิปติยา ตีณิ, นอนนฺตเร ตีณิ, นสมนนฺตเร ตีณิ, นอญฺญมญฺเญ เอกํ, นฺเอาปนิสฺสเย ตีณิ ฯเปฯ นมคฺเค ตีณิ, นสมฺปยุตฺเต เอกํ, นวิปฺปยุตฺเต ตีณิ, โนนตฺถิยา ตีณิ, โนวิคเต ตีณิ.}
\proseitemwithcloser{44}{\csromanfolio{351}นเหตุปจฺจยา อารมฺมเณ นว, อธิปติยา นว, (อนุโลมมาติกา) ฯเปฯ อวิคเต ปญฺจ.}{ปรามาสทุกํ นิฏฺฐิตํ.}
\csromanheader{2}{51. \textbf{ปรามฏฺฐทุก 1-7. ปฏิจฺจวาร}}
\csromanmarkh{1. ปจฺจยานุโลม}\csromanmarkhii{1. วิภงฺควาร}\csromanheadernomark{2}{1. \textbf{ปจฺจยานุโลม 1. วิภงฺควาร}}
\vspace{\tipitakaparskip}
\csromancenter{\textbf{เหตุ}}
\proseitem{45}{ปรามฏฺฐํ ธมฺมํ ปฏิจฺจ ปรามฏฺโฐ ธมฺโม อุปฺปชฺชติ เหตุปจฺจยา. ปรามฏฺฐํ เอกํ ขนฺธํ ปฏิจฺจ ตโย ขนฺธา จิตฺต\-สมุฏฺฐานญฺจ รูปํ ฯเปฯ เทฺว ขนฺเธ ฯเปฯ. ปฏิสนฺธิ\-กฺขเณ ฯเปฯ. (ยาว อชฺฌตฺติกา มหาภูตา.)}

\prose{ปรามฏฺฐํ ธมฺมํ ปฏิจฺจ ฯเปฯ.}

\vspace{\tipitakaparskip}
\csromancenter{(ปรามฏฺฐทุกํ, ยถา จูฬนฺตรทุเก โลกิยทุกํ, เอวํ กาตพฺพํ, นินฺนานากรณํ.)}
\nitthitam{ปรามฏฺฐทุกํ นิฏฺฐิตํ.}
\csromanmarkh{52. ปรามาส\-สมฺปยุตฺตทุก}\csromanmarkhii{1. ปฏิจฺจวาร}\csromanheadernomark{2}{52. ปรามาส\-สมฺปยุตฺตทุก 1. ปฏิจฺจวาร}
\csromanmarkh{1. ปจฺจยานุโลม}\csromanmarkhii{1. วิภงฺควาร}\csromanheadernomark{2}{1. \textbf{ปจฺจยานุโลม 1. วิภงฺควาร}}
\vspace{\tipitakaparskip}
\csromancenter{\textbf{เหตุ}}
\proseitem{46}{ปรามาส\-สมฺปยุตฺตํ ธมฺมํ ปฏิจฺจ ปรามาส\-สมฺปยุตฺโต ธมฺโม อุปฺปชฺชติ เหตุปจฺจยา. ปรามาส\-สมฺปยุตฺตํ เอกํ ขนฺธํ ปฏิจฺจ ตโย ขนฺธา ฯเปฯ เทฺว ขนฺเธ ฯเปฯ. (1)}

\prose{\csromanfolio{352}ปรามาส\-สมฺปยุตฺตํ ธมฺมํ ปฏิจฺจ ปรามาส\-วิปฺปยุตฺโต ธมฺโม อุปฺปชฺชติ เหตุปจฺจยา. ปรามาส\-สมฺปยุตฺเต ขนฺเธ ปฏิจฺจ จิตฺต\-สมุฏฺฐานํ รูปํ. (2)}

\prose{ปรามาส\-สมฺปยุตฺตํ ธมฺมํ ปฏิจฺจ ปรามาส\-สมฺปยุตฺโต จ ปรามาส\-วิปฺปยุตฺโต จ ธมฺมา อุปฺปชฺชนฺติ เหตุปจฺจยา. ปรามาส\-สมฺปยุตฺตํ เอกํ ขนฺธํ ปฏิจฺจ ตโย ขนฺธา จิตฺต\-สมุฏฺฐานญฺจ รูปํ. เทฺว ขนฺเธ ฯเปฯ. (3)}

\proseitem{47}{ปรามาส\-วิปฺปยุตฺตํ ธมฺมํ ปฏิจฺจ ปรามาส\-วิปฺปยุตฺโต ธมฺโม อุปฺปชฺชติ เหตุปจฺจยา. ปรามาส\-วิปฺปยุตฺตํ เอกํ ขนฺธํ ปฏิจฺจ ตโย ขนฺธา จิตฺต\-สมุฏฺฐานญฺจ รูปํ ฯเปฯ เทฺว ขนฺเธ ฯเปฯ. ปฏิสนฺธิ\-กฺขเณ ฯเปฯ. ขนฺเธ ปฏิจฺจ วตฺถุ, วตฺถุํ ปฏิจฺจ ขนฺธา. เอกํ มหาภูตํ ฯเปฯ. (1)}

\prose{ปรามาส\-สมฺปยุตฺตญฺจ ปรามาส\-วิปฺปยุตฺตญฺจ ธมฺมํ ปฏิจฺจ ปรามาส\-วิปฺปยุตฺโต ธมฺโม อุปฺปชฺชติ เหตุปจฺจยา. ปรามาส\-สมฺปยุตฺเต ขนฺเธ จ มหาภูเต จ ปฏิจฺจ จิตฺต\-สมุฏฺฐานํ รูปํ. (สํขิตฺตํ.) (1)}

\csromanmarkh{1. ปจฺจยานุโลม}\csromanmarkhii{2. สงฺขฺยาวาร}\csromanheadernomark{2}{1. \textbf{ปจฺจยานุโลม 2. สงฺขฺยาวาร}}
\proseitem{48}{เหตุยา ปญฺจ, อารมฺมเณ เทฺว, อธิปติยา ปญฺจ, อนนฺตเร เทฺว, สมนนฺตเร เทฺว, สหชาเต ปญฺจ, อญฺญมญฺเญ เทฺว, นิสฺสเย ปญฺจ, อุปนิสฺสเย เทฺว, ปุเรชาเต เทฺว, อาเสวเน เทฺว, กมฺเม ปญฺจ, วิปาเก เอกํ, อาหาเร ปญฺจ ฯเปฯ มคฺเค ปญฺจ, สมฺปยุตฺเต เทฺว, วิปฺปยุตฺเต ปญฺจ, อตฺถิยา ปญฺจ, นตฺถิยา เทฺว, วิคเต เทฺว, อวิคเต ปญฺจ.}

\vspace{\tipitakaparskip}
\csromancenter{\textbf{อนุโลมํ.}}
\csromanmarkh{2. ปจฺจย\-ปจฺจนีย}\csromanmarkhii{1. วิภงฺควาร}\csromanheadernomark{2}{2. ปจฺจย\-ปจฺจนีย 1. วิภงฺควาร}
\proseitem{49}{ปรามาส\-วิปฺปยุตฺตํ ธมฺมํ ปฏิจฺจ ปรามาส\-วิปฺปยุตฺโต ธมฺโม อุปฺปชฺชติ นเหตุปจฺจยา. อเหตุกํ ปรามาส\-วิปฺปยุตฺตํ เอกํ ขนฺธํ ปฏิจฺจ ตโย \csromanfolio{353}ขนฺธา จิตฺต\-สมุฏฺฐานญฺจ รูปํ ฯเปฯ เทฺว ขนฺเธ ฯเปฯ. อเหตุก\-ปฏิสนฺธิ\-กฺขเณ ฯเปฯ. (ยาว อสญฺญสตฺตา.) วิจิกิจฺฉา\-สหคเต อุทฺธจฺจ\-สหคเต ขนฺเธ ปฏิจฺจ วิจิกิจฺฉา\-สหคโต อุทฺธจฺจ\-สหคโต โมโห. (1)}

\prose{นอารมฺมณ}

\proseitem{49}{ปรามาส\-สมฺปยุตฺตํ ธมฺมํ ปฏิจฺจ ปรามาส\-วิปฺปยุตฺโต ธมฺโม อุปฺปชฺชติ นอารมฺมณ\-ปจฺจยา. ปรามาส\-สมฺปยุตฺเต ขนฺเธ ปฏิจฺจ จิตฺต\-สมุฏฺฐานํ รูปํ. (สํขิตฺตํ.) (1)}

\prose{ปรามาส\-วิปฺปยุตฺตํ ธมฺมํ ปฏิจฺจ ปรามาส\-วิปฺปยุตฺโต ธมฺโม อุปฺปชฺชติ นอารมฺมณ\-ปจฺจยา. ปรามาส\-วิปฺปยุตฺเต ขนฺเธ ปฏิจฺจ จิตฺต\-สมุฏฺฐานํ รูปํ. (ยาว อสญฺญสตฺตา.) (1)}

\prose{ปรามาส\-สมฺปยุตฺตญฺจ ปรามาส\-วิปฺปยุตฺตญฺจ ธมฺมํ ปฏิจฺจ ปรามาส\-วิปฺปยุตฺโต ธมฺโม อุปฺปชฺชติ นอารมฺมณ\-ปจฺจยา. ปรามาส\-สมฺปยุตฺเต ขนฺเธ จ มหาภูเต จ ปฏิจฺจ จิตฺต\-สมุฏฺฐานํ รูปํ. (สํขิตฺตํ.) (1)}

\csromanmarkh{2. ปจฺจย\-ปจฺจนีย}\csromanmarkhii{2. สงฺขฺยาวาร}\csromanheadernomark{2}{2. ปจฺจย\-ปจฺจนีย 2. สงฺขฺยาวาร}
\proseitem{51}{นเหตุยา เอกํ, นอารมฺมเณ ตีณิ, นอธิปติยา ปญฺจ, นอนนฺตเร ตีณิ, นสมนนฺตเร ตีณิ, นอญฺญมญฺเญ ตีณิ, นอุปนิสฺสเย ตีณิ, นปุเรชาเต จตฺตาริ, นปจฺฉาชาเต ปญฺจ, นอาเสวเน ปญฺจ, นกมฺเม เทฺว, นวิปาเก ปญฺจ, นอาหาเร เอกํ, นอินฺทฺริเย เอกํ, นฌาเน เอกํ, นมคฺเค เอกํ, นสมฺปยุตฺเต ตีณิ, นวิปฺปยุตฺเต เทฺว, โนนตฺถิยา ตีณิ, โนวิคเต ตีณิ.}

\proseitem{49}{เหตุปจฺจยา นอารมฺมเณ ตีณิ, นอธิปติยา ปญฺจ. (เอวํ คเณตพฺพํ.)}

\vspace{\tipitakaparskip}
\csromancenter{นเหตุปจฺจยา อารมฺมเณ เอกํ. (สํขิตฺตํ) อวิคเต เอกํ.}
\csromancenter{(สหชาตวาโรปิ ปฏิจฺจ\-วาร\-สทิโส.)}
\csromanmarkh{52. ปรามาส\-สมฺปยุตฺตทุก}\csromanmarkhii{3. ปจฺจยวาร}\csromanheadernomark{2}{52. ปรามาส\-สมฺปยุตฺตทุก 3. ปจฺจยวาร}
\csromanmarkh{1. ปจฺจยานุโลม}\csromanmarkhii{1. วิภงฺควาร}\csromanheadernomark{2}{1. \textbf{ปจฺจยานุโลม 1. วิภงฺควาร}}
\csromanheader{2}{\textbf{เหตุ}}
\proseitem{54}{\csromanfolio{354}ปรามาส\-สมฺปยุตฺตํ ธมฺมํ ปจฺจยา ปรามาส\-สมฺปยุตฺโต ธมฺโม อุปฺปชฺชติ เหตุปจฺจยา. ตีณิ. (ปฏิจฺจ\-วาร\-สทิโส.)}

\prose{ปรามาส\-วิปฺปยุตฺตํ ธมฺมํ ปจฺจยา ปรามาส\-วิปฺปยุตฺโต ธมฺโม อุปฺปชฺชติ เหตุปจฺจยา. ปรามาส\-วิปฺปยุตฺตํ เอกํ ขนฺธํ ปจฺจยา ตโย ขนฺธา จิตฺต\-สมุฏฺฐานญฺจ รูปํ ฯเปฯ เทฺว ขนฺเธ ฯเปฯ. ปฏิสนฺธิ\-กฺขเณ ฯเปฯ. (ยาว อชฺฌตฺติกา มหาภูตา.) วตฺถุํ ปจฺจยา ปรามาส\-วิปฺปยุตฺตา ขนฺธา. (1)}

\prose{ปรามาส\-วิปฺปยุตฺตํ ธมฺมํ ปจฺจยา ปรามาส\-สมฺปยุตฺโต ธมฺโม อุปฺปชฺชติ เหตุปจฺจยา. วตฺถุํ ปจฺจยา ปรามาส\-สมฺปยุตฺตกา ขนฺธา. (2)}

\prose{ปรามาส\-วิปฺปยุตฺตํ ธมฺมํ ปจฺจยา ปรามาส\-สมฺปยุตฺโต จ ปรามาส\-วิปฺปยุตฺโต จ ธมฺมา อุปฺปชฺชนฺติ เหตุปจฺจยา. วตฺถุํ ปจฺจยา ปรามาส\-สมฺปยุตฺตกา ขนฺธา. มหาภูเต ปจฺจยา จิตฺต\-สมุฏฺฐานํ รูปํ. (3)}

\proseitem{55}{ปรามาส\-สมฺปยุตฺตญฺจ ปรามาส\-วิปฺปยุตฺตญฺจ ธมฺมํ ปจฺจยา ปรามาส\-สมฺปยุตฺโต ธมฺโม อุปฺปชฺชติ เหตุปจฺจยา. ปรามาส\-สมฺปยุตฺตํ เอกํ ขนฺธญฺจ วตฺถุญฺจ ปจฺจยา ตโย ขนฺธา ฯเปฯ เทฺว ขนฺเธ ฯเปฯ. (1)}

\prose{ปรามาส\-สมฺปยุตฺตญฺจ ปรามาส\-วิปฺปยุตฺตญฺจ ธมฺมํ ปจฺจยา ปรามาส\-วิปฺปยุตฺโต ธมฺโม อุปฺปชฺชติ เหตุปจฺจยา. ปรามาส\-สมฺปยุตฺเต ขนฺเธ จ มหาภูเต จ ปจฺจยา จิตฺต\-สมุฏฺฐานํ รูปํ. (2)}

\prose{ปรามาส\-สมฺปยุตฺตญฺจ ปรามาส\-วิปฺปยุตฺตญฺจ ธมฺมํ ปจฺจยา ปรามาส\-สมฺปยุตฺโต จ ปรามาส\-วิปฺปยุตฺโต จ ธมฺมา อุปฺปชฺชนฺติ เหตุปจฺจยา. ปรามาส\-สมฺปยุตฺตํ เอกํ ขนฺธญฺจ วตฺถุญฺจ ปจฺจยา ตโย ขนฺธา ฯเปฯ เทฺว ขนฺเธ ฯเปฯ. ปรามาส\-สมฺปยุตฺเต ขนฺเธ จ มหาภูเต จ ปจฺจยา จิตฺต\-สมุฏฺฐานํ รูปํ. (สํขิตฺตํ.) (3)}

\csromanmarkh{52. ปรามาส\-สมฺปยุตฺตทุก}\csromanmarkhii{1. ปจฺจยานุโลม 2. สงฺขฺยาวาร}\csromanheadernomark{2}{52. ปรามาส\-สมฺปยุตฺตทุก 1. ปจฺจยานุโลม 2. สงฺขฺยาวาร}
\proseitem{56}{\csromanfolio{355}เหตุยา นว, อารมฺมเณ จตฺตาริ, อธิปติยา นว, อนนฺตเร จตฺตาริ, สมนนฺตเร จตฺตาริ, สหชาเต นว, อญฺญมญฺเญ จตฺตริ, นิสฺสเย นว, อุปนิสฺสเย จตฺตาริ, ปุเรชาเต จตฺตาริ, อาเสวเน จตฺตาริ, กมฺเม นว, วิปาเก เอกํ, อาหาเร นว ฯเปฯ อวิคเต นว.}

\vspace{\tipitakaparskip}
\csromancenter{อนุโลมํ.}
\csromanmarkh{2. ปจฺจย\-ปจฺจนีย}\csromanmarkhii{1. วิภงฺควาร}\csromanheadernomark{2}{2. ปจฺจย\-ปจฺจนีย 1. วิภงฺควาร}
\proseitem{57}{ปรามาส\-วิปฺปยุตฺตํ ธมฺมํ ปจฺจยา ปรามาส\-วิปฺปยุตฺโต ธมฺโม อุปฺปชฺชติ นเหตุปจฺจยา. อเหตุกํ ปรามาส\-วิปฺปยุตฺตํ เอกํ ขนฺธํ ปจฺจยา ตโย ขนฺธา จิตฺต\-สมุฏฺฐานญฺจ รูปํ ฯเปฯ เทฺว ขนฺเธ ฯเปฯ. อเหตุก\-ปฏิสนฺธิ\-กฺขเณ ฯเปฯ. (ยาว อสญฺญสตฺตา.) จกฺขายตนํ ปจฺจยา จกฺขุ วิญฺญาณํ ฯเปฯ. กายายตนํ ปจฺจยา กายวิญฺญาณํ. วตฺถุํ ปจฺจยา อเหตุก\-ปรามาส\-วิปฺปยุตฺตา ขนฺธา. วิจิกิจฺฉา\-สหคเต อุทฺธจฺจ\-สหคเต ขนฺเธ จ วตฺถุญฺจ ปจฺจยา วิจิกิจฺฉา\-สหคโต อุทฺธจฺจ\-สหคโต โมโห.}

\csromanmarkh{2. ปจฺจย\-ปจฺจนีย}\csromanmarkhii{2. สงฺขฺยาวาร}\csromanheadernomark{2}{2. ปจฺจย\-ปจฺจนีย 2. สงฺขฺยาวาร}
\proseitem{58}{นเหตุยา เอกํ, นอารมฺมเณ ตีณิ, นอธิปติยา นว, นอนนฺตเร ตีณิ, นสมนนฺตเร ตีณิ, นอญฺญมญฺเญ ตีณิ, นอุปนิสฺสเย ตีณิ, นปุเรชาเต จตฺตาริ, นปจฺฉาชาเต นว, นอาเสวเน นว, นกมฺเม จตฺตาริ, นวิปาเก นว, นอาหาเร เอกํ, นอินฺทฺริเย เอกํ, นฌาเน เอกํ, นมคฺเค เอกํ, นสมฺปยุตฺเต ตีณิ, นวิปฺปยุตฺเต เทฺว, โนนตฺถิยา ตีณิ โนวิคเต ตีณิ.}

\vspace{\tipitakaparskip}
\csromancenter{\textbf{เหตุ}ปจฺจยา นอารมฺมเณ ตีณิ, นอธิปติยา นว. (เอวํ คเณตพฺพํ.)}
\csromancenter{นเหตุปจฺจยา อารมฺมเณ เอกํ ฯเปฯ อวิคเต เอกํ.}
\csromancenter{(นิสฺสยวาโร ปจฺจยวาร\-สทิโส.)}
\csromanmarkh{52. ปรามาส\-สมฺปยุตฺตทุก}\csromanmarkhii{5. สํสฏฺฐวาร}\csromanheadernomark{2}{52. ปรามาส\-สมฺปยุตฺตทุก 5. สํสฏฺฐวาร}
\proseitem{61}{\csromanfolio{356}ปรามาส\-สมฺปยุตฺตํ ธมฺมํ สํสฏฺโฐ ปรามาส\-สมฺปยุตฺโต ธมฺโม อุปฺปชฺชติ เหตุปจฺจยา. (สํขิตฺตํ.)}

\proseitem{62}{เหตุยา เทฺว, (สพฺพตฺถ เทฺว,) วิปาเก เอกํ ฯเปฯ อวิคเต เทฺว.}

\prose{นเหตุยา เอกํ, นอธิปติยา เทฺว, นปุเรชาเต เทฺว, นปจฺฉาชาเต เทฺว, นอาเสวเน เทฺว, นกมฺเม เทฺว, นวิปาเก เทฺว, นฌาเน เอกํ, นมคฺเค เอกํ, นวิปฺปยุตฺเต เทฺว.}

\vspace{\tipitakaparskip}
\csromancenter{(เอวํ อิตเร เทฺว คณนาปิ สมฺปยุตฺต\-วาราปิ กาตพฺพา.)}
\csromanmarkh{52. ปรามาส\-สมฺปยุตฺตทุก}\csromanmarkhii{7. ปญฺหาวาร}\csromanheadernomark{2}{52. ปรามาส\-สมฺปยุตฺตทุก 7. ปญฺหาวาร}
\csromanmarkh{1. ปจฺจยานุโลม}\csromanmarkhii{1. วิภงฺควาร}\csromanheadernomark{2}{1. \textbf{ปจฺจยานุโลม 1. วิภงฺควาร}}
\vspace{\tipitakaparskip}
\csromancenter{\textbf{เหตุ}}
\proseitem{63}{ปรามาส\-สมฺปยุตฺโต ธมฺโม ปรามาส\-สมฺปยุตฺตสฺส ธมฺมสฺส เหตุปจฺจเยน ปจฺจโย. ปรามาส\-สมฺปยุตฺตา เหตู สมฺปยุตฺตกานํ \csromanfolio{357}ขนฺธานํ เหตุปจฺจเยน ปจฺจโย. (มูลํ กาตพฺพํ.) ปรามาส\-สมฺปยุตฺตา เหตู จิตฺต\-สมุฏฺฐานานํ รูปานํ เหตุปจฺจเยน ปจฺจโย. (มูลํ กาตพฺพํ.) ปรามาส\-สมฺปยุตฺตา เหตู สมฺปยุตฺตกานํ ขนฺธานํ จิตฺตสมุฏฺฐานานญฺจ รูปานํ เหตุปจฺจเยน ปจฺจโย. (3)}

\prose{ปรามาส\-วิปฺปยุตฺโต ธมฺโม ปรามาส\-วิปฺปยุตฺตสฺส ธมฺมสฺส เหตุปจฺจเยน ปจฺจโย. ปรามาส\-วิปฺปยุตฺโต เหตู สมฺปยุตฺตกานํ ขนฺธานํ จิตฺตสมุฏฺฐานานญฺจ รูปานํ เหตุปจฺจเยน ปจฺจโย. ปฏิสนฺธิ\-กฺขเณ ฯเปฯ. (1)}

\proseitem{64}{ปรามาส\-สมฺปยุตฺโต ธมฺโม ปรามาส\-สมฺปยุตฺตสฺส ธมฺมสฺส อารมฺมณ\-ปจฺจเยน ปจฺจโย. ราคํ อสฺสาเทติ อภินนฺทติ, ตํ อารพฺภ ปรามาส\-สมฺปยุตฺโต ราโค อุปฺปชฺชติ. ปรามาส\-สมฺปยุตฺเต ขนฺเธ อสฺสาเทติ อภินนฺทติ, ตํ อารพฺภ ปรามาส\-สมฺปยุตฺโต ราโค อุปฺปชฺชติ. (1)}

\prose{ปรามาส\-สมฺปยุตฺโต ธมฺโม ปรามาส\-วิปฺปยุตฺตสฺส ธมฺมสฺส อารมฺมณ\-ปจฺจเยน ปจฺจโย. อริยา ปรามาส\-สมฺปยุตฺเต ปหีเน กิเลเส ปจฺจเวกฺขนฺติ, ปุพฺเพ สมุทาจิณฺเณ กิเลเส ชานนฺติ. ปรามาส\-สมฺปยุตฺเต ขนฺเธ อนิจฺจโต ฯเปฯ วิปสฺสติ, อสฺสาเทติ อภินนฺทติ, ตํ อารพฺภ ปรามาส\-วิปฺปยุตฺโต ราโค ฯเปฯ. วิจิกิจฺฉา ฯเปฯ อุทฺธจฺจํ ฯเปฯ โทมนสฺสํ อุปฺปชฺชติ. เจโตปริย\-ญาเณน ปรามาส\-สมฺปยุตฺตจิตฺต\-สมงฺคิสฺส จิตฺตํ ชานาติ. ปรามาส\-สมฺปยุตฺตา ขนฺธา เจโตปริย\-ญาณสฺส, ปุพฺเพ\-นิวาสา\-นุสฺสติ\-ญาณสฺส, ยถากมฺมูปคญาณสฺส, อนาคตํสญาณสฺส, อาวชฺชนาย อารมฺมณ\-ปจฺจเยน ปจฺจโย. (2)}

\proseitem{65}{ปรามาส\-วิปฺปยุตฺโต ธมฺโม ปรามาส\-วิปฺปยุตฺตสฺส ธมฺมสฺส อารมฺมณ\-ปจฺจเยน ปจฺจโย. ทานํ ฯเปฯ สีลํ ฯเปฯ อุโปสถกมฺมํ กตฺวา ตํ ปจฺจเวกฺขติ, อสฺสาเทติ อภินนฺทติ, ตํ อารพฺภ ปรามาส\-วิปฺปยุตฺโต ราโค ฯเปฯ วิจิกิจฺฉา ฯเปฯ อุทฺธจฺจํ ฯเปฯ โทมนสฺสํ อุปฺปชฺชติ. ปุพฺเพ สุจิณฺณานิ ฯเปฯ. ฌานา ฯเปฯ. อริยา มคฺคา วุฏฺฐหิตฺวา มคฺคํ ปจฺจเวกฺขนฺติ. ผลํ ฯเปฯ. นิพฺพานํ ฯเปฯ. นิพฺพานํ โคตฺรภุสฺส, โวทานสฺส, มคฺคสฺส, ผลสฺส, อาวชฺชนาย อารมฺมณ\-ปจฺจเยน ปจฺจโย. อริยา ปรามาส\-วิปฺปยุตฺเต ปหีเน กิเลเส ฯเปฯ. วิกฺขมฺภิเต กิเลเส ฯเปฯ. ปุพฺเพ ฯเปฯ. จกฺขุํ ฯเปฯ วตฺถุํ ปรามาส\-วิปฺปยุตฺเต \csromanfolio{358}ขนฺเธ อนิจฺจโต ฯเปฯ วิปสฺสติ, อสฺสาเทติ อภินนฺทติ, ตํ อารพฺภ ปรามาส\-วิปฺปยุตฺโต ราโค ฯเปฯ วิจิกิจฺฉา ฯเปฯ อุทฺธจฺจํ ฯเปฯ โทมนสฺสํ อุปฺปชฺชติ. ทิพฺเพน จกฺขุนา รูปํ ปสฺสติ. ทิพฺพาย โสตธาตุยา สทฺทํ สุณาติ. เจโตปริย\-ญาเณน ปรามาส\-วิปฺปยุตฺต\-จิตฺต\-สมงฺคิสฺส จิตฺตํ ชานาติ. อากาสา\-นญฺจา\-ยตนํ วิญฺญาณญฺจา\-ยตนสฺส ฯเปฯ. อากิญฺจญฺญา\-ยตนํ เนว\-สญฺญา\-นา\-สญฺญา\-ยตนสฺส ฯเปฯ. รูปายตนํ จกฺขุ\-วิญฺญาณสฺส ฯเปฯ. โผฏฺฐพฺพา\-ยตนํ กายวิญฺญาณสฺส ฯเปฯ. ปรามาส\-วิปฺปยุตฺตา ขนฺธา อิทฺธิวิธ\-ญาณสฺส, เจโตปริย\-ญาณสฺส, ปุพฺเพ\-นิวาสา\-นุสฺสติ\-ญาณสฺส, ยถากมฺมูปคญาณสฺส, อนาคตํสญาณสฺส, อาวชฺชนาย อารมฺมณ\-ปจฺจเยน ปจฺจโย. (1)}

\prose{ปรามาส\-วิปฺปยุตฺโต ธมฺโม ปรามาส\-สมฺปยุตฺตสฺส ธมฺมสฺส อารมฺมณ\-ปจฺจเยน ปจฺจโย. ทานํ ฯเปฯ สีลํ ฯเปฯ อุโปสถกมฺมํ ฯเปฯ. ปุพฺเพ ฯเปฯ. ฌานา ฯเปฯ. จกฺขุํ ฯเปฯ. วตฺถุํ ปรามาส\-วิปฺปยุตฺเต ขนฺเธ อสฺสาเทติ อภินนฺทติ ตํ อารพฺภ ปรามาส\-สมฺปยุตฺโต ราโค อุปฺปชฺชติ. (2)}

\prose{ปรามาส\-สมฺปยุตฺโต ธมฺโม ปรามาส\-สมฺปยุตฺตสฺส ธมฺมสฺส อธิปติ\-ปจฺจเยน ปจฺจโย. อารมฺมณา\-ธิปติ, สหชาตา\-ธิปติ.  อารมฺมณา\-ธิปติ\csromandash{}ราคํ ครุํ กตฺวา อสฺสาเทติ อภินนฺทติ, ตํ ครุํ กตฺวา ปรามาส\-สมฺปยุตฺโต ราโค อุปฺปชฺชติ. ปรามาส\-สมฺปยุตฺเต ขนฺเธ ครุํ กตฺวา อสฺสาเทติ อภินนฺทติ, ตํ ครุํ กตฺวา ปรามาส\-สมฺปยุตฺโต ราโค อุปฺปชฺชติ. สหชาตา\-ธิปติ\csromandash{}ปรามาส\-สมฺปยุตฺตา\-ธิปติ สมฺปยุตฺตกานํ ขนฺธานํ อธิปติ\-ปจฺจเยน ปจฺจโย. (1)}

\prose{ปรามาส\-สมฺปยุตฺโต ธมฺโม ปรามาส\-วิปฺปยุตฺตสฺส ธมฺมสฺส อธิปติ\-ปจฺจเยน ปจฺจโย. อารมฺมณา\-ธิปติ, สหชาตา\-ธิปติ.  อารมฺมณา\-ธิปติ\csromandash{}ปรามาส\-สมฺปยุตฺเต ขนฺเธ ครุํ กตฺวา อสฺสาเทติ อภินนฺทติ, ตํ ครุํ กตฺวา ปรามาส\-วิปฺปยุตฺโต ราโค อุปฺปชฺชติ. สหชาตา\-ธิปติ\csromandash{}ปรามาส\-สมฺปยุตฺตา\-ธิปติ จิตฺต\-สมุฏฺฐานานํ รูปานํ อธิปติ\-ปจฺจเยน ปจฺจโย. (2)}

\prose{ปรามาส\-สมฺปยุตฺโต ธมฺโม ปรามาส\-สมฺปยุตฺตสฺส จ ปรามาส\-วิปฺปยุตฺตสฺส จ ธมฺมสฺส อธิปติ\-ปจฺจเยน ปจฺจโย. สหชาตา\-ธิปติ\csromandash{} \csromanfolio{359}ปรามาส\-สมฺปยุตฺตา\-ธิปติ สมฺปยุตฺตกานํ ขนฺธานํ จิตฺตสมุฏฺฐานานญฺจ รูปานํ อธิปติ\-ปจฺจเยน ปจฺจโย. (3)}

\proseitem{66}{ปรามาส\-วิปฺปยุตฺโต ธมฺโม ปรามาส\-วิปฺปยุตฺตสฺส ธมฺมสฺส อธิปติ\-ปจฺจเยน ปจฺจโย. อารมฺมณา\-ธิปติ, สหชาตา\-ธิปติ.  อารมฺมณา\-ธิปติ\csromandash{}ทานํ ฯเปฯ สีลํ ฯเปฯ อุโปสถกมฺมํ กตฺวา ตํ ครุํ กตฺวา ปจฺจเวกฺขติ, อสฺสาเทติ อภินนฺทติ, ตํ ครุํ กตฺวา ปรามาส\-วิปฺปยุตฺโต ราโค อุปฺปชฺชติ. ปุพฺเพ ฯเปฯ. ฌานา ฯเปฯ. อริยา มคฺคา วุฏฺฐหิตฺวา มคฺคํ ครุํ กตฺวา ฯเปฯ. ผลํ ฯเปฯ นิพฺพานํ ฯเปฯ นิพฺพานํ โคตฺรภุสฺส, โวทานสฺส, มคฺคสฺส, ผลสฺส อธิปติ\-ปจฺจเยน ปจฺจโย. จกฺขุํ ฯเปฯ วตฺถุํ ปรามาส\-วิปฺปยุตฺเต ขนฺเธ ครุํ กตฺวา อสฺสาเทติ อภินนฺทติ, ตํ ครุํ กตฺวา ปรามาส\-วิปฺปยุตฺโต ราโค อุปฺปชฺชติ. สหชาตา\-ธิปติ\csromandash{}ปรามาส\-วิปฺปยุตฺตา\-ธิปติ สมฺปยุตฺตกานํ ขนฺธานํ จิตฺตสมุฏฺฐานานญฺจ รูปานํ อธิปติ\-ปจฺจเยน ปจฺจโย. (1)}

\prose{ปรามาส\-วิปฺปยุตฺโต ธมฺโม ปรามาส\-สมฺปยุตฺตสฺส ธมฺมสฺส อธิปติ\-ปจฺจเยน ปจฺจโย. อารมฺมณา\-ธิปติ\csromandash{}ทานํ ฯเปฯ สีลํ ฯเปฯ อุโปสถกมฺมํ กตฺวา ตํ ครุํ กตฺวา อสฺสาเทติ อภินนฺทติ, ตํ ครุํ กตฺวา ปรามาส\-สมฺปยุตฺโต ราโค อุปฺปชฺชติ. ปุพฺเพ ฯเปฯ. ฌานา ฯเปฯ. จกฺขุํ ฯเปฯ วตฺถุํ ปรามาส\-วิปฺปยุตฺเต ขนฺเธ ครุํ กตฺวา อสฺสาเทติ อภินนฺทติ, ตํ ครุํ กตฺวา ปรามาส\-สมฺปยุตฺโต ราโค อุปฺปชฺชติ. (2)}

\csromanheader{2}{อนนฺตร\-สมนนฺตร}
\proseitem{66}{ปรามาส\-สมฺปยุตฺโต ธมฺโม ปรามาส\-สมฺปยุตฺตสฺส ธมฺมสฺส อนนฺตร\-ปจฺจเยน ปจฺจโย. ปุริมา ปุริมา ปรามาส\-สมฺปยุตฺตา ขนฺธา ปจฺฉิมานํ ปจฺฉิมานํ ปรามาส\-สมฺปยุตฺตกานํ ขนฺธานํ อนนฺตร\-ปจฺจเยน ปจฺจโย. (1)}

\prose{ปรามาส\-สมฺปยุตฺโต ธมฺโม ปรามาส\-วิปฺปยุตฺตสฺส ธมฺมสฺส อนนฺตร\-ปจฺจเยน ปจฺจโย. ปรามาส\-สมฺปยุตฺตา ขนฺธา วุฏฺฐานสฺส อนนฺตร\-ปจฺจเยน ปจฺจโย. (2)}

\prose{ปรามาส\-วิปฺปยุตฺโต ธมฺโม ปรามาส\-วิปฺปยุตฺตสฺส ธมฺมสฺส อนนฺตร\-ปจฺจเยน ปจฺจโย. ปุริมา ปุริมา ปรามาส\-วิปฺปยุตฺตา ขนฺธา \csromanfolio{360}ปจฺฉิมานํ ปจฺฉิมานํ ปรามาส\-วิปฺปยุตฺตานํ ขนฺธานํ อนนฺตร\-ปจฺจเยน ปจฺจโย. อนุโลมํ ฯเปฯ ผลสมาปตฺติยา อนนฺตร\-ปจฺจเยน ปจฺจโย. (1)}

\proseitem{69}{ปรามาส\-วิปฺปยุตฺโต ธมฺโม ปรามาส\-สมฺปยุตฺตสฺส ธมฺมสฺส อนนฺตร\-ปจฺจเยน ปจฺจโย. อาวชฺชนา ปรามาส\-สมฺปยุตฺตกานํ ขนฺธานํ อนนฺตร\-ปจฺจเยน ปจฺจโย. สมนนฺตร\-ปจฺจเยน ปจฺจโย.}

\csromanheader{2}{\textbf{สหชาตาทิ}}
\proseitem{70}{ปรามาส\-สมฺปยุตฺโต ธมฺโม ปรามาส\-สมฺปยุตฺตสฺส ธมฺมสฺส สหชาต\-ปจฺจเยน ปจฺจโย. ปญฺจ. อญฺญมญฺญ\-ปจฺจเยน ปจฺจโย. เทฺว. นิสฺสย\-ปจฺจเยน ปจฺจโย. สตฺต.}

\csromanheader{2}{\textbf{อุปนิสฺสย}}
\proseitem{69}{ปรามาส\-สมฺปยุตฺโต ธมฺโม ปรามาส\-สมฺปยุตฺตสฺส ธมฺมสฺส อุปนิสฺสย\-ปจฺจเยน ปจฺจโย. อารมฺมณู\-ปนิสฺสโย, อนนฺตรู\-ปนิสฺสโย, ปกตู\-ปนิสฺสโย ฯเปฯ. ปกตู\-ปนิสฺสโย\csromandash{}ปรามาส\-สมฺปยุตฺโต ราโค. โมโห. ปตฺถนา ปรามาส\-สมฺปยุตฺตสฺส ราคสฺส. โมหสฺส. ปตฺถนาย อุปนิสฺสย\-ปจฺจเยน ปจฺจโย. (1)}

\prose{ปรามาส\-สมฺปยุตฺโต ธมฺโม ปรามาส\-วิปฺปยุตฺตสฺส ธมฺมสฺส อุปนิสฺสย\-ปจฺจเยน ปจฺจโย. อารมฺมณู\-ปนิสฺสโย, อนนฺตรู\-ปนิสฺสโย, ปกตู\-ปนิสฺสโย ฯเปฯ. ปกตู\-ปนิสฺสโย\csromandash{}ปรามาส\-สมฺปยุตฺตํ ราคํ อุปนิสฺสาย ทานํ เทติ ฯเปฯ สมาปตฺติํ อุปฺปาเทติ. มานํ ชปฺเปติ. ปรามาส\-สมฺปยุตฺตํ โมหํ. ปตฺถนํ อุปนิสฺสาย ทานํ เทติ ฯเปฯ สมาปตฺติํ อุปฺปาเทติ. มานํ ชปฺเปติ. ปาณํ หนติ ฯเปฯ สํฆํ ภินฺทติ. ปรามาส\-สมฺปยุตฺโต ราโค โมโห. ปตฺถนา สทฺธาย ฯเปฯ. ปญฺญาย. ราคสฺส. โทสสฺส. โมหสฺส. มานสฺส. ปตฺถนาย. กายิกสฺส สุขสฺส ฯเปฯ ผลสมาปตฺติยา อุปนิสฺสย\-ปจฺจเยน ปจฺจโย. (2)}

\proseitem{72}{ปรามาส\-วิปฺปยุตฺโต ธมฺโม ปรามาส\-วิปฺปยุตฺตสฺส ธมฺมสฺส อุปนิสฺสย\-ปจฺจเยน ปจฺจโย. อารมฺมณู\-ปนิสฺสโย, อนนฺตรู\-ปนิสฺสโย, ปกตู\-ปนิสฺสโย ฯเปฯ. ปกตู\-ปนิสฺสโย\csromandash{}สทฺธํ อุปนิสฺสาย ทานํ \csromanfolio{361}เทติ ฯเปฯ สมาปตฺติํ อุปฺปาเทติ. มานํ ชปฺเปติ. สีลํ ฯเปฯ. ปญฺญํ. ราคํ ฯเปฯ. มานํ ฯเปฯ. ปตฺถนํ อุปนิสฺสาย ทานํ เทติ ฯเปฯ สมาปตฺติํ อุปฺปาเทติ. ปาณํ หนติ ฯเปฯ สํฆํ ภินฺทติ. กายิกํ สุขํ ฯเปฯ. เสนาสนํ อุปนิสฺสาย ทานํ เทติ ฯเปฯ สํฆํ ภินฺทติ. สทฺธา ฯเปฯ. ปญฺญา. ราโค ฯเปฯ. มาโน. ปตฺถนา. กายิกํ สุขํ ฯเปฯ. เสนาสนํ สทฺธาย ฯเปฯ ปญฺญาย. ราคสฺส ฯเปฯ มานสฺส. ปตฺถนาย. กายิกสฺส สุขสฺส ฯเปฯ ผลสมปตฺติยา อุปนิสฺสย\-ปจฺจเยน ปจฺจโย. (1)}

\prose{ปรามาส\-วิปฺปยุตฺโต ธมฺโม ปรามาส\-สมฺปยุตฺตสฺส ธมฺมสฺส อุปนิสฺสย\-ปจฺจเยน ปจฺจโย. อารมฺมณู\-ปนิสฺสโย, อนนฺตรู\-ปนิสฺสโย, ปกตู\-ปนิสฺสโย ฯเปฯ. ปกตู\-ปนิสฺสโย\csromandash{}สทฺธํ อุปนิสฺสาย ปรามาส\-สมฺปยุตฺโต ราโค อุปฺปชฺชติ. สีลํ ฯเปฯ. เสนาสนํ อุปนิสฺสาย ปรามาส\-สมฺปยุตฺโต ราโค อุปฺปชฺชติ. สทฺธา ฯเปฯ. เสนาสนํ ปรามาส\-สมฺปยุตฺตสฺส ราคสฺส ฯเปฯ ปตฺถนาย อุปนิสฺสย\-ปจฺจเยน ปจฺจโย. (2)}

\proseitem{73}{ปรามาส\-วิปฺปยุตฺโต ธมฺโม ปรามาส\-วิปฺปยุตฺตสฺส ธมฺมสฺส ปุเรชาต\-ปจฺจเยน ปจฺจโย. อารมฺมณ\-ปุเรชาตํ, วตฺถุ\-ปุเรชาตํ.  อารมฺมณ\-ปุเรชาตํ\csromandash{}จกฺขุํ ฯเปฯ วตฺถุํ อนิจฺจโต ฯเปฯ วิปสฺสติ, อสฺสาเทติ อภินนฺทติ, ตํ อารพฺภ ปรามาส\-วิปฺปยุตฺโต ราโค ฯเปฯ วิจิกิจฺฉา ฯเปฯ อุทฺธจฺจํ ฯเปฯ โทมนสฺสํ อุปฺปชฺชติ. ทิพฺเพน จกฺขุนา รูปํ ปสฺสติ. ทิพฺพาย โสตธาตุยา สทฺทํ สุณาติ. รูปายตนํ จกฺขุ\-วิญฺญาณสฺส ฯเปฯ. โผฏฺฐพฺพา\-ยตนํ กายวิญฺญาณสฺส ปุเรชาต\-ปจฺจเยน ปจฺจโย. วตฺถุ\-ปุเรชาตํ\csromandash{}จกฺขายตนํ จกฺขุ\-วิญฺญาณสฺส ฯเปฯ กายายตนํ กายวิญฺญาณสฺส ฯเปฯ. วตฺถุ ปรามาส\-วิปฺปยุตฺตานํ ขนฺธานํ ปุเรชาต\-ปจฺจเยน ปจฺจโย. (1)}

\prose{ปรามาส\-วิปฺปยุตฺโต ธมฺโม ปรามาส\-สมฺปยุตฺตสฺส ธมฺมสฺส ปุเรชาต\-ปจฺจเยน ปจฺจโย. อารมฺมณ\-ปุเรชาตํ, วตฺถุ\-ปุเรชาตํ.  อารมฺมณ\-ปุเรชาตํ\csromandash{}จกฺขุํ ฯเปฯ วตฺถุํ อสฺสาเทติ อภินนฺทติ, ตํ อารพฺภ ปรามาส\-สมฺปยุตฺโต ราโค อุปฺปชฺชติ. วตฺถุ\-ปุเรชาตํ\csromandash{}วตฺถุ ปรามาส\-สมฺปยุตฺตกานํ ขนฺธานํ ปุเรชาต\-ปจฺจเยน ปจฺจโย. (2)}

\proseitem{74}{\csromanfolio{362}ปรามาส\-สมฺปยุตฺโต ธมฺโม ปรามาส\-วิปฺปยุตฺตสฺส ธมฺมสฺส ปจฺฉาชาต\-ปจฺจเยน ปจฺจโย. (สํขิตฺตํ.) (1)}

\prose{ปรามาส\-วิปฺปยุตฺโต ธมฺโม ปรามาส\-วิปฺปยุตฺตสฺส ธมฺมสฺส ปจฺฉาชาต\-ปจฺจเยน ปจฺจโย. (สํขิตฺตํ.) (1)}

\prose{ปรามาส\-สมฺปยุตฺโต ธมฺโม ปรามาส\-สมฺปยุตฺตสฺส ธมฺมสฺส อาเสวน\-ปจฺจเยน ปจฺจโย. เทฺว.}

\csromanheader{2}{\textbf{กมฺมาทิ}}
\proseitem{75}{ปรามาส\-สมฺปยุตฺโต ธมฺโม ปรามาส\-สมฺปยุตฺตสฺส ธมฺมสฺส กมฺมปจฺจเยน ปจฺจโย. ปรามาส\-สมฺปยุตฺตา เจตนา สมฺปยุตฺตกานํ ขนฺธานํ กมฺมปจฺจเยน ปจฺจโย. (1)}

\prose{ปรามาส\-สมฺปยุตฺโต ธมฺโม ปรามาส\-วิปฺปยุตฺตสฺส ธมฺมสฺส กมฺมปจฺจเยน ปจฺจโย. สหชาตา, นานากฺขณิกา.  สหชาตา ปรามาส\-สมฺปยุตฺตา เจตนา จิตฺต\-สมุฏฺฐานานํ รูปานํ กมฺมปจฺจเยน ปจฺจโย. นานากฺขณิกา ปรามาส\-สมฺปยุตฺตา เจตนา วิปากานํ ขนฺธานํ กฏตฺตา จ รูปานํ กมฺมปจฺจเยน ปจฺจโย. (มูลํ กาตพฺพํ.) ปรามาส\-สมฺปยุตฺตา เจตนา สมฺปยุตฺตกานํ ขนฺธานํ จิตฺตสมุฏฺฐานานญฺจ รูปานํ กมฺมปจฺจเยน ปจฺจโย. (3)}

\prose{ปรามาส\-วิปฺปยุตฺโต ธมฺโม ปรามาส\-วิปฺปยุตฺตสฺส ธมฺมสฺส กมฺมปจฺจเยน ปจฺจโย. สหชาตา, นานากฺขณิกา.  สหชาตา ปรามาส\-วิปฺปยุตฺตา เจตนา สมฺปยุตฺตกานํ ขนฺธานํ จิตฺตสมุฏฺฐานานญฺจ รูปานํ กมฺมปจฺจเยน ปจฺจโย. ปฏิสนฺธิ\-กฺขเณ ฯเปฯ. นานากฺขณิกา ปรามาส\-วิปฺปยุตฺตา เจตนา วิปากานํ ขนฺธานํ กฏตฺตา จ รูปานํ กมฺมปจฺจเยน ปจฺจโย. (1)}

\prose{วิปาก\-ปจฺจเยน ปจฺจโย. เอกํ. อาหารปจฺจเยน ปจฺจโย. จตฺตาริ. อินฺทฺริย\-ปจฺจเยน ปจฺจโย. จตฺตาริ. ฌานปจฺจเยน ปจฺจโย. จตฺตาริ. มคฺคปจฺจเยน ปจฺจโย. จตฺตาริ. สมฺปยุตฺต\-ปจฺจเยน ปจฺจโย. เทฺว.}

\csromanheader{2}{52. ปรามาส\-สมฺปยุตฺตทุก วิปฺปยุตฺต}
\proseitem{76}{\csromanfolio{363}ปรามาส\-สมฺปยุตฺโต ธมฺโม ปรามาส\-วิปฺปยุตฺตสฺส ธมฺมสฺส วิปฺปยุตฺต\-ปจฺจเยน ปจฺจโย. สหชาตํ, ปจฺฉาชาตํ. (สํขิตฺตํ.) (1)}

\prose{ปรามาส\-วิปฺปยุตฺโต ธมฺโม ปรามาส\-วิปฺปยุตฺตสฺส ธมฺมสฺส วิปฺปยุตฺต\-ปจฺจเยน ปจฺจโย. สหชาตํ, ปุเรชาตํ, ปจฺฉาชาตํ. (สํขิตฺตํ.) (1)}

\prose{ปรามาส\-วิปฺปยุตฺโต ธมฺโม ปรามาส\-สมฺปยุตฺตสฺส ธมฺมสฺส วิปฺปยุตฺต\-ปจฺจเยน ปจฺจโย. ปุเรชาตํ วตฺถุ ปรามาส\-สมฺปยุตฺตกานํ ขนฺธานํ วิปฺปยุตฺต\-ปจฺจเยน ปจฺจโย. (1)}

\proseitem{77}{ปรามาส\-สมฺปยุตฺโต ธมฺโม ปรามาส\-สมฺปยุตฺตสฺส ธมฺมสฺส อตฺถิปจฺจเยน ปจฺจโย. ปรามาส\-สมฺปยุตฺตา เอโก ขนฺโธ ติณฺณนฺนํ ขนฺธานํ อตฺถิปจฺจเยน ปจฺจโย ฯเปฯ เทฺว ขนฺธา ฯเปฯ. (1)}

\prose{ปรามาส\-สมฺปยุตฺโต ธมฺโม ปรามาส\-วิปฺปยุตฺตสฺส ธมฺมสฺส อตฺถิปจฺจเยน ปจฺจโย. ปรามาส\-สมฺปยุตฺตา ขนฺธา จิตฺต\-สมุฏฺฐานานํ รูปานํ อตฺถิปจฺจเยน ปจฺจโย. (มูลํ กาตพฺพํ.) (2)}

\prose{ปรามาส\-สมฺปยุตฺโต ธมฺโม ปรามาส\-สมฺปยุตฺตสฺส จ ปรามาส\-วิปฺปยุตฺตสฺส จ ธมฺมสฺส อตฺถิปจฺจเยน ปจฺจโย. ปรามาส\-สมฺปยุตฺโต เอโก ขนฺโธ ติณฺณนฺนํ ขนฺธานํ จิตฺตสมุฏฺฐานานญฺจ รูปานํ อตฺถิปจฺจเยน ปจฺจโย ฯเปฯ เทฺว ขนฺธา ฯเปฯ. (3)}

\proseitem{78}{ปรามาส\-วิปฺปยุตฺโต ธมฺโม ปรามาส\-วิปฺปยุตฺตสฺส ธมฺมสฺส อตฺถิปจฺจเยน ปจฺจโย. สหชาตํ, ปุเรชาตํ, ปจฺฉาชาตํ, อาหารํ, อินฺทฺริยํ. (สํขิตฺตํ.) (1)}

\prose{ปรามาส\-วิปฺปยุตฺโต ธมฺโม ปรามาส\-สมฺปยุตฺตสฺส ธมฺมสฺส อตฺถิปจฺจเยน ปจฺจโย. ปุเรชาตํ จกฺขุํ ฯเปฯ วตฺถุํ อสฺสาเทติ อภินนฺทติ, ตํ อารพฺภ ปรามาส\-สมฺปยุตฺโต ราโค อุปฺปชฺชติ. วตฺถุ ปรามาส\-สมฺปยุตฺตกานํ ขนฺธานํ อตฺถิปจฺจเยน ปจฺจโย. (2)}

\proseitem{79}{\csromanfolio{364}ปรามาส\-สมฺปยุตฺโต จ ปรามาส\-วิปฺปยุตฺโต จ ธมฺมา ปรามาส\-สมฺปยุตฺตสฺส ธมฺมสฺส อตฺถิปจฺจเยน ปจฺจโย. สหชาตํ, ปุเรชาตํ.  สหชาโต ปรามาส\-สมฺปยุตฺโต เอโก ขนฺโธ จ วตฺถุ จ ติณฺณนฺนํ ขนฺธานํ อตฺถิปจฺจเยน ปจฺจโย ฯเปฯ เทฺว ขนฺธา จ ฯเปฯ. (1)}

\prose{ปรามาส\-สมฺปยุตฺโต จ ปรามาส\-วิปฺปยุตฺโต จ ธมฺมา ปรามาส\-วิปฺปยุตฺตสฺส ธมฺมสฺส อตฺถิปจฺจเยน ปจฺจโย. สหชาตํ, ปจฺฉาชาตํ, อาหารํ, อินฺทฺริยํ.  สหชาตา ปรามาส\-สมฺปยุตฺตา ขนฺธา จ มหาราภูตา จ จิตฺต\-สมุฏฺฐานานํ รูปานํ อตฺถิปจฺจเยน ปจฺจโย. ปจฺฉาชาตา ปรามาส\-สมฺปยุตฺตา ขนฺธา จ กพฬีกาโร อาหาโร จ อิมสฺส กายสฺส อตฺถิปจฺจเยน ปจฺจโย. ปจฺฉาชาตา ปรามาส\-สมฺปยุตฺตา ขนฺธา จ รูปชีวิตินฺทฺริยญฺจ กฏตฺตารูปานํ อตฺถิปจฺจเยน ปจฺจโย. (2)}

\csromanmarkh{1. ปจฺจยานุโลม}\csromanmarkhii{2. สงฺขฺยาวาร}\csromanheadernomark{2}{1. \textbf{ปจฺจยานุโลม 2. สงฺขฺยาวาร}}
\proseitem{80}{เหตุยา จตฺตาริ, อารมฺมเณ จตฺตาริ, อธิปติยา ปญฺจ, อนนฺตเร จตฺตาริ, สมนนฺตเร จตฺตาริ, สหชาเต ปญฺจ, อญฺญมญฺเญ เทฺว, นิสฺสเย สตฺต, อุปนิสฺสเย จตฺตาริ, ปุเรชาเต เทฺว, ปจฺฉาชาเต เทฺว, อาเสวเน เทฺว, กมฺเม จตฺตาริ, วิปาเก เอกํ, อาหาเร จตฺตาริ, อินฺทฺริเย จตฺตาริ, ฌาเน จตฺตาริ, มคฺเค จตฺตาริ, สมฺปยุตฺเต เทฺว, วิปฺปยุตฺเต ตีณิ, อตฺถิยา สตฺต, นตฺถิยา จตฺตาริ, วิคเต จตฺตาริ, อวิคเต สตฺต.}

\vspace{\tipitakaparskip}
\csromancenter{\textbf{อนุโลมํ.}}
\proseitem{81}{ปรามาส\-สมฺปยุตฺโต ธมฺโม ปรามาส\-สมฺปยุตฺตสฺส ธมฺมสฺส อารมฺมณ\-ปจฺจเยน ปจฺจโย. สหชาต\-ปจฺจเยน ปจฺจโย. อุปนิสฺสย\-ปจฺจเยน ปจฺจโย. (1)}

\prose{ปรามาส\-สมฺปยุตฺโต ธมฺโม ปรามาส\-วิปฺปยุตฺตสฺส ธมฺมสฺส อารมฺมณ\-ปจฺจเยน ปจฺจโย. สหชาต\-ปจฺจเยน ปจฺจโย. อุปนิสฺสย\-ปจฺจเยน ปจฺจโย. ปจฺฉาชาต\-ปจฺจเยน ปจฺจโย. กมฺมปจฺจเยน ปจฺจโย. (2)}

\prose{\csromanfolio{365}ปรามาส\-สมฺปยุตฺโต ธมฺโม ปรามาส\-สมฺปยุตฺตสฺส จ ปรามาส\-วิปฺปยุตฺตสฺส จ ธมฺมสฺส สหชาต\-ปจฺจเยน ปจฺจโย. (3)}

\proseitem{82}{ปรามาส\-วิปฺปยุตฺโต ธมฺโม ปรามาส\-วิปฺปยุตฺตสฺส ธมฺมสฺส อารมฺมณ\-ปจฺจเยน ปจฺจโย. สหชาต\-ปจฺจเยน ปจฺจโย. อุปนิสฺสย\-ปจฺจเยน ปจฺจโย. ปุเรชาต\-ปจฺจเยน ปจฺจโย. ปจฺฉาชาต\-ปจฺจเยน ปจฺจโย. กมฺมปจฺจเยน ปจฺจโย. อาหารปจฺจเยน ปจฺจโย. อินฺทฺริย\-ปจฺจเยน ปจฺจโย. (1)}

\prose{ปรามาส\-วิปฺปยุตฺโต ธมฺโม ปรามาส\-สมฺปยุตฺตสฺส ธมฺมสฺส อารมฺมณ\-ปจฺจเยน ปจฺจโย. อุปนิสฺสย\-ปจฺจเยน ปจฺจโย. ปุเรชาต\-ปจฺจเยน ปจฺจโย. (2)}

\prose{ปรามาส\-สมฺปยุตฺโต จ ปรามาส\-วิปฺปยุตฺโต จ ธมฺมา ปรามาส\-สมฺปยุตฺตสฺส ธมฺมสฺส สหชาตํ, ปุเรชาตํ. (1)}

\prose{ปรามาส\-สมฺปยุตฺโต จ ปรามาส\-วิปฺปยุตฺโต จ ธมฺมา ปรามาส\-วิปฺปยุตฺตสฺส ธมฺมสฺส สหชาตํ, ปจฺฉาชาตํ, อาหารํ, อินฺทฺริยํ. (2)}

\csromanmarkh{2. ปจฺจย\-ปจฺจนีย}\csromanmarkhii{2. สงฺขฺยาวาร}\csromanheadernomark{2}{2. ปจฺจย\-ปจฺจนีย 2. สงฺขฺยาวาร}
\proseitem{83}{นเหตุยา สตฺต, นอารมฺมเณ สตฺต, นอธิปติยา สตฺต, นอนนฺตเร สตฺต, นสมนนฺตเร สตฺต, นสหชาเต ปญฺจ, นอญฺญมญฺเญ ปญฺจ, นนิสฺสเย ปญฺจ, นอุปนิสฺสเย สตฺต, นปุเรชาเต ฉ, นปจฺฉาชาเต สตฺต, (สพฺพตฺถ สตฺต,) นมคฺเค สตฺต, นสมฺปยุตฺเต ปญฺจ, นวิปฺปยุตฺเต จตฺตาริ, โนอตฺถิยา จตฺตาริ, โนนตฺถิยา สตฺต, โนวิคเต สตฺต, โนอวิคเต จตฺตาริ.}

\proseitem{84}{เหตุปจฺจยา นอารมฺมเณ จตฺตาริ, นอธิปติยา จตฺตาริ, นอนนฺตเร จตฺตาริ, นสมนนฺตเร จตฺตาริ, นอญฺญมญฺเญ เทฺว, นอุปนิสฺสเย \csromanfolio{366}จตฺตาริ, (สพฺพตฺถ จตฺตาริ,) นมคฺเค จตฺตาริ, นสมฺปยุตฺเต เทฺว, นวิปฺปยุตฺเต เทฺว, โนนตฺถิยา จตฺตาริ, โนวิคเต จตฺตาริ.}

\proseitemwithcloser{85}{นเหตุปจฺจยา อารมฺมเณ จตฺตาริ, อธิปติยา ปญฺจ, (อนุโลมมาติกา กาตพฺพา.) ฯเปฯ อวิคเต สตฺต.}{ปรามาส\-สมฺปยุตฺตทุกํ นิฏฺฐิตํ.}
\csromanmatikaanchor{mtk.51}
\csromantocmarkref{subsection}{53. ปรามาสปรามฏฺฐทุก}{mtk.51}
\bookmark[dest={mtk.51},level=2]{53. ปรามาสปรามฏฺฐทุก}
\csromanmarkh{53. ปรามาส\-ปรามฏฺฐทุก}\csromanmarkhii{1. ปฏิจฺจวาร}\csromanheadernomark{2}{53. ปรามาส\-ปรามฏฺฐทุก 1. ปฏิจฺจวาร}
\csromanmarkh{1. ปจฺจยานุโลม}\csromanmarkhii{1. วิภงฺควาร}\csromanheadernomark{2}{1. \textbf{ปจฺจยานุโลม 1. วิภงฺควาร}}
\csromanheader{2}{\textbf{เหตุ}}
\proseitem{86}{ปรามาสญฺเจว ปรามฏฺฐญฺจ ธมฺมํ ปฏิจฺจ ปรามฏฺโฐ เจว โน จ ปรามาโส ธมฺโม อุปฺปชฺชติ เหตุปจฺจยา. ปรามาสํ ปฏิจฺจ สมฺปยุตฺตกา ขนฺธา จิตฺต\-สมุฏฺฐานญฺจ รูปํ. (1)}

\prose{ปรามฏฺฐ\-ญฺเจว โน จ ปรามาสํ ธมฺมํ ปฏิจฺจ ปรามฏฺโฐ เจว โน จ ปรามาโส ธมฺโม อุปฺปชฺชติ เหตุปจฺจยา. ปรามฏฺฐ\-ญฺเจว โน จ ปรามาสํ เอกํ ขนฺธํ ปฏิจฺจ ตโย ขนฺธา จิตฺต\-สมุฏฺฐานญฺจ รูปํ ฯเปฯ เทฺว ขนฺเธ ฯเปฯ. ปฏิสนฺธิ\-กฺขเณ ฯเปฯ. (ยาว อชฺฌตฺติกา มหาภูตา.) (1)}

\prose{ปรามฏฺฐ\-ญฺเจว โน จ ปรามาสํ ธมฺมํ ปฏิจฺจ ปรามาโส เจว ปรามฏฺโฐ จ ธมฺโม อุปฺปชฺชติ เหตุปจฺจยา. ปรามฏฺเฐ เจว โน จ ปรามาเส ขนฺเธ ปฏิจฺจ ปรามาโส. (2)}

\prose{ปรามฏฺฐ\-ญฺเจว โน จ ปรามาสํ ธมฺมํ ปฏิจฺจ ปรามาโส เจว ปรามฏฺโฐ จ ปรามฏฺโฐ เจว โน จ ปรามาโส จ ธมฺมา อุปฺปชฺชนฺติ เหตุปจฺจยา. ปรามฏฺฐ\-ญฺเจว โน จ ปรามาสํ เอกํ ขนฺธํ ปฏิจฺจ ตโย ขนฺธา ปรามาโส จ จิตฺต\-สมุฏฺฐานญฺจ รูปํ ฯเปฯ เทฺว ขนฺเธ ฯเปฯ. (3)}

\prose{ปรามาสญฺเจว ปรามฏฺฐญฺจ ปรามฏฺฐ\-ญฺเจว โน จ ปรามาสญฺจ ธมฺมํ ปฏิจฺจ ปรามฏฺโฐ เจว โน จ ปรามาโส ธมฺโม อุปฺปชฺชติ เหตุปจฺจยา.}

\vspace{\tipitakaparskip}
\csromancenter{ปรามาส\-ปรามฏฺฐทุก ปรามฏฺฐ\-ญฺเจว โน จ ปรามาสํ เอกํ ขนฺธญฺจ ปรามาสญฺจ ปฏิจฺจ ตโย ขนฺธา จิตฺต\-สมุฏฺฐานญฺจ รูปํ ฯเปฯ เทฺว ขนฺเธ จ ฯเปฯ. (สํขิตฺตํ.)}
\csromancenter{(สพฺเพ วารา ยถา ปรามาสทุกํ, เอวํ กาตพฺพํ, นินฺนานากรณํ.)}
\csromanmarkh{53. ปรามาส\-ปรามฏฺฐทุก}\csromanmarkhii{7. ปญฺหาวาร}\csromanheadernomark{2}{53. ปรามาส\-ปรามฏฺฐทุก 7. ปญฺหาวาร}
\csromanmarkh{1. ปจฺจยานุโลม}\csromanmarkhii{1. วิภงฺควาร}\csromanheadernomark{2}{1. \textbf{ปจฺจยานุโลม 1. วิภงฺควาร}}
\csromanheader{2}{\textbf{เหตุ}}
\proseitem{87}{\csromanfolio{367}ปรามฏฺโฐ เจว โน จ ปรามาโส ธมฺโม ปรามฏฺฐสฺส เจว โน จ ปรามาสสฺส ธมฺมสฺส เหตุปจฺจเยน ปจฺจโย. ปรามฏฺฐา เจว โน จ ปรามาสา เหตู สมฺปยุตฺตกานํ ขนฺธานํ จิตฺตสมุฏฺฐานานญฺจ รูปานํ เหตุปจฺจเยน ปจฺจโย. ปฏิสนฺธิ\-กฺขเณ ฯเปฯ. (1)}

\prose{ปรามฏฺโฐ เจว โน จ ปรามาโส ธมฺโม ปรามาสสฺส เจว ปรามฏฺฐสฺส จ ธมฺมสฺส เหตุปจฺจเยน ปจฺจโย. ปรามฏฺฐา เจว โน จ ปรามาสา เหตู ปรามาสสฺส เหตุปจฺจเยน ปจฺจโย. (2)}

\prose{ปรามฏฺโฐ เจว โน จ ปรามาโส ธมฺโม ปรามาสสฺส เจว ปรามฏฺฐสฺส จ ปรามฏฺฐสฺส เจว โน จ ปรามาสสฺส จ ธมฺมสฺส เหตุปจฺจเยน ปจฺจโย. ปรามฏฺฐา เจว โน จ ปรามาสา เหตู สมฺปยุตฺตกานํ ขนฺธานํ ปรามาสสฺส จ จิตฺตสมุฏฺฐานานญฺจ รูปานํ เหตุปจฺจเยน ปจฺจโย. (3)}

\proseitem{88}{ปรามาโส เจว ปรามฏฺโฐ จ ธมฺโม ปรามาสสฺส เจว ปรามฏฺฐสฺส จ ธมฺมสฺส อารมฺมณ\-ปจฺจเยน ปจฺจโย. ตีณิ. (อารพฺภ กาตพฺพานิ. ปรามาส\-ทุก\-สทิสํ.)}

\proseitem{89}{ปรามฏฺโฐ เจว โน จ ปรามาโส ธมฺโม ปรามฏฺฐสฺส เจว โน จ ปรามาสสฺส ธมฺมสฺส อารมฺมณ\-ปจฺจเยน ปจฺจโย. ทานํ ฯเปฯ สีลํ ฯเปฯ อุโปสถกมฺมํ กตฺวา ตํ ปจฺจเวกฺขติ, อสฺสาเทติ อภินนฺทติ, ตํ อารพฺภ ราโค ฯเปฯ. วิจิกิจฺฉา. อุทฺธจฺจํ. โทมนสฺสํ อุปฺปชฺชติ. ปุพฺเพ ฯเปฯ. ฌานา ฯเปฯ. อริยา โคตฺรภุํ ปจฺจเวกฺขนฺติ, โวทานํ ปจฺจเวกฺขนฺติ. ปหีเน กิเลเส ฯเปฯ. \csromanfolio{368}วิกฺขมฺภิเต กิเลเส ฯเปฯ. ปุพฺเพ ฯเปฯ. จกฺขุํ ฯเปฯ วตฺถุํ ปรามฏฺเฐ เจว โน จ ปรามาเส ขนฺเธ อนิจฺจโต ฯเปฯ โทมนสฺสํ อุปฺปชฺชติ. ทิพฺเพน \textbf{จกฺขุนา รูปํ ปสฺสติ. (ยาว อาวชฺชนา สพฺพํ กาตพฺพํ.) (1)}}

\prose{ปรามฏฺโฐ เจว โน จ ปรามาโส ธมฺโม ปรามาสสฺส เจว ปรามฏฺฐสฺส จ ธมฺมสฺส อารมฺมณ\-ปจฺจเยน ปจฺจโย. ทานํ ฯเปฯ สีลํ ฯเปฯ อุโปสถกมฺมํ ฯเปฯ. ปุพฺเพ ฯเปฯ. ฌานา ฯเปฯ. จกฺขุํ ฯเปฯ. วตฺถุํ ปรามฏฺเฐ เจว โน จ ปรามาเส ขนฺเธ อสฺสาเทติ อภินนฺทติ, ตํ อารพฺภ ทิฏฺฐิ อุปฺปชฺชติ. (2)}

\prose{ปรามฏฺโฐ เจว โน จ ปรามาโส ธมฺโม ปรามาสสฺส เจว ปรามฏฺฐสฺส จ ปรามฏฺฐสฺส เจว โน จ ปรามาสสฺส จ ธมฺมสฺส อารมฺมณ\-ปจฺจเยน ปจฺจโย. ทานํ ฯเปฯ สีลํ ฯเปฯ อุโปสถกมฺมํ ฯเปฯ. ปุพฺเพ ฯเปฯ. ฌานา ฯเปฯ. จกฺขุํ ฯเปฯ วตฺถุํ ปรามฏฺเฐ เจว โน จ ปรามาเส ขนฺเธ อสฺสาเทติ อภินนฺทติ, ตํ อารพฺภ ปรามาโส จ สมฺปยุตฺตกา จ ขนฺธา อุปฺปชฺชนฺติ. (3)}

\prose{(เอวํ อิตเรปิ ตีณิ. อารพฺภ กาตพฺพานิ. อิมํ ทุกํ ปรามาส\-ทุก\-สทิสํ. โลกุตฺตรํ ยหิํ น ลพฺภติ, ตหิํ น กาตพฺพํ.)}

\csromanmatikaanchor{mtk.52}
\csromantocmarkref{subsection}{54. ปรามาสวิปฺปยุตฺตปรามฏฺฐทุก}{mtk.52}
\bookmark[dest={mtk.52},level=2]{54. ปรามาสวิปฺปยุตฺตปรามฏฺฐทุก}
\csromanmarkh{54. ปรามาส\-วิปฺปยุตฺต\-ปรามฏฺฐทุก}\csromanmarkhii{1. ปฏิจฺจวาร}\csromanheadernomark{2}{54. ปรามาส\-วิปฺปยุตฺต\-ปรามฏฺฐทุก 1. ปฏิจฺจวาร}
\csromanmarkh{1. ปจฺจยานุโลม}\csromanmarkhii{1. วิภงฺควาร}\csromanheadernomark{2}{1. \textbf{ปจฺจยานุโลม 1. วิภงฺควาร}}
\csromanheader{2}{\textbf{เหตุ}}
\proseitem{90}{ปรามาส\-วิปฺปยุตฺตํ ปรามฏฺฐํ ธมฺมํ ปฏิจฺจ ปรามาส\-วิปฺปยุตฺโต ปรามฏฺโฐ ธมฺโม อุปฺปชฺชติ เหตุปจฺจยา. ปรามาส\-วิปฺปยุตฺตํ ปรามฏฺฐํ เอกํ ขนฺธํ ปฏิจฺจ ตโย ขนฺธา จิตฺต\-สมุฏฺฐานญฺจ รูปํ ฯเปฯ เทฺว ขนฺเธ ฯเปฯ.}

\prose{ปรามาส\-วิปฺปยุตฺตํ อปรามฏฺฐํ ธมฺมํ ปฏิจฺจ ปรามาส\-วิปฺปยุตฺโต อปรามฏฺโฐ ธมฺโม อุปฺปชฺชติ เหตุปจฺจยา. (สํขิตฺตํ.)}

\prosewithcloser{(ยถา จูฬนฺตรทุเก โลกิยทุกํ, เอวํ กาตพฺพํ, นินฺนานากรณํ.)}{ปรามาส\-วิปฺปยุตฺต\-ปรามฏฺฐทุกํ นิฏฺฐิตํ.}
\nitthitam{ปรามาส\-โคจฺฉกํ นิฏฺฐิตํ.}
\csromanmatikaanchor{mtk.53}
\csromanmatikaanchor{mtk.54}
\csromantocmarknopage{section}{10. มหนฺตรทุก}{mtk.53}
\bookmark[dest={mtk.53},level=1]{10. มหนฺตรทุก}
\csromantocmarkref{subsection}{55. สารมฺมณทุก}{mtk.54}
\bookmark[dest={mtk.54},level=2]{55. สารมฺมณทุก}
\csromanmarkh{55. สารมฺมณทุก}\csromanmarkhii{10. มหนฺตรทุก}\csromanheadernomark{2}{55. สารมฺมณทุก \textbf{10. มหนฺตรทุก}}
\csromanmarkh{55. สารมฺมณทุก}\csromanmarkhii{1. ปฏิจฺจวาร}\csromanheadernomark{2}{55. สารมฺมณทุก 1. ปฏิจฺจวาร}
\csromanmarkh{1. ปจฺจยานุโลม}\csromanmarkhii{1. วิภงฺควาร}\csromanheadernomark{2}{1. \textbf{ปจฺจยานุโลม 1. วิภงฺควาร}}
\csromanheader{2}{\textbf{เหตุ}}
\proseitem{1}{\csromanfolio{369}สารมฺมณํ ธมฺมํ ปฏิจฺจ สารมฺมโณ ธมฺโม อุปฺปชฺชติ เหตุปจฺจยา. สารมฺมณํ เอกํ ขนฺธํ ปฏิจฺจ ตโย ขนฺธา ฯเปฯ เทฺว ขนฺเธ ฯเปฯ. ปฏิสนฺธิ\-กฺขเณ ฯเปฯ. (1)}

\prose{สารมฺมณํ ธมฺมํ ปฏิจฺจ อนารมฺมโณ ธมฺโม อุปฺปชฺชติ เหตุปจฺจยา. สารมฺมเณ ขนฺเธ ปฏิจฺจ จิตฺต\-สมุฏฺฐานํ รูปํ. ปฏิสนฺธิ\-กฺขเณ ฯเปฯ. (2)}

\prose{สารมฺมณํ ธมฺมํ ปฏิจฺจ สารมฺมโณ จ อนารมฺมโณ จ ธมฺมา อุปฺปชฺชนฺติ เหตุปจฺจยา. สารมฺมณํ เอกํ ขนฺธํ ปฏิจฺจ ตโย ขนฺธา จิตฺต\-สมุฏฺฐานญฺจ รูปํ ฯเปฯ เทฺว ขนฺเธ ฯเปฯ. ปฏิสนฺธิ\-กฺขเณ ฯเปฯ. (3)}

\proseitem{2}{อนารมฺมณํ ธมฺมํ ปฏิจฺจ อนารมฺมโณ ธมฺโม อุปฺปชฺชติ เหตุปจฺจยา. เอกํ มหาภูตํ ฯเปฯ มหาภูเต ปฏิจฺจ จิตฺต\-สมุฏฺฐานํ รูปํ, กฏตฺตารูปํ, อุปาทารูปํ. (1)}

\prose{อนารมฺมณํ ธมฺมํ ปฏิจฺจ สารมฺมโณ ธมฺโม อุปฺปชฺชติ เหตุปจฺจยา. ปฏิสนฺธิ\-กฺขเณ วตฺถุํ ปฏิจฺจ สารมฺมณา ขนฺธา. (2)}

\prose{อนารมฺมณํ ธมฺมํ ปฏิจฺจ สารมฺมโณ จ อนารมฺมโณ จ ธมฺมา อุปฺปชฺชนฺติ เหตุปจฺจยา. ปฏิสนฺธิ\-กฺขเณ วตฺถุํ ปฏิจฺจ สารมฺมณา ขนฺธา, มหาภูเต ปฏิจฺจ กฏตฺตารูปํ. (3)}

\proseitem{3}{สารมฺมณญฺจ อนารมฺมณญฺจ ธมฺมํ ปฏิจฺจ สารมฺมโณ ธมฺโม อุปฺปชฺชติ เหตุปจฺจยา. ปฏิสนฺธิ\-กฺขเณ สารมฺมณํ เอกํ ขนฺธญฺจ วตฺถุญฺจ ปฏิจฺจ ตโย ขนฺธา ฯเปฯ เทฺว ขนฺเธ จ ฯเปฯ. (1)}

\prose{\csromanfolio{370}สารมฺมณญฺจ อนารมฺมณญฺจ ธมฺมํ ปฏิจฺจ อนารมฺมโณ ธมฺโม อุปฺปชฺชติ เหตุปจฺจยา. สารมฺมเณ ขนฺเธ จ มหาภูเต จ ปฏิจฺจ จิตฺต\-สมุฏฺฐานํ รูปํ. ปฏิสนฺธิ\-กฺขเณ ฯเปฯ. (2)}

\prose{สารมฺมณญฺจ อนารมฺมณญฺจ ธมฺมํ ปฏิจฺจ สารมฺมโณ จ อนารมฺมโณ จ ธมฺมา อุปฺปชฺชนฺติ เหตุปจฺจยา. ปฏิสนฺธิ\-กฺขเณ สารมฺมณํ เอกํ ขนฺธญฺจ วตฺถุญฺจ ปฏิจฺจ ตโย ขนฺธา ฯเปฯ เทฺว ขนฺเธ จ ฯเปฯ. สารมฺมเณ ขนฺเธ จ มหาภูเต จ ปฏิจฺจ กฏตฺตารูปํ. (3)}

\proseitem{4}{สารมฺมณํ ธมฺมํ ปฏิจฺจ สารมฺมโณ ธมฺโม อุปฺปชฺชติ อารมฺมณ\-ปจฺจยา. สารมฺมณํ เอกํ ขนฺธํ ปฏิจฺจ ตโย ขนฺธา ฯเปฯ เทฺว ขนฺเธ ฯเปฯ. ปฏิสนฺธิ\-กฺขเณ ฯเปฯ.}

\prose{อนารมฺมณํ ธมฺมํ ปฏิจฺจ สารมฺมโณ ธมฺโม อุปฺปชฺชติ อารมฺมณ\-ปจฺจยา. ปฏิสนฺธิ\-กฺขเณ วตฺถุํ ปฏิจฺจ สารมฺมณา ขนฺธา.}

\prose{สารมฺมณญฺจ อนารมฺมณญฺจ ธมฺมํ ปฏิจฺจ สารมฺมโณ ธมฺโม อุปฺปชฺชติ อารมฺมณ\-ปจฺจยา. ปฏิสนฺธิ\-กฺขเณ สารมฺมณํ เอกํ ขนฺธญฺจ วตฺถุญฺจ ปฏิจฺจ ตโย ขนฺธา ฯเปฯ เทฺว ขนฺเธ จ ฯเปฯ. (สํขิตฺตํ.)}

\csromanmarkh{1. ปจฺจยานุโลม}\csromanmarkhii{2. สญฺขฺยาวาร}\csromanheadernomark{2}{1. \textbf{ปจฺจยานุโลม 2. สญฺขฺยาวาร}}
\proseitem{5}{เหตุยา นว, อารมฺมเณ ตีณิ, อธิปติยา ปญฺจ, อนนฺตเร ตีณิ, สมนนฺตเร ตีณิ, สหชาเต นว, อญฺญมญฺเญ ฉ, นิสฺสเย นว, อุปนิสฺสเย ตีณิ, ปุเรชาเต เอกํ, อาเสวเน เอกํ, กมฺเม นว, วิปาเก นว, อาหาเร นว, อินฺทฺริเย นว, ฌาเน นว, มคฺเค นว, สมฺปยุตฺเต ตีณิ, วิปฺปยุตฺเต นว, อตฺถิยา นว, นตฺถิยา ตีณิ, วิคเต ตีณิ, อวิคเต นว.}

\csromanmarkh{55. สารมฺมณทุก}\csromanmarkhii{2. ปจฺจย\-ปจฺจนีย 1. วิภงฺควาร}\csromanheadernomark{2}{55. สารมฺมณทุก 2. ปจฺจย\-ปจฺจนีย 1. วิภงฺควาร}
\proseitem{6}{\csromanfolio{371}สารมฺมณํ ธมฺมํ ปฏิจฺจ สารมฺมโณ ธมฺโม อุปฺปชฺชติ นเหตุปจฺจยา. อเหตุกํ สารมฺมณํ เอกํ ขนฺธํ ปฏิจฺจ ตโย ขนฺธา ฯเปฯ เทฺว ขนฺเธ ฯเปฯ. อเหตุก\-ปฏิสนฺธิ\-กฺขเณ ฯเปฯ. วิจิกิจฺฉา\-สหคเต อุทฺธจฺจ\-สหคเต ขนฺเธ ปฏิจฺจ วิจิกิจฺฉา\-สหคโต อุทฺธจฺจ\-สหคโต โมโห. (1)}

\prose{สารมฺมณํ ธมฺมํ ปฏิจฺจ อนารมฺมโณ ธมฺโม อุปฺปชฺชติ นเหตุปจฺจยา. อเหตุเก สารมฺมเณ ขนฺเธ ปฏิจฺจ จิตฺต\-สมุฏฺฐานํ รูปํ. อเหตุก\-ปฏิสนฺธิ\-กฺขเณ ฯเปฯ. (2)}

\prose{สารมฺมณํ ธมฺมํ ปฏิจฺจ สารมฺมโณ จ อนารมฺมโณ จ ธมฺมา อุปฺปชฺชนฺติ นเหตุปจฺจยา. อเหตุกํ สารมฺมณํ เอกํ ขนฺธํ ปฏิจฺจ ตโย ขนฺธา จิตฺต\-สมุฏฺฐานญฺจ รูปํ ฯเปฯ เทฺว ขนฺเธ ฯเปฯ. อเหตุก\-ปฏิสนฺธิ\-กฺขเณ ฯเปฯ. (3)}

\proseitem{7}{อนารมฺมณํ ธมฺมํ ปฏิจฺจ อนารมฺมโณ ธมฺโม อุปฺปชฺชติ นเหตุปจฺจยา. เอกํ มหาภูตํ ฯเปฯ. อสญฺญสตฺตานํ เอกํ มหาภูตํ ฯเปฯ. (1)}

\prose{อนารมฺมณํ ธมฺมํ ปฏิจฺจ สารมฺมโณ ธมฺโม อุปฺปชฺชติ นเหตุปจฺจยา. อเหตุก\-ปฏิสนฺธิ\-กฺขเณ วตฺถุํ ปฏิจฺจ สารมฺมณา ขนฺธา. (2)}

\prose{อนารมฺมณํ ธมฺมํ ปฏิจฺจ สารมฺมโณ จ อนารมฺมโณ จ ธมฺมา อุปฺปชฺชนฺติ นเหตุปจฺจยา. อเหตุก\-ปฏิสนฺธิ\-กฺขเณ วตฺถุํ ปฏิจฺจ สารมฺมณา ขนฺธา. มหาภูเต ปฏิจฺจ กฏตฺตารูปํ. (3)}

\proseitem{8}{สารมฺมณญฺจ อนารมฺมณญฺจ ธมฺมํ ปฏิจฺจ สารมฺมโณ ธมฺโม อุปฺปชฺชติ นเหตุปจฺจยา. อเหตุก\-ปฏิสนฺธิ\-กฺขเณ สารมฺมณํ เอกํ ขนฺธญฺจ วตฺถุญฺจ ปฏิจฺจ ตโย ขนฺธา ฯเปฯ เทฺว ขนฺเธ จ ฯเปฯ. (1)}

\prose{สารมฺมณญฺจ อนารมฺมณญฺจ ธมฺมํ ปฏิจฺจ อนารมฺมโณ ธมฺโม อุปฺปชฺชติ นเหตุปจฺจยา. อเหตุเก สารมฺมเณ ขนฺเธ จ มหาภูเต จ ปฏิจฺจ จิตฺต\-สมุฏฺฐานํ รูปํ. อเหตุก\-ปฏิสนฺธิ\-กฺขเณ ฯเปฯ. (2)}

\prose{สารมฺมณญฺจ อนารมฺมณญฺจ ธมฺมํ ปฏิจฺจ สารมฺมโณ จ อนารมฺมโณ จ ธมฺมา อุปฺปชฺชนฺติ นเหตุปจฺจยา. อเหตุก\-ปฏิสนฺธิ\-กฺขเณ สารมฺมณํ เอกํ \csromanfolio{372}\textbf{ขนฺธญฺจ วตฺถุญฺจ ปฏิจฺจ ตโย ขนฺธา ฯเปฯ เทฺว ขนฺเธ จ ฯเปฯ.} \textbf{สารมฺมเณ ขนฺเธ จ มหาภูเต จ ปฏิจฺจ กฏตฺตารูปํ. (3)}}

\prose{นอารมฺมณ}

\proseitem{8}{สารมฺมณํ ธมฺมํ ปฏิจฺจ อนารมฺมโณ ธมฺโม อุปฺปชฺชติ นอารมฺมณ\-ปจฺจยา. สารมฺมเณ ขนฺเธ ปฏิจฺจ จิตฺต\-สมุฏฺฐานํ รูปํ. ปฏิสนฺธิ\-กฺขเณ ฯเปฯ. (1)}

\prose{อนารมฺมณํ ธมฺมํ ปฏิจฺจ อนารมฺมโณ ธมฺโม อุปฺปชฺชติ นอารมฺมณ\-ปจฺจยา. (ยาว อสญฺญสตฺตา.) (1)}

\prose{สารมฺมณญฺจ อนารมฺมณญฺจ ธมฺมํ ปฏิจฺจ อนารมฺมโณ ธมฺโม อุปฺปชฺชติ นอารมฺมณ\-ปจฺจยา. สารมฺมเณ ขนฺเธ จ มหาภูเต จ ปฏิจฺจ จิตฺต\-สมุฏฺฐานํ รูปํ. ปฏิสนฺธิ\-กฺขเณ ฯเปฯ. (สํขิตฺตํ.) (1)}

\csromanmarkh{2. ปจฺจย\-ปจฺจนีย}\csromanmarkhii{2. สงฺขฺยาวาร}\csromanheadernomark{2}{2. ปจฺจย\-ปจฺจนีย 2. สงฺขฺยาวาร}
\proseitem{8}{นเหตุยา นว, นอารมฺมเณ ตีณิ, นอธิปติยา นว, นอนนฺตเร ตีณิ, นสมนนฺตเร ตีณิ, นอญฺญมญฺเญ ตีณิ, นอุปนิสฺสเย ตีณิ, นปุเรชาเต นว, นปจฺฉาชาเต นว, นอาเสวเน นว, นกมฺเม เทฺว, นวิปาเก ปญฺจ, นอาหาเร เอกํ, นอินฺทฺริเย เอกํ, นฌาเน เทฺว, นมคฺเค นว, นสมฺปยุตฺเต ตีณิ, นวิปฺปยุตฺเต เทฺว, โนนตฺถิยา ตีณิ, โนวิคเต ตีณิ.}
\csromansectionrule

\prose{เหตุปจฺจยา นอารมฺมเณ ตีณิ, นอธิปติยา นว ฯเปฯ นกมฺเม เอกํ ฯเปฯ นวิปฺปยุตฺเต เอกํ. (สํขิตฺตํ.)}

\prose{นเหตุปจฺจยา อารมฺมเณ ตีณิ ฯเปฯ สหชาเต นว ฯเปฯ มคฺเค เอกํ ฯเปฯ อวิคเต นว.}

\vspace{\tipitakaparskip}
\csromancenter{(สหชาตวาโรปิ ปฏิจฺจ\-วาร\-สทิโส.)}
\csromanmarkh{55. สารมฺมณทุก}\csromanmarkhii{55. สารมฺมณทุก 3. ปจฺจยวาร}\csromanheadernomark{2}{55. สารมฺมณทุก \textbf{55. สารมฺมณทุก 3. ปจฺจยวาร}}
\csromanmarkh{1. ปจฺจยานุโลม}\csromanmarkhii{1. วิภงฺควาร}\csromanheadernomark{2}{1. \textbf{ปจฺจยานุโลม 1. วิภงฺควาร}}
\csromanheader{2}{\textbf{เหตุ}}
\proseitem{13}{\csromanfolio{373}สารมฺมณํ ธมฺมํ ปจฺจยา สารมฺมโณ ธมฺโม อุปฺปชฺชติ เหตุปจฺจยา. ตีณิ. (ปฏิจฺจสทิสา.)}

\proseitem{14}{อนารมฺมณํ ธมฺมํ ปจฺจยา อนารมฺมโณ ธมฺโม อุปฺปชฺชติ เหตุปจฺจยา. เอกํ มหาภูตํ ฯเปฯ มหาภูเต ปจฺจยา จิตฺต\-สมุฏฺฐานํ รูปํ, กฏตฺตารูปํ, อุปาทารูปํ. (1)}

\prose{อนารมฺมณํ ธมฺมํ ปจฺจยา สารมฺมโณ ธมฺโม อุปฺปชฺชติ เหตุปจฺจยา. วตฺถุํ ปจฺจยา สารมฺมณา ขนฺธา. ปฏิสนฺธิ\-กฺขเณ วตฺถุํ ปจฺจยา สารมฺมณา ขนฺธา. (2)}

\prose{อนารมฺมณํ ธมฺมํ ปจฺจยา สารมฺมโณ จ อนารมฺมโณ จ ธมฺมา อุปฺปชฺชนฺติ เหตุปจฺจยา. วตฺถุํ ปจฺจยา สารมฺมณา ขนฺธา. มหาภูเต ปจฺจยา จิตฺต\-สมุฏฺฐานํ รูปํ. ปฏิสนฺธิ\-กฺขเณ ฯเปฯ. (3)}

\proseitem{15}{สารมฺมณญฺจ อนารมฺมณญฺจ ธมฺมํ ปจฺจยา สารมฺมโณ ธมฺโม อุปฺปชฺชติ เหตุปจฺจยา. สารมฺมณํ เอกํ ขนฺธญฺจ วตฺถุญฺจ ปจฺจยา ตโย ขนฺธา ฯเปฯ เทฺว ขนฺเธ จ ฯเปฯ. ปฏิสนฺธิ\-กฺขเณ ฯเปฯ. (1)}

\prose{สารมฺมณญฺจ อนารมฺมณญฺจ ธมฺมํ ปจฺจยา อนารมฺมโณ ธมฺโม อุปฺปชฺชติ เหตุปจฺจยา. สารมฺมเณ ขนฺเธ จ มหาภูเต จ ปจฺจยา จิตฺต\-สมุฏฺฐานํ รูปํ. ปฏิสนฺธิ\-กฺขเณ ฯเปฯ. (2)}

\prose{สารมฺมณญฺจ อนารมฺมณญฺจ ธมฺมํ ปจฺจยา สารมฺมโณ จ อนารมฺมโณ จ ธมฺมา อุปฺปชฺชนฺติ เหตุปจฺจยา. สารมฺมณํ เอกํ ขนฺธญฺจ วตฺถุญฺจ ปจฺจยา ตโย ขนฺธา ฯเปฯ เทฺว ขนฺเธ จ ฯเปฯ. สารมฺมเณ ขนฺเธ จ มหาภูเต จ ปจฺจยา จิตฺต\-สมุฏฺฐานํ รูปํ. ปฏิสนฺธิ\-กฺขเณ ฯเปฯ. (3)}

\proseitem{16}{สารมฺมณํ ธมฺมํ ปจฺจยา สารมฺมโณ ธมฺโม อุปฺปชฺชติ อารมฺมณ\-ปจฺจยา. สารมฺมณํ เอกํ ขนฺธํ ปจฺจยา ตโย ขนฺธา ฯเปฯ เทฺว ขนฺเธ ฯเปฯ. ปฏิสนฺธิ\-กฺขเณ ฯเปฯ. (1)}

\prose{\csromanfolio{374}อนารมฺมณํ ธมฺมํ ปจฺจยา สารมฺมโณ ธมฺโม อุปฺปชฺชติ อารมฺมณ\-ปจฺจยา. จกฺขายตนํ ปจฺจยา จกฺขุวิญฺญาณํ ฯเปฯ. กายายตนํ ปจฺจยา กายวิญฺญาณํ. วตฺถุํ ปจฺจยา สารมฺมณา ขนฺธา. ปฏิสนฺธิ\-กฺขเณ ฯเปฯ. (1)}

\prose{สารมฺมณญฺจ อนารมฺมณญฺจ ธมฺมํ ปจฺจยา สารมฺมโณ ธมฺโม อุปฺปชฺชติ อารมฺมณ\-ปจฺจยา. จกฺขุ\-วิญฺญาณสหคตํ เอกํ ขนฺธญฺจ จกฺขายตนญฺจ ปจฺจยา ตโย ขนฺธา ฯเปฯ เทฺว ขนฺเธ จ ฯเปฯ. กาย\-วิญฺญาณสหคตํ ฯเปฯ. สารมฺมณํ เอกํ ขนฺธญฺจ วตฺถุญฺจ ปจฺจยา ตโย ขนฺธา ฯเปฯ เทฺว ขนฺเธ จ ฯเปฯ. ปฏิสนฺธิ\-กฺขเณ ฯเปฯ. (สํขิตฺตํ.) (1)}

\csromanmarkh{1. ปจฺจยานุโลม}\csromanmarkhii{2. สงฺขฺยาวาร}\csromanheadernomark{2}{1. \textbf{ปจฺจยานุโลม 2. สงฺขฺยาวาร}}
\proseitem{17}{เหตุยา นว, อารมฺมเณ ตีณิ, อธิปติยา นว, อนนฺตเร ตีณิ, สมนนฺตเร ตีณิ, สหชาเต นว, อญฺญมญฺเญ ฉ, นิสฺสเย นว, อุปนิสฺสเย ตีณิ, ปุเรชาเต ตีณิ, อาเสวเน ตีณิ, กมฺเม นว, วิปาเก นว, อาหาเร นว, อินฺทฺริเย นว, ฌาเน นว, มคฺเค นว, สมฺปยุตฺเต ตีณิ, วิปฺปยุตฺเต นว, อตฺถิยา นว, นตฺถิยา ตีณิ, วิคเต ตีณิ, อวิคเต นว.}

\csromanmarkh{2. ปจฺจย\-ปจฺจนีย}\csromanmarkhii{1. วิภงฺควาร}\csromanheadernomark{2}{2. ปจฺจย\-ปจฺจนีย 1. วิภงฺควาร}
\proseitem{18}{สารมฺมณํ ธมฺมํ ปจฺจยา สารมฺมโณ ธมฺโม อุปฺปชฺชติ นเหตุปจฺจยา. ตีณิ. (ปฏิจฺจ\-วาร\-สทิสํ.)}

\prose{อนารมฺมณํ ธมฺมํ ปจฺจยา อนารมฺมโณ ธมฺโม อุปฺปชฺชติ นเหตุปจฺจยา. เอกํ มหาภูตํ ฯเปฯ. อสญฺญสตฺตานํ เอกํ มหาภูตํ ฯเปฯ. (1)}

\prose{อนารมฺมณํ ธมฺมํ ปจฺจยา สารมฺมโณ ธมฺโม อุปฺปชฺชติ นเหตุปจฺจยา. จกฺขายตนํ ปจฺจยา จกฺขุวิญฺญาณํ ฯเปฯ. กายายตนํ ปจฺจยา กายวิญฺญาณํ. วตฺถุํ ปจฺจยา อเหตุกา สารมฺมณา ขนฺธา. ปฏิสนฺธิ\-กฺขเณ ฯเปฯ. วตฺถุํ ปจฺจยา วิจิกิจฺฉา\-สหคโต อุทฺธจฺจ\-สหคโต โมโห. (2)}

\prose{\csromanfolio{375}อนารมฺมณํ ธมฺมํ ปจฺจยา สารมฺมโณ จ อนารมฺมโณ จ ธมฺมา อุปฺปชฺชนฺติ นเหตุปจฺจยา. วตฺถุํ ปจฺจยา อเหตุกา สารมฺมณา ขนฺธา. มหาภูเต ปจฺจยา จิตฺต\-สมุฏฺฐานํ รูปํ. ปฏิสนฺธิ\-กฺขเณ ฯเปฯ. (3)}

\proseitem{19}{สารมฺมณญฺจ อนารมฺมณญฺจ ธมฺมํ ปจฺจยา สารมฺมโณ ธมฺโม อุปฺปชฺชติ นเหตุปจฺจยา. จกฺขุ\-วิญฺญาณสหคตํ เอกํ ขนฺธญฺจ จกฺขายตนญฺจ ปจฺจยา ตโย ขนฺธา ฯเปฯ เทฺว ขนฺเธ จ ฯเปฯ. กาย\-วิญฺญาณสหคตํ ฯเปฯ. อเหตุกํ สารมฺมณํ เอกํ ขนฺธญฺจ วตฺถุญฺจ ปจฺจยา ตโย ขนฺธา ฯเปฯ เทฺว ขนฺเธ จ ฯเปฯ. อเหตุก\-ปฏิสนฺธิ\-กฺขเณ ฯเปฯ. วิจิกิจฺฉา\-สหคเต อุทฺธจฺจ\-สหคเต ขนฺเธ จ วตฺถุญฺจ ปจฺจยา วิจิกิจฺฉา\-สหคโต อุทฺธจฺจ\-สหคโต โมโห. (1)}

\prose{สารมฺมณญฺจ อนารมฺมณญฺจ ธมฺมํ ปจฺจยา อนารมฺมโณ ธมฺโม อุปฺปชฺชติ นเหตุปจฺจยา. อเหตุเก สารมฺมเณ ขนฺเธ จ มหาภูเต จ ปจฺจยา จิตฺต\-สมุฏฺฐานํ รูปํ. ปฏิสนฺธิ\-กฺขเณ ฯเปฯ. (2)}

\prose{สารมฺมณญฺจ อนารมฺมณญฺจ ธมฺมํ ปจฺจยา สารมฺมโณ จ อนารมฺมโณ จ ธมฺมา อุปฺปชฺชนฺติ นเหตุปจฺจยา. อเหตุกํ สารมฺมณํ เอกํ ขนฺธญฺจ วตฺถุญฺจ ปจฺจยา ตโย ขนฺธา ฯเปฯ เทฺว ขนฺเธ จ ฯเปฯ. อเหตุเก สารมฺมเณ ขนฺเธ จ มหาภูเต จ ปจฺจยา จิตฺต\-สมุฏฺฐานํ รูปํ. อเหตุก\-ปฏิสนฺธิ\-กฺขเณ อเหตุกํ สารมฺมณํ เอกํ ขนฺธญฺจ วตฺถุญฺจ ปจฺจยา ตโย ขนฺธา ฯเปฯ เทฺว ขนฺเธ จ ฯเปฯ. อเหตุเก สารมฺมเณ ขนฺเธ จ มหาภูเต จ ปจฺจยา กฏตฺตารูปํ ฯเปฯ. (สํขิตฺตํ.) (3)}

\csromanmarkh{2. ปจฺจย\-ปจฺจนีย}\csromanmarkhii{2. สงฺขฺยาวาร}\csromanheadernomark{2}{2. ปจฺจย\-ปจฺจนีย 2. สงฺขฺยาวาร}
\proseitem{20}{นเหตุยา นว, นอารมฺมเณ ตีณิ, นอธิปติยา นว, นอนนฺตเร ตีณิ, นสมนนฺตเร ตีณิ, นอญฺญมญฺเญ ตีณิ, นอุปนิสฺสเย ตีณิ, นปุเรชาเต นว, นปจฺฉาชาเต นว, นอาเสวเน นว, นกมฺเม จตฺตาริ, นวิปาเก นว, นอาหาเร เอกํ, นอินฺทฺริเย เอกํ, นฌาเน จตฺตาริ, นมคฺเค นว, นสมฺปยุตฺเต ตีณิ, นวิปฺปยุตฺเต เทฺว, โนนตฺถิยา ตีณิ, โนวิคเต ตีณิ.}

\proseitem{21}{\csromanfolio{376}เหตุปจฺจยา นอารมฺมเณ ตีณิ, นอธิปติยา นว ฯเปฯ นกมฺเม เอกํ ฯเปฯ นวิปฺปยุตฺเต เอกํ. (สํขิตฺตํ.)}

\proseitem{22}{นเหตุปจฺจยา อารมฺมเณ ตีณิ, อนนฺตเร ตีณิ, สมนนฺตเร ตีณิ, สหชาเต นว ฯเปฯ มคฺเค ตีณิ ฯเปฯ อวิคเต นว.}

\vspace{\tipitakaparskip}
\csromancenter{(นิสฺสยวาโร ปจฺจยวาร\-สทิโส.)}
\csromanmarkh{55. สารมฺมณทุก}\csromanmarkhii{5. สํสฏฺฐวาร}\csromanheadernomark{2}{55. \textbf{สารมฺมณทุก 5. สํสฏฺฐวาร}}
\proseitem{23}{สารมฺมณํ ธมฺมํ สํสฏฺโฐ สารมฺมโณ ธมฺโม อุปฺปชฺชติ เหตุปจฺจยา. สารมฺมณํ เอกํ ขนฺธํ สํสฏฺฐา ตโย ขนฺธา ฯเปฯ เทฺว ขนฺเธ สํสฏฺฐา เทฺว ขนฺธา. ปฏิสนฺธิ\-กฺขเณ ฯเปฯ.}

\proseitem{24}{เหตุยา เอกํ, อารมฺมเณ เอกํ, อธิปติยา เอกํ, (สพฺพตฺถ เอกํ.) อวิคเต เอกํ. อนุโลมํ.}

\proseitem{25}{สารมฺมณํ ธมฺมํ สํสฏฺโฐ สารมฺมโณ ธมฺโม อุปฺปชฺชติ นเหตุปจฺจยา. อเหตุกํ สารมฺมณํ เอกํ ขนฺธํ สํสฏฺฐา ตโย ขนฺธา ฯเปฯ เทฺว ขนฺเธ ฯเปฯ. วิจิกิจฺฉา\-สหคเต อุทฺธจฺจ\-สหคเต ขนฺเธ สํสฏฺโฐ วิจิกิจฺฉา\-สหคโต อุทฺธจฺจ\-สหคโต โมโห. (สํขิตฺตํ.)}

\proseitem{26}{นเหตุยา เอกํ, นอธิปติยา เอกํ, นปุเรชาเต เอกํ, นปจฺฉาชาเต เอกํ, นอาเสวเน เอกํ, นกมฺเม เอกํ, นวิปาเก เอกํ, นฌาเน เอกํ, นมคฺเค เอกํ, นวิปฺปยุตฺเต เอกํ.}

\vspace{\tipitakaparskip}
\csromancenter{ปจฺจนียํ.}
\csromancenter{(เอวํ อิตเร เทฺว คณนาปิ สมฺปยุตฺตวาโรปิ กาตพฺพา.)}
\csromanmarkh{55. สารมฺมณทุก}\csromanmarkhii{55. สารมฺมณทุก 7. ปญฺหาวาร}\csromanheadernomark{2}{55. สารมฺมณทุก \textbf{55. สารมฺมณทุก 7. ปญฺหาวาร}}
\csromanmarkh{1. ปจฺจยานุโลม}\csromanmarkhii{1. วิภงฺควาร}\csromanheadernomark{2}{1. \textbf{ปจฺจยานุโลม 1. วิภงฺควาร}}
\csromanheader{2}{\textbf{เหตุ}}
\proseitem{27}{\csromanfolio{377}สารมฺมโณ ธมฺโม สารมฺมณสฺส ธมฺมสฺส เหตุปจฺจเยน ปจฺจโย. สารมฺมณา เหตู สมฺปยุตฺตกานํ ขนฺธานํ เหตุปจฺจเยน ปจฺจโย. ปฏิสนฺธิ\-กฺขเณ ฯเปฯ. (1)}

\prose{สารมฺมโณ ธมฺโม อนารมฺมณสฺส ธมฺมสฺส เหตุปจฺจเยน ปจฺจโย. สารมฺมณา เหตู จิตฺต\-สมุฏฺฐานานํ รูปานํ เหตุปจฺจเยน ปจฺจโย. ปฏิสนฺธิ\-กฺขเณ ฯเปฯ. (2)}

\prose{สารมฺมโณ ธมฺโม สารมฺมณสฺส จ อนารมฺมณสฺส จ ธมฺมสฺส เหตุปจฺจเยน ปจฺจโย. สารมฺมณา เหตู สมฺปยุตฺตกานํ ขนฺธานํ จิตฺตสมุฏฺฐานานญฺจ รูปานํ เหตุปจฺจเยน ปจฺจโย. ปฏิสนฺธิ\-กฺขเณ. (3)}

\proseitem{28}{สารมฺมโณ ธมฺโม สารมฺมณสฺส ธมฺมสฺส อารมฺมณ\-ปจฺจเยน ปจฺจโย. ทานํ ฯเปฯ สีลํ ฯเปฯ อุโปสถกมฺมํ กตฺวา ตํ ปจฺจเวกฺขติ, อสฺสาเทติ อภินนฺทติ, ตํ อารพฺภ ราโค อุปฺปชฺชติ ฯเปฯ โทมนสฺสํ อุปฺปชฺชติ. ปุพฺเพ สุจิณฺณานิ ฯเปฯ. ฌานา ฯเปฯ. อริยา มคฺคา วุฏฺฐหิตฺวา มคฺคํ ปจฺจเวกฺขนฺติ, ผลํ ปจฺจเวกฺขนฺติ. ปหีเน กิเลเส ฯเปฯ. วิกฺขมฺภิเต กิเลเส ฯเปฯ. ปุพฺเพ สมุทาจิณฺเณ กิเสเล ชานนฺติ. สารมฺมเณ ขนฺเธ อนิจฺจโต ฯเปฯ โทมนสฺสํ อุปฺปชฺชติ. เจโตปริย\-ญาเณน สารมฺมณ\-จิตฺต\-สมงฺคิสฺส จิตฺตํ ชานาติ. อากาสา\-นญฺจา\-ยตนํ วิญฺญาณญฺจา\-ยตนสฺส ฯเปฯ. อากิญฺจญฺญา\-ยตนํ เนว\-สญฺญา\-นา\-สญฺญา\-ยตนสฺส ฯเปฯ. สารมฺมณา ขนฺธา อิทฺธิวิธ\-ญาณสฺส, เจโตปริย\-ญาณสฺส, ปุพฺเพ\-นิวาสา\-นุสฺสติ\-ญาณสฺส, ยถากมฺมูปคญาณสฺส, อนาคตํสญาณสฺส, อาวชฺชนาย อารมฺมณ\-ปจฺจเยน ปจฺจโย. (1)}

\proseitem{29}{อนารมฺมโณ ธมฺโม สารมฺมณสฺส ธมฺมสฺส อารมฺมณ\-ปจฺจเยน ปจฺจโย. อริยา นิพฺพานํ ปจฺจเวกฺขนฺติ. นิพฺพานํ โคตฺรภุสฺส, โวทานสฺส, \csromanfolio{378}มคฺคสฺส, ผลสฺส, อาวชฺชนาย อารมฺมณ\-ปจฺจเยน ปจฺจโย. จกฺขุํ ฯเปฯ วตฺถุํ อนิจฺจโต ฯเปฯ โทมนสฺสํ อุปฺปชฺชติ. ทิพฺเพน จกฺขุนา รูปํ ปสฺสติ. ทิพฺพาย โสตธาตุยา สทฺทํ สุณาติ. รูปายตนํ จกฺขุ\-วิญฺญาณสฺส ฯเปฯ. โผฏฺฐพฺพา\-ยตนํ กายวิญฺญาณสฺส ฯเปฯ. อนารมฺมณา ขนฺธา อิทฺธิวิธ\-ญาณสฺส, ปุพฺเพ\-นิวาสา\-นุสฺสติ\-ญาณสฺส, อนาคตํสญาณสฺส, อาวชฺชนาย อารมฺมณ\-ปจฺจเยน ปจฺจโย. (1)}

\proseitem{30}{สารมฺมโณ ธมฺโม สารมฺมณสฺส ธมฺมสฺส อธิปติ\-ปจฺจเยน ปจฺจโย. อารมฺมณา\-ธิปติ, สหชาตา\-ธิปติ.  อารมฺมณา\-ธิปติ\csromandash{}ทานํ ฯเปฯ สีลํ ฯเปฯ อุโปสถกมฺมํ กตฺวา ตํ ครุํ กตฺวา ปจฺจเวกฺขติ, อสฺสาเทติ อภินนฺทติ, ตํ ครุํ กตฺวา ราโค อุปฺปชฺชติ. ทิฏฺฐิ อุปฺปชฺชติ. ปุพฺเพ สุจิณฺณานิ ฯเปฯ. ฌานา วุฏฺฐหิตฺวา ฌานํ ฯเปฯ. อริยา มคฺคา วุฏฺฐหิตฺวา มคฺคํ ครุํ ฯเปฯ. ผลํ ครุํ ฯเปฯ. สารมฺมเณ ขนฺเธ ครุํ กตฺวา อสฺสาเทติ อภินนฺทติ, ตํ ครุํ กตฺวา ราโค อุปฺปชฺชติ. ทิฏฺฐิ อุปฺปชฺชติ. สหชาตา\-ธิปติ\csromandash{}สารมฺมณา\-ธิปติ สมฺปยุตฺตกานํ ขนฺธานํ อธิปติ\-ปจฺจเยน ปจฺจโย. (1)}

\prose{สารมฺมโณ ธมฺโม อนารมฺมณสฺส ธมฺมสฺส อธิปติ\-ปจฺจเยน ปจฺจโย. สหชาตา\-ธิปติ\csromandash{}สารมฺมณา\-ธิปติ จิตฺต\-สมุฏฺฐานานํ รูปานํ อธิปติ\-ปจฺจเยน ปจฺจโย. (2)}

\prose{สารมฺมโณ ธมฺโม สารมฺมณสฺส จ อนารมฺมณสฺส จ ธมฺมสฺส อธิปติ\-ปจฺจเยน ปจฺจโย. สหชาตา\-ธิปติ\csromandash{}สารมฺมณา\-ธิปติ สมฺปยุตฺตกานํ ขนฺธานํ จิตฺตสมุฏฺฐานานญฺจ รูปานํ อธิปติ\-ปจฺจเยน ปจฺจโย. (3)}

\proseitem{31}{อนารมฺมโณ ธมฺโม สารมฺมณสฺส ธมฺมสฺส อธิปติ\-ปจฺจเยน ปจฺจโย. อารมฺมณา\-ธิปติ\csromandash{}อริยา นิพฺพานํ ครุํ กตฺวา ปจฺจเวกฺขนฺติ. นิพฺพานํ โคตฺรภุสฺส, โวทานสฺส, มคฺคสฺส, ผลสฺส อธิปติ\-ปจฺจเยน ปจฺจโย. จกฺขุํ ฯเปฯ วตฺถุํ ครุํ กตฺวา อสฺสาเทติ อภินนฺทติ, ตํ ครุํ กตฺวา ราโค อุปฺปชฺชติ. ทิฏฺฐิ อุปฺปชฺชติ. (1)}

\csromanheader{2}{55. สารมฺมณทุก \textbf{อนนฺตราทิ}}
\proseitem{32}{\csromanfolio{379}สารมฺมโณ ธมฺโม สารมฺมณสฺส ธมฺมสฺส อนนฺตร\-ปจฺจเยน ปจฺจโย. ปุริมา ปุริมา สารมฺมณา ฯเปฯ ผลสมาปตฺติยา อนนฺตร\-ปจฺจเยน ปจฺจโย. สมนนฺตร\-ปจฺจเยน ปจฺจโย. เอกํ. สหชาต\-ปจฺจเยน ปจฺจโย. สตฺต. (ปฏิจฺจวาเร สหชาตสทิสา.) อญฺญมญฺญ\-ปจฺจเยน ปจฺจโย. ฉ. (ปฏิจฺจวาร อญฺญมญฺญ\-สทิสํ.) นิสฺสย\-ปจฺจเยน ปจฺจโย. สตฺต. (ปจฺจยวาเร นิสฺสยสทิโส.)}

\csromanheader{2}{\textbf{อุปนิสฺสย}}
\proseitem{33}{สารมฺมโณ ธมฺโม สารมฺมณสฺส ธมฺมสฺส อุปนิสฺสย\-ปจฺจเยน ปจฺจโย. อารมฺมณู\-ปนิสฺสโย, อนนฺตรู\-ปนิสฺสโย, ปกตู\-ปนิสฺสโย ฯเปฯ. ปกตู\-ปนิสฺสโย\csromandash{}สทฺธํ อุปนิสฺสาย ทานํ เทติ ฯเปฯ. มานํ ชปฺเปติ. ทิฏฺฐิํ คณฺหาติ. สีลํ ฯเปฯ. ปญฺญํ. ราคํ ฯเปฯ. ปตฺถนํ. กายิกํ สุขํ. กายิกํ ทุกฺขํ อุปนิสฺสาย ทานํ เทติ ฯเปฯ สมาปตฺติํ อุปฺปาเทติ. ปาณํ หนติ ฯเปฯ สํฆํ ภินฺทติ. สทฺธา ฯเปฯ. ปญฺญา. ราโค ฯเปฯ. ปตฺถนา. กายิกํ สุขํ. กายิกํ ทุกฺข สทฺธาย ฯเปฯ ปญฺญาย. ราคสฺส ฯเปฯ ปตฺถนาย. กายิกสฺส สุขสฺส. กายิกสฺส ทุกฺขสฺส. มคฺคสฺส, ผลสมาปตฺติยา อุปนิสฺสย\-ปจฺจเยน ปจฺจโย. (1)}

\prose{อนารมฺมโณ ธมฺโม สารมฺมณสฺส ธมฺมสฺส อุปนิสฺสย\-ปจฺจเยน ปจฺจโย. อารมฺมณู\-ปนิสฺสโย, ปกตู\-ปนิสฺสโย ฯเปฯ. ปกตู\-ปนิสฺสโย\csromandash{} อุตุํ. โภชนํ. เสนาสนํ อุปนิสฺสาย ทานํ เทติ ฯเปฯ สมาปตฺติํ อุปฺปาเทติ. ปาณํ หนติ ฯเปฯ สํฆํ ภินฺทติ. อุตุ. โภชนํ. เสนาสนํ สทฺธาย ฯเปฯ ปตฺถนาย. กายิกสฺส สุขสฺส. กายิกสฺส ทุกฺขสฺส. มคฺคสฺส, ผลสมาปตฺติยา อุปนิสฺสย\-ปจฺจเยน ปจฺจโย. (1)}

\proseitem{34}{อนารมฺมโณ ธมฺโม สารมฺมณสฺส ธมฺมสฺส ปุเรชาต\-ปจฺจเยน ปจฺจโย. อารมฺมณ\-ปุเรชาตํ, วตฺถุ\-ปุเรชาตํ.  อารมฺมณ\-ปุเรชาตํ\csromandash{}จกฺขุํ ฯเปฯ วตฺถุํ อนิจฺจโต ฯเปฯ โทมนสฺสํ อุปฺปชฺชติ. ทิพฺเพน จกฺขุนา รูปํ \csromanfolio{380}ปสฺสติ. ทิพฺพาย โสตธาตุยา สทฺทํ สุณาติ. รูปายตนํ จกฺขุ\-วิญฺญาณสฺส ฯเปฯ. โผฏฺฐพฺพา\-ยตนํ กายวิญฺญาณสฺส ฯเปฯ. วตฺถุ\-ปุเรชาตํ\csromandash{}จกฺขายตนํ จกฺขุ\-วิญฺญาณสฺส ฯเปฯ กายายตนํ กายวิญฺญาณสฺส ฯเปฯ. วตฺถุ สารมฺมณานํ ขนฺธานํ ปุเรชาต\-ปจฺจเยน ปจฺจโย. (1)}

\proseitem{35}{สารมฺมโณ ธมฺโม อนารมฺมณสฺส ธมฺมสฺส ปจฺฉาชาต\-ปจฺจเยน ปจฺจโย. เอกํ.}

\prose{สารมฺมโณ ธมฺโม สารมฺมณสฺส ธมฺมสฺส อาเสวน\-ปจฺจเยน ปจฺจโย. เอกํ.}

\proseitem{36}{สารมฺมโณ ธมฺโม สารมฺมณสฺส ธมฺมสฺส กมฺมปจฺจเยน ปจฺจโย. สหชาตา, นานากฺขณิกา.  สหชาตา สารมฺมณา เจตนา สมฺปยุตฺตกานํ ขนฺธานํ กมฺมปจฺจเยน ปจฺจโย. ปฏิสนฺธิ\-กฺขเณ ฯเปฯ. นานากฺขณิกา สารมฺมณา เจตนา วิปากานํ ขนฺธานํ กมฺมปจฺจเยน ปจฺจโย. (1)}

\prose{สารมฺมโณ ธมฺโม อนารมฺมณสฺส ธมฺมสฺส กมฺมปจฺจเยน ปจฺจโย. สหชาตา, นานากฺขณิกา.  สหชาตา สารมฺมณา เจตนา จิตฺต\-สมุฏฺฐานานํ รูปานํ กมฺมปจฺจเยน ปจฺจโย. ปฏิสนฺธิ\-กฺขเณ ฯเปฯ. นานากฺขณิกา สารมฺมณา เจตนา กฏตฺตารูปานํ กมฺมปจฺจเยน ปจฺจโย. (2)}

\prose{สารมฺมโณ ธมฺโม สารมฺมณสฺส จ อนารมฺมณสฺส จ ธมฺมสฺส กมฺมปจฺจเยน ปจฺจโย. สหชาตา, นานากฺขณิกา.  สหชาตา สารมฺมณา เจตนา สมฺปยุตฺตกานํ ขนฺธานํ จิตฺตสมุฏฺฐานานญฺจ รูปานํ กมฺมปจฺจเยน ปจฺจโย. ปฏิสนฺธิ\-กฺขเณ ฯเปฯ. นานากฺขณิกา สารมฺมณา เจตนา วิปากานํ ขนฺธานํ กฏตฺตา จ รูปานํ กมฺมปจฺจเยน ปจฺจโย. (3)}

\csromanheader{2}{\textbf{วิปาก อาหาร}}
\vspace{\tipitakaparskip}
\csromancenter{สารมฺมโณ ธมฺโม สารมฺมณสฺส ธมฺมสฺส วิปาก\-ปจฺจเยน ปจฺจโย. ตีณิ.}
\prose{\csromanfolio{381}สารมฺมโณ ธมฺโม สารมฺมณสฺส ธมฺมสฺส อาหารปจฺจเยน ปจฺจโย. ตีณิ.}

\prose{อนารมฺมโณ ธมฺโม อนารมฺมณสฺส ธมฺมสฺส อาหารปจฺจเยน ปจฺจโย. กพฬีกาโร อาหาโร อิมสฺส กายสฺส อาหารปจฺจเยน ปจฺจโย. (1)}

\csromanheader{2}{\textbf{อินฺทฺริย}}
\proseitem{38}{สารมฺมโณ ธมฺโม สารมฺมณสฺส ธมฺมสฺส อินฺทฺริย\-ปจฺจเยน ปจฺจโย. ตีณิ.}

\prose{อนารมฺมโณ ธมฺโม อนารมฺมณสฺส ธมฺมสฺส อินฺทฺริย\-ปจฺจเยน ปจฺจโย. รูปชีวิตินฺทฺริยํ กฏตฺตารูปานํ อินฺทฺริย\-ปจฺจเยน ปจฺจโย. (1)}

\prose{อนารมฺมโณ ธมฺโม สารมฺมณสฺส ธมฺมสฺส อินฺทฺริย\-ปจฺจเยน ปจฺจโย. จกฺขุนฺทฺริยํ จกฺขุ\-วิญฺญาณสฺส ฯเปฯ. กายินฺทฺริยํ กายวิญฺญาณสฺส อินฺทฺริย\-ปจฺจเยน ปจฺจโย. (1)}

\prose{สารมฺมโณ จ อนารมฺมโณ จ ธมฺมา สารมฺมณสฺส ธมฺมสฺส อินฺทฺริย\-ปจฺจเยน ปจฺจโย. จกฺขุนฺทฺริยญฺจ จกฺขุวิญฺญาณญฺจ จกฺขุ\-วิญฺญาณสหคตานํ ขนฺธานํ อินฺทฺริย\-ปจฺจเยน ปจฺจโย ฯเปฯ. กายินฺทฺริยญฺจ กายวิญฺญาณญฺจ กาย\-วิญฺญาณสหคตานํ ขนฺธานํ อินฺทฺริย\-ปจฺจเยน ปจฺจโย. (1)}

\csromanheader{2}{\textbf{ฌานาทิ}}
\proseitem{39}{สารมฺมโณ ธมฺโม สารมฺมณสฺส ธมฺมสฺส ฌานปจฺจเยน ปจฺจโย. ตีณิ. มคฺคปจฺจเยน ปจฺจโย. ตีณิ. สมฺปยุตฺต\-ปจฺจเยน ปจฺจโย. เอกํ.}

\proseitem{40}{สารมฺมโณ ธมฺโม อนารมฺมณสฺส ธมฺมสฺส วิปฺปยุตฺต\-ปจฺจเยน ปจฺจโย. สหชาตํ, ปจฺฉาชาตํ. (สํขิตฺตํ.) (1)}

\prose{อนารมฺมโณ ธมฺโม สารมฺมณสฺส ธมฺมสฺส วิปฺปยุตฺต\-ปจฺจเยน ปจฺจโย. สหชาตํ, ปุเรชาตํ.  สหชาตํ ปฏิสนฺธิ\-กฺขเณ วตฺถุ สารมฺมณานํ \csromanfolio{382}ขนฺธานํ วิปฺปยุตฺต\-ปจฺจเยน ปจฺจโย. ปุเรชาตํ จกฺขายตนํ จกฺขุ\-วิญฺญาณสฺส ฯเปฯ กายายตนํ กายวิญฺญาณสฺส วิปฺปยุตฺต\-ปจฺจเยน ปจฺจโย. วตฺถุ สารมฺมณานํ ขนฺธานํ วิปฺปยุตฺต\-ปจฺจเยน ปจฺจโย. (1)}

\proseitem{41}{สารมฺมโณ ธมฺโม สารมฺมณสฺส ธมฺมสฺส อตฺถิปจฺจเยน ปจฺจโย. สารมฺมโณ เอโก ขนฺโธ ติณฺณนฺนํ ขนฺธานํ อตฺถิปจฺจเยน ปจฺจโย ฯเปฯ เทฺว ขนฺธา ฯเปฯ. ปฏิสนฺธิ\-กฺขเณ ฯเปฯ. (1)}

\prose{สารมฺมโณ ธมฺโม อนารมฺมณสฺส ธมฺมสฺส อตฺถิปจฺจเยน ปจฺจโย. สหชาตํ, ปจฺฉาชาตํ. (สํขิตฺตํ.) (2)}

\prose{สารมฺมโณ ธมฺโม สารมฺมณสฺส จ อนารมฺมณสฺส จ ธมฺมสฺส อตฺถิปจฺจเยน ปจฺจโย. (ปฏิจฺจสทิสํ.) (3)}

\proseitem{42}{อนารมฺมโณ ธมฺโม อนารมฺมณสฺส ธมฺมสฺส อตฺถิปจฺจเยน ปจฺจโย. \textbf{สหชาตํ}, อาหารํ, อินฺทฺริยํ.  สหชาตํ เอกํ มหาภูตํ ฯเปฯ. (ยาว อสญฺญสตฺตา.) (1)}

\prose{อนารมฺมโณ ธมฺโม สารมฺมณสฺส ธมฺมสฺส อตฺถิปจฺจเยน ปจฺจโย. สหชาตํ, ปุเรชาตํ.  สหชาตํ ปฏิสนฺธิ\-กฺขเณ วตฺถุ สารมฺมณานํ ขนฺธานํ อตฺถิปจฺจเยน ปจฺจโย. ปุเรชาตํ จกฺขุํ ฯเปฯ วตฺถุํ อนิจฺจโต ฯเปฯ โทมนสฺสํ อุปฺปชฺชติ. ทิพฺเพน จกฺขุนา รูปํ ปสฺสติ. ทิพฺพาย โสตธาตุยา สทฺทํ สุณาติ. รูปายตนํ จกฺขุ\-วิญฺญาณสฺส ฯเปฯ. โผฏฺฐพฺพา\-ยตนํ กายวิญฺญาณสฺส ฯเปฯ. จกฺขายตนํ จกฺขุ\-วิญฺญาณสฺส ฯเปฯ. กายายตนํ กายวิญฺญาณสฺส ฯเปฯ. วตฺถุ สารมฺมณานํ ขนฺธานํ อตฺถิปจฺจเยน ปจฺจโย. (2)}

\proseitem{43}{สารมฺมโณ จ อนารมฺมโณ จ ธมฺมา สารมฺมณสฺส ธมฺมสฺส อตฺถิปจฺจเยน ปจฺจโย. สหชาตํ, ปุเรชาตํ.  สหชาโต จกฺขุ\-วิญฺญาณ\-สหคโต เอโก ขนฺโธ จ จกฺขายตนญฺจ ติณฺณนฺนํ ขนฺธานํ อตฺถิปจฺจเยน ปจฺจโย ฯเปฯ เทฺว ขนฺธา จ ฯเปฯ. กาย\-วิญฺญาณ\-สหคโต เอโก ขนฺโธ จ กายายตนญฺจ ติณฺณนฺนํ ขนฺธานํ อตฺถิปจฺจเยน ปจฺจโย ฯเปฯ เทฺว ขนฺธา จ ฯเปฯ. สารมฺมโณ เอโก ขนฺโธ จ วตฺถุ จ ติณฺณนฺนํ ขนฺธานํ อตฺถิปจฺจเยน ปจฺจโย ฯเปฯ เทฺว ขนฺธา จ ฯเปฯ. ปฏิสนฺธิ\-กฺขเณ ฯเปฯ. (1)}

\prose{\csromanfolio{383}สารมฺมโณ จ อนารมฺมโณ จ ธมฺมา อนารมฺมณสฺส ธมฺมสฺส อตฺถิปจฺจเยน ปจฺจโย. สหชาตํ, ปจฺฉาชาตํ, อาหารํ, อินฺทฺริยํ.  สหชาตา สารมฺมณา ขนฺธา จ มหาภูตา จ จิตฺต\-สมุฏฺฐานานํ รูปานํ อตฺถิปจฺจเยน ปจฺจโย. ปฏิสนฺธิ\-กฺขเณ ฯเปฯ. ปจฺฉาชาตา สารมฺมณา ขนฺธา จ กพฬีกาโร อาหาโร จ อิมสฺส กายสฺส อตฺถิปจฺจเยน ปจฺจโย. ปจฺฉาชาตา สารมฺมณา ขนฺธา จ รูปชีวิตินฺทฺริยญฺจ กฏตฺตารูปานํ อตฺถิปจฺจเยน ปจฺจโย. (2)}

\csromanmarkh{1. ปจฺจยานุโลม}\csromanmarkhii{2. สงฺขฺยาวาร}\csromanheadernomark{2}{1. \textbf{ปจฺจยานุโลม 2. สงฺขฺยาวาร}}
\proseitem{44}{เหตุยา ตีณิ, อารมฺมเณ เทฺว, อธิปติยา จตฺตาริ, อนนฺตเร เอกํ, สมนนฺตเร เอกํ, สหชาเต สตฺต, อญฺญมญฺเญ ฉ, นิสฺสเย สตฺต, อุปนิสฺสเย เทฺว, ปุเรชาเต เอกํ, ปจฺฉาชาเต เอกํ, อาเสวเน เอกํ, กมฺเม ตีณิ, วิปาเก ตีณิ, อาหาเร จตฺตาริ, อินฺทฺริเย ฉ, ฌาเน ตีณิ, มคฺเค ตีณิ, สมฺปยุตฺเต เอกํ, วิปฺปยุตฺเต เทฺว, อตฺถิยา สตฺต, นตฺถิยา เอกํ, วิคเต เอกํ, อวิคเต สตฺต.}

\proseitem{45}{สารมฺมโณ ธมฺโม สารมฺมณสฺส ธมฺมสฺส อารมฺมณ\-ปจฺจเยน ปจฺจโย. สหชาต\-ปจฺจเยน ปจฺจโย. อุปนิสฺสย\-ปจฺจเยน ปจฺจโย. กมฺมปจฺจเยน ปจฺจโย. (1)}

\prose{สารมฺมโณ ธมฺโม อนารมฺมณสฺส ธมฺมสฺส สหชาต\-ปจฺจเยน ปจฺจโย. ปจฺฉาชาต\-ปจฺจเยน ปจฺจโย. กมฺมปจฺจเยน ปจฺจโย. (2)}

\prose{สารมฺมโณ ธมฺโม สารมฺมณสฺส จ อนารมฺมณสฺส จ ธมฺมสฺส สหชาต\-ปจฺจเยน ปจฺจโย. กมฺมปจฺจเยน ปจฺจโย. (3)}

\proseitem{46}{อนารมฺมโณ ธมฺโม อนารมฺมณสฺส ธมฺมสฺส สหชาต\-ปจฺจเยน ปจฺจโย. อาหารปจฺจเยน ปจฺจโย. อินฺทฺริย\-ปจฺจเยน ปจฺจโย. (1)}

\prose{\csromanfolio{384}อนารมฺมโณ ธมฺโม สารมฺมณสฺส ธมฺมสฺส อารมฺมณ\-ปจฺจเยน ปจฺจโย. สหชาต\-ปจฺจเยน ปจฺจโย. อุปนิสฺสย\-ปจฺจเยน ปจฺจโย. ปุเรชาต\-ปจฺจเยน ปจฺจโย. (2)}

\prose{สารมฺมโณ จ อนารมฺมโณ จ ธมฺมา สารมฺมณสฺส ธมฺมสฺส สหชาตํ, ปุเรชาตํ. (1)}

\prose{สารมฺมโณ จ อนารมฺมโณ จ ธมฺมา อนารมฺมณสฺส ธมฺมสฺส สหชาตํ, ปจฺฉาชาตํ, อาหารํ, อินฺทฺริยํ. (2)}

\csromanmarkh{2. ปจฺจย\-ปจฺจนีย}\csromanmarkhii{2. สงฺขฺยาวาร}\csromanheadernomark{2}{2. ปจฺจย\-ปจฺจนีย 2. สงฺขฺยาวาร}
\proseitem{47}{นเหตุยา สตฺต, นอารมฺมเณ สตฺต ฯเปฯ นสมนนฺตเร สตฺต, นสหชาเต ฉ, นอญฺญมญฺเญ ฉ, นนิสฺสเย ฉ, นอุปนิสฺสเย สตฺต, นปุเรชาเต สตฺต, นปจฺฉาชาเต สตฺต ฯเปฯ นมคฺเค สตฺต, นสมฺปยุตฺเต ฉ, นวิปฺปยุตฺเต ปญฺจ, โนอตฺถิยา จตฺตาริ, โนนตฺถิยา สตฺต, โนวิคเต สตฺต, โนอวิคเต จตฺตาริ.}

\proseitem{48}{เหตุปจฺจยา นอารมฺมเณ ตีณิ, นอธิปติยา ฯเปฯ นสมนนฺตเร ตีณิ, นอญฺญมญฺเญ เอกํ, นอุปนิสฺสเย ตีณิ ฯเปฯ นมคฺเค ตีณิ, นสมฺปยุตฺเต เอกํ, นวิปฺปยุตฺเต เอกํ, โนนตฺถิยา ตีณิ, โนวิคเต ตีณิ.}

\proseitemwithcloser{49}{นเหตุปจฺจยา อารมฺมเณ เทฺว, อธิปติยา จตฺตาริ, อนนฺตเร เอกํ, (อนุโลมมาติกา กาตพฺพา.) ฯเปฯ อวิคเต สตฺต.}{สารมฺมณทุกํ นิฏฺฐิตํ.}
\csromanmatikaanchor{mtk.55}
\csromantocmarkref{subsection}{56. จิตฺตทุก}{mtk.55}
\bookmark[dest={mtk.55},level=2]{56. จิตฺตทุก}
\csromanmarkh{56. จิตฺตทุก}\csromanmarkhii{56. จิตฺตทุก 1. ปฏิจฺจวาร}\csromanheadernomark{2}{56. จิตฺตทุก \textbf{56. จิตฺตทุก 1. ปฏิจฺจวาร}}
\csromanmarkh{1. ปจฺจยานุโลม}\csromanmarkhii{1. วิภงฺควาร}\csromanheadernomark{2}{1. \textbf{ปจฺจยานุโลม 1. วิภงฺควาร}}
\csromanheader{2}{\textbf{เหตุ}}
\proseitem{50}{\csromanfolio{385}จิตฺตํ ธมฺมํ ปฏิจฺจ โนจิตฺโต ธมฺโม อุปฺปชฺชติ เหตุปจฺจยา. จิตฺตํ ปฏิจฺจ สมฺปยุตฺตกา ขนฺธา จิตฺต\-สมุฏฺฐานญฺจ รูปํ. ปฏิสนฺธิ\-กฺขเณ จิตฺตํ ปฏิจฺจ สมฺปยุตฺตกา ขนฺธา กฏตฺตา จ รูปํ. (1)}

\prose{โนจิตฺตํ ธมฺมํ ปฏิจฺจ โนจิตฺโต ธมฺโม อุปฺปชฺชติ เหตุปจฺจยา. โนจิตฺตํ เอกํ ขนฺธํ ปฏิจฺจ เทฺว ขนฺธา จิตฺต\-สมุฏฺฐานญฺจ รูปํ. เทฺว ขนฺเธ ปฏิจฺจ เอโก ขนฺโธ จิตฺต\-สมุฏฺฐานญฺจ รูปํ. ปฏิสนฺธิ\-กฺขเณ ฯเปฯ. ขนฺเธ ปฏิจฺจ วตฺถุ, วตฺถุํ ปฏิจฺจ ขนฺธา. เอกํ มหาภูตํ ปฏิจฺจ ฯเปฯ. (1)}

\prose{โนจิตฺตํ ธมฺมํ ปฏิจฺจ จิตฺโต ธมฺโม อุปฺปชฺชติ เหตุปจฺจยา. โนจิตฺเต ขนฺเธ ปฏิจฺจ จิตฺตํ. ปฏิสนฺธิ\-กฺขเณ โนจิตฺเต ขนฺเธ ปฏิจฺจ จิตฺตํ. ปฏิสนฺธิ\-กฺขเณ วตฺถุํ ปฏิจฺจ จิตฺตํ. (2)}

\prose{โนจิตฺตํ ธมฺมํ ปฏิจฺจ จิตฺโต จ โนจิตฺโต จ ธมฺมา อุปฺปชฺชนฺติ เหตุปจฺจยา. โนจิตฺตํ เอกํ ขนฺธํ ปฏิจฺจ เทฺว ขนฺธา จิตฺตญฺจ จิตฺต\-สมุฏฺฐานญฺจ รูปํ. เทฺว ขนฺเธ ฯเปฯ. ปฏิสนฺธิ\-กฺขเณ โนจิตฺตํ เอกํ ขนฺธํ ปฏิจฺจ เทฺว ขนฺธา จิตฺตญฺจ กฏตฺตา จ รูปํ. เทฺว ขนฺเธ ปฏิจฺจ เอโก ขนฺโธ จิตฺตญฺจ กฏตฺตา จ รูปํ. ปฏิสนฺธิ\-กฺขเณ วตฺถุํ ปฏิจฺจ จิตฺตญฺจ สมฺปยุตฺตกา จ ขนฺธา. (3)}

\prose{จิตฺตญฺจ โนจิตฺตญฺจ ธมฺมํ ปฏิจฺจ โนจิตฺโต ธมฺโม อุปฺปชฺชติ เหตุปจฺจยา. โนจิตฺตํ เอกํ ขนฺธญฺจ จิตฺตญฺจ ปฏิจฺจ เทฺว ขนฺธา จิตฺต\-สมุฏฺฐานญฺจ รูปํ. เทฺว ขนฺเธ จ ฯเปฯ. จิตฺตญฺจ มหาภูเต จ ปฏิจฺจ จิตฺต\-สมุฏฺฐานํ รูปํ. ปฏิสนฺธิ\-กฺขเณ โนจิตฺตํ เอกํ ขนฺธญฺจ จิตฺตญฺจ ปฏิจฺจ เทฺว ขนฺธา กฏตฺตา จ รูปํ. เทฺว ขนฺเธ ฯเปฯ. ปฏิสนฺธิ\-กฺขเณ จิตฺตญฺจ วตฺถุญฺจ ปฏิจฺจ โนจิตฺตา ขนฺธา. ปฏิสนฺธิ\-กฺขเณ จิตฺตญฺจ มหาภูเต จ ปฏิจฺจ กฏตฺตารูปํ. (1)}

\proseitem{51}{จิตฺตํ ธมฺมํ ปฏิจฺจ โนจิตฺโต ธมฺโม อุปฺปชฺชติ อารมฺมณ\-ปจฺจยา. จิตฺตํ ปฏิจฺจ สมฺปยุตฺตกา ขนฺธา. ปฏิสนฺธิ\-กฺขเณ ฯเปฯ. (1)}

\prose{\csromanfolio{386}โนจิตฺตํ ธมฺมํ ปฏิจฺจ โนจิตฺโต ธมฺโม อุปฺปชฺชติ อารมฺมณ\-ปจฺจยา. โนจิตฺตํ เอกํ ขนฺธํ ปฏิจฺจ เทฺว ขนฺธา. เทฺว ขนฺเธ ฯเปฯ. ปฏิสนฺธิ\-กฺขเณ ฯเปฯ. ปฏิสนฺธิ\-กฺขเณ วตฺถุํ ปฏิจฺจ ขนฺธา. (1)}

\prose{โนจิตฺตํ ธมฺมํ ปฏิจฺจ จิตฺโต ธมฺโม อุปฺปชฺชติ อารมฺมณ\-ปจฺจยา. โนจิตฺเต ขนฺเธ ปฏิจฺจ จิตฺตํ. ปฏิสนฺธิ\-กฺขเณ โนจิตฺเต ขนฺเธ ปฏิจฺจ จิตฺตํ. ปฏิสนฺธิ\-กฺขเณ วตฺถุํ ปฏิจฺจ จิตฺตํ. (2)}

\prose{โนจิตฺตํ ธมฺมํ ปฏิจฺจ จิตฺโต จ โนจิตฺโต จ ธมฺมา อุปฺปชฺชนฺติ อารมฺมณ\-ปจฺจยา. โนจิตฺตํ เอกํ ขนฺธํ ปฏิจฺจ เทฺว ขนฺธา จิตฺตญฺจ. เทฺว ขนฺเธ ปฏิจฺจ ฯเปฯ. ปฏิสนฺธิ\-กฺขเณ โนจิตฺตํ เอกํ ขนฺธํ ปฏิจฺจ เทฺว ขนฺธา จิตฺตญฺจ. เทฺว ขนฺเธ ฯเปฯ. ปฏิสนฺธิ\-กฺขเณ วตฺถุํ ปฏิจฺจ จิตฺตญฺจ สมฺปยุตฺตกา จ ขนฺธา. (3)}

\prose{จิตฺตญฺจ โนจิตฺตญฺจ ธมฺมํ ปฏิจฺจ โนจิตฺโต ธมฺโม อุปฺปชฺชติ อารมฺมณ\-ปจฺจยา. โนจิตฺตํ เอกํ ขนฺธญฺจ จิตฺตญฺจ ปฏิจฺจ เทฺว ขนฺธา. เทฺว ขนฺเธ จ ฯเปฯ. ปฏิสนฺธิ\-กฺขเณ ฯเปฯ. ปฏิสนฺธิ\-กฺขเณ จิตฺตญฺจ วตฺถุญฺจ ปฏิจฺจ โนจิตฺตา ขนฺธา. (1) (สํขิตฺตํ.)}

\csromanmarkh{1. ปจฺจยานุโลม}\csromanmarkhii{2. สงฺขฺยาวาร}\csromanheadernomark{2}{1. \textbf{ปจฺจยานุโลม 2. สงฺขฺยาวาร}}
\proseitem{52}{เหตุยา ปญฺจ, อารมฺมเณ ปญฺจ, อธิปติยา ปญฺจ, อนนฺตเร ปญฺจ, สมนนฺตเร ปญฺจ, สหชาเต ปญฺจ, อญฺญมญฺเญ ปญฺจ, นิสฺสเย ปญฺจ, อุปนิสฺสเย ปญฺจ, ปุเรชาเต ปญฺจ, อาเสวเน ปญฺจ, กมฺเม ปญฺจ, วิปาเก ปญฺจ, อาหาเร ปญฺจ, อินฺทฺริเย ปญฺจ, ฌาเน ปญฺจ, มคฺเค ปญฺจ, สมฺปยุตฺเต ปญฺจ, วิปฺปยุตฺเต ปญฺจ, อตฺถิยา ปญฺจ, นตฺถิยา ปญฺจ, วิคเต ปญฺจ, อวิคเต ปญฺจ.}

\csromanmarkh{2. ปจฺจย\-ปจฺจนีย}\csromanmarkhii{1. วิภงฺควาร}\csromanheadernomark{2}{2. ปจฺจย\-ปจฺจนีย 1. วิภงฺควาร}
\proseitem{53}{จิตฺตํ ธมฺมํ ปฏิจฺจ โนจิตฺโต ธมฺโม อุปฺปชฺชติ นเหตุปจฺจยา. อเหตุกํ จิตฺตํ ปฏิจฺจ สมฺปยุตฺตกา ขนฺธา จิตฺต\-สมุฏฺฐานญฺจ รูปํ. อเหตุก\-ปฏิสนฺธิ\-กฺขเณ จิตฺตํ ฯเปฯ. วิจิกิจฺฉา\-สหคตํ อุทฺธจฺจ\-สหคตํ จิตฺตํ ปฏิจฺจ วิจิกิจฺฉา\-สหคโต อุทฺธจฺจ\-สหคโต โมโห. (1)}

\prose{\csromanfolio{387}โนจิตฺตํ ธมฺมํ ปฏิจฺจ โนจิตฺโต ธมฺโม อุปฺปชฺชติ นเหตุปจฺจยา. อเหตุกํ โนจิตฺตํ เอกํ ขนฺธํ ปฏิจฺจ เทฺว ขนฺธา จิตฺต\-สมุฏฺฐานญฺจ รูปํ. เทฺว ขนฺเธ ฯเปฯ. อเหตุก\-ปฏิสนฺธิ\-กฺขเณ ฯเปฯ. (ยาว อสญฺญสตฺตา.) วิจิกิจฺฉา\-สหคเต อุทฺธจฺจ\-สหคเต ขนฺเธ ปฏิจฺจ วิจิกิจฺฉา\-สหคโต อุทฺธจฺจ\-สหคโต โมโห. (1)}

\prose{โนจิตฺตํ ธมฺมํ ปฏิจฺจ จิตฺโต ธมฺโม อุปฺปชฺชติ นเหตุปจฺจยา. อเหตุเก โนจิตฺเต ขนฺเธ ปฏิจฺจ จิตฺตํ. อเหตุก\-ปฏิสนฺธิ\-กฺขเณ โนจิตฺเต ขนฺเธ ปฏิจฺจ จิตฺตํ. อเหตุก\-ปฏิสนฺธิ\-กฺขเณ วตฺถุํ ปฏิจฺจ จิตฺตํ.}

\prose{โนจิตฺตํ ธมฺมํ ปฏิจฺจ จิตฺโต จ โนจิตฺโต จ ธมฺมา อุปฺปชฺชนฺติ นเหตุปจฺจยา. อเหตุกํ โนจิตฺตํ เอกํ ขนฺธํ ปฏิจฺจ เทฺว ขนฺธา จิตฺตญฺจ จิตฺต\-สมุฏฺฐานญฺจ รูปํ. เทฺว ขนฺเธ ฯเปฯ. อเหตุก\-ปฏิสนฺธิ\-กฺขเณ ฯเปฯ. อเหตุก\-ปฏิสนฺธิ\-กฺขเณ วตฺถุํ ปฏิจฺจ จิตฺตญฺจ สมฺปยุตฺตกา จ ขนฺธา. (3)}

\prose{จิตฺตญฺจ โนจิตฺตญฺจ ธมฺมํ ปฏิจฺจ โนจิตฺโต ธมฺโม อุปฺปชฺชติ นเหตุปจฺจยา. อเหตุกํ โนจิตฺตํ เอกํ ขนฺธญฺจ จิตฺตญฺจ ปฏิจฺจ เทฺว ขนฺธา จิตฺต\-สมุฏฺฐานญฺจ รูปํ. เทฺว ขนฺเธ จ ฯเปฯ. อเหตุก\-ปฏิสนฺธิ\-กฺขเณ โนจิตฺตํ เอกํ ขนฺธญฺจ จิตฺตญฺจ ปฏิจฺจ เทฺว ขนฺธา กฏตฺตา จ รูปํ. เทฺว ขนฺเธ จ ฯเปฯ. อเหตุก\-ปฏิสนฺธิ\-กฺขเณ จิตฺตญฺจ วตฺถุญฺจ ปฏิจฺจ โนจิตฺตา ขนฺธา. วิจิกิจฺฉา\-สหคตํ อุทฺธจฺจ\-สหคตํ จิตฺตญฺจ สมฺปยุตฺตเก จ ขนฺเธ ปฏิจฺจ วิจิกิจฺฉา\-สหคโต อุทฺธจฺจ\-สหคโต โมโห. (1)}

\prose{นอารมฺมณ}

\proseitem{54}{จิตฺตํ ธมฺมํ ปฏิจฺจ โนจิตฺโต ธมฺโม อุปฺปชฺชติ นอารมฺมณ\-ปจฺจยา. จิตฺตํ ปฏิจฺจ จิตฺต\-สมุฏฺฐานํ รูปํ. ปฏิสนฺธิ\-กฺขเณ ฯเปฯ. (1)}

\prose{โนจิตฺตํ ธมฺมํ ปฏิจฺจ โนจิตฺโต ธมฺโม อุปฺปชฺชติ นอารมฺมณ\-ปจฺจยา. โนจิตฺเต ขนฺเธ ปฏิจฺจ จิตฺต\-สมุฏฺฐานํ รูปํ. ปฏิสนฺธิ\-กฺขเณ ฯเปฯ. (ยาว อสญฺญสตฺตา.) (2)}

\prose{จิตฺตญฺจ โนจิตฺตญฺจ ธมฺมํ ปฏิจฺจ โนจิตฺโต ธมฺโม อุปฺปชฺชติ นอารมฺมณ\-ปจฺจยา. จิตฺตญฺจ สมฺปยุตฺตเก จ ขนฺเธ ปฏิจฺจ จิตฺต\-สมุฏฺฐานํ รูปํ. จิตฺตญฺจ มหาภูเต จ ปฏิจฺจ จิตฺต\-สมุฏฺฐานํ รูปํ. ปฏิสนฺธิ\-กฺขเณ จิตฺตญฺจ สมฺปยุตฺตเก จ ขนฺเธ ปฏิจฺจ กฏตฺตารูปํ. จิตฺตญฺจ มหาภูเต จ ปฏิจฺจ กฏตฺตารูปํ. (1)}

\prose{\csromanfolio{388}นอธิปตฺยาทิ}

\proseitem{55}{จิตฺตํ ธมฺมํ ปฏิจฺจ โนจิตฺโต ธมฺโม อุปฺปชฺชติ นอธิปติ\-ปจฺจยา. ปญฺจ. นอนนฺตร\-ปจฺจยา ฯเปฯ. นอุปนิสฺสยปจฺจยา. ตีณิ.}

\csromanheader{2}{\textbf{นปุเรชาตาทิ}}
\proseitem{56}{จิตฺตํ ธมฺมํ ปฏิจฺจ โนจิตฺโต ธมฺโม อุปฺปชฺชติ นปุเรชาต\-ปจฺจยา. อรูเป จิตฺตํ ปฏิจฺจ สมฺปยุตฺตกา ขนฺธา. จิตฺตํ ปฏิจฺจ จิตฺต\-สมุฏฺฐานํ รูปํ. ปฏิสนฺธิ\-กฺขเณ จิตฺตํ ปฏิจฺจ สมฺปยุตฺตกา ขนฺธา กฏตฺตา จ รูปํ. (1)}

\prose{โนจิตฺตํ ธมฺมํ ปฏิจฺจ โนจิตฺโต ธมฺโม อุปฺปชฺชติ นปุเรชาต\-ปจฺจยา. อรูเป โนจิตฺตํ เอกํ ขนฺธํ ปฏิจฺจ เทฺว ขนฺธา. เทฺว ขนฺเธ ฯเปฯ. โนจิตฺเต ขนฺเธ ปฏิจฺจ จิตฺต\-สมุฏฺฐานํ รูปํ. ปฏิสนฺธิ\-กฺขเณ ฯเปฯ. (ยาว อสญฺญสตฺตา.) (1)}

\prose{โนจิตฺตํ ธมฺมํ ปฏิจฺจ จิตฺโต ธมฺโม อุปฺปชฺชติ นปุเรชาต\-ปจฺจยา. อรูเป โนจิตฺเต ขนฺเธ ปฏิจฺจ จิตฺตํ. ปฏิสนฺธิ\-กฺขเณ. วตฺถุํ ปฏิจฺจ จิตฺตํ. (2)}

\prose{โนจิตฺตํ ธมฺมํ ปฏิจฺจ จิตฺโต จ โนจิตฺโต จ ธมฺมา อุปฺปชฺชนฺติ นปุเรชาต\-ปจฺจยา. อรูเป โนจิตฺตํ เอกํ ขนฺธํ ปฏิจฺจ เทฺว ขนฺธา จิตฺตญฺจ. เทฺว ขนฺเธ ฯเปฯ. ปฏิสนฺธิ\-กฺขเณ วตฺถุํ ปฏิจฺจ จิตฺตญฺจ สมฺปยุตฺตกา จ ขนฺธา. (3)}

\prose{จิตฺตญฺจ โนจิตฺตญฺจ ธมฺมํ ปฏิจฺจ โนจิตฺโต ธมฺโม อุปฺปชฺชติ นปุเรชาต\-ปจฺจยา. อรูเป โนจิตฺตํ เอกํ ขนฺธญฺจ จิตฺตญฺจ ปฏิจฺจ เทฺว ขนฺธา. เทฺว ขนฺเธ จ ฯเปฯ. โนจิตฺเต ขนฺเธ จ จิตฺตญฺจ ปฏิจฺจ จิตฺต\-สมุฏฺฐานํ รูปํ. จิตฺตญฺจ มหาภูเต จ ปฏิจฺจ จิตฺต\-สมุฏฺฐานํ รูปํ. ปฏิสนฺธิ\-กฺขเณ จิตฺตญฺจ วตฺถุญฺจ ปฏิจฺจ โนจิตฺตา ขนฺธา. ปฏิสนฺธิ\-กฺขเณ จิตฺตญฺจ สมฺปยุตฺตเก จ ขนฺเธ ปฏิจฺจ กฏตฺตารูปํ. จิตฺตญฺจ มหาภูเต จ ปฏิจฺจ กฏตฺตารูปํ. (1)}

\prose{นปจฺฉาชาต\-ปจฺจยา. นอาเสวน\-ปจฺจยา.}

\csromanheader{2}{\textbf{นกมฺม}}
\proseitem{57}{จิตฺตํ ธมฺมํ ปฏิจฺจ โนจิตฺโต ธมฺโม อุปฺปชฺชติ นกมฺมปจฺจยา. จิตฺตํ ปฏิจฺจ สมฺปยุตฺตกา เจตนา. (1)}

\prose{\csromanfolio{389}โนจิตฺตํ ธมฺมํ ปฏิจฺจ โนจิตฺโต ธมฺโม อุปฺปชฺชติ นกมฺมปจฺจยา. โนจิตฺเต ขนฺเธ ปฏิจฺจ สมฺปยุตฺตกา เจตนา. พาหิรํ. อาหารสมุฏฺฐานํ. อุตุสมุฏฺฐานํ ฯเปฯ. (1)}

\prose{จิตฺตญฺจ โนจิตฺตญฺจ ธมฺมํ ปฏิจฺจ โนจิตฺโต ธมฺโม อุปฺปชฺชติ นกมฺมปจฺจยา. โนจิตฺเต ขนฺเธ จ จิตฺตญฺจ ปฏิจฺจ สมฺปยุตฺตกา เจตนา. (สํขิตฺตํ.) (1)}

\csromanmarkh{2. ปจฺจย\-ปจฺจนีย}\csromanmarkhii{2. สงฺขฺยาวาร}\csromanheadernomark{2}{2. ปจฺจย\-ปจฺจนีย 2. สงฺขฺยาวาร}
\proseitem{58}{นเหตุยา ปญฺจ, นอารมฺมเณ ตีณิ, นอธิปติยา ปญฺจ, นอนนฺตเร ตีณิ, นสมนนฺตเร ตีณิ, นอญฺญมญฺเญ ตีณิ, นอุปนิสฺสเย ตีณิ, นปุเรชาเต ปญฺจ, นปจฺฉาชาเต ปญฺจ, นอาเสวเน ปญฺจ, นกมฺเม ตีณิ, นวิปาเก ปญฺจ, นอาหาเร เอกํ, นอินฺทฺริเย เอกํ, นฌาเน ปญฺจ, นมคฺเค ปญฺจ, นสมฺปยุตฺเต ตีณิ, นวิปฺปยุตฺเต ปญฺจ, โนนตฺถิยา ตีณิ, โนวิคเต ตีณิ.}

\vspace{\tipitakaparskip}
\csromancenter{\textbf{เหตุ}ปจฺจยา นอารมฺมเณ ตีณิ, นอธิปติยา ปญฺจ. (สํขิตฺตํ.)}
\proseitem{60}{นเหตุปจฺจยา อารมฺมเณ ปญฺจ, อนนฺตเร ปญฺจ, (สพฺพตฺถ ปญฺจ,) มคฺเค ตีณิ ฯเปฯ อวิคเต ปญฺจ.}

\vspace{\tipitakaparskip}
\csromancenter{(สหชาตวาโร ปฏิจฺจ\-วาร\-สทิโส.)}
\csromanmarkh{56. จิตฺตทุก}\csromanmarkhii{3. ปจฺจยวาร}\csromanheadernomark{2}{56. จิตฺตทุก 3. ปจฺจยวาร}
\csromanmarkh{1. ปจฺจยานุโลม}\csromanmarkhii{1. วิภงฺควาร}\csromanheadernomark{2}{1. \textbf{ปจฺจยานุโลม 1. วิภงฺควาร}}
\vspace{\tipitakaparskip}
\csromancenter{\textbf{เหตุ}}
\proseitem{61}{จิตฺตํ ธมฺมํ ปจฺจยา โนจิตฺโต ธมฺโม อุปฺปชฺชติ เหตุปจฺจยา. จิตฺตํ ปจฺจยา สมฺปยุตฺตกา ขนฺธา จิตฺต\-สมุฏฺฐานญฺจ รูปํ. ปฏิสนฺธิ\-กฺขเณ ฯเปฯ. (1)}

\prose{\csromanfolio{390}โนจิตฺตํ ธมฺมํ ปจฺจยา โนจิตฺโต ธมฺโม อุปฺปชฺชติ เหตุปจฺจยา. โนจิตฺตํ เอกํ ขนฺธํ ปจฺจยา เทฺว ขนฺธา จิตฺต\-สมุฏฺฐานญฺจ รูปํ. เทฺว ขนฺเธ ฯเปฯ. ปฏิสนฺธิ\-กฺขเณ ฯเปฯ. (ยาว มหาภูตา.) วตฺถุํ ปจฺจยา โนจิตฺตา ขนฺธา. (1)}

\prose{โนจิตฺตํ ธมฺมํ ปจฺจยา จิตฺโต ธมฺโม อุปฺปชฺชติ เหตุปจฺจยา. โนจิตฺเต ขนฺเธ ปจฺจยา จิตฺตํ. วตฺถุํ ปจฺจยา จิตฺตํ. ปฏิสนฺธิ\-กฺขเณ โนจิตฺเต ขนฺเธ ปจฺจยา จิตฺตํ. ปฏิสนฺธิ\-กฺขเณ วตฺถุํ ปจฺจยา จิตฺตํ.}

\prose{โนจิตฺตํ ธมฺมํ ปจฺจยา จิตฺโต จ โนจิตฺโต จ ธมฺมา อุปฺปชฺชนฺติ เหตุปจฺจยา. โนจิตฺตํ เอกํ ขนฺธํ ปจฺจยา เทฺว ขนฺธา จิตฺตญฺจ จิตฺต\-สมุฏฺฐานญฺจ รูปํ. เทฺว ขนฺเธ ฯเปฯ. วตฺถุํ ปจฺจยา จิตฺตญฺจ สมฺปยุตฺตกา จ ขนฺธา. ปฏิสนฺธิ\-กฺขเณ โนจิตฺตํ เอกํ ขนฺธํ ปจฺจยา เทฺว ขนฺธา จิตฺตญฺจ กฏตฺตา จ รูปํ. เทฺว ขนฺเธ ฯเปฯ. ปฏิสนฺธิ\-กฺขเณ วตฺถุํ ปจฺจยา จิตฺตญฺจ สมฺปยุตฺตกา จ ขนฺธา. (3)}

\prose{จิตฺตญฺจ โนจิตฺตญฺจ ธมฺมํ ปจฺจยา โนจิตฺโต ธมฺโม อุปฺปชฺชติ เหตุปจฺจยา. โนจิตฺตํ เอกํ ขนฺธญฺจ จิตฺตญฺจ ปจฺจยา เทฺว ขนฺธา จิตฺต\-สมุฏฺฐานญฺจ รูปํ. เทฺว ขนฺเธ จ ฯเปฯ. จิตฺตญฺจ มหาภูเต จ ปจฺจยา จิตฺต\-สมุฏฺฐานํ รูปํ. จิตฺตญฺจ วตฺถุญฺจ ปจฺจยา โนจิตฺตา ขนฺธา. ปฏิสนฺธิ\-กฺขเณ โนจิตฺตํ เอกํ ขนฺธญฺจ จิตฺตญฺจ ปจฺจยา เทฺว ขนฺธา กฏตฺตา จ รูปํ. เทฺว ขนฺเธ จ ฯเปฯ. ปฏิสนฺธิ\-กฺขเณ จิตฺตญฺจ มหาภูเต จ ปจฺจยา กฏตฺตารูปํ. ปฏิสนฺธิ\-กฺขเณ จิตฺตญฺจ วตฺถุญฺจ ปจฺจยา โนจิตฺตา ขนฺธา. (1)}

\proseitem{62}{จิตฺตํ ธมฺมํ ปจฺจยา โนจิตฺโต ธมฺโม อุปฺปชฺชติ อารมฺมณ\-ปจฺจยา. เอกํ. (ปฏิจฺจ\-วาร\-สทิสํ.)}

\prose{โนจิตฺตํ ธมฺมํ ปจฺจยา โนจิตฺโต ธมฺโม อุปฺปชฺชติ อารมฺมณ\-ปจฺจยา. โนจิตฺตํ เอกํ ขนฺธํ ปจฺจยา เทฺว ขนฺธา. เทฺว ขนฺเธ ฯเปฯ. ปฏิสนฺธิ\-กฺขเณ ฯเปฯ. ปฏิสนฺธิ\-กฺขเณ วตฺถุํ ปจฺจยา ขนฺธา. จกฺขายตนํ ปจฺจยา จกฺขุ\-วิญฺญาณสหคตา ขนฺธา ฯเปฯ. กายายตนํ ปจฺจยา กาย\-วิญฺญาณสหคตา ขนฺธา. วตฺถุํ ปจฺจยา โนจิตฺตา ขนฺธา. (1)}

\prose{โนจิตฺตํ ธมฺมํ ปจฺจยา จิตฺโต ธมฺโม อุปฺปชฺชติ อารมฺมณ\-ปจฺจยา. โนจิตฺเต ขนฺเธ ปจฺจยา จิตฺตํ. วตฺถุํ ปจฺจยา จิตฺตํ. ปฏิสนฺธิ\-กฺขเณ \csromanfolio{391}โนจิตฺเต ขนฺเธ ปจฺจยา จิตฺตํ. ปฏิสนฺธิ\-กฺขเณ วตฺถุํ ปจฺจยา จิตฺตํ. จกฺขายตนํ ปจฺจยา จกฺขุวิญฺญาณํ ฯเปฯ. กายายตนํ ปจฺจยา กายวิญฺญาณํ. (2)}

\prose{โนจิตฺตํ ธมฺมํ ปจฺจยา จิตฺโต จ โนจิตฺโต จ ธมฺมา อุปฺปชฺชนฺติ อารมฺมณ\-ปจฺจยา. โนจิตฺตํ เอกํ ขนฺธํ ปจฺจยา เทฺว ขนฺธา จิตฺตญฺจ. เทฺว ขนฺเธ ฯเปฯ. วตฺถุํ ปจฺจยา จิตฺตญฺจ สมฺปยุตฺตกา จ ขนฺธา. ปฏิสนฺธิ\-กฺขเณ ฯเปฯ. ปฏิสนฺธิ\-กฺขเณ วตฺถุํ ปจฺจยา จิตฺตญฺจ สมฺปยุตฺตกา จ ขนฺธา. จกฺขายตนํ ปจฺจยา จกฺขุวิญฺญาณํ สมฺปยุตฺตกา จ ขนฺธา ฯเปฯ. กายายตนํ ปจฺจยา กายวิญฺญาณํ สมฺปยุตฺตกา จ ขนฺธา. (3)}

\prose{จิตฺตญฺจ โนจิตฺตญฺจ ธมฺมํ ปจฺจยา โนจิตฺโต ธมฺโม อุปฺปชฺชติ อารมฺมณ\-ปจฺจยา. โนจิตฺตํ เอกํ ขนฺธญฺจ จิตฺตญฺจ ปจฺจยา เทฺว ขนฺธา. เทฺว ขนฺเธ จ ฯเปฯ. จิตฺตญฺจ วตฺถุญฺจ ปจฺจยา โนจิตฺตา ขนฺธา. ปฏิสนฺธิ\-กฺขเณ ฯเปฯ. ปฏิสนฺธิ\-กฺขเณ จิตฺตญฺจ วตฺถุญฺจ ปจฺจยา โนจิตฺตา ขนฺธา. จกฺขายตนญฺจ จกฺขุวิญฺญาณญฺจ ปจฺจยา จกฺขุ\-วิญฺญาณสหคตา ขนฺธา ฯเปฯ. กายายตนญฺจ. (สํขิตฺตํ.) (1)}

\csromanmarkh{1. ปจฺจยานุโลม}\csromanmarkhii{2. สงฺขฺยาวาร}\csromanheadernomark{2}{1. \textbf{ปจฺจยานุโลม 2. สงฺขฺยาวาร}}
\proseitem{62}{เหตุยา ปญฺจ, อารมฺมเณ ปญฺจ, อธิปติยา ปญฺจ, (สพฺพตฺถ ปญฺจ,) อวิคเต ปญฺจ.}

\csromanmarkh{2. ปจฺจย\-ปจฺจนีย}\csromanmarkhii{1. วิภงฺควาร}\csromanheadernomark{2}{2. ปจฺจย\-ปจฺจนีย 1. วิภงฺควาร}
\proseitem{64}{จิตฺตํ ธมฺมํ ปจฺจยา โนจิตฺโต ธมฺโม อุปฺปชฺชติ นเหตุปจฺจยา. อเหตุกํ จิตฺตํ ปจฺจยา สมฺปยุตฺตกา ขนฺธา จิตฺต\-สมุฏฺฐานญฺจ รูปํ. อเหตุก\-ปฏิสนฺธิ\-กฺขเณ ฯเปฯ. วิจิกิจฺฉา\-สหคตํ อุทฺธจฺจ\-สหคตํ จิตฺตํ ปจฺจยา วิจิกิจฺฉา\-สหคโต อุทฺธจฺจ\-สหคโต โมโห. (1)}

\proseitem{65}{โนจิตฺตํ ธมฺมํ ปจฺจยา โนจิตฺโต ธมฺโม อุปฺปชฺชติ นเหตุปจฺจยา. อเหตุกํ โนจิตฺตํ เอกํ ขนฺธํ ปจฺจยา เทฺว ขนฺธา จิตฺต\-สมุฏฺฐานญฺจ รูปํ. เทฺว ขนฺเธ ฯเปฯ. อเหตุก\-ปฏิสนฺธิ\-กฺขเณ ฯเปฯ. (ยาว อสญฺญสตฺตา.) จกฺขายตนํ ปจฺจยา จกฺขุ\-วิญฺญาณสหคตา ขนฺธา ฯเปฯ. กายายตนํ ปจฺจยา กาย\-วิญฺญาณสหคตา ขนฺธา. วตฺถุํ ปจฺจยา อเหตุกา \csromanfolio{392}โนจิตฺตา ขนฺธา. วิจิกิจฺฉา\-สหคเต อุทฺธจฺจ\-สหคเต ขนฺเธ จ วตฺถุญฺจ ปจฺจยา วิจิกิจฺฉา\-สหคโต อุทฺธจฺจ\-สหคโต โมโห. (1)}

\prose{โนจิตฺตํ ธมฺมํ ปจฺจยา จิตฺโต ธมฺโม อุปฺปชฺชติ นเหตุปจฺจยา. อเหตุเก โนจิตฺเต ขนฺเธ ปจฺจยา จิตฺตํ. วตฺถุํ ปจฺจยา จิตฺตํ. อเหตุก\-ปฏิสนฺธิ\-กฺขเณ ฯเปฯ. อเหตุก\-ปฏิสนฺธิ\-กฺขเณ วตฺถุํ ปจฺจยา จิตฺตํ. จกฺขายตนํ ปจฺจยา จกฺขุวิญฺญาณํ ฯเปฯ. กายายตนํ ปจฺจยา กายวิญฺญาณํ. (2)}

\prose{โนจิตฺตํ ธมฺมํ ปจฺจยา จิตฺโต จ โนจิตฺโต จ ธมฺมา อุปฺปชฺชนฺติ นเหตุปจฺจยา. อเหตุกํ โนจิตฺตํ เอกํ ขนฺธํ ปจฺจยา เทฺว ขนฺธา จิตฺตญฺจ จิตฺต\-สมุฏฺฐานญฺจ รูปํ. เทฺว ขนฺเธ ฯเปฯ. วตฺถุํ ปจฺจยา จิตฺตญฺจ สมฺปยุตฺตกา จ ขนฺธา. อเหตุก\-ปฏิสนฺธิ\-กฺขเณ ฯเปฯ. อเหตุก\-ปฏิสนฺธิ\-กฺขเณ วตฺถุํ ปจฺจยา จิตฺตญฺจ สมฺปยุตฺตกา จ ขนฺธา. จกฺขายตนํ ปจฺจยา จกฺขุวิญฺญาณํ สมฺปยุตฺตกา จ ขนฺธา ฯเปฯ. กายายตนํ ปจฺจยา กายวิญฺญาณํ ฯเปฯ. (3)}

\prose{จิตฺตญฺจ โนจิตฺตญฺจ ธมฺมํ ปจฺจยา โนจิตฺโต ธมฺโม อุปฺปชฺชติ นเหตุปจฺจยา. อเหตุกํ โนจิตฺตํ เอกํ ขนฺธญฺจ จิตฺตญฺจ ปจฺจยา เทฺว ขนฺธา จิตฺต\-สมุฏฺฐานญฺจ รูปํ. เทฺว ขนฺเธ จ ฯเปฯ. จิตฺตญฺจ วตฺถุญฺจ ปจฺจยา โนจิตฺตา ขนฺธา. อเหตุก\-ปฏิสนฺธิ\-กฺขเณ ฯเปฯ. อเหตุก\-ปฏิสนฺธิ\-กฺขเณ จิตฺตญฺจ วตฺถุญฺจ ปจฺจยา โนจิตฺตา ขนฺธา. จกฺขายตนญฺจ จกฺขุวิญฺญาณญฺจ ปจฺจยา จกฺขุ\-วิญฺญาณสหคตา ขนฺธา ฯเปฯ. กายายตนญฺจ ฯเปฯ. วิจิกิจฺฉา\-สหคเต อุทฺธจฺจ\-สหคเต ขนฺเธ จ จิตฺตญฺจ ปจฺจยา วิจิกิจฺฉา\-สหคโต อุทฺธจฺจ\-สหคโต โมโห. (1) (สํขิตฺตํ.)}

\csromanmarkh{2. ปจฺจย\-ปจฺจนีย}\csromanmarkhii{2. สงฺขฺยาวาร}\csromanheadernomark{2}{2. ปจฺจย\-ปจฺจนีย 2. สงฺขฺยาวาร}
\proseitem{65}{นเหตุยา ปญฺจ, นอารมฺมเณ ตีณิ, นอธิปติยา ปญฺจ, นอนนฺตเร ตีณิ, นสมนนฺตเร ตีณิ, นอญฺญมญฺเญ ตีณิ, นอุปนิสฺสเย ตีณิ, นปุเรชาเต ปญฺจ, นปจฺฉาชาเต ปญฺจ, นอาเสวเน ปญฺจ, นกมฺเม ตีณิ, นวิปาเก ปญฺจ, นอาหาเร เอกํ, นอินฺทฺริเย เอกํ, นฌาเน ปญฺจ, นมคฺเค ปญฺจ, นสมฺปยุตฺเต ตีณิ, นวิปฺปยุตฺเต ปญฺจ, โนนตฺถิยา ตีณิ, โนวิคเต ตีณิ.}

\csromanmarkh{56. จิตฺตทุก}\csromanmarkhii{3. ปจฺจยา\-นุโลม\-ปจฺจนีย}\csromanheadernomark{2}{56. จิตฺตทุก 3. ปจฺจยา\-นุโลม\-ปจฺจนีย}
\vspace{\tipitakaparskip}
\csromancenter{\textbf{เหตุ}ปจฺจยา นอารมฺมเณ ตีณิ, นอธิปติยา ปญฺจ. (สํขิตฺตํ.)}
\proseitem{68}{\csromanfolio{393}นเหตุปจฺจยา อารมฺมเณ ปญฺจ, อนนฺตเร ปญฺจ, (สพฺพตฺถ ปญฺจ,) มคฺเค ตีณิ ฯเปฯ อวิคเต ปญฺจ.}

\vspace{\tipitakaparskip}
\csromancenter{(นิสฺสยวาโร ปจฺจยวาร\-สทิโส.)}
\csromanmarkh{56. จิตฺตทุก}\csromanmarkhii{5. สํสฏฺฐวาร}\csromanheadernomark{2}{56. จิตฺตทุก 5. สํสฏฺฐวาร}
\vspace{\tipitakaparskip}
\csromancenter{\textbf{เหตุ}}
\proseitem{69}{จิตฺตํ ธมฺมํ สํสฏฺโฐ โนจิตฺโต ธมฺโม อุปฺปชฺชติ เหตุปจฺจยา. จิตฺตํ สํสฏฺฐา สมฺปยุตฺตกา ขนฺธา. ปฏิสนฺธิ\-กฺขเณ ฯเปฯ. (1)}

\prose{โนจิตฺตํ ธมฺมํ สํสฏฺโฐ โนจิตฺโต ธมฺโม อุปฺปชฺชติ เหตุปจฺจยา. โนจิตฺตํ เอกํ ขนฺธํ สํสฏฺฐา เทฺว ขนฺธา. เทฺว ขนฺเธ สํสฏฺโฐ เอโก ขนฺโธ. ปฏิสนฺธิ\-กฺขเณ ฯเปฯ. (1)}

\prose{โนจิตฺตํ ธมฺมํ สํสฏฺโฐ จิตฺโต ธมฺโม อุปฺปชฺชติ เหตุปจฺจยา. โนจิตฺเต ขนฺเธ สํสฏฺฐํ จิตฺตํ. ปฏิสนฺธิเณ ฯเปฯ. (2)}

\prose{โนจิตฺตํ ธมฺมํ สํสฏฺโฐ จิตฺโต จ โนจิตฺโต จ ธมฺมา อุปฺปชฺชนฺติ เหตุปจฺจยา. โนจิตฺตํ เอกํ ขนฺธํ สํสฏฺฐา เทฺว ขนฺธา จิตฺตญฺจ. เทฺว ขนฺเธ ฯเปฯ. ปฏิสนฺธิ\-กฺขเณ ฯเปฯ. (3)}

\prose{จิตฺตญฺจ โนจิตฺตญฺจ ธมฺมํ สํสฏฺโฐ โนจิตฺโต ธมฺโม อุปฺปชฺชติ เหตุปจฺจยา. โนจิตฺตํ เอกํ ขนฺธญฺจ จิตฺตญฺจ สํสฏฺฐา เทฺว ขนฺธา. เทฺว ขนฺเธ จ ฯเปฯ. ปฏิสนฺธิ\-กฺขเณ ฯเปฯ. (สํขิตฺตํ.)}

\prose{\textbf{เหตุ}ยา ปญฺจ, อารมฺมเณ ปญฺจ, อธิปติยา ปญฺจ, (สพฺพตฺถ ปญฺจ,) อวิคเต ปญฺจ. (สํขิตฺตํ.)}

\prose{\csromanfolio{394}นเหตุยา ปญฺจ, นอธิปติยา ปญฺจ, นปุเรชาเต ปญฺจ, นปจฺฉาชาเต ปญฺจ, นอาเสวเน ปญฺจ, นกมฺเม ตีณิ, นวิปาเก ปญฺจ, นฌาเน ปญฺจ, นมคฺเค ปญฺจ, นวิปฺปยุตฺเต ปญฺจ.}

\prose{(เอวํ อิตเร เทฺว คณนาปิ สมฺปยุตฺตวาโรปิ สพฺเพ กาตพฺพา.)}

\csromanmarkh{56. จิตฺตทุก}\csromanmarkhii{7. ปญฺหาวาร}\csromanheadernomark{2}{56. \textbf{จิตฺตทุก 7. ปญฺหาวาร}}
\csromanmarkh{1. ปจฺจยานุโลม}\csromanmarkhii{1. วิภงฺควาร}\csromanheadernomark{2}{1. \textbf{ปจฺจยานุโลม 1. วิภงฺควาร}}
\csromanheader{2}{\textbf{เหตุ}}
\proseitem{70}{โนจิตฺโต ธมฺโม โนจิตฺตสฺส ธมฺมสฺส เหตุปจฺจเยน ปจฺจโย. โนจิตฺตา เหตู สมฺปยุตฺตกานํ ขนฺธานํ จิตฺตสมุฏฺฐานานญฺจ รูปานํ เหตุปจฺจเยน ปจฺจโย. ปฏิสนฺธิ\-กฺขเณ ฯเปฯ. (1)}

\prose{โนจิตฺโต ธมฺโม จิตฺตสฺส ธมฺมสฺส เหตุปจฺจเยน ปจฺจโย. โนจิตฺตา เหตู จิตฺตสฺส เหตุปจฺจเยน ปจฺจโย. ปฏิสนฺธิ\-กฺขเณ ฯเปฯ. (2)}

\prose{โนจิตฺโต ธมฺโม จิตฺตสฺส จ โนจิตฺตสฺส จ ธมฺมสฺส เหตุปจฺจเยน ปจฺจโย. โนจิตฺตา เหตู สมฺปยุตฺตกานํ ขนฺธานํ จิตฺตสฺส จ จิตฺตสมุฏฺฐานานญฺจ รูปานํ เหตุปจฺจเยน ปจฺจโย. ปฏิสนฺธิ\-กฺขเณ ฯเปฯ. (3)}

\proseitem{71}{จิตฺโต ธมฺโม จิตฺตสฺส ธมฺมสฺส อารมฺมณ\-ปจฺจเยน ปจฺจโย. จิตฺตํ อารพฺภ จิตฺตํ อุปฺปชฺชติ. (มูลํ กาตพฺพํ.) จิตฺตํ อารพฺภ โนจิตฺตา ขนฺธา อุปฺปชฺชนฺติ. (มูลํ กาตพฺพํ.) จิตฺตํ อารพฺภ จิตฺตญฺจ สมฺปยุตฺตกา จ ขนฺธา อุปฺปชฺชนฺติ. (3)}

\proseitem{72}{โนจิตฺโต ธมฺโม โนจิตฺตสฺส ธมฺมสฺส อารมฺมณ\-ปจฺจเยน ปจฺจโย. ทานํ, ทตฺวา, สีลํ ฯเปฯ อุโปสถกมฺมํ กตฺวา ตํ ปจฺจเวกฺขติ, อสฺสาเทติ อภินนฺทติ, ตํ อารพฺภ ราโค อุปฺปชฺชติ ฯเปฯ โทมนสฺสํ อุปฺปชฺชติ. ปุพฺเพ สุจิณฺณานิ ฯเปฯ. ฌานา วุฏฺฐหิตฺวา ฌานํ ฯเปฯ. อริยา มคฺคา \csromanfolio{395}วุฏฺฐหิตฺวา มคฺคํ ปจฺจเวกฺขนฺติ, ผลํ ปจฺจเวกฺขนฺติ, นิพฺพานํ ปจฺจเวกฺขนฺติ. นิพฺพานํ โคตฺรภุสฺส, โวทานสฺส, มคฺคสฺส, ผลสฺส อาวชฺชนาย อารมฺมณ\-ปจฺจเยน ปจฺจโย. อริยา โนจิตฺเต ปหีเน กิเลเส ปจฺจเวกฺขนฺติ, วิกฺขมฺภิเต กิเลเส ปจฺจเวกฺขนฺติ. ปุพฺเพ สมุทาจิณฺเณ กิเลเส ชานนฺติ. จกฺขุํ ฯเปฯ วตฺถุํ โนจิตฺเต ขนฺเธ อนิจฺจโต ฯเปฯ โทมนสฺสํ อุปฺปชฺชติ. ทิพฺเพน จกฺขุนา รูปํ ปสฺสติ. ทิพฺพาย โสตธาตุยา สทฺทํ สุณาติ. เจโตปริย\-ญาเณน โนจิตฺต\-สมงฺคิสฺส จิตฺตํ ชานาติ. อากาสา\-นญฺจา\-ยตนํ ฯเปฯ. อากิญฺจญฺญา\-ยตนํ ฯเปฯ. รูปายตนํ จกฺขุ\-วิญฺญาณสหคตานํ ขนฺธานํ ฯเปฯ. โผฏฺฐพฺพา\-ยตนํ กาย\-วิญฺญาณสหคตานํ ขนฺธานํ ฯเปฯ. โนจิตฺตา ขนฺธา อิทฺธิวิธ\-ญาณสฺส, เจโตปริย\-ญาณสฺส, ปุพฺเพ\-นิวาสา\-นุสฺสติ\-ญาณสฺส, ยถากมฺมูปคญาณสฺส, อนาคตํสญาณสฺส, อาวชฺชนาย อารมฺมณ\-ปจฺจเยน ปจฺจโย. (1)}

\prose{โนจิตฺโต ธมฺโม จิตฺตสฺส ธมฺมสฺส อารมฺมณ\-ปจฺจเยน ปจฺจโย. ทานํ ทตฺวา ฯเปฯ. (ปฐม\-คมน\-สทิสํ. นินฺนานากรณํ. อิมํ นานํ) รูปายตนํ จกฺขุ\-วิญฺญาณสฺส ฯเปฯ. โผฏฺฐพฺพา\-ยตนํ กายวิญฺญาณสฺส ฯเปฯ. โนจิตฺตา ขนฺธา อิทฺธิวิธ\-ญาณสฺส ฯเปฯ อาวชฺชนาย อารมฺมณ\-ปจฺจเยน ปจฺจโย. (2)}

\prose{โนจิตฺโต ธมฺโม จิตฺตสฺส จ โนจิตฺตสฺส จ ธมฺมสฺส อารมฺมณ\-ปจฺจเยน ปจฺจโย. ทานํ ทตฺวา ฯเปฯ. (ปฐม\-คมน\-สทิสํ. นินฺนานากรณํ. อิมํ นานํ.) รูปายตนํ จกฺขุ\-วิญฺญาณสฺส สมฺปยุตฺตกานญฺจ ขนฺธานํ ฯเปฯ. โผฏฺฐพฺพา\-ยตนํ กายวิญฺญาณสฺส สมฺปยุตฺตกานญฺจ ขนฺธานํ ฯเปฯ. โนจิตฺตา ขนฺธา อิทฺธิวิธ\-ญาณสฺส ฯเปฯ อาวชฺชนาย อารมฺมณ\-ปจฺจเยน ปจฺจโย. (3)}

\prose{จิตฺโต จ โนจิตฺโต จ ธมฺมา จิตฺตสฺส ธมฺมสฺส อารมฺมณ\-ปจฺจเยน ปจฺจโย. จิตฺตญฺจ สมฺปยุตฺตเก จ ขนฺเธ อารพฺภ จิตฺตํ อุปฺปชฺชติ. ตีณิ.}

\proseitem{73}{จิตฺโต ธมฺโม จิตฺตสฺส ธมฺมสฺส อธิปติ\-ปจฺจเยน ปจฺจโย. อารมฺมณา\-ธิปติ\csromandash{}จิตฺตํ ครุํ กตฺวา จิตฺตํ อุปฺปชฺชติ. (1)}

\prose{จิตฺโต ธมฺโม โนจิตฺตสฺส ธมฺมสฺส อธิปติ\-ปจฺจเยน ปจฺจโย. อารมฺมณา\-ธิปติ, สหชาตา\-ธิปติ.  อารมฺมณา\-ธิปติ\csromandash{}จิตฺตํ ครุํ \csromanfolio{396}กตฺวา โนจิตฺตา ขนฺธา อุปฺปชฺชนฺติ. สหชาตา\-ธิปติ\csromandash{}จิตฺตาธิปติ สมฺปยุตฺตกานํ ขนฺธานํ จิตฺตสมุฏฺฐานานญฺจ รูปานํ อธิปติ\-ปจฺจเยน ปจฺจโย. (2)}

\prose{จิตฺโต ธมฺโม จิตฺตสฺส จ โนจิตฺตสฺส จ ธมฺมสฺส อธิปติ\-ปจฺจเยน ปจฺจโย. อารมฺมณา\-ธิปติ\csromandash{}จิตฺตํ ครุํ กตฺวา จิตฺตญฺจ สมฺปยุตฺตกา จ ขนฺธา อุปฺปชฺชนฺติ. (3)}

\proseitem{74}{โนจิตฺโต ธมฺโม โนจิตฺตสฺส ธมฺมสฺส อธิปติ\-ปจฺจเยน ปจฺจโย. อารมฺมณา\-ธิปติ, สหชาตา\-ธิปติ.  อารมฺมณา\-ธิปติ\csromandash{}ทานํ ฯเปฯ สีลํ ฯเปฯ อุโปสถกมฺมํ กตฺวา ตํ ครุํ กตฺวา ปจฺจเวกฺขติ, อสฺสาเทติ อภินนฺทติ, ตํ ครุํ กตฺวา ราโค อุปฺปชฺชติ. ทิฏฺฐิ อุปฺปชฺชติ. ปุพฺเพ ฯเปฯ. ฌานา ฯเปฯ. อริยา มคฺคา วุฏฺฐหิตฺวา มคฺคํ ครุํ กตฺวา ปจฺจเวกฺขนฺติ, ผลํ ครุํ กตฺวา ปจฺจเวกฺขนฺติ, นิพฺพานํ ครุํ กตฺวา ปจฺจเวกฺขนฺติ. นิพฺพานํ โคตฺรภุสฺส, โวทานสฺส, มคฺคสฺส, ผลสฺส อธิปติ\-ปจฺจเยน ปจฺจโย. จกฺขุํ ฯเปฯ วตฺถุํ โนจิตฺเต ขนฺเธ ครุํ กตฺวา อสฺสาเทติ อภินนฺทติ, ตํ ครุํ กตฺวา ราโค อุปฺปชฺชติ. ทิฏฺฐิ อุปฺปชฺชติ. สหชาตา\-ธิปติ\csromandash{}โนจิตฺตาธิปติ สมฺปยุตฺตกานํ ขนฺธานํ จิตฺตสมุฏฺฐานานญฺจ รูปานํ อธิปติ\-ปจฺจเยน ปจฺจโย. (1)}

\prose{โนจิตฺโต ธมฺโม จิตฺตสฺส ธมฺมสฺส อธิปติ\-ปจฺจเยน ปจฺจโย. (เทฺวปิ คมนา ปฐม\-คมน\-สทิสํ. นินฺนานากรณํ. อารมฺมณา\-ธิปติ สหชาตา\-ธิปติ กาตพฺพา.) (2)}

\prose{จิตฺโต จ โนจิตฺโต จ ธมฺมา จิตฺตสฺส ธมฺมสฺส อธิปติ\-ปจฺจเยน ปจฺจโย. อารมฺมณา\-ธิปติ. (ตีณิปิ ครุการมฺมณา กาตพฺพา. อารมฺมณา\-ธิปติเยว.)}

\proseitem{75}{จิตฺโต ธมฺโม จิตฺตสฺส ธมฺมสฺส อนนฺตร\-ปจฺจเยน ปจฺจโย. ปุริมํ ปุริมํ จิตฺตํ ปจฺฉิมสฺส ปจฺฉิมสฺส จิตฺตสฺส อนนฺตร\-ปจฺจเยน ปจฺจโย. (1)}

\prose{จิตฺโต ธมฺโม โนจิตฺตสฺส ธมฺมสฺส อนนฺตร\-ปจฺจเยน ปจฺจโย. ปุริมํ ปุริมํ จิตฺตํ ปจฺฉิมานํ ปจฺฉิมานํ โนจิตฺตานํ ขนฺธานํ อนนฺตร\-ปจฺจเยน ปจฺจโย. จิตฺตํ วุฏฺฐานสฺส อนนฺตร\-ปจฺจเยน ปจฺจโย. (2)}

\prose{\csromanfolio{397}จิตฺโต ธมฺโม จิตฺตสฺส จ โนจิตฺตสฺส จ ธมฺมสฺส อนนฺตร\-ปจฺจเยน ปจฺจโย. ปุริมํ ปุริมํ จิตฺตํ ปจฺฉิมสฺส ปจฺฉิมสฺส จิตฺตสฺส สมฺปยุตฺตกานญฺจ ขนฺธานํ อนนฺตร\-ปจฺจเยน ปจฺจโย. (3)}

\proseitem{76}{โนจิตฺโต ธมฺโม โนจิตฺตสฺส ธมฺมสฺส อนนฺตร\-ปจฺจเยน ปจฺจโย. ปุริมา ปุริมา โนจิตฺตา ขนฺธา ฯเปฯ ผลสมาปตฺติยา อนนฺตร\-ปจฺจเยน ปจฺจโย.}

\vspace{\tipitakaparskip}
\csromancenter{(อิเม เทฺว ปูเรตุกาเมน กาตพฺพา ปุริมคมน\-สทิสํ.)}
\prose{จิตฺโต จ โนจิตฺโต จ ธมฺมา จิตฺตสฺส ธมฺมสฺส อนนฺตร\-ปจฺจเยน ปจฺจโย. ปุริมํ ปุริมํ จิตฺตญฺจ สมฺปยุตฺตกา จ ขนฺธา ปจฺฉิมสฺส ปจฺฉิมสฺส จิตฺตสฺส อนนฺตร\-ปจฺจเยน ปจฺจโย. (มูลํ ปุจฺฉิตพฺพํ.) ปุริมํ ปุริมํ จิตฺตญฺจ สมฺปยุตฺตกา จ ขนฺธา ปจฺฉิมานํ ปจฺฉิมานํ โนจิตฺตานํ ขนฺธานํ อนนฺตร\-ปจฺจเยน ปจฺจโย. จิตฺตญฺจ สมฺปยุตฺตกา จ ขนฺธา วุฏฺฐานสฺส อนนฺตร\-ปจฺจเยน ปจฺจโย. (มูลํ ปุจฺฉิตพฺพํ.) ปุริมํ ปุริมํ จิตฺตญฺจ สมฺปยุตฺตกา จ ขนฺธา ปจฺฉิมสฺส ปจฺฉิมสฺส จิตฺตสฺส สมฺปยุตฺตกานญฺจ ขนฺธานํ อนนฺตร\-ปจฺจเยน ปจฺจโย. (3)}

\prose{สมนนฺตร\-ปจฺจเยน ปจฺจโย. สหชาต\-ปจฺจเยน ปจฺจโย. อญฺญมญฺญ\-ปจฺจเยน ปจฺจโย. (ปฏิจฺจ\-วาร\-สทิสา.) นิสฺสย\-ปจฺจเยน ปจฺจโย. ปญฺจ. (ปจฺจยวาร\-สทิสา.)}

\csromanheader{2}{\textbf{อุปนิสฺสย}}
\proseitem{77}{จิตฺโต ธมฺโม จิตฺตสฺส ธมฺมสฺส อุปนิสฺสย\-ปจฺจเยน ปจฺจโย. อารมฺมณู\-ปนิสฺสโย, อนนฺตรู\-ปนิสฺสโย, ปกตู\-ปนิสฺสโย ฯเปฯ. ปกตู\-ปนิสฺสโย\csromandash{}จิตฺตํ จิตฺตสฺส อุปนิสฺสย\-ปจฺจเยน ปจฺจโย. ตีณิ.}

\prose{โนจิตฺโต ธมฺโม โนจิตฺตสฺส ธมฺมสฺส อุปนิสฺสย\-ปจฺจเยน ปจฺจโย. อารมฺมณู\-ปนิสฺสโย, อนนฺตรู\-ปนิสฺสโย, ปกตู\-ปนิสฺสโย ฯเปฯ. ปกตู\-ปนิสฺสโย\csromandash{}สทฺธํ อุปนิสฺสาย ทานํ เทติ ฯเปฯ. มานํ ชปฺเปติ. ทิฏฺฐิํ คณฺหาติ. สีลํ ฯเปฯ. เสนาสนํ อุปนิสฺสาย ทานํ เทติ ฯเปฯ สํฆํ ภินฺทติ. สทฺธา ฯเปฯ. เสนาสนํ สทฺธาย ฯเปฯ มคฺคสฺส, ผลสมาปตฺติยา อุปนิสฺสย\-ปจฺจเยน ปจฺจโย. (1)}

\prose{โนจิตฺโต ธมฺโม จิตฺตสฺส ธมฺมสฺส อุปนิสฺสย\-ปจฺจเยน ปจฺจโย. (อิเม เทฺวปิ ปุเรตุกาเมน สพฺพตฺถ กาตพฺพา. ปฐม\-คมน\-สทิสํ. นินฺนานากรณํ.)}

\prose{\csromanfolio{398}จิตฺโต จ โนจิตฺโต จ ธมฺมา จิตฺตสฺส ธมฺมสฺส อุปนิสฺสย\-ปจฺจเยน ปจฺจโย. อารมฺมณู\-ปนิสฺสโย, อนนฺตรู\-ปนิสฺสโย, ปกตู\-ปนิสฺสโย ฯเปฯ. ปกตู\-ปนิสฺสโย\csromandash{}จิตฺตญฺจ สมฺปยุตฺตกา จ ขนฺธา จิตฺตสฺส อุปนิสฺสย\-ปจฺจเยน ปจฺจโย. ตีณิ.}

\proseitem{78}{โนจิตฺโต ธมฺโม โนจิตฺตสฺส ธมฺมสฺส ปุเรชาต\-ปจฺจเยน ปจฺจโย. อารมฺมณ\-ปุเรชาตํ, วตฺถุ\-ปุเรชาตํ.  อารมฺมณ\-ปุเรชาตํ\csromandash{}จกฺขุํ ฯเปฯ วตฺถุํ อนิจฺจโต ฯเปฯ โทมนสฺสํ อุปฺปชฺชติ. ทิพฺเพน จกฺขุนา รูปํ ปสฺสติ. ทิพฺพาย โสตธาตุยา สทฺทํ สุณาติ. รูปายตนํ จกฺขุ\-วิญฺญาณสหคตานํ ขนฺธานํ ฯเปฯ. โผฏฺฐพฺพา\-ยตนํ กาย\-วิญฺญาณสหคตานํ ขนฺธานํ ปุเรชาต\-ปจฺจเยน ปจฺจโย. วตฺถุ\-ปุเรชาตํ\csromandash{}จกฺขายตนํ จกฺขุ\-วิญฺญาณสหคตานํ ขนฺธานํ ฯเปฯ กายายตนํ ฯเปฯ. วตฺถุ โนจิตฺตานํ ขนฺธานํ ปุเรชาต\-ปจฺจเยน ปจฺจโย. (1)}

\prose{โนจิตฺโต ธมฺโม จิตฺตสฺส ธมฺมสฺส ปุเรชาต\-ปจฺจเยน ปจฺจโย. อารมฺมณ\-ปุเรชาตํ, วตฺถุ\-ปุเรชาตํ.  อารมฺมณ\-ปุเรชาตํ\csromandash{}จกฺขุํ ฯเปฯ วตฺถุํ อนิจฺจโต ฯเปฯ โทมนสฺสํ อุปฺปชฺชติ. ทิพฺเพน จกฺขุนา รูปํ ปสฺสติ. ทิพฺพาย โสตธาตุยา สทฺทํ สุณาติ. รูปายตนํ จกฺขุ\-วิญฺญาณสฺส ฯเปฯ. โผฏฺฐพฺพา\-ยตนํ กายวิญฺญาณสฺส ฯเปฯ. วตฺถุ\-ปุเรชาตํ\csromandash{}จกฺขายตนํ จกฺขุ\-วิญฺญาณสฺส ฯเปฯ กายายตนํ กายวิญฺญาณสฺส ฯเปฯ. วตฺถุ จิตฺตสฺส ปุเรชาต\-ปจฺจเยน ปจฺจโย. (2)}

\prose{โนจิตฺโต ธมฺโม จิตฺตสฺส จ โนจิตฺตสฺส จ ธมฺมสฺส ปุเรชาต\-ปจฺจเยน ปจฺจโย. อารมฺมณ\-ปุเรชาตํ, วตฺถุ\-ปุเรชาตํ.  อารมฺมณ\-ปุเรชาตํ\csromandash{} จกฺขุํ ฯเปฯ วตฺถุํ อนิจฺจโต ฯเปฯ โทมนสฺสํ อุปฺปชฺชติ. ทิพฺเพน จกฺขุนา รูปํ ปสฺสติ. ทิพฺพาย โสตธาตุยา สทฺทํ สุณาติ. รูปายตนํ จกฺขุ\-วิญฺญาณสฺส จ สมฺปยุตฺตกานญฺจ ขนฺธานํ ฯเปฯ. โผฏฺฐพฺพา\-ยตนํ ฯเปฯ. วตฺถุ\-ปุเรชาตํ\csromandash{}จกฺขายตนํ จกฺขุ\-วิญฺญาณสฺส สมฺปยุตฺตกานญฺจ ขนฺธานํ ปุเรชาต\-ปจฺจเยน ปจฺจโย ฯเปฯ กายายตนํ กายวิญฺญาณสฺส สมฺปยุตฺตกานญฺจ ขนฺธานํ ฯเปฯ. วตฺถุ จิตฺตสฺส จ สมฺปยุตฺตกานญฺจ ขนฺธานํ ปุเรชาต\-ปจฺจเยน ปจฺจโย. (3)}

\csromanheader{2}{56. จิตฺตทุก ปจฺฉาชาตา\-เสวน}
\proseitem{79}{\csromanfolio{399}จิตฺโต ธมฺโม โนจิตฺตสฺส ธมฺมสฺส ปจฺฉาชาต\-ปจฺจเยน ปจฺจโย. (สํขิตฺตํ.) (1)}

\prose{โนจิตฺโต ธมฺโม โนจิตฺตสฺส ธมฺมสฺส ปจฺฉาชาต\-ปจฺจเยน ปจฺจโย. (สํขิตฺตํ.) (1)}

\prose{จิตฺโต จ โนจิตฺโต จ ธมฺมา โนจิตฺตสฺส ธมฺมสฺส ปจฺฉาชาต\-ปจฺจเยน ปจฺจโย. (สํขิตฺตํ.) (1)}

\prose{จิตฺโต ธมฺโม จิตฺตสฺส ธมฺมสฺส อาเสวน\-ปจฺจเยน ปจฺจโย. นว.}

\proseitem{80}{โนจิตฺโต ธมฺโม โนจิตฺตสฺส ธมฺมสฺส กมฺมปจฺจเยน ปจฺจโย. สหชาตา, นานากฺขณิกา.  สหชาตา โนจิตฺตา เจตนา สมฺปยุตฺตกานํ ขนฺธานํ จิตฺตสมุฏฺฐานานญฺจ รูปานํ กมฺมปจฺจเยน ปจฺจโย. ปฏิสนฺธิ\-กฺขเณ ฯเปฯ. นานากฺขณิกา โนจิตฺตา เจตนา วิปากานํ ขนฺธานํ กฏตฺตา จ รูปานํ กมฺมปจฺจเยน ปจฺจโย. (1)}

\prose{โนจิตฺโต ธมฺโม จิตฺตสฺส ธมฺมสฺส กมฺมปจฺจเยน ปจฺจโย. สหชาตา, นานากฺขณิกา.  สหชาตา โนจิตฺตา เจตนา จิตฺตสฺส กมฺมปจฺจเยน ปจฺจโย. ปฏิสนฺธิ\-กฺขเณ ฯเปฯ. นานากฺขณิกา โนจิตฺตา เจตนา วิปากสฺส จิตฺตสฺส กมฺมปจฺจเยน ปจฺจโย. (2)}

\prose{โนจิตฺโต ธมฺโม จิตฺตสฺส จ โนจิตฺตสฺส จ ธมฺมสฺส กมฺมปจฺจเยน ปจฺจโย. สหชาตา, นานากฺขณิกา.  สหชาตา โนจิตฺตา เจตนา สมฺปยุตฺตกานํ ขนฺธานํ จิตฺตสฺส จ จิตฺตสมุฏฺฐานานญฺจ รูปานํ กมฺมปจฺจเยน ปจฺจโย. ปฏิสนฺธิ\-กฺขเณ ฯเปฯ. นานากฺขณิกา โนจิตฺตา เจตนา วิปากานํ ขนฺธานํ จิตฺตสฺส จ กฏตฺตา จ รูปานํ กมฺมปจฺจเยน ปจฺจโย. (3)}

\csromanheader{2}{\textbf{วิปากาทิ}}
\proseitem{81}{จิตฺโต ธมฺโม โนจิตฺตสฺส ธมฺมสฺส วิปาก\-ปจฺจเยน ปจฺจโย. ปญฺจ. อาหารปจฺจเยน ปจฺจโย. ปญฺจ. อินฺทฺริย\-ปจฺจเยน ปจฺจโย. ปญฺจ. โนจิตฺโต ธมฺโม โนจิตฺตสฺส ธมฺมสฺส ฌานปจฺจเยน ปจฺจโย. ตีณิ. มคฺคปจฺจเยน ปจฺจโย. ตีณิ. สมฺปยุตฺต\-ปจฺจเยน ปจฺจโย. ปญฺจ.}

\proseitem{82}{\csromanfolio{400}จิตฺโต ธมฺโม โนจิตฺตสฺส ธมฺมสฺส วิปฺปยุตฺต\-ปจฺจเยน ปจฺจโย. สหชาตํ, ปจฺฉาชาตํ. (สํขิตฺตํ.) (1)}

\prose{โนจิตฺโต ธมฺโม โนจิตฺตสฺส ธมฺมสฺส วิปฺปยุตฺต\-ปจฺจเยน ปจฺจโย. สหชาตํ, ปุเรชาตํ, ปจฺฉาชาตํ.  สหชาตา โนจิตฺตา ขนฺธา จิตฺต\-สมุฏฺฐานานํ รูปานํ วิปฺปยุตฺต\-ปจฺจเยน ปจฺจโย. ปฏิสนฺธิ\-กฺขเณ โนจิตฺตา ขนฺธา กฏตฺตารูปานํ ฯเปฯ. โนจิตฺตา ขนฺธา วตฺถุสฺส วิปฺปยุตฺต\-ปจฺจเยน ปจฺจโย. วตฺถุ ขนฺธานํ วิปฺปยุตฺต\-ปจฺจเยน ปจฺจโย. ปุเรชาตํ จกฺขายตนํ จกฺขุ\-วิญฺญาณสหคตานํ ขนฺธานํ ฯเปฯ กายายตนํ กาย\-วิญฺญาณสหคตานํ ขนฺธานํ ฯเปฯ. วตฺถุ โนจิตฺตานํ ขนฺธานํ วิปฺปยุตฺต\-ปจฺจเยน ปจฺจโย. ปจฺฉาชาตา โนจิตฺตา ขนฺธา ปุเรชาตสฺส อิมสฺส กายสฺส วิปฺปยุตฺต\-ปจฺจเยน ปจฺจโย. (1)}

\prose{โนจิตฺโต ธมฺโม จิตฺตสฺส ธมฺมสฺส วิปฺปยุตฺต\-ปจฺจเยน ปจฺจโย. สหชาตํ, ปุเรชาตํ.  สหชาตํ ปฏิสนฺธิ\-กฺขเณ วตฺถุ จิตฺตสฺส วิปฺปยุตฺต\-ปจฺจเยน ปจฺจโย. ปุเรชาตํ จกฺขายตนํ จกฺขุ\-วิญฺญาณสฺส ฯเปฯ กายายตนํ กายวิญฺญาณสฺส ฯเปฯ. วตฺถุ จิตฺตสฺส วิปฺปยุตฺต\-ปจฺจเยน ปจฺจโย. (2)}

\prose{โนจิตฺโต ธมฺโม จิตฺตสฺส จ โนจิตฺตสฺส จ ธมฺมสฺส วิปฺปยุตฺต\-ปจฺจเยน ปจฺจโย. สหชาตํ, ปุเรชาตํ.  สหชาตํ ปฏิสนฺธิ\-กฺขเณ วตฺถุ จิตฺตสฺส สมฺปยุตฺตกานญฺจ ขนฺธานํ วิปฺปยุตฺต\-ปจฺจเยน ปจฺจโย. ปุเรชาตํ จกฺขายตนํ จกฺขุ\-วิญฺญาณสฺส สมฺปยุตฺตกานญฺจ ขนฺธานํ วิปฺปยุตฺต\-ปจฺจเยน ปจฺจโย ฯเปฯ กายายตนํ กายวิญฺญาณสฺส สมฺปยุตฺตกานญฺจ ขนฺธานํ ฯเปฯ. วตฺถุ จิตฺตสฺส จ สมฺปยุตฺตกานญฺจ ขนฺธานํ วิปฺปยุตฺต\-ปจฺจเยน ปจฺจโย. (3)}

\prose{จิตฺโต จ โนจิตฺโต จ ธมฺมา โนจิตฺตสฺส ธมฺมสฺส วิปฺปยุตฺต\-ปจฺจเยน ปจฺจโย. สหชาตํ, ปจฺฉาชาตํ. (สํขิตฺตํ.) (1)}

\proseitem{83}{จิตฺโต ธมฺโม โนจิตฺตสฺส ธมฺมสฺส อตฺถิปจฺจเยน ปจฺจโย. สหชาตํ, ปจฺฉาชาตํ. (สํขิตฺตํ.) (1)}

\prose{โนจิตฺโต ธมฺโม โนจิตฺตสฺส ธมฺมสฺส อตฺถิปจฺจเยน ปจฺจโย. สหชาตํ, ปุเรชาตํ, ปจฺฉาชาตํ, อาหารํ, อินฺทฺริยํ. (สํขิตฺตํ.) (1)}

\prose{\csromanfolio{401}โนจิตฺโต ธมฺโม จิตฺตสฺส ธมฺมสฺส อตฺถิปจฺจเยน ปจฺจโย. สหชาตํ, ปุเรชาตํ.  สหชาตา โนจิตฺตา ขนฺธา จิตฺตสฺส อตฺถิปจฺจเยน ปจฺจโย. ปฏิสนฺธิ\-กฺขเณ ฯเปฯ. ปฏิสนฺธิ\-กฺขเณ วตฺถุ จิตฺตสฺส อตฺถิปจฺจเยน ปจฺจโย. ปุเรชาตํ จกฺขุํ ฯเปฯ วตฺถุํ อนิจฺจโต ฯเปฯ. (ปุเรชาต\-สทิสํ, สํขิตฺตํ.) (2)}

\prose{โนจิตฺโต ธมฺโม จิตฺตสฺส จ โนจิตฺตสฺส จ ธมฺมสฺส อตฺถิปจฺจเยน ปจฺจโย. สหชาตํ, ปุเรชาตํ.  สหชาโต โนจิตฺโต เอโก ขนฺโธ ทฺวินฺนํ ขนฺธานํ จิตฺตสฺส จ จิตฺตสมุฏฺฐานานญฺจ รูปานํ อตฺถิปจฺจเยน ปจฺจโย. ปฏิสนฺธิ\-กฺขเณ โนจิตฺโต เอโก ขนฺโธ ฯเปฯ. ปฏิสนฺธิ\-กฺขเณ วตฺถุ จิตฺตสฺส สมฺปยุตฺตกานญฺจ ขนฺธานํ อตฺถิปจฺจเยน ปจฺจโย. ปุเรชาตํ จกฺขุํ ฯเปฯ วตฺถุํ ฯเปฯ. (ปุเรชาต\-สทิสํ, สํขิตฺตํ.) (3)}

\proseitem{84}{จิตฺโต จ โนจิตฺโต จ ธมฺมา โนจิตฺตสฺส ธมฺมสฺส อตฺถิปจฺจเยน ปจฺจโย. สหชาตํ, ปุเรชาตํ, ปจฺฉาชาตํ, อาหารํ, อินฺทฺริยํ.  สหชาโต โนจิตฺโต เอโก ขนฺโธ จ จิตฺตญฺจ ทฺวินฺนํ ขนฺธานํ จิตฺตสมุฏฺฐานานญฺจ รูปานํ อตฺถิปจฺจเยน ปจฺจโย. เทฺว ขนฺธา จ ฯเปฯ. สหชาตํ จิตฺตญฺจ วตฺถุญฺจ โนจิตฺตานํ ขนฺธานํ อตฺถิปจฺจเยน ปจฺจโย. (ปฏิสนฺธิ\-กฺขเณปิ เทฺว.) สหชาตํ จิตฺตญฺจ สมฺปยุตฺตกา จ ขนฺธา จิตฺต\-สมุฏฺฐานานํ รูปานํ อตฺถิปจฺจเยน ปจฺจโย. สหชาตํ จิตฺตญฺจ มหาภูตา จ จิตฺต\-สมุฏฺฐานานํ รูปานํ อตฺถิปจฺจเยน ปจฺจโย. (ปฏิสนฺธิ\-กฺขเณปิ เทฺว.) ปจฺฉาชาตํ จิตฺตญฺจ สมฺปยุตฺตกา จ ขนฺธา ปุเรชาตสฺส อิมสฺส กายสฺส อตฺถิปจฺจเยน ปจฺจโย. ปจฺฉาชาตํ จิตฺตญฺจ สมฺปยุตฺตกา จ ขนฺธา กพฬีกาโร อาหาโร จ อิมสฺส กายสฺส อตฺถิปจฺจเยน ปจฺจโย. ปจฺฉาชาตํ จิตฺตญฺจ สมฺปยุตฺตกา จ ขนฺธา รูปชีวิตินฺทฺริยญฺจ กฏตฺตารูปานํ อตฺถิปจฺจเยน ปจฺจโย. (1)}

\csromanmarkh{1. ปจฺจยานุโลม}\csromanmarkhii{2. สงฺขฺยาวาร}\csromanheadernomark{2}{1. \textbf{ปจฺจยานุโลม 2. สงฺขฺยาวาร}}
\proseitem{85}{เหตุยา ตีณิ, อารมฺมเณ นว, อธิปติยา นว, อนนฺตเร นว, สมนนฺตเร นว, สหชาเต ปญฺจ, อญฺญมญฺเญ ปญฺจ, นิสฺสเย ปญฺจ, อุปนิสฺสเย นว, ปุเรชาเต ตีณิ, ปจฺฉาชาเต ตีณิ, อาเสวเน \csromanfolio{402}นว, กมฺเม ตีณิ, วิปาเก ปญฺจ, อาหาเร ปญฺจ, อินฺทฺริเย ปญฺจ, ฌาเน ตีณิ, มคฺเค ตีณิ, สมฺปยุตฺเต ปญฺจ, วิปฺปยุตฺเต ปญฺจ, อตฺถิยา ปญฺจ, นตฺถิยา นว, วิคเต นว, อวิคเต ปญฺจ.}

\proseitem{86}{จิตฺโต ธมฺโม จิตฺตสฺส ธมฺมสฺส อารมฺมณ\-ปจฺจเยน ปจฺจโย. อุปนิสฺสย\-ปจฺจเยน ปจฺจโย. (1)}

\prose{จิตฺโต ธมฺโม โนจิตฺตสฺส ธมฺมสฺส อารมฺมณ\-ปจฺจเยน ปจฺจโย. สหชาต\-ปจฺจเยน ปจฺจโย. อุปนิสฺสย\-ปจฺจเยน ปจฺจโย. ปจฺฉาชาต\-ปจฺจเยน ปจฺจโย. (2)}

\prose{จิตฺโต ธมฺโม จิตฺตสฺส จ โนจิตฺตสฺส จ ธมฺมสฺส อารมฺมณ\-ปจฺจเยน ปจฺจโย. อุปนิสฺสย\-ปจฺจเยน ปจฺจโย. (3)}

\proseitem{87}{โนจิตฺโต ธมฺโม โนจิตฺตสฺส ธมฺมสฺส อารมฺมณ\-ปจฺจเยน ปจฺจโย. สหชาต\-ปจฺจเยน ปจฺจโย. อุปนิสฺสย\-ปจฺจเยน ปจฺจโย. ปุเรชาต\-ปจฺจเยน ปจฺจโย. ปจฺฉาชาต\-ปจฺจเยน ปจฺจโย. กมฺมปจฺจเยน ปจฺจโย. อาหารปจฺจเยน ปจฺจโย. อินฺทฺริย\-ปจฺจเยน ปจฺจโย. (1)}

\prose{โนจิตฺโต ธมฺโม จิตฺตสฺส ธมฺมสฺส อารมฺมณ\-ปจฺจเยน ปจฺจโย. สหชาต\-ปจฺจเยน ปจฺจโย. อุปนิสฺสย\-ปจฺจเยน ปจฺจโย. ปุเรชาต\-ปจฺจเยน ปจฺจโย. กมฺมปจฺจเยน ปจฺจโย. (2)}

\prose{โนจิตฺโต ธมฺโม จิตฺตสฺส จ โนจิตฺตสฺส จ ธมฺมสฺส อารมฺมณ\-ปจฺจเยน ปจฺจโย. สหชาต\-ปจฺจเยน ปจฺจโย. อุปนิสฺสย\-ปจฺจเยน ปจฺจโย. ปุเรชาต\-ปจฺจเยน ปจฺจโย. กมฺมปจฺจเยน ปจฺจโย. (3)}

\proseitem{88}{จิตฺโต จ โนจิตฺโต จ ธมฺมา จิตฺตสฺส ธมฺมสฺส อารมฺมณ\-ปจฺจเยน ปจฺจโย. อุปนิสฺสย\-ปจฺจเยน ปจฺจโย. (1)}

\prose{จิตฺโต จ โนจิตฺโต จ ธมฺมา โนจิตฺตสฺส ธมฺมสฺส อารมฺมณ\-ปจฺจเยน ปจฺจโย. สหชาต\-ปจฺจเยน ปจฺจเยน ปจฺจโย. อุปนิสฺสย\-ปจฺจเยน ปจฺจโย. ปจฺฉาชาต\-ปจฺจเยน ปจฺจโย. (2)}

\prose{\csromanfolio{403}จิตฺโต จ โนจิตฺโต จ ธมฺมา จิตฺตสฺส จ โนจิตฺตสฺส จ ธมฺมสฺส อารมฺมณ\-ปจฺจเยน ปจฺจโย. อุปนิสฺสย\-ปจฺจเยน ปจฺจโย. (3)}

\csromanmarkh{2. ปจฺจย\-ปจฺจนีย}\csromanmarkhii{2. สงฺขฺยาวาร}\csromanheadernomark{2}{2. ปจฺจย\-ปจฺจนีย 2. สงฺขฺยาวาร}
\vspace{\tipitakaparskip}
\csromancenter{นเหตุยา นว, นอารมฺมเณ นว, (สพฺพตฺถ นว.) โนอวิคเต นว.}
\csromancenter{\textbf{เหตุ}ปจฺจยา นอารมฺมเณ ตีณิ ฯเปฯ นสมนนฺตเร ตีณิ, นอญฺญมญฺเญ เอกํ, นฺเอาปนิสฺสเย ตีณิ, (สพฺพตฺถ ตีณิ.) นมคฺเค ตีณิ, นสมฺปยุตฺเต เอกํ, นวิปฺปยุตฺเต ตีณิ, โนนตฺถิยา ตีณิ, โนวิคเต ตีณิ.}
\csromanheader{2}{91. นเหตุปจฺจยา อารมฺมเณ นว, อธิปติยา นว. (อนุโลมมาติกา กาตพฺพา.)}
\nitthitam{จิตฺตทุกํ นิฏฺฐิตํ.}
\csromanmatikaanchor{mtk.56}
\csromantocmarkref{subsection}{57. เจตสิกทุก}{mtk.56}
\bookmark[dest={mtk.56},level=2]{57. เจตสิกทุก}
\csromanmarkh{57. เจตสิกทุก}\csromanmarkhii{1. ปฏิจฺจวาร}\csromanheadernomark{2}{57. เจตสิกทุก 1. ปฏิจฺจวาร}
\csromanmarkh{1. ปจฺจยานุโลม}\csromanmarkhii{1. วิภงฺควาร}\csromanheadernomark{2}{1. \textbf{ปจฺจยานุโลม 1. วิภงฺควาร}}
\vspace{\tipitakaparskip}
\csromancenter{\textbf{เหตุ}}
\proseitem{92}{เจตสิกํ ธมฺมํ ปฏิจฺจ เจตสิโก ธมฺโม อุปฺปชฺชติ เหตุปจฺจยา. เจตสิกํ เอกํ ขนฺธํ ปฏิจฺจ เทฺว ขนฺธา. เทฺว ขนฺเธ ปฏิจฺจ เอโก ขนฺโธ. ปฏิสนฺธิ\-กฺขเณ ฯเปฯ. (1)}

\prose{เจตสิกํ ธมฺมํ ปฏิจฺจ อเจตสิโก ธมฺโม อุปฺปชฺชติ เหตุปจฺจยา. เจตสิเก ขนฺเธ ปฏิจฺจ จิตฺตญฺจ จิตฺต\-สมุฏฺฐานญฺจ รูปํ. ปฏิสนฺธิ\-กฺขเณ เจตสิเก ขนฺเธ ปฏิจฺจ จิตฺตญฺจ กฏตฺตา จ รูปํ. (2)}

\prose{\csromanfolio{404}เจตสิกํ ธมฺมํ ปฏิจฺจ เจตสิโก จ อเจตสิโก จ ธมฺมา อุปฺปชฺชนฺติ เหตุปจฺจยา. เจตสิกํ เอกํ ขนฺธํ ปฏิจฺจ เทฺว ขนฺธา จิตฺตญฺจ จิตฺต\-สมุฏฺฐานญฺจ รูปํ. เทฺว ขนฺเธ ฯเปฯ. ปฏิสนฺธิ\-กฺขเณ ฯเปฯ. (3)}

\proseitem{93}{อเจตสิกํ ธมฺมํ ปฏิจฺจ อเจตสิโก ธมฺโม อุปฺปชฺชติ เหตุปจฺจยา. จิตฺตํ ปฏิจฺจ จิตฺต\-สมุฏฺฐานํ รูปํ. ปฏิสนฺธิ\-กฺขเณ จิตฺตํ ปฏิจฺจ กฏตฺตารูปํ. จิตฺตํ ปฏิจฺจ วตฺถุ, วตฺถุํ ปฏิจฺจ จิตฺตํ. เอกํ มหาภูตํ ฯเปฯ มหาภูเต ปฏิจฺจ จิตฺต\-สมุฏฺฐานํ รูปํ, กฏตฺตารูปํ, อุปาทารูปํ. (1)}

\prose{อเจตสิกํ ธมฺมํ ปฏิจฺจ เจตสิโก ธมฺโม อุปฺปชฺชติ เหตุปจฺจยา, จิตฺตํ ปฏิจฺจ สมฺปยุตฺตกา ขนฺธา. ปฏิสนฺธิ\-กฺขเณ จิตฺตํ ฯเปฯ. ปฏิสนฺธิ\-กฺขเณ วตฺถุํ ปฏิจฺจ เจตสิกา ขนฺธา. (2)}

\prose{อเจตสิกํ ธมฺมํ ปฏิจฺจ เจตสิโก จ อเจตสิโก จ ธมฺมา อุปฺปชฺชนฺติ เหตุปจฺจยา. จิตฺตํ ปฏิจฺจ สมฺปยุตฺตกา ขนฺธา จิตฺต\-สมุฏฺฐานญฺจ รูปํ. ปฏิสนฺธิ\-กฺขเณ จิตฺตํ ฯเปฯ. ปฏิสนฺธิ\-กฺขเณ วตฺถุํ ปฏิจฺจ จิตฺตญฺจ สมฺปยุตฺตกา จ ขนฺธา. (3)}

\proseitem{94}{เจตสิกญฺจ อเจตสิกญฺจ ธมฺมํ ปฏิจฺจ เจตสิโก ธมฺโม อุปฺปชฺชติ เหตุปจฺจยา. เจตสิกํ เอกํ ขนฺธญฺจ จิตฺตญฺจ ปฏิจฺจ เทฺว ขนฺธา. เทฺว ขนฺเธ จ ฯเปฯ. ปฏิสนฺธิ\-กฺขเณ เจตสิกํ เอกํ ขนฺธญฺจ จิตฺตญฺจ ปฏิจฺจ เทฺว ขนฺธา. เทฺว ขนฺเธ จ ฯเปฯ. ปฏิสนฺธิ\-กฺขเณ เจตสิกํ เอกํ ขนฺธญฺจ วตฺถุญฺจ ปฏิจฺจ เทฺว ขนฺธา. เทฺว ขนฺเธ จ ฯเปฯ. (1)}

\prose{เจตสิกญฺจ อเจตสิกญฺจ ธมฺมํ ปฏิจฺจ อเจตสิโก ธมฺโม อุปฺปชฺชติ เหตุปจฺจยา. เจตสิเก ขนฺเธ จ จิตฺตญฺจ ปฏิจฺจ จิตฺต\-สมุฏฺฐานํ รูปํ. เจตสิเก ขนฺเธ จ มหาภูเต จ ปฏิจฺจ จิตฺต\-สมุฏฺฐานํ รูปํ. ปฏิสนฺธิ\-กฺขเณ เจตสิเก ขนฺเธ จ จิตฺตญฺจ ปฏิจฺจ กฏตฺตารูปํ. ปฏิสนฺธิ\-กฺขเณ เจตสิเก ขนฺเธ จ มหาภูเต จ ปฏิจฺจ กฏตฺตารูปํ. ปฏิสนฺธิ\-กฺขเณ เจตสิเก ขนฺเธ จ วตฺถุญฺจ ปฏิจฺจ จิตฺตํ. (2)}

\prose{เจตสิกญฺจ อเจตสิกญฺจ ธมฺมํ ปฏิจฺจ เจตสิโก จ อเจตสิโก จ ธมฺมา อุปฺปชฺชนฺติ เหตุปจฺจยา. เจตสิกํ เอกํ ขนฺธญฺจ จิตฺตญฺจ ปฏิจฺจ เทฺว ขนฺธา จิตฺต\-สมุฏฺฐานญฺจ รูปํ. เทฺว ขนฺเธ จ ฯเปฯ. ปฏิสนฺธิ\-กฺขเณ เจตสิกํ เอกํ \csromanfolio{405}ขนฺธญฺจ จิตฺตญฺจ ปฏิจฺจ เทฺว ขนฺธา กฏตฺตา จ รูปํ. เทฺว ขนฺเธ จ ฯเปฯ. ปฏิสนฺธิ\-กฺขเณ เจตสิกํ เอกํ ขนฺธญฺจ วตฺถุญฺจ ปฏิจฺจ เทฺว ขนฺธา จิตฺตญฺจ. เทฺว ขนฺเธ จ วตฺถุญฺจ ปฏิจฺจ เอโก ขนฺโธ จิตฺตญฺจ. (3)}

\proseitem{95}{เจตสิกํ ธมฺมํ ปฏิจฺจ เจตสิโก ธมฺโม อุปฺปชฺชติ อารมฺมณ\-ปจฺจยา. เจตสิกํ เอกํ ขนฺธํ ปฏิจฺจ เทฺว ขนฺธา. เทฺว ขนฺเธ ฯเปฯ. ปฏิสนฺธิ\-กฺขเณ ฯเปฯ. (1)}

\prose{เจตสิกํ ธมฺมํ ปฏิจฺจ อเจตสิโก ธมฺโม อุปฺปชฺชติ อารมฺมณ\-ปจฺจยา. เจตสิเก ขนฺเธ ปฏิจฺจ จิตฺตํ. ปฏิสนฺธิ\-กฺขเณ ฯเปฯ.}

\prose{เจตสิกํ ธมฺมํ ปฏิจฺจ เจตสิโก จ อเจตสิโก จ ธมฺมา อุปฺปชฺชนฺติ อารมฺมณ\-ปจฺจยา. เจตสิกํ เอกํ ขนฺธํ ปฏิจฺจ เทฺว ขนฺธา จิตฺตญฺจ. เทฺว ขนฺเธ ฯเปฯ. ปฏิสนฺธิ\-กฺขเณ ฯเปฯ. (3)}

\proseitem{96}{อเจตสิกํ ธมฺมํ ปฏิจฺจ อเจตสิโก ธมฺโม อุปฺปชฺชติ อารมฺมณ\-ปจฺจยา. ปฏิสนฺธิ\-กฺขเณ วตฺถุํ ปฏิจฺจ จิตฺตํ. (1)}

\prose{อเจตสิกํ ธมฺมํ ปฏิจฺจ เจตสิโก ธมฺโม อุปฺปชฺชติ อารมฺมณ\-ปจฺจยา. จิตฺตํ ปฏิจฺจ สมฺปยุตฺตกา ขนฺธา. ปฏิสนฺธิ\-กฺขเณ จิตฺตํ ปฏิจฺจ สมฺปยุตฺตกา ขนฺธา. ปฏิสนฺธิ\-กฺขเณ วตฺถุํ ปฏิจฺจ เจตสิกา ขนฺธา. (2)}

\prose{อเจตสิกํ ธมฺมํ ปฏิจฺจ เจตสิโก จ อเจตสิโก จ ธมฺมา อุปฺปชฺชนฺติ อารมฺมณ\-ปจฺจยา. ปฏิสนฺธิ\-กฺขเณ วตฺถุํ ปฏิจฺจ จิตฺตญฺจ สมฺปยุตฺตกา จ ขนฺธา. (3)}

\prose{เจตสิกญฺจ อเจตสิกญฺจ ธมฺมํ ปฏิจฺจ เจตสิโก ธมฺโม อุปฺปชฺชติ อารมฺมณ\-ปจฺจยา. เจตสิกํ เอกํ ขนฺธญฺจ จิตฺตญฺจ ปฏิจฺจ เทฺว ขนฺธา. เทฺว ขนฺเธ จ ฯเปฯ. ปฏิสนฺธิ\-กฺขเณ เจตสิกํ เอกํ ขนฺธญฺจ จิตฺตญฺจ ปฏิจฺจ เทฺว ขนฺธา. เทฺว ขนฺเธ จ ฯเปฯ. ปฏิสนฺธิ\-กฺขเณ เจตสิกํ เอกํ ขนฺธญฺจ วตฺถุญฺจ ปฏิจฺจ เทฺว ขนฺธา. เทฺว ขนฺเธ จ ฯเปฯ. (1)}

\prose{เจตสิกญฺจ อเจตสิกญฺจ ธมฺมํ ปฏิจฺจ อเจตสิโก ธมฺโม อุปฺปชฺชติ อารมฺมณ\-ปจฺจยา. ปฏิสนฺธิ\-กฺขเณ เจตสิเก ขนฺเธ จ วตฺถุญฺจ ปฏิจฺจ จิตฺตํ. (2)}

\prose{\csromanfolio{406}เจตสิกญฺจ อเจตสิกญฺจ ธมฺมํ ปฏิจฺจ เจตสิโก จ อเจตสิโก จ ธมฺมา อุปฺปชฺชนฺติ อารมฺมณ\-ปจฺจยา. ปฏิสนฺธิ\-กฺขเณ เจตสิกํ เอกํ ขนฺธญฺจ วตฺถุญฺจ ปฏิจฺจ เทฺว ขนฺธา จิตฺตญฺจ. เทฺว ขนฺเธ จ ฯเปฯ. (3)}

\proseitem{97}{เจตสิกํ ธมฺมํ ปฏิจฺจ เจตสิโก ธมฺโม อุปฺปชฺชติ อธิปติ\-ปจฺจยา. (สํขิตฺตํ.)}

\csromanmarkh{1. ปจฺจยานุโลม}\csromanmarkhii{2. สงฺขฺยาวาร}\csromanheadernomark{2}{1. \textbf{ปจฺจยานุโลม 2. สงฺขฺยาวาร}}
\proseitem{98}{เหตุยา นว, อารมฺมเณ นว, อธิปติยา นว, อนนฺตเร นว, สมนนฺตเร นว, สหชาเต นว, อญฺญมญฺเญ นว, นิสฺสเย นว, อุปนิสฺสเย นว, ปุเรชาเต ปญฺจ, อาเสวเน ปญฺจ, กมฺเม นว, (สพฺพตฺถ นว.) อวิคเต นว.}

\csromanmarkh{2. ปจฺจย\-ปจฺจนีย}\csromanmarkhii{1. วิภงฺควาร}\csromanheadernomark{2}{2. ปจฺจย\-ปจฺจนีย 1. วิภงฺควาร}
\proseitem{99}{เจตสิกํ ธมฺมํ ปฏิจฺจ เจตสิโก ธมฺโม อุปฺปชฺชติ นเหตุปจฺจยา. อเหตุกํ เจตสิกํ เอกํ ขนฺธํ ปฏิจฺจ เทฺว ขนฺธา. เทฺว ขนฺเธ ฯเปฯ. อเหตุก\-ปฏิสนฺธิ\-กฺขเณ ฯเปฯ. วิจิกิจฺฉา\-สหคเต อุทฺธจฺจ\-สหคเต ขนฺเธ ปฏิจฺจ วิจิกิจฺฉา\-สหคโต อุทฺธจฺจ\-สหคโต โมโห. (1)}

\prose{เจตสิกํ ธมฺมํ ปฏิจฺจ อเจตสิโก ธมฺโม อุปฺปชฺชติ นเหตุปจฺจยา. อเหตุเก เจตสิเก ขนฺเธ ปฏิจฺจ จิตฺตญฺจ จิตฺต\-สมุฏฺฐานญฺจ รูปํ. อเหตุก\-ปฏิสนฺธิ\-กฺขเณ ฯเปฯ. (2)}

\prose{เจตสิกํ ธมฺมํ ปฏิจฺจ เจตสิโก จ อเจตสิโก จ ธมฺมา อุปฺปชฺชนฺติ นเหตุปจฺจยา. อเหตุกํ เจตสิกํ เอกํ ขนฺธํ ปฏิจฺจ เทฺว ขนฺธา จิตฺตญฺจ จิตฺต\-สมุฏฺฐานญฺจ รูปํ. เทฺว ขนฺเธ ฯเปฯ. อเหตุก\-ปฏิสนฺธิ\-กฺขเณ ฯเปฯ. (3)}

\proseitem{100}{\csromanfolio{407}อเจตสิกํ ธมฺมํ ปฏิจฺจ อเจตสิโก ธมฺโม อุปฺปชฺชติ นเหตุปจฺจยา. อเหตุกํ จิตฺตํ ปฏิจฺจ จิตฺต\-สมุฏฺฐานํ รูปํ. อเหตุก\-ปฏิสนฺธิ\-กฺขเณ จิตฺตํ ปฏิจฺจ กฏตฺตารูปํ. จิตฺตํ ปฏิจฺจ วตฺถุ, วตฺถุํ ปฏิจฺจ จิตฺตํ. เอกํ มหาภูตํ ฯเปฯ. อสญฺญสตฺตานํ เอกํ มหาภูตํ ฯเปฯ. (1)}

\prose{อเจตสิกํ ธมฺมํ ปฏิจฺจ เจตสิโก ธมฺโม อุปฺปชฺชติ นเหตุปจฺจยา. อเหตุกํ จิตฺตํ ปฏิจฺจ สมฺปยุตฺตกํ ขนฺธา. อเหตุก\-ปฏิสนฺธิ\-กฺขเณ จิตฺตํ ปฏิจฺจ สมฺปยุตฺตกา ขนฺธา. อเหตุก\-ปฏิสนฺธิ\-กฺขเณ วตฺถุํ ปฏิจฺจ เจตสิกา ขนฺธา. วิจิกิจฺฉา\-สหคตํ อุทฺธจฺจ\-สหคตํ จิตฺตํ ปฏิจฺจ วิจิกิจฺฉา\-สหคโต อุทฺธจฺจ\-สหคโต โมโห. (2)}

\prose{อเจตสิกํ ธมฺมํ ปฏิจฺจ เจตสิโก จ อเจตสิโก จ ธมฺมา อุปฺปชฺชนฺติ นเหตุปจฺจยา. อเหตุกํ จิตฺตํ ปฏิจฺจ สมฺปยุตฺตกา ขนฺธา จิตฺต\-สมุฏฺฐานญฺจ รูปํ. อเหตุก\-ปฏิสนฺธิ\-กฺขเณ จิตฺตํ ปฏิจฺจ สมฺปยุตฺตกา ขนฺธา กฏตฺตา จ รูปํ. อเหตุก\-ปฏิสนฺธิ\-กฺขเณ วตฺถุํ ปฏิจฺจ จิตฺตญฺจ สมฺปยุตฺตกา จ ขนฺธา. (3)}

\proseitem{101}{เจตสิกญฺจ อเจตสิกญฺจ ธมฺมํ ปฏิจฺจ เจตสิโก ธมฺโม อุปฺปชฺชติ นเหตุปจฺจยา. อเหตุกํ เจตสิกํ เอกํ ขนฺธญฺจ จิตฺตญฺจ ปฏิจฺจ เทฺว ขนฺธา. เทฺว ขนฺเธ จ ฯเปฯ. อเหตุก\-ปฏิสนฺธิ\-กฺขเณ เจตสิกํ เอกํ ขนฺธญฺจ จิตฺตญฺจ ปฏิจฺจ เทฺว ขนฺธา. เทฺว ขนฺเธ จ ฯเปฯ. อเหตุก\-ปฏิสนฺธิ\-กฺขเณ เจตสิกํ เอกํ ขนฺธญฺจ วตฺถุญฺจ ปฏิจฺจ เทฺว ขนฺธา. เทฺว ขนฺเธ จ ฯเปฯ. วิจิกิจฺฉา\-สหคเต อุทฺธจฺจ\-สหคเต ขนฺเธ จ จิตฺตญฺจ ปฏิจฺจ วิจิกิจฺฉา\-สหคโต อุทฺธจฺจ\-สหคโต โมโห. (1)}

\prose{เจตสิกญฺจ อเจตสิกญฺจ ธมฺมํ ปฏิจฺจ อเจตสิโก ธมฺโม อุปฺปชฺชติ นเหตุปจฺจยา. อเหตุเก เจตสิเก ขนฺเธ จ จิตฺตญฺจ ปฏิจฺจ จิตฺต\-สมุฏฺฐานํ รูปํ. อเหตุเก เจตสิเก ขนฺเธ จ มหาภูเต จ ปฏิจฺจ จิตฺต\-สมุฏฺฐานํ รูปํ. อเหตุก\-ปฏิสนฺธิ\-กฺขเณ เจตสิเก ขนฺเธ จ จิตฺตญฺจ ปฏิจฺจ กฏตฺตารูปํ. อเหตุก\-ปฏิสนฺธิ\-กฺขเณ เจตสิเก ขนฺเธ จ มหาภูเต จ ปฏิจฺจ กฏตฺตารูปํ. อเหตุก\-ปฏิสนฺธิ\-กฺขเณ เจตสิเก ขนฺเธ จ วตฺถุญฺจ ปฏิจฺจ จิตฺตํ. (2)}

\prose{เจตสิกญฺจ อเจตสิกญฺจ ธมฺมํ ปฏิจฺจ เจตสิโก จ อเจตสิโก จ ธมฺมา อุปฺปชฺชนฺติ นเหตุปจฺจยา. อเหตุกํ เจตสิกํ เอกํ ขนฺธญฺจ จิตฺตญฺจ \csromanfolio{408}ปฏิจฺจ เทฺว ขนฺธา จิตฺต\-สมุฏฺฐานญฺจ รูปํ. อเหตุก\-ปฏิสนฺธิ\-กฺขเณ เจตสิกํ เอกํ ขนฺธญฺจ จิตฺตญฺจ ปฏิจฺจ เทฺว ขนฺธา กฏตฺตา จ รูปํ. อเหตุก\-ปฏิสนฺธิ\-กฺขเณ เจตสิกํ เอกํ ขนฺธญฺจ วตฺถุญฺจ ปฏิจฺจ เทฺว ขนฺธา จิตฺตญฺจ. เทฺว ขนฺเธ จ ฯเปฯ. (3)}

\prose{นอารมฺมณ}

\proseitem{102}{เจตสิกํ ธมฺมํ ปฏิจฺจ อเจตสิโก ธมฺโม อุปฺปชฺชติ นอารมฺมณ\-ปจฺจยา. เจตสิเก ขนฺเธ ปฏิจฺจ จิตฺต\-สมุฏฺฐานํ รูปํ. ปฏิสนฺธิ\-กฺขเณ ฯเปฯ. (1)}

\prose{อเจตสิกํ ธมฺมํ ปฏิจฺจ อเจตสิโก ธมฺโม อุปฺปชฺชติ นอารมฺมณ\-ปจฺจยา. จิตฺตํ ปฏิจฺจ จิตฺต\-สมุฏฺฐานํ รูปํ. ปฏิสนฺธิ\-กฺขเณ ฯเปฯ. (ยาว อสญฺญสตฺตา.) (1)}

\prose{เจตสิกญฺจ อเจตสิกญฺจ ธมฺมํ ปฏิจฺจ อเจตสิโก ธมฺโม อุปฺปชฺชติ นอารมฺมณ\-ปจฺจยา. เจตสิเก ขนฺเธ จ จิตฺตญฺจ ปฏิจฺจ จิตฺต\-สมุฏฺฐานํ รูปํ. เจตสิเก ขนฺเธ จ มหาภูเต จ ปฏิจฺจ จิตฺต\-สมุฏฺฐานํ รูปํ. (ปฏิสนฺธิกฺขเณ เทฺวปิ กาตพฺพา, สํขิตฺตํ.)}

\vspace{\tipitakaparskip}
\csromancenter{(ปฏิสนฺธิ\-กฺขเณ เทฺวปิ กาตพฺพา, สํขิตฺตํ.)}
\csromanmarkh{2. ปจฺจย\-ปจฺจนีย}\csromanmarkhii{2. สงฺขฺยาวาร}\csromanheadernomark{2}{2. ปจฺจย\-ปจฺจนีย 2. สงฺขฺยาวาร}
\proseitem{103}{นเหตุยา นว, นอารมฺมเณ ตีณิ, นอธิปติยา นว, นอนนฺตเร ตีณิ, นสมนนฺตเร ตีณิ, นอญฺญมญฺเญ ตีณิ, นอุปนิสฺสเย ตีณิ, นปุเรชาเต นว, นปจฺฉาชาเต นว, นอาเสวเน นว, นกมฺเม จตฺตาริ, นวิปาเก นว, นอาหาเร เอกํ, นอินฺทฺริเย เอกํ, นฌาเน ฉ, นมคฺเค นว, นสมฺปยุตฺเต ตีณิ, นวิปฺปยุตฺเต ฉ, โนนตฺถิยา ตีณิ, โนวิคเต ตีณิ.}

\proseitem{104}{เหตุปจฺจยา นอารมฺมเณ ตีณิ, นอธิปติยา นว. (สํขิตฺตํ.)}

\csromanmarkh{57. เจตสิกทุก}\csromanmarkhii{4. ปจฺจย\-ปจฺจนียา\-นุโลม}\csromanheadernomark{2}{57. เจตสิกทุก 4. ปจฺจย\-ปจฺจนียา\-นุโลม}
\proseitem{105}{\csromanfolio{409}นเหตุปจฺจยา อารมฺมเณ นว, อนนฺตเร นว ฯเปฯ ปุเรชาเต ปญฺจ, อาเสวเน ปญฺจ, กมฺเม นว, (สพฺพตฺถ นว.) มคฺเค ตีณิ, อวิคเต นว.}

\vspace{\tipitakaparskip}
\csromancenter{(สหชาตวาโรปิ ปฏิจฺจ\-วาร\-สทิโส.)}
\csromanmarkh{57. เจตสิกทุก}\csromanmarkhii{3. ปจฺจยวาร}\csromanheadernomark{2}{57. เจตสิกทุก 3. ปจฺจยวาร}
\csromanmarkh{1. ปจฺจยานุโลม}\csromanmarkhii{1. วิภงฺควาร}\csromanheadernomark{2}{1. \textbf{ปจฺจยานุโลม 1. วิภงฺควาร}}
\csromanheader{2}{\textbf{เหตุ}}
\proseitem{106}{เจตสิกํ ธมฺมํ ปจฺจยา เจตสิโก ธมฺโม อุปฺปชฺชติ เหตุปจฺจยา. ตีณิ. (ปฏิจฺจสทิสา.)}

\prose{อเจตสิกํ ธมฺมํ ปจฺจยา อเจตสิโก ธมฺโม อุปฺปชฺชติ เหตุปจฺจยา. จิตฺตํ ปจฺจยา จิตฺต\-สมุฏฺฐานํ รูปํ. วตฺถุํ ปจฺจยา จิตฺตํ. ปฏิสนฺธิ\-กฺขเณ จิตฺตํ ปจฺจยา กฏตฺตารูปํ. จิตฺตํ ปจฺจยา วตฺถุ, วตฺถุํ ปจฺจยา จิตฺตํ. เอกํ มหาภูตํ ฯเปฯ. (1)}

\prose{อเจตสิกํ ธมฺมํ ปจฺจยา เจตสิโก ธมฺโม อุปฺปชฺชติ เหตุปจฺจยา. จิตฺตํ ปจฺจยา สมฺปยุตฺตกา ขนฺธา, วตฺถุํ ปจฺจยา เจตสิกา ขนฺธา. (ปฏิสนฺธิ\-กฺขเณ เทฺวปิ กาตพฺพา.) (2)}

\prose{อเจตสิกํ ธมฺมํ ปจฺจยา เจตสิโก จ อเจตสิโก จ ธมฺมา อุปฺปชฺชนฺติ เหตุปจฺจยา. จิตฺตํ ปจฺจยา สมฺปยุตฺตกา ขนฺธา จิตฺต\-สมุฏฺฐานญฺจ รูปํ. วตฺถุํ ปจฺจยา จิตฺตญฺจ สมฺปยุตฺตกา จ ขนฺธา. (ปฏิสนฺธิ\-กฺขเณ เทฺวปิ กาตพฺพา.) (3)}

\proseitem{107}{เจตสิกญฺจ อเจตสิกญฺจ ธมฺมํ ปจฺจยา เจตสิโก ธมฺโม อุปฺปชฺชติ เหตุปจฺจยา. เจตสิกํ เอกํ ขนฺธญฺจ จิตฺตญฺจ ปจฺจยา เทฺว ขนฺธา. เทฺว ขนฺเธ จ ฯเปฯ. เจตสิกํ เอกํ ขนฺธญฺจ วตฺถุญฺจ ปจฺจยา เทฺว ขนฺธา. เทฺว ขนฺเธ จ ฯเปฯ. (ปฏิสนฺธิ\-กฺขเณ เทฺวปิ กาตพฺพา.) (1)}

\prose{\csromanfolio{410}เจตสิกญฺจ อเจตสิกญฺจ ธมฺมํ ปจฺจยา อเจตสิโก ธมฺโม อุปฺปชฺชติ เหตุปจฺจยา. เจตสิเก ขนฺเธ จ จิตฺตญฺจ ปจฺจยา จิตฺต\-สมุฏฺฐานํ รูปํ. เจตสิเก ขนฺเธ จ มหาภูเต จ ปจฺจยา จิตฺต\-สมุฏฺฐานํ รูปํ. เจตสิเก ขนฺเธ จ วตฺถุญฺจ ปจฺจยา จิตฺตํ. (ปฏิสนฺธิ\-กฺขเณ ตีณิปิ กาตพฺพา.) (2)}

\prose{เจตสิกญฺจ อเจตสิกญฺจ ธมฺมํ ปจฺจยา เจตสิโก จ อเจตสิโก จ ธมฺมา อุปฺปชฺชนฺติ เหตุปจฺจยา. เจตสิกํ เอกํ ขนฺธญฺจ จิตฺตญฺจ ปจฺจยา เทฺว ขนฺธา จิตฺต\-สมุฏฺฐานญฺจ รูปํ. เทฺว ขนฺเธ จ ฯเปฯ. เจตสิกํ เอกํ ขนฺธญฺจ วตฺถุญฺจ ปจฺจยา เทฺว ขนฺธา จิตฺตญฺจ. เทฺว ขนฺเธ จ ฯเปฯ. (ปฏิสนฺธิ\-กฺขเณ เทฺวปิ กาตพฺพา.) (3)}

\csromanheader{2}{\textbf{อรมฺมณ}}
\proseitem{108}{เจตสิกํ ธมฺมํ ปจฺจยา เจตสิโก ธมฺโม อุปฺปชฺชติ อารมฺมณ\-ปจฺจยา. ตีณิ. (ปฏิจฺจสทิสา.)}

\prose{อเจตสิกํ ธมฺมํ ปจฺจยา อเจตสิโก ธมฺโม อุปฺปชฺชติ อารมฺมณ\-ปจฺจยา. จกฺขายตนํ ปจฺจยา จกฺขุวิญฺญาณํ ฯเปฯ. กายายตนํ ปจฺจยา กายวิญฺญาณํ. วตฺถุํ ปจฺจยา จิตฺตํ. ปฏิสนฺธิ\-กฺขเณ ฯเปฯ. (1)}

\prose{อเจตสิกํ ธมฺมํ ปจฺจยา เจตสิโก ธมฺโม อุปฺปชฺชติ อารมฺมณ\-ปจฺจยา. จกฺขายตนํ ปจฺจยา จกฺขุ\-วิญฺญาณสหคตา ขนฺธา ฯเปฯ. กายายตนํ ปจฺจยา ฯเปฯ. จิตฺตํ ปจฺจยา สมฺปยุตฺตกา ขนฺธา. วตฺถุํ ปจฺจยา เจตสิกา ขนฺธา. (ปฏิสนฺธิ\-กฺขเณ เทฺวปิ กาตพฺพา.) (2)}

\prose{อเจตสิกํ ธมฺมํ ปจฺจยา เจตสิโก จ อเจตสิโก จ ธมฺมา อุปฺปชฺชนฺติ อารมฺมณ\-ปจฺจยา. จกฺขายตนํ ปจฺจยา จกฺขุวิญฺญาณํ สมฺปยุตฺตกา จ ขนฺธา ฯเปฯ. กายายตนํ ปจฺจยา. วตฺถุํ ปจฺจยา จิตฺตํ สมฺปยุตฺตกา จ ขนฺธา. (ปฏิสนฺธิ\-กฺขเณ เอกํ.) (3)}

\proseitem{109}{เจตสิกญฺจ อเจตสิกญฺจ ธมฺมํ ปจฺจยา เจตสิโก ธมฺโม อุปฺปชฺชติ อารมฺมณ\-ปจฺจยา. จกฺขุ\-วิญฺญาณสหคตํ เอกํ ขนฺธญฺจ จกฺขุวิญฺญาณญฺจ ปจฺจยา เทฺว ขนฺธา. เทฺว ขนฺเธ จ ฯเปฯ. จกฺขุ\-วิญฺญาณสหคตํ เอกํ ขนฺธญฺจ จกฺขายตนญฺจ ปจฺจยา เทฺว ขนฺธา. เทฺว ขนฺเธ จ ฯเปฯ. กาย\-วิญฺญาณสหคตํ ฯเปฯ. เจตสิกํ เอกํ ขนฺธญฺจ จิตฺตญฺจ ปจฺจยา เทฺว ขนฺธา. เทฺว ขนฺเธ จ ฯเปฯ. เจตสิกํ เอกํ ขนฺธญฺจ วตฺถุญฺจ ปจฺจยา เทฺว ขนฺธา. เทฺว ขนฺเธ จ ฯเปฯ. (ปฏิสนฺธิ\-กฺขเณ เทฺว.) (1)}

\prose{\csromanfolio{411}เจตสิกญฺจ อเจตสิกญฺจ ธมฺมํ ปจฺจยา อเจตสิโก ธมฺโม อุปฺปชฺชติ อารมฺมณ\-ปจฺจยา. จกฺขุ\-วิญฺญาณสหคเต ขนฺเธ จ จกฺขายตนญฺจ ปจฺจยา จกฺขุวิญฺญาณํ ฯเปฯ. กาย\-วิญฺญาณสหคเต ฯเปฯ. เจตสิเก ขนฺเธ จ วตฺถุญฺจ ปจฺจยา จิตฺตํ. (ปฏิสนฺธิ\-กฺขเณ เอกํ.)}

\prose{เจตสิกญฺจ อเจตสิกญฺจ ธมฺมํ ปจฺจยา เจตสิโก จ อเจตสิโก จ ธมฺมา อุปฺปชฺชนฺติ อารมฺมณ\-ปจฺจยา. จกฺขุ\-วิญฺญาณสหคตํ เอกํ ขนฺธญฺจ จกฺขายตนญฺจ ปจฺจยา เทฺว ขนฺธา จกฺขุวิญฺญาณญฺจ. เทฺว ขนฺเธ จ ฯเปฯ. เจตสิกํ เอกํ ขนฺธญฺจ วตฺถุญฺจ ปจฺจยา เทฺว ขนฺธา จิตฺตญฺจ. เทฺว ขนฺเธ จ ฯเปฯ. (ปฏิสนฺธิ\-กฺขเณ เอกํ. สํ ขิตฺตํ.) (3)}

\csromanmarkh{1. ปจฺจยานุโลม}\csromanmarkhii{2. สงฺขฺยาวาร}\csromanheadernomark{2}{1. \textbf{ปจฺจยานุโลม 2. สงฺขฺยาวาร}}
\proseitem{110}{เหตุยา นว, อารมฺมเณ นว, อธิปติยา นว, (สพฺพตฺถ นว.) ปุเรชาเต นว, อาเสวเน นว ฯเปฯ อวิคเต นว.}

\csromanmarkh{2. ปจฺจย\-ปจฺจนีย}\csromanmarkhii{1. วิภงฺควาร}\csromanheadernomark{2}{2. ปจฺจย\-ปจฺจนีย 1. วิภงฺควาร}
\proseitem{111}{เจตสิกํ ธมฺมํ ปจฺจยา เจตสิโก ธมฺโม อุปฺปชฺชติ นเหตุปจฺจยา. อเหตุกํ เจตสิกํ. (สํขิตฺตํ.)}

\prose{(นว ปญฺหา, ปญฺจวิญฺญาณมฺปิ ยถา อารมฺมณ\-ปจฺจยา เอวํ กาตพฺพํ. ตีสุเยว โมโห. สพฺเพ ปญฺหา ปวตฺติ\-ปฏิสนฺธิยา กาตพฺพา อสมฺโมหนฺเตน.)}

\csromanmarkh{2. ปจฺจย\-ปจฺจนีย}\csromanmarkhii{2. สงฺขฺยาวาร}\csromanheadernomark{2}{2. ปจฺจย\-ปจฺจนีย 2. สงฺขฺยาวาร}
\vspace{\tipitakaparskip}
\csromancenter{\textbf{นเหตุ}ยา นว, นอารมฺมเณ ตีณิ, นอธิปติยา นว, นอนนฺตเร ตีณิ, นสมนนฺตเร ตีณิ, นอญฺญมญฺเญ ตีณิ, นฺเอาปนิสฺสเย ตีณิ, นปุเรชาเต นว, นปจฺฉาชาเต นว, นอาเสวเน นว, นกมฺเม}
\prosecont{\csromanfolio{412}จตฺตาริ, นวิปาเก นว, นอาหาเร เอกํ, นอินฺทฺริเย เอกํ, นฌาเน นว, นมคฺเค นว, นสมฺปยุตฺเต ตีณิ, นวิปฺปยุตฺเต ฉ, โนนตฺถิยา ตีณิ, โนวิคเต ตีณิ.}

\vspace{\tipitakaparskip}
\csromancenter{\textbf{เหตุ}ปจฺจยา นอารมฺมเณ ตีณิ, นอธิปติยา นว. (สํขิตฺตํ.)}
\proseitem{114}{นเหตุปจฺจยา อารมฺมเณ นว, อนนฺตเร นว, (สพฺพตฺถ นว.) มคฺเค ตีณิ ฯเปฯ อวิคเต นว.}

\vspace{\tipitakaparskip}
\csromancenter{(นิสฺสยวาโร ปจฺจยวาร\-สทิโส.)}
\csromanmarkh{57. เจตสิกทุก}\csromanmarkhii{5. สํสฏฺฐวาร}\csromanheadernomark{2}{57. \textbf{เจตสิกทุก 5. สํสฏฺฐวาร}}
\csromanmarkh{1. ปจฺจยานุโลม}\csromanmarkhii{1. วิภงฺควาร}\csromanheadernomark{2}{1. \textbf{ปจฺจยานุโลม 1. วิภงฺควาร}}
\vspace{\tipitakaparskip}
\csromancenter{\textbf{เหตุ}}
\proseitem{115}{เจตสิกํ ธมฺมํ สํสฏฺโฐ เจตสิโก ธมฺโม อุปฺปชฺชติ เหตุปจฺจยา. เจตสิกํ เอกํ ขนฺธํ สํสฏฺฐา เทฺว ขนฺธา. เทฺว ขนฺเธ ฯเปฯ. ปฏิสนฺธิ\-กฺขเณ ฯเปฯ. (1)}

\prose{เจตสิกํ ธมฺมํ สํสฏฺโฐ อเจตสิโก ธมฺโม อุปฺปชฺชติ เหตุปจฺจยา. เจตสิเก ขนฺเธ สํสฏฺฐํ จิตฺตํ. ปฏิสนฺธิ\-กฺขเณ ฯเปฯ.}

\prose{เจตสิกํ ธมฺมํ สํสฏฺโฐ เจตสิโก จ อเจตสิโก จ ธมฺมา อุปฺปชฺชนฺติ เหตุปจฺจยา. เจตสิกํ เอกํ ขนฺธํ สํสฏฺฐา เทฺว ขนฺธา จิตฺตญฺจ. เทฺว ขนฺเธ ฯเปฯ. ปฏิสนฺธิ\-กฺขเณ ฯเปฯ. (3)}

\prose{อเจตสิกํ ธมฺมํ สํสฏฺโฐ เจตสิโก ธมฺโม อุปฺปชฺชติ เหตุปจฺจยา. จิตฺตํ สํสฏฺฐา สมฺปยุตฺตกา ขนฺธา. ปฏิสนฺธิ\-กฺขเณ ฯเปฯ. (1)}

\prose{\csromanfolio{413}เจตสิกญฺจ อเจตสิกญฺจ ธมฺมํ สํสฏฺโฐ เจตสิโก ธมฺโม อุปฺปชฺชติ เหตุปจฺจยา. เจตสิกํ เอกํ ขนฺธญฺจ จิตฺตญฺจ สํสฏฺฐา เทฺว ขนฺธา. เทฺว ขนฺเธ จ ฯเปฯ. ปฏิสนฺธิ\-กฺขเณ ฯเปฯ. (1) (สํขิตฺตํ.)}

\csromanmarkh{1. ปจฺจยานุโลม}\csromanmarkhii{2. สงฺขฺยาวาร}\csromanheadernomark{2}{1. \textbf{ปจฺจยานุโลม 2. สงฺขฺยาวาร}}
\proseitem{116}{เหตุยา ปญฺจ, อารมฺมเณ ปญฺจ, อธิปติยา ปญฺจ, (สพฺพตฺถ ปญฺจ.) อวิคเต ปญฺจ.}

\csromanmarkh{2. ปจฺจย\-ปจฺจนีย}\csromanmarkhii{1. วิภงฺควาร}\csromanheadernomark{2}{2. ปจฺจย\-ปจฺจนีย 1. วิภงฺควาร}
\proseitem{117}{เจตสิกํ ธมฺมํ สํสฏฺโฐ เจตสิโก ธมฺโม อุปฺปชฺชติ นเหตุปจฺจยา. (เอวํ ปญฺจปิ ปญฺหา กาตพฺพา, ตีณิเยว โมโห. สํขิตฺตํ.)}

\csromanmarkh{2. ปจฺจย\-ปจฺจนีย}\csromanmarkhii{2. สงฺขฺยาวาร}\csromanheadernomark{2}{2. ปจฺจย\-ปจฺจนีย 2. สงฺขฺยาวาร}
\proseitem{118}{นเหตุยา ปญฺจ, นอธิปติยา ปญฺจ, นปุเรชาเต ปญฺจ, นปจฺฉาชาเต ปญฺจ, นอาเสวเน ปญฺจ, นกมฺเม ตีณิ, นวิปาเก ปญฺจ, นฌาเน ปญฺจ, นมคฺเค ปญฺจ, นวิปฺปยุตฺเต ปญฺจ.}

\vspace{\tipitakaparskip}
\csromancenter{(เอวํ อิตเร เทฺว คณนาปิ สมฺปยุตฺตวาโรปิ กาตพฺโพ.)}
\csromanmarkh{57. เจตสิกทุก}\csromanmarkhii{7. ปญฺหาวาร}\csromanheadernomark{2}{57. เจตสิกทุก 7. ปญฺหาวาร}
\csromanmarkh{1. ปจฺจยานุโลม}\csromanmarkhii{1. วิภงฺควาร}\csromanheadernomark{2}{1. \textbf{ปจฺจยานุโลม 1. วิภงฺควาร}}
\vspace{\tipitakaparskip}
\csromancenter{\textbf{เหตุ}}
\proseitem{119}{เจตสิโก ธมฺโม เจตสิกสฺส ธมฺมสฺส เหตุปจฺจเยน ปจฺจโย. เจตสิกา เหตู สมฺปยุตฺตกานํ ขนฺธานํ เหตุปจฺจเยน ปจฺจโย. ปฏิสนฺธิ\-กฺขเณ ฯเปฯ. (1)}

\prose{เจตสิโก ธมฺโม อเจตสิกสฺส ธมฺมสฺส เหตุปจฺจเยน ปจฺจโย. เจตสิกา เหตู จิตฺตสฺส จิตฺตสมุฏฺฐานานญฺจ รูปานํ เหตุปจฺจเยน ปจฺจโย. ปฏิสนฺธิ\-กฺขเณ ฯเปฯ. (2)}

\prose{\csromanfolio{414}เจตสิโก ธมฺโม เจตสิกสฺส จ อเจตสิกสฺส จ ธมฺมสฺส เหตุปจฺจเยน ปจฺจโย. เจตสิกา เหตู สมฺปยุตฺตกานํ ขนฺธานํ จิตฺตสฺส จ จิตฺตสมุฏฺฐานาญฺจ รูปานํ เหตุปจฺจเยน ปจฺจโย. ปฏิสนฺธิ\-กฺขเณ ฯเปฯ. (3)}

\proseitem{120}{เจตสิโก ธมฺโม เจตสิกสฺส ธมฺมสฺส อารมฺมณ\-ปจฺจเยน ปจฺจโย. เจตสิเก ขนฺเธ อารพฺภ เจตสิกา ขนฺธา อุปฺปชฺชนฺติ. (มูลํ ปุจฺฉิตพฺพํ.) เจตสิเก ขนฺเธ อารพฺภ จิตฺตํ อุปฺปชฺชติ. (มูลํ ปุจฺฉิตพฺพํ.) เจตสิเก ขนฺเธ อารพฺภ เจตสิกา ขนฺธา จ จิตฺตญฺจ อุปฺปชฺชนฺติ. (3)}

\proseitem{121}{อเจตสิโก ธมฺโม อเจตสิกสฺส ธมฺมสฺส อารมฺมณ\-ปจฺจเยน ปจฺจโย. อริยา มคฺคา วุฏฺฐหิตฺวา ฯเปฯ. ผลํ ปจฺจเวกฺขนฺติ, นิพฺพานํ ปจฺจเวกฺขนฺติ. นิพฺพานํ โคตฺรภุสฺส, โวทานสฺส, มคฺคสฺส, ผลสฺส, อาวชฺชนาย อารมฺมณ\-ปจฺจเยน ปจฺจโย. จกฺขุํ ฯเปฯ วตฺถุํ อเจตสิเก ขนฺเธ อนิจฺจโต ฯเปฯ วิปสฺสติ, อสฺสาเทติ อภินนฺทติ, ตํ อารพฺภ จิตฺตํ อุปฺปชฺชติ. ทิพฺเพน จกฺขุนา รูปํ ปสฺสติ. ทิพฺพาย โสตธาตุยา สทฺทํ สุณาติ. เจโตปริย\-ญาเณน อเจตสิก\-จิตฺต\-สมงฺคิสฺส จิตฺตํ ชานาติ. อากาสา\-นญฺจา\-ยตนํ ฯเปฯ. อากิญฺจญฺญา\-ยตนํ ฯเปฯ. รูปายตนํ จกฺขุ\-วิญฺญาณสฺส ฯเปฯ. โผฏฺฐพฺพา\-ยตนํ กายวิญฺญาณสฺส ฯเปฯ. อเจตสิกา ขนฺธา อิทฺธิวิธ\-ญาณสฺส, เจโตปริย\-ญาณสฺส, ปุพฺเพ\-นิวาสา\-นุสฺสติ\-ญาณสฺส, อนาคตํสญาณสฺส, อาวชฺชนาย อารมฺมณ\-ปจฺจเยน ปจฺจโย. (1)}

\prose{อเจตสิโก ธมฺโม เจตสิกสฺส ธมฺมสฺส อารมฺมณ\-ปจฺจเยน ปจฺจโย. อริยา มคฺคา วุฏฺฐหิตฺวา ฯเปฯ. นิพฺพานํ ปจฺจเวกฺขนฺติ. (ปฐม\-คมน\-สทิสํ.) จกฺขุํ ฯเปฯ วตฺถุํ อเจตสิเก ขนฺเธ อนิจฺจโต ฯเปฯ โทมนสฺสํ อุปฺปชฺชติ. ทิพฺเพน จกฺขุนา รูปํ ปสฺสติ. ทิพฺพาย โสตธาตุยา สทฺทํ สุณาติ. เจโตปริย\-ญาเณน อเจตสิก\-จิตฺต\-สมงฺคิสฺส จิตฺตํ ชานาติ. อากาสา\-นญฺจา\-ยตนํ ฯเปฯ. อากิญฺจญฺญา\-ยตนํ เนว\-สญฺญา\-นา\-สญฺญา\-ยตนสฺส ฯเปฯ. รูปายตนํ จกฺขุ\-วิญฺญาณสหคตานํ ขนฺธานํ ฯเปฯ. โผฏฺฐพฺพา\-ยตนํ กาย\-วิญฺญาณสหคตานํ ขนฺธานํ ฯเปฯ. อเจตสิกา ขนฺธา อิทฺธิวิธ\-ญาณสฺส, เจโตปริย\-ญาณสฺส, ปุพฺเพ\-นิวาสา\-นุสฺสติ\-ญาณสฺส, อนาคตํสญาณสฺส, อาวชฺชนาย อารมฺมณ\-ปจฺจเยน ปจฺจโย. (2)}

\prose{\csromanfolio{415}อเจตสิโก ธมฺโม เจตสิกสฺส จ อเจตสิกสฺส จ ธมฺมสฺส อารมฺมณ\-ปจฺจเยน ปจฺจโย. อริยา มคฺคา วุฏฺฐหิตฺวา ฯเปฯ. นิพฺพานํ ปจฺจเวกฺขนฺติ. (ปฐม\-คมน\-สทิสํ.) จกฺขุํ ฯเปฯ วตฺถุํ อเจตสิเก ขนฺเธ อนิจฺจโต ฯเปฯ วิปสฺสติ, อสฺสาเทติ อภินนฺทติ, ตํ อารพฺภ จิตฺตญฺจ สมฺปยุตฺตกา จ ขนฺธา อุปฺปชฺชนฺติ. ทิพฺเพน จกฺขุนา รูปํ ปสฺสติ ฯเปฯ. รูปายตนํ จกฺขุ\-วิญฺญาณสฺส สมฺปยุตฺตกานญฺจ ขนฺธานํ อารมฺมณ\-ปจฺจเยน ปจฺจโย ฯเปฯ. โผฏฺฐพฺพา\-ยตนํ กายวิญฺญาณสฺส สมฺปยุตฺตกานญฺจ ขนฺธานํ ฯเปฯ. อเจตสิกา ขนฺธา อิทฺธิวิธ\-ญาณสฺส, เจโตปริย\-ญาณสฺส, ปุพฺเพ\-นิวาสา\-นุสฺสติ\-ญาณสฺส, อนาคตํสญาณสฺส, อาวชฺชนาย อารมฺมณ\-ปจฺจเยน ปจฺจโย. (3)}

\proseitem{112}{เจตสิโก จ อเจตสิโก จ ธมฺมา เจตสิกสฺส ธมฺมสฺส อารมฺมณ\-ปจฺจเยน ปจฺจโย. เจตสิเก ขนฺเธ จ จิตฺตญฺจ อารพฺภ เจตสิกา ขนฺธา อุปฺปชฺชนฺติ. (มูลํ ปุจฺฉิตพฺพํ.) เจตสิเก ขนฺเธ จ จิตฺตญฺจ อารพฺภ จิตฺตํ อุปฺปชฺชติ. (มูลํ ปุจฺฉิตพฺพํ.) เจตสิเก ขนฺเธ จ จิตฺตญฺจ อารพฺภ เจตสิกา ขนฺธา จ จิตฺตญฺจ อุปฺปชฺชนฺติ. (3)}

\proseitem{123}{เจตสิโก ธมฺโม เจตสิกสฺส ธมฺมสฺส อธิปติ\-ปจฺจเยน ปจฺจโย. อารมฺมณา\-ธิปติ, สหชาตา\-ธิปติ.  อารมฺมณ\-ธิปติ\csromandash{}เจตสิเก ขนฺเธ ครุํ กตฺวา เจตสิกา ขนฺธา อุปฺปชฺชนฺติ. สหชาตา\-ธิปติ\csromandash{} เจตสิกา\-ธิปติ สมฺปยุตฺตกานํ ขนฺธานํ อธิปติ\-ปจฺจเยน ปจฺจโย. (มูลํ ปุจฺฉิตพฺพํ.) อารมฺมณา\-ธิปติ\csromandash{}เจตสิเก ขนฺเธ ครุํ กตฺวา จิตฺตํ อุปฺปชฺชติ. สหชาตา\-ธิปติ\csromandash{}เจตสิกา\-ธิปติ จิตฺตสฺส จิตฺตสมุฏฺฐานานญฺจ รูปานํ อธิปติ\-ปจฺจเยน ปจฺจโย. (มูลํ ปุจฺฉิตพฺพํ.) อารมฺมณา\-ธิปติ\csromandash{}เจตสิเก ขนฺเธ ครุํ กตฺวา เจตสิกา ขนฺธา จ จิตฺตญฺจ อุปฺปชฺชนฺติ. สหชาตา\-ธิปติ\csromandash{}เจตสิกา\-ธิปติ สมฺปยุตฺตกานํ ขนฺธานํ จิตฺตสฺส จ จิตฺตสมุฏฺฐานานญฺจ รูปานํ อธิปติ\-ปจฺจเยน ปจฺจโย. (3)}

\proseitem{124}{อเจตสิโก ธมฺโม อเจตสิกสฺส ธมฺมสฺส อธิปติ\-ปจฺจเยน ปจฺจโย. อารมฺมณา\-ธิปติ, สหชาตา\-ธิปติ.  อารมฺมณา\-ธิปติ\csromandash{}อริยา มคฺคา วุฏฺฐหิตฺวา ฯเปฯ. นิพฺพานํ ครุํ กตฺวา ปจฺจเวกฺขนฺติ. นิพฺพานํ โคตฺรภุสฺส, โวทานสฺส, มคฺคสฺส, ผลสฺส อธิปติ\-ปจฺจเยน ปจฺจโย. \csromanfolio{416}จกฺขุํ ฯเปฯ วตฺถุํ อเจตสิเก ขนฺเธ ครุํ กตฺวา อสฺสาเทติ อภินนฺทติ, ตํ ครุํ กตฺวา จิตฺตํ อุปฺปชฺชติ. สหชาตา\-ธิปติ\csromandash{}อเจตสิกา\-ธิปติ จิตฺต\-สมุฏฺฐานานํ รูปานํ อธิปติ\-ปจฺจเยน ปจฺจโย. (1)}

\prose{อเจตสิโก ธมฺโม เจตสิกสฺส ธมฺมสฺส อธิปติ\-ปจฺจเยน ปจฺจโย. อารมฺมณา\-ธิปติ, สหชาตา\-ธิปติ.  อารมฺมณา\-ธิปติ\csromandash{}อริยา มคฺคา วุฏฺฐหิตฺวา ฯเปฯ. นิพฺพานํ ครุํ กตฺวา ฯเปฯ. (ปฐม\-คมน\-สทิสํ.) จกฺขุํ ฯเปฯ วตฺถุํ อเจตสิเก ขนฺเธ ครุํ กตฺวา อสฺสาเทติ อภินนฺทติ, ตํ ครุํ กตฺวา ราโค อุปฺปชฺชติ, ทิฏฺฐิ อุปฺปชฺชติ. สหชาตา\-ธิปติ\csromandash{} อเจตสิกา\-ธิปติ สมฺปยุตฺตกานํ ขนฺธานํ อธิปติ\-ปจฺจเยน ปจฺจโย.}

\prose{อเจตสิโก ธมฺโม เจตสิกสฺส จ อเจตสิกสฺส จ ธมฺมสฺส อธิปติ\-ปจฺจเยน ปจฺจโย. อารมฺมณา\-ธิปติ, สหชาตา\-ธิปติ.  อารมฺมณา\-ธิปติ\csromandash{}อริยา มคฺคา วุฏฺฐหิตฺวา ฯเปฯ. นิพฺพานํ ฯเปฯ. (ปฐมคมนํ.) อเจตสิเก ขนฺเธ ครุํ กตฺวา อสฺสาเทติ อภินนฺทติ, ตํ ครุํ กตฺวา เจตสิกา ขนฺธา จ จิตฺตญฺจ อุปฺปชฺชนฺติ. สหชาตา\-ธิปติ\csromandash{} อเจตสิกา\-ธิปติ สมฺปยุตฺตกานํ ขนฺธานํ จิตฺตสมุฏฺฐานานญฺจ รูปานํ อธิปติ\-ปจฺจเยน ปจฺจโย. (3)}

\prose{เจตสิโก จ อเจตสิโก จ ธมฺมา เจตสิกสฺส ธมฺมสฺส อธิปติ\-ปจฺจเยน ปจฺจโย. อารมฺมณา\-ธิปติ. ตีณิ. (อารมฺมณา\-ธิปติเยว.)}

\proseitem{125}{เจตสิโก ธมฺโม เจตสิกสฺส ธมฺมสฺส อนนฺตร\-ปจฺจเยน ปจฺจโย. ปุริมา ปุริมา เจตสิกา ขนฺธา ปจฺฉิมานํ ปจฺฉิมานํ เจตสิกานํ ขนฺธานํ อนนฺตร\-ปจฺจเยน ปจฺจโย. (มูลํ ปุจฺฉิตพฺพํ.) ปุริมา ปุริมา เจตสิกา ขนฺธา ปจฺฉิมสฺส ปจฺฉิมสฺส จิตฺตสฺส อนนฺตร\-ปจฺจเยน ปจฺจโย. (มูลํ ปุจฺฉิตพฺพํ.) ปุริมา ปุริมา เจตสิกา ขนฺธา ปจฺฉิมานํ ปจฺฉิมานํ เจตสิกานํ ขนฺธานํ จิตฺตสฺส จ นนฺตรปจฺจเยน ปจฺจโย. (3)}

\proseitem{126}{อเจตสิโก ธมฺโม อเจตสิกสฺส ธมฺมสฺส อนนฺตร\-ปจฺจเยน ปจฺจโย. ปุริมํ ปุริมํ จิตฺตํ ปจฺฉิมสฺส ปจฺฉิมสฺส จิตฺตสฺส อนนฺตร\-ปจฺจเยน ปจฺจโย. อนุโลมํ โคตฺรภุสฺส ฯเปฯ ผลสมาปตฺติยา อนนฺตร\-ปจฺจเยน ปจฺจโย. (1)}

\prose{\csromanfolio{417}อเจตสิโก ธมฺโม เจตสิกสฺส ธมฺมสฺส อนนฺตร\-ปจฺจเยน ปจฺจโย. (เอวํ ตีณิ กาตพฺพานิ, ปฐม\-คมน\-สทิสํ. ปูริตฺวา กาตพฺพํ, นินฺนานากรณํ.)}

\prose{เจตสิโก จ อเจตสิโก จ ธมฺมา เจตสิกสฺส ธมฺมสฺส อนนฺตร\-ปจฺจเยน ปจฺจโย. ตีณิ. (อาวชฺชนาปิ วุฏฺฐานมฺปิ นตฺถิ.)}

\prose{สมนนฺตร\-ปจฺจเยน ปจฺจโย. นว. สหชาต\-ปจฺจเยน ปจฺจโย. นว. (ปฏิจฺจ\-วาร\-สทิสํ.) อญฺญมญฺญ\-ปจฺจเยน ปจฺจโย. นว. (ปฏิจฺจ\-วาร\-สทิสํ.) นิสฺสย\-ปจฺจเยน ปจฺจโย. นว. (ปจฺจยวาร\-สทิสํ.)}

\csromanheader{2}{\textbf{อุปนิสฺสย}}
\proseitem{127}{เจตสิโก ธมฺโม เจตสิกสฺส ธมฺมสฺส อุปนิสฺสย\-ปจฺจเยน ปจฺจโย. อารมฺมณู\-ปนิสฺสโย, อนนฺตรู\-ปนิสฺสโย, ปกตู\-ปนิสฺสโย ฯเปฯ. ปกตู\-ปนิสฺสโย\csromandash{}เจตสิกา ขนฺธา เจตสิกานํ ขนฺธานํ อุปนิสฺสย\-ปจฺจเยน ปจฺจโย. (มูลํ ปุจฺฉิตพฺพํ. ตีณิ อุปนิสฺสยา.) เจตสิกา ขนฺธา จิตฺตสฺส อุปนิสฺสย\-ปจฺจเยน ปจฺจโย. (มูลํ ปุจฺฉิตพฺพํ. ตีณิ อุปนิสฺสยา.) เจตสิกา ขนฺธา เจตสิกานํ ขนฺธานํ จิตฺตสฺส จ อุปนิสฺสย\-ปจฺจเยน ปจฺจโย. (3)}

\proseitem{128}{อเจตสิโก ธมฺโม อเจตสิกสฺส ธมฺมสฺส อุปนิสฺสย\-ปจฺจเยน ปจฺจโย. อารมฺมณู\-ปนิสฺสโย, อนนฺตรู\-ปนิสฺสโย, ปกตู\-ปนิสฺสโย ฯเปฯ. ปกตู\-ปนิสฺสโย\csromandash{}อุตุํ. โภชนํ. เสนาสนํ จิตฺตํ อุปนิสฺสาย ทานํ เทติ ฯเปฯ สํฆํ ภินฺทติ. อุตุํ. โภชนํ. เสนาสนํ จิตฺตํ สทฺธาย ฯเปฯ ปญฺญาย. ราคสฺส ฯเปฯ ปตฺถนาย กายิกสฺส สุขสฺส, กายิกสฺส ทุกฺขสฺส, มคฺคสฺส, ผลสมาปตฺติยา อุปนิสฺสย\-ปจฺจเยน ปจฺจโย. (1)}

\prose{อเจตสิโก ธมฺโม เจตสิกสฺส ธมฺมสฺส อุปนิสฺสย\-ปจฺจเยน ปจฺจโย. อารมฺมณู\-ปนิสฺสโย, อนนฺตรู\-ปนิสฺสโย, ปกตู\-ปนิสฺสโย ฯเปฯ. ปกตู\-ปนิสฺสโย\csromandash{}อุตุํ. โภชนํ. เสนาสนํ จิตฺตํ อุปนิสฺสาย ทานํ เทติ ฯเปฯ. (ตีณิ. ปฐม\-คมน\-สทิสํ. นินฺนานากรณํ.)}

\prose{เจตสิโก จ อเจตสิโก จ ธมฺมา เจตสิกสฺส ธมฺมสฺส อุปนิสฺสย\-ปจฺจเยน ปจฺจโย. ตีณิ.}

\proseitem{129}{\csromanfolio{418}อเจตสิโก ธมฺโม อเจตสิกสฺส ธมฺมสฺส ปุเรชาต\-ปจฺจเยน ปจฺจโย. อารมฺมณ\-ปุเรชาตํ, วตฺถุ\-ปุเรชาตํ.  อารมฺมณ\-ปุเรชาตํ\csromandash{} จกฺขุํ ฯเปฯ วตฺถุํ อนิจฺจโต ฯเปฯ วิปสฺสติ, อสฺสาเทติ อภินนฺทติ, ตํ อารพฺภ จิตฺตํ อุปฺปชฺชติ. ทิพฺเพน จกฺขุนา รูปํ ปสฺสติ. ทิพฺพาย โสตธาตุยา สทฺทํ สุณาติ. รูปายตนํ จกฺขุ\-วิญฺญาณสฺส ฯเปฯ. โผฏฺฐพฺพา\-ยตนํ กายวิญฺญาณสฺส ฯเปฯ. วตฺถุ\-ปุเรชาตํ\csromandash{}จกฺขายตนํ จกฺขุ\-วิญฺญาณสฺส ฯเปฯ กายายตนํ กายวิญฺญาณสฺส ฯเปฯ. วตฺถุ จิตฺตสฺส ปุเรชาต\-ปจฺจเยน ปจฺจโย. (1)}

\prose{อเจตสิโก ธมฺโม เจตสิกสฺส ธมฺมสฺส ปุเรชาต\-ปจฺจเยน ปจฺจโย. อารมฺมณ\-ปุเรชาตํ, วตฺถุ\-ปุเรชาตํ.  อารมฺมณ\-ปุเรชาตํ\csromandash{}จกฺขุํ ฯเปฯ. วตฺถุํ อนิจฺจโต ฯเปฯ โทมนสฺสํ อุปฺปชฺชติ. ทิพฺเพน จกฺขุนา รูปํ ปสฺสติ. ทิพฺพาย โสตธาตุยา สทฺทํ สุณาติ. รูปายตนํ จกฺขุ\-วิญฺญาณสหคตานํ ขนฺธานํ ฯเปฯ. โผฏฺฐพฺพา\-ยตนํ กาย\-วิญฺญาณสหคตานํ ขนฺธานํ ฯเปฯ. วตฺถุ\-ปุเรชาตํ\csromandash{} จกฺขายตนํ จกฺขุ\-วิญฺญาณสหคตานํ ขนฺธานํ ฯเปฯ กายายตนํ ฯเปฯ. วตฺถุ เจตสิกานํ ขนฺธานํ ปุเรชาต\-ปจฺจเยน ปจฺจโย. (2)}

\prose{อเจตสิโก ธมฺโม เจตสิกสฺส จ อเจตสิกสฺส จ ธมฺมสฺส ปุเรชาต\-ปจฺจเยน ปจฺจโย. อารมฺมณ\-ปุเรชาตํ, วตฺถุ\-ปุเรชาตํ.  อารมฺมณ\-ปุเรชาตํ\csromandash{}จกฺขุํ ฯเปฯ วตฺถุํ อนิจฺจโต ฯเปฯ วิปสฺสติ, อสฺสาเทติ อภินนฺทติ, ตํ อารพฺภ จิตฺตญฺจ สมฺปยุตฺตกา จ ขนฺธา อุปฺปชฺชนฺติ. ทิพฺเพน จกฺขุนา รูปํ ปสฺสติ. ทิพฺพาย โสตธาตุยา สทฺทํ สุณาติ. รูปายตนํ จกฺขุ\-วิญฺญาณสฺส สมฺปยุตฺตกานญฺจ ขนฺธานํ ฯเปฯ. โผฏฺฐพฺพา\-ยตนํ ฯเปฯ. วตฺถุ\-ปุเรชาตํ\csromandash{}จกฺขายตนํ จกฺขุ\-วิญฺญาณสฺส สมฺปยุตฺตกานญฺจ ขนฺธานํ ฯเปฯ กายายตนํ ฯเปฯ. วตฺถุ จิตฺตสฺส สมฺปยุตฺตกานญฺจ ขนฺธานํ ปุเรชาต\-ปจฺจเยน ปจฺจโย. (3)}

\proseitem{130}{เจตสิโก ธมฺโม อเจตสิกสฺส ธมฺมสฺส ปจฺฉาชาต\-ปจฺจเยน ปจฺจโย. (สํขิตฺตํ.)}

\prose{\csromanfolio{419}อเจตสิโก ธมฺโม อเจตสิกสฺส ธมฺมสฺส ปจฺฉาชาต\-ปจฺจเยน ปจฺจโย. (สํขิตฺตํ.)}

\prose{เจตสิโก จ อเจตสิโก จ ธมฺมา อเจตสิกสฺส ธมฺมสฺส ปจฺฉาชาต\-ปจฺจเยน ปจฺจโย. (สํขิตฺตํ.)}

\prose{เจตสิโก ธมฺโม เจตสิกสฺส ธมฺมสฺส อาเสวน\-ปจฺจเยน ปจฺจโย. (สํขิตฺตํ.)}

\proseitem{131}{เจตสิโก ธมฺโม เจตสิกสฺส ธมฺมสฺส กมฺมปจฺจเยน ปจฺจโย. สหชาตา, นานากฺขณิกา.  สหชาตา เจตสิกา เจตนา สมฺปยุตฺตกานํ ขนฺธานํ กมฺมปจฺจเยน ปจฺจโย. ปฏิสนฺธิ\-กฺขเณ ฯเปฯ. นานากฺขณิกา เจตสิกา เจตนา วิปากานํ ขนฺธานํ กมฺมปจฺจเยน ปจฺจโย. (1)}

\prose{เจตสิโก ธมฺโม อเจตสิกสฺส ธมฺมสฺส กมฺมปจฺจเยน ปจฺจโย. สหชาตา, นานากฺขณิกา.  สหชาตา เจตสิกา เจตนา จิตฺตสฺส จิตฺตสมุฏฺฐานานญฺจ รูปานํ กมฺมปจฺจเยน ปจฺจโย. ปฏิสนฺธิ\-กฺขเณ ฯเปฯ. นานากฺขณิกา เจตสิกา เจตนา วิปากสฺส จิตฺตสฺส กฏตฺตา จ รูปานํ กมฺมปจฺจเยน ปจฺจโย. (2)}

\prose{เจตสิโก ธมฺโม เจตสิกสฺส จ อเจตสิกสฺส จ ธมฺมสฺส กมฺมปจฺจเยน ปจฺจโย. สหชาตา, นานากฺขณิกา.  สหชาตา เจตสิกา เจตนา สมฺปยุตฺตกานํ ขนฺธานํ จิตฺตสฺส จ จิตฺต\-สมุฏฺฐานญฺจ รูปานํ กมฺมปจฺจเยน ปจฺจโย. ปฏิสนฺธิ\-กฺขเณ ฯเปฯ. นานากฺขณิกา เจตสิกา เจตนา วิปากานํ ขนฺธานํ จิตฺตสฺส กฏตฺตา จ รูปานํ กมฺมปจฺจเยน ปจฺจโย. (3)}

\csromanheader{2}{\textbf{วิปากาทิ}}
\proseitem{132}{เจตสิโก ธมฺโม เจตสิกสฺส ธมฺมสฺส วิปาก\-ปจฺจเยน ปจฺจโย. นว. อาหารปจฺจเยน ปจฺจโย. นว. อินฺทฺริย\-ปจฺจเยน ปจฺจโย. นว. ฌานปจฺจเยน ปจฺจโย. ตีณิ. มคฺคปจฺจเยน ปจฺจโย. ตีณิ. สมฺปยุตฺต\-ปจฺจเยน ปจฺจโย. ปญฺจ.}

\proseitem{133}{\csromanfolio{420}เจตสิโก ธมฺโม อเจตสิกสฺส ธมฺมสฺส วิปฺปยุตฺต\-ปจฺจเยน ปจฺจโย. สหชาตํ, ปจฺฉาชาตํ. (สํขิตฺตํ.) (1)}

\prose{อเจตสิโก ธมฺโม อเจตสิกสฺส ธมฺมสฺส วิปฺปยุตฺต\-ปจฺจเยน ปจฺจโย. สหชาตํ, ปุเรชาตํ, ปจฺฉาชาตํ.  สหชาตํ จิตฺตํ จิตฺต\-สมุฏฺฐานานํ รูปานํ วิปฺปยุตฺต\-ปจฺจเยน ปจฺจโย. ปฏิสนฺธิ\-กฺขเณ จิตฺตํ กฏตฺตารูปานํ วิปฺปยุตฺต\-ปจฺจเยน ปจฺจโย. จิตฺตํ วตฺถุสฺส วิปฺปยุตฺต\-ปจฺจเยน ปจฺจโย. วตฺถุ จิตฺตสฺส วิปฺปยุตฺต\-ปจฺจเยน ปจฺจโย. ปุเรชาตํ จกฺขายตนํ จกฺขุ\-วิญฺญาณสฺส วิปฺปยุตฺต\-ปจฺจเยน ปจฺจโย ฯเปฯ กายายตนํ กายวิญฺญาณสฺส ฯเปฯ. วตฺถุ จิตฺตสฺส วิปฺปยุตฺต\-ปจฺจเยน ปจฺจโย. ปจฺฉาชาตํ จิตฺตํ ปุเรชาตสฺส อิมสฺส กายสฺส วิปฺปยุตฺต\-ปจฺจเยน ปจฺจโย. (1)}

\prose{อเจตสิโก ธมฺโม เจตสิกสฺส ธมฺมสฺส วิปฺปยุตฺต\-ปจฺจเยน ปจฺจโย. สหชาตํ, ปุเรชาตํ.  สหชาตํ ปฏิสนฺธิ\-กฺขเณ วตฺถุ เจตสิกานํ ขนฺธานํ วิปฺปยุตฺต\-ปจฺจเยน ปจฺจโย. ปุเรชาตํ จกฺขายตนํ จกฺขุ\-วิญฺญาณสหคตานํ ขนฺธานํ ฯเปฯ กายายตนํ กาย\-วิญฺญาณสหคตานํ ขนฺธานํ ฯเปฯ. วตฺถุ เจตสิกานํ ขนฺธานํ วิปฺปยุตฺต\-ปจฺจเยน ปจฺจโย. (2)}

\prose{อเจตสิโก ธมฺโม เจตสิกสฺส จ อเจตสิกสฺส จ ธมฺมสฺส วิปฺปยุตฺต\-ปจฺจเยน ปจฺจโย. สหชาตํ, ปุเรชาตํ.  สหชาตํ ปฏิสนฺธิ\-กฺขเณ วตฺถุ เจตสิกานํ ขนฺธานํ จิตฺตสฺส จ วิปฺปยุตฺต\-ปจฺจเยน ปจฺจโย. ปุเรชาตํ จกฺขายตนํ จกฺขุ\-วิญฺญาณสฺส สมฺปยุตฺตกานญฺจ ขนฺธานํ ฯเปฯ กายายตนํ กายวิญฺญาณสฺส สมฺปยุตฺตกานญฺจ ขนฺธานํ ฯเปฯ. วตฺถุ จิตฺตสฺส สมฺปยุตฺตกานญฺจ ขนฺธานํ วิปฺปยุตฺต\-ปจฺจเยน ปจฺจโย. (3)}

\prose{เจตสิโก จ อเจตสิโก จ ธมฺมา อเจตสิกสฺส ธมฺมสฺส วิปฺปยุตฺต\-ปจฺจเยน ปจฺจโย. สหชาตํ, ปจฺฉาชาตํ. (สํขิตฺตํ.)}

\proseitem{134}{เจตสิโก ธมฺโม เจตสิกสฺส ธมฺมสฺส อตฺถิปจฺจเยน ปจฺจโย. เอกํ. (ปฏิจฺจ\-วาร\-สทิสํ.)}

\prose{เจตสิโก ธมฺโม อเจตสิกสฺส ธมฺมสฺส อตฺถิปจฺจเยน ปจฺจโย. สหชาตํ, ปจฺฉาชาตํ. (สํขิตฺตํ.) (2)}

\prose{\csromanfolio{421}เจตสิโก ธมฺโม เจตสิกสฺส จ อเจตสิกสฺส จ ธมฺมสฺส อตฺถิปจฺจเยน ปจฺจโย. เอกํ. (ปฏิจฺจ\-วาร\-สทิสํ.) (3)}

\proseitem{135}{อเจตสิโก ธมฺโม อเจตสิกสฺส ธมฺมสฺส อตฺถิปจฺจเยน ปจฺจโย. สหชาตํ, ปุเรชาตํ, ปจฺฉาชาตํ, อาหารํ, อินฺทฺริยํ. (สํขิตฺตํ.) (1)}

\prose{อเจตสิโก ธมฺโม เจตสิกสฺส ธมฺมสฺส อตฺถิปจฺจเยน ปจฺจโย. สหชาตํ, ปุเรชาตํ.  สหชาตํ จิตฺตํ เจตสิกานํ ขนฺธานํ อตฺถิปจฺจเยน ปจฺจโย. ปฏิสนฺธิ\-กฺขเณ จิตฺตํ ฯเปฯ. ปฏิสนฺธิ\-กฺขเณ วตฺถุ เจตสิกา ขนฺธานํ อตฺถิปจฺจเยน ปจฺจโย. ปุเรชาตํ จกฺขุํ ฯเปฯ วตฺถุํ อนิจฺจโต ฯเปฯ. (ปุเรชาต\-สทิสํ, นินฺนานากรณํ.) (2)}

\prose{อเจตสิโก ธมฺโม เจตสิกสฺส จ อเจตสิกสฺส จ ธมฺมสฺส อตฺถิปจฺจเยน ปจฺจโย. สหชาตํ, ปุเรชาตํ.  สหชาตํ จิตฺตํ สมฺปยุตฺตกานํ ขนฺธานํ จิตฺตสมุฏฺฐานานญฺจ รูปานํ อตฺถิปจฺจเยน ปจฺจโย. ปฏิสนฺธิ\-กฺขเณ จิตฺตํ ฯเปฯ. ปฏิสนฺธิ\-กฺขเณ วตฺถุ เจตสิกานํ ขนฺธานํ จิตฺตสฺส จ อตฺถิปจฺจเยน ปจฺจโย. ปุเรชาตํ จกฺขุํ ฯเปฯ วตฺถุํ อนิจฺจโต ฯเปฯ. (ปุเรชาต\-สทิสํ นินฺนานํ.) (3)}

\proseitem{136}{เจตสิโก จ อเจตสิโก จ ธมฺมา เจตสิกสฺส ธมฺมสฺส อตฺถิปจฺจเยน ปจฺจโย. สหชาตํ, ปุเรชาตํ.  สหชาโต จกฺขุ\-วิญฺญาณ\-สหคโต เอโก ขนฺโธ จ จกฺขายตนญฺจ ทฺวินฺนํ ขนฺธานํ ฯเปฯ. กาย\-วิญฺญาณ\-สหคโต ฯเปฯ. เจตสิโก เอโก ขนฺโธ จ วตฺถุ จ ทฺวินฺนํ ขนฺธานํ อตฺถิปจฺจเยน ปจฺจโย. เทฺว ขนฺธา จ ฯเปฯ. สหชาโต เจตสิโก เอโก ขนฺโธ จ จิตฺตญฺจ ทฺวินฺนํ ขนฺธานํ อตฺถิปจฺจเยน ปจฺจโย. เทฺว ขนฺธา จ ฯเปฯ. (ปฏิสนฺธิ\-กฺขเณ เทฺวปิ กาตพฺพา.) (1)}

\prose{เจตสิโก จ อเจตสิโก จ ธมฺมา อเจตสิกสฺส ธมฺมสฺส อตฺถิปจฺจเยน ปจฺจโย. สหชาตํ, ปุเรชาตํ, ปจฺฉาชาตํ, อาหารํ, อินฺทฺริยํ.  สหชาตา จกฺขุ\-วิญฺญาณสหคตา ขนฺธา จ จกฺขายตนญฺจ จกฺขุ\-วิญฺญาณสฺส ฯเปฯ. กายวิญฺญาณสฺส ฯเปฯ. เจตสิกา ขนฺธา จ จิตฺตญฺจ จิตฺต\-สมุฏฺฐานานํ รูปานํ อตฺถิปจฺจเยน ปจฺจโย. เจตสิกา ขนฺธา จ จิตฺตญฺจ มหาภูตา จ จิตฺต\-สมุฏฺฐานานํ รูปานํ อตฺถิปจฺจเยน ปจฺจโย. \csromanfolio{422}สหชาตา เจตสิกา ขนฺธา จ วตฺถุ จ จิตฺตสฺส อตฺถิปจฺจเยน ปจฺจโย. (ปฏิสนฺธิ\-กฺขเณ ตีณิปิ กาตพฺพา.) ปจฺฉาชาตา เจตสิกา ขนฺธา จ จิตฺตญฺจ ปุเรชาตสฺส อิมสฺส กายสฺส อตฺถิปจฺจเยน ปจฺจโย. ปจฺฉาชาตา เจตสิกา ขนฺธา จ จิตฺตญฺจ กพฬีกาโร อาหาโร จ อิมสฺส กายสฺส อตฺถิปจฺจเยน ปจฺจโย. ปจฺฉาชาตา เจตสิกา ขนฺธา จ จิตฺตญฺจ รูปชีวิตินฺทฺริยญฺจ กฏตฺตารูปานํ อตฺถิปจฺจเยน ปจฺจโย. (2)}

\prose{เจตสิโก จ อเจตสิโก จ ธมฺมา เจตสิกสฺส จ อเจตสิกสฺส จ ธมฺมสฺส อตฺถิปจฺจเยน ปจฺจโย. สหชาตํ, ปุเรชาตํ.  สหชาโต จกฺขุ\-วิญฺญาณ\-สหคโต เอโก ขนฺโธ จ จกฺขายตนญฺจ ทฺวินฺนํ ขนฺธานํ จกฺขุ\-วิญฺญาณสฺส จ อตฺถิปจฺจเยน ปจฺจโย ฯเปฯ. กาย\-วิญฺญาณ\-สหคโต เอโก ขนฺโธ จ กายายตนญฺจ ทฺวินฺนํ ขนฺธานํ กายวิญฺญาณสฺส จ อตฺถิปจฺจเยน ปจฺจโย. เทฺว ขนฺธา จ ฯเปฯ. สหชาโต เจตสิโก เอโก ขนฺโธ จ วตฺถุ จ ทฺวินฺนํ ขนฺธานํ จิตฺตสฺส จ อตฺถิปจฺจเยน ปจฺจโย. เทฺว ขนฺธา จ ฯเปฯ. (ปฏิสนฺธิยา เทฺว กาตพฺพา.) (3)}

\csromanmarkh{1. ปจฺจยานุโลม}\csromanmarkhii{2. สงฺขฺยาวาร}\csromanheadernomark{2}{1. \textbf{ปจฺจยานุโลม 2. สงฺขฺยาวาร}}
\proseitem{137}{เหตุยา ตีณิ, อารมฺมเณ นว, อธิปติยา นว, อนนฺตเร นว, สมนนฺตเร นว, สหชาเต นว, อญฺญมญฺเญ นว, นิสฺสเย นว, อุปนิสฺสเย อนว, ปุเรชาเต ตีณิ, ปจฺฉาชาเต ตีณิ, อาเสวเน นว, กมฺเม ตีณิ, วิปาเก นว, อาหาเร นว, อินฺทฺริเย นว, ฌาเน ตีณิ, มคฺเค ตีณิ, สมฺปยุตฺเต ปญฺจ, วิปฺปยุตฺเต ปญฺจ, อตฺถิยา นว, นตฺถิยา นว, วิคเต นว, อวิคเต นว.}

\vspace{\tipitakaparskip}
\csromancenter{อนุโลมํ.}
\proseitem{138}{เจตสิโก ธมฺโม เจตสิกสฺส ธมฺมสฺส อารมฺมณ\-ปจฺจเยน ปจฺจโย. สหชาต\-ปจฺจเยน ปจฺจโย. อุปนิสฺสย\-ปจฺจเยน ปจฺจโย. กมฺมปจฺจเยน ปจฺจโย. (1)}

\prose{\csromanfolio{423}เจตสิโก ธมฺโม อเจตสิกสฺส ธมฺมสฺส อารมฺมณ\-ปจฺจเยน ปจฺจโย. สหชาต\-ปจฺจเยน ปจฺจโย. อุปนิสฺสย\-ปจฺจเยน ปจฺจโย. ปจฺฉาชาต\-ปจฺจเยน ปจฺจโย. กมฺมปจฺจเยน ปจฺจโย. (2)}

\prose{เจตสิโก ธมฺโม เจตสิกสฺส จ อเจตสิกสฺส จ ธมฺมสฺส อารมฺมณ\-ปจฺจเยน ปจฺจโย. สหชาต\-ปจฺจเยน ปจฺจโย. อุปนิสฺสย\-ปจฺจเยน ปจฺจโย. กมฺมปจฺจเยน ปจฺจโย. (3)}

\proseitem{139}{อเจตสิโก ธมฺโม อเจตสิกสฺส ธมฺมสฺส อารมฺมณ\-ปจฺจเยน ปจฺจโย. สหชาต\-ปจฺจเยน ปจฺจโย. อุปนิสฺสย\-ปจฺจเยน ปจฺจโย. ปุเรชาต\-ปจฺจเยน ปจฺจโย. ปจฺฉาชาต\-ปจฺจเยน ปจฺจโย. อาหารปจฺจเยน ปจฺจโย. อินฺทฺริย\-ปจฺจเยน ปจฺจโย. (1)}

\prose{อเจตสิโก ธมฺโม เจตสิกสฺส ธมฺมสฺส อารมฺมณ\-ปจฺจเยน ปจฺจโย. สหชาต\-ปจฺจเยน ปจฺจโย. อุปนิสฺสย\-ปจฺจเยน ปจฺจโย. ปุเรชาต\-ปจฺจเยน ปจฺจโย. (2)}

\prose{อเจตสิโก ธมฺโม เจตสิกสฺส จ อเจตสิกสฺส จ ธมฺมสฺส อารมฺมณ\-ปจฺจเยน ปจฺจโย. สหชาต\-ปจฺจเยน ปจฺจโย. อุปนิสฺสย\-ปจฺจเยน ปจฺจโย. ปุเรชาต\-ปจฺจเยน ปจฺจโย. (3)}

\proseitem{140}{เจตสิโก จ อเจตสิโก จ ธมฺมา เจตสิกสฺส ธมฺมสฺส อารมฺมณ\-ปจฺจเยน ปจฺจโย. สหชาต\-ปจฺจเยน ปจฺจโย. อุปนิสฺสย\-ปจฺจเยน ปจฺจโย. (1)}

\prose{เจตสิโก จ อเจตสิโก จ ธมฺมา อเจตสิกสฺส ธมฺมสฺส อารมฺมณ\-ปจฺจเยน ปจฺจโย. สหชาต\-ปจฺจเยน ปจฺจโย. อุปนิสฺสย\-ปจฺจเยน ปจฺจโย. ปจฺฉาชาต\-ปจฺจเยน ปจฺจโย. (2)}

\prose{เจตสิโก จ อเจตสิโก จ ธมฺมา เจตสิกสฺส จ อเจตสิกสฺส จ ธมฺมสฺส อารมฺมณ\-ปจฺจเยน ปจฺจโย. สหชาต\-ปจฺจเยน ปจฺจโย. อุปนิสฺสย\-ปจฺจเยน ปจฺจโย. (3)}

\csromanmarkh{2. ปจฺจย\-ปจฺจนีย}\csromanmarkhii{2. สงฺขฺยาวาร}\csromanheadernomark{2}{2. ปจฺจย\-ปจฺจนีย 2. สงฺขฺยาวาร}
\vspace{\tipitakaparskip}
\csromancenter{นเหตุยา นว, นอารมฺมเณ นว, (สพฺพตฺถ นว,) โนอวิคเต นว.}
\csromancenter{\textbf{เหตุ}ปจฺจยา นอารมฺมเณ ตีณิ, นอธิปติยา ตีณิ, นอนนฺตเร ตีณิ, นสมนนฺตเร ตีณิ, นอญฺญมญฺเญ เอกํ, นฺเอาปนิสฺสเย ตีณิ, (สพฺพตฺถ ตีณิ,) นมคฺเค ตีณิ, นสมฺปยุตฺเต เอกํ, นวิปฺปยุตฺเต ตีณิ, โนนตฺถิยา ตีณิ, โนวิคเต ตีณิ.}
\proseitemwithcloser{143}{\csromanfolio{424}นเหตุปจฺจยา อารมฺมเณ นว, อธิปติยา นว, (อนุโลมมาติกา กาตพฺพา.) ฯเปฯ อวิคเต นว.}{เจตสิกทุกํ นิฏฺฐิตํ.}
\csromanmatikaanchor{mtk.57}
\csromantocmarkref{subsection}{58. จิตฺตสมฺปยุตฺตทุก}{mtk.57}
\bookmark[dest={mtk.57},level=2]{58. จิตฺตสมฺปยุตฺตทุก}
\csromanmarkh{58. จิตฺต\-สมฺปยุตฺตทุก}\csromanmarkhii{1. ปฏิจฺจวาร}\csromanheadernomark{2}{58. จิตฺต\-สมฺปยุตฺตทุก 1. ปฏิจฺจวาร}
\csromanmarkh{1. ปจฺจยานุโลม}\csromanmarkhii{1. วิภงฺควาร}\csromanheadernomark{2}{1. \textbf{ปจฺจยานุโลม 1. วิภงฺควาร}}
\vspace{\tipitakaparskip}
\csromancenter{\textbf{เหตุ}}
\proseitem{144}{จิตฺต\-สมฺปยุตฺตํ ธมฺมํ ปฏิจฺจ จิตฺต\-สมฺปยุตฺโต ธมฺโม อุปฺปชฺชติ เหตุปจฺจยา. จิตฺต\-สมฺปยุตฺตํ เอกํ ขนฺธํ ปฏิจฺจ เทฺว ขนฺธา, เทฺว ขนฺเธ ปฏิจฺจ เอโก ขนฺโธ ปฏิสนฺธิ\-กฺขเณ ฯเปฯ. (1)}

\prose{จิตฺต\-สมฺปยุตฺตํ ธมฺมํ ปฏิจฺจ จิตฺต\-วิปฺปยุตฺโต ธมฺโม อุปฺปชฺชติ เหตุปจฺจยา. จิตฺต\-สมฺปยุตฺเต ขนฺเธ จิตฺต\-สมุฏฺฐานํ รูปํ. ปฏิสนฺธิ\-กฺขเณ ฯเปฯ. (2)}

\prose{จิตฺต\-สมฺปยุตฺตํ ธมฺมํ ปฏิจฺจ จิตฺต\-สมฺปยุตฺโต จ จิตฺต\-วิปฺปยุตฺโต จ ธมฺมา อุปฺปชฺชนฺติ เหตุปจฺจยา. จิตฺต\-สมฺปยุตฺตํ เอกํ ขนฺธํ ปฏิจฺจ เทฺว ขนฺธา จิตฺต\-สมุฏฺฐานญฺจ รูปํ. เทฺว ขนฺเธ ฯเปฯ. ปฏิสนฺธิ\-กฺขเณ ฯเปฯ. (3)}

\proseitem{145}{จิตฺต\-วิปฺปยุตฺตํ ธมฺมํ ปฏิจฺจ จิตฺต\-วิปฺปยุตฺโต ธมฺโม อุปฺปชฺชติ เหตุปจฺจยา. เอกํ มหาภูตํ ฯเปฯ มหาภูเต ปฏิจฺจ จิตฺต\-สมุฏฺฐานํ รูปํ ฯเปฯ. ปฏิสนฺธิ\-กฺขเณ ฯเปฯ. เอกํ มหาภูตํ ฯเปฯ มหาภูเต ปฏิจฺจ กฏตฺตารูปํ, อุปาทารูปํ. (1)}

\prose{\csromanfolio{425}จิตฺต\-วิปฺปยุตฺตํ ธมฺมํ ปฏิจฺจ จิตฺต\-สมฺปยุตฺโต ธมฺโม อุปฺปชฺชติ เหตุปจฺจยา. ปฏิสนฺธิ\-กฺขเณ วตฺถุํ ปฏิจฺจ จิตฺต\-สมฺปยุตฺตกา ขนฺธา.}

\prose{จิตฺต\-วิปฺปยุตฺตํ ธมฺมํ ปฏิจฺจ จิตฺต\-สมฺปยุตฺโต จ จิตฺต\-วิปฺปยุตฺโต จ ธมฺมา อุปฺปชฺชนฺติ เหตุปจฺจยา. ปฏิสนฺธิ\-กฺขเณ วตฺถุํ ปฏิจฺจ จิตฺต\-สมฺปยุตฺตกา ขนฺธา. มหาภูเต ปฏิจฺจ กฏตฺตารูปํ. (3)}

\proseitem{146}{จิตฺต\-สมฺปยุตฺตญฺจ จิตฺต\-วิปฺปยุตฺตญฺจ ธมฺมํ ปฏิจฺจ จิตฺต\-สมฺปยุตฺโต ธมฺโม อุปฺปชฺชติ เหตุปจฺจยา. ปฏิสนฺธิ\-กฺขเณ จิตฺต\-สมฺปยุตฺตํ เอกํ ขนฺธญฺจ วตฺถุญฺจ ปฏิจฺจ เทฺว ขนฺธา. เทฺว ขนฺเธ จ ฯเปฯ. (1)}

\prose{จิตฺต\-สมฺปยุตฺตญฺจ จิตฺต\-วิปฺปยุตฺตญฺจ ธมฺมํ ปฏิจฺจ จิตฺต\-วิปฺปยุตฺโต ธมฺโม อุปฺปชฺชติ เหตุปจฺจยา. จิตฺต\-สมฺปยุตฺเต ขนฺเธ จ มหาภูเต จ ปฏิจฺจ จิตฺต\-สมุฏฺฐานํ รูปํ. ปฏิสนฺธิ\-กฺขเณ จิตฺต\-สมฺปยุตฺเต ขนฺเธ จ มหาภูเต จ ปฏิจฺจ กฏตฺตารูปํ. (2)}

\prose{จิตฺต\-สมฺปยุตฺตญฺจ จิตฺต\-วิปฺปยุตฺตญฺจ ธมฺมํ ปฏิจฺจ จิตฺต\-สมฺปยุตฺโต จ จิตฺต\-วิปฺปยุตฺโต จ ธมฺมา อุปฺปชฺชนฺติ เหตุปจฺจยา. ปฏิสนฺธิ\-กฺขเณ จิตฺต\-สมฺปยุตฺตํ เอกํ ขนฺธญฺจ วตฺถุญฺจ ปฏิจฺจ เทฺว ขนฺธา. เทฺว ขนฺเธ จ ฯเปฯ. จิตฺต\-สมฺปยุตฺเต ขนฺเธ จ มหาภูเต จ ปฏิจฺจ กฏตฺตารูปํ. (3)}

\proseitem{147}{จิตฺต\-สมฺปยุตฺตํ ธมฺมํ ปฏิจฺจ จิตฺต\-สมฺปยุตฺโต ธมฺโม อุปฺปชฺชติ อารมฺมณ\-ปจฺจยา. จิตฺต\-สมฺปยุตฺตํ เอกํ ขนฺธํ ปฏิจฺจ เทฺว ขนฺธา. เทฺว ขนฺเธ ฯเปฯ. ปฏิสนฺธิ\-กฺขเณ ฯเปฯ. (1)}

\prose{จิตฺต\-วิปฺปยุตฺตํ ธมฺมํ ปฏิจฺจ จิตฺต\-สมฺปยุตฺโต ธมฺโม อุปฺปชฺชติ อารมฺมณ\-ปจฺจยา. ปฏิสนฺธิ\-กฺขเณ วตฺถุํ ปฏิจฺจ จิตฺต\-สมฺปยุตฺตกา ขนฺธา. (1)}

\prose{จิตฺต\-สมฺปยุตฺตญฺจ จิตฺต\-วิปฺปยุตฺตญฺจ ธมฺมํ ปฏิจฺจ จิตฺต\-สมฺปยุตฺโต ธมฺโม อุปฺปชฺชติ อารมฺมณ\-ปจฺจยา. ปฏิสนฺธิ\-กฺขเณ จิตฺต\-สมฺปยุตฺตํ เอกํ ขนฺธญฺจ วตฺถุญฺจ ปฏิจฺจ เทฺว ขนฺธา. เทฺว ขนฺเธ จ ฯเปฯ. (สํขิตฺตํ.) (1)}

\csromanmarkh{1. ปจฺจยานุโลม}\csromanmarkhii{2. สงฺขฺยาวาร}\csromanheadernomark{2}{1. \textbf{ปจฺจยานุโลม 2. สงฺขฺยาวาร}}
\proseitem{148}{\csromanfolio{426}เหตุยา นว, อารมฺมเณ ตีณิ, อธิปติยา ปญฺจ, อนนฺตเร ตีณิ, สมนนฺตเร ตีณิ, สหชาเต นว, อญฺญมญฺเญ ฉ, นิสฺสเย นว, อุปนิสฺสเย ตีณิ, ปุเรชาเต เอกํ, อาเสวเน เอกํ, กมฺเม นว, วิปาเก นว, อาหาเร นว, อินฺทฺริเย นว, ฌาเน นว, มคฺเค นว, สมฺปยุตฺเต ตีณิ, วิปฺปยุตฺเต นว, อตฺถิยา นว, นตฺถิยา ตีณิ, วิคเต ตีณิ, อวิคเต นว.}

\csromanmarkh{2. ปจฺจย\-ปจฺจนีย}\csromanmarkhii{1. วิภงฺควาร}\csromanheadernomark{2}{2. ปจฺจย\-ปจฺจนีย 1. วิภงฺควาร}
\proseitem{149}{จิตฺต\-สมฺปยุตฺตํ ธมฺมํ ปฏิจฺจ จิตฺต\-สมฺปยุตฺโต ธมฺโม อุปฺปชฺชติ นเหตุปจฺจยา. อเหตุกํ จิตฺต\-สมฺปยุตฺตํ เอกํ ขนฺธํ ปฏิจฺจ เทฺว ขนฺธา. เทฺว ขนฺเธ ฯเปฯ. อเหตุก\-ปฏิสนฺธิ\-กฺขเณ ฯเปฯ. วิจิกิจฺฉา\-สหคเต อุทฺธจฺจ\-สหคเต ขนฺเธ ปฏิจฺจ วิจิกิจฺฉา\-สหคโต อุทฺธจฺจ\-สหคโต โมโห. (1)}

\prose{(เอวํ นวปิ ปญฺหา กาตพฺพา. “อเหตุกนฺ”ติ สพฺพตฺถ นิยาเมตพฺพํ, เอกํเยว โมหํ มูลปเท.)}

\prose{นอารมฺมณ}

\proseitem{150}{จิตฺต\-สมฺปยุตฺตํ ธมฺมํ ปฏิจฺจ จิตฺต\-วิปฺปยุตฺโต ธมฺโม อุปฺปชฺชติ นอารมฺมณ\-ปจฺจยา. จิตฺต\-สมฺปยุตฺเต ขนฺเธ ปฏิจฺจ จิตฺต\-สมุฏฺฐานํ รูปํ. ปฏิสนฺธิ\-กฺขเณ ฯเปฯ. (1)}

\prose{จิตฺต\-วิปฺปยุตฺตํ ธมฺมํ ปฏิจฺจ จิตฺต\-วิปฺปยุตฺโต ธมฺโม อุปฺปชฺชติ นอารมฺมณ\-ปจฺจยา. เอกํ มหาภูตํ ฯเปฯ. (ยาว อสญฺญสตฺตา.) (1)}

\prose{จิตฺต\-สมฺปยุตฺตญฺจ จิตฺต\-วิปฺปยุตฺตญฺจ ธมฺมํ ปฏิจฺจ จิตฺต\-วิปฺปยุตฺโต ธมฺโม อุปฺปชฺชติ นอารมฺมณ\-ปจฺจยา. จิตฺต\-สมฺปยุตฺเต ขนฺเธ จ มหาภูเต จ ปฏิจฺจ จิตฺต\-สมุฏฺฐานํ รูปํ. ปฏิสนฺธิ\-กฺขเณ เอกํ. (สํขิตฺตํ.) (1)}

\csromanmarkh{58. จิตฺต\-สมฺปยุตฺตทุก}\csromanmarkhii{2. ปจฺจย\-ปจฺจนีย 2. สงฺขฺยาวาร}\csromanheadernomark{2}{58. จิตฺต\-สมฺปยุตฺตทุก 2. ปจฺจย\-ปจฺจนีย 2. สงฺขฺยาวาร}
\proseitem{151}{\csromanfolio{427}นเหตุยา นว, นอารมฺมเณ ตีณิ, นอธิปติยา นว, นอนนฺตเร ตีณิ, นสมนนฺตเร ตีณิ, นอญฺญมญฺเญ ตีณิ, นอุปนิสฺสเย ตีณิ, นปุเรชาเต นว, นปจฺฉาชาเต นว, นอาเสวเน นว, นกมฺเม เทฺว, นวิปาเก ปญฺจ, นอาหาเร เอกํ, นอินฺทฺริเย เอกํ, นฌาเน เทฺว, นมคฺเค นว, นสมฺปยุตฺเต ตีณิ, นวิปฺปยุตฺเต เทฺว, โนนตฺถิยา ตีณิ, โนวิคเต ตีณิ.}

\vspace{\tipitakaparskip}
\csromancenter{\textbf{เหตุ}ปจฺจยา นอารมฺมเณ ตีณิ, นอธิปติยา นว ฯเปฯ นกมฺเม เอกํ, นวิปาเก ปญฺจ, นสมฺปยุตฺเต ตีณิ, นวิปฺปยุตฺเต เอกํ, โนนตฺถิยา ตีณิ, โนวิคเต ตีณิ.}
\proseitem{153}{นเหตุปจฺจยา อารมฺมเณ ตีณิ, อนนฺตเร ตีณิ, สมนนฺตเร ตีณิ ฯเปฯ อญฺญมญฺเญ ฉ ฯเปฯ ปุเรชาเต เอกํ, อาเสวเน เอกํ, กมฺเม นว ฯเปฯ มคฺเค เอกํ ฯเปฯ อวิคเต นว.}

\vspace{\tipitakaparskip}
\csromancenter{(สหชาตวาโร ปฏิจฺจ\-วาร\-สทิโส.)}
\csromanmarkh{58. จิตฺต\-สมฺปยุตฺตทุก}\csromanmarkhii{3. ปจฺจยวาร}\csromanheadernomark{2}{58. จิตฺต\-สมฺปยุตฺตทุก 3. ปจฺจยวาร}
\csromanmarkh{1. ปจฺจยานุโลม}\csromanmarkhii{1. วิภงฺควาร}\csromanheadernomark{2}{1. \textbf{ปจฺจยานุโลม 1. วิภงฺควาร}}
\vspace{\tipitakaparskip}
\csromancenter{\textbf{เหตุ}}
\proseitem{154}{จิตฺต\-สมฺปยุตฺตํ ธมฺมํ ปจฺจยา จิตฺต\-สมฺปยุตฺโต ธมฺโม อุปฺปชฺชติ เหตุปจฺจยา. ตีณิ. (ปฏิจฺจสทิสา.)}

\prose{จิตฺต\-วิปฺปยุตฺตํ ธมฺมํ ปจฺจยา จิตฺต\-วิปฺปยุตฺโต ธมฺโม อุปฺปชฺชติ เหตุปจฺจยา. เอกํ. (ปฏิจฺจสทิสํ) (1)}

\prose{\csromanfolio{428}จิตฺต\-วิปฺปยุตฺตํ ธมฺมํ ปจฺจยา จิตฺต\-สมฺปยุตฺโต ธมฺโม อุปฺปชฺชติ เหตุปจฺจยา. วตฺถุํ ปจฺจยา จิตฺต\-สมฺปยุตฺตกา ขนฺธา. ปฏิสนฺธิ\-กฺขเณ ฯเปฯ. (2)}

\prose{จิตฺต\-วิปฺปยุตฺตํ ธมฺมํ ปจฺจยา จิตฺต\-สมฺปยุตฺโต จ จิตฺต\-วิปฺปยุตฺโต จ ธมฺมา อุปฺปชฺชนฺติ เหตุปจฺจยา. วตฺถุํ ปจฺจยา จิตฺต\-สมฺปยุตฺตกา ขนฺธา. มหาภูเต ปจฺจยา จิตฺต\-สมุฏฺฐานํ รูปํ. ปฏิสนฺธิ\-กฺขเณ ฯเปฯ. (3)}

\proseitem{155}{จิตฺต\-สมฺปยุตฺตญฺจ จิตฺต\-วิปฺปยุตฺตญฺจ ธมฺมํ ปจฺจยา จิตฺต\-สมฺปยุตฺโต ธมฺโม อุปฺปชฺชติ เหตุปจฺจยา. จิตฺต\-สมฺปยุตฺตํ เอกํ ขนฺธญฺจ วตฺถุญฺจ ปจฺจยา เทฺว ขนฺธา. เทฺว ขนฺเธ จ ฯเปฯ. ปฏิสนฺธิ\-กฺขเณ ฯเปฯ. (1)}

\prose{จิตฺต\-สมฺปยุตฺตญฺจ จิตฺต\-วิปฺปยุตฺตญฺจ ธมฺมํ ปจฺจยา จิตฺต\-วิปฺปยุตฺโต ธมฺโม อุปฺปชฺชติ เหตุปจฺจยา. จิตฺต\-สมฺปยุตฺเต ขนฺเธ จ มหาภูเต จ ปจฺจยา จิตฺต\-สมุฏฺฐานํ รูปํ. ปฏิสนฺธิ\-กฺขเณ ฯเปฯ. (2)}

\prose{จิตฺต\-สมฺปยุตฺตญฺจ จิตฺต\-วิปฺปยุตฺตญฺจ ธมฺมํ ปจฺจยา จิตฺต\-สมฺปยุตฺโต จ จิตฺต\-วิปฺปยุตฺโต จ ธมฺมา อุปฺปชฺชนฺติ เหตุปจฺจยา. จิตฺต\-สมฺปยุตฺตํ เอกํ ขนฺธญฺจ วตฺถุญฺจ ปจฺจยา เทฺว ขนฺธา. เทฺว ขนฺเธ จ ฯเปฯ. จิตฺต\-สมฺปยุตฺเต ขนฺเธ จ มหาภูเต จ ปจฺจยา จิตฺต\-สมุฏฺฐานํ รูปํ. ปฏิสนฺธิ\-กฺขเณ ฯเปฯ. (3)}

\proseitem{156}{จิตฺต\-สมฺปยุตฺตํ ธมฺมํ ปจฺจยา จิตฺต\-สมฺปยุตฺโต ธมฺโม อุปฺปชฺชติ อารมฺมณ\-ปจฺจยา. เอกํ. (ปฏิจฺจสทิสํ.) (1)}

\prose{จิตฺต\-วิปฺปยุตฺตํ ธมฺมํ ปจฺจยา จิตฺต\-สมฺปยุตฺโต ธมฺโม อุปฺปชฺชติ อารมฺมณ\-ปจฺจยา. จกฺขายตนํ ปจฺจยา จกฺขุ\-วิญฺญาณสหคตา ขนฺธา ฯเปฯ. กายายตนํ ปจฺจยา กาย\-วิญฺญาณสหคตา ขนฺธา ฯเปฯ. วตฺถุํ ปจฺจยา จิตฺต\-สมฺปยุตฺตกา ขนฺธา. ปฏิสนฺธิ\-กฺขเณ ฯเปฯ. (1)}

\prose{จิตฺต\-สมฺปยุตฺตญฺจ จิตฺต\-วิปฺปยุตฺตญฺจ ธมฺมํ ปจฺจยา จิตฺต\-สมฺปยุตฺโต ธมฺโม อุปฺปชฺชติ อารมฺมณ\-ปจฺจยา. จกฺขุ\-วิญฺญาณสหคตํ เอกํ ขนฺธญฺจ จกฺขายตนญฺจ ปจฺจยา เทฺว ขนฺธา. เทฺว ขนฺเธ จ ฯเปฯ. กาย\-วิญฺญาณสหคตํ เอกํ ขนฺธญฺจ กายายตนญฺจ ปจฺจยา เทฺว ขนฺธา. เทฺว ขนฺเธ จ ฯเปฯ. จิตฺต\-สมฺปยุตฺตํ เอกํ ขนฺธญฺจ วตฺถุญฺจ ปจฺจยา เทฺว ขนฺธา. เทฺว ขนฺเธ จ ฯเปฯ. ปฏิสนฺธิ\-กฺขเณ. (สํขิตฺตํ.) (1)}

\csromanmarkh{58. จิตฺต\-สมฺปยุตฺตทุก}\csromanmarkhii{1. ปจฺจยานุโลม 2. สงฺขฺยาวาร}\csromanheadernomark{2}{58. จิตฺต\-สมฺปยุตฺตทุก 1. ปจฺจยานุโลม 2. สงฺขฺยาวาร}
\proseitem{157}{\csromanfolio{429}เหตุยา นว, อารมฺมเณ ตีณิ, อธิปติยา นว, อนนฺตเร ตีณิ, สมนนฺตเร ตีณิ, สหชาเต นว, อญฺญมญฺเญ ฉ, นิสฺสเย นว, อุปนิสฺสเย ตีณิ, ปุเรชาเต ตีณิ, อาเสวเน ตีณิ, กมฺเม นว ฯเปฯ อวิคเต นว.}

\csromanmarkh{2. ปจฺจย\-ปจฺจนีย}\csromanmarkhii{1. วิภงฺควาร}\csromanheadernomark{2}{2. ปจฺจย\-ปจฺจนีย 1. วิภงฺควาร}
\proseitem{158}{จิตฺต\-สมฺปยุตฺตํ ธมฺมํ ปจฺจยา จิตฺต\-สมฺปยุตฺโต ธมฺโม อุปฺปชฺชติ นเหตุปจฺจยา. ตีณิ. (ปฏิจฺจสทิสา.)}

\prose{จิตฺต\-วิปฺปยุตฺตํ ธมฺมํ ปจฺจยา จิตฺต\-วิปฺปยุตฺโต ธมฺโม อุปฺปชฺชติ นเหตุปจฺจยา. เอกํ มหาภูตํ ฯเปฯ. (ยาว อสญฺญสตฺตา.) (1)}

\prose{จิตฺต\-วิปฺปยุตฺตํ ธมฺมํ ปจฺจยา จิตฺต\-สมฺปยุตฺโต ธมฺโม อุปฺปชฺชติ นเหตุปจฺจยา. จกฺขายตนํ ปจฺจยา จกฺขุ\-วิญฺญาณสหคตา ขนฺธา ฯเปฯ. กายายตนํ ฯเปฯ. วตฺถุํ ปจฺจยา อเหตุกา จิตฺต\-สมฺปยุตฺตกา ขนฺธา. ปฏิสนฺธิ\-กฺขเณ ฯเปฯ. วตฺถุํ ปจฺจยา วิจิกิจฺฉา\-สหคโต อุทฺธจฺจ\-สหคโต โมโห. (2)}

\prose{จิตฺต\-วิปฺปยุตฺตํ ธมฺมํ ปจฺจยา จิตฺต\-สมฺปยุตฺโต จ จิตฺต\-วิปฺปยุตฺโต จ ธมฺมา อุปฺปชฺชนฺติ นเหตุปจฺจยา. วตฺถุํ ปจฺจยา อเหตุกา จิตฺต\-สมฺปยุตฺตกา ขนฺธา. มหาภูเต ปจฺจยา จิตฺต\-สมุฏฺฐานํ รูปํ. ปฏิสนฺธิ\-กฺขเณ ฯเปฯ. (3)}

\prose{จิตฺต\-สมฺปยุตฺตญฺจ จิตฺต\-วิปฺปยุตฺตญฺจ ธมฺมํ ปจฺจยา จิตฺต\-สมฺปยุตฺโต ธมฺโม อุปฺปชฺชติ นเหตุปจฺจยา. จกฺขุ\-วิญฺญาณสหคตํ เอกํ ขนฺธญฺจ จกฺขายตนญฺจ ปจฺจยา เทฺว ขนฺธา. เทฺว ขนฺเธ จ ฯเปฯ. กาย\-วิญฺญาณสหคตํ เอกํ ขนฺธญฺจ กายายตนญฺจ ปจฺจยา เทฺว ขนฺธา. เทฺว ขนฺเธ จ ฯเปฯ. อเหตุกํ จิตฺต\-สมฺปยุตฺตํ เอกํ ขนฺธญฺจ วตฺถุญฺจ ปจฺจยา เทฺว ขนฺธา. เทฺว ขนฺเธ จ ฯเปฯ. อเหตุก\-ปฏิสนฺธิ\-กฺขเณ ฯเปฯ. วิจิกิจฺฉา\-สหคเต อุทฺธจฺจ\-สหคเต ขนฺเธ จ วตฺถุญฺจ ปจฺจยา วิจิกิจฺฉา\-สหคโต อุทฺธจฺจ\-สหคโต โมโห. (เอวํ เทฺว ปญฺหา ปวตฺติ\-ปฏิสนฺธิ กาตพฺพา. สํขิตฺตํ.)}

\csromanmarkh{2. ปจฺจย\-ปจฺจนีย}\csromanmarkhii{2. สงฺขฺยาวาร}\csromanheadernomark{2}{2. ปจฺจย\-ปจฺจนีย 2. สงฺขฺยาวาร}
\proseitem{159}{\csromanfolio{430}นเหตุยา นว, นอารมฺมเณ ตีณิ, นอธิปติยา นว, นอนนฺตเร ตีณิ, นสมนนฺตเร ตีณิ, นอญฺญมญฺเญ ตีณิ, นอุปนิสฺสเย ตีณิ, นปุเรชาเต นว, นปจฺฉาชาเต นว, นอาเสวเน นว, นกมฺเม จตฺตาริ, นวิปาเก นว, นอาหาเร เอกํ, นอินฺทฺริเย เอกํ, นฌาเน จตฺตาริ, นมคฺเค นว, นสมฺปยุตฺเต ตีณิ, นวิปฺปยุตฺเต เทฺว, โนนตฺถิยา ตีณิ, โนวิคเต ตีณิ.}

\proseitem{160}{เหตุปจฺจยา นอารมฺมเณ ตีณิ ฯเปฯ นกมฺเม ตีณิ ฯเปฯ นวิปฺปยุตฺเต เอกํ, โนนตฺถิยา ตีณิ, โนวิคเต ตีณิ.}

\vspace{\tipitakaparskip}
\csromancenter{นเหตุปจฺจยา อารมฺมเณ ตีณิ ฯเปฯ มคฺเค ตีณิ ฯเปฯ อวิคเต นว.}
\csromancenter{(นิสฺสยวาโร ปจฺจยวาร\-สทิโส.)}
\csromanmarkh{58. จิตฺต\-สมฺปยุตฺตทุก}\csromanmarkhii{5. สํสฏฺฐวาร}\csromanheadernomark{2}{58. จิตฺต\-สมฺปยุตฺตทุก 5. สํสฏฺฐวาร}
\prose{1. ปจฺจยา\-นุโลมาทิ}

\proseitem{162}{จิตฺต\-สมฺปยุตฺตํ ธมฺมํ สํสฏฺโฐ จิตฺต\-สมฺปยุตฺโต ธมฺโม อุปฺปชฺชติ เหตุปจฺจยา. จิตฺต\-สมฺปยุตฺตํ เอกํ ขนฺธํ สํสฏฺฐา เทฺว ขนฺธา. เทฺว ขนฺเธ ฯเปฯ. ปฏิสนฺธิ\-กฺขเณ ฯเปฯ.}

\prose{เหตุยา เอกํ, อารมฺมเณ เอกํ, อธิปติยา เอกํ, (สพฺพตฺถ เอกํ,) อวิคเต เอกํ.}

\prose{\csromanfolio{431}นเหตุยา เอกํ, นอธิปติยา เอกํ, นปุเรชาเต เอกํ, นปจฺฉาชาเต เอกํ, นอาเสวเน เอกํ, นกมฺเม เอกํ, นวิปาเก เอกํ, นฌาเน เอกํ, นมคฺเค เอกํ, นวิปฺปยุตฺเต เอกํ.}

\prose{(เอวํ อิตเร เทฺว คณนาปิ สมฺปยุตฺตวาโรปิ กาตพฺโพ.)}

\csromanmarkh{58. จิตฺต\-สมฺปยุตฺตทุก}\csromanmarkhii{7. ปญฺหาวาร}\csromanheadernomark{2}{58. จิตฺต\-สมฺปยุตฺตทุก 7. ปญฺหาวาร}
\csromanmarkh{1. ปจฺจยานุโลม}\csromanmarkhii{1. วิภงฺควาร}\csromanheadernomark{2}{1. \textbf{ปจฺจยานุโลม 1. วิภงฺควาร}}
\csromanheader{2}{\textbf{เหตุ}}
\proseitem{163}{จิตฺต\-สมฺปยุตฺโต ธมฺโม จิตฺต\-สมฺปยุตฺตสฺส ธมฺมสฺส เหตุปจฺจเยน ปจฺจโย. จิตฺต\-สมฺปยุตฺตา เหตู สมฺปยุตฺตกานํ ขนฺธานํ เหตุปจฺจเยน ปจฺจโย. ปฏิสนฺธิ\-กฺขเณ. (1)}

\prose{จิตฺต\-สมฺปยุตฺโต ธมฺโม จิตฺต\-วิปฺปยุตฺตสฺส ธมฺมสฺส เหตุปจฺจเยน ปจฺจโย. จิตฺต\-สมฺปยุตฺตา เหตู จิตฺต\-สมุฏฺฐานานํ รูปานํ เหตุปจฺจเยน ปจฺจโย. ปฏิสนฺธิ\-กฺขเณ ฯเปฯ. (2)}

\prose{จิตฺต\-สมฺปยุตฺโต ธมฺโม จิตฺต\-สมฺปยุตฺตสฺส จ จิตฺต\-วิปฺปยุตฺตสฺส จ ธมฺมสฺส เหตุปจฺจเยน ปจฺจโย. จิตฺต\-สมฺปยุตฺตา เหตู สมฺปยุตฺตกานํ ขนฺธานํ จิตฺตสมุฏฺฐานานญฺจ รูปานํ เหตุปจฺจเยน ปจฺจโย. ปฏิสนฺธิ\-กฺขเณ ฯเปฯ. (3)}

\proseitem{164}{จิตฺต\-สมฺปยุตฺโต ธมฺโม จิตฺต\-สมฺปยุตฺตสฺส ธมฺมสฺส อารมฺมณ\-ปจฺจเยน ปจฺจโย. ทานํ ฯเปฯ สีลํ ฯเปฯ อุโปสถกมฺมํ กตฺวา ตํ ปจฺจเวกฺขติ, อสฺสาเทติ อภินนฺทติ, ตํ อารพฺภ ราโค อุปฺปชฺชติ ฯเปฯ โทมนสฺสํ อุปฺปชฺชติ. ปุพฺเพ สุจิณฺณานิ ฯเปฯ. ฌานา วุฏฺฐหิตฺวา ฌานํ ฯเปฯ. อริยา มคฺคา วุฏฺฐหิตฺวา มคฺคํ ปจฺจเวกฺขนฺติ, ผลํ ปจฺจเวกฺขนฺติ. ปหีเน กิเลเส ฯเปฯ. วิกฺขมฺภิเต กิเลเส ปจฺจเวกฺขนฺติ. ปุพฺเพ สมุทาจิณฺเณ กิเลเส ชานนฺติ. จิตฺต\-สมฺปยุตฺเต ขนฺเธ อนิจฺจโต ฯเปฯ โทมนสฺสํ อุปฺปชฺชติ. เจโตปริย\-ญาเณน จิตฺต\-สมฺปยุตฺต\-สมงฺคิสฺส จิตฺตํ ชานนฺติ. อากาสา\-นญฺจา\-ยตนํ วิญฺญาณญฺจา\-ยตนสฺส ฯเปฯ. \csromanfolio{432}อากิญฺจญฺญา\-ยตนํ เนว\-สญฺญา\-นา\-สญฺญา\-ยตนสฺส ฯเปฯ. จิตฺต\-สมฺปยุตฺตา ขนฺธา อิทฺธิวิธ\-ญาณสฺส, เจโตปริย\-ญาณสฺส, ปุพฺเพ\-นิวาสา\-นุสฺสติ\-ญาณสฺส, ยถากมฺมูปคญาณสฺส, อนาคตํสญาณสฺส, อาวชฺชนาย อารมฺมณ\-ปจฺจเยน ปจฺจโย. (1)}

\proseitem{165}{จิตฺต\-วิปฺปยุตฺโต ธมฺโม จิตฺต\-สมฺปยุตฺตสฺส ธมฺมสฺส อารมฺมณ\-ปจฺจเยน ปจฺจโย. อริยา นิพฺพานํ ปจฺจเวกฺขนฺติ. นิพฺพานํ โคตฺรภุสฺส, โวทานสฺส, มคฺคสฺส, ผลสฺส, อาวชฺชนาย อารมฺมณ\-ปจฺจเยน ปจฺจโย. จกฺขุํ ฯเปฯ วตฺถุํ จิตฺต\-วิปฺปยุตฺเต ขนฺเธ อนิจฺจโต ฯเปฯ โทมนสฺสํ อุปฺปชฺชติ. ทิพฺเพน จกฺขุนา รูปํ ปสฺสติ. ทิพฺพาย โสตธาตุยา สทฺทํ สุณาติ. รูปายตนํ จกฺขุ\-วิญฺญาณสหคตานํ ขนฺธานํ ฯเปฯ. โผฏฺฐพฺพา\-ยตนํ ฯเปฯ. จิตฺต\-วิปฺปยุตฺตา ขนฺธา อิทฺธิวิธ\-ญาณสฺส ปุพฺเพ\-นิวาสา\-นุสฺสติ\-ญาณสฺส, อนาคตํสญาณสฺส, อาวชฺชนาย อารมฺมณ\-ปจฺจเยน ปจฺจโย. (1)}

\proseitem{166}{จิตฺต\-สมฺปยุตฺโต ธมฺโม จิตฺต\-สมฺปยุตฺตสฺส ธมฺมสฺส อธิปติ\-ปจฺจเยน ปจฺจโย. อารมฺมณา\-ธิปติ, สหชาตา\-ธิปติ.  อารมฺมณา\-ธิปติ\csromandash{}ทานํ ฯเปฯ สีลํ ฯเปฯ อุโปสถกมฺมํ กตฺวา ตํ ครุํ กตฺวา ปจฺจเวกฺขติ, อสฺสาเทติ อภินนฺทติ, ตํ ครุํ กตฺวา ราโค อุปฺปชฺชติ. ทิฏฺฐิ อุปฺปชฺชติ. ปุพฺเพ ฯเปฯ. ฌานา ฯเปฯ. อริยา มคฺคา วุฏฺฐหิตฺวา มคฺคํ ครุํ กตฺวา ปจฺจเวกฺขนฺติ. ผลํ ครุํ กตฺวา ปจฺจเวกฺขนฺติ. จิตฺต\-สมฺปยุตฺเต ขนฺเธ ครุํ กตฺวา อสฺสาเทติ อภินนฺทติ, ตํ ครุํ กตฺวา ราโค อุปฺปชฺชติ. ทิฏฺฐิ อุปฺปชฺชติ. สหชาตา\-ธิปติ\csromandash{} จิตฺต\-สมฺปยุตฺตา\-ธิปติ สมฺปยุตฺตกานํ ขนฺธานํ อธิปติ\-ปจฺจเยน ปจฺจโย. (1)}

\prose{จิตฺต\-สมฺปยุตฺโต ธมฺโม จิตฺต\-วิปฺปยุตฺตสฺส ธมฺมสฺส อธิปติ\-ปจฺจเยน ปจฺจโย. สหชาตา\-ธิปติ\csromandash{}จิตฺต\-สมฺปยุตฺตา\-ธิปติ จิตฺต\-สมุฏฺฐานานํ รูปานํ อธิปติ\-ปจฺจเยน ปจฺจโย. (2)}

\prose{จิตฺต\-สมฺปยุตฺโต ธมฺโม จิตฺต\-สมฺปยุตฺตสฺส จ จิตฺต\-วิปฺปยุตฺตสฺส จ ธมฺมสฺส อธิปติ\-ปจฺจเยน ปจฺจโย. สหชาตา\-ธิปติ\csromandash{} จิตฺต\-สมฺปยุตฺตา\-ธิปติ สมฺปยุตฺตกานํ ขนฺธานํ จิตฺตสมุฏฺฐานานญฺจ รูปานํ อธิปติ\-ปจฺจเยน ปจฺจโย. (3)}

\proseitem{167}{\csromanfolio{433}จิตฺต\-วิปฺปยุตฺโต ธมฺโม จิตฺต\-สมฺปยุตฺตสฺส ธมฺมสฺส อธิปติ\-ปจฺจเยน ปจฺจโย. อารมฺมณา\-ธิปติ\csromandash{}อริยา นิพฺพานํ ครุํ กตฺวา ปจฺจเวกฺขนฺติ. นิพฺพานํ โคตฺรภุสฺส, โวทานสฺส, มคฺคสฺส, ผลสฺส อธิปติ\-ปจฺจเยน ปจฺจโย. จกฺขุํ ฯเปฯ วตฺถุํ จิตฺต\-วิปฺปยุตฺเต ขนฺเธ ครุํ กตฺวา อสฺสาเทติ อภินนฺทติ, ตํ ครุํ กตฺวา ราโค อุปฺปชฺชติ. ทิฏฺฐิ อุปฺปชฺชติ. (1)}

\proseitem{168}{จิตฺต\-สมฺปยุตฺโต ธมฺโม จิตฺต\-สมฺปยุตฺตสฺส ธมฺมสฺส อนนฺตร\-ปจฺจเยน ปจฺจโย. ปุริมา ปุริมา จิตฺต\-สมฺปยุตฺตา ขนฺธา ฯเปฯ ผลสมาปตฺติยา อนนฺตร\-ปจฺจเยน ปจฺจโย. สมนนฺตร\-ปจฺจเยน ปจฺจโย. สหชาต\-ปจฺจเยน ปจฺจโย. สตฺต. (ปฏิจฺจสทิสา, ปญฺหาฆฏนา นตฺถิ.) อญฺญมญฺญ\-ปจฺจเยน ปจฺจโย. ฉ. (ปฏิจฺจสทิสา.) นิสฺสย\-ปจฺจเยน ปจฺจโย. สตฺต. (ปจฺจยวาร\-สทิสา, ปญฺหาฆฏนา นตฺถิ.)}

\csromanheader{2}{\textbf{อุปนิสฺสย}}
\proseitem{169}{จิตฺต\-สมฺปยุตฺโต ธมฺโม จิตฺต\-สมฺปยุตฺตสฺส ธมฺมสฺส อุปนิสฺสย\-ปจฺจเยน ปจฺจโย. อารมฺมณู\-ปนิสฺสโย, อนนฺตรู\-ปนิสฺสโย, ปกตู\-ปนิสฺสโย ฯเปฯ. ปกตู\-ปนิสฺสโย\csromandash{}สทฺธํ อุปนิสฺสาย ทานํ เทติ ฯเปฯ มานํ ชปฺเปติ, ทิฏฺฐิํ คณฺหาติ, สีลํ ฯเปฯ. ปตฺถนํ. กายิกํ สุขํ. กายิกํ ทุกฺขํ อุปนิสฺสาย ทานํ เทติ ฯเปฯ. ปาณํ หนติ ฯเปฯ สํฆํ ภินฺทติ. สทฺธา ฯเปฯ. กายิกํ ทุกฺขํ สทฺธาย ฯเปฯ มคฺคสฺส, ผลสมาปตฺติยา อุปนิสฺสย\-ปจฺจเยน ปจฺจโย. (1)}

\prose{จิตฺต\-วิปฺปยุตฺโต ธมฺโม จิตฺต\-สมฺปยุตฺตสฺส ธมฺมสฺส อุปนิสฺสย\-ปจฺจเยน ปจฺจโย. อารมฺมณู\-ปนิสฺสโย, ปกตู\-ปนิสฺสโย ฯเปฯ. ปกตู\-ปนิสฺสโย\csromandash{}อุตุํ. โภชนํ. เสนาสนํ อุปนิสฺสาย ทานํ เทติ ฯเปฯ สํฆํ ภินฺทติ. อุตุ. โภชนํ. เสนาสนํ สทฺธาย ฯเปฯ ผลสมาปตฺติยา อุปนิสฺสย\-ปจฺจเยน ปจฺจโย. (1)}

\proseitem{170}{จิตฺต\-วิปฺปยุตฺโต ธมฺโม จิตฺต\-สมฺปยุตฺตสฺส ธมฺมสฺส ปุเรชาต\-ปจฺจเยน ปจฺจโย. อารมฺมณ\-ปุเรชาตํ, วตฺถุ\-ปุเรชาตํ.  อารมฺมณ\-ปุเรชาตํ\csromandash{}จกฺขุํ ฯเปฯ วตฺถุํ อนิจฺจโต ฯเปฯ โทมนสฺสํ อุปฺปชฺชติ. ทิพฺเพน \csromanfolio{434}จกฺขุนา รูปํ ปสฺสติ. ทิพฺพาย โสตธาตุยา สทฺทํ สุณาติ. รูปายตนํ จกฺขุ\-วิญฺญาณสหคตานํ ขนฺธานํ ฯเปฯ โผฏฺฐพฺพา\-ยตนํ ฯเปฯ. วตฺถุ\-ปุเรชาตํ\csromandash{}จกฺขายตนํ จกฺขุ\-วิญฺญาณสหคตานํ ขนฺธานํ ฯเปฯ กายายตนํ ฯเปฯ. วตฺถุ จิตฺต\-สมฺปยุตฺตกานํ ขนฺธานํ ปุเรชาต\-ปจฺจเยน ปจฺจโย. (1)}

\proseitem{171}{จิตฺต\-สมฺปยุตฺโต ธมฺโม จิตฺต\-วิปฺปยุตฺตสฺส ธมฺมสฺส ปจฺฉาชาต\-ปจฺจเยน ปจฺจโย. (สํขิตฺตํ.) เอกํ. อาเสวน\-ปจฺจเยน ปจฺจโย. เอกํ.}

\proseitem{172}{จิตฺต\-สมฺปยุตฺโต ธมฺโม จิตฺต\-สมฺปยุตฺตสฺส ธมฺมสฺส กมฺมปจฺจเยน ปจฺจโย. สหชาตา, นานากฺขณิกา.  สหชาตา จิตฺต\-สมฺปยุตฺตา เจตนา สมฺปยุตฺตกานํ ขนฺธานํ กมฺมปจฺจเยน ปจฺจโย. ปฏิสนฺธิ\-กฺขเณ ฯเปฯ. นานากฺขณิกา จิตฺต\-สมฺปยุตฺตา เจตนา วิปากานํ ขนฺธานํ กมฺมปจฺจเยน ปจฺจโย. (1)}

\prose{จิตฺต\-สมฺปยุตฺโต ธมฺโม จิตฺต\-วิปฺปยุตฺตสฺส ธมฺมสฺส กมฺมปจฺจเยน ปจฺจโย. สหชาตา, นานากฺขณิกา.  สหชาตา จิตฺต\-สมฺปยุตฺตา เจตนา จิตฺต\-สมุฏฺฐานานํ รูปานํ กมฺมปจฺจเยน ปจฺจโย. ปฏิสนฺธิ\-กฺขเณ ฯเปฯ. นานากฺขณิกา จิตฺต\-สมฺปยุตฺตา เจตนา กฏตฺตารูปานํ กมฺมปจฺจเยน ปจฺจโย. (2)}

\prose{จิตฺต\-สมฺปยุตฺโต ธมฺโม จิตฺต\-สมฺปยุตฺตสฺส จ จิตฺต\-วิปฺปยุตฺตสฺส จ ธมฺมสฺส กมฺมปจฺจเยน ปจฺจโย. สหชาตา, นานากฺขณิกา.  สหชาตา จิตฺต\-สมฺปยุตฺตา เจตนา สมฺปยุตฺตกานํ ขนฺธานํ จิตฺตสมุฏฺฐานานญฺจ รูปานํ กมฺมปจฺจเยน ปจฺจโย. ปฏิสนฺธิ\-กฺขเณ ฯเปฯ. นานากฺขณิกา จิตฺต\-สมฺปยุตฺตา เจตนา วิปากานํ ขนฺธานํ กฏตฺตา จ รูปานํ กมฺมปจฺจเยน ปจฺจโย. (3)}

\csromanheader{2}{\textbf{วิปากาหาร}}
\proseitem{173}{จิตฺต\-สมฺปยุตฺโต ธมฺโม จิตฺต\-สมฺปยุตฺตสฺส ธมฺมสฺส วิปาก\-ปจฺจเยน ปจฺจโย. ตีณิ.}

\prose{จิตฺต\-สมฺปยุตฺโต ธมฺโม จิตฺต\-สมฺปยุตฺตสฺส ธมฺมสฺส อาหารปจฺจเยน ปจฺจโย. ตีณิ.}

\prose{\csromanfolio{435}จิตฺต\-วิปฺปยุตฺโต ธมฺโม จิตฺต\-วิปฺปยุตฺตสฺส ธมฺมสฺส อาหารปจฺจเยน ปจฺจโย. กพฬีกาโร อาหาโร อิมสฺส กายสฺส อาหารปจฺจเยน ปจฺจโย. (1)}

\csromanheader{2}{\textbf{อินฺทฺริยาทิ}}
\proseitem{174}{จิตฺต\-สมฺปยุตฺโต ธมฺโม จิตฺต\-สมฺปยุตฺตสฺส ธมฺมสฺส อินฺทฺริย\-ปจฺจเยน ปจฺจโย. ตีณิ.}

\prose{จิตฺต\-วิปฺปยุตฺโต ธมฺโม จิตฺต\-วิปฺปยุตฺตสฺส ธมฺมสฺส อินฺทฺริย\-ปจฺจเยน ปจฺจโย. รูปชีวิตินฺทฺริยํ กฏตฺตารูปานํ อินฺทฺริย\-ปจฺจเยน ปจฺจโย. (1)}

\prose{จิตฺต\-วิปฺปยุตฺโต ธมฺโม จิตฺต\-สมฺปยุตฺตสฺส ธมฺมสฺส อินฺทฺริย\-ปจฺจเยน ปจฺจโย. จกฺขุนฺทฺริยํ จกฺขุ\-วิญฺญาณสหคตานํ ขนฺธานํ อินฺทฺริย\-ปจฺจเยน ปจฺจโย ฯเปฯ. กายินฺทฺริยํ. (2)}

\prose{จิตฺต\-สมฺปยุตฺโต จ จิตฺต\-วิปฺปยุตฺโต จ ธมฺมา จิตฺต\-สมฺปยุตฺตสฺส ธมฺมสฺส อินฺทฺริย\-ปจฺจเยน ปจฺจโย. จกฺขุนฺทฺริยญฺจ อุเปกฺขินฺทฺริยญฺจ จกฺขุ\-วิญฺญาณสหคตานํ ขนฺธานํ อินฺทฺริย\-ปจฺจเยน ปจฺจโย ฯเปฯ. กายินฺทฺริยญฺจ สุขินฺทฺริยญฺจ ฯเปฯ. กายินฺทฺริยญฺจ ทุกฺขินฺทฺริยญฺจ กาย\-วิญฺญาณสหคตานํ ขนฺธานํ อินฺทฺริย\-ปจฺจเยน ปจฺจโย. (1)}

\prose{ฌานปจฺจเยน ปจฺจโย. ตีณิ. มคฺคปจฺจเยน ปจฺจโย. ตีณิ. สมฺปยุตฺต\-ปจฺจเยน ปจฺจโย. เอกํ.}

\prose{\textbf{วิปฺปยุตฺต}}

\proseitem{175}{จิตฺต\-สมฺปยุตฺโต ธมฺโม จิตฺต\-วิปฺปยุตฺตสฺส ธมฺมสฺส วิปฺปยุตฺต\-ปจฺจเยน ปจฺจโย. สหชาตํ, ปจฺฉาชาตํ. (สํขิตฺตํ.) (1)}

\prose{จิตฺต\-วิปฺปยุตฺโต ธมฺโม จิตฺต\-สมฺปยุตฺตสฺส ธมฺมสฺส วิปฺปยุตฺต\-ปจฺจเยน ปจฺจโย. สหชาตํ, ปุเรชาตํ.  สหชาตํ ปฏิสนฺธิ\-กฺขเณ วตฺถุ จิตฺต\-สมฺปยุตฺตกานํ ขนฺธานํ วิปฺปยุตฺต\-ปจฺจเยน ปจฺจโย. ปุเรชาตํ จกฺขายตนํ จกฺขุ\-วิญฺญาณสหคตานํ ขนฺธานํ วิปฺปยุตฺต\-ปจฺจเยน ปจฺจโย ฯเปฯ. กายายตนํ ฯเปฯ. วตฺถุ จิตฺต\-สมฺปยุตฺตกานํ ขนฺธานํ วิปฺปยุตฺต\-ปจฺจเยน ปจฺจโย. (1)}

\proseitem{176}{\csromanfolio{436}จิตฺต\-สมฺปยุตฺโต ธมฺโม จิตฺต\-สมฺปยุตฺตสฺส ธมฺมสฺส อตฺถิปจฺจเยน ปจฺจโย. เอกํ. (ปฏิจฺจสทิสํ.) (1)}

\prose{จิตฺต\-สมฺปยุตฺโต ธมฺโม จิตฺต\-วิปฺปยุตฺตสฺส ธมฺมสฺส อตฺถิปจฺจเยน ปจฺจโย. สหชาตํ, ปจฺฉาชาตํ. (สํขิตฺตํ.) (2)}

\prose{จิตฺต\-สมฺปยุตฺโต ธมฺโม จิตฺต\-สมฺปยุตฺตสฺส จ จิตฺต\-วิปฺปยุตฺตสฺส จ ธมฺมสฺส อตฺถิปจฺจเยน ปจฺจโย. (ปฏิจฺจสทิสํ.) (3)}

\proseitem{177}{จิตฺต\-วิปฺปยุตฺโต ธมฺโม จิตฺต\-วิปฺปยุตฺตสฺส ธมฺมสฺส อตฺถิปจฺจเยน ปจฺจโย. สหชาตํ, อาหารํ, อินฺทฺริยํ. (สํขิตฺตํ.) (1)}

\prose{จิตฺต\-วิปฺปยุตฺโต ธมฺโม จิตฺต\-สมฺปยุตฺตสฺส ธมฺมสฺส อตฺถิปจฺจเยน ปจฺจโย. สหชาตํ, ปุเรชาตํ.  สหชาตํ ปฏิสนฺธิ\-กฺขเณ วตฺถุ จิตฺต\-สมฺปยุตฺตกานํ ขนฺธานํ อตฺถิปจฺจเยน ปจฺจโย. ปุเรชาตํ จกฺขุํ ฯเปฯ วตฺถุํ อนิจฺจโต ฯเปฯ. (ปุเรชาต\-สทิสํ.) (2)}

\proseitem{178}{จิตฺต\-สมฺปยุตฺโต จ จิตฺต\-วิปฺปยุตฺโต จ ธมฺมา จิตฺต\-สมฺปยุตฺตสฺส ธมฺมสฺส อตฺถิปจฺจเยน ปจฺจโย. สหชาตํ, ปุเรชาตํ.  สหชาโต จกฺขุ\-วิญฺญาณ\-สหคโต เอโก ขนฺโธ จ จกฺขายตนญฺจ ทฺวินฺนํ ขนฺธานํ ฯเปฯ. กาย\-วิญฺญาณ\-สหคโต เอโก ขนฺโธ จ กายายตนญฺจ ทฺวินฺนํ ขนฺธานํ ฯเปฯ. จิตฺต\-สมฺปยุตฺโต เอโก ขนฺโธ จ วตฺถุ จ ทฺวินฺนํ ขนฺธานํ ฯเปฯ เทฺว ขนฺธา จ ฯเปฯ. ปฏิสนฺธิ\-กฺขเณ ฯเปฯ. (1)}

\prose{จิตฺต\-สมฺปยุตฺโต จ จิตฺต\-วิปฺปยุตฺโต จ ธมฺมา จิตฺต\-วิปฺปยุตฺตสฺส ธมฺมสฺส อตฺถิปจฺจเยน ปจฺจโย. สหชาตํ, ปจฺฉาชาตํ, อาหารํ, อินฺทฺริยํ.  สหชาตา จิตฺต\-สมฺปยุตฺตา ขนฺธา จ มหาภูตา จ จิตฺต\-สมุฏฺฐานานํ รูปานํ อตฺถิปจฺจเยน ปจฺจโย. ปฏิสนฺธิ\-กฺขเณ ฯเปฯ. ปจฺฉาชาตา จิตฺต\-สมฺปยุตฺตกา ขนฺธา จ กพฬีกาโร อาหาโร จ อิมสฺส กายสฺส อตฺถิปจฺจเยน ปจฺจโย. ปจฺฉาชาตา จิตฺต\-สมฺปยุตฺตกา ขนฺธา จ รูปชีวิตินฺทฺริยญฺจ กฏตฺตารูปานํ อตฺถิปจฺจเยน ปจฺจโย. (2)}

\csromanmarkh{58. จิตฺต\-สมฺปยุตฺตทุก}\csromanmarkhii{1. ปจฺจยานุโลม 2. สงฺขฺยาวาร}\csromanheadernomark{2}{58. จิตฺต\-สมฺปยุตฺตทุก 1. ปจฺจยานุโลม 2. สงฺขฺยาวาร}
\proseitem{179}{\csromanfolio{437}เหตุยา ตีณิ, อารมฺมเณ เทฺว, อธิปติยา จตฺตาริ, อนนฺตเร เอกํ, สมนนฺตเร เอกํ, สหชาเต สตฺต, อญฺญมญฺเญ ฉ, นิสฺสเย สตฺต, อุปนิสฺสเย เทฺว, ปุเรชาเต เอกํ, ปจฺฉาชาเต เอกํ, อาเสวเน เอกํ, กมฺเม ตีณิ, วิปาเก ตีณิ, อาหาเร จตฺตาริ, อินฺทฺริเย ฉ, ฌาเน ตีณิ, มคฺเค ตีณิ, สมฺปยุตฺเต เอกํ, วิปฺปยุตฺเต เทฺว, อตฺถิยา สตฺต, นตฺถิยา เอกํ, วิคเต เอกํ, อวิคเต สตฺต.}

\vspace{\tipitakaparskip}
\csromancenter{\textbf{อนุโลมํ.}}
\proseitem{180}{จิตฺต\-สมฺปยุตฺโต ธมฺโม จิตฺต\-สมฺปยุตฺตสฺส ธมฺมสฺส อารมฺมณ\-ปจฺจเยน ปจฺจโย. สหชาต\-ปจฺจเยน ปจฺจโย. อุปนิสฺสย\-ปจฺจเยน ปจฺจโย. กมฺมปจฺจเยน ปจฺจโย. (1)}

\prose{จิตฺต\-สมฺปยุตฺโต ธมฺโม จิตฺต\-วิปฺปยุตฺตสฺส ธมฺมสฺส สหชาต\-ปจฺจเยน ปจฺจโย. ปจฺฉาชาต\-ปจฺจเยน ปจฺจโย. กมฺมปจฺจเยน ปจฺจโย. (2)}

\prose{จิตฺต\-สมฺปยุตฺโต ธมฺโม จิตฺต\-สมฺปยุตฺตสฺส จ จิตฺต\-วิปฺปยุตฺตสฺส จ ธมฺมสฺส สหชาต\-ปจฺจเยน ปจฺจโย. กมฺมปจฺจเยน ปจฺจโย. (3)}

\proseitem{181}{จิตฺต\-วิปฺปยุตฺโต ธมฺโม จิตฺต\-วิปฺปยุตฺตสฺส ธมฺมสฺส สหชาต\-ปจฺจเยน ปจฺจโย. อาหารปจฺจเยน ปจฺจโย. อินฺทฺริย\-ปจฺจเยน ปจฺจโย. (1)}

\prose{จิตฺต\-วิปฺปยุตฺโต ธมฺโม จิตฺต\-สมฺปยุตฺตสฺส ธมฺมสฺส อารมฺมณ\-ปจฺจเยน ปจฺจโย. สหชาต\-ปจฺจเยน ปจฺจโย. อุปนิสฺสย\-ปจฺจเยน ปจฺจโย. ปุเรชาต\-ปจฺจเยน ปจฺจโย. (2)}

\prose{จิตฺต\-สมฺปยุตฺโต จ จิตฺต\-วิปฺปยุตฺโต จ ธมฺมา จิตฺต\-สมฺปยุตฺตสฺส ธมฺมสฺส สหชาตํ, ปุเรชาตํ. (1)}

\prose{จิตฺต\-สมฺปยุตฺโต จ จิตฺต\-วิปฺปยุตฺโต จ ธมฺมา จิตฺต\-วิปฺปยุตฺตสฺส ธมฺมสฺส สหชาตํ, ปจฺฉาชาตํ, อาหารํ, อินฺทฺริยํ. (2)}

\csromanmarkh{2. ปจฺจย\-ปจฺจนีย}\csromanmarkhii{2. สงฺขฺยาวาร}\csromanheadernomark{2}{2. ปจฺจย\-ปจฺจนีย 2. สงฺขฺยาวาร}
\proseitem{182}{\csromanfolio{438}นเหตุยา สตฺต, นอารมฺมเณ สตฺต, นอธิปติยา สตฺต, นอนนฺตเร สตฺต, นสมนนฺตเร สตฺต, นสหชาเต ฉ, นอญฺญมญฺเญ ฉ, นนิสฺสเย ฉ, นอุปนิสฺสเย สตฺต, นปุเรชาเต สตฺต, (สพฺพตฺถ สตฺต.) นสมฺปยุตฺเต ฉ, นวิปฺปยุตฺเต ปญฺจ, โนอตฺถิยา จตฺตาริ, โนนตฺถิยา สตฺต, โนวิคเต สตฺต, โนอวิคเต จตฺตาริ.}

\proseitem{183}{เหตุปจฺจยา นอารมฺมเณ ตีณิ, นอธิปติยา ตีณิ, นอนนฺตเร ตีณิ, นสมนนฺตเร ตีณิ, นอญฺญมญฺเญ เอกํ, นอุปนิสฺสเย ตีณิ, (สพฺพตฺถ ตีณิ.) นสมฺปยุตฺเต เอกํ, นวิปฺปยุตฺเต เอกํ, โนนตฺถิยา ตีณิ, โนวิคเต ตีณิ.}

\proseitemwithcloser{184}{นเหตุปจฺจยา อารมฺมเณ เทฺว, อธิปติยา จตฺตาริ, (อนุโลมมาติกา กาตพฺพา.) ฯเปฯ อวิคเต สตฺต.}{จิตฺต\-สมฺปยุตฺตทุกํ นิฏฺฐิตํ.}
\csromanmatikaanchor{mtk.58}
\csromantocmarkref{subsection}{59. จิตฺตสํสฏฺฐทุก}{mtk.58}
\bookmark[dest={mtk.58},level=2]{59. จิตฺตสํสฏฺฐทุก}
\csromanmarkh{59. จิตฺต\-สํสฏฺฐทุก}\csromanmarkhii{1. ปฏิจฺจวาร}\csromanheadernomark{2}{59. จิตฺต\-สํสฏฺฐทุก 1. ปฏิจฺจวาร}
\csromanmarkh{1. ปจฺจยานุโลม}\csromanmarkhii{1. วิภงฺควาร}\csromanheadernomark{2}{1. \textbf{ปจฺจยานุโลม 1. วิภงฺควาร}}
\proseitem{185}{จิตฺต\-สํสฏฺฐํ ธมฺมํ ปฏิจฺจ จิตฺตสํสฏฺโฐ ธมฺโม อุปฺปชฺชติ เหตุปจฺจยา. จิตฺต\-สํสฏฺฐํ เอกํ ขนฺธํ ปฏิจฺจ เทฺว ขนฺธา. เทฺว ขนฺเธ ฯเปฯ. ปฏิสนฺธิ\-กฺขเณ ฯเปฯ. (1)}

\prose{จิตฺต\-สํสฏฺฐํ ธมฺมํ ปฏิจฺจ จิตฺต\-วิสํสฏฺโฐ ธมฺโม อุปฺปชฺชติ เหตุปจฺจยา. จิตฺตสํสฏฺเฐ ขนฺเธ ปฏิจฺจ จิตฺต\-สมุฏฺฐานํ รูปํ. ปฏิสนฺธิ\-กฺขเณ ฯเปฯ. (2)}

\vspace{\tipitakaparskip}
\csromancenter{(จิตฺต\-สํสฏฺฐทุกํ ยถา จิตฺต\-สมฺปยุตฺตทุกํ, เอวํ กาตพฺพํ, นินฺนานากรณํ.)}
\nitthitam{จิตฺต\-สํสฏฺฐทุกํ นิฏฺฐิตํ.}
\csromanmatikaanchor{mtk.59}
\csromantocmarkref{subsection}{60. จิตฺตสมุฏฺฐานทุก}{mtk.59}
\bookmark[dest={mtk.59},level=2]{60. จิตฺตสมุฏฺฐานทุก}
\csromanmarkh{60. จิตฺต\-สมุฏฺฐานทุก}\csromanmarkhii{1. ปฏิจฺจวาร}\csromanheadernomark{2}{60. จิตฺต\-สมุฏฺฐานทุก 1. ปฏิจฺจวาร}
\csromanmarkh{1. ปจฺจยานุโลม}\csromanmarkhii{1. วิภงฺควาร}\csromanheadernomark{2}{1. \textbf{ปจฺจยานุโลม 1. วิภงฺควาร}}
\csromanheader{2}{\textbf{เหตุ}}
\proseitem{186}{\csromanfolio{439}จิตฺต\-สมุฏฺฐานํ ธมฺมํ ปฏิจฺจ จิตฺต\-สมุฏฺฐาโน ธมฺโม อุปฺปชฺชติ เหตุปจฺจยา. จิตฺต\-สมุฏฺฐานํ เอกํ ขนฺธํ ปฏิจฺจ เทฺว ขนฺธา จิตฺต\-สมุฏฺฐานญฺจ รูปํ. เทฺว ขนฺเธ ฯเปฯ. ปฏิสนฺธิ\-กฺขเณ จิตฺต\-สมุฏฺฐานํ เอกํ ขนฺธํ ปฏิจฺจ เทฺว ขนฺธา. เทฺว ขนฺเธ ฯเปฯ. เอกํ มหาภูตํ ฯเปฯ มหาภูเต ปฏิจฺจ จิตฺต\-สมุฏฺฐานํ รูปํ, อุปาทารูปํ. (1)}

\prose{จิตฺต\-สมุฏฺฐานํ ธมฺมํ ปฏิจฺจ โนจิตฺต\-สมุฏฺฐาโน ธมฺโม อุปฺปชฺชติ เหตุปจฺจยา. จิตฺต\-สมุฏฺฐาเน ขนฺเธ ปฏิจฺจ จิตฺตํ. ปฏิสนฺธิ\-กฺขเณ จิตฺต\-สมุฏฺฐาเน ขนฺเธ ปฏิจฺจ จิตฺตํ กฏตฺตา จ รูปํ.}

\prose{จิตฺต\-สมุฏฺฐานํ ธมฺมํ ปฏิจฺจ จิตฺต\-สมุฏฺฐาโน จ โนจิตฺต\-สมุฏฺฐาโน จ ธมฺมา อุปฺปชฺชนฺติ เหตุปจฺจยา. จิตฺต\-สมุฏฺฐานํ เอกํ ขนฺธํ ปฏิจฺจ เทฺว ขนฺธา จิตฺตญฺจ จิตฺต\-สมุฏฺฐานญฺจ รูปํ. เทฺว ขนฺเธ ฯเปฯ. ปฏิสนฺธิ\-กฺขเณ จิตฺต\-สมุฏฺฐานํ เอกํ ขนฺธํ ปฏิจฺจ เทฺว ขนฺธา จิตฺตญฺจ กฏตฺตา จ รูปํ. เทฺว ขนฺเธ ฯเปฯ. (3)}

\proseitem{187}{โนจิตฺต\-สมุฏฺฐานํ ธมฺมํ ปฏิจฺจ โนจิตฺต\-สมุฏฺฐาโน ธมฺโม อุปฺปชฺชติ เหตุปจฺจยา. ปฏิสนฺธิ\-กฺขเณ จิตฺตํ ปฏิจฺจ กฏตฺตารูปํ. จิตฺตํ ปฏิจฺจ วตฺถุ, วตฺถุ ปฏิจฺจ จิตฺตํ. เอกํ มหาภูตํ ฯเปฯ มหาภูเต ปฏิจฺจ กฏตฺตารูปํ, อุปาทารูปํ. (1)}

\prose{โนจิตฺต\-สมุฏฺฐานํ ธมฺมํ ปฏิจฺจ จิตฺต\-สมุฏฺฐาโน ธมฺโม อุปฺปชฺชติ เหตุปจฺจยา. จิตฺตํ ปฏิจฺจ สมฺปยุตฺตกา ขนฺธา จิตฺต\-สมุฏฺฐานญฺจ รูปํ. ปฏิสนฺธิ\-กฺขเณ จิตฺตํ ปฏิจฺจ สมฺปยุตฺตกา ขนฺธา. ปฏิสนฺธิ\-กฺขเณ วตฺถุํ ปฏิจฺจ จิตฺต\-สมุฏฺฐานา ขนฺธา. (2)}

\prose{\csromanfolio{440}โนจิตฺต\-สมุฏฺฐานํ ธมฺมํ ปฏิจฺจ จิตฺต\-สมุฏฺฐาโน จ โนจิตฺต\-สมุฏฺฐาโน จ ธมฺมา อุปฺปชฺชนฺติ เหตุปจฺจยา. ปฏิสนฺธิ\-กฺขเณ จิตฺตํ ปฏิจฺจ สมฺปยุตฺตกา ขนฺธา กฏตฺตา จ รูปํ. ปฏิสนฺธิ\-กฺขเณ วตฺถุํ ปฏิจฺจ จิตฺตํ สมฺปยุตฺตกา จ ขนฺธา. (3)}

\proseitem{188}{จิตฺต\-สมุฏฺฐานญฺจ โนจิตฺต\-สมุฏฺฐานญฺจ ธมฺมํ ปฏิจฺจ จิตฺต\-สมุฏฺฐาโน ธมฺโม อุปฺปชฺชติ เหตุปจฺจยา. จิตฺต\-สมุฏฺฐานํ เอกํ ขนฺธญฺจ จิตฺตญฺจ ปฏิจฺจ เทฺว ขนฺธา จิตฺต\-สมุฏฺฐานญฺจ รูปํ. เทฺว ขนฺเธ จ ฯเปฯ. ปฏิสนฺธิ\-กฺขเณ จิตฺต\-สมุฏฺฐานํ เอกํ ขนฺธญฺจ จิตฺตญฺจ ปฏิจฺจ เทฺว ขนฺธา. เทฺว ขนฺเธ จ ฯเปฯ. ปฏิสนฺธิ\-กฺขเณ จิตฺต\-สมุฏฺฐานํ เอกํ ขนฺธญฺจ วตฺถุญฺจ ปฏิจฺจ เทฺว ขนฺธา. เทฺว ขนฺเธ จ ฯเปฯ. (1)}

\prose{จิตฺต\-สมุฏฺฐานญฺจ โนจิตฺต\-สมุฏฺฐานญฺจ ธมฺมํ ปฏิจฺจ โนจิตฺต\-สมุฏฺฐาโน ธมฺโม อุปฺปชฺชติ เหตุปจฺจยา. ปฏิสนฺธิ\-กฺขเณ จิตฺต\-สมุฏฺฐาเน ขนฺเธ จ จิตฺตญฺจ ปฏิจฺจ กฏตฺตารูปํ. ปฏิสนฺธิ\-กฺขเณ จิตฺต\-สมุฏฺฐาเน ขนฺเธ จ มหาภูเต จ ปฏิจฺจ กฏตฺตารูปํ. ปฏิสนฺธิ\-กฺขเณ จิตฺต\-สมุฏฺฐาเน ขนฺเธ จ วตฺถุญฺจ ปฏิจฺจ จิตฺตํ. (2)}

\prose{จิตฺต\-สมุฏฺฐานญฺจ โนจิตฺต\-สมุฏฺฐานญฺจ ธมฺมํ ปฏิจฺจ จิตฺต\-สมุฏฺฐาโน จ โนจิตฺต\-สมุฏฺฐาโน จ ธมฺมา อุปฺปชฺชนฺติ เหตุปจฺจยา. ปฏิสนฺธิ\-กฺขเณ จิตฺต\-สมุฏฺฐานํ เอกํ ขนฺธญฺจ จิตฺตญฺจ ปฏิจฺจ เทฺว ขนฺธา กฏตฺตา จ รูปํ. เทฺว ขนฺเธ จ ฯเปฯ. ปฏิสนฺธิ\-กฺขเณ จิตฺต\-สมุฏฺฐานํ เอกํ ขนฺธญฺจ วตฺถุญฺจ ปฏิจฺจ เทฺว ขนฺธา จิตฺตญฺจ. เทฺว ขนฺเธ จ ฯเปฯ. (3)}

\proseitem{189}{จิตฺต\-สมุฏฺฐานํ ธมฺมํ ปฏิจฺจ จิตฺต\-สมุฏฺฐาโน ธมฺโม อุปฺปชฺชติ อารมฺมณ\-ปจฺจยา. จิตฺต\-สมุฏฺฐานํ เอกํ ขนฺธํ ปฏิจฺจ เทฺว ขนฺธา. เทฺว ขนฺเธ ฯเปฯ. ปฏิสนฺธิ\-กฺขเณ ฯเปฯ. (1)}

\prose{จิตฺต\-สมุฏฺฐานํ ธมฺมํ ปฏิจฺจ โนจิตฺต\-สมุฏฺฐาโน ธมฺโม อุปฺปชฺชติ อารมฺมณ\-ปจฺจยา. จิตฺต\-สมุฏฺฐาเน ขนฺเธ ปฏิจฺจ จิตฺตํ. ปฏิสนฺธิ\-กฺขเณ ฯเปฯ. (2)}

\prose{จิตฺต\-สมุฏฺฐานํ ธมฺมํ ปฏิจฺจ จิตฺต\-สมุฏฺฐาโน จ โนจิตฺต\-สมุฏฺฐาโน จ ธมฺมา อุปฺปชฺชนฺติ อารมฺมณ\-ปจฺจยา. จิตฺต\-สมุฏฺฐานํ เอกํ ขนฺธํ ปฏิจฺจ เทฺว ขนฺธา จิตฺตญฺจ. เทฺว ขนฺเธ ฯเปฯ. ปฏิสนฺธิ\-กฺขเณ ฯเปฯ. (3)}

\proseitem{190}{\csromanfolio{441}โนจิตฺต\-สมุฏฺฐานํ ธมฺมํ ปฏิจฺจ โนจิตฺต\-สมุฏฺฐาโน ธมฺโม อุปฺปชฺชติ อารมฺมณ\-ปจฺจยา. ปฏิสนฺธิ\-กฺขเณ วตฺถุํ ปฏิจฺจ จิตฺตํ. (1)}

\prose{โนจิตฺต\-สมุฏฺฐานํ ธมฺมํ ปฏิจฺจ จิตฺต\-สมุฏฺฐาโน ธมฺโม อุปฺปชฺชติ อารมฺมณ\-ปจฺจยา. จิตฺตํ ปฏิจฺจ สมฺปยุตฺตกา ขนฺธา. ปฏิสนฺธิ\-กฺขเณ จิตฺตํ ปฏิจฺจ สมฺปยุตฺตกา ขนฺธา. ปฏิสนฺธิ\-กฺขเณ วตฺถุํ ปฏิจฺจ จิตฺต\-สมุฏฺฐานา ขนฺธา. (2)}

\prose{โนจิตฺต\-สมุฏฺฐานํ ธมฺมํ ปฏิจฺจ จิตฺต\-สมุฏฺฐาโน จ โนจิตฺต\-สมุฏฺฐาโน จ ธมฺมา อุปฺปชฺชนฺติ อารมฺมณ\-ปจฺจยา. ปฏิสนฺธิ\-กฺขเณ วตฺถุํ ปฏิจฺจ จิตฺตํ สมฺปยุตฺตกา จ ขนฺธา. (3)}

\proseitem{191}{จิตฺต\-สมุฏฺฐานญฺจ โนจิตฺต\-สมุฏฺฐานญฺจ ธมฺมํ ปฏิจฺจ จิตฺต\-สมุฏฺฐาโน ธมฺโม อุปฺปชฺชติ อารมฺมณ\-ปจฺจยา. จิตฺต\-สมุฏฺฐานํ เอกํ ขนฺธญฺจ จิตฺตญฺจ ปฏิจฺจ เทฺว ขนฺธา. เทฺว ขนฺเธ ฯเปฯ. ปฏิสนฺธิ\-กฺขเณ จิตฺต\-สมุฏฺฐานํ เอกํ ขนฺธญฺจ จิตฺตญฺจ ปฏิจฺจ เทฺว ขนฺธา. เทฺว ขนฺเธ ฯเปฯ. ปฏิสนฺธิ\-กฺขเณ จิตฺต\-สมุฏฺฐานํ เอกํ ขนฺธญฺจ วตฺถุญฺจ ปฏิจฺจ เทฺว ขนฺธา. เทฺว ขนฺเธ ฯเปฯ. (1)}

\prose{จิตฺต\-สมุฏฺฐานญฺจ โนจิตฺต\-สมุฏฺฐานญฺจ ธมฺมํ ปฏิจฺจ โนจิตฺต\-สมุฏฺฐาโน ธมฺโม อุปฺปชฺชติ อารมฺมณ\-ปจฺจยา. ปฏิสนฺธิ\-กฺขเณ จิตฺต\-สมุฏฺฐาเน ขนฺเธ จ วตฺถุญฺจ ปฏิจฺจ จิตฺตํ. (2)}

\prose{จิตฺต\-สมุฏฺฐานญฺจ โนจิตฺต\-สมุฏฺฐานญฺจ ธมฺมํ ปฏิจฺจ จิตฺต\-สมุฏฺฐาโน จ โนจิตฺต\-สมุฏฺฐาโน จ ธมฺมา อุปฺปชฺชนฺติ อารมฺมณ\-ปจฺจยา. ปฏิสนฺธิ\-กฺขเณ จิตฺต\-สมุฏฺฐานํ เอกํ ขนฺธญฺจ วตฺถุญฺจ ปฏิจฺจ เทฺว ขนฺธา จิตฺตญฺจ. เทฺว ขนฺเธ ฯเปฯ. (3)}

\csromanmarkh{1. ปจฺจยานุโลม}\csromanmarkhii{2. สงฺขฺยาวาร}\csromanheadernomark{2}{1. \textbf{ปจฺจยานุโลม 2. สงฺขฺยาวาร}}
\proseitem{192}{เหตุยา นว, อารมฺมเณ นว, อธิปติยา ปญฺจ, อนนฺตเร นว, สมนนฺตเร นว, สหชาเต นว, อญฺญมญฺเญ นว, นิสฺสเย นว, อุปนิสฺสเย นว, ปุเรชาเต ปญฺจ, อาเสวเน ปญฺจ, กมฺเม นว, วิปาเก นว, อาหาเร นว, อินฺทฺริเย นว, ฌาเน นว, มคฺเค นว, สมฺปยุตฺเต นว, (สพฺพตฺถ นว.) อวิคเต นว.}

\vspace{\tipitakaparskip}
\csromancenter{อนุโลมํ.}
\csromanmarkh{2. ปจฺจย\-ปจฺจนีย}\csromanmarkhii{1. วิภงฺควาร}\csromanheadernomark{2}{2. ปจฺจย\-ปจฺจนีย 1. วิภงฺควาร}
\proseitem{193}{\csromanfolio{442}จิตฺต\-สมุฏฺฐานํ ธมฺมํ ปฏิจฺจ จิตฺต\-สมุฏฺฐาโน ธมฺโม อุปฺปชฺชติ นเหตุปจฺจยา. อเหตุกํ จิตฺต\-สมุฏฺฐานํ เอกํ ขนฺธํ ปฏิจฺจ เทฺว ขนฺธา จิตฺต\-สมุฏฺฐานญฺจ รูปํ. เทฺว ขนฺเธ ฯเปฯ. อเหตุก\-ปฏิสนฺธิ\-กฺขเณ จิตฺต\-สมุฏฺฐานํ เอกํ ขนฺธํ ปฏิจฺจ เทฺว ขนฺธา. เทฺว ขนฺเธ ฯเปฯ. จิตฺต\-สมุฏฺฐานํ เอกํ มหาภูตํ ฯเปฯ. มหาภูเต ปฏิจฺจ จิตฺต\-สมุฏฺฐานํ รูปํ, อุปาทารูปํ. วิจิกิจฺฉา\-สหคเต อุทฺธจฺจ\-สหคเต ขนฺเธ ปฏิจฺจ วิจิกิจฺฉา\-สหคโต อุทฺธจฺจ\-สหคโต โมโห. (1)}

\prose{จิตฺต\-สมุฏฺฐานํ ธมฺมํ ปฏิจฺจ โนจิตฺต\-สมุฏฺฐาโน ธมฺโม อุปฺปชฺชติ นเหตุปจฺจยา. อเหตุเก จิตฺต\-สมุฏฺฐาเน ขนฺเธ ปฏิจฺจ จิตฺตํ. อเหตุก\-ปฏิสนฺธิ\-กฺขเณ จิตฺต\-สมุฏฺฐาเน ขนฺเธ ปฏิจฺจ จิตฺตํ กฏตฺตา จ รูปํ. (2)}

\prose{จิตฺต\-สมุฏฺฐานํ ธมฺมํ ปฏิจฺจ จิตฺต\-สมุฏฺฐาโน จ โนจิตฺต\-สมุฏฺฐาโน จ ธมฺมา อุปฺปชฺชนฺติ นเหตุปจฺจยา. อเหตุกํ จิตฺต\-สมุฏฺฐานํ เอกํ ขนฺธํ ปฏิจฺจ เทฺว ขนฺธา จิตฺตญฺจ จิตฺต\-สมุฏฺฐานญฺจ รูปํ. เทฺว ขนฺเธ ฯเปฯ. อเหตุก\-ปฏิสนฺธิ\-กฺขเณ ฯเปฯ. (3)}

\proseitem{194}{โนจิตฺต\-สมุฏฺฐานํ ธมฺมํ ปฏิจฺจ โนจิตฺต\-สมุฏฺฐาโน ธมฺโม อุปฺปชฺชติ นเหตุปจฺจยา. อเหตุก\-ปฏิสนฺธิ\-กฺขเณ จิตฺตํ ปฏิจฺจ กฏตฺตารูปํ. จิตฺตํ ปฏิจฺจ วตฺถุ, วตฺถุํ ปฏิจฺจ จิตฺตํ. เอกํ มหาภูตํ ฯเปฯ. (ยาว อสญฺญสตฺตา กาตพฺพา.) (1)}

\prose{โนจิตฺต\-สมุฏฺฐานํ ธมฺมํ ปฏิจฺจ จิตฺต\-สมุฏฺฐาโน ธมฺโม อุปฺปชฺชติ นเหตุปจฺจยา. อเหตุกํ จิตฺตํ ปฏิจฺจ สมฺปยุตฺตกา ขนฺธา จิตฺต\-สมุฏฺฐานญฺจ รูปํ. อเหตุก\-ปฏิสนฺธิ\-กฺขเณ จิตฺตํ ฯเปฯ. อเหตุก\-ปฏิสนฺธิ\-กฺขเณ วตฺถุํ ปฏิจฺจ จิตฺต\-สมุฏฺฐานา ขนฺธา. วิจิกิจฺฉา\-สหคเต อุทฺธจฺจ\-สหคเต ขนฺเธ ปฏิจฺจ วิจิกิจฺฉา\-สหคโต อุทฺธจฺจ\-สหคโต โมโห. (2)}

\prose{โนจิตฺต\-สมุฏฺฐานํ ธมฺมํ ปฏิจฺจ จิตฺต\-สมุฏฺฐาโน จ โนจิตฺต\-สมุฏฺฐาโน จ ธมฺมา อุปฺปชฺชนฺติ นเหตุปจฺจยา. อเหตุก\-ปฏิสนฺธิ\-กฺขเณ จิตฺตํ ปฏิจฺจ สมฺปยุตฺตกา ขนฺธา กฏตฺตา จ รูปํ. อเหตุก\-ปฏิสนฺธิ\-กฺขเณ วตฺถุํ ปฏิจฺจ จิตฺตํ สมฺปยุตฺตกา ขนฺธา จ. (3)}

\proseitem{195}{\csromanfolio{443}จิตฺต\-สมุฏฺฐานญฺจ โนจิตฺต\-สมุฏฺฐานญฺจ ธมฺมํ ปฏิจฺจ จิตฺต\-สมุฏฺฐาโน ธมฺโม อุปฺปชฺชติ นเหตุปจฺจยา. อเหตุกํ จิตฺต\-สมุฏฺฐานํ เอกํ ขนฺธญฺจ จิตฺตญฺจ ปฏิจฺจ เทฺว ขนฺธา จิตฺต\-สมุฏฺฐานญฺจ รูปํ. เทฺว ขนฺเธ ฯเปฯ. อเหตุก\-ปฏิสนฺธิ\-กฺขเณ จิตฺต\-สมุฏฺฐานํ เอกํ ขนฺธญฺจ จิตฺตญฺจ ปฏิจฺจ เทฺว ขนฺธา. เทฺว ขนฺเธ ฯเปฯ. อเหตุก\-ปฏิสนฺธิ\-กฺขเณ จิตฺต\-สมุฏฺฐานํ เอกํ ขนฺธญฺจ วตฺถุญฺจ ปฏิจฺจ เทฺว ขนฺธา. เทฺว ขนฺเธ ฯเปฯ. วิจิกิจฺฉา\-สหคเต อุทฺธจฺจ\-สหคเต ขนฺเธ จ จิตฺตญฺจ ปฏิจฺจ วิจิกิจฺฉา\-สหคโต อุทฺธจฺจ\-สหคโต โมโห. (1)}

\prose{จิตฺต\-สมุฏฺฐานญฺจ โนจิตฺต\-สมุฏฺฐานญฺจ ธมฺมํ ปฏิจฺจ โนจิตฺต\-สมุฏฺฐาโน ธมฺโม อุปฺปชฺชติ นเหตุปจฺจยา. อเหตุก\-ปฏิสนฺธิ\-กฺขเณ จิตฺต\-สมุฏฺฐาเน ขนฺเธ จ จิตฺตญฺจ ปฏิจฺจ กฏตฺตารูปํ. อเหตุก\-ปฏิสนฺธิ\-กฺขเณ จิตฺต\-สมุฏฺฐาเน ขนฺเธ จ มหาภูเต จ ปฏิจฺจ กฏตฺตารูปํ. อเหตุก\-ปฏิสนฺธิ\-กฺขเณ จิตฺต\-สมุฏฺฐาเน ขนฺเธ จ วตฺถุญฺจ ปฏิจฺจ จิตฺตํ. (2)}

\prose{จิตฺต\-สมุฏฺฐานญฺจ โนจิตฺต\-สมุฏฺฐานญฺจ ธมฺมํ ปฏิจฺจ จิตฺต\-สมุฏฺฐาโน จ โนจิตฺต\-สมุฏฺฐาโน จ ธมฺมา อุปฺปชฺชนฺติ นเหตุปจฺจยา. อเหตุก\-ปฏิสนฺธิ\-กฺขเณ จิตฺต\-สมุฏฺฐานํ เอกํ ขนฺธญฺจ จิตฺตญฺจ ปฏิจฺจ เทฺว ขนฺธา กฏตฺตา จ รูปํ. เทฺว ขนฺเธ ฯเปฯ. อเหตุก\-ปฏิสนฺธิ\-กฺขเณ จิตฺต\-สมุฏฺฐานํ เอกํ ขนฺธญฺจ วตฺถุญฺจ ปฏิจฺจ เทฺว ขนฺธา จิตฺตญฺจ. เทฺว ขนฺเธ ฯเปฯ. (3)}

\prose{นอารมฺมณ}

\proseitem{196}{จิตฺต\-สมุฏฺฐานํ ธมฺมํ ปฏิจฺจ จิตฺต\-สมุฏฺฐาโน ธมฺโม อุปฺปชฺชติ นอารมฺมณ\-ปจฺจยา. จิตฺต\-สมุฏฺฐาเน ขนฺเธ ปฏิจฺจ จิตฺต\-สมุฏฺฐานํ รูปํ. จิตฺต\-สมุฏฺฐานํ เอกํ มหาภูตํ ฯเปฯ. มหาภูเต ปฏิจฺจ จิตฺต\-สมุฏฺฐานํ รูปํ, อุปาทารูปํ. (1)}

\prose{จิตฺต\-สมุฏฺฐานํ ธมฺมํ ปฏิจฺจ โนจิตฺต\-สมุฏฺฐาโน ธมฺโม อุปฺปชฺชติ นอารมฺมณ\-ปจฺจยา. ปฏิสนฺธิ\-กฺขเณ จิตฺต\-สมุฏฺฐาเน ขนฺเธ ปฏิจฺจ กฏตฺตารูปํ. (2)}

\proseitem{197}{โนจิตฺต\-สมุฏฺฐานํ ธมฺมํ ปฏิจฺจ โนจิตฺต\-สมุฏฺฐาโน ธมฺโม อุปฺปชฺชติ นอารมฺมณ\-ปจฺจยา. ปฏิสนฺธิ\-กฺขเณ จิตฺตํ ปฏิจฺจ กฏตฺตารูปํ. จิตฺตํ ปฏิจฺจ วตฺถุ. เอกํ มหาภูตํ ฯเปฯ. (ยาว อสญฺญสตฺตา.) (1)}

\prose{โนจิตฺต\-สมุฏฺฐานํ ธมฺมํ ปฏิจฺจ จิตฺต\-สมุฏฺฐาโน ธมฺโม อุปฺปชฺชติ นอารมฺมณ\-ปจฺจยา. จิตฺตํ ปฏิจฺจ จิตฺต\-สมุฏฺฐานํ รูปํ. (2)}

\prose{\csromanfolio{444}จิตฺต\-สมุฏฺฐานญฺจ โนจิตฺต\-สมุฏฺฐานญฺจ ธมฺมํ ปฏิจฺจ จิตฺต\-สมุฏฺฐาโน ธมฺโม อุปฺปชฺชติ นอารมฺมณ\-ปจฺจยา. จิตฺต\-สมุฏฺฐาเน ขนฺเธ จ จิตฺตญฺจ ปฏิจฺจ จิตฺต\-สมุฏฺฐานํ รูปํ. จิตฺตญฺจ มหาภูเต จ ปฏิจฺจ จิตฺต\-สมุฏฺฐานํ รูปํ. (1)}

\prose{จิตฺต\-สมุฏฺฐานญฺจ โนจิตฺต\-สมุฏฺฐานญฺจ ธมฺมํ ปฏิจฺจ โนจิตฺต\-สมุฏฺฐาโน ธมฺโม อุปฺปชฺชติ นอารมฺมณ\-ปจฺจยา. ปฏิสนฺธิ\-กฺขเณ จิตฺต\-สมุฏฺฐาเน ขนฺเธ จ จิตฺตญฺจ ปฏิจฺจ กฏตฺตารูปํ. ปฏิสนฺธิ\-กฺขเณ จิตฺต\-สมุฏฺฐาเน ขนฺเธ จ มหาภูเต จ ปฏิจฺจ กฏตฺตารูปํ. (สํขิตฺตํ.) (2)}

\csromanmarkh{2. ปจฺจย\-ปจฺจนีย}\csromanmarkhii{2. สงฺขฺยาวาร}\csromanheadernomark{2}{2. ปจฺจย\-ปจฺจนีย 2. สงฺขฺยาวาร}
\proseitem{198}{นเหตุยา นว, นอารมฺมเณ ฉ, นอธิปติยา นว, นอนนฺตเร ฉ, นสมนนฺตเร ฉ, นอญฺญมญฺเญ ฉ, นอุปนิสฺสเย ฉ, นปุเรชาเต นว, นปจฺฉาชาเต นว, นอาเสวเน นว, นกมฺเม จตฺตาริ, นวิปาเก ปญฺจ, นอาหาเร เอกํ, นอินฺทฺริเย เอกํ, นฌาเน ฉ, นมคฺเค นว, นสมฺปยุตฺเต ฉ, นวิปฺปยุตฺเต ฉ, โนนตฺถิยา ฉ, โนวิคเต ฉ.}

\proseitem{199}{เหตุปจฺจยา นอารมฺมเณ ฉ, นอธิปติยา นว ฯเปฯ นกมฺเม ตีณิ, นวิปาเก ปญฺจ, นสมฺปยุตฺเต ฉ, นวิปฺปยุตฺเต ปญฺจ, โนนตฺถิยา ฉ, โนวิคเต ฉ.}

\proseitem{200}{นเหตุปจฺจยา อารมฺมเณ นว, อนนฺตเร นว, สมนนฺตเร นว ฯเปฯ ปุเรชาเต ปญฺจ, อาเสวเน ปญฺจ ฯเปฯ ฌาเน นว, มคฺเค ตีณิ ฯเปฯ อวิคเต นว.}

\vspace{\tipitakaparskip}
\csromancenter{(สหชาตวาโร ปฏิจฺจ\-วาร\-สทิโส.)}
\csromanmarkh{60. จิตฺต\-สมุฏฺฐานทุก}\csromanmarkhii{60. จิตฺต\-สมุฏฺฐานทุก 3. ปจฺจยวาร}\csromanheadernomark{2}{60. จิตฺต\-สมุฏฺฐานทุก 60. จิตฺต\-สมุฏฺฐานทุก 3. ปจฺจยวาร}
\csromanmarkh{1. ปจฺจยานุโลม}\csromanmarkhii{1. วิภงฺควาร}\csromanheadernomark{2}{1. \textbf{ปจฺจยานุโลม 1. วิภงฺควาร}}
\csromanheader{2}{\textbf{เหตุ}}
\proseitem{201}{\csromanfolio{445}จิตฺต\-สมุฏฺฐานํ ธมฺมํ ปจฺจยา จิตฺต\-สมุฏฺฐาโน ธมฺโม อุปฺปชฺชติ เหตุปจฺจยา. ตีณิ. (ปฏิจฺจ\-วาร\-สทิสํ.)}

\prose{โนจิตฺต\-สมุฏฺฐานํ ธมฺมํ ปจฺจยา โนจิตฺต\-สมุฏฺฐาโน ธมฺโม อุปฺปชฺชติ เหตุปจฺจยา. วตฺถุํ ปจฺจยา จิตฺตํ. ปฏิสนฺธิ\-กฺขเณ ฯเปฯ. (ปฏิจฺจ\-วาร\-สทิสา.) (1)}

\prose{โนจิตฺต\-สมุฏฺฐานํ ธมฺมํ ปจฺจยา จิตฺต\-สมุฏฺฐาโน ธมฺโม อุปฺปชฺชติ เหตุปจฺจยา. จิตฺตํ ปจฺจยา สมฺปยุตฺตกา ขนฺธา จิตฺต\-สมุฏฺฐานญฺจ รูปํ. วตฺถุํ ปจฺจยา จิตฺต\-สมุฏฺฐานา ขนฺธา. (ปฏิสนฺธิ\-กฺขเณ เทฺวปิ กาตพฺพา.) (2)}

\prose{โนจิตฺต\-สมุฏฺฐานํ ธมฺมํ ปจฺจยา จิตฺต\-สมุฏฺฐาโน จ โนจิตฺต\-สมุฏฺฐาโน จ ธมฺมา อุปฺปชฺชนฺติ เหตุปจฺจยา. วตฺถุํ ปจฺจยา จิตฺตํ สมฺปยุตฺตกา จ ขนฺธา. ปฏิสนฺธิ\-กฺขเณ เทฺว. (ปฏิจฺจ\-วาร\-สทิสํ.) (3)}

\proseitem{202}{จิตฺต\-สมุฏฺฐานญฺจ โนจิตฺต\-สมุฏฺฐานญฺจ ธมฺมํ ปจฺจยา จิตฺต\-สมุฏฺฐาโน ธมฺโม อุปฺปชฺชติ เหตุปจฺจยา. จิตฺต\-สมุฏฺฐานํ เอกํ ขนฺธญฺจ จิตฺตญฺจ ปจฺจยา เทฺว ขนฺธา. เทฺว ขนฺเธ ฯเปฯ. จิตฺต\-สมุฏฺฐานํ เอกํ ขนฺธญฺจ วตฺถุญฺจ ปจฺจยา เทฺว ขนฺธา. เทฺว ขนฺเธ ฯเปฯ. (ปฏิสนฺธิ\-กฺขเณ เทฺวปิ ปฏิจฺจ\-วาร\-สทิสา.) (1)}

\prose{จิตฺต\-สมุฏฺฐานญฺจ โนจิตฺต\-สมุฏฺฐานญฺจ ธมฺมํ ปจฺจยา โนจิตฺต\-สมุฏฺฐาโน ธมฺโม อุปฺปชฺชติ เหตุปจฺจยา. จิตฺต\-สมุฏฺฐาเน ขนฺเธ จ วตฺถุญฺจ ปจฺจยา จิตฺตํ. (ปฏิสนฺธิ\-กฺขเณ ตีณิปิ กาตพฺพา. ปฏิจฺจ\-วาร\-สทิสา.) (2)}

\prose{จิตฺต\-สมุฏฺฐานญฺจ โนจิตฺต\-สมุฏฺฐานญฺจ ธมฺมํ ปจฺจยา จิตฺต\-สมุฏฺฐาโน จ โนจิตฺต\-สมุฏฺฐาโน จ ธมฺมา อุปฺปชฺชนฺติ เหตุปจฺจยา. จิตฺต\-สมุฏฺฐานํ เอกํ ขนฺธญฺจ วตฺถุญฺจ ปจฺจยา เทฺว ขนฺธา จิตฺตญฺจ. เทฺว ขนฺเธ ฯเปฯ. (ปฏิสนฺธิ\-กฺขเณ เทฺวปิ กาตพฺพา. ปฏิจฺจ\-วาร\-สทิสา.) (3)}

\proseitem{203}{จิตฺต\-สมุฏฺฐานํ ธมฺมํ ปจฺจยา จิตฺต\-สมุฏฺฐาโน ธมฺโม อุปฺปชฺชติ อารมฺมณ\-ปจฺจยา. ตีณิ. (ปฏิจฺจ\-วาร\-สทิสา.)}

\prose{\csromanfolio{446}โนจิตฺต\-สมุฏฺฐานํ ธมฺมํ ปจฺจยา โนจิตฺต\-สมุฏฺฐาโน ธมฺโม อุปฺปชฺชติ อารมฺมณ\-ปจฺจยา. จกฺขายตนํ ปจฺจยา จกฺขุวิญฺญาณํ ฯเปฯ. กายายตนํ ปจฺจยา กายวิญฺญาณํ ฯเปฯ. วตฺถุํ ปจฺจยา จิตฺตํ. (1)}

\prose{โนจิตฺต\-สมุฏฺฐานํ ธมฺมํ ปจฺจยา จิตฺต\-สมุฏฺฐาโน ธมฺโม อุปฺปชฺชติ อารมฺมณ\-ปจฺจยา. จกฺขายตนํ ปจฺจยา จกฺขุ\-วิญฺญาณสหคตา ขนฺธา ฯเปฯ. กายายตนํ ปจฺจยา กาย\-วิญฺญาณสหคตา ขนฺธา. จิตฺตํ ปจฺจยา สมฺปยุตฺตกา ขนฺธา. วตฺถุํ ปจฺจยา จิตฺต\-สมุฏฺฐานา ขนฺธา. (ปฏิสนฺธิ\-กฺขเณ เทฺวปิ.) (2)}

\prose{โนจิตฺต\-สมุฏฺฐานํ ธมฺมํ ปจฺจยา จิตฺต\-สมุฏฺฐาโน จ โนจิตฺต\-สมุฏฺฐาโน จ ธมฺมา อุปฺปชฺชนฺติ อารมฺมณ\-ปจฺจยา. จกฺขายตนํ ปจฺจยา จกฺขุวิญฺญาณํ สมฺปยุตฺตกา จ ขนฺธา ฯเปฯ. กายายตนํ ฯเปฯ. วตฺถุํ ปจฺจยา จิตฺตํ สมฺปยุตฺตกา จ ขนฺธา. ปฏิสนฺธิ\-กฺขเณ ฯเปฯ. (3)}

\proseitem{204}{จิตฺต\-สมุฏฺฐานญฺจ โนจิตฺต\-สมุฏฺฐานญฺจ ธมฺมํ ปจฺจยา จิตฺต\-สมุฏฺฐาโน ธมฺโม อุปฺปชฺชติ อารมฺมณ\-ปจฺจยา. จิตฺต\-สมุฏฺฐานํ เอกํ ขนฺธญฺจ จิตฺตญฺจ ปจฺจยา เทฺว ขนฺธา. เทฺว ขนฺเธ ฯเปฯ. จิตฺต\-สมุฏฺฐานํ เอกํ ขนฺธญฺจ วตฺถุญฺจ ปจฺจยา เทฺว ขนฺธา. เทฺว ขนฺเธ ฯเปฯ. (ปฏิสนฺธิ\-กฺขเณ เทฺวปิ กาตพฺพา.) (1)}

\prose{จิตฺต\-สมุฏฺฐานญฺจ โนจิตฺต\-สมุฏฺฐานญฺจ ธมฺมํ ปจฺจยา โนจิตฺต\-สมุฏฺฐาโน ธมฺโม อุปฺปชฺชติ อารมฺมณ\-ปจฺจยา. จกฺขุ\-วิญฺญาณสหคเต ขนฺเธ จ จกฺขายตนญฺจ ปจฺจยา จกฺขุวิญฺญาณํ ฯเปฯ. กาย\-วิญฺญาณสหคเต ฯเปฯ. จิตฺต\-สมุฏฺฐาเน ขนฺเธ จ วตฺถุญฺจ ปจฺจยา จิตฺตํ. ปฏิสนฺธิ\-กฺขเณ ฯเปฯ. (2)}

\prose{จิตฺต\-สมุฏฺฐานญฺจ โนจิตฺต\-สมุฏฺฐานญฺจ ธมฺมํ ปจฺจยา จิตฺต\-สมุฏฺฐาโน จ โนจิตฺต\-สมุฏฺฐาโน จ ธมฺมา อุปฺปชฺชนฺติ อารมฺมณ\-ปจฺจยา. จกฺขุ\-วิญฺญาณสหคตํ เอกํ ขนฺธญฺจ จกฺขายตนญฺจ ปจฺจยา เทฺว ขนฺธา จกฺขุวิญฺญาณญฺจ. เทฺว ขนฺเธ ฯเปฯ. กาย\-วิญฺญาณสหคตํ ฯเปฯ. จิตฺต\-สมุฏฺฐานํ เอกํ ขนฺธญฺจ วตฺถุญฺจ ปจฺจยา เทฺว ขนฺธา จิตฺตญฺจ. เทฺว ขนฺเธ ฯเปฯ. ปฏิสนฺธิ\-กฺขเณ ฯเปฯ. (3) (สํขิตฺตํ.)}

\csromanmarkh{1. ปจฺจยานุโลม}\csromanmarkhii{2. สงฺขฺยาวาร}\csromanheadernomark{2}{1. \textbf{ปจฺจยานุโลม 2. สงฺขฺยาวาร}}
\vspace{\tipitakaparskip}
\csromancenter{เหตุยา นว, อารมฺมเณ นว, อธิปติยา นว, (สพฺพตฺถ นว.) อวิคเต นว.}
\csromanmarkh{60. จิตฺต\-สมุฏฺฐานทุก}\csromanmarkhii{2. ปจฺจย\-ปจฺจนีย 1. วิภงฺควาร}\csromanheadernomark{2}{60. จิตฺต\-สมุฏฺฐานทุก 2. ปจฺจย\-ปจฺจนีย 1. วิภงฺควาร}
\proseitem{206}{\csromanfolio{447}จิตฺต\-สมุฏฺฐานํ ธมฺมํ ปจฺจยา จิตฺต\-สมุฏฺฐาโน ธมฺโม อุปฺปชฺชติ นเหตุปจฺจยา. (สพฺเพ นวปิ ปญฺหา กาตพฺพา, ปฏิจฺจ\-วาร\-สทิสา. ปญฺจวิญฺญาณมฺปิ กาตพฺพํ. ตีณิเยว, โมโห.)}

\csromanmarkh{2. ปจฺจย\-ปจฺจนีย}\csromanmarkhii{2. สงฺขฺยาวาร}\csromanheadernomark{2}{2. ปจฺจย\-ปจฺจนีย 2. สงฺขฺยาวาร}
\vspace{\tipitakaparskip}
\csromancenter{\textbf{นเหตุ}ยา นว, นอารมฺมเณ ฉ, นอธิปติยา นว, นอนนฺตเร ฉ, นสมนนฺตเร ฉ, นอญฺญมญฺเญ ฉ, นฺเอาปนิสฺสเย ฉ, นปุเรชาเต นว, นปจฺฉาชาเต นว, นอาเสวเน นว, นกมฺเม จตฺตาริ, นวิปาเก นว, นอาหาเร เอกํ, ไนนฺทฺริเย เอกํ, นฌาเน นว, นมคฺเค นว, นสมฺปยุตฺเต ฉ, นวิปฺปยุตฺเต ฉ, โนนตฺถิยา ฉ, โนวิคเต ฉ.}
\csromancenter{(เอวํ อิตเร เทฺว คณนาปิ นิสฺสยวาโรปิ กาตพฺโพ.)}
\csromanmarkh{60. จิตฺต\-สมุฏฺฐานทุก}\csromanmarkhii{5. สํสฏฺฐวาร}\csromanheadernomark{2}{60. จิตฺต\-สมุฏฺฐานทุก 5. สํสฏฺฐวาร}
\proseitem{208}{จิตฺต\-สมุฏฺฐานํ ธมฺมํ สํสฏฺโฐ จิตฺต\-สมุฏฺฐาโน ธมฺโม อุปฺปชฺชติ เหตุปจฺจยา. จิตฺต\-สมุฏฺฐานํ เอกํ ขนฺธํ สํสฏฺฐา เทฺว ขนฺธา. เทฺว ขนฺเธ ฯเปฯ. ปฏิสนฺธิ\-กฺขเณ ฯเปฯ. (1)}

\prose{จิตฺต\-สมุฏฺฐานํ ธมฺมํ สํสฏฺโฐ โนจิตฺต\-สมุฏฺฐาโน ธมฺโม อุปฺปชฺชติ เหตุปจฺจยา. จิตฺต\-สมุฏฺฐาเน ขนฺเธ สํสฏฺฐํ จิตฺตํ. ปฏิสนฺธิ\-กฺขเณ ฯเปฯ. (2)}

\prose{จิตฺต\-สมุฏฺฐานํ ธมฺมํ สํสฏฺโฐ จิตฺต\-สมุฏฺฐาโน จ โนจิตฺต\-สมุฏฺฐาโน จ ธมฺมา อุปฺปชฺชนฺติ เหตุปจฺจยา. จิตฺต\-สมุฏฺฐานํ เอกํ ขนฺธํ สํสฏฺฐา เทฺว ขนฺธา จิตฺตญฺจ. เทฺว ขนฺเธ ฯเปฯ. ปฏิสนฺธิ\-กฺขเณ ฯเปฯ. (3)}

\proseitem{209}{\csromanfolio{448}โนจิตฺต\-สมุฏฺฐานํ ธมฺมํ สํสฏฺโฐ จิตฺต\-สมุฏฺฐาโน ธมฺโม อุปฺปชฺชติ เหตุปจฺจยา. จิตฺตสํสฏฺฐา สมฺปยุตฺตกา ขนฺธา. ปฏิสนฺธิ\-กฺขเณ ฯเปฯ. (1)}

\prose{จิตฺต\-สมุฏฺฐานญฺจ โนจิตฺต\-สมุฏฺฐานญฺจ ธมฺมํ สํสฏฺโฐ จิตฺต\-สมุฏฺฐาโน ธมฺโม อุปฺปชฺชติ เหตุปจฺจยา. จิตฺต\-สมุฏฺฐานํ เอกํ ขนฺธญฺจ จิตฺตญฺจ สํสฏฺฐา เทฺว ขนฺธา. เทฺว ขนฺเธ ฯเปฯ. ปฏิสนฺธิ\-กฺขเณ ฯเปฯ. (สํขิตฺตํ.) (1)}

\prose{เหตุยา ปญฺจ, อารมฺมเณ ปญฺจ, (สพฺพตฺถ ปญฺจ.) อวิคเต ปญฺจ.}

\prose{จิตฺต\-สมุฏฺฐานํ ธมฺมํ สํสฏฺโฐ จิตฺต\-สมุฏฺฐาโน ธมฺโม อุปฺปชฺชติ นเหตุปจฺจยา. (ปญฺจ ปญฺหา กาตพฺพา. ตีณิ. โมโห.)}

\prose{นเหตุยา ปญฺจ, นอธิปติยา ปญฺจ, นปุเรชาเต ปญฺจ, นปจฺฉาชาเต ปญฺจ, นอาเสวเน ปญฺจ, นกมฺเม ตีณิ, นวิปาเก ปญฺจ, นฌาเน ปญฺจ, นมคฺเค ปญฺจ, นวิปฺปยุตฺเต ปญฺจ.}

\vspace{\tipitakaparskip}
\csromancenter{(เอวํ อิตเร เทฺว คณนาปิ สมฺปยุตฺตวาโรปิ กาตพฺโพ.)}
\csromanmarkh{60. จิตฺต\-สมุฏฺฐานทุก}\csromanmarkhii{7. ปญฺหาวาร}\csromanheadernomark{2}{60. จิตฺต\-สมุฏฺฐานทุก 7. ปญฺหาวาร}
\csromanmarkh{1. ปจฺจยานุโลม}\csromanmarkhii{1. วิภงฺควาร}\csromanheadernomark{2}{1. \textbf{ปจฺจยานุโลม 1. วิภงฺควาร}}
\vspace{\tipitakaparskip}
\csromancenter{\textbf{เหตุ}}
\proseitem{210}{จิตฺต\-สมุฏฺฐาโน ธมฺโม จิตฺต\-สมุฏฺฐานสฺส ธมฺมสฺส เหตุปจฺจเยน ปจฺจโย. จิตฺต\-สมุฏฺฐานา เหตู สมฺปยุตฺตกานํ ขนฺธานํ จิตฺต\-สมุฏฺฐานานญฺจ รูปานํ เหตุปจฺจเยน ปจฺจโย. ปฏิสนฺธิ\-กฺขเณ ฯเปฯ. (1)}

\prose{จิตฺต\-สมุฏฺฐาโน ธมฺโม โนจิตฺต\-สมุฏฺฐานสฺส ธมฺมสฺส เหตุปจฺจเยน ปจฺจโย. จิตฺต\-สมุฏฺฐานา เหตู จิตฺตสฺส เหตุปจฺจเยน ปจฺจโย. ปฏิสนฺธิ\-กฺขเณ จิตฺต\-สมุฏฺฐานา เหตู จิตฺตสฺส กฏตฺตา จ รูปานํ เหตุปจฺจเยน ปจฺจโย. (2)}

\prose{จิตฺต\-สมุฏฺฐาโน ธมฺโม จิตฺต\-สมุฏฺฐานสฺส จ โนจิตฺต\-สมุฏฺฐานสฺส จ ธมฺมสฺส เหตุปจฺจเยน ปจฺจโย. จิตฺต\-สมุฏฺฐานา เหตู สมฺปยุตฺตกานํ ขนฺธานํ จิตฺตสฺส จ จิตฺต\-สมุฏฺฐานานญฺจ รูปานํ เหตุปจฺจเยน ปจฺจโย. ปฏิสนฺธิ\-กฺขเณ ฯเปฯ. (3)}

\csromanheader{2}{60. จิตฺต\-สมุฏฺฐานทุก อารมฺมณ}
\proseitem{211}{\csromanfolio{449}จิตฺต\-สมุฏฺฐาโน ธมฺโม จิตฺต\-สมุฏฺฐานสฺส ธมฺมสฺส อารมฺมณ\-ปจฺจเยน ปจฺจโย. จิตฺต\-สมุฏฺฐาเน ขนฺเธ อารพฺภ จิตฺต\-สมุฏฺฐานา ขนฺธา อุปฺปชฺชนฺติ. (มูลํ กาตพฺพํ.) จิตฺต\-สมุฏฺฐาเน ขนฺเธ อารพฺภ จิตฺตํ อุปฺปชฺชติ. (มูลํ กาตพฺพํ.) จิตฺต\-สมุฏฺฐาเน ขนฺเธ อารพฺภ จิตฺต\-สมุฏฺฐานา ขนฺธา จ จิตฺตญฺจ อุปฺปชฺชนฺติ. (3)}

\proseitem{212}{โนจิตฺต\-สมุฏฺฐาโน ธมฺโม โนจิตฺต\-สมุฏฺฐานสฺส ธมฺมสฺส อารมฺมณ\-ปจฺจเยน ปจฺจโย. อริยา มคฺคา วุฏฺฐหิตฺวา มคฺคํ ปจฺจเวกฺขนฺติ. ผลํ ฯเปฯ. นิพฺพานํ ปจฺจเวกฺขนฺติ. นิพฺพานํ โคตฺรภุสฺส, โวทานสฺส, มคฺคสฺส, ผลสฺส, อาวชฺชนาย อารมฺมณ\-ปจฺจเยน ปจฺจโย. จกฺขุํ ฯเปฯ. วตฺถุํ โนจิตฺต\-สมุฏฺฐาเน ขนฺเธ อนิจฺจโต ฯเปฯ วิปสฺสติ, อสฺสาเทติ อภินนฺทติ, ตํ อารพฺภ จิตฺตํ อุปฺปชฺชติ. ทิพฺเพน จกฺขุนา รูปํ ปสฺสติ. ทิพฺพาย โสตธาตุยา สทฺทํ สุณาติ. เจโตปริย\-ญาเณน โนจิตฺตสมุฏฺฐานจิตฺตสมงฺคิสฺส จิตฺตํ ชานาติ. อากาสา\-นญฺจา\-ยตนํ ฯเปฯ. อากิญฺจญฺญา\-ยตนํ ฯเปฯ. รูปายตนํ จกฺขุ\-วิญฺญาณสฺส ฯเปฯ. โผฏฺฐพฺพา\-ยตนํ ฯเปฯ. โนจิตฺต\-สมุฏฺฐานา ขนฺธา อิทฺธิวิธ\-ญาณสฺส, เจโตปริย\-ญาณสฺส, ปุพฺเพ\-นิวาสา\-นุสฺสติ\-ญาณสฺส, อนาคตํสญาณสฺส, อาวชฺชนาย อารมฺมณ\-ปจฺจเยน ปจฺจโย. (1)}

\prose{โนจิตฺต\-สมุฏฺฐาโน ธมฺโม จิตฺต\-สมุฏฺฐานสฺส ธมฺมสฺส อารมฺมณ\-ปจฺจเยน ปจฺจโย. อริยา มคฺคา ฯเปฯ. นิพฺพานํ ปจฺจเวกฺขนฺติ. (ปฐม\-คมน\-สทิสํ.) จกฺขุํ ฯเปฯ. วตฺถุํ โนจิตฺต\-สมุฏฺฐาเน ขนฺเธ อนิจฺจโต ฯเปฯ โทมนสฺสํ อุปฺปชฺชติ. ทิพฺเพน จกฺขุนา รูปํ ปสฺสติ ฯเปฯ. รูปายตนํ จกฺขุ\-วิญฺญาณสหคตานํ ขนฺธานํ ฯเปฯ. โผฏฺฐพฺพา\-ยตนํ ฯเปฯ. โนจิตฺต\-สมุฏฺฐานา ขนฺธา อิทฺธิวิธ\-ญาณสฺส, เจโตปริย\-ญาณสฺส, ปุพฺเพ\-นิวาสา\-นุสฺสติ\-ญาณสฺส, อนาคตํสญาณสฺส, อาวชฺชนาย อารมฺมณ\-ปจฺจเยน ปจฺจโย. (2)}

\prose{โนจิตฺต\-สมุฏฺฐาโน ธมฺโม จิตฺต\-สมุฏฺฐานสฺส จ โนจิตฺต\-สมุฏฺฐานสฺส จ ธมฺมสฺส อารมฺมณ\-ปจฺจเยน ปจฺจโย. อริยา มคฺคา ฯเปฯ. นิพฺพานํ ปจฺจเวกฺขนฺติ. (ปฐม\-คมน\-สทิสํ.) โนจิตฺต\-สมุฏฺฐาเน ขนฺเธ อนิจฺจโต ฯเปฯ วิปสฺสติ, อสฺสาเทติ, อภินนฺทติ, ตํ อารพฺภ จิตฺตญฺจ สมฺปยุตฺตกา จ ขนฺธา อุปฺปชฺชนฺติ. ทิพฺเพน จกฺขุนา รูปํ ปสฺสติ. ทิพฺพาย โสตธาตุยา สทฺทํ สุณาติ. รูปายตนํ จกฺขุ\-วิญฺญาณสฺส สมฺปยุตฺตกานญฺจ ขนฺธานํ ฯเปฯ. โผฏฺฐพฺพา\-ยตนํ ฯเปฯ. \csromanfolio{450}โนจิตฺต\-สมุฏฺฐานา ขนฺธา อิทฺธิวิธ\-ญาณสฺส, เจโตปริย\-ญาณสฺส, ปุพฺเพ\-นิวาสา\-นุสฺสติ\-ญาณสฺส, อนาคตํสญาณสฺส, อาวชฺชนาย อารมฺมณ\-ปจฺจเยน ปจฺจโย. (3)}

\prose{จิตฺต\-สมุฏฺฐาโน จ โนจิตฺต\-สมุฏฺฐาโน จ ธมฺมา จิตฺต\-สมุฏฺฐานสฺส ธมฺมสฺส อารมฺมณ\-ปจฺจเยน ปจฺจโย. ตีณิ. (อารพฺภ กาตพฺพา.)}

\proseitem{213}{จิตฺต\-สมุฏฺฐาโน ธมฺโม จิตฺต\-สมุฏฺฐานสฺส ธมฺมสฺส อธิปติ\-ปจฺจเยน ปจฺจโย. อารมฺมณา\-ธิปติ, สหชาตา\-ธิปติ.  อารมฺมณา\-ธิปติ\csromandash{}จิตฺต\-สมุฏฺฐาเน ขนฺเธ ครุํ กตฺวา จิตฺต\-สมุฏฺฐานา ขนฺธา อุปฺปชฺชนฺติ. สหชาตา\-ธิปติ\csromandash{}จิตฺต\-สมุฏฺฐานา\-ธิปติ สมฺปยุตฺตกานํ ขนฺธานํ จิตฺต\-สมุฏฺฐานานญฺจ รูปานํ อธิปติ\-ปจฺจเยน ปจฺจโย. (ตีณิปิ. อารมฺมณา\-ธิปติ สหชาตา\-ธิปติปิ กาตพฺพา.) (3)}

\prose{โนจิตฺต\-สมุฏฺฐาโน ธมฺโม โนจิตฺต\-สมุฏฺฐานสฺส ธมฺมสฺส อธิปติ\-ปจฺจเยน ปจฺจโย. อารมฺมณา\-ธิปติ\csromandash{}อริยา มคฺคา ฯเปฯ. นิพฺพานํ ครุํ กตฺวา ฯเปฯ. โนจิตฺต\-สมุฏฺฐาเน ขนฺเธ ครุํ กตฺวา อสฺสาเทติ อภินนฺทติ, ตํ ครุํ กตฺวา จิตฺตํ อุปฺปชฺชติ. (1)}

\prose{โนจิตฺต\-สมุฏฺฐาโน ธมฺโม จิตฺต\-สมุฏฺฐานสฺส ธมฺมสฺส อธิปติ\-ปจฺจเยน ปจฺจโย. อารมฺมณา\-ธิปติ, สหชาตา\-ธิปติ.  อารมฺมณา\-ธิปติ\csromandash{}อริยา มคฺคา ฯเปฯ. นิพฺพานํ ครุํ กตฺวา ปจฺจเวกฺขนฺติ ฯเปฯ. โนจิตฺต\-สมุฏฺฐาเน ขนฺเธ ครุํ กตฺวา อสฺสาเทติ อภินนฺทติ, ตํ ครุํ กตฺวา ราโค อุปฺปชฺชติ. ทิฏฺฐิ อุปฺปชฺชติ. สหชาตา\-ธิปติ\csromandash{}โนจิตฺตสมุฏฺฐานา\-ธิปติ สมฺปยุตฺตกานํ ขนฺธานํ จิตฺตสมุฏฺฐานานญฺจ รูปานํ อธิปติ\-ปจฺจเยน ปจฺจโย. (2)}

\prose{โนจิตฺต\-สมุฏฺฐาโน ธมฺโม จิตฺต\-สมุฏฺฐานสฺส จ โนจิตฺต\-สมุฏฺฐานสฺส จ ธมฺมสฺส อธิปติ\-ปจฺจเยน ปจฺจโย. อารมฺมณา\-ธิปติ\csromandash{}อริยา มคฺคา ฯเปฯ. นิพฺพานํ ครุํ กตฺวา ปจฺจเวกฺขนฺติ ฯเปฯ. โนจิตฺต\-สมุฏฺฐาเน ขนฺเธ ครุํ กตฺวา ฯเปฯ. จิตฺตญฺจ สมฺปยุตฺตกา จ ขนฺธา อุปฺปชฺชนฺติ. (3)}

\prose{จิตฺต\-สมุฏฺฐาโน จ โนจิตฺต\-สมุฏฺฐาโน จ ธมฺมา จิตฺต\-สมุฏฺฐานสฺส ธมฺมสฺส อธิปติ\-ปจฺจเยน ปจฺจโย. อารมฺมณา\-ธิปติ. ตีณิ. (อารมฺมณา\-ธิปติเยว.)}

\csromanheader{2}{60. จิตฺต\-สมุฏฺฐานทุก อนนฺตร\-สมนนฺตร}
\proseitem{214}{\csromanfolio{451}จิตฺต\-สมุฏฺฐาโน ธมฺโม จิตฺต\-สมุฏฺฐานสฺส ธมฺมสฺส อนนฺตร\-ปจฺจเยน ปจฺจโย. ตีณิ. (วุฏฺฐานํ นตฺถิ.)}

\prose{โนจิตฺต\-สมุฏฺฐาโน ธมฺโม โนจิตฺต\-สมุฏฺฐานสฺส ธมฺมสฺส อนนฺตร\-ปจฺจเยน ปจฺจโย. ปุริมํ ปุริมํ จิตฺตํ ปจฺฉิมสฺส ปจฺฉิมสฺส ฯเปฯ. นิโรธา วุฏฺฐหนฺตสฺส เนว\-สญฺญา\-นา\-สญฺญา\-ยตนํ ผลสมาปตฺติยา อนนฺตร\-ปจฺจเยน ปจฺจโย.}

\vspace{\tipitakaparskip}
\csromancenter{(อิตเร เทฺว คณนา, อิมสฺส สทิสาเยว กาตพฺพา.)}
\prose{จิตฺต\-สมุฏฺฐาโน จ โนจิตฺต\-สมุฏฺฐาโน จ ธมฺมา จิตฺต\-สมุฏฺฐานสฺส ธมฺมสฺส อนนฺตร\-ปจฺจเยน ปจฺจโย. (ตีณิ กาตพฺพา. วุฏฺฐานํ นตฺถิ.) สมนนฺตร\-ปจฺจเยน ปจฺจโย.}

\csromanheader{2}{\textbf{สหชาตาทิ}}
\proseitem{215}{จิตฺต\-สมุฏฺฐาโน ธมฺโม จิตฺต\-สมุฏฺฐานสฺส ธมฺมสฺส สหชาต\-ปจฺจเยน ปจฺจโย. (ปฏิจฺจ\-วาร\-สทิสํ.) อญฺญมญฺญ\-ปจฺจเยน ปจฺจโย. (ปฏิจฺจ\-วาร\-สทิสํ.) นิสฺสย\-ปจฺจเยน ปจฺจโย. (ปจฺจยวาร\-สทิสํ)}

\csromanheader{2}{\textbf{อุปนิสฺสย}}
\proseitem{216}{จิตฺต\-สมุฏฺฐาโน ธมฺโม จิตฺต\-สมุฏฺฐานสฺส ธมฺมสฺส อุปนิสฺสย\-ปจฺจเยน ปจฺจโย. อารมฺมณู\-ปนิสฺสโย, อนนฺตรู\-ปนิสฺสโย, ปกตู\-ปนิสฺสโย ฯเปฯ. ปกตู\-ปนิสฺสโย\csromandash{}(ตีณิ ปญฺหา กาตพฺพา.) (3)}

\prose{โนจิตฺต\-สมุฏฺฐาโน ธมฺโม โนจิตฺต\-สมุฏฺฐานสฺส ธมฺมสฺส อุปนิสฺสย\-ปจฺจเยน ปจฺจโย. อารมฺมณู\-ปนิสฺสโย, อนนฺตรู\-ปนิสฺสโย, ปกตู\-ปนิสฺสโย ฯเปฯ. ปกตู\-ปนิสฺสโย\csromandash{}อุตุํ. โภชนํ. เสนาสนํ จิตฺตํ อุปนิสฺสาย ทานํ เทติ ฯเปฯ สํฆํ ภินฺทติ. อุตุ. โภชนํ. เสนาสนํ จิตฺตํ จิตฺตสฺส อุปนิสฺสย\-ปจฺจเยน ปจฺจโย. (1)}

\prose{โนจิตฺต\-สมุฏฺฐาโน ธมฺโม จิตฺต\-สมุฏฺฐานสฺส ธมฺมสฺส อุปนิสฺสย\-ปจฺจเยน ปจฺจโย. อารมฺมณู\-ปนิสฺสโย, อนนฺตรู\-ปนิสฺสโย, ปกตู\-ปนิสฺสโย ฯเปฯ. ปกตู\-ปนิสฺสโย\csromandash{}อุตุํ. โภชนํ. เสนาสนํ จิตฺตํ \csromanfolio{452}อุปนิสฺสาย ทานํ เทติ ฯเปฯ สํฆํ ภินฺทติ. อุตุ. โภชนํ. เสนาสนํ จิตฺตํ สทฺธาย ฯเปฯ มคฺคสฺส ผลสมาปตฺติยา อุปนิสฺสย\-ปจฺจเยน ปจฺจโย. (2)}

\prose{โนจิตฺต\-สมุฏฺฐาโน ธมฺโม จิตฺต\-สมุฏฺฐานสฺส จ โนจิตฺต\-สมุฏฺฐานสฺส จ ธมฺมสฺส อุปนิสฺสย\-ปจฺจเยน ปจฺจโย. อารมฺมณู\-ปนิสฺสโย, อนนฺตรู\-ปนิสฺสโย, ปกตู\-ปนิสฺสโย ฯเปฯ. ปกตู\-ปนิสฺสโย\csromandash{}อุตุํ. โภชนํ. เสนาสนํ จิตฺตํ อุปนิสฺสาย ทานํ เทติ ฯเปฯ สํฆํ ภินฺทติ. อุตุ. โภชนํ. เสนาสนํ จิตฺตํ จิตฺต\-สมุฏฺฐานานํ ขนฺธานํ จิตฺตสฺส จ อุปนิสฺสย\-ปจฺจเยน ปจฺจโย. (3)}

\prose{จิตฺต\-สมุฏฺฐาโน จ โนจิตฺต\-สมุฏฺฐาโน จ ธมฺมา จิตฺต\-สมุฏฺฐานสฺส ธมฺมสฺส อุปนิสฺสย\-ปจฺจเยน ปจฺจโย. อารมฺมณู\-ปนิสฺสโย, อนนฺตรู\-ปนิสฺสโย, ปกตู\-ปนิสฺสโย ฯเปฯ. ปกตู\-ปนิสฺสโย\csromandash{}ตีณิ.}

\proseitem{217}{จิตฺต\-สมุฏฺฐาโน ธมฺโม จิตฺต\-สมุฏฺฐานสฺส ธมฺมสฺส ปุเรชาต\-ปจฺจเยน ปจฺจโย. อารมฺมณ\-ปุเรชาตํ\csromandash{}จิตฺต\-สมุฏฺฐาเน รูเป ฯเปฯ. โผฏฺฐพฺเพ อนิจฺจโต ฯเปฯ วิปสฺสติ ฯเปฯ โทมนสฺสํ อุปฺปชฺชติ. ทิพฺเพน จกฺขุนา รูปํ ปสฺสติ. ทิพฺพาย โสตธาตุยา สทฺทํ สุณาติ. รูปายตนํ จกฺขุ\-วิญฺญาณสหคตานํ ขนฺธานํ ปุเรชาต\-ปจฺจเยน ปจฺจโย ฯเปฯ. โผฏฺฐพฺพา\-ยตนํ กาย\-วิญฺญาณสหคตานํ ขนฺธานํ ปุเรชาต\-ปจฺจเยน ปจฺจโย. (1)}

\prose{จิตฺต\-สมุฏฺฐาโน ธมฺโม โนจิตฺต\-สมุฏฺฐานสฺส ธมฺมสฺส ปุเรชาต\-ปจฺจเยน ปจฺจโย. อารมฺมณ\-ปุเรชาตํ\csromandash{}จิตฺต\-สมุฏฺฐาเน รูเป ฯเปฯ โผฏฺฐพฺเพ อนิจฺจโต ฯเปฯ วิปสฺสติ, อสฺสาเทติ อภินนฺทติ, ตํ อารพฺภ จิตฺตํ อุปฺปชฺชติ. ทิพฺเพน จกฺขุนา รูปํ ปสฺสติ. ทิพฺพาย โสตธาตุยา สทฺทํ สุณาติ. รูปายตนํ จกฺขุ\-วิญฺญาณสฺส ฯเปฯ. โผฏฺฐพฺพา\-ยตนํ กายวิญฺญาณสฺส ปุเรชาต\-ปจฺจเยน ปจฺจโย. (2)}

\prose{จิตฺต\-สมุฏฺฐาโน ธมฺโม จิตฺต\-สมุฏฺฐานสฺส จ โนจิตฺต\-สมุฏฺฐานสฺส จ ธมฺมสฺส ปุเรชาต\-ปจฺจเยน ปจฺจโย. อารมฺมณ\-ปุเรชาตํ\csromandash{} จิตฺต\-สมุฏฺฐาเน รูเป ฯเปฯ โผฏฺฐพฺเพ อนิจฺจโต ฯเปฯ วิปสฺสติ, อสฺสาเทติ อภินนฺทติ, ตํ อารพฺภ จิตฺตญฺจ สมฺปยุตฺตกา จ ขนฺธา อุปฺปชฺชนฺติ. ทิพฺเพน จกฺขุนา รูปํ ปสฺสติ. \csromanfolio{453}ทิพฺพาย โสตธาตุยา สทฺทํ สุณาติ. รูปายตนํ จกฺขุ\-วิญฺญาณสฺส สมฺปยุตฺตกานญฺจ ขนฺธานํ ปุเรชาต\-ปจฺจเยน ปจฺจโย ฯเปฯ. โผฏฺฐพฺพา\-ยตนํ. (3)}

\proseitem{218}{โนจิตฺต\-สมุฏฺฐาโน ธมฺโม โนจิตฺต\-สมุฏฺฐานสฺส ธมฺมสฺส ปุเรชาต\-ปจฺจเยน ปจฺจโย. อารมฺมณ\-ปุเรชาตํ, วตฺถุ\-ปุเรชาตํ.  อารมฺมณ\-ปุเรชาตํ\csromandash{}จกฺขุํ ฯเปฯ กายํ. รูเป ฯเปฯ เผฏฺฐพฺเพ วตฺถุํ อนิจฺจโต ฯเปฯ ตํ อารพฺภ จิตฺตํ อุปฺปชฺชติ. ทิพฺเพน จกฺขุนา รูปํ ปสฺสติ. ทิพฺพาย โสตธาตุยา สทฺทํ สุณาติ. รูปายตนํ จกฺขุ\-วิญฺญาณสฺส ฯเปฯ. โผฏฺฐพฺพา\-ยตนํ ฯเปฯ. วตฺถุ\-ปุเรชาตํ\csromandash{} จกฺขายตนํ จกฺขุ\-วิญฺญาณสฺส ฯเปฯ กายายตนํ ฯเปฯ. วตฺถุ จิตฺตสฺส ปุเรชาต\-ปจฺจเยน ปจฺจโย. (1)}

\prose{โนจิตฺต\-สมุฏฺฐาโน ธมฺโม จิตฺต\-สมุฏฺฐานสฺส ธมฺมสฺส ปุเรชาต\-ปจฺจเยน ปจฺจโย. อารมฺมณ\-ปุเรชาตํ, วตฺถุ\-ปุเรชาตํ.  อารมฺมณ\-ปุเรชาตํ\csromandash{}จกฺขุํ ฯเปฯ วตฺถุํ อนิจฺจโต ฯเปฯ โทมนสฺสํ อุปชฺชติ. ทิพฺเพน จกฺขุนา รูปํ ปสฺสติ. ทิพฺพาย โสตธาตุยา สทฺทํ สุณาติ. รูปายตนํ จกฺขุ\-วิญฺญาณสหคตานํ ขนฺธานํ ฯเปฯ โผฏฺฐพฺพา\-ยตนํ ฯเปฯ. วตฺถุ\-ปุเรชาตํ\csromandash{}จกฺขายตนํ จกฺขุ\-วิญฺญาณสหคตานํ ขนฺธานํ ฯเปฯ กายายตนํ ฯเปฯ. วตฺถุ จิตฺต\-สมุฏฺฐานานํ ขนฺธานํ ปุเรชาต\-ปจฺจเยน ปจฺจโย. (2)}

\prose{โนจิตฺต\-สมุฏฺฐาโน ธมฺโม จิตฺต\-สมุฏฺฐานสฺส จ โนจิตฺต\-สมุฏฺฐานสฺส จ ธมฺมสฺส ปุเรชาต\-ปจฺจเยน ปจฺจโย. อารมฺมณ\-ปุเรชาตํ, วตฺถุ\-ปุเรชาตํ.  อารมฺมณ\-ปุเรชาตํ\csromandash{}จกฺขุํ ฯเปฯ วตฺถุํ อนิจฺจโต ฯเปฯ ตํ อารพฺภ จิตฺตญฺจ สมฺปยุตฺตกา จ ขนฺธา อุปฺปชฺชนฺติ. ทิพฺเพน จกฺขุนา รูปํ ปสฺสติ ฯเปฯ. ทิพฺพาย โสตธาตุยา สทฺทํ สุณาติ. รูปายตนํ จกฺขุ\-วิญฺญาณสฺส สมฺปยุตฺตกานญฺจ ขนฺธานํ ฯเปฯ. โผฏฺฐพฺพา\-ยตนํ ฯเปฯ. วตฺถุ\-ปุเรชาตํ\csromandash{}จกฺขายตนํ จกฺขุ\-วิญฺญาณสฺส สมฺปยุตฺตกานญฺจ ขนฺธานํ ฯเปฯ กายายตนํ ฯเปฯ. วตฺถุ จิตฺตสฺส สมฺปยุตฺตกานญฺจ ขนฺธานํ ปุเรชาต\-ปจฺจเยน ปจฺจโย. (3)}

\proseitem{219}{จิตฺต\-สมุฏฺฐาโน จ โนจิตฺต\-สมุฏฺฐาโน จ ธมฺมา จิตฺต\-สมุฏฺฐานสฺส ธมฺมสฺส ปุเรชาต\-ปจฺจเยน ปจฺจโย. อารมฺมณ\-ปุเรชาตํ, วตฺถุ\-ปุเรชาตํ. จิตฺต\-สมุฏฺฐานํ รูปายตนญฺจ จกฺขายตนญฺจ จกฺขุ\-วิญฺญาณสหคตานํ \csromanfolio{454}ขนฺธานํ ฯเปฯ. จิตฺต\-สมุฏฺฐานํ โผฏฺฐพฺพายตนญฺจ กายายตนญฺจ กาย\-วิญฺญาณสหคตานํ ขนฺธานํ ปุเรชาต\-ปจฺจเยน ปจฺจโย. จิตฺต\-สมุฏฺฐานํ รูปายตนญฺจ วตฺถุ จ จิตฺต\-สมุฏฺฐานานํ ขนฺธานํ ปุเรชาต\-ปจฺจเยน ปจฺจโย ฯเปฯ. จิตฺต\-สมุฏฺฐานํ โผฏฺฐพฺพายตนญฺจ วตฺถุ จ ฯเปฯ. (1)}

\prose{จิตฺต\-สมุฏฺฐาโน จ โนจิตฺต\-สมุฏฺฐาโน จ ธมฺมา โนจิตฺต\-สมุฏฺฐานสฺส ธมฺมสฺส ปุเรชาต\-ปจฺจเยน ปจฺจโย. อารมฺมณ\-ปุเรชาตํ, วตฺถุ\-ปุเรชาตํ. จิตฺต\-สมุฏฺฐานํ รูปายตนญฺจ จกฺขายตนญฺจ จกฺขุ\-วิญฺญาณสฺส ฯเปฯ. จิตฺต\-สมุฏฺฐานํ โผฏฺฐพฺพายตนญฺจ กายายตนญฺจ กายวิญฺญาณสฺส ปุเรชาต\-ปจฺจเยน ปจฺจโย. จิตฺต\-สมุฏฺฐานํ รูปายตนญฺจ วตฺถุ จ จิตฺตสฺส ปุเรชาต\-ปจฺจเยน ปจฺจโย ฯเปฯ. จิตฺต\-สมุฏฺฐานํ โผฏฺฐพฺพายตนญฺจ วตฺถุ จ ฯเปฯ. (2)}

\prose{จิตฺต\-สมุฏฺฐาโน จ โนจิตฺต\-สมุฏฺฐาโน จ ธมฺมา จิตฺต\-สมุฏฺฐานสฺส จ โนจิตฺต\-สมุฏฺฐานสฺส จ ธมฺมสฺส ปุเรชาต\-ปจฺจเยน ปจฺจโย. อารมฺมณ\-ปุเรชาตํ, วตฺถุ\-ปุเรชาตํ. จิตฺต\-สมุฏฺฐานํ รูปายตนญฺจ จกฺขายตนญฺจ จกฺขุ\-วิญฺญาณสฺส สมฺปยุตฺตกานญฺจ ขนฺธานํ ปุเรชาต\-ปจฺจเยน ปจฺจโย ฯเปฯ. จิตฺต\-สมุฏฺฐานํ โผฏฺฐพฺพายตนญฺจ ฯเปฯ. จิตฺต\-สมุฏฺฐานํ รูปายตนญฺจ วตฺถุ จ จิตฺตสฺส สมฺปยุตฺตกานญฺจ ขนฺธานํ ปุเรชาต\-ปจฺจเยน ปจฺจโย ฯเปฯ. จิตฺต\-สมุฏฺฐานํ โผฏฺฐพฺพายตนญฺจ วตฺถุ จ ฯเปฯ. (3)}

\proseitem{220}{จิตฺต\-สมุฏฺฐาโน ธมฺโม จิตฺต\-สมุฏฺฐานสฺส ธมฺมสฺส ปจฺฉาชาต\-ปจฺจเยน ปจฺจโย. ปจฺฉาชาตา จิตฺต\-สมุฏฺฐานา ขนฺธา ปุเรชาตสฺส อิมสฺส จิตฺต\-สมุฏฺฐานสฺส กายสฺส ปจฺฉาชาปจฺจเยน ปจฺจโย. (อิมินากาเรเนว ปจฺฉาชาโต วิตฺถาเรตพฺโพ.) อาเสวน\-ปจฺจเยน ปจฺจโย. นว.}

\csromanheader{2}{\textbf{กมฺมวิปาก}}
\proseitem{221}{จิตฺต\-สมุฏฺฐาโน ธมฺโม จิตฺต\-สมุฏฺฐานสฺส ธมฺมสฺส กมฺมปจฺจเยน ปจฺจโย. สหชาตา, นานากฺขณิกา.  สหชาตา จิตฺต\-สมุฏฺฐานา เจตนา สมฺปยุตฺตกานํ ขนฺธานํ จิตฺต\-สมุฏฺฐานานญฺจ รูปานํ กมฺมปจฺจเยน ปจฺจโย. ปฏิสนฺธิ\-กฺขเณ ฯเปฯ. นานากฺขณิกา จิตฺต\-สมุฏฺฐานา เจตนา วิปากานํ ขนฺธานํ กมฺมปจฺจเยน ปจฺจโย. (1)}

\prose{\csromanfolio{455}จิตฺต\-สมุฏฺฐาโน ธมฺโม โนจิตฺต\-สมุฏฺฐานสฺส ธมฺมสฺส กมฺมปจฺจเยน ปจฺจโย. สหชาตา, นานากฺขณิกา.  สหชาตา จิตฺต\-สมุฏฺฐานา เจตนา จิตฺตสฺส กมฺมปจฺจเยน ปจฺจโย. ปฏิสนฺธิ\-กฺขเณ ฯเปฯ. นานากฺขณิกา จิตฺต\-สมุฏฺฐานา เจตนา วิปากสฺส จิตฺตสฺส กฏตฺตา จ รูปานํ กมฺมปจฺจเยน ปจฺจโย. (2)}

\prose{จิตฺต\-สมุฏฺฐาโน ธมฺโม จิตฺต\-สมุฏฺฐานสฺส จ โนจิตฺต\-สมุฏฺฐานสฺส จ ธมฺมสฺส กมฺมปจฺจเยน ปจฺจโย. สหชาตา, นานากฺขณิกา.  สหชาตา จิตฺต\-สมุฏฺฐานา เจตนา สมฺปยุตฺตกานํ ขนฺธานํ จิตฺตสฺส จ จิตฺต\-สมุฏฺฐานานญฺจ รูปานํ กมฺมปจฺจเยน ปจฺจโย. ปฏิสนฺธิ\-กฺขเณ ฯเปฯ. นานากฺขณิกา จิตฺต\-สมุฏฺฐานา เจตนา วิปากานํ ขนฺธานํ จิตฺตสฺส จ กฏตฺตา จ รูปานํ กมฺมปจฺจเยน ปจฺจโย. (3)}

\prose{จิตฺต\-สมุฏฺฐาโน ธมฺโม จิตฺต\-สมุฏฺฐานสฺส ธมฺมสฺส วิปาก\-ปจฺจเยน ปจฺจโย. นว.}

\csromanheader{2}{\textbf{อาหาร}}
\proseitem{222}{จิตฺต\-สมุฏฺฐาโน ธมฺโม จิตฺต\-สมุฏฺฐานสฺส ธมฺมสฺส อาหารปจฺจเยน ปจฺจโย. จิตฺต\-สมุฏฺฐานา อาหารา สมฺปยุตฺตกานํ ขนฺธานํ จิตฺต\-สมุฏฺฐานานญฺจ รูปานํ อาหารปจฺจเยน ปจฺจโย. ปฏิสนฺธิ\-กฺขเณ ฯเปฯ. (มูลํ กาตพฺพํ.) จิตฺต\-สมุฏฺฐานา อาหารา จิตฺตสฺส อาหารปจฺจเยน ปจฺจโย. ปฏิสนฺธิ\-กฺขเณ ฯเปฯ. จิตฺต\-สมุฏฺฐาโน กพฬีกาโร อาหาโร อิมสฺส โนจิตฺต\-สมุฏฺฐานสฺส กายสฺส อาหารปจฺจเยน ปจฺจโย. (มูลํ กาตพฺพํ.) จิตฺต\-สมุฏฺฐานา อาหารา สมฺปยุตฺตกานํ ขนฺธานํ จิตฺตสฺส จ จิตฺต\-สมุฏฺฐานานญฺจ รูปานํ อาหารปจฺจเยน ปจฺจโย. ปฏิสนฺธิ\-กฺขเณ ฯเปฯ. (3)}

\proseitem{223}{โนจิตฺต\-สมุฏฺฐาโน ธมฺโม โนจิตฺต\-สมุฏฺฐานสฺส ธมฺมสฺส อาหารปจฺจเยน ปจฺจโย. ปฏิสนฺธิ\-กฺขเณ โนจิตฺต\-สมุฏฺฐานา อาหารา กฏตฺตารูปานํ อาหารปจฺจเยน ปจฺจโย. โนจิตฺต\-สมุฏฺฐาโน กพฬีกาโร อาหาโร อิมสฺส โนจิตฺต\-สมุฏฺฐานสฺส กายสฺส อาหารปจฺจเยน ปจฺจโย. (มูลํ กาตพฺพํ.) โนจิตฺต\-สมุฏฺฐานา อาหารา สมฺปยุตฺตกานํ ขนฺธานํ จิตฺตสมุฏฺฐานานญฺจ รูปานํ อาหารปจฺจเยน ปจฺจโย. ปฏิสนฺธิ\-กฺขเณ ฯเปฯ. (มูลํ กาตพฺพํ.) ปฏิสนฺธิ\-กฺขเณ โนจิตฺต\-สมุฏฺฐานา \csromanfolio{456}อาหารา สมฺปยุตฺตกานํ ขนฺธานํ กฏตฺตา จ รูปานํ อาหารปจฺจเยน ปจฺจโย. (3)}

\proseitem{224}{จิตฺต\-สมุฏฺฐาโน จ โนจิตฺต\-สมุฏฺฐาโน จ ธมฺมา จิตฺต\-สมุฏฺฐานสฺส ธมฺมสฺส อาหารปจฺจเยน ปจฺจโย. จิตฺต\-สมุฏฺฐานา จ โนจิตฺต\-สมุฏฺฐานา จ อาหารา สมฺปยุตฺตกานํ ขนฺธานํ จิตฺต\-สมุฏฺฐานานญฺจ รูปานํ อาหารปจฺจเยน ปจฺจโย. ปฏิสนฺธิ\-กฺขเณ ฯเปฯ. (มูลํ กาตพฺพํ.) ปฏิสนฺธิ\-กฺขเณ จิตฺต\-สมุฏฺฐานา จ โนจิตฺต\-สมุฏฺฐานา จ อาหารา กฏตฺตารูปานํ อาหารปจฺจเยน ปจฺจโย. จิตฺต\-สมุฏฺฐาโน จ โนจิตฺต\-สมุฏฺฐาโน จ กพฬีกาโร อาหาโร อิมสฺส โนจิตฺต\-สมุฏฺฐานสฺส กายสฺส อาหารปจฺจเยน ปจฺจโย. (มูลํ กาตพฺพํ.) ปฏิสนฺธิ\-กฺขเณ จิตฺต\-สมุฏฺฐานา จ โนจิตฺต\-สมุฏฺฐานา จ อาหารา สมฺปยุตฺตกานํ ขนฺธานํ กฏตฺตา จ รูปานํ อาหารปจฺจเยน ปจฺจโย. (3)}

\csromanheader{2}{\textbf{อินฺทฺริยาทิ}}
\proseitem{225}{จิตฺต\-สมุฏฺฐาโน ธมฺโม จิตฺต\-สมุฏฺฐานสฺส ธมฺมสฺส อินฺทฺริย\-ปจฺจเยน ปจฺจโย. ตีณิ.}

\prose{โนจิตฺต\-สมุฏฺฐาโน ธมฺโม โนจิตฺต\-สมุฏฺฐานสฺส ธมฺมสฺส อินฺทฺริย\-ปจฺจเยน ปจฺจโย. ปฏิสนฺธิ\-กฺขเณ โนจิตฺต\-สมุฏฺฐานา อินฺทฺริยา กฏตฺตารูปานํ อินฺทฺริย\-ปจฺจเยน ปจฺจโย. ปฏิสนฺธิ\-กฺขเณ ฯเปฯ. จกฺขุนฺทฺริยํ จกฺขุ\-วิญฺญาณสฺส ฯเปฯ. กายินฺทฺริยํ กายวิญฺญาณสฺส ฯเปฯ. รูปชีวิตินฺทฺริยํ กฏตฺตารูปานํ อินฺทฺริย\-ปจฺจเยน ปจฺจโย. (มูลํ กาตพฺพํ.) โนจิตฺต\-สมุฏฺฐานา อินฺทฺริยา สมฺปยุตฺตกานํ ขนฺธานํ จิตฺตสมุฏฺฐานานญฺจ รูปานํ อินฺทฺริย\-ปจฺจเยน ปจฺจโย. ปฏิสนฺธิ\-กฺขเณ ฯเปฯ. จกฺขุนฺทฺริยํ จกฺขุ\-วิญฺญาณสหคตานํ ขนฺธานํ ฯเปฯ. กายินฺทฺริยํ ฯเปฯ. (มูลํ กาตพฺพํ.) ปฏิสนฺธิ\-กฺขเณ โนจิตฺต\-สมุฏฺฐานา อินฺทฺริยา สมฺปยุตฺตกานํ ขนฺธานํ กฏตฺตา จ รูปานํ อินฺทฺริย\-ปจฺจเยน ปจฺจโย. จกฺขุนฺทฺริยํ จกฺขุ\-วิญฺญาณสฺส จกฺขุ\-วิญฺญาณสหคตานญฺจ ขนฺธานํ ฯเปฯ. กายินฺทฺริยํ ฯเปฯ. (3)}

\proseitem{226}{จิตฺต\-สมุฏฺฐาโน จ โนจิตฺต\-สมุฏฺฐาโน จ ธมฺมา จิตฺตสมุฏฺฐานุสฺส ธมฺมสฺส อินฺทฺริย\-ปจฺจเยน ปจฺจโย. จิตฺต\-สมุฏฺฐานา จ โนจิตฺต\-สมุฏฺฐานา จ อินฺทฺริยา สมฺปยุตฺตกานํ ขนฺธานํ จิตฺต\-สมุฏฺฐานานญฺจ รูปานํ อินฺทฺริย\-ปจฺจเยน \csromanfolio{457}ปจฺจโย. ปฏิสนฺธิ\-กฺขเณ ฯเปฯ. จกฺขุนฺทฺริยญฺจ อุเปกฺขินฺทฺริยญฺจ จกฺขุ\-วิญฺญาณสหคตานํ ขนฺธานํ ฯเปฯ. กายินฺทฺริยญฺจ สุขินฺทฺริยญฺจ ฯเปฯ. กายินฺทฺริยญฺจ ทุกฺขินฺทฺริยญฺจ กาย\-วิญฺญาณสหคตานํ ขนฺธานํ อินฺทฺริย\-ปจฺจเยน ปจฺจโย. (มูลํ กาตพฺพํ.) ปฏิสนฺธิ\-กฺขเณ จิตฺต\-สมุฏฺฐานา จ โนจิตฺต\-สมุฏฺฐานา จ อินฺทฺริยา กฏตฺตารูปานํ อินฺทฺริย\-ปจฺจเยน ปจฺจโย. จกฺขุนฺทฺริยญฺจ อุเปกฺขินฺทฺริยญฺจ จกฺขุ\-วิญฺญาณสฺส ฯเปฯ. กายินฺทฺริยญฺจ ฯเปฯ. (มูลํ กาตพฺพํ.) ปฏิสนฺธิ\-กฺขเณ จิตฺต\-สมุฏฺฐานา จ โนจิตฺต\-สมุฏฺฐานา จ อินฺทฺริยา สมฺปยุตฺตกานํ ขนฺธานํ กฏตฺตา จ รูปานํ อินฺทฺริย\-ปจฺจเยน ปจฺจโย. จกฺขุนฺทฺริยญฺจ อุเปกฺขินฺทฺริยญฺจ จกฺขุ\-วิญฺญาณสฺส สมฺปยุตฺตกานญฺจ ขนฺธานํ อินฺทฺริย\-ปจฺจเยน ปจฺจโย. กายินฺทฺริยญฺจ ฯเปฯ. (3)}

\prose{ฌานปจฺจเยน ปจฺจโย. ตีณิ, มคฺคปจฺจเยน ปจฺจโย. ตีณิ, สมฺปยุตฺต\-ปจฺจเยน ปจฺจโย. ปญฺจ ฯเปฯ.}

\proseitem{227}{จิตฺต\-สมุฏฺฐาโน ธมฺโม จิตฺต\-สมุฏฺฐานสฺส ธมฺมสฺส วิปฺปยุตฺต\-ปจฺจเยน ปจฺจโย. สหชาตํ, ปจฺฉาชาตํ. (สํขิตฺตํ.) (1)}

\prose{จิตฺต\-สมุฏฺฐาโน ธมฺโม โนจิตฺต\-สมุฏฺฐานสฺส ธมฺมสฺส วิปฺปยุตฺต\-ปจฺจเยน ปจฺจโย. สหชาตํ, ปจฺฉาชาตํ.  สหชาตํ ปฏิสนฺธิ\-กฺขเณ. (สํขิตฺตํ.) (2)}

\prose{จิตฺต\-สมุฏฺฐาโน ธมฺโม จิตฺต\-สมุฏฺฐานสฺส จ โนจิตฺต\-สมุฏฺฐานสฺส จ ธมฺมสฺส วิปฺปยุตฺต\-ปจฺจเยน ปจฺจโย. ปจฺฉาชาตํ. (สํขิตฺตํ.) (3)}

\proseitem{228}{โนจิตฺต\-สมุฏฺฐาโน ธมฺโม โนจิตฺต\-สมุฏฺฐานสฺส ธมฺมสฺส วิปฺปยุตฺต\-ปจฺจเยน ปจฺจโย. สหชาตํ, ปุเรชาตํ, ปจฺฉาชาตํ.  สหชาตํ ปฏิสนฺธิ\-กฺขเณ จิตฺตํ กฏตฺตารูปานํ วิปฺปยุตฺต\-ปจฺจเยน ปจฺจโย. จิตฺตํ วตฺถุสฺส วิปฺปยุตฺต\-ปจฺจเยน ปจฺจโย. วตฺถุ จิตฺตสฺส วิปฺปยุตฺต\-ปจฺจเยน ปจฺจโย. ปุเรชาตํ จกฺขายตนํ จกฺขุ\-วิญฺญาณสฺส ฯเปฯ. กายายตนํ กายวิญฺญาณสฺส ฯเปฯ. วตฺถุ จิตฺตสฺส วิปฺปยุตฺต\-ปจฺจเยน ปจฺจโย. ปจฺฉาชาตา โนจิตฺต\-สมุฏฺฐานา ขนฺธา ปุเรชาตสฺส อิมสฺส โนจิตฺต\-สมุฏฺฐานสฺส กายสฺส วิปฺปยุตฺต\-ปจฺจเยน ปจฺจโย. (1)}

\prose{โนจิตฺต\-สมุฏฺฐาโน ธมฺโม จิตฺต\-สมุฏฺฐานสฺส ธมฺมสฺส วิปฺปยุตฺต\-ปจฺจเยน ปจฺจโย. สหชาตํ, ปุเรชาตํ, ปจฺฉาชาตํ. (สํขิตฺตํ.) (2)}

\prose{\csromanfolio{458}โนจิตฺต\-สมุฏฺฐาโน ธมฺโม จิตฺต\-สมุฏฺฐานสฺส จ โนจิตฺต\-สมุฏฺฐานสฺส จ ธมฺมสฺส วิปฺปยุตฺต\-ปจฺจเยน ปจฺจโย. สหชาตํ, ปุเรชาตํ, ปจฺฉาชาตํ. (สํขิตฺตํ.) (3)}

\proseitem{229}{จิตฺต\-สมุฏฺฐาโน จ โนจิตฺต\-สมุฏฺฐาโน จ ธมฺมา จิตฺต\-สมุฏฺฐานสฺส ธมฺมสฺส วิปฺปยุตฺต\-ปจฺจเยน ปจฺจโย. สหชาตํ, ปจฺฉาชาตํ. (สํขิตฺตํ.) (1)}

\prose{จิตฺต\-สมุฏฺฐาโน จ โนจิตฺต\-สมุฏฺฐาโน จ ธมฺมา โนจิตฺต\-สมุฏฺฐานสฺส ธมฺมสฺส วิปฺปยุตฺต\-ปจฺจเยน ปจฺจโย. สหชาตํ, ปจฺฉาชาตํ. (สํขิตฺตํ.) (2)}

\prose{จิตฺต\-สมุฏฺฐาโน จ โนจิตฺต\-สมุฏฺฐาโน จ ธมฺมา จิตฺต\-สมุฏฺฐานสฺส จ โนจิตฺต\-สมุฏฺฐานสฺส จ ธมฺมสฺส วิปฺปยุตฺต\-ปจฺจเยน ปจฺจโย. ปจฺฉาชาตํ. (สํขิตฺตํ.) (3)}

\proseitem{230}{จิตฺต\-สมุฏฺฐาโน ธมฺโม จิตฺต\-สมุฏฺฐานสฺส ธมฺมสฺส อตฺถิปจฺจเยน ปจฺจโย. สหชาตํ, ปุเรชาตํ, ปจฺฉาชาตํ. (สํขิตฺตํ.) (1)}

\prose{จิตฺต\-สมุฏฺฐาโน ธมฺโม โนจิตฺต\-สมุฏฺฐานสฺส ธมฺมสฺส อตฺถิปจฺจเยน ปจฺจโย. สหชาตํ, ปุเรชาตํ, ปจฺฉาชาตํ, อาหารํ. (สํขิตฺตํ.) (2)}

\prose{จิตฺต\-สมุฏฺฐาโน ธมฺโม จิตฺต\-สมุฏฺฐานสฺส จ โนจิตฺต\-สมุฏฺฐานสฺส จ ธมฺมสฺส อตฺถิปจฺจเยน ปจฺจโย. สหชาตํ, ปุเรชาตํ, ปจฺฉาชาตํ. (สํขิตฺตํ.) (3)}

\prose{โนจิตฺต\-สมุฏฺฐาโน ธมฺโม โนจิตฺต\-สมุฏฺฐานสฺส ธมฺมสฺส อตฺถิปจฺจเยน ปจฺจโย. สหชาตํ, ปุเรชาตํ, ปจฺฉาชาตํ, อาหารํ, อินฺทฺริยํ. (สํขิตฺตํ.) (1)}

\prose{โนจิตฺต\-สมุฏฺฐาโน ธมฺโม จิตฺต\-สมุฏฺฐานสฺส ธมฺมสฺส อตฺถิปจฺจเยน ปจฺจโย. สหชาตํ, ปุเรชาตํ, ปจฺฉาชาตํ. (สํขิตฺตํ.)}

\prose{โนจิตฺต\-สมุฏฺฐาโน ธมฺโม จิตฺต\-สมุฏฺฐานสฺส จ โนจิตฺต\-สมุฏฺฐานสฺส จ ธมฺมสฺส อตฺถิปจฺจเยน ปจฺจโย. สหชาตํ, ปุเรชาตํ, ปจฺฉาชาตํ, อาหารํ. (สํขิตฺตํ.) (3)}

\proseitem{231}{\csromanfolio{459}จิตฺต\-สมุฏฺฐาโน จ โนจิตฺต\-สมุฏฺฐาโน จ ธมฺมา จิตฺต\-สมุฏฺฐานสฺส ธมฺมสฺส อตฺถิปจฺจเยน ปจฺจโย. สหชาตํ, ปุเรชาตํ, ปจฺฉาชาตํ.  สหชาโต จกฺขุ\-วิญฺญาณ\-สหคโต เอโก ขนฺโธ ฯเปฯ. (สํขิตฺตํ.) (1)}

\prose{จิตฺต\-สมุฏฺฐาโน จ โนจิตฺต\-สมุฏฺฐาโน จ ธมฺมา โนจิตฺต\-สมุฏฺฐานสฺส ธมฺมสฺส อตฺถิปจฺจเยน ปจฺจโย. สหชาตํ, ปุเรชาตํ, ปจฺฉาชาตํ, อาหารํ, อินฺทฺริยํ.  สหชาตา จกฺขุ\-วิญฺญาณสหคตา ขนฺธา จ จกฺขายตนญฺจ จกฺขุ\-วิญฺญาณสฺส อตฺถิปจฺจเยน ปจฺจโย ฯเปฯ. กาย\-วิญฺญาณสหคตา ฯเปฯ. สหชาตา จิตฺต\-สมุฏฺฐานา ฯเปฯ. (ปจฺจยวาร\-สทิสํ ปฏิสนฺธิปิ ปวตฺติปิ กาตพฺพา สพฺเพสมฺปิ ปญฺหานํ.) ปจฺฉาชาตา จิตฺต\-สมุฏฺฐานา ขนฺธา จ จิตฺตญฺจ ปุเรชาตสฺส อิมสฺส โนจิตฺต\-สมุฏฺฐานสฺส กายสฺส อตฺถิปจฺจเยน ปจฺจโย. ปจฺฉาชาตา จิตฺต\-สมุฏฺฐานา ขนฺธา จ จิตฺตญฺจ กพฬีกาโร อาหาโร จ อิมสฺส โนจิตฺต\-สมุฏฺฐานสฺส กายสฺส อตฺถิปจฺจเยน ปจฺจโย. ปจฺฉาชาตา จิตฺต\-สมุฏฺฐานา ขนฺธา จ จิตฺตญฺจ รูปชีวิตินฺทฺริยญฺจ กฏตฺตารูปานํ อตฺถิปจฺจเยน ปจฺจโย. (2)}

\prose{จิตฺต\-สมุฏฺฐาโน จ โนจิตฺต\-สมุฏฺฐาโน จ ธมฺมา จิตฺต\-สมุฏฺฐานสฺส จ โนจิตฺต\-สมุฏฺฐานสฺส จ ธมฺมสฺส อตฺถิปจฺจเยน ปจฺจโย. สหชาตํ, ปุเรชาตํ, ปจฺฉาชาตํ.  สหชาโต จกฺขุ\-วิญฺญาณ\-สหคโต ฯเปฯ. (สํขิตฺตํ.) (3) นตฺถิ\-ปจฺจเยน ปจฺจโย. วิคต\-ปจฺจเยน ปจฺจโย. อวิคต\-ปจฺจเยน ปจฺจโย.}

\csromanmarkh{1. ปจฺจยานุโลม}\csromanmarkhii{2. สงฺขฺยาวาร}\csromanheadernomark{2}{1. \textbf{ปจฺจยานุโลม 2. สงฺขฺยาวาร}}
\proseitem{232}{เหตุยา ตีณิ, อารมฺมเณ นว, อธิปติยา นว, อนนฺตเร นว, สมนนฺตเร นว, สหชาเต นว, อญฺญมญฺเญ นว, นิสฺสเย นว, อุปนิสฺสเย นว, ปุเรชาเต นว, ปจฺฉาชาเต นว, อาเสวเน นว, กมฺเม ตีณิ, วิปาเก นว, (สพฺพตฺถ นว,) อินฺทฺริเย นว, ฌาเน ตีณิ, มคฺเค ตีณิ, สมฺปยุตฺเต ปญฺจ, วิปฺปยุตฺเต นว ฯเปฯ อวิคเต นว.}

\vspace{\tipitakaparskip}
\csromancenter{อนุโลมํ.}
\proseitem{233}{\csromanfolio{460}จิตฺต\-สมุฏฺฐาโน ธมฺโม จิตฺต\-สมุฏฺฐานสฺส ธมฺมสฺส อารมฺมณ\-ปจฺจเยน ปจฺจโย. สหชาต\-ปจฺจเยน ปจฺจโย. อุปนิสฺสย\-ปจฺจเยน ปจฺจโย. ปุเรชาต\-ปจฺจเยน ปจฺจโย. ปจฺฉาชาต\-ปจฺจเยน ปจฺจโย. กมฺมปจฺจเยน ปจฺจโย. (1)}

\prose{จิตฺต\-สมุฏฺฐาโน ธมฺโม โนจิตฺต\-สมุฏฺฐานสฺส ธมฺมสฺส อารมฺมณ\-ปจฺจเยน ปจฺจโย. สหชาต\-ปจฺจเยน ปจฺจโย. อุปนิสฺสย\-ปจฺจเยน ปจฺจโย. ปุเรชาต\-ปจฺจเยน ปจฺจโย. ปจฺฉาชาต\-ปจฺจเยน ปจฺจโย. กมฺมปจฺจเยน ปจฺจโย. อาหารปจฺจเยน ปจฺจโย. (2)}

\prose{จิตฺต\-สมุฏฺฐาโน ธมฺโม จิตฺต\-สมุฏฺฐานสฺส จ โนจิตฺต\-สมุฏฺฐานสฺส จ ธมฺมสฺส อารมฺมณ\-ปจฺจเยน ปจฺจโย. สหชาต\-ปจฺจเยน ปจฺจโย. อุปนิสฺสย\-ปจฺจเยน ปจฺจโย. ปุเรชาต\-ปจฺจเยน ปจฺจโย. ปจฺฉาชาต\-ปจฺจเยน ปจฺจโย. กมฺมปจฺจเยน ปจฺจโย. (3)}

\proseitem{234}{โนจิตฺต\-สมุฏฺฐาโน ธมฺโม โนจิตฺต\-สมุฏฺฐานสฺส ธมฺมสฺส อารมฺมณ\-ปจฺจเยน ปจฺจโย. สหชาต\-ปจฺจเยน ปจฺจโย. อุปนิสฺสย\-ปจฺจเยน ปจฺจโย. ปุเรชาต\-ปจฺจเยน ปจฺจโย. ปจฺฉาชาต\-ปจฺจเยน ปจฺจโย. อาหารปจฺจเยน ปจฺจโย. อินฺทฺริย\-ปจฺจเยน ปจฺจโย. (1)}

\prose{โนจิตฺต\-สมุฏฺฐาโน ธมฺโม จิตฺต\-สมุฏฺฐานสฺส ธมฺมสฺส อารมฺมณ\-ปจฺจเยน ปจฺจโย. สหชาต\-ปจฺจเยน ปจฺจโย. อุปนิสฺสย\-ปจฺจเยน ปจฺจโย. ปุเรชาต\-ปจฺจเยน ปจฺจโย. ปจฺฉาชาต\-ปจฺจเยน ปจฺจโย. (2)}

\prose{โนจิตฺต\-สมุฏฺฐาโน ธมฺโม จิตฺต\-สมุฏฺฐานสฺส จ โนจิตฺต\-สมุฏฺฐานสฺส จ ธมฺมสฺส อารมฺมณ\-ปจฺจเยน ปจฺจโย. สหชาต\-ปจฺจเยน ปจฺจโย. อุปนิสฺสย\-ปจฺจเยน ปจฺจโย. ปุเรชาต\-ปจฺจเยน ปจฺจโย. ปจฺฉาชาต\-ปจฺจเยน ปจฺจโย. (3)}

\proseitem{235}{จิตฺต\-สมุฏฺฐาโน จ โนจิตฺต\-สมุฏฺฐาโน จ ธมฺมา จิตฺต\-สมุฏฺฐานสฺส ธมฺมสฺส อารมฺมณ\-ปจฺจเยน ปจฺจโย. สหชาต\-ปจฺจเยน ปจฺจโย. อุปนิสฺสย\-ปจฺจเยน ปจฺจโย. ปุเรชาต\-ปจฺจเยน ปจฺจโย. ปจฺฉาชาต\-ปจฺจเยน ปจฺจโย. (1)}

\prose{\csromanfolio{461}จิตฺต\-สมุฏฺฐาโน จ โนจิตฺต\-สมุฏฺฐาโน จ ธมฺมา โนจิตฺต\-สมุฏฺฐานสฺส ธมฺมสฺส อารมฺมณ\-ปจฺจเยน ปจฺจโย. สหชาต\-ปจฺจเยน ปจฺจโย. อุปนิสฺสย\-ปจฺจเยน ปจฺจโย. ปุเรชาต\-ปจฺจเยน ปจฺจโย. ปจฺฉาชาต\-ปจฺจเยน ปจฺจโย. อาหารปจฺจเยน ปจฺจโย. (2)}

\prose{จิตฺต\-สมุฏฺฐาโน จ โนจิตฺต\-สมุฏฺฐาโน จ ธมฺมา จิตฺต\-สมุฏฺฐานสฺส จ โนจิตฺต\-สมุฏฺฐานสฺส จ ธมฺมสฺส อารมฺมณ\-ปจฺจเยน ปจฺจโย. สหชาต\-ปจฺจเยน ปจฺจโย. อุปนิสฺสย\-ปจฺจเยน ปจฺจโย. ปุเรชาต\-ปจฺจเยน ปจฺจโย. ปจฺฉาชาต\-ปจฺจเยน ปจฺจโย. (3)}

\csromanmarkh{2. ปจฺจย\-ปจฺจนีย}\csromanmarkhii{2. สงฺขฺยาวาร}\csromanheadernomark{2}{2. ปจฺจย\-ปจฺจนีย 2. สงฺขฺยาวาร}
\vspace{\tipitakaparskip}
\csromancenter{นเหตุยา นว, นอารมฺมเณ นว, (สพฺพตฺถ นว,) โนอวิคเต นว.}
\csromancenter{\textbf{เหตุ}ปจฺจยา นอารมฺมเณ ตีณิ, นอธิปติยา ตีณิ, นอนนฺตเร ตีณิ, นสมนนฺตเร ตีณิ, นอญฺญมญฺเญ เทฺว, นฺเอาปนิสฺสเย ตีณิ ฯเปฯ นมคฺเค ตีณิ, นสมฺปยุตฺเต เทฺว, นวิปฺปยุตฺเต ตีณิ, โนนตฺถิยา ตีณิ, โนวิคเต ตีณิ.}
\proseitemwithcloser{238}{นเหตุปจฺจยา อารมฺมเณ นว, อธิปติยา นว, (อนุโลมคณนา กาตพฺพา.) ฯเปฯ อวิคเต นว.}{จิตฺต\-สมุฏฺฐานทุกํ นิฏฺฐิตํ.}
\csromanmatikaanchor{mtk.60}
\csromantocmarkref{subsection}{61. จิตฺตสหภูทุก}{mtk.60}
\bookmark[dest={mtk.60},level=2]{61. จิตฺตสหภูทุก}
\csromanmarkh{61. จิตฺต\-สหภูทุก}\csromanmarkhii{1. ปฏิจฺจวาร}\csromanheadernomark{2}{61. จิตฺต\-สหภูทุก 1. ปฏิจฺจวาร}
\csromanmarkh{1. ปจฺจยานุโลม}\csromanmarkhii{1. วิภงฺควาร}\csromanheadernomark{2}{1. \textbf{ปจฺจยานุโลม 1. วิภงฺควาร}}
\vspace{\tipitakaparskip}
\csromancenter{\textbf{เหตุ}}
\proseitem{239}{จิตฺตสหภุํ ธมฺมํ ปฏิจฺจ จิตฺตสหภู ธมฺโม อุปฺปชฺชติ เหตุปจฺจยา. จิตฺตสหภุํ เอกํ ขนฺธํ ปฏิจฺจ เทฺว ขนฺธา จิตฺตสหภุ จิตฺต\-สมุฏฺฐานญฺจ \csromanfolio{462}รูปํ. เทฺว ขนฺเธ ฯเปฯ. ปฏิสนฺธิ\-กฺขเณ จิตฺตสหภุํ เอกํ ขนฺธํ ปฏิจฺจ เทฺว ขนฺธา. เทฺว ขนฺเธ ฯเปฯ. (1)}

\prose{จิตฺตสหภุํ ธมฺมํ ปฏิจฺจ โนจิตฺตสหภู ธมฺโม อุปฺปชฺชติ เหตุปจฺจยา. จิตฺตสหภู ขนฺเธ ปฏิจฺจ จิตฺตํ โนจิตฺตสหภุ จิตฺต\-สมุฏฺฐานญฺจ รูปํ. ปฏิสนฺธิ\-กฺขเณ จิตฺตสหภู ขนฺเธ ปฏิจฺจ จิตฺตํ กฏตฺตา จ รูปํ. (2)}

\prose{จิตฺตสหภุํ ธมฺมํ ปฏิจฺจ จิตฺตสหภู จ โนจิตฺตสหภู จ ธมฺมา อุปฺปชฺชนฺติ เหตุปจฺจยา. จิตฺตสหภุํ เอกํ ขนฺธํ ปฏิจฺจ เทฺว ขนฺธา จิตฺตญฺจ จิตฺตสหภุ จ โนจิตฺตสหภุ จ จิตฺต\-สมุฏฺฐานํ รูปํ. เทฺว ขนฺเธ ฯเปฯ. ปฏิสนฺธิ\-กฺขเณ จิตฺตสหภุํ เอกํ ขนฺธํ ปฏิจฺจ เทฺว ขนฺธา จิตฺตญฺจ กฏตฺตา จ รูปํ. เทฺว ขนฺเธ ฯเปฯ. (3)}

\proseitem{240}{โนจิตฺต\-สหภุํ ธมฺมํ ปฏิจฺจ โนจิตฺตสหภู ธมฺโม อุปฺปชฺชติ เหตุปจฺจยา. จิตฺตํ ปฏิจฺจ โนจิตฺตสหภุ จิตฺต\-สมุฏฺฐานํ รูปํ. ปฏิสนฺธิ\-กฺขเณ จิตฺตํ ปฏิจฺจ กฏตฺตารูปํ. จิตฺตํ ปฏิจฺจ วตฺถุ, วตฺถุํ ปฏิจฺจ จิตฺตํ. เอกํ มหาภูตํ ฯเปฯ. มหาภูเต ปฏิจฺจ โนจิตฺตสหภุ จิตฺต\-สมุฏฺฐานํ รูปํ, กฏตฺตารูปํ, อุปาทารูปํ. (1)}

\prose{โนจิตฺต\-สหภุํ ธมฺมํ ปฏิจฺจ จิตฺตสหภู ธมฺโม อุปฺปชฺชติ เหตุปจฺจยา. จิตฺตํ ปฏิจฺจ สมฺปยุตฺตกา ขนฺธา จิตฺตสหภุ จิตฺต\-สมุฏฺฐานญฺจ รูปํ. ปฏิสนฺธิ\-กฺขเณ จิตฺตํ ปฏิจฺจ สมฺปยุตฺตกา ขนฺธา. ปฏิสนฺธิ\-กฺขเณ วตฺถุํ ปฏิจฺจ จิตฺตสหภู ขนฺธา. มหาภูเต ปฏิจฺจ จิตฺตสหภุ จิตฺต\-สมุฏฺฐานํ รูปํ, อุปาทารูปํ. (2)}

\prose{โนจิตฺต\-สหภุํ ธมฺมํ ปฏิจฺจ จิตฺตสหภู จ โนจิตฺตสหภู จ ธมฺมา อุปฺปชฺชนฺติ เหตุปจฺจยา. จิตฺตํ ปฏิจฺจ สมฺปยุตฺตกา ขนฺธา จิตฺตสหภุ จ โนจิตฺตสหภุ จ จิตฺต\-สมุฏฺฐานํ รูปํ. ปฏิสนฺธิ\-กฺขเณ จิตฺตํ ปฏิจฺจ สมฺปยุตฺตกา ขนฺธา กฏตฺตา จ รูปํ. ปฏิสนฺธิ\-กฺขเณ วตฺถุํ ปฏิจฺจ จิตฺตญฺจ สมฺปยุตฺตกา จ ขนฺธา. มหาภูเต ปฏิจฺจ จิตฺตสหภุ จ โนจิตฺตสหภุ จ จิตฺต\-สมุฏฺฐานํ รูปํ, อุปาทารูปํ. (3)}

\proseitem{241}{จิตฺต\-สหภุญฺจ โนจิตฺต\-สหภุญฺจ ธมฺมํ ปฏิจฺจ จิตฺตสหภู ธมฺโม อุปฺปชฺชติ เหตุปจฺจยา. จิตฺตสหภุํ เอกํ ขนฺธญฺจ จิตฺตญฺจ ปฏิจฺจ เทฺว ขนฺธา จิตฺตสหภุ จิตฺต\-สมุฏฺฐานญฺจ รูปํ. เทฺว ขนฺเธ จ ฯเปฯ. ปฏิสนฺธิ\-กฺขเณ จิตฺตสหภุํ เอกํ ขนฺธญฺจ จิตฺตญฺจ ปฏิจฺจ เทฺว ขนฺธา. เทฺว ขนฺเธ จ ฯเปฯ. ปฏิสนฺธิ\-กฺขเณ จิตฺตสหภุํ \csromanfolio{463}เอกํ ขนฺธญฺจ วตฺถุญฺจ ปฏิจฺจ เทฺว ขนฺธา. เทฺว ขนฺเธ จ ฯเปฯ. จิตฺตสหภู ขนฺเธ จ มหาภูเต จ ปฏิจฺจ จิตฺตสหภุ จิตฺต\-สมุฏฺฐานํ รูปํ, อุปาทารูปํ. (1)}

\prose{จิตฺต\-สหภุญฺจ โนจิตฺต\-สหภุญฺจ ธมฺมํ ปฏิจฺจ โนจิตฺตสหภู ธมฺโม อุปฺปชฺชติ เหตุปจฺจยา. จิตฺตสหภู ขนฺเธ จ จิตฺตญฺจ ปฏิจฺจ โนจิตฺตสหภุ จิตฺต\-สมุฏฺฐานํ รูปํ. ปฏิสนฺธิ\-กฺขเณ จิตฺตสหภู ขนฺเธ จ จิตฺตญฺจ ปฏิจฺจ กฏตฺตารูปํ. ปฏิสนฺธิ\-กฺขเณ จิตฺตสหภู ขนฺเธ จ วตฺถุญฺจ ปฏิจฺจ จิตฺตํ. จิตฺตสหภู ขนฺเธ จ มหาภูเต จ ปฏิจฺจ โนจิตฺตสหภุ จิตฺต\-สมุฏฺฐานํ รูปํ, กฏตฺตารูปํ, อุปาทารูปํ.}

\prose{จิตฺต\-สหภุญฺจ โนจิตฺต\-สหภุญฺจ ธมฺมํ ปฏิจฺจ จิตฺตสหภู จ โนจิตฺตสหภู จ ธมฺมา อุปฺปชฺชนฺติ เหตุปจฺจยา. จิตฺตสหภุํ เอกํ ขนฺธญฺจ จิตฺตญฺจ ปฏิจฺจ เทฺว ขนฺธา จิตฺตสหภุ จ โนจิตฺตสหภุ จ จิตฺต\-สมุฏฺฐานํ รูปํ. เทฺว ขนฺเธ จ ฯเปฯ. ปฏิสนฺธิ\-กฺขเณ จิตฺตสหภุํ เอกํ ขนฺธญฺจ จิตฺตญฺจ ปฏิจฺจ เทฺว ขนฺธา กฏตฺตา จ รูปํ. เทฺว ขนฺเธ จ ฯเปฯ. ปฏิสนฺธิ\-กฺขเณ จิตฺตสหภุํ เอกํ ขนฺธญฺจ วตฺถุญฺจ ปฏิจฺจ เทฺว ขนฺธา จิตฺตญฺจ. เทฺว ขนฺเธ จ ฯเปฯ. จิตฺตสหภู ขนฺเธ จ มหาภูเต จ ปฏิจฺจ กฏตฺตารูปํ, อุปาทารูปํ. (สํขิตฺตํ.) (3)}

\csromanmarkh{1. ปจฺจยานุโลม}\csromanmarkhii{2. สงฺขฺยาวาร}\csromanheadernomark{2}{1. \textbf{ปจฺจยานุโลม 2. สงฺขฺยาวาร}}
\proseitem{242}{เหตุยา นว, อารมฺมเณ นว, (อรูปํ สพฺพํ อุทฺธริตพฺพํ. จิตฺตสมุฏฺฐาน\-ทุก\-สทิสํ.) อธิปติยา นว, (มหาภูตา ฉสุปิ ปเญฺหสุ กาตพฺพา. อธิปติยา ตีสุ นตฺถิ.) อนนฺตเร นว, สมนนฺตเร นว, สหชาเต นว, อญฺญมญฺเญ นว, นิสฺสเย นว, อุปนิสฺสเย นว, ปุเรชาเต ปญฺจ, อาเสวเน ปญฺจ, กมฺเม นว, (สพฺพตฺถ นว,) อวิคเต นว.}

\csromanmarkh{2. ปจฺจย\-ปจฺจนีย}\csromanmarkhii{2. วิภงฺควาร}\csromanheadernomark{2}{2. ปจฺจย\-ปจฺจนีย 2. วิภงฺควาร}
\proseitem{243}{จิตฺตสหภุํ ธมฺมํ ปฏิจฺจ จิตฺตสหภู ธมฺโม อุปฺปชฺชติ นเหตุปจฺจยา. อเหตุกํ จิตฺตสหภุํ เอกํ ขนฺธํ ปฏิจฺจ เทฺว ขนฺธา จิตฺตสหภุ จิตฺต\-สมุฏฺฐานญฺจ รูปํ. เทฺว ขนฺเธ ฯเปฯ. อเหตุก\-ปฏิสนฺธิ\-กฺขเณ ฯเปฯ. วิจิกิจฺฉา\-สหคเต \csromanfolio{464}อุทฺธจฺจ\-สหคเต ขนฺเธ ปฏิจฺจ วิจิกิจฺฉา\-สหคโต อุทฺธจฺจ\-สหคโต โมโห. (เอวํ นวปิ ปญฺหา กาตพฺพา. “อเหตุกนฺ”ติ นิยาเมตพฺพํ. ยถา อนุโลเม ลพฺภติ, เอวํ กาตพฺพํ. ตีณิ, โมโห. ยถา จิตฺตสมุฏฺฐาน\-ทุเก, เอวเมว กาตพฺพา.)}

\csromanheader{2}{\textbf{นกมฺม}}
\proseitem{244}{จิตฺตสหภุํ ธมฺมํ ปฏิจฺจ จิตฺตสหภู ธมฺโม อุปฺปชฺชติ นกมฺมปจฺจยา. จิตฺตสหภู ขนฺเธ ปฏิจฺจ จิตฺตสหภู เจตนา.}

\prose{โนจิตฺต\-สหภุํ ธมฺมํ ปฏิจฺจ โนจิตฺตสหภู ธมฺโม อุปฺปชฺชติ นกมฺมปจฺจยา. พาหิรํ. อาหารสมุฏฺฐานํ. อุตุสมุฏฺฐานํ ฯเปฯ.}

\prose{โนจิตฺต\-สหภุํ ธมฺมํ ปฏิจฺจ จิตฺตสหภู ธมฺโม อุปฺปชฺชติ นกมฺมปจฺจยา. จิตฺตํ ปฏิจฺจ สมฺปยุตฺตกา เจตนา.}

\prose{จิตฺต\-สหภุญฺจ โนจิตฺต\-สหภุญฺจ ธมฺมํ ปฏิจฺจ จิตฺตสหภู ธมฺโม อุปฺปชฺชติ นกมฺมปจฺจยา. จิตฺตสหภู ขนฺเธ จ จิตฺตญฺจ ปฏิจฺจ สมฺปยุตฺตกา เจตนา.}

\csromanheader{2}{\textbf{นฌาน}}
\proseitem{245}{จิตฺตสหภุํ ธมฺมํ ปฏิจฺจ จิตฺตสหภู ธมฺโม อุปฺปชฺชติ นฌานปจฺจยา. ปญฺจ\-วิญฺญาณสหคตํ ฯเปฯ.}

\csromanmarkh{2. ปจฺจย\-ปจฺจนีย}\csromanmarkhii{2. สงฺขฺยาวาร}\csromanheadernomark{2}{2. ปจฺจย\-ปจฺจนีย 2. สงฺขฺยาวาร}
\proseitem{246}{นเหตุยา นว, นอารมฺมเณ นว, นอธิปติยา นว, นอนนฺตเร นว, นสมนนฺตเร นว, นอญฺญมญฺเญ นว, นอุปนิสฺสเย นว, นปุเรชาเต นว, นปจฺฉาชาเต นว, นอาเสวเน นว. นกมฺเม จตฺตาริ, นวิปาเก นว, นอาหาเร เอกํ, นอินฺทฺริเย เอกํ, นฌาเน ฉ, นมคฺเค นว, นสมฺปยุตฺเต นว, นวิปฺปยุตฺเต ฉ, โนนตฺถิยา นว, โนวิคเต นว. (เอวํ อิตเร เทฺวปิ คณนา กาตพฺพา.)}

\vspace{\tipitakaparskip}
\csromancenter{(สหชาตวาโรปิ ปฏิจฺจ\-วาร\-สทิโส.)}
\csromanmarkh{61. จิตฺต\-สหภูทุก}\csromanmarkhii{61. จิตฺต\-สหภูทุก 3. ปจฺจยวาร}\csromanheadernomark{2}{61. จิตฺต\-สหภูทุก 61. จิตฺต\-สหภูทุก 3. ปจฺจยวาร}
\csromanmarkh{1. ปจฺจยานุโลม}\csromanmarkhii{1. วิภงฺควาร}\csromanheadernomark{2}{1. \textbf{ปจฺจยานุโลม 1. วิภงฺควาร}}
\csromanheader{2}{\textbf{เหตุ}}
\proseitem{247}{\csromanfolio{465}จิตฺตสหภุํ ธมฺมํ ปจฺจยา จิตฺตสหภู ธมฺโม อุปฺปชฺชติ เหตุปจฺจยา. ตีณิ. (ปฏิจฺจสทิสา.)}

\prose{โนจิตฺต\-สหภุํ ธมฺมํ ปจฺจยา โนจิตฺตสหภู ธมฺโม อุปฺปชฺชติ เหตุปจฺจยา. วตฺถุํ ปจฺจยา จิตฺตํ. จิตฺตํ ปจฺจยา โนจิตฺตสหภุ จิตฺต\-สมุฏฺฐานํ รูปํ. ปฏิสนฺธิ\-กฺขเณ ฯเปฯ. (ปฏิจฺจ\-วาร\-สทิสํ. สพฺเพ มหาภูตา.) (1)}

\prose{โนจิตฺต\-สหภุํ ธมฺมํ ปจฺจยา จิตฺตสหภู ธมฺโม อุปฺปชฺชติ เหตุปจฺจยา. จิตฺตํ ปจฺจยา สมฺปยุตฺตกา ขนฺธา จิตฺตสหภุ จิตฺต\-สมุฏฺฐานญฺจ รูปํ. วตฺถุํ ปจฺจยา จิตฺตสหภู ขนฺธา. ปฏิสนฺธิ\-กฺขเณ ฯเปฯ. (สพฺเพ มหาภูตา. ปฏิจฺจสทิสํ.) (2)}

\prose{โนจิตฺต\-สหภุํ ธมฺมํ ปจฺจยา จิตฺตสหภู จ โนจิตฺตสหภู จ ธมฺมา อุปฺปชฺชนฺติ เหตุปจฺจยา. จิตฺตํ ปจฺจยา สมฺปยุตฺตกา ขนฺธา จิตฺตสหภุ จ โนจิตฺตสหภุ จ จิตฺต\-สมุฏฺฐานํ รูปํ. วตฺถุํ ปจฺจยา จิตฺตญฺจ สมฺปยุตฺตกา จ ขนฺธา. ปฏิสนฺธิ\-กฺขเณ ฯเปฯ. (ปฏิจฺจสทิสํ. สพฺเพ มหาภูตา.) (3)}

\proseitem{248}{จิตฺต\-สหภุญฺจ โนจิตฺต\-สหภุญฺจ ธมฺมํ ปจฺจยา จิตฺตสหภู ธมฺโม อุปฺปชฺชติ เหตุปจฺจยา. จิตฺตสหตุํ เอกํ ขนฺธญฺจ จิตฺตญฺจ ปจฺจยา เทฺว ขนฺธา จิตฺตสหภุ จิตฺต\-สมุฏฺฐานญฺจ รูปํ. เทฺว ขนฺเธ จ ฯเปฯ. จิตฺตสหภุํ เอกํ ขนฺธญฺจ วตฺถุญฺจ ปจฺจยา เทฺว ขนฺธา. เทฺว ขนฺเธ จ ฯเปฯ. ปฏิสนฺธิ\-กฺขเณ ฯเปฯ. (ปฏิจฺจสทิสํ. สพฺเพ มหาภูตา.) (1)}

\prose{จิตฺต\-สหภุญฺจ โนจิตฺต\-สหภุญฺจ ธมฺมํ ปจฺจยา โนจิตฺตสหภู ธมฺโม อุปฺปชฺชติ เหตุปจฺจยา. จิตฺตสหภู ขนฺเธ จ จิตฺตญฺจ ปจฺจยา โนจิตฺตสหภุ จิตฺต\-สมุฏฺฐานํ รูปํ. จิตฺตสหภู ขนฺเธ จ วตฺถุญฺจ ปจฺจยา จิตฺตํ. จิตฺตสหภู ขนฺเธ จ มหาภูเต จ ปจฺจยา โนจิตฺตสหภุ จิตฺต\-สมุฏฺฐานํ รูปํ. ปฏิสนฺธิ\-กฺขเณ ฯเปฯ. (ปฏิจฺจสทิสํ. สพฺเพ มหาภูตา.) (2)}

\prose{จิตฺต\-สหภุญฺจ โนจิตฺต\-สหภุญฺจ ธมฺมํ ปจฺจยา จิตฺตสหภู จ โนจิตฺตสหภู จ ธมฺมา อุปฺปชฺชนฺติ เหตุปจฺจยา. จิตฺตสหภุํ เอกํ ขนฺธญฺจ จิตฺตญฺจ \csromanfolio{466}ปจฺจยา เทฺว ขนฺธา จิตฺตสหภุ จ โนจิตฺตสหภุ จ จิตฺต\-สมุฏฺฐานํ รูปํ. จิตฺตสหภุํ เอกํ ขนฺธญฺจ วตฺถุญฺจ ปจฺจยา เทฺว ขนฺธา จิตฺตญฺจ. เทฺว ขนฺเธ จ ฯเปฯ. ปฏิสนฺธิ\-กฺขเณ ฯเปฯ. (ปฏิจฺจสทิสํ. สพฺเพ มหาภูตา.) (3)}

\proseitem{249}{จิตฺตสหภุํ ธมฺมํ ปจฺจยา จิตฺตสหภู ธมฺโม อุปฺปชฺชติ อารมฺมณ\-ปจฺจยา. ตีณิ. (ปฏิจฺจสทิสํ.)}

\prose{โนจิตฺต\-สหภุํ ธมฺมํ ปจฺจยา โนจิตฺตสหภู ธมฺโม อุปฺปชฺชติ อารมฺมณ\-ปจฺจยา ฯเปฯ. จกฺขายตนํ ปจฺจยา จกฺขุวิญฺญาณํ ฯเปฯ. กายายตนํ ฯเปฯ. (อิมํ จิตฺต\-สมุฏฺฐานทุกํ ปจฺจยวาเร อารมฺมณ\-สทิสํ ฉนฺนมฺปิ. อิเมสํ ปญฺจวิญฺญาณ\-มูลา กาตพฺพา. สํขิตฺตํ.)}

\csromanmarkh{1. ปจฺจยานุโลม}\csromanmarkhii{2. สงฺขฺยาวาร}\csromanheadernomark{2}{1. \textbf{ปจฺจยานุโลม 2. สงฺขฺยาวาร}}
\vspace{\tipitakaparskip}
\csromancenter{เหตุยา นว, อารมฺมเณ นว, อธิปติยา นว, (สพฺพตฺถ นว.) อวิคเต นว.}
\csromanmarkh{2. ปจฺจย\-ปจฺจนีย}\csromanmarkhii{1. วิภงฺควาร}\csromanheadernomark{2}{2. ปจฺจย\-ปจฺจนีย 1. วิภงฺควาร}
\proseitem{251}{จิตฺตสหภุํ ธมฺมํ ปจฺจยา จิตฺตสหภู ธมฺโม อุปฺปชฺชติ นเหตุปจฺจยา. อเหตุกํ จิตฺตสหภุํ เอกํ ขนฺธํ ฯเปฯ. (สํขิตฺตํ. สพฺพํ กาตพฺพํ.) (ปจฺจยวารสฺส ปญฺจวิญฺญาณํ ฉนฺนมฺปิ มูลา กาตพฺพา. สพฺเพ มหาภูเต ตีณิเยว. โมโห. สํขิตฺตํ.)}

\csromanmarkh{2. ปจฺจย\-ปจฺจนีย}\csromanmarkhii{2. สงฺขฺยาวาร}\csromanheadernomark{2}{2. ปจฺจย\-ปจฺจนีย 2. สงฺขฺยาวาร}
\vspace{\tipitakaparskip}
\csromancenter{\textbf{นเหตุ}ยา นว, นอารมฺมเณ นว, นอธิปติยา นว, นอนนฺตเร นว, นสมนนฺตเร นว, นอญฺญมญฺเญ นว, นฺเอาปนิสฺสเย นว, นปุเรชาเต}
\prose{\csromanfolio{467}นว, นปจฺฉาชาเต นว, นอาเสวเน นว, นกมฺเม จตฺตาริ, นวิปาเก นว, นอาหาเร เอกํ, ไนนฺทฺริเย เอกํ, นฌาเน นว, นมคฺเค นว, นสมฺปยุตฺเต นว, นวิปฺปยุตฺเต ฉ, โนนตฺถิยา นว, โนวิคเต นว.}

\proseitem{253}{เหตุปจฺจยา นอารมฺมเณ นว, (สพฺพตฺถ นว.) นกมฺเม ตีณิ, นวิปาเก นว, นสมฺปยุตฺเต นว, นวิปฺปยุตฺเต ปญฺจ, โนนตฺถิยา นว, โนวิคเต นว.}

\proseitem{254}{นเหตุปจฺจยา อารมฺมเณ นว, อนนฺตเร นว, (สพฺพตฺถ นว.) มคฺเค ตีณิ ฯเปฯ อวิคเต นว.}

\vspace{\tipitakaparskip}
\csromancenter{(นิสฺสยวาโร ปจฺจยวาร\-สทิโส.)}
\csromanmarkh{61. จิตฺต\-สหภูทุก}\csromanmarkhii{5. สํสฏฺฐวาร}\csromanheadernomark{2}{61. จิตฺต\-สหภูทุก 5. สํสฏฺฐวาร}
\proseitem{255}{จิตฺตสหภุํ ธมฺมํ สํสฏฺโฐ จิตฺตสหภู ธมฺโม อุปฺปชฺชติ เหตุปจฺจยา. จิตฺตสหภุํ เอกํ ขนฺธํ สํสฏฺฐา เทฺว ขนฺธา. เทฺว ขนฺเธ ฯเปฯ. ปฏิสนฺธิ\-กฺขเณ ฯเปฯ. (1)}

\prose{จิตฺตสหภุํ ธมฺมํ สํสฏฺโฐ โนจิตฺตสหภู ธมฺโม อุปฺปชฺชติ เหตุปจฺจยา. จิตฺตสหภู ขนฺเธ สํสฏฺฐํ จิตฺตํ. ปฏิสนฺธิ\-กฺขเณ ฯเปฯ. (2)}

\prose{จิตฺตสหภุํ ธมฺมํ สํสฏฺโฐ จิตฺตสหภู จ โนจิตฺตสหภู จ ธมฺมา อุปฺปชฺชนฺติ เหตุปจฺจยา. จิตฺตสหภุํ เอกํ ขนฺธํ สํสฏฺฐา เทฺว ขนฺธา จิตฺตญฺจ. เทฺว ขนฺเธ ฯเปฯ. ปฏิสนฺธิ\-กฺขเณ ฯเปฯ. (3)}

\proseitem{256}{โนจิตฺต\-สหภุํ ธมฺมํ สํสฏฺโฐ จิตฺตสหภู ธมฺโม อุปฺปชฺชติ เหตุปจฺจยา. จิตฺตํ สํสฏฺฐา สมฺปยุตฺตกา ขนฺธา. ปฏิสนฺธิ\-กฺขเณ ฯเปฯ. (1)}

\prose{\csromanfolio{468}จิตฺต\-สหภุญฺจ โนจิตฺต\-สหภุญฺจ ธมฺมํ สํสฏฺโฐ จิตฺตสหภู ธมฺโม อุปฺปชฺชติ เหตุปจฺจยา. จิตฺตสหภุํ เอกํ ขนฺธญฺจ จิตฺตญฺจ สํสฏฺฐา เทฺว ขนฺธา. เทฺว ขนฺเธ จ ฯเปฯ. ปฏิสนฺธิ\-กฺขเณ ฯเปฯ. (สํขิตฺตํ.) (2)}

\prose{เหตุยา ปญฺจ, อารมฺมเณ ปญฺจ, (สพฺพตฺถ ปญฺจ.) อวิคเต ปญฺจ. อนุโลมํ.}

\prose{นเหตุยา ปญฺจ, (ตีณิ, โมโห.) นอธิปติยา ปญฺจ, นปุเรชาเต ปญฺจ, นปจฺฉาชาเต ปญฺจ, นอาเสวเน ปญฺจ, นกมฺเม ตีณิ, นวิปาเก ปญฺจ, นฌาเน ปญฺจ, นมคฺเค ปญฺจ, นวิปฺปยุตฺเต ปญฺจ.}

\vspace{\tipitakaparskip}
\csromancenter{(อิตเร เทฺว คณนาปิ สมฺปยุตฺตวาโรปิ สพฺพํ กาตพฺพํ.)}
\csromanmarkh{61. จิตฺต\-สหภูทุก}\csromanmarkhii{7. ปญฺหาวาร}\csromanheadernomark{2}{61. จิตฺต\-สหภูทุก 7. ปญฺหาวาร}
\csromanmarkh{1. ปจฺจยานุโลม}\csromanmarkhii{1. วิภงฺควาร}\csromanheadernomark{2}{1. \textbf{ปจฺจยานุโลม 1. วิภงฺควาร}}
\vspace{\tipitakaparskip}
\csromancenter{\textbf{เหตุ}}
\proseitem{257}{จิตฺตสหภู ธมฺโม จิตฺต\-สหภุสฺส ธมฺมสฺส เหตุปจฺจเยน ปจฺจโย. จิตฺตสหภู เหตู สมฺปยุตฺตกานํ ขนฺธานํ จิตฺต\-สหภูนํ จิตฺตสมุฏฺฐานานญฺจ รูปานํ เหตุปจฺจเยน ปจฺจโย. ปฏิสนฺธิ\-กฺขเณ ฯเปฯ. (1)}

\prose{จิตฺตสหภู ธมฺโม โนจิตฺต\-สหภุสฺส ธมฺมสฺส เหตุปจฺจเยน ปจฺจโย. จิตฺตสหภู เหตู จิตฺตสฺส โนจิตฺต\-สหภูนํ จิตฺตสมุฏฺฐานานญฺจ รูปานํ เหตุปจฺจเยน ปจฺจโย. ปฏิสนฺธิ\-กฺขเณ จิตฺตสหภู เหตู จิตฺตสฺส กฏตฺตา จ รูปานํ เหตุปจฺจเยน ปจฺจโย. (2)}

\prose{จิตฺตสหภู ธมฺโม จิตฺต\-สหภุสฺส จ โนจิตฺต\-สหภุสฺส จ ธมฺมสฺส เหตุปจฺจเยน ปจฺจโย. จิตฺตสหภู เหตู สมฺปยุตฺตกานํ ขนฺธานํ จิตฺตสฺส จ จิตฺต\-สหภูนญฺจ โนจิตฺต\-สหภูนญฺจ จิตฺต\-สมุฏฺฐานานํ รูปานํ เหตุปจฺจเยน ปจฺจโย. (3)}

\csromanheader{2}{\textbf{อรมฺมณ}}
\proseitem{258}{จิตฺตสหภู ธมฺโม จิตฺต\-สหภุสฺส ธมฺมสฺส อารมฺมณ\-ปจฺจเยน ปจฺจโย. นว. (จิตฺตสมุฏฺฐานทุกสทิสํ นินฺนานากรณํ.)}

\vspace{\tipitakaparskip}
\csromancenter{(จิตฺตสมุฏฺฐาน\-ทุก\-สทิสํ นินฺนานากรณํ.)}
\csromanheader{2}{61. จิตฺต\-สหภูทุก อธิปติ}
\proseitem{259}{\csromanfolio{469}จิตฺตสหภู ธมฺโม จิตฺต\-สหภุสฺส ธมฺมสฺส อธิปติ\-ปจฺจเยน ปจฺจโย. ตีณิ. (อารมฺมณาธิปติปิ สหชาตาธิปติปิ กาตพฺพา.)}

\vspace{\tipitakaparskip}
\csromancenter{(อารมฺมณาธิปติปิ สหชาตาธิปติปิ กาตพฺพา.)}
\prose{โนจิตฺตสหภู ธมฺโม โนจิตฺต\-สหภุสฺส ธมฺมสฺส อธิปติ\-ปจฺจเยน ปจฺจโย. ตีณิ. (อารมฺมณาธิปติปิ สหชาตาธิปติปิ อิเมสมฺปิ ติณฺณํ กาตพฺพา. นวปิ ปญฺหา จิตฺตสมุฏฺฐาน\-ทุก\-สทิสา. อนฺเต ตีณิ อารมฺมณา\-ธิปติเยว.)}

\proseitem{260}{จิตฺตสหภู ธมฺโม จิตฺต\-สหภุสฺส ธมฺมสฺส อนนฺตร\-ปจฺจเยน ปจฺจโย. นว. (จิตฺตสมุฏฺฐาน\-ทุก\-สทิสํ นินฺนานากรณํ.) สมนนฺตร\-ปจฺจเยน ปจฺจโย. นว. (ปฏิจฺจสทิสา.) สหชาต\-ปจฺจเยน ปจฺจโย. นว. (ปฏิจฺจสทิสา.) อญฺญมญฺญ\-ปจฺจเยน ปจฺจโย. นว. (ปฏิจฺจสทิสา.) นิสฺสย\-ปจฺจเยน ปจฺจโย. นว. (ปจฺจยวาร\-สทิสา.) อุปนิสฺสย\-ปจฺจเยน ปจฺจโย. นว. (จิตฺตสมุฏฺฐาน\-ทุก\-สทิสา.)}

\proseitem{261}{โนจิตฺตสหภู ธมฺโม โนจิตฺต\-สหภุสฺส ธมฺมสฺส ปุเรชาต\-ปจฺจเยน ปจฺจโย. อารมฺมณ\-ปุเรชาตํ, วตฺถุ\-ปุเรชาตํ. ตีณิ. (โนจิตฺตสหภุมูลํเยว ลพฺภติ. จิตฺตสมุฏฺฐาน\-ทุก\-สทิสา. ตีณิปิ นินฺนานากรณํ.)}

\proseitem{262}{จิตฺตสหภู ธมฺโม โนจิตฺต\-สหภุสฺส ธมฺมสฺส ปจฺฉาชาต\-ปจฺจเยน ปจฺจโย. ปจฺฉาชาตา จิตฺตสหภู ขนฺธา ปุเรชาตสฺส อิมสฺส โนจิตฺต\-สหภุสฺส กายสฺส ปจฺฉาชาต\-ปจฺจเยน ปจฺจโย. (1)}

\prose{โนจิตฺตสหภู ธมฺโม โนจิตฺต\-สหภุสฺส ธมฺมสฺส ปจฺฉาชาต\-ปจฺจเยน ปจฺจโย ฯเปฯ. (1)}

\prose{จิตฺตสหภู จ โนจิตฺตสหภู จ ธมฺมา โนจิตฺต\-สหภุสฺส ธมฺมสฺส ปจฺฉาชาต\-ปจฺจเยน ปจฺจโย. (สํขิตฺตํ.) อาเสวน\-ปจฺจเยน ปจฺจโย. นว.}

\proseitem{263}{\csromanfolio{470}จิตฺตสหภู ธมฺโม จิตฺต\-สหภุสฺส ธมฺมสฺส กมฺมปจฺจเยน ปจฺจโย. สหชาตา, นานากฺขณิกา.  สหชาตา จิตฺตสหภู เจตนา สมฺปยุตฺตกานํ ขนฺธานํ จิตฺต\-สหภูนํ จิตฺตสมุฏฺฐานานญฺจ รูปานํ กมฺมปจฺจเยน ปจฺจโย. ปฏิสนฺธิ\-กฺขเณ ฯเปฯ. นานากฺขณิกา จิตฺตสหภู เจตนา วิปากานํ จิตฺต\-สหภูนํ ขนฺธานํ กมฺมปจฺจเยน ปจฺจโย. (1)}

\prose{จิตฺตสหภู ธมฺโม โนจิตฺต\-สหภุสฺส ธมฺมสฺส กมฺมปจฺจเยน ปจฺจโย. สหชาตา, นานากฺขณิกา.  สหชาตา จิตฺตสหภู เจตนา จิตฺตสฺส โนจิตฺต\-สหภูนํ จิตฺตสมุฏฺฐานานญฺจ รูปานํ กมฺมปจฺจเยน ปจฺจโย. ปฏิสนฺธิ\-กฺขเณ ฯเปฯ. นานากฺขณิกา จิตฺตสหภู เจตนา วิปากสฺส จิตฺตสฺส กฏตฺตา จ รูปานํ กมฺมปจฺจเยน ปจฺจโย. (2)}

\prose{จิตฺตสหภู ธมฺโม จิตฺต\-สหภุสฺส จ โนจิตฺต\-สหภุสฺส จ ธมฺมสฺส กมฺมปจฺจเยน ปจฺจโย. สหชาตา, นานากฺขณิกา.  สหชาตา จิตฺตสหภู เจตนา สมฺปยุตฺตกานํ ขนฺธานํ จิตฺตสฺส จ จิตฺต\-สหภูนญฺจ โนจิตฺต\-สหภูนญฺจ จิตฺตสมุฏฺฐานานญฺจ รูปานํ กมฺมปจฺจเยน ปจฺจโย. ปฏิสนฺธิ\-กฺขเณ ฯเปฯ. นานากฺขณิกา จิตฺตสหภู เจตนา วิปากานํ ขนฺธานํ จิตฺตสฺส กฏตฺตา จ รูปานํ กมฺมปจฺจเยน ปจฺจโย. (3)}

\csromanheader{2}{\textbf{วิปากาทิ}}
\proseitem{264}{จิตฺตสหภู ธมฺโม จิตฺต\-สหภุสฺส ธมฺมสฺส วิปาก\-ปจฺจเยน ปจฺจโย. (จิตฺตสมุฏฺฐาน\-ทุก\-สทิสํ.) อาหารปจฺจเยน ปจฺจโย. นว. (จิตฺตสมุฏฺฐาน\-ทุก\-สทิสา. อิมมฺปิ เอกํ กพฬีการ\-อาหาร\-สทิสํ.) จิตฺตสหภู ธมฺโม จิตฺต\-สหภุสฺส ธมฺมสฺส อินฺทฺริย\-ปจฺจเยน ปจฺจโย. นว. (จิตฺตสมุฏฺฐาน\-ทุก\-สทิสํ นินฺนานากรณํ.) ฌานปจฺจเยน ปจฺจโย. ตีณิ. มคฺคปจฺจเยน ปจฺจโย. ตีณิ. สมฺปยุตฺต\-ปจฺจเยน ปจฺจโย. ปญฺจ.}

\proseitem{265}{จิตฺตสหภู ธมฺโม จิตฺต\-สหภุสฺส ธมฺมสฺส วิปฺปยุตฺต\-ปจฺจเยน ปจฺจโย. สหชาตา จิตฺตสหภู ขนฺธา จิตฺต\-สหภูนํ จิตฺต\-สมุฏฺฐานานํ รูปานํ วิปฺปยุตฺต\-ปจฺจเยน ปจฺจโย. (1)}

\prose{\csromanfolio{471}จิตฺตสหภู ธมฺโม โนจิตฺต\-สหภุสฺส ธมฺมสฺส วิปฺปยุตฺต\-ปจฺจเยน ปจฺจโย. สหชาตํ, ปจฺฉาชาตํ.  สหชาตา จิตฺตสหภู ขนฺธา โนจิตฺต\-สหภูนํ จิตฺต\-สมุฏฺฐานานํ รูปานํ วิปฺปยุตฺต\-ปจฺจเยน ปจฺจโย. ปฏิสนฺธิ\-กฺขเณ ฯเปฯ. ปจฺฉาชาตา จิตฺตสหภู ขนฺธา ปุเรชาตสฺส อิมสฺส โนจิตฺต\-สหภุสฺส กายสฺส วิปฺปยุตฺต\-ปจฺจเยน ปจฺจโย. (2)}

\prose{จิตฺตสหภู ธมฺโม จิตฺต\-สหภุสฺส จ โนจิตฺต\-สหภุสฺส จ ธมฺมสฺส วิปฺปยุตฺต\-ปจฺจเยน ปจฺจโย. สหชาตา จิตฺตสหภู ขนฺธา จิตฺต\-สหภูนญฺจ โนจิตฺต\-สหภูนญฺจ จิตฺต\-สมุฏฺฐานานํ รูปานํ วิปฺปยุตฺต\-ปจฺจเยน ปจฺจโย. (3)}

\proseitem{266}{โนจิตฺตสหภู ธมฺโม โนจิตฺต\-สหภุสฺส ธมฺมสฺส วิปฺปยุตฺต\-ปจฺจเยน ปจฺจโย. สหชาตํ, ปุเรชาตํ, ปจฺฉาชาตํ.  สหชาตํ จิตฺตํ โนจิตฺต\-สหภูนํ จิตฺต\-สมุฏฺฐานานํ รูปานํ วิปฺปยุตฺต\-ปจฺจเยน ปจฺจโย. ปฏิสนฺธิ\-กฺขเณ จิตฺตํ กฏตฺตารูปานํ วิปฺปยุตฺต\-ปจฺจเยน ปจฺจโย. จิตฺตํ วตฺถุสฺส วิปฺปยุตฺต\-ปจฺจเยน ปจฺจโย. วตฺถุ จิตฺตสฺส วิปฺปยุตฺต\-ปจฺจเยน ปจฺจโย. ปุเรชาตํ จกฺขายตนํ จกฺขุ\-วิญฺญาณสฺส วิปฺปยุตฺต\-ปจฺจเยน ปจฺจโย ฯเปฯ. กายายตนํ ฯเปฯ. วตฺถุ จิตฺตสฺส วิปฺปยุตฺต\-ปจฺจเยน ปจฺจโย. ปจฺฉาชาตํ จิตฺตํ ปุเรชาตสฺส อิมสฺส โนจิตฺต\-สหภุสฺส กายสฺส วิปฺปยุตฺต\-ปจฺจเยน ปจฺจโย. (1)}

\prose{โนจิตฺตสหภู ธมฺโม จิตฺต\-สหภุสฺส ธมฺมสฺส วิปฺปยุตฺต\-ปจฺจเยน ปจฺจโย. สหชาตํ, ปุเรชาตํ.  สหชาตํ จิตฺตํ จิตฺต\-สหภูนํ จิตฺต\-สมุฏฺฐานานํ รูปานํ วิปฺปยุตฺต\-ปจฺจเยน ปจฺจโย. ปุเรชาตํ จกฺขายตนํ จกฺขุ\-วิญฺญาณสหคตานํ ขนฺธานํ วิปฺปยุตฺต\-ปจฺจเยน ปจฺจโย ฯเปฯ. กายายตนํ ฯเปฯ. วตฺถุ จิตฺต\-สหภูนํ ขนฺธานํ วิปฺปยุตฺต\-ปจฺจเยน ปจฺจโย. (2)}

\prose{โนจิตฺตสหภู ธมฺโม จิตฺต\-สหภุสฺส จ โนจิตฺต\-สหภุสฺส จ ธมฺมสฺส วิปฺปยุตฺต\-ปจฺจเยน ปจฺจโย. สหชาตํ, ปุเรชาตํ.  สหชาตํ จิตฺตํ จิตฺต\-สหภูนญฺจ โนจิตฺต\-สหภูนญฺจ จิตฺต\-สมุฏฺฐานานํ รูปานํ วิปฺปยุตฺต\-ปจฺจเยน ปจฺจโย. ปฏิสนฺธิ\-กฺขเณ ฯเปฯ. วตฺถุ\-ปุเรชาตํ จกฺขายตนํ จกฺขุ\-วิญฺญาณสฺส สมฺปยุตฺตกานญฺจ ขนฺธานํ วิปฺปยุตฺต\-ปจฺจเยน ปจฺจโย ฯเปฯ. กายายตนํ ฯเปฯ. วตฺถุ จิตฺตสฺส สมฺปยุตฺตกานญฺจ ขนฺธานํ วิปฺปยุตฺต\-ปจฺจเยน ปจฺจโย. (3)}

\proseitem{267}{\csromanfolio{472}จิตฺตสหภู จ โนจิตฺตสหภู จ ธมฺมา จิตฺต\-สหภุสฺส ธมฺมสฺส วิปฺปยุตฺต\-ปจฺจเยน ปจฺจโย. สหชาตา จิตฺตสหภู ขนฺธา จ จิตฺตญฺจ จิตฺต\-สหภูนํ จิตฺต\-สมุฏฺฐานานํ รูปานํ วิปฺปยุตฺต\-ปจฺจเยน ปจฺจโย. (1)}

\prose{จิตฺตสหภู จ โนจิตฺตสหภู จ ธมฺมา โนจิตฺต\-สหภุสฺส ธมฺมสฺส วิปฺปยุตฺต\-ปจฺจเยน ปจฺจโย. สหชาตํ, ปจฺฉาชาตํ.  สหชาตา จิตฺตสหภู ขนฺธา จ จิตฺตญฺจ โนจิตฺต\-สหภูนํ จิตฺต\-สมุฏฺฐานานํ รูปานํ วิปฺปยุตฺต\-ปจฺจเยน ปจฺจโย. ปฏิสนฺธิ\-กฺขเณ ฯเปฯ. ปจฺฉาชาตา จิตฺตสหภู ขนฺธา จ จิตฺตญฺจ ปุเรชาตสฺส อิมสฺส โนจิตฺต\-สหภุสฺส กายสฺส วิปฺปยุตฺต\-ปจฺจเยน ปจฺจโย.}

\prose{จิตฺตสหภู จ โนจิตฺตสหภู จ ธมฺมา จิตฺต\-สหภุสฺส จ โนจิตฺต\-สหภุสฺส จ ธมฺมสฺส วิปฺปยุตฺต\-ปจฺจเยน ปจฺจโย. จิตฺตสหภู ขนฺธา จ จิตฺตญฺจ จิตฺต\-สหภูนญฺจ โนจิตฺต\-สหภูนญฺจ จิตฺต\-สมุฏฺฐานานํ รูปานํ วิปฺปยุตฺต\-ปจฺจเยน ปจฺจโย. (3)}

\proseitem{268}{จิตฺตสหภู ธมฺโม จิตฺต\-สหภุสฺส ธมฺมสฺส อตฺถิปจฺจเยน ปจฺจโย. จิตฺตสหภู เอโก ขนฺโธ ฯเปฯ. (ปฏิจฺจสทิสํ.) (1)}

\prose{จิตฺตสหภู ธมฺโม โนจิตฺต\-สหภุสฺส ธมฺมสฺส อตฺถิปจฺจเยน ปจฺจโย. สหชาตํ, ปจฺฉาชาตํ. (สํขิตฺตํ.) (2)}

\prose{จิตฺตสหภู ธมฺโม จิตฺต\-สหภุสฺส จ โนจิตฺต\-สหภุสฺส จ ธมฺมสฺส อตฺถิปจฺจเยน ปจฺจโย. จิตฺตสหภู เอโก ขนฺโธ ฯเปฯ. (ปฏิจฺจสทิสํ.) (3)}

\prose{โนจิตฺตสหภู ธมฺโม โนจิตฺต\-สหภุสฺส ธมฺมสฺส อตฺถิปจฺจเยน ปจฺจโย. สหชาตํ, ปุเรชาตํ, ปจฺฉาชาตํ, อาหารํ, อินฺทฺริยํ. (สํขิตฺตํ.) (1)}

\prose{โนจิตฺตสหภู ธมฺโม โนจิตฺต\-สหภุสฺส ธมฺมสฺส อตฺถิปจฺจเยน ปจฺจโย. สหชาตํ, ปุเรชาตํ. (สํขิตฺตํ.) (2)}

\prose{โนจิตฺตสหภู ธมฺโม จิตฺต\-สหภุสฺส จ โนจิตฺต\-สหภุสฺส จ ธมฺมสฺส อตฺถิปจฺจเยน ปจฺจโย. สหชาตํ, ปุเรชาตํ. (สํขิตฺตํ.) (3)}

\proseitem{269}{\csromanfolio{473}จิตฺตสหภู จ โนจิตฺตสหภู จ ธมฺมา จิตฺต\-สหภุสฺส ธมฺมสฺส อตฺถิปจฺจเยน ปจฺจโย. สหชาตํ, ปุเรชาตํ.  สหชาโต จกฺขุ\-วิญฺญาณ\-สหคโต เอโก ขนฺโธ จ จกฺขายตนญฺจ จกฺขุวิญฺญาณญฺจ ทฺวินฺนํ ขนฺธานํ ฯเปฯ. เทฺว ขนฺธา จ ฯเปฯ. (สพฺพํ ปฏิสนฺธิยํ กาตพฺพํ. สหชาตํ ปุเรชาตมฺปิ.) (1)}

\prose{จิตฺตสหภู จ โนจิตฺตสหภู จ ธมฺมา โนจิตฺต\-สหภุสฺส ธมฺมสฺส อตฺถิปจฺจเยน ปจฺจโย. สหชาตํ, ปุเรชาตํ, ปจฺฉาชาตํ, อาหารํ, อินฺทฺริยํ.  สหชาตา จกฺขุ\-วิญฺญาณสหคตา ขนฺธา จ จกฺขายตนญฺจ จกฺขุ\-วิญฺญาณสฺส อตฺถิปจฺจเยน ปจฺจโย ฯเปฯ. กาย\-วิญฺญาณสหคตา ฯเปฯ. จิตฺตสหภู ขนฺธา จ จิตฺตญฺจ โนจิตฺต\-สหภูนํ จิตฺต\-สมุฏฺฐานานํ รูปานํ อตฺถิปจฺจเยน ปจฺจโย. สหชาตา จิตฺตสหภู ขนฺธา จ วตฺถุ จ จิตฺตสฺส อตฺถิปจฺจเยน ปจฺจโย. สหชาตา จิตฺตสหภู ขนฺธา จ มหาภูตา จ โนจิตฺต\-สหภูนํ จิตฺต\-สมุฏฺฐานานํ รูปานํ อตฺถิปจฺจเยน ปจฺจโย. (ปฏิสนฺธิ\-กฺขเณ ตีณิปิ กาตพฺพา.) ปจฺฉาชาตา จิตฺตสหภู ขนฺธา จ จิตฺตญฺจ ปุเรชาตสฺส อิมสฺส โนจิตฺต\-สหภุสฺส กายสฺส อตฺถิปจฺจเยน ปจฺจโย. ปจฺฉาชาตา จิตฺตสหภู ขนฺธา จ กพฬีกาโร อาหาโร จ อิมสฺส โนจิตฺต\-สหภุสฺส กายสฺส อตฺถิปจฺจเยน ปจฺจโย. ปจฺฉาชาตา จิตฺตสหภู ขนฺธา จ รูปชีวิตินฺทฺริยญฺจ กฏตฺตารูปานํ อตฺถิปจฺจเยน ปจฺจโย. (2)}

\prose{จิตฺตสหภู จ โนจิตฺตสหภู จ ธมฺมา จิตฺต\-สหภุสฺส จ โนจิตฺต\-สหภุสฺส จ ธมฺมสฺส อตฺถิปจฺจเยน ปจฺจโย. สหชาตํ, ปุเรชาตํ.  สหชาโต จกฺขุ\-วิญฺญาณ\-สหคโต เอโก ขนฺโธ จ จกฺขายตนญฺจ ฯเปฯ. (ปจฺจยวาร\-สทิสํ.) (3)}

\csromanmarkh{1. ปจฺจยานุโลม}\csromanmarkhii{2. สงฺขฺยาวาร}\csromanheadernomark{2}{1. \textbf{ปจฺจยานุโลม 2. สงฺขฺยาวาร}}
\proseitem{270}{เหตุยา ตีณิ, อารมฺมเณ นว, อธิปติยา นว, อนนฺตเร นว, สมนนฺตเร นว, สหชาเต นว, อญฺญมญฺเญ นว, นิสฺสเย นว, อุปนิสฺสเย นว, ปุเรชาเต ตีณิ, ปจฺฉาชาเต ตีณิ, อาเสวเน นว, \csromanfolio{474}กมฺเม ตีณิ, วิปาเก นว, อาหาเร นว, อินฺทฺริเย นว, ฌาเน ตีณิ, มคฺเค ตีณิ, สมฺปยุตฺเต ปญฺจ, วิปฺปยุตฺเต นว, อตฺถิยา นว, นตฺถิยา นว, วิคเต นว, อวิคเต นว.}

\vspace{\tipitakaparskip}
\csromancenter{อนุโลมํ.}
\proseitem{271}{จิตฺตสหภู ธมฺโม จิตฺต\-สหภุสฺส ธมฺมสฺส อารมฺมณ\-ปจฺจเยน ปจฺจโย. สหชาต\-ปจฺจเยน ปจฺจโย. อุปนิสฺสย\-ปจฺจเยน ปจฺจโย. กมฺมปจฺจเยน ปจฺจโย. (1)}

\prose{จิตฺตสหภู ธมฺโม โนจิตฺต\-สหภุสฺส ธมฺมสฺส อารมฺมณ\-ปจฺจเยน ปจฺจโย. สหชาต\-ปจฺจเยน ปจฺจโย. อุปนิสฺสย\-ปจฺจเยน ปจฺจโย. ปจฺฉาชาต\-ปจฺจเยน ปจฺจโย. กมฺมปจฺจเยน ปจฺจโย. (2)}

\prose{จิตฺตสหภู ธมฺโม จิตฺต\-สหภุสฺส จ โนจิตฺต\-สหภุสฺส จ ธมฺมสฺส อารมฺมณ\-ปจฺจเยน ปจฺจโย. สหชาต\-ปจฺจเยน ปจฺจโย. อุปนิสฺสย\-ปจฺจเยน ปจฺจโย. กมฺมปจฺจเยน ปจฺจโย. (3)}

\proseitem{272}{โนจิตฺตสหภู ธมฺโม โนจิตฺต\-สหภุสฺส ธมฺมสฺส อารมฺมณ\-ปจฺจเยน ปจฺจโย. สหชาต\-ปจฺจเยน ปจฺจโย. อุปนิสฺสย\-ปจฺจเยน ปจฺจโย. ปุเรชาต\-ปจฺจเยน ปจฺจโย. ปจฺฉาชาต\-ปจฺจเยน ปจฺจโย. อาหารปจฺจเยน ปจฺจโย. อินฺทฺริย\-ปจฺจเยน ปจฺจโย. (1)}

\prose{โนจิตฺตสหภู ธมฺโม จิตฺต\-สหภุสฺส ธมฺมสฺส อารมฺมณ\-ปจฺจเยน ปจฺจโย. สหชาต\-ปจฺจเยน ปจฺจโย. อุปนิสฺสย\-ปจฺจเยน ปจฺจโย. ปุเรชาต\-ปจฺจเยน ปจฺจโย. (2)}

\prose{โนจิตฺตสหภู ธมฺโม จิตฺต\-สหภุสฺส จ โนจิตฺต\-สหภุสฺส จ ธมฺมสฺส อารมฺมณ\-ปจฺจเยน ปจฺจโย. สหชาต\-ปจฺจเยน ปจฺจโย. อุปนิสฺสย\-ปจฺจเยน ปจฺจโย. ปุเรชาต\-ปจฺจเยน ปจฺจโย. (3)}

\csromanmatikaanchor{mtk.61}
\csromantocmarkref{subsection}{62. จิตฺตานุปริวตฺติทุก}{mtk.61}
\bookmark[dest={mtk.61},level=2]{62. จิตฺตานุปริวตฺติทุก}
\csromanheader{2}{62. จิตฺตา\-นุปริวตฺติทุก}
\proseitem{273}{\csromanfolio{475}จิตฺตสหภู จ โนจิตฺตสหภู จ ธมฺมา จิตฺต\-สหภุสฺส ธมฺมสฺส อารมฺมณ\-ปจฺจเยน ปจฺจโย. สหชาต\-ปจฺจเยน ปจฺจโย. อุปนิสฺสย\-ปจฺจเยน ปจฺจโย. (1)}

\prose{จิตฺตสหภู จ โนจิตฺตสหภู จ ธมฺมา โนจิตฺต\-สหภุสฺส ธมฺมสฺส อารมฺมณ\-ปจฺจเยน ปจฺจโย. สหชาต\-ปจฺจเยน ปจฺจโย. อุปนิสฺสย\-ปจฺจเยน ปจฺจโย. ปจฺฉาชาต\-ปจฺจเยน ปจฺจโย. (2)}

\prose{จิตฺตสหภู จ โนจิตฺตสหภู จ ธมฺมา จิตฺต\-สหภุสฺส จ โนจิตฺต\-สหภุสฺส จ ธมฺมสฺส อารมฺมณ\-ปจฺจเยน ปจฺจโย. สหชาต\-ปจฺจเยน ปจฺจโย. อุปนิสฺสย\-ปจฺจเยน ปจฺจโย. (3)}

\csromanmarkh{2. ปจฺจย\-ปจฺจนีย}\csromanmarkhii{2. สงฺขฺยาวาร}\csromanheadernomark{2}{2. ปจฺจย\-ปจฺจนีย 2. สงฺขฺยาวาร}
\vspace{\tipitakaparskip}
\csromancenter{นเหตุยา นว, นอารมฺมเณ นว, (สพฺพตฺถ นว.) โนอวิคเต นว.}
\proseitem{275}{เหตุปจฺจยา นอารมฺมเณ ตีณิ, นอธิปติยา ตีณิ, นอนนฺตเร ตีณิ, นสมนนฺตเร ตีณิ, นอญฺญมญฺเญ ตีณิ, นอุปนิสฺสเย ตีณิ, (สพฺพตฺถ ตีณิ.) นสมฺปยุตฺเต ตีณิ, นวิปฺปยุตฺเต ตีณิ, โนนตฺถิยา ตีณิ, โนวิคเต ตีณิ.}

\csromanheader{2}{276. นเหตุปจฺจยา อารมฺมเณ นว, อธิปติยา นว. (อนุโลมมาติกา กาตพฺพา.)}
\nitthitam{จิตฺต\-สหภูทุกํ นิฏฺฐิตํ.}
\csromanmarkh{62. จิตฺตา\-นุปริวตฺติทุก}\csromanmarkhii{1. ปฏิจฺจวาร}\csromanheadernomark{2}{62. จิตฺตา\-นุปริวตฺติทุก 1. ปฏิจฺจวาร}
\proseitem{277}{จิตฺตา\-นุปริวตฺติํ ธมฺมํ ปฏิจฺจ จิตฺตา\-นุปริวตฺตี ธมฺโม อุปฺปชฺชติ เหตุปจฺจยา. จิตฺตา\-นุปริวตฺติํ เอกํ ขนฺธํ ปฏิจฺจ เทฺว ขนฺธา จิตฺตา\-นุปริวตฺติํ \csromanfolio{476}จิตฺต\-สมุฏฺฐานญฺจ รูปํ. เทฺว ขนฺเธ ฯเปฯ. ปฏิสนฺธิ\-กฺขเณ ฯเปฯ.}

\vspace{\tipitakaparskip}
\csromancenter{(ยถา จิตฺต\-สหภูทุกํ, เอวํ อิมํ ทุกํ กาตพฺพํ, นินฺนานากรณํ.)}
\nitthitam{จิตฺตา\-นุปริวตฺติทุกํ นิฏฺฐิตํ.}
\csromanmatikaanchor{mtk.62}
\csromantocmarkref{subsection}{63. จิตฺตสํสฏฺฐสมุฏฺฐานทุก}{mtk.62}
\bookmark[dest={mtk.62},level=2]{63. จิตฺตสํสฏฺฐสมุฏฺฐานทุก}
\csromanmarkh{63. จิตฺต\-สํสฏฺฐ\-สมุฏฺฐานทุก}\csromanmarkhii{1. ปฏิจฺจวาร}\csromanheadernomark{2}{63. จิตฺต\-สํสฏฺฐ\-สมุฏฺฐานทุก 1. ปฏิจฺจวาร}
\csromanmarkh{1. ปจฺจยานุโลม}\csromanmarkhii{1. วิภงฺควาร}\csromanheadernomark{2}{1. \textbf{ปจฺจยานุโลม 1. วิภงฺควาร}}
\csromanheader{2}{\textbf{เหตุ}}
\proseitem{278}{จิตฺต\-สํสฏฺฐ\-สมุฏฺฐานํ ธมฺมํ ปฏิจฺจ จิตฺต\-สํสฏฺฐ\-สมุฏฺฐาโน ธมฺโม อุปฺปชฺชติ เหตุปจฺจยา. จิตฺต\-สํสฏฺฐ\-สมุฏฺฐานํ เอกํ ขนฺธํ ปฏิจฺจ เทฺว ขนฺธา. เทฺว ขนฺเธ ปฏิจฺจ เอโก ขนฺโธ. ปฏิสนฺธิ\-กฺขเณ ฯเปฯ. (1)}

\prose{จิตฺต\-สํสฏฺฐ\-สมุฏฺฐานํ ธมฺมํ ปฏิจฺจ โนจิตฺต\-สํสฏฺฐ\-สมุฏฺฐาโน ธมฺโม อุปฺปชฺชติ เหตุปจฺจยา. จิตฺต\-สํสฏฺฐ\-สมุฏฺฐาเน ขนฺเธ ปฏิจฺจ จิตฺตํ จิตฺต\-สมุฏฺฐานญฺจ รูปํ. ปฏิสนฺธิ\-กฺขเณ จิตฺต\-สํสฏฺฐ\-สมุฏฺฐาเน ขนฺเธ ปฏิจฺจ จิตฺตํ กฏตฺตา จ รูปํ. (2)}

\prose{จิตฺต\-สํสฏฺฐ\-สมุฏฺฐานํ ธมฺมํ ปฏิจฺจ จิตฺต\-สํสฏฺฐ\-สมุฏฺฐาโน จ โนจิตฺต\-สํสฏฺฐ\-สมุฏฺฐาโน จ ธมฺมา อุปฺปชฺชนฺติ เหตุปจฺจยา. จิตฺต\-สํสฏฺฐ\-สมุฏฺฐานํ เอกํ ขนฺธํ ปฏิจฺจ เทฺว ขนฺธา จิตฺตญฺจ จิตฺต\-สมุฏฺฐานญฺจ รูปํ. เทฺว ขนฺเธ ฯเปฯ. ปฏิสนฺธิ\-กฺขเณ ฯเปฯ. (3)}

\proseitem{279}{โนจิตฺต\-สํสฏฺฐ\-สมุฏฺฐานํ ธมฺมํ ปฏิจฺจ โนจิตฺต\-สํสฏฺฐ\-สมุฏฺฐาโน ธมฺโม อุปฺปชฺชติ เหตุปจฺจยา. จิตฺตํ ปฏิจฺจ จิตฺต\-สมุฏฺฐานํ รูปํ. ปฏิสนฺธิ\-กฺขเณ จิตฺตํ ปฏิจฺจ กฏตฺตารูปํ. จิตฺตํ ปฏิจฺจ วตฺถุ, วตฺถุํ ปฏิจฺจ จิตฺตํ. เอกํ มหาภูตํ ฯเปฯ มหาภูเต ปฏิจฺจ จิตฺต\-สมุฏฺฐานํ รูปํ, กฏตฺตารูปํ, อุปาทารูปํ. (1)}

\prose{โนจิตฺต\-สํสฏฺฐ\-สมุฏฺฐานํ ธมฺมํ ปฏิจฺจ จิตฺต\-สํสฏฺฐ\-สมุฏฺฐาโน ธมฺโม อุปฺปชฺชติ เหตุปจฺจยา. จิตฺตํ ปฏิจฺจ สมฺปยุตฺตกา ขนฺธา. ปฏิสนฺธิ\-กฺขเณ จิตฺตํ ปฏิจฺจ สมฺปยุตฺตกา ขนฺธา. ปฏิสนฺธิ\-กฺขเณ วตฺถุํ ปฏิจฺจ จิตฺต\-สํสฏฺฐ\-สมุฏฺฐานา ขนฺธา. (2)}

\prose{\csromanfolio{477}โนจิตฺต\-สํสฏฺฐ\-สมุฏฺฐานํ ธมฺมํ ปฏิจฺจ จิตฺต\-สํสฏฺฐ\-สมุฏฺฐาโน จ โนจิตฺต\-สํสฏฺฐ\-สมุฏฺฐาโน จ ธมฺมา อุปฺปชฺชนฺติ เหตุปจฺจยา. จิตฺตํ ปฏิจฺจ สมฺปยุตฺตกา ขนฺธา จิตฺต\-สมุฏฺฐานญฺจ รูปํ. ปฏิสนฺธิ\-กฺขเณ จิตฺตํ ปฏิจฺจ สมฺปยุตฺตกา ขนฺธา กฏตฺตา จ รูปํ. ปฏิสนฺธิ\-กฺขเณ วตฺถุํ ปฏิจฺจ จิตฺตํ สมฺปยุตฺตกา จ ขนฺธา. (3)}

\proseitem{280}{จิตฺต\-สํสฏฺฐ\-สมุฏฺฐานญฺจ โนจิตฺต\-สํสฏฺฐ\-สมุฏฺฐานญฺจ ธมฺมํ ปฏิจฺจ จิตฺต\-สํสฏฺฐ\-สมุฏฺฐาโน ธมฺโม อุปฺปชฺชติ เหตุปจฺจยา. จิตฺต\-สํสฏฺฐ\-สมุฏฺฐานํ เอกํ ขนฺธญฺจ จิตฺตญฺจ ปฏิจฺจ เทฺว ขนฺธา. เทฺว ขนฺเธ ฯเปฯ. ปฏิสนฺธิ\-กฺขเณ จิตฺต\-สํสฏฺฐ\-สมุฏฺฐานํ เอกํ ขนฺธญฺจ จิตฺตญฺจ ปฏิจฺจ เทฺว ขนฺธา. เทฺว ขนฺเธ ฯเปฯ. ปฏิสนฺธิ\-กฺขเณ จิตฺต\-สํสฏฺฐ\-สมุฏฺฐานํ เอกํ ขนฺธญฺจ วตฺถุญฺจ ปฏิจฺจ เทฺว ขนฺธา. เทฺว ขนฺเธ ฯเปฯ. (1)}

\prose{จิตฺต\-สํสฏฺฐ\-สมุฏฺฐานญฺจ โนจิตฺต\-สํสฏฺฐ\-สมุฏฺฐานญฺจ ธมฺมํ ปฏิจฺจ โนจิตฺต\-สํสฏฺฐ\-สมุฏฺฐาโน ธมฺโม อุปฺปชฺชติ เหตุปจฺจยา. จิตฺต\-สํสฏฺฐ\-สมุฏฺฐาเน ขนฺเธ จ จิตฺตญฺจ ปฏิจฺจ จิตฺต\-สมุฏฺฐานํ รูปํ. จิตฺต\-สํสฏฺฐ\-สมุฏฺฐาเน ขนฺเธ จ มหาภูเต จ ปฏิจฺจ จิตฺต\-สมุฏฺฐานํ รูปํ. ปฏิสนฺธิ\-กฺขเณ จิตฺต\-สํสฏฺฐ\-สมุฏฺฐาเน ขนฺเธ จ จิตฺตญฺจ ปฏิจฺจ กฏตฺตารูปํ. ปฏิสนฺธิ\-กฺขเณ จิตฺต\-สํสฏฺฐ\-สมุฏฺฐาเน ขนฺเธ จ มหาภูเต จ ปฏิจฺจ กฏตฺตารูปํ. ปฏิสนฺธิ\-กฺขเณ จิตฺต\-สํสฏฺฐ\-สมุฏฺฐาเน ขนฺเธ จ วตฺถุญฺจ ปฏิจฺจ จิตฺตํ. (2)}

\prose{จิตฺต\-สํสฏฺฐ\-สมุฏฺฐานญฺจ โนจิตฺต\-สํสฏฺฐ\-สมุฏฺฐานญฺจ ธมฺมํ ปฏิจฺจ จิตฺต\-สํสฏฺฐ\-สมุฏฺฐาโน จ โนจิตฺต\-สํสฏฺฐ\-สมุฏฺฐาโน จ ธมฺมา อุปฺปชฺชนฺติ เหตุปจฺจยา. จิตฺต\-สํสฏฺฐ\-สมุฏฺฐานํ เอกํ ขนฺธญฺจ จิตฺตญฺจ ปฏิจฺจ เทฺว ขนฺธา จิตฺต\-สมุฏฺฐานญฺจ รูปํ. เทฺว ขนฺเธ ฯเปฯ. ปฏิสนฺธิ\-กฺขเณ จิตฺต\-สํสฏฺฐ\-สมุฏฺฐานํ เอกํ ขนฺธญฺจ จิตฺตญฺจ ปฏิจฺจ เทฺว ขนฺธา กฏตฺตา จ รูปํ. เทฺว ขนฺเธ ฯเปฯ. ปฏิสนฺธิ\-กฺขเณ จิตฺต\-สํสฏฺฐ\-สมุฏฺฐานํ เอกํ ขนฺธญฺจ วตฺถุญฺจ ปฏิจฺจ เทฺว ขนฺธา จิตฺตญฺจ. เทฺว ขนฺเธ ฯเปฯ. (สํขิตฺตํ.) (3)}

\csromanmarkh{1. ปจฺจยานุโลม}\csromanmarkhii{2. สงฺขฺยาวาร}\csromanheadernomark{2}{1. \textbf{ปจฺจยานุโลม 2. สงฺขฺยาวาร}}
\proseitem{281}{เหตุยา นว, อารมฺมเณ นว, อธิปติยา นว, อนนฺตเร นว, สมนนฺตเร นว, สหชาเต นว, อญฺญมญฺเญ นว, นิสฺสเย นว, อุปนิสฺสเย นว, ปุเรชาเต ปญฺจ, อาเสวเน ปญฺจ, กมฺเม นว, วิปาเก นว, (สพฺพตฺถ นว.) อวิคเต นว.}

\vspace{\tipitakaparskip}
\csromancenter{อนุโลมํ.}
\csromanmarkh{2. ปจฺจย\-ปจฺจนีย}\csromanmarkhii{1. วิภงฺควาร}\csromanheadernomark{2}{2. ปจฺจย\-ปจฺจนีย 1. วิภงฺควาร}
\proseitem{282}{\csromanfolio{478}จิตฺต\-สํสฏฺฐ\-สมุฏฺฐานํ ธมฺมํ ปฏิจฺจ จิตฺต\-สํสฏฺฐ\-สมุฏฺฐาโน ธมฺโม อุปฺปชฺชติ นเหตุปจฺจยา. อเหตุกํ จิตฺต\-สํสฏฺฐ\-สมุฏฺฐานํ เอกํ ขนฺธํ ปฏิจฺจ เทฺว ขนฺธา. เทฺว ขนฺเธ ฯเปฯ. อเหตุก\-ปฏิสนฺธิ\-กฺขเณ ฯเปฯ. วิจิกิจฺฉา\-สหคเต อุทฺธจฺจ\-สหคเต ขนฺเธ ปฏิจฺจ วิจิกิจฺฉา\-สหคโต อุทฺธจฺจ\-สหคโต โมโห. (1)}

\prose{จิตฺต\-สํสฏฺฐ\-สมุฏฺฐานํ ธมฺมํ ปฏิจฺจ โนจิตฺต\-สํสฏฺฐ\-สมุฏฺฐาโน ธมฺโม อุปฺปชฺชติ นเหตุปจฺจยา. อเหตุเก จิตฺต\-สํสฏฺฐ\-สมุฏฺฐาเน ขนฺเธ ปฏิจฺจ จิตฺตํ จิตฺต\-สมุฏฺฐานญฺจ รูปํ. อเหตุก\-ปฏิสนฺธิ\-กฺขเณ ฯเปฯ. (2)}

\prose{จิตฺต\-สํสฏฺฐ\-สมุฏฺฐานํ ธมฺมํ ปฏิจฺจ จิตฺต\-สํสฏฺฐ\-สมุฏฺฐาโน จ โนจิตฺต\-สํสฏฺฐ\-สมุฏฺฐาโน จ ธมฺมา อุปฺปชฺชนฺติ นเหตุปจฺจยา. อเหตุกํ จิตฺต\-สํสฏฺฐ\-สมุฏฺฐานํ เอกํ ขนฺธํ ปฏิจฺจ เทฺว ขนฺธา จิตฺตญฺจ จิตฺต\-สมุฏฺฐานญฺจ รูปํ. เทฺว ขนฺเธ ฯเปฯ. อเหตุก\-ปฏิสนฺธิ\-กฺขเณ ฯเปฯ. (3)}

\proseitem{283}{โนจิตฺต\-สํสฏฺฐ\-สมุฏฺฐานํ ธมฺมํ ปฏิจฺจ โนจิตฺต\-สํสฏฺฐ\-สมุฏฺฐาโน ธมฺโม อุปฺปชฺชติ นเหตุปจฺจยา. อเหตุกํ จิตฺตํ ปฏิจฺจ จิตฺต\-สมุฏฺฐานํ รูปํ. อเหตุก\-ปฏิสนฺธิ\-กฺขเณ จิตฺตํ ปฏิจฺจ กฏตฺตารูปํ. จิตฺตํ ปฏิจฺจ วตฺถุ, วตฺถุํ ปฏิจฺจ จิตฺตํ. เอกํ มหาภูตํ ฯเปฯ. (ยาว อสญฺญสตฺตา.) (1)}

\prose{โนจิตฺต\-สํสฏฺฐ\-สมุฏฺฐานํ ธมฺมํ ปฏิจฺจ จิตฺต\-สํสฏฺฐ\-สมุฏฺฐาโน ธมฺโม อุปฺปชฺชติ นเหตุปจฺจยา. อเหตุกํ จิตฺตํ ปฏิจฺจ สมฺปยุตฺตกา ขนฺธา. อเหตุก\-ปฏิสนฺธิ\-กฺขเณ จิตฺตํ ปฏิจฺจ สมฺปยุตฺตกา ขนฺธา. อเหตุก\-ปฏิสนฺธิ\-กฺขเณ วตฺถุํ ปฏิจฺจ จิตฺต\-สํสฏฺฐ\-สมุฏฺฐานา ขนฺธา. วิจิกิจฺฉา\-สหคตํ อุทฺธจฺจ\-สหคตํ จิตฺตํ ปฏิจฺจ วิจิกิจฺฉา\-สหคโต อุทฺธจฺจ\-สหคโต โมโห. (2)}

\prose{โนจิตฺต\-สํสฏฺฐ\-สมุฏฺฐานํ ธมฺมํ ปฏิจฺจ จิตฺต\-สํสฏฺฐ\-สมุฏฺฐาโน จ โนจิตฺต\-สํสฏฺฐ\-สมุฏฺฐาโน จ ธมฺมา อุปฺปชฺชนฺติ นเหตุปจฺจยา. อเหตุกํ จิตฺตํ ปฏิจฺจ สมฺปยุตฺตกา ขนฺธา จิตฺต\-สมุฏฺฐานญฺจ รูปํ. อเหตุก\-ปฏิสนฺธิ\-กฺขเณ จิตฺตํ ปฏิจฺจ สมฺปยุตฺตกา ขนฺธา กฏตฺตา จ รูปํ. อเหตุก\-ปฏิสนฺธิ\-กฺขเณ วตฺถุํ ปฏิจฺจ จิตฺตํ สมฺปยุตฺตกา จ ขนฺธา. (3)}

\proseitem{284}{\csromanfolio{479}จิตฺต\-สํสฏฺฐ\-สมุฏฺฐานญฺจ โนจิตฺต\-สํสฏฺฐ\-สมุฏฺฐานญฺจ ธมฺมํ ปฏิจฺจ จิตฺต\-สํสฏฺฐ\-สมุฏฺฐาโน ธมฺโม อุปฺปชฺชติ นเหตุปจฺจยา. อเหตุกํ จิตฺต\-สํสฏฺฐ\-สมุฏฺฐานํ เอกํ ขนฺธญฺจ จิตฺตญฺจ ปฏิจฺจ เทฺว ขนฺธา. เทฺว ขนฺเธ ฯเปฯ. อเหตุก\-ปฏิสนฺธิ\-กฺขเณ จิตฺต\-สํสฏฺฐ\-สมุฏฺฐานํ เอกํ ขนฺธญฺจ จิตฺตญฺจ ฯเปฯ. อเหตุก\-ปฏิสนฺธิ\-กฺขเณ จิตฺต\-สํสฏฺฐ\-สมุฏฺฐานํ เอกํ ขนฺธญฺจ วตฺถุญฺจ ปฏิจฺจ เทฺว ขนฺธา. เทฺว ขนฺเธ ฯเปฯ. วิจิกิจฺฉา\-สหคเต อุทฺธจฺจ\-สหคเต ขนฺเธ จ จิตฺตญฺจ ปฏิจฺจ วิจิกิจฺฉาสคหโต อุทฺธจฺจ\-สหคโต โมโห. (1)}

\prose{จิตฺต\-สํสฏฺฐ\-สมุฏฺฐานญฺจ โนจิตฺต\-สํสฏฺฐ\-สมุฏฺฐานญฺจ ธมฺมํ ปฏิจฺจ โนจิตฺต\-สํสฏฺฐ\-สมุฏฺฐาโน ธมฺโม อุปฺปชฺชติ นเหตุปจฺจยา. อเหตุเก จิตฺต\-สํสฏฺฐ\-สมุฏฺฐาเน ขนฺเธ จ จิตฺตญฺจ ปฏิจฺจ จิตฺต\-สมุฏฺฐานํ รูปํ. อเหตุเก จิตฺต\-สํสฏฺฐ\-สมุฏฺฐาเน ขนฺเธ จ มหาภูเต จ ปฏิจฺจ จิตฺต\-สมุฏฺฐานํ รูปํ. อเหตุก\-ปฏิสนฺธิ\-กฺขเณ จิตฺต\-สํสฏฺฐ\-สมุฏฺฐาเน ขนฺเธ จ จิตฺตญฺจ ปฏิจฺจ กฏตฺตารูปํ. อเหตุก\-ปฏิสนฺธิ\-กฺขเณ จิตฺต\-สํสฏฺฐ\-สมุฏฺฐาเน ขนฺเธ จ มหาภูเต จ ปฏิจฺจ กฏตฺตารูปํ. อเหตุก\-ปฏิสนฺธิ\-กฺขเณ จิตฺต\-สํสฏฺฐ\-สมุฏฺฐาเน ขนฺเธ จ วตฺถุญฺจ ปฏิจฺจ จิตฺตํ. (2)}

\prose{จิตฺต\-สํสฏฺฐ\-สมุฏฺฐานญฺจ โนจิตฺต\-สํสฏฺฐ\-สมุฏฺฐานญฺจ ธมฺมํ ปฏิจฺจ จิตฺต\-สํสฏฺฐ\-สมุฏฺฐาโน จ โนจิตฺต\-สํสฏฺฐ\-สมุฏฺฐาโน จ ธมฺมา อุปฺปชฺชนฺติ นเหตุปจฺจยา. อเหตุกํ จิตฺต\-สํสฏฺฐ\-สมุฏฺฐานํ เอกํ ขนฺธญฺจ จิตฺตญฺจ ปฏิจฺจ เทฺว ขนฺธา จิตฺต\-สมุฏฺฐานญฺจ รูปํ. เทฺว ขนฺเธ ฯเปฯ. (อเหตุก\-ปฏิสนฺธิ\-กฺขเณ เทฺวปิ กาตพฺพา.) (3)}

\prose{นอารมฺมณ}

\proseitem{285}{จิตฺต\-สํสฏฺฐ\-สมุฏฺฐานํ ธมฺมํ ปฏิจฺจ โนจิตฺต\-สํสฏฺฐ\-สมุฏฺฐาโน ธมฺโม อุปฺปชฺชติ นอารมฺมณ\-ปจฺจยา. จิตฺต\-สํสฏฺฐ\-สมุฏฺฐาเน ขนฺเธ ปฏิจฺจ จิตฺต\-สมุฏฺฐานํ รูปํ. ปฏิสนฺธิ\-กฺขเณ ฯเปฯ. (1)}

\prose{โนจิตฺต\-สํสฏฺฐ\-สมุฏฺฐานํ ธมฺมํ ปฏิจฺจ โนจิตฺต\-สํสฏฺฐ\-สมุฏฺฐาโน ธมฺโม อุปฺปชฺชติ นอารมฺมณ\-ปจฺจยา. จิตฺตํ ปฏิจฺจ จิตฺต\-สมุฏฺฐานํ รูปํ. ปฏิสนฺธิ\-กฺขเณ จิตฺตํ ปฏิจฺจ กฏตฺตารูปํ. จิตฺตํ ปฏิจฺจ วตฺถุ. เอกํ มหาภูตํ ฯเปฯ. (ยาว อสญฺญสตฺตา.) (1)}

\prose{จิตฺต\-สํสฏฺฐ\-สมุฏฺฐานญฺจ โนจิตฺต\-สํสฏฺฐ\-สมุฏฺฐานญฺจ ธมฺมํ ปฏิจฺจ โนจิตฺต\-สํสฏฺฐ\-สมุฏฺฐาโน ธมฺโม อุปฺปชฺชติ นอารมฺมณ\-ปจฺจยา. จิตฺต\-สํสฏฺฐ\-สมุฏฺฐาเน \csromanfolio{480}ขนฺเธ จ จิตฺตญฺจ ปฏิจฺจ จิตฺต\-สมุฏฺฐานํ รูปํ. จิตฺต\-สํสฏฺฐ\-สมุฏฺฐาเน ขนฺเธ จ มหาภูเต จ ปฏิจฺจ จิตฺต\-สมุฏฺฐานํ รูปํ. (ปฏิสนฺธิ\-กฺขเณ เทฺว. สํขิตฺตํ.) (1)}

\csromanmarkh{2. ปจฺจย\-ปจฺจนีย}\csromanmarkhii{2. สงฺขฺยาวาร}\csromanheadernomark{2}{2. ปจฺจย\-ปจฺจนีย 2. สงฺขฺยาวาร}
\proseitem{286}{นเหตุยา นว, นอารมฺมเณ ตีณิ, นอธิปติยา นว, นอนนฺตเร ตีณิ, นสมนนฺตเร ตีณิ, นอญฺญมญฺเญ ตีณิ, นอุปนิสฺสเย ตีณิ, นปุเรชาเต นว, นปจฺฉาชาเต นว, นอาเสวเน นว, นกมฺเม จตฺตาริ, นวิปาเก นว, นอาหาเร เอกํ, นอินฺทฺริเย เอกํ, นฌาเน ฉ, นมคฺเค นว, นสมฺปยุตฺเต ตีณิ, นวิปฺปยุตฺเต ฉ, โนนตฺถิยา ตีณิ, โนวิคเต ตีณิ.}

\vspace{\tipitakaparskip}
\csromancenter{287. \textbf{เหตุ}ปจฺจยา นอารมฺมเณ ตีณิ. (สํขิตฺตํ.)}
\proseitem{288}{นเหตุปจฺจยา อารมฺมเณ นว, อนนฺตเร นว. (สํขิตฺตํ.)}

\vspace{\tipitakaparskip}
\csromancenter{(สหชาตวาโร ปฏิจฺจ\-วาร\-สทิโส.)}
\csromanmarkh{63. จิตฺต\-สํสฏฺฐ\-สมุฏฺฐานทุก}\csromanmarkhii{3. ปจฺจยวาร}\csromanheadernomark{2}{63. จิตฺต\-สํสฏฺฐ\-สมุฏฺฐานทุก 3. ปจฺจยวาร}
\csromanmarkh{1. ปจฺจยานุโลม}\csromanmarkhii{1. วิภงฺควาร}\csromanheadernomark{2}{1. \textbf{ปจฺจยานุโลม 1. วิภงฺควาร}}
\vspace{\tipitakaparskip}
\csromancenter{\textbf{เหตุ}}
\proseitem{289}{จิตฺต\-สํสฏฺฐ\-สมุฏฺฐานํ ธมฺมํ ปจฺจยา จิตฺต\-สํสฏฺฐ\-สมุฏฺฐาโน ธมฺโม อุปฺปชฺชติ เหตุปจฺจยา. (สํขิตฺตํ.) ตีณิ. (ปฏิจฺจ\-วาร\-สทิสา.)}

\prose{\csromanfolio{481}โนจิตฺต\-สํสฏฺฐ\-สมุฏฺฐานํ ธมฺมํ ปจฺจยา โนจิตฺต\-สํสฏฺฐ\-สมุฏฺฐาโน ธมฺโม อุปฺปชฺชติ เหตุปจฺจยา. จิตฺตํ ปจฺจยา จิตฺต\-สมุฏฺฐานํ รูปํ. วตฺถุํ ปจฺจยา จิตฺตํ. ปฏิสนฺธิ\-กฺขเณ ฯเปฯ. (ยาว มหาภูตา.) (1)}

\prose{โนจิตฺต\-สํสฏฺฐ\-สมุฏฺฐานํ ธมฺมํ ปจฺจยา จิตฺต\-สํสฏฺฐ\-สมุฏฺฐาโน ธมฺโม อุปฺปชฺชติ เหตุปจฺจยา. จิตฺตํ ปจฺจยา สมฺปยุตฺตกา ขนฺธา. วตฺถุํ ปจฺจยา จิตฺต\-สํสฏฺฐ\-สมุฏฺฐานา ขนฺธา. (ปฏิสนฺธิ\-กฺขเณ เทฺวปิ กาตพฺพา.) ตีณิ. (3)}

\proseitem{290}{จิตฺต\-สํสฏฺฐ\-สมุฏฺฐานญฺจ โนจิตฺต\-สํสฏฺฐ\-สมุฏฺฐานญฺจ ธมฺมํ ปจฺจยา จิตฺต\-สํสฏฺฐ\-สมุฏฺฐาโน ธมฺโม อุปฺปชฺชติ เหตุปจฺจยา. จิตฺต\-สํสฏฺฐ\-สมุฏฺฐานํ เอกํ ขนฺธญฺจ จิตฺตญฺจ ปจฺจยา เทฺว ขนฺธา. เทฺว ขนฺเธ จ ฯเปฯ. จิตฺต\-สํสฏฺฐ\-สมุฏฺฐานํ เอกํ ขนฺธญฺจ วตฺถุญฺจ ปจฺจยา เทฺว ขนฺธา. เทฺว ขนฺเธ จ ฯเปฯ. (ปฏิสนฺธิ\-กฺขเณ เทฺวปิ กาตพฺพา.) (1)}

\prose{จิตฺต\-สํสฏฺฐ\-สมุฏฺฐานญฺจ โนจิตฺต\-สํสฏฺฐ\-สมุฏฺฐานญฺจ ธมฺมํ ปจฺจยา โนจิตฺต\-สํสฏฺฐ\-สมุฏฺฐาโน ธมฺโม อุปฺปชฺชติ เหตุปจฺจยา. จิตฺต\-สํสฏฺฐ\-สมุฏฺฐาเน ขนฺเธ จ จิตฺตญฺจ ปจฺจยา โนจิตฺต\-สํสฏฺฐ\-สมุฏฺฐานํ รูปํ. จิตฺต\-สํสฏฺฐ\-สมุฏฺฐาเน ขนฺเธ จ มหาภูเต จ ปจฺจยา โนจิตฺต\-สํสฏฺฐ\-สมุฏฺฐานํ รูปํ. จิตฺต\-สํสฏฺฐ\-สมุฏฺฐาเน ขนฺเธ จ วตฺถุญฺจ ปจฺจยา จิตฺตํ. (ปฏิสนฺธิ\-กฺขเณ ตีณิปิ กาตพฺพา.) (2)}

\prose{จิตฺต\-สํสฏฺฐ\-สมุฏฺฐานญฺจ โนจิตฺต\-สํสฏฺฐ\-สมุฏฺฐานญฺจ ธมฺมํ ปจฺจยา จิตฺต\-สํสฏฺฐ\-สมุฏฺฐาโน จ โนจิตฺต\-สํสฏฺฐ\-สมุฏฺฐาโน จ ธมฺมา อุปฺปชฺชนฺติ เหตุปจฺจยา. จิตฺต\-สํสฏฺฐ\-สมุฏฺฐานํ เอกํ ขนฺธญฺจ จิตฺตญฺจ ปจฺจยา เทฺว ขนฺธา จิตฺต\-สํสฏฺฐ\-สมุฏฺฐานญฺจ รูปํ. เทฺว ขนฺเธ ฯเปฯ. จิตฺต\-สํสฏฺฐ\-สมุฏฺฐานํ เอกํ ขนฺธญฺจ วตฺถุญฺจ ปจฺจยา เทฺว ขนฺธา จิตฺตญฺจ. เทฺว ขนฺเธ ฯเปฯ. (ปฏิสนฺธิ\-กฺขเณ เทฺวปิ กาตพฺพา. (3)}

\csromanheader{2}{\textbf{อรมฺมณ}}
\proseitem{291}{จิตฺต\-สํสฏฺฐ\-สมุฏฺฐานํ ธมฺมํ ปจฺจยา จิตฺต\-สํสฏฺฐ\-สมุฏฺฐาโน ธมฺโม อุปฺปชฺชติ อารมฺมณ\-ปจฺจยา. ตีณิ. (ปฏิจฺจสทิสา.)}

\prose{โนจิตฺต\-สํสฏฺฐ\-สมุฏฺฐานํ ธมฺมํ ปจฺจยา โนจิตฺต\-สํสฏฺฐ\-สมุฏฺฐาโน ธมฺโม อุปฺปชฺชติ อารมฺมณ\-ปจฺจยา. จกฺขายตนํ ปจฺจยา จกฺขุวิญฺญาณํ ฯเปฯ. กายายตนํ ฯเปฯ. วตฺถุํ ปจฺจยา จิตฺตํ. ปฏิสนฺธิ\-กฺขเณ ฯเปฯ. (1)}

\prose{\csromanfolio{482}โนจิตฺต\-สํสฏฺฐ\-สมุฏฺฐานํ ธมฺมํ ปจฺจยา จิตฺต\-สํสฏฺฐ\-สมุฏฺฐาโน ธมฺโม อุปฺปชฺชติ อารมฺมณ\-ปจฺจยา. จกฺขายตนํ ปจฺจยา จกฺขุ\-วิญฺญาณสหคตา ขนฺธา ฯเปฯ. กายายตนํ ฯเปฯ. จิตฺตํ ปจฺจยา สมฺปยุตฺตกา ขนฺธา. วตฺถุํ ปจฺจยา จิตฺต\-สํสฏฺฐ\-สมุฏฺฐานา ขนฺธา. (ปฏิสนฺธิ\-กฺขเณ เทฺวปิ กาตพฺพา.) (2)}

\prose{โนจิตฺต\-สํสฏฺฐ\-สมุฏฺฐานํ ธมฺมํ ปจฺจยา จิตฺต\-สํสฏฺฐ\-สมุฏฺฐาโน จ โนจิตฺต\-สํสฏฺฐ\-สมุฏฺฐาโน จ ธมฺมา อุปฺปชฺชนฺติ อารมฺมณ\-ปจฺจยา. จกฺขายตนํ ปจฺจยา จกฺขุวิญฺญาณํ สมฺปยุตฺตกา จ ขนฺธา ฯเปฯ. กายายตนํ ฯเปฯ. วตฺถุํ ปจฺจยา จิตฺตํ สมฺปยุตฺตกา จ ขนฺธา. (ปฏิสนฺธิ\-กฺขเณ เอกํ กาตพฺพํ.) (3)}

\proseitem{292}{จิตฺต\-สํสฏฺฐ\-สมุฏฺฐานญฺจ โนจิตฺต\-สํสฏฺฐ\-สมุฏฺฐานญฺจ ธมฺมํ ปจฺจยา จิตฺต\-สํสฏฺฐ\-สมุฏฺฐาโน ธมฺโม อุปฺปชฺชติ อารมฺมณ\-ปจฺจยา. จกฺขุ\-วิญฺญาณสหคตํ เอกํ ขนฺธญฺจ จกฺขายตนญฺจ จกฺขุวิญฺญาณญฺจ ปจฺจยา เทฺว ขนฺธา. เทฺว ขนฺเธ ฯเปฯ. กาย\-วิญฺญาณสหคตํ ฯเปฯ. จิตฺต\-สํสฏฺฐ\-สมุฏฺฐานํ เอกํ ขนฺธญฺจ จิตฺตญฺจ ปจฺจยา เทฺว ขนฺธา. เทฺว ขนฺเธ ฯเปฯ. จิตฺต\-สํสฏฺฐ\-สมุฏฺฐานํ เอกํ ขนฺธญฺจ วตฺถุญฺจ ปจฺจยา เทฺว ขนฺธา. เทฺว ขนฺเธ ฯเปฯ. (ปฏิสนฺธิ\-กฺขเณ เทฺว กาตพฺพา.) (1)}

\prose{จิตฺต\-สํสฏฺฐ\-สมุฏฺฐานญฺจ โนจิตฺต\-สํสฏฺฐ\-สมุฏฺฐานญฺจ ธมฺมํ ปจฺจยา โนจิตฺต\-สํสฏฺฐ\-สมุฏฺฐาโน ธมฺโม อุปฺปชฺชติ อารมฺมณ\-ปจฺจยา. จกฺขุ\-วิญฺญาณสหคเต ขนฺเธ จ จกฺขายตนญฺจ ปจฺจยา จกฺขุวิญฺญาณํ ฯเปฯ. กาย\-วิญฺญาณสหคเต ฯเปฯ. จิตฺต\-สํสฏฺฐ\-สมุฏฺฐาเน ขนฺเธ จ วตฺถุญฺจ ปจฺจยา จิตฺตํ. (ปฏิสนฺธิ\-กฺขเณ เอกํ กาตพฺพํ.)}

\prose{จิตฺต\-สํสฏฺฐ\-สมุฏฺฐานญฺจ โนจิตฺต\-สํสฏฺฐ\-สมุฏฺฐานญฺจ ธมฺมํ ปจฺจยา จิตฺต\-สํสฏฺฐ\-สมุฏฺฐาโน จ โนจิตฺต\-สํสฏฺฐ\-สมุฏฺฐาโน จ ธมฺมา อุปฺปชฺชนฺติ อารมฺมณ\-ปจฺจยา. จิตฺต\-สํสฏฺฐ\-สมุฏฺฐานํ เอกํ ขนฺธญฺจ วตฺถุญฺจ ปจฺจยา เทฺว ขนฺธา จิตฺตญฺจ. เทฺว ขนฺเธ ฯเปฯ. (ปฏิสนฺธิ\-กฺขเณ เอกํ กาตพฺพํ. สํขิตฺตํ.) (3)}

\csromanmarkh{1. ปจฺจยานุโลม}\csromanmarkhii{2. สงฺขฺยาวาร}\csromanheadernomark{2}{1. \textbf{ปจฺจยานุโลม 2. สงฺขฺยาวาร}}
\vspace{\tipitakaparskip}
\csromancenter{เหตุยา นว, อารมฺมเณ นว, อธิปติยา นว, (สพฺพตฺถ นว.) อวิคเต นว.}
\csromancenter{2. ปจฺจย\-ปจฺจนีย 1. วิภงฺควาร}
\proseitem{294}{\csromanfolio{483}จิตฺต\-สํสฏฺฐ\-สมุฏฺฐานํ ธมฺมํ ปจฺจยา จิตฺต\-สํสฏฺฐ\-สมุฏฺฐาโน ธมฺโม อุปฺปชฺชติ นเหตุปจฺจยา. (เอวํ นว ปญฺหา กาตพฺพา. ปจฺจยวาเร ปญฺจวิญฺญาณมฺปิ กาตพฺพํ. ตีณิเยว, โมโห.)}

\csromanmarkh{2. ปจฺจย\-ปจฺจนีย}\csromanmarkhii{2. สงฺขฺยาวาร}\csromanheadernomark{2}{2. ปจฺจย\-ปจฺจนีย 2. สงฺขฺยาวาร}
\proseitem{295}{นเหตุยา นว, นอารมฺมเณ ตีณิ, นอธิปติยา นว, นอนนฺตเร ตีณิ, นสมนนฺตเร ตีณิ, นอญฺญมญฺเญ ตีณิ, นอุปนิสฺสเย ตีณิ, นปุเรชาเต นว, นปจฺฉาชาเต นว, นอาเสวเน นว, นกมฺเม จตฺตาริ, นวิปาเก นว, นอาหาเร เอกํ, นอินฺทฺริเย เอกํ, นฌาเน นว, นมคฺเค นว, นสมฺปยุตฺเต ตีณิ, นวิปฺปยุตฺเต ฉ, โนนตฺถิยา ตีณิ, โนวิคเต ตีณิ.}

\vspace{\tipitakaparskip}
\csromancenter{(เอวํ อิตเร เทฺว คณนาปิ นิสฺสยวาโรปิ กาตพฺโพ.)}
\csromanmarkh{63. จิตฺต\-สํสฏฺฐ\-สมุฏฺฐานทุก}\csromanmarkhii{5. สํสฏฺฐวาร}\csromanheadernomark{2}{63. จิตฺต\-สํสฏฺฐ\-สมุฏฺฐานทุก 5. สํสฏฺฐวาร}
\csromanmarkh{1. ปจฺจยานุโลม}\csromanmarkhii{1. วิภงฺควาร}\csromanheadernomark{2}{1. \textbf{ปจฺจยานุโลม 1. วิภงฺควาร}}
\csromanheader{2}{\textbf{เหตุ}}
\proseitem{296}{จิตฺต\-สํสฏฺฐ\-สมุฏฺฐานํ ธมฺมํ สํสฏฺโฐ จิตฺต\-สํสฏฺฐ\-สมุฏฺฐาโน ธมฺโม อุปฺปชฺชติ เหตุปจฺจยา. จิตฺต\-สํสฏฺฐ\-สมุฏฺฐานํ เอกํ ขนฺธํ สํสฏฺฐา เทฺว ขนฺธา. เทฺว ขนฺเธ ฯเปฯ. ปฏิสนฺธิ\-กฺขเณ ฯเปฯ. (1)}

\prose{จิตฺต\-สํสฏฺฐ\-สมุฏฺฐานํ ธมฺมํ สํสฏฺโฐ โนจิตฺต\-สํสฏฺฐ\-สมุฏฺฐาโน ธมฺโม อุปฺปชฺชติ เหตุปจฺจยา. จิตฺต\-สํสฏฺฐ\-สมุฏฺฐาเน ขนฺเธ สํสฏฺฐํ จิตฺตํ. ปฏิสนฺธิ\-กฺขเณ ฯเปฯ. (2)}

\prose{จิตฺต\-สํสฏฺฐ\-สมุฏฺฐานํ ธมฺมํ สํสฏฺโฐ จิตฺต\-สํสฏฺฐ\-สมุฏฺฐาโน จ โนจิตฺต\-สํสฏฺฐ\-สมุฏฺฐาโน จ ธมฺมา อุปฺปชฺชนฺติ เหตุปจฺจยา. จิตฺต\-สํสฏฺฐ\-สมุฏฺฐานํ เอกํ ขนฺธํ สํสฏฺฐา เทฺว ขนฺธา จิตฺตญฺจ. เทฺว ขนฺเธ ฯเปฯ. ปฏิสนฺธิ\-กฺขเณ ฯเปฯ. (3)}

\prose{\csromanfolio{484}โนจิตฺต\-สํสฏฺฐ\-สมุฏฺฐานํ ธมฺมํ สํสฏฺโฐ จิตฺต\-สํสฏฺฐ\-สมุฏฺฐาโน ธมฺโม อุปฺปชฺชติ เหตุปจฺจยา. จิตฺตํ สํสฏฺฐา สมฺปยุตฺตกา ขนฺธา. ปฏิสนฺธิ\-กฺขเณ ฯเปฯ. (1)}

\prose{จิตฺต\-สํสฏฺฐ\-สมุฏฺฐานญฺจ โนจิตฺต\-สํสฏฺฐ\-สมุฏฺฐานญฺจ ธมฺมํ สํสฏฺโฐ จิตฺต\-สํสฏฺฐ\-สมุฏฺฐาโน ธมฺโม อุปฺปชฺชติ เหตุปจฺจยา. จิตฺต\-สํสฏฺฐ\-สมุฏฺฐานํ เอกํ ขนฺธญฺจ จิตฺตญฺจ สํสฏฺฐา เทฺว ขนฺธา. เทฺว ขนฺเธ ฯเปฯ. ปฏิสนฺธิ\-กฺขเณ ฯเปฯ. (1) (สํขิตฺตํ.)}

\csromanmarkh{1. ปจฺจยานุโลม}\csromanmarkhii{2. สงฺขฺยาวาร}\csromanheadernomark{2}{1. \textbf{ปจฺจยานุโลม 2. สงฺขฺยาวาร}}
\proseitem{297}{เหตุยา ปญฺจ, อารมฺมเณ ปญฺจ, อธิปติยา ปญฺจ, (สพฺพตฺถ ปญฺจ,) อวิคเต ปญฺจ.}

\csromanmarkh{2. ปจฺจย\-ปจฺจนีย}\csromanmarkhii{1. วิภงฺควาร}\csromanheadernomark{2}{2. ปจฺจย\-ปจฺจนีย 1. วิภงฺควาร}
\proseitem{298}{จิตฺต\-สํสฏฺฐ\-สมุฏฺฐานํ ธมฺมํ สํสฏฺโฐ จิตฺต\-สํสฏฺฐ\-สมุฏฺฐาโน ธมฺโม อุปฺปชฺชติ นเหตุปจฺจยา. (สํขิตฺตํ. ตีณิเยว, โมโห.)}

\csromanmarkh{2. ปจฺจย\-ปจฺจนีย}\csromanmarkhii{1. สงฺขฺยาวาร}\csromanheadernomark{2}{2. ปจฺจย\-ปจฺจนีย 1. สงฺขฺยาวาร}
\proseitem{299}{นเหตุยา ปญฺจ, นอธิปติยา ปญฺจ, นปุเรชาเต ปญฺจ, นปจฺฉาชาเต ปญฺจ, นอาเสวเน ปญฺจ, นกมฺเม ตีณิ, นวิปาเก ปญฺจ, นฌาเน ปญฺจ, นมคฺเค ปญฺจ, นวิปฺปยุตฺเต ปญฺจ.}

\vspace{\tipitakaparskip}
\csromancenter{(เอวํ อิตเร เทฺว คณนาปิ สมฺปยุตฺตวาโรปิ กาตพฺโพ.)}
\csromanmarkh{63. จิตฺต\-สํสฏฺฐ\-สมุฏฺฐานทุก}\csromanmarkhii{7. ปญฺหาวาร}\csromanheadernomark{2}{63. จิตฺต\-สํสฏฺฐ\-สมุฏฺฐานทุก 7. ปญฺหาวาร}
\csromanmarkh{1. ปจฺจยานุโลม}\csromanmarkhii{1. วิภงฺควาร}\csromanheadernomark{2}{1. \textbf{ปจฺจยานุโลม 1. วิภงฺควาร}}
\vspace{\tipitakaparskip}
\csromancenter{\textbf{เหตุ}}
\proseitem{300}{จิตฺต\-สํสฏฺฐ\-สมุฏฺฐาโน ธมฺโม จิตฺต\-สํสฏฺฐ\-สมุฏฺฐานสฺส ธมฺมสฺส เหตุปจฺจเยน ปจฺจโย. จิตฺต\-สํสฏฺฐ\-สมุฏฺฐานา เหตู สมฺปยุตฺตกานํ \csromanfolio{485}ขนฺธานํ เหตุปจฺจเยน ปจฺจโย. ปฏิสนฺธิ\-กฺขเณ ฯเปฯ. (มูลํ กาตพฺพํ.) จิตฺต\-สํสฏฺฐ\-สมุฏฺฐานา เหตู จิตฺตสฺส จิตฺตสมุฏฺฐานานญฺจ รูปานํ เหตุปจฺจเยน ปจฺจโย. ปฏิสนฺธิ\-กฺขเณ ฯเปฯ. (มูลํ กาตพฺพํ.) จิตฺต\-สํสฏฺฐ\-สมุฏฺฐานา เหตู สมฺปยุตฺตกานํ ขนฺธานํ จิตฺตสฺส จ จิตฺตสมุฏฺฐานานญฺจ รูปานํ เหตุปจฺจเยน ปจฺจโย. ปฏิสนฺธิ\-กฺขเณ ฯเปฯ. (3)}

\proseitem{301}{จิตฺต\-สํสฏฺฐ\-สมุฏฺฐาโน ธมฺโม จิตฺต\-สํสฏฺฐ\-สมุฏฺฐานสฺส ธมฺมสฺส อารมฺมณ\-ปจฺจเยน ปจฺจโย. จิตฺต\-สํสฏฺฐ\-สมุฏฺฐาเน ขนฺเธ อารพฺภ จิตฺต\-สํสฏฺฐ\-สมุฏฺฐานา ขนฺธา อุปฺปชฺชนฺติ. (มูลํ กาตพฺพํ.) จิตฺต\-สํสฏฺฐ\-สมุฏฺฐาเน ขนฺเธ อารพฺภ จิตฺตํ อุปฺปชฺชติ. (มูลํ กาตพฺพํ.) จิตฺต\-สํสฏฺฐ\-สมุฏฺฐาเน ขนฺเธ อารพฺภ จิตฺตญฺจ สมฺปยุตฺตกา จ ขนฺธา อุปฺปชฺชนฺติ. (3)}

\prose{โนจิตฺต\-สํสฏฺฐา\-สมุฏฺฐาโน ธมฺโม โนจิตฺต\-สํสฏฺฐ\-สมุฏฺฐานสฺส ธมฺมสฺส อารมฺมณ\-ปจฺจเยน ปจฺจโย. อริยา มคฺคา วุฏฺฐหิตฺวา มคฺคํ ปจฺจเวกฺขนฺติ ฯเปฯ. นิพฺพานํ ปจฺจเวกฺขนฺติ. นิพฺพานํ โคตฺรภุสฺส ฯเปฯ. (สํขิตฺตํ. ยถา จิตฺต\-สหภู\-ทุเก อารมฺมณํ, เอวํ กาตพฺพํ. นินฺนานากรณํ. นวปิ ปญฺหา.)}

\proseitem{302}{จิตฺต\-สํสฏฺฐ\-สมุฏฺฐาโน ธมฺโม จิตฺต\-สํสฏฺฐ\-สมุฏฺฐานสฺส ธมฺมสฺส อธิปติ\-ปจฺจเยน ปจฺจโย. อารมฺมณา\-ธิปติ, สหชาตา\-ธิปติ. ตีณิ. (เทฺวปิ อธิปตี กาตพฺพา.) (3)}

\prose{โนจิตฺต\-สํสฏฺฐ\-สมุฏฺฐาโน ธมฺโม โนจิตฺต\-สํสฏฺฐ\-สมุฏฺฐานสฺส ธมฺมสฺส อธิปติ\-ปจฺจเยน ปจฺจโย. อารมฺมณา\-ธิปติ, สหชาตา\-ธิปติ. ตีณิ. (เทฺวปิ อธิปตี กาตพฺพา.) (3)}

\prose{จิตฺต\-สํสฏฺฐ\-สมุฏฺฐาโน จ โนจิตฺต\-สํสฏฺฐ\-สมุฏฺฐาโน จ ธมฺมา จิตฺต\-สํสฏฺฐ\-สมุฏฺฐานสฺส ธมฺมสฺส อธิปติ\-ปจฺจเยน ปจฺจโย. อารมฺมณา\-ธิปติ. (เอกาเยว อธิปติ กาตพฺพา, นวปิ ปญฺหา. ยถา จิตฺต\-สหภูทุกํ, เอวํ กาตพฺพํ. นินฺนานากรณํ.)}

\proseitem{303}{\csromanfolio{486}จิตฺต\-สํสฏฺฐ\-สมุฏฺฐาโน ธมฺโม จิตฺต\-สํสฏฺฐ\-สมุฏฺฐานสฺส ธมฺมสฺส อนนฺตร\-ปจฺจเยน ปจฺจโย. (นวปิ ปญฺหา. จิตฺต\-สหภู\-ทุก\-สทิสา.) สมนนฺตร\-ปจฺจเยน ปจฺจโย. นว. สหชาต\-ปจฺจเยน ปจฺจโย. นว. (ปฏิจฺจสทิสา.) อญฺญมญฺญ\-ปจฺจเยน ปจฺจโย. นว. (ปฏิจฺจสทิสา.) นิสฺสย\-ปจฺจเยน ปจฺจโย. นว. (ปจฺจยสทิสา.) อุปนิสฺสย\-ปจฺจเยน ปจฺจโย. (นวปิ ปญฺหา. จิตฺต\-สหภู\-ทุก\-สทิสา. นินฺนานากรณํ.)}

\csromanheader{2}{\textbf{ปุเรชาตาทิ}}
\proseitem{304}{โนจิตฺต\-สํสฏฺฐ\-สมุฏฺฐาโน ธมฺโม โนจิตฺต\-สํสฏฺฐ\-สมุฏฺฐานสฺส ธมฺมสฺส ปุเรชาต\-ปจฺจเยน ปจฺจโย. อารมฺมณ\-ปุเรชาตํ, วตฺถุ\-ปุเรชาตํ. ตีณิ. (จิตฺต\-สหภู\-ทุก\-สทิสา. นินฺนานากรณํ.)}

\prose{จิตฺต\-สํสฏฺฐ\-สมุฏฺฐาโน ธมฺโม โนจิตฺต\-สํสฏฺฐ\-สมุฏฺฐานสฺส ธมฺมสฺส ปจฺฉาชาต\-ปจฺจเยน ปจฺจโย. (จิตฺต\-สหภู\-ทุก\-สทิสา. นินฺนานากรณํ. ตีณิปิ ปจฺฉาชาตา. เทฺว. เอกมูลานํ เอกา ฆฏนา.) อาเสวน\-ปจฺจเยน ปจฺจโย. นว.}

\csromanheader{2}{\textbf{กมฺมาทิ}}
\proseitem{305}{จิตฺต\-สํสฏฺฐ\-สมุฏฺฐาโน ธมฺโม จิตฺต\-สํสฏฺฐ\-สมุฏฺฐานสฺส ธมฺมสฺส กมฺมปจฺจเยน ปจฺจโย. ตีณิ. (จิตฺต\-สหภู\-ทุก\-สทิสา. นินฺนานากรณา. ตีณิปิ สหชาตา, นานากฺขณิกา.) วิปาก\-ปจฺจเยน ปจฺจโย. นว. อาหารปจฺจเยน ปจฺจโย. นว. จิตฺต\-สหภู\-คมน\-สทิสา. เอกํเยว กพฬีการํ อาหารํ.) อินฺทฺริย\-ปจฺจเยน ปจฺจโย. นว. ฌานปจฺจเยน ปจฺจโย. ตีณิ. มคฺคปจฺจเยน ปจฺจโย. ตีณิ. สมฺปยุตฺต\-ปจฺจเยน ปจฺจโย. ปญฺจ.}

\proseitem{306}{จิตฺต\-สํสฏฺฐ\-สมุฏฺฐาโน ธมฺโม โนจิตฺต\-สํสฏฺฐ\-สมุฏฺฐานสฺส ธมฺมสฺส วิปฺปยุตฺต\-ปจฺจเยน ปจฺจโย. สหชาตํ, ปจฺฉาชาตํ. (สํขิตฺตํ.) (1)}

\prose{\csromanfolio{487}โนจิตฺต\-สํสฏฺฐ\-สมุฏฺฐาโน ธมฺโม โนจิตฺต\-สํสฏฺฐ\-สมุฏฺฐานสฺส ธมฺมสฺส วิปฺปยุตฺต\-ปจฺจเยน ปจฺจโย. สหชาตํ, ปุเรชาตํ, ปจฺฉาชาตํ. (สํขิตฺตํ.) (1)}

\prose{โนจิตฺต\-สํสฏฺฐ\-สมุฏฺฐาโน ธมฺโม จิตฺต\-สํสฏฺฐ\-สมุฏฺฐานสฺส ธมฺมสฺส วิปฺปยุตฺต\-ปจฺจเยน ปจฺจโย. สหชาตํ, ปุเรชาตํ.  สหชาตํ ปฏิสนฺธิ\-กฺขเณ วตฺถุ จิตฺต\-สํสฏฺฐ\-สมุฏฺฐานานํ ขนฺธานํ วิปฺปยุตฺต\-ปจฺจเยน ปจฺจโย. ปุเรชาตํ จกฺขายตนํ จกฺขุ\-วิญฺญาณสหคตานํ ขนฺธานํ วิปฺปยุตฺต\-ปจฺจเยน ปจฺจโย ฯเปฯ. กายายตนํ ฯเปฯ. วตฺถุ จิตฺต\-สํสฏฺฐ\-สมุฏฺฐานานํ ขนฺธานํ วิปฺปยุตฺต\-ปจฺจเยน ปจฺจโย. (2)}

\prose{โนจิตฺต\-สํสฏฺฐ\-สมุฏฺฐาโน ธมฺโม จิตฺต\-สํสฏฺฐ\-สมุฏฺฐานสฺส จ โนจิตฺต\-สํสฏฺฐ\-สมุฏฺฐานสฺส จ ธมฺมสฺส วิปฺปยุตฺต\-ปจฺจเยน ปจฺจโย. สหชาตํ, ปุเรชาตํ.  สหชาตํ ปฏิสนฺธิ\-กฺขเณ วตฺถุ จิตฺตสฺส สมฺปยุตฺตกานญฺจ ขนฺธานํ วิปฺปยุตฺต\-ปจฺจเยน ปจฺจโย. ปุเรชาตํ จกฺขายตนํ จกฺขุ\-วิญฺญาณสฺส สมฺปยุตฺตกานญฺจ ขนฺธานํ วิปฺปยุตฺต\-ปจฺจเยน ปจฺจโย ฯเปฯ. กายายตนํ ฯเปฯ. วตฺถุ จิตฺตสฺส สมฺปยุตฺตกานญฺจ ขนฺธานํ วิปฺปยุตฺต\-ปจฺจเยน ปจฺจโย. (3)}

\prose{จิตฺต\-สํสฏฺฐ\-สมุฏฺฐาโน จ โนจิตฺต\-สํสฏฺฐ\-สมุฏฺฐาโน จ ธมฺมา โนจิตฺต\-สํสฏฺฐ\-สมุฏฺฐานสฺส ธมฺมสฺส วิปฺปยุตฺต\-ปจฺจเยน ปจฺจโย. สหชาตํ, ปจฺฉาชาตํ. (สํขิตฺตํ.)}

\proseitem{307}{จิตฺต\-สํสฏฺฐ\-สมุฏฺฐาโน ธมฺโม จิตฺต\-สํสฏฺฐ\-สมุฏฺฐานสฺส ธมฺมสฺส อตฺถิปจฺจเยน ปจฺจโย. (ปฏิจฺจสทิสา.) จิตฺต\-สํสฏฺฐ\-สมุฏฺฐาโน ธมฺโม โนจิตฺต\-สํสฏฺฐ\-สมุฏฺฐานสฺส ธมฺมสฺส อตฺถิปจฺจเยน ปจฺจโย. สหชาตํ, ปจฺฉาชาตํ. (สํขิตฺตํ.) จิตฺต\-สํสฏฺฐ\-สมุฏฺฐาโน ธมฺโม จิตฺต\-สํสฏฺฐ\-สมุฏฺฐานสฺส จ โนจิตฺต\-สํสฏฺฐ\-สมุฏฺฐานสฺส จ ธมฺมสฺส อตฺถิปจฺจเยน ปจฺจโย. (ปฏิจฺจสทิสา.) (3)}

\prose{โนจิตฺต\-สํสฏฺฐ\-สมุฏฺฐาโน ธมฺโม โนจิตฺต\-สํสฏฺฐ\-สมุฏฺฐานสฺส ธมฺมสฺส อตฺถิปจฺจเยน ปจฺจโย, สหชาตํ, ปุเรชาตํ, ปจฺฉาชาตํ, อาหารํ, อินฺทฺริยํ. (ปุเรชาต\-สทิสํ, ปุเรชาตํ กาตพฺพํ. สพฺพํ สํขิตฺตํ. วิตฺถาเรตพฺพํ. (1)}

\prose{\csromanfolio{488}โนจิตฺต\-สํสฏฺฐ\-สมุฏฺฐาโน ธมฺโม จิตฺต\-สํสฏฺฐ\-สมุฏฺฐานสฺส ธมฺมสฺส อตฺถิปจฺจเยน ปจฺจโย. สหชาตํ, ปุเรชาตํ.  สหชาตํ จิตฺตํ สมฺปยุตฺตกานํ ขนฺธานํ อตฺถิปจฺจเยน ปจฺจโย. ปฏิสนฺธิ\-กฺขเณ จิตฺตํ สมฺปยุตฺตกานํ ขนฺธานํ อตฺถิปจฺจเยน ปจฺจโย. ปฏิสนฺธิ\-กฺขเณ วตฺถุ จิตฺต\-สํสฏฺฐ\-สมุฏฺฐานานํ ขนฺธานํ อตฺถิปจฺจเยน ปจฺจโย. (ปุเรชาตํ ปุเรชาต\-สทิสํ. นินฺนานากรณํ.) (2)}

\prose{โนจิตฺต\-สํสฏฺฐ\-สมุฏฺฐาโน ธมฺโม จิตฺต\-สํสฏฺฐ\-สมุฏฺฐานสฺส จ โนจิตฺต\-สํสฏฺฐ\-สมุฏฺฐานสฺส จ ธมฺมสฺส อตฺถิปจฺจเยน ปจฺจโย. สหชาตํ, ปุเรชาตํ.  สหชาตํ จิตฺตํ สมฺปยุตฺตกานํ ขนฺธานํ จิตฺตสมุฏฺฐานานญฺจ รูปานํ อตฺถิปจฺจเยน ปจฺจโย. ปฏิสนฺธิ\-กฺขเณ จิตฺตํ สมฺปยุตฺตกานํ ขนฺธานํ กฏตฺตา จ รูปานํ อตฺถิปจฺจเยน ปจฺจโย. ปฏิสนฺธิ\-กฺขเณ วตฺถุ จิตฺตสฺส สมฺปยุตฺตกานญฺจ ขนฺธานํ อตฺถิปจฺจเยน ปจฺจโย. (ปุเรชาตํ ปุเรชาต\-สทิสํ.) (3)}

\proseitem{308}{จิตฺต\-สํสฏฺฐ\-สมุฏฺฐาโน จ โนจิตฺต\-สํสฏฺฐ\-สมุฏฺฐาโน จ ธมฺมา จิตฺต\-สํสฏฺฐ\-สมุฏฺฐานสฺส ธมฺมสฺส อตฺถิปจฺจเยน ปจฺจโย. สหชาตํ, ปุเรชาตํ.  สหชาโต จกฺขุ\-วิญฺญาณ\-สหคโต เอโก ขนฺโธ จ จกฺขุวิญฺญาณญฺจ ทฺวินฺนํ ขนฺธานํ ฯเปฯ เทฺว ขนฺธา จ ฯเปฯ. จกฺขุ\-วิญฺญาณ\-สหคโต เอโก ขนฺโธ จ จกฺขายตนญฺจ ทฺวินฺนํ ขนฺธานํ ฯเปฯ เทฺว ขนฺธา จ ฯเปฯ. กาย\-วิญฺญาณ\-สหคโต ฯเปฯ. จิตฺต\-สํสฏฺฐ\-สมุฏฺฐาโน เอโก ขนฺโธ จ จิตฺตญฺจ ทฺวินฺนํ ขนฺธานํ อตฺถิปจฺจเยน ปจฺจโย. เทฺว ขนฺธา จ ฯเปฯ. จิตฺต\-สํสฏฺฐ\-สมุฏฺฐาโน เอโก ขนฺโธ จ วตฺถุ จ ทฺวินฺนํ ขนฺธานํ อตฺถิปจฺจเยน ปจฺจโย ฯเปฯ เทฺว ขนฺธา จ ฯเปฯ. (ปฏิสนฺธิ\-กฺขเณ เทฺวปิ กาตพฺพา.) (1)}

\prose{จิตฺต\-สํสฏฺฐ\-สมุฏฺฐาโน จ โนจิตฺต\-สํสฏฺฐ\-สมุฏฺฐาโน จ ธมฺมา โนจิตฺต\-สํสฏฺฐ\-สมุฏฺฐานสฺส ธมฺมสฺส อตฺถิปจฺจเยน ปจฺจโย. สหชาตํ, ปุเรชาตํ, ปจฺฉาชาตํ, อาหารํ, อินฺทฺริยํ.  สหชาตา จกฺขุ\-วิญฺญาณสหคตา ขนฺธา จ จกฺขายตนญฺจ จกฺขุ\-วิญฺญาณสฺส อตฺถิปจฺจเยน ปจฺจโย ฯเปฯ. กาย\-วิญฺญาณสหคตา ฯเปฯ. จิตฺต\-สํสฏฺฐ\-สมุฏฺฐานา ขนฺธา จ จิตฺตญฺจ จิตฺต\-สํสฏฺฐ\-สมุฏฺฐานานํ รูปานํ อตฺถิปจฺจเยน ปจฺจโย. จิตฺต\-สํสฏฺฐ\-สมุฏฺฐานา ขนฺธา จ มหาภูตา จ จิตฺต\-สํสฏฺฐ\-สมุฏฺฐานานํ รูปานํ อตฺถิปจฺจเยน ปจฺจโย. จิตฺต\-สํสฏฺฐ\-สมุฏฺฐานา ขนฺธา จ วตฺถุ จ จิตฺตสฺส อตฺถิปจฺจเยน ปจฺจโย. \csromanfolio{489}(ปฏิสนฺธิ\-กฺขเณ ตีณิ กาตพฺพา.) ปจฺฉาชาตา จิตฺต\-สํสฏฺฐ\-สมุฏฺฐานา ขนฺธา จ จิตฺตญฺจ ปุเรชาตสฺส อิมสฺส โนจิตฺต\-สํสฏฺฐ\-สมุฏฺฐานสฺส กายสฺส อตฺถิปจฺจเยน ปจฺจโย. ปจฺฉาชาตา จิตฺต\-สํสฏฺฐ\-สมุฏฺฐานา ขนฺธา จ จิตฺตญฺจ กพฬีกาโร อาหาโร จ อิมสฺส โนจิตฺต\-สํสฏฺฐ\-สมุฏฺฐานสฺส กายสฺส อตฺถิปจฺจเยน ปจฺจโย. ปจฺฉาชาตา จิตฺต\-สํสฏฺฐ\-สมุฏฺฐานา ขนฺธา จ จิตฺตญฺจ รูปชีวิตินฺทฺริยญฺจ กฏตฺตารูปานํ อตฺถิปจฺจเยน ปจฺจโย. (2)}

\prose{จิตฺต\-สํสฏฺฐ\-สมุฏฺฐาโน จ โนจิตฺต\-สํสฏฺฐ\-สมุฏฺฐาโน จ ธมฺมา จิตฺต\-สํสฏฺฐ\-สมุฏฺฐานสฺส จ โนจิตฺต\-สํสฏฺฐ\-สมุฏฺฐานสฺส จ ธมฺมสฺส อตฺถิปจฺจเยน ปจฺจโย. สหชาตํ, ปุเรชาตํ.  สหชาโต จกฺขุ\-วิญฺญาณ\-สหคโต เอโก ขนฺโธ จ จกฺขายตนญฺจ ทฺวินฺนํ ขนฺธานํ จกฺขุ\-วิญฺญาณสฺส จ อตฺถิปจฺจเยน ปจฺจโย. เทฺว ขนฺธา จ ฯเปฯ. กาย\-วิญฺญาณ\-สหคโต ฯเปฯ. สหชาโต จิตฺต\-สํสฏฺฐ\-สมุฏฺฐาโน เอโก ขนฺโธ จ จิตฺตญฺจ ทฺวินฺนํ ขนฺธานํ จิตฺตสมุฏฺฐานานญฺจ รูปานํ อตฺถิปจฺจเยน ปจฺจโย. เทฺว ขนฺธา จ ฯเปฯ. จิตฺตํสํสฏฺฐสมุฏฺฐาโน เอโก ขนฺโธ จ วตฺถุ จ ทฺวินฺนํ ขนฺธานํ จิตฺตสฺส จ อตฺถิปจฺจเยน ปจฺจโย. เทฺว ขนฺธา จ ฯเปฯ. (ปฏิสนฺธิ\-กฺขเณ เทฺวปิ กาตพฺพา.) (3)}

\csromanmarkh{1. ปจฺจยานุโลม}\csromanmarkhii{2. สงฺขฺยาวาร}\csromanheadernomark{2}{1. \textbf{ปจฺจยานุโลม 2. สงฺขฺยาวาร}}
\proseitem{309}{เหตุยา ตีณิ, อารมฺมเณ นว, อธิปติยา นว, อนนฺตเร นว, สมนนฺตเร นว, สหชาเต นว, อญฺญมญฺเญ นว, นิสฺสเย นว, อุปนิสฺสเย นว, ปุเรชาเต ตีณิ, ปจฺฉาชาเต ตีณิ, อาเสวเน นว, กมฺเม ตีณิ, วิปาเก นว, อาหาเร นว, อินฺทฺริเย นว, ฌาเน ตีณิ, มคฺเค ตีณิ, สมฺปยุตฺเต ปญฺจ, วิปฺปยุตฺเต ปญฺจ, อตฺถิยา นว, นตฺถิยา นว, วิคเต นว, อวิคเต นว.}

\vspace{\tipitakaparskip}
\csromancenter{\textbf{อนุโลมํ.}}
\proseitem{310}{จิตฺต\-สํสฏฺฐ\-สมุฏฺฐาโน ธมฺโม จิตฺต\-สํสฏฺฐ\-สมุฏฺฐานสฺส ธมฺมสฺส อารมฺมณ\-ปจฺจเยน ปจฺจโย. สหชาต\-ปจฺจเยน ปจฺจโย. อุปนิสฺสย\-ปจฺจเยน ปจฺจโย. กมฺมปจฺจเยน ปจฺจโย. (1)}

\prose{\csromanfolio{490}จิตฺต\-สํสฏฺฐ\-สมุฏฺฐาโน ธมฺโม โนจิตฺต\-สํสฏฺฐ\-สมุฏฺฐานสฺส ธมฺมสฺส อารมฺมณ\-ปจฺจเยน ปจฺจโย. สหชาต\-ปจฺจเยน ปจฺจโย. อุปนิสฺสย\-ปจฺจเยน ปจฺจโย. ปจฺฉาชาต\-ปจฺจเยน ปจฺจโย. กมฺมปจฺจเยน ปจฺจโย. (2)}

\prose{จิตฺต\-สํสฏฺฐ\-สมุฏฺฐาโน ธมฺโม จิตฺต\-สํสฏฺฐ\-สมุฏฺฐานสฺส จ โนจิตฺต\-สํสฏฺฐ\-สมุฏฺฐานสฺส จ ธมฺมสฺส อารมฺมณ\-ปจฺจเยน ปจฺจโย. สหชาต\-ปจฺจเยน ปจฺจโย. อุปนิสฺสย\-ปจฺจเยน ปจฺจโย. กมฺมปจฺจเยน ปจฺจโย. (3)}

\proseitem{311}{โนจิตฺต\-สํสฏฺฐ\-สมุฏฺฐาโน ธมฺโม โนจิตฺต\-สํสฏฺฐ\-สมุฏฺฐานสฺส ธมฺมสฺส อารมฺมณ\-ปจฺจเยน ปจฺจโย. สหชาต\-ปจฺจเยน ปจฺจโย. อุปนิสฺสย\-ปจฺจเยน ปจฺจโย. ปุเรชาต\-ปจฺจเยน ปจฺจโย. ปจฺฉาชาต\-ปจฺจเยน ปจฺจโย. อาหารปจฺจเยน ปจฺจโย. อินฺทฺริย\-ปจฺจเยน ปจฺจโย. (1)}

\prose{โนจิตฺต\-สํสฏฺฐ\-สมุฏฺฐาโน ธมฺโม จิตฺต\-สํสฏฺฐ\-สมุฏฺฐานสฺส ธมฺมสฺส อารมฺมณ\-ปจฺจเยน ปจฺจโย. สหชาต\-ปจฺจเยน ปจฺจโย. อุปนิสฺสย\-ปจฺจเยน ปจฺจโย. ปุเรชาต\-ปจฺจเยน ปจฺจโย. (2)}

\prose{โนจิตฺต\-สํสฏฺฐ\-สมุฏฺฐาโน ธมฺโม จิตฺต\-สํสฏฺฐ\-สมุฏฺฐานสฺส จ โนจิตฺต\-สํสฏฺฐ\-สมุฏฺฐานสฺส จ ธมฺมสฺส อารมฺมณ\-ปจฺจเยน ปจฺจโย. สหชาต\-ปจฺจเยน ปจฺจโย. อุปนิสฺสย\-ปจฺจเยน ปจฺจโย. ปุเรชาต\-ปจฺจเยน ปจฺจโย. (3)}

\proseitem{312}{จิตฺต\-สํสฏฺฐ\-สมุฏฺฐาโน จ โนจิตฺต\-สํสฏฺฐ\-สมุฏฺฐาโน จ ธมฺมา จิตฺต\-สํสฏฺฐ\-สมุฏฺฐานสฺส ธมฺมสฺส อารมฺมณ\-ปจฺจเยน ปจฺจโย. สหชาต\-ปจฺจเยน ปจฺจโย. อุปนิสฺสย\-ปจฺจเยน ปจฺจโย. (1)}

\prose{จิตฺต\-สํสฏฺฐ\-สมุฏฺฐาโน จ โนจิตฺต\-สํสฏฺฐ\-สมุฏฺฐาโน จ ธมฺมา โนจิตฺต\-สํสฏฺฐ\-สมุฏฺฐานสฺส ธมฺมสฺส อารมฺมณ\-ปจฺจเยน ปจฺจโย. สหชาต\-ปจฺจเยน ปจฺจโย. อุปนิสฺสย\-ปจฺจเยน ปจฺจโย. ปจฺฉาชาต\-ปจฺจเยน ปจฺจโย. (2)}

\prose{จิตฺต\-สํสฏฺฐ\-สมุฏฺฐาโน จ โนจิตฺต\-สํสฏฺฐ\-สมุฏฺฐาโน จ ธมฺมา จิตฺต\-สํสฏฺฐ\-สมุฏฺฐานสฺส จ โนจิตฺต\-สํสฏฺฐ\-สมุฏฺฐานสฺส จ ธมฺมสฺส อารมฺมณ\-ปจฺจเยน ปจฺจโย. สหชาต\-ปจฺจเยน ปจฺจโย. อุปนิสฺสย\-ปจฺจเยน ปจฺจโย. (3)}

\proseitem{65}{\csromanfolio{491}จิตฺต\-สํสฏฺฐ\-สมุฏฺฐานา\-นุปริวตฺติทุก 2. ปจฺจย\-ปจฺจนีย 2. สงฺขฺยาวาร}

\proseitem{313}{นเหตุยา นว, นอารมฺมเณ นว, (สพฺพตฺถ นว,) โนอวิคเต นว.}

\vspace{\tipitakaparskip}
\csromancenter{\textbf{เหตุ}ปจฺจยา นอารมฺมเณ ตีณิ, นอธิปติยา นว, นอนนฺตเร ตีณิ, นสมนนฺตเร ตีณิ, นอญฺญมญฺเญ เอกํ, นฺเอาปนิสฺสเย ตีณิ, (สพฺพตฺถ ตีณิ,) นสมฺปยุตฺเต เอกํ, นวิปฺปยุตฺเต ตีณิ, โนนตฺถิยา ตีณิ, โนวิคเต ตีณิ.}
\csromanheader{2}{315. นเหตุปจฺจยา อารมฺมเณ นว, อธิปติยา นว. (สพฺพตฺถ นว, อนุโลมมาติกา.)}
\nitthitam{จิตฺต\-สํสฏฺฐ\-สมุฏฺฐานทุกํ นิฏฺฐิตํ.}
\csromanmatikaanchor{mtk.63}
\csromantocmarkref{subsection}{64. จิตฺตสํสฏฺฐสมุฏฺฐานสหภูทุก}{mtk.63}
\bookmark[dest={mtk.63},level=2]{64. จิตฺตสํสฏฺฐสมุฏฺฐานสหภูทุก}
\csromanmarkh{64. จิตฺต\-สํสฏฺฐ\-สมุฏฺฐาน\-สหภูทุก}\csromanmarkhii{1. ปฏิจฺจวาร}\csromanheadernomark{2}{64. จิตฺต\-สํสฏฺฐ\-สมุฏฺฐาน\-สหภูทุก 1. ปฏิจฺจวาร}
\vspace{\tipitakaparskip}
\csromancenter{\textbf{เหตุ}}
\proseitem{316}{จิตฺต\-สํสฏฺฐ\-สมุฏฺฐาน\-สหภุํ ธมฺมํ ปฏิจฺจ จิตฺต\-สํสฏฺฐ\-สมุฏฺฐาน\-สหภู ธมฺโม อุปฺปชฺชติ เหตุปจฺจยา. จิตฺต\-สํสฏฺฐ\-สมุฏฺฐาน\-สหภุํ เอกํ ขนฺธํ ปฏิจฺจ เทฺว ขนฺธา. เทฺว ขนฺเธ ฯเปฯ. ปฏิสนฺธิ\-กฺขเณ ฯเปฯ.}

\vspace{\tipitakaparskip}
\csromancenter{(ยถา จิตฺต\-สํสฏฺฐ\-สมุฏฺฐานทุกํ, เอวํ อิมมฺปิ ทุกํ, นินฺนานากรณํ.)}
\nitthitam{จิตฺต\-สํสฏฺฐ\-สมุฏฺฐาน\-สหภูทุกํ นิฏฺฐิตํ.}
\csromanmatikaanchor{mtk.64}
\csromantocmarkref{subsection}{65. จิตฺตสํสฏฺฐสมุฏฺฐานานุปริวตฺติทุก}{mtk.64}
\bookmark[dest={mtk.64},level=2]{65. จิตฺตสํสฏฺฐสมุฏฺฐานานุปริวตฺติทุก}
\csromanmarkh{65. จิตฺต\-สํสฏฺฐ\-สมุฏฺฐานา\-นุปริวตฺติทุก}\csromanmarkhii{1. ปฏิจฺจวาร}\csromanheadernomark{2}{65. จิตฺต\-สํสฏฺฐ\-สมุฏฺฐานา\-นุปริวตฺติทุก 1. ปฏิจฺจวาร}
\csromanheader{2}{\textbf{เหตุ}}
\proseitem{317}{จิตฺต\-สํสฏฺฐ\-สมุฏฺฐานา\-นุปริวตฺติํ ธมฺมํ ปฏิจฺจ จิตฺต\-สํสฏฺฐ\-สมุฏฺฐานา\-นุปริวตฺตี ธมฺโม อุปฺปชฺชติ เหตุปจฺจยา. จิตฺต\-สํสฏฺฐ\-สมุฏฺฐานา\-นุปริวตฺติํ \csromanfolio{492}เอกํ ขนฺธํ ปฏิจฺจ เทฺว ขนฺธา, เทฺว ขนฺเธ ฯเปฯ. ปฏิสนฺธิ\-กฺขเณ ฯเปฯ.}

\vspace{\tipitakaparskip}
\csromancenter{(ยถา จิตฺต\-สํสฏฺฐ\-สมุฏฺฐาน\-ทุก\-สทิสํ, นินฺนานากรณํ.)}
\nitthitam{จิตฺต\-สํสฏฺฐ\-สมุฏฺฐานา\-นุปริวตฺติทุกํ นิฏฺฐิตํ.}
\csromanmatikaanchor{mtk.65}
\csromantocmarkref{subsection}{66. อชฺฌตฺติกทุก}{mtk.65}
\bookmark[dest={mtk.65},level=2]{66. อชฺฌตฺติกทุก}
\csromanmarkh{66. อชฺฌตฺติกทุก}\csromanmarkhii{1. ปฏิจฺจวาร}\csromanheadernomark{2}{66. \textbf{อชฺฌตฺติกทุก 1. ปฏิจฺจวาร}}
\csromanmarkh{1. ปจฺจยานุโลม}\csromanmarkhii{1. วิภงฺควาร}\csromanheadernomark{2}{1. \textbf{ปจฺจยานุโลม 1. วิภงฺควาร}}
\csromanheader{2}{\textbf{เหตุ}}
\proseitem{318}{อชฺฌตฺติกํ ธมฺมํ ปฏิจฺจ อชฺฌตฺติโก ธมฺโม อุปฺปชฺชติ เหตุปจฺจยา. ปฏิสนฺธิ\-กฺขเณ จิตฺตํ ปฏิจฺจ อชฺฌตฺติกํ กฏตฺตารูปํ. (1)}

\prose{อชฺฌตฺติกํ ธมฺมํ ปฏิจฺจ พาหิโร ธมฺโม อุปฺปชฺชติ เหตุปจฺจยา. จิตฺตํ ปฏิจฺจ สมฺปยุตฺตกา ขนฺธา จิตฺต\-สมุฏฺฐานญฺจ รูปํ. ปฏิสนฺธิ\-กฺขเณ จิตฺตํ ปฏิจฺจ สมฺปยุตฺตกา ขนฺธา พาหิรํ กฏตฺตา จ รูปํ. (2)}

\prose{อชฺฌตฺติกํ ธมฺมํ ปฏิจฺจ อชฺฌตฺติโก จ พาหิโร จ ธมฺมา อุปฺปชฺชนฺติ เหตุปจฺจยา. ปฏิสนฺธิ\-กฺขเณ จิตฺตํ ปฏิจฺจ สมฺปยุตฺตกา ขนฺธา อชฺฌตฺติกญฺจ พาหิรญฺจ กฏตฺตารูปํ. (3)}

\proseitem{319}{พาหิรํ ธมฺมํ ปฏิจฺจ พาหิโร ธมฺโม อุปฺปชฺชติ เหตุปจฺจยา. พาหิรํ เอกํ ขนฺธํ ปฏิจฺจ เทฺว ขนฺธา จิตฺต\-สมุฏฺฐานญฺจ รูปํ. เทฺว ขนฺเธ ฯเปฯ. ปฏิสนฺธิ\-กฺขเณ พาหิรํ เอกํ ขนฺธํ ปฏิจฺจ เทฺว ขนฺธา พาหิรํ กฏตฺตา จ รูปํ. เทฺว ขนฺเธ ฯเปฯ. ขนฺเธ ปฏิจฺจ วตฺถุ, วตฺถุํ ปฏิจฺจ ขนฺธา. เอกํ มหาภูตํ ฯเปฯ มหาภูเต ปฏิจฺจ จิตฺต\-สมุฏฺฐานํ รูปํ, กฏตฺตารูปํ, อุปาทารูปํ. (1)}

\prose{พาหิรํ ธมฺมํ ปฏิจฺจ อชฺฌตฺติโก ธมฺโม อุปฺปชฺชติ เหตุปจฺจยา. พาหิเร ขนฺเธ ปฏิจฺจ จิตฺตํ. ปฏิสนฺธิ\-กฺขเณ พาหิเร ขนฺเธ ปฏิจฺจ จิตฺตํ อชฺฌตฺติกํ กฏตฺตา จ รูปํ. ปฏิสนฺธิ\-กฺขเณ พาหิรํ วตฺถุํ ปฏิจฺจ จิตฺตํ. (2)}

\prose{พาหิรํ ธมฺมํ ปฏิจฺจ อชฺฌตฺติโก จ พาหิโร จ ธมฺมา อุปฺปชฺชนฺติ เหตุปจฺจยา. พาหิรํ เอกํ ขนฺธํ ปฏิจฺจ เทฺว ขนฺธา จิตฺตญฺจ จิตฺต\-สมุฏฺฐานญฺจ รูปํ. เทฺว ขนฺเธ ฯเปฯ. ปฏิสนฺธิ\-กฺขเณ พาหิรํ เอกํ ขนฺธํ ปฏิจฺจ เทฺว ขนฺธา จิตฺตญฺจ อชฺฌตฺติกญฺจ พาหิรญฺจ \csromanfolio{493}กฏตฺตารูปํ. เทฺว ขนฺเธ ฯเปฯ. ปฏิสนฺธิ\-กฺขเณ วตฺถุํ ปฏิจฺจ จิตฺตํ สมฺปยุตฺตกา จ ขนฺธา. (3)}

\proseitem{320}{อชฺฌตฺติกญฺจ พาหิรญฺจ ธมฺมํ ปฏิจฺจ อชฺฌตฺติโก ธมฺโม อุปฺปชฺชติ เหตุปจฺจยา. ปฏิสนฺธิ\-กฺขเณ จิตฺตญฺจ สมฺปยุตฺตเก จ ขนฺเธ ปฏิจฺจ อชฺฌตฺติกํ กฏตฺตารูปํ. (1)}

\prose{อชฺฌตฺติกญฺจ พาหิรญฺจ ธมฺมํ ปฏิจฺจ พาหิโร ธมฺโม อุปฺปชฺชติ เหตุปจฺจยา. พาหิรํ เอกํ ขนฺธญฺจ จิตฺตญฺจ ปฏิจฺจ เทฺว ขนฺธา จิตฺต\-สมุฏฺฐานญฺจ รูปํ. เทฺว ขนฺเธ จ ฯเปฯ. จิตฺตญฺจ มหาภูเต จ ปฏิจฺจ จิตฺต\-สมุฏฺฐานํ รูปํ. ปฏิสนฺธิ\-กฺขเณ พาหิรํ เอกํ ขนฺธญฺจ จิตฺตญฺจ ปฏิจฺจ เทฺว ขนฺธา พาหิรํ กฏตฺตา จ รูปํ. เทฺว ขนฺเธ จ ฯเปฯ. จิตฺตญฺจ มหาภูเต จ ปฏิจฺจ พาหิรํ กฏตฺตารูปํ. ปฏิสนฺธิ\-กฺขเณ จิตฺตญฺจ วตฺถุญฺจ ปฏิจฺจ พาหิรา ขนฺธา. (2)}

\prose{อชฺฌตฺติกญฺจ พาหิรญฺจ ธมฺมํ ปฏิจฺจ อชฺฌตฺติโก จ พาหิโร จ ธมฺมา อุปฺปชฺชนฺติ เหตุปจฺจยา. ปฏิสนฺธิ\-กฺขเณ พาหิรํ เอกํ ขนฺธญฺจ จิตฺตญฺจ ปฏิจฺจ เทฺว ขนฺธา อชฺฌตฺติกญฺจ พาหิรญฺจ กฏตฺตารูปํ. เทฺว ขนฺเธ จ ฯเปฯ. (3)}

\proseitem{321}{อชฺฌตฺติกํ ธมฺมํ ปฏิจฺจ พาหิโร ธมฺโม อุปฺปชฺชติ อารมฺมณ\-ปจฺจยา. จิตฺตํ ปฏิจฺจ สมฺปยุตฺตกา ขนฺธา. ปฏิสนฺธิ\-กฺขเณ จิตฺตํ ปฏิจฺจ สมฺปยุตฺตกา ขนฺธา. (1)}

\prose{พาหิรํ ธมฺมํ ปฏิจฺจ พาหิโร ธมฺโม อุปฺปชฺชติ อารมฺมณ\-ปจฺจยา. พาหิรํ เอกํ ขนฺธํ ปฏิจฺจ เทฺว ขนฺธา. เทฺว ขนฺเธ ฯเปฯ. ปฏิสนฺธิ\-กฺขเณ ฯเปฯ. วตฺถุํ ปฏิจฺจ ขนฺธา. (1)}

\prose{พาหิรํ ธมฺมํ ปฏิจฺจ อชฺฌตฺติโก ธมฺโม อุปฺปชฺชติ อารมฺมณ\-ปจฺจยา. พาหิเร ขนฺเธ ปฏิจฺจ จิตฺตํ. ปฏิสนฺธิ\-กฺขเณ พาหิเร ขนฺเธ ปฏิจฺจ จิตฺตํ. ปฏิสนฺธิ\-กฺขเณ วตฺถุํ ปฏิจฺจ จิตฺตํ. (2)}

\prose{พาหิรํ ธมฺมํ ปฏิจฺจ อชฺฌตฺติโก จ พาหิโร จ ธมฺมา อุปฺปชฺชนฺติ อารมฺมณ\-ปจฺจยา. พาหิรํ เอกํ ขนฺธํ ปฏิจฺจ เทฺว ขนฺธา จิตฺตญฺจ. เทฺว ขนฺเธ ฯเปฯ. ปฏิสนฺธิ\-กฺขเณ พาหิรํ เอกํ ขนฺธํ ปฏิจฺจ เทฺว ขนฺธา จิตฺตญฺจ. เทฺว ขนฺเธ ฯเปฯ. ปฏิสนฺธิ\-กฺขเณ วตฺถุํ ปฏิจฺจ จิตฺตํ สมฺปยุตฺตกา จ ขนฺธา. (3)}

\prose{\csromanfolio{494}อชฺฌตฺติกญฺจ พาหิรญฺจ ธมฺมํ ปฏิจฺจ พาหิโร ธมฺโม อุปฺปชฺชติ อารมฺมณ\-ปจฺจยา. พาหิรํ เอกํ ขนฺธญฺจ จิตฺตญฺจ ปฏิจฺจ เทฺว ขนฺธา. เทฺว ขนฺเธ จ ฯเปฯ. ปฏิสนฺธิ\-กฺขเณ พาหิรํ เอกํ ขนฺธญฺจ จิตฺตญฺจ ปฏิจฺจ เทฺว ขนฺธา. เทฺว ขนฺเธ จ ฯเปฯ. ปฏิสนฺธิ\-กฺขเณ พาหิรํ เอกํ ขนฺธญฺจ วตฺถุญฺจ ปฏิจฺจ เทฺว ขนฺธา. เทฺว ขนฺเธ จ ฯเปฯ. (สํขิตฺตํ.) (1)}

\csromanmarkh{1. ปจฺจยานุโลม}\csromanmarkhii{2. สงฺขฺยาวาร}\csromanheadernomark{2}{1. \textbf{ปจฺจยานุโลม 2. สงฺขฺยาวาร}}
\proseitem{322}{เหตุยา นว, อารมฺมเณ ปญฺจ, อธิปติยา ปญฺจ, อนนฺตเร ปญฺจ, สมนนฺตเร ปญฺจ, สหชาเต นว, อญฺญมญฺเญ ปญฺจ, นิสฺสเย นว, อุปนิสฺสเย ปญฺจ, ปุเรชาเต ปญฺจ, อาเสวเน ปญฺจ, กมฺเม นว, วิปาเก นว, (สพฺพตฺถ นว,) สมฺปยุตฺเต ปญฺจ, วิปฺปยุตฺเต นว, อตฺถิยา นว, นตฺถิยา ปญฺจ, วิคเต ปญฺจ, อวิคเต นว.}

\csromanmarkh{2. ปจฺจย\-ปจฺจนีย}\csromanmarkhii{1. วิภงฺควาร}\csromanheadernomark{2}{2. ปจฺจย\-ปจฺจนีย 1. วิภงฺควาร}
\proseitem{323}{อชฺฌตฺติกํ ธมฺมํ ปฏิจฺจ อชฺฌตฺติโก ธมฺโม อุปฺปชฺชติ นเหตุปจฺจยา. อเหตุก\-ปฏิสนฺธิ\-กฺขเณ จิตฺตํ ปฏิจฺจ อชฺฌตฺติกํ กฏตฺตารูปํ. (1)}

\prose{อชฺฌตฺติกํ ธมฺมํ ปฏิจฺจ พาหิโร ธมฺโม อุปฺปชฺชติ นเหตุปจฺจยา. อเหตุกํ จิตฺตํ ปฏิจฺจ สมฺปยุตฺตกา ขนฺธา จิตฺต\-สมุฏฺฐานญฺจ รูปํ. อเหตุก\-ปฏิสนฺธิ\-กฺขเณ จิตฺตํ ปฏิจฺจ สมฺปยุตฺตกา ขนฺธา พาหิรํ กฏตฺตา จ รูปํ. วิจิกิจฺฉา\-สหคตํ อุทฺธจฺจ\-สหคตํ จิตฺตํ ปฏิจฺจ วิจิกิจฺฉา\-สหคโต อุทฺธจฺจ\-สหคโต โมโห. (2)}

\prose{อชฺฌตฺติกํ ธมฺมํ ปฏิจฺจ อชฺฌตฺติโก จ พาหิโร จ ธมฺมา อุปฺปชฺชนฺติ นเหตุปจฺจยา. อเหตุก\-ปฏิสนฺธิ\-กฺขเณ จิตฺตํ ปฏิจฺจ สมฺปยุตฺตกา ขนฺธา อชฺฌตฺติกญฺจ พาหิรญฺจ กฏตฺตารูปํ. (3)}

\proseitem{324}{พาหิรํ ธมฺมํ ปฏิจฺจ พาหิโร ธมฺโม อุปฺปชฺชติ นเหตุปจฺจยา. อเหตุกํ พาหิรํ เอกํ ขนฺธํ ปฏิจฺจ เทฺว ขนฺธา จิตฺต\-สมุฏฺฐานญฺจ รูปํ. เทฺว ขนฺเธ ฯเปฯ. \csromanfolio{495}อเหตุก\-ปฏิสนฺธิ\-กฺขเณ ฯเปฯ. (ยาว อสญฺญสตฺตา.) วิจิกิจฺฉา\-สหคเต อุทฺธจฺจ\-สหคเต ขนฺเธ ปฏิจฺจ วิจิกิจฺฉา\-สหคโต อุทฺธจฺจ\-สหคโต โมโห. (1)}

\prose{พาหิรํ ธมฺมํ ปฏิจฺจ อชฺฌตฺติโก ธมฺโม อุปฺปชฺชติ นเหตุปจฺจยา. อเหตุเก พาหิเร ขนฺเธ ปฏิจฺจ จิตฺตํ. อเหตุก\-ปฏิสนฺธิ\-กฺขเณ พาหิเร ขนฺเธ ปฏิจฺจ จิตฺตํ อชฺฌตฺติกํ กฏตฺตา จ รูปํ. อเหตุก\-ปฏิสนฺธิ\-กฺขเณ วตฺถุํ ปฏิจฺจ จิตฺตํ. (2)}

\prose{พาหิรํ ธมฺมํ ปฏิจฺจ อชฺฌตฺติโก จ พาหิโร จ ธมฺมา อุปฺปชฺชนฺติ นเหตุปจฺจยา. อเหตุกํ พาหิรํ เอกํ ขนฺธํ ปฏิจฺจ เทฺว ขนฺธา จิตฺตญฺจ จิตฺต\-สมุฏฺฐานญฺจ รูปํ. เทฺว ขนฺเธ ฯเปฯ. อเหตุก\-ปฏิสนฺธิ\-กฺขเณ พาหิรํ เอกํ ขนฺธํ ปฏิจฺจ เทฺว ขนฺธา จิตฺตญฺจ อชฺฌตฺติกญฺจ พาหิรญฺจ กฏตฺตารูปํ. อเหตุก\-ปฏิสนฺธิ\-กฺขเณ วตฺถุํ ปฏิจฺจ จิตฺตํ สมฺปยุตฺตกา จ ขนฺธา. (3)}

\proseitem{325}{อชฺฌตฺติกญฺจ พาหิรญฺจ ธมฺมํ ปฏิจฺจ อชฺฌตฺติโก ธมฺโม อุปฺปชฺชติ นเหตุปจฺจยา. อเหตุก\-ปฏิสนฺธิ\-กฺขเณ จิตฺตญฺจ สมฺปยุตฺตเก จ ขนฺเธ ปฏิจฺจ อชฺฌตฺติกํ กฏตฺตารูปํ. (1)}

\prose{อชฺฌตฺติกญฺจ พาหิรญฺจ ธมฺมํ ปฏิจฺจ พาหิโร ธมฺโม อุปฺปชฺชติ นเหตุปจฺจยา. อเหตุกํ พาหิรํ เอกํ ขนฺธญฺจ จิตฺตญฺจ ปฏิจฺจ เทฺว ขนฺธา จิตฺต\-สมุฏฺฐานญฺจ รูปํ. เทฺว ขนฺเธ ฯเปฯ. อเหตุกํ จิตฺตญฺจ มหาภูเต จ ปฏิจฺจ จิตฺต\-สมุฏฺฐานํ รูปํ. อเหตุก\-ปฏิสนฺธิ\-กฺขเณ พาหิรํ เอกํ ขนฺธญฺจ จิตฺตญฺจ ปฏิจฺจ เทฺว ขนฺธา พาหิรํ กฏตฺตา จ รูปํ. อเหตุก\-ปฏิสนฺธิ\-กฺขเณ จิตฺตญฺจ มหาภูเต จ ปฏิจฺจ พาหิรํ กฏตฺตารูปํ. อเหตุก\-ปฏิสนฺธิ\-กฺขเณ จิตฺตญฺจ วตฺถุญฺจ ปฏิจฺจ พาหิรา ขนฺธา. วิจิกิจฺฉา\-สหคเต อุทฺธจฺจ\-สหคเต ขนฺเธ จ จิตฺตญฺจ ปฏิจฺจ วิจิกิจฺฉา\-สหคโต อุทฺธจฺจ\-สหคโต โมโห. (2)}

\prose{อชฺฌตฺติกญฺจ พาหิรญฺจ ธมฺมํ ปฏิจฺจ อชฺฌตฺติโก จ พาหิโร จ ธมฺมา อุปฺปชฺชนฺติ นเหตุปจฺจยา. อเหตุก\-ปฏิสนฺธิ\-กฺขเณ พาหิรํ เอกํ ขนฺธญฺจ จิตฺตญฺจ ปฏิจฺจ เทฺว ขนฺธา อชฺฌตฺติกญฺจ พาหิรญฺจ กฏตฺตารูปํ. เทฺว ขนฺเธ จ ฯเปฯ. (3)}

\prose{\csromanfolio{496}นอารมฺมณ}

\proseitem{326}{อชฺฌตฺติกํ ธมฺมํ ปฏิจฺจ อชฺฌตฺติโก ธมฺโม อุปฺปชฺชติ นอารมฺมณ\-ปจฺจยา. ปฏิสนฺธิ\-กฺขเณ จิตฺตํ ปฏิจฺจ อชฺฌตฺติกํ กฏตฺตารูปํ. (1)}

\prose{อชฺฌตฺติกํ ธมฺมํ ปฏิจฺจ พาหิโร ธมฺโม อุปฺปชฺชติ นอารมฺมณ\-ปจฺจยา. จิตฺตํ ปฏิจฺจ จิตฺต\-สมุฏฺฐานํ รูปํ. ปฏิสนฺธิ\-กฺขเณ จิตฺตํ ปฏิจฺจ พาหิรํ กฏตฺตารูปํ. (2)}

\prose{อชฺฌตฺติกํ ธมฺมํ ปฏิจฺจ อชฺฌตฺติโก จ พาหิโร จ ธมฺมา อุปฺปชฺชนฺติ นอารมฺมณ\-ปจฺจยา. ปฏิสนฺธิ\-กฺขเณ จิตฺตํ ปฏิจฺจ อชฺฌตฺติกญฺจ พาหิรญฺจ กฏตฺตารูปํ. (3)}

\prose{พาหิรํ ธมฺมํ ปฏิจฺจ พาหิโร ธมฺโม อุปฺปชฺชติ นอารมฺมณ\-ปจฺจยา. พาหิเร ขนฺเธ ปฏิจฺจ จิตฺต\-สมุฏฺฐานํ รูปํ. ปฏิสนฺธิ\-กฺขเณ พาหิเร ขนฺเธ ปฏิจฺจ พาหิรํ กฏตฺตารูปํ. พาหิเร ขนฺเธ ปฏิจฺจ วตฺถุ. เอกํ มหาภูตํ ฯเปฯ. (ยาว อสญฺญสตฺตา.) (1)}

\prose{พาหิรํ ธมฺมํ ปฏิจฺจ อชฺฌตฺติโก ธมฺโม อุปฺปชฺชติ นอารมฺมณ\-ปจฺจยา. ปฏิสนฺธิ\-กฺขเณ พาหิเร ขนฺเธ ปฏิจฺจ อชฺฌตฺติกํ กฏตฺตารูปํ. (2)}

\prose{พาหิรํ ธมฺมํ ปฏิจฺจ อชฺฌตฺติโก จ พาหิโร จ ธมฺมา อุปฺปชฺชนฺติ นอารมฺมณ\-ปจฺจยา. ปฏิสนฺธิ\-กฺขเณ พาหิเร ขนฺเธ ปฏิจฺจ อชฺฌตฺติกญฺจ พาหิรญฺจ กฏตฺตารูปํ. (3)}

\proseitem{327}{อชฺฌตฺติกญฺจ พาหิรญฺจ ธมฺมํ ปฏิจฺจ อชฺฌตฺติโก ธมฺโม อุปฺปชฺชติ นอารมฺมณ\-ปจฺจยา. ปฏิสนฺธิ\-กฺขเณ จิตฺตญฺจ สมฺปยุตฺตเก จ ขนฺเธ ปฏิจฺจ อชฺฌตฺติกํ กฏตฺตารูปํ. (1)}

\prose{อชฺฌตฺติกญฺจ พาหิรญฺจ ธมฺมํ ปฏิจฺจ พาหิโร ธมฺโม อุปฺปชฺชติ นอารมฺมณ\-ปจฺจยา. พาหิเร ขนฺเธ จ จิตฺตญฺจ ปฏิจฺจ จิตฺต\-สมุฏฺฐานํ รูปํ. จิตฺตญฺจ มหาภูเต จ ปฏิจฺจ จิตฺต\-สมุฏฺฐานํ รูปํ. (ปฏิสนฺธิ\-กฺขเณ เทฺว กาตพฺพา.)}

\prose{อชฺฌตฺติกญฺจ พาหิรญฺจ ธมฺมํ ปฏิจฺจ อชฺฌตฺติโก จ พาหิโร จ ธมฺมา อุปฺปชฺชนฺติ นอารมฺมณ\-ปจฺจยา. ปฏิสนฺธิ\-กฺขเณ จิตฺตญฺจ สมฺปยุตฺตเก จ ขนฺเธ ปฏิจฺจ อชฺฌตฺติกญฺจ พาหิรญฺจ กฏตฺตารูปํ. (สํขิตฺตํ.) (3)}

\csromanheader{2}{66. อชฺฌตฺติกทุก \textbf{นฌาน}}
\proseitem{328}{\csromanfolio{497}อชฺฌตฺติกํ ธมฺมํ ปฏิจฺจ พาหิโร ธมฺโม อุปฺปชฺชติ นฌานปจฺจยา. จกฺขุวิญฺญาณํ ปฏิจฺจ สมฺปยุตฺตกา ขนฺธา ฯเปฯ. กายวิญฺญาณํ ฯเปฯ.}

\prose{พาหิรํ ธมฺมํ ปฏิจฺจ พาหิโร ธมฺโม อุปฺปชฺชติ นฌานปจฺจยา. จกฺขุ\-วิญฺญาณสหคตํ เอกํ ขนฺธํ ปฏิจฺจ เทฺว ขนฺธา. เทฺว ขนฺเธ ฯเปฯ. กาย\-วิญฺญาณสหคตํ ฯเปฯ. พาหิรํ. อาหารสมุฏฺฐานํ. อุตุสมุฏฺฐานํ. อสญฺญสตฺตานํ ฯเปฯ. พาหิรํ ธมฺมํ ปฏิจฺจ อชฺฌตฺติโก ธมฺโม อุปฺปชฺชติ นฌานปจฺจยา. จกฺขุ\-วิญฺญาณสหคเต ขนฺเธ ปฏิจฺจ จกฺขุวิญฺญาณํ ฯเปฯ. กาย\-วิญฺญาณสหคเต ขนฺเธ ปฏิจฺจ กายวิญฺญาณํ.  พาหิรํ ธมฺมํ ปฏิจฺจ อชฺฌตฺติโก จ พาหิโร จ ธมฺมา อุปฺปชฺชนฺติ นฌานปจฺจยา. จกฺขุ\-วิญฺญาณสหคตํ เอกํ ขนฺธํ ปฏิจฺจ เทฺว ขนฺธา จกฺขุวิญฺญาณญฺจ. เทฺว ขนฺเธ ฯเปฯ. กาย\-วิญฺญาณสหคตํ เอกํ ขนฺธํ ฯเปฯ. (3)}

\prose{อชฺฌตฺติกญฺจ พาหิรญฺจ ธมฺมํ ปฏิจฺจ พาหิโร ธมฺโม อุปฺปชฺชติ นฌานปจฺจยา. จกฺขุ\-วิญฺญาณสหคตํ เอกํ ขนฺธญฺจ จกฺขุวิญฺญาณญฺจ ปฏิจฺจ เทฺว ขนฺธา. เทฺว ขนฺเธ ฯเปฯ. กายวิญฺญาณํ. (จกฺกํ.)}

\csromanmarkh{2. ปจฺจย\-ปจฺจนีย}\csromanmarkhii{2. สงฺขฺยาวาร}\csromanheadernomark{2}{2. ปจฺจย\-ปจฺจนีย 2. สงฺขฺยาวาร}
\proseitem{329}{นเหตุยา นว, นอารมฺมเณ นว, นอธิปติยา นว, นอนนฺตเร นว, (สพฺพตฺถ นว.) นกมฺเม ตีณิ, นวิปาเก ปญฺจ, นอาหาเร เอกํ, นอินฺทฺริเย เอกํ, นฌาเน ปญฺจ, นมคฺเค นว, นสมฺปยุตฺเต นว, นวิปฺปยุตฺเต ปญฺจ, โนนตฺถิยา นว, โนวิคเต นว.}

\proseitem{330}{เหตุปจฺจยา นอารมฺมเณ นว, นอธิปติยา นว. (สํขิตฺตํ.)}

\proseitem{331}{นเหตุปจฺจยา อารมฺมเณ ปญฺจ, อนนฺตเร ปญฺจ, สมนนฺตเร ปญฺจ, สหชาเต นว ฯเปฯ มคฺเค ตีณิ. (สํขิตฺตํ.)}

\vspace{\tipitakaparskip}
\csromancenter{(สหชาตวาโรปิ ปฏิจฺจ\-วาร\-สทิโส.)}
\csromanmarkh{66. อชฺฌตฺติกทุก}\csromanmarkhii{3. ปจฺจยวาร}\csromanheadernomark{2}{66. \textbf{อชฺฌตฺติกทุก 3. ปจฺจยวาร}}
\csromanmarkh{1. ปจฺจยานุโลม}\csromanmarkhii{1. วิภงฺควาร}\csromanheadernomark{2}{1. \textbf{ปจฺจยานุโลม 1. วิภงฺควาร}}
\csromanheader{2}{\textbf{เหตุ}}
\proseitem{332}{\csromanfolio{498}อชฺฌตฺติกํ ธมฺมํ ปจฺจยา อชฺฌตฺติโก ธมฺโม อุปฺปชฺชติ เหตุปจฺจยา. ตีณิ. (ปฏิจฺจสทิสา.)}

\prose{พาหิรํ ธมฺมํ ปจฺจยา พาหิโร ธมฺโม อุปฺปชฺชติ เหตุปจฺจยา. พาหิรํ เอกํ ขนฺธํ ปจฺจยา เทฺว ขนฺธา จิตฺต\-สมุฏฺฐานญฺจ รูปํ. เทฺว ขนฺเธ ฯเปฯ. (ปฏิสนฺธิ\-กฺขเณ เทฺวปิ กาตพฺพา, ยาว อชฺฌตฺติกา มหาภูตา.) วตฺถุํ ปจฺจยา พาหิรา ขนฺธา.  พาหิรํ ธมฺมํ ปจฺจยา อชฺฌตฺติโก ธมฺโม อุปฺปชฺชติ เหตุปจฺจยา. พาหิเร ขนฺเธ ปจฺจยา จิตฺตํ. วตฺถุํ ปจฺจยา จิตฺตํ. (ปฏิสนฺธิ\-กฺขเณ เทฺวปิ กาตพฺพา.). พาหิรํ ธมฺมํ ปจฺจยา อชฺฌตฺติโก จ พาหิโร จ ธมฺมา อุปฺปชฺชนฺติ เหตุปจฺจยา. พาหิรํ เอกํ ขนฺธํ ปจฺจยา เทฺว ขนฺธา จิตฺตญฺจ จิตฺต\-สมุฏฺฐานญฺจ รูปํ. เทฺว ขนฺเธ ฯเปฯ. วตฺถุํ ปจฺจยา จิตฺตํ สมฺปยุตฺตกา จ ขนฺธา. (ปฏิสนฺธิ\-กฺขเณ เทฺวปิ กาตพฺพา.) (3)}

\proseitem{333}{อชฺฌตฺติกญฺจ พาหิรญฺจ ธมฺมํ ปจฺจยา อชฺฌตฺติโก ธมฺโม อุปฺปชฺชติ เหตุปจฺจยา. ปฏิสนฺธิ\-กฺขเณ จิตฺตญฺจ สมฺปยุตฺตเก จ ขนฺเธ ปจฺจยา อชฺฌตฺติกํ กฏตฺตารูปํ.  อชฺฌตฺติกญฺจ พาหิรญฺจ ธมฺมํ ปจฺจยา พาหิโร ธมฺโม อุปฺปชฺชติ เหตุปจฺจยา. พาหิรํ เอกํ ขนฺธญฺจ จิตฺตญฺจ ปจฺจยา เทฺว ขนฺธา จิตฺต\-สมุฏฺฐานญฺจ รูปํ. เทฺว ขนฺเธ ฯเปฯ. จิตฺตญฺจ มหาภูเต จ ปจฺจยา จิตฺต\-สมุฏฺฐานํ รูปํ. จิตฺตญฺจ วตฺถุญฺจ ปจฺจยา พาหิรา ขนฺธา. (ปฏิสนฺธิ\-กฺขเณ ตีณิปิ กาตพฺพา.). อชฺฌตฺติกญฺจ พาหิรญฺจ ธมฺมํ ปจฺจยา อชฺฌตฺติโก จ พาหิโร จ ธมฺมา อุปฺปชฺชนฺติ เหตุปจฺจยา. ปฏิสนฺธิ\-กฺขเณ พาหิรํ เอกํ ขนฺธญฺจ จิตฺตญฺจ ปจฺจยา เทฺว ขนฺธา อชฺฌตฺติกญฺจ พาหิรญฺจ กฏตฺตารูปํ. เทฺว ขนฺเธ ฯเปฯ. (3)}

\csromanheader{2}{\textbf{อรมฺมณ}}
\proseitem{334}{อชฺฌตฺติกํ ธมฺมํ ปจฺจยา อชฺฌตฺติโก ธมฺโม อุปฺปชฺชติ อารมฺมณ\-ปจฺจยา. จกฺขายตนํ ปจฺจยา จกฺขุวิญฺญาณํ ฯเปฯ. กายายตนํ ฯเปฯ.  อชฺฌตฺติกํ ธมฺมํ ปจฺจยา พาหิโร ธมฺโม อุปฺปชฺชติ อารมฺมณ\-ปจฺจโย. \csromanfolio{499}จกฺขายตนญฺจ จกฺขุวิญฺญาณญฺจ ปจฺจยา จกฺขุ\-วิญฺญาณสหคตา ขนฺธา ฯเปฯ. กายายตนญฺจ ฯเปฯ. จิตฺตํ ปจฺจยา สมฺปยุตฺตกา ขนฺธา. ปฏิสนฺธิ\-กฺขเณ ฯเปฯ. อชฺฌตฺติกํ ธมฺมํ ปจฺจยา อชฺฌตฺติโก จ พาหิโร จ ธมฺมา อุปฺปชฺชนฺติ อารมฺมณ\-ปจฺจยา. จกฺขายตนํ ปจฺจยา จกฺขุวิญฺญาณํ สมฺปยุตฺตกา จ ขนฺธา ฯเปฯ. กายายตนํ. (3)}

\prose{พาหิรํ ธมฺมํ ปจฺจยา พาหิโร ธมฺโม อุปฺปชฺชติ อารมฺมณ\-ปจฺจยา. พาหิรํ เอกํ ขนฺธํ ปจฺจยา เทฺว ขนฺธา. เทฺว ขนฺเธ ฯเปฯ. ปฏิสนฺธิ\-กฺขเณ ฯเปฯ. วตฺถุํ ปจฺจยา พาหิรา ขนฺธา.  พาหิรํ ธมฺมํ ปจฺจยา อชฺฌตฺติโก ธมฺโม อุปฺปชฺชติ อารมฺมณ\-ปจฺจยา. พาหิเร ขนฺเธ ปจฺจยา จิตฺตํ. วตฺถุํ ปจฺจยา จิตฺตํ. (ปฏิสนฺธิ\-กฺขเณ เทฺวปิ กาตพฺพา.). พาหิรํ ธมฺมํ ปจฺจยา อชฺฌตฺติโก จ พาหิโร จ ธมฺมา อุปฺปชฺชนฺติ อารมฺมณ\-ปจฺจยา. วตฺถุํ ปจฺจยา จิตฺตญฺจ สมฺปยุตฺตกา จ ขนฺธา. (ปฏิสนฺธิ\-กฺขเณ เอกํ กาตพฺพํ.) (3)}

\proseitem{335}{อชฺฌตฺติกญฺจ พาหิรญฺจ ธมฺมํ ปจฺจยา อชฺฌตฺติโก ธมฺโม อุปฺปชฺชติ อารมฺมณ\-ปจฺจยา. จกฺขุ\-วิญฺญาณสหคเต ขนฺเธ จ จกฺขายตนญฺจ ปจฺจยา จกฺขุวิญฺญาณํ ฯเปฯ. กาย\-วิญฺญาณสหคเต ฯเปฯ. อชฺฌตฺติกญฺจ พาหิรญฺจ ธมฺมํ ปจฺจยา พาหิโร ธมฺโม อุปฺปชฺชติ อารมฺมณ\-ปจฺจยา. จกฺขุ\-วิญฺญาณสหคตํ เอกํ ขนฺธญฺจ จกฺขายตนญฺจ จกฺขุวิญฺญาณญฺจ ปจฺจยา เทฺว ขนฺธา. เทฺว ขนฺเธ ฯเปฯ. กาย\-วิญฺญาณสหคตํ ฯเปฯ. พาหิรํ เอกํ ขนฺธญฺจ จิตฺตญฺจ ปจฺจยา เทฺว ขนฺธา. เทฺว ขนฺเธ ฯเปฯ. จิตฺตญฺจ วตฺถุญฺจ ปจฺจยา พาหิรา ขนฺธา. (ปฏิสนฺธิ\-กฺขเณ เทฺวปิ กาตพฺพา.) อชฺฌตฺติกญฺจ พาหิรญฺจ ธมฺมํ ปจฺจยา อชฺฌตฺติโก จ พาหิโร จ ธมฺมา อุปฺปชฺชนฺติ อารมฺมณ\-ปจฺจยา. จกฺขุ\-วิญฺญาณสหคตํ เอกํ ขนฺธญฺจ จกฺขายตนญฺจ ปจฺจยา เทฺว ขนฺธา จกฺขุวิญฺญาณญฺจ. เทฺว ขนฺเธ ฯเปฯ. (สํขิตฺตํ.) (3)}

\csromanmarkh{1. ปจฺจยานุโลม}\csromanmarkhii{2. สงฺขฺยาวาร}\csromanheadernomark{2}{1. \textbf{ปจฺจยานุโลม 2. สงฺขฺยาวาร}}
\proseitem{336}{เหตุยา นว, อารมฺมเณ นว, อธิปติยา ปญฺจ, อนนฺตเร นว, สมนนฺตเร นว, สหชาเต นว, อญฺญมญฺเญ นว, นิสฺสเย นว, อุปนิสฺสเย นว, ปุเรชาเต นว, อาเสวเน นว, กมฺเม นว, (สพฺพตฺถ นว.) อวิคเต นว.}

\csromanmarkh{2. ปจฺจย\-ปจฺจนีย}\csromanmarkhii{1. วิภงฺควาร}\csromanheadernomark{2}{2. ปจฺจย\-ปจฺจนีย 1. วิภงฺควาร}
\proseitem{337}{\csromanfolio{500}อชฺฌตฺติกํ ธมฺมํ ปจฺจยา อชฺฌตฺติโก ธมฺโม อุปฺปชฺชติ นเหตุปจฺจยา. อเหตุก\-ปฏิสนฺธิ\-กฺขเณ จิตฺตํ ปจฺจยา อชฺฌตฺติกํ กฏตฺตารูปํ. จกฺขายตนํ ปจฺจยา จกฺขุวิญฺญาณํ. (สํขิตฺตํ.) (เอวํ นวปิ ปญฺหา กาตพฺพา, ปญฺจวิญฺญาณมฺปิ ปเวเสตฺวา ตีณิเยว โมโห.)}

\csromanmarkh{2. ปจฺจย\-ปจฺจนีย}\csromanmarkhii{2. สงฺขฺยาวาร}\csromanheadernomark{2}{2. ปจฺจย\-ปจฺจนีย 2. สงฺขฺยาวาร}
\vspace{\tipitakaparskip}
\csromancenter{\textbf{นเหตุ}ยา นว, นอารมฺมเณ นว, นอธิปติยา นว, นอนนฺตเร นว, นสมนนฺตเร นว, นอญฺญมญฺเญ นว, นฺเอาปนิสฺสเย นว, นปุเรชาเต นว, นปจฺฉาชาเต นว, นอาเสวเน นว, นกมฺเม ตีณิ, นวิปาเก ปญฺจ, นอาหาเร เอกํ, ไนนฺทฺริเย เอกํ, นฌาเน นว, นมคฺเค นว, นสมฺปยุตฺเต นว, นวิปฺปยุตฺเต ปญฺจ, โนนตฺถิยา นว, โนวิคเต นว.}
\csromancenter{(เอวํ อิตเร เทฺว คณนาปิ นิสฺสยวาโรปิ กาตพฺโพ.)}
\csromanmarkh{66. อชฺฌตฺติกทุก}\csromanmarkhii{5. สํสฏฺฐวาร}\csromanheadernomark{2}{66. \textbf{อชฺฌตฺติกทุก 5. สํสฏฺฐวาร}}
\proseitem{339}{อชฺฌตฺติกํ ธมฺมํ สํสฏฺโฐ พาหิโร ธมฺโม อุปฺปชฺชติ เหตุปจฺจยา. จิตฺตํ สํสฏฺฐา สมฺปยุตฺตกา ขนฺธา. ปฏิสนฺธิ\-กฺขเณ จิตฺตํ สํสฏฺฐา สมฺปยุตฺตกา ขนฺธา. (1)}

\prose{พาหิรํ ธมฺมํ สํสฏฺโฐ พาหิโร ธมฺโม อุปฺปชฺชติ เหตุปจฺจยา. พาหิรํ เอกํ ขนฺธํ สํสฏฺฐา เทฺว ขนฺธา. เทฺว ขนฺเธ ฯเปฯ. ปฏิสนฺธิ\-กฺขเณ ฯเปฯ. พาหิรํ ธมฺมํ สํสฏฺโฐ อชฺฌตฺติโก ธมฺโม อุปฺปชฺชติ เหตุปจฺจยา. พาหิเร ขนฺเธ สํสฏฺฐํ จิตฺตํ. ปฏิสนฺธิ\-กฺขเณ ฯเปฯ. พาหิรํ ธมฺมํ สํสฏฺโฐ อชฺฌตฺติโก จ พาหิโร จ ธมฺมา อุปฺปชฺชนฺติ เหตุปจฺจยา. พาหิรํ เอกํ ขนฺธํ สํสฏฺฐา เทฺว ขนฺธา จิตฺตญฺจ. เทฺว ขนฺเธ ฯเปฯ. ปฏิสนฺธิ\-กฺขเณ ฯเปฯ. (3)}

\prose{\csromanfolio{501}อชฺฌตฺติกญฺจ พาหิรญฺจ ธมฺมํ สํสฏฺโฐ พาหิโร ธมฺโม อุปฺปชฺชติ เหตุปจฺจยา. พาหิรํ เอกํ ขนฺธญฺจ จิตฺตญฺจ สํสฏฺฐา เทฺว ขนฺธา. เทฺว ขนฺเธ ฯเปฯ. ปฏิสนฺธิ\-กฺขเณ ฯเปฯ. (สํขิตฺตํ.)}

\prose{เหตุยา ปญฺจ, อารมฺมเณ ปญฺจ, อธิปติยา ปญฺจ, (สพฺพตฺถ ปญฺจ.) อวิคเต ปญฺจ. (อนุโลมํ.)}

\prose{อชฺฌตฺติกํ ธมฺมํ สํสฏฺโฐ พาหิโร ธมฺโม อุปฺปชฺชติ นเหตุปจฺจยา. (เอวํ ปญฺจ กาตพฺพา, ตีณิเยว โมโห.)}

\prose{นเหตุยา ปญฺจ, นอธิปติยา ปญฺจ, นปุเรชาเต ปญฺจ, นปจฺฉาชาเต ปญฺจ, นอาเสวเน ปญฺจ, นกมฺเม ตีณิ, นวิปาเก ปญฺจ, นฌาเน ปญฺจ, นมคฺเค ปญฺจ, นวิปฺปยุตฺเต ปญฺจ. (ปจฺจนียํ.)}

\prose{(เอวํ อิตเร เทฺว ค คณนาปิ สมฺปยุตฺตวาโรปิ กาตพฺโพ.)}

\csromanmarkh{66. อชฺฌตฺติกทุก}\csromanmarkhii{7. ปญฺหาวาร}\csromanheadernomark{2}{66. อชฺฌตฺติกทุก 7. ปญฺหาวาร}
\csromanmarkh{1. ปจฺจยานุโลม}\csromanmarkhii{1. วิภงฺควาร}\csromanheadernomark{2}{1. \textbf{ปจฺจยานุโลม 1. วิภงฺควาร}}
\vspace{\tipitakaparskip}
\csromancenter{\textbf{เหตุ}}
\proseitem{340}{พาหิโร ธมฺโม พาหิรสฺส ธมฺมสฺส เหตุปจฺจเยน ปจฺจโย. พาหิรา เหตู สมฺปยุตฺตกานํ ขนฺธานํ จิตฺตสมุฏฺฐานานญฺจ รูปานํ เหตุปจฺจเยน ปจฺจโย. ปฏิสนฺธิ\-กฺขเณ พาหิรา เหตู สมฺปยุตฺตกานํ ขนฺธานํ พาหิรานญฺจ กฏตฺตารูปานํ เหตุปจฺจเยน ปจฺจโย. (1)}

\prose{พาหิโร ธมฺโม อชฺฌตฺติกสฺส ธมฺมสฺส เหตุปจฺจเยน ปจฺจโย. พาหิรา เหตู จิตฺตสฺส เหตุปจฺจเยน ปจฺจโย. ปฏิสนฺธิ\-กฺขเณ พาหิรา เหตู จิตฺตสฺส อชฺฌตฺติกานญฺจ กฏตฺตารูปานํ เหตุปจฺจเยน ปจฺจโย. (2)}

\prose{พาหิโร ธมฺโม อชฺฌตฺติกสฺส จ พาหิรสฺส จ ธมฺมสฺส เหตุปจฺจเยน ปจฺจโย. พาหิรา เหตู สมฺปยุตฺตกานํ ขนฺธานํ จิตฺตสฺส จ จิตฺตสมุฏฺฐานานญฺจ รูปานํ เหตุปจฺจเยน ปจฺจโย. ปฏิสนฺธิ\-กฺขเณ พาหิรา เหตู สมฺปยุตฺตกานํ ขนฺธานํ จิตฺตสฺส จ อชฺฌตฺติกานญฺจ พาหิรานญฺจ กฏตฺตารูปานํ เหตุปจฺจเยน ปจฺจโย. (3)}

\proseitem{341}{\csromanfolio{502}อชฺฌตฺติโก ธมฺโม อชฺฌตฺติกสฺส ธมฺมสฺส อารมฺมณ\-ปจฺจเยน ปจฺจโย. จิตฺตํ อารพฺภ จิตฺตํ อุปฺปชฺชติ. (มูลํ ปุจฺฉิตพฺพํ.) จิตฺตํ อารพฺภ พาหิรา ขนฺธา อุปฺปชฺชนฺติ. (มูลํ ปุจฺฉิตพฺพํ.) จิตฺตํ อารพฺภ จิตฺตญฺจ สมฺปยุตฺตกา จ ขนฺธา อุปฺปชฺชนฺติ. (3)}

\prose{พาหิโร ธมฺโม พาหิรสฺส ธมฺมสฺส อารมฺมณ\-ปจฺจเยน ปจฺจโย. ทานํ ฯเปฯ สีลํ ฯเปฯ อุโปสถกมฺมํ กตฺวา ตํ ปจฺจเวกฺขติ, อสฺสาเทติ อภินนฺทติ, ตํ อารพฺภ ราโค ฯเปฯ โทมนสฺสํ อุปฺปชฺชติ. ปุพฺเพ สุจิณฺณานิ ปจฺจเวกฺขติ. ฌานา วุฏฺฐหิตฺวา ฌานํ ฯเปฯ. อริยา มคฺคา วุฏฺฐหิตฺวา มคฺคํ ฯเปฯ. ผลํ ฯเปฯ. นิพฺพานํ ปจฺจเวกฺขนฺติ. นิพฺพานํ โคตฺรภุสฺส, โวทานสฺส, มคฺคสฺส, ผลสฺส, อาวชฺชนาย อารมฺมณ\-ปจฺจเยน ปจฺจโย. อริยา พาหิเร ปหีเน กิเลเส ปจฺจเวกฺขนฺติ. วิกฺขมฺภิเต กิเลเส ฯเปฯ. ปุพฺเพ สมุทาจิณฺเณ กิเลเส ชานนฺติ. รูเป ฯเปฯ. วตฺถุํ พาหิเร ขนฺเธ อนิจฺจโต ฯเปฯ โทมนสฺสํ อุปฺปชฺชติ. ทิพฺเพน จกฺขุนา รูปํ ปสฺสติ. ทิพฺพาย โสตธาตุยา สทฺทํ สุณาติ. เจโตปริย\-ญาเณน พาหิรจิตฺตสมงฺคสฺส จิตฺตํ ชานาติ. อากาสา\-นญฺจา\-ยตนํ วิญฺญาณญฺจา\-ยตนสฺส ฯเปฯ. อากิญฺจญฺญา\-ยตนํ เนว\-สญฺญา\-นา\-สญฺญา\-ยตนสฺส ฯเปฯ. รูปายตนํ จกฺขุ\-วิญฺญาณสหคตานํ ขนฺธานํ อารมฺมณ\-ปจฺจเยน ปจฺจโย ฯเปฯ. โผฏฺฐพฺพา\-ยตนํ. พาหิรา ขนฺธา อิทฺธิวิธ\-ญาณสฺส, เจโตปริย\-ญาณสฺส, ปุพฺเพ\-นิวาสา\-นุสฺสติ\-ญาณสฺส, ยถากมฺมูปคญาณสฺส, อนาคตํสญาณสฺส, อาวชฺชนาย อารมฺมณ\-ปจฺจเยน ปจฺจโย. (1)}

\prose{พาหิโร ธมฺโม อชฺฌตฺติกสฺส ธมฺมสฺส อารมฺมณ\-ปจฺจเยน ปจฺจโย. ทานํ ฯเปฯ สีลํ ฯเปฯ อุโปสถกมฺมํ กตฺวา ตํ ปจฺจเวกฺขติ, อสฺสาเทติ อภินนฺทติ, ตํ อารพฺภ จิตฺตํ อุปฺปชฺชติ. ปุพฺเพ สุจิณฺณานิ ฯเปฯ. ฌานํ ฯเปฯ. (สํขิตฺตํ. สพฺพํ กาตพฺพํ.) ปุพฺเพ สมุทาจิณฺเณ ฯเปฯ. รูเป ฯเปฯ. วตฺถุํ พาหิเร ขนฺเธ อนิจฺจโต ฯเปฯ วิปสฺสติ, อสฺสาเทติ อภินนฺทติ, ตํ อารพฺภ จิตฺตํ อุปฺปชฺชติ. ทิพฺเพน จกฺขุนา รูปํ ปสฺสติ ฯเปฯ. รูปายตนํ จกฺขุ\-วิญฺญาณสฺส ฯเปฯ. โผฏฺฐพฺพา\-ยตนํ ฯเปฯ. พาหิรา ขนฺธา อิทฺธิวิธ\-ญาณสฺส, เจโตปริย\-ญาณสฺส, ปุพฺเพ\-นิวาสา\-นุสฺสติ\-ญาณสฺส, ยถากมฺมูปคญาณสฺส, อนาคตํสญาณสฺส, อาวชฺชนาย อารมฺมณ\-ปจฺจเยน ปจฺจโย. (2)}

\prose{\csromanfolio{503}พาหิโร ธมฺโม อชฺฌตฺติกสฺส จ พาหิรสฺส จ ธมฺมสฺส อารมฺมณ\-ปจฺจเยน ปจฺจโย. ทานํ ฯเปฯ สีลํ ฯเปฯ อุโปสถกมฺมํ กตฺวา ตํ ปจฺจเวกฺขติ, อสฺสาเทติ อภินนฺทติ, ตํ อารพฺภ จิตฺตญฺจ สมฺปยุตฺตกา จ ขนฺธา อุปฺปชฺชนฺติ. (สํขิตฺตํ. สพฺพํ กาตพฺพํ.) พาหิเร ขนฺเธ อนิจฺจโต ฯเปฯ วิปสฺสติ, อสฺสาเทติ อภินนฺทติ, ตํ อารพฺภ จิตฺตญฺจ สมฺปยุตฺตกา จ ขนฺธา อุปฺปชฺชนฺติ. ทิพฺเพน จกฺขุนา รูปํ ปสฺสติ ฯเปฯ. รูปายตนํ จกฺขุ\-วิญฺญาณสฺส จ สมฺปยุตฺตกานญฺจ ขนฺธานํ ฯเปฯ. โผฏฺฐพฺพา\-ยตนํ ฯเปฯ. พาหิรา ขนฺธา อิทฺธิวิธ\-ญาณสฺส, เจโตปริย\-ญาณสฺส, ปุพฺเพ\-นิวาสา\-นุสฺสติ\-ญาณสฺส, ยถากมฺมูปคญาณสฺส, อนาคตํสญาณสฺส, อาวชฺชนาย อารมฺมณ\-ปจฺจเยน ปจฺจโย. (3)}

\prose{อชฺฌตฺติโก จ พาหิโร จ ธมฺมา อชฺฌตฺติกสฺส ธมฺมสฺส อารมฺมณ\-ปจฺจเยน ปจฺจโย. ตีณิ.}

\proseitem{342}{อชฺฌตฺติโก ธมฺโม อชฺฌตฺติกสฺส ธมฺมสฺส อธิปติ\-ปจฺจเยน ปจฺจโย. อารมฺมณา\-ธิปติ\csromandash{}จิตฺตํ ครุํ กตฺวา จิตฺตํ อุปฺปชฺชติ. (1)}

\prose{อชฺฌตฺติโก ธมฺโม พาหิรสฺส ธมฺมสฺส อธิปติ\-ปจฺจเยน ปจฺจโย. อารมฺมณา\-ธิปติ, สหชาตา\-ธิปติ.  อารมฺมณา\-ธิปติ\csromandash{}จิตฺตํ ครุํ กตฺวา พาหิรา ขนฺธา อุปฺปชฺชนฺติ. สหชาตา\-ธิปติ\csromandash{}จิตฺตาธิปติ สมฺปยุตฺตกานํ ขนฺธานํ จิตฺตสมุฏฺฐานานญฺจ รูปานํ อธิปติ\-ปจฺจเยน ปจฺจโย. (มูลํ.) อารมฺมณา\-ธิปติ\csromandash{}อชฺฌตฺติกํ จิตฺตํ ครุํ กตฺวา จิตฺตญฺจ สมฺปยุตฺตกา จ ขนฺธา อุปฺปชฺชนฺติ. (3)}

\prose{พาหิโร ธมฺโม พาหิรสฺส ธมฺมสฺส อธิปติ\-ปจฺจเยน ปจฺจโย. อารมฺมณา\-ธิปติ, สหชาตา\-ธิปติ.  อารมฺมณา\-ธิปติ\csromandash{}ทานํ ทตฺวา ฯเปฯ. ตีณิ. (เทฺว อธิปตี, ติณฺณมฺปิ กาตพฺพา.) (3)}

\prose{อชฺฌตฺติโก จ พาหิโร จ ธมฺมา อชฺฌตฺติกสฺส ธมฺมสฺส อธิปติ\-ปจฺจเยน ปจฺจโย. ตีณิ. (ติณฺณมฺปิ เอกาเยว อธิปติ.)}

\proseitem{343}{\csromanfolio{504}อชฺฌตฺติโก ธมฺโม อชฺฌตฺติกสฺส ธมฺมสฺส อนนฺตร\-ปจฺจเยน ปจฺจโย. ปุริมํ ปุริมํ จิตฺตํ ปจฺฉิมสฺส ปจฺฉิมสฺส จิตฺตสฺส อนนฺตร\-ปจฺจเยน ปจฺจโย. ตีณิ.}

\prose{พาหิโร ธมฺโม พาหิรสฺส ธมฺมสฺส อนนฺตร\-ปจฺจเยน ปจฺจโย. ปุริมา ปุริมา พาหิรา ขนฺธา ปจฺฉิมานํ ปจฺฉิมานํ ขนฺธานํ อนนฺตร\-ปจฺจเยน ปจฺจโย. อนุโลมํ โคตฺรภุสฺส ฯเปฯ. ตีณิ. (ติณฺณมฺปิ เอกสทิสา.)}

\prose{อชฺฌตฺติโก จ พาหิโร จ ธมฺมา อชฺฌตฺติกสฺส ธมฺมสฺส อนนฺตร\-ปจฺจเยน ปจฺจโย. ตีณิ. สมนนฺตร\-ปจฺจเยน ปจฺจโย. นว. สหชาต\-ปจฺจเยน ปจฺจโย. นว. (ปฏิจฺจสทิสา.) อญฺญมญฺญ\-ปจฺจเยน ปจฺจโย. ปญฺจ. (ปฏิจฺจสทิสา.) นิสฺสย\-ปจฺจเยน ปจฺจโย. นว. (ปจฺจยวาร\-สทิสา.)}

\csromanheader{2}{\textbf{อุปนิสฺสย}}
\proseitem{344}{อชฺฌตฺติโก ธมฺโม อชฺฌตฺติกสฺส ธมฺมสฺส อุปนิสฺสย\-ปจฺจเยน ปจฺจโย. อารมฺมณู\-ปนิสฺสโย, อนนฺตรู\-ปนิสฺสโย. ปกตู\-ปนิสฺสโย ฯเปฯ. ปกตู\-ปนิสฺสโย\csromandash{}จิตฺตํ จิตฺตสฺส อุปนิสฺสย\-ปจฺจเยน ปจฺจโย. ตีณิ.}

\prose{พาหิโร ธมฺโม พาหิรสฺส ธมฺมสฺส อุปนิสฺสย\-ปจฺจเยน ปจฺจโย. อารมฺมณู\-ปนิสฺสโย, อนนฺตรู\-ปนิสฺสโย, ปกตู\-ปนิสฺสโย ฯเปฯ. ปกตู\-ปนิสฺสโย\csromandash{}สทฺธํ อุปนิสฺสาย ทานํ เทติ ฯเปฯ มานํ ชปฺเปติ, ทิฏฺฐิํ คณฺหาติ. สีลํ ฯเปฯ. เสนาสนํ อุปนิสฺสาย ทานํ เทติ ฯเปฯ สํฆํ ภินฺทติ. สทฺธา ฯเปฯ. เสนาสนํ สทฺธาย ฯเปฯ ผลสมาปตฺติยา อุปนิสฺสย\-ปจฺจเยน ปจฺจโย.}

\vspace{\tipitakaparskip}
\csromancenter{(ตีณิปิ ปูเรตฺวา กาตพฺพา. “จิตฺตสฺสา”ติ กาตพฺพา. “สมฺปยุตฺตกานญฺจา”ติ กาตพฺพา.)}
\prose{อชฺฌตฺติโก จ พาหิโร จ ธมฺมา อชฺฌตฺติกสฺส ธมฺมสฺส อุปนิสฺสย\-ปจฺจเยน ปจฺจโย. ตีณิ.}

\proseitem{345}{อชฺฌตฺติโก ธมฺโม อชฺฌตฺติกสฺส ธมฺมสฺส ปุเรชาต\-ปจฺจเยน ปจฺจโย. อารมฺมณ\-ปุเรชาตํ, วตฺถุ\-ปุเรชาตํ.  อารมฺมณ\-ปุเรชาตํ\csromandash{} \csromanfolio{505}จกฺขุํ ฯเปฯ กายํ อนิจฺจโต ฯเปฯ วิปสฺสติ, อสฺสาเทติ อภินนฺทติ, ตํ อารพฺภ จิตฺตํ อุปฺปชฺชติ. วตฺถุ\-ปุเรชาตํ\csromandash{}จกฺขายตนํ จกฺขุ\-วิญฺญาณสฺส ฯเปฯ กายายตนํ กายวิญฺญาณสฺส ปุเรชาต\-ปจฺจเยน ปจฺจโย. (1)}

\prose{อชฺฌตฺติโก ธมฺโม พาหิรสฺส ธมฺมสฺส ปุเรชาต\-ปจฺจเยน ปจฺจโย. อารมฺมณ\-ปุเรชาตํ, วตฺถุ\-ปุเรชาตํ.  อารมฺมณ\-ปุเรชาตํ\csromandash{}จกฺขุํ ฯเปฯ กายํ อนิจฺจโต ฯเปฯ วิปสฺสติ, อสฺสาเทติ ฯเปฯ โทมนสฺสํ อุปฺปชฺชติ. วตฺถุ\-ปุเรชาตํ\csromandash{}จกฺขายตนํ จกฺขุ\-วิญฺญาณสหคตานํ ขนฺธานํ ฯเปฯ กายายตนํ กาย\-วิญฺญาณสหคตานํ ขนฺธานํ ปุเรชาต\-ปจฺจเยน ปจฺจโย. (2)}

\prose{อชฺฌตฺติโก ธมฺโม อชฺฌตฺติกสฺส จ พาหิรสฺส จ ธมฺมสฺส ปุเรชาต\-ปจฺจเยน ปจฺจโย. อารมฺมณ\-ปุเรชาตํ, วตฺถุ\-ปุเรชาตํ.  อารมฺมณ\-ปุเรชาตํ\csromandash{}จกฺขุํ ฯเปฯ กายํ อนิจฺจโต ฯเปฯ วิปสฺสติ, ตํ อารพฺภ จิตฺตญฺจ สมฺปยุตฺตกา ขนฺธา จ อุปฺปชฺชนฺติ. วตฺถุ\-ปุเรชาตํ\csromandash{} จกฺขายตนํ จกฺขุ\-วิญฺญาณสฺส สมฺปยุตฺตกานญฺจ ขนฺธานํ ฯเปฯ กายายตนํ กายวิญฺญาณสฺส สมฺปยุตฺตกานญฺจ ขนฺธานํ ปุเรชาต\-ปจฺจเยน ปจฺจโย. (3)}

\proseitem{346}{พาหิโร ธมฺโม พาหิรสฺส ธมฺมสฺส ปุเรชาต\-ปจฺจเยน ปจฺจโย. อารมฺมณ\-ปุเรชาตํ, วตฺถุ\-ปุเรชาตํ.  อารมฺมณ\-ปุเรชาตํ\csromandash{}รูเป ฯเปฯ โผฏฺฐพฺเพ. วตฺถุํ อนิจฺจโต ฯเปฯ โทมนสฺสํ อุปฺปชฺชติ. ทิพฺเพน จกฺขุนา รูปํ ปสฺสติ. ทิพฺพาย โสตธาตุยา สทฺทํ สุณาติ. วตฺถุ\-ปุเรชาตํ\csromandash{}วตฺถุ พาหิรานํ ขนฺธานํ ปุเรชาต\-ปจฺจเยน ปจฺจโย. (1)}

\prose{พาหิโร ธมฺโม อชฺฌตฺติกสฺส ธมฺมสฺส ปุเรชาต\-ปจฺจเยน ปจฺจโย. อารมฺมณ\-ปุเรชาตํ, วตฺถุ\-ปุเรชาตํ.  อารมฺมณ\-ปุเรชาตํ\csromandash{}รูเป ฯเปฯ โผฏฺฐพฺเพ. วตฺถุํ อนิจฺจโต ฯเปฯ ตํ อารพฺภ จิตฺตํ อุปฺปชฺชติ. ทิพฺเพน จกฺขุนา รูปํ ปสฺสติ. ทิพฺพาย โสตธาตุยา สทฺทํ สุณาติ. วตฺถุ\-ปุเรชาตํ\csromandash{}วตฺถุ จิตฺตสฺส ปุเรชาต\-ปจฺจเยน ปจฺจโย. (2)}

\prose{พาหิโร ธมฺโม อชฺฌตฺติกสฺส จ พาหิรสฺส จ ธมฺมสฺส ปุเรชาต\-ปจฺจเยน ปจฺจโย. อารมฺมณ\-ปุเรชาตํ, วตฺถุ\-ปุเรชาตํ.  อารมฺมณ\-ปุเรชาตํ\csromandash{}รูเป ฯเปฯ โผฏฺฐพฺเพ. วตฺถุํ อนิจฺจโต ฯเปฯ ตํ อารพฺภ จิตฺตญฺจ สมฺปยุตฺตกา ขนฺธา จ อุปฺปชฺชนฺติ. ทิพฺเพน จกฺขุนา รูปํ ปสฺสติ. ทิพฺพาย โสตธาตุยา \csromanfolio{506}สทฺทํ สุณาติ. วตฺถุ\-ปุเรชาตํ\csromandash{}วตฺถุ จิตฺตสฺส สมฺปยุตฺตกานญฺจ ขนฺธานํ ปุเรชาต\-ปจฺจเยน ปจฺจโย. (3)}

\proseitem{347}{อชฺฌตฺติโก จ พาหิโร จ ธมฺมา อชฺฌตฺติกสฺส ธมฺมสฺส ปุเรชาต\-ปจฺจเยน ปจฺจโย. อารมฺมณ\-ปุเรชาตํ, วตฺถุ\-ปุเรชาตํ. จกฺขายตนญฺจ วตฺถุ จ จิตฺตสฺส ฯเปฯ. กายายตนญฺจ วตฺถุ จ จิตฺตสฺส ปุเรชาต\-ปจฺจเยน ปจฺจโย. รูปายตนญฺจ จกฺขายตนญฺจ จกฺขุ\-วิญฺญาณสฺส ฯเปฯ. โผฏฺฐพฺพายตนญฺจ กายายตนญฺจ กายวิญฺญาณสฺส ปุเรชาต\-ปจฺจเยน ปจฺจโย. (1)}

\prose{อชฺฌตฺติโก จ พาหิโร จ ธมฺมา พาหิรสฺส ธมฺมสฺส ปุเรชาต\-ปจฺจเยน ปจฺจโย. อารมฺมณ\-ปุเรชาตํ, วตฺถุ\-ปุเรชาตํ. จกฺขายตนญฺจ วตฺถุ จ พาหิรานํ ขนฺธานํ ปุเรชาต\-ปจฺจเยน ปจฺจโย ฯเปฯ. กายายตนญฺจ วตฺถุ จ พาหิรานํ ขนฺธานํ ปุเรชาต\-ปจฺจเยน ปจฺจโย. รูปายตนญฺจ จกฺขายตนญฺจ จกฺขุ\-วิญฺญาณสหคตานํ ขนฺธานํ ปุเรชาต\-ปจฺจเยน ปจฺจโย ฯเปฯ. โผฏฺฐพฺพายตนญฺจ กายายตนญฺจ กาย\-วิญฺญาณสหคตานํ ขนฺธานํ ปุเรชาต\-ปจฺจเยน ปจฺจโย. (2)}

\prose{อชฺฌตฺติโก จ พาหิโร จ ธมฺมา อชฺฌตฺติกสฺส จ พาหิรสฺส จ ธมฺมสฺส ปุเรชาต\-ปจฺจเยน ปจฺจโย. อารมฺมณ\-ปุเรชาตํ, วตฺถุ\-ปุเรชาตํ. จกฺขายตนญฺจ วตฺถุ จ จิตฺตสฺส สมฺปยุตฺตกานญฺจ ขนฺธานํ ปุเรชาต\-ปจฺจเยน ปจฺจโย ฯเปฯ. กายายตนญฺจ วตฺถุ จ ฯเปฯ. รูปายตนญฺจ จกฺขายตนญฺจ จกฺขุ\-วิญฺญาณสฺส สมฺปยุตฺตกานญฺจ ขนฺธานํ ปุเรชาต\-ปจฺจเยน ปจฺจโย ฯเปฯ. โผฏฺฐพฺพายตนญฺจ ฯเปฯ. (3)}

\proseitem{348}{อชฺฌตฺติโก ธมฺโม อชฺฌตฺติกสฺส ธมฺมสฺส ปจฺฉาชาต\-ปจฺจเยน ปจฺจโย. ปจฺฉาชาตา อชฺฌตฺติกา ขนฺธา ปุเรชาตสฺส อิมสฺส อชฺฌตฺติกสฺส กายสฺส ปจฺฉาชาต\-ปจฺจเยน ปจฺจโย. (มูลํ กาตพฺพํ.) ปจฺฉาชาตา อชฺฌตฺติกา ขนฺธา ปุเรชาตสฺส อิมสฺส พาหิรสฺส กายสฺส ปจฺฉาชาต\-ปจฺจเยน ปจฺจโย. (มูลํ กาตพฺพํ.) ปจฺฉาชาตา อชฺฌตฺติกา ขนฺธา ปุเรชาตสฺส อิมสฺส อชฺฌตฺติกสฺส จ พาหิรสฺส จ \csromanfolio{507}กายสฺส ปจฺฉาชาต\-ปจฺจเยน ปจฺจโย. (เอวํ นวปิ ปญฺหา กาตพฺพา.) อาเสวน\-ปจฺจเยน ปจฺจโย. (นว ปญฺหา กาตพฺพา.)}

\csromanheader{2}{\textbf{กมฺม วิปาก}}
\proseitem{349}{พาหิโร ธมฺโม พาหิรสฺส ธมฺมสฺส กมฺมปจฺจเยน ปจฺจโย. สหชาตา, นานากฺขณิกา.  สหชาตา พาหิรา เจตนา สมฺปยุตฺตกานํ ขนฺธานํ จิตฺตสมุฏฺฐานานญฺจ รูปานํ กมฺมปจฺจเยน ปจฺจโย. ปฏิสนฺธิ\-กฺขเณ ฯเปฯ. นานากฺขณิกา พาหิรา เจตนา วิปากานํ พาหิรานํ ขนฺธานํ กฏตฺตา จ รูปานํ กมฺมปจฺจเยน ปจฺจโย. (1)}

\prose{พาหิโร ธมฺโม อชฺฌตฺติกสฺส ธมฺมสฺส กมฺมปจฺจเยน ปจฺจโย. สหชาตา, นานากฺขณิกา.  สหชาตา พาหิรา เจตนา จิตฺตสฺส กมฺมปจฺจเยน ปจฺจโย. ปฏิสนฺธิ\-กฺขเณ ฯเปฯ. นานากฺขณิกา พาหิรา เจตนา วิปากสฺส จิตฺตสฺส อชฺฌตฺติกานญฺจ กฏตฺตารูปานํ กมฺมปจฺจเยน ปจฺจโย. (2)}

\prose{พาหิโร ธมฺโม อชฺฌตฺติกสฺส จ พาหิรสฺส จ ธมฺมสฺส กมฺมปจฺจเยน ปจฺจโย. สหชาตา, นานากฺขณิกา.  สหชาตา พาหิรา เจตนา สมฺปยุตฺตกานํ ขนฺธานํ จิตฺตสฺส จ จิตฺตสมุฏฺฐานานญฺจ รูปานํ กมฺมปจฺจเยน ปจฺจโย. ปฏิสนฺธิ\-กฺขเณ ฯเปฯ. นานากฺขณิกา พาหิรา เจตนา วิปากานํ ขนฺธานํ จิตฺตสฺส จ อชฺฌตฺติกานญฺจ พาหิรานญฺจ กฏตฺตารูปานํ กมฺมปจฺจเยน ปจฺจโย. (3) วิปาก\-ปจฺจเยน ปจฺจโย. นว.}

\csromanheader{2}{\textbf{อาหาร}}
\proseitem{350}{อชฺฌตฺติโก ธมฺโม อชฺฌตฺติกสฺส ธมฺมสฺส อาหารปจฺจเยน ปจฺจโย. ปฏิสนฺธิ\-กฺขเณ อชฺฌตฺติกา อาหารา อชฺฌตฺติกานํ กฏตฺตารูปานํ อาหารปจฺจเยน ปจฺจโย. (1)}

\prose{อชฺฌตฺติโก ธมฺโม พาหิรสฺส ธมฺมสฺส อาหารปจฺจเยน ปจฺจโย. อชฺฌตฺติกา อาหารา สมฺปยุตฺตกานํ ขนฺธานํ จิตฺตสมุฏฺฐานานญฺจ รูปานํ อาหารปจฺจเยน ปจฺจโย. ปฏิสนฺธิ\-กฺขเณ อชฺฌตฺติกา อาหารา สมฺปยุตฺตกานํ ขนฺธานํ พาหิรานญฺจ กฏตฺตารูปานํ อาหารปจฺจเยน ปจฺจโย. (มูลํ กาตพฺพํ.) ปฏิสนฺธิ\-กฺขเณ อชฺฌตฺติกา อาหารา \csromanfolio{508}สมฺปยุตฺตกานํ ขนฺธานํ อชฺฌตฺติกานญฺจ พาหิรานญฺจ กฏตฺตารูปานํ อาหารปจฺจเยน ปจฺจโย. (3)}

\proseitem{351}{พาหิโร ธมฺโม พาหิรสฺส ธมฺมสฺส อาหารปจฺจเยน ปจฺจโย. พาหิรา อาหารา สมฺปยุตฺตกานํ ขนฺธานํ จิตฺตสมุฏฺฐานานญฺจ รูปานํ อาหารปจฺจเยน ปจฺจโย. ปฏิสนฺธิ\-กฺขเณ ฯเปฯ. พาหิโร กพฬีกาโร อาหาโร พาหิรสฺส กายสฺส อาหารปจฺจเยน ปจฺจโย. (1)}

\prose{พาหิโร ธมฺโม อชฺฌตฺติกสฺส ธมฺมสฺส อาหารปจฺจเยน ปจฺจโย. พาหิรา อาหารา จิตฺตสฺส อาหารปจฺจเยน ปจฺจโย. ปฏิสนฺธิ\-กฺขเณ พาหิรา อาหารา จิตฺตสฺส อชฺฌตฺติกานญฺจ กฏตฺตารูปานํ อาหารปจฺจเยน ปจฺจโย. พาหิโร กพฬีกาโร อาหาโร อชฺฌตฺติกสฺส กายสฺส อาหารปจฺจเยน ปจฺจโย. (2)}

\prose{พาหิโร ธมฺโม อชฺฌตฺติกสฺส จ พาหิรสฺส จ ธมฺมสฺส อาหารปจฺจเยน ปจฺจโย. พาหิรา อาหารา สมฺปยุตฺตกานํ ขนฺธานํ จิตฺตสฺส จ จิตฺตสมุฏฺฐานานญฺจ รูปานํ อาหารปจฺจเยน ปจฺจโย. ปฏิสนฺธิ\-กฺขเณ พาหิรา อาหารา สมฺปยุตฺตกานํ ขนฺธานํ จิตฺตสฺส จ อชฺฌตฺติกานญฺจ พาหิรานญฺจ กฏตฺตารูปานํ อาหารปจฺจเยน ปจฺจโย. พาหิโร กพฬีกาโร อาหาโร อชฺฌตฺติกสฺส จ พาหิรสฺส จ กายสฺส อาหารปจฺจเยน ปจฺจโย. (3)}

\proseitem{352}{อชฺฌตฺติโก จ พาหิโร จ ธมฺมา อชฺฌตฺติกสฺส ธมฺมสฺส อาหารปจฺจเยน ปจฺจโย. ปฏิสนฺธิ\-กฺขเณ อชฺฌตฺติกา จ พาหิรา จ อาหารา อชฺฌตฺติกานํ กฏตฺตารูปานํ อาหารปจฺจเยน ปจฺจโย. (1)}

\prose{อชฺฌตฺติโก จ พาหิโร จ ธมฺมา พาหิรสฺส ธมฺมสฺส อาหารปจฺจเยน ปจฺจโย. อชฺฌตฺติกา จ พาหิรา จ อาหารา สมฺปยุตฺตกานํ ขนฺธานํ จิตฺตสมุฏฺฐานานญฺจ รูปานํ อาหารปจฺจเยน ปจฺจโย. ปฏิสนฺธิ\-กฺขเณ อชฺฌตฺติกา จ พาหิรา จ อาหารา สมฺปยุตฺตกานํ ขนฺธานํ พาหิรานญฺจ กฏตฺตารูปานํ อาหารปจฺจเยน ปจฺจโย. (2)}

\prose{อชฺฌตฺติโก จ พาหิโร จ ธมฺมา อชฺฌตฺติกสฺส จ พาหิรสฺส จ ธมฺมสฺส อาหารปจฺจเยน ปจฺจโย. ปฏิสนฺธิ\-กฺขเณ อชฺฌตฺติกา จ พาหิรา จ \csromanfolio{509}อาหารา สมฺปยุตฺตกานํ ขนฺธานํ อชฺฌตฺติกานญฺจ พาหิรานญฺจ กฏตฺตารูปานํ อาหารปจฺจเยน ปจฺจโย. (3)}

\csromanheader{2}{\textbf{อินฺทฺริย}}
\proseitem{353}{อชฺฌตฺติโก ธมฺโม อชฺฌตฺติกสฺส ธมฺมสฺส อินฺทฺริย\-ปจฺจเยน ปจฺจโย. ปฏิสนฺธิ\-กฺขเณ อชฺฌตฺติกา อินฺทฺริยา อชฺฌตฺติกานํ กฏตฺตารูปานํ อินฺทฺริย\-ปจฺจเยน ปจฺจโย. จกฺขุนฺทฺริยํ จกฺขุ\-วิญฺญาณสฺส ฯเปฯ. กายินฺทฺริยํ กายวิญฺญาณสฺส อินฺทฺริย\-ปจฺจเยน ปจฺจโย. (1)}

\prose{อชฺฌตฺติโก ธมฺโม พาหิรสฺส ธมฺมสฺส อินฺทฺริย\-ปจฺจเยน ปจฺจโย. อชฺฌตฺติกา อินฺทฺริยา สมฺปยุตฺตกานํ ขนฺธานํ จิตฺตสมุฏฺฐานานญฺจ รูปานํ อินฺทฺริย\-ปจฺจเยน ปจฺจโย. ปฏิสนฺธิ\-กฺขเณ อชฺฌตฺติกา อินฺทฺริยา สมฺปยุตฺตกานํ ขนฺธานํ พาหิรานญฺจ กฏตฺตารูปานํ อินฺทฺริย\-ปจฺจเยน ปจฺจโย. จกฺขุนฺทฺริยํ จกฺขุ\-วิญฺญาณสหคตานํ ขนฺธานํ ฯเปฯ. กายินฺทฺริยํ กาย\-วิญฺญาณสหคตานํ ขนฺธานํ อินฺทฺริย\-ปจฺจเยน ปจฺจโย. (2)}

\prose{อชฺฌตฺติโก ธมฺโม อชฺฌตฺติกสฺส จ พาหิรสฺส จ ธมฺมสฺส อินฺทฺริย\-ปจฺจเยน ปจฺจโย. ปฏิสนฺธิ\-กฺขเณ อชฺฌตฺติกา อินฺทฺริยา สมฺปยุตฺตกานํ ขนฺธานํ อชฺฌตฺติกานญฺจ พาหิรานญฺจ กฏตฺตารูปานํ อินฺทฺริย\-ปจฺจเยน ปจฺจโย. จกฺขุนฺทฺริยํ จกฺขุ\-วิญฺญาณสฺส สมฺปยุตฺตกานญฺจ ขนฺธานํ ฯเปฯ. กายินฺทฺริยํ กายวิญฺญาณสฺส สมฺปยุตฺตกานญฺจ ขนฺธานํ อินฺทฺริย\-ปจฺจเยน ปจฺจโย. (3)}

\proseitem{354}{พาหิโร ธมฺโม พาหิรสฺส ธมฺมสฺส อินฺทฺริย\-ปจฺจเยน ปจฺจโย. พาหิรา อินฺทฺริยา สมฺปยุตฺตกานํ ขนฺธานํ จิตฺตสมุฏฺฐานานญฺจ รูปานํ อินฺทฺริย\-ปจฺจเยน ปจฺจโย. ปฏิสนฺธิ\-กฺขเณ พาหิรา อินฺทฺริยา สมฺปยุตฺตกานํ ขนฺธานํ พาหิรานญฺจ กฏตฺตารูปานํ อินฺทฺริย\-ปจฺจเยน ปจฺจโย. รูปชีวิตินฺทฺริยํ พาหิรานํ กฏตฺตารูปานํ อินฺทฺริย\-ปจฺจเยน ปจฺจโย. (1)}

\prose{พาหิโร ธมฺโม อชฺฌตฺติกสฺส ธมฺมสฺส อินฺทฺริย\-ปจฺจเยน ปจฺจโย. พาหิรา อินฺทฺริยา จิตฺตสฺส อินฺทฺริย\-ปจฺจเยน ปจฺจโย. ปฏิสนฺธิ\-กฺขเณ พาหิรา อินฺทฺริยา จิตฺตสฺส อชฺฌตฺติกานญฺจ กฏตฺตารูปานํ อินฺทฺริย\-ปจฺจเยน ปจฺจโย. รูปชีวิตินฺทฺริยํ อชฺฌตฺติกานํ กฏตฺตารูปานํ อินฺทฺริย\-ปจฺจเยน ปจฺจโย.}

\prose{\csromanfolio{510}พาหิโร ธมฺโม อชฺฌตฺติกสฺส จ พาหิรสฺส จ ธมฺมสฺส อินฺทฺริย\-ปจฺจเยน ปจฺจโย. พาหิรา อินฺทฺริยา สมฺปยุตฺตกานํ ขนฺธานํ จิตฺตสฺส จ จิตฺตสมุฏฺฐานานญฺจ รูปานํ อินฺทฺริย\-ปจฺจเยน ปจฺจโย. ปฏิสนฺธิ\-กฺขเณ พาหิรา อินฺทฺริยา สมฺปยุตฺตกานํ ขนฺธานํ จิตฺตสฺส จ อชฺฌตฺติกานญฺจ พาหิรานญฺจ กฏตฺตารูปานํ อินฺทฺริย\-ปจฺจเยน ปจฺจโย. รูปชีวิตินฺทฺริยํ อชฺฌตฺติกานญฺจ พาหิรานญฺจ กฏตฺตารูปานํ อินฺทฺริย\-ปจฺจเยน ปจฺจโย. (3)}

\proseitem{355}{อชฺฌตฺติโก จ พาหิโร จ ธมฺมา อชฺฌตฺติกสฺส ธมฺมสฺส อินฺทฺริย\-ปจฺจเยน ปจฺจโย. ปฏิสนฺธิ\-กฺขเณ อชฺฌตฺติกา จ พาหิรา จ อินฺทฺริยา อชฺฌตฺติกานํ กฏตฺตารูปานํ อินฺทฺริย\-ปจฺจเยน ปจฺจโย. จกฺขุนฺทฺริยญฺจ อุเปกฺขินฺทฺริยญฺจ จกฺขุ\-วิญฺญาณสฺส อินฺทฺริย\-ปจฺจเยน ปจฺจโย ฯเปฯ. กายินฺทฺริยญฺจ สุขินฺทฺริยญฺจ ฯเปฯ. กายินฺทฺริยญฺจ ทุกฺขินฺทฺริยญฺจ กายวิญฺญาณสฺส อินฺทฺริย\-ปจฺจเยน ปจฺจโย. (1)}

\prose{อชฺฌตฺติโก จ พาหิโร จ ธมฺมา พาหิรสฺส ธมฺมสฺส อินฺทฺริย\-ปจฺจเยน ปจฺจโย. อชฺฌตฺติกา จ พาหิรา จ อินฺทฺริยา สมฺปยุตฺตกานํ ขนฺธานํ จิตฺตสมุฏฺฐานานญฺจ รูปานํ อินฺทฺริย\-ปจฺจเยน ปจฺจโย. ปฏิสนฺธิ\-กฺขเณ อชฺฌตฺติกา จ พาหิรา จ อินฺทฺริยา สมฺปยุตฺตกานํ ขนฺธานํ พาหิรานญฺจ กฏตฺตารูปานํ อินฺทฺริย\-ปจฺจเยน ปจฺจโย. จกฺขุนฺทฺริยญฺจ อุเปกฺขินฺทฺริยญฺจ จกฺขุ\-วิญฺญาณสหคตานํ ขนฺธานํ อินฺทฺริย\-ปจฺจเยน ปจฺจโย ฯเปฯ. กายินฺทฺริยญฺจ สุขินฺทฺริยญฺจ ฯเปฯ. กายินฺทฺริยญฺจ ทุกฺขินฺทฺริยญฺจ กาย\-วิญฺญาณสหคตานํ ขนฺธานํ อินฺทฺริย\-ปจฺจเยน ปจฺจโย. (2)}

\prose{อชฺฌตฺติโก จ พาหิโร จ ธมฺมา อชฺฌตฺติกสฺส จ พาหิรสฺส จ ธมฺมสฺส อินฺทฺริย\-ปจฺจเยน ปจฺจโย. ปฏิสนฺธิ\-กฺขเณ อชฺฌตฺติกา จ พาหิรา จ อินฺทฺริยา สมฺปยุตฺตกานํ ขนฺธานํ อชฺฌตฺติกานญฺจ พาหิรานญฺจ กฏตฺตารูปานํ อินฺทฺริย\-ปจฺจเยน ปจฺจโย. จกฺขุนฺทฺริยญฺจ อุเปกฺขินฺทฺริยญฺจ จกฺขุ\-วิญฺญาณสฺส สมฺปยุตฺตกานญฺจ ขนฺธานํ อินฺทฺริย\-ปจฺจเยน ปจฺจโย. กายินฺทฺริยญฺจ ฯเปฯ. (3)}

\csromanheader{2}{\textbf{ฌานาทิ}}
\proseitem{356}{พาหิโร ธมฺโม พาหิรสฺส ธมฺมสฺส ฌานปจฺจเยน ปจฺจโย. ตีณิ. มคฺคปจฺจเยน ปจฺจโย. ตีณิ. สมฺปยุตฺต\-ปจฺจเยน ปจฺจโย. ปญฺจ.}

\csromanheader{2}{66. อชฺฌตฺติกทุก \textbf{วิปฺปยุตฺต}}
\proseitem{357}{\csromanfolio{511}อชฺฌตฺติโก ธมฺโม อชฺฌตฺติกสฺส ธมฺมสฺส วิปฺปยุตฺต\-ปจฺจเยน ปจฺจโย. สหชาตํ, ปุเรชาตํ, ปจฺฉาชาตํ.  สหชาตํ ปฏิสนฺธิ\-กฺขเณ จิตฺตํ อชฺฌตฺติกานํ กฏตฺตารูปานํ วิปฺปยุตฺต\-ปจฺจเยน ปจฺจโย. ปุเรชาตํ จกฺขายตนํ จกฺขุ\-วิญฺญาณสฺส ฯเปฯ. กายายตนํ กายวิญฺญาณสฺส วิปฺปยุตฺต\-ปจฺจเยน ปจฺจโย. ปจฺฉาชาตา อชฺฌตฺติกา ขนฺธา ปุเรชาตสฺส อิมสฺส อชฺฌตฺติกสฺส กายสฺส วิปฺปยุตฺต\-ปจฺจเยน ปจฺจโย. (1)}

\prose{อชฺฌตฺติโก ธมฺโม พาหิรสฺส ธมฺมสฺส วิปฺปยุตฺต\-ปจฺจเยน ปจฺจโย. สหชาตํ, ปุเรชาตํ, ปจฺฉาชาตํ.  สหชาตา อชฺฌตฺติกา ขนฺธา จิตฺต\-สมุฏฺฐานานํ รูปานํ วิปฺปยุตฺต\-ปจฺจเยน ปจฺจโย. ปฏิสนฺธิ\-กฺขเณ ฯเปฯ. ปุเรชาตํ จกฺขายตนํ จกฺขุ\-วิญฺญาณสหคตานํ ขนฺธานํ ฯเปฯ กายายตนํ กาย\-วิญฺญาณสหคตานํ ขนฺธานํ วิปฺปยุตฺต\-ปจฺจเยน ปจฺจโย. ปจฺฉาชาตา อชฺฌตฺติกา ขนฺธา ปุเรชาตสฺส อิมสฺส พาหิรสฺส กายสฺส วิปฺปยุตฺต\-ปจฺจเยน ปจฺจโย. (2)}

\prose{อชฺฌตฺติโก ธมฺโม อชฺฌตฺติกสฺส จ พาหิรสฺส จ ธมฺมสฺส วิปฺปยุตฺต\-ปจฺจเยน ปจฺจโย. สหชาตํ, ปุเรชาตํ, ปจฺฉาชาตํ.  สหชาตา ปฏิสนฺธิ\-กฺขเณ อชฺฌตฺติกา ขนฺธา อชฺฌตฺติกานญฺจ พาหิรานญฺจ กฏตฺตารูปานํ วิปฺปยุตฺต\-ปจฺจเยน ปจฺจโย. ปุเรชาตํ จกฺขายตนํ จกฺขุ\-วิญฺญาณสฺส สมฺปยุตฺตกานญฺจ ขนฺธานํ ฯเปฯ. กายายตนํ กายวิญฺญาณสฺส สมฺปยุตฺตกานํ ขนฺธานํ วิปฺปยุตฺต\-ปจฺจเยน ปจฺจโย. ปจฺฉาชาตา อชฺฌตฺติกา ขนฺธา ปุเรชาตสฺส อิมสฺส อชฺฌตฺติกสฺส จ พาหิรสฺส จ กายสฺส วิปฺปยุตฺต\-ปจฺจเยน ปจฺจโย. (3)}

\proseitem{358}{พาหิโร ธมฺโม พาหิรสฺส ธมฺมสฺส วิปฺปยุตฺต\-ปจฺจเยน ปจฺจโย. สหชาตํ, ปุเรชาตํ, ปจฺฉาชาตํ.  สหชาตา พาหิรา ขนฺธา จิตฺต\-สมุฏฺฐานานํ รูปานํ วิปฺปยุตฺต\-ปจฺจเยน ปจฺจโย. ปฏิสนฺธิ\-กฺขเณ ฯเปฯ. ขนฺธา วตฺถุสฺส วิปฺปยุตฺต\-ปจฺจเยน ปจฺจโย. วตฺถุ ขนฺธานํ วิปฺปยุตฺต\-ปจฺจเยน ปจฺจโย. ปุเรชาตํ วตฺถุ พาหิรานํ ขนฺธานํ วิปฺปยุตฺต\-ปจฺจเยน ปจฺจโย. ปจฺฉาชาตา พาหิรา ขนฺธา ปุเรชาตสฺส อิมสฺส พาหิรสฺส กายสฺส วิปฺปยุตฺต\-ปจฺจเยน ปจฺจโย. (1)}

\prose{\csromanfolio{512}พาหิโร ธมฺโม อชฺฌตฺติกสฺส ธมฺมสฺส วิปฺปยุตฺต\-ปจฺจเยน ปจฺจโย. สหชาตํ, ปุเรชาตํ, ปจฺฉาชาตํ.  สหชาตา ปฏิสนฺธิ\-กฺขเณ พาหิรา ขนฺธา อชฺฌตฺติกานํ กฏตฺตารูปานํ วิปฺปยุตฺต\-ปจฺจเยน ปจฺจโย. ปฏิสนฺธิ\-กฺขเณ วตฺถุ จิตฺตสฺส วิปฺปยุตฺต\-ปจฺจเยน ปจฺจโย. ปุเรชาตํ วตฺถุ จิตฺตสฺส วิปฺปยุตฺต\-ปจฺจเยน ปจฺจโย. ปจฺฉาชาตา พาหิรา ขนฺธา ปุเรชาตสฺส อิมสฺส อชฺฌตฺติกสฺส กายสฺส วิปฺปยุตฺต\-ปจฺจเยน ปจฺจโย. (2)}

\prose{พาหิโร ธมฺโม อชฺฌตฺติกสฺส จ พาหิรสฺส จ ธมฺมสฺส วิปฺปยุตฺต\-ปจฺจเยน ปจฺจโย. สหชาตํ, ปุเรชาตํ, ปจฺฉาชาตํ. (สํขิตฺตํ.) (3)}

\proseitem{359}{อชฺฌตฺติโก จ พาหิโร จ ธมฺมา อชฺฌตฺติกสฺส ธมฺมสฺส วิปฺปยุตฺต\-ปจฺจเยน ปจฺจโย. สหชาตํ, ปจฺฉาชาตํ.  สหชาตา ปฏิสนฺธิ\-กฺขเณ อชฺฌตฺติกา จ พาหิรา จ ขนฺธา อชฺฌตฺติกานํ กฏตฺตารูปานํ วิปฺปยุตฺต\-ปจฺจเยน ปจฺจโย. ปจฺฉาชาตํ\csromandash{}ปจฺฉาชาตา ฯเปฯ. (สํขิตฺตํ.) (1)}

\prose{อชฺฌตฺติโก จ พาหิโร จ ธมฺมา พาหิรสฺส ธมฺมสฺส วิปฺปยุตฺต\-ปจฺจเยน ปจฺจโย. สหชาตํ, ปจฺฉาชาตํ. (สํขิตฺตํ.) (2)}

\prose{อชฺฌตฺติโก จ พาหิโร จ ธมฺมา อชฺฌตฺติกสฺส จ พาหิรสฺส จ ธมฺมสฺส วิปฺปยุตฺต\-ปจฺจเยน ปจฺจโย. สหชาตํ, ปจฺฉาชาตํ.  สหชาตา ปฏิสนฺธิ\-กฺขเณ อชฺฌตฺติกา จ พาหิรา จ ขนฺธา ฯเปฯ. (สํขิตฺตํ.) (3)}

\prose{อตฺถิอาทิ}

\proseitem{360}{อชฺฌตฺติโก ธมฺโม อชฺฌตฺติกสฺส ธมฺมสฺส อตฺถิปจฺจเยน ปจฺจโย. สหชาตํ, ปุเรชาตํ, ปจฺฉาชาตํ.  สหชาตํ ปฏิสนฺธิ\-กฺขเณ จิตฺตํ อชฺฌตฺติกานํ กฏตฺตารูปานํ อตฺถิปจฺจเยน ปจฺจโย. ปุเรชาตํ จกฺขุํ ฯเปฯ กายํ อนิจฺจโต ฯเปฯ. (ปุเรชาต\-สทิสํ นินฺนานากรณํ.) ปจฺฉาชาตํ\csromandash{}(ปจฺฉาชาต\-สทิสํ กาตพฺพํ.) (1)}

\prose{อชฺฌตฺติโก ธมฺโม พาหิรสฺส ธมฺมสฺส อตฺถิปจฺจเยน ปจฺจโย. สหชาตํ, ปุเรชาตํ, ปจฺฉาชาตํ.  \textbf{สหชาตํ}\csromandash{}สหชาตา อชฺฌตฺติกา ขนฺธา สมฺปยุตฺตกานํ ขนฺธานํ จิตฺตสมุฏฺฐานานญฺจ รูปานํ อตฺถิปจฺจเยน ปจฺจโย. (สํขิตฺตํ.) (2)}

\prose{\csromanfolio{513}(อิธ อตฺถิ สพฺพฏฺฐาเน สหชาตํ ปจฺจยวาร\-สทิสํ, ปุเรชาตํ ปุเรชาต\-สทิสํ, ปจฺฉาชาตํ ปจฺฉาชาต\-สทิสํ กาตพฺพํ. นินฺนานากรณํ.)}

\prose{อชฺฌตฺติโก ธมฺโม อชฺฌตฺติกสฺส จ พาหิรสฺส จ ธมฺมสฺส อตฺถิปจฺจเยน ปจฺจโย. สหชาตํ, ปุเรชาตํ, ปจฺฉาชาตํ. (สํขิตฺตํ.) (3)}

\proseitem{361}{พาหิโร ธมฺโม พาหิรสฺส ธมฺมสฺส อตฺถิปจฺจเยน ปจฺจโย. สหชาตํ, ปุเรชาตํ, ปจฺฉาชาตํ, อาหารํ, อินฺทฺริยํ. (สพฺพํ วิตฺถาเรตพฺพํ.) (1)}

\prose{พาหิโร ธมฺโม อชฺฌตฺติกสฺส ธมฺมสฺส อตฺถิปจฺจเยน ปจฺจโย. สหชาตํ, ปุเรชาตํ, ปจฺฉาชาตํ, อาหารํ, อินฺทฺริยํ. (สํขิตฺตํ.) (2)}

\prose{พาหิโร ธมฺโม อชฺฌตฺติกสฺส จ พาหิรสฺส จ ธมฺมสฺส อตฺถิปจฺจเยน ปจฺจโย. สหชาตํ, ปุเรชาตํ, ปจฺฉาชาตํ, อาหารํ, อินฺทฺริยํ. (สํขิตฺตํ.) (3)}

\proseitem{362}{อชฺฌตฺติโก จ พาหิโร จ ธมฺมา อชฺฌตฺติกสฺส ธมฺมสฺส อตฺถิปจฺจเยน ปจฺจโย. สหชาตํ, ปุเรชาตํ, ปจฺฉาชาตํ, อาหารํ, อินฺทฺริยํ.  สหชาตา จกฺขุ\-วิญฺญาณสหคตา ขนฺธา จ จกฺขายตนญฺจ จกฺขุ\-วิญฺญาณสฺส อตฺถิปจฺจเยน ปจฺจโย ฯเปฯ. กาย\-วิญฺญาณสหคตา ขนฺธา จ ฯเปฯ. ปฏิสนฺธิ\-กฺขเณ อชฺฌตฺติกา จ พาหิรา จ ขนฺธา อชฺฌตฺติกานํ กฏตฺตารูปานํ อตฺถิปจฺจเยน ปจฺจโย. ปุเรชาตํ จกฺขายตนญฺจ วตฺถุ จ จิตฺตสฺส อตฺถิปจฺจเยน ปจฺจโย ฯเปฯ. กายายตนญฺจ วตฺถุ จ จิตฺตสฺส อตฺถิปจฺจเยน ปจฺจโย. รูปายตนญฺจ จกฺขายตนญฺจ จกฺขุ\-วิญฺญาณสฺส ฯเปฯ. โผฏฺฐพฺพายตนญฺจ กายายตนญฺจ กายวิญฺญาณสฺส อตฺถิปจฺจเยน ปจฺจโย. ปจฺฉาชาตา อชฺฌตฺติกา จ พาหิรา จ ขนฺธา ปุเรชาตสฺส อิมสฺส อชฺฌตฺติกสฺส กายสฺส อตฺถิปจฺจเยน ปจฺจโย. ปจฺฉาชาตา อชฺฌตฺติกา จ พาหิรา จ ขนฺธา กพฬีกาโร อาหาโร จ อิมสฺส อชฺฌตฺติกสฺส กายสฺส อตฺถิปจฺจเยน ปจฺจโย. ปจฺฉาชาตา อชฺฌตฺติกา จ พาหิรา จ ขนฺธา รูปชีวิตินฺทฺริยญฺจ อชฺฌตฺติกานํ กฏตฺตารูปานํ อตฺถิปจฺจเยน ปจฺจโย. (1)}

\prose{อชฺฌตฺติโก จ พาหิโร จ ธมฺมา พาหิรสฺส ธมฺมสฺส อตฺถิปจฺจเยน ปจฺจโย. สหชาตํ, ปุเรชาตํ, อาหารํ, อินฺทฺริยํ.   \csromanfolio{514}สหชาโต จกฺขุ\-วิญฺญาณ\-สหคโต เอโก ขนฺโธ จ จกฺขายตนญฺจ จกฺขุวิญฺญาณญฺจ ทฺวินฺนํ ขนฺธานํ อตฺถิปจฺจเยน ปจฺจโย ฯเปฯ. (สหชาต\-ปจฺจยวาร\-สทิสํ นินฺนานากรณํ, ปฐมคมนสทิสํเยว. สพฺเพ ปทา ปฐม\-ฆฏนา\-นเยน วิภชิตพฺพา.) (2)}

\prose{อชฺฌตฺติโก จ พาหิโร จ ธมฺมา อชฺฌตฺติกสฺส จ พาหิรสฺส จ ธมฺมสฺส อตฺถิปจฺจเยน ปจฺจโย. สหชาตํ, ปุเรชาตํ, ปจฺฉาชาตํ, อาหารํ, อินฺทฺริยํ.  สหชาโต จกฺขุ\-วิญฺญาณ\-สหคโต เอโก ขนฺโธ จ จกฺขายตนญฺจ ทฺวินฺนํ ขนฺธานํ จกฺขุ\-วิญฺญาณสฺส จ อตฺถิปจฺจเยน ปจฺจโย. (สํขิตฺตํ. สพฺเพ ปทา วิภชิตพฺพา, ปฐม\-ฆฏนา\-นเยน.) (3)}

\prose{นตฺถิ\-ปจฺจเยน ปจฺจโย. วิคต\-ปจฺจเยน ปจฺจโย. อวิคต\-ปจฺจเยน ปจฺจโย.}

\csromanmarkh{1. ปจฺจยานุโลม}\csromanmarkhii{2. สงฺขฺยาวาร}\csromanheadernomark{2}{1. \textbf{ปจฺจยานุโลม 2. สงฺขฺยาวาร}}
\proseitem{363}{เหตุยา ตีณิ, อารมฺมเณ นว, อธิปติยา นว, อนนฺตเร นว, สมนนฺตเร นว, สหชาเต นว, อญฺญมญฺเญ ปญฺจ, นิสฺสเย นว, อุปนิสฺสเย นว, ปุเรชาเต นว, ปจฺฉาชาเต นว, อาเสวเน นว, กมฺเม ตีณิ, วิปาเก นว, อาหาเร นว, อินฺทฺริเย นว, ฌาเน ตีณิ, มคฺเค ตีณิ, สมฺปยุตฺเต ปญฺจ, วิปฺปยุตฺเต นว, อตฺถิยา นว, นตฺถิยา นว, วิคเต นว, อวิคเต นว.}

\vspace{\tipitakaparskip}
\csromancenter{\textbf{อนุโลมํ.}}
\proseitem{364}{อชฺฌตฺติโก ธมฺโม อชฺฌตฺติกสฺส ธมฺมสฺส อารมฺมณ\-ปจฺจเยน ปจฺจโย. สหชาต\-ปจฺจเยน ปจฺจโย. อุปนิสฺสย\-ปจฺจเยน ปจฺจโย. ปุเรชาต\-ปจฺจเยน ปจฺจโย. ปจฺฉาชาต\-ปจฺจเยน ปจฺจโย. (1)}

\prose{อชฺฌตฺติโก ธมฺโม พาหิรสฺส ธมฺมสฺส อารมฺมณ\-ปจฺจเยน ปจฺจโย. สหชาต\-ปจฺจเยน ปจฺจโย. อุปนิสฺสย\-ปจฺจเยน ปจฺจโย. ปุเรชาต\-ปจฺจเยน ปจฺจโย. ปจฺฉาชาต\-ปจฺจเยน ปจฺจโย. (2)}

\prose{\csromanfolio{515}อชฺฌตฺติโก ธมฺโม อชฺฌตฺติกสฺส จ พาหิรสฺส จ ธมฺมสฺส อารมฺมณ\-ปจฺจเยน ปจฺจโย. สหชาต\-ปจฺจเยน ปจฺจโย. อุปนิสฺสย\-ปจฺจเยน ปจฺจโย. ปุเรชาต\-ปจฺจเยน ปจฺจโย. ปจฺฉาชาต\-ปจฺจเยน ปจฺจโย. (3)}

\proseitem{365}{พาหิโร ธมฺโม พาหิรสฺส ธมฺมสฺส อารมฺมณ\-ปจฺจเยน ปจฺจโย. สหชาต\-ปจฺจเยน ปจฺจโย. อุปนิสฺสย\-ปจฺจเยน ปจฺจโย. ปุเรชาต\-ปจฺจเยน ปจฺจโย. ปจฺฉาชาต\-ปจฺจเยน ปจฺจโย. กมฺมปจฺจเยน ปจฺจโย. อาหารปจฺจเยน ปจฺจโย. อินฺทฺริย\-ปจฺจเยน ปจฺจโย. (1)}

\prose{พาหิโร ธมฺโม อชฺฌตฺติกสฺส ธมฺมสฺส อารมฺมณ\-ปจฺจเยน ปจฺจโย. สหชาต\-ปจฺจเยน ปจฺจโย. อุปนิสฺสย\-ปจฺจเยน ปจฺจโย. ปุเรชาต\-ปจฺจเยน ปจฺจโย. ปจฺฉาชาต\-ปจฺจเยน ปจฺจโย. กมฺมปจฺจเยน ปจฺจโย. อาหารปจฺจเยน ปจฺจโย. อินฺทฺริย\-ปจฺจเยน ปจฺจโย. (2)}

\prose{พาหิโร ธมฺโม อชฺฌตฺติกสฺส จ พาหิรสฺส จ ธมฺมสฺส อารมฺมณ\-ปจฺจเยน ปจฺจโย. สหชาต\-ปจฺจเยน ปจฺจโย. อุปนิสฺสย\-ปจฺจเยน ปจฺจโย. ปุเรชาต\-ปจฺจเยน ปจฺจโย. ปจฺฉาชาต\-ปจฺจเยน ปจฺจโย. กมฺมปจฺจเยน ปจฺจโย. (3)}

\proseitem{366}{อชฺฌตฺติโก จ พาหิโร จ ธมฺมา อชฺฌตฺติกสฺส ธมฺมสฺส อารมฺมณ\-ปจฺจเยน ปจฺจโย. สหชาต\-ปจฺจเยน ปจฺจโย. อุปนิสฺสย\-ปจฺจเยน ปจฺจโย. ปุเรชาต\-ปจฺจเยน ปจฺจโย. ปจฺฉาชาต\-ปจฺจเยน ปจฺจโย. (1)}

\prose{อชฺฌตฺติโก จ พาหิโร จ ธมฺมา พาหิรสฺส ธมฺมสฺส อารมฺมณ\-ปจฺจเยน ปจฺจโย. สหชาต\-ปจฺจเยน ปจฺจโย. อุปนิสฺสย\-ปจฺจเยน ปจฺจโย. ปุเรชาต\-ปจฺจเยน ปจฺจโย. ปจฺฉาชาต\-ปจฺจเยน ปจฺจโย. (2)}

\prose{อชฺฌตฺติโก จ พาหิโร จ ธมฺมา อชฺฌตฺติตสฺส จ พาหิรสฺส จ ธมฺมสฺส อารมฺมณ\-ปจฺจเยน ปจฺจโย. สหชาต\-ปจฺจเยน ปจฺจโย. อุปนิสฺสย\-ปจฺจเยน ปจฺจโย. ปุเรชาต\-ปจฺจเยน ปจฺจโย. ปจฺฉาชาต\-ปจฺจเยน ปจฺจโย. (3)}

\csromanmarkh{2. ปจฺจย\-ปจฺจนีย}\csromanmarkhii{2. สงฺขฺยาวาร}\csromanheadernomark{2}{2. ปจฺจย\-ปจฺจนีย 2. สงฺขฺยาวาร}
\vspace{\tipitakaparskip}
\csromancenter{นเหตุยา นว, นอารมฺมเณ นว, (สพฺพตฺถ นว,) โนอวิคเต นว.}
\csromancenter{\textbf{เหตุ}ปจฺจยา นอารมฺมเณ ตีณิ, นอธิปติยา ตีณิ, นอนนฺตเร ตีณิ, นสมนนฺตเร ตีณิ, นอญฺญมญฺเญ ตีณิ, นฺเอาปนิสฺสเย ตีณิ, (สพฺพตฺถ ตีณิ,) นสมฺปยุตฺเต ตีณิ, นวิปฺปยุตฺเต ตีณิ, โนนตฺถิยา ตีณิ, โนวิคเต ตีณิ.}
\proseitemwithcloser{369}{\csromanfolio{516}นเหตุปจฺจยา อารมฺมเณ นว, อธิปติยา นว, (อนุโลมมาติกา กาตพฺพา.) ฯเปฯ อวิคเต นว.}{อชฺฌตฺติกทุกํ นิฏฺฐิตํ.}
\csromanmatikaanchor{mtk.66}
\csromantocmarkref{subsection}{67. อุปาทาทุก}{mtk.66}
\bookmark[dest={mtk.66},level=2]{67. อุปาทาทุก}
\csromanmarkh{67. อุปาทาทุก}\csromanmarkhii{1. ปฏิจฺจวาร}\csromanheadernomark{2}{67. \textbf{อุปาทาทุก 1. ปฏิจฺจวาร}}
\csromanmarkh{1. ปจฺจยานุโลม}\csromanmarkhii{1. วิภงฺควาร}\csromanheadernomark{2}{1. \textbf{ปจฺจยานุโลม 1. วิภงฺควาร}}
\vspace{\tipitakaparskip}
\csromancenter{\textbf{เหตุ}}
\proseitem{370}{อุปาทา ธมฺมํ ปฏิจฺจ โนอุปาทา ธมฺโม อุปฺปชฺชติ เหตุปจฺจยา. ปฏิสนฺธิ\-กฺขเณ วตฺถุํ ปฏิจฺจ โนอุปาทา ขนฺธา. (1)}

\prose{โนอุปาทา ธมฺมํ ปฏิจฺจ โนอุปาทา ธมฺโม อุปฺปชฺชติ เหตุปจฺจยา. โนอุปาทา เอกํ ขนฺธํ ปฏิจฺจ ตโย ขนฺธา โนอุปาทา จ จิตฺต\-สมุฏฺฐานํ รูปํ ฯเปฯ. เทฺว ขนฺเธ ฯเปฯ. ปฏิสนฺธิ\-กฺขเณ โนอุปาทา เอกํ ขนฺธํ ปฏิจฺจ ตโย ขนฺธา โนอุปาทา จ กฏตฺตารูปํ ฯเปฯ เทฺว ขนฺเธ ฯเปฯ. เอกํ มหาภูตํ ฯเปฯ เทฺว มหาภูเต ปฏิจฺจ เทฺว มหาภูตา. (1)}

\prose{โนอุปาทา ธมฺมํ ปฏิจฺจ อุปาทา ธมฺโม อุปฺปชฺชติ เหตุปจฺจยา. โนอุปาทา ขนฺเธ ปฏิจฺจ อุปาทา จิตฺต\-สมุฏฺฐานํ รูปํ. ปฏิสนฺธิ\-กฺขเณ ฯเปฯ มหภูเต ปฏิจฺจ อุปาทา จิตฺต\-สมุฏฺฐานํ รูปํ, กฏตฺตารูปํ, อุปาทารูปํ. (2)}

\prose{โนอุปาทา ธมฺมํ ปฏิจฺจ อุปาทา จ โนอุปาทา จ ธมฺมา อุปฺปชฺชนฺติ เหตุปจฺจยา. โนอุปาทา เอกํ ขนฺธํ ปฏิจฺจ ตโย ขนฺธา อุปาทา จ โนอุปาทา จ จิตฺต\-สมุฏฺฐานํ รูปํ ฯเปฯ เทฺว ขนฺเธ ฯเปฯ. ปฏิสนฺธิ\-กฺขเณ ฯเปฯ. (3)}

\prose{\csromanfolio{517}อุปาทา จ โนอุปาทา จ ธมฺมํ ปฏิจฺจ โนอุปาทา ธมฺโม อุปฺปชฺชติ เหตุปจฺจยา. ปฏิสนฺธิ\-กฺขเณ โนอุปาทา เอกํ ขนฺธญฺจ วตฺถุญฺจ ปฏิจฺจ ตโย ขนฺธา. เทฺว ขนฺเธ จ ฯเปฯ. (1)}

\proseitem{371}{อุปาทา ธมฺมํ ปฏิจฺจ โนอุปาทา ธมฺโม อุปฺปชฺชติ อารมฺมณ\-ปจฺจยา. ปฏิสนฺธิ\-กฺขเณ วตฺถุํ ปฏิจฺจ โนอุปาทา ขนฺธา. (1)}

\prose{โนอุปาทา ธมฺมํ ปฏิจฺจ โนอุปาทา ธมฺโม อุปฺปชฺชติ อารมฺมณ\-ปจฺจยา. โนอุปาทา เอกํ ขนฺธํ ปฏิจฺจ ตโย ขนฺธา ฯเปฯ เทฺว ขนฺเธ ฯเปฯ. ปฏิสนฺธิ\-กฺขเณ ฯเปฯ. (1)}

\prose{อุปาทา จ โนอุปาทา จ ธมฺมํ ปฏิจฺจ โนอุปาทา ธมฺโม อุปฺปชฺชติ อารมฺมณ\-ปจฺจยา. ปฏิสนฺธิ\-กฺขเณ โนอุปาทา เอกํ ขนฺธญฺจ วตฺถุญฺจ ปฏิจฺจ ตโย ขนฺธา ฯเปฯ เทฺว ขนฺเธ จ ฯเปฯ. (1)}

\csromanheader{2}{\textbf{อธิปตฺยาทิ}}
\proseitem{372}{โนอุปาทา ธมฺมํ ปฏิจฺจ โนอุปาทา ธมฺโม อุปฺปชฺชติ อธิปติ\-ปจฺจยา. โนอุปาทา เอกํ ขนฺธํ ปฏิจฺจ ตโย ขนฺธา โนอุปาทา จ จิตฺต\-สมุฏฺฐานํ รูปํ ฯเปฯ เทฺว ขนฺเธ ฯเปฯ. (1)}

\prose{โนอุปาทา ธมฺมํ ปฏิจฺจ อุปาทา ธมฺโม อุปฺปชฺชติ อธิปติ\-ปจฺจยา. โนอุปาทา ขนฺเธ ปฏิจฺจ อุปาทา จิตฺต\-สมุฏฺฐานํ รูปํ. มหาภูเต ปฏิจฺจ อุปาทา จิตฺต\-สมุฏฺฐานํ รูปํ. (2)}

\prose{โนอุปาทา ธมฺมํ ปฏิจฺจ อุปาทา จ โนอุปาทา จ ธมฺมา อุปฺปชฺชนฺติ อธิปติ\-ปจฺจยา. โนอุปาทา เอกํ ขนฺธํ ปฏิจฺจ ตโย ขนฺธา อุปาทา จ โนอุปาทา จ จิตฺต\-สมุฏฺฐานํ รูปํ ฯเปฯ เทฺว ขนฺเธ ฯเปฯ. (3)}

\prose{อนนฺตร\-ปจฺจยา ตีณิ. สมนนฺตร\-ปจฺจยา ตีณิ. สหชาต\-ปจฺจยา ปญฺจ.}

\csromanheader{2}{\textbf{อญฺญมญฺญ}}
\proseitem{373}{อุปาทา ธมฺมํ ปฏิจฺจ โนอุปาทา ธมฺโม อุปฺปชฺชติ อญฺญมญฺญ\-ปจฺจยา. ปฏิสนฺธิ\-กฺขเณ วตฺถุํ ปฏิจฺจ โนอุปาทา ขนฺธา. (1)}

\prose{โนอุปาทา ธมฺมํ ปฏิจฺจ โนอุปาทา ธมฺโม อุปฺปชฺชติ อญฺญมญฺญ\-ปจฺจยา. โนอุปาทา เอกํ ขนฺธํ ปฏิจฺจ ตโย ขนฺธา ฯเปฯ เทฺว ขนฺเธ ฯเปฯ. ปฏิสนฺธิ\-กฺขเณ ฯเปฯ. \csromanfolio{518}เอกํ มหาภูตํ ฯเปฯ. อสญฺญสตฺตานํ เอกํ มหาภูตํ ปฏิจฺจ ฯเปฯ เทฺว มหาภูเต ปฏิจฺจ เทฺว มหาภูตา. (1)}

\prose{โนอุปาทา ธมฺมํ ปฏิจฺจ อุปาทา ธมฺโม อุปฺปชฺชติ อญฺญมญฺญ\-ปจฺจยา. ปฏิสนฺธิ\-กฺขเณ โนอุปาทา ขนฺเธ ปฏิจฺจ วตฺถุ. (2)}

\prose{โนอุปาทา ธมฺมํ ปฏิจฺจ อุปาทา จ โนอุปาทา จ ธมฺมา อุปฺปชฺชนฺติ อญฺญมญฺญ\-ปจฺจยา. ปฏิสนฺธิ\-กฺขเณ โนอุปาทา เอกํ ขนฺธํ ปฏิจฺจ ตโย ขนฺธา วตฺถุ จ ฯเปฯ เทฺว ขนฺเธ ฯเปฯ. (3)}

\prose{อุปาทา จ โนอุปาทา จ ธมฺมํ ปฏิจฺจ โนอุปาทา ธมฺโม อุปฺปชฺชติ อญฺญมญฺญ\-ปจฺจยา. ปฏิสนฺธิ\-กฺขเณ โนอุปาทา เอกํ ขนฺธญฺจ วตฺถุญฺจ ปฏิจฺจ ตโย ขนฺธา ฯเปฯ เทฺว ขนฺเธ จ ฯเปฯ. (สํขิตฺตํ.) (1)}

\csromanmarkh{1. ปจฺจยานุโลม}\csromanmarkhii{2. สงฺขฺยาวาร}\csromanheadernomark{2}{1. \textbf{ปจฺจยานุโลม 2. สงฺขฺยาวาร}}
\proseitem{374}{เหตุยา ปญฺจ, อารมฺมเณ ตีณิ, อธิปติยา ตีณิ, อนนฺตเร ตีณิ, สมนนฺตเร ตีณิ, สหชาเต ปญฺจ, อญฺญมญฺเญ ปญฺจ, นิสฺสเย ปญฺจ, อุปนิสฺสเย ตีณิ, ปุเรชาเต เอกํ, อาเสวเน เอกํ, กมฺเม ปญฺจ, วิปาเก ปญฺจ, (สพฺพตฺถ ปญฺจ,) สมฺปยุตฺเต ตีณิ, วิปฺปยุตฺเต ปญฺจ, อตฺถิยา ปญฺจ, นตฺถิยา ตีณิ, วิคเต ตีณิ, อวิคเต ปญฺจ.}

\csromanmarkh{2. ปจฺจย\-ปจฺจนีย}\csromanmarkhii{1. วิภงฺควาร}\csromanheadernomark{2}{2. ปจฺจย\-ปจฺจนีย 1. วิภงฺควาร}
\proseitem{375}{อุปาทา ธมฺมํ ปฏิจฺจ โนอุปาทา ธมฺโม อุปฺปชฺชติ นเหตุปจฺจยา. อเหตุก\-ปฏิสนฺธิ\-กฺขเณ วตฺถุํ ปฏิจฺจ โนอุปาทา ขนฺธา. (1)}

\prose{โนอุปาทา ธมฺมํ ปฏิจฺจ โนอุปาทา ธมฺโม อุปฺปชฺชติ นเหตุปจฺจยา. อเหตุกํ โนอุปาทา เอกํ ขนฺธํ ปฏิจฺจ ตโย ขนฺธา โนอุปาทา จ จิตฺต\-สมุฏฺฐานํ รูปํ ฯเปฯ เทฺว ขนฺเธ ฯเปฯ. อเหตุก\-ปฏิสนฺธิ\-กฺขเณ ฯเปฯ. เอกํ มหาภูตํ ฯเปฯ. อสญฺญสตฺตานํ เอกํ มหาภูตํ ปฏิจฺจ ตโย มหาภูตา ฯเปฯ เทฺว มหาภูเต ปฏิจฺจ เทฺว มหาภูตา. วิจิกิจฺฉา\-สหคเต \csromanfolio{519}อุทฺธจฺจ\-สหคเต ขนฺเธ ปฏิจฺจ วิจิกิจฺฉา\-สหคโต อุทฺธจฺจ\-สหคโต โมโห. (1)}

\prose{โนอุปาทา ธมฺมํ ปฏิจฺจ อุปาทา ธมฺโม อุปฺปชฺชติ นเหตุปจฺจยา. อเหตุเก โนอุปาทา ขนฺเธ ปฏิจฺจ อุปาทา จิตฺต\-สมุฏฺฐานํ รูปํ. อเหตุก\-ปฏิสนฺธิ\-กฺขเณ ฯเปฯ มหาภูเต ปฏิจฺจ อุปาทา จิตฺต\-สมุฏฺฐานํ รูปํ, กฏตฺตารูปํ อุปาทารูปํ ฯเปฯ. (ยาว อสญฺญสตฺตา.) (2)}

\prose{โนอุปาทา ธมฺมํ ปฏิจฺจ อุปาทา จ โนอุปาทา จ ธมฺมา อุปฺปชฺชนฺติ นเหตุปจฺจยา. อเหตุกํ โนอุปาทา เอกํ ขนฺธํ ปฏิจฺจ ตโย ขนฺธา อุปาทา จ โนอุปาทา จ จิตฺต\-สมุฏฺฐานํ รูปํ ฯเปฯ เทฺว ขนฺเธ ฯเปฯ. อเหตุก\-ปฏิสนฺธิ\-กฺขเณ ฯเปฯ. (3)}

\prose{อุปาทา จ โนอุปาทา จ ธมฺมํ ปฏิจฺจ โนอุปาทา ธมฺโม อุปฺปชฺชติ นเหตุปจฺจยา. อเหตุก\-ปฏิสนฺธิ\-กฺขเณ โนอุปาทา เอกํ ขนฺธญฺจ วตฺถุญฺจ ปฏิจฺจ ตโย ขนฺธา ฯเปฯ เทฺว ขนฺเธ จ ฯเปฯ. (1)}

\prose{นอารมฺมณาทิ}

\proseitem{376}{โนอุปาทา ธมฺมํ ปฏิจฺจ โนอุปาทา ธมฺโม อุปฺปชฺชติ นอารมฺมณ\-ปจฺจยา. โนอุปาทา ขนฺเธ ปฏิจฺจ โนอุปาทา จิตฺต\-สมุฏฺฐานํ รูปํ. ปฏิสนฺธิ\-กฺขเณ ฯเปฯ. เอกํ มหาภูตํ ฯเปฯ. (ยาว อสญฺญสตฺตา.) เทฺว มหาภูเต ปฏิจฺจ เทฺว มหาภูตา. (1)}

\prose{โนอุปาทา ธมฺมํ ปฏิจฺจ อุปาทา ธมฺโม อุปฺปชฺชติ นอารมฺมณ\-ปจฺจยา. โนอุปาทา ขนฺเธ ปฏิจฺจ อุปาทา จิตฺต\-สมุฏฺฐานํ รูปํ. ปฏิสนฺธิ\-กฺขเณ ฯเปฯ มหาภูเต ปฏิจฺจ อุปาทา จิตฺต\-สมุฏฺฐานํ รูปํ, กฏตฺตารูปํ, อุปาทารูปํ. (ยาว อสญฺญสตฺตา.) (2)}

\prose{โนอุปาทา ธมฺมํ ปฏิจฺจ อุปาทา จ โนอุปาทา จ ธมฺมา อุปฺปชฺชนฺติ นอารมฺมณ\-ปจฺจยา. โนอุปาทา ขนฺเธ ปฏิจฺจ อุปาทา จ โนอุปาทา จ จิตฺต\-สมุฏฺฐานํ รูปํ. ปฏิสนฺธิ\-กฺขเณ ฯเปฯ. (3)}

\prose{นอธิปติ\-ปจฺจยา ปญฺจ. นอนนฺตร\-ปจฺจยา ตีณิ ฯเปฯ. นอุปนิสฺสยปจฺจยา ปญฺจ. นปุเรชาต\-ปจฺจยา ปญฺจ. นปจฺฉาชาต\-ปจฺจยา ปญฺจ. นอาเสวน\-ปจฺจยา ปญฺจ.}

\csromanheader{2}{\textbf{นกมฺม}}
\proseitem{377}{\csromanfolio{520}โนอุปาทา ธมฺมํ ปฏิจฺจ โนอุปาทา ธมฺโม อุปฺปชฺชติ นกมฺมปจฺจยา. โนอุปาทา ขนฺเธ ปฏิจฺจ สมฺปยุตฺตกา เจตนา. พาหิรํ อาหารสมุฏฺฐานํ อุตุสมุฏฺฐานํ ฯเปฯ เทฺว มหาภูเต ปฏิจฺจ เทฺว มหาภูตา. (1)}

\prose{โนอุปาทา ธมฺมํ ปฏิจฺจ อุปาทา ธมฺโม อุปฺปชฺชติ นกมฺมปจฺจยา. พาหิเร อาหารสมุฏฺฐาเน อุตุสมุฏฺฐาเน มหาภูเต ปฏิจฺจ อุปาทารูปํ. (2)}

\prose{โนอุปาทา ธมฺมํ ปฏิจฺจ อุปาทา จ โนอุปาทา จ ธมฺมา อุปฺปชฺชนฺติ นกมฺมปจฺจยา. พาหิรํ อาหารสมุฏฺฐานํ อุตุสมุฏฺฐานํ เอกํ มหาภูตํ ปฏิจฺจ ตโย มหาภูตา อุปาทา จ รูปํ ฯเปฯ เทฺว มหาภูเต ฯเปฯ. (3)}

\csromanheader{2}{\textbf{นวิปาก}}
\proseitem{378}{โนอุปาทา ธมฺมํ ปฏิจฺจ โนอุปาทา ธมฺโม อุปฺปชฺชติ นวิปาก\-ปจฺจยา. โนอุปาทา เอกํ ขนฺธํ ปฏิจฺจ ตโย ขนฺธา โนอุปาทา จ จิตฺต\-สมุฏฺฐานํ รูปํ ฯเปฯ เทฺว ขนฺเธ ฯเปฯ. เอกํ มหาภูตํ ฯเปฯ.}

\vspace{\tipitakaparskip}
\csromancenter{(ยาว อสญฺญสตฺตา, เอวํ ตีณิ, โนอุปาทามูลเก.)}
\prose{นอาหาร}

\proseitem{379}{โนอุปาทา ธมฺมํ ปฏิจฺจ โนอุปาทา ธมฺโม อุปฺปชฺชติ นอาหารปจฺจยา. พาหิรํ อุตุสมุฏฺฐานํ อสญฺญสตฺตานํ เอกํ มหาภูตํ ฯเปฯ เทฺว มหาภูเต ปฏิจฺจ มหาภูตา. (1)}

\prose{โนอุปาทา ธมฺมํ ปฏิจฺจ อุปาทา ธมฺโม อุปฺปชฺชติ นอาหารปจฺจยา. พาหิรํ อุตุสมุฏฺฐานํ อสญฺญสตฺตานํ มหาภูเต ปฏิจฺจ อุปาทารูปํ.}

\prose{โนอุปาทา ธมฺมํ ปฏิจฺจ อุปาทา จ โนอุปาทา จ ธมฺมา อุปฺปชฺชนฺติ นอาหารปจฺจยา. พาหิรํ อุตุสมุฏฺฐานํ อสญฺญสตฺตานํ เอกํ มหาภูตํ ปฏิจฺจ ตโย มหาภูตา อุปาทา จ รูปํ ฯเปฯ เทฺว มหาภูเต ฯเปฯ. (3)}

\prose{นอินฺทฺริย}

\proseitem{380}{โนอุปาทา ธมฺมํ ปฏิจฺจ โนอุปาทา ธมฺโม อุปฺปชฺชติ นอินฺทฺริยปจฺจยา. พาหิรํ อาหารสมุฏฺฐานํ อุตุสมุฏฺฐานํ เอกํ มหาภูตํ ฯเปฯ. (1)}

\prose{\csromanfolio{521}โนอุปาทา ธมฺมํ ปฏิจฺจ อุปาทา ธมฺโม อุปฺปชฺชติ นอินฺทฺริยปจฺจยา. พาหิเร อาหารสมุฏฺฐาเน อุตุสมุฏฺฐาเน มหาภูเต ปฏิจฺจ อุปาทารูปํ. อสญฺญสตฺตานํ มหาภูเต ปฏิจฺจ รูปชีวิตินฺทฺริยํ.}

\prose{โนอุปาทา ธมฺมํ ปฏิจฺจ อุปาทา จ โนอุปาทา จ ธมฺมา อุปฺปชฺชนฺติ นอินฺทฺริยปจฺจยา. พาหิรํ อาหารสมุฏฺฐานํ อุตุสมุฏฺฐานํ เอกํ มหาภูตํ ฯเปฯ. (3)}

\csromanheader{2}{\textbf{นฌานาทิ}}
\proseitem{381}{โนอุปาทา ธมฺมํ ปฏิจฺจ โนอุปาทา ธมฺโม อุปฺปชฺชติ นฌานปจฺจยา. ปญฺจ\-วิญฺญาณสหคตํ เอกํ ขนฺธํ ปฏิจฺจ ตโย ขนฺธา ฯเปฯ เทฺว ขนฺเธ ฯเปฯ. พาหิรํ อาหารสมุฏฺฐานํ อุตุสมุฏฺฐานํ อสญฺญสตฺตานํ ฯเปฯ เทฺว มหาภูเต ปฏิจฺจ เทฺว มหาภูตา. (1)}

\prose{โนอุปาทา ธมฺมํ ปฏิจฺจ อุปาทา ธมฺโม อุปฺปชฺชติ นฌานปจฺจยา. พาหิเร อาหารสมุฏฺฐาเน อุตุสมุฏฺฐาเน อสญฺญสตฺตานํ มหาภูเต ปฏิจฺจ อุปาทา กฏตฺตารูปํ. (2)}

\prose{โนอุปาทา ธมฺมํ ปฏิจฺจ อุปาทา จ โนอุปาทา จ ธมฺมา อุปฺปชฺชนฺติ นฌานปจฺจยา. พาหิรํ อาหารสมุฏฺฐานํ อุตุสมุฏฺฐานํ อสญฺญสตฺตานํ เอกํ มหาภูตํ ปฏิจฺจ ตโย มหาภูตา ฯเปฯ เทฺว มหาภูเต ฯเปฯ มหาภูเต ปฏิจฺจ อุปาทา กฏตฺตารูปํ, อุปาทารูปํ. (3)}

\prose{อุปาทา ธมฺมํ ปฏิจฺจ โนอุปาทา ธมฺโม อุปฺปชฺชติ นมคฺคปจฺจยา ปญฺจ. นสมฺปยุตฺต\-ปจฺจยา. ตีณิ.}

\csromanheader{2}{\textbf{นวิปฺปยุตฺตาทิ}}
\proseitem{382}{โนอุปาทา ธมฺมํ ปฏิจฺจ โนอุปาทา ธมฺโม อุปฺปชฺชติ นวิปฺปยุตฺต\-ปจฺจยา. อรูเป โนอุปาทา เอกํ ขนฺธํ ปฏิจฺจ ตโย ขนฺธา ฯเปฯ เทฺว ขนฺเธ ฯเปฯ. พาหิรํ อาหารสมุฏฺฐานํ อุตุสมุฏฺฐานํ อสญฺญสตฺตานํ เอกํ มหาภูตํ ฯเปฯ. (1)}

\prose{โนอุปาทา ธมฺมํ ปฏิจฺจ อุปาทา ธมฺโม อุปฺปชฺชติ นวิปฺปยุตฺต\-ปจฺจยา. พาหิเร อาหารสมุฏฺฐาเน อุตุสมุฏฺฐาเน อสญฺญสตฺตานํ มหาภูเต ปฏิจฺจ อุปาทา กฏตฺตารูปํ, อุปาทารูปํ. (2)}

\prose{\csromanfolio{522}โนอุปาทา ธมฺมํ ปฏิจฺจ อุปาทา จ โนอุปาทา จ ธมฺมา อุปฺปชฺชนฺติ นวิปฺปยุตฺต\-ปจฺจยา. พาหิรํ อาหารสมุฏฺฐานํ อุตุสมุฏฺฐานํ เอกํ มหาภูตํ ปฏิจฺจ ตโย มหาภูตา อุปาทา จ รูปํ ฯเปฯ เทฺว มหาภูเต ปฏิจฺจ เทฺว มหาภูตา อุปาทา จ รูปํ. อสญฺญสตฺตานํ เอกํ มหาภูตํ ปฏิจฺจ ตโย มหาภูตา กฏตฺตา จ รูปํ, อุปาทารูปํ ฯเปฯ เทฺว มหาภูเต ปฏิจฺจ เทฺว มหาภูตา กฏตฺตา จ รูปํ อุปาทารูปํ. โนนตฺถิปจฺจยา. โนวิคต\-ปจฺจยา. (3)}

\csromanmarkh{2. ปจฺจย\-ปจฺจนีย}\csromanmarkhii{2. สงฺขฺยาวาร}\csromanheadernomark{2}{2. ปจฺจย\-ปจฺจนีย 2. สงฺขฺยาวาร}
\proseitem{383}{นเหตุยา ปญฺจ, นอารมฺมเณ ตีณิ, นอธิปติยา ปญฺจ, นอนนฺตเร ตีณิ, นสมนนฺตเร ตีณิ, นอญฺญมญฺเญ ตีณิ, นอุปนิสฺสเย ตีณิ, นปุเรชาเต ปญฺจ, นปจฺฉาชาเต ปญฺจ, นอาเสวเน ปญฺจ, นกมฺเม ตีณิ, นวิปาเก ตีณิ, นอาหาเร ตีณิ, นอินฺทฺริเย ตีณิ, นฌาเน ตีณิ, นมคฺเค ปญฺจ, นสมฺปยุตฺเต ตีณิ, นวิปฺปยุตฺเต ตีณิ, โนนตฺถิยา ตีณิ, โนวิคเต ตีณิ.}

\proseitem{384}{เหตุปจฺจยา นอารมฺมเณ ตีณิ, นอธิปติยา ปญฺจ ฯเปฯ นกมฺเม เอกํ, นวิปาเก ตีณิ ฯเปฯ นสมฺปยุตฺเต ตีณิ, นวิปฺปยุตฺเต เอกํ, โนนตฺถิยา ตีณิ, โนวิคเต ตีณิ.}

\proseitem{385}{นเหตุปจฺจยา อารมฺมเณ ตีณิ, อนนฺตเร ตีณิ, สมนนฺตเร ตีณิ, สหชาเต ปญฺจ, อญฺญมญฺเญ ปญฺจ, นิสฺสเย ปญฺจ, อุปนิสฺสเย ตีณิ, ปุเรชาเต เอกํ, อาเสวเน เอกํ ฯเปฯ มคฺเค เอกํ ฯเปฯ อวิคเต ปญฺจ.}

\vspace{\tipitakaparskip}
\csromancenter{(สหชาตวาโร ปฏิจฺจ\-วาร\-สทิโส.)}
\csromanmarkh{67. อุปาทาทุก}\csromanmarkhii{67. อุปาทาทุก 3. ปจฺจยวาร}\csromanheadernomark{2}{67. อุปาทาทุก \textbf{67. อุปาทาทุก 3. ปจฺจยวาร}}
\csromanmarkh{1. ปจฺจยานุโลม}\csromanmarkhii{1. วิภงฺควาร}\csromanheadernomark{2}{1. \textbf{ปจฺจยานุโลม 1. วิภงฺควาร}}
\csromanheader{2}{\textbf{เหตุ}}
\proseitem{386}{\csromanfolio{523}อุปาทา ธมฺมํ ปจฺจยา โนอุปาทา ธมฺโม อุปฺปชฺชติ เหตุปจฺจยา. วตฺถุํ ปจฺจยา โนอุปาทา ขนฺธา. ปฏิสนฺธิ\-กฺขเณ ฯเปฯ. (1)}

\prose{โนอุปาทา ธมฺมํ ปจฺจยา โนอุปาทา ธมฺโม อุปฺปชฺชติ เหตุปจฺจยา. ตีณิ. (โนอุปาทามูลเก ตีณิปิ ปฏิจฺจสทิสา, นินฺนานากรณา.) (3)}

\prose{อุปาทา จ โนอุปาทา จ ธมฺมํ ปจฺจยา โนอุปาทา ธมฺโม อุปฺปชฺชติ เหตุปจฺจยา. โนอุปาทา เอกํ ขนฺธญฺจ วตฺถุญฺจ ปจฺจยา ตโย ขนฺธา ฯเปฯ เทฺว ขนฺเธ จ ฯเปฯ. ปฏิสนฺธิ\-กฺขเณ ฯเปฯ. (1)}

\proseitem{387}{อุปาทา ธมฺมํ ปจฺจยา โนอุปาทา ธมฺโม อุปฺปชฺชติ อารมฺมณ\-ปจฺจยา. จกฺขายตนํ ปจฺจยา จกฺขุวิญฺญาณํ ฯเปฯ. กายายตนํ ปจฺจยา กายวิญฺญาณํ. วตฺถุํ ปจฺจยา โนอุปาทา ขนฺธา. ปฏิสนฺธิ\-กฺขเณ ฯเปฯ. (1)}

\prose{โนอุปาทา ธมฺมํ ปจฺจยา โนอุปาทา ธมฺโม อุปฺปชฺชติ อารมฺมณ\-ปจฺจยา. เอกํ. (ปฏิจฺจสทิสํ.) (1)}

\prose{อุปาทา ธมฺมญฺจ โนอุปาทา ธมฺมญฺจ ธมฺมํ ปจฺจยา โนอุปาทา ธมฺโม อุปฺปชฺชติ อารมฺมณ\-ปจฺจยา. จกฺขุ\-วิญฺญาณสหคตํ เอกํ ขนฺธญฺจ จกฺขายตนญฺจ ปจฺจยา ตโย ขนฺธา ฯเปฯ เทฺว ขนฺเธ จ ฯเปฯ. กาย\-วิญฺญาณสหคตํ ฯเปฯ. โนอุปาทา เอกํ ขนฺธญฺจ วตฺถุญฺจ ปจฺจยา ตโย ขนฺธา ฯเปฯ เทฺว ขนฺเธ จ ฯเปฯ. ปฏิสนฺธิ\-กฺขเณ ฯเปฯ. (สํขิตฺตํ.) (1)}

\csromanmarkh{1. ปจฺจยานุโลม}\csromanmarkhii{2. สงฺขฺยาวาร}\csromanheadernomark{2}{1. \textbf{ปจฺจยานุโลม 2. สงฺขฺยาวาร}}
\proseitem{388}{\textbf{เหตุ}ยา ปญฺจ, อารมฺมเณ ตีณิ, อธิปติยา ปญฺจ, อนนฺตเร ตีณิ, สมนนฺตเร ตีณิ, สหชาเต ปญฺจ, อญฺญมญฺเญ ปญฺจ, นิสฺสเย ปญฺจ, อุปนิสฺสเย ตีณิ, ปุเรชาเต ตีณิ, อาเสวเน ตีณิ, กมฺเม ปญฺจ ฯเปฯ อวิคเต ปญฺจ. (เอวํ คเณตพฺพํ.)}

\csromanmarkh{2. ปจฺจย\-ปจฺจนีย}\csromanmarkhii{1. วิภงฺควาร}\csromanheadernomark{2}{2. ปจฺจย\-ปจฺจนีย 1. วิภงฺควาร}
\proseitem{389}{\csromanfolio{524}อุปาทา ธมฺมํ ปจฺจยา โนอุปาทา ธมฺโม อุปฺปชฺชติ นเหตุปจฺจยา. จกฺขายตนํ ปจฺจยา จกฺขุวิญฺญาณํ ฯเปฯ. กายายตนํ ปจฺจยา กายวิญฺญาณํ. วตฺถุํ ปจฺจยา อเหตุกา โนอุปาทา ขนฺธา. อเหตุก\-ปฏิสนฺธิ\-กฺขเณ ฯเปฯ. วตฺถุํ ปจฺจยา วิจิกิจฺฉา\-สหคโต อุทฺธจฺจ\-สหคโต โมโห. (1)}

\prose{โนอุปาทา ธมฺมํ ปจฺจยา โนอุปาทา ธมฺโม อุปฺปชฺชติ นเหตุปจฺจยา. อเหตุกํ โนอุปาทา เอกํ ขนฺธํ ปจฺจยา ตโย ขนฺธา โนอุปาทา จ จิตฺต\-สมุฏฺฐานํ รูปํ ฯเปฯ เทฺว ขนฺเธ ฯเปฯ. อเหตุก\-ปฏิสนฺธิ\-กฺขเณ ฯเปฯ. เอกํ มหาภูตํ ฯเปฯ. (ยาว อสญฺญสตฺตา.) วิจิกิจฺฉา\-สหคเต อุทฺธจฺจ\-สหคเต ขนฺเธ ปจฺจยา วิจิกิจฺฉา\-สหคโต อุทฺธจฺจ\-สหคโต โมโห. (ตีณิปิ. ปฏิจฺจสทิสา นินฺนานากรณา.) (3)}

\prose{อุปาทา จ โนอุปาทา จ ธมฺมํ ปจฺจยา โนอุปาทา ธมฺโม อุปฺปชฺชติ นเหตุปจฺจยา. จกฺขุ\-วิญฺญาณสหคตํ เอกํ ขนฺธญฺจ จกฺขายตนญฺจ ปจฺจยา ตโย ขนฺธา ฯเปฯ เทฺว ขนฺเธ จ ฯเปฯ. กาย\-วิญฺญาณสหคตํ ฯเปฯ. อเหตุกํ โนอุปาทา เอกํ ขนฺธญฺจ วตฺถุญฺจ ปจฺจยา ตโย ขนฺธา ฯเปฯ เทฺว ขนฺเธ จ ฯเปฯ. ปฏิสนฺธิ\-กฺขเณ ฯเปฯ. วิจิกิจฺฉา\-สหคเต อุทฺธจฺจ\-สหคเต ขนฺเธ จ วตฺถุญฺจ ปจฺจยา วิจิกิจฺฉา\-สหคโต อุทฺธจฺจ\-สหคโต โมโห. (1)}

\prose{นอารมฺมณ\-ปจฺจยา ตีณิ. นอาเสวน\-ปจฺจยา ปญฺจ.}

\csromanheader{2}{\textbf{นกมฺมาทิ}}
\proseitem{390}{อุปาทา ธมฺมํ ปจฺจยา โนอุปาทา ธมฺโม อุปฺปชฺชติ นกมฺมปจฺจยา. วตฺถุํ ปจฺจยา โนอุปาทา เจตนา. (1)}

\prose{โนอุปาทา ธมฺมํ ปจฺจยา โนอุปาทา ธมฺโม อุปฺปชฺชติ นกมฺมปจฺจยา. โนอุปาทา ขนฺเธ ปจฺจยา สมฺปยุตฺตกา เจตนา. พาหิรํ อาหารสมุฏฺฐานํ อุตุสมุฏฺฐานํ ฯเปฯ เทฺว มหาภูเต ปจฺจยา เทฺว มหาภูตา. (1)}

\prose{โนอุปาทา ธมฺมํ ปจฺจยา อุปาทา ธมฺโม อุปฺปชฺชติ นกมฺมปจฺจยา. พาหิเร อาหารสมุฏฺฐาเน อุตุสมุฏฺฐาเน มหาภูเต ปจฺจยา อุปาทารูปํ. (2)}

\prose{\csromanfolio{525}โนอุปาทา ธมฺมํ ปจฺจยา อุปาทา จ โนอุปาทา จ ธมฺมา อุปฺปชฺชนฺติ นกมฺมปจฺจยา. พาหิรํ อาหารสมุฏฺฐานํ อุตุสมุฏฺฐานํ เอกํ มหาภูตํ ปจฺจยา ตโย มหาภูตา อุปาทา จ รูปํ ฯเปฯ เทฺว มหาภูเต ฯเปฯ. (3)}

\prose{อุปาทาจ โนอุปาทา จ ธมฺมํ ปจฺจยา โนอุปาทา ธมฺโม อุปฺปชฺชติ นกมฺมปจฺจยา. โนอุปาทา ขนฺเธ จ วตฺถุญฺจ ปจฺจยา โนอุปาทา เจตนา. (1)}

\prose{นวิปาก\-ปจฺจยา ปญฺจ. นอาหารปจฺจยา ตีณิ. นอินฺทฺริยปจฺจยา ตีณิ.}

\csromanheader{2}{\textbf{นฌานาทิ}}
\proseitem{391}{อุปาทา ธมฺมํ ปจฺจยา โนอุปาทา ธมฺโม อุปฺปชฺชติ นฌานปจฺจยา. จกฺขายตนํ ปจฺจยา จกฺขุวิญฺญาณํ ฯเปฯ. กายายตนํ ปจฺจยา กายวิญฺญาณํ. (1)}

\prose{โนอุปาทา ธมฺมํ ปจฺจยา โนอุปาทา ธมฺโม อุปฺปชฺชติ นฌานปจฺจยา. ปญฺจ\-วิญฺญาณสหคตํ เอกํ ขนฺธํ ปจฺจยา ตโย ขนฺธา ฯเปฯ เทฺว ขนฺเธ ฯเปฯ. พาหิรํ อาหารสมุฏฺฐานํ อุตุสมุฏฺฐานํ อสญฺญสตฺตานํ ฯเปฯ เทฺว มหาภูเต ปจฺจยา เทฺว มหาภูตา. (1)}

\prose{โนอุปาทา ธมฺมํ ปจฺจยา อุปาทา ธมฺโม อุปฺปชฺชติ นฌานปจฺจยา. พาหิเร อาหารสมุฏฺฐาเน อุตุสมุฏฺฐาเน อสญฺญสตฺตานํ มหาภูเต ปจฺจยา อุปาทา กฏตฺตารูปํ. (2)}

\prose{โนอุปาทา ธมฺมํ ปจฺจยา อุปาทา จ โนอุปาทา จ ธมฺมา อุปฺปชฺชนฺติ นฌานปจฺจยา. พาหิรํ อาหารสมุฏฺฐานํ อุตุสมุฏฺฐานํ อสญฺญสตฺตานํ เอกํ มหาภูตํ ปจฺจยา ตโย มหาภูตา อุปาทา จ กฏตฺตารูปํ ฯเปฯ เทฺว ฯเปฯ. (3)}

\prose{อุปาทา จ โนอุปาทา จ ธมฺมํ ปจฺจยา โนอุปาทา ธมฺโม อุปฺปชฺชติ นฌานปจฺจยา. จกฺขุ\-วิญฺญาณสหคตํ เอกํ ขนฺธญฺจ จกฺขายตนญฺจ ปจฺจยา ตโย ขนฺธา ฯเปฯ เทฺว ขนฺเธ จ ฯเปฯ. (1)}

\prose{นมคฺคปจฺจยา ปญฺจ. นสมฺปยุตฺต\-ปจฺจยา ตีณิ. นวิปฺปยุตฺต\-ปจฺจยา ตีณิ. โนนตฺถิปจฺจยา ตีณิ. โนวิคต\-ปจฺจยา ตีณิ.}

\csromanmarkh{2. ปจฺจย\-ปจฺจนีย}\csromanmarkhii{2. สงฺขฺยาวาร}\csromanheadernomark{2}{2. ปจฺจย\-ปจฺจนีย 2. สงฺขฺยาวาร}
\proseitem{392}{\csromanfolio{526}นเหตุยา ปญฺจ, นอารมฺมเณ ตีณิ, นอธิปติยา ปญฺจ ฯเปฯ นกมฺเม ปญฺจ, นวิปาเก ปญฺจ, นอาหาเร ตีณิ, นอินฺทฺริเย ตีณิ, นฌาเน ปญฺจ, นมคฺเค ปญฺจ, นสมฺปยุตฺเต ตีณิ, นวิปฺปยุตฺเต ตีณิ, โนนตฺถิยา ตีณิ, โนวิคเต ตีณิ.}

\vspace{\tipitakaparskip}
\csromancenter{(เอวํ อิตเร เทฺว คณนาปิ นิสฺสวาโรปิ กาตพฺโพ.)}
\csromanmarkh{67. อุปาทาทุก}\csromanmarkhii{5. สํสฏฺฐวาร}\csromanheadernomark{2}{67. \textbf{อุปาทาทุก 5. สํสฏฺฐวาร}}
\proseitem{393}{โนอุปาทา ธมฺมํ สํสฏฺโฐ โนอุปาทา ธมฺโม อุปฺปชฺชติ เหตุปจฺจยา. โนอุปาทา เอกํ ขนฺธํ สํสฏฺฐา ตโย ขนฺธา ฯเปฯ เทฺว ขนฺเธ สํสฏฺฐา เทฺว ขนฺธา. ปฏิสนฺธิ\-กฺขเณ ฯเปฯ. (สํขิตฺตํ.)}

\prose{เหตุยา เอกํ, อารมฺมเณ เอกํ, อธิปติยา เอกํ, (สพฺพตฺถ เอกํ.) อวิคเต เอกํ. อนุโลมํ.}

\prose{โนอุปาทา ธมฺมํ สํสฏฺโฐ โนอุปาทา ธมฺโม อุปฺปชฺชติ นเหตุปจฺจยา. อเหตุกํ โนอุปาทา เอกํ ขนฺธํ สํสฏฺฐา ตโย ขนฺธา ฯเปฯ เทฺว ขนฺเธ ฯเปฯ. อเหตุก\-ปฏิสนฺธิ\-กฺขเณ ฯเปฯ. วิจิกิจฺฉา\-สหคเต อุทฺธจฺจ\-สหคเต ขนฺเธ สํสฏฺโฐ วิจิกิจฺฉา\-สหคโต อุทฺธจฺจ\-สหคโต โมโห. (สํขิตฺตํ.)}

\prose{นเหตุยา เอกํ, นอธิปติยา เอกํ, นปุเรชาเต เอกํ, นปจฺฉาชาเต เอกํ, นอาเสวเน เอกํ, นกมฺเม เอกํ, นวิปาเก เอกํ, นฌาเน เอกํ, นมคฺเค เอกํ, นวิปฺปยุตฺเต เอกํ.}

\vspace{\tipitakaparskip}
\csromancenter{ปจฺจนียํ.}
\csromancenter{(อิตเร เทฺว คณนาปิ สมฺปยุตฺตวาโรปิ กาตพฺโพ.)}
\csromanmarkh{67. อุปาทาทุก}\csromanmarkhii{67. อุปาทาทุก 7. ปญฺหาวาร}\csromanheadernomark{2}{67. อุปาทาทุก \textbf{67. อุปาทาทุก 7. ปญฺหาวาร}}
\csromanmarkh{1. ปจฺจยานุโลม}\csromanmarkhii{1. วิภงฺควาร}\csromanheadernomark{2}{1. \textbf{ปจฺจยานุโลม 1. วิภงฺควาร}}
\csromanheader{2}{\textbf{เหตุ}}
\proseitem{394}{\csromanfolio{527}โนอุปาทา ธมฺโม โนอุปาทา ธมฺมสฺส เหตุปจฺจเยน ปจฺจโย. โนอุปาทา เหตู สมฺปยุตฺตกานํ ขนฺธานํ โนอุปาทา จ จิตฺต\-สมุฏฺฐานานํ รูปานํ เหตุปจฺจเยน ปจฺจโย. ปฏิสนฺธิ\-กฺขเณ ฯเปฯ. (1)}

\prose{โนอุปาทา ธมฺโม อุปาทา ธมฺมสฺส เหตุปจฺจเยน ปจฺจโย. โนอุปาทา เหตู อุปาทา จิตฺต\-สมุฏฺฐานานํ รูปานํ เหตุปจฺจเยน ปจฺจโย. ปฏิสนฺธิ\-กฺขเณ ฯเปฯ. (2)}

\prose{โนอุปาทา ธมฺโม อุปาทา จ โนอุปาทา จ ธมฺมสฺส เหตุปจฺจเยน ปจฺจโย. โนอุปาทา เหตู สมฺปยุตฺตกานํ ขนฺธานํ อุปาทา จ โนอุปาทา จ จิตฺต\-สมุฏฺฐานานํ รูปานํ เหตุปจฺจเยน ปจฺจโย. ปฏิสนฺธิ\-กฺขเณ ฯเปฯ. (3)}

\proseitem{395}{อุปาทา ธมฺโม โนอุปาทา ธมฺมสฺส อารมฺมณ\-ปจฺจเยน ปจฺจโย. จกฺขุํ ฯเปฯ กายํ. รูเป ฯเปฯ รเส วตฺถุํ อนิจฺจโต ฯเปฯ โทมนสฺสํ อุปฺปชฺชติ. ทิพฺเพน จกฺขุนา รูปํ ปสฺสติ. ทิพฺพาย โสตธาตุยา สทฺทํ สุณาติ. รูปายตนํ จกฺขุ\-วิญฺญาณสฺส ฯเปฯ. รสายตนํ ชิวฺหา\-วิญฺญาณสฺส อารมฺมณ\-ปจฺจเยน ปจฺจโย. อุปาทา ขนฺธา อิทฺธิวิธ\-ญาณสฺส, ปุพฺเพ\-นิวาสา\-นุสฺสติ\-ญาณสฺส, อนาคตํสญาณสฺส, อาวชฺชนาย อารมฺมณ\-ปจฺจเยน ปจฺจโย. (1)}

\prose{โนอุปาทา ธมฺโม โนอุปาทา ธมฺมสฺส อารมฺมณ\-ปจฺจเยน ปจฺจโย. ทานํ ฯเปฯ สีลํ ฯเปฯ อุโปสถกมฺมํ กตฺวา ตํ ปจฺจเวกฺขติ, อสฺสาเทติ อภินนฺทติ, ตํ อารพฺภ ราโค อุปฺปชฺชติ ฯเปฯ โทมนสฺสํ อุปฺปชฺชติ. ปุพฺเพ ฯเปฯ. ฌานา ฯเปฯ. อริยา มคฺคา วุฏฺฐหิตฺวา มคฺคํ ปจฺจเวกฺขนฺติ. ผลํ ฯเปฯ. นิพฺพานํ ฯเปฯ. นิพฺพานํ โคตฺรภุสฺส, โวทานสฺส, มคฺคสฺส, ผลสฺส, อาวชฺชนาย อารมฺมณ\-ปจฺจเยน ปจฺจโย. อริยา ปหีเน กิเลเส ปจฺจเวกฺขนฺติ. วิกฺขมฺภิเต กิเลเส ฯเปฯ. ปุพฺเพ ฯเปฯ. โผฏฺฐพฺเพ โนอุปาทา ขนฺเธ อนิจฺจโต ฯเปฯ โทมนสฺสํ อุปฺปชฺชติ. เจโตปริย\-ญาเณน โนอุปาทาจิตฺตสมงฺคิสฺส จิตฺตํ \csromanfolio{528}ชานาติ. อากาสา\-นญฺจา\-ยตนํ วิญฺญาณญฺจา\-ยตนสฺส ฯเปฯ. อากิญฺจญฺญา\-ยตนํ เนว\-สญฺญา\-นา\-สญฺญา\-ยตนสฺส ฯเปฯ. โผฏฺฐพฺพา\-ยตนํ กายวิญฺญาณสฺส, อารมฺมณ\-ปจฺจเยน ปจฺจโย. โนอุปาทา ขนฺธา อิทฺธิวิธ\-ญาณสฺส เจโตปริย\-ญาณสฺส, ปุพฺเพ\-นิวาสา\-นุสฺสติ\-ญาณสฺส, ยถากมฺมูปคญาณสฺส, อนาคตํสญาณสฺส, อาวชฺชนาย อารมฺมณ\-ปจฺจเยน ปจฺจโย. (1)}

\proseitem{396}{อุปาทา ธมฺโม โนอุปาทา ธมฺมสฺส อธิปติ\-ปจฺจเยน ปจฺจโย. อารมฺมณา\-ธิปติ\csromandash{}จกฺขุํ ฯเปฯ กายํ. รูเป ฯเปฯ รเส วตฺถุํ ครุํ กตฺวา อสฺสาเทติ, อภินนฺทติ, ตํ ครุํ กตฺวา ราโค อุปฺปชฺชติ. ทิฏฺฐิ อุปฺปชฺชติ. (1)}

\prose{โนอุปาทา ธมฺโม โนอุปาทา ธมฺมสฺส อธิปติ\-ปจฺจเยน ปจฺจโย. อารมฺมณา\-ธิปติ, สหชาตา\-ธิปติ.  อารมฺมณา\-ธิปติ\csromandash{}ทานํ ฯเปฯ สีลํ ฯเปฯ อุโปสถกมฺมํ กตฺวา ตํ ครุํ กตฺวา ปจฺจเวกฺขติ, อสฺสาเทติ อภินนฺทติ, ตํ ครุํ กตฺวา ราโค อุปฺปชฺชติ. ทิฏฺฐิ อุปฺปชฺชติ. ปุพฺเพ สุจิณฺณานิ ฯเปฯ. ฌานา วุฏฺฐหิตฺวา ฌานํ ฯเปฯ. อริยา มคฺคา วุฏฺฐหิตฺวา มคฺคํ ครุํ กตฺวา ฯเปฯ. ผลํ ฯเปฯ. นิพฺพานํ ครุํ กตฺวา ปจฺจเวกฺขนฺติ. นิพฺพานํ โคตฺรภุสฺส, โวทานสฺส, มคฺคสฺส, ผลสฺส อธิปติ\-ปจฺจเยน ปจฺจโย. โผฏฺฐพฺเพ โนอุปาทา ขนฺเธ ครุํ กตฺวา อสฺสาเทติ อภินนฺทติ, ตํ ครุํ กตฺวา ราโค อุปฺปชฺชติ. ทิฏฺฐิ อุปฺปชฺชติ. สหชาตา\-ธิปติ\csromandash{}โนอุปาทาธิปติ สมฺปยุตฺตกานํ ขนฺธานํ โนอุปาทา จ จิตฺต\-สมุฏฺฐานานํ รูปานํ อธิปติ\-ปจฺจเยน ปจฺจโย. (1)}

\prose{โนอุปาทา ธมฺโม อุปาทา ธมฺมสฺส อธิปติ\-ปจฺจเยน ปจฺจโย. สหชาตา\-ธิปติ\csromandash{}โนอุปาทาธิปติ อุปาทา จิตฺต\-สมุฏฺฐานานํ รูปานํ อธิปติ\-ปจฺจเยน ปจฺจโย. (2)}

\prose{โนอุปาทา ธมฺโม อุปาทา จ โนอุปาทา จ ธมฺมสฺส อธิปติ\-ปจฺจเยน ปจฺจโย. สหชาตา\-ธิปติ\csromandash{}โนอุปาทาธิปติ สมฺปยุตฺตกานํ ขนฺธานํ อุปาทา จ โนอุปาทา จ จิตฺต\-สมุฏฺฐานานํ รูปานํ อธิปติ\-ปจฺจเยน ปจฺจโย. (3)}

\csromanheader{2}{67. อุปาทาทุก \textbf{อนนฺตราทิ}}
\proseitem{397}{\csromanfolio{529}โนอุปาทา ธมฺโม โนอุปาทา ธมฺมสฺส อนนฺตร\-ปจฺจเยน ปจฺจโย. ปุริมา ปุริมา โนอุปาทา ขนฺธา ปจฺฉิมานํ ปจฺฉิมานํ โนอุปาทา ขนฺธานํ อนนฺตร\-ปจฺจเยน ปจฺจโย. อนุโลมํ โคตฺรภุสฺส ฯเปฯ ผลสมาปตฺติยา อนนฺตร\-ปจฺจเยน ปจฺจโย. (1)}

\prose{สมนนฺตร\-ปจฺจเยน ปจฺจโย. สหชาต\-ปจฺจเยน ปจฺจโย. (ปฏิจฺจสทิสํ.) อญฺญมญฺญ\-ปจฺจเยน ปจฺจโย. (ปฏิจฺจสทิสํ.) นิสฺสย\-ปจฺจเยน ปจฺจโย. (ปจฺจวาเร นิสฺสยสทิสํ.)}

\csromanheader{2}{\textbf{อุปนิสฺสย}}
\proseitem{398}{อุปาทา ธมฺโม โนอุปาทา ธมฺมสฺส อุปนิสฺสย\-ปจฺจเยน ปจฺจโย. อารมฺมณู\-ปนิสฺสโย, ปกตู\-ปนิสฺสโย ฯเปฯ. ปกตู\-ปนิสฺสโย\csromandash{} จกฺขุสมฺปทํ ฯเปฯ กายสมฺปทํ. วณฺณสมฺปทํ. สทฺทสมฺปทํ. คนฺธสมฺปทํ. รสสมฺปทํ. โภชนํ อุปนิสฺสาย ทานํ เทติ ฯเปฯ สํฆํ ภินฺทติ. จกฺขุสมฺปทา ฯเปฯ. กายสมฺปทา. วณฺณสมฺปทา. สทฺทสมฺปทา. คนฺธสมฺปทา. รสสมฺปทา. โภชนํ สทฺธาย ฯเปฯ ผลสมาปตฺติยา อุปนิสฺสย\-ปจฺจเยน ปจฺจโย. (1)}

\prose{โนอุปาทา ธมฺโม โนอุปาทา ธมฺมสฺส อุปนิสฺสย\-ปจฺจเยน ปจฺจโย. อารมฺมณู\-ปนิสฺสโย, อนนฺตรู\-ปนิสฺสโย, ปกตู\-ปนิสฺสโย ฯเปฯ. ปกตู\-ปนิสฺสโย\csromandash{}สทฺธํ อุปนิสฺสาย ทานํ เทติ ฯเปฯ สมาปตฺติํ อุปฺปาเทติ. มานํ ชปฺเปติ. ทิฏฺฐิํ คณฺหาติ. สีลํ ฯเปฯ. ปญฺญํ. ราคํ ฯเปฯ. ปตฺถนํ. กายิกํ สุขํ. กายิกํ ทุกฺขํ. อุตุํ. เสนาสนํ อุปนิสฺสาย ทานํ เทติ ฯเปฯ สํฆํ ภินฺทติ. สทฺธา ฯเปฯ. เสนาสนํ สทฺธาย ฯเปฯ ผลสมาปตฺติยา อุปนิสฺสย\-ปจฺจเยน ปจฺจโย. (1)}

\proseitem{399}{อุปาทา ธมฺโม โนอุปาทา ธมฺมสฺส ปุเรชาต\-ปจฺจเยน ปจฺจโย. อารมฺมณ\-ปุเรชาตํ, วตฺถุ\-ปุเรชาตํ.  อารมฺมณ\-ปุเรชาตํ\csromandash{}จกฺขุํ ฯเปฯ. วตฺถุํ อนิจฺจโต ฯเปฯ โทมนสฺสํ อุปฺปชฺชติ. ทิพฺเพน จกฺขุนา รูปํ ปสฺสติ. ทิพฺพาย โสตธาตุยา สทฺทํ สุณาติ. รูปายตนํ จกฺขุ\-วิญฺญาณสฺส ฯเปฯ. \csromanfolio{530}รสายตนํ ชิวฺหา\-วิญฺญาณสฺส ฯเปฯ. วตฺถุ\-ปุเรชาตํ\csromandash{}จกฺขายตนํ จกฺขุ\-วิญฺญาณสฺส ฯเปฯ กายายตนํ ฯเปฯ. วตฺถุ โนอุปาทา ขนฺธานํ ปุเรชาต\-ปจฺจเยน ปจฺจโย. (1)}

\prose{โนอุปาทา ธมฺโม โนอุปาทา ธมฺมสฺส ปุเรชาต\-ปจฺจเยน ปจฺจโย. อารมฺมณ\-ปุเรชาตํ\csromandash{}โผฏฺฐพฺเพ อนิจฺจโต ฯเปฯ โทมนสฺสํ อุปฺปชฺชติ. โผฏฺฐพฺพา\-ยตนํ กายวิญฺญาณสฺส ปุเรชาต\-ปจฺจเยน ปจฺจโย. (1)}

\prose{อุปาทา จ โนอุปาทา จ ธมฺมา โนอุปาทา ธมฺมสฺส ปุเรชาต\-ปจฺจเยน ปจฺจโย. อารมฺมณ\-ปุเรชาตํ, วตฺถุ\-ปุเรชาตํ. โผฏฺฐพฺพายตนญฺจ กายายตนญฺจ กายวิญฺญาณสฺส ปุเรชาต\-ปจฺจเยน ปจฺจโย. โผฏฺฐพฺพายตนญฺจ วตฺถุ จ โนอุปาทา ขนฺธานํ ปุเรชาต\-ปจฺจเยน ปจฺจโย. (1)}

\csromanheader{2}{\textbf{ปจฺฉาชาตเสวน}}
\proseitem{400}{โนอุปาทา ธมฺโม โนอุปาทา ธมฺมสฺส ปจฺฉาชาต\-ปจฺจเยน ปจฺจโย. ปจฺฉาชาตา โนอุปาทา ขนฺธา ปุเรชาตสฺส อิมสฺส โนอุปาทา กายสฺส ปจฺฉาชาต\-ปจฺจเยน ปจฺจโย. (มูลํ ปุจฺฉิตพฺพํ.) ปจฺฉาชาตา โนอุปาทา ขนฺธา ปุเรชาตสฺส อิมสฺส อุปาทา กายสฺส ปจฺฉาชาต\-ปจฺจเยน ปจฺจโย. (มูลํ ปุจฺฉิตพฺพํ.) ปจฺฉาชาตา โนอุปาทา ขนฺธา ปุเรชาตสฺส อิมสฺส อุปาทา กายสฺส จ โนอุปาทา กายสฺส จ ปจฺฉาชาต\-ปจฺจเยน ปจฺจโย. (3) อาเสวน\-ปจฺจเยน ปจฺจโย. เอกํ.}

\proseitem{401}{โนอุปาทา ธมฺโม โนอุปาทา ธมฺมสฺส กมฺมปจฺจเยน ปจฺจโย. สหชาตา, นานากฺขณิกา.  สหชาตา โนอุปาทา เจตนา สมฺปยุตฺตกานํ ขนฺธานํ โนอุปาทา จ จิตฺต\-สมุฏฺฐานานํ รูปานํ กมฺมปจฺจเยน ปจฺจโย. ปฏิสนฺธิ\-กฺขเณ ฯเปฯ. นานากฺขณิกา โนอุปาทา เจตนา วิปากานํ ขนฺธานํ โนอุปาทา จ กฏตฺตารูปานํ กมฺมปจฺจเยน ปจฺจโย. (1)}

\prose{โนอุปาทา ธมฺโม อุปาทา ธมฺมสฺส กมฺมปจฺจเยน ปจฺจโย. สหชาตา, นานากฺขณิกา.  สหชาตา โนอุปาทา เจตนา อุปาทา จิตฺต\-สมุฏฺฐานานํ รูปานํ กมฺมปจฺจเยน ปจฺจโย. ปฏิสนฺธิ\-กฺขเณ ฯเปฯ. นานากฺขณิกา โนอุปาทา เจตนา อุปาทา กฏตฺตารูปานํ กมฺมปจฺจเยน ปจฺจโย. (2)}

\prose{\csromanfolio{531}โนอุปาทา ธมฺโม อุปาทา จ โนอุปาทา จ ธมฺมสฺส กมฺมปจฺจเยน ปจฺจโย. สหชาตา, นานากฺขณิกา.  สหชาตา โนอุปาทา เจตนา สมฺปยุตฺตกานํ ขนฺธานํ อุปาทา จ โนอุปาทา จ จิตฺต\-สมุฏฺฐานานํ รูปานํ กมฺมปจฺจเยน ปจฺจโย. ปฏิสนฺธิ\-กฺขเณ ฯเปฯ. นานากฺขณิกา โนอุปาทา เจตนา วิปากานํ ขนฺธานํ อุปาทา จ โนอุปาทา จ กฏตฺตารูปานํ กมฺมปจฺจเยน ปจฺจโย. (3)}

\csromanheader{2}{\textbf{วิปาก}}
\proseitem{402}{โนอุปาทา ธมฺโม โนอุปาทา ธมฺมสฺส วิปาก\-ปจฺจเยน ปจฺจโย. วิปาโก โนอุปาทา เอโก ขนฺโธ ติณฺณนฺนํ ขนฺธานํ ฯเปฯ เทฺว ขนฺธา ฯเปฯ. ปฏิสนฺธิ\-กฺขเณ ฯเปฯ. ตีณิ.}

\csromanheader{2}{\textbf{อาหาร}}
\proseitem{403}{อุปาทา ธมฺโม อุปาทา ธมฺมสฺส อาหารปจฺจเยน ปจฺจโย. กพฬีกาโร อาหาโร อิมสฺส อุปาทา กายสฺส อาหารปจฺจเยน ปจฺจโย. (มูลํ ปุจฺฉิตพฺพํ.) กพฬีกาโร อาหาโร อิมสฺส โนอุปาทา กายสฺส อาหารปจฺจเยน ปจฺจโย. (มูลํ ปุจฺฉิตพฺพํ.) กพฬีกาโร อาหาโร อิมสฺส อุปาทา กายสฺส จ โนอุปาทา กายสฺส จ อาหารปจฺจเยน ปจฺจโย. (3)}

\prose{โนอุปาทา ธมฺโม โนอุปาทา ธมฺมสฺส อาหารปจฺจเยน ปจฺจโย. โนอุปาทา อาหารา สมฺปยุตฺตกานํ ขนฺธานํ โนอุปาทา จ จิตฺต\-สมุฏฺฐานานํ รูปานํ อาหารปจฺจเยน ปจฺจโย. ปฏิสนฺธิ\-กฺขเณ ฯเปฯ. (โนอุปาทามูลเก ตีณิ. ปฏิสนฺธิ\-กฺขเณ ตีณิปิ กาตพฺพา.) (3)}

\csromanheader{2}{\textbf{อินฺทฺริย}}
\proseitem{404}{อุปาทา ธมฺโม อุปาทา ธมฺมสฺส อินฺทฺริย\-ปจฺจเยน ปจฺจโย. รูปชีวิตินฺทฺริยํ อุปาทา กฏตฺตารูปานํ อินฺทฺริย\-ปจฺจเยน ปจฺจโย. (มูลํ กาตพฺพํ.) จกฺขุนฺทฺริยํ จกฺขุ\-วิญฺญาณสฺส ฯเปฯ. กายินฺทฺริยํ กายวิญฺญาณสฺส รูปชีวิตินฺทฺริยํ โนอุปาทา กฏตฺตารูปานํ อินฺทฺริย\-ปจฺจเยน ปจฺจโย. (มูลํ กาตพฺพํ.) รูปชีวิตินฺทฺริยํ อุปาทา จ โนอุปาทา จ กฏตฺตารูปานํ อินฺทฺริย\-ปจฺจเยน ปจฺจโย. (3)}

\prose{\csromanfolio{532}โนอุปาทา ธมฺโม โนอุปาทา ธมฺมสฺส อินฺทฺริย\-ปจฺจเยน ปจฺจโย. ตีณิ. ปฏิสนฺธิ\-กฺขเณ ฯเปฯ.}

\prose{อุปาทา จ โนอุปาทา จ ธมฺมา โนอุปาทา ธมฺมสฺส อินฺทฺริย\-ปจฺจเยน ปจฺจโย. จกฺขุนฺทฺริยญฺจ จกฺขุวิญฺญาณญฺจ จกฺขุ\-วิญฺญาณสหคตานํ ขนฺธานํ อินฺทฺริย\-ปจฺจเยน ปจฺจโย ฯเปฯ. กยินฺทฺริยญฺจ ฯเปฯ.}

\csromanheader{2}{\textbf{ฌานาทิ}}
\proseitem{405}{โน อุปาทา ธมฺโม โนอุปาทา ธมฺมสฺส ฌานปจฺจเยน ปจฺจโย. ตีณิ. มคฺคปจฺจเยน ปจฺจโย. ตีณิ. ปฏิสนฺธิ\-กฺขเณ ฯเปฯ. สมฺปยุตฺต\-ปจฺจเยน ปจฺจโย. เอกํ.}

\proseitem{406}{อุปาทา ธมฺโม โนอุปาทา ธมฺมสฺส วิปฺปยุตฺต\-ปจฺจเยน ปจฺจโย. สหชาตํ, ปุเรชาตํ.  สหชาตํ ปฏิสนฺธิ\-กฺขเณ วตฺถุ โนอุปาทา ขนฺธานํ วิปฺปยุตฺต\-ปจฺจเยน ปจฺจโย. ปุเรชาตํ จกฺขายตนํ จกฺขุ\-วิญฺญาณสฺส ฯเปฯ กายายตนํ ฯเปฯ. วตฺถุ โนอุปาทา ขนฺธานํ วิปฺปยุตฺต\-ปจฺจเยน ปจฺจโย. (1)}

\prose{โนอุปาทา ธมฺโม โนอุปาทา ธมฺมสฺส วิปฺปยุตฺต\-ปจฺจเยน ปจฺจโย. สหชาตํ, ปจฺฉาชาตํ.  สหชาตา โนอุปาทา ขนฺธา โนอุปาทา จิตฺต\-สมุฏฺฐานานํ รูปานํ วิปฺปยุตฺต\-ปจฺจเยน ปจฺจโย. ปฏิสนฺธิ\-กฺขเณ โนอุปาทา ขนฺธา โนอุปาทา กฏตฺตารูปานํ วิปฺปยุตฺต\-ปจฺจเยน ปจฺจโย. ปจฺฉาชาตา โนอุปาทา ขนฺธา ปุเรชาตสฺส อิมสฺส โนอุปาทา กายสฺส วิปฺปยุตฺต\-ปจฺจเยน ปจฺจโย. (1)}

\prose{โนอุปาทา ธมฺโม อุปาทา ธมฺมสฺส วิปฺปยุตฺต\-ปจฺจเยน ปจฺจโย. สหชาตํ, ปจฺฉาชาตํ.  สหชาตา โนอุปาทา ขนฺธา อุปาทา จิตฺต\-สมุฏฺฐานานํ รูปานํ วิปฺปยุตฺต\-ปจฺจเยน ปจฺจโย. ปฏิสนฺธิ\-กฺขเณ ฯเปฯ. ปจฺฉาชาตา โนอุปาทา ขนฺธา ปุเรชาตสฺส อิมสฺส อุปาทา กายสฺส วิปฺปยุตฺต\-ปจฺจเยน ปจฺจโย. (2)}

\prose{โนอุปาทา ธมฺโม อุปาทา ธมฺมสฺส จ โนอุปาทา ธมฺมสฺส จ วิปฺปยุตฺต\-ปจฺจเยน ปจฺจโย. สหชาตํ, ปจฺฉาชาตํ.  สหชาตา โนอุปาทา ขนฺธา}

\prose{\csromanfolio{533}อุปาทา จ โนอุปาทา จ จิตฺต\-สมุฏฺฐานานํ รูปานํ วิปฺปยุตฺต\-ปจฺจเยน ปจฺจโย. ปฏิสนฺธิ\-กฺขเณ ฯเปฯ. ปจฺฉาชาตา โนอุปาทา ขนฺธา ปุเรชาตสฺส อิมสฺส อุปาทา กายสฺส จ โนอุปาทา กายสฺส จ วิปฺปยุตฺต\-ปจฺจเยน ปจฺจโย. (3)}

\proseitem{407}{อุปาทา ธมฺโม อุปาทา ธมฺมสฺส อตฺถิปจฺจเยน ปจฺจโย. อาหารํ, อินฺทฺริยํ. \textbf{กพฬีกาโร อาหาโร} อิมสฺส อุปาทา กายสฺส อตฺถิปจฺจเยน ปจฺจโย. \textbf{รูปชีวิตินฺทฺริยํ} อุปาทา กฏตฺตารูปานํ อตฺถิปจฺจเยน ปจฺจโย. (1)}

\prose{อุปาทา ธมฺโม โนอุปาทา ธมฺมสฺส อตฺถิปจฺจเยน ปจฺจโย. สหชาตํ, ปุเรชาตํ, อาหารํ, อินฺทฺริยํ.  สหชาตํ ปฏิสนฺธิ\-กฺขเณ วตฺถุ โนอุปาทา ขนฺธานํ อตฺถิปจฺจเยน ปจฺจโย. ปุเรชาตํ จกฺขุํ อนิจฺจโต ฯเปฯ. (สํขิตฺตํ. ปุเรชาต\-สทิสํ. นินฺนานากรณํ.) กพฬีกาโร อาหาโร อิมสฺส โนอุปาทา กายสฺส อตฺถิปจฺจเยน ปจฺจโย. รูปชีวิตินฺทฺริยํ โนอุปาทา กฏตฺตารูปานํ อตฺถิปจฺจเยน ปจฺจโย. (2)}

\prose{อุปาทา ธมฺโม อุปาทา ธมฺมสฺส จ โนอุปาทา ธมฺมสฺส จ อตฺถิปจฺจเยน ปจฺจโย. อาหารํ, อินฺทฺริยํ. \textbf{กพฬีกาโร อาหาโร} อิมสฺส อุปาทา กายสฺส จ โนอุปาทา กายสฺส จ อตฺถิปจฺจเยน ปจฺจโย. \textbf{รูปชีวิตินฺทฺริยํ} อุปาทา จ โนอุปาทา จ กฏตฺตารูปานํ อตฺถิปจฺจเยน ปจฺจโย. (3)}

\proseitem{408}{โนอุปาทา ธมฺโม โนอุปาทา ธมฺมสฺส อตฺถิปจฺจเยน ปจฺจโย. สหชาตํ, ปุเรชาตํ, ปจฺฉาชาตํ.  สหชาโต โนอุปาทา เอโก ขนฺโธ ติณฺณนฺนํ ขนฺธานํ โนอุปาทา จ จิตฺต\-สมุฏฺฐานานํ รูปานํ อตฺถิปจฺจเยน ปจฺจโย ฯเปฯ เทฺว ขนฺธา ฯเปฯ. ปฏิสนฺธิ\-กฺขเณ ฯเปฯ. เอกํ มหาภูตํ ติณฺณนฺนํ มหาภูตานํ อตฺถิปจฺจเยน ปจฺจโย ฯเปฯ เทฺว มหาภูตา ทฺวินฺนํ มหาภูตานํ อตฺถิปจฺจเยน ปจฺจโย. (ยาว อสญฺญสตฺตา.) ปุเรชาตํ\csromandash{}โผฏฺฐพฺเพ อนิจฺจโต ฯเปฯ โทมนสฺสํ อุปฺปชฺชติ. โผฏฺฐพฺพา\-ยตนํ กายวิญฺญาณสฺส อตฺถิปจฺจเยน ปจฺจโย. ปจฺฉาชาตา โนอุปาทา ขนฺธา ปุเรชาตสฺส อิมสฺส โนอุปาทา กายสฺส อตฺถิปจฺจเยน ปจฺจโย. (1)}

\prose{โนอุปาทา ธมฺโม อุปาทา ธมฺมสฺส อตฺถิปจฺจเยน ปจฺจโย. สหชาตํ, ปจฺฉาชาตํ.  สหชาตา โนอุปาทา ขนฺธา อุปาทา จิตฺต\-สมุฏฺฐานานํ \csromanfolio{534}รูปานํ อตฺถิปจฺจเยน ปจฺจโย. ปฏิสนฺธิ\-กฺขเณ ฯเปฯ. ปจฺฉาชาตา โนอุปาทา ขนฺธา ปุเรชาตสฺส อิมสฺส อุปาทา กายสฺส อตฺถิปจฺจเยน ปจฺจโย. (2)}

\prose{โนอุปาทา ธมฺโม อุปาทา จ โนอุปาทา จ ธมฺมสฺส อตฺถิปจฺจเยน ปจฺจโย. สหชาตํ, ปจฺฉาชาตํ.  สหชาโต โนอุปาทา เอโก ขนฺโธ ติณฺณนฺนํ ขนฺธานํ อุปาทา จ โนอุปาทา จ จิตฺต\-สมุฏฺฐานานํ รูปานํ อตฺถิปจฺจเยน ปจฺจโย. เทฺว ขนฺธา ฯเปฯ. ปฏิสนฺธิ\-กฺขเณ ฯเปฯ. ปจฺฉาชาตา โนอุปาทา ขนฺธา ปุเรชาตสฺส อิมสฺส อุปาทา กายสฺส จ โนอุปาทา กายสฺส จ อตฺถิปจฺจเยน ปจฺจโย. (3)}

\proseitem{409}{อุปาทา จ โนอุปาทา จ ธมฺมา อุปาทา ธมฺมสฺส อตฺถิปจฺจเยน ปจฺจโย. ปจฺฉาชาตํ, อาหารํ, อินฺทฺริยํ.  ปจฺฉาชาตา โนอุปาทา ขนฺธา จ กพฬีกาโร อาหาโร จ อิมสฺส อุปาทา กายสฺส อตฺถิปจฺจเยน ปจฺจโย. ปจฺฉาชาตา โนอุปาทา ขนฺธา จ รูปชีวิตินฺทฺริยญฺจ อุปาทา กฏตฺตารูปานํ อตฺถิปจฺจเยน ปจฺจโย. (1)}

\prose{อุปาทา จ โนอุปาทา จ ธมฺมา โนอุปาทา ธมฺมสฺส อตฺถิปจฺจเยน ปจฺจโย. สหชาตํ, ปุเรชาตํ, ปจฺฉาชาตํ, อาหารํ, อินฺทฺริยํ.  สหชาโต จกฺขุ\-วิญฺญาณ\-สหคโต เอโก ขนฺโธ จ จกฺขายตนญฺจ ติณฺณนฺนํ ขนฺธานํ อตฺถิปจฺจเยน ปจฺจโย ฯเปฯ เทฺว ขนฺธา จ ฯเปฯ. โนอุปาทา เอโก ขนฺโธ จ วตฺถุ จ ติณฺณนฺนํ ขนฺธานํ อตฺถิปจฺจเยน ปจฺจโย ฯเปฯ เทฺว ขนฺธา จ ฯเปฯ. ปฏิสนฺธิ\-กฺขเณ โนอุปาทา เอโก ขนฺโธ จ วตฺถุ จ ติณฺณนฺนํ ขนฺธานํ อตฺถิปจฺจเยน ปจฺจโย ฯเปฯ เทฺว ขนฺธา จ ฯเปฯ. ปุเรชาตํ โผฏฺฐพฺพายตนญฺจ กายายตนญฺจ กายวิญฺญาณสฺส อตฺถิปจฺจเยน ปจฺจโย. โผฏฺฐพฺพายตนญฺจ วตฺถุ จ โนอุปาทา ขนฺธานํ อตฺถิปจฺจเยน ปจฺจโย. ปจฺฉาชาตา โนอุปาทา ขนฺธา จ กพฬีกาโร อาหาโร จ อิมสฺส โนอุปาทา กายสฺส อตฺถิปจฺจเยน ปจฺจโย. ปจฺฉาชาตา โนอุปาทา ขนฺธา จ รูปชีวิตินฺทฺริยญฺจ โนอุปาทา กฏตฺตารูปานํ อตฺถิปจฺจเยน ปจฺจโย. (2)}

\prose{อุปาทา จ โนอุปาทา จ ธมฺมา อุปาทา ธมฺมสฺส จ โนอุปาทา ธมฺมสฺส จ อตฺถิปจฺจเยน ปจฺจโย. ปจฺฉาชาตํ, อาหารํ, อินฺทฺริยํ.  \textbf{ปจฺฉาชาตา} โนอุปาทา ขนฺธา จ กพฬีกาโร อาหาโร จ อิมสฺส อุปาทา กายสฺส จ โนอุปาทา กายสฺส จ อตฺถิปจฺจเยน ปจฺจโย.}

\prose{\csromanfolio{535}ปจฺฉาชาตา โนอุปาทา ขนฺธา จ รูปชีวิตินฺทฺริยญฺจ อุปาทา จ โนอุปาทา จ กฏตฺตารูปานํ อตฺถิปจฺจเยน ปจฺจโย. (3)}

\csromanmarkh{2. ปจฺจยานุโลม}\csromanmarkhii{2. สงฺขฺยาวาร}\csromanheadernomark{2}{2. \textbf{ปจฺจยานุโลม 2. สงฺขฺยาวาร}}
\proseitem{410}{เหตุยา ตีณิ, อารมฺมเณ เทฺว, อธิปติยา จตฺตาริ, อนนฺตเร เอกํ, สมนนฺตเร เอกํ, สหชาเต ปญฺจ, อญฺญมญฺเญ ปญฺจ, นิสฺสเย ปญฺจ, อุปนิสฺสเย เทฺว, ปุเรชาเต ตีณิ, ปจฺฉาชาเต ตีณิ, อาเสวเน เอกํ, กมฺเม ตีณิ, วิปาเก ตีณิ, อาหาเร ฉ, อินฺทฺริเย สตฺต, ฌาเน ตีณิ, มคฺเค ตีณิ, สมฺปยุตฺเต เอกํ, วิปฺปยุตฺเต จตฺตาริ, อตฺถิยา นว, นตฺถิยา เอกํ, วิคเต เอกํ, อวิคเต นว.}

\vspace{\tipitakaparskip}
\csromancenter{อนุโลมํ.}
\proseitem{411}{อุปาทา ธมฺโม อุปาทา ธมฺมสฺส อาหารปจฺจเยน ปจฺจโย. อินฺทฺริย\-ปจฺจเยน ปจฺจโย. (1)}

\prose{อุปาทา ธมฺโม โนอุปาทา ธมฺมสฺส อารมฺมณ\-ปจฺจเยน ปจฺจโย. สหชาต\-ปจฺจเยน ปจฺจโย. อุปนิสฺสย\-ปจฺจเยน ปจฺจโย. ปุเรชาต\-ปจฺจเยน ปจฺจโย. อาหารปจฺจเยน ปจฺจโย. อินฺทฺริย\-ปจฺจเยน ปจฺจโย. (2)}

\prose{อุปาทา ธมฺโม อุปาทา ธมฺมสฺส จ โนอุปาทา ธมฺมสฺส จ อาหารปจฺจเยน ปจฺจโย. อินฺทฺริย\-ปจฺจเยน ปจฺจโย. (3)}

\proseitem{412}{โนอุปาทา ธมฺโม โนอุปาทา ธมฺมสฺส อารมฺมณ\-ปจฺจเยน ปจฺจโย. สหชาต\-ปจฺจเยน ปจฺจโย. อุปนิสฺสย\-ปจฺจเยน ปจฺจโย. ปุเรชาต\-ปจฺจเยน ปจฺจโย. ปจฺฉาชาต\-ปจฺจเยน ปจฺจโย. กมฺมปจฺจเยน ปจฺจโย. (1)}

\prose{โนอุปาทา ธมฺโม อุปาทา ธมฺมสฺส สหชาต\-ปจฺจเยน ปจฺจโย. ปจฺฉาชาต\-ปจฺจเยน ปจฺจโย. กมฺมปจฺจเยน ปจฺจโย. (2)}

\prose{\csromanfolio{536}โนอุปาทา ธมฺโม อุปาทา ธมฺมสฺส จ โนอุปาทา ธมฺมสฺส จ สหชาต\-ปจฺจเยน ปจฺจโย. ปจฺฉาชาต\-ปจฺจเยน ปจฺจโย. กมฺมปจฺจเยน ปจฺจโย. (3)}

\prose{อุปาทา จ โนอุปาทา จ ธมฺมา อุปาทา ธมฺมสฺส ปจฺฉาชาตํ, อาหารํ, อินฺทฺริยํ. (1)}

\prose{อุปาทา จ โนอุปาทา จ ธมฺมา โนอุปาทา ธมฺมสฺส สหชาตํ, ปุเรชาตํ, ปจฺฉาชาตํ, อาหารํ, อินฺทฺริยํ. (2)}

\prose{อุปาทา จ โนอุปาทา จ ธมฺมา อุปาทา ธมฺมสฺส จ โนอุปาทา ธมฺมสฺส จ ปจฺฉาชาตํ, อาหารํ, อินฺทฺริยํ. (3)}

\csromanmarkh{2. ปจฺจย\-ปจฺจนีย}\csromanmarkhii{2. สงฺขฺยาวาร}\csromanheadernomark{2}{2. ปจฺจย\-ปจฺจนีย 2. สงฺขฺยาวาร}
\proseitem{413}{นเหตุยา นว, นอารมฺมเณ นว, (สพฺพตฺถ นว.) นสมฺปยุตฺเต นว, นวิปฺปยุตฺเต ฉ, โนอตฺถิยา จตฺตาริ, โนนตฺถิยา นว, โนวิคเต นว, โนอวิคเต จตฺตาริ.}

\proseitem{414}{เหตุปจฺจยา นอารมฺมเณ ตีณิ, นอธิปติยา ตีณิ, นอนนฺตเร ตีณิ, นสมนนฺตเร ตีณิ, นอญฺญมญฺเญ ตีณิ, นอุปนิสฺสเย ตีณิ, (สพฺพตฺถ ตีณิ.) นสมฺปยุตฺเต ตีณิ, นวิปฺปยุตฺเต เอกํ, โนนตฺถิยา ตีณิ, โนวิคเต ตีณิ.}

\proseitemwithcloser{415}{นเหตุปจฺจยา อารมฺมเณ เทฺว, อธิปติยา จตฺตาริ ฯเปฯ. (อนุโลมมาติกา วิตฺถาเรตพฺพา.) ฯเปฯ อวิคเต นว.}{อุปาทาทุกํ นิฏฺฐิตํ.}
\csromanmatikaanchor{mtk.67}
\csromantocmarkref{subsection}{68. อุปาทินฺนทุก}{mtk.67}
\bookmark[dest={mtk.67},level=2]{68. อุปาทินฺนทุก}
\csromanmarkh{68. อุปาทินฺนทุก}\csromanmarkhii{68. อุปาทินฺนทุก 1. ปฏิจฺจวาร}\csromanheadernomark{2}{68. อุปาทินฺนทุก \textbf{68. อุปาทินฺนทุก 1. ปฏิจฺจวาร}}
\csromanmarkh{1. ปจฺจยานุโลม}\csromanmarkhii{1. วิภงฺควาร}\csromanheadernomark{2}{1. \textbf{ปจฺจยานุโลม 1. วิภงฺควาร}}
\csromanheader{2}{\textbf{เหตุ}}
\proseitem{416}{\csromanfolio{537}อุปาทินฺนํ ธมฺมํ ปฏิจฺจ อุปาทินฺโน ธมฺโม อุปฺปชฺชติ เหตุปจฺจยา. อุปาทินฺนํ เอกํ ขนฺธํ ปฏิจฺจ ตโย ขนฺธา ฯเปฯ เทฺว ขนฺเธ ฯเปฯ. ปฏิสนฺธิ\-กฺขเณ อุปาทินฺนํ เอกํ ขนฺธํ ปฏิจฺจ ตโย ขนฺธา กฏตฺตา จ รูปํ ฯเปฯ เทฺว ขนฺเธ ฯเปฯ. ขนฺเธ ปฏิจฺจ วตฺถุ, วตฺถุํ ปฏิจฺจ ขนฺธา. เอกํ มหาภูตํ ฯเปฯ มหาภูเต ปฏิจฺจ กฏตฺตารูปํ, อุปาทารูปํ. (1)}

\prose{อุปาทินฺนํ ธมฺมํ ปฏิจฺจ อนุปาทินฺโน ธมฺโม อุปฺปชฺชติ เหตุปจฺจยา. อุปาทินฺเน ขนฺเธ ปฏิจฺจ จิตฺต\-สมุฏฺฐานํ รูปํ. (2)}

\prose{อุปาทินฺนํ ธมฺมํ ปฏิจฺจ อุปาทินฺโน จ อนุปาทินฺโน จ ธมฺมา อุปฺปชฺชนฺติ เหตุปจฺจยา. อุปาทินฺนํ เอกํ ขนฺธํ ปฏิจฺจ ตโย ขนฺธา จิตฺต\-สมุฏฺฐานญฺจ รูปํ ฯเปฯ เทฺว ขนฺเธ ฯเปฯ. (3)}

\prose{อนุปาทินฺนํ ธมฺมํ ปฏิจฺจ อนุปาทินฺโน ธมฺโม อุปฺปชฺชติ เหตุปจฺจยา. อนุปาทินฺนํ เอกํ ขนฺธํ ปฏิจฺจ ตโย ขนฺธา จิตฺต\-สมุฏฺฐานญฺจ รูปํ ฯเปฯ เทฺว ขนฺเธ ฯเปฯ. เอกํ มหาภูตํ ปฏิจฺจ ฯเปฯ มหาภูเต ปฏิจฺจ จิตฺต\-สมุฏฺฐานํ รูปํ, อุปาทารูปํ. (1)}

\prose{อุปาทินฺนญฺจ อนุปาทินฺนญฺจ ธมฺมํ ปฏิจฺจ อนุปาทินฺโน ธมฺโม อุปฺปชฺชติ เหตุปจฺจยา. อุปาทินฺเน ขนฺเธ จ มหาภูเต จ ปฏิจฺจ จิตฺต\-สมุฏฺฐานํ รูปํ. (1)}

\proseitem{417}{อุปาทินฺนํ ธมฺมํ ปฏิจฺจ อุปาทินฺโน ธมฺโม อุปฺปชฺชติ อารมฺมณ\-ปจฺจยา. อุปาทินฺนํ เอกํ ขนฺธํ ปฏิจฺจ ตโย ขนฺธา ฯเปฯ เทฺว ขนฺเธ ฯเปฯ. ปฏิสนฺธิ\-กฺขเณ ฯเปฯ. วตฺถุํ ปฏิจฺจ ขนฺธา. (1)}

\prose{อนุปาทินฺนํ ธมฺมํ ปฏิจฺจ อนุปาทินฺโน ธมฺโม อุปฺปชฺชติ อารมฺมณ\-ปจฺจยา. อนุปาทินฺนํ เอกํ ขนฺธํ ปฏิจฺจ ตโย ขนฺธา ฯเปฯ เทฺว ขนฺเธ ฯเปฯ. (1)}

\proseitem{418}{\csromanfolio{538}อนุปาทินฺนํ ธมฺมํ ปฏิจฺจ อนุปาทินฺโน ธมฺโม อุปฺปชฺชติ อธิปติ\-ปจฺจยา. อนุปาทินฺนํ เอกํ ขนฺธํ ปฏิจฺจ ตโย ขนฺธา จิตฺต\-สมุฏฺฐานญฺจ รูปํ ฯเปฯ เทฺว ขนฺเธ ฯเปฯ. เอกํ มหาภูตํ ปฏิจฺจ ฯเปฯ มหาภูเต ปฏิจฺจ จิตฺต\-สมุฏฺฐานํ รูปํ, อุปาทารูปํ. (สํขิตฺตํ.)}

\csromanmarkh{1. ปจฺจยานุโลม}\csromanmarkhii{2. สงฺขฺยาวาร}\csromanheadernomark{2}{1. \textbf{ปจฺจยานุโลม 2. สงฺขฺยาวาร}}
\proseitem{419}{เหตุยา ปญฺจ, อารมฺมเณ เทฺว, อธิปติยา เอกํ, อนนฺตเร เทฺว, สมนนฺตเร เทฺว, สหชาเต ปญฺจ, อญฺญมญฺเญ เทฺว, นิสฺสเย ปญฺจ, อุปนิสฺสเย เทฺว, ปุเรชาเต เทฺว, อาเสวเน เอกํ, กมฺเม ปญฺจ, วิปาเก ปญฺจ, อาหาเร ปญฺจ, อินฺทฺริเย ปญฺจ, ฌาเน ปญฺจ, มคฺเค ปญฺจ, สมฺปยุตฺเต เทฺว, วิปฺปยุตฺเต ปญฺจ, อตฺถิยา ปญฺจ, นตฺถิยา เทฺว, วิคเต เทฺว, อวิคเต ปญฺจ.}

\csromanmarkh{2. ปจฺจย\-ปจฺจนีย}\csromanmarkhii{1. วิภงฺควาร}\csromanheadernomark{2}{2. ปจฺจย\-ปจฺจนีย 1. วิภงฺควาร}
\proseitem{420}{อุปาทินฺนํ ธมฺมํ ปฏิจฺจ อุปาทินฺโน ธมฺโม อุปฺปชฺชติ นเหตุปจฺจยา. อเหตุกํ อุปาทินฺนํ เอกํ ขนฺธํ ปฏิจฺจ ตโย ขนฺธา ฯเปฯ เทฺว ขนฺเธ ฯเปฯ. อเหตุก\-ปฏิสนฺธิ\-กฺขเณ อุปาทินฺนํ เอกํ ขนฺธํ ปฏิจฺจ ตโย ขนฺธา กฏตฺตา จ รูปํ ฯเปฯ เทฺว ขนฺเธ ฯเปฯ. ขนฺเธ ปฏิจฺจ วตฺถุ, วตฺถุํ ปฏิจฺจ ขนฺธา. เอกํ มหาภูตํ ปฏิจฺจ ฯเปฯ มหาภูเต ปฏิจฺจ กฏตฺตารูปํ, อุปาทารูปํ. อสญฺญสตฺตานํ เอกํ มหาภูตํ ปฏิจฺจ ฯเปฯ มหาภูเต ปฏิจฺจ กฏตฺตารูปํ, อุปาทารูปํ. (1)}

\prose{อุปาทินฺนํ ธมฺมํ ปฏิจฺจ อนุปาทินฺโน ธมฺโม อุปฺปชฺชติ นเหตุปจฺจยา. อเหตุเก อุปาทินฺเน ขนฺเธ ปฏิจฺจ จิตฺต\-สมุฏฺฐานํ รูปํ.}

\prose{อุปาทินฺนํ ธมฺมํ ปฏิจฺจ อุปาทินฺโน จ อนุปาทินฺโน จ ธมฺมา อุปฺปชฺชนฺติ นเหตุปจฺจยา. อเหตุกํ อุปาทินฺนํ เอกํ ขนฺธํ ปฏิจฺจ ตโย ขนฺธา จิตฺต\-สมุฏฺฐานญฺจ รูปํ ฯเปฯ เทฺว ขนฺเธ ฯเปฯ. (3)}

\proseitem{421}{\csromanfolio{539}อนุปาทินฺนํ ธมฺมํ ปฏิจฺจ อนุปาทินฺโน ธมฺโม อุปฺปชฺชติ นเหตุปจฺจยา. อเหตุกํ อนุปาทินฺนํ เอกํ ขนฺธํ ปฏิจฺจ ตโย ขนฺธา จิตฺต\-สมุฏฺฐานญฺจ รูปํ ฯเปฯ เทฺว ขนฺเธ ฯเปฯ. เอกํ มหาภูตํ ฯเปฯ มหาภูเต ปฏิจฺจ จิตฺต\-สมุฏฺฐานํ รูปํ, อุปาทารูปํ. พาหิรํ. อาหารสมุฏฺฐานํ. อุตุสมุฏฺฐานํ เอกํ มหาภูตํ ฯเปฯ มหาภูเต ปฏิจฺจ อุปาทารูปํ. วิจิกิจฺฉา\-สหคเต อุทฺธจฺจ\-สหคเต ขนฺเธ ปฏิจฺจ วิจิกิจฺฉา\-สหคโต อุทฺธจฺจ\-สหคโต โมโห. (1)}

\prose{อุปาทินฺนญฺจ อนุปาทินฺนญฺจ ธมฺมํ ปฏิจฺจ อนุปาทินฺโน ธมฺโม อุปฺปชฺชติ นเหตุปจฺจยา. อเหตุเก อุปาทินฺเน ขนฺเธ จ มหาภูเต จ ปฏิจฺจ จิตฺต\-สมุฏฺฐานํ รูปํ. (1)}

\prose{นอารมฺมณ}

\proseitem{422}{อุปาทินฺนํ ธมฺมํ ปฏิจฺจ อุปาทินฺโน ธมฺโม อุปฺปชฺชติ นอารมฺมณ\-ปจฺจยา. ปฏิสนฺธิ\-กฺขเณ อุปาทินฺเน ขนฺเธ ปฏิจฺจ กฏตฺตารูปํ. ขนฺเธ ปฏิจฺจ วตฺถุ. เอกํ มหาภูตํ ฯเปฯ มหาภูเต ปฏิจฺจ กฏตฺตารูปํ, อุปาทารูปํ. อสญฺญสตฺตานํ เอกํ มหาภูตํ ฯเปฯ มหาภูเต ปฏิจฺจ กฏตฺตารูปํ, อุปาทารูปํ. (1)}

\prose{อุปาทินฺนํ ธมฺมํ ปฏิจฺจ อนุปาทินฺโน ธมฺโม อุปฺปชฺชติ นอารมฺมณ\-ปจฺจยา. อุปาทินฺเน ขนฺเธ ปฏิจฺจ จิตฺต\-สมุฏฺฐานํ รูปํ.}

\prose{อนุปาทินฺนํ ธมฺมํ ปฏิจฺจ อนุปาทินฺโน ธมฺโม อุปฺปชฺชติ นอารมฺมณ\-ปจฺจยา. อนุปาทินฺเน ขนฺเธ ปฏิจฺจ จิตฺต\-สมุฏฺฐานํ รูปํ. เอกํ มหาภูตํ ฯเปฯ. พาหิรํ. อาหารสมุฏฺฐานํ. อุตุสมุฏฺฐานํ เอกํ มหาภูตํ ฯเปฯ มหาภูเต ปฏิจฺจ อุปาทารูปํ. (1)}

\prose{อุปาทินฺนญฺจ อนุปาทินฺนญฺจ ธมฺมํ ปฏิจฺจ อนุปาทินฺโน ธมฺโม อุปฺปชฺชติ นอารมฺมณ\-ปจฺจยา. อุปาทินฺเน ขนฺเธ จ มหาภูเต จ ปฏิจฺจ จิตฺต\-สมุฏฺฐานํ รูปํ. (1)}

\prose{นอธิปตฺยาทิ}

\proseitem{423}{อุปาทินฺนํ ธมฺมํ ปฏิจฺจ อุปาทินฺโน ธมฺโม อุปฺปชฺชติ นอธิปติ\-ปจฺจยา. นอนนฺตร\-ปจฺจยา. นสมนนฺตร\-ปจฺจยา. นอญฺญมญฺญ\-ปจฺจยา. นอุปนิสฺสยปจฺจยา.}

\csromanheader{2}{\textbf{นปุเรชาตาทิ}}
\proseitem{424}{อุปาทินฺนํ ธมฺมํ ปฏิจฺจ อุปาทินฺโน ธมฺโม อุปฺปชฺชติ นปุเรชาต\-ปจฺจยา. อรูเป อุปาทินฺนํ เอกํ ขนฺธํ ปฏิจฺจ ตโย ขนฺธา ฯเปฯ เทฺว ขนฺเธ ฯเปฯ. \csromanfolio{540}ปฏิสนฺธิ\-กฺขเณ อุปาทินฺนํ เอกํ ขนฺธํ ปฏิจฺจ ตโย ขนฺธา กฏตฺตา จ รูปํ ฯเปฯ เทฺว ขนฺเธ ฯเปฯ. (ยาว อสญฺญสตฺตา.) (1)}

\prose{อุปาทินฺนํ ธมฺมํ ปฏิจฺจ อนุปาทินฺโน ธมฺโม อุปฺปชฺชติ นปุเรชาต\-ปจฺจยา. อุปาทินฺเน ขนฺเธ ปฏิจฺจ จิตฺต\-สมุฏฺฐานํ รูปํ. (2)}

\prose{อนุปาทินฺนํ ธมฺมํ ปฏิจฺจ อนุปาทินฺโน ธมฺโม อุปฺปชฺชติ นปุเรชาต\-ปจฺจยา. อรูเป อนุปาทินฺนํ เอกํ ขนฺธํ ปฏิจฺจ ตโย ขนฺธา ฯเปฯ เทฺว ขนฺเธ ฯเปฯ. อนุปาทินฺเน ขนฺเธ ปฏิจฺจ จิตฺต\-สมุฏฺฐานํ รูปํ. เอกํ มหาภูตํ ฯเปฯ. (ยาว อสญฺญสตฺตา.) (1)}

\prose{อุปาทินฺนญฺจ อนุปาทินฺนญฺจ ธมฺมํ ปฏิจฺจ อนุปาทินฺโน ธมฺโม อุปฺปชฺชติ นปุเรชาต\-ปจฺจยา. อุปาทินฺเน ขนฺเธ จ มหาภูเต จ ปฏิจฺจ จิตฺต\-สมุฏฺฐานํ รูปํ. นปจฺฉาชาต\-ปจฺจยา. นอาเสวน\-ปจฺจยา. (1)}

\csromanheader{2}{\textbf{นกมฺม}}
\proseitem{425}{อนุปาทินฺนํ ธมฺมํ ปฏิจฺจ อนุปาทินฺโน ธมฺโม อุปฺปชฺชติ นกมฺมปจฺจยา. อนุปาทินฺเน ขนฺเธ ปฏิจฺจ อนุปาทินฺนา เจตนา. พาหิรํ. อาหารสมุฏฺฐานํ อุตุสมุฏฺฐานํ ฯเปฯ. มหาภูเต ปฏิจฺจ อุปาทารูปํ. (1)}

\csromanheader{2}{\textbf{นวิปาก}}
\proseitem{426}{อุปาทินฺนํ ธมฺมํ ปฏิจฺจ อุปาทินฺโน ธมฺโม อุปฺปชฺชติ นวิปาก\-ปจฺจยา. อสญฺญสตฺตานํ เอกํ มหาภูตํ ฯเปฯ มหาภูเต ปฏิจฺจ กฏตฺตารูปํ, อุปาทารูปํ. (1)}

\prose{อนุปาทินฺนํ ธมฺมํ ปฏิจฺจ อนุปาทินฺโน ธมฺโม อุปฺปชฺชติ นวิปาก\-ปจฺจยา. อนุปาทินฺนํ เอกํ ขนฺธํ ปฏิจฺจ ตโย ขนฺธา จิตฺต\-สมุฏฺฐานญฺจ รูปํ ฯเปฯ เทฺว ขนฺเธ ฯเปฯ. เอกํ มหาภูตํ ฯเปฯ. (ยาว อุตุสมุฏฺฐานํ.)}

\prose{นอาหาร}

\proseitem{427}{อุปาทินฺนํ ธมฺมํ ปฏิจฺจ อุปาทินฺโน ธมฺโม อุปฺปชฺชติ นอาหารปจฺจยา. อสญฺญสตฺตานํ เอกํ มหาภูตํ. (1)}

\prose{อนุปาทินฺนํ ธมฺมํ ปฏิจฺจ อนุปาทินฺโน ธมฺโม อุปฺปชฺชติ นอาหารปจฺจยา. พาหิรํ. อุตุสมุฏฺฐานํ ฯเปฯ. (1)}

\csromanheader{2}{68. อุปาทินฺนทุก นอินฺทฺริย}
\proseitem{428}{\csromanfolio{541}อุปาทินฺนํ ธมฺมํ ปฏิจฺจ อุปาทินฺโน ธมฺโม อุปฺปชฺชติ นอินฺทฺริยปจฺจยา. อสญฺญสตฺตานํ มหาภูเต ปฏิจฺจ รูปชีวิตินฺทฺริยํ. (1)}

\prose{อนุปาทินฺนํ ธมฺมํ ปฏิจฺจ อนุปาทินฺโน ธมฺโม อุปฺปชฺชติ นอินฺทฺริยปจฺจยา. พาหิรํ. อาหารสมุฏฺฐานํ, อุตุสมุฏฺฐานํ ฯเปฯ. (1)}

\csromanheader{2}{\textbf{นฌานาทิ}}
\proseitem{429}{อุปาทินฺนํ ธมฺมํ ปฏิจฺจ อุปาทินฺโน ธมฺโม อุปฺปชฺชติ นฌานปจฺจยา. ปญฺจ\-วิญฺญาณสหคตํ เอกํ ขนฺธํ ปฏิจฺจ ตโย ขนฺธา ฯเปฯ เทฺว ขนฺเธ ฯเปฯ. อสญฺญสตฺตานํ ฯเปฯ.}

\prose{อนุปาทินฺนํ ธมฺมํ ปฏิจฺจ อนุปาทินฺโน ธมฺโม อุปฺปชฺชติ นฌานปจฺจยา. พาหิรํ. อาหารสมุฏฺฐานํ. อุตุสมุฏฺฐานํ ฯเปฯ.}

\prose{อุปาทินฺนํ ธมฺมํ ปฏิจฺจ อุปาทินฺโน ธมฺโม อุปฺปชฺชติ นมคฺคปจฺจยา. (นเหตุสทิสํ. โมโห นตฺถิ.) นสมฺปยุตฺต\-ปจฺจยา.}

\csromanheader{2}{\textbf{นวิปฺปยุตฺตาทิ}}
\proseitem{430}{อุปาทินฺนํ ธมฺมํ ปฏิจฺจ อุปาทินฺโน ธมฺโม อุปฺปชฺชติ นวิปฺปยุตฺต\-ปจฺจยา. อรูเป อุปาทินฺนํ เอกํ ขนฺธํ ปฏิจฺจ ตโย ขนฺธา ฯเปฯ เทฺว ขนฺเธ ฯเปฯ. อสญฺญสตฺตานํ เอกํ มหาภูตํ ปฏิจฺจ ฯเปฯ.}

\prose{อนุปาทินฺนํ ธมฺมํ ปฏิจฺจ อนุปาทินฺโน ธมฺโม อุปฺปชฺชติ นวิปฺปยุตฺต\-ปจฺจยา. อรูเป อนุปาทินฺนํ เอกํ ขนฺธํ ปฏิจฺจ ตโย ขนฺธา ฯเปฯ เทฺว ขนฺเธ ฯเปฯ. พาหิรํ. อาหารสมุฏฺฐานํ. อุตุสมุฏฺฐานํ ฯเปฯ. โนนตฺถิปจฺจยา. โนวิคต\-ปจฺจยา.}

\csromanmarkh{2. ปจฺจย\-ปจฺจนีย}\csromanmarkhii{2. สงฺขฺยาวาร}\csromanheadernomark{2}{2. ปจฺจย\-ปจฺจนีย 2. สงฺขฺยาวาร}
\proseitem{431}{นเหตุยา ปญฺจ, นอารมฺมเณ จตฺตาริ, นอธิปติยา ปญฺจ, นอนนฺตเร จตฺตาริ, นสมนนฺตเร จตฺตาริ, นอญฺญมญฺเญ จตฺตาริ, นอุปนิสฺสเย จตฺตาริ, นปุเรชาเต จตฺตาริ, นปจฺฉาชาเต ปญฺจ, นอาเสวเน ปญฺจ, \csromanfolio{542}นกมฺเม เอกํ, นวิปาเก เทฺว, นอาหาเร เทฺว, นอินฺทฺริเย เทฺว, นฌาเน เทฺว, นมคฺเค ปญฺจ, นสมฺปยุตฺเต จตฺตาริ, นวิปฺปยุตฺเต เทฺว, โนนตฺถิยา จตฺตาริ, โนวิคเต จตฺตาริ.}

\vspace{\tipitakaparskip}
\csromancenter{\textbf{เหตุ}ปจฺจยา นอารมฺมเณ จตฺตาริ, นอธิปติยา ปญฺจ ฯเปฯ นปุเรชาเต จตฺตาริ, นปจฺฉาชาเต ปญฺจ, นอาเสวเน ปญฺจ, นกมฺเม เอกํ, นวิปาเก เอกํ ฯเปฯ นสมฺปยุตฺเต จตฺตาริ, นวิปฺปยุตฺเต เทฺว, โนนตฺถิยา จตฺตาริ, โนวิคเต จตฺตาริ.}
\proseitem{433}{นเหตุปจฺจยา อารมฺมเณ เทฺว, อนนฺตเร เทฺว, สมนนฺตเร เทฺว, สหชาเต ปญฺจ, อญฺญมญฺเญ เทฺว, นิสฺสเย ปญฺจ, อุปนิสฺสเย เทฺว, ปุเรชาเต เทฺว, อาเสวเน เอกํ, กมฺเม ปญฺจ, วิปาเก ปญฺจ ฯเปฯ มคฺเค เอกํ, สมฺปยุตฺเต เทฺว ฯเปฯ อวิคเต ปญฺจ.}

\vspace{\tipitakaparskip}
\csromancenter{(สหชาตวาโร ปฏิจฺจ\-วาร\-สทิโส.)}
\csromanmarkh{68. อุปาทินฺนทุก}\csromanmarkhii{3. ปจฺจยวาร}\csromanheadernomark{2}{68. \textbf{อุปาทินฺนทุก 3. ปจฺจยวาร}}
\csromanmarkh{1. ปจฺจยานุโลม}\csromanmarkhii{1. วิภงฺควาร}\csromanheadernomark{2}{1. \textbf{ปจฺจยานุโลม 1. วิภงฺควาร}}
\vspace{\tipitakaparskip}
\csromancenter{\textbf{เหตุ}}
\proseitem{434}{อุปาทินฺนํ ธมฺมํ ปจฺจยา อุปาทินฺโน ธมฺโม อุปฺปชฺชติ เหตุปจฺจยา. อุปาทินฺนํ เอขํ ขนฺธํ ปจฺจยา ตโย ขนฺธา ฯเปฯ เทฺว ขนฺเธ ฯเปฯ. ปฏิสนฺธิ\-กฺขเณ อุปาทินฺนํ เอกํ ขนฺธํ ปจฺจยา ตโย ขนฺธา กฏตฺตา จ รูปํ ฯเปฯ เทฺว ขนฺเธ ฯเปฯ. ขนฺเธ ปจฺจยา วตฺถุ, วตฺถุํ ปจฺจยา ขนฺธา. เอกํ มหาภูตํ ฯเปฯ มหาภูเต ปจฺจยา กฏตฺตารูปํ, อุปาทารูปํ. วตฺถุํ ปจฺจยา อุปาทินฺนา ขนฺธา. (1)}

\prose{อุปาทินฺนํ ธมฺมํ ปจฺจยา อนุปาทินฺโน ธมฺโม อุปฺปชฺชติ เหตุปจฺจยา. อุปาทินฺเน ขนฺเธ ปจฺจยา จิตฺต\-สมุฏฺฐานํ รูปํ. วตฺถุํ ปจฺจยา อนุปาทินฺนา ขนฺธา. (2)}

\prose{\csromanfolio{543}อุปาทินฺนํ ธมฺมํ ปจฺจยา อุปาทินฺโน จ อนุปาทินฺโน จ ธมฺมา อุปฺปชฺชนฺติ เหตุปจฺจยา. อุปาทินฺนํ เอกํ ขนฺธํ ปจฺจยา ตโย ขนฺธา จิตฺต\-สมุฏฺฐานญฺจ รูปํ ฯเปฯ เทฺว ขนฺเธ ฯเปฯ. (3)}

\prose{อนุปาทินฺนํ ธมฺมํ ปจฺจยา อนุปาทินฺโน ธมฺโม อุปฺปชฺชติ เหตุปจฺจยา. อนุปาทินฺนํ เอกํ ขนฺธํ ปจฺจยา ตโย ขนฺธา จิตฺต\-สมุฏฺฐานญฺจ รูปํ ฯเปฯ เทฺว ขนฺเธ ฯเปฯ. เอกํ มหาภูตํ ฯเปฯ มหาภูเต ปจฺจยา จิตฺต\-สมุฏฺฐานํ รูปํ, อุปาทารูปํ. (1)}

\prose{อุปาทินฺนญฺจ อนุปาทินฺนญฺจ ธมฺมํ ปจฺจยา อนุปาทินฺโน ธมฺโม อุปฺปชฺชติ เหตุปจฺจยา. อุปาทินฺเน ขนฺเธ จ มหาภูเต จ ปจฺจยา จิตฺต\-สมุฏฺฐานํ รูปํ. อนุปาทินฺนํ เอกํ ขนฺธญฺจ วตฺถุญฺจ ปจฺจยา ตโย ขนฺธา ฯเปฯ เทฺว ขนฺเธ จ ฯเปฯ. (1)}

\proseitem{435}{อุปาทินฺนํ ธมฺมํ ปจฺจยา อุปาทินฺโน ธมฺโม อุปฺปชฺชติ อารมฺมณ\-ปจฺจยา. อุปาทินฺนํ เอกํ ขนฺธํ ปจฺจยา ตโย ขนฺธา ฯเปฯ เทฺว ขนฺเธ ฯเปฯ. ปฏิสนฺธิ\-กฺขเณ ฯเปฯ. วตฺถุํ ปจฺจยา ขนฺธา. จกฺขายตนํ ปจฺจยา จกฺขุวิญฺญาณํ ฯเปฯ. กายายตนํ ฯเปฯ. วตฺถุํ ปจฺจยา อุปาทินฺนา ขนฺธา. (1)}

\prose{อุปาทินฺนํ ธมฺมํ ปจฺจยา อนุปาทินฺโน ธมฺโม อุปฺปชฺชติ อารมฺมณ\-ปจฺจยา. วตฺถุํ ปจฺจยา อนุปาทินฺนา ขนฺธา. (2)}

\prose{อนุปาทินฺนํ ธมฺมํ ปจฺจยา อนุปาทินฺโน ธมฺโม อุปฺปชฺชติ อารมฺมณ\-ปจฺจยา. อนุปาทินฺนํ เอกํ ขนฺธํ ปจฺจยา ตโย ขนฺธา ฯเปฯ เทฺว ขนฺเธ ฯเปฯ. (1)}

\prose{อุปาทินฺนญฺจ อนุปาทินฺนญฺจ ธมฺมํ ปจฺจยา อนุปาทินฺโน ธมฺโม อุปฺปชฺชติ อารมฺมณ\-ปจฺจยา. อนุปาทินฺนํ เอกํ ขนฺธญฺจ วตฺถุญฺจ ปจฺจยา ตโย ขนฺธา ฯเปฯ เทฺว ขนฺเธ จ ฯเปฯ. (1)}

\proseitem{436}{อุปาทินฺนํ ธมฺมํ ปจฺจยา อนุปาทินฺโน ธมฺโม อุปฺปชฺชติ อธิปติ\-ปจฺจยา. วตฺถุํ ปจฺจยา อนุปาทินฺนา ขนฺธา. (1)}

\prose{อนุปาทินฺนํ ธมฺมํ ปจฺจยา อนุปาทินฺโน ธมฺโม อุปฺปชฺชติ อธิปติ\-ปจฺจยา. อนุปาทินฺนํ เอกํ ขนฺธํ ปจฺจยา ตโย ขนฺธา จิตฺต\-สมุฏฺฐานญฺจ รูปํ ฯเปฯ เทฺว ขนฺเธ ฯเปฯ. เอกํ มหาภูตํ ฯเปฯ มหาภูเต ปจฺจยา จิตฺต\-สมุฏฺฐานํ รูปํ, อุปาทารูปํ. (1)}

\prose{\csromanfolio{544}อุปาทินฺนญฺจ อนุปาทินฺนญฺจ ธมฺมํ ปจฺจยา อนุปาทินฺโน ธมฺโม อุปฺปชฺชติ อธิปติ\-ปจฺจยา. อนุปาทินฺนํ เอกํ ขนฺธญฺจ วตฺถุญฺจ ปจฺจยา ตโย ขนฺธา ฯเปฯ เทฺว ขนฺเธ จ ฯเปฯ. (1) (สํขิตฺตํ.)}

\csromanmarkh{1. ปจฺจยานุโลม}\csromanmarkhii{2. สงฺขฺยาวาร}\csromanheadernomark{2}{1. \textbf{ปจฺจยานุโลม 2. สงฺขฺยาวาร}}
\proseitem{437}{เหตุยา ปญฺจ, อารมฺมเณ จตฺตาริ, อธิปติยา ตีณิ, อนนฺตเร จตฺตาริ, สมนนฺตเร จตฺตาริ, สหชาเต ปญฺจ, อญฺญมญฺเญ จตฺตาริ, นิสฺสเย ปญฺจ, อุปนิสฺสเย จตฺตาริ, ปุเรชาเต จตฺตาริ, อาเสวเน ตีณิ, กมฺเม ปญฺจ, วิปาเก ปญฺจ ฯเปฯ อวิคเต ปญฺจ.}

\csromanmarkh{2. ปจฺจย\-ปจฺจนีย}\csromanmarkhii{1. วิภงฺควาร}\csromanheadernomark{2}{2. ปจฺจย\-ปจฺจนีย 1. วิภงฺควาร}
\proseitem{438}{อุปาทินฺนํ ธมฺมํ ปจฺจยา อุปาทินฺโน ธมฺโม อุปฺปชฺชติ นเหตุปจฺจยา. อเหตุกํ อุปาทินฺนํ เอกํ ขนฺธํ ปจฺจยา ตโย ขนฺธา ฯเปฯ เทฺว ขนฺเธ ฯเปฯ. อเหตุก\-ปฏิสนฺธิ\-กฺขเณ ฯเปฯ. ขนฺเธ ปจฺจยา วตฺถุ, วตฺถุํ ปจฺจยา ขนฺธา. เอกํ มหาภูตํ ฯเปฯ มหาภูเต ปจฺจยา กฏตฺตารูปํ, อุปาทารูปํ. อสญฺญสตฺตานํ เอกํ มหาภูตํ ฯเปฯ. จกฺขายตนํ ปจฺจยา จกฺขุวิญฺญาณํ ฯเปฯ. กายายตนํ ฯเปฯ. วตฺถุํ ปจฺจยา อเหตุกา อุปาทินฺนา ขนฺธา. (1)}

\prose{อุปาทินฺนํ ธมฺมํ ปจฺจยา อนุปาทินฺโน ธมฺโม อุปฺปชฺชติ นเหตุปจฺจยา. อเหตุเก อุปาทินฺเน ขนฺเธ ปจฺจยา จิตฺต\-สมุฏฺฐานํ รูปํ. วตฺถุํ ปจฺจยา อเหตุกา อนุปาทินฺนา ขนฺธา. วตฺถุํ ปจฺจยา วิจิกิจฺฉา\-สหคโต อุทฺธจฺจ\-สหคโต โมโห. (2)}

\prose{อุปาทินฺนํ ธมฺมํ ปจฺจยา อุปาทินฺโน จ อนุปาทินฺโน จ ธมฺมา อุปฺปชฺชนฺติ นเหตุปจฺจยา. อเหตุกํ อุปาทินฺนํ เอกํ ขนฺธํ ปจฺจยา ตโย ขนฺธา จิตฺต\-สมุฏฺฐานญฺจ รูปํ ฯเปฯ เทฺว ขนฺเธ ฯเปฯ. (3)}

\proseitem{439}{อนุปาทินฺนํ ธมฺมํ ปจฺจยา อนุปาทินฺโน ธมฺโม อุปฺปชฺชติ นเหตุปจฺจยา. อเหตุกํ อนุปาทินฺนํ เอกํ ขนฺธํ ปจฺจยา ตโย ขนฺธา จิตฺต\-สมุฏฺฐานญฺจ รูปํ ฯเปฯ \csromanfolio{545}เทฺว ขนฺเธ ฯเปฯ. เอกํ มหาภูตํ ฯเปฯ. พาหิรํ. อาหารสมุฏฺฐานํ. อุตุสมุฏฺฐานํ ฯเปฯ. วิจิกิจฺฉา\-สหคเต อุทฺธจฺจ\-สหคเต ขนฺเธ ปจฺจยา วิจิกิจฺฉา\-สหคโต อุทฺธจฺจ\-สหคโต โมโห. (1)}

\prose{อุปาทินฺนญฺจ อนุปาทินญฺจ ธมฺมํ ปจฺจยา อนุปาทินฺโน ธมฺโม อุปฺปชฺชติ นเหตุปจฺจยา. อเหตุเก อุปาทินฺเน ขนฺเธ จ มหาภูเต จ ปจฺจยา จิตฺต\-สมุฏฺฐานํ รูปํ. อเหตุกํ อนุปาทินฺนํ เอกํ ขนฺธญฺจ วตฺถุญฺจ ปจฺจยา ตโย ขนฺธา ฯเปฯ เทฺว ขนฺเธ จ ฯเปฯ. วิจิกิจฺฉา\-สหคเต อุทฺธจฺจ\-สหคเต ขนฺเธ จ วตฺถุญฺจ ปจฺจยา วิจิกิจฺฉา\-สหคโต อุทฺธจฺจ\-สหคโต โมโห. (1) (สํขิตฺตํ.)}

\csromanmarkh{2. ปจฺจย\-ปจฺจนีย}\csromanmarkhii{2. สงฺขฺยาวาร}\csromanheadernomark{2}{2. ปจฺจย\-ปจฺจนีย 2. สงฺขฺยาวาร}
\proseitem{440}{นเหตุยา ปญฺจ, นอารมฺมเณ จตฺตาริ, นอธิปติยา ปญฺจ, นอนนฺตเร จตฺตาริ, นสมนนฺตเร จตฺตาริ, นอญฺญมญฺเญ จตฺตาริ, นอุปนิสฺสเย จตฺตาริ, นปุเรชาเต จตฺตาริ, นปจฺฉาชาเต ปญฺจ, นอาเสวเน ปญฺจ, นกมฺเม ตีณิ, นวิปาเก จตฺตาริ, นอาหาเร เทฺว, นอินฺทฺริเย เทฺว, นฌาเน เทฺว, นมคฺเค ปญฺจ, นสมฺปยุตฺเต จตฺตาริ, นวิปฺปยุตฺเต เทฺว, โนนตฺถิยา จตฺตาริ, โนวิคเต จตฺตาริ.}

\proseitem{441}{เหตุปจฺจยา นอารมฺมเณ จตฺตาริ, นอธิปติยา ปญฺจ ฯเปฯ นปุเรชาเต จตฺตาริ, นปจฺฉาชาเต ปญฺจ, นอาเสวเน ปญฺจ, นกมฺเม ตีณิ, นวิปาเก ตีณิ ฯเปฯ นสมฺปยุตฺเต จตฺตาริ, นวิปฺปยุตฺเต เทฺว, โนนตฺถิยา จตฺตาริ, โนวิคเต จตฺตาริ.}

\csromanheader{2}{4. \textbf{ปจฺจยาปจฺจนียานุโลม}}
\proseitem{442}{นเหตุปจฺจยา อารมฺมเณ จตฺตาริ, อนนฺตเร จตฺตาริ ฯเปฯ มคฺเค ตีณิ ฯเปฯ อวิคเต ปญฺจ.}

\vspace{\tipitakaparskip}
\csromancenter{(นิสฺสยวาโร ปจฺจยวาร\-สทิโส.)}
\csromanmarkh{68. อุปาทินฺนทุก}\csromanmarkhii{5. สํสฏฺฐวาร}\csromanheadernomark{2}{68. \textbf{อุปาทินฺนทุก 5. สํสฏฺฐวาร}}
\proseitem{443}{\csromanfolio{546}อุปาทินฺนํ ธมฺมํ สํสฏฺโฐ อุปาทินฺโน ธมฺโม อุปฺปชฺชติ เหตุปจฺจยา. อุปาทินฺนํ เอกํ ขนฺธํ สํสฏฺฐา ตโย ขนฺธา ฯเปฯ เทฺว ขนฺเธ ฯเปฯ. ปฏิสนฺธิ\-กฺขเณ ฯเปฯ.}

\prose{อนุปาทินฺนํ ธมฺมํ สํสฏฺโฐ อนุปาทินฺโน ธมฺโม อุปฺปชฺชติ เหตุปจฺจยา. อนุปาทินฺนํ เอกํ ขนฺธํ สํสฏฺฐา ตโย ขนฺธา ฯเปฯ เทฺว ขนฺเธ ฯเปฯ. (1)}

\prose{เหตุยา เทฺว, อารมฺมเณ เทฺว, อธิปติยา เอกํ, อนนฺตเร เทฺว, สมนนฺตเร เทฺว, สหชาเต เทฺว, อญฺญมญฺเญ เทฺว, นิสฺสเย เทฺว, อุปนิสฺสเย เทฺว, ปุเรชาเต เทฺว, อาเสวเน เอกํ, กมฺเม เทฺว, วิปาเก เทฺว ฯเปฯ อวิคเต เทฺว. อนุโลมํ.}

\proseitem{444}{อุปาทินฺนํ ธมฺมํ สํสฏฺโฐ อุปาทินฺโน ธมฺโม อุปฺปชฺชติ นเหตุปจฺจยา. อเหตุกํ อุปาทินฺนํ เอกํ ขนฺธํ สํสฏฺฐา ตโย ขนฺธา ฯเปฯ เทฺว ขนฺเธ ฯเปฯ. อเหตุก\-ปฏิสนฺธิ\-กฺขเณ ฯเปฯ. (1)}

\prose{อนุปาทินฺนํ ธมฺมํ สํสฏฺโฐ อนุปาทินฺโน ธมฺโม อุปฺปชฺชติ นเหตุปจฺจยา. อเหตุกํ อนุปาทินฺนํ เอกํ ขนฺธํ สํสฏฺฐา ตโย ขนฺธา ฯเปฯ เทฺว ขนฺเธ ฯเปฯ. วิจิกิจฺฉา\-สหคเต อุทฺธจฺจ\-สหคเต ขนฺเธ สํสฏฺโฐ วิจิกิจฺฉา\-สหคโต อุทฺธจฺจ\-สหคโต โมโห. (1)}

\prose{นเหตุยา เทฺว, นอธิปติยา เทฺว, นปุเรชาเต เทฺว, นปจฺฉาชาเต เทฺว, นอาเสวเน เทฺว, นกมฺเม เอกํ, นวิปาเก เอกํ, นฌาเน เอกํ, นมคฺเค เทฺว, นวิปฺปยุตฺเต เทฺว.}

\vspace{\tipitakaparskip}
\csromancenter{ปจฺจนียํ.}
\csromancenter{(เอวํ อิตเร เทฺว คณนาปิ สมฺปยุตฺตวาโรปิ กาตพฺโพ.)}
\csromanmarkh{68. อุปาทินฺนทุก}\csromanmarkhii{68. อุปาทินฺนทุก 7. ปญฺหาวาร}\csromanheadernomark{2}{68. อุปาทินฺนทุก \textbf{68. อุปาทินฺนทุก 7. ปญฺหาวาร}}
\csromanmarkh{1. ปจฺจยานุโลม}\csromanmarkhii{1. วิภงฺควาร}\csromanheadernomark{2}{1. \textbf{ปจฺจยานุโลม 1. วิภงฺควาร}}
\csromanheader{2}{\textbf{เหตุ}}
\proseitem{445}{\csromanfolio{547}อุปาทินฺโน ธมฺโม อุปาทินฺนสฺส ธมฺมสฺส เหตุปจฺจเยน ปจฺจโย. อุปาทินฺนา เหตู สมฺปยุตฺตกานํ ขนฺธานํ เหตุปจฺจเยน ปจฺจโย. ปฏิสนฺธิ\-กฺขเณ อุปาทินฺนา เหตู สมฺปยุตฺตกานํ ขนฺธานํ กฏตฺตา จ รูปานํ เหตุปจฺจเยน ปจฺจโย. (1)}

\prose{อุปาทินฺโน ธมฺโม อนุปาทินฺนสฺส ธมฺมสฺส เหตุปจฺจเยน ปจฺจโย. อุปาทินฺนา เหตู จิตฺต\-สมุฏฺฐานานํ รูปานํ เหตุปจฺจเยน ปจฺจโย. (2)}

\prose{อุปาทินฺโน ธมฺโม อุปาทินฺนสฺส จ อนุปาทินฺนสฺส จ ธมฺมสฺส เหตุปจฺจเยน ปจฺจโย. อุปาทินฺนา เหตู สมฺปยุตฺตกานํ ขนฺธานํ จิตฺตสมุฏฺฐานานญฺจ รูปานํ เหตุปจฺจเยน ปจฺจโย. (3)}

\prose{อนุปาทินฺโน ธมฺโม อนุปาทินฺนสฺส ธมฺมสฺส เหตุปจฺจเยน ปจฺจโย. อนุปาทินฺนา เหตู สมฺปยุตฺตกานํ ขนฺธานํ จิตฺตสมุฏฺฐานานญฺจ รูปานํ เหตุปจฺจเยน ปจฺจโย. (1)}

\proseitem{446}{อุปาทินฺโน ธมฺโม อุปาทินฺนสฺส ธมฺมสฺส อารมฺมณ\-ปจฺจเยน ปจฺจโย. จกฺขุํ ฯเปฯ. กายํ อุปาทินฺเน รูเป. คนฺเธ. รเส. โผฏฺฐพฺเพ. วตฺถุํ อุปาทินฺเน ขนฺเธ อนิจฺจโต ฯเปฯ โทมนสฺสํ อุปฺปชฺชติ. กุสลากุสเล นิรุทฺเธ วิปาโก ตทารมฺมณตา อุปฺปชฺชติ. อุปาทินฺนํ รูปายตนํ จกฺขุ\-วิญฺญาณสฺส ฯเปฯ. อุปาทินฺนํ คนฺธายตนํ ฆาน\-วิญฺญาณสฺส ฯเปฯ. ผฏฺฐพฺพายตนํ กายวิญฺญาณสฺส อารมฺมณ\-ปจฺจเยน ปจฺจโย. (1)}

\prose{อุปาทินฺโน ธมฺโม อนุปาทินฺนสฺส ธมฺมสฺส อารมฺมณ\-ปจฺจเยน ปจฺจโย. จกฺขุํ ฯเปฯ. กายํ อุปาทินฺเน รูเป. คนฺเธ. รเส. โผฏฺฐพฺเพ. วตฺถุํ อุปาทินฺเน ขนฺเธ อนิจฺจโต ฯเปฯ โทมนสฺสํ อุปฺปชฺชติ. ทิพฺเพน จกฺขุนา อุปาทินฺนํ รูปํ ปสฺสติ. เจโตปริย\-ญาเณน อุปาทินฺน\-จิตฺต\-สมงฺคิสฺส จิตฺตํ ชานาติ. \csromanfolio{548}อุปาทินฺนา ขนฺธา อิทฺธิวิธ\-ญาณสฺส, เจโตปริย\-ญาณสฺส, ปุพฺเพ\-นิวาสา\-นุสฺสติ\-ญาณสฺส, อนาคตํสญาณสฺส, อาวชฺชนาย อารมฺมณ\-ปจฺจเยน ปจฺจโย. (2)}

\proseitem{447}{อนุปาทินฺโน ธมฺโม อนุปาทินฺนสฺส ธมฺมสฺส อารมฺมณ\-ปจฺจเยน ปจฺจโย. ทานํ ฯเปฯ สีลํ ฯเปฯ อุโปสถกมฺมํ กตฺวา ตํ ปจฺจเวกฺขติ อสฺสาเทติ อภินนฺทติ, ตํ อารพฺภ ราโค ฯเปฯ โทมนสฺสํ อุปฺปชฺชติ. ปุพฺเพ สุจิณฺณานิ ฯเปฯ. ฌานา ฯเปฯ. อริยา มคฺคา วุฏฺฐหิตฺวา มคฺคํ ปจฺจเวกฺขนฺติ. ผลํ ฯเปฯ. นิพฺพานํ ฯเปฯ. นิพฺพานํ โคตฺรภุสฺส, โวทานสฺส, มคฺคสฺส, ผลสฺส, อาวชฺชนาย อารมฺมณ\-ปจฺจเยน ปจฺจโย. อริยา ปหีเน กิเลเส ฯเปฯ. วิกฺขมฺภิเต กิเลเส ฯเปฯ. ปุพฺเพ ฯเปฯ. อนุปาทินฺเน รูเป. สทฺเท. คนฺเธ. รเส. โผฏฺฐพฺเพ อนุปาทินฺเน ขนฺเธ อนิจฺจโต ฯเปฯ โทมนสฺสํ อุปฺปชฺชติ. ทิพฺเพน จกฺขุนา อนุปาทินฺนํ รูปํ ปสฺสติ. ทิพฺพาย โสตธาตุยา สทฺทํ สุณาติ. เจโตปริย\-ญาเณน อนุปาทินฺน\-จิตฺต\-สมงฺคิสฺส จิตฺตํ ชานาติ. อากาสา\-นญฺจา\-ยตนํ ฯเปฯ. อากิญฺจญฺญา\-ยตนํ ฯเปฯ. อนุปาทินฺนา ขนฺธา อิทฺธิวิธ\-ญาณสฺส, เจโตปริย\-ญาณสฺส, ปุพฺเพ\-นิวาสา\-นุสฺสติ\-ญาณสฺส, ยถากมฺมูปคญาณสฺส, อนาคตํสญาณสฺส, อาวชฺชนาย อารมฺมณ\-ปจฺจเยน ปจฺจโย. (1)}

\prose{อนุปาทินฺโน ธมฺโม อุปาทินฺนสฺส ธมฺมสฺส อารมฺมณ\-ปจฺจเยน ปจฺจโย. อนุปาทินฺเน รูเป. สทฺเท. คนฺเธ. รเส. โผฏฺฐพฺเพ อนุปาทินฺเน ขนฺเธ อนิจฺจโต ฯเปฯ โทมนสฺสํ อุปฺปชฺชติ. กุสลากุสเล นิรุทฺเธ วิปาโก ตทารมฺมณตา อุปฺปชฺชติ. อากาสา\-นญฺจา\-ยตนกุสลํ วิญฺญาณญฺจา\-ยตนวิปากสฺส ฯเปฯ. อากิญฺจญฺญายตน\-กุสลํ ฯเปฯ. อนุปาทินฺนํ รูปายตนํ จกฺขุ\-วิญฺญาณสฺส ฯเปฯ. สทฺทายตนํ ฯเปฯ. โผฏฺฐพฺพา\-ยตนํ กายวิญฺญาณสฺส อารมฺมณ\-ปจฺจเยน ปจฺจโย. (2)}

\proseitem{448}{อุปาทินฺโน ธมฺโม อนุปาทินฺนสฺส ธมฺมสฺส อธิปติ\-ปจฺจเยน ปจฺจโย. อารมฺมณา\-ธิปติ\csromandash{}จกฺขุํ ฯเปฯ กายํ. อุปาทินฺเน รูเป. คนฺเธ. รเส. โผฏฺฐพฺเพ. วตฺถุํ อุปาทินฺเน ขนฺเธ ครุํ กตฺวา อสฺสาเทติ อภินนฺทติ, ตํ ครุํ กตฺวา ราโค อุปฺปชฺชติ. ทิฏฺฐิ อุปฺปชฺชติ. (1)}

\prose{\csromanfolio{549}อนุปาทินฺโน ธมฺโม อนุปาทินฺนสฺส ธมฺมสฺส อธิปติ\-ปจฺจเยน ปจฺจโย. อารมฺมณา\-ธิปติ, สหชาตา\-ธิปติ.  อารมฺมณา\-ธิปติ\csromandash{}ทานํ ฯเปฯ สีลํ ฯเปฯ อุโปสถกมฺมํ กตฺวา ตํ ครุํ กตฺวา ปจฺจเวกฺขติ, อสฺสาเทติ อภินนฺทติ, ตํ ครุํ กตฺวา ราโค อุปฺปชฺชติ. ทิฏฺฐิ อุปฺปชฺชติ. ปุพฺเพ ฯเปฯ. ฌานา ฯเปฯ. อริยา มคฺคา วุฏฺฐหิตฺวา มคฺคํ ครุํ กตฺวา ปจฺจเวกฺขนฺติ, ผลํ ฯเปฯ. นิพฺพานํ ครุํ กตฺวา ปจฺจเวกฺขนฺติ. นิพฺพานํ โคตฺรภุสฺส, โวทานสฺส, มคฺคสฺส, ผลสฺส อธิปติ\-ปจฺจเยน ปจฺจโย. อนุปาทินฺเน รูเป. สทฺเท. คนฺเธ. รเส. โผฏฺฐพฺเพ อนุปาทินฺเน ขนฺเธ ครุํ กตฺวา อสฺสาเทติ อภินนฺทติ, ตํ ครุํ กตฺวา ราโค อุปฺปชฺชติ. ทิฏฺฐิ อุปฺปชฺชติ. สหชาตา\-ธิปติ\csromandash{}อนุปาทินฺนา\-ธิปติ สมฺปยุตฺตกานํ ขนฺธานํ จิตฺตสมุฏฺฐานานญฺจ รูปานํ อธิปติ\-ปจฺจเยน ปจฺจโย. (1)}

\proseitem{449}{อุปาทินฺโน ธมฺโม อุปาทินฺนสฺส ธมฺมสฺส อนนฺตร\-ปจฺจเยน ปจฺจโย. ปุริมา ปุริมา อุปาทินฺนา ขนฺธา ปจฺฉิมานํ ปจฺฉิมานํ อุปาทินฺนานํ ขนฺธานํ อนนฺตร\-ปจฺจเยน ปจฺจโย. ปญฺจวิญฺญาณํ วิปาก\-มโนธาตุยา อนนฺตร\-ปจฺจเยน ปจฺจโย. วิปาก\-มโนธาตุ วิปาก\-มโนวิญฺญาณ\-ธาตุยา อนนฺตร\-ปจฺจเยน ปจฺจโย. (1)}

\prose{อุปาทินฺโน ธมฺโม อนุปาทินฺนสฺส ธมฺมสฺส อนนฺตร\-ปจฺจเยน ปจฺจโย. ภวงฺคํ อาวชฺชนาย. วิปาก\-มโนวิญฺญาณ\-ธาตุ กิริย\-มโนวิญฺญาณ\-ธาตุยา อนนฺตร\-ปจฺจเยน ปจฺจโย. (2)}

\proseitem{450}{อนุปาทินฺโน ธมฺโม อนุปาทินฺนสฺส ธมฺมสฺส อนนฺตร\-ปจฺจเยน ปจฺจโย. ปุริมา ปุริมา อนุปาทินฺนา ขนฺธา ปจฺฉิมานํ ปจฺฉิมานํ อนุปาทินฺนานํ ขนฺธานํ อนนฺตร\-ปจฺจเยน ปจฺจโย. อนุโลมํ โคตฺรภุสฺส ฯเปฯ. ผลสมาปตฺติยา อนนฺตร\-ปจฺจเยน ปจฺจโย. (1)}

\prose{อนุปาทินฺโน ธมฺโม อุปาทินฺนสฺส ธมฺมสฺส อนนฺตร\-ปจฺจเยน ปจฺจโย. อาวชฺชนา ปญฺจนฺนํ วิญฺญาณาณํ ฯเปฯ. อนุปาทินฺนา ขนฺธา วุฏฺฐานสฺส อนนฺตร\-ปจฺจเยน ปจฺจโย. สมนนฺตร\-ปจฺจเยน ปจฺจโย. สหชาต\-ปจฺจเยน ปจฺจโย. ปญฺจ. (ปฏิจฺจสทิสา.) อญฺญมญฺญ\-ปจฺจเยน ปจฺจโย. เทฺว. (ปฏิจฺจสทิสา.) นิสฺสย\-ปจฺจเยน ปจฺจโย. (ปจฺจยวาเร นิสฺสยสทิสา.) ปญฺจ.}

\csromanheader{2}{\textbf{อุปนิสฺสย}}
\proseitem{451}{\csromanfolio{550}อุปาทินฺโน ธมฺโม อุปาทินฺนสฺส ธมฺมสฺส อุปนิสฺสย\-ปจฺจเยน ปจฺจโย. อนนฺตรู\-ปนิสฺสโย, ปกตู\-ปนิสฺสโย ฯเปฯ. ปกตู\-ปนิสฺสโย\csromandash{} กายิกํ สุขํ กายิกสฺส สุขสฺส. กายิกสฺส ทุกฺขสฺส อุปนิสฺสย\-ปจฺจเยน ปจฺจโย. กายิกํ ทุกฺขํ. อุปาทินฺนํ อุตุ. โภชนํ กายิกสฺส สุขสฺส. กายิกสฺส ทุกฺขสฺส อุปนิสฺสย\-ปจฺจเยน ปจฺจโย. กายิกํ สุขํ. กายิกํ ทุกฺขํ. อุตุ. โภชนํ กายิกสฺส สุขสฺส. กายิกสฺส ทุกฺขสฺส อุปนิสฺสย\-ปจฺจเยน ปจฺจโย. (1)}

\prose{อุปาทินฺโน ธมฺโม อนุปาทินฺนสฺส ธมฺมสฺส อุปนิสฺสย\-ปจฺจเยน ปจฺจโย. อารมฺมณู\-ปนิสฺสโย, อนนฺตรู\-ปนิสฺสโย, ปกตู\-ปนิสฺสโย ฯเปฯ. ปกตู\-ปนิสฺสโย\csromandash{}กายิกํ สุขํ อุปนิสฺสาย ทานํ เทติ ฯเปฯ สํฆํ ภินฺทติ. กายิกํ ทุกฺขํ. อุปาทินฺนํ อุตุํ. โภชนํ อุปนิสฺสาย ทานํ เทติ ฯเปฯ สมาปตฺติํ อุปฺปาเทติ. ปาณํ หนติ ฯเปฯ สํฆํ ภินฺทติ. กายิกํ สุขํ. กายิกํ ทุกฺขํ. อุตุ. โภชนํ สทฺธาย ฯเปฯ ปตฺถนาย มคฺคสฺส, ผลสมาปตฺติยา อุปนิสฺสย\-ปจฺจเยน ปจฺจโย. (2)}

\proseitem{452}{อนุปาทินฺโน ธมฺโม อนุปาทินฺนสฺส ธมฺมสฺส อุปนิสฺสย\-ปจฺจเยน ปจฺจโย. อารมฺมณู\-ปนิสฺสโย, อนนฺตรู\-ปนิสฺสโย, ปกตู\-ปนิสฺสโย ฯเปฯ. ปกตู\-ปนิสฺสโย\csromandash{}สทฺธํ อุปนิสฺสาย ทานํ เทติ. มานํ ชปฺเปติ, ทิฏฺฐิํ คณฺหาติ. สีลํ ฯเปฯ. ปตฺถนํ. อุตุํ. โภชนํ. เสนาสนํ อุปนิสฺสาย ทานํ เทติ ฯเปฯ สํฆํ ภินฺทติ. สทฺธา ฯเปฯ. ปตฺถนา. อุตุ. โภชนํ. เสนาสนํ สทฺธาย ฯเปฯ ปตฺถนาย มคฺคสฺส, ผลสมาปตฺติยา อุปนิสฺสย\-ปจฺจเยน ปจฺจโย. (1)}

\prose{อนุปาทินฺโน ธมฺโม อุปาทินฺนสฺส ธมฺมสฺส อุปนิสฺสย\-ปจฺจเยน ปจฺจโย. อนนฺตรู\-ปนิสฺสโย, ปกตู\-ปนิสฺสโย ฯเปฯ. ปกตู\-ปนิสฺสโย\csromandash{} สทฺธํ อุปนิสฺสาย อตฺตานํ อาตาเปติ ปริตาเปติ, ปริยิฏฺฐิมูลกํ ทุกฺขํ ปจฺจนุโภติ. สีลํ ฯเปฯ. ปตฺถนํ. อุตุํ. โภชนํ. เสนาสนํ อุปนิสฺสาย อตฺตานํ อาตาเปติ ปริตาเปติ, ปริยิฏฺฐิมูลกํ ทุกฺขํ ปจฺจนุโภติ. สทฺธา ฯเปฯ. เสนาสนํ กายิกสฺส สุขสฺส. กายิกสฺส ทุกฺขสฺส อุปนิสฺสย\-ปจฺจเยน ปจฺจโย. กุสลากุสลํ ธมฺมํ วิปากสฺส อุปนิสฺสย\-ปจฺจเยน ปจฺจโย. (2)}

\csromanheader{2}{68. อุปาทินฺนทุก \textbf{ปุเรชาต}}
\proseitem{453}{\csromanfolio{551}อุปาทินฺโน ธมฺโม อุปาทินฺนสฺส ธมฺมสฺส ปุเรชาต\-ปจฺจเยน ปจฺจโย. อารมฺมณ\-ปุเรชาตํ, วตฺถุ\-ปุเรชาตํ.  อารมฺมณ\-ปุเรชาตํ\csromandash{} จกฺขุํ ฯเปฯ กายํ. อุปาทินฺเน รูเป. คนฺเธ. รเส. โผฏฺฐพฺเพ. วตฺถุํ อนิจฺจโต ฯเปฯ โทมนสฺสํ อุปฺปชฺชติ. กุสลากุสเล นิรุทฺเธ วิปาโก ตทารมฺมณตา อุปฺปชฺชติ. อุปาทินฺนํ รูปายตนํ จกฺขุ\-วิญฺญาณสฺส ฯเปฯ คนฺธายตนํ ฯเปฯ โผฏฺฐพฺพา\-ยตนํ กายวิญฺญาณสฺส ฯเปฯ. วตฺถุ\-ปุเรชาตํ\csromandash{}จกฺขายตนํ จกฺขุ\-วิญฺญาณสฺส ฯเปฯ กายายตนํ กายวิญฺญาณสฺส ฯเปฯ. วตฺถุ อุปาทินฺนานํ ขนฺธานํ ปุเรชาต\-ปจฺจเยน ปจฺจโย. (1)}

\prose{อุปาทินฺโน ธมฺโม อนุปาทินฺนสฺส ธมฺมสฺส ปุเรชาต\-ปจฺจเยน ปจฺจโย. อารมฺมณ\-ปุเรชาตํ, วตฺถุ\-ปุเรชาตํ.  อารมฺมณ\-ปุเรชาตํ\csromandash{} จกฺขุํ ฯเปฯ กายํ. อุปาทินฺเน รูเป. คนฺเธ. รเส. โผฏฺฐพฺเพ วตฺถุํ อนิจฺจโต ฯเปฯ โทมนสฺสํ อุปฺปชฺชติ. ทิพฺเพน จกฺขุนา อุปาทินฺนํ รูปํ ปสฺสติ. วตฺถุ\-ปุเรชาตํ\csromandash{}วตฺถุ อนุปาทินฺนานํ ขนฺธานํ ปุเรชาต\-ปจฺจเยน ปจฺจโย. (2)}

\proseitem{454}{อนุปาทินฺโน ธมฺโม อนุปาทินฺนสฺส ธมฺมสฺส ปุเรชาต\-ปจฺจเยน ปจฺจโย. อารมฺมณ\-ปุเรชาตํ\csromandash{}อนุปาทินฺเน รูเป. สทฺเท. คนฺเธ. รเส. โผฏฺฐพฺเพ อนิจฺจโต ฯเปฯ โทมนสฺสํ อุปฺปชฺชติ. ทิพฺเพน จกฺขุนา อนุปาทินฺนํ รูปํ ปสฺสติ. ทิพฺพาย โสตธาตุยา สทฺทํ สุณาติ. (1)}

\prose{อนุปาทินฺโน ธมฺโม อุปาทินฺนสฺส ธมฺมสฺส ปุเรชาต\-ปจฺจเยน ปจฺจโย. อารมฺมณ\-ปุเรชาตํ\csromandash{}อนุปาทินฺเน รูเป. สทฺเท. คนฺเธ. รเส. โผฏฺฐพฺเพ อนิจฺจโต ฯเปฯ โทมนสฺสํ อุปฺปชฺชติ. กุสลากุสเล นิรุทฺเธ วิปาโก ตทารมฺมณตา อุปฺปชฺชติ. อนุปาทินฺนํ รูปายตนํ จกฺขุ\-วิญฺญาณสฺส ฯเปฯ สทฺทายตนํ ฯเปฯ โผฏฺฐพฺพา\-ยตนํ กายวิญฺญาณสฺส ปุเรชาต\-ปจฺจเยน ปจฺจโย. (2)}

\prose{อุปาทินฺโน จ อนุปาทินฺโน จ ธมฺมา อุปาทินฺนสฺส ธมฺมสฺส ปุเรชาต\-ปจฺจเยน ปจฺจโย. อารมฺมณ\-ปุเรชาตํ, วตฺถุ\-ปุเรชาตํ.  อนุปาทินฺนํ รูปายตนญฺจ จกฺขายตนญฺจ จกฺขุ\-วิญฺญาณสฺส ฯเปฯ สทฺทายตนญฺจ โสตายตนญฺจ ฯเปฯ โผฏฺฐพฺพายตนญฺจ กายายตนญฺจ กายวิญฺญาณสฺส ปุเรชาต\-ปจฺจเยน ปจฺจโย. อนุปาทินฺนํ รูปายตนญฺจ วตฺถุ จ ฯเปฯ โผฏฺฐพฺพายตนญฺจ วตฺถุ จ อุปาทินฺนานํ ขนฺธานํ ปุเรชาต\-ปจฺจเยน ปจฺจโย. (1)}

\prose{\csromanfolio{552}อุปาทินฺโน จ อนุปาทินฺโน จ ธมฺมา อนุปาทินฺนสฺส ธมฺมสฺส ปุเรชาต\-ปจฺจเยน ปจฺจโย. อารมฺมณ\-ปุเรชาตํ, วตฺถุ\-ปุเรชาตํ. อนุปาทินฺนํ รูปายตนญฺจ วตฺถุ จ ฯเปฯ โผฏฺฐพฺพายตนญฺจ วตฺถุ จ อนุปาทินฺนานํ ขนฺธานํ ปุเรชาต\-ปจฺจเยน ปจฺจโย. (2)}

\proseitem{455}{อุปาทินฺโน ธมฺโม อุปาทินฺนสฺส ธมฺมสฺส ปจฺฉาชาต\-ปจฺจเยน ปจฺจโย. ปจฺฉาชาตา อุปาทินฺนา ขนฺธา ปุเรชาตสฺส อิมสฺส อุปาทินฺนสฺส กายสฺส ปจฺฉาชาต\-ปจฺจเยน ปจฺจโย. (1)}

\prose{อุปาทินฺโน ธมฺโม อนุปาทินฺนสฺส ธมฺมสฺส ปจฺฉาชาต\-ปจฺจเยน ปจฺจโย. ปจฺฉาชาตา อุปาทินฺนา ขนฺธา ปุเรชาตสฺส อิมสฺส อนุปาทินฺนสฺส กายสฺส ปจฺฉาชาต\-ปจฺจเยน ปจฺจโย. (2)}

\prose{อุปาทินฺโน ธมฺโม อุปาทินฺนสฺส จ อนุปาทินฺนสฺส จ ธมฺมสฺส ปจฺฉาชาต\-ปจฺจเยน ปจฺจโย. (สํขิตฺตํ.) (3)}

\prose{อนุปาทินฺโน ธมฺโม อนุปาทินฺนสฺส ธมฺมสฺส ปจฺฉาชาต\-ปจฺจเยน ปจฺจโย. (สํขิตฺตํ.) (1)}

\prose{อนุปาทินฺโน ธมฺโม อุปาทินฺนสฺส ธมฺมสฺส ปจฺฉาชาต\-ปจฺจเยน ปจฺจโย. (สํขิตฺตํ.) (2)}

\prose{อนุปาทินฺโน ธมฺโม อุปาทินฺนสฺส จ อนุปาทินฺนสฺส จ ธมฺมสฺส ปจฺฉาชาต\-ปจฺจเยน ปจฺจโย. (สํขิตฺตํ.) (3)}

\proseitem{456}{อนุปาทินฺโน ธมฺโม อนุปาทินฺนสฺส ธมฺมสฺส อาเสวน\-ปจฺจเยน ปจฺจโย. เอกํ. (สํขิตฺตํ.)}

\proseitem{457}{อุปาทินฺโน ธมฺโม อุปาทินฺนสฺส ธมฺมสฺส กมฺมปจฺจเยน ปจฺจโย. อุปาทินฺนา เจตนา สมฺปยุตฺตกานํ ขนฺธานํ กมฺมปจฺจเยน ปจฺจโย. ปฏิสนฺธิ\-กฺขเณ อุปาทินฺนา เจตนา สมฺปยุตฺตกานํ ขนฺธานํ กฏตฺตา จ รูปานํ กมฺมปจฺจเยน ปจฺจโย. (1)}

\prose{\csromanfolio{553}อุปาทินฺโน ธมฺโม อนุปาทินฺนสฺส ธมฺมสฺส กมฺมปจฺจเยน ปจฺจโย. อุปาทินฺนา เจตนา จิตฺต\-สมุฏฺฐานานํ รูปานํ กมฺมปจฺจเยน ปจฺจโย. (2)}

\prose{อุปาทินฺโน ธมฺโม อุปาทินฺนสฺส จ อนุปาทินฺนสฺส จ ธมฺมสฺส กมฺมปจฺจเยน ปจฺจโย. อุปาทินฺนา เจตนา สมฺปยุตฺตกานํ ขนฺธานํ จิตฺตสมุฏฺฐานานญฺจ รูปานํ กมฺมปจฺจเยน ปจฺจโย. (3)}

\prose{อนุปาทินฺโน ธมฺโม อนุปาทินฺนสฺส ธมฺมสฺส กมฺมปจฺจเยน ปจฺจโย. อนุปาทินฺนา เจตนา สมฺปยุตฺตกานํ ขนฺธานํ จิตฺตสมุฏฺฐานานญฺจ รูปานํ กมฺมปจฺจเยน ปจฺจโย. (1)}

\prose{อนุปาทินฺโน ธมฺโม อุปาทินฺนสฺส ธมฺมสฺส กมฺมปจฺจเยน ปจฺจโย. \textbf{นานากฺขณิกา} อนุปาทินฺนา เจตนา วิปากานํ ขนฺธานํ กฏตฺตา จ รูปานํ กมฺมปจฺจเยน ปจฺจโย. (2)}

\csromanheader{2}{\textbf{วิปาก}}
\proseitem{458}{อุปาทินฺโน ธมฺโม อุปาทินฺนสฺส ธมฺมสฺส วิปาก\-ปจฺจเยน ปจฺจโย. อุปาทินฺโน เอโก ขนฺโธ. (สํขิตฺตํ. อุปาทินฺนมูลเก ตีณิ.)}

\prose{อนุปาทินฺโน ธมฺโม อนุปาทินฺนสฺส ธมฺมสฺส วิปาก\-ปจฺจเยน ปจฺจโย. วิปาโก อนุปาทินฺโน เอโก ขนฺโธ ติณฺณนฺนํ ขนฺธานํ จิตฺตสมุฏฺฐานาญฺจ รูปานํ วิปาก\-ปจฺจเยน ปจฺจโย ฯเปฯ เทฺว ขนฺธา ฯเปฯ. (1)}

\csromanheader{2}{\textbf{อาหาร}}
\proseitem{459}{อุปาทินฺโน ธมฺโม อุปาทินฺนสฺส ธมฺมสฺส อาหารปจฺจเยน ปจฺจโย. อุปาทินฺนา อาหารา สมฺปยุตฺตกานํ ขนฺธานํ อาหารปจฺจเยน ปจฺจโย. ปฏิสนฺธิ\-กฺขเณ ฯเปฯ. อุปาทินฺโน กพฬีกาโร อาหาโร อิมสฺส อุปาทินฺนสฺส กายสฺส อาหารปจฺจเยน ปจฺจโย. (1)}

\prose{อุปาทินฺโน ธมฺโม อนุปาทินฺนสฺส ธมฺมสฺส อาหารปจฺจเยน ปจฺจโย. อุปาทินฺนา อาหารา จิตฺต\-สมุฏฺฐานานํ รูปานํ อาหารปจฺจเยน ปจฺจโย. อุปาทินฺโน กพฬีกาโร อาหาโร อิมสฺส อนุปาทินฺนสฺส กายสฺส อาหารปจฺจเยน ปจฺจโย. (2)}

\prose{\csromanfolio{554}อุปาทินฺโน ธมฺโม อุปาทินฺนสฺส จ อนุปาทินฺนสฺส จ ธมฺมสฺส อาหารปจฺจเยน ปจฺจโย. อุปาทินฺนา อาหารา สมฺปยุตฺตกานํ ขนฺธานํ จิตฺตสมุฏฺฐานานญฺจ รูปานํ อาหารปจฺจเยน ปจฺจโย. อุปาทินฺโน กพฬีกาโร อาหาโร อิมสฺส อุปาทินฺนสฺส จ อนุปาทินฺนสฺส จ กายสฺส อาหารปจฺจเยน ปจฺจโย. (3)}

\proseitem{460}{อนุปาทินฺโน ธมฺโม อนุปาทินฺนสฺส ธมฺมสฺส อาหารปจฺจเยน ปจฺจโย. อนุปาทินฺนา อาหารา สมฺปยุตฺตกานํ ขนฺธานํ จิตฺตสมุฏฺฐานานญฺจ รูปานํ อาหารปจฺจเยน ปจฺจโย. อนุปาทินฺโน กพฬีกาโร อาหาโร อิมสฺส อนุปาทินฺนสฺส กายสฺส อาหารปจฺจเยน ปจฺจโย. (1)}

\prose{อนุปาทินฺโน ธมฺโม อุปาทินฺนสฺส ธมฺมสฺส อาหารปจฺจเยน ปจฺจโย. อนุปาทินฺโน กพฬีกาโร อาหาโร อิมสฺส อุปาทินฺนสฺส กายสฺส อาหารปจฺจเยน ปจฺจโย. (2)}

\prose{อนุปาทินฺโน ธมฺโม อุปาทินฺนสฺส จ อนุปาทินฺนสฺส จ ธมฺมสฺส อาหารปจฺจเยน ปจฺจโย. อนุปาทินฺโน กพฬีกาโร อาหาโร อิมสฺส อุปาทินฺนสฺส จ อนุปาทินฺนสฺส จ กายสฺส อาหารปจฺจเยน ปจฺจโย. (3)}

\prose{อุปาทินฺโน จ อนุปาทินฺโน จ ธมฺมา อุปาทินฺนสฺส ธมฺมสฺส อาหารปจฺจเยน ปจฺจโย. อุปาทินฺโน จ อนุปาทินฺโน จ กพฬีกาโร อาหาโร อิมสฺส อุปาทินฺนสฺส กายสฺส อาหารปจฺจเยน ปจฺจโย. (1)}

\prose{อุปาทินฺโน จ อนุปาทินฺโน จ ธมฺมา อนุปาทินฺนสฺส ธมฺมสฺส อาหารปจฺจเยน ปจฺจโย. อุปาทินฺโน จ อนุปาทินฺโน จ กพฬีกาโร อาหาโร อิมสฺส อนุปาทินฺนสฺส กายสฺส อาหารปจฺจเยน ปจฺจโย. (2)}

\prose{อุปาทินฺโน จ อนุปาทินฺโน จ ธมฺมา อุปาทินฺนสฺส จ อนุปาทินฺนสฺส จ ธมฺมสฺส อาหารปจฺจเยน ปจฺจโย. อุปาทินฺโน จ อนุปาทินฺโน จ กพฬีกาโร อาหาโร อิมสฺส อุปาทินฺนสฺส จ อนุปาทินฺนสฺส จ กายสฺส อาหารปจฺจเยน ปจฺจโย. (3)}

\csromanheader{2}{\textbf{อินฺทฺริยาทิ}}
\proseitem{461}{อุปาทินฺโน ธมฺโม อุปาทินฺนสฺส ธมฺมสฺส อินฺทฺริย\-ปจฺจเยน ปจฺจโย. อุปาทินฺนา อินฺทฺริยา สมฺปยุตฺตกานํ ขนฺธานํ อินฺทฺริย\-ปจฺจเยน ปจฺจโย. (1) \csromanfolio{555}ปฏิสนฺธิ\-กฺขเณ อุปาทินฺนา อินฺทฺริยา สมฺปยุตฺตกานํ ขนฺธานํ กฏตฺตา จ รูปานํ อินฺทฺริย\-ปจฺจเยน ปจฺจโย. จกฺขุนฺทฺริยํ จกฺขุ\-วิญฺญาณสฺส ฯเปฯ. กายินฺทฺริยํ กายวิญฺญาณสฺส ฯเปฯ. รูปชีวิตินฺทฺริยํ กฏตฺตารูปานํ อินฺทฺริย\-ปจฺจเยน ปจฺจโย. (2)}

\prose{อุปาทินฺโน ธมฺโม อนุปาทินฺนสฺส ธมฺมสฺส ฯเปฯ. (อุปาทินฺนมูลเก ตีณิ, ปฐมสฺเสว รูปชีวิตินฺทฺริยํ, อิตเรสุ นตฺถิ.) (3)}

\prose{อนุปาทินฺโน ธมฺโม อนุปาทินฺนสฺส ธมฺมสฺส อินฺทฺริย\-ปจฺจเยน ปจฺจโย. อนุปาทินฺนา อินฺทฺริยา สมฺปยุตฺตกานํ ขนฺธานํ จิตฺตสมุฏฺฐานานญฺจ รูปานํ อินฺทฺริย\-ปจฺจเยน ปจฺจโย. (1)}

\prose{อุปาทินฺโน ธมฺโม อุปาทินฺนสฺส ธมฺมสฺส ฌานปจฺจเยน ปจฺจโย. จตฺตาริ. มคฺคปจฺจเยน ปจฺจโย. จตฺตาริ. สมฺปยุตฺต\-ปจฺจเยน ปจฺจโย. เทฺว.}

\proseitem{462}{อุปาทินฺโน ธมฺโม อุปาทินฺนสฺส ธมฺมสฺส วิปฺปยุตฺต\-ปจฺจเยน ปจฺจโย. สหชาตํ, ปุเรชาตํ, ปจฺฉาชาตํ.  สหชาตา ปฏิสนฺธิ\-กฺขเณ อุปาทินฺนา ขนฺธา กฏตฺตารูปานํ วิปฺปยุตฺต\-ปจฺจเยน ปจฺจโย. ขนฺธา วตฺถุสฺส วิปฺปยุตฺต\-ปจฺจเยน ปจฺจโย. วตฺถุ ขนฺธานํ วิปฺปยุตฺต\-ปจฺจเยน ปจฺจโย. ปุเรชาตํ จกฺขายตนํ จกฺขุ\-วิญฺญาณสฺส ฯเปฯ กายายตนํ กายวิญฺญาณสฺส ฯเปฯ. วตฺถุ อุปาทินฺนานํ ขนฺธานํ วิปฺปยุตฺต\-ปจฺจเยน ปจฺจโย. ปจฺฉาชาตา อุปาทินฺนา ขนฺธา ปุเรชาตสฺส อิมสฺส อุปาทินฺนสฺส กายสฺส วิปฺปยุตฺต\-ปจฺจเยน ปจฺจโย. (1)}

\prose{อุปาทินฺโน ธมฺโม อนุปาทินฺนสฺส ธมฺมสฺส วิปฺปยุตฺต\-ปจฺจเยน ปจฺจโย. สหชาตํ, ปุเรชาตํ, ปจฺฉาชาตํ.  สหชาตา อุปาทินฺนา ขนฺธา จิตฺต\-สมุฏฺฐานานํ รูปานํ วิปฺปยุตฺต\-ปจฺจเยน ปจฺจโย. ปุเรชาตํ วตฺถุ อนุปาทินฺนานํ ขนฺธานํ วิปฺปยุตฺต\-ปจฺจเยน ปจฺจโย. ปจฺฉาชาตา อุปาทินฺนา ขนฺธา ปุเรชาตสฺส อิมสฺส อนุปาทินฺนสฺส กายสฺส วิปฺปยุตฺต\-ปจฺจเยน ปจฺจโย. (2)}

\prose{อุปาทินฺโน ธมฺโม อุปาทินฺนสฺส จ อนุปาทินฺนสฺส จ ธมฺมสฺส วิปฺปยุตฺต\-ปจฺจเยน ปจฺจโย. ปจฺฉาชาตา อุปาทินฺนา ขนฺธา ปุเรชาตสฺส อิมสฺส อุปาทินฺนสฺส จ อนุปาทินฺนสฺส จ กายสฺส วิปฺปยุตฺต\-ปจฺจเยน ปจฺจโย. (3)}

\proseitem{463}{\csromanfolio{556}อนุปาทินฺโน ธมฺโม อนุปาทินฺนสฺส ธมฺมสฺส วิปฺปยุตฺต\-ปจฺจเยน ปจฺจโย. สหชาตํ, ปจฺฉาชาตํ.  สหชาตา อนุปาทินฺนา ขนฺธา จิตฺต\-สมุฏฺฐานานํ รูปานํ วิปฺปยุตฺต\-ปจฺจเยน ปจฺจโย. ปจฺฉาชาตา อนุปาทินฺนา ขนฺธา ปุเรชาตสฺส อิมสฺส อนุปาทินฺนสฺส กายสฺส วิปฺปยุตฺต\-ปจฺจเยน ปจฺจโย. (1)}

\prose{อนุปาทินฺโน ธมฺโม อุปาทินฺนสฺส ธมฺมสฺส วิปฺปยุตฺต\-ปจฺจเยน ปจฺจโย. ปจฺฉาชาตา อนุปาทินฺนา ขนฺธา ปุเรชาตสฺส อิมสฺส อุปาทินฺนสฺส กายสฺส วิปฺปยุตฺต\-ปจฺจเยน ปจฺจโย. (2)}

\prose{อนุปาทินฺโน ธมฺโม อุปาทินฺนสฺส จ อนุปาทินฺนสฺส จ ธมฺมสฺส วิปฺปยุตฺต\-ปจฺจเยน ปจฺจโย. ปจฺฉาชาตา อนุปาทินฺนา ขนฺธา ปุเรชาตสฺส อิมสฺส อุปาทินฺนสฺส จ อนุปาทินฺนสฺส จ กายสฺส วิปฺปยุตฺต\-ปจฺจเยน ปจฺจโย. (3)}

\prose{อตฺถิอาทิ}

\proseitem{464}{อุปาทินฺโน ธมฺโม อุปาทินฺนสฺส ธมฺมสฺส อตฺถิปจฺจเยน ปจฺจโย. สหชาตํ, ปุเรชาตํ, ปจฺฉาชาตํ, อาหารํ, อินฺทฺริยํ. (สํขิตฺตํ. ยถา นิกฺขิตฺต\-ปทานิ วิภชิตพฺพานิ ปริปุณฺณานิ.) (1)}

\prose{อุปาทินฺโน ธมฺโม อนุปาทินฺนสฺส ธมฺมสฺส อตฺถิปจฺจเยน ปจฺจโย. สหชาตํ, ปุเรชาตํ, ปจฺฉาชาตํ, อาหารํ. (สํขิตฺตํ. ยถา นิกฺขิตฺต\-ปทานิ วิตฺถาเร\-ตพฺพานิ.) (2)}

\prose{อุปาทินฺโน ธมฺโม อุปาทินฺนสฺส จ อนุปาทินฺนสฺส จ ธมฺมสฺส อตฺถิปจฺจเยน ปจฺจโย. สหชาตํ, ปจฺฉาชาตํ, อาหารํ. (สํขิตฺตํ. ยถา นิกฺขิตฺต\-ปทานิ วิตฺถาเร\-ตพฺพานิ.) (3)}

\prose{อนุปาทินฺโน ธมฺโม อนุปาทินฺนสฺส ธมฺมสฺส อตฺถิปจฺจเยน ปจฺจโย. สหชาตํ, ปุเรชาตํ, ปจฺฉาชาตํ, อาหารํ. (สํขิตฺตํ. ยถา นิกฺขิตฺต\-ปทานิ วิภชิตพฺพานิ.) (1)}

\prose{อนุปาทินฺโน ธมฺโม อุปาทินฺนสฺส ธมฺมสฺส อตฺถิปจฺจเยน ปจฺจโย. ปุเรชาตํ, ปจฺฉาชาตํ, อาหารํ.  ปุเรชาตํ\csromandash{}อนุปาทินฺเน รูเป. สทฺเท. คนฺเธ. รเส. โผฏฺฐพฺเพ อนิจฺจโต ฯเปฯ โทมนสฺสํ อุปฺปชฺชติ. กุสลากุสเล นิรุทฺเธ วิปาโก ตทารมฺมณตา อุปฺปชฺชติ. อนุปาทินฺนํ \csromanfolio{557}รูปายตนํ จกฺขุ\-วิญฺญาณสฺส ฯเปฯ โผฏฺฐพฺพา\-ยตนํ กายวิญฺญาณสฺส ฯเปฯ. ปจฺฉาชาตา อนุปาทินฺนา ขนฺธา ปุเรชาตสฺส อิมสฺส อุปาทินฺนสฺส กายสฺส อตฺถิปจฺจเยน ปจฺจโย. อนุปาทินฺโน กพฬีกาโร อาหาโร อิมสฺส อุปาทินฺนสฺส กายสฺส อตฺถิปจฺจเยน ปจฺจโย. (2)}

\prose{อนุปาทินฺโน ธมฺโม อุปาทินฺนสฺส จ อนุปาทินฺนสฺส จ ธมฺมสฺส อตฺถิปจฺจเยน ปจฺจโย. ปจฺฉาชาตํ, อาหารํ.  \textbf{ปจฺฉาชาตา} อนุปาทินฺนา ขนฺธา ปุเรชาตสฺส อิมสฺส อุปาทินฺนสฺส จ อนุปาทินฺนสฺส จ กายสฺส อตฺถิปจฺจเยน ปจฺจโย. อนุปาทินฺโน \textbf{กพฬีกาโร อาหาโร} อิมสฺส อุปาทินฺนสฺส จ อนุปาทินฺนสฺส จ กายสฺส อตฺถิปจฺจเยน ปจฺจโย. (3)}

\proseitem{465}{อุปาทินฺโน จ อนุปาทินฺโน จ ธมฺมา อุปาทินฺนสฺส ธมฺมสฺส อตฺถิปจฺจเยน ปจฺจโย. ปุเรชาตํ, ปจฺฉาชาตํ, อาหารํ, อินฺทฺริยํ.  ปุเรชาตํ อนุปาทินฺนํ รูปายตนญฺจ จกฺขายตนญฺจ จกฺขุ\-วิญฺญาณสฺส ฯเปฯ. อนุปาทินฺนํ โผฏฺฐพฺพายตนญฺจ กายายตนญฺจ กายวิญฺญาณสฺส อตฺถิปจฺจเยน ปจฺจโย. อนุปาทินฺนํ รูปายตนญฺจ วตฺถุ จ ฯเปฯ โผฏฺฐพฺพายตนญฺจ วตฺถุ จ อุปาทินฺนานํ ขนฺธานํ อตฺถิปจฺจเยน ปจฺจโย. ปจฺฉาชาตา อุปาทินฺนา ขนฺธา จ อนุปาทินฺโน กพฬีกาโร อาหาโร จ อิมสฺส อุปาทินฺนสฺส กายสฺส อตฺถิปจฺจเยน ปจฺจโย. ปจฺฉาชาตา อนุปาทินฺนา ขนฺธา จ รูปชีวิตินฺทฺริยญฺจ กฏตฺตารูปานํ อตฺถิปจฺจเยน ปจฺจโย. (1)}

\prose{อุปาทินฺโน จ อนุปาทินฺโน จ ธมฺมา อนุปาทินฺนสฺส ธมฺมสฺส อตฺถิปจฺจเยน ปจฺจโย. สหชาตํ, ปุเรชาตํ, ปจฺฉาชาตํ, อาหารํ.  สหชาตา อุปาทินฺนา ขนฺธา จ มหาภูตา จ จิตฺต\-สมุฏฺฐานานํ รูปานํ อตฺถิปจฺจเยน ปจฺจโย. สหชาโต อนุปาทินฺโน เอโก ขนฺโธ จ วตฺถุ จ ติณฺณนฺนํ ขนฺธานํ อตฺถิปจฺจเยน ปจฺจโย ฯเปฯ เทฺว ขนฺธา จ ฯเปฯ. ปุเรชาตํ อนุปาทินฺนํ รูปายตนญฺจ วตฺถุ จ อนุปาทินฺนานํ ขนฺธานํ อตฺถิปจฺจเยน ปจฺจโย ฯเปฯ โผฏฺฐพฺพายตนญฺจ วตฺถุ จ อนุปาทินฺนานํ ขนฺธานํ อตฺถิปจฺจเยน ปจฺจโย. ปจฺฉาชาตา อุปาทินฺนา ขนฺธา จ อนุปาทินฺโน กพฬีกาโร อาหาโร จ อิมสฺส อนุปาทินฺนสฺส กายสฺส อตฺถิปจฺจเยน ปจฺจโย. (2)}

\prose{อุปาทินฺโน จ อนุปาทินฺโน จ ธมฺมา อุปาทินฺนสฺส จ อนุปาทินฺนสฺส จ ธมฺมสฺส อตฺถิปจฺจเยน ปจฺจโย. อาหารํ\csromandash{}อุปาทินฺโน จ อนุปาทินฺโน จ \csromanfolio{558}กพฬีกาโร อาหาโร อิมสฺส อุปาทินฺนสฺส จ อนุปาทินฺนสฺส จ กายสฺส อตฺถิปจฺจเยน ปจฺจโย. (3)}

\prose{นตฺถิ\-ปจฺจเยน ปจฺจโย. วิคต\-ปจฺจเยน ปจฺจโย. อวิคต\-ปจฺจเยน ปจฺจโย.}

\csromanmarkh{1. ปจฺจยานุโลม}\csromanmarkhii{2. สงฺขฺยาวาร}\csromanheadernomark{2}{1. \textbf{ปจฺจยานุโลม 2. สงฺขฺยาวาร}}
\proseitem{466}{เหตุยา จตฺตาริ, อารมฺมเณ จตฺตาริ, อธิปติยา เทฺว, อนนฺตเร จตฺตาริ, สมนนฺตเร จตฺตาริ, สหชาเต ปญฺจ, อญฺญมญฺเญ เทฺว, นิสฺสเย ปญฺจ, อุปนิสฺสเย จตฺตาริ, ปุเรชาเต ฉ, ปจฺฉาชาเต ฉ, อาเสวเน เอกํ, กมฺเม ปญฺจ, วิปาเก จตฺตาริ, อาหาเร นว, อินฺทฺริเย จตฺตาริ, ฌาเน จตฺตาริ, มคฺเค จตฺตาริ, สมฺปยุตฺเต เทฺว, วิปฺปยุตฺเต ฉ, อตฺถิยา นว, นตฺถิยา จตฺตาริ, วิคเต จตฺตาริ, อวิคเต นว.}

\proseitem{467}{อุปาทินฺโน ธมฺโม อุปาทินฺนสฺส ธมฺมสฺส อารมฺมณ\-ปจฺจเยน ปจฺจโย. สหชาต\-ปจฺจเยน ปจฺจโย. อุปนิสฺสย\-ปจฺจเยน ปจฺจโย, ปุเรชาต\-ปจฺจเยน ปจฺจโย. ปจฺฉาชาต\-ปจฺจเยน ปจฺจโย. อาหารปจฺจเยน ปจฺจโย. อินฺทฺริย\-ปจฺจเยน ปจฺจโย. (1)}

\prose{อุปาทินฺโน ธมฺโม อนุปาทินฺนสฺส ธมฺมสฺส อารมฺมณ\-ปจฺจเยน ปจฺจโย. สหชาต\-ปจฺจเยน ปจฺจโย. อุปนิสฺสย\-ปจฺจเยน ปจฺจโย. ปุเรชาต\-ปจฺจเยน ปจฺจโย. ปจฺฉาชาต\-ปจฺจเยน ปจฺจโย. อาหารปจฺจเยน ปจฺจโย. (2)}

\prose{อุปาทินฺโน ธมฺโม อุปาทินฺนสฺส จ อนุปาทินฺนสฺส จ ธมฺมสฺส สหชาต\-ปจฺจเยน ปจฺจโย. ปจฺฉาชาต\-ปจฺจเยน ปจฺจโย. อาหารปจฺจเยน ปจฺจโย. (3)}

\proseitem{468}{อนุปาทินฺโน ธมฺโม อนุปาทินฺนสฺส ธมฺมสฺส อารมฺมณ\-ปจฺจเยน ปจฺจโย. สหชาต\-ปจฺจเยน ปจฺจโย, อุปนิสฺสย\-ปจฺจเยน ปจฺจโย. ปจฺฉาชาต\-ปจฺจเยน ปจฺจโย. อาหารปจฺจเยน ปจฺจโย. (1)}

\prose{\csromanfolio{559}อนุปาทินฺโน ธมฺโม อุปาทินฺนสฺส ธมฺมสฺส อารมฺมณ\-ปจฺจเยน ปจฺจโย. อุปนิสฺสย\-ปจฺจเยน ปจฺจโย. ปจฺฉาชาต\-ปจฺจเยน ปจฺจโย. กมฺมปจฺจเยน ปจฺจโย. อาหารปจฺจเยน ปจฺจโย. (2)}

\prose{อนุปาทินฺโน ธมฺโม อุปาทินฺนสฺส จ อนุปาทินฺนสฺส จ ธมฺมสฺส ปจฺฉาชาต\-ปจฺจเยน ปจฺจโย. อาหารปจฺจเยน ปจฺจโย. (3)}

\prose{อุปาทินฺโน จ อนุปาทินฺโน จ ธมฺมา อุปาทินฺนสฺส ธมฺมสฺส ปุเรชาตํ, ปจฺฉาชาตํ, อาหารํ, อินฺทฺริยํ. (1)}

\prose{อุปาทินฺโน จ อนุปาทินฺโน จ ธมฺมา อนุปาทินฺนสฺส ธมฺมสฺส สหชาตํ, ปุเรชาตํ, ปจฺฉาชาตํ, อาหารํ. (2)}

\prose{อุปาทินฺโน จ อนุปาทินฺโน จ ธมฺมา อุปาทินฺนสฺส จ อนุปาทินฺนสฺส จ ธมฺมสฺส อาหารปจฺจเยน ปจฺจโย. (3)}

\csromanmarkh{2. ปจฺจย\-ปจฺจนีย}\csromanmarkhii{2. สงฺขฺยาวาร}\csromanheadernomark{2}{2. ปจฺจย\-ปจฺจนีย 2. สงฺขฺยาวาร}
\proseitem{469}{นเหตุยา นว, นอารมฺมเณ นว, (สพฺพตฺถ นว,) นอาหาเร อฏฺฐ ฯเปฯ นสมฺปยุตฺเต นว, นวิปฺปยุตฺเต นว, โนอตฺถิยา จตฺตาริ, โนนตฺถิยา นว, โนวิคเต นว, โนอวิคเต จตฺตาริ.}

\proseitem{470}{เหตุปจฺจยา นอารมฺมเณ จตฺตาริ ฯเปฯ นอญฺญมญฺเญ ตีณิ, นอุปนิสฺสเย จตฺตาริ, นสมฺปยุตฺเต ตีณิ, นวิปฺปยุตฺเต เทฺว, โนนตฺถิยา จตฺตาริ, โนวิคเต จตฺตาริ.}

\proseitemwithcloser{471}{นเหตุปจฺจยา อารมฺมเณ จตฺตาริ, อธิปติยา เทฺว, (อนุโลมมาติกา กาตพฺพา.) ฯเปฯ อวิคเต นว.}{อุปาทินฺนทุกํ นิฏฺฐิตํ.}
\nitthitam{มหนฺตรทุกํ นิฏฺฐิตํ.}
\csromanmatikaanchor{mtk.68}
\csromantocmarknopage{section}{11. อุปาทานโคจฺฉก}{mtk.68}
\bookmark[dest={mtk.68},level=1]{11. อุปาทานโคจฺฉก}
\csromanheader{1}{11. อุปาทาน\-โคจฺฉก}
\csromanmatikaanchor{mtk.69}
\csromantocmarkref{subsection}{69. อุปาทานทุก}{mtk.69}
\bookmark[dest={mtk.69},level=2]{69. อุปาทานทุก}
\csromanmarkh{69. อุปาทานทุก}\csromanmarkhii{1. ปฏิจฺจวาร}\csromanheadernomark{2}{69. \textbf{อุปาทานทุก 1. ปฏิจฺจวาร}}
\csromanmarkh{1. ปจฺจยานุโลม}\csromanmarkhii{1. วิภงฺควาร}\csromanheadernomark{2}{1. \textbf{ปจฺจยานุโลม 1. วิภงฺควาร}}
\csromanheader{2}{\textbf{เหตุ}}
\proseitem{1}{\csromanfolio{560}อุปาทานํ ธมฺมํ ปฏิจฺจ อุปาทาโน ธมฺโม อุปฺปชฺชติ เหตุปจฺจยา. ทิฏฺฐุปาทานํ ปฏิจฺจ กามุปาทานํ. กามุปาทานํ ปฏิจฺจ ทิฏฺฐุปาทานํ. สีลพฺพตุ\-ปาทานํ ปฏิจฺจ กามุปาทานํ. กามุปาทานํ ปฏิจฺจ สีลพฺพตุ\-ปาทานํ. อตฺตวาทุ\-ปาทานํ ปฏิจฺจ กามุปาทานํ. กามุปาทานํ ปฏิจฺจ อตฺตวาทุ\-ปาทานํ. (1)}

\prose{อุปาทานํ ธมฺมํ ปฏิจฺจ โนอุปาทาโน ธมฺโม อุปฺปชฺชติ เหตุปจฺจยา. อุปาทาเน ปฏิจฺจ สมฺปยุตฺตกา ขนฺธา จิตฺต\-สมุฏฺฐานญฺจ รูปํ. (2)}

\prose{อุปาทานํ ธมฺมํ ปฏิจฺจ อุปาทาโน จ โนอุปาทาโน จ ธมฺมา อุปฺปชฺชนฺติ เหตุปจฺจยา. ทิฏฺฐุปาทานํ ปฏิจฺจ กามุปาทานํ สมฺปยุตฺตกา จ ขนฺธา จิตฺต\-สมุฏฺฐานญฺจ รูปํ. กามุปาทานํ ฯเปฯ. (สพฺพํ จกฺกํ กาตพฺพํ.) (3)}

\proseitem{2}{โนอุปาทานํ ธมฺมํ ปฏิจฺจ โนอุปาทาโน ธมฺโม อุปฺปชฺชติ เหตุปจฺจยา. โนอุปาทานํ เอกํ ขนฺธํ ปฏิจฺจ ตโย ขนฺธา จิตฺต\-สมุฏฺฐานญฺจ รูปํ ฯเปฯ เทฺว ขนฺเธ ฯเปฯ. ปฏิสนฺธิ\-กฺขเณ โนอุปาทานํ เอกํ ขนฺธํ ปฏิจฺจ ตโย ขนฺธา กฏตฺตา จ รูปํ ฯเปฯ เทฺว ขนฺเธ ฯเปฯ. ขนฺเธ ปฏิจฺจ วตฺถุ, วตฺถุํ ปฏิจฺจ ขนฺธา. เอกํ มหาภูตํ ฯเปฯ มหาภูเต ปฏิจฺจ จิตฺต\-สมุฏฺฐานํ รูปํ, กฏตฺตารูปํ, อุปาทารูปํ. (1)}

\prose{โนอุปาทานํ ธมฺมํ ปฏิจฺจ อุปาทาโน ธมฺโม อุปฺปชฺชติ เหตุปจฺจยา. โนอุปาทาเน ขนฺเธ ปฏิจฺจ อุปาทานา. (2)}

\prose{โนอุปาทานํ ธมฺมํ ปฏิจฺจ อุปาทาโน จ โนอุปาทาโน จ ธมฺมา อุปฺปชฺชนฺติ เหตุปจฺจยา. โนอุปาทานํ เอกํ ขนฺธํ ปฏิจฺจ ตโย ขนฺธา อุปาทานา จ จิตฺต\-สมุฏฺฐานญฺจ รูปํ ฯเปฯ เทฺว ขนฺเธ ปฏิจฺจ เทฺว ขนฺธา อุปาทานา จ จิตฺต\-สมุฏฺฐานญฺจ รูปํ. (3)}

\proseitem{3}{\csromanfolio{561}อุปาทานญฺจ โนอุปาทานญฺจ ธมฺมํ ปฏิจฺจ อุปาทาโน ธมฺโม อุปฺปชฺชติ เหตุปจฺจยา. ทิฏฺฐุปาทานญฺจ สมฺปยุตฺตเก จ ขนฺเธ ปฏิจฺจ กามุปาทานํ. (สพฺเพ จกฺกา กาตพฺพา.) (1)}

\prose{อุปาทานญฺจ โนอุปาทานญฺจ ธมฺมํ ปฏิจฺจ โนอุปาทาโน ธมฺโม อุปฺปชฺชติ เหตุปจฺจยา. โนอุปาทานํ เอกํ ขนฺธญฺจ อุปาทานญฺจ ปฏิจฺจ ตโย ขนฺธา จิตฺต\-สมุฏฺฐานญฺจ รูปํ ฯเปฯ เทฺว ขนฺเธ จ ฯเปฯ. อุปาทานญฺจ มหาภูเต จ ปฏิจฺจ จิตฺต\-สมุฏฺฐานํ รูปํ. (2)}

\prose{อุปาทานญฺจ โนอุปาทานญฺจ ธมฺมํ ปฏิจฺจ อุปาทาโน จ โนอุปาทาโน จ ธมฺมา อุปฺปชฺชนฺติ เหตุปจฺจยา. โนอุปาทานํ เอกํ ขนฺธญฺจ ทิฏฺฐุปาทานญฺจ ปฏิจฺจ ตโย ขนฺธา กามุปาทานญฺจ จิตฺต\-สมุฏฺฐานํ รูปํ. เทฺว ขนฺเธ จ ฯเปฯ. (จกฺกํ กาตพฺพํ.) (3)}

\proseitem{4}{อุปาทานํ ธมฺมํ ปฏิจฺจ อุปาทาโน ธมฺโม อุปฺปชฺชติ อารมฺมณ\-ปจฺจยา. (นวปิ ปญฺหา กาตพฺพา. รูปํ ฉฑฺเฑตพฺพํ.)}

\csromanmarkh{1. ปจฺจยานุโลม}\csromanmarkhii{2. สงฺขฺยาวาร}\csromanheadernomark{2}{1. \textbf{ปจฺจยานุโลม 2. สงฺขฺยาวาร}}
\proseitem{5}{เหตุยา นว, อารมฺมเณ นว, อธิปติยา นว, (สพฺพตฺถ นว.) วิปาเก เอกํ ฯเปฯ อวิคเต นว.}

\csromanmarkh{2. ปจฺจย\-ปจฺจนีย}\csromanmarkhii{1. วิภงฺควาร}\csromanheadernomark{2}{2. ปจฺจย\-ปจฺจนีย 1. วิภงฺควาร}
\proseitem{6}{โนอุปาทานํ ธมฺมํ ปฏิจฺจ โนอุปาทาโน ธมฺโม อุปฺปชฺชติ นเหตุปจฺจยา. อเหตุกํ โนอุปาทานํ เอกํ ขนฺธํ ปฏิจฺจ ตโย ขนฺธา จิตฺต\-สมุฏฺฐานญฺจ รูปํ ฯเปฯ เทฺว ขนฺเธ ฯเปฯ. อเหตุก\-ปฏิสนฺธิ\-กฺขเณ ฯเปฯ. ขนฺเธ ปฏิจฺจ วตฺถุ, วตฺถุํ ปฏิจฺจ ขนฺธา. เอกํ มหาภูตํ ฯเปฯ มหาภูเต ปฏิจฺจ จิตฺต\-สมุฏฺฐานํ รูปํ, กฏตฺตารูปํ, อุปาทารูปํ. พาหิรํ. อาหารสมุฏฺฐานํ. อุตุสมุฏฺฐานํ. อสญฺญสตฺตานํ ฯเปฯ. วิจิกิจฺฉา\-สหคเต อุทฺธจฺจ\-สหคเต ขนฺเธ ปฏิจฺจ วิจิกิจฺฉา\-สหคโต อุทฺธจฺจ\-สหคโต โมโห. (1)}

\prose{\csromanfolio{562}นอารมฺมณาทิ}

\proseitem{7}{อุปาทานํ ธมฺมํ ปฏิจฺจ โนอุปาทาโน ธมฺโม อุปฺปชฺชติ นอารมฺมณ\-ปจฺจยา. อุปาทาเน ปฏิจฺจ จิตฺต\-สมุฏฺฐานํ รูปํ. (1)}

\prose{โนอุปาทานํ ธมฺมํ ปฏิจฺจ โนอุปาทาโน ธมฺโม อุปฺปชฺชติ นอารมฺมณ\-ปจฺจยา. โนอุปาทาเน ขนฺเธ ปฏิจฺจ จิตฺต\-สมุฏฺฐานํ รูปํ. ปฏิสนฺธิ\-กฺขเณ โนอุปาทาเน ขนฺเธ ปฏิจฺจ กฏตฺตารูปํ. ขนฺเธ ปฏิจฺจ วตฺถุ. เอกํ มหาภูตํ ฯเปฯ. (ยาว อสญฺญสตฺตา.) (1)}

\prose{อุปาทานญฺจ โนอุปาทานญฺจ ธมฺมํ ปฏิจฺจ โนอุปาทาโน ธมฺโม อุปฺปชฺชติ นอารมฺมณ\-ปจฺจยา. อุปาทาเน จ สมฺปยุตฺตเก จ ขนฺเธ ปฏิจฺจ จิตฺต\-สมุฏฺฐานํ รูปํ. อุปาทาเน จ มหาภูเต จ ปฏิจฺจ จิตฺต\-สมุฏฺฐานํ รูปํ. (1)}

\prose{นอธิปติ\-ปจฺจยา. นอนนฺตร\-ปจฺจยา ฯเปฯ. นอุปนิสฺสยปจฺจยา.}

\csromanheader{2}{\textbf{นปุเรชาต}}
\proseitem{8}{อุปาทานํ ธมฺมํ ปฏิจฺจ อุปาทาโน ธมฺโม อุปฺปชฺชติ นปุเรชาต\-ปจฺจยา. อรูเป อตฺตวาทุ\-ปาทานํ ปฏิจฺจ กามุปาทานํ. กามุปาทานํ ปฏิจฺจ อตฺตวาทุ\-ปาทานํ. (1)}

\prose{อุปาทานํ ธมฺมํ ปฏิจฺจ โนอุปาทาโน ธมฺโม อุปฺปชฺชติ นปุเรชาต\-ปจฺจยา. อรูเป อุปาทาเน ปฏิจฺจ สมฺปยุตฺตกา ขนฺธา. อุปาทาเน ปฏิจฺจ จิตฺต\-สมุฏฺฐานํ รูปํ. (สํขิตฺตํ. นวปิ ปญฺหา อรูเป เทฺว อุปาทานา.)}

\csromanmarkh{2. ปจฺจย\-ปจฺจนีย}\csromanmarkhii{2. สงฺขฺยาวาร}\csromanheadernomark{2}{2. ปจฺจย\-ปจฺจนีย 2. สงฺขฺยาวาร}
\proseitem{9}{นเหตุยา เอกํ, นอารมฺมเณ ตีณิ, นอธิปติยา นว, นอนนฺตเร ตีณิ ฯเปฯ นอุปนิสฺสเย ตีณิ, นปุเรชาเต นว, นปจฺฉาชาเต นว, นอาเสวเน นว, นกมฺเม ตีณิ, นวิปาเก นว, นอาหาเร เอกํ, นอินฺทฺริเย เอกํ, นฌาเน เอกํ, นมคฺเค เอกํ, นสมฺปยุตฺเต ตีณิ, นวิปฺปยุตฺเต นว, โนนตฺถิยา ตีณิ, โนวิคเต ตีณิ.}

\csromanmarkh{69. อุปาทานทุก}\csromanmarkhii{3. ปจฺจยา\-นุโลม\-ปจฺจนีย}\csromanheadernomark{2}{69. อุปาทานทุก 3. ปจฺจยา\-นุโลม\-ปจฺจนีย}
\vspace{\tipitakaparskip}
\csromancenter{\textbf{เหตุ}ปจฺจยา นอารมฺมเณ ตีณิ, นอธิปติยา นว ฯเปฯ นฺเอาปนิสฺสเย ตีณิ, นปุเรชาเต นว, นปจฺฉาชาเต นว, นอาเสวเน นว, นกมฺเม ตีณิ, นวิปาเก นว, นสมฺปยุตฺเต ตีณิ, นวิปฺปยุตฺเต นว, โนนตฺถิยา ตีณิ, โนวิคเต ตีณิ.}
\proseitem{11}{\csromanfolio{563}นเหตุปจฺจยา อารมฺมเณ เอกํ, (สพฺพตฺถ เอกํ.) มคฺเค เอกํ ฯเปฯ อวิคเต เอกํ.}

\prose{(สหชาตวาโร ปฏิจฺจ\-วาร\-สทิโส. วิภชนฺเตน ทิฏฺฐุปาทานํ “สหชาตํ กามุปาทานนฺ”ติ กาตพฺพํ.)}

\csromanmarkh{69. อุปาทานทุก}\csromanmarkhii{3. ปจฺจยวาร}\csromanheadernomark{2}{69. อุปาทานทุก 3. ปจฺจยวาร}
\vspace{\tipitakaparskip}
\csromancenter{\textbf{เหตุ}}
\proseitem{12}{อุปาทานํ ธมฺมํ ปจฺจยา อุปาทาโน ธมฺโม อุปฺปชฺชติ เหตุปจฺจยา. ทิฏฺฐุปาทานํ ปจฺจยา กามุปาทานํ. ตีณิ. (ปฏิจฺจสทิสา.)}

\prose{โนอุปาทานํ ธมฺมํ ปจฺจยา โนอุปาทาโน ธมฺโม อุปฺปชฺชติ เหตุปจฺจยา. โนอุปาทานํ เอกํ ขนฺธํ ปจฺจยา ตโย ขนฺธา จิตฺต\-สมุฏฺฐานญฺจ รูปํ ฯเปฯ. เทฺว ขนฺเธ ฯเปฯ. ปฏิสนฺธิ\-กฺขเณ ฯเปฯ. ขนฺเธ ปจฺจยา วตฺถุ ฯเปฯ. (ยาว อชฺฌตฺติกา มหาภูตา.) วตฺถุํ ปจฺจยา โนอุปาทานา ขนฺธา. (1)}

\prose{โนอุปาทานํ ธมฺมํ ปจฺจยา อุปาทาโน ธมฺโม อุปฺปชฺชติ เหตุปจฺจยา. โนอุปาทาเน ขนฺเธ ปจฺจยา อุปาทานา. วตฺถุํ ปจฺจยา อุปาทานา. (2)}

\prose{โนอุปาทานํ ธมฺมํ ปจฺจยา อุปาทาโน จ โนอุปาทาโน จ ธมฺมา อุปฺปชฺชนฺติ เหตุปจฺจยา. โนอุปาทานํ เอกํ ขนฺธํ ปจฺจยา ตโย ขนฺธา อุปาทานา จ จิตฺต\-สมุฏฺฐานญฺจ รูปํ ฯเปฯ เทฺว ขนฺเธ ฯเปฯ. วตฺถุํ ปจฺจยา อุปาทานา. มหาภูเต ปจฺจยา จิตฺต\-สมุฏฺฐานํ รูปํ. วตฺถุํ ปจฺจยา อุปาทานา สมฺปยุตฺตกา จ ขนฺธา. (3)}

\proseitem{13}{\csromanfolio{564}อุปาทานญฺจ โนอุปาทานญฺจ ธมฺมํ ปจฺจยา อุปาทาโน ธมฺโม อุปฺปชฺชติ เหตุปจฺจยา. ทิฏฺฐุปาทานญฺจ สมฺปยุตฺตเก จ ขนฺเธ ปจฺจยา กามุปาทานํ. กามุปาทานญฺจ สมฺปยุตฺตเก จ ขนฺเธ ปจฺจยา ทิฏฺฐุปาทานํ. (จกฺกํ กาตพฺพํ.) ทิฏฺฐุปาทานญฺจ วตฺถุญฺจ ปจฺจยา กามุปาทานํ. (จกฺกํ กาตพฺพํ.) (1)}

\prose{อุปาทานญฺจ โนอุปาทานญฺจ ธมฺมํ ปจฺจยา โนอุปาทาโน ธมฺโม อุปฺปชฺชติ เหตุปจฺจยา. โนอุปาทานํ เอกํ ขนฺธญฺจ อุปาทานญฺจ ปจฺจยา ตโย ขนฺธา จิตฺต\-สมุฏฺฐานญฺจ รูปํ ฯเปฯ เทฺว ขนฺเธ จ ฯเปฯ. (จกฺกํ กาตพฺพํ.) อุปาทาเน จ มหาภูเต จ ปจฺจยา จิตฺต\-สมุฏฺฐานํ รูปํ. อุปาทานญฺจ วตฺถุญฺจ ปจฺจยา โนอุปาทาโน ขนฺธา. (2)}

\prose{อุปาทานญฺจ โนอุปาทานญฺจ ธมฺมํ ปจฺจยา อุปาทาโน จ โนอุปาทาโน จ ธมฺมา อุปฺปชฺชนฺติ เหตุปจฺจยา. โนอุปาทานํ เอกํ ขนฺธญฺจ ทิฏฺฐุปาทานญฺจ ปจฺจยา ตโย ขนฺธา กามุปาทานํ จิตฺต\-สมุฏฺฐานญฺจ รูปํ ฯเปฯ. (จกฺกํ.) ทิฏฺฐุปาทานญฺจ วตฺถุญฺจ ปจฺจยา กามุปาทานํ สมฺปยุตฺตกา จ ขนฺธา ฯเปฯ. (จกฺกํ.) (3)}

\prose{อารมฺมณ\-ปจฺจยา.}

\vspace{\tipitakaparskip}
\csromancenter{(อารมฺมเณ โนอุปาทาน\-มูลเก ปญฺจายตนญฺจ วตฺถุญฺจ กาตพฺพา.)}
\prose{เหตุยา นว, อารมฺมเณ นว, อธิปติยา นว, (สพฺพตฺถ นว,) วิปาเก เอกํ ฯเปฯ อวิคเต นว. อนุโลมํ.}

\proseitem{14}{โนอุปาทานํ ธมฺมํ ปจฺจยา โนอุปาทาโน ธมฺโม อุปฺปชฺชติ นเหตุปจฺจยา. อเหตุกํ โนอุปาทานํ เอกํ ขนฺธํ ปจฺจยา ฯเปฯ. อเหตุก\-ปฏิสนฺธิ\-กฺขเณ ฯเปฯ. (ยาว อสญฺญสตฺตา.) จกฺขายตนํ ปจฺจยา จกฺขุวิญฺญาณํ ฯเปฯ. กายายตนํ ปจฺจยา กายวิญฺญาณํ. วตฺถุํ ปจฺจยา อเหตุกา โนอุปาทานา ขนฺธา. วิจิกิจฺฉา\-สหคเต อุทฺธจฺจ\-สหคเต ขนฺเธ จ วตฺถุญฺจ ปจฺจยา วิจิกิจฺฉา\-สหคโต อุทฺธจฺจ\-สหคโต โมโห.}

\prose{นเหตุยา เอกํ, นอารมฺมเณ ตีณิ, นอธิปติยา นว, นอนนฺตเร ตีณิ, นสมนนฺตเร ตีณิ ฯเปฯ นปุเรชาเต นว ฯเปฯ นกมฺเม ตีณิ, นวิปาเก นว, (ปฏิจฺจสทิสํ.) ฯเปฯ โนวิคเต ตีณิ.}

\vspace{\tipitakaparskip}
\csromancenter{ปจฺจนียํ.}
\csromancenter{(เอวํ อิตเร เทฺว คณนาปิ นิสฺสยวาโรปิ กาตพฺโพ.)}
\csromanmarkh{69. อุปาทานทุก}\csromanmarkhii{69. อุปาทานทุก 5. สํสฏฺฐวาร}\csromanheadernomark{2}{69. อุปาทานทุก \textbf{69. อุปาทานทุก 5. สํสฏฺฐวาร}}
\csromanheader{2}{\textbf{เหตุ}}
\proseitem{15}{\csromanfolio{565}อุปาทานํ ธมฺมํ สํสฏฺโฐ อุปาทาโน ธมฺโม อุปฺปชฺชติ เหตุปจฺจยา. ทิฏฺฐุปาทานํ สํสฏฺฐํ กามุปาทานํ. กามุปาทานํ สํสฏฺฐํ ทิฏฺฐุปาทานํ. (จกฺกํ. เอวํ นวปิ ปญฺหา กาตพฺพา.)}

\prose{\textbf{เหตุ}ยา นว, อารมฺมเณ นว, อธิปติยา นว, (สพฺพตฺถ นว.) วิปาเก เอกํ ฯเปฯ วิคเต นว, อวิคเต นว. อนุโลมํ.}

\proseitem{16}{โนอุปาทานํ ธมฺมํ สํสฏฺโฐ โนอุปาทาโน ธมฺโม อุปฺปชฺชติ นเหตุปจฺจยา. อเหตุกํ โนอุปาทานํ เอกํ ขนฺธํ สํสฏฺฐา ตโย ขนฺธา ฯเปฯ. เทฺว ขนฺเธ ฯเปฯ. อเหตุก\-ปฏิสนฺธิ\-กฺขเณ ฯเปฯ. วิจิกิจฺฉา\-สหคเต อุทฺธจฺจ\-สหคเต ขนฺเธ สํสฏฺโฐ วิจิกิจฺฉา\-สหคโต อุทฺธจฺจ\-สหคโต โมโห.}

\prose{นเหตุยา เอกํ, นอธิปติยา นว, นปุเรชาเต นว, นปจฺฉาชาเต นว, นอาเสวเน นว, นกมฺเม ตีณิ, นวิปาเก นว, นฌาเน เอกํ, นมคฺเค เอกํ, นวิปฺปยุตฺเต นว.}

\vspace{\tipitakaparskip}
\csromancenter{ปจฺจนียํ.}
\csromancenter{(เอวํ อิตเร เทฺว คณนาปิ สมฺปยุตฺตวาโรปิ กาตพฺโพ.)}
\csromanmarkh{69. อุปาทานทุก}\csromanmarkhii{7. ปญฺหาวาร}\csromanheadernomark{2}{69. อุปาทานทุก 7. ปญฺหาวาร}
\csromanmarkh{1. ปจฺจยานุโลม}\csromanmarkhii{1. วิภงฺควาร}\csromanheadernomark{2}{1. \textbf{ปจฺจยานุโลม 1. วิภงฺควาร}}
\csromanheader{2}{\textbf{เหตุ}}
\proseitem{17}{อุปาทาโน ธมฺโม อุปาทานสฺส ธมฺมสฺส เหตุปจฺจเยน ปจฺจโย. อุปาทานา เหตู สมฺปยุตฺตกานํ อุปาทานานํ เหตุปจฺจเยน ปจฺจโย. (มูลํ กาตพฺพํ.) อุปาทานา เหตู สมฺปยุตฺตกานํ ขนฺธานํ จิตฺตสมุฏฺฐานานญฺจ \csromanfolio{566}รูปานํ เหตุปจฺจเยน ปจฺจโย. (มูลํ กาตพฺพํ.) อุปาทานา เหตู สมฺปยุตฺตกานํ ขนฺธานํ อุปาทานานญฺจ จิตฺต\-สมุฏฺฐานานํ รูปานํ เหตุปจฺจเยน ปจฺจโย. (3)}

\prose{โนอุปาทาโน ธมฺโม โนอุปาทานสฺส ธมฺมสฺส เหตุปจฺจเยน ปจฺจโย. โนอุปาทานา เหตู สมฺปยุตฺตกานํ ขนฺธานํ จิตฺตสมุฏฺฐานานญฺจ รูปานํ เหตุปจฺจเยน ปจฺจโย. ปฏิสนฺธิ\-กฺขเณ ฯเปฯ. (1)}

\prose{โนอุปาทาโน ธมฺโม อุปาทานสฺส ธมฺมสฺส เหตุปจฺจเยน ปจฺจโย. โนอุปาทานา เหตู สมฺปยุตฺตกานํ อุปาทานานํ เหตุปจฺจเยน ปจฺจโย. (มูลํ กาตพฺพํ.) โนอุปาทานา เหตู สมฺปยุตฺตกานํ ขนฺธานํ อุปาทานานญฺจ จิตฺต\-สมุฏฺฐานานํ รูปานํ เหตุปจฺจเยน ปจฺจโย. (3)}

\proseitem{19}{อุปาทาโน จ โนอุปาทาโน จ ธมฺมา อุปาทานสฺส ธมฺมสฺส เหตุปจฺจเยน ปจฺจโย. อุปาทานา จ โนอุปาทานา จ เหตู สมฺปยุตฺตกานํ อุปาทานานํ เหตุปจฺจเยน ปจฺจโย. (มูลํ กาตพฺพํ.) อุปาทานา จ โนอุปาทานา จ เหตู สมฺปยุตฺตกานํ ขนฺธานํ จิตฺต\-สมุฏฺฐานญฺจ รูปานํ เหตุปจฺจเยน ปจฺจโย. (มูลํ กาตพฺพํ.) อุปาทานา จ โนอุปาทานา จ เหตู สมฺปยุตฺตกานํ ขนฺธานํ อุปาทานานญฺจ จิตฺต\-สมุฏฺฐานานํ รูปานํ เหตุปจฺจเยน ปจฺจโย. (3)}

\proseitem{17}{อุปาทาโน ธมฺโม อุปาทานสฺส ธมฺมสฺส อารมฺมณ\-ปจฺจเยน ปจฺจโย. อุปาทาเน อารพฺภ อุปาทานา อุปฺปชฺชนฺติ. ตีณิ. (อารพฺภ กาตพฺพา.)}

\prose{โนอุปาทาโน ธมฺโม โนอุปาทานสฺส ธมฺมสฺส อารมฺมณ\-ปจฺจเยน ปจฺจโย. ทานํ ทตฺวา, สีลํ ฯเปฯ อุโปสถกมฺมํ กตฺวา ตํ ปจฺจเวกฺขติ, อสฺสาเทติ อภินนฺทติ, ตํ อารพฺภ ราโค อุปฺปชฺชติ. ทิฏฺฐิ อุปฺปชฺชติ. วิจิกิจฺฉา ฯเปฯ. อุทฺธจฺจํ ฯเปฯ. โทมนสฺสํ อุปฺปชฺชติ. ปุพฺเพ สุจิณฺณานิ ฯเปฯ. ฌานา วุฏฺฐหิตฺวา ฌานํ ปจฺจเวกฺขติ. อริยา มคฺคา วุฏฺฐหิตฺวา มคฺคํ ปจฺจเวกฺขนฺติ. ผลํ ฯเปฯ. นิพฺพานํ ปจฺจเวกฺขนฺติ. นิพฺพานํ โคตฺรภุสฺส, โวทานสฺส, มคฺคสฺส, ผลสฺส, อาวชฺชนาย อารมฺมณ\-ปจฺจเยน ปจฺจโย. อริยา โนอุปาทาเน ปหีเน \csromanfolio{567}กิเลเส ฯเปฯ. วิกฺขมฺภิเต กิเลเส ปจฺจเวกฺขนฺติ, ปุพฺเพ สมุทาจิณฺเณ ฯเปฯ. จกฺขุํ ฯเปฯ วตฺถุํ โนอุปาทาเน ขนฺเธ อนิจฺจโต ฯเปฯ โทมนสฺสํ อุปฺปชฺชติ. ทิพฺเพน จกฺขุนา รูปํ ปสฺสติ. ทิพฺพาย โสตธาตุยา สทฺทํ สุณาติ. เจโตปริย\-ญาเณน โนอุปาทานจิตฺตสมงฺคิสฺส จิตฺตํ ชานาติ. อากาสา\-นญฺจา\-ยตนํ วิญฺญาณญฺจา\-ยตนสฺส ฯเปฯ. อากิญฺจญฺญา\-ยตนํ เนว\-สญฺญา\-นา\-สญฺญา\-ยตนสฺส ฯเปฯ. รูปายตนํ จกฺขุ\-วิญฺญาณสฺส ฯเปฯ. โผฏฺฐพฺพา\-ยตนํ กายวิญฺญาณสฺส ฯเปฯ. โนอุปาทานา ขนฺธา อิทฺธิวิธ\-ญาณสฺส, เจโตปริย\-ญาณสฺส, ปุพฺเพ\-นิวาสา\-นุสฺสติ\-ญาณสฺส, ยถากมฺมูปคญาณสฺส, อนาคตํสญาณสฺส, อาวชฺชนาย อารมฺมณ\-ปจฺจเยน ปจฺจโย. (1)}

\proseitem{20}{โนอุปาทาโน ธมฺโม อุปาทานสฺส ธมฺมสฺส อารมฺมณ\-ปจฺจเยน ปจฺจโย. ทานํ ฯเปฯ สีลํ ฯเปฯ อุโปสถกมฺมํ กตฺวา ตํ อสฺสาเทติ อภินนฺทติ, ตํ อารพฺภ ราโค อุปฺปชฺชติ. ทิฏฺฐิ อุปฺปชฺชติ. ปุพฺเพ สุจิณฺณานิ ฯเปฯ. ฌานา ฯเปฯ. จกฺขุํ ฯเปฯ วตฺถุํ โนอุปาทาเน ขนฺเธ อสฺสาเทติ อภินนฺทติ, ตํ อารพฺภ ราโค อุปฺปชฺชติ. ทิฏฺฐิ อุปฺปชฺชติ. (2)}

\prose{โนอุปาทาโน ธมฺโม อุปาทานสฺส จ โนอุปาทานสฺส จ ธมฺมสฺส อารมฺมณ\-ปจฺจเยน ปจฺจโย. ทานํ ฯเปฯ สีลํ ฯเปฯ อุโปสถกมฺมํ กตฺวา ฯเปฯ. ปุพฺเพ สุจิณฺณานิ ฯเปฯ. ฌานา วุฏฺฐหิตฺวา ฌานํ ฯเปฯ. จกฺขุํ ฯเปฯ วตฺถุํ โนอุปาทาเน ขนฺเธ อสฺสาเทติ อภินนฺทติ, ตํ อารพฺภ อุปาทานา จ สมฺปยุตฺตกา จ ขนฺธา อุปฺปชฺชนฺติ. (3)}

\prose{อุปาทาโน จ โนอุปาทาโน จ ธมฺมา อุปาทานสฺส ธมฺมสฺส อารมฺมณ\-ปจฺจเยน ปจฺจโย. ตีณิ. (อารพฺภ กาตพฺพา.)}

\proseitem{21}{อุปาทาโน ธมฺโม อุปาทานสฺส ธมฺมสฺส อธิปติ\-ปจฺจเยน ปจฺจโย. อารมฺมณา\-ธิปติ\csromandash{}อุปาทาเน ครุํ กตฺวา อุปาทานา อุปฺปชฺชนฺติ. ตีณิ. (อารมฺมณา\-ธิปติเยว.)}

\proseitem{20}{โนอุปาทาโน ธมฺโม โนอุปาทานสฺส ธมฺมสฺส อธิปติ\-ปจฺจเยน ปจฺจโย. อารมฺมณา\-ธิปติ, สหชาตา\-ธิปติ.  อารมฺมณา\-ธิปติ\csromandash{}ทานํ ทตฺวา, สีลํ ฯเปฯ อุโปสถกมฺมํ ฯเปฯ. ปุพฺเพ ฯเปฯ. ฌานา วุฏฺฐหิตฺวา ฌานํ ครุํ กตฺวา ปจฺจเวกฺขติ, อสฺสาเทติ อภินนฺทติ, อริยา มคฺคา \csromanfolio{568}วุฏฺฐหิตฺวา มคฺคํ ครุํ กตฺวา ปจฺจเวกฺขนฺติ ฯเปฯ ผลสฺส อธิปติ\-ปจฺจเยน ปจฺจโย. จกฺขุํ ฯเปฯ วตฺถุํ โนอุปาทาเน ขนฺเธ ครุํ กตฺวา อสฺสาเทติ อภินนฺทติ, ตํ ครุํ กตฺวา ราโค อุปฺปชฺชติ. ทิฏฺฐิ อุปฺปชฺชติ. สหชาตา\-ธิปติ\csromandash{}โนอุปาทานา\-ธิปติ สมฺปยุตฺตกานํ ขนฺธานํ จิตฺตสมุฏฺฐานานญฺจ รูปานํ อธิปติ\-ปจฺจเยน ปจฺจโย. (1)}

\proseitem{22}{โนอุปาทาโน ธมฺโม อุปาทานสฺส ธมฺมสฺส อธิปติ\-ปจฺจเยน ปจฺจโย. อารมฺมณา\-ธิปติ, สหชาตา\-ธิปติ.  อารมฺมณา\-ธิปติ\csromandash{}ทานํ ฯเปฯ สีลํ ฯเปฯ อุโปสถกมฺมํ กตฺวา ตํ ครุํ กตฺวา อสฺสาเทติ อภินนฺทติ, ตํ ครุํ กตฺวา ราโค อุปฺปชฺชติ. ทิฏฺฐิ อุปฺปชฺชติ. ปุพฺเพ ฯเปฯ. ฌานา ฯเปฯ. จกฺขุํ ฯเปฯ วตฺถุํ โนอุปาทาเน ขนฺเธ ครุํ กตฺวา อสฺสาเทติ อภินนฺทติ, ตํ ครุํ กตฺวา ราโค อุปฺปชฺชติ. ทิฏฺฐิ อุปฺปชฺชติ. สหชาตา\-ธิปติ\csromandash{}โนอุปาทานา\-ธิปติ สมฺปยุตฺตกานํ อุปาทานานํ อธิปติ\-ปจฺจเยน ปจฺจโย. (2)}

\prose{โนอุปาทาโน ธมฺโม อุปาทานสฺส จ โนอุปาทานสฺส จ ธมฺมสฺส อธิปติ\-ปจฺจเยน ปจฺจโย. อารมฺมณา\-ธิปติ, สหชาตา\-ธิปติ.  อารมฺมณา\-ธิปติ\csromandash{}ทานํ ฯเปฯ สีลํ ฯเปฯ อุโปสถกมฺมํ ฯเปฯ. ปุพฺเพ ฯเปฯ. ฌานํ ฯเปฯ. จกฺขุํ ฯเปฯ วตฺถุํ โนอุปาทาเน ขนฺเธ ครุํ กตฺวา อสฺสาเทติ อภินนฺทติ, ตํ ครุํ กตฺวา อุปาทานา จ สมฺปยุตฺตกา จ ขนฺธา อุปฺปชฺชนฺติ. สหชาตา\-ธิปติ\csromandash{}โนอุปาทานา\-ธิปติ สมฺปยุตฺตกานํ ขนฺธานํ อุปาทานานญฺจ จิตฺต\-สมุฏฺฐานานํ รูปานํ อธิปติ\-ปจฺจเยน ปจฺจโย. (3)}

\prose{อุปาทาโน จ โนอุปาทาโน จ ธมฺมา อุปาทานสฺส ธมฺมสฺส อธิปติ\-ปจฺจเยน ปจฺจโย. ตีณิ. อารมฺมณา\-ธิปติ. ตีณิ. (อารพฺภ กาตพฺพา. อารมฺมณา\-ธิปติเยว.)}

\proseitem{23}{อุปาทาโน ธมฺโม อุปาทานสฺส ธมฺมสฺส อนนฺตร\-ปจฺจเยน ปจฺจโย. ปุริมา ปุริมา อุปาทานา ปจฺฉิมานํ ปจฺฉิมานํ อุปาทานานํ อนนฺตร\-ปจฺจเยน ปจฺจโย. (มูลํ กาตพฺพํ.) ปุริมา ปุริมา อุปาทานา ปจฺฉิมานํ ปจฺฉิมานํ โนอุปาทานานํ ขนฺธานํ อนนฺตร\-ปจฺจเยน ปจฺจโย. อุปาทานํ วุฏฺฐานสฺส อนนฺตร\-ปจฺจเยน ปจฺจโย. (มูลํ กาตพฺพํ.) ปุริมา ปุริมา อุปาทานา ปจฺฉิมานํ ปจฺฉิมานํ อุปาทานานํ สมฺปยุตฺตกานญฺจ ขนฺธานํ อนนฺตร\-ปจฺจเยน ปจฺจโย. (3)}

\proseitem{24}{\csromanfolio{569}โนอุปาทาโน ธมฺโม โนอุปาทานสฺส ธมฺมสฺส อนนฺตร\-ปจฺจเยน ปจฺจโย. ปุริมา ปุริมา โนอุปาทานา ขนฺธา ปจฺฉิมานํ ปจฺฉิมานํ โนอุปาทานานํ ขนฺธานํ อนนฺตร\-ปจฺจเยน ปจฺจโย. อนุโลมํ โคตฺรภุสฺส ฯเปฯ. ผลสมาปตฺติยา อนนฺตร\-ปจฺจเยน ปจฺจโย. (1)}

\prose{โนอุปาทาโน ธมฺโม อุปาทานสฺส ธมฺมสฺส อนนฺตร\-ปจฺจเยน ปจฺจโย. ปุริมา ปุริมา โนอุปาทานา ขนฺธา ปจฺฉิมานํ ปจฺฉิมานํ อุปาทานานํ ขนฺธานํ อนนฺตร\-ปจฺจเยน ปจฺจโย. อาวชฺชนา อุปาทานานํ อนนฺตร\-ปจฺจเยน ปจฺจโย. (2)}

\prose{โนอุปาทาโน ธมฺโม อุปาทานสฺส จ โนอุปาทานสฺส จ ธมฺมสฺส อนนฺตร\-ปจฺจเยน ปจฺจโย. ปุริมา ปุริมา โนอุปาทานา ขนฺธา ปจฺฉิมานํ ปจฺฉิมานํ อุปาทานานํ สมฺปยุตฺตกานญฺจ ขนฺธานํ อนนฺตร\-ปจฺจเยน ปจฺจโย. อาวชฺชนา อุปาทานานํ สมฺปยุตฺตกานญฺจ ขนฺธานํ อนนฺตร\-ปจฺจเยน ปจฺจโย. (3)}

\proseitem{25}{อุปาทาโน จ โนอุปาทาโน จ ธมฺมา อุปาทานสฺส ธมฺมสฺส อนนฺตร\-ปจฺจเยน ปจฺจโย. ปุริมา ปุริมา อุปาทานา จ สมฺปยุตฺตกา จ ขนฺธา ปจฺฉิมานํ ปจฺฉิมานํ อุปาทานานํ อนนฺตร\-ปจฺจเยน ปจฺจโย. (มูลํ กาตพฺพํ.) ปุริมา ปุริมา อุปาทานา จ สมฺปยุตฺตกา จ ขนฺธา ปจฺฉิมานํ ปจฺฉิมานํ โนอุปาทานานํ ขนฺธานํ อนนฺตร\-ปจฺจเยน ปจฺจโย. อุปาทานา จ สมฺปยุตฺตกา จ ขนฺธา วุฏฺฐานสฺส อนนฺตร\-ปจฺจเยน ปจฺจโย. (มูลํ กาตพฺพํ.) ปุริมา ปุริมา อุปาทานา จ สมฺปยุตฺตกา จ ขนฺธา ปจฺฉิมานํ ปจฺฉิมานํ อุปาทานานํ สมฺปยุตฺตกานญฺจ ขนฺธานํ อนนฺตร\-ปจฺจเยน ปจฺจโย. (3)}

\prose{สมนนฺตร\-ปจฺจเยน ปจฺจโย. สหชาต\-ปจฺจเยน ปจฺจโย. (ปฏิจฺจสทิสํ.) อญฺญมญฺญ\-ปจฺจเยน ปจฺจโย. (ปฏิจฺจสทิสํ.) นิสฺสย\-ปจฺจเยน ปจฺจโย. (ปจฺจยสทิสํ.)}

\csromanheader{2}{\textbf{อุปนิสฺสย}}
\proseitem{26}{อุปาทาโน ธมฺโม อุปาทานสฺส ธมฺมสฺส อุปนิสฺสย\-ปจฺจเยน ปจฺจโย. อารมฺมณู\-ปนิสฺสโย, อนนฺตรู\-ปนิสฺสโย, ปกตู\-ปนิสฺสโย ฯเปฯ. ปกตู\-ปนิสฺสโย\csromandash{}อุปาทานา อุปาทานานํ อุปนิสฺสย\-ปจฺจเยน ปจฺจโย. ตีณิ.}

\proseitem{27}{\csromanfolio{570}โนอุปาทาโน ธมฺโม โนอุปาทานสฺส ธมฺมสฺส อุปนิสฺสย\-ปจฺจเยน ปจฺจโย. อารมฺมณู\-ปนิสฺสโย, อนนฺตรู\-ปนิสฺสโย, ปกตู\-ปนิสฺสโย ฯเปฯ. ปกตู\-ปนิสฺสโย\csromandash{}สทฺธํ อุปนิสฺสาย ทานํ เทติ ฯเปฯ สมาปตฺติํ อุปฺปาเทติ. มานํ ชปฺเปติ, ทิฏฺฐิํ คณฺหาติ. สีลํ ฯเปฯ. ปญฺญํ. ราคํ ฯเปฯ. ปตฺถนํ. กายิกํ สุขํ ฯเปฯ. เสนาสนํ อุปนิสฺสาย ทานํ เทติ ฯเปฯ สมาปตฺติํ อุปฺปาเทติ. มานํ ชปฺเปติ. ทิฏฺฐิํ คณฺหาติ. ปาณํ หนติ ฯเปฯ สํฆํ ภินฺทติ. สทฺธา ฯเปฯ. เสนาสนํ สทฺธาย ฯเปฯ ผลสมาปตฺติยา อุปนิสฺสย\-ปจฺจเยน ปจฺจโย. (1)}

\prose{โนอุปาทาโน ธมฺโม อุปาทานสฺส ธมฺมสฺส อุปนิสฺสย\-ปจฺจเยน ปจฺจโย. ปกตู\-ปนิสฺสโย\csromandash{}สทฺธํ อุปนิสฺสาย มานํ ชปฺเปติ, ทิฏฺฐิํ คณฺหาติ. สีลํ ฯเปฯ. เสนาสนํ อุปนิสฺสาย อทินฺนํ ฯเปฯ มุสา ฯเปฯ ปิสุณํ ฯเปฯ สมฺผํ ฯเปฯ สนฺธิํ ฯเปฯ นิลฺโลปํ ฯเปฯ เอกาคาริกํ ฯเปฯ ปริปนฺเถ ฯเปฯ ปรทารํ ฯเปฯ คามฆาตํ ฯเปฯ นิคมฆาตํ กโรติ. สทฺธา ฯเปฯ. เสนาสนํ อุปาทานานํ อุปนิสฺสย\-ปจฺจเยน ปจฺจโย. (มูลํ กาตพฺพํ.) สทฺธํ อุปนิสฺสาย มานํ ชปฺเปติ, ทิฏฺฐิํ คณฺหาติ. สีลํ ฯเปฯ. เสนาสนํ อุปนิสฺสาย อทินฺนํ อาทิยติ, มุสา ฯเปฯ ปิสุณํ ฯเปฯ สมฺผํ ฯเปฯ สนฺธิํ ฯเปฯ นิลฺโลปํ ฯเปฯ เอกาคาริกํ ฯเปฯ ปริปนฺเถ ฯเปฯ ปรทารํ ฯเปฯ คามฆาตํ ฯเปฯ นิคมฆาตํ กโรติ. สทฺธา ฯเปฯ. เสนาสนํ อุปาทานานํ สมฺปยุตฺตกานญฺจ ขนฺธานํ อุปนิสฺสย\-ปจฺจเยน ปจฺจโย. (3)}

\prose{อุปาทาโน จ โนอุปาทาโน จ ธมฺมา อุปาทานสฺส ธมฺมสฺส อุปนิสฺสย\-ปจฺจเยน ปจฺจโย. อารมฺมณู\-ปนิสฺสโย, อนนฺตรู\-ปนิสฺสโย, ปกตู\-ปนิสฺสโย ฯเปฯ. ปกตู\-ปนิสฺสโย\csromandash{}ตีณิ.}

\proseitem{28}{โนอุปาทาโน ธมฺโม โนอุปาทานสฺส ธมฺมสฺส ปุเรชาต\-ปจฺจเยน ปจฺจโย. อารมฺมณ\-ปุเรชาตํ, วตฺถุ\-ปุเรชาตํ.  อารมฺมณ\-ปุเรชาตํ\csromandash{} จกฺขุํ ฯเปฯ วตฺถุํ อนิจฺจโต ฯเปฯ โทมนสฺสํ อุปฺปชฺชติ. ทิพฺเพน จกฺขุนา รูปํ ปสฺสติ. ทิพฺพาย โสตธาตุยา สทฺทํ สุณาติ. รูปายตนํ จกฺขุ\-วิญฺญาณสฺส ฯเปฯ. โผฏฺฐพฺพา\-ยตนํ กายวิญฺญาณสฺส ฯเปฯ. วตฺถุ\-ปุเรชาตํ\csromandash{}จกฺขายตนํ จกฺขุ\-วิญฺญาณสฺส ฯเปฯ กายายตนํ กายวิญฺญาณสฺส ฯเปฯ. วตฺถุ โนอุปาทานานํ ขนฺธานํ ปุเรชาต\-ปจฺจเยน ปจฺจโย. (1)}

\prose{\csromanfolio{571}โนอุปาทาโน ธมฺโม อุปาทานสฺส ธมฺมสฺส ปุเรชาต\-ปจฺจเยน ปจฺจโย. อารมฺมณ\-ปุเรชาตํ, วตฺถุ\-ปุเรชาตํ.  อารมฺมณ\-ปุเรชาตํ\csromandash{} จกฺขุํ ฯเปฯ วตฺถุํ อสฺสาเทติ อภินนฺทติ, ตํ อารพฺภ ราโค อุปฺปชฺชติ. ทิฏฺฐิ อุปฺปชฺชติ. วตฺถุ\-ปุเรชาตํ\csromandash{}วตฺถุ อุปาทานานํ ปุเรชาต\-ปจฺจเยน ปจฺจโย. (2)}

\prose{โนอุปาทาโน ธมฺโม อุปาทานสฺส จ โนอุปาทานสฺส จ ธมฺมสฺส ปุเรชาต\-ปจฺจเยน ปจฺจโย. อารมฺมณ\-ปุเรชาตํ, วตฺถุ\-ปุเรชาตํ.  อารมฺมณ\-ปุเรชาตํ\csromandash{}จกฺขุํ ฯเปฯ วตฺถุํ อสฺสาเทติ อภินนฺทติ, ตํ อารพฺภ อุปาทานา จ สมฺปยุตฺตกา จ ขนฺธา อุปฺปชฺชนฺติ. วตฺถุ\-ปุเรชาตํ\csromandash{}วตฺถุ อุปาทานานํ สมฺปยุตฺตกานญฺจ ขนฺธานํ ปุเรชาต\-ปจฺจเยน ปจฺจโย. (3)}

\proseitem{29}{อุปาทาโน ธมฺโม โนอุปาทานสฺส ธมฺมสฺส ปจฺฉาชาต\-ปจฺจเยน ปจฺจโย. (สํขิตฺตํ.) (1)}

\prose{โนอุปาทาโน ธมฺโม โนอุปาทานสฺส ธมฺมสฺส ปจฺฉาชาต\-ปจฺจเยน ปจฺจโย. (สํขิตฺตํ.) (1)}

\prose{อุปาทาโน จ โนอุปาทาโน จ ธมฺมา โนอุปาทานสฺส ธมฺมสฺส ปจฺฉาชาต\-ปจฺจเยน ปจฺจโย. (สํขิตฺตํ.) (1) อาเสวน\-ปจฺจเยน ปจฺจโย.}

\proseitem{30}{โนอุปาทาโน ธมฺโม โนอุปาทานสฺส ธมฺมสฺส กมฺมปจฺจเยน ปจฺจโย. สหชาตา, นานากฺขณิกา.  สหชาตา โนอุปาทานา เจตนา สมฺปยุตฺตกานํ ขนฺธานํ จิตฺตสมุฏฺฐานานญฺจ รูปานํ กมฺมปจฺจเยน ปจฺจโย. ปฏิสนฺธิ\-กฺขเณ ฯเปฯ. นานากฺขณิกา โนอุปาทานา เจตนา วิปากานํ ขนฺธานํ กฏตฺตา จ รูปานํ กมฺมปจฺจเยน ปจฺจโย. (1)}

\prose{โนอุปาทาโน ธมฺโม อุปาทานสฺส ธมฺมสฺส กมฺมปจฺจเยน ปจฺจโย. โนอุปาทานา เจตนา สมฺปยุตฺตกานํ อุปาทานานํ กมฺมปจฺจเยน ปจฺจโย. (2)}

\prose{โนอุปาทาโน ธมฺโม อุปาทานสฺส จ โนอุปาทานสฺส จ ธมฺมสฺส กมฺมปจฺจเยน ปจฺจโย. โนอุปาทานา เจตนา สมฺปยุตฺตกานํ ขนฺธานํ อุปาทานานญฺจ จิตฺต\-สมุฏฺฐานานํ รูปานํ กมฺมปจฺจเยน ปจฺจโย. (3)}

\csromanheader{2}{\textbf{วิปากาทิ}}
\proseitem{31}{\csromanfolio{572}โนอุปาทาโน ธมฺโม โนอุปาทานสฺส ธมฺมสฺส วิปาก\-ปจฺจเยน ปจฺจโย. วิปาโก โนอุปาทาโน เอโก ขนฺโธ ฯเปฯ. เอกํ.}

\prose{โนอุปาทาโน ธมฺโม โนอุปาทานสฺส ธมฺมสฺส อาหารปจฺจเยน ปจฺจโย. ตีณิ. (เอโกเยว กพฬีกาโร อาหาโร.) อินฺทฺริย\-ปจฺจเยน ปจฺจโย. ตีณิ. (รูปชีวิตินฺทฺริย เอกํเยว.) ฌานปจฺจเยน ปจฺจโย. ตีณิ.}

\prose{อุปาทาโน ธมฺโม อุปาทานสฺส ธมฺมสฺส มคฺคปจฺจเยน ปจฺจโย. อุปาทานานิ มคฺคงฺคานิ สมฺปยุตฺตกานํ อุปาทานานํ มคฺคปจฺจเยน ปจฺจโย. (อิมินา การเณน นว ปญฺหา กาตพฺพา.) สมฺปยุตฺต\-ปจฺจเยน ปจฺจโย. นว.}

\proseitem{32}{อุปาทาโน ธมฺโม โนอุปาทานสฺส ธมฺมสฺส วิปฺปยุตฺต\-ปจฺจเยน ปจฺจโย. สหชาตํ, ปจฺฉาชาตํ.  สหชาตา อุปาทานา จิตฺต\-สมุฏฺฐานานํ รูปานํ วิปฺปยุตฺต\-ปจฺจเยน ปจฺจโย. ปจฺฉาชาตํ\csromandash{}ปจฺฉาชาตา อุปาทานา ปุเรชาตสฺส อิมสฺส กายสฺส วิปฺปยุตฺต\-ปจฺจเยน ปจฺจโย. (1)}

\prose{โนอุปาทาโน ธมฺโม โนอุปาทานสฺส ธมฺมสฺส วิปฺปยุตฺต\-ปจฺจเยน ปจฺจโย. สหชาตํ, ปุเรชาตํ, ปจฺฉาชาตํ. (สํขิตฺตํ.) (1)}

\prose{โนอุปาทาโน ธมฺโม อุปาทานสฺส ธมฺมสฺส วิปฺปยุตฺต\-ปจฺจเยน ปจฺจโย. ปุเรชาตํ วตฺถุ อุปาทานานํ วิปฺปยุตฺต\-ปจฺจเยน ปจฺจโย. (2)}

\prose{โนอุปาทาโน ธมฺโม อุปาทานสฺส จ โนอุปาทานสฺส จ ธมฺมสฺส วิปฺปยุตฺต\-ปจฺจเยน ปจฺจโย. ปุเรชาตํ วตฺถุ อุปาทานานํ สมฺปยุตฺตกานญฺจ ขนฺธานํ วิปฺปยุตฺต\-ปจฺจเยน ปจฺจโย. (3)}

\proseitem{33}{อุปาทาโน จ โนอุปาทาโน จ ธมฺมา โนอุปาทานสฺส ธมฺมสฺส วิปฺปยุตฺต\-ปจฺจเยน ปจฺจโย. สหชาตํ, ปจฺฉาชาตํ.  สหชาตา อุปาทานา จ สมฺปยุตฺตกา จ ขนฺธา จิตฺต\-สมุฏฺฐานานํ รูปานํ วิปฺปยุตฺต\-ปจฺจเยน ปจฺจโย. ปจฺฉาชาตา อุปาทานา จ สมฺปยุตฺตกา จ ขนฺธา ปุเรชาตสฺส อิมสฺส กายสฺส วิปฺปยุตฺต\-ปจฺจเยน ปจฺจโย. (1)}

\csromanheader{2}{69. อุปาทานทุก \textbf{อตฺถิ}}
\proseitem{34}{\csromanfolio{573}อุปาทาโน ธมฺโม อุปาทานสฺส ธมฺมสฺส อตฺถิปจฺจเยน ปจฺจโย. ทิฏฺฐุปาทานํ กามุปาทานสฺส อตฺถิปจฺจเยน ปจฺจโย. (จกฺกํ.) (1)}

\prose{อุปาทาโน ธมฺโม โนอุปาทานสฺส ธมฺมสฺส อตฺถิปจฺจเยน ปจฺจโย. \textbf{สหชาตํ}, ปจฺฉาชาตํ.  \textbf{สหชาตา} อุปาทานา สมฺปยุตฺตกานํ ขนฺธานํ จิตฺตสมุฏฺฐานานญฺจ รูปานํ อตฺถิปจฺจเยน ปจฺจโย. \textbf{ปจฺฉาชาตา} อุปาทานา ปุเรชาตสฺส อิมสฺส กายสฺส อตฺถิปจฺจเยน ปจฺจโย. (2)}

\prose{อุปาทาโน ธมฺโม อุปาทานสฺส จ โนอุปาทานสฺส จ ธมฺมสฺส อตฺถิปจฺจเยน ปจฺจโย. (สํขิตฺตํ. ปฏิจฺจสทิสํ.) (3)}

\proseitem{35}{โนอุปาทาโน ธมฺโม โนอุปาทานสฺส ธมฺมสฺส อตฺถิปจฺจเยน ปจฺจโย. \textbf{สหชาตํ}, ปุเรชาตํ, ปจฺฉาชาตํ, อาหารํ, อินฺทฺริยํ. (สํขิตฺตํ. วิตฺถาเรตพฺพํ.) (1)}

\prose{โนอุปาทาโน ธมฺโม อุปาทานสฺส ธมฺมสฺส อตฺถิปจฺจเยน ปจฺจโย. สหชาตํ, ปุเรชาตํ.  สหชาตํ\csromandash{}(สหชาต\-สทิสํ.) ปุเรชาตํ\csromandash{} (ปุเรชาต\-สทิสํ.) (2)}

\prose{โนอุปาทาโน ธมฺโม อุปาทานสฺส จ โนอุปาทานสฺส จ ธมฺมสฺส อตฺถิปจฺจเยน ปจฺจโย. สหชาตํ, ปุเรชาตํ. (สํขิตฺตํ. สหชาต\-สทิสํ. สหชาตํ วิภชิตพฺพํ. ปุเรชาต\-สทิสํ ปุเรชาตํ.) (3)}

\proseitem{36}{อุปาทาโน จ โนอุปาทาโน จ ธมฺมา อุปาทานสฺส ธมฺมสฺส อตฺถิปจฺจเยน ปจฺจโย. สหชาตํ, ปุเรชาตํ.  สหชาตํ ทิฏฺฐุปาทานญฺจ สมฺปยุตฺตกา จ ขนฺธา กามุปาทานสฺส อตฺถิปจฺจเยน ปจฺจโย. (จกฺกํ.) สหชาตํ ทิฏฺฐุปาทานญฺจ วตฺถุ จ กามุปาทานสฺส อตฺถิปจฺจเยน ปจฺจโย. (จกฺกํ.) (1)}

\prose{อุปาทาโน จ โนอุปาทาโน จ ธมฺมา โนอุปาทานสฺส ธมฺมสฺส อตฺถิปจฺจเยน ปจฺจโย. สหชาตํ, ปุเรชาตํ, ปจฺฉาชาตํ, อาหารํ, อินฺทฺริยํ.  สหชาโต โนอุปาทาโน เอโก ขนฺโธ จ อุปาทาโน จ ติณฺณนฺนํ ขนฺธานํ จิตฺตสมุฏฺฐานานญฺจ รูปานํ อตฺถิปจฺจเยน ปจฺจโย ฯเปฯ เทฺว ขนฺธา จ ฯเปฯ. สหชาตา อุปาทานา จ มหาภูตา จ จิตฺต\-สมุฏฺฐานานํ \csromanfolio{574}รูปานํ อตฺถิปจฺจเยน ปจฺจโย. \textbf{สหชาตา} อุปาทานา จ วตฺถุ จ โนอุปาทานานํ ขนฺธานํ อตฺถิปจฺจเยน ปจฺจโย. ปจฺฉาชาตา อุปาทานา จ สมฺปยุตฺตกา จ ขนฺธา ปุเรชาตสฺส อิมสฺส กายสฺส อตฺถิปจฺจเยน ปจฺจโย. ปจฺฉาชาตา อุปาทานา จ สมฺปยุตฺตกา จ ขนฺธา กพฬีกาโร อาหาโร อิมสฺส กายสฺส อตฺถิปจฺจเยน ปจฺจโย. ปจฺฉาชาตา อุปาทานา จ สมฺปยุตฺตกา จ ขนฺธา รูปชีวิตินฺทฺริยญฺจ กฏตฺตารูปานํ อตฺถิปจฺจเยน ปจฺจโย. (2)}

\prose{อุปาทาโน จ โนอุปาทาโน จ ธมฺมา อุปาทานสฺส จ โนอุปาทานสฺส จ ธมฺมสฺส อตฺถิปจฺจเยน ปจฺจโย. สหชาตํ, ปุเรชาตํ.  สหชาโต โนอุปาทาโน เอโก ขนฺโธ จ ทิฏฺฐุปาทานญฺจ ติณฺณนฺนํ ขนฺธานํ กามุปาทานสฺส จ จิตฺต\-สมุฏฺฐานานํ รูปานํ อตฺถิปจฺจเยน ปจฺจโย ฯเปฯ เทฺว ขนฺธา ฯเปฯ. (จกฺกํ.) สหชาตํ ทิฏฺฐุปาทานญฺจ วตฺถุ จ กามุปาทานสฺส สมฺปยุตฺตกานญฺจ ขนฺธานํ อตฺถิปจฺจเยน ปจฺจโย. (จกฺกํ.) (3)}

\csromanmarkh{1. ปจฺจยานุโลม}\csromanmarkhii{2. สงฺขฺยาวาร}\csromanheadernomark{2}{1. \textbf{ปจฺจยานุโลม 2. สงฺขฺยาวาร}}
\proseitem{37}{เหตุยา นว, อารมฺมเณ นว, อธิปติยา นว, อนนฺตเร นว, สมนนฺตเร นว, สหชาเต นว, อญฺญมญฺเญ นว, นิสฺสเย นว, อุปนิสฺสเย นว, ปุเรชาเต ตีณิ, ปจฺฉาชาเต ตีณิ, อาเสวเน นว, กมฺเม ตีณิ, วิปาเก เอกํ, อาหาเร ตีณิ, อินฺทฺริเย ตีณิ, ฌาเน ตีณิ, มคฺเค นว, สมฺปยุตฺเต นว, วิปฺปยุตฺเต ปญฺจ, อตฺถิยา นว, นตฺถิยา นว, วิคเต นว, อวิคเต นว.}

\vspace{\tipitakaparskip}
\csromancenter{อนุโลมํ.}
\proseitem{38}{อุปาทาโน ธมฺโม อุปาทานสฺส ธมฺมสฺส อารมฺมณ\-ปจฺจเยน ปจฺจโย. สหชาต\-ปจฺจเยน ปจฺจโย. อุปนิสฺสย\-ปจฺจเยน ปจฺจโย. (1)}

\prose{\csromanfolio{575}อุปาทาโน ธมฺโม โนอุปาทานสฺส ธมฺมสฺส อารมฺมณ\-ปจฺจเยน ปจฺจโย. สหชาต\-ปจฺจเยน ปจฺจโย. อุปนิสฺสย\-ปจฺจเยน ปจฺจโย. ปจฺฉาชาต\-ปจฺจเยน ปจฺจโย. (2)}

\prose{อุปาทาโน ธมฺโม อุปาทานสฺส จ โนอุปาทานสฺส จ ธมฺมสฺส อารมฺมณ\-ปจฺจเยน ปจฺจโย. สหชาต\-ปจฺจเยน ปจฺจโย. อุปนิสฺสย\-ปจฺจเยน ปจฺจโย. (3)}

\proseitem{39}{โนอุปาทาโน ธมฺโม โนอุปาทานสฺส ธมฺมสฺส อารมฺมณ\-ปจฺจเยน ปจฺจโย. สหชาต\-ปจฺจเยน ปจฺจโย. อุปนิสฺสย\-ปจฺจเยน ปจฺจโย. ปุเรชาต\-ปจฺจเยน ปจฺจโย. ปจฺฉาชาต\-ปจฺจเยน ปจฺจโย. กมฺมปจฺจเยน ปจฺจโย. อาหารปจฺจเยน ปจฺจโย. อินฺทฺริย\-ปจฺจเยน ปจฺจโย. (1)}

\prose{โนอุปาทาโน ธมฺโม อุปาทานสฺส ธมฺมสฺส อารมฺมณ\-ปจฺจเยน ปจฺจโย. สหชาต\-ปจฺจเยน ปจฺจโย. อุปนิสฺสย\-ปจฺจเยน ปจฺจโย. ปุเรชาต\-ปจฺจเยน ปจฺจโย. (2)}

\prose{โนอุปาทาโน ธมฺโม อุปาทานสฺส จ โนอุปาทานสฺส จ ธมฺมสฺส อารมฺมณ\-ปจฺจเยน ปจฺจโย. สหชาต\-ปจฺจเยน ปจฺจโย. อุปนิสฺสย\-ปจฺจเยน ปจฺจโย. ปุเรชาต\-ปจฺจเยน ปจฺจโย. (3)}

\proseitem{40}{อุปาทาโน จ โนอุปาทาโน จ ธมฺมา อุปาทานสฺส ธมฺมสฺส อารมฺมณ\-ปจฺจเยน ปจฺจโย. สหชาต\-ปจฺจเยน ปจฺจโย. อุปนิสฺสย\-ปจฺจเยน ปจฺจโย. (1)}

\prose{อุปาทาโน จ โนอุปาทาโน จ ธมฺมา โนอุปาทานสฺส ธมฺมสฺส อารมฺมณ\-ปจฺจเยน ปจฺจโย. สหชาต\-ปจฺจเยน ปจฺจโย. อุปนิสฺสย\-ปจฺจเยน ปจฺจโย. ปจฺฉาชาต\-ปจฺจเยน ปจฺจโย. (2)}

\prose{อุปาทาโน จ โนอุปาทาโน จ ธมฺมา อุปาทานสฺส จ โนอุปาทานสฺส จ ธมฺมสฺส อารมฺมณ\-ปจฺจเยน ปจฺจโย. สหชาต\-ปจฺจเยน ปจฺจโย. อุปนิสฺสย\-ปจฺจเยน ปจฺจโย. (3)}

\csromanmarkh{2. ปจฺจย\-ปจฺจนีย}\csromanmarkhii{2. สงฺขฺยาวาร}\csromanheadernomark{2}{2. ปจฺจย\-ปจฺจนีย 2. สงฺขฺยาวาร}
\vspace{\tipitakaparskip}
\csromancenter{นเหตุยา นว, นอารมฺมเณ นว, (สพฺพตฺถ นว.) โนอวิคเต นว.}
\csromancenter{\textbf{เหตุ}ปจฺจยา นอารมฺมเณ นว, นอธิปติยา นว ฯเปฯ นอญฺญมญฺเญ ตีณิ, นฺเอาปนิสฺสเย นว, (สพฺพตฺถ นว.) นสมฺปยุตฺเต ตีณิ, นวิปฺปยุตฺเต นว, โนนตฺถิยา นว, โนวิคเต นว.}
\proseitemwithcloser{43}{\csromanfolio{576}นเหตุปจฺจยา อารมฺมเณ นว, อธิปติยา นว, (อนุโลมมาติกา กาตพฺพา.) ฯเปฯ อวิคเต นว.}{อุปาทานทุกํ นิฏฺฐิตํ.}
\csromanmatikaanchor{mtk.70}
\csromantocmarkref{subsection}{70. อุปาทานิยทุก}{mtk.70}
\bookmark[dest={mtk.70},level=2]{70. อุปาทานิยทุก}
\csromanmarkh{70. อุปาทานิยทุก}\csromanmarkhii{1. ปฏิจฺจวาร}\csromanheadernomark{2}{70. \textbf{อุปาทานิยทุก 1. ปฏิจฺจวาร}}
\proseitemwithcloser{44}{อุปาทานิยํ ธมฺมํ ปฏิจฺจ อุปาทานิโย ธมฺโม อุปฺปชฺชติ เหตุปจฺจยา. อุปาทานิยํ เอกํ ขนฺธํ ปฏิจฺจ ตโย ขนฺธา จิตฺต\-สมุฏฺฐานญฺจ รูปํ ฯเปฯ เทฺว ขนฺเธ ฯเปฯ. ปฏิสนฺธิ\-กฺขเณ ฯเปฯ. ขนฺเธ ปฏิจฺจ วตฺถุ, วตฺถุํ ปฏิจฺจ ขนฺธา. เอกํ มหาภูตํ ฯเปฯ. (ยถา โลกิยทุกํ, เอวํ กาตพฺพํ. นินฺนานากรณํ.)}{อุปาทานิยทุกํ นิฏฺฐิตํ.}
\csromanmatikaanchor{mtk.71}
\csromantocmarkref{subsection}{71. อุปาทานสมฺปยุตฺตทุก}{mtk.71}
\bookmark[dest={mtk.71},level=2]{71. อุปาทานสมฺปยุตฺตทุก}
\csromanmarkh{71. อุปาทาน\-สมฺปยุตฺตทุก}\csromanmarkhii{1. ปฏิจฺจวาร}\csromanheadernomark{2}{71. อุปาทาน\-สมฺปยุตฺตทุก 1. ปฏิจฺจวาร}
\csromanmarkh{1. ปจฺจยานุโลม}\csromanmarkhii{1. วิภงฺควาร}\csromanheadernomark{2}{1. \textbf{ปจฺจยานุโลม 1. วิภงฺควาร}}
\vspace{\tipitakaparskip}
\csromancenter{\textbf{เหตุ}}
\proseitem{45}{อุปาทาน\-สมฺปยุตฺตํ ธมฺมํ ปฏิจฺจ อุปาทาน\-สมฺปยุตฺโต ธมฺโม อุปฺปชฺชติ เหตุปจฺจยา. อุปาทาน\-สมฺปยุตฺตํ เอกํ ขนฺธํ ปฏิจฺจ ตโย ขนฺธา ฯเปฯ เทฺว ขนฺเธ ฯเปฯ. (1)}

\prose{อุปาทาน\-สมฺปยุตฺตํ ธมฺมํ ปฏิจฺจ อุปาทาน\-วิปฺปยุตฺโต ธมฺโม อุปฺปชฺชติ เหตุปจฺจยา. อุปาทาน\-สมฺปยุตฺเต ขนฺเธ ปฏิจฺจ จิตฺต\-สมุฏฺฐานํ รูปํ. ทิฏฺฐิคต\-วิปฺปยุตฺต\-โลภ\-สหคเต ขนฺเธ ปฏิจฺจ โลโภ จิตฺต\-สมุฏฺฐานญฺจ รูปํ. (2)}

\prose{\csromanfolio{577}อุปาทาน\-สมฺปยุตฺตํ ธมฺมํ ปฏิจฺจ อุปาทาน\-สมฺปยุตฺโต จ อุปาทาน\-วิปฺปยุตฺโต จ ธมฺมา อุปฺปชฺชนฺติ เหตุปจฺจยา. อุปาทาน\-สมฺปยุตฺตํ เอกํ ขนฺธํ ปฏิจฺจ ตโย ขนฺธา จิตฺต\-สมุฏฺฐานญฺจ รูปํ ฯเปฯ เทฺว ขนฺเธ ฯเปฯ. ทิฏฺฐิคต\-วิปฺปยุตฺต\-โลภ\-สหคตํ เอกํ ขนฺธํ ปฏิจฺจ ตโย ขนฺธา โลโภ จ จิตฺต\-สมุฏฺฐานญฺจ รูปํ ฯเปฯ เทฺว ขนฺเธ ฯเปฯ. (3)}

\proseitem{46}{อุปาทาน\-วิปฺปยุตฺตํ ธมฺมํ ปฏิจฺจ อุปาทาน\-วิปฺปยุตฺโต ธมฺโม อุปฺปชฺชติ เหตุปจฺจยา. อุปาทาน\-วิปฺปยุตฺตํ เอกํ ขนฺธํ ปฏิจฺจ ตโย ขนฺธา จิตฺต\-สมุฏฺฐานญฺจ รูปํ ฯเปฯ เทฺว ขนฺเธ ฯเปฯ. ทิฏฺฐิคต\-วิปฺปยุตฺตํ โลภํ ปฏิจฺจ จิตฺต\-สมุฏฺฐานํ รูปํ. ปฏิสนฺธิ\-กฺขเณ อุปาทาน\-วิปฺปยุตฺตํ เอกํ ขนฺธํ ปฏิจฺจ ตโย ขนฺธา กฏตฺตา จ รูปํ ฯเปฯ เทฺว ขนฺเธ ฯเปฯ. ขนฺเธ ปฏิจฺจ วตฺถุ, วตฺถุํ ปฏิจฺจ ขนฺธา. เอกํ มหาภูตํ ฯเปฯ. (1)}

\prose{อุปาทาน\-วิปฺปยุตฺตํ ธมฺมํ ปฏิจฺจ อุปาทาน\-สมฺปยุตฺโต ธมฺโม อุปฺปชฺชติ เหตุปจฺจยา. ทิฏฺฐิคต\-วิปฺปยุตฺตํ โลภํ ปฏิจฺจ สมฺปยุตฺตกา ขนฺธา. (2)}

\prose{อุปาทาน\-วิปฺปยุตฺตํ ธมฺมํ ปฏิจฺจ อุปาทาน\-สมฺปยุตฺโต จ อุปาทาน\-วิปฺปยุตฺโต จ ธมฺมา อุปฺปชฺชนฺติ เหตุปจฺจยา. ทิฏฺฐิคต\-วิปฺปยุตฺตํ โลภํ ปฏิจฺจ สมฺปยุตฺตกา ขนฺธา จิตฺต\-สมุฏฺฐานญฺจ รูปํ. (3)}

\proseitem{47}{อุปาทาน\-สมฺปยุตฺตญฺจ อุปาทาน\-วิปฺปยุตฺตญฺจ ธมฺมํ ปฏิจฺจ อุปาทาน\-สมฺปยุตฺโต ธมฺโม อุปฺปชฺชติ เหตุปจฺจยา. ทิฏฺฐิคต\-วิปฺปยุตฺต\-โลภ\-สหคตํ เอกํ ขนฺธญฺจ โลภญฺจ ปฏิจฺจ ตโย ขนฺธา ฯเปฯ เทฺว ขนฺเธ จ ฯเปฯ. (1)}

\prose{อุปาทาน\-สมฺปยุตฺตญฺจ อุปาทาน\-วิปฺปยุตฺตญฺจ ธมฺมํ ปฏิจฺจ อุปาทาน\-วิปฺปยุตฺโต ธมฺโม อุปฺปชฺชติ เหตุปจฺจยา. อุปาทาน\-สมฺปยุตฺเต ขนฺเธ จ มหาภูเต จ ปฏิจฺจ จิตฺต\-สมุฏฺฐานํ รูปํ. ทิฏฺฐิคต\-วิปฺปยุตฺต\-โลภ\-สหคเต ขนฺเธ จ โลภญฺจ ปฏิจฺจ จิตฺต\-สมุฏฺฐานํ รูปํ. (2)}

\prose{อุปาทาน\-สมฺปยุตฺตญฺจ อุปาทาน\-วิปฺปยุตฺตญฺจ ธมฺมํ ปฏิจฺจ อุปาทาน\-สมฺปยุตฺโต จ อุปาทาน\-วิปฺปยุตฺโต จ ธมฺมา อุปฺปชฺชนฺติ เหตุปจฺจยา. ทิฏฺฐิคต\-วิปฺปยุตฺต\-โลภ\-สหคตํ เอกํ ขนฺธญฺจ โลภญฺจ ปฏิจฺจ ตโย ขนฺธา จิตฺต\-สมุฏฺฐานญฺจ รูปํ ฯเปฯ เทฺว ขนฺเธ จ ฯเปฯ. (3)}

\proseitem{48}{\csromanfolio{578}อุปาทาน\-สมฺปยุตฺตํ ธมฺมํ ปฏิจฺจ อุปาทาน\-สมฺปยุตฺโต ธมฺโม อุปฺปชฺชติ อารมฺมณ\-ปจฺจยา. อุปาทาน\-สมฺปยุตฺตํ เอกํ ขนฺธํ ปฏิจฺจ ตโย ขนฺธา ฯเปฯ เทฺว ขนฺเธ ฯเปฯ. (1)}

\prose{อุปาทาน\-สมฺปยุตฺตํ ธมฺมํ ปฏิจฺจ อุปาทาน\-วิปฺปยุตฺโต ธมฺโม อุปฺปชฺชติ อารมฺมณ\-ปจฺจยา. ทิฏฺฐิคต\-วิปฺปยุตฺต\-โลภ\-สหคเต ขนฺเธ ปฏิจฺจ โลโภ. (2)}

\prose{อุปาทาน\-สมฺปยุตฺตํ ธมฺมํ ปฏิจฺจ อุปาทาน\-สมฺปยุตฺโต จ อุปาทาน\-วิปฺปยุตฺโต จ ธมฺมา อุปฺปชฺชนฺติ อารมฺมณ\-ปจฺจยา. ทิฏฺฐิคต\-วิปฺปยุตฺต\-โลภ\-สหคตํ เอกํ ขนฺธํ ปฏิจฺจ ตโย ขนฺธา โลโภ จ ฯเปฯ เทฺว ขนฺเธ ฯเปฯ. (3)}

\proseitem{49}{อุปาทาน\-วิปฺปยุตฺตํ ธมฺมํ ปฏิจฺจ อุปาทาน\-วิปฺปยุตฺโต ธมฺโม อุปฺปชฺชติ อารมฺมณ\-ปจฺจยา. อุปาทาน\-วิปฺปยุตฺตํ เอกํ ขนฺธํ ฯเปฯ เทฺว ขนฺเธ ฯเปฯ. ปฏิสนฺธิ\-กฺขเณ ฯเปฯ. วตฺถุํ ปฏิจฺจ ขนฺธา. (1)}

\prose{อุปาทาน\-วิปฺปยุตฺตํ ธมฺมํ ปฏิจฺจ อุปาทาน\-สมฺปยุตฺโต ธมฺโม อุปฺปชฺชติ อารมฺมณ\-ปจฺจยา. ทิฏฺฐิคต\-วิปฺปยุตฺตํ โลภํ ปฏิจฺจ สมฺปยุตฺตกา ขนฺธา. (2)}

\prose{อุปาทาน\-สมฺปยุตฺตญฺจ อุปาทาน\-วิปฺปยุตฺตญฺจ ธมฺมํ ปฏิจฺจ อุปาทาน\-สมฺปยุตฺโต ธมฺโม อุปฺปชฺชติ อารมฺมณ\-ปจฺจยา. ทิฏฺฐิคต\-วิปฺปยุตฺต\-โลภ\-สหคตํ เอกํ ขนฺธญฺจ โลภญฺจ ปฏิจฺจ ตโย ขนฺธา ฯเปฯ เทฺว ขนฺเธ จ ฯเปฯ. (3) (สํขิตฺตํ.)}

\csromanmarkh{1. ปจฺจยานุโลม}\csromanmarkhii{2. สงฺขฺยาวาร}\csromanheadernomark{2}{1. \textbf{ปจฺจยานุโลม 2. สงฺขฺยาวาร}}
\proseitem{50}{เหตุยา นว, อารมฺมเณ ฉ, อธิปติยา นว, อนนฺตเร ฉ, สมนนฺตเร ฉ, สหชาเต นว, อญฺญมญฺเญ ฉ, นิสฺสเย นว, อุปนิสฺสเย ฉ, ปุเรชาเต ฉ, อาเสวเน ฉ, กมฺเม นว, วิปาเก เอกํ, อาหาเร นว, (สพฺพตฺถ นว,) มคฺเค นว, สมฺปยุตฺเต ฉ, วิปฺปยุตฺเต นว, อตฺถิยา นว, นตฺถิยา ฉ, วิคเต ฉ, อวิคเต นว.}

\csromanmarkh{71. อุปาทาน\-สมฺปยุตฺตทุก}\csromanmarkhii{2. ปจฺจย\-ปจฺจนีย 1. วิภงฺควาร}\csromanheadernomark{2}{71. อุปาทาน\-สมฺปยุตฺตทุก 2. ปจฺจย\-ปจฺจนีย 1. วิภงฺควาร}
\proseitem{51}{\csromanfolio{579}อุปาทาน\-วิปฺปยุตฺตํ ธมฺมํ ปฏิจฺจ อุปาทาน\-วิปฺปยุตฺโต ธมฺโม อุปฺปชฺชติ นเหตุปจฺจยา. อเหตุกํ อุปาทาน\-วิปฺปยุตฺตํ เอกํ ขนฺธํ ปฏิจฺจ ตโย ขนฺธา จิตฺต\-สมุฏฺฐานญฺจ รูปํ ฯเปฯ เทฺว ขนฺเธ ฯเปฯ. อเหตุก\-ปฏิสนฺธิ\-กฺขเณ ฯเปฯ. (ยาว อสญฺญสตฺตา.) วิจิกิจฺฉา\-สหคเต อุทฺธจฺจ\-สหคเต ขนฺเธ ปฏิจฺจ วิจิกิจฺฉา\-สหคโต อุทฺธจฺจ\-สหคโต โมโห. (1)}

\prose{นอารมฺมณาทิ}

\proseitem{52}{อุปาทาน\-สมฺปยุตฺตํ ธมฺมํ ปฏิจฺจ อุปาทาน\-วิปฺปยุตฺโต ธมฺโม อุปฺปชฺชติ นอารมฺมณ\-ปจฺจยา. อุปาทาน\-สมฺปยุตฺเต ขนฺเธ ปฏิจฺจ จิตฺต\-สมุฏฺฐานํ รูปํ. (1)}

\prose{อุปาทาน\-วิปฺปยุตฺตํ ธมฺมํ ปฏิจฺจ อุปาทาน\-วิปฺปยุตฺโต ธมฺโม อุปฺปชฺชติ นอารมฺมณ\-ปจฺจยา. อุปาทาน\-วิปฺปยุตฺเต ขนฺเธ ปฏิจฺจ จิตฺต\-สมุฏฺฐานํ รูปํ. ทิฏฺฐิคต\-วิปฺปยุตฺตํ โลภํ ปฏิจฺจ จิตฺต\-สมุฏฺฐานํ รูปํ. ปฏิสนฺธิ\-กฺขเณ ฯเปฯ. (ยาว อสญฺญสตฺตา.) (1)}

\prose{อุปาทาน\-สมฺปยุตฺตญฺจ อุปาทาน\-วิปฺปยุตฺตญฺจ ธมฺมํ ปฏิจฺจ อุปาทาน\-วิปฺปยุตฺโต ธมฺโม อุปฺปชฺชติ นอารมฺมณ\-ปจฺจยา. อุปาทาน\-สมฺปยุตฺเต ขนฺเธ จ มหาภูเต จ ปฏิจฺจ จิตฺต\-สมุฏฺฐานํ รูปํ. ทิฏฺฐิคต\-วิปฺปยุตฺต\-โลภ\-สหคเต ขนฺเธ จ โลภญฺจ ปฏิจฺจ จิตฺต\-สมุฏฺฐานํ รูปํ. (1)}

\prose{นอธิปติ\-ปจฺจยา. นอนนฺตร\-ปจฺจยา. นสมนนฺตร\-ปจฺจยา. นอุปนิสฺสยปจฺจยา.}

\csromanheader{2}{\textbf{นปุเรชาตาทิ}}
\proseitem{53}{อุปาทาน\-สมฺปยุตฺตํ ธมฺมํ ปฏิจฺจ อุปาทาน\-สมฺปยุตฺโต ธมฺโม อุปฺปชฺชติ นปุเรชาต\-ปจฺจยา. อรูเป อุปาทาน\-สมฺปยุตฺตํ เอกํ ขนฺธํ ปฏิจฺจ ตโย ขนฺธา ฯเปฯ เทฺว ขนฺเธ ฯเปฯ. (1)}

\prose{อุปาทาน\-สมฺปยุตฺตํ ธมฺมํ ปฏิจฺจ อุปาทาน\-วิปฺปยุตฺโต ธมฺโม อุปฺปชฺชติ นปุเรชาต\-ปจฺจยา. อรูเป ทิฏฺฐิคต\-วิปฺปยุตฺต\-โลภ\-สหคเต ขนฺเธ ปฏิจฺจ โลโภ. อุปาทาน\-สมฺปยุตฺเต ขนฺเธ ปฏิจฺจ จิตฺต\-สมุฏฺฐานํ รูปํ. (2)}

\prose{\csromanfolio{580}อุปาทาน\-สมฺปยุตฺตํ ธมฺมํ ปฏิจฺจ อุปาทาน\-สมฺปยุตฺโต จ อุปาทาน\-วิปฺปยุตฺโต จ ธมฺมา อุปฺปชฺชนฺติ นปุเรชาต\-ปจฺจยา. อรูเป ทิฏฺฐิคต\-วิปฺปยุตฺต\-โลภ\-สหคตํ เอกํ ขนฺทํ ปฏิจฺจ ตโย ขนฺธา โลโภ จ ฯเปฯ เทฺว ขนฺเธ ฯเปฯ. (3)}

\proseitem{54}{อุปาทาน\-วิปฺปยุตฺตํ ธมฺมํ ปฏิจฺจ อุปาทาน\-วิปฺปยุตฺโต ธมฺโม อุปฺปชฺชติ นปุเรชาต\-ปจฺจยา. อรูเป อุปาทาน\-วิปฺปยุตฺตํ เอกํ ขนฺธํ ปฏิจฺจ ตโย ขนฺธา ฯเปฯ เทฺว ขนฺเธ ฯเปฯ. อุปาทาน\-วิปฺปยุตฺเต ขนฺเธ ปฏิจฺจ จิตฺต\-สมุฏฺฐานํ รูปํ. ทิฏฺฐิคต\-วิปฺปยุตฺตํ โลภํ ปฏิจฺจ จิตฺต\-สมุฏฺฐานํ รูปํ. ปฏิสนฺธิ\-กฺขเณ ฯเปฯ. (ยาว อสญฺญสตฺตา.) (1)}

\prose{อุปาทาน\-วิปฺปยุตฺตํ ธมฺมํ ปฏิจฺจ อุปาทาน\-สมฺปยุตฺโต ธมฺโม อุปฺปชฺชติ นปุเรชาต\-ปจฺจยา. อรูเป ทิฏฺฐิคต\-วิปฺปยุตฺตํ โลภํ ปฏิจฺจ สมฺปยุตฺตกา ขนฺธา. (2)}

\prose{อุปาทาน\-สมฺปยุตฺตญฺจ อุปาทาน\-วิปฺปยุตฺตญฺจ ธมฺมํ ปฏิจฺจ อุปาทาน\-สมฺปยุตฺโต ธมฺโม อุปฺปชฺชติ นปุเรชาต\-ปจฺจยา. อรูเป ทิฏฺฐิคต\-วิปฺปยุตฺต\-โลภ\-สหคตํ เอกํ ขนฺธญฺจ โลภญฺจ ปฏิจฺจ ตโย ขนฺธา ฯเปฯ เทฺว ขนฺเธ จ ฯเปฯ. (1)}

\prose{อุปาทาน\-สมฺปยุตฺตญฺจ อุปาทาน\-วิปฺปยุตฺตญฺจ ธมฺมํ ปฏิจฺจ อุปาทาน\-วิปฺปยุตฺโต ธมฺโม อุปฺปชฺชติ นปุเรชาต\-ปจฺจยา. อุปาทาน\-สมฺปยุตฺเต ขนฺเธ จ มหาภูเต จ ปฏิจฺจ จิตฺต\-สมุฏฺฐานํ รูปํ. ทิฏฺฐิคต\-วิปฺปยุตฺต\-โลภ\-สหคเต ขนฺเธ จ โลภญฺจ ปฏิจฺจ จิตฺต\-สมุฏฺฐานํ รูปํ. (2)}

\prose{นปจฺฉาชาต\-ปจฺจยา. นอาเสวน\-ปจฺจยา.}

\csromanheader{2}{\textbf{นกมฺม}}
\proseitem{55}{อุปาทาน\-สมฺปยุตฺตํ ธมฺมํ ปฏิจฺจ อุปาทาน\-สมฺปยุตฺโต ธมฺโม อุปฺปชฺชติ นกมฺมปจฺจยา. อุปาทาน\-สมฺปยุตฺเต ขนฺเธ ปฏิจฺจ สมฺปยุตฺตกา เจตนา. (1)}

\prose{อุปาทาน\-วิปฺปยุตฺตํ ธมฺมํ ปฏิจฺจ อุปาทาน\-วิปฺปยุตฺโต ธมฺโม อุปฺปชฺชติ นกมฺมปจฺจยา. อุปาทาน\-วิปฺปยุตฺเต ขนฺเธ ปฏิจฺจ สมฺปยุตฺตกา เจตนา. พาหิรํ. อาหารสมุฏฺฐานํ. อุตุสมุฏฺฐานํ ฯเปฯ. (1)}

\prose{อุปาทาน\-วิปฺปยุตฺตํ ธมฺมํ ปฏิจฺจ อุปาทาน\-สมฺปยุตฺโต ธมฺโม อุปฺปชฺชติ นกมฺมปจฺจยา. ทิฏฺฐิคต\-วิปฺปยุตฺตํ โลภํ ปฏิจฺจ สมฺปยุตฺตกา เจตนา. (2)}

\prose{\csromanfolio{581}อุปาทาน\-สมฺปยุตฺตญฺจ อุปาทาน\-วิปฺปยุตฺตญฺจ ธมฺมํ ปฏิจฺจ อุปาทาน\-สมฺปยุตฺโต ธมฺโม อุปฺปชฺชติ นกมฺมปจฺจยา. ทิฏฺฐิคต\-วิปฺปยุตฺต\-โลภ\-สหคเต ขนฺเธ จ โลภญฺจ ปฏิจฺจ สมฺปยุตฺตกา เจตนา. (1) (สํขิตฺตํ.)}

\csromanmarkh{2. ปจฺจย\-ปจฺจนีย}\csromanmarkhii{2. สงฺขฺยาวาร}\csromanheadernomark{2}{2. ปจฺจย\-ปจฺจนีย 2. สงฺขฺยาวาร}
\proseitem{56}{นเหตุยา เอกํ, นอารมฺมเณ ตีณิ, นอธิปติยา นว, นอนนฺตเร ตีณิ, นสมนนฺตเร ตีณิ, นอญฺญมญฺเญ ตีณิ, นอุปนิสฺสเย ตีณิ, นปุเรชาเต สตฺต, นปจฺฉาชาเต นว, นอาเสวเน นว, นกมฺเม จตฺตาริ, นวิปาเก นว, นอาหาเร เอกํ, นอินฺทฺริเย เอกํ, นฌาเน เอกํ, นมคฺเค เอกํ, นสมฺปยุตฺเต ตีณิ, นวิปฺปยุตฺเต ฉ, โนนตฺถิยา ตีณิ. โนวิคเต ตีณิ.}

\prose{(อิตเร เทฺว คณนาปิ กาตพฺพา. สหชาตวาโร ปฏิจฺจ\-วาร\-สทิโส.)}

\csromanmarkh{71. อุปาทาน\-สมฺปยุตฺตทุก}\csromanmarkhii{3. ปจฺจยวาร}\csromanheadernomark{2}{71. อุปาทาน\-สมฺปยุตฺตทุก 3. ปจฺจยวาร}
\csromanmarkh{1. ปจฺจยานุโลม}\csromanmarkhii{1. วิภงฺควาร}\csromanheadernomark{2}{1. \textbf{ปจฺจยานุโลม 1. วิภงฺควาร}}
\csromanheader{2}{\textbf{เหตุ}}
\proseitem{57}{อุปาทาน\-สมฺปยุตฺตํ ธมฺมํ ปจฺจยา อุปาทาน\-สมฺปยุตฺโต ธมฺโม อุปฺปชฺชติ เหตุปจฺจยา. ตีณิ. (ปฏิจฺจสทิสํ.)}

\prose{อุปาทาน\-วิปฺปยุตฺตํ ธมฺมํ ปจฺจยา อุปาทาน\-วิปฺปยุตฺโต ธมฺโม อุปฺปชฺชติ เหตุปจฺจยา. อุปาทาน\-วิปฺปยุตฺตํ เอกํ ขนฺธํ ปจฺจยา ฯเปฯ. (ยาว อชฺฌตฺติกา มหาภูตา.) วตฺถุํ ปจฺจยา อุปาทาน\-วิปฺปยุตฺตา ขนฺธา. วตฺถุํ ปจฺจยา ทิฏฺฐิคต\-วิปฺปยุตฺโต โลโภ. (1)}

\prose{อุปาทาน\-วิปฺปยุตฺตํ ธมฺมํ ปจฺจยา อุปาทาน\-สมฺปยุตฺโต ธมฺโม อุปฺปชฺชติ เหตุปจฺจยา. วตฺถุํ ปจฺจยา อุปาทาน\-สมฺปยุตฺตา ขนฺธา. ทิฏฺฐิคต\-วิปฺปยุตฺตํ โลภํ ปจฺจยา สมฺปยุตฺตกา ขนฺธา. (2)}

\prose{\csromanfolio{582}อุปาทาน\-วิปฺปยุตฺตํ ธมฺมํ ปจฺจยา อุปาทาน\-สมฺปยุตฺโต จ อุปาทาน\-วิปฺปยุตฺโต จ ธมฺมา อุปฺปชฺชนฺติ เหตุปจฺจยา. วตฺถุํ ปจฺจยา อุปาทาน\-สมฺปยุตฺตกา ขนฺธา. มหาภูเต ปจฺจยา จิตฺต\-สมุฏฺฐานํ รูปํ. ทิฏฺฐิคต\-วิปฺปยุตฺตํ โลภํ ปจฺจยา สมฺปยุตฺตกา ขนฺธา จิตฺต\-สมุฏฺฐานญฺจ รูปํ. วตฺถุํ ปจฺจยา ทิฏฺฐิคต\-วิปฺปยุตฺต\-โลภ\-สหคตา ขนฺธา จ โลโภ จ. (3)}

\proseitem{58}{อุปาทาน\-สมฺปยุตฺตญฺจ อุปาทาน\-วิปฺปยุตฺตญฺจ ธมฺมํ ปจฺจยา อุปาทาน\-สมฺปยุตฺโต ธมฺโม อุปฺปชฺชติ เหตุปจฺจยา. อุปาทาน\-สมฺปยุตฺตํ เอกํ ขนฺธญฺจ วตฺถุญฺจ ปจฺจยา ตโย ขนฺธา ฯเปฯ เทฺว ขนฺเธ จ ฯเปฯ. ทิฏฺฐิคต\-วิปฺปยุตฺต\-โลภ\-สหคตํ เอกํ ขนฺธญฺจ โลภญฺจ ปจฺจยา ตโย ขนฺธา ฯเปฯ เทฺว ขนฺเธ จ ฯเปฯ. (1)}

\prose{อุปาทาน\-สมฺปยุตฺตญฺจ อุปาทาน\-วิปฺปยุตฺตญฺจ ธมฺมํ ปจฺจยา อุปาทาน\-วิปฺปยุตฺโต ธมฺโม อุปฺปชฺชติ เหตุปจฺจยา. อุปาทาน\-สมฺปยุตฺเต ขนฺเธ จ มหาภูเต จ ปจฺจยา จิตฺต\-สมุฏฺฐานํ รูปํ. ทิฏฺฐิคต\-วิปฺปยุตฺต\-โลภ\-สหคเต ขนฺเธ จ โลภญฺจ ปจฺจยา จิตฺต\-สมุฏฺฐานํ รูปํ. ทิฏฺฐิคต\-วิปฺปยุตฺต\-โลภ\-สหคเต ขนฺเธ จ วตฺถุญฺจ ปจฺจยา โลโภ. (2)}

\prose{อุปาทาน\-สมฺปยุตฺตญฺจ อุปาทาน\-วิปฺปยุตฺตญฺจ ธมฺมํ ปจฺจยา อุปาทาน\-สมฺปยุตฺโต จ อุปาทาน\-วิปฺปยุตฺโต จ ธมฺมา อุปฺปชฺชนฺติ เหตุปจฺจยา. อุปาทาน\-สมฺปยุตฺตํ เอกํ ขนฺธญฺจ วตฺถุญฺจ ปจฺจยา ตโย ขนฺธา ฯเปฯ เทฺว ขนฺเธ จ ฯเปฯ. อุปาทาน\-สมฺปยุตฺเต ขนฺเธ จ มหาภูเต จ ปจฺจยา จิตฺต\-สมุฏฺฐานํ รูปํ. ทิฏฺฐิคต\-วิปฺปยุตฺต\-โลภ\-สหคตํ เอกํ ขนฺธญฺจ วตฺถุญฺจ ปจฺจยา ตโย ขนฺธา โลโภ จ ฯเปฯ เทฺว ขนฺเธ จ ฯเปฯ. (3) (สํขิตฺตํ. อารมฺมณ\-ปจฺจยมฺหิ ปญฺจ วิญฺญาณา กาตพฺพา.)}

\csromanmarkh{1. ปจฺจยานุโลม}\csromanmarkhii{2. สงฺขฺยาวาร}\csromanheadernomark{2}{1. \textbf{ปจฺจยานุโลม 2. สงฺขฺยาวาร}}
\proseitem{59}{เหตุยา นว, อารมฺมเณ นว, อธิปติยา นว, อนนฺตเร นว, (สพฺพตฺถ นว,) วิปาเก เอกํ ฯเปฯ อวิคเต นว.}

\csromanmarkh{2. ปจฺจย\-ปจฺจนีย}\csromanmarkhii{1. วิภงฺควาร}\csromanheadernomark{2}{2. ปจฺจย\-ปจฺจนีย 1. วิภงฺควาร}
\proseitem{60}{อุปาทาน\-วิปฺปยุตฺตํ ธมฺมํ ปจฺจยา อุปาทาน\-วิปฺปยุตฺโต ธมฺโม อุปฺปชฺชติ นเหตุปจฺจยา. อเหตุกํ อุปาทาน\-วิปฺปยุตฺตํ เอกํ ขนฺธํ ปจฺจยา ฯเปฯ. (ยาว \csromanfolio{583}อสญฺญสตฺตา.) จกฺขายตนํ ปจฺจยา จกฺขุวิญฺญาณํ ฯเปฯ. กายายตนํ ปจฺจยา กายวิญฺญาณํ ฯเปฯ. วตฺถุํ ปจฺจยา อเหตุกา อุปาทาน\-วิปฺปยุตฺตา ขนฺธา. วิจิกิจฺฉา\-สหคเต อุทฺธจฺจ\-สหคเต ขนฺเธ จ วตฺถุญฺจ ปจฺจยา วิจิกิจฺฉา\-สหคโต อุทฺธจฺจ\-สหคโต โมโห. (สํขิตฺตํ.)}

\csromanmarkh{2. ปจฺจย\-ปจฺจนีย}\csromanmarkhii{2. สงฺขฺยาวาร}\csromanheadernomark{2}{2. ปจฺจย\-ปจฺจนีย 2. สงฺขฺยาวาร}
\proseitem{60}{นเหตุยา เอกํ, นอารมฺมเณ ตีณิ, นอธิปติยา นว, นอนนฺตเร ตีณิ, นสมนนฺตเร ตีณิ, นอญฺญมญฺเญ ตีณิ, นอุปนิสฺสเย ตีณิ, นปุเรชาเต สตฺต, นปจฺฉาชาเต นว, นอาเสวเน นว, นกมฺเม จตฺตาริ, นวิปาเก นว, นอาหาเร เอกํ, นอินฺทฺริเย เอกํ, นฌาเน เอกํ, นมคฺเค เอกํ, นสมฺปยุตฺเต ตีณิ, นวิปฺปยุตฺเต ฉ, โนนตฺถิยา ตีณิ, โนวิคเต ตีณิ.}

\vspace{\tipitakaparskip}
\csromancenter{(เอวํ อิตเร เทฺว คณนาปิ นิสฺสยวาโรปิ กาตพฺโพ.)}
\csromanmarkh{71. อุปาทาน\-สมฺปยุตฺตทุก}\csromanmarkhii{5. สํสฏฺฐวาร}\csromanheadernomark{2}{71. อุปาทาน\-สมฺปยุตฺตทุก 5. สํสฏฺฐวาร}
\proseitem{62}{อุปาทาน\-สมฺปยุตฺตํ ธมฺมํ สํสฏฺโฐ อุปาทาน\-สมฺปยุตฺโต ธมฺโม อุปฺปชฺชติ เหตุปจฺจยา. อุปาทาน\-สมฺปยุตฺตํ เอกํ ขนฺธํ สํสฏฺฐา. ตีณิ.}

\prose{อุปาทาน\-วิปฺปยุตฺตํ ธมฺมํ สํสฏฺโฐ อุปาทาน\-วิปฺปยุตฺโต ธมฺโม อุปฺปชฺชติ เหตุปจฺจยา. (ปฏิจฺจสทิสํ. อรูปํเยว กาตพฺพํ.)}

\prose{อุปาทาน\-วิปฺปยุตฺตํ ธมฺมํ สํสฏฺโฐ อุปาทาน\-สมฺปยุตฺโต ธมฺโม อุปฺปชฺชติ เหตุปจฺจยา. (ปฏิจฺจสทิสํ. อรูปํเยว กาตพฺพํ.)}

\prose{อุปาทาน\-สมฺปยุตฺตญฺจ อุปาทาน\-วิปฺปยุตฺตญฺจ ธมฺมํ สํสฏฺโฐ อุปาทาน\-สมฺปยุตฺโต ธมฺโม อุปฺปชฺชติ เหตุปจฺจยา. (ปฏิจฺจสทิสํ. อรูปํเยว กาตพฺพํ.)}

\prose{เหตุยา ฉ, อารมฺมเณ ฉ, อธิปติยา ฉ, (สพฺพตฺถ ฉ,) วิปาเก เอกํ ฯเปฯ อวิคเต ฉ. อนุโลมํ.}

\prose{อุปาทาน\-วิปฺปยุตฺตํ ธมฺมํ สํสฏฺโฐ อุปาทาน\-วิปฺปยุตฺโต ธมฺโม อุปฺปชฺชติ นเหตุปจฺจยา. (สํขิตฺตํ.)}

\prose{\csromanfolio{584}นเหตุยา เอกํ, นอธิปติยา ฉ, นปุเรชาเต ฉ, นปจฺฉาชาเต ฉ, นอาเสวเน ฉ, นกมฺเม จตฺตาริ, นวิปาเก ฉ, นฌาเน เอกํ, นมคฺเค เอกํ, นวิปฺปยุตฺเต ฉ.}

\vspace{\tipitakaparskip}
\csromancenter{ปจฺจนียํ.}
\csromancenter{(เอวํ อิตเร เทฺว คณนาปิ สมฺปยุตฺตวาโรปิ กาตพฺโพ.)}
\csromanmarkh{71. อุปาทาน\-สมฺปยุตฺตทุก}\csromanmarkhii{7. ปญฺหาวาร}\csromanheadernomark{2}{71. อุปาทาน\-สมฺปยุตฺตทุก 7. ปญฺหาวาร}
\csromanmarkh{1. ปจฺจยานุโลม}\csromanmarkhii{1. วิภงฺควาร}\csromanheadernomark{2}{1. \textbf{ปจฺจยานุโลม 1. วิภงฺควาร}}
\csromanheader{2}{\textbf{เหตุ}}
\proseitem{63}{อุปาทาน\-สมฺปยุตฺโต ธมฺโม อุปาทาน\-สมฺปยุตฺตสฺส ธมฺมสฺส เหตุปจฺจเยน ปจฺจโย. อุปาทาน\-สมฺปยุตฺตา เหตู สมฺปยุตฺตกานํ ขนฺธานํ เหตุปจฺจเยน ปจฺจโย. (1)}

\prose{อุปาทาน\-สมฺปยุตฺโต ธมฺโม อุปาทาน\-วิปฺปยุตฺตสฺส ธมฺมสฺส เหตุปจฺจเยน ปจฺจโย. อุปาทาน\-สมฺปยุตฺตา เหตู จิตฺต\-สมุฏฺฐานํ รูปานํ เหตุปจฺจเยน ปจฺจโย. ทิฏฺฐิคต\-วิปฺปยุตฺต\-โลภ\-สหคตา เหตู ทิฏฺฐิคต\-วิปฺปยุตฺต\-โลภสฺส จิตฺตสมุฏฺฐานานญฺจ รูปานํ เหตุปจฺจเยน ปจฺจโย. (มูลํ กาตพฺพํ.) อุปาทาน\-สมฺปยุตฺตา เหตู สมฺปยุตฺตกานํ ขนฺธานํ จิตฺตสมุฏฺฐานานญฺจ รูปานํ เหตุปจฺจเยน ปจฺจโย. ทิฏฺฐิคต\-วิปฺปยุตฺต\-โลภ\-สหคตา เหตู สมฺปยุตฺตกานํ ขนฺธานํ โลภสฺส จ จิตฺตสมุฏฺฐานานญฺจ รูปานํ เหตุปจฺจเยน ปจฺจโย. (3)}

\proseitem{64}{อุปาทาน\-วิปฺปยุตฺโต ธมฺโม อุปาทาน\-วิปฺปยุตฺตสฺส ธมฺมสฺส เหตุปจฺจเยน ปจฺจโย. อุปาทาน\-วิปฺปยุตฺตา เหตู สมฺปยุตฺตกานํ ขนฺธานํ จิตฺตสมุฏฺฐานานญฺจ รูปานํ เหตุปจฺจเยน ปจฺจโย. ทิฏฺฐิคต\-วิปฺปยุตฺโต โลโภ จิตฺต\-สมุฏฺฐานานํ รูปานํ เหตุปจฺจเยน ปจฺจโย. ปฏิสนฺธิ\-กฺขเณ ฯเปฯ. (มูลํ กาตพฺพํ.) ทิฏฺฐิคต\-วิปฺปยุตฺโต โลโภ สมฺปยุตฺตกานํ ขนฺธานํ เหตุปจฺจเยน ปจฺจโย. (มูลํ กาตพฺพํ.) ทิฏฺฐิคต\-วิปฺปยุตฺโต \csromanfolio{585}โลโภ สมฺปยุตฺตกานํ ขนฺธานํ จิตฺตสมุฏฺฐานานญฺจ รูปานํ เหตุปจฺจเยน ปจฺจโย. (3)}

\proseitem{65}{อุปาทาน\-สมฺปยุตฺโต จ อุปาทานวิปฺปยุตฺโตจ ธมฺมา อุปาทาน\-สมฺปยุตฺตสฺส ธมฺมสฺส เหตุปจฺจเยน ปจฺจโย. ทิฏฺฐิคต\-วิปฺปยุตฺต\-โลภ\-สหคโต โมโห จ โลโภ จ สมฺปยุตฺตกานํ ขนฺธานํ เหตุปจฺจเยน ปจฺจโย. (มูลํ กาตพฺพํ.) ทิฏฺฐิคต\-วิปฺปยุตฺต\-โลภ\-สหคโต โมโห จ โลโภ จ จิตฺต\-สมุฏฺฐานานํ รูปานํ เหตุปจฺจเยน ปจฺจโย. (2)}

\prose{อุปาทาน\-สมฺปยุตฺโต จ อุปาทาน\-วิปฺปยุตฺโต จ ธมฺมา อุปาทาน\-สมฺปยุตฺตสฺส จ อุปาทาน\-วิปฺปยุตฺตสฺส จ ธมฺมสฺส เหตุปจฺจเยน ปจฺจโย. ทิฏฺฐิคต\-วิปฺปยุตฺต\-โลภ\-สหคโต โมโห จ โลโภ จ สมฺปยุตฺตกานํ ขนฺธานํ จิตฺตสมุฏฺฐานานญฺจ รูปานํ เหตุปจฺจเยน ปจฺจโย. (3)}

\proseitem{64}{อุปาทาน\-สมฺปยุตฺโต ธมฺโม อุปาทาน\-สมฺปยุตฺตสฺส ธมฺมสฺส อารมฺมณ\-ปจฺจเยน ปจฺจโย. อุปาทาน\-สมฺปยุตฺเต ขนฺเธ อารพฺภ อุปาทาน\-สมฺปยุตฺตา ขนฺธา อุปฺปชฺชนฺติ. (มูลํ กาตพฺพํ.) อุปาทาน\-สมฺปยุตฺเต ขนฺเธ อารพฺภ อุปาทาน\-วิปฺปยุตฺตา ขนฺธา จ โลโภ จ อุปฺปชฺชนฺติ. (มูลํ กาตพฺพํ.) อุปาทาน\-สมฺปยุตฺเต ขนฺเธ อารพฺภ ทิฏฺฐิคต\-วิปฺปยุตฺต\-โลภ\-สหคตา ขนฺธา จ โลโภ จ อุปฺปชฺชนฺติ. (3)}

\prose{อุปาทาน\-วิปฺปยุตฺโต ธมฺโม อุปาทาน\-วิปฺปยุตฺตสฺส ธมฺมสฺส อารมฺมณ\-ปจฺจเยน ปจฺจโย. ทานํ ฯเปฯ สีลํ ฯเปฯ อุโปสถกมฺมํ ฯเปฯ. ปุพฺเพ สุจิณฺณานิ ฯเปฯ. ฌานา วุฏฺฐหิตฺวา ฌานํ ปจฺจเวกฺขติ, อสฺสาเทติ อภินนฺทติ, ตํ อารพฺภ ทิฏฺฐิคต\-วิปฺปยุตฺโต ราโค ฯเปฯ. วิจิกิจฺฉา ฯเปฯ อุทฺธจฺจํ อุปฺปชฺชติ. ฌาเน ปริหีเน วิปฺปฏิสาริสฺส โทมนสฺสํ อุปฺปชฺชติ. อริยา มคฺคา วุฏฺฐหิตฺวา มคฺคํ ปจฺจเวกฺขนฺติ. ผลํ ฯเปฯ. นิพฺพานํ ปจฺจเวกฺขนฺติ, นิพฺพานํ โคตฺรภุสฺส, โวทานสฺส, มคฺคสฺส, ผลสฺส, อาวชฺชนาย อารมฺมณ\-ปจฺจเยน ปจฺจโย. อริยา อุปาทาน\-วิปฺปยุตฺเต ปหีเน กิเลเส ฯเปฯ. วิกฺขมฺภิเต กิเลเส ฯเปฯ. ปุพฺเพ สมุทาจิณฺเณ กิเลเส ฯเปฯ. จกฺขุํ ฯเปฯ วตฺถุํ อุปาทาน\-วิปฺปยุตฺเต \csromanfolio{586}ขนฺเธ จ โลภญฺจ อนิจฺจโต ฯเปฯ วิปสฺสติ, อสฺสาเทติ อภินนฺทติ, ตํ อารพฺภ ทิฏฺฐิคต\-วิปฺปยุตฺโต ราโค ฯเปฯ วิจิกิจฺฉา ฯเปฯ อุทฺธจฺจํ ฯเปฯ โทมนสฺสํ อุปฺปชฺชติ ฯเปฯ. (สพฺพํ ปริปุณฺณํ.) ทิพฺเพน จกฺขุนา. (ยาว กายวิญฺญาณํ.) อุปาทาน\-วิปฺปยุตฺตา ขนฺธา อิทฺธิวิธ\-ญาณสฺส, เจโตปริย\-ญาณสฺส, ปุพฺเพ\-นิวาสา\-นุสฺสติ\-ญาณสฺส, ยถากมฺมูปคญาณสฺส, อนาคตํสญาณสฺส, อาวชฺชนาย อารมฺมณ\-ปจฺจเยน ปจฺจโย. (1)}

\proseitem{67}{อุปาทาน\-วิปฺปยุตฺโต ธมฺโม อุปาทาน\-สมฺปยุตฺตสฺส ธมฺมสฺส อารมฺมณ\-ปจฺจเยน ปจฺจโย. ทานํ ฯเปฯ. ฌานํ ฯเปฯ. จกฺขุํ ฯเปฯ. วตฺถุํ อุปาทาน\-วิปฺปยุตฺเต ขนฺเธ จ โลภญฺจ อสฺสาเทติ อภินนฺทติ, ตํ อารพฺภ ราโค อุปฺปชฺชติ. ทิฏฺฐิ อุปฺปชฺชติ. (2)}

\prose{อุปาทาน\-วิปฺปยุตฺโต ธมฺโม อุปาทาน\-สมฺปยุตฺตสฺส จ อุปาทาน\-วิปฺปยุตฺตสฺส จ ธมฺมสฺส อารมฺมณ\-ปจฺจเยน ปจฺจโย. จกฺขุํ ฯเปฯ. วตฺถุํ อุปาทาน\-วิปฺปยุตฺเต ขนฺเธ จ โลภญฺจ อารพฺภ ทิฏฺฐิคต\-วิปฺปยุตฺต\-โลภ\-สหคตา ขนฺธา จ โลโภ จ อุปฺปชฺชนฺติ. (3)}

\prose{อุปาทาน\-สมฺปยุตฺโต จ อุปาทาน\-วิปฺปยุตฺโต จ ธมฺมา อุปาทาน\-สมฺปยุตฺตสฺส ธมฺมสฺส อารมฺมณ\-ปจฺจเยน ปจฺจโย. ทิฏฺฐิคต\-วิปฺปยุตฺต\-โลภ สหคเต ขนฺเธ จ โลภญฺจ อารพฺภ อุปาทาน\-สมฺปยุตฺตา ขนฺธา อุปฺปชฺชนฺติ. (มูลํ กาตพฺพํ.) ทิฏฺฐิคต\-วิปฺปยุตฺต\-โลภ\-สหคเต ขนฺเธ จ โลภญฺจ อารพฺภ อุปาทาน\-วิปฺปยุตฺโต ขนฺธา จ โลโภ จ อุปฺปชฺชนฺติ. (มูลํ กาตพฺพํ.) ทิฏฺฐิคต\-วิปฺปยุตฺต\-โลภ\-สหคเต ขนฺเธ จ โลภญฺจ อารพฺภ ทิฏฺฐิคต\-วิปฺปยุตฺต\-โลภ\-สหคเต ขนฺธา จ โลโภ จ อุปฺปชฺชนฺติ. (3)}

\prose{อุปาทาน\-สมฺปยุตฺโต ธมฺโม อุปาทาน\-สมฺปยุตฺตสฺส ธมฺมสฺส อธิปติ\-ปจฺจเยน ปจฺจโย. อารมฺมณา\-ธิปติ, สหชาตา\-ธิปติ.  อารมฺมณา\-ธิปติ\csromandash{}อุปาทาน\-สมฺปยุตฺเต ขนฺเธ ครุํ กตฺวา อุปาทาน\-สมฺปยุตฺตา ขนฺธา อุปฺปชฺชนฺติ. สหชาตา\-ธิปติ\csromandash{} อุปาทาน\-สมฺปยุตฺตา\-ธิปติ สมฺปยุตฺตกานํ ขนฺธานํ \csromanfolio{587}อธิปติ\-ปจฺจเยน ปจฺจโย. (มูลํ กาตพฺพํ.) อารมฺมณา\-ธิปติ, สหชาตา\-ธิปติ.  อารมฺมณา\-ธิปติ\csromandash{}อุปาทาน\-สมฺปยุตฺเต ขนฺเธ ครุํ กตฺวา ทิฏฺฐิคต\-วิปฺปยุตฺโต โลโภ อุปฺปชฺชติ. สหชาตา\-ธิปติ\csromandash{} อุปาทาน\-สมฺปยุตฺตา\-ธิปติ จิตฺต\-สมุฏฺฐานานํ รูปานํ อธิปติ\-ปจฺจเยน ปจฺจโย. ทิฏฺฐิคต\-วิปฺปยุตฺต\-โลภ\-สหคตา\-ธิปติ โลภสฺส จ จิตฺตสมุฏฺฐานานญฺจ รูปานํ อธิปติ\-ปจฺจเยน ปจฺจโย. (มูลํ กาตพฺพํ.) อารมฺมณา\-ธิปติ, สหชาตา\-ธิปติ.  อารมฺมณา\-ธิปติ\csromandash{} อุปาทาน\-สมฺปยุตฺเต ขนฺเธ ครุํ กตฺวา ทิฏฺฐิคต\-วิปฺปยุตฺต\-โลภ\-สหคตา ขนฺธา จ โลโภ จ อุปฺปชฺชนฺติ. สหชาตา\-ธิปติ\csromandash{}อุปาทาน\-สมฺปยุตฺตา\-ธิปติ สมฺปยุตฺตกานํ ขนฺธานํ จิตฺตสมุฏฺฐานานญฺจ รูปานํ อธิปติ\-ปจฺจเยน ปจฺจโย. ทิฏฺฐิคต\-วิปฺปยุตฺต\-โลภ\-สหคตา\-ธิปติ สมฺปยุตฺตกานํ ขนฺธานํ โลภสฺส จ จิตฺตสมุฏฺฐานานญฺจ รูปานํ อธิปติ\-ปจฺจเยน ปจฺจโย. (3)}

\proseitem{70}{อุปาทาน\-วิปฺปยุตฺโต ธมฺโม อุปาทาน\-วิปฺปยุตฺตสฺส ธมฺมสฺส อธิปติ\-ปจฺจเยน ปจฺจโย. อารมฺมณา\-ธิปติ, สหชาตา\-ธิปติ.  อารมฺมณา\-ธิปติ\csromandash{}ทานํ ฯเปฯ. สีลํ ฯเปฯ. ฌานา วุฏฺฐหิตฺวา ฌานํ ครุํ กตฺวา ปจฺจเวกฺขติ, อสฺสาเทติ อภินนฺทติ, ตํ ครุํ กตฺวา ทิฏฺฐิคต\-วิปฺปยุตฺโต ราโค อุปฺปชฺชติ ฯเปฯ. อริยา มคฺคา วุฏฺฐหิตฺวา ฯเปฯ ผลสฺส อธิปติ\-ปจฺจเยน ปจฺจโย. จกฺขุํ ฯเปฯ. วตฺถุํ อุปาทาน\-วิปฺปยุตฺเต ขนฺเธ จ โลภญฺจ ครุํ กตฺวา อสฺสาเทติ อภินนฺทติ, ตํ ครุํ กตฺวา ทิฏฺฐิคต\-วิปฺปยุตฺโต ราโค อุปฺปชฺชติ. สหชาตา\-ธิปติ\csromandash{} อุปาทาน\-วิปฺปยุตฺตา\-ธิปติ สมฺปยุตฺตกานํ ขนฺธานํ จิตฺตสมุฏฺฐานานญฺจ รูปานํ อธิปติ\-ปจฺจเยน ปจฺจโย. (1)}

\prose{อุปาทาน\-วิปฺปยุตฺโต ธมฺโม อุปาทาน\-สมฺปยุตฺตสฺส ธมฺมสฺส อธิปติ\-ปจฺจเยน ปจฺจโย.  อารมฺมณา\-ธิปติ\csromandash{}ทานํ ฯเปฯ. ฌานํ ฯเปฯ. จกฺขุํ ฯเปฯ. วตฺถุํ อุปาทาน\-วิปฺปยุตฺเต ขนฺเธ จ โลภญฺจ ครุํ กตฺวา อสฺสาเทติ อภินนฺทติ, ตํ ครุํ กตฺวา ราโค อุปฺปชฺชติ. ทิฏฺฐิ อุปฺปชฺชติ.}

\prose{อุปาทาน\-วิปฺปยุตฺโต ธมฺโม อุปาทาน\-สมฺปยุตฺตสฺส จ อุปาทาน\-วิปฺปยุตฺตสฺส จ ธมฺมสฺส อธิปติ\-ปจฺจเยน ปจฺจโย.  อารมฺมณา\-ธิปติ\csromandash{}จกฺขุํ ฯเปฯ วตฺถุํ อุปาทาน\-วิปฺปยุตฺเต ขนฺเธ จ โลภญฺจ ครุํ กตฺวา ทิฏฺฐิคต\-วิปฺปยุตฺต\-โลภ\-สหคตา ขนฺธา จ โลโภ จ อุปฺปชฺชนฺติ. (3)}

\proseitem{71}{\csromanfolio{588}อุปาทาน\-สมฺปยุตฺโต จ อุปาทาน\-วิปฺปยุตฺโต จ ธมฺมา อุปาทาน\-สมฺปยุตฺตสฺส ธมฺมสฺส อธิปติ\-ปจฺจเยน ปจฺจโย.  อารมฺมณา\-ธิปติ\csromandash{}ทิฏฺฐิคต\-วิปฺปยุตฺต\-โลภ\-สหคเต ขนฺเธ จ โลภญฺจ ครุํ กตฺวา อุปาทาน\-สมฺปยุตฺตา ขนฺธา อุปฺปชฺชนฺติ. (มูลํ กาตพฺพํ.) ทิฏฺฐิคต\-วิปฺปยุตฺต\-โลภ\-สหคเต ขนฺเธ จ โลภญฺจ ครุํ กตฺวา ทิฏฺฐิคต\-วิปฺปยุตฺโต โลโภ อุปฺปชฺชติ. (มูลํ กาตพฺพํ.) ทิฏฺฐิคต\-วิปฺปยุตฺต\-โลภ\-สหคเต ขนฺเธ จ โลภญฺจ ครุํ กตฺวา ทิฏฺฐิคต\-วิปฺปยุตฺต\-โลภ\-สหคตา ขนฺธา จ โลโภ จ อุปฺปชฺชนฺติ. (3)}

\proseitem{72}{อุปาทาน\-สมฺปยุตฺโต ธมฺโม อุปาทาน\-สมฺปยุตฺตสฺส ธมฺมสฺส อนนฺตร\-ปจฺจเยน ปจฺจโย. ปุริมา ปุริมา อุปาทาน\-สมฺปยุตฺตา ขนฺธา ปจฺฉิมานํ ปจฺฉิมานํ อุปาทาน\-สมฺปยุตฺตกานํ ขนฺธานํ อนนฺตร\-ปจฺจเยน ปจฺจโย. (มูลํ กาตพฺพํ.) ปุริมา ปุริมา ทิฏฺฐิคต\-วิปฺปยุตฺต\-โลภ\-สหคตา ขนฺธา ปจฺฉิมสฺส ปจฺฉิมสฺส ทิฏฺฐิคต\-วิปฺปยุตฺตสฺส โลภสฺส อนนฺตร\-ปจฺจเยน ปจฺจโย. อุปาทาน\-สมฺปยุตฺตา ขนฺธา วุฏฺฐานสฺส อนนฺตร\-ปจฺจเยน ปจฺจโย. (มูลํ กาตพฺพํ.) ปุริมา ปุริมา ทิฏฺฐิคต\-วิปฺปยุตฺต\-โลภ\-สหคตา ขนฺธา ปจฺฉิมานํ ปจฺฉิมานํ ทิฏฺฐิคต\-วิปฺปยุตฺต\-โลภ\-สหคตานํ ขนฺธานํ โลภสฺส จ อนนฺตร\-ปจฺจเยน ปจฺจโย. (3)}

\proseitem{73}{อุปาทาน\-วิปฺปยุตฺโต ธมฺโม อุปาทาน\-วิปฺปยุตฺตสฺส ธมฺมสฺส อนนฺตร\-ปจฺจเยน ปจฺจโย. ปุริโม ปุริโม ทิฏฺฐิคต\-วิปฺปยุตฺโต โลโภ ปจฺฉิมสฺส ปจฺฉิมสฺส ทิฏฺฐิคต\-วิปฺปยุตฺตสฺส โลภสฺส อนนฺตร\-ปจฺจเยน ปจฺจโย. ปุริมา ปุริมา อุปาทาน\-วิปฺปยุตฺตา ขนฺธา ปจฺฉิมานํ ปจฺฉิมานํ อุปาทาน\-วิปฺปยุตฺตานํ ขนฺธานํ อนนฺตร\-ปจฺจเยน ปจฺจโย. ทิฏฺฐิคต\-วิปฺปยุตฺโต โลโภ วุฏฺฐานสฺส อนนฺตร\-ปจฺจเยน ปจฺจโย. อนุโลมํ โคตฺรภุสฺส ฯเปฯ ผลสมาปตฺติยา อนนฺตร\-ปจฺจเยน ปจฺจโย. (มูลํ กาตพฺพํ.) ปุริโม ปุริโม ทิฏฺฐิคต\-วิปฺปยุตฺโต โลโภ ปจฺฉิมานํ ปจฺฉิมานํ ทิฏฺฐิคต\-วิปฺปยุตฺต\-โลภ\-สหคตานํ ขนฺธานํ อนนฺตร\-ปจฺจเยน ปจฺจโย. อาวชฺชนา อุปาทาน\-สมฺปยุตฺตกานํ ขนฺธานํ อนนฺตร\-ปจฺจเยน ปจฺจโย. (มูลํ กาตพฺพํ.) ปุริโม ปุริโม ทิฏฺฐิคต\-วิปฺปยุตฺโต โลโภ ปจฺฉิมานํ ปจฺฉิมานํ ทิฏฺฐิคต\-วิปฺปยุตฺต\-โลภ\-สหคตานํ ขนฺธานํ โลภสฺส จ อนนฺตร\-ปจฺจเยน ปจฺจโย. \csromanfolio{589}อาวชฺชนา ทิฏฺฐิคต\-วิปฺปยุตฺต\-โลภ\-สหคตานํ ขนฺธานํ โลภสฺส จ อนนฺตร\-ปจฺจเยน ปจฺจโย. (3)}

\proseitem{74}{อุปาทาน\-สมฺปยุตฺโต จ อุปาทาน\-วิปฺปยุตฺโต จ ธมฺมา อุปาทาน\-สมฺปยุตฺตสฺส ธมฺมสฺส อนนฺตร\-ปจฺจเยน ปจฺจโย. ปุริมา ปุริมา ทิฏฺฐิคต\-วิปฺปยุตฺต\-โลภ\-สหคตา ขนฺธา จ โลโภ จ ปจฺฉิมานํ ปจฺฉิมานํ ทิฏฺฐิคต\-วิปฺปยุตฺต\-โลภ\-สหคตานํ ขนฺธานํ อนนฺตร\-ปจฺจเยน ปจฺจโย. (มูลํ กาตพฺพํ.) ปุริมา ปุริมา ทิฏฺฐิคต\-วิปฺปยุตฺต\-โลภ\-สหคตา ขนฺธา จ โลโภ จ ปจฺฉิมสฺส ปจฺฉิมสฺส ทิฏฺฐิคต\-วิปฺปยุตฺตสฺส โลภสฺส อนนฺตร\-ปจฺจเยน ปจฺจโย. ทิฏฺฐิคต\-วิปฺปยุตฺต\-โลภ\-สหคตา ขนฺธา จ โลโภ จ วุฏฺฐานสฺส อนนฺตร\-ปจฺจเยน ปจฺจโย. (มูลํ กาตพฺพํ.) ปุริมา ปุริมา ทิฏฺฐิคต\-วิปฺปยุตฺต\-โลภ\-สหคตา ขนฺธา จ โลโภ จ ปจฺฉิมานํ ปจฺฉิมานํ ทิฏฺฐิคต\-วิปฺปยุตฺต\-โลภ\-สหคตานํ ขนฺธานํ โลภสฺส จ อนนฺตร\-ปจฺจเยน ปจฺจโย. (3)}

\csromanheader{2}{\textbf{สมนนฺตราทิ}}
\proseitem{75}{อุปาทาน\-สมฺปยุตฺโต ธมฺโม อุปาทาน\-สมฺปยุตฺตสฺส ธมฺมสฺส สมนนฺตร\-ปจฺจเยน ปจฺจโย. สหชาต\-ปจฺจเยน ปจฺจโย. (ปฏิจฺจสทิสํ.) นว. อญฺญมญฺญ\-ปจฺจเยน ปจฺจโย. (ปฏิจฺจสทิสํ.) ฉ. นิสฺสย\-ปจฺจเยน ปจฺจโย. (ปจฺจยวาร\-สทิสํ.) นว.}

\csromanheader{2}{\textbf{อุปนิสฺสย}}
\proseitem{76}{อุปาทาน\-สมฺปยุตฺโต ธมฺโม อุปาทาน\-สมฺปยุตฺตสฺส ธมฺมสฺส อุปนิสฺสย\-ปจฺจเยน ปจฺจโย. อารมฺมณู\-ปนิสฺสโย, อนนฺตรู\-ปนิสฺสโย, ปกตู\-ปนิสฺสโย ฯเปฯ. ปกตู\-ปนิสฺสโย\csromandash{}อุปทานสมฺปยุตฺตา ขนฺธา อุปาทาน\-สมฺปยุตฺตกานํ ขนฺธานํ อุปนิสฺสย\-ปจฺจเยน ปจฺจโย. (มูลํ กาตพฺพํ.) อุปาทาน\-สมฺปยุตฺตา ขนฺธา อุปาทาน\-วิปฺป\-ยุตฺตกานํ ขนฺธานํ โลภสฺส จ อุปนิสฺสย\-ปจฺจเยน ปจฺจโย. (มูลํ กาตพฺพํ.) อุปาทาน\-สมฺปยุตฺตา ขนฺธา ทิฏฺฐิคต\-วิปฺปยุตฺต\-โลภ\-สหคตานํ ขนฺธานํ โลภสฺส จ อุปนิสฺสย\-ปจฺจเยน ปจฺจโย. (3)}

\proseitem{77}{\csromanfolio{590}อุปาทาน\-วิปฺปยุตฺโต ธมฺโม อุปาทาน\-วิปฺปยุตฺตสฺส ธมฺมสฺส อุปนิสฺสย\-ปจฺจเยน ปจฺจโย. อารมฺมณู\-ปนิสฺสโย, อนนฺตรู\-ปนิสฺสโย, ปกตู\-ปนิสฺสโย ฯเปฯ. ปกตู\-ปนิสฺสโย\csromandash{}สทฺธํ อุปนิสฺสาย ทานํ เทติ ฯเปฯ สมาปตฺติํ อุปฺปาเทติ, มานํ ชปฺเปติ. สีลํ ฯเปฯ. ปญฺญํ. อุปาทาน\-วิปฺปยุตฺตํ ราคํ. มานํ. ปตฺถนํ ฯเปฯ. เสนาสนํ อุปนิสฺสาย ทานํ เทติ ฯเปฯ สํฆํ ภินฺทติ. สทฺธา ฯเปฯ. เสนาสนํ สทฺธาย ฯเปฯ ปญฺญาย อุปาทาน\-วิปฺปยุตฺตสฺส ราคสฺส. มานสฺส. ปตฺถนาย. กายิกสฺส สุขสฺส ฯเปฯ ผลสมาปตฺติยา อุปนิสฺสย\-ปจฺจเยน ปจฺจโย. (1)}

\prose{อุปาทาน\-วิปฺปยุตฺโต ธมฺโม อุปาทาน\-สมฺปยุตฺตสฺส ธมฺมสฺส อุปนิสฺสย\-ปจฺจเยน ปจฺจโย. (ตีณิ อุปนิสฺสยา.) สทฺธํ อุปนิสฺสาย มานํ ชปฺเปติ. ทิฏฺฐิํ คณฺหาติ. สีลํ ฯเปฯ. ปญฺญํ. อุปาทาน\-วิปฺปยุตฺตํ ราคํ. มานํ. ปตฺถนํ. เสนาสนํ อุปนิสฺสาย อทินฺนํ ฯเปฯ มุสา ฯเปฯ ปิสุณํ ฯเปฯ สมฺผํ ฯเปฯ สนฺธิํ ฯเปฯ นิลฺโลปํ ฯเปฯ เอกาคาริกํ ฯเปฯ ปริปนฺเถ ฯเปฯ ปรทารํ ฯเปฯ คามฆาตํ ฯเปฯ นิคฆาตํ กโรติ. สทฺธา ฯเปฯ. เสนาสนํ อุปาทาน\-สมฺปยุตฺตสฺส ราคสฺส. โมหสฺส. มานสฺส. ทิฏฺฐิยา. ปตฺถนาย อุปนิสฺสย\-ปจฺจเยน ปจฺจโย. (2)}

\prose{อุปาทาน\-วิปฺปยุตฺโต ธมฺโม อุปาทาน\-สมฺปยุตฺตสฺส จ อุปาทาน\-วิปฺปยุตฺตสฺส จ ธมฺมสฺส อุปนิสฺสย\-ปจฺจเยน ปจฺจโย. อารมฺมณู\-ปนิสฺสโย, อนนฺตรู\-ปนิสฺสโย, ปกตู\-ปนิสฺสโย ฯเปฯ. ปกตู\-ปนิสฺสโย\csromandash{}สทฺธํ อุปนิสฺสาย มานํ ชปฺเปติ. สีลํ ฯเปฯ. ปญฺญํ. อุปาทาน\-วิปฺปยุตฺตํ ราคํ ฯเปฯ. เสนาสนํ อุปนิสฺสาย อทินฺนํ อาทิยติ ฯเปฯ. คามฆาตํ กโรติ. นิคมฆาตํ กโรติ. สทฺธา ฯเปฯ. เสนาสนํ ทิฏฺฐิคต\-วิปฺปยุตฺต\-โลภ\-สหคตานํ ขนฺธานํ โลภสฺส จ อุปนิสฺสย\-ปจฺจเยน ปจฺจโย. (3)}

\proseitem{78}{อุปาทาน\-สมฺปยุตฺโต จ อุปาทาน\-วิปฺปยุตฺโต จ ธมฺมา อุปาทาน\-สมฺปยุตฺตสฺส ธมฺมสฺส อุปนิสฺสย\-ปจฺจเยน ปจฺจโย. ทิฏฺฐิคต\-วิปฺปยุตฺต\-โลภ\-สหคตา ขนฺธา จ โลโภ จ อุปาทาน\-สมฺปยุตฺตกานํ ขนฺธานํ อุปนิสฺสย\-ปจฺจเยน ปจฺจโย. (มูลํ กาตพฺพํ.) ทิฏฺฐิคต\-วิปฺปยุตฺต\-โลภ\-สหคตา ขนฺธา จ โลโภ จ อุปาทาน\-วิปฺปยุตฺตานํ ขนฺธานํ โลภสฺส จ อุปนิสฺสย\-ปจฺจเยน ปจฺจโย. (มูลํ กาตพฺพํ.) ทิฏฺฐิคต\-วิปฺปยุตฺต\-โลภ\-สหคตา ขนฺธา จ โลโภ จ ทิฏฺฐิคต\-วิปฺปยุตฺต\-โลภ\-สหคตานํ ขนฺธานํ โลภสฺส จ อุปนิสฺสย\-ปจฺจเยน ปจฺจโย. (3)}

\csromanheader{2}{71. อุปาทาน\-สมฺปยุตฺตทุก ปุเรชาต}
\proseitem{79}{\csromanfolio{591}อุปาทาน\-วิปฺปยุตฺโต ธมฺโม อุปาทาน\-วิปฺปยุตฺตสฺส ธมฺมสฺส ปุเรชาต\-ปจฺจเยน ปจฺจโย. อารมฺมณ\-ปุเรชาตํ, วตฺถุ\-ปุเรชาตํ.  อารมฺมณ\-ปุเรชาตํ\csromandash{}จกฺขุํ ฯเปฯ วตฺถุํ อนิจฺจโต ฯเปฯ โทมนสฺสํ อุปฺปชฺชติ. ทิพฺเพน จกฺขุนา รูปํ ปสฺสติ. ทิพฺพาย โสตธาตุยา สทฺทํ สุณาติ. รูปายตนํ จกฺขุ\-วิญฺญาณสฺส ฯเปฯ. โผฏฺฐพฺพา\-ยตนํ ฯเปฯ. วตฺถุ\-ปุเรชาตํ\csromandash{}จกฺขายตนํ จกฺขุ\-วิญฺญาณสฺส ฯเปฯ กายายตนํ ฯเปฯ. วตฺถุ อุปาทาน\-วิปฺปยุตฺตานํ ขนฺธานํ โลภสฺส จ ปุเรชาต\-ปจฺจเยน ปจฺจโย. (1)}

\prose{อุปาทาน\-วิปฺปยุตฺโต ธมฺโม อุปาทาน\-สมฺปยุตฺตสฺส ธมฺมสฺส ปุเรชาต\-ปจฺจเยน ปจฺจโย. อารมฺมณ\-ปุเรชาตํ, วตฺถุ\-ปุเรชาตํ.  อารมฺมณ\-ปุเรชาตํ\csromandash{}จกฺขุํ ฯเปฯ วตฺถุํ อสฺสาเทติ อภินนฺทติ, ตํ อารพฺภ ราโค อุปฺปชฺชติ. ทิฏฺฐิ อุปฺปชฺชติ. วตฺถุ\-ปุเรชาตํ\csromandash{}วตฺถุ อุปาทาน\-สมฺปยุตฺตกานํ ขนฺธานํ ปุเรชาต\-ปจฺจเยน ปจฺจโย. (2)}

\prose{อุปาทาน\-วิปฺปยุตฺโต ธมฺโม อุปาทาน\-สมฺปยุตฺตสฺส จ อุปาทาน\-วิปฺปยุตฺตสฺส จ ธมฺมสฺส ปุเรชาต\-ปจฺจเยน ปจฺจโย. อารมฺมณ\-ปุเรชาตํ, วตฺถุ\-ปุเรชาตํ.  อารมฺมณ\-ปุเรชาตํ\csromandash{}จกฺขุํ ฯเปฯ วตฺถุํ อารพฺภ ทิฏฺฐิคต\-วิปฺปยุตฺต\-โลภ\-สหคตา ขนฺธา จ โลโภ จ อุปฺปชฺชนฺติ. วตฺถุ\-ปุเรชาตํ\csromandash{}วตฺถุ ทิฏฺฐิคต\-วิปฺปยุตฺต\-โลภ\-สหคตานํ ขนฺธานํ โลภสฺส จ ปุเรชาต\-ปจฺจเยน ปจฺจโย. (3)}

\proseitem{80}{อุปาทาน\-สมฺปยุตฺโต ธมฺโม อุปาทาน\-วิปฺปยุตฺตสฺส ธมฺมสฺส ปจฺฉาชาต\-ปจฺจเยน ปจฺจโย. (สํขิตฺตํ.) (1)}

\prose{อุปาทาน\-วิปฺปยุตฺโต ธมฺโม อุปาทาน\-วิปฺปยุตฺตสฺส ธมฺมสฺส ปจฺฉาชาต\-ปจฺจเยน ปจฺจโย. (สํขิตฺตํ.) (1)}

\prose{อุปาทาน\-สมฺปยุตฺโต จ อุปาทาน\-วิปฺปยุตฺโต จ ธมฺมา อุปาทาน\-วิปฺปยุตฺตสฺส ธมฺมสฺส ปจฺฉาชาต\-ปจฺจเยน ปจฺจโย. (สํขิตฺตํ.) (1) อาเสวน\-ปจฺจเยน ปจฺจโย.}

\csromanheader{2}{\textbf{กมฺมวิปาก}}
\proseitem{81}{\csromanfolio{592}อุปาทาน\-สมฺปยุตฺโต ธมฺโม อุปาทาน\-สมฺปยุตฺตสฺส ธมฺมสฺส กมฺมปจฺจเยน ปจฺจโย. อุปาทาน\-สมฺปยุตฺตา เจตนา สมฺปยุตฺตกานํ ขนฺธานํ กมฺมปจฺจเยน ปจฺจโย. (1)}

\prose{อุปาทาน\-สมฺปยุตฺโต ธมฺโม อุปาทาน\-วิปฺปยุตฺตสฺส ธมฺมสฺส กมฺมปจฺจเยน ปจฺจโย. สหชาตา, นานากฺขณิกา.  สหชาตา อุปาทาน\-สมฺปยุตฺตา เจตนา จิตฺต\-สมุฏฺฐานานํ รูปานํ กมฺมปจฺจเยน ปจฺจโย. ทิฏฺฐิคต\-วิปฺปยุตฺต\-โลภ\-สหคตา เจตนา โลภสฺส จิตฺตสมุฏฺฐานานญฺจ รูปานํ กมฺมปจฺจเยน ปจฺจโย. นานากฺขณิกา อุปาทาน\-สมฺปยุตฺตา เจตนา วิปากานํ ขนฺธานํ กฏตฺตา จ รูปานํ กมฺมปจฺจเยน ปจฺจโย. (2)}

\prose{อุปาทาน\-สมฺปยุตฺโต ธมฺโม อุปาทาน\-สมฺปยุตฺตสฺส จ อุปาทาน\-วิปฺปยุตฺตสฺส จ ธมฺมสฺส กมฺมปจฺจเยน ปจฺจโย. อุปาทาน\-สมฺปยุตฺตา เจตนา สมฺปยุตฺตกานํ ขนฺธานํ จิตฺตสมุฏฺฐานานญฺจ รูปานํ กมฺมปจฺจเยน ปจฺจโย. ทิฏฺฐิคต\-วิปฺปยุตฺต\-โลภ\-สหคตา เจตนา สมฺปยุตฺตกานํ ขนฺธานํ โลภสฺส จ จิตฺตสมุฏฺฐานานญฺจ รูปานํ กมฺมปจฺจเยน ปจฺจโย. (3)}

\proseitem{82}{อุปาทาน\-วิปฺปยุตฺโต ธมฺโม อุปาทาน\-วิปฺปยุตฺตสฺส ธมฺมสฺส กมฺมปจฺจเยน ปจฺจโย. สหชาตา, นานากฺขณิกา.  สหชาตา อุปาทาน\-วิปฺปยุตฺตา เจตนา สมฺปยุตฺตกานํ ขนฺธานํ จิตฺตสมุฏฺฐานานญฺจ รูปานํ กมฺมปจฺจเยน ปจฺจโย. ปฏิสนฺธิ\-กฺขเณ ฯเปฯ. นานากฺขณิกา อุปาทาน\-วิปฺปยุตฺตา เจตนา วิปากานํ ขนฺธานํ กฏตฺตา จ รูปานํ กมฺมปจฺจเยน ปจฺจโย. (1)}

\prose{อุปาทาน\-วิปฺปยุตฺโต ธมฺโม อุปาทาน\-วิปฺปยุตฺตสฺส ธมฺมสฺส วิปาก\-ปจฺจเยน ปจฺจโย. วิปาโก อุปาทาน\-วิปฺปยุตฺโต เอโก ฯเปฯ. เอกํ.}

\csromanheader{2}{\textbf{อาหาราทิ}}
\proseitem{83}{อุปาทาน\-สมฺปยุตฺโต ธมฺโม อุปาทาน\-สมฺปยุตฺตสฺส ธมฺมสฺส อาหารปจฺจเยน ปจฺจโย. อินฺทฺริย\-ปจฺจเยน ปจฺจโย. ฌานปจฺจเยน ปจฺจโย. มคฺคปจฺจเยน ปจฺจโย. (อิเมสุ จตูสุปิ ยถา กมฺมปจฺจเย ทิฏฺฐิคต\-วิปฺปยุตฺโต โลโภ ทสฺสิโต, เอวํ ทสฺเสตพฺโพ. จตฺตาริ จตฺตาริ ปญฺหา.) สมฺปยุตฺต\-ปจฺจเยน ปจฺจโย. ฉ.}

\csromanheader{2}{71. อุปาทาน\-สมฺปยุตฺตทุก วิปฺปยุตฺต}
\proseitem{84}{\csromanfolio{593}อุปาทาน\-สมฺปยุตฺโต ธมฺโม อุปาทาน\-วิปฺปยุตฺตสฺส ธมฺมสฺส วิปฺปยุตฺต\-ปจฺจเยน ปจฺจโย. สหชาตํ, ปจฺฉาชาตํ. (สํขิตฺตํ.) (1)}

\prose{อุปาทาน\-วิปฺปยุตฺโต ธมฺโม อุปาทาน\-วิปฺปยุตฺตสฺส ธมฺมสฺส วิปฺปยุตฺต\-ปจฺจเยน ปจฺจโย. สหชาตํ, ปุเรชาตํ, ปจฺฉาชาตํ. (สํขิตฺตํ.) (1)}

\prose{อุปาทาน\-วิปฺปยุตฺโต ธมฺโม อุปาทาน\-สมฺปยุตฺตสฺส ธมฺมสฺส วิปฺปยุตฺต\-ปจฺจเยนปจฺจโย. ปุเรชาตํ วตฺถุ อุปาทาน\-สมฺปยุตฺตกานํ ขนฺธานํ วิปฺปยุตฺต\-ปจฺจเยน ปจฺจโย. (2)}

\prose{อุปาทาน\-วิปฺปยุตฺโต ธมฺโม อุปาทาน\-สมฺปยุตฺตสฺส จ อุปาทาน\-วิปฺปยุตฺตสฺส จ ธมฺมสฺส วิปฺปยุตฺต\-ปจฺจเยน ปจฺจโย. ปุเรชาตํ วตฺถุ ทิฏฺฐิคต\-วิปฺปยุตฺต\-โลภ\-สหคตานํ ขนฺธานํ โลภสฺส จ วิปฺปยุตฺต\-ปจฺจเยน ปจฺจโย. (3)}

\prose{อุปาทาน\-สมฺปยุตฺโต จ อุปาทาน\-วิปฺปยุตฺโต จ ธมฺมา อุปาทาน\-วิปฺปยุตฺตสฺส ธมฺมสฺส วิปฺปยุตฺต\-ปจฺจเยน ปจฺจโย. สหชาตํ, ปจฺฉาชาตํ.  สหชาตา ทิฏฺฐิคต\-วิปฺปยุตฺต\-โลภ\-สหคตา ขนฺธา จ โลโภ จ จิตฺต\-สมุฏฺฐานานํ รูปานํ วิปฺปยุตฺต\-ปจฺจเยน ปจฺจโย. ปจฺฉาชาตา\csromandash{} ฯเปฯ. (1)}

\prose{อตฺถิอาทิ}

\proseitem{85}{อุปาทาน\-สมฺปยุตฺโต ธมฺโม อุปาทาน\-สมฺปยุตฺตสฺส ธมฺมสฺส อตฺถิปจฺจเยน ปจฺจโย. เอกํ. (ปฏิจฺจสทิสํ.) (1)}

\prose{อุปาทาน\-สมฺปยุตฺโต ธมฺโม อุปาทาน\-วิปฺปยุตฺตสฺส ธมฺมสฺส อตฺถิปจฺจเยน ปจฺจโย. สหชาตํ, ปจฺฉาชาตํ. (สํขิตฺตํ.) (2)}

\prose{อุปาทาน\-สมฺปยุตฺโต ธมฺโม อุปาทาน\-สมฺปยุตฺตสฺส จ อุปาทาน\-วิปฺปยุตฺตสฺส จ ธมฺมสฺส อตฺถิปจฺจเยน ปจฺจโย. (ปฏิจฺจสทิสํ.) (3)}

\proseitem{86}{อุปาทาน\-วิปฺปยุตฺโต ธมฺโม อุปาทาน\-วิปฺปยุตฺตสฺส ธมฺมสฺส อตฺถิปจฺจเยน ปจฺจโย. สหชาตํ, ปุเรชาตํ, ปจฺฉาชาตํ, อาหารํ. อินฺทฺริยํ. (สํขิตฺตํ.) (1)}

\prose{อุปาทาน\-วิปฺปยุตฺโต ธมฺโม อุปาทาน\-สมฺปยุตฺตสฺส ธมฺมสฺส อตฺถิปจฺจเยน ปจฺจโย, สหชาตํ, ปุเรชาตํ. (สํขิตฺตํ.) (2)}

\prose{\csromanfolio{594}อุปาทาน\-วิปฺปยุตฺโต ธมฺโม อุปาทาน\-สมฺปยุตฺตสฺส จ อุปาทาน\-วิปฺปยุตฺตสฺส จ ธมฺมสฺส อตฺถิปจฺจเยน ปจฺจโย. สหชาตํ, ปุเรชาตํ. (อิเมสุ สหชาตํ สหชาต\-สทิสํ, ปุเรชาตํ ปุเรชาต\-สทิสํ.) (3)}

\proseitem{87}{อุปาทาน\-สมฺปยุตฺโต จ อุปาทาน\-วิปฺปยุตฺโต จ ธมฺมา อุปาทาน\-สมฺปยุตฺตสฺส ธมฺมสฺส อตฺถิปจฺจเยน ปจฺจโย. สหชาตํ, ปุเรชาตํ.  สหชาโต อุปาทาน\-สมฺปยุตฺโต เอโก ขนฺโธ จ วตฺถุ จ ติณฺณนฺนํ ขนฺธานํ อตฺถิปจฺจเยน ปจฺจโย ฯเปฯ. เทฺว ขนฺธา จ ฯเปฯ. สหชาโต ทิฏฺฐิคต\-วิปฺปยุตฺต\-โลภ\-สหคโต เอโก ขนฺโธ จ โลโภ จ ติณฺณนฺนํ ขนฺธานํ อตฺถิปจฺจเยน ปจฺจโย ฯเปฯ. เทฺว ขนฺธา จ ฯเปฯ. (1)}

\prose{อุปาทาน\-สมฺปยุตฺโต จ อุปาทาน\-วิปฺปยุตฺโต จ ธมฺมา อุปาทาน\-วิปฺปยุตฺตสฺส ธมฺมสฺส อตฺถิปจฺจเยน ปจฺจโย. สหชาตํ, ปุเรชาตํ, ปจฺฉาชาตํ, อาหารํ, อินฺทฺริยํ.  สหชาตา อุปาทาน\-สมฺปยุตฺตา ขนฺธา จ มหาภูตา จ จิตฺต\-สมุฏฺฐานานํ รูปานํ อตฺถิปจฺจเยน ปจฺจโย. สหชาตา ทิฏฺฐิคต\-วิปฺปยุตฺต\-โลภ\-สหคตา ขนฺธา จ โลโภ จ จิตฺต\-สมุฏฺฐานานํ รูปานํ อตฺถิปจฺจเยน ปจฺจโย. สหชาตา ทิฏฺฐิคต\-วิปฺปยุตฺต\-โลภ\-สหคตา ขนฺธา จ วตฺถุ จ โลภสฺส อตฺถิปจฺจเยน ปจฺจโย. ปจฺฉาชาตา ทิฏฺฐิคต\-วิปฺปยุตฺต\-โลภ\-สหคตา ขนฺธา จ โลโภ จ ปุเรชาตสฺส อิมสฺส กายสฺส อตฺถิปจฺจเยน ปจฺจโย. ปจฺฉาชาตา ทิฏฺฐิคต\-วิปฺปยุตฺต\-โลภ\-สหคตา ขนฺธา จ โลโภ จ กพฬีกาโร อาหาโร จ อิมสฺส กายสฺส อตฺถิปจฺจเยน ปจฺจโย. ปจฺฉาชาตา ทิฏฺฐิคต\-วิปฺปยุตฺต\-โลภ\-สหคตา ขนฺธา จ โลโภ จ รูปชีวิตินฺทฺริยญฺจ กฏตฺตารูปานํ อตฺถิปจฺจเยน ปจฺจโย. (2)}

\prose{อุปาทาน\-สมฺปยุตฺโต จ อุปาทาน\-วิปฺปยุตฺโต จ ธมฺมา อุปาทาน\-สมฺปยุตฺตสฺส จ อุปาทาน\-วิปฺปยุตฺตสฺส จ ธมฺมสฺส อตฺถิปจฺจเยน ปจฺจโย. สหชาตํ, ปุเรชาตํ.  สหชาโต ทิฏฺฐิคต\-วิปฺปยุตฺต\-โลภ\-สหคโต เอโก ขนฺโธ จ โลโภ จ ติณฺณนฺนํ ขนฺธานํ จิตฺตสมุฏฺฐานานญฺจ รูปานํ อตฺถิปจฺจเยน ปจฺจโย ฯเปฯ. เทฺว ขนฺธา จ ฯเปฯ. สหชาโต ทิฏฺฐิคต\-วิปฺปยุตฺต\-โลภ\-สหคโต เอโก ขนฺโธ จ วตฺถุ จ ติณฺณนฺนํ ขนฺธานํ โลภสฺส จ อตฺถิปจฺจเยน ปจฺจโย ฯเปฯ. เทฺว ขนฺธา จ ฯเปฯ. (3)}

\prose{\csromanfolio{595}นตฺถิ\-ปจฺจเยน ปจฺจโย. วิคต\-ปจฺจเยน ปจฺจโย. อวิคต\-ปจฺจเยน ปจฺจโย.}

\csromanmarkh{1. ปจฺจยานุโลม}\csromanmarkhii{2. สงฺขฺยาวาร}\csromanheadernomark{2}{1. \textbf{ปจฺจยานุโลม 2. สงฺขฺยาวาร}}
\proseitem{88}{เหตุยา นว, อารมฺมเณ นว, อธิปติยา นว, อนนฺตเร นว, สมนนฺตเร นว, สหชาเต นว, อญฺญมญฺเญ ฉ, นิสฺสเย นว, อุปนิสฺสเย นว, ปุเรชาเต ตีณิ, ปจฺฉาชาเต ตีณิ, อาเสวเน นว, กมฺเม จตฺตาริ, วิปาเก เอกํ, อาหาเร จตฺตาริ, อินฺทฺริเย จตฺตาริ, ฌาเน จตฺตาริ, มคฺเค จตฺตาริ, สมฺปยุตฺเต ฉ, วิปฺปยุตฺเต ปญฺจ, อตฺถิยา นว, นตฺถิยา นว, วิคเต นว, อวิคเต นว.}

\proseitem{89}{อุปาทาน\-สมฺปยุตฺโต ธมฺโม อุปาทาน\-สมฺปยุตฺตสฺส ธมฺมสฺส อารมฺมณ\-ปจฺจเยน ปจฺจโย. สหชาต\-ปจฺจเยน ปจฺจโย. อุปนิสฺสย\-ปจฺจเยน ปจฺจโย. (1)}

\prose{อุปาทาน\-สมฺปยุตฺโต ธมฺโม อุปาทาน\-วิปฺปยุตฺตสฺส ธมฺมสฺส อารมฺมณ\-ปจฺจเยน ปจฺจโย. สหชาต\-ปจฺจเยน ปจฺจโย. อุปนิสฺสย\-ปจฺจเยน ปจฺจโย. ปจฺฉาชาต\-ปจฺจเยน ปจฺจโย. กมฺมปจฺจเยน ปจฺจโย. (2)}

\prose{อุปาทาน\-สมฺปยุตฺโต ธมฺโม อุปาทาน\-สมฺปยุตฺตสฺส จ อุปาทาน\-วิปฺปยุตฺตสฺส จ ธมฺมสฺส อารมฺมณ\-ปจฺจเยน ปจฺจโย. สหชาต\-ปจฺจเยน ปจฺจโย. อุปนิสฺสย\-ปจฺจเยน ปจฺจโย. (3)}

\proseitem{90}{อุปาทาน\-วิปฺปยุตฺโต ธมฺโม อุปาทาน\-วิปฺปยุตฺตสฺส ธมฺมสฺส อารมฺมณ\-ปจฺจเยน ปจฺจโย. สหชาต\-ปจฺจเยน ปจฺจโย. อุปนิสฺสย\-ปจฺจเยน ปจฺจโย. ปุเรชาต\-ปจฺจเยน ปจฺจโย. ปจฺฉาชาต\-ปจฺจเยน ปจฺจโย. กมฺมปจฺจเยน ปจฺจโย. อาหารปจฺจเยน ปจฺจโย. อินฺทฺริย\-ปจฺจเยน ปจฺจโย. (1)}

\prose{\csromanfolio{596}อุปาทาน\-วิปฺปยุตฺโต ธมฺโม อุปาทาน\-สมฺปยุตฺตสฺส ธมฺมสฺส อารมฺมณ\-ปจฺจเยน ปจฺจโย. สหชาต\-ปจฺจเยน ปจฺจโย. อุปนิสฺสย\-ปจฺจเยน ปจฺจโย. ปุเรชาต\-ปจฺจเยน ปจฺจโย. (2)}

\prose{อุปาทาน\-วิปฺปยุตฺโต ธมฺโม อุปาทาน\-สมฺปยุตฺตสฺส จ อุปาทาน\-วิปฺปยุตฺตสฺส จ ธมฺมสฺส อารมฺมณ\-ปจฺจเยน ปจฺจโย. สหชาต\-ปจฺจเยน ปจฺจโย. อุปนิสฺสย\-ปจฺจเยน ปจฺจโย. ปุเรชาต\-ปจฺจเยน ปจฺจโย. (3)}

\proseitem{91}{อุปาทาน\-สมฺปยุตฺโต จ อุปาทาน\-วิปฺปยุตฺโต จ ธมฺมา อุปาทาน\-สมฺปยุตฺตสฺส ธมฺมสฺส อารมฺมณ\-ปจฺจเยน ปจฺจโย. สหชาต\-ปจฺจเยน ปจฺจโย. อุปนิสฺสย\-ปจฺจเยน ปจฺจโย. (1)}

\prose{อุปาทาน\-สมฺปยุตฺโต จ อุปาทาน\-วิปฺปยุตฺโต จ ธมฺมา อุปาทาน\-วิปฺปยุตฺตสฺส ธมฺมสฺส อารมฺมณ\-ปจฺจเยน ปจฺจโย. สหชาต\-ปจฺจเยน ปจฺจโย. อุปนิสฺสย\-ปจฺจเยน ปจฺจโย. ปจฺฉาชาต\-ปจฺจเยน ปจฺจโย. (2)}

\prose{อุปาทาน\-สมฺปยุตฺโต จ อุปาทาน\-วิปฺปยุตฺโต จ ธมฺมา อุปาทาน\-สมฺปยุตฺตสฺส จ อุปาทาน\-วิปฺปยุตฺตสฺส จ ธมฺมสฺส อารมฺมณ\-ปจฺจเยน ปจฺจโย. สหชาต\-ปจฺจเยน ปจฺจโย. อุปนิสฺสย\-ปจฺจเยน ปจฺจโย. (3)}

\csromanmarkh{2. ปจฺจย\-ปจฺจนีย}\csromanmarkhii{2. สงฺขฺยาวาร}\csromanheadernomark{2}{2. ปจฺจย\-ปจฺจนีย 2. สงฺขฺยาวาร}
\proseitem{92}{นเหตุยา นว, นอารมฺมเณ นว, นอธิปติยา นว, (สพฺพตฺถ นว.) โนอวิคเต นว.}

\proseitem{93}{เหตุปจฺจยา นอารมฺมเณ นว, นอธิปติยา นว, นอนนฺตเร นว, นสมนนฺตเร นว, นอญฺญมญฺเญ ตีณิ, นอุปนิสฺสเย นว, (สพฺพตฺถ นว.) นสมฺปยุตฺเต ตีณิ, นวิปฺปยุตฺเต ฉ, โนนตฺถิยา นว, โนวิคเต นว.}

\proseitemwithcloser{94}{นเหตุปจฺจยา อารมฺมเณ นว, อธิปติยา นว, (อนุโลมมาติกา.) ฯเปฯ อวิคเต นว.}{อุปาทาน\-สมฺปยุตฺตทุกํ นิฏฺฐิตํ.}
\csromanmatikaanchor{mtk.72}
\csromantocmarkref{subsection}{72. อุปาทานอุปาทานิยทุก}{mtk.72}
\bookmark[dest={mtk.72},level=2]{72. อุปาทานอุปาทานิยทุก}
\proseitem{72}{\csromanfolio{597}อุปาทานอุปาทานิยทุก 72. อุปาทานอุปาทานิยทุก 1. ปฏิจฺจวาร}

\csromanmarkh{1. ปจฺจยานุโลม}\csromanmarkhii{1. วิภงฺควาร}\csromanheadernomark{2}{1. \textbf{ปจฺจยานุโลม 1. วิภงฺควาร}}
\csromanheader{2}{\textbf{เหตุ}}
\proseitem{95}{อุปาทานญฺเจว อุปาทานิยญฺจ ธมฺมํ ปฏิจฺจ อุปาทาโน เจว อุปาทานิโย จ ธมฺโม อุปฺปชฺชติ เหตุปจฺจยา. ทิฏฺฐุปาทานํ ปฏิจฺจ กามุปาทานํ. กามุปาทานํ ปฏิจฺจ ทิฏฺฐุปาทานํ. (จกฺกํ.) (1)}

\prose{อุปาทานญฺเจว อุปาทานิยญฺจ ธมฺมํ ปฏิจฺจ อุปาทานิโย เจว โน จ อุปาทาโน ธมฺโม อุปฺปชฺชติ เหตุปจฺจยา. อุปาทาเน ปฏิจฺจ สมฺปยุตฺตกา ขนฺธา จิตฺต\-สมุฏฺฐานญฺจ รูปํ. (2)}

\prose{อุปาทานญฺเจว อุปาทานิยญฺจ ธมฺมํ ปฏิจฺจ อุปาทาโน เจว อุปาทานิโย จ อุปาทานิโย เจว โน จ อุปาทาโน จ ธมฺมา อุปฺปชฺชนฺติ เหตุปจฺจยา. ทิฏฺฐุปาทานํ ปฏิจฺจ กามุปาทานํ สมฺปยุตฺตกา จ ขนฺธา จิตฺต\-สมุฏฺฐานญฺจ รูปํ. (จกฺกํ.) (3)}

\proseitem{96}{อุปาทานิยญฺเจว โน จ อุปาทานํ ธมฺมํ ปฏิจฺจ อุปาทานิโย เจว โน จ อุปาทาโน ธมฺโม อุปฺปชฺชติ เหตุปจฺจยา. อุปาทานิยญฺเจว โน จ อุปาทานํ เอกํ ขนฺธํ ปฏิจฺจ ตโย ขนฺธา จิตฺต\-สมุฏฺฐานญฺจ รูปํ ฯเปฯ. เทฺว ขนฺเธ ฯเปฯ. ปฏิสนฺธิ\-กฺขเณ ฯเปฯ. (ยาว อชฺฌตฺติกา มหาภูตา.) (1)}

\prose{อุปาทานิยญฺเจว โน จ อุปาทานํ ธมฺมํ ปฏิจฺจ อุปาทาโน เจว อุปาทานิโย จ ธมฺโม อุปฺปชฺชติ เหตุปจฺจยา. อุปาทานิเย เจว โน จ อุปาทาเน ขนฺเธ ปฏิจฺจ อุปาทานา. (2)}

\prose{อุปาทานิยญฺเจว โน จ อุปาทานํ ธมฺมํ ปฏิจฺจ อุปาทาโน เจว อุปาทานิโย จ อุปาทานิโย เจว โน จ อุปาทาโน จ ธมฺมา อุปฺปชฺชนฺติ เหตุปจฺจยา. อุปาทานิยญฺเจว โน จ อุปาทานํ เอกํ ขนฺธํ ปฏิจฺจ ตโย ขนฺธา อุปาทานา จ จิตฺต\-สมุฏฺฐานญฺจ รูปํ ฯเปฯ. เทฺว ขนฺเธ ฯเปฯ. (3)}

\proseitem{97}{อุปาทานญฺเจว อุปาทานิยญฺจ อุปาทานิยญฺเจว โน จ อุปาทานญฺจ ธมฺมํ ปฏิจฺจ อุปาทาโน เจว อุปาทานิโย จ ธมฺโม อุปฺปชฺชติ เหตุปจฺจยา. ทิฏฺฐุปาทานญฺจ สมฺปยุตฺตเก จ ขนฺเธ ปฏิจฺจ กามุปาทานํ. (จกฺกํ.) (1)}

\prose{\csromanfolio{598}อุปาทานญฺเจว อุปาทานิยญฺจ อุปาทานิยญฺเจว โน จ อุปาทานญฺจ ธมฺมํ ปฏิจฺจ อุปาทานิโย เจว โน จ อุปาทาโน ธมฺโม อุปฺปชฺชติ เหตุปจฺจยา. อุปาทานิยญฺเจว โน จ อุปาทานํ เอกํ ขนฺธญฺจ อุปาทาเน จ ปฏิจฺจ ตโย ขนฺธา จิตฺต\-สมุฏฺฐานญฺจ รูปํ ฯเปฯ เทฺว ขนฺเธ จ ฯเปฯ. อุปาทาเน จ มหาภูเต จ ปฏิจฺจ จิตฺต\-สมุฏฺฐานํ รูปํ. (2)}

\prose{อุปาทานญฺเจว อุปาทานิยญฺจ อุปาทานิยญฺเจว โน จ อุปาทานญฺจ ธมฺมํ ปฏิจฺจ อุปาทาโน เจว อุปาทานิโย จ อุปาทานิโย เจว โน จ อุปาทาโน จ ธมฺมา อุปฺปชฺชนฺติ เหตุปจฺจยา. อุปาทานิยญฺเจว โน จ อุปาทานํ เอกํ ขนฺธญฺจ ทิฏฺฐุปาทานญฺจ ปฏิจฺจ ตโย ขนฺธา กามุปาทานญฺจ จิตฺต\-สมุฏฺฐานญฺจ รูปํ ฯเปฯ. เทฺว ขนฺเธ จ ฯเปฯ. (จกฺกํ.) (3)}

\prose{(ยถา อุปาทานทุกํ. เอวํ ปฏิจฺจวาโรปิ สหชาตวาโรปิ ปจฺจยวาโรปิ นิสฺสยวาโรปิ สํสฏฺฐ\-วาโรปิ สมฺปยุตฺตวาโรปิ กาตพฺพา. นินฺนานากรณา, อามสนํ นานากรณํ.)}

\prose{72. อุปาทานอุปาทานิยทุก 7. ปญฺหาวาร}

\csromanmarkh{1. ปจฺจยานุโลม}\csromanmarkhii{1. วิภงฺควาร}\csromanheadernomark{2}{1. \textbf{ปจฺจยานุโลม 1. วิภงฺควาร}}
\csromanheader{2}{\textbf{เหตุ}}
\proseitem{98}{อุปาทาโน เจว อุปาทานิโย จ ธมฺโม อุปาทานสฺส เจว อุปาทานิยสฺส จ ธมฺมสฺส เหตุปจฺจเยน ปจฺจโย. อุปาทานา เจว อุปาทานิยา จ เหตู สมฺปยุตฺตกานํ อุปาทานานํ เหตุปจฺจเยน ปจฺจโย. (1)}

\prose{อุปาทาโน เจว อุปาทานิโย จ ธมฺโม อุปาทานิยสฺส เจว โน จ อุปาทานสฺส ธมฺมสฺส เหตุปจฺจเยน ปจฺจโย. อุปาทานา เจว อุปาทานิยา จ เหตู สมฺปยุตฺตกานํ ขนฺธานํ จิตฺตสมุฏฺฐานานญฺจ รูปานํ เหตุปจฺจเยน ปจฺจโย. (อุปาทาน\-ทุก\-สทิสา, นินฺนานา. นว ปญฺหา.)}

\prose{\csromanfolio{599}72. อุปาทานฺ\-เอาปาทานิยทุก อารมฺมณ}

\proseitem{99}{อุปาทาโน เจว อุปาทานิโย จ ธมฺโม อุปาทานสฺส เจว อุปาทานิยสฺส จ ธมฺมสฺส อารมฺมณ\-ปจฺจเยน ปจฺจโย. อุปาทาเน อารพฺภ อุปาทานา อุปฺปชฺชนฺติ. ตีณิ.}

\prose{อุปาทานิโย เจว โน จ อุปาทาโน ธมฺโม อุปาทานิยสฺส เจว โน จ อุปาทานสฺส ธมฺมสฺส อารมฺมณ\-ปจฺจเยน ปจฺจโย. ทานํ ฯเปฯ. ฌานา วุฏฺฐหิตฺวา ฌานํ ปจฺจเวกฺขติ, อสฺสาเทติ อภินนฺทติ, ตํ อารพฺภ ราโค อุปฺปชฺชติ. ทิฏฺฐิ อุปฺปชฺชติ. วิจิกิจฺฉา. อุทฺธจฺจํ. โทมนสฺสํ อุปฺปชฺชติ. อริยา โคตฺรภุํ ปจฺจเวกฺขนฺติ. โวทานํ ปจฺจเวกฺขนฺติ. ปหีเน กิเลเส ฯเปฯ. วิกฺขมฺภิเต กิเลเส ฯเปฯ. ปุพฺเพ สมุทาจิณฺเณ ฯเปฯ. จกฺขุํ ฯเปฯ. วตฺถุํ ฯเปฯ. (สํขิตฺตํ.) อนาคตํสญาณสฺส, อาวชฺชนาย อารมฺมณ\-ปจฺจเยน ปจฺจโย. (1)}

\prose{อุปาทานิโย เจว โน จ อุปาทาโน ธมฺโม อุปาทานสฺส เจว อุปาทานิยสฺส จ ธมฺมสฺส อารมฺมณ\-ปจฺจเยน ปจฺจโย. (สํขิตฺตํ. อิตเร เทฺว อุปาทาน\-ทุก\-สทิสา.) (3)}

\prose{อุปาทาโน เจว อุปาทานิโย จ อุปาทานิโย เจว โน จ อุปาทาโน จ ธมฺมา อุปาทานสฺส เจว อุปาทานิยสฺส จ ธมฺมสฺส อารมฺมณ\-ปจฺจเยน ปจฺจโย. ตีณิ. (เหฏฺฐา อธิปติ. ตีณิ. อุปาทาน\-ทุก\-สทิสา.)}

\proseitem{100}{อุปาทานิโย เจว โน จ อุปาทาโน ธมฺโม อุปาทานิยสฺส เจว โน จ อุปาทานสฺส จ ธมฺมสฺส อธิปติ\-ปจฺจเยน ปจฺจโย. อารมฺมณา\-ธิปติ, สหชาตา\-ธิปติ.  อารมฺมณา\-ธิปติ\csromandash{}ทานํ. ฌานา วุฏฺฐหิตฺวา ฌานํ ครุํ กตฺวา ปจฺจเวกฺขติ, อสฺสาเทติ อภินนฺทติ, ตํ ครุํ กตฺวา ราโค อุปฺปชฺชติ. ทิฏฺฐิ อุปฺปชฺชติ. เสกฺขา โคตฺรภุํ ครุํ กตฺวา ฯเปฯ. โวทานํ ฯเปฯ. จกฺขุํ ฯเปฯ วตฺถุํ อุปาทานิเย เจว โน จ อุปาทาเน ขนฺเธ ครุํ กตฺวา อุปาทานิยา เจว โน จ อุปาทานา ขนฺธา อุปฺปชฺชนฺติ. สหชาตา\-ธิปติ\csromandash{} อุปาทานิยา เจว โน จ อุปาทานาธิปติ สมฺปยุตฺตกานํ ขนฺธานํ จิตฺตสมุฏฺฐานานญฺจ รูปานํ อธิปติ\-ปจฺจเยน ปจฺจโย. (อวเสสา เทฺวปิ อารมฺมณา\-ธิปติปิ สหชาตา\-ธิปติปิ อุปาทาน\-ทุก\-สทิสา.) (3)}

\prosewithcloser{\csromanfolio{600}(ฆฏนา อธิปติปิ ตีณิ. อุปาทาน\-ทุก\-สทิสา. สพฺเพ ปจฺจยา อุปาทาน\-ทุก\-สทิสา. อุปาทานิเย โลกุตฺตรํ นตฺถิ, ปจฺจนียมฺปิ อิตเร เทฺว คณนาปิ อุปาทาน\-ทุก\-สทิสํ.)}{อุปาทานอุปาทานิยทุกํ นิฏฺฐิตํ.}
\csromanmatikaanchor{mtk.73}
\csromantocmarkref{subsection}{73. อุปาทานอุปาทานสมฺปยุตฺตทุก}{mtk.73}
\bookmark[dest={mtk.73},level=2]{73. อุปาทานอุปาทานสมฺปยุตฺตทุก}
\csromanmarkh{73. อุปาทานอุปาทานสมฺปยุตฺตทุก}\csromanmarkhii{1. ปฏิจฺจวาร}\csromanheadernomark{2}{73. อุปาทานอุปาทานสมฺปยุตฺตทุก 1. ปฏิจฺจวาร}
\csromanmarkh{1. ปจฺจยานุโลม}\csromanmarkhii{1. วิภงฺควาร}\csromanheadernomark{2}{1. \textbf{ปจฺจยานุโลม 1. วิภงฺควาร}}
\csromanheader{2}{\textbf{เหตุ}}
\proseitem{101}{อุปาทานญฺเจว อุปาทาน\-สมฺปยุตฺตญฺจ ธมฺมํ ปฏิจฺจ อุปาทาโน เจว อุปาทาน\-สมฺปยุตฺโต จ ธมฺโม อุปฺปชฺชติ เหตุปจฺจยา. ทิฏฺฐุปาทานํ ปฏิจฺจ กามุปาทานํ. (จกฺกํ กาตพฺพํ.) (1)}

\prose{อุปาทานญฺเจว อุปาทาน\-สมฺปยุตฺตญฺจ ธมฺมํ ปฏิจฺจ อุปาทาน\-สมฺปยุตฺโต เจว โน จ อุปาทาโน ธมฺโม อุปฺปชฺชติ เหตุปจฺจยา. อุปาทาเน ปฏิจฺจ สมฺปยุตฺตกา ขนฺธา. (2)}

\prose{อุปาทานญฺเจว อุปาทาน\-สมฺปยุตฺตญฺจ ธมฺมํ ปฏิจฺจ อุปาทาโน เจว อุปาทาน\-สมฺปยุตฺโต จ อุปาทาน\-สมฺปยุตฺโต เจว โน จ อุปาทาโน จ ธมฺมา อุปฺปชฺชนฺติ เหตุปจฺจยา. ทิฏฺฐุปาทานํ ปฏิจฺจ กามุปาทานํ สมฺปยุตฺตกา จ ขนฺธา. (จกฺกํ.) (3)}

\prose{อุปาทาน\-สมฺปยุตฺต\-ญฺเจว โน จ อุปาทานํ ธมฺมํ ปฏิจฺจ ฯเปฯ.}

\prose{(สํขิตฺตํ. อามสนํ นานากรณํ. อุปาทาน\-ทุก\-สทิสํ. นว ปญฺหา. รูปํ นตฺถิ. เอวํ สพฺเพปิ วารา วิตฺถาเรตพฺพา. อรูปํเยว.)}

\prose{73. อุปาทานอุปาทานสมฺปยุตฺตทุก 7. ปญฺหาวาร}

\csromanmarkh{1. ปจฺจยานุโลม}\csromanmarkhii{1. วิภงฺควาร}\csromanheadernomark{2}{1. \textbf{ปจฺจยานุโลม 1. วิภงฺควาร}}
\csromanheader{2}{\textbf{เหตุ}}
\proseitem{102}{อุปาทาโน เจว อุปาทาน\-สมฺปยุตฺโต จ ธมฺโม อุปาทานสฺส เจว อุปาทาน\-สมฺปยุตฺตสฺส จ ธมฺมสฺส เหตุปจฺจเยน ปจฺจโย. อุปาทานา เจว}

\proseitem{73}{\csromanfolio{601}อุปาทานอุปาทานสมฺปยุตฺตทุก อุปาทาน\-สมฺปยุตฺตา จ เหตู สมฺปยุตฺตกานํ อุปาทานานํ เหตุปจฺจเยน ปจฺจโย. (มูลํ ปุจฺฉิตพฺพํ.) อุปาทานา เจว อุปาทาน\-สมฺปยุตฺตา จ เหตู สมฺปยุตฺตกานํ ขนฺธานํ เหตุปจฺจเยน ปจฺจโย. (มูลํ ปุจฺฉิตพฺพํ.) อุปาทานา เจว อุปาทาน\-สมฺปยุตฺตา จ เหตู สมฺปยุตฺตกานํ ขนฺธานํ อุปาทานานญฺจ เหตุปจฺจเยน ปจฺจโย. (3)}

\proseitem{103}{อุปาทาน\-สมฺปยุตฺโต เจว โน จ อุปาทาโน ธมฺโม อุปาทาน\-สมฺปยุตฺตสฺส เจว โน จ อุปาทานสฺส ธมฺมสฺส เหตุปจฺจเยน ปจฺจโย. อุปาทาน\-สมฺปยุตฺตา เจว โน จ อุปาทานา เหตู สมฺปยุตฺตกานํ ขนฺธานํ เหตุปจฺจเยน ปจฺจโย. (มูลํ ปุจฺฉิตพฺพํ.) อุปาทาน\-สมฺปยุตฺตา เจว โน จ อุปาทานา เหตู สมฺปยุตฺตกานํ อุปาทานานํ เหตุปจฺจเยน ปจฺจโย. (มูลํ ปุจฺฉิตพฺพํ.) อุปาทาน\-สมฺปยุตฺตา เจว โน จ อุปาทานา เหตู สมฺปยุตฺตกานํ ขนฺธานํ อุปาทานานญฺจ เหตุปจฺจเยน ปจฺจโย. (3)}

\proseitem{104}{อุปาทาโน เจว อุปาทาน\-สมฺปยุตฺโต จ อุปาทาน\-สมฺปยุตฺโต เจว โน จ อุปาทาโน จ ธมฺมา อุปาทานสฺส เจว อุปาทาน\-สมฺปยุตฺตสฺส จ ธมฺมสฺส เหตุปจฺจเยน ปจฺจโย. อุปาทานา เจว อุปาทาน\-สมฺปยุตฺตา จ อุปาทาน\-สมฺปยุตฺตา เจว โน จ อุปาทานา จ เหตู สมฺปยุตฺตกานํ อุปาทานานํ เหตุปจฺจเยน ปจฺจโย. (มูลํ ปุจฺฉิตพฺพํ.) อุปาทานา เจว อุปาทาน\-สมฺปยุตฺตา จ อุปาทาน\-สมฺปยุตฺตา เจว โน จ อุปาทานา จ เหตู สมฺปยุตฺตกานํ ขนฺธานํ เหตุปจฺจเยน ปจฺจโย. (มูลํ ปุจฺฉิตพฺพํ.) อุปาทานา เจว อุปาทาน\-สมฺปยุตฺตา จ อุปาทาน\-สมฺปยุตฺตา เจว โน จ อุปาทานา จ เหตู สมฺปยุตฺตกานํ ขนฺธานํ อุปาทานานญฺจ เหตุปจฺจเยน ปจฺจโย. (3)}

\proseitem{105}{อุปาทาโน เจว อุปาทาน\-สมฺปยุตฺโต จ ธมฺโม อุปาทานสฺส เจว อุปาทาน\-สมฺปยุตฺตสฺส จ ธมฺมสฺส อารมฺมณ\-ปจฺจเยน ปจฺจโย. อุปาทาเน อารพฺภ อุปาทานา อุปฺปชฺชนฺติ. (มูลํ ปุจฺฉิตพฺพํ.) อุปาทาเน อารพฺภ อุปาทาน\-สมฺปยุตฺตา เจว โน จ อุปาทานา ขนฺธา อุปฺปชฺชนฺติ. (มูลํ ปุจฺฉิตพฺพํ.) อุปาทาเน อารพฺภ อุปาทานา จ สมฺปยุตฺตกา จ ขนฺธา อุปฺปชฺชนฺติ. (3)}

\prose{\csromanfolio{602}อุปาทาน\-สมฺปยุตฺโต เจว โน จ อุปาทาโน ธมฺโม อุปาทาน\-สมฺปยุตฺตสฺส เจว โน จ อุปาทานสฺส ธมฺมสฺส อารมฺมณ\-ปจฺจเยน ปจฺจโย. อุปาทาน\-สมฺปยุตฺเต เจว โน จ อุปาทาเน ขนฺเธ อารพฺภ อุปาทาน\-สมฺปยุตฺตา เจว โน จ อุปาทานา ขนฺธา อุปฺปชฺชนฺติ. (ตีณิปิ กาตพฺพา, ฆฏเน ตีณิปิ กาตพฺพา.)}

\proseitem{106}{อุปาทาโน เจว อุปาทาน\-สมฺปยุตฺโต จ ธมฺโม อุปาทานสฺส เจว อุปาทาน\-สมฺปยุตฺตสฺส จ ธมฺมสฺส อธิปติ\-ปจฺจเยน ปจฺจโย. ตีณิ. (3)}

\prose{อุปาทาน\-สมฺปยุตฺโต เจว โน จ อุปาทาโน ธมฺโม อุปาทาน\-สมฺปยุตฺตสฺส เจว โน จ อุปาทานสฺส ธมฺมสฺส อธิปติ\-ปจฺจเยน ปจฺจโย. อารมฺมณา\-ธิปติ, สหชาตา\-ธิปติ.}

\vspace{\tipitakaparskip}
\csromancenter{(ตีณิปิ, ตีสุปิ เทฺวปิ อธิปตี กาตพฺพา, ฆฏนาธิปติปิ ตีณิ.)}
\proseitem{107}{อุปาทาโน เจว อุปาทาน\-สมฺปยุตฺโต จ ธมฺโม อุปาทานสฺส เจว อุปาทาน\-สมฺปยุตฺตสฺส จ ธมฺมสฺส อนนฺตร\-ปจฺจเยน ปจฺจโย. ปุริมา ปุริมา อุปาทานา เจว อุปาทาน\-สมฺปยุตฺตา จ ขนฺธา ปจฺฉิมานํ ปจฺฉิมานํ อุปาทานานํ อนนฺตร\-ปจฺจเยน ปจฺจโย. (เอวํ นวปิ ปญฺหา กาตพฺพา, อาวชฺชนาปิ วุฏฺฐานมฺปิ นตฺถิ.}

\csromanheader{2}{\textbf{สมนนฺตราทิ}}
\proseitem{108}{อุปาทาโน เจว อุปาทาน\-สมฺปยุตฺโต จ ธมฺโม อุปาทานสฺส เจว อุปาทาน\-สมฺปยุตฺตสฺส จ ธมฺมสฺส สมนนฺตร\-ปจฺจเยน ปจฺจโย. นว. สหชาต\-ปจฺจเยน ปจฺจโย. นว. อญฺญมญฺญ\-ปจฺจเยน ปจฺจโย. นว. นิสฺสย\-ปจฺจเยน ปจฺจโย. นว.}

\csromanheader{2}{\textbf{อุปนิสฺสย}}
\proseitem{109}{อุปาทาโน เจว อุปาทาน\-สมฺปยุตฺโต จ ธมฺโม อุปาทานสฺส เจว อุปาทาน\-สมฺปยุตฺตสฺส จ ธมฺมสฺส อุปนิสฺสย\-ปจฺจเยน ปจฺจโย ฯเปฯ. ตีณิ.}

\prose{\csromanfolio{603}73. อุปาทานอุปาทานสมฺปยุตฺตทุก}

\prose{อุปาทาน\-สมฺปยุตฺโต เจว โน จ อุปาทาโน ธมฺโม อุปาทาน\-สมฺปยุตฺตสฺส เจว โน จ อุปาทานสฺส ธมฺมสฺส อุปนิสฺสย\-ปจฺจเยน ปจฺจโย. อารมฺมณู\-ปนิสฺสโย, อนนฺตรู\-ปนิสฺสโย, ปกตู\-ปนิสฺสโย ฯเปฯ. ปกตู\-ปนิสฺสโย\csromandash{}อุปาทาน\-สมฺปยุตฺตา เจว โน จ อุปาทานา ขนฺธา อุปาทาน\-สมฺปยุตฺตาน\-ญฺเจว โน จ อุปาทานานํ ขนฺธานํ อุปนิสฺสย\-ปจฺจเยน ปจฺจโย. (ตีณิ ฆฏนุปนิสฺสเยปิ ตีณิ.) อาเสวนปจฺจเยนปจฺจโย. นว.}

\csromanheader{2}{\textbf{กมฺมาทิ}}
\proseitem{110}{อุปาทาน\-สมฺปยุตฺโต เจว โน จ อุปาทาโน ธมฺโม อุปาทาน\-สมฺปยุตฺตสฺส เจว โน จ อุปาทานสฺส ธมฺมสฺส กมฺมปจฺจเยน ปจฺจโย. ตีณิ. อาหารปจฺจเยน ปจฺจโย. ตีณิ. อินฺทฺริย\-ปจฺจเยน ปจฺจโย. ตีณิ. ฌานปจฺจเยน ปจฺจโย. ตีณิ. มคฺคปจฺจเยน ปจฺจโย. นว. สมฺปยุตฺต\-ปจฺจเยน ปจฺจโย. นว. อตฺถิปจฺจเยน ปจฺจโย. นว. นตฺถิ\-ปจฺจเยน ปจฺจโย. นว. วิคต\-ปจฺจเยน ปจฺจโย. นว. อวิคต\-ปจฺจเยน ปจฺจโย. นว.}

\csromanmarkh{1. ปจฺจยานุโลม}\csromanmarkhii{2. สงฺขฺยาวาร}\csromanheadernomark{2}{1. \textbf{ปจฺจยานุโลม 2. สงฺขฺยาวาร}}
\proseitem{111}{เหตุยา นว, อารมฺมเณ นว, อธิปติยา นว, อนนฺตเร นว, สมนนฺตเร นว, สหชาเต นว, อญฺญมญฺเญ นว, นิสฺสเย นว, อุปนิสฺสเย นว, อาเสวเน นว, กมฺเม ตีณิ, อาหาเร ตีณิ, อินฺทฺริเย ตีณิ, ฌาเน ตีณิ, มคฺเค นว, สมฺปยุตฺเต นว, อตฺถิยา นว, นตฺถิยา นว, วิคเต นว, อวิคเต นว.}

\proseitem{112}{อุปาทาโน เจว อุปาทาน\-สมฺปยุตฺโต จ ธมฺโม อุปาทานสฺส เจว อุปาทาน\-สมฺปยุตฺตสฺส จ ธมฺมสฺส อารมฺมณ\-ปจฺจเยน ปจฺจโย. สหชาต\-ปจฺจเยน ปจฺจโย. อุปนิสฺสย\-ปจฺจเยน ปจฺจโย. (1)}

\prose{\csromanfolio{604}อุปาทาโน เจว อุปาทาน\-สมฺปยุตฺโต จ ธมฺโม อุปาทาน\-สมฺปยุตฺตสฺส เจว โน จ อุปาทานสฺส ธมฺมสฺส อารมฺมณ\-ปจฺจเยน ปจฺจโย. สหชาต\-ปจฺจเยน ปจฺจโย. อุปนิสฺสย\-ปจฺจเยน ปจฺจโย. (2)}

\prose{อุปาทาโน เจว อุปาทาน\-สมฺปยุตฺโต จ ธมฺโม อุปาทานสฺส เจว อุปาทาน\-สมฺปยุตฺตสฺส จ อุปาทาน\-สมฺปยุตฺตสฺส เจว โน จ อุปาทานสฺส จ ธมฺมสฺส อารมฺมณ\-ปจฺจเยน ปจฺจโย. สหชาต\-ปจฺจเยน ปจฺจโย. อุปนิสฺสย\-ปจฺจเยน ปจฺจโย.}

\prose{(เอวํ นวปิ กาตพฺพา, เอเกกสฺส มูเล ตีณิ ตีณิ ปญฺหา.)}

\csromanmarkh{2. ปจฺจย\-ปจฺจนีย}\csromanmarkhii{2. สงฺขฺยาวาร}\csromanheadernomark{2}{2. ปจฺจย\-ปจฺจนีย 2. สงฺขฺยาวาร}
\proseitem{113}{นเหตุยา นว, นอารมฺมเณ นว, นอธิปติยา นว, (สพฺพตฺถ นว.) โนอวิคเต นว.}

\vspace{\tipitakaparskip}
\csromancenter{\textbf{เหตุ}ปจฺจยา นอารมฺมเณ นว, นอธิปติยา นว, นอนนฺตเร นว, นสมนนฺตเร นว, นฺเอาปนิสฺสเย นว, (สพฺพตฺถ นว.) นมคฺเค นว, นสมฺปยุตฺเต นว, โนนตฺถิยา นว, โนวิคเต นว.}
\csromanheader{2}{115. นเหตุปจฺจยา อารมฺมเณ นว, อธิปติยา นว ฯเปฯ. (อนุโลมมาติกา กาตพฺพา.)}
\nitthitam{อุปาทานอุปาทานสมฺปยุตฺตทุกํ นิฏฺฐิตํ.}
\csromanmatikaanchor{mtk.74}
\csromantocmarkref{subsection}{74. อุปาทานวิปฺปยุตฺตอุปาทานิยทุก}{mtk.74}
\bookmark[dest={mtk.74},level=2]{74. อุปาทานวิปฺปยุตฺตอุปาทานิยทุก}
\csromanmarkh{74. อุปาทานวิปฺปยุตฺตอุปาทานิยทุก}\csromanmarkhii{1. ปฏิจฺจวาร}\csromanheadernomark{2}{74. อุปาทานวิปฺปยุตฺตอุปาทานิยทุก 1. ปฏิจฺจวาร}
\vspace{\tipitakaparskip}
\csromancenter{\textbf{เหตุ}}
\proseitem{116}{อุปาทาน\-วิปฺปยุตฺตํ อุปาทานิยํ ธมฺมํ ปฏิจฺจ อุปาทาน\-วิปฺปยุตฺโต อุปาทานิโย ธมฺโม อุปฺปชฺชติ เหตุปจฺจยา. อุปาทาน\-วิปฺปยุตฺตํ อุปาทานิยํ}

\proseitem{74}{\csromanfolio{605}อุปาทานวิปฺปยุตฺตอุปาทานิยทุก เอกํ ขนฺธํ ปฏิจฺจ ตโย ขนฺธา จิตฺต\-สมุฏฺฐานญฺจ รูปํ ฯเปฯ เทฺว ขนฺเธ ฯเปฯ. ปฏิสนฺธิ\-กฺขเณ ฯเปฯ. เอกํ มหาภูตํ ฯเปฯ มหาภูเต ปฏิจฺจ จิตฺต\-สมุฏฺฐานํ รูปํ, กฏตฺตารูปํ, อุปาทารูปํ. (1)}

\prose{อุปาทาน\-วิปฺปยุตฺตํ อนุปาทานิยํ ธมฺมํ ปฏิจฺจ อุปาทาน\-วิปฺปยุตฺโต อนุปาทานิโย ธมฺโม อุปฺปชฺชติ เหตุปจฺจยา. อุปาทาน\-วิปฺปยุตฺตํ อนุปาทานิยํ เอกํ ขนฺธํ ปฏิจฺจ ตโย ขนฺธา ฯเปฯ เทฺว ขนฺเธ ฯเปฯ. (1)}

\prose{อุปาทาน\-วิปฺปยุตฺตํ อนุปาทานิยํ ธมฺมํ ปฏิจฺจ อุปาทาน\-วิปฺปยุตฺโต อุปาทานิโย ธมฺโม อุปฺปชฺชติ เหตุปจฺจยา. อุปาทาน\-วิปฺปยุตฺเต อนุปาทานิเย ขนฺเธ ปฏิจฺจ จิตฺต\-สมุฏฺฐานํ รูปํ.}

\prose{อุปาทาน\-วิปฺปยุตฺตํ อนุปาทานิยํ ธมฺมํ ปฏิจฺจ อุปาทาน\-วิปฺปยุตฺโต อุปาทานิโย จ อุปาทาน\-วิปฺปยุตฺโต อนุปาทานิโย จ ธมฺมา อุปฺปชฺชนฺติ เหตุปจฺจยา. อุปาทาน\-วิปฺปยุตฺตํ อนุปาทานิยํ เอกํ ขนฺธํ ปฏิจฺจ ตโย ขนฺธา จิตฺต\-สมุฏฺฐานญฺจ รูปํ ฯเปฯ เทฺว ขนฺเธ ฯเปฯ. (3)}

\prose{อุปาทาน\-วิปฺปยุตฺตํ อุปาทานิยญฺจ อุปาทาน\-วิปฺปยุตฺตํ อนุปาทานิยญฺจ ธมฺมํ ปฏิจฺจ อุปาทาน\-วิปฺปยุตฺโต อุปาทานิโย ธมฺโม อุปฺปชฺชติ เหตุปจฺจยา. อุปาทาน\-วิปฺปยุตฺเต อนุปาทานิเย ขนฺเธ จ มหาภูเต จ ปฏิจฺจ จิตฺต\-สมุฏฺฐานํ รูปํ. (1)}

\prose{เหตุยา ปญฺจ, อารมฺมเณ เทฺว ฯเปฯ อวิคเต ปญฺจ.}

\prosewithcloser{(อิมํ ทุกํ จูฬนฺตรทุเก โลกิย\-ทุก\-สทิสํ. นินฺนานากรณํ.)}{อุปาทานวิปฺปยุตฺตอุปาทานิยทุกํ นิฏฺฐิตํ.}
\nitthitam{อุปาทาน\-โคจฺฉกํ นิฏฺฐิตํ.}
